%% Generated by Sphinx.
\def\sphinxdocclass{report}
\documentclass[a4,10pt,english]{sphinxmanual}
\ifdefined\pdfpxdimen
   \let\sphinxpxdimen\pdfpxdimen\else\newdimen\sphinxpxdimen
\fi \sphinxpxdimen=.75bp\relax
\ifdefined\pdfimageresolution
    \pdfimageresolution= \numexpr \dimexpr1in\relax/\sphinxpxdimen\relax
\fi
\newdimen\sphinxremdimen\sphinxremdimen = 10pt
%% let collapsible pdf bookmarks panel have high depth per default
\PassOptionsToPackage{bookmarksdepth=5}{hyperref}
%% turn off hyperref patch of \index as sphinx.xdy xindy module takes care of
%% suitable \hyperpage mark-up, working around hyperref-xindy incompatibility
\PassOptionsToPackage{hyperindex=false}{hyperref}
%% memoir class requires extra handling
\makeatletter\@ifclassloaded{memoir}
{\ifdefined\memhyperindexfalse\memhyperindexfalse\fi}{}\makeatother

\PassOptionsToPackage{booktabs}{sphinx}
\PassOptionsToPackage{colorrows}{sphinx}

\PassOptionsToPackage{warn}{textcomp}

\catcode`^^^^00a0\active\protected\def^^^^00a0{\leavevmode\nobreak\ }
\usepackage{cmap}
\usepackage{xeCJK}
\usepackage{amsmath,amssymb,amstext}
\usepackage{babel}



\setmainfont{FreeSerif}[
  Extension      = .otf,
  UprightFont    = *,
  ItalicFont     = *Italic,
  BoldFont       = *Bold,
  BoldItalicFont = *BoldItalic
]
\setsansfont{FreeSans}[
  Extension      = .otf,
  UprightFont    = *,
  ItalicFont     = *Oblique,
  BoldFont       = *Bold,
  BoldItalicFont = *BoldOblique,
]
\setmonofont{FreeMono}[Scale=0.9,
  Extension      = .otf,
  UprightFont    = *,
  ItalicFont     = *Oblique,
  BoldFont       = *Bold,
  BoldItalicFont = *BoldOblique,
]



\usepackage[Sonny]{fncychap}
\ChNameVar{\Large\normalfont\sffamily}
\ChTitleVar{\Large\normalfont\sffamily}
\usepackage{sphinx}

\fvset{fontsize=\small,formatcom=\xeCJKVerbAddon}
\usepackage{geometry}


% Include hyperref last.
\usepackage{hyperref}
% Fix anchor placement for figures with captions.
\usepackage{hypcap}% it must be loaded after hyperref.
% Set up styles of URL: it should be placed after hyperref.
\urlstyle{same}

\addto\captionsenglish{\renewcommand{\contentsname}{MTDataPro 入门}}

\usepackage{sphinxmessages}
\setcounter{tocdepth}{1}


\usepackage{xeCJK}
\setCJKmainfont{FandolSong}
\setCJKsansfont{FandolHei}
\setCJKmonofont{FandolFang}


\title{MTDataPro 中文手册}
\date{2026 年 03 月 29 日}
\release{1.9.5}
\author{王培杰}
\newcommand{\sphinxlogo}{\vbox{}}
\renewcommand{\releasename}{发行版本}
\makeindex
\begin{document}

\ifdefined\shorthandoff
  \ifnum\catcode`\=\string=\active\shorthandoff{=}\fi
  \ifnum\catcode`\"=\active\shorthandoff{"}\fi
\fi

\pagestyle{empty}
\sphinxmaketitle
\pagestyle{plain}
\sphinxtableofcontents
\pagestyle{normal}
\phantomsection\label{\detokenize{index::doc}}


\sphinxAtStartPar
欢迎来到 \sphinxstylestrong{MTDataPro} —— 大地电磁（MT）数据处理软件。

\sphinxAtStartPar
MTDataPro 是地球科学领域广泛使用的专业 MT 数据处理软件之一。本手册由 王培杰 维护整理，详尽介绍了 MTDataPro 的用法并提供了大量实用示例，既可以作为初学者的入门读物，也可以作为日常使用的参考书。


\bigskip\hrule\bigskip



\chapter{手册学习指南}
\label{\detokenize{index:id1}}
\sphinxAtStartPar
本手册主要包含如下三个部分：
\begin{enumerate}
\sphinxsetlistlabels{\arabic}{enumi}{enumii}{}{.}%
\item {} 
\sphinxAtStartPar
\sphinxstylestrong{MTDataPro 入门}：介绍 MTDataPro 的功能特点、安装方法，并为初学者提供快速入门教程。初学者应完整阅读"入门"章节，并通过练习掌握 MTDataPro 的基本用法。

\item {} 
\sphinxAtStartPar
\sphinxstylestrong{MTDataPro 进阶}：详细介绍各功能模块的使用方法和参数配置，可以作为参考书查阅。

\item {} 
\sphinxAtStartPar
\sphinxstylestrong{MTDataPro 实例}：包含丰富的实用脚本和处理案例，可以作为日常数据处理参考。

\end{enumerate}


\bigskip\hrule\bigskip



\chapter{快速链接}
\label{\detokenize{index:id2}}

\begin{savenotes}\sphinxattablestart
\sphinxthistablewithglobalstyle
\centering
\begin{tabulary}{\linewidth}[t]{TTT}
\sphinxtoprule
\begin{varwidth}[t]{\sphinxcolwidth{1}{3}}
\sphinxstyletheadfamily \sphinxAtStartPar
快速入门
\sphinxbeforeendvarwidth
\end{varwidth}%
&\begin{varwidth}[t]{\sphinxcolwidth{1}{3}}
\sphinxstyletheadfamily \sphinxAtStartPar
安装指南
\sphinxbeforeendvarwidth
\end{varwidth}%
&\begin{varwidth}[t]{\sphinxcolwidth{1}{3}}
\sphinxstyletheadfamily \sphinxAtStartPar
处理流程
\sphinxbeforeendvarwidth
\end{varwidth}%
\\
\sphinxmidrule
\sphinxtableatstartofbodyhook\begin{varwidth}[t]{\sphinxcolwidth{1}{3}}
\sphinxAtStartPar
{\hyperref[\detokenize{intro/quickstart::doc}]{\sphinxcrossref{\DUrole{doc}{\DUrole{std}{\DUrole{std-doc}{10 分钟完成第一次数据处理}}}}}}
\sphinxbeforeendvarwidth
\end{varwidth}%
&\begin{varwidth}[t]{\sphinxcolwidth{1}{3}}
\sphinxAtStartPar
{\hyperref[\detokenize{intro/install::doc}]{\sphinxcrossref{\DUrole{doc}{\DUrole{std}{\DUrole{std-doc}{安装、配置与授权}}}}}}
\sphinxbeforeendvarwidth
\end{varwidth}%
&\begin{varwidth}[t]{\sphinxcolwidth{1}{3}}
\sphinxAtStartPar
{\hyperref[\detokenize{chapters/chapter4::doc}]{\sphinxcrossref{\DUrole{doc}{\DUrole{std}{\DUrole{std-doc}{时间序列处理详解}}}}}}
\sphinxbeforeendvarwidth
\end{varwidth}%
\\
\sphinxbottomrule
\end{tabulary}
\sphinxtableafterendhook\par
\sphinxattableend\end{savenotes}


\begin{savenotes}\sphinxattablestart
\sphinxthistablewithglobalstyle
\centering
\begin{tabulary}{\linewidth}[t]{TTT}
\sphinxtoprule
\begin{varwidth}[t]{\sphinxcolwidth{1}{3}}
\sphinxstyletheadfamily \sphinxAtStartPar
仪器支持
\sphinxbeforeendvarwidth
\end{varwidth}%
&\begin{varwidth}[t]{\sphinxcolwidth{1}{3}}
\sphinxstyletheadfamily \sphinxAtStartPar
数据导出
\sphinxbeforeendvarwidth
\end{varwidth}%
&\begin{varwidth}[t]{\sphinxcolwidth{1}{3}}
\sphinxstyletheadfamily \sphinxAtStartPar
常见问题
\sphinxbeforeendvarwidth
\end{varwidth}%
\\
\sphinxmidrule
\sphinxtableatstartofbodyhook\begin{varwidth}[t]{\sphinxcolwidth{1}{3}}
\sphinxAtStartPar
{\hyperref[\detokenize{chapters/chapter3::doc}]{\sphinxcrossref{\DUrole{doc}{\DUrole{std}{\DUrole{std-doc}{Phoenix、Metronix、LEMI}}}}}}
\sphinxbeforeendvarwidth
\end{varwidth}%
&\begin{varwidth}[t]{\sphinxcolwidth{1}{3}}
\sphinxAtStartPar
{\hyperref[\detokenize{chapters/chapter6::doc}]{\sphinxcrossref{\DUrole{doc}{\DUrole{std}{\DUrole{std-doc}{EDI 格式}}}}}}
\sphinxbeforeendvarwidth
\end{varwidth}%
&\begin{varwidth}[t]{\sphinxcolwidth{1}{3}}
\sphinxAtStartPar
{\hyperref[\detokenize{appendices/appendixD::doc}]{\sphinxcrossref{\DUrole{doc}{\DUrole{std}{\DUrole{std-doc}{FAQ 与故障排除}}}}}}
\sphinxbeforeendvarwidth
\end{varwidth}%
\\
\sphinxbottomrule
\end{tabulary}
\sphinxtableafterendhook\par
\sphinxattableend\end{savenotes}


\bigskip\hrule\bigskip



\chapter{快速开始}
\label{\detokenize{index:id3}}\begin{enumerate}
\sphinxsetlistlabels{\arabic}{enumi}{enumii}{}{.}%
\item {} 
\sphinxAtStartPar
安装软件并完成授权

\item {} 
\sphinxAtStartPar
创建或打开工程

\item {} 
\sphinxAtStartPar
导入仪器数据

\item {} 
\sphinxAtStartPar
执行 FFT 处理

\item {} 
\sphinxAtStartPar
数据筛选与阻抗估计

\item {} 
\sphinxAtStartPar
导出 EDI 结果

\end{enumerate}


\bigskip\hrule\bigskip



\chapter{文档下载}
\label{\detokenize{index:id4}}\begin{itemize}
\item {} 
\sphinxAtStartPar
\sphinxhref{https://github.com/EMWPJ/MTDataPro-docs}{MTDataPro 中文手册源码} (https://github.com/EMWPJ/MTDataPro\sphinxhyphen{}docs)

\item {} 
\sphinxAtStartPar
\sphinxhref{https://github.com/EMWPJ/MTDataPro-docs/releases}{MTDataPro 中文手册 PDF} (https://github.com/EMWPJ/MTDataPro\sphinxhyphen{}docs/releases)（即将提供）

\end{itemize}


\bigskip\hrule\bigskip



\chapter{外部链接}
\label{\detokenize{index:id5}}\begin{itemize}
\item {} 
\sphinxAtStartPar
\sphinxhref{https://github.com/EMWPJ/MTDataPro-docs}{MTDataPro 文档仓库} (https://github.com/EMWPJ/MTDataPro\sphinxhyphen{}docs)

\item {} 
\sphinxAtStartPar
\sphinxhref{https://github.com/EMWPJ/MTDataPro-docs/discussions}{参与讨论} (https://github.com/EMWPJ/MTDataPro\sphinxhyphen{}docs/discussions)

\end{itemize}


\bigskip\hrule\bigskip



\chapter{引用 MTDataPro}
\label{\detokenize{index:id6}}
\sphinxAtStartPar
如果您在科研工作中使用了 MTDataPro，请引用：

\begin{sphinxVerbatim}[commandchars=\\\{\}]
\PYG{n+nc}{@software}\PYG{p}{\PYGZob{}}\PYG{n+nl}{mtdp2026}\PYG{p}{,}
\PYG{+w}{  }\PYG{n+na}{title}\PYG{+w}{ }\PYG{p}{=}\PYG{+w}{ }\PYG{l+s}{\PYGZob{}}\PYG{l+s}{MTDataPro: A Professional Magnetotelluric Data Processing Software}\PYG{l+s}{\PYGZcb{}}\PYG{p}{,}
\PYG{+w}{  }\PYG{n+na}{author}\PYG{+w}{ }\PYG{p}{=}\PYG{+w}{ }\PYG{l+s}{\PYGZob{}}\PYG{l+s}{王培杰}\PYG{l+s}{\PYGZcb{}}\PYG{p}{,}
\PYG{+w}{  }\PYG{n+na}{year}\PYG{+w}{ }\PYG{p}{=}\PYG{+w}{ }\PYG{l+s}{\PYGZob{}}\PYG{l+s}{2026}\PYG{l+s}{\PYGZcb{}}\PYG{p}{,}
\PYG{+w}{  }\PYG{n+na}{version}\PYG{+w}{ }\PYG{p}{=}\PYG{+w}{ }\PYG{l+s}{\PYGZob{}}\PYG{l+s}{1.9.5}\PYG{l+s}{\PYGZcb{}}\PYG{p}{,}
\PYG{+w}{  }\PYG{n+na}{url}\PYG{+w}{ }\PYG{p}{=}\PYG{+w}{ }\PYG{l+s}{\PYGZob{}}\PYG{l+s}{https://github.com/EMWPJ/MTDataPro\PYGZhy{}docs}\PYG{l+s}{\PYGZcb{}}
\PYG{p}{\PYGZcb{}}
\end{sphinxVerbatim}


\bigskip\hrule\bigskip


\sphinxstepscope


\section{简介}
\label{\detokenize{intro/index:id1}}\label{\detokenize{intro/index::doc}}
\sphinxAtStartPar
MTDataPro（Magnetotelluric Data Processing）是一款专业的 MT 数据处理软件，用于大地电磁测深数据的处理和分析。


\subsection{什么是 MTDataPro？}
\label{\detokenize{intro/index:mtdatapro}}
\sphinxAtStartPar
MTDataPro 是由 王培杰 开发和维护的大地电磁数据处理软件。该软件提供了从原始时间序列数据导入、处理、分析到结果导出的完整工作流程，支持多种主流 MT 仪器数据格式。


\subsection{主要特点}
\label{\detokenize{intro/index:id2}}\begin{itemize}
\item {} 
\sphinxAtStartPar
\sphinxstylestrong{多仪器支持}：Phoenix、Metronix、LEMI、ATTS、RMT 等主流仪器格式

\item {} 
\sphinxAtStartPar
\sphinxstylestrong{完整处理流程}：时间序列预处理、FFT 频谱分析、传递函数估计

\item {} 
\sphinxAtStartPar
\sphinxstylestrong{智能数据筛选}：自动质量评估、Robust 估计、参考道处理

\item {} 
\sphinxAtStartPar
\sphinxstylestrong{正演建模}：1D/2D/3D 大地电磁正演计算

\item {} 
\sphinxAtStartPar
\sphinxstylestrong{可视化输出}：专业图表绘制、地理信息展示

\item {} 
\sphinxAtStartPar
\sphinxstylestrong{标准格式导出}：EDI、J\sphinxhyphen{}format 等国际标准格式

\end{itemize}


\subsection{版本信息}
\label{\detokenize{intro/index:id3}}\begin{itemize}
\item {} 
\sphinxAtStartPar
\sphinxstylestrong{当前版本}：v1.9.5

\item {} 
\sphinxAtStartPar
\sphinxstylestrong{发布日期}：2026 年 2 月

\item {} 
\sphinxAtStartPar
\sphinxstylestrong{开发环境}：Embarcadero Delphi 10.4 Sydney

\item {} 
\sphinxAtStartPar
\sphinxstylestrong{运行平台}：Windows 64 位

\end{itemize}


\subsection{系统要求}
\label{\detokenize{intro/index:id4}}\begin{itemize}
\item {} 
\sphinxAtStartPar
\sphinxstylestrong{操作系统}：Windows 10/11 64 位

\item {} 
\sphinxAtStartPar
\sphinxstylestrong{内存}：建议 16 GB RAM（最低 8 GB）

\item {} 
\sphinxAtStartPar
\sphinxstylestrong{硬盘}：2 GB 可用空间

\item {} 
\sphinxAtStartPar
\sphinxstylestrong{显示器}：1920×1080 或更高分辨率

\end{itemize}


\subsection{学习资源}
\label{\detokenize{intro/index:id5}}\begin{itemize}
\item {} 
\sphinxAtStartPar
{\hyperref[\detokenize{intro/quickstart::doc}]{\sphinxcrossref{\DUrole{std}{\DUrole{std-doc}{快速入门}}}}} — 10 分钟完成第一次数据处理

\item {} 
\sphinxAtStartPar
{\hyperref[\detokenize{intro/install::doc}]{\sphinxcrossref{\DUrole{std}{\DUrole{std-doc}{安装指南}}}}} — 安装与配置

\item {} 
\sphinxAtStartPar
\sphinxhref{https://www.bilibili.com/}{视频教程} (https://www.bilibili.com/) — 配套视频教程（即将上线）

\end{itemize}


\subsection{引用 MTDataPro}
\label{\detokenize{intro/index:id6}}
\sphinxAtStartPar
如果您在科研工作中使用了 MTDataPro，请引用：

\begin{sphinxVerbatim}[commandchars=\\\{\}]
\PYG{n+nc}{@software}\PYG{p}{\PYGZob{}}\PYG{n+nl}{mtdp2026}\PYG{p}{,}
\PYG{+w}{  }\PYG{n+na}{title}\PYG{+w}{ }\PYG{p}{=}\PYG{+w}{ }\PYG{l+s}{\PYGZob{}}\PYG{l+s}{MTDataPro: A Professional Magnetotelluric Data Processing Software}\PYG{l+s}{\PYGZcb{}}\PYG{p}{,}
\PYG{+w}{  }\PYG{n+na}{author}\PYG{+w}{ }\PYG{p}{=}\PYG{+w}{ }\PYG{l+s}{\PYGZob{}}\PYG{l+s}{王培杰}\PYG{l+s}{\PYGZcb{}}\PYG{p}{,}
\PYG{+w}{  }\PYG{n+na}{year}\PYG{+w}{ }\PYG{p}{=}\PYG{+w}{ }\PYG{l+s}{\PYGZob{}}\PYG{l+s}{2026}\PYG{l+s}{\PYGZcb{}}\PYG{p}{,}
\PYG{+w}{  }\PYG{n+na}{version}\PYG{+w}{ }\PYG{p}{=}\PYG{+w}{ }\PYG{l+s}{\PYGZob{}}\PYG{l+s}{1.9.5}\PYG{l+s}{\PYGZcb{}}\PYG{p}{,}
\PYG{+w}{  }\PYG{n+na}{url}\PYG{+w}{ }\PYG{p}{=}\PYG{+w}{ }\PYG{l+s}{\PYGZob{}}\PYG{l+s}{https://github.com/EMWPJ/MTDataPro\PYGZhy{}docs}\PYG{l+s}{\PYGZcb{}}
\PYG{p}{\PYGZcb{}}
\end{sphinxVerbatim}


\subsection{联系方式}
\label{\detokenize{intro/index:id7}}\begin{itemize}
\item {} 
\sphinxAtStartPar
\sphinxstylestrong{技术支持}：wangpj@yangtzeu.edu.cn

\item {} 
\sphinxAtStartPar
\sphinxstylestrong{商务合作}：691195217@qq.com

\item {} 
\sphinxAtStartPar
\sphinxstylestrong{Issue 反馈}：\sphinxhref{https://github.com/EMWPJ/MTDataPro-docs/issues}{GitHub Issues} (https://github.com/EMWPJ/MTDataPro\sphinxhyphen{}docs/issues)

\end{itemize}

\sphinxstepscope


\section{安装}
\label{\detokenize{intro/install:id1}}\label{\detokenize{intro/install::doc}}
\sphinxAtStartPar
MTDataPro 是免安装软件，无需运行安装程序。下载压缩包后解压即可使用。


\subsection{下载软件}
\label{\detokenize{intro/install:id2}}
\sphinxAtStartPar
从以下渠道获取 MTDataPro：
\begin{itemize}
\item {} 
\sphinxAtStartPar
\sphinxstylestrong{官方发布}：联系开发团队获取最新版本

\item {} 
\sphinxAtStartPar
\sphinxstylestrong{版本更新}：关注更新通知获取新版本

\end{itemize}


\subsection{解压软件}
\label{\detokenize{intro/install:id3}}\begin{enumerate}
\sphinxsetlistlabels{\arabic}{enumi}{enumii}{}{.}%
\item {} 
\sphinxAtStartPar
下载 MTDataPro 压缩包（ZIP格式）

\item {} 
\sphinxAtStartPar
解压到目标目录（如 \sphinxcode{\sphinxupquote{D:\textbackslash{}Software\textbackslash{}MTDataPro}} 或 \sphinxcode{\sphinxupquote{C:\textbackslash{}MTDataPro}}）

\item {} 
\sphinxAtStartPar
确保目录路径不包含中文字符或特殊符号

\end{enumerate}


\subsection{目录结构}
\label{\detokenize{intro/install:id4}}
\sphinxAtStartPar
解压后的目录结构如下：

\begin{sphinxVerbatim}[commandchars=\\\{\}]
MTDataPro/
├── MTDP.exe          \PYGZsh{} 主程序
├── MTDP.chm          \PYGZsh{} 帮助文件
├── CalibrationFiles/  \PYGZsh{} 标定文件目录
├── FFTWLibraries/    \PYGZsh{} FFTW库
├── MKLLibraries/     \PYGZsh{} Intel MKL库
├── Samples/          \PYGZsh{} 示例数据
└── docs/             \PYGZsh{} 文档目录
\end{sphinxVerbatim}


\subsection{软件注册}
\label{\detokenize{intro/install:id5}}
\sphinxAtStartPar
首次运行 MTDataPro 需要进行注册授权：


\subsubsection{输入序列号}
\label{\detokenize{intro/install:id6}}\begin{enumerate}
\sphinxsetlistlabels{\arabic}{enumi}{enumii}{}{.}%
\item {} 
\sphinxAtStartPar
双击 \sphinxcode{\sphinxupquote{MTDP.exe}} 启动 MTDataPro

\item {} 
\sphinxAtStartPar
在授权对话框中输入序列号

\item {} 
\sphinxAtStartPar
点击"验证"完成注册

\end{enumerate}

\sphinxAtStartPar
序列号与电脑硬件绑定，每个序列号只能在授权的电脑上使用。


\subsubsection{配置数据路径}
\label{\detokenize{intro/install:id7}}
\sphinxAtStartPar
首次运行时建议设置以下路径：
\begin{itemize}
\item {} 
\sphinxAtStartPar
\sphinxstylestrong{工程目录}：存储工程文件的默认位置

\item {} 
\sphinxAtStartPar
\sphinxstylestrong{数据目录}：原始数据文件的位置

\item {} 
\sphinxAtStartPar
\sphinxstylestrong{临时目录}：处理过程中的临时文件

\end{itemize}


\subsubsection{选择界面语言}
\label{\detokenize{intro/install:id8}}
\sphinxAtStartPar
MTDataPro 支持中英文界面，可在"设置"中切换。


\subsection{外部依赖}
\label{\detokenize{intro/install:id9}}
\sphinxAtStartPar
MTDataPro 依赖以下外部库（已包含在压缩包中）：


\begin{savenotes}\sphinxattablestart
\sphinxthistablewithglobalstyle
\centering
\begin{tabulary}{\linewidth}[t]{TT}
\sphinxtoprule
\begin{varwidth}[t]{\sphinxcolwidth{1}{2}}
\sphinxstyletheadfamily \sphinxAtStartPar
依赖
\sphinxbeforeendvarwidth
\end{varwidth}%
&\begin{varwidth}[t]{\sphinxcolwidth{1}{2}}
\sphinxstyletheadfamily \sphinxAtStartPar
用途
\sphinxbeforeendvarwidth
\end{varwidth}%
\\
\sphinxmidrule
\sphinxtableatstartofbodyhook\begin{varwidth}[t]{\sphinxcolwidth{1}{2}}
\sphinxAtStartPar
Intel MKL
\sphinxbeforeendvarwidth
\end{varwidth}%
&\begin{varwidth}[t]{\sphinxcolwidth{1}{2}}
\sphinxAtStartPar
数学核心库，用于 FFT 计算
\sphinxbeforeendvarwidth
\end{varwidth}%
\\
\sphinxhline\begin{varwidth}[t]{\sphinxcolwidth{1}{2}}
\sphinxAtStartPar
FFTW
\sphinxbeforeendvarwidth
\end{varwidth}%
&\begin{varwidth}[t]{\sphinxcolwidth{1}{2}}
\sphinxAtStartPar
快速傅里叶变换库
\sphinxbeforeendvarwidth
\end{varwidth}%
\\
\sphinxhline\begin{varwidth}[t]{\sphinxcolwidth{1}{2}}
\sphinxAtStartPar
TeeChart
\sphinxbeforeendvarwidth
\end{varwidth}%
&\begin{varwidth}[t]{\sphinxcolwidth{1}{2}}
\sphinxAtStartPar
图表绑制组件
\sphinxbeforeendvarwidth
\end{varwidth}%
\\
\sphinxhline\begin{varwidth}[t]{\sphinxcolwidth{1}{2}}
\sphinxAtStartPar
仪器 DLL
\sphinxbeforeendvarwidth
\end{varwidth}%
&\begin{varwidth}[t]{\sphinxcolwidth{1}{2}}
\sphinxAtStartPar
各仪器专用通信库
\sphinxbeforeendvarwidth
\end{varwidth}%
\\
\sphinxbottomrule
\end{tabulary}
\sphinxtableafterendhook\par
\sphinxattableend\end{savenotes}


\subsection{常见问题}
\label{\detokenize{intro/install:id10}}

\subsubsection{无法启动}
\label{\detokenize{intro/install:id11}}\begin{itemize}
\item {} 
\sphinxAtStartPar
检查是否缺少运行库（VC++ Redistributable）

\item {} 
\sphinxAtStartPar
确保解压路径不包含中文字符

\item {} 
\sphinxAtStartPar
查看日志文件获取错误信息

\end{itemize}


\subsubsection{注册失败}
\label{\detokenize{intro/install:id12}}\begin{itemize}
\item {} 
\sphinxAtStartPar
检查序列号是否正确输入

\item {} 
\sphinxAtStartPar
确认是在授权的电脑上使用

\item {} 
\sphinxAtStartPar
联系技术支持获取帮助

\end{itemize}


\subsection{软件更新}
\label{\detokenize{intro/install:id13}}\begin{enumerate}
\sphinxsetlistlabels{\arabic}{enumi}{enumii}{}{.}%
\item {} 
\sphinxAtStartPar
下载新版本压缩包

\item {} 
\sphinxAtStartPar
关闭正在运行的 MTDataPro

\item {} 
\sphinxAtStartPar
解压到原安装目录，覆盖旧文件

\item {} 
\sphinxAtStartPar
原有的工程文件和授权通常可以继续使用

\end{enumerate}

\sphinxstepscope


\section{快速入门}
\label{\detokenize{intro/quickstart:id1}}\label{\detokenize{intro/quickstart::doc}}
\sphinxAtStartPar
本教程将引导您在 10 分钟内完成第一次 MT 数据处理。


\subsection{基本流程}
\label{\detokenize{intro/quickstart:id2}}
\sphinxAtStartPar
MTDataPro 的数据处理流程如下：


\subsection{步骤 1：创建工程}
\label{\detokenize{intro/quickstart:id3}}\begin{enumerate}
\sphinxsetlistlabels{\arabic}{enumi}{enumii}{}{.}%
\item {} 
\sphinxAtStartPar
启动 MTDataPro

\item {} 
\sphinxAtStartPar
点击 \sphinxstylestrong{文件 → 新建工程}

\item {} 
\sphinxAtStartPar
输入工程名称，选择保存位置

\item {} 
\sphinxAtStartPar
点击"创建"

\end{enumerate}


\subsection{步骤 2：导入数据}
\label{\detokenize{intro/quickstart:id4}}\begin{enumerate}
\sphinxsetlistlabels{\arabic}{enumi}{enumii}{}{.}%
\item {} 
\sphinxAtStartPar
在工程面板中右键点击"测点"

\item {} 
\sphinxAtStartPar
选择"导入数据"

\item {} 
\sphinxAtStartPar
选择仪器类型（如 Phoenix）

\item {} 
\sphinxAtStartPar
浏览并选择数据文件

\item {} 
\sphinxAtStartPar
点击"打开"

\end{enumerate}

\sphinxAtStartPar
MTDataPro 将自动识别采样率、时间范围和通道配置。


\subsection{步骤 3：查看时间序列}
\label{\detokenize{intro/quickstart:id5}}\begin{enumerate}
\sphinxsetlistlabels{\arabic}{enumi}{enumii}{}{.}%
\item {} 
\sphinxAtStartPar
双击工程面板中的测点

\item {} 
\sphinxAtStartPar
时间序列视图将显示原始数据波形

\item {} 
\sphinxAtStartPar
使用工具栏按钮进行缩放、平移

\item {} 
\sphinxAtStartPar
右键菜单可切换通道显示

\end{enumerate}


\subsection{步骤 4：执行 FFT 处理}
\label{\detokenize{intro/quickstart:fft}}\begin{enumerate}
\sphinxsetlistlabels{\arabic}{enumi}{enumii}{}{.}%
\item {} 
\sphinxAtStartPar
选择要处理的时间序列

\item {} 
\sphinxAtStartPar
点击 \sphinxstylestrong{处理 → 新建处理方案}

\item {} 
\sphinxAtStartPar
配置处理参数：
\begin{itemize}
\item {} 
\sphinxAtStartPar
窗口长度

\item {} 
\sphinxAtStartPar
重叠率

\item {} 
\sphinxAtStartPar
窗函数类型

\end{itemize}

\item {} 
\sphinxAtStartPar
点击"运行"开始处理

\end{enumerate}


\subsection{步骤 5：数据筛选}
\label{\detokenize{intro/quickstart:id6}}\begin{enumerate}
\sphinxsetlistlabels{\arabic}{enumi}{enumii}{}{.}%
\item {} 
\sphinxAtStartPar
在傅里叶系数视图中查看处理结果

\item {} 
\sphinxAtStartPar
使用筛选工具移除质量差的数据点

\item {} 
\sphinxAtStartPar
可选：启用参考道处理

\item {} 
\sphinxAtStartPar
可选：调整筛选阈值

\end{enumerate}


\subsection{步骤 6：导出结果}
\label{\detokenize{intro/quickstart:id7}}\begin{enumerate}
\sphinxsetlistlabels{\arabic}{enumi}{enumii}{}{.}%
\item {} 
\sphinxAtStartPar
选择要导出的测点或数据组

\item {} 
\sphinxAtStartPar
点击 \sphinxstylestrong{文件 → 导出}

\item {} 
\sphinxAtStartPar
选择导出格式（EDI、J\sphinxhyphen{}format 等）

\item {} 
\sphinxAtStartPar
设置导出选项

\item {} 
\sphinxAtStartPar
点击"导出"

\end{enumerate}


\subsection{下一步}
\label{\detokenize{intro/quickstart:id8}}
\sphinxAtStartPar
完成快速入门后，您可以：
\begin{itemize}
\item {} 
\sphinxAtStartPar
阅读 {\hyperref[\detokenize{chapters/chapter4::doc}]{\sphinxcrossref{\DUrole{std}{\DUrole{std-doc}{第 4 章：时间序列处理}}}}} 了解更多处理细节

\item {} 
\sphinxAtStartPar
学习 {\hyperref[\detokenize{chapters/chapter5::doc}]{\sphinxcrossref{\DUrole{std}{\DUrole{std-doc}{处理方案配置}}}}} 进行高级处理

\item {} 
\sphinxAtStartPar
查看 {\hyperref[\detokenize{gallery/index::doc}]{\sphinxcrossref{\DUrole{std}{\DUrole{std-doc}{实例展示}}}}} 了解典型应用场景

\end{itemize}

\sphinxstepscope


\section{第1章 软件概述}
\label{\detokenize{chapters/chapter1:id1}}\label{\detokenize{chapters/chapter1::doc}}

\subsection{1.1 MTDP简介}
\label{\detokenize{chapters/chapter1:mtdp}}

\subsubsection{1.1.1 软件定位}
\label{\detokenize{chapters/chapter1:id2}}
\sphinxAtStartPar
MTDP（Magnetotelluric Data Processing）是一款专业的MT数据处理软件，用于大地电磁测深数据的处理和分析。

\sphinxAtStartPar
\sphinxstylestrong{核心功能：}
\begin{itemize}
\item {} 
\sphinxAtStartPar
多仪器数据格式支持

\item {} 
\sphinxAtStartPar
时间序列处理与频谱分析

\item {} 
\sphinxAtStartPar
智能数据筛选与质量评估

\item {} 
\sphinxAtStartPar
MT传递函数估计

\item {} 
\sphinxAtStartPar
数据导出与可视化

\end{itemize}


\subsubsection{1.1.2 版本信息}
\label{\detokenize{chapters/chapter1:id3}}\begin{itemize}
\item {} 
\sphinxAtStartPar
\sphinxstylestrong{当前版本}: v1.9.4

\item {} 
\sphinxAtStartPar
\sphinxstylestrong{发布日期}: 2026年2月2日

\item {} 
\sphinxAtStartPar
\sphinxstylestrong{开发环境}: Embarcadero Delphi 10.4 Sydney

\item {} 
\sphinxAtStartPar
\sphinxstylestrong{运行平台}: Windows 64位

\end{itemize}


\subsubsection{1.1.3 主要特点}
\label{\detokenize{chapters/chapter1:id4}}\begin{enumerate}
\sphinxsetlistlabels{\arabic}{enumi}{enumii}{}{.}%
\item {} 
\sphinxAtStartPar
\sphinxstylestrong{多仪器支持}
\begin{itemize}
\item {} 
\sphinxAtStartPar
Phoenix MTU\sphinxhyphen{}5/5A/5C/8A

\item {} 
\sphinxAtStartPar
Metronix ADU\sphinxhyphen{}06/ADU\sphinxhyphen{}07/MMS

\item {} 
\sphinxAtStartPar
LEMI长周期仪器

\item {} 
\sphinxAtStartPar
RMT/CASRMT格式

\item {} 
\sphinxAtStartPar
Aether (ATTS) 格式

\item {} 
\sphinxAtStartPar
EDI格式导入导出

\end{itemize}

\item {} 
\sphinxAtStartPar
\sphinxstylestrong{智能处理}
\begin{itemize}
\item {} 
\sphinxAtStartPar
Robust稳健估计

\item {} 
\sphinxAtStartPar
多目标遗传算法筛选

\item {} 
\sphinxAtStartPar
马氏距离异常检测

\item {} 
\sphinxAtStartPar
AI参考曲线

\end{itemize}

\item {} 
\sphinxAtStartPar
\sphinxstylestrong{专业工具}
\begin{itemize}
\item {} 
\sphinxAtStartPar
完整标定系统

\item {} 
\sphinxAtStartPar
工频滤波

\item {} 
\sphinxAtStartPar
降采样处理

\item {} 
\sphinxAtStartPar
批量处理

\end{itemize}

\item {} 
\sphinxAtStartPar
\sphinxstylestrong{辅助工具}
\begin{itemize}
\item {} 
\sphinxAtStartPar
时间序列格式转换

\item {} 
\sphinxAtStartPar
标定文件格式转换

\end{itemize}

\end{enumerate}


\bigskip\hrule\bigskip



\subsection{1.2 系统要求与安装}
\label{\detokenize{chapters/chapter1:id5}}

\subsubsection{1.2.1 硬件要求}
\label{\detokenize{chapters/chapter1:id6}}

\begin{savenotes}\sphinxattablestart
\sphinxthistablewithglobalstyle
\centering
\begin{tabulary}{\linewidth}[t]{TTT}
\sphinxtoprule
\begin{varwidth}[t]{\sphinxcolwidth{1}{3}}
\sphinxstyletheadfamily \sphinxAtStartPar
组件
\sphinxbeforeendvarwidth
\end{varwidth}%
&\begin{varwidth}[t]{\sphinxcolwidth{1}{3}}
\sphinxstyletheadfamily \sphinxAtStartPar
最低配置
\sphinxbeforeendvarwidth
\end{varwidth}%
&\begin{varwidth}[t]{\sphinxcolwidth{1}{3}}
\sphinxstyletheadfamily \sphinxAtStartPar
推荐配置
\sphinxbeforeendvarwidth
\end{varwidth}%
\\
\sphinxmidrule
\sphinxtableatstartofbodyhook\begin{varwidth}[t]{\sphinxcolwidth{1}{3}}
\sphinxAtStartPar
CPU
\sphinxbeforeendvarwidth
\end{varwidth}%
&\begin{varwidth}[t]{\sphinxcolwidth{1}{3}}
\sphinxAtStartPar
4核心
\sphinxbeforeendvarwidth
\end{varwidth}%
&\begin{varwidth}[t]{\sphinxcolwidth{1}{3}}
\sphinxAtStartPar
8核心以上
\sphinxbeforeendvarwidth
\end{varwidth}%
\\
\sphinxhline\begin{varwidth}[t]{\sphinxcolwidth{1}{3}}
\sphinxAtStartPar
内存
\sphinxbeforeendvarwidth
\end{varwidth}%
&\begin{varwidth}[t]{\sphinxcolwidth{1}{3}}
\sphinxAtStartPar
8GB
\sphinxbeforeendvarwidth
\end{varwidth}%
&\begin{varwidth}[t]{\sphinxcolwidth{1}{3}}
\sphinxAtStartPar
16GB以上
\sphinxbeforeendvarwidth
\end{varwidth}%
\\
\sphinxhline\begin{varwidth}[t]{\sphinxcolwidth{1}{3}}
\sphinxAtStartPar
硬盘
\sphinxbeforeendvarwidth
\end{varwidth}%
&\begin{varwidth}[t]{\sphinxcolwidth{1}{3}}
\sphinxAtStartPar
10GB可用空间
\sphinxbeforeendvarwidth
\end{varwidth}%
&\begin{varwidth}[t]{\sphinxcolwidth{1}{3}}
\sphinxAtStartPar
SSD, 50GB以上
\sphinxbeforeendvarwidth
\end{varwidth}%
\\
\sphinxhline\begin{varwidth}[t]{\sphinxcolwidth{1}{3}}
\sphinxAtStartPar
显卡
\sphinxbeforeendvarwidth
\end{varwidth}%
&\begin{varwidth}[t]{\sphinxcolwidth{1}{3}}
\sphinxAtStartPar
支持OpenGL 3.0
\sphinxbeforeendvarwidth
\end{varwidth}%
&\begin{varwidth}[t]{\sphinxcolwidth{1}{3}}
\sphinxAtStartPar
独立显卡
\sphinxbeforeendvarwidth
\end{varwidth}%
\\
\sphinxbottomrule
\end{tabulary}
\sphinxtableafterendhook\par
\sphinxattableend\end{savenotes}


\subsubsection{1.2.2 软件要求}
\label{\detokenize{chapters/chapter1:id7}}\begin{itemize}
\item {} 
\sphinxAtStartPar
操作系统：Windows 10/11 64位

\item {} 
\sphinxAtStartPar
运行时库：随软件安装

\end{itemize}


\subsubsection{1.2.3 安装步骤}
\label{\detokenize{chapters/chapter1:id8}}\begin{enumerate}
\sphinxsetlistlabels{\arabic}{enumi}{enumii}{}{.}%
\item {} 
\sphinxAtStartPar
运行安装程序 \sphinxcode{\sphinxupquote{MTDPSetup.exe}}

\item {} 
\sphinxAtStartPar
选择安装目录（建议使用默认路径）

\item {} 
\sphinxAtStartPar
完成安装

\item {} 
\sphinxAtStartPar
首次运行时进行授权

\end{enumerate}


\subsubsection{1.2.4 授权验证}
\label{\detokenize{chapters/chapter1:id9}}
\sphinxAtStartPar
MTDP采用序列号授权机制：
\begin{enumerate}
\sphinxsetlistlabels{\arabic}{enumi}{enumii}{}{.}%
\item {} 
\sphinxAtStartPar
首次启动时，软件会显示硬件码

\item {} 
\sphinxAtStartPar
将硬件码发送给开发者获取授权码

\item {} 
\sphinxAtStartPar
输入授权码完成激活

\end{enumerate}
\begin{quote}

\sphinxAtStartPar
\sphinxstylestrong{注意}: 授权与计算机硬件绑定，更换计算机需要重新获取授权。
\end{quote}


\bigskip\hrule\bigskip



\subsection{1.3 软件界面导览}
\label{\detokenize{chapters/chapter1:id10}}

\subsubsection{1.3.1 主窗口布局}
\label{\detokenize{chapters/chapter1:id11}}
\sphinxAtStartPar
MTDP主界面采用树状结构组织数据，配合多选项卡显示不同视图：

\sphinxAtStartPar
\sphinxstylestrong{主要界面元素：}
\begin{itemize}
\item {} 
\sphinxAtStartPar
\sphinxstylestrong{菜单栏}：位于窗口顶部，包含工程、工具、设置、帮助等主菜单

\item {} 
\sphinxAtStartPar
\sphinxstylestrong{工程树}：左侧面板，显示当前工程的结构层次

\item {} 
\sphinxAtStartPar
\sphinxstylestrong{数据视图}：右侧主区域，显示选中对象的内容

\item {} 
\sphinxAtStartPar
\sphinxstylestrong{状态栏}：底部显示当前状态信息

\end{itemize}


\subsubsection{1.3.2 数据层次结构}
\label{\detokenize{chapters/chapter1:id12}}
\sphinxAtStartPar
MTDP采用四级层次结构组织数据：

\begin{sphinxVerbatim}[commandchars=\\\{\}]
工程 (Project)
└── 工区 (WorkArea)
    └── 测段 (Group)
        └── 测点 (Site)
\end{sphinxVerbatim}

\sphinxAtStartPar
\sphinxstylestrong{各层级功能：}


\begin{savenotes}\sphinxattablestart
\sphinxthistablewithglobalstyle
\centering
\begin{tabulary}{\linewidth}[t]{TTT}
\sphinxtoprule
\begin{varwidth}[t]{\sphinxcolwidth{1}{3}}
\sphinxstyletheadfamily \sphinxAtStartPar
层级
\sphinxbeforeendvarwidth
\end{varwidth}%
&\begin{varwidth}[t]{\sphinxcolwidth{1}{3}}
\sphinxstyletheadfamily \sphinxAtStartPar
说明
\sphinxbeforeendvarwidth
\end{varwidth}%
&\begin{varwidth}[t]{\sphinxcolwidth{1}{3}}
\sphinxstyletheadfamily \sphinxAtStartPar
典型用途
\sphinxbeforeendvarwidth
\end{varwidth}%
\\
\sphinxmidrule
\sphinxtableatstartofbodyhook\begin{varwidth}[t]{\sphinxcolwidth{1}{3}}
\sphinxAtStartPar
工程
\sphinxbeforeendvarwidth
\end{varwidth}%
&\begin{varwidth}[t]{\sphinxcolwidth{1}{3}}
\sphinxAtStartPar
最顶层容器
\sphinxbeforeendvarwidth
\end{varwidth}%
&\begin{varwidth}[t]{\sphinxcolwidth{1}{3}}
\sphinxAtStartPar
整个勘探项目
\sphinxbeforeendvarwidth
\end{varwidth}%
\\
\sphinxhline\begin{varwidth}[t]{\sphinxcolwidth{1}{3}}
\sphinxAtStartPar
工区
\sphinxbeforeendvarwidth
\end{varwidth}%
&\begin{varwidth}[t]{\sphinxcolwidth{1}{3}}
\sphinxAtStartPar
按区域划分
\sphinxbeforeendvarwidth
\end{varwidth}%
&\begin{varwidth}[t]{\sphinxcolwidth{1}{3}}
\sphinxAtStartPar
同一区域的测点集合
\sphinxbeforeendvarwidth
\end{varwidth}%
\\
\sphinxhline\begin{varwidth}[t]{\sphinxcolwidth{1}{3}}
\sphinxAtStartPar
测段
\sphinxbeforeendvarwidth
\end{varwidth}%
&\begin{varwidth}[t]{\sphinxcolwidth{1}{3}}
\sphinxAtStartPar
按时间/任务划分
\sphinxbeforeendvarwidth
\end{varwidth}%
&\begin{varwidth}[t]{\sphinxcolwidth{1}{3}}
\sphinxAtStartPar
同一批处理的测点
\sphinxbeforeendvarwidth
\end{varwidth}%
\\
\sphinxhline\begin{varwidth}[t]{\sphinxcolwidth{1}{3}}
\sphinxAtStartPar
测点
\sphinxbeforeendvarwidth
\end{varwidth}%
&\begin{varwidth}[t]{\sphinxcolwidth{1}{3}}
\sphinxAtStartPar
最小数据单元
\sphinxbeforeendvarwidth
\end{varwidth}%
&\begin{varwidth}[t]{\sphinxcolwidth{1}{3}}
\sphinxAtStartPar
单个观测点数据
\sphinxbeforeendvarwidth
\end{varwidth}%
\\
\sphinxbottomrule
\end{tabulary}
\sphinxtableafterendhook\par
\sphinxattableend\end{savenotes}


\subsubsection{1.3.3 菜单结构}
\label{\detokenize{chapters/chapter1:id13}}
\sphinxAtStartPar
MTDP的菜单系统包含主菜单和右键菜单两部分，涵盖了数据处理的所有功能。


\paragraph{主菜单}
\label{\detokenize{chapters/chapter1:id14}}

\begin{savenotes}
\sphinxatlongtablestart
\sphinxthistablewithglobalstyle
\makeatletter
  \LTleft \@totalleftmargin plus1fill
  \LTright\dimexpr\columnwidth-\@totalleftmargin-\linewidth\relax plus1fill
\makeatother
\begin{longtable}{lll}
\sphinxtoprule
\begin{varwidth}[t]{\sphinxcolwidth{1}{3}}
\sphinxstyletheadfamily \sphinxAtStartPar
菜单
\sphinxbeforeendvarwidth
\end{varwidth}%
&\begin{varwidth}[t]{\sphinxcolwidth{1}{3}}
\sphinxstyletheadfamily \sphinxAtStartPar
功能项
\sphinxbeforeendvarwidth
\end{varwidth}%
&\begin{varwidth}[t]{\sphinxcolwidth{1}{3}}
\sphinxstyletheadfamily \sphinxAtStartPar
说明
\sphinxbeforeendvarwidth
\end{varwidth}%
\\
\sphinxmidrule
\endfirsthead

\multicolumn{3}{c}{\sphinxnorowcolor
    \makebox[0pt]{\sphinxtablecontinued{\tablename\ \thetable{} \textendash{} 接上页}}%
}\\
\sphinxtoprule
\begin{varwidth}[t]{\sphinxcolwidth{1}{3}}
\sphinxstyletheadfamily \sphinxAtStartPar
菜单
\sphinxbeforeendvarwidth
\end{varwidth}%
&\begin{varwidth}[t]{\sphinxcolwidth{1}{3}}
\sphinxstyletheadfamily \sphinxAtStartPar
功能项
\sphinxbeforeendvarwidth
\end{varwidth}%
&\begin{varwidth}[t]{\sphinxcolwidth{1}{3}}
\sphinxstyletheadfamily \sphinxAtStartPar
说明
\sphinxbeforeendvarwidth
\end{varwidth}%
\\
\sphinxmidrule
\endhead

\sphinxbottomrule
\multicolumn{3}{r}{\sphinxnorowcolor
    \makebox[0pt][r]{\sphinxtablecontinued{续下页}}%
}\\
\endfoot

\endlastfoot
\sphinxtableatstartofbodyhook
\begin{varwidth}[t]{\sphinxcolwidth{1}{3}}
\sphinxAtStartPar
\sphinxstylestrong{工程}
\sphinxbeforeendvarwidth
\end{varwidth}%
&\begin{varwidth}[t]{\sphinxcolwidth{1}{3}}
\sphinxAtStartPar
新建工程
\sphinxbeforeendvarwidth
\end{varwidth}%
&\begin{varwidth}[t]{\sphinxcolwidth{1}{3}}
\sphinxAtStartPar
创建新的MTDP工程
\sphinxbeforeendvarwidth
\end{varwidth}%
\\
\sphinxhline\begin{varwidth}[t]{\sphinxcolwidth{1}{3}}
\sphinxAtStartPar

\sphinxbeforeendvarwidth
\end{varwidth}%
&\begin{varwidth}[t]{\sphinxcolwidth{1}{3}}
\sphinxAtStartPar
打开工程
\sphinxbeforeendvarwidth
\end{varwidth}%
&\begin{varwidth}[t]{\sphinxcolwidth{1}{3}}
\sphinxAtStartPar
打开已存在的工程文件
\sphinxbeforeendvarwidth
\end{varwidth}%
\\
\sphinxhline\begin{varwidth}[t]{\sphinxcolwidth{1}{3}}
\sphinxAtStartPar

\sphinxbeforeendvarwidth
\end{varwidth}%
&\begin{varwidth}[t]{\sphinxcolwidth{1}{3}}
\sphinxAtStartPar
重新打开
\sphinxbeforeendvarwidth
\end{varwidth}%
&\begin{varwidth}[t]{\sphinxcolwidth{1}{3}}
\sphinxAtStartPar
打开最近使用的工程
\sphinxbeforeendvarwidth
\end{varwidth}%
\\
\sphinxhline\begin{varwidth}[t]{\sphinxcolwidth{1}{3}}
\sphinxAtStartPar

\sphinxbeforeendvarwidth
\end{varwidth}%
&\begin{varwidth}[t]{\sphinxcolwidth{1}{3}}
\sphinxAtStartPar
保存工程
\sphinxbeforeendvarwidth
\end{varwidth}%
&\begin{varwidth}[t]{\sphinxcolwidth{1}{3}}
\sphinxAtStartPar
保存当前工程
\sphinxbeforeendvarwidth
\end{varwidth}%
\\
\sphinxhline\begin{varwidth}[t]{\sphinxcolwidth{1}{3}}
\sphinxAtStartPar

\sphinxbeforeendvarwidth
\end{varwidth}%
&\begin{varwidth}[t]{\sphinxcolwidth{1}{3}}
\sphinxAtStartPar
另存为
\sphinxbeforeendvarwidth
\end{varwidth}%
&\begin{varwidth}[t]{\sphinxcolwidth{1}{3}}
\sphinxAtStartPar
将工程保存为新文件
\sphinxbeforeendvarwidth
\end{varwidth}%
\\
\sphinxhline\begin{varwidth}[t]{\sphinxcolwidth{1}{3}}
\sphinxAtStartPar

\sphinxbeforeendvarwidth
\end{varwidth}%
&\begin{varwidth}[t]{\sphinxcolwidth{1}{3}}
\sphinxAtStartPar
关闭工程
\sphinxbeforeendvarwidth
\end{varwidth}%
&\begin{varwidth}[t]{\sphinxcolwidth{1}{3}}
\sphinxAtStartPar
关闭当前工程
\sphinxbeforeendvarwidth
\end{varwidth}%
\\
\sphinxhline\begin{varwidth}[t]{\sphinxcolwidth{1}{3}}
\sphinxAtStartPar

\sphinxbeforeendvarwidth
\end{varwidth}%
&\begin{varwidth}[t]{\sphinxcolwidth{1}{3}}
\sphinxAtStartPar
导出版本
\sphinxbeforeendvarwidth
\end{varwidth}%
&\begin{varwidth}[t]{\sphinxcolwidth{1}{3}}
\sphinxAtStartPar
导出工程版本信息
\sphinxbeforeendvarwidth
\end{varwidth}%
\\
\sphinxhline\begin{varwidth}[t]{\sphinxcolwidth{1}{3}}
\sphinxAtStartPar

\sphinxbeforeendvarwidth
\end{varwidth}%
&\begin{varwidth}[t]{\sphinxcolwidth{1}{3}}
\sphinxAtStartPar
导出选中项
\sphinxbeforeendvarwidth
\end{varwidth}%
&\begin{varwidth}[t]{\sphinxcolwidth{1}{3}}
\sphinxAtStartPar
导出选中的测点/测段
\sphinxbeforeendvarwidth
\end{varwidth}%
\\
\sphinxhline\begin{varwidth}[t]{\sphinxcolwidth{1}{3}}
\sphinxAtStartPar

\sphinxbeforeendvarwidth
\end{varwidth}%
&\begin{varwidth}[t]{\sphinxcolwidth{1}{3}}
\sphinxAtStartPar
插入工程
\sphinxbeforeendvarwidth
\end{varwidth}%
&\begin{varwidth}[t]{\sphinxcolwidth{1}{3}}
\sphinxAtStartPar
将其他工程合并到当前工程
\sphinxbeforeendvarwidth
\end{varwidth}%
\\
\sphinxhline\begin{varwidth}[t]{\sphinxcolwidth{1}{3}}
\sphinxAtStartPar

\sphinxbeforeendvarwidth
\end{varwidth}%
&\begin{varwidth}[t]{\sphinxcolwidth{1}{3}}
\sphinxAtStartPar
退出
\sphinxbeforeendvarwidth
\end{varwidth}%
&\begin{varwidth}[t]{\sphinxcolwidth{1}{3}}
\sphinxAtStartPar
退出程序
\sphinxbeforeendvarwidth
\end{varwidth}%
\\
\sphinxhline\begin{varwidth}[t]{\sphinxcolwidth{1}{3}}
\sphinxAtStartPar
\sphinxstylestrong{🔧 工具}
\sphinxbeforeendvarwidth
\end{varwidth}%
&\begin{varwidth}[t]{\sphinxcolwidth{1}{3}}
\sphinxAtStartPar
标定查看
\sphinxbeforeendvarwidth
\end{varwidth}%
&\begin{varwidth}[t]{\sphinxcolwidth{1}{3}}
\sphinxAtStartPar
查看标定文件内容
\sphinxbeforeendvarwidth
\end{varwidth}%
\\
\sphinxhline\begin{varwidth}[t]{\sphinxcolwidth{1}{3}}
\sphinxAtStartPar

\sphinxbeforeendvarwidth
\end{varwidth}%
&\begin{varwidth}[t]{\sphinxcolwidth{1}{3}}
\sphinxAtStartPar
Phoenix TS分割
\sphinxbeforeendvarwidth
\end{varwidth}%
&\begin{varwidth}[t]{\sphinxcolwidth{1}{3}}
\sphinxAtStartPar
分割Phoenix时间序列
\sphinxbeforeendvarwidth
\end{varwidth}%
\\
\sphinxhline\begin{varwidth}[t]{\sphinxcolwidth{1}{3}}
\sphinxAtStartPar

\sphinxbeforeendvarwidth
\end{varwidth}%
&\begin{varwidth}[t]{\sphinxcolwidth{1}{3}}
\sphinxAtStartPar
Phoenix TS修复
\sphinxbeforeendvarwidth
\end{varwidth}%
&\begin{varwidth}[t]{\sphinxcolwidth{1}{3}}
\sphinxAtStartPar
修复损坏的Phoenix时间序列
\sphinxbeforeendvarwidth
\end{varwidth}%
\\
\sphinxhline\begin{varwidth}[t]{\sphinxcolwidth{1}{3}}
\sphinxAtStartPar

\sphinxbeforeendvarwidth
\end{varwidth}%
&\begin{varwidth}[t]{\sphinxcolwidth{1}{3}}
\sphinxAtStartPar
Phoenix TS合并
\sphinxbeforeendvarwidth
\end{varwidth}%
&\begin{varwidth}[t]{\sphinxcolwidth{1}{3}}
\sphinxAtStartPar
合并多个Phoenix时间序列文件
\sphinxbeforeendvarwidth
\end{varwidth}%
\\
\sphinxhline\begin{varwidth}[t]{\sphinxcolwidth{1}{3}}
\sphinxAtStartPar

\sphinxbeforeendvarwidth
\end{varwidth}%
&\begin{varwidth}[t]{\sphinxcolwidth{1}{3}}
\sphinxAtStartPar
Phoenix通道重排
\sphinxbeforeendvarwidth
\end{varwidth}%
&\begin{varwidth}[t]{\sphinxcolwidth{1}{3}}
\sphinxAtStartPar
重新排列Phoenix通道顺序
\sphinxbeforeendvarwidth
\end{varwidth}%
\\
\sphinxhline\begin{varwidth}[t]{\sphinxcolwidth{1}{3}}
\sphinxAtStartPar

\sphinxbeforeendvarwidth
\end{varwidth}%
&\begin{varwidth}[t]{\sphinxcolwidth{1}{3}}
\sphinxAtStartPar
TBL坐标修复
\sphinxbeforeendvarwidth
\end{varwidth}%
&\begin{varwidth}[t]{\sphinxcolwidth{1}{3}}
\sphinxAtStartPar
修复TBL文件中的坐标问题
\sphinxbeforeendvarwidth
\end{varwidth}%
\\
\sphinxhline\begin{varwidth}[t]{\sphinxcolwidth{1}{3}}
\sphinxAtStartPar

\sphinxbeforeendvarwidth
\end{varwidth}%
&\begin{varwidth}[t]{\sphinxcolwidth{1}{3}}
\sphinxAtStartPar
TBL参数替换
\sphinxbeforeendvarwidth
\end{varwidth}%
&\begin{varwidth}[t]{\sphinxcolwidth{1}{3}}
\sphinxAtStartPar
批量替换TBL参数
\sphinxbeforeendvarwidth
\end{varwidth}%
\\
\sphinxhline\begin{varwidth}[t]{\sphinxcolwidth{1}{3}}
\sphinxAtStartPar

\sphinxbeforeendvarwidth
\end{varwidth}%
&\begin{varwidth}[t]{\sphinxcolwidth{1}{3}}
\sphinxAtStartPar
TS查看
\sphinxbeforeendvarwidth
\end{varwidth}%
&\begin{varwidth}[t]{\sphinxcolwidth{1}{3}}
\sphinxAtStartPar
查看时间序列数据
\sphinxbeforeendvarwidth
\end{varwidth}%
\\
\sphinxhline\begin{varwidth}[t]{\sphinxcolwidth{1}{3}}
\sphinxAtStartPar

\sphinxbeforeendvarwidth
\end{varwidth}%
&\begin{varwidth}[t]{\sphinxcolwidth{1}{3}}
\sphinxAtStartPar
时间序列转换
\sphinxbeforeendvarwidth
\end{varwidth}%
&\begin{varwidth}[t]{\sphinxcolwidth{1}{3}}
\sphinxAtStartPar
时间序列格式转换
\sphinxbeforeendvarwidth
\end{varwidth}%
\\
\sphinxhline\begin{varwidth}[t]{\sphinxcolwidth{1}{3}}
\sphinxAtStartPar

\sphinxbeforeendvarwidth
\end{varwidth}%
&\begin{varwidth}[t]{\sphinxcolwidth{1}{3}}
\sphinxAtStartPar
⊢ TXT→TSN
\sphinxbeforeendvarwidth
\end{varwidth}%
&\begin{varwidth}[t]{\sphinxcolwidth{1}{3}}
\sphinxAtStartPar
Phoenix文本格式转TSN
\sphinxbeforeendvarwidth
\end{varwidth}%
\\
\sphinxhline\begin{varwidth}[t]{\sphinxcolwidth{1}{3}}
\sphinxAtStartPar

\sphinxbeforeendvarwidth
\end{varwidth}%
&\begin{varwidth}[t]{\sphinxcolwidth{1}{3}}
\sphinxAtStartPar
⊢ TSN→DAT
\sphinxbeforeendvarwidth
\end{varwidth}%
&\begin{varwidth}[t]{\sphinxcolwidth{1}{3}}
\sphinxAtStartPar
TSN转DAT格式
\sphinxbeforeendvarwidth
\end{varwidth}%
\\
\sphinxhline\begin{varwidth}[t]{\sphinxcolwidth{1}{3}}
\sphinxAtStartPar

\sphinxbeforeendvarwidth
\end{varwidth}%
&\begin{varwidth}[t]{\sphinxcolwidth{1}{3}}
\sphinxAtStartPar
⊢ DAT→TSN
\sphinxbeforeendvarwidth
\end{varwidth}%
&\begin{varwidth}[t]{\sphinxcolwidth{1}{3}}
\sphinxAtStartPar
DAT转TSN格式
\sphinxbeforeendvarwidth
\end{varwidth}%
\\
\sphinxhline\begin{varwidth}[t]{\sphinxcolwidth{1}{3}}
\sphinxAtStartPar

\sphinxbeforeendvarwidth
\end{varwidth}%
&\begin{varwidth}[t]{\sphinxcolwidth{1}{3}}
\sphinxAtStartPar
⊢ TSN→GMT
\sphinxbeforeendvarwidth
\end{varwidth}%
&\begin{varwidth}[t]{\sphinxcolwidth{1}{3}}
\sphinxAtStartPar
TSN转GMT格式
\sphinxbeforeendvarwidth
\end{varwidth}%
\\
\sphinxhline\begin{varwidth}[t]{\sphinxcolwidth{1}{3}}
\sphinxAtStartPar

\sphinxbeforeendvarwidth
\end{varwidth}%
&\begin{varwidth}[t]{\sphinxcolwidth{1}{3}}
\sphinxAtStartPar
⊢ 时间序列→GMT
\sphinxbeforeendvarwidth
\end{varwidth}%
&\begin{varwidth}[t]{\sphinxcolwidth{1}{3}}
\sphinxAtStartPar
通用时间序列转GMT
\sphinxbeforeendvarwidth
\end{varwidth}%
\\
\sphinxhline\begin{varwidth}[t]{\sphinxcolwidth{1}{3}}
\sphinxAtStartPar

\sphinxbeforeendvarwidth
\end{varwidth}%
&\begin{varwidth}[t]{\sphinxcolwidth{1}{3}}
\sphinxAtStartPar
⊢ 时间序列→ATGMT
\sphinxbeforeendvarwidth
\end{varwidth}%
&\begin{varwidth}[t]{\sphinxcolwidth{1}{3}}
\sphinxAtStartPar
时间序列转ATGMT格式
\sphinxbeforeendvarwidth
\end{varwidth}%
\\
\sphinxhline\begin{varwidth}[t]{\sphinxcolwidth{1}{3}}
\sphinxAtStartPar

\sphinxbeforeendvarwidth
\end{varwidth}%
&\begin{varwidth}[t]{\sphinxcolwidth{1}{3}}
\sphinxAtStartPar
Metronix ATS→DAT
\sphinxbeforeendvarwidth
\end{varwidth}%
&\begin{varwidth}[t]{\sphinxcolwidth{1}{3}}
\sphinxAtStartPar
Metronix ATS格式转DAT
\sphinxbeforeendvarwidth
\end{varwidth}%
\\
\sphinxhline\begin{varwidth}[t]{\sphinxcolwidth{1}{3}}
\sphinxAtStartPar

\sphinxbeforeendvarwidth
\end{varwidth}%
&\begin{varwidth}[t]{\sphinxcolwidth{1}{3}}
\sphinxAtStartPar
LEMI测点合并
\sphinxbeforeendvarwidth
\end{varwidth}%
&\begin{varwidth}[t]{\sphinxcolwidth{1}{3}}
\sphinxAtStartPar
合并LEMI格式测点
\sphinxbeforeendvarwidth
\end{varwidth}%
\\
\sphinxhline\begin{varwidth}[t]{\sphinxcolwidth{1}{3}}
\sphinxAtStartPar

\sphinxbeforeendvarwidth
\end{varwidth}%
&\begin{varwidth}[t]{\sphinxcolwidth{1}{3}}
\sphinxAtStartPar
LEMI H自动恢复
\sphinxbeforeendvarwidth
\end{varwidth}%
&\begin{varwidth}[t]{\sphinxcolwidth{1}{3}}
\sphinxAtStartPar
自动恢复LEMI H通道
\sphinxbeforeendvarwidth
\end{varwidth}%
\\
\sphinxhline\begin{varwidth}[t]{\sphinxcolwidth{1}{3}}
\sphinxAtStartPar

\sphinxbeforeendvarwidth
\end{varwidth}%
&\begin{varwidth}[t]{\sphinxcolwidth{1}{3}}
\sphinxAtStartPar
Metronix测点合并
\sphinxbeforeendvarwidth
\end{varwidth}%
&\begin{varwidth}[t]{\sphinxcolwidth{1}{3}}
\sphinxAtStartPar
合并Metronix格式测点
\sphinxbeforeendvarwidth
\end{varwidth}%
\\
\sphinxhline\begin{varwidth}[t]{\sphinxcolwidth{1}{3}}
\sphinxAtStartPar

\sphinxbeforeendvarwidth
\end{varwidth}%
&\begin{varwidth}[t]{\sphinxcolwidth{1}{3}}
\sphinxAtStartPar
WEM数据处理
\sphinxbeforeendvarwidth
\end{varwidth}%
&\begin{varwidth}[t]{\sphinxcolwidth{1}{3}}
\sphinxAtStartPar
WEM格式数据处理
\sphinxbeforeendvarwidth
\end{varwidth}%
\\
\sphinxhline\begin{varwidth}[t]{\sphinxcolwidth{1}{3}}
\sphinxAtStartPar

\sphinxbeforeendvarwidth
\end{varwidth}%
&\begin{varwidth}[t]{\sphinxcolwidth{1}{3}}
\sphinxAtStartPar
EDI修复
\sphinxbeforeendvarwidth
\end{varwidth}%
&\begin{varwidth}[t]{\sphinxcolwidth{1}{3}}
\sphinxAtStartPar
修复损坏的EDI文件
\sphinxbeforeendvarwidth
\end{varwidth}%
\\
\sphinxhline\begin{varwidth}[t]{\sphinxcolwidth{1}{3}}
\sphinxAtStartPar

\sphinxbeforeendvarwidth
\end{varwidth}%
&\begin{varwidth}[t]{\sphinxcolwidth{1}{3}}
\sphinxAtStartPar
EDI转换相干度
\sphinxbeforeendvarwidth
\end{varwidth}%
&\begin{varwidth}[t]{\sphinxcolwidth{1}{3}}
\sphinxAtStartPar
EDI格式相干度转换
\sphinxbeforeendvarwidth
\end{varwidth}%
\\
\sphinxhline\begin{varwidth}[t]{\sphinxcolwidth{1}{3}}
\sphinxAtStartPar

\sphinxbeforeendvarwidth
\end{varwidth}%
&\begin{varwidth}[t]{\sphinxcolwidth{1}{3}}
\sphinxAtStartPar
旧版MT→EDI
\sphinxbeforeendvarwidth
\end{varwidth}%
&\begin{varwidth}[t]{\sphinxcolwidth{1}{3}}
\sphinxAtStartPar
旧版MT数据转EDI格式
\sphinxbeforeendvarwidth
\end{varwidth}%
\\
\sphinxhline\begin{varwidth}[t]{\sphinxcolwidth{1}{3}}
\sphinxAtStartPar

\sphinxbeforeendvarwidth
\end{varwidth}%
&\begin{varwidth}[t]{\sphinxcolwidth{1}{3}}
\sphinxAtStartPar
Phoenix格式转换
\sphinxbeforeendvarwidth
\end{varwidth}%
&\begin{varwidth}[t]{\sphinxcolwidth{1}{3}}
\sphinxAtStartPar
Phoenix数据格式互转
\sphinxbeforeendvarwidth
\end{varwidth}%
\\
\sphinxhline\begin{varwidth}[t]{\sphinxcolwidth{1}{3}}
\sphinxAtStartPar

\sphinxbeforeendvarwidth
\end{varwidth}%
&\begin{varwidth}[t]{\sphinxcolwidth{1}{3}}
\sphinxAtStartPar
⊢ CLB→JSON
\sphinxbeforeendvarwidth
\end{varwidth}%
&\begin{varwidth}[t]{\sphinxcolwidth{1}{3}}
\sphinxAtStartPar
CLB标定转JSON
\sphinxbeforeendvarwidth
\end{varwidth}%
\\
\sphinxhline\begin{varwidth}[t]{\sphinxcolwidth{1}{3}}
\sphinxAtStartPar

\sphinxbeforeendvarwidth
\end{varwidth}%
&\begin{varwidth}[t]{\sphinxcolwidth{1}{3}}
\sphinxAtStartPar
⊢ CLC→JSON
\sphinxbeforeendvarwidth
\end{varwidth}%
&\begin{varwidth}[t]{\sphinxcolwidth{1}{3}}
\sphinxAtStartPar
CLC标定转JSON
\sphinxbeforeendvarwidth
\end{varwidth}%
\\
\sphinxhline\begin{varwidth}[t]{\sphinxcolwidth{1}{3}}
\sphinxAtStartPar

\sphinxbeforeendvarwidth
\end{varwidth}%
&\begin{varwidth}[t]{\sphinxcolwidth{1}{3}}
\sphinxAtStartPar
⊢ JSON→CLB
\sphinxbeforeendvarwidth
\end{varwidth}%
&\begin{varwidth}[t]{\sphinxcolwidth{1}{3}}
\sphinxAtStartPar
JSON转CLB标定
\sphinxbeforeendvarwidth
\end{varwidth}%
\\
\sphinxhline\begin{varwidth}[t]{\sphinxcolwidth{1}{3}}
\sphinxAtStartPar

\sphinxbeforeendvarwidth
\end{varwidth}%
&\begin{varwidth}[t]{\sphinxcolwidth{1}{3}}
\sphinxAtStartPar
⊢ JSON→CLC
\sphinxbeforeendvarwidth
\end{varwidth}%
&\begin{varwidth}[t]{\sphinxcolwidth{1}{3}}
\sphinxAtStartPar
JSON转CLC标定
\sphinxbeforeendvarwidth
\end{varwidth}%
\\
\sphinxhline\begin{varwidth}[t]{\sphinxcolwidth{1}{3}}
\sphinxAtStartPar

\sphinxbeforeendvarwidth
\end{varwidth}%
&\begin{varwidth}[t]{\sphinxcolwidth{1}{3}}
\sphinxAtStartPar
⊢ TBL→JSON
\sphinxbeforeendvarwidth
\end{varwidth}%
&\begin{varwidth}[t]{\sphinxcolwidth{1}{3}}
\sphinxAtStartPar
TBL头文件转JSON
\sphinxbeforeendvarwidth
\end{varwidth}%
\\
\sphinxhline\begin{varwidth}[t]{\sphinxcolwidth{1}{3}}
\sphinxAtStartPar

\sphinxbeforeendvarwidth
\end{varwidth}%
&\begin{varwidth}[t]{\sphinxcolwidth{1}{3}}
\sphinxAtStartPar
⊢ JSON→TBL
\sphinxbeforeendvarwidth
\end{varwidth}%
&\begin{varwidth}[t]{\sphinxcolwidth{1}{3}}
\sphinxAtStartPar
JSON转TBL头文件
\sphinxbeforeendvarwidth
\end{varwidth}%
\\
\sphinxhline\begin{varwidth}[t]{\sphinxcolwidth{1}{3}}
\sphinxAtStartPar
\sphinxstylestrong{时间序列处理}
\sphinxbeforeendvarwidth
\end{varwidth}%
&\begin{varwidth}[t]{\sphinxcolwidth{1}{3}}
\sphinxAtStartPar
SSA/MSSA
\sphinxbeforeendvarwidth
\end{varwidth}%
&\begin{varwidth}[t]{\sphinxcolwidth{1}{3}}
\sphinxAtStartPar
奇异谱分析多窗口
\sphinxbeforeendvarwidth
\end{varwidth}%
\\
\sphinxhline\begin{varwidth}[t]{\sphinxcolwidth{1}{3}}
\sphinxAtStartPar

\sphinxbeforeendvarwidth
\end{varwidth}%
&\begin{varwidth}[t]{\sphinxcolwidth{1}{3}}
\sphinxAtStartPar
ATTS降采样
\sphinxbeforeendvarwidth
\end{varwidth}%
&\begin{varwidth}[t]{\sphinxcolwidth{1}{3}}
\sphinxAtStartPar
ATTS数据降采样
\sphinxbeforeendvarwidth
\end{varwidth}%
\\
\sphinxhline\begin{varwidth}[t]{\sphinxcolwidth{1}{3}}
\sphinxAtStartPar

\sphinxbeforeendvarwidth
\end{varwidth}%
&\begin{varwidth}[t]{\sphinxcolwidth{1}{3}}
\sphinxAtStartPar
MTU降采样
\sphinxbeforeendvarwidth
\end{varwidth}%
&\begin{varwidth}[t]{\sphinxcolwidth{1}{3}}
\sphinxAtStartPar
MTU数据降采样
\sphinxbeforeendvarwidth
\end{varwidth}%
\\
\sphinxhline\begin{varwidth}[t]{\sphinxcolwidth{1}{3}}
\sphinxAtStartPar
\sphinxstylestrong{设置}
\sphinxbeforeendvarwidth
\end{varwidth}%
&\begin{varwidth}[t]{\sphinxcolwidth{1}{3}}
\sphinxAtStartPar
FFT参数
\sphinxbeforeendvarwidth
\end{varwidth}%
&\begin{varwidth}[t]{\sphinxcolwidth{1}{3}}
\sphinxAtStartPar
设置FFT处理参数
\sphinxbeforeendvarwidth
\end{varwidth}%
\\
\sphinxhline\begin{varwidth}[t]{\sphinxcolwidth{1}{3}}
\sphinxAtStartPar

\sphinxbeforeendvarwidth
\end{varwidth}%
&\begin{varwidth}[t]{\sphinxcolwidth{1}{3}}
\sphinxAtStartPar
Phoenix FFT设置
\sphinxbeforeendvarwidth
\end{varwidth}%
&\begin{varwidth}[t]{\sphinxcolwidth{1}{3}}
\sphinxAtStartPar
Phoenix专用FFT参数
\sphinxbeforeendvarwidth
\end{varwidth}%
\\
\sphinxhline\begin{varwidth}[t]{\sphinxcolwidth{1}{3}}
\sphinxAtStartPar

\sphinxbeforeendvarwidth
\end{varwidth}%
&\begin{varwidth}[t]{\sphinxcolwidth{1}{3}}
\sphinxAtStartPar
临时文件路径
\sphinxbeforeendvarwidth
\end{varwidth}%
&\begin{varwidth}[t]{\sphinxcolwidth{1}{3}}
\sphinxAtStartPar
设置临时文件存储位置
\sphinxbeforeendvarwidth
\end{varwidth}%
\\
\sphinxhline\begin{varwidth}[t]{\sphinxcolwidth{1}{3}}
\sphinxAtStartPar

\sphinxbeforeendvarwidth
\end{varwidth}%
&\begin{varwidth}[t]{\sphinxcolwidth{1}{3}}
\sphinxAtStartPar
清理临时文件
\sphinxbeforeendvarwidth
\end{varwidth}%
&\begin{varwidth}[t]{\sphinxcolwidth{1}{3}}
\sphinxAtStartPar
清理缓存的临时文件
\sphinxbeforeendvarwidth
\end{varwidth}%
\\
\sphinxhline\begin{varwidth}[t]{\sphinxcolwidth{1}{3}}
\sphinxAtStartPar

\sphinxbeforeendvarwidth
\end{varwidth}%
&\begin{varwidth}[t]{\sphinxcolwidth{1}{3}}
\sphinxAtStartPar
最大线程数
\sphinxbeforeendvarwidth
\end{varwidth}%
&\begin{varwidth}[t]{\sphinxcolwidth{1}{3}}
\sphinxAtStartPar
设置并行处理线程数
\sphinxbeforeendvarwidth
\end{varwidth}%
\\
\sphinxhline\begin{varwidth}[t]{\sphinxcolwidth{1}{3}}
\sphinxAtStartPar

\sphinxbeforeendvarwidth
\end{varwidth}%
&\begin{varwidth}[t]{\sphinxcolwidth{1}{3}}
\sphinxAtStartPar
设置采集盒标定目录
\sphinxbeforeendvarwidth
\end{varwidth}%
&\begin{varwidth}[t]{\sphinxcolwidth{1}{3}}
\sphinxAtStartPar
配置采集盒标定文件路径
\sphinxbeforeendvarwidth
\end{varwidth}%
\\
\sphinxhline\begin{varwidth}[t]{\sphinxcolwidth{1}{3}}
\sphinxAtStartPar

\sphinxbeforeendvarwidth
\end{varwidth}%
&\begin{varwidth}[t]{\sphinxcolwidth{1}{3}}
\sphinxAtStartPar
设置传感器标定目录
\sphinxbeforeendvarwidth
\end{varwidth}%
&\begin{varwidth}[t]{\sphinxcolwidth{1}{3}}
\sphinxAtStartPar
配置传感器标定文件路径
\sphinxbeforeendvarwidth
\end{varwidth}%
\\
\sphinxhline\begin{varwidth}[t]{\sphinxcolwidth{1}{3}}
\sphinxAtStartPar

\sphinxbeforeendvarwidth
\end{varwidth}%
&\begin{varwidth}[t]{\sphinxcolwidth{1}{3}}
\sphinxAtStartPar
管理标定路径
\sphinxbeforeendvarwidth
\end{varwidth}%
&\begin{varwidth}[t]{\sphinxcolwidth{1}{3}}
\sphinxAtStartPar
管理标定文件搜索路径列表
\sphinxbeforeendvarwidth
\end{varwidth}%
\\
\sphinxhline\begin{varwidth}[t]{\sphinxcolwidth{1}{3}}
\sphinxAtStartPar

\sphinxbeforeendvarwidth
\end{varwidth}%
&\begin{varwidth}[t]{\sphinxcolwidth{1}{3}}
\sphinxAtStartPar
设置系统标定文件
\sphinxbeforeendvarwidth
\end{varwidth}%
&\begin{varwidth}[t]{\sphinxcolwidth{1}{3}}
\sphinxAtStartPar
配置系统级标定文件
\sphinxbeforeendvarwidth
\end{varwidth}%
\\
\sphinxhline\begin{varwidth}[t]{\sphinxcolwidth{1}{3}}
\sphinxAtStartPar

\sphinxbeforeendvarwidth
\end{varwidth}%
&\begin{varwidth}[t]{\sphinxcolwidth{1}{3}}
\sphinxAtStartPar
记录时间可见
\sphinxbeforeendvarwidth
\end{varwidth}%
&\begin{varwidth}[t]{\sphinxcolwidth{1}{3}}
\sphinxAtStartPar
设置时间标记显示
\sphinxbeforeendvarwidth
\end{varwidth}%
\\
\sphinxhline\begin{varwidth}[t]{\sphinxcolwidth{1}{3}}
\sphinxAtStartPar

\sphinxbeforeendvarwidth
\end{varwidth}%
&\begin{varwidth}[t]{\sphinxcolwidth{1}{3}}
\sphinxAtStartPar
语言切换
\sphinxbeforeendvarwidth
\end{varwidth}%
&\begin{varwidth}[t]{\sphinxcolwidth{1}{3}}
\sphinxAtStartPar
切换界面语言
\sphinxbeforeendvarwidth
\end{varwidth}%
\\
\sphinxhline\begin{varwidth}[t]{\sphinxcolwidth{1}{3}}
\sphinxAtStartPar

\sphinxbeforeendvarwidth
\end{varwidth}%
&\begin{varwidth}[t]{\sphinxcolwidth{1}{3}}
\sphinxAtStartPar
⊢ 中文
\sphinxbeforeendvarwidth
\end{varwidth}%
&\begin{varwidth}[t]{\sphinxcolwidth{1}{3}}
\sphinxAtStartPar
简体中文界面
\sphinxbeforeendvarwidth
\end{varwidth}%
\\
\sphinxhline\begin{varwidth}[t]{\sphinxcolwidth{1}{3}}
\sphinxAtStartPar

\sphinxbeforeendvarwidth
\end{varwidth}%
&\begin{varwidth}[t]{\sphinxcolwidth{1}{3}}
\sphinxAtStartPar
⊢ English
\sphinxbeforeendvarwidth
\end{varwidth}%
&\begin{varwidth}[t]{\sphinxcolwidth{1}{3}}
\sphinxAtStartPar
英文界面
\sphinxbeforeendvarwidth
\end{varwidth}%
\\
\sphinxhline\begin{varwidth}[t]{\sphinxcolwidth{1}{3}}
\sphinxAtStartPar

\sphinxbeforeendvarwidth
\end{varwidth}%
&\begin{varwidth}[t]{\sphinxcolwidth{1}{3}}
\sphinxAtStartPar
预测模型
\sphinxbeforeendvarwidth
\end{varwidth}%
&\begin{varwidth}[t]{\sphinxcolwidth{1}{3}}
\sphinxAtStartPar
AI预测模型设置
\sphinxbeforeendvarwidth
\end{varwidth}%
\\
\sphinxhline\begin{varwidth}[t]{\sphinxcolwidth{1}{3}}
\sphinxAtStartPar

\sphinxbeforeendvarwidth
\end{varwidth}%
&\begin{varwidth}[t]{\sphinxcolwidth{1}{3}}
\sphinxAtStartPar
使用FFTW3
\sphinxbeforeendvarwidth
\end{varwidth}%
&\begin{varwidth}[t]{\sphinxcolwidth{1}{3}}
\sphinxAtStartPar
启用FFTW3加速库
\sphinxbeforeendvarwidth
\end{varwidth}%
\\
\sphinxhline\begin{varwidth}[t]{\sphinxcolwidth{1}{3}}
\sphinxAtStartPar

\sphinxbeforeendvarwidth
\end{varwidth}%
&\begin{varwidth}[t]{\sphinxcolwidth{1}{3}}
\sphinxAtStartPar
坐标轴设置
\sphinxbeforeendvarwidth
\end{varwidth}%
&\begin{varwidth}[t]{\sphinxcolwidth{1}{3}}
\sphinxAtStartPar
设置图表坐标轴范围
\sphinxbeforeendvarwidth
\end{varwidth}%
\\
\sphinxhline\begin{varwidth}[t]{\sphinxcolwidth{1}{3}}
\sphinxAtStartPar
\sphinxstylestrong{帮助}
\sphinxbeforeendvarwidth
\end{varwidth}%
&\begin{varwidth}[t]{\sphinxcolwidth{1}{3}}
\sphinxAtStartPar
关于
\sphinxbeforeendvarwidth
\end{varwidth}%
&\begin{varwidth}[t]{\sphinxcolwidth{1}{3}}
\sphinxAtStartPar
查看软件版本和授权信息
\sphinxbeforeendvarwidth
\end{varwidth}%
\\
\sphinxbottomrule
\end{longtable}
\sphinxtableafterendhook
\sphinxatlongtableend
\end{savenotes}


\bigskip\hrule\bigskip



\paragraph{右键菜单详细说明}
\label{\detokenize{chapters/chapter1:id15}}
\sphinxAtStartPar
右键菜单是MTDP中最常用的操作方式，根据选中的对象类型显示不同的操作选项。MTDP提供了四个层级的右键菜单：\sphinxstylestrong{项目}、\sphinxstylestrong{工区}、\sphinxstylestrong{分组}和\sphinxstylestrong{测点}。


\subparagraph{项目右键菜单}
\label{\detokenize{chapters/chapter1:id16}}
\sphinxAtStartPar
项目右键菜单用于管理整个工程项目的显示方式。


\begin{savenotes}\sphinxattablestart
\sphinxthistablewithglobalstyle
\centering
\begin{tabulary}{\linewidth}[t]{TTTT}
\sphinxtoprule
\begin{varwidth}[t]{\sphinxcolwidth{1}{4}}
\sphinxstyletheadfamily \sphinxAtStartPar
功能
\sphinxbeforeendvarwidth
\end{varwidth}%
&\begin{varwidth}[t]{\sphinxcolwidth{1}{4}}
\sphinxstyletheadfamily \sphinxAtStartPar
说明
\sphinxbeforeendvarwidth
\end{varwidth}%
&\begin{varwidth}[t]{\sphinxcolwidth{1}{4}}
\sphinxstyletheadfamily \sphinxAtStartPar
使用场景
\sphinxbeforeendvarwidth
\end{varwidth}%
&\begin{varwidth}[t]{\sphinxcolwidth{1}{4}}
\sphinxstyletheadfamily \sphinxAtStartPar
操作步骤
\sphinxbeforeendvarwidth
\end{varwidth}%
\\
\sphinxmidrule
\sphinxtableatstartofbodyhook\begin{varwidth}[t]{\sphinxcolwidth{1}{4}}
\sphinxAtStartPar
复选框显示
\sphinxbeforeendvarwidth
\end{varwidth}%
&\begin{varwidth}[t]{\sphinxcolwidth{1}{4}}
\sphinxAtStartPar
切换项目树的复选框显示状态
\sphinxbeforeendvarwidth
\end{varwidth}%
&\begin{varwidth}[t]{\sphinxcolwidth{1}{4}}
\sphinxAtStartPar
当需要批量选择多个测点进行操作时
\sphinxbeforeendvarwidth
\end{varwidth}%
&\begin{varwidth}[t]{\sphinxcolwidth{1}{4}}
\sphinxAtStartPar
1. 右键项目名称2. 点击"复选框显示"3. 项目树中会出现复选框
\sphinxbeforeendvarwidth
\end{varwidth}%
\\
\sphinxhline\begin{varwidth}[t]{\sphinxcolwidth{1}{4}}
\sphinxAtStartPar
全部展开
\sphinxbeforeendvarwidth
\end{varwidth}%
&\begin{varwidth}[t]{\sphinxcolwidth{1}{4}}
\sphinxAtStartPar
展开所有节点，显示完整项目结构
\sphinxbeforeendvarwidth
\end{varwidth}%
&\begin{varwidth}[t]{\sphinxcolwidth{1}{4}}
\sphinxAtStartPar
当需要查看所有工区、分组和测点时
\sphinxbeforeendvarwidth
\end{varwidth}%
&\begin{varwidth}[t]{\sphinxcolwidth{1}{4}}
\sphinxAtStartPar
1. 右键项目名称2. 点击"全部展开"3. 所有节点都会展开
\sphinxbeforeendvarwidth
\end{varwidth}%
\\
\sphinxhline\begin{varwidth}[t]{\sphinxcolwidth{1}{4}}
\sphinxAtStartPar
全部折叠
\sphinxbeforeendvarwidth
\end{varwidth}%
&\begin{varwidth}[t]{\sphinxcolwidth{1}{4}}
\sphinxAtStartPar
折叠所有节点，只显示顶级项目
\sphinxbeforeendvarwidth
\end{varwidth}%
&\begin{varwidth}[t]{\sphinxcolwidth{1}{4}}
\sphinxAtStartPar
当项目结构复杂，需要简化视图时
\sphinxbeforeendvarwidth
\end{varwidth}%
&\begin{varwidth}[t]{\sphinxcolwidth{1}{4}}
\sphinxAtStartPar
1. 右键项目名称2. 点击"全部折叠"3. 所有子节点都会隐藏
\sphinxbeforeendvarwidth
\end{varwidth}%
\\
\sphinxbottomrule
\end{tabulary}
\sphinxtableafterendhook\par
\sphinxattableend\end{savenotes}
\begin{quote}

\sphinxAtStartPar
\sphinxstylestrong{提示}：项目树的展开/折叠状态会被保存，下次打开工程时会恢复。
\end{quote}


\subparagraph{工区右键菜单}
\label{\detokenize{chapters/chapter1:id17}}
\sphinxAtStartPar
工区（WorkArea）是同一区域内测点的集合，工区右键菜单包含以下功能分类：

\sphinxAtStartPar
\sphinxstylestrong{1. 测段管理}


\begin{savenotes}\sphinxattablestart
\sphinxthistablewithglobalstyle
\centering
\begin{tabulary}{\linewidth}[t]{TTTTT}
\sphinxtoprule
\begin{varwidth}[t]{\sphinxcolwidth{1}{5}}
\sphinxstyletheadfamily \sphinxAtStartPar
功能
\sphinxbeforeendvarwidth
\end{varwidth}%
&\begin{varwidth}[t]{\sphinxcolwidth{1}{5}}
\sphinxstyletheadfamily \sphinxAtStartPar
说明
\sphinxbeforeendvarwidth
\end{varwidth}%
&\begin{varwidth}[t]{\sphinxcolwidth{1}{5}}
\sphinxstyletheadfamily \sphinxAtStartPar
使用场景
\sphinxbeforeendvarwidth
\end{varwidth}%
&\begin{varwidth}[t]{\sphinxcolwidth{1}{5}}
\sphinxstyletheadfamily \sphinxAtStartPar
操作步骤
\sphinxbeforeendvarwidth
\end{varwidth}%
&\begin{varwidth}[t]{\sphinxcolwidth{1}{5}}
\sphinxstyletheadfamily \sphinxAtStartPar
注意事项
\sphinxbeforeendvarwidth
\end{varwidth}%
\\
\sphinxmidrule
\sphinxtableatstartofbodyhook\begin{varwidth}[t]{\sphinxcolwidth{1}{5}}
\sphinxAtStartPar
新建测段
\sphinxbeforeendvarwidth
\end{varwidth}%
&\begin{varwidth}[t]{\sphinxcolwidth{1}{5}}
\sphinxAtStartPar
在工区下创建新的测段
\sphinxbeforeendvarwidth
\end{varwidth}%
&\begin{varwidth}[t]{\sphinxcolwidth{1}{5}}
\sphinxAtStartPar
当需要按时间、任务或测线组织测点时
\sphinxbeforeendvarwidth
\end{varwidth}%
&\begin{varwidth}[t]{\sphinxcolwidth{1}{5}}
\sphinxAtStartPar
1. 右键工区名称2. 点击"新建测段"3. 输入测段名称4. 设置测段属性
\sphinxbeforeendvarwidth
\end{varwidth}%
&\begin{varwidth}[t]{\sphinxcolwidth{1}{5}}
\sphinxAtStartPar
测段名称建议包含日期或测线编号
\sphinxbeforeendvarwidth
\end{varwidth}%
\\
\sphinxhline\begin{varwidth}[t]{\sphinxcolwidth{1}{5}}
\sphinxAtStartPar
加载测段
\sphinxbeforeendvarwidth
\end{varwidth}%
&\begin{varwidth}[t]{\sphinxcolwidth{1}{5}}
\sphinxAtStartPar
从已有文件加载测段配置
\sphinxbeforeendvarwidth
\end{varwidth}%
&\begin{varwidth}[t]{\sphinxcolwidth{1}{5}}
\sphinxAtStartPar
当需要复用之前的测段设置时
\sphinxbeforeendvarwidth
\end{varwidth}%
&\begin{varwidth}[t]{\sphinxcolwidth{1}{5}}
\sphinxAtStartPar
1. 右键工区名称2. 点击"加载测段"3. 选择测段配置文件
\sphinxbeforeendvarwidth
\end{varwidth}%
&\begin{varwidth}[t]{\sphinxcolwidth{1}{5}}
\sphinxAtStartPar
仅加载配置，不包含测点数据
\sphinxbeforeendvarwidth
\end{varwidth}%
\\
\sphinxhline\begin{varwidth}[t]{\sphinxcolwidth{1}{5}}
\sphinxAtStartPar
粘贴测段
\sphinxbeforeendvarwidth
\end{varwidth}%
&\begin{varwidth}[t]{\sphinxcolwidth{1}{5}}
\sphinxAtStartPar
粘贴已复制的测段
\sphinxbeforeendvarwidth
\end{varwidth}%
&\begin{varwidth}[t]{\sphinxcolwidth{1}{5}}
\sphinxAtStartPar
当需要在工区之间复制测段时
\sphinxbeforeendvarwidth
\end{varwidth}%
&\begin{varwidth}[t]{\sphinxcolwidth{1}{5}}
\sphinxAtStartPar
1. 先在源工区复制测段2. 右键目标工区3. 点击"粘贴测段"
\sphinxbeforeendvarwidth
\end{varwidth}%
&\begin{varwidth}[t]{\sphinxcolwidth{1}{5}}
\sphinxAtStartPar
测点数据会同时复制
\sphinxbeforeendvarwidth
\end{varwidth}%
\\
\sphinxbottomrule
\end{tabulary}
\sphinxtableafterendhook\par
\sphinxattableend\end{savenotes}

\sphinxAtStartPar
\sphinxstylestrong{2. 工区管理}


\begin{savenotes}\sphinxattablestart
\sphinxthistablewithglobalstyle
\centering
\begin{tabulary}{\linewidth}[t]{TTTTT}
\sphinxtoprule
\begin{varwidth}[t]{\sphinxcolwidth{1}{5}}
\sphinxstyletheadfamily \sphinxAtStartPar
功能
\sphinxbeforeendvarwidth
\end{varwidth}%
&\begin{varwidth}[t]{\sphinxcolwidth{1}{5}}
\sphinxstyletheadfamily \sphinxAtStartPar
说明
\sphinxbeforeendvarwidth
\end{varwidth}%
&\begin{varwidth}[t]{\sphinxcolwidth{1}{5}}
\sphinxstyletheadfamily \sphinxAtStartPar
使用场景
\sphinxbeforeendvarwidth
\end{varwidth}%
&\begin{varwidth}[t]{\sphinxcolwidth{1}{5}}
\sphinxstyletheadfamily \sphinxAtStartPar
操作步骤
\sphinxbeforeendvarwidth
\end{varwidth}%
&\begin{varwidth}[t]{\sphinxcolwidth{1}{5}}
\sphinxstyletheadfamily \sphinxAtStartPar
注意事项
\sphinxbeforeendvarwidth
\end{varwidth}%
\\
\sphinxmidrule
\sphinxtableatstartofbodyhook\begin{varwidth}[t]{\sphinxcolwidth{1}{5}}
\sphinxAtStartPar
编辑
\sphinxbeforeendvarwidth
\end{varwidth}%
&\begin{varwidth}[t]{\sphinxcolwidth{1}{5}}
\sphinxAtStartPar
编辑工区属性（名称、描述等）
\sphinxbeforeendvarwidth
\end{varwidth}%
&\begin{varwidth}[t]{\sphinxcolwidth{1}{5}}
\sphinxAtStartPar
当需要修改工区基本信息时
\sphinxbeforeendvarwidth
\end{varwidth}%
&\begin{varwidth}[t]{\sphinxcolwidth{1}{5}}
\sphinxAtStartPar
1. 右键工区名称2. 点击"编辑"3. 在属性面板中修改信息
\sphinxbeforeendvarwidth
\end{varwidth}%
&\begin{varwidth}[t]{\sphinxcolwidth{1}{5}}
\sphinxAtStartPar
修改后需要保存工程
\sphinxbeforeendvarwidth
\end{varwidth}%
\\
\sphinxhline\begin{varwidth}[t]{\sphinxcolwidth{1}{5}}
\sphinxAtStartPar
删除
\sphinxbeforeendvarwidth
\end{varwidth}%
&\begin{varwidth}[t]{\sphinxcolwidth{1}{5}}
\sphinxAtStartPar
删除整个工区及其所有数据
\sphinxbeforeendvarwidth
\end{varwidth}%
&\begin{varwidth}[t]{\sphinxcolwidth{1}{5}}
\sphinxAtStartPar
当工区不再需要时
\sphinxbeforeendvarwidth
\end{varwidth}%
&\begin{varwidth}[t]{\sphinxcolwidth{1}{5}}
\sphinxAtStartPar
1. 右键工区名称2. 点击"删除"3. 确认删除操作
\sphinxbeforeendvarwidth
\end{varwidth}%
&\begin{varwidth}[t]{\sphinxcolwidth{1}{5}}
\sphinxAtStartPar
⚠️ 删除后无法恢复，请谨慎操作
\sphinxbeforeendvarwidth
\end{varwidth}%
\\
\sphinxhline\begin{varwidth}[t]{\sphinxcolwidth{1}{5}}
\sphinxAtStartPar
复制
\sphinxbeforeendvarwidth
\end{varwidth}%
&\begin{varwidth}[t]{\sphinxcolwidth{1}{5}}
\sphinxAtStartPar
复制工区（包含所有测段和测点）
\sphinxbeforeendvarwidth
\end{varwidth}%
&\begin{varwidth}[t]{\sphinxcolwidth{1}{5}}
\sphinxAtStartPar
当需要创建相似的工区时
\sphinxbeforeendvarwidth
\end{varwidth}%
&\begin{varwidth}[t]{\sphinxcolwidth{1}{5}}
\sphinxAtStartPar
1. 右键工区名称2. 点击"复制"3. 右键目标位置粘贴
\sphinxbeforeendvarwidth
\end{varwidth}%
&\begin{varwidth}[t]{\sphinxcolwidth{1}{5}}
\sphinxAtStartPar
复制后建议重命名以区分
\sphinxbeforeendvarwidth
\end{varwidth}%
\\
\sphinxbottomrule
\end{tabulary}
\sphinxtableafterendhook\par
\sphinxattableend\end{savenotes}

\sphinxAtStartPar
\sphinxstylestrong{3. 批量处理}


\begin{savenotes}\sphinxattablestart
\sphinxthistablewithglobalstyle
\centering
\begin{tabulary}{\linewidth}[t]{TTTTT}
\sphinxtoprule
\begin{varwidth}[t]{\sphinxcolwidth{1}{5}}
\sphinxstyletheadfamily \sphinxAtStartPar
功能
\sphinxbeforeendvarwidth
\end{varwidth}%
&\begin{varwidth}[t]{\sphinxcolwidth{1}{5}}
\sphinxstyletheadfamily \sphinxAtStartPar
说明
\sphinxbeforeendvarwidth
\end{varwidth}%
&\begin{varwidth}[t]{\sphinxcolwidth{1}{5}}
\sphinxstyletheadfamily \sphinxAtStartPar
使用场景
\sphinxbeforeendvarwidth
\end{varwidth}%
&\begin{varwidth}[t]{\sphinxcolwidth{1}{5}}
\sphinxstyletheadfamily \sphinxAtStartPar
操作步骤
\sphinxbeforeendvarwidth
\end{varwidth}%
&\begin{varwidth}[t]{\sphinxcolwidth{1}{5}}
\sphinxstyletheadfamily \sphinxAtStartPar
注意事项
\sphinxbeforeendvarwidth
\end{varwidth}%
\\
\sphinxmidrule
\sphinxtableatstartofbodyhook\begin{varwidth}[t]{\sphinxcolwidth{1}{5}}
\sphinxAtStartPar
处理
\sphinxbeforeendvarwidth
\end{varwidth}%
&\begin{varwidth}[t]{\sphinxcolwidth{1}{5}}
\sphinxAtStartPar
批量处理工区内所有测段的数据
\sphinxbeforeendvarwidth
\end{varwidth}%
&\begin{varwidth}[t]{\sphinxcolwidth{1}{5}}
\sphinxAtStartPar
当需要对整个工区进行统一处理时
\sphinxbeforeendvarwidth
\end{varwidth}%
&\begin{varwidth}[t]{\sphinxcolwidth{1}{5}}
\sphinxAtStartPar
1. 右键工区名称2. 点击"处理"3. 设置处理参数4. 执行处理
\sphinxbeforeendvarwidth
\end{varwidth}%
&\begin{varwidth}[t]{\sphinxcolwidth{1}{5}}
\sphinxAtStartPar
处理时间取决于测点数量和计算机性能
\sphinxbeforeendvarwidth
\end{varwidth}%
\\
\sphinxhline\begin{varwidth}[t]{\sphinxcolwidth{1}{5}}
\sphinxAtStartPar
Phoenix TS转FT
\sphinxbeforeendvarwidth
\end{varwidth}%
&\begin{varwidth}[t]{\sphinxcolwidth{1}{5}}
\sphinxAtStartPar
将Phoenix时间序列转换为傅里叶系数
\sphinxbeforeendvarwidth
\end{varwidth}%
&\begin{varwidth}[t]{\sphinxcolwidth{1}{5}}
\sphinxAtStartPar
当使用Phoenix仪器且需要频谱数据时
\sphinxbeforeendvarwidth
\end{varwidth}%
&\begin{varwidth}[t]{\sphinxcolwidth{1}{5}}
\sphinxAtStartPar
1. 右键工区名称2. 点击"Phoenix TS转FT"3. 设置FFT参数4. 执行转换
\sphinxbeforeendvarwidth
\end{varwidth}%
&\begin{varwidth}[t]{\sphinxcolwidth{1}{5}}
\sphinxAtStartPar
需要先导入Phoenix时间序列数据
\sphinxbeforeendvarwidth
\end{varwidth}%
\\
\sphinxbottomrule
\end{tabulary}
\sphinxtableafterendhook\par
\sphinxattableend\end{savenotes}

\sphinxAtStartPar
\sphinxstylestrong{4. 数据导出}


\begin{savenotes}\sphinxattablestart
\sphinxthistablewithglobalstyle
\centering
\begin{tabulary}{\linewidth}[t]{TTTTT}
\sphinxtoprule
\begin{varwidth}[t]{\sphinxcolwidth{1}{5}}
\sphinxstyletheadfamily \sphinxAtStartPar
功能
\sphinxbeforeendvarwidth
\end{varwidth}%
&\begin{varwidth}[t]{\sphinxcolwidth{1}{5}}
\sphinxstyletheadfamily \sphinxAtStartPar
说明
\sphinxbeforeendvarwidth
\end{varwidth}%
&\begin{varwidth}[t]{\sphinxcolwidth{1}{5}}
\sphinxstyletheadfamily \sphinxAtStartPar
使用场景
\sphinxbeforeendvarwidth
\end{varwidth}%
&\begin{varwidth}[t]{\sphinxcolwidth{1}{5}}
\sphinxstyletheadfamily \sphinxAtStartPar
输出文件
\sphinxbeforeendvarwidth
\end{varwidth}%
&\begin{varwidth}[t]{\sphinxcolwidth{1}{5}}
\sphinxstyletheadfamily \sphinxAtStartPar
注意事项
\sphinxbeforeendvarwidth
\end{varwidth}%
\\
\sphinxmidrule
\sphinxtableatstartofbodyhook\begin{varwidth}[t]{\sphinxcolwidth{1}{5}}
\sphinxAtStartPar
导出参数
\sphinxbeforeendvarwidth
\end{varwidth}%
&\begin{varwidth}[t]{\sphinxcolwidth{1}{5}}
\sphinxAtStartPar
导出工区内所有测点的处理参数
\sphinxbeforeendvarwidth
\end{varwidth}%
&\begin{varwidth}[t]{\sphinxcolwidth{1}{5}}
\sphinxAtStartPar
当需要汇总所有测点的MT参数时
\sphinxbeforeendvarwidth
\end{varwidth}%
&\begin{varwidth}[t]{\sphinxcolwidth{1}{5}}
\sphinxAtStartPar
文本格式的参数文件
\sphinxbeforeendvarwidth
\end{varwidth}%
&\begin{varwidth}[t]{\sphinxcolwidth{1}{5}}
\sphinxAtStartPar
包含视电阻率、阻抗、相位等
\sphinxbeforeendvarwidth
\end{varwidth}%
\\
\sphinxhline\begin{varwidth}[t]{\sphinxcolwidth{1}{5}}
\sphinxAtStartPar
导出测点坐标
\sphinxbeforeendvarwidth
\end{varwidth}%
&\begin{varwidth}[t]{\sphinxcolwidth{1}{5}}
\sphinxAtStartPar
导出所有测点的地理坐标
\sphinxbeforeendvarwidth
\end{varwidth}%
&\begin{varwidth}[t]{\sphinxcolwidth{1}{5}}
\sphinxAtStartPar
当需要制作测点分布图或导入GIS时
\sphinxbeforeendvarwidth
\end{varwidth}%
&\begin{varwidth}[t]{\sphinxcolwidth{1}{5}}
\sphinxAtStartPar
TXT或CSV格式的坐标文件
\sphinxbeforeendvarwidth
\end{varwidth}%
&\begin{varwidth}[t]{\sphinxcolwidth{1}{5}}
\sphinxAtStartPar
包含经度、纬度、高程
\sphinxbeforeendvarwidth
\end{varwidth}%
\\
\sphinxhline\begin{varwidth}[t]{\sphinxcolwidth{1}{5}}
\sphinxAtStartPar
导出测点信息
\sphinxbeforeendvarwidth
\end{varwidth}%
&\begin{varwidth}[t]{\sphinxcolwidth{1}{5}}
\sphinxAtStartPar
导出测点的详细信息（含采集参数）
\sphinxbeforeendvarwidth
\end{varwidth}%
&\begin{varwidth}[t]{\sphinxcolwidth{1}{5}}
\sphinxAtStartPar
当需要完整的测点元数据时
\sphinxbeforeendvarwidth
\end{varwidth}%
&\begin{varwidth}[t]{\sphinxcolwidth{1}{5}}
\sphinxAtStartPar
文本格式的信息文件
\sphinxbeforeendvarwidth
\end{varwidth}%
&\begin{varwidth}[t]{\sphinxcolwidth{1}{5}}
\sphinxAtStartPar
包含测点名称、坐标、时间、仪器等
\sphinxbeforeendvarwidth
\end{varwidth}%
\\
\sphinxhline\begin{varwidth}[t]{\sphinxcolwidth{1}{5}}
\sphinxAtStartPar
导出到项目
\sphinxbeforeendvarwidth
\end{varwidth}%
&\begin{varwidth}[t]{\sphinxcolwidth{1}{5}}
\sphinxAtStartPar
将工区导出为独立的MTDP项目
\sphinxbeforeendvarwidth
\end{varwidth}%
&\begin{varwidth}[t]{\sphinxcolwidth{1}{5}}
\sphinxAtStartPar
当需要将工区独立出来作为新项目时
\sphinxbeforeendvarwidth
\end{varwidth}%
&\begin{varwidth}[t]{\sphinxcolwidth{1}{5}}
\sphinxAtStartPar
新的.MTDPE工程文件
\sphinxbeforeendvarwidth
\end{varwidth}%
&\begin{varwidth}[t]{\sphinxcolwidth{1}{5}}
\sphinxAtStartPar
导出后可单独打开和管理
\sphinxbeforeendvarwidth
\end{varwidth}%
\\
\sphinxhline\begin{varwidth}[t]{\sphinxcolwidth{1}{5}}
\sphinxAtStartPar
导出KMZ
\sphinxbeforeendvarwidth
\end{varwidth}%
&\begin{varwidth}[t]{\sphinxcolwidth{1}{5}}
\sphinxAtStartPar
导出Google Earth KMZ格式
\sphinxbeforeendvarwidth
\end{varwidth}%
&\begin{varwidth}[t]{\sphinxcolwidth{1}{5}}
\sphinxAtStartPar
当需要在Google Earth中查看测点分布时
\sphinxbeforeendvarwidth
\end{varwidth}%
&\begin{varwidth}[t]{\sphinxcolwidth{1}{5}}
\sphinxAtStartPar
.kmz压缩文件
\sphinxbeforeendvarwidth
\end{varwidth}%
&\begin{varwidth}[t]{\sphinxcolwidth{1}{5}}
\sphinxAtStartPar
包含测点位置、名称和基本信息
\sphinxbeforeendvarwidth
\end{varwidth}%
\\
\sphinxhline\begin{varwidth}[t]{\sphinxcolwidth{1}{5}}
\sphinxAtStartPar
导出OmapKML
\sphinxbeforeendvarwidth
\end{varwidth}%
&\begin{varwidth}[t]{\sphinxcolwidth{1}{5}}
\sphinxAtStartPar
导出Omap专用KML格式
\sphinxbeforeendvarwidth
\end{varwidth}%
&\begin{varwidth}[t]{\sphinxcolwidth{1}{5}}
\sphinxAtStartPar
当使用Omap软件进行成图时
\sphinxbeforeendvarwidth
\end{varwidth}%
&\begin{varwidth}[t]{\sphinxcolwidth{1}{5}}
\sphinxAtStartPar
.kml文件
\sphinxbeforeendvarwidth
\end{varwidth}%
&\begin{varwidth}[t]{\sphinxcolwidth{1}{5}}
\sphinxAtStartPar
包含Omap专用的样式和属性
\sphinxbeforeendvarwidth
\end{varwidth}%
\\
\sphinxhline\begin{varwidth}[t]{\sphinxcolwidth{1}{5}}
\sphinxAtStartPar
导出普通KML
\sphinxbeforeendvarwidth
\end{varwidth}%
&\begin{varwidth}[t]{\sphinxcolwidth{1}{5}}
\sphinxAtStartPar
导出标准KML格式
\sphinxbeforeendvarwidth
\end{varwidth}%
&\begin{varwidth}[t]{\sphinxcolwidth{1}{5}}
\sphinxAtStartPar
当需要与其他GIS软件交换数据时
\sphinxbeforeendvarwidth
\end{varwidth}%
&\begin{varwidth}[t]{\sphinxcolwidth{1}{5}}
\sphinxAtStartPar
.kml文件
\sphinxbeforeendvarwidth
\end{varwidth}%
&\begin{varwidth}[t]{\sphinxcolwidth{1}{5}}
\sphinxAtStartPar
标准KML格式，兼容性好
\sphinxbeforeendvarwidth
\end{varwidth}%
\\
\sphinxhline\begin{varwidth}[t]{\sphinxcolwidth{1}{5}}
\sphinxAtStartPar
导出EDI
\sphinxbeforeendvarwidth
\end{varwidth}%
&\begin{varwidth}[t]{\sphinxcolwidth{1}{5}}
\sphinxAtStartPar
批量导出所有测点的EDI文件
\sphinxbeforeendvarwidth
\end{varwidth}%
&\begin{varwidth}[t]{\sphinxcolwidth{1}{5}}
\sphinxAtStartPar
当需要将处理结果提供给反演软件时
\sphinxbeforeendvarwidth
\end{varwidth}%
&\begin{varwidth}[t]{\sphinxcolwidth{1}{5}}
\sphinxAtStartPar
.edi文件（每个测点一个）
\sphinxbeforeendvarwidth
\end{varwidth}%
&\begin{varwidth}[t]{\sphinxcolwidth{1}{5}}
\sphinxAtStartPar
建议先检查数据质量再导出
\sphinxbeforeendvarwidth
\end{varwidth}%
\\
\sphinxbottomrule
\end{tabulary}
\sphinxtableafterendhook\par
\sphinxattableend\end{savenotes}

\sphinxAtStartPar
\sphinxstylestrong{5. 其他功能}


\begin{savenotes}\sphinxattablestart
\sphinxthistablewithglobalstyle
\centering
\begin{tabulary}{\linewidth}[t]{TTTT}
\sphinxtoprule
\begin{varwidth}[t]{\sphinxcolwidth{1}{4}}
\sphinxstyletheadfamily \sphinxAtStartPar
功能
\sphinxbeforeendvarwidth
\end{varwidth}%
&\begin{varwidth}[t]{\sphinxcolwidth{1}{4}}
\sphinxstyletheadfamily \sphinxAtStartPar
说明
\sphinxbeforeendvarwidth
\end{varwidth}%
&\begin{varwidth}[t]{\sphinxcolwidth{1}{4}}
\sphinxstyletheadfamily \sphinxAtStartPar
使用场景
\sphinxbeforeendvarwidth
\end{varwidth}%
&\begin{varwidth}[t]{\sphinxcolwidth{1}{4}}
\sphinxstyletheadfamily \sphinxAtStartPar
操作步骤
\sphinxbeforeendvarwidth
\end{varwidth}%
\\
\sphinxmidrule
\sphinxtableatstartofbodyhook\begin{varwidth}[t]{\sphinxcolwidth{1}{4}}
\sphinxAtStartPar
更新TBL
\sphinxbeforeendvarwidth
\end{varwidth}%
&\begin{varwidth}[t]{\sphinxcolwidth{1}{4}}
\sphinxAtStartPar
更新工区内所有测点的TBL文件
\sphinxbeforeendvarwidth
\end{varwidth}%
&\begin{varwidth}[t]{\sphinxcolwidth{1}{4}}
\sphinxAtStartPar
当TBL文件有更新或需要重新生成时
\sphinxbeforeendvarwidth
\end{varwidth}%
&\begin{varwidth}[t]{\sphinxcolwidth{1}{4}}
\sphinxAtStartPar
1. 右键工区名称2. 点击"更新TBL"3. 等待更新完成
\sphinxbeforeendvarwidth
\end{varwidth}%
\\
\sphinxbottomrule
\end{tabulary}
\sphinxtableafterendhook\par
\sphinxattableend\end{savenotes}


\subparagraph{📁 分组（测段）右键菜单}
\label{\detokenize{chapters/chapter1:id18}}
\sphinxAtStartPar
分组（Group/测段）是同一批处理测点的集合，分组右键菜单功能最为丰富，是日常数据处理中最常用的菜单。

\sphinxAtStartPar
\sphinxstylestrong{1. 测点管理}


\begin{savenotes}\sphinxattablestart
\sphinxthistablewithglobalstyle
\centering
\begin{tabulary}{\linewidth}[t]{TTTTT}
\sphinxtoprule
\begin{varwidth}[t]{\sphinxcolwidth{1}{5}}
\sphinxstyletheadfamily \sphinxAtStartPar
功能
\sphinxbeforeendvarwidth
\end{varwidth}%
&\begin{varwidth}[t]{\sphinxcolwidth{1}{5}}
\sphinxstyletheadfamily \sphinxAtStartPar
说明
\sphinxbeforeendvarwidth
\end{varwidth}%
&\begin{varwidth}[t]{\sphinxcolwidth{1}{5}}
\sphinxstyletheadfamily \sphinxAtStartPar
使用场景
\sphinxbeforeendvarwidth
\end{varwidth}%
&\begin{varwidth}[t]{\sphinxcolwidth{1}{5}}
\sphinxstyletheadfamily \sphinxAtStartPar
操作步骤
\sphinxbeforeendvarwidth
\end{varwidth}%
&\begin{varwidth}[t]{\sphinxcolwidth{1}{5}}
\sphinxstyletheadfamily \sphinxAtStartPar
注意事项
\sphinxbeforeendvarwidth
\end{varwidth}%
\\
\sphinxmidrule
\sphinxtableatstartofbodyhook\begin{varwidth}[t]{\sphinxcolwidth{1}{5}}
\sphinxAtStartPar
加载测点
\sphinxbeforeendvarwidth
\end{varwidth}%
&\begin{varwidth}[t]{\sphinxcolwidth{1}{5}}
\sphinxAtStartPar
从目录批量加载测点数据
\sphinxbeforeendvarwidth
\end{varwidth}%
&\begin{varwidth}[t]{\sphinxcolwidth{1}{5}}
\sphinxAtStartPar
当需要导入一批新采集的测点时
\sphinxbeforeendvarwidth
\end{varwidth}%
&\begin{varwidth}[t]{\sphinxcolwidth{1}{5}}
\sphinxAtStartPar
1. 右键分组名称2. 点击"加载测点"3. 选择数据目录4. 选择要加载的测点
\sphinxbeforeendvarwidth
\end{varwidth}%
&\begin{varwidth}[t]{\sphinxcolwidth{1}{5}}
\sphinxAtStartPar
会清除分组内现有测点
\sphinxbeforeendvarwidth
\end{varwidth}%
\\
\sphinxhline\begin{varwidth}[t]{\sphinxcolwidth{1}{5}}
\sphinxAtStartPar
追加加载
\sphinxbeforeendvarwidth
\end{varwidth}%
&\begin{varwidth}[t]{\sphinxcolwidth{1}{5}}
\sphinxAtStartPar
追加加载测点（不清除现有测点）
\sphinxbeforeendvarwidth
\end{varwidth}%
&\begin{varwidth}[t]{\sphinxcolwidth{1}{5}}
\sphinxAtStartPar
当需要向分组中添加新测点时
\sphinxbeforeendvarwidth
\end{varwidth}%
&\begin{varwidth}[t]{\sphinxcolwidth{1}{5}}
\sphinxAtStartPar
1. 右键分组名称2. 点击"追加加载"3. 选择数据目录4. 选择要添加的测点
\sphinxbeforeendvarwidth
\end{varwidth}%
&\begin{varwidth}[t]{\sphinxcolwidth{1}{5}}
\sphinxAtStartPar
不会影响已存在的测点
\sphinxbeforeendvarwidth
\end{varwidth}%
\\
\sphinxhline\begin{varwidth}[t]{\sphinxcolwidth{1}{5}}
\sphinxAtStartPar
从EDI加载
\sphinxbeforeendvarwidth
\end{varwidth}%
&\begin{varwidth}[t]{\sphinxcolwidth{1}{5}}
\sphinxAtStartPar
从EDI文件加载测点
\sphinxbeforeendvarwidth
\end{varwidth}%
&\begin{varwidth}[t]{\sphinxcolwidth{1}{5}}
\sphinxAtStartPar
当需要导入已处理的EDI数据时
\sphinxbeforeendvarwidth
\end{varwidth}%
&\begin{varwidth}[t]{\sphinxcolwidth{1}{5}}
\sphinxAtStartPar
1. 右键分组名称2. 点击"从EDI加载"3. 选择EDI文件或目录
\sphinxbeforeendvarwidth
\end{varwidth}%
&\begin{varwidth}[t]{\sphinxcolwidth{1}{5}}
\sphinxAtStartPar
适合导入历史数据或他人数据
\sphinxbeforeendvarwidth
\end{varwidth}%
\\
\sphinxhline\begin{varwidth}[t]{\sphinxcolwidth{1}{5}}
\sphinxAtStartPar
加载源测点
\sphinxbeforeendvarwidth
\end{varwidth}%
&\begin{varwidth}[t]{\sphinxcolwidth{1}{5}}
\sphinxAtStartPar
加载源场测点（用于远参考处理）
\sphinxbeforeendvarwidth
\end{varwidth}%
&\begin{varwidth}[t]{\sphinxcolwidth{1}{5}}
\sphinxAtStartPar
当需要进行远参考处理时
\sphinxbeforeendvarwidth
\end{varwidth}%
&\begin{varwidth}[t]{\sphinxcolwidth{1}{5}}
\sphinxAtStartPar
1. 右键分组名称2. 点击"加载源测点"3. 选择源场测点数据
\sphinxbeforeendvarwidth
\end{varwidth}%
&\begin{varwidth}[t]{\sphinxcolwidth{1}{5}}
\sphinxAtStartPar
源测点需要与测站时间同步
\sphinxbeforeendvarwidth
\end{varwidth}%
\\
\sphinxhline\begin{varwidth}[t]{\sphinxcolwidth{1}{5}}
\sphinxAtStartPar
加载多文件测点
\sphinxbeforeendvarwidth
\end{varwidth}%
&\begin{varwidth}[t]{\sphinxcolwidth{1}{5}}
\sphinxAtStartPar
从多个不同文件加载测点
\sphinxbeforeendvarwidth
\end{varwidth}%
&\begin{varwidth}[t]{\sphinxcolwidth{1}{5}}
\sphinxAtStartPar
当测点数据分散在多个文件时
\sphinxbeforeendvarwidth
\end{varwidth}%
&\begin{varwidth}[t]{\sphinxcolwidth{1}{5}}
\sphinxAtStartPar
1. 右键分组名称2. 点击"加载多文件测点"3. 选择多个文件
\sphinxbeforeendvarwidth
\end{varwidth}%
&\begin{varwidth}[t]{\sphinxcolwidth{1}{5}}
\sphinxAtStartPar
支持不同格式混合加载
\sphinxbeforeendvarwidth
\end{varwidth}%
\\
\sphinxhline\begin{varwidth}[t]{\sphinxcolwidth{1}{5}}
\sphinxAtStartPar
从目录加载Phoenix
\sphinxbeforeendvarwidth
\end{varwidth}%
&\begin{varwidth}[t]{\sphinxcolwidth{1}{5}}
\sphinxAtStartPar
从目录批量加载Phoenix格式测点
\sphinxbeforeendvarwidth
\end{varwidth}%
&\begin{varwidth}[t]{\sphinxcolwidth{1}{5}}
\sphinxAtStartPar
当使用Phoenix仪器且数据按目录组织时
\sphinxbeforeendvarwidth
\end{varwidth}%
&\begin{varwidth}[t]{\sphinxcolwidth{1}{5}}
\sphinxAtStartPar
1. 右键分组名称2. 点击"从目录加载Phoenix"3. 选择Phoenix数据目录
\sphinxbeforeendvarwidth
\end{varwidth}%
&\begin{varwidth}[t]{\sphinxcolwidth{1}{5}}
\sphinxAtStartPar
自动识别Phoenix数据格式
\sphinxbeforeendvarwidth
\end{varwidth}%
\\
\sphinxhline\begin{varwidth}[t]{\sphinxcolwidth{1}{5}}
\sphinxAtStartPar
加载MTU8
\sphinxbeforeendvarwidth
\end{varwidth}%
&\begin{varwidth}[t]{\sphinxcolwidth{1}{5}}
\sphinxAtStartPar
加载MTU\sphinxhyphen{}8格式数据
\sphinxbeforeendvarwidth
\end{varwidth}%
&\begin{varwidth}[t]{\sphinxcolwidth{1}{5}}
\sphinxAtStartPar
当使用MTU\sphinxhyphen{}8仪器采集数据时
\sphinxbeforeendvarwidth
\end{varwidth}%
&\begin{varwidth}[t]{\sphinxcolwidth{1}{5}}
\sphinxAtStartPar
1. 右键分组名称2. 点击"加载MTU8"3. 选择MTU\sphinxhyphen{}8数据文件
\sphinxbeforeendvarwidth
\end{varwidth}%
&\begin{varwidth}[t]{\sphinxcolwidth{1}{5}}
\sphinxAtStartPar
需要recmeta.json文件
\sphinxbeforeendvarwidth
\end{varwidth}%
\\
\sphinxhline\begin{varwidth}[t]{\sphinxcolwidth{1}{5}}
\sphinxAtStartPar
加载KML坐标
\sphinxbeforeendvarwidth
\end{varwidth}%
&\begin{varwidth}[t]{\sphinxcolwidth{1}{5}}
\sphinxAtStartPar
从KML文件导入测点坐标
\sphinxbeforeendvarwidth
\end{varwidth}%
&\begin{varwidth}[t]{\sphinxcolwidth{1}{5}}
\sphinxAtStartPar
当测点坐标已用KML整理好时
\sphinxbeforeendvarwidth
\end{varwidth}%
&\begin{varwidth}[t]{\sphinxcolwidth{1}{5}}
\sphinxAtStartPar
1. 右键分组名称2. 点击"加载KML坐标"3. 选择KML文件
\sphinxbeforeendvarwidth
\end{varwidth}%
&\begin{varwidth}[t]{\sphinxcolwidth{1}{5}}
\sphinxAtStartPar
仅导入坐标，不导入时间序列
\sphinxbeforeendvarwidth
\end{varwidth}%
\\
\sphinxhline\begin{varwidth}[t]{\sphinxcolwidth{1}{5}}
\sphinxAtStartPar
粘贴测点
\sphinxbeforeendvarwidth
\end{varwidth}%
&\begin{varwidth}[t]{\sphinxcolwidth{1}{5}}
\sphinxAtStartPar
粘贴已复制的单个或多个测点
\sphinxbeforeendvarwidth
\end{varwidth}%
&\begin{varwidth}[t]{\sphinxcolwidth{1}{5}}
\sphinxAtStartPar
当需要在分组之间移动测点时
\sphinxbeforeendvarwidth
\end{varwidth}%
&\begin{varwidth}[t]{\sphinxcolwidth{1}{5}}
\sphinxAtStartPar
1. 先在源分组复制测点2. 右键目标分组3. 点击"粘贴测点"
\sphinxbeforeendvarwidth
\end{varwidth}%
&\begin{varwidth}[t]{\sphinxcolwidth{1}{5}}
\sphinxAtStartPar
测点的所有数据都会复制
\sphinxbeforeendvarwidth
\end{varwidth}%
\\
\sphinxhline\begin{varwidth}[t]{\sphinxcolwidth{1}{5}}
\sphinxAtStartPar
粘贴分组测点
\sphinxbeforeendvarwidth
\end{varwidth}%
&\begin{varwidth}[t]{\sphinxcolwidth{1}{5}}
\sphinxAtStartPar
粘贴整组测点
\sphinxbeforeendvarwidth
\end{varwidth}%
&\begin{varwidth}[t]{\sphinxcolwidth{1}{5}}
\sphinxAtStartPar
当需要合并两个分组的测点时
\sphinxbeforeendvarwidth
\end{varwidth}%
&\begin{varwidth}[t]{\sphinxcolwidth{1}{5}}
\sphinxAtStartPar
1. 先复制源分组2. 右键目标分组3. 点击"粘贴分组测点"
\sphinxbeforeendvarwidth
\end{varwidth}%
&\begin{varwidth}[t]{\sphinxcolwidth{1}{5}}
\sphinxAtStartPar
会追加所有测点
\sphinxbeforeendvarwidth
\end{varwidth}%
\\
\sphinxbottomrule
\end{tabulary}
\sphinxtableafterendhook\par
\sphinxattableend\end{savenotes}
\begin{quote}

\sphinxAtStartPar
\sphinxstylestrong{提示}：加载测点时，MTDP会自动识别仪器类型（Phoenix、LEMI、Metronix等），无需手动指定。
\end{quote}

\sphinxAtStartPar
\sphinxstylestrong{2. 测点批量操作}


\begin{savenotes}\sphinxattablestart
\sphinxthistablewithglobalstyle
\centering
\begin{tabulary}{\linewidth}[t]{TTTTT}
\sphinxtoprule
\begin{varwidth}[t]{\sphinxcolwidth{1}{5}}
\sphinxstyletheadfamily \sphinxAtStartPar
功能
\sphinxbeforeendvarwidth
\end{varwidth}%
&\begin{varwidth}[t]{\sphinxcolwidth{1}{5}}
\sphinxstyletheadfamily \sphinxAtStartPar
说明
\sphinxbeforeendvarwidth
\end{varwidth}%
&\begin{varwidth}[t]{\sphinxcolwidth{1}{5}}
\sphinxstyletheadfamily \sphinxAtStartPar
使用场景
\sphinxbeforeendvarwidth
\end{varwidth}%
&\begin{varwidth}[t]{\sphinxcolwidth{1}{5}}
\sphinxstyletheadfamily \sphinxAtStartPar
操作步骤
\sphinxbeforeendvarwidth
\end{varwidth}%
&\begin{varwidth}[t]{\sphinxcolwidth{1}{5}}
\sphinxstyletheadfamily \sphinxAtStartPar
注意事项
\sphinxbeforeendvarwidth
\end{varwidth}%
\\
\sphinxmidrule
\sphinxtableatstartofbodyhook\begin{varwidth}[t]{\sphinxcolwidth{1}{5}}
\sphinxAtStartPar
测点重命名
\sphinxbeforeendvarwidth
\end{varwidth}%
&\begin{varwidth}[t]{\sphinxcolwidth{1}{5}}
\sphinxAtStartPar
批量重命名分组内测点
\sphinxbeforeendvarwidth
\end{varwidth}%
&\begin{varwidth}[t]{\sphinxcolwidth{1}{5}}
\sphinxAtStartPar
当测点命名需要统一调整时
\sphinxbeforeendvarwidth
\end{varwidth}%
&\begin{varwidth}[t]{\sphinxcolwidth{1}{5}}
\sphinxAtStartPar
1. 右键分组名称2. 点击"测点重命名"3. 设置重命名规则
\sphinxbeforeendvarwidth
\end{varwidth}%
&\begin{varwidth}[t]{\sphinxcolwidth{1}{5}}
\sphinxAtStartPar
支持查找替换、添加前后缀等
\sphinxbeforeendvarwidth
\end{varwidth}%
\\
\sphinxhline\begin{varwidth}[t]{\sphinxcolwidth{1}{5}}
\sphinxAtStartPar
添加前缀
\sphinxbeforeendvarwidth
\end{varwidth}%
&\begin{varwidth}[t]{\sphinxcolwidth{1}{5}}
\sphinxAtStartPar
为所有测点名添加前缀
\sphinxbeforeendvarwidth
\end{varwidth}%
&\begin{varwidth}[t]{\sphinxcolwidth{1}{5}}
\sphinxAtStartPar
当需要标识测点来源或批次时
\sphinxbeforeendvarwidth
\end{varwidth}%
&\begin{varwidth}[t]{\sphinxcolwidth{1}{5}}
\sphinxAtStartPar
1. 右键分组名称2. 点击"添加前缀"3. 输入前缀内容
\sphinxbeforeendvarwidth
\end{varwidth}%
&\begin{varwidth}[t]{\sphinxcolwidth{1}{5}}
\sphinxAtStartPar
例如：添加"Line1\_"前缀
\sphinxbeforeendvarwidth
\end{varwidth}%
\\
\sphinxhline\begin{varwidth}[t]{\sphinxcolwidth{1}{5}}
\sphinxAtStartPar
添加后缀
\sphinxbeforeendvarwidth
\end{varwidth}%
&\begin{varwidth}[t]{\sphinxcolwidth{1}{5}}
\sphinxAtStartPar
为所有测点名添加后缀
\sphinxbeforeendvarwidth
\end{varwidth}%
&\begin{varwidth}[t]{\sphinxcolwidth{1}{5}}
\sphinxAtStartPar
当需要标识测点处理版本时
\sphinxbeforeendvarwidth
\end{varwidth}%
&\begin{varwidth}[t]{\sphinxcolwidth{1}{5}}
\sphinxAtStartPar
1. 右键分组名称2. 点击"添加后缀"3. 输入后缀内容
\sphinxbeforeendvarwidth
\end{varwidth}%
&\begin{varwidth}[t]{\sphinxcolwidth{1}{5}}
\sphinxAtStartPar
例如：添加"\_v2"后缀
\sphinxbeforeendvarwidth
\end{varwidth}%
\\
\sphinxhline\begin{varwidth}[t]{\sphinxcolwidth{1}{5}}
\sphinxAtStartPar
测点排序
\sphinxbeforeendvarwidth
\end{varwidth}%
&\begin{varwidth}[t]{\sphinxcolwidth{1}{5}}
\sphinxAtStartPar
按纬度/经度/名称/时间排序测点
\sphinxbeforeendvarwidth
\end{varwidth}%
&\begin{varwidth}[t]{\sphinxcolwidth{1}{5}}
\sphinxAtStartPar
当需要按特定顺序组织测点时
\sphinxbeforeendvarwidth
\end{varwidth}%
&\begin{varwidth}[t]{\sphinxcolwidth{1}{5}}
\sphinxAtStartPar
1. 右键分组名称2. 点击"测点排序"3. 选择排序方式4. 选择升序或降序
\sphinxbeforeendvarwidth
\end{varwidth}%
&\begin{varwidth}[t]{\sphinxcolwidth{1}{5}}
\sphinxAtStartPar
排序后影响测点在树中的显示顺序
\sphinxbeforeendvarwidth
\end{varwidth}%
\\
\sphinxhline\begin{varwidth}[t]{\sphinxcolwidth{1}{5}}
\sphinxAtStartPar
设置通道为首个
\sphinxbeforeendvarwidth
\end{varwidth}%
&\begin{varwidth}[t]{\sphinxcolwidth{1}{5}}
\sphinxAtStartPar
将指定通道设置为首个通道
\sphinxbeforeendvarwidth
\end{varwidth}%
&\begin{varwidth}[t]{\sphinxcolwidth{1}{5}}
\sphinxAtStartPar
当通道顺序需要调整时
\sphinxbeforeendvarwidth
\end{varwidth}%
&\begin{varwidth}[t]{\sphinxcolwidth{1}{5}}
\sphinxAtStartPar
1. 右键分组名称2. 点击"设置通道为首个"3. 选择通道
\sphinxbeforeendvarwidth
\end{varwidth}%
&\begin{varwidth}[t]{\sphinxcolwidth{1}{5}}
\sphinxAtStartPar
影响后续处理的通道选择
\sphinxbeforeendvarwidth
\end{varwidth}%
\\
\sphinxhline\begin{varwidth}[t]{\sphinxcolwidth{1}{5}}
\sphinxAtStartPar
设置电极为首个
\sphinxbeforeendvarwidth
\end{varwidth}%
&\begin{varwidth}[t]{\sphinxcolwidth{1}{5}}
\sphinxAtStartPar
将电极长度设置为首个
\sphinxbeforeendvarwidth
\end{varwidth}%
&\begin{varwidth}[t]{\sphinxcolwidth{1}{5}}
\sphinxAtStartPar
当需要统一电极长度设置时
\sphinxbeforeendvarwidth
\end{varwidth}%
&\begin{varwidth}[t]{\sphinxcolwidth{1}{5}}
\sphinxAtStartPar
1. 右键分组名称2. 点击"设置电极为首个"
\sphinxbeforeendvarwidth
\end{varwidth}%
&\begin{varwidth}[t]{\sphinxcolwidth{1}{5}}
\sphinxAtStartPar
以第一个测点的电极长度为准
\sphinxbeforeendvarwidth
\end{varwidth}%
\\
\sphinxhline\begin{varwidth}[t]{\sphinxcolwidth{1}{5}}
\sphinxAtStartPar
设置旋转角为首个
\sphinxbeforeendvarwidth
\end{varwidth}%
&\begin{varwidth}[t]{\sphinxcolwidth{1}{5}}
\sphinxAtStartPar
将旋转角设置为首个
\sphinxbeforeendvarwidth
\end{varwidth}%
&\begin{varwidth}[t]{\sphinxcolwidth{1}{5}}
\sphinxAtStartPar
当需要统一旋转角度时
\sphinxbeforeendvarwidth
\end{varwidth}%
&\begin{varwidth}[t]{\sphinxcolwidth{1}{5}}
\sphinxAtStartPar
1. 右键分组名称2. 点击"设置旋转角为首个"
\sphinxbeforeendvarwidth
\end{varwidth}%
&\begin{varwidth}[t]{\sphinxcolwidth{1}{5}}
\sphinxAtStartPar
以第一个测点的旋转角为准
\sphinxbeforeendvarwidth
\end{varwidth}%
\\
\sphinxbottomrule
\end{tabulary}
\sphinxtableafterendhook\par
\sphinxattableend\end{savenotes}

\sphinxAtStartPar
\sphinxstylestrong{3. 数据处理}


\begin{savenotes}\sphinxattablestart
\sphinxthistablewithglobalstyle
\centering
\begin{tabulary}{\linewidth}[t]{TTTTT}
\sphinxtoprule
\begin{varwidth}[t]{\sphinxcolwidth{1}{5}}
\sphinxstyletheadfamily \sphinxAtStartPar
功能
\sphinxbeforeendvarwidth
\end{varwidth}%
&\begin{varwidth}[t]{\sphinxcolwidth{1}{5}}
\sphinxstyletheadfamily \sphinxAtStartPar
说明
\sphinxbeforeendvarwidth
\end{varwidth}%
&\begin{varwidth}[t]{\sphinxcolwidth{1}{5}}
\sphinxstyletheadfamily \sphinxAtStartPar
使用场景
\sphinxbeforeendvarwidth
\end{varwidth}%
&\begin{varwidth}[t]{\sphinxcolwidth{1}{5}}
\sphinxstyletheadfamily \sphinxAtStartPar
操作步骤
\sphinxbeforeendvarwidth
\end{varwidth}%
&\begin{varwidth}[t]{\sphinxcolwidth{1}{5}}
\sphinxstyletheadfamily \sphinxAtStartPar
注意事项
\sphinxbeforeendvarwidth
\end{varwidth}%
\\
\sphinxmidrule
\sphinxtableatstartofbodyhook\begin{varwidth}[t]{\sphinxcolwidth{1}{5}}
\sphinxAtStartPar
处理全部
\sphinxbeforeendvarwidth
\end{varwidth}%
&\begin{varwidth}[t]{\sphinxcolwidth{1}{5}}
\sphinxAtStartPar
处理分组内所有测点
\sphinxbeforeendvarwidth
\end{varwidth}%
&\begin{varwidth}[t]{\sphinxcolwidth{1}{5}}
\sphinxAtStartPar
当首次处理或需要重新处理所有测点时
\sphinxbeforeendvarwidth
\end{varwidth}%
&\begin{varwidth}[t]{\sphinxcolwidth{1}{5}}
\sphinxAtStartPar
1. 右键分组名称2. 点击"处理全部"3. 设置FFT参数4. 执行处理
\sphinxbeforeendvarwidth
\end{varwidth}%
&\begin{varwidth}[t]{\sphinxcolwidth{1}{5}}
\sphinxAtStartPar
处理时间取决于测点数量
\sphinxbeforeendvarwidth
\end{varwidth}%
\\
\sphinxhline\begin{varwidth}[t]{\sphinxcolwidth{1}{5}}
\sphinxAtStartPar
处理剩余
\sphinxbeforeendvarwidth
\end{varwidth}%
&\begin{varwidth}[t]{\sphinxcolwidth{1}{5}}
\sphinxAtStartPar
仅处理未处理的测点
\sphinxbeforeendvarwidth
\end{varwidth}%
&\begin{varwidth}[t]{\sphinxcolwidth{1}{5}}
\sphinxAtStartPar
当部分测点已处理，需要处理剩余测点时
\sphinxbeforeendvarwidth
\end{varwidth}%
&\begin{varwidth}[t]{\sphinxcolwidth{1}{5}}
\sphinxAtStartPar
1. 右键分组名称2. 点击"处理剩余"3. 确认处理
\sphinxbeforeendvarwidth
\end{varwidth}%
&\begin{varwidth}[t]{\sphinxcolwidth{1}{5}}
\sphinxAtStartPar
已处理的测点不会被重新处理
\sphinxbeforeendvarwidth
\end{varwidth}%
\\
\sphinxhline\begin{varwidth}[t]{\sphinxcolwidth{1}{5}}
\sphinxAtStartPar
同标定处理
\sphinxbeforeendvarwidth
\end{varwidth}%
&\begin{varwidth}[t]{\sphinxcolwidth{1}{5}}
\sphinxAtStartPar
使用相同标定文件处理所有测点
\sphinxbeforeendvarwidth
\end{varwidth}%
&\begin{varwidth}[t]{\sphinxcolwidth{1}{5}}
\sphinxAtStartPar
当所有测点使用同一套标定文件时
\sphinxbeforeendvarwidth
\end{varwidth}%
&\begin{varwidth}[t]{\sphinxcolwidth{1}{5}}
\sphinxAtStartPar
1. 右键分组名称2. 点击"同标定处理"3. 选择标定文件
\sphinxbeforeendvarwidth
\end{varwidth}%
&\begin{varwidth}[t]{\sphinxcolwidth{1}{5}}
\sphinxAtStartPar
适用于使用相同仪器的情况
\sphinxbeforeendvarwidth
\end{varwidth}%
\\
\sphinxhline\begin{varwidth}[t]{\sphinxcolwidth{1}{5}}
\sphinxAtStartPar
Phoenix Robust处理
\sphinxbeforeendvarwidth
\end{varwidth}%
&\begin{varwidth}[t]{\sphinxcolwidth{1}{5}}
\sphinxAtStartPar
Phoenix Robust稳健估计
\sphinxbeforeendvarwidth
\end{varwidth}%
&\begin{varwidth}[t]{\sphinxcolwidth{1}{5}}
\sphinxAtStartPar
当需要使用Phoenix专用Robust算法时
\sphinxbeforeendvarwidth
\end{varwidth}%
&\begin{varwidth}[t]{\sphinxcolwidth{1}{5}}
\sphinxAtStartPar
1. 右键分组名称2. 点击"Phoenix Robust处理"3. 设置参数
\sphinxbeforeendvarwidth
\end{varwidth}%
&\begin{varwidth}[t]{\sphinxcolwidth{1}{5}}
\sphinxAtStartPar
仅适用于Phoenix数据
\sphinxbeforeendvarwidth
\end{varwidth}%
\\
\sphinxhline\begin{varwidth}[t]{\sphinxcolwidth{1}{5}}
\sphinxAtStartPar
Phoenix TS转FT
\sphinxbeforeendvarwidth
\end{varwidth}%
&\begin{varwidth}[t]{\sphinxcolwidth{1}{5}}
\sphinxAtStartPar
Phoenix时间序列转傅里叶系数
\sphinxbeforeendvarwidth
\end{varwidth}%
&\begin{varwidth}[t]{\sphinxcolwidth{1}{5}}
\sphinxAtStartPar
当需要将时域数据转换为频域时
\sphinxbeforeendvarwidth
\end{varwidth}%
&\begin{varwidth}[t]{\sphinxcolwidth{1}{5}}
\sphinxAtStartPar
1. 右键分组名称2. 点击"Phoenix TS转FT"3. 设置FFT参数
\sphinxbeforeendvarwidth
\end{varwidth}%
&\begin{varwidth}[t]{\sphinxcolwidth{1}{5}}
\sphinxAtStartPar
生成.FT文件供后续处理使用
\sphinxbeforeendvarwidth
\end{varwidth}%
\\
\sphinxhline\begin{varwidth}[t]{\sphinxcolwidth{1}{5}}
\sphinxAtStartPar
计算标定
\sphinxbeforeendvarwidth
\end{varwidth}%
&\begin{varwidth}[t]{\sphinxcolwidth{1}{5}}
\sphinxAtStartPar
计算测点的标定系数
\sphinxbeforeendvarwidth
\end{varwidth}%
&\begin{varwidth}[t]{\sphinxcolwidth{1}{5}}
\sphinxAtStartPar
当需要重新计算或更新标定系数时
\sphinxbeforeendvarwidth
\end{varwidth}%
&\begin{varwidth}[t]{\sphinxcolwidth{1}{5}}
\sphinxAtStartPar
1. 右键分组名称2. 点击"计算标定"3. 选择标定文件
\sphinxbeforeendvarwidth
\end{varwidth}%
&\begin{varwidth}[t]{\sphinxcolwidth{1}{5}}
\sphinxAtStartPar
需要先设置标定文件路径
\sphinxbeforeendvarwidth
\end{varwidth}%
\\
\sphinxhline\begin{varwidth}[t]{\sphinxcolwidth{1}{5}}
\sphinxAtStartPar
计算远参考
\sphinxbeforeendvarwidth
\end{varwidth}%
&\begin{varwidth}[t]{\sphinxcolwidth{1}{5}}
\sphinxAtStartPar
计算远参考传递函数
\sphinxbeforeendvarwidth
\end{varwidth}%
&\begin{varwidth}[t]{\sphinxcolwidth{1}{5}}
\sphinxAtStartPar
当分组内包含远参考站数据时
\sphinxbeforeendvarwidth
\end{varwidth}%
&\begin{varwidth}[t]{\sphinxcolwidth{1}{5}}
\sphinxAtStartPar
1. 右键分组名称2. 点击"计算远参考"3. 选择远参考站
\sphinxbeforeendvarwidth
\end{varwidth}%
&\begin{varwidth}[t]{\sphinxcolwidth{1}{5}}
\sphinxAtStartPar
所有测点之间会互为远参考
\sphinxbeforeendvarwidth
\end{varwidth}%
\\
\sphinxbottomrule
\end{tabulary}
\sphinxtableafterendhook\par
\sphinxattableend\end{savenotes}

\sphinxAtStartPar
\sphinxstylestrong{4. 数据导出}


\begin{savenotes}\sphinxattablestart
\sphinxthistablewithglobalstyle
\centering
\begin{tabulary}{\linewidth}[t]{TTTTT}
\sphinxtoprule
\begin{varwidth}[t]{\sphinxcolwidth{1}{5}}
\sphinxstyletheadfamily \sphinxAtStartPar
功能
\sphinxbeforeendvarwidth
\end{varwidth}%
&\begin{varwidth}[t]{\sphinxcolwidth{1}{5}}
\sphinxstyletheadfamily \sphinxAtStartPar
说明
\sphinxbeforeendvarwidth
\end{varwidth}%
&\begin{varwidth}[t]{\sphinxcolwidth{1}{5}}
\sphinxstyletheadfamily \sphinxAtStartPar
使用场景
\sphinxbeforeendvarwidth
\end{varwidth}%
&\begin{varwidth}[t]{\sphinxcolwidth{1}{5}}
\sphinxstyletheadfamily \sphinxAtStartPar
输出文件
\sphinxbeforeendvarwidth
\end{varwidth}%
&\begin{varwidth}[t]{\sphinxcolwidth{1}{5}}
\sphinxstyletheadfamily \sphinxAtStartPar
注意事项
\sphinxbeforeendvarwidth
\end{varwidth}%
\\
\sphinxmidrule
\sphinxtableatstartofbodyhook\begin{varwidth}[t]{\sphinxcolwidth{1}{5}}
\sphinxAtStartPar
\sphinxstylestrong{📊 导出绘图数据}
\sphinxbeforeendvarwidth
\end{varwidth}%
&\begin{varwidth}[t]{\sphinxcolwidth{1}{5}}
\sphinxAtStartPar
导出参数/相位张量/坐标供第三方软件绘图
\sphinxbeforeendvarwidth
\end{varwidth}%
&\begin{varwidth}[t]{\sphinxcolwidth{1}{5}}
\sphinxAtStartPar
当需要使用Golden绘图、Surfer等软件成图时
\sphinxbeforeendvarwidth
\end{varwidth}%
&\begin{varwidth}[t]{\sphinxcolwidth{1}{5}}
\sphinxAtStartPar
\sphinxhyphen{} \{组名\}\sphinxhyphen{}Paint.dat\sphinxhyphen{} \{频率\}\sphinxhyphen{}PhaseTensor.dat\sphinxhyphen{} 各测点参数文件
\sphinxbeforeendvarwidth
\end{varwidth}%
&\begin{varwidth}[t]{\sphinxcolwidth{1}{5}}
\sphinxAtStartPar
推荐用于批量绘图需求
\sphinxbeforeendvarwidth
\end{varwidth}%
\\
\sphinxhline\begin{varwidth}[t]{\sphinxcolwidth{1}{5}}
\sphinxAtStartPar
导出SpeEDI
\sphinxbeforeendvarwidth
\end{varwidth}%
&\begin{varwidth}[t]{\sphinxcolwidth{1}{5}}
\sphinxAtStartPar
批量导出SpeEDI格式
\sphinxbeforeendvarwidth
\end{varwidth}%
&\begin{varwidth}[t]{\sphinxcolwidth{1}{5}}
\sphinxAtStartPar
当需要与支持SpeEDI的软件交换数据时
\sphinxbeforeendvarwidth
\end{varwidth}%
&\begin{varwidth}[t]{\sphinxcolwidth{1}{5}}
\sphinxAtStartPar
.edi文件（SpeEDI格式）
\sphinxbeforeendvarwidth
\end{varwidth}%
&\begin{varwidth}[t]{\sphinxcolwidth{1}{5}}
\sphinxAtStartPar
标准EDI格式，广泛兼容
\sphinxbeforeendvarwidth
\end{varwidth}%
\\
\sphinxhline\begin{varwidth}[t]{\sphinxcolwidth{1}{5}}
\sphinxAtStartPar
导出ZTEDI
\sphinxbeforeendvarwidth
\end{varwidth}%
&\begin{varwidth}[t]{\sphinxcolwidth{1}{5}}
\sphinxAtStartPar
批量导出ZTEDI格式
\sphinxbeforeendvarwidth
\end{varwidth}%
&\begin{varwidth}[t]{\sphinxcolwidth{1}{5}}
\sphinxAtStartPar
当需要包含完整张量信息的EDI时
\sphinxbeforeendvarwidth
\end{varwidth}%
&\begin{varwidth}[t]{\sphinxcolwidth{1}{5}}
\sphinxAtStartPar
.edi文件（ZTEDI格式）
\sphinxbeforeendvarwidth
\end{varwidth}%
&\begin{varwidth}[t]{\sphinxcolwidth{1}{5}}
\sphinxAtStartPar
包含阻抗和倾子张量
\sphinxbeforeendvarwidth
\end{varwidth}%
\\
\sphinxhline\begin{varwidth}[t]{\sphinxcolwidth{1}{5}}
\sphinxAtStartPar
导出MTpyEDI
\sphinxbeforeendvarwidth
\end{varwidth}%
&\begin{varwidth}[t]{\sphinxcolwidth{1}{5}}
\sphinxAtStartPar
批量导出MTpyEDI格式
\sphinxbeforeendvarwidth
\end{varwidth}%
&\begin{varwidth}[t]{\sphinxcolwidth{1}{5}}
\sphinxAtStartPar
当需要使用Python MTpy库处理时
\sphinxbeforeendvarwidth
\end{varwidth}%
&\begin{varwidth}[t]{\sphinxcolwidth{1}{5}}
\sphinxAtStartPar
.edi文件（MTpyEDI格式）
\sphinxbeforeendvarwidth
\end{varwidth}%
&\begin{varwidth}[t]{\sphinxcolwidth{1}{5}}
\sphinxAtStartPar
Python兼容格式
\sphinxbeforeendvarwidth
\end{varwidth}%
\\
\sphinxhline\begin{varwidth}[t]{\sphinxcolwidth{1}{5}}
\sphinxAtStartPar
导出PLT
\sphinxbeforeendvarwidth
\end{varwidth}%
&\begin{varwidth}[t]{\sphinxcolwidth{1}{5}}
\sphinxAtStartPar
批量导出PLT格式
\sphinxbeforeendvarwidth
\end{varwidth}%
&\begin{varwidth}[t]{\sphinxcolwidth{1}{5}}
\sphinxAtStartPar
当需要绘图软件直接使用时
\sphinxbeforeendvarwidth
\end{varwidth}%
&\begin{varwidth}[t]{\sphinxcolwidth{1}{5}}
\sphinxAtStartPar
.plt文件
\sphinxbeforeendvarwidth
\end{varwidth}%
&\begin{varwidth}[t]{\sphinxcolwidth{1}{5}}
\sphinxAtStartPar
包含视电阻率和相位曲线数据
\sphinxbeforeendvarwidth
\end{varwidth}%
\\
\sphinxhline\begin{varwidth}[t]{\sphinxcolwidth{1}{5}}
\sphinxAtStartPar
导出数据质量
\sphinxbeforeendvarwidth
\end{varwidth}%
&\begin{varwidth}[t]{\sphinxcolwidth{1}{5}}
\sphinxAtStartPar
导出测点质量评级
\sphinxbeforeendvarwidth
\end{varwidth}%
&\begin{varwidth}[t]{\sphinxcolwidth{1}{5}}
\sphinxAtStartPar
当需要汇总数据质量信息时
\sphinxbeforeendvarwidth
\end{varwidth}%
&\begin{varwidth}[t]{\sphinxcolwidth{1}{5}}
\sphinxAtStartPar
.dat文件
\sphinxbeforeendvarwidth
\end{varwidth}%
&\begin{varwidth}[t]{\sphinxcolwidth{1}{5}}
\sphinxAtStartPar
包含每个测点的质量等级
\sphinxbeforeendvarwidth
\end{varwidth}%
\\
\sphinxhline\begin{varwidth}[t]{\sphinxcolwidth{1}{5}}
\sphinxAtStartPar
导出测点信息
\sphinxbeforeendvarwidth
\end{varwidth}%
&\begin{varwidth}[t]{\sphinxcolwidth{1}{5}}
\sphinxAtStartPar
导出测点详细信息
\sphinxbeforeendvarwidth
\end{varwidth}%
&\begin{varwidth}[t]{\sphinxcolwidth{1}{5}}
\sphinxAtStartPar
当需要完整的测点元数据时
\sphinxbeforeendvarwidth
\end{varwidth}%
&\begin{varwidth}[t]{\sphinxcolwidth{1}{5}}
\sphinxAtStartPar
.dat文件
\sphinxbeforeendvarwidth
\end{varwidth}%
&\begin{varwidth}[t]{\sphinxcolwidth{1}{5}}
\sphinxAtStartPar
包含坐标、时间、仪器等信息
\sphinxbeforeendvarwidth
\end{varwidth}%
\\
\sphinxhline\begin{varwidth}[t]{\sphinxcolwidth{1}{5}}
\sphinxAtStartPar
导出KMZ
\sphinxbeforeendvarwidth
\end{varwidth}%
&\begin{varwidth}[t]{\sphinxcolwidth{1}{5}}
\sphinxAtStartPar
导出Google Earth KMZ
\sphinxbeforeendvarwidth
\end{varwidth}%
&\begin{varwidth}[t]{\sphinxcolwidth{1}{5}}
\sphinxAtStartPar
当需要在Google Earth中查看测点分布时
\sphinxbeforeendvarwidth
\end{varwidth}%
&\begin{varwidth}[t]{\sphinxcolwidth{1}{5}}
\sphinxAtStartPar
.kmz文件
\sphinxbeforeendvarwidth
\end{varwidth}%
&\begin{varwidth}[t]{\sphinxcolwidth{1}{5}}
\sphinxAtStartPar
压缩格式，包含图标和样式
\sphinxbeforeendvarwidth
\end{varwidth}%
\\
\sphinxhline\begin{varwidth}[t]{\sphinxcolwidth{1}{5}}
\sphinxAtStartPar
导出OmapKML
\sphinxbeforeendvarwidth
\end{varwidth}%
&\begin{varwidth}[t]{\sphinxcolwidth{1}{5}}
\sphinxAtStartPar
导出Omap专用KML
\sphinxbeforeendvarwidth
\end{varwidth}%
&\begin{varwidth}[t]{\sphinxcolwidth{1}{5}}
\sphinxAtStartPar
当使用Omap成图时
\sphinxbeforeendvarwidth
\end{varwidth}%
&\begin{varwidth}[t]{\sphinxcolwidth{1}{5}}
\sphinxAtStartPar
.kml文件
\sphinxbeforeendvarwidth
\end{varwidth}%
&\begin{varwidth}[t]{\sphinxcolwidth{1}{5}}
\sphinxAtStartPar
包含Omap专用属性
\sphinxbeforeendvarwidth
\end{varwidth}%
\\
\sphinxhline\begin{varwidth}[t]{\sphinxcolwidth{1}{5}}
\sphinxAtStartPar
导出普通KML
\sphinxbeforeendvarwidth
\end{varwidth}%
&\begin{varwidth}[t]{\sphinxcolwidth{1}{5}}
\sphinxAtStartPar
导出标准KML
\sphinxbeforeendvarwidth
\end{varwidth}%
&\begin{varwidth}[t]{\sphinxcolwidth{1}{5}}
\sphinxAtStartPar
当需要与其他GIS软件交换数据时
\sphinxbeforeendvarwidth
\end{varwidth}%
&\begin{varwidth}[t]{\sphinxcolwidth{1}{5}}
\sphinxAtStartPar
.kml文件
\sphinxbeforeendvarwidth
\end{varwidth}%
&\begin{varwidth}[t]{\sphinxcolwidth{1}{5}}
\sphinxAtStartPar
标准格式，兼容性好
\sphinxbeforeendvarwidth
\end{varwidth}%
\\
\sphinxhline\begin{varwidth}[t]{\sphinxcolwidth{1}{5}}
\sphinxAtStartPar
追加EDI
\sphinxbeforeendvarwidth
\end{varwidth}%
&\begin{varwidth}[t]{\sphinxcolwidth{1}{5}}
\sphinxAtStartPar
将EDI数据追加到测点
\sphinxbeforeendvarwidth
\end{varwidth}%
&\begin{varwidth}[t]{\sphinxcolwidth{1}{5}}
\sphinxAtStartPar
当需要将处理结果合并到测点时
\sphinxbeforeendvarwidth
\end{varwidth}%
&\begin{varwidth}[t]{\sphinxcolwidth{1}{5}}
\sphinxAtStartPar
更新测点数据
\sphinxbeforeendvarwidth
\end{varwidth}%
&\begin{varwidth}[t]{\sphinxcolwidth{1}{5}}
\sphinxAtStartPar
不会覆盖原有数据
\sphinxbeforeendvarwidth
\end{varwidth}%
\\
\sphinxbottomrule
\end{tabulary}
\sphinxtableafterendhook\par
\sphinxattableend\end{savenotes}

\sphinxAtStartPar
\sphinxstylestrong{5. 数据显示}


\begin{savenotes}\sphinxattablestart
\sphinxthistablewithglobalstyle
\centering
\begin{tabulary}{\linewidth}[t]{TTTT}
\sphinxtoprule
\begin{varwidth}[t]{\sphinxcolwidth{1}{4}}
\sphinxstyletheadfamily \sphinxAtStartPar
功能
\sphinxbeforeendvarwidth
\end{varwidth}%
&\begin{varwidth}[t]{\sphinxcolwidth{1}{4}}
\sphinxstyletheadfamily \sphinxAtStartPar
说明
\sphinxbeforeendvarwidth
\end{varwidth}%
&\begin{varwidth}[t]{\sphinxcolwidth{1}{4}}
\sphinxstyletheadfamily \sphinxAtStartPar
使用场景
\sphinxbeforeendvarwidth
\end{varwidth}%
&\begin{varwidth}[t]{\sphinxcolwidth{1}{4}}
\sphinxstyletheadfamily \sphinxAtStartPar
操作步骤
\sphinxbeforeendvarwidth
\end{varwidth}%
\\
\sphinxmidrule
\sphinxtableatstartofbodyhook\begin{varwidth}[t]{\sphinxcolwidth{1}{4}}
\sphinxAtStartPar
添加时序显示
\sphinxbeforeendvarwidth
\end{varwidth}%
&\begin{varwidth}[t]{\sphinxcolwidth{1}{4}}
\sphinxAtStartPar
在时序窗口显示测点时间序列
\sphinxbeforeendvarwidth
\end{varwidth}%
&\begin{varwidth}[t]{\sphinxcolwidth{1}{4}}
\sphinxAtStartPar
当需要检查原始时间序列数据时
\sphinxbeforeendvarwidth
\end{varwidth}%
&\begin{varwidth}[t]{\sphinxcolwidth{1}{4}}
\sphinxAtStartPar
1. 右键分组名称2. 点击"添加时序显示"3. 时序窗口显示所有测点时序
\sphinxbeforeendvarwidth
\end{varwidth}%
\\
\sphinxhline\begin{varwidth}[t]{\sphinxcolwidth{1}{4}}
\sphinxAtStartPar
添加标定显示
\sphinxbeforeendvarwidth
\end{varwidth}%
&\begin{varwidth}[t]{\sphinxcolwidth{1}{4}}
\sphinxAtStartPar
在标定窗口显示标定曲线
\sphinxbeforeendvarwidth
\end{varwidth}%
&\begin{varwidth}[t]{\sphinxcolwidth{1}{4}}
\sphinxAtStartPar
当需要查看标定文件的频率响应时
\sphinxbeforeendvarwidth
\end{varwidth}%
&\begin{varwidth}[t]{\sphinxcolwidth{1}{4}}
\sphinxAtStartPar
1. 右键分组名称2. 点击"添加标定显示"3. 标定窗口显示标定曲线
\sphinxbeforeendvarwidth
\end{varwidth}%
\\
\sphinxhline\begin{varwidth}[t]{\sphinxcolwidth{1}{4}}
\sphinxAtStartPar
添加到地图视图
\sphinxbeforeendvarwidth
\end{varwidth}%
&\begin{varwidth}[t]{\sphinxcolwidth{1}{4}}
\sphinxAtStartPar
将测点添加到地图视图
\sphinxbeforeendvarwidth
\end{varwidth}%
&\begin{varwidth}[t]{\sphinxcolwidth{1}{4}}
\sphinxAtStartPar
当需要在地图上查看测点位置时
\sphinxbeforeendvarwidth
\end{varwidth}%
&\begin{varwidth}[t]{\sphinxcolwidth{1}{4}}
\sphinxAtStartPar
1. 右键分组名称2. 点击"添加到地图视图"3. 地图视图显示测点位置
\sphinxbeforeendvarwidth
\end{varwidth}%
\\
\sphinxbottomrule
\end{tabulary}
\sphinxtableafterendhook\par
\sphinxattableend\end{savenotes}

\sphinxAtStartPar
\sphinxstylestrong{6. 其他功能}


\begin{savenotes}\sphinxattablestart
\sphinxthistablewithglobalstyle
\centering
\begin{tabulary}{\linewidth}[t]{TTTT}
\sphinxtoprule
\begin{varwidth}[t]{\sphinxcolwidth{1}{4}}
\sphinxstyletheadfamily \sphinxAtStartPar
功能
\sphinxbeforeendvarwidth
\end{varwidth}%
&\begin{varwidth}[t]{\sphinxcolwidth{1}{4}}
\sphinxstyletheadfamily \sphinxAtStartPar
说明
\sphinxbeforeendvarwidth
\end{varwidth}%
&\begin{varwidth}[t]{\sphinxcolwidth{1}{4}}
\sphinxstyletheadfamily \sphinxAtStartPar
使用场景
\sphinxbeforeendvarwidth
\end{varwidth}%
&\begin{varwidth}[t]{\sphinxcolwidth{1}{4}}
\sphinxstyletheadfamily \sphinxAtStartPar
操作步骤
\sphinxbeforeendvarwidth
\end{varwidth}%
\\
\sphinxmidrule
\sphinxtableatstartofbodyhook\begin{varwidth}[t]{\sphinxcolwidth{1}{4}}
\sphinxAtStartPar
编辑
\sphinxbeforeendvarwidth
\end{varwidth}%
&\begin{varwidth}[t]{\sphinxcolwidth{1}{4}}
\sphinxAtStartPar
编辑分组属性
\sphinxbeforeendvarwidth
\end{varwidth}%
&\begin{varwidth}[t]{\sphinxcolwidth{1}{4}}
\sphinxAtStartPar
当需要修改分组名称或描述时
\sphinxbeforeendvarwidth
\end{varwidth}%
&\begin{varwidth}[t]{\sphinxcolwidth{1}{4}}
\sphinxAtStartPar
1. 右键分组名称2. 点击"编辑"3. 修改属性
\sphinxbeforeendvarwidth
\end{varwidth}%
\\
\sphinxhline\begin{varwidth}[t]{\sphinxcolwidth{1}{4}}
\sphinxAtStartPar
删除
\sphinxbeforeendvarwidth
\end{varwidth}%
&\begin{varwidth}[t]{\sphinxcolwidth{1}{4}}
\sphinxAtStartPar
删除分组及所有测点
\sphinxbeforeendvarwidth
\end{varwidth}%
&\begin{varwidth}[t]{\sphinxcolwidth{1}{4}}
\sphinxAtStartPar
当分组不再需要时
\sphinxbeforeendvarwidth
\end{varwidth}%
&\begin{varwidth}[t]{\sphinxcolwidth{1}{4}}
\sphinxAtStartPar
1. 右键分组名称2. 点击"删除"3. 确认删除
\sphinxbeforeendvarwidth
\end{varwidth}%
\\
\sphinxhline\begin{varwidth}[t]{\sphinxcolwidth{1}{4}}
\sphinxAtStartPar
复制
\sphinxbeforeendvarwidth
\end{varwidth}%
&\begin{varwidth}[t]{\sphinxcolwidth{1}{4}}
\sphinxAtStartPar
复制分组
\sphinxbeforeendvarwidth
\end{varwidth}%
&\begin{varwidth}[t]{\sphinxcolwidth{1}{4}}
\sphinxAtStartPar
当需要创建相似的分组时
\sphinxbeforeendvarwidth
\end{varwidth}%
&\begin{varwidth}[t]{\sphinxcolwidth{1}{4}}
\sphinxAtStartPar
1. 右键分组名称2. 点击"复制"
\sphinxbeforeendvarwidth
\end{varwidth}%
\\
\sphinxhline\begin{varwidth}[t]{\sphinxcolwidth{1}{4}}
\sphinxAtStartPar
导入绘图用远点
\sphinxbeforeendvarwidth
\end{varwidth}%
&\begin{varwidth}[t]{\sphinxcolwidth{1}{4}}
\sphinxAtStartPar
导入用于绘图的远点数据
\sphinxbeforeendvarwidth
\end{varwidth}%
&\begin{varwidth}[t]{\sphinxcolwidth{1}{4}}
\sphinxAtStartPar
当需要绘制带远点的图件时
\sphinxbeforeendvarwidth
\end{varwidth}%
&\begin{varwidth}[t]{\sphinxcolwidth{1}{4}}
\sphinxAtStartPar
1. 右键分组名称2. 点击"导入绘图用远点"3. 选择远点数据
\sphinxbeforeendvarwidth
\end{varwidth}%
\\
\sphinxbottomrule
\end{tabulary}
\sphinxtableafterendhook\par
\sphinxattableend\end{savenotes}


\subparagraph{测点右键菜单}
\label{\detokenize{chapters/chapter1:id19}}
\sphinxAtStartPar
测点（Site）是数据处理的最小单元，测点右键菜单用于单点操作。

\sphinxAtStartPar
\sphinxstylestrong{1. 测点管理}


\begin{savenotes}\sphinxattablestart
\sphinxthistablewithglobalstyle
\centering
\begin{tabulary}{\linewidth}[t]{TTTTT}
\sphinxtoprule
\begin{varwidth}[t]{\sphinxcolwidth{1}{5}}
\sphinxstyletheadfamily \sphinxAtStartPar
功能
\sphinxbeforeendvarwidth
\end{varwidth}%
&\begin{varwidth}[t]{\sphinxcolwidth{1}{5}}
\sphinxstyletheadfamily \sphinxAtStartPar
说明
\sphinxbeforeendvarwidth
\end{varwidth}%
&\begin{varwidth}[t]{\sphinxcolwidth{1}{5}}
\sphinxstyletheadfamily \sphinxAtStartPar
使用场景
\sphinxbeforeendvarwidth
\end{varwidth}%
&\begin{varwidth}[t]{\sphinxcolwidth{1}{5}}
\sphinxstyletheadfamily \sphinxAtStartPar
操作步骤
\sphinxbeforeendvarwidth
\end{varwidth}%
&\begin{varwidth}[t]{\sphinxcolwidth{1}{5}}
\sphinxstyletheadfamily \sphinxAtStartPar
注意事项
\sphinxbeforeendvarwidth
\end{varwidth}%
\\
\sphinxmidrule
\sphinxtableatstartofbodyhook\begin{varwidth}[t]{\sphinxcolwidth{1}{5}}
\sphinxAtStartPar
编辑
\sphinxbeforeendvarwidth
\end{varwidth}%
&\begin{varwidth}[t]{\sphinxcolwidth{1}{5}}
\sphinxAtStartPar
编辑测点属性（名称、坐标、电极等）
\sphinxbeforeendvarwidth
\end{varwidth}%
&\begin{varwidth}[t]{\sphinxcolwidth{1}{5}}
\sphinxAtStartPar
当需要修改测点基本信息时
\sphinxbeforeendvarwidth
\end{varwidth}%
&\begin{varwidth}[t]{\sphinxcolwidth{1}{5}}
\sphinxAtStartPar
1. 右键测点名称2. 点击"编辑"3. 在属性面板中修改
\sphinxbeforeendvarwidth
\end{varwidth}%
&\begin{varwidth}[t]{\sphinxcolwidth{1}{5}}
\sphinxAtStartPar
修改后需要保存工程
\sphinxbeforeendvarwidth
\end{varwidth}%
\\
\sphinxhline\begin{varwidth}[t]{\sphinxcolwidth{1}{5}}
\sphinxAtStartPar
重命名
\sphinxbeforeendvarwidth
\end{varwidth}%
&\begin{varwidth}[t]{\sphinxcolwidth{1}{5}}
\sphinxAtStartPar
重命名测点
\sphinxbeforeendvarwidth
\end{varwidth}%
&\begin{varwidth}[t]{\sphinxcolwidth{1}{5}}
\sphinxAtStartPar
当测点命名需要调整时
\sphinxbeforeendvarwidth
\end{varwidth}%
&\begin{varwidth}[t]{\sphinxcolwidth{1}{5}}
\sphinxAtStartPar
1. 右键测点名称2. 点击"重命名"3. 输入新名称
\sphinxbeforeendvarwidth
\end{varwidth}%
&\begin{varwidth}[t]{\sphinxcolwidth{1}{5}}
\sphinxAtStartPar
名称应唯一且易于识别
\sphinxbeforeendvarwidth
\end{varwidth}%
\\
\sphinxhline\begin{varwidth}[t]{\sphinxcolwidth{1}{5}}
\sphinxAtStartPar
复制
\sphinxbeforeendvarwidth
\end{varwidth}%
&\begin{varwidth}[t]{\sphinxcolwidth{1}{5}}
\sphinxAtStartPar
复制测点
\sphinxbeforeendvarwidth
\end{varwidth}%
&\begin{varwidth}[t]{\sphinxcolwidth{1}{5}}
\sphinxAtStartPar
当需要在分组之间移动测点时
\sphinxbeforeendvarwidth
\end{varwidth}%
&\begin{varwidth}[t]{\sphinxcolwidth{1}{5}}
\sphinxAtStartPar
1. 右键测点名称2. 点击"复制"3. 右键目标分组粘贴
\sphinxbeforeendvarwidth
\end{varwidth}%
&\begin{varwidth}[t]{\sphinxcolwidth{1}{5}}
\sphinxAtStartPar
复制所有数据
\sphinxbeforeendvarwidth
\end{varwidth}%
\\
\sphinxhline\begin{varwidth}[t]{\sphinxcolwidth{1}{5}}
\sphinxAtStartPar
删除
\sphinxbeforeendvarwidth
\end{varwidth}%
&\begin{varwidth}[t]{\sphinxcolwidth{1}{5}}
\sphinxAtStartPar
删除测点
\sphinxbeforeendvarwidth
\end{varwidth}%
&\begin{varwidth}[t]{\sphinxcolwidth{1}{5}}
\sphinxAtStartPar
当测点数据不再需要时
\sphinxbeforeendvarwidth
\end{varwidth}%
&\begin{varwidth}[t]{\sphinxcolwidth{1}{5}}
\sphinxAtStartPar
1. 右键测点名称2. 点击"删除"3. 确认删除
\sphinxbeforeendvarwidth
\end{varwidth}%
&\begin{varwidth}[t]{\sphinxcolwidth{1}{5}}
\sphinxAtStartPar
⚠️ 删除后无法恢复
\sphinxbeforeendvarwidth
\end{varwidth}%
\\
\sphinxhline\begin{varwidth}[t]{\sphinxcolwidth{1}{5}}
\sphinxAtStartPar
标记完成
\sphinxbeforeendvarwidth
\end{varwidth}%
&\begin{varwidth}[t]{\sphinxcolwidth{1}{5}}
\sphinxAtStartPar
标记测点已完成处理
\sphinxbeforeendvarwidth
\end{varwidth}%
&\begin{varwidth}[t]{\sphinxcolwidth{1}{5}}
\sphinxAtStartPar
当测点处理完成并已检查时
\sphinxbeforeendvarwidth
\end{varwidth}%
&\begin{varwidth}[t]{\sphinxcolwidth{1}{5}}
\sphinxAtStartPar
1. 右键测点名称2. 点击"标记完成"
\sphinxbeforeendvarwidth
\end{varwidth}%
&\begin{varwidth}[t]{\sphinxcolwidth{1}{5}}
\sphinxAtStartPar
完成的测点会显示特殊标记
\sphinxbeforeendvarwidth
\end{varwidth}%
\\
\sphinxbottomrule
\end{tabulary}
\sphinxtableafterendhook\par
\sphinxattableend\end{savenotes}

\sphinxAtStartPar
\sphinxstylestrong{2. 测点设置}


\begin{savenotes}\sphinxattablestart
\sphinxthistablewithglobalstyle
\centering
\begin{tabulary}{\linewidth}[t]{TTTTT}
\sphinxtoprule
\begin{varwidth}[t]{\sphinxcolwidth{1}{5}}
\sphinxstyletheadfamily \sphinxAtStartPar
功能
\sphinxbeforeendvarwidth
\end{varwidth}%
&\begin{varwidth}[t]{\sphinxcolwidth{1}{5}}
\sphinxstyletheadfamily \sphinxAtStartPar
说明
\sphinxbeforeendvarwidth
\end{varwidth}%
&\begin{varwidth}[t]{\sphinxcolwidth{1}{5}}
\sphinxstyletheadfamily \sphinxAtStartPar
使用场景
\sphinxbeforeendvarwidth
\end{varwidth}%
&\begin{varwidth}[t]{\sphinxcolwidth{1}{5}}
\sphinxstyletheadfamily \sphinxAtStartPar
操作步骤
\sphinxbeforeendvarwidth
\end{varwidth}%
&\begin{varwidth}[t]{\sphinxcolwidth{1}{5}}
\sphinxstyletheadfamily \sphinxAtStartPar
注意事项
\sphinxbeforeendvarwidth
\end{varwidth}%
\\
\sphinxmidrule
\sphinxtableatstartofbodyhook\begin{varwidth}[t]{\sphinxcolwidth{1}{5}}
\sphinxAtStartPar
EH类型
\sphinxbeforeendvarwidth
\end{varwidth}%
&\begin{varwidth}[t]{\sphinxcolwidth{1}{5}}
\sphinxAtStartPar
设置电磁场类型（0\sphinxhyphen{}3）
\sphinxbeforeendvarwidth
\end{varwidth}%
&\begin{varwidth}[t]{\sphinxcolwidth{1}{5}}
\sphinxAtStartPar
当需要指定电场和磁场的组合方式时
\sphinxbeforeendvarwidth
\end{varwidth}%
&\begin{varwidth}[t]{\sphinxcolwidth{1}{5}}
\sphinxAtStartPar
1. 右键测点名称2. 选择"EH类型"3. 选择0/1/2/3
\sphinxbeforeendvarwidth
\end{varwidth}%
&\begin{varwidth}[t]{\sphinxcolwidth{1}{5}}
\sphinxAtStartPar
0=ExEyHxHyHz1=HxHyHz2=ExEy3=ExEyHxHy
\sphinxbeforeendvarwidth
\end{varwidth}%
\\
\sphinxhline\begin{varwidth}[t]{\sphinxcolwidth{1}{5}}
\sphinxAtStartPar
设置H测点
\sphinxbeforeendvarwidth
\end{varwidth}%
&\begin{varwidth}[t]{\sphinxcolwidth{1}{5}}
\sphinxAtStartPar
设置为H测点
\sphinxbeforeendvarwidth
\end{varwidth}%
&\begin{varwidth}[t]{\sphinxcolwidth{1}{5}}
\sphinxAtStartPar
当此测点作为H测点使用时
\sphinxbeforeendvarwidth
\end{varwidth}%
&\begin{varwidth}[t]{\sphinxcolwidth{1}{5}}
\sphinxAtStartPar
1. 右键测点名称2. 点击"设置H测点"
\sphinxbeforeendvarwidth
\end{varwidth}%
&\begin{varwidth}[t]{\sphinxcolwidth{1}{5}}
\sphinxAtStartPar
用于特定处理流程
\sphinxbeforeendvarwidth
\end{varwidth}%
\\
\sphinxhline\begin{varwidth}[t]{\sphinxcolwidth{1}{5}}
\sphinxAtStartPar
设置为H组
\sphinxbeforeendvarwidth
\end{varwidth}%
&\begin{varwidth}[t]{\sphinxcolwidth{1}{5}}
\sphinxAtStartPar
设置为H组
\sphinxbeforeendvarwidth
\end{varwidth}%
&\begin{varwidth}[t]{\sphinxcolwidth{1}{5}}
\sphinxAtStartPar
当此测点作为H组的参考时
\sphinxbeforeendvarwidth
\end{varwidth}%
&\begin{varwidth}[t]{\sphinxcolwidth{1}{5}}
\sphinxAtStartPar
1. 右键测点名称2. 点击"设置为H组"
\sphinxbeforeendvarwidth
\end{varwidth}%
&\begin{varwidth}[t]{\sphinxcolwidth{1}{5}}
\sphinxAtStartPar
会影响同组其他测点
\sphinxbeforeendvarwidth
\end{varwidth}%
\\
\sphinxhline\begin{varwidth}[t]{\sphinxcolwidth{1}{5}}
\sphinxAtStartPar
设置为远参考组
\sphinxbeforeendvarwidth
\end{varwidth}%
&\begin{varwidth}[t]{\sphinxcolwidth{1}{5}}
\sphinxAtStartPar
设置为远参考组
\sphinxbeforeendvarwidth
\end{varwidth}%
&\begin{varwidth}[t]{\sphinxcolwidth{1}{5}}
\sphinxAtStartPar
当此测点作为远参考站时
\sphinxbeforeendvarwidth
\end{varwidth}%
&\begin{varwidth}[t]{\sphinxcolwidth{1}{5}}
\sphinxAtStartPar
1. 右键测点名称2. 点击"设置为远参考组"
\sphinxbeforeendvarwidth
\end{varwidth}%
&\begin{varwidth}[t]{\sphinxcolwidth{1}{5}}
\sphinxAtStartPar
远参考站应远离测区
\sphinxbeforeendvarwidth
\end{varwidth}%
\\
\sphinxhline\begin{varwidth}[t]{\sphinxcolwidth{1}{5}}
\sphinxAtStartPar
旋转恢复
\sphinxbeforeendvarwidth
\end{varwidth}%
&\begin{varwidth}[t]{\sphinxcolwidth{1}{5}}
\sphinxAtStartPar
恢复原始旋转角度
\sphinxbeforeendvarwidth
\end{varwidth}%
&\begin{varwidth}[t]{\sphinxcolwidth{1}{5}}
\sphinxAtStartPar
当需要取消坐标旋转时
\sphinxbeforeendvarwidth
\end{varwidth}%
&\begin{varwidth}[t]{\sphinxcolwidth{1}{5}}
\sphinxAtStartPar
1. 右键测点名称2. 点击"旋转恢复"
\sphinxbeforeendvarwidth
\end{varwidth}%
&\begin{varwidth}[t]{\sphinxcolwidth{1}{5}}
\sphinxAtStartPar
恢复到采集时的方向
\sphinxbeforeendvarwidth
\end{varwidth}%
\\
\sphinxbottomrule
\end{tabulary}
\sphinxtableafterendhook\par
\sphinxattableend\end{savenotes}

\sphinxAtStartPar
\sphinxstylestrong{3. 数据质量}


\begin{savenotes}\sphinxattablestart
\sphinxthistablewithglobalstyle
\centering
\begin{tabulary}{\linewidth}[t]{TTTT}
\sphinxtoprule
\begin{varwidth}[t]{\sphinxcolwidth{1}{4}}
\sphinxstyletheadfamily \sphinxAtStartPar
功能
\sphinxbeforeendvarwidth
\end{varwidth}%
&\begin{varwidth}[t]{\sphinxcolwidth{1}{4}}
\sphinxstyletheadfamily \sphinxAtStartPar
说明
\sphinxbeforeendvarwidth
\end{varwidth}%
&\begin{varwidth}[t]{\sphinxcolwidth{1}{4}}
\sphinxstyletheadfamily \sphinxAtStartPar
使用场景
\sphinxbeforeendvarwidth
\end{varwidth}%
&\begin{varwidth}[t]{\sphinxcolwidth{1}{4}}
\sphinxstyletheadfamily \sphinxAtStartPar
质量标准
\sphinxbeforeendvarwidth
\end{varwidth}%
\\
\sphinxmidrule
\sphinxtableatstartofbodyhook\begin{varwidth}[t]{\sphinxcolwidth{1}{4}}
\sphinxAtStartPar
数据质量标记
\sphinxbeforeendvarwidth
\end{varwidth}%
&\begin{varwidth}[t]{\sphinxcolwidth{1}{4}}
\sphinxAtStartPar
设置质量等级
\sphinxbeforeendvarwidth
\end{varwidth}%
&\begin{varwidth}[t]{\sphinxcolwidth{1}{4}}
\sphinxAtStartPar
当需要标记测点数据质量时
\sphinxbeforeendvarwidth
\end{varwidth}%
&\begin{varwidth}[t]{\sphinxcolwidth{1}{4}}
\sphinxAtStartPar
见下表
\sphinxbeforeendvarwidth
\end{varwidth}%
\\
\sphinxhline\begin{varwidth}[t]{\sphinxcolwidth{1}{4}}
\sphinxAtStartPar
质量0
\sphinxbeforeendvarwidth
\end{varwidth}%
&\begin{varwidth}[t]{\sphinxcolwidth{1}{4}}
\sphinxAtStartPar
未评估
\sphinxbeforeendvarwidth
\end{varwidth}%
&\begin{varwidth}[t]{\sphinxcolwidth{1}{4}}
\sphinxAtStartPar
默认状态
\sphinxbeforeendvarwidth
\end{varwidth}%
&\begin{varwidth}[t]{\sphinxcolwidth{1}{4}}
\sphinxAtStartPar
数据尚未评估
\sphinxbeforeendvarwidth
\end{varwidth}%
\\
\sphinxhline\begin{varwidth}[t]{\sphinxcolwidth{1}{4}}
\sphinxAtStartPar
质量1
\sphinxbeforeendvarwidth
\end{varwidth}%
&\begin{varwidth}[t]{\sphinxcolwidth{1}{4}}
\sphinxAtStartPar
优秀
\sphinxbeforeendvarwidth
\end{varwidth}%
&\begin{varwidth}[t]{\sphinxcolwidth{1}{4}}
\sphinxAtStartPar
相干性>0.8，曲线平滑，无异常点
\sphinxbeforeendvarwidth
\end{varwidth}%
&\begin{varwidth}[t]{\sphinxcolwidth{1}{4}}
\sphinxAtStartPar
可直接用于反演
\sphinxbeforeendvarwidth
\end{varwidth}%
\\
\sphinxhline\begin{varwidth}[t]{\sphinxcolwidth{1}{4}}
\sphinxAtStartPar
质量2
\sphinxbeforeendvarwidth
\end{varwidth}%
&\begin{varwidth}[t]{\sphinxcolwidth{1}{4}}
\sphinxAtStartPar
良好
\sphinxbeforeendvarwidth
\end{varwidth}%
&\begin{varwidth}[t]{\sphinxcolwidth{1}{4}}
\sphinxAtStartPar
相干性0.6\sphinxhyphen{}0.8，曲线基本平滑
\sphinxbeforeendvarwidth
\end{varwidth}%
&\begin{varwidth}[t]{\sphinxcolwidth{1}{4}}
\sphinxAtStartPar
可用于反演，需注意
\sphinxbeforeendvarwidth
\end{varwidth}%
\\
\sphinxhline\begin{varwidth}[t]{\sphinxcolwidth{1}{4}}
\sphinxAtStartPar
质量3
\sphinxbeforeendvarwidth
\end{varwidth}%
&\begin{varwidth}[t]{\sphinxcolwidth{1}{4}}
\sphinxAtStartPar
一般
\sphinxbeforeendvarwidth
\end{varwidth}%
&\begin{varwidth}[t]{\sphinxcolwidth{1}{4}}
\sphinxAtStartPar
相干性0.5\sphinxhyphen{}0.6，曲线有波动
\sphinxbeforeendvarwidth
\end{varwidth}%
&\begin{varwidth}[t]{\sphinxcolwidth{1}{4}}
\sphinxAtStartPar
谨慎使用，建议复查
\sphinxbeforeendvarwidth
\end{varwidth}%
\\
\sphinxhline\begin{varwidth}[t]{\sphinxcolwidth{1}{4}}
\sphinxAtStartPar
质量4
\sphinxbeforeendvarwidth
\end{varwidth}%
&\begin{varwidth}[t]{\sphinxcolwidth{1}{4}}
\sphinxAtStartPar
较差
\sphinxbeforeendvarwidth
\end{varwidth}%
&\begin{varwidth}[t]{\sphinxcolwidth{1}{4}}
\sphinxAtStartPar
相干性<0.5，曲线畸变明显
\sphinxbeforeendvarwidth
\end{varwidth}%
&\begin{varwidth}[t]{\sphinxcolwidth{1}{4}}
\sphinxAtStartPar
不建议使用，需重新处理
\sphinxbeforeendvarwidth
\end{varwidth}%
\\
\sphinxbottomrule
\end{tabulary}
\sphinxtableafterendhook\par
\sphinxattableend\end{savenotes}
\begin{quote}

\sphinxAtStartPar
\sphinxstylestrong{建议}：处理完成后立即标记数据质量，方便后续筛选和批量导出。
\end{quote}

\sphinxAtStartPar
\sphinxstylestrong{4. 数据处理}


\begin{savenotes}\sphinxattablestart
\sphinxthistablewithglobalstyle
\centering
\begin{tabulary}{\linewidth}[t]{TTTTT}
\sphinxtoprule
\begin{varwidth}[t]{\sphinxcolwidth{1}{5}}
\sphinxstyletheadfamily \sphinxAtStartPar
功能
\sphinxbeforeendvarwidth
\end{varwidth}%
&\begin{varwidth}[t]{\sphinxcolwidth{1}{5}}
\sphinxstyletheadfamily \sphinxAtStartPar
说明
\sphinxbeforeendvarwidth
\end{varwidth}%
&\begin{varwidth}[t]{\sphinxcolwidth{1}{5}}
\sphinxstyletheadfamily \sphinxAtStartPar
使用场景
\sphinxbeforeendvarwidth
\end{varwidth}%
&\begin{varwidth}[t]{\sphinxcolwidth{1}{5}}
\sphinxstyletheadfamily \sphinxAtStartPar
操作步骤
\sphinxbeforeendvarwidth
\end{varwidth}%
&\begin{varwidth}[t]{\sphinxcolwidth{1}{5}}
\sphinxstyletheadfamily \sphinxAtStartPar
注意事项
\sphinxbeforeendvarwidth
\end{varwidth}%
\\
\sphinxmidrule
\sphinxtableatstartofbodyhook\begin{varwidth}[t]{\sphinxcolwidth{1}{5}}
\sphinxAtStartPar
处理
\sphinxbeforeendvarwidth
\end{varwidth}%
&\begin{varwidth}[t]{\sphinxcolwidth{1}{5}}
\sphinxAtStartPar
单点FFT处理
\sphinxbeforeendvarwidth
\end{varwidth}%
&\begin{varwidth}[t]{\sphinxcolwidth{1}{5}}
\sphinxAtStartPar
当需要对单个测点进行处理时
\sphinxbeforeendvarwidth
\end{varwidth}%
&\begin{varwidth}[t]{\sphinxcolwidth{1}{5}}
\sphinxAtStartPar
1. 右键测点名称2. 点击"处理"3. 设置FFT参数4. 执行处理
\sphinxbeforeendvarwidth
\end{varwidth}%
&\begin{varwidth}[t]{\sphinxcolwidth{1}{5}}
\sphinxAtStartPar
适合测试或特殊处理
\sphinxbeforeendvarwidth
\end{varwidth}%
\\
\sphinxhline\begin{varwidth}[t]{\sphinxcolwidth{1}{5}}
\sphinxAtStartPar
远参考计算
\sphinxbeforeendvarwidth
\end{varwidth}%
&\begin{varwidth}[t]{\sphinxcolwidth{1}{5}}
\sphinxAtStartPar
设置远参考站
\sphinxbeforeendvarwidth
\end{varwidth}%
&\begin{varwidth}[t]{\sphinxcolwidth{1}{5}}
\sphinxAtStartPar
当需要指定远参考站时
\sphinxbeforeendvarwidth
\end{varwidth}%
&\begin{varwidth}[t]{\sphinxcolwidth{1}{5}}
\sphinxAtStartPar
1. 右键测点名称2. 点击"远参考计算"3. 选择远参考站
\sphinxbeforeendvarwidth
\end{varwidth}%
&\begin{varwidth}[t]{\sphinxcolwidth{1}{5}}
\sphinxAtStartPar
远参考站需时间同步
\sphinxbeforeendvarwidth
\end{varwidth}%
\\
\sphinxhline\begin{varwidth}[t]{\sphinxcolwidth{1}{5}}
\sphinxAtStartPar
频谱组编辑
\sphinxbeforeendvarwidth
\end{varwidth}%
&\begin{varwidth}[t]{\sphinxcolwidth{1}{5}}
\sphinxAtStartPar
编辑频谱数据
\sphinxbeforeendvarwidth
\end{varwidth}%
&\begin{varwidth}[t]{\sphinxcolwidth{1}{5}}
\sphinxAtStartPar
当需要手动编辑频谱点时
\sphinxbeforeendvarwidth
\end{varwidth}%
&\begin{varwidth}[t]{\sphinxcolwidth{1}{5}}
\sphinxAtStartPar
1. 右键测点名称2. 点击"频谱组编辑"3. 在编辑器中修改
\sphinxbeforeendvarwidth
\end{varwidth}%
&\begin{varwidth}[t]{\sphinxcolwidth{1}{5}}
\sphinxAtStartPar
可删除异常频点
\sphinxbeforeendvarwidth
\end{varwidth}%
\\
\sphinxhline\begin{varwidth}[t]{\sphinxcolwidth{1}{5}}
\sphinxAtStartPar
加载FC
\sphinxbeforeendvarwidth
\end{varwidth}%
&\begin{varwidth}[t]{\sphinxcolwidth{1}{5}}
\sphinxAtStartPar
加载傅里叶系数
\sphinxbeforeendvarwidth
\end{varwidth}%
&\begin{varwidth}[t]{\sphinxcolwidth{1}{5}}
\sphinxAtStartPar
当需要加载已有的FC文件时
\sphinxbeforeendvarwidth
\end{varwidth}%
&\begin{varwidth}[t]{\sphinxcolwidth{1}{5}}
\sphinxAtStartPar
1. 右键测点名称2. 点击"加载FC"3. 选择FC文件
\sphinxbeforeendvarwidth
\end{varwidth}%
&\begin{varwidth}[t]{\sphinxcolwidth{1}{5}}
\sphinxAtStartPar
.FC或.FT文件
\sphinxbeforeendvarwidth
\end{varwidth}%
\\
\sphinxbottomrule
\end{tabulary}
\sphinxtableafterendhook\par
\sphinxattableend\end{savenotes}

\sphinxAtStartPar
\sphinxstylestrong{5. 数据导出}


\begin{savenotes}\sphinxattablestart
\sphinxthistablewithglobalstyle
\centering
\begin{tabulary}{\linewidth}[t]{TTTTT}
\sphinxtoprule
\begin{varwidth}[t]{\sphinxcolwidth{1}{5}}
\sphinxstyletheadfamily \sphinxAtStartPar
功能
\sphinxbeforeendvarwidth
\end{varwidth}%
&\begin{varwidth}[t]{\sphinxcolwidth{1}{5}}
\sphinxstyletheadfamily \sphinxAtStartPar
说明
\sphinxbeforeendvarwidth
\end{varwidth}%
&\begin{varwidth}[t]{\sphinxcolwidth{1}{5}}
\sphinxstyletheadfamily \sphinxAtStartPar
使用场景
\sphinxbeforeendvarwidth
\end{varwidth}%
&\begin{varwidth}[t]{\sphinxcolwidth{1}{5}}
\sphinxstyletheadfamily \sphinxAtStartPar
输出文件
\sphinxbeforeendvarwidth
\end{varwidth}%
&\begin{varwidth}[t]{\sphinxcolwidth{1}{5}}
\sphinxstyletheadfamily \sphinxAtStartPar
注意事项
\sphinxbeforeendvarwidth
\end{varwidth}%
\\
\sphinxmidrule
\sphinxtableatstartofbodyhook\begin{varwidth}[t]{\sphinxcolwidth{1}{5}}
\sphinxAtStartPar
\sphinxstylestrong{📊 导出绘图数据}
\sphinxbeforeendvarwidth
\end{varwidth}%
&\begin{varwidth}[t]{\sphinxcolwidth{1}{5}}
\sphinxAtStartPar
导出测点参数数据供第三方软件绘图
\sphinxbeforeendvarwidth
\end{varwidth}%
&\begin{varwidth}[t]{\sphinxcolwidth{1}{5}}
\sphinxAtStartPar
当需要使用Golden绘图、Surfer等软件成图时
\sphinxbeforeendvarwidth
\end{varwidth}%
&\begin{varwidth}[t]{\sphinxcolwidth{1}{5}}
\sphinxAtStartPar
\sphinxhyphen{} \{测点名\}\sphinxhyphen{}\{参数名\}.dat\sphinxhyphen{} \{测点名\}\sphinxhyphen{}lgRP.dat
\sphinxbeforeendvarwidth
\end{varwidth}%
&\begin{varwidth}[t]{\sphinxcolwidth{1}{5}}
\sphinxAtStartPar
推荐用于单点详细绘图
\sphinxbeforeendvarwidth
\end{varwidth}%
\\
\sphinxhline\begin{varwidth}[t]{\sphinxcolwidth{1}{5}}
\sphinxAtStartPar
导出EDI
\sphinxbeforeendvarwidth
\end{varwidth}%
&\begin{varwidth}[t]{\sphinxcolwidth{1}{5}}
\sphinxAtStartPar
导出标准EDI文件
\sphinxbeforeendvarwidth
\end{varwidth}%
&\begin{varwidth}[t]{\sphinxcolwidth{1}{5}}
\sphinxAtStartPar
当需要将处理结果提供给反演软件时
\sphinxbeforeendvarwidth
\end{varwidth}%
&\begin{varwidth}[t]{\sphinxcolwidth{1}{5}}
\sphinxAtStartPar
.edi文件
\sphinxbeforeendvarwidth
\end{varwidth}%
&\begin{varwidth}[t]{\sphinxcolwidth{1}{5}}
\sphinxAtStartPar
建议先检查数据质量
\sphinxbeforeendvarwidth
\end{varwidth}%
\\
\sphinxhline\begin{varwidth}[t]{\sphinxcolwidth{1}{5}}
\sphinxAtStartPar
导出SpeEDI
\sphinxbeforeendvarwidth
\end{varwidth}%
&\begin{varwidth}[t]{\sphinxcolwidth{1}{5}}
\sphinxAtStartPar
导出SpeEDI格式
\sphinxbeforeendvarwidth
\end{varwidth}%
&\begin{varwidth}[t]{\sphinxcolwidth{1}{5}}
\sphinxAtStartPar
当需要与支持SpeEDI的软件交换数据时
\sphinxbeforeendvarwidth
\end{varwidth}%
&\begin{varwidth}[t]{\sphinxcolwidth{1}{5}}
\sphinxAtStartPar
.edi文件（SpeEDI格式）
\sphinxbeforeendvarwidth
\end{varwidth}%
&\begin{varwidth}[t]{\sphinxcolwidth{1}{5}}
\sphinxAtStartPar
标准格式
\sphinxbeforeendvarwidth
\end{varwidth}%
\\
\sphinxhline\begin{varwidth}[t]{\sphinxcolwidth{1}{5}}
\sphinxAtStartPar
导出ZTEDI
\sphinxbeforeendvarwidth
\end{varwidth}%
&\begin{varwidth}[t]{\sphinxcolwidth{1}{5}}
\sphinxAtStartPar
导出ZTEDI格式
\sphinxbeforeendvarwidth
\end{varwidth}%
&\begin{varwidth}[t]{\sphinxcolwidth{1}{5}}
\sphinxAtStartPar
当需要完整张量信息时
\sphinxbeforeendvarwidth
\end{varwidth}%
&\begin{varwidth}[t]{\sphinxcolwidth{1}{5}}
\sphinxAtStartPar
.edi文件（ZTEDI格式）
\sphinxbeforeendvarwidth
\end{varwidth}%
&\begin{varwidth}[t]{\sphinxcolwidth{1}{5}}
\sphinxAtStartPar
含阻抗和倾子
\sphinxbeforeendvarwidth
\end{varwidth}%
\\
\sphinxhline\begin{varwidth}[t]{\sphinxcolwidth{1}{5}}
\sphinxAtStartPar
导出MTpyEDI
\sphinxbeforeendvarwidth
\end{varwidth}%
&\begin{varwidth}[t]{\sphinxcolwidth{1}{5}}
\sphinxAtStartPar
导出MTpyEDI格式
\sphinxbeforeendvarwidth
\end{varwidth}%
&\begin{varwidth}[t]{\sphinxcolwidth{1}{5}}
\sphinxAtStartPar
当使用Python处理时
\sphinxbeforeendvarwidth
\end{varwidth}%
&\begin{varwidth}[t]{\sphinxcolwidth{1}{5}}
\sphinxAtStartPar
.edi文件（MTpyEDI格式）
\sphinxbeforeendvarwidth
\end{varwidth}%
&\begin{varwidth}[t]{\sphinxcolwidth{1}{5}}
\sphinxAtStartPar
Python兼容
\sphinxbeforeendvarwidth
\end{varwidth}%
\\
\sphinxhline\begin{varwidth}[t]{\sphinxcolwidth{1}{5}}
\sphinxAtStartPar
导出PLTEDI
\sphinxbeforeendvarwidth
\end{varwidth}%
&\begin{varwidth}[t]{\sphinxcolwidth{1}{5}}
\sphinxAtStartPar
导出PLT格式
\sphinxbeforeendvarwidth
\end{varwidth}%
&\begin{varwidth}[t]{\sphinxcolwidth{1}{5}}
\sphinxAtStartPar
当绘图软件需要时
\sphinxbeforeendvarwidth
\end{varwidth}%
&\begin{varwidth}[t]{\sphinxcolwidth{1}{5}}
\sphinxAtStartPar
.plt文件
\sphinxbeforeendvarwidth
\end{varwidth}%
&\begin{varwidth}[t]{\sphinxcolwidth{1}{5}}
\sphinxAtStartPar
含视电阻率/相位
\sphinxbeforeendvarwidth
\end{varwidth}%
\\
\sphinxhline\begin{varwidth}[t]{\sphinxcolwidth{1}{5}}
\sphinxAtStartPar
导出RhoPlus预测
\sphinxbeforeendvarwidth
\end{varwidth}%
&\begin{varwidth}[t]{\sphinxcolwidth{1}{5}}
\sphinxAtStartPar
导出RhoPlus预测数据
\sphinxbeforeendvarwidth
\end{varwidth}%
&\begin{varwidth}[t]{\sphinxcolwidth{1}{5}}
\sphinxAtStartPar
当需要验证数据一致性时
\sphinxbeforeendvarwidth
\end{varwidth}%
&\begin{varwidth}[t]{\sphinxcolwidth{1}{5}}
\sphinxAtStartPar
\sphinxhyphen{} \{测点名\}\sphinxhyphen{}RhoPlus.dat\sphinxhyphen{} \{测点名\}\sphinxhyphen{}Predict.dat
\sphinxbeforeendvarwidth
\end{varwidth}%
&\begin{varwidth}[t]{\sphinxcolwidth{1}{5}}
\sphinxAtStartPar
用于1D模型验证
\sphinxbeforeendvarwidth
\end{varwidth}%
\\
\sphinxhline\begin{varwidth}[t]{\sphinxcolwidth{1}{5}}
\sphinxAtStartPar
导出GMT时序
\sphinxbeforeendvarwidth
\end{varwidth}%
&\begin{varwidth}[t]{\sphinxcolwidth{1}{5}}
\sphinxAtStartPar
导出GMT格式时间序列
\sphinxbeforeendvarwidth
\end{varwidth}%
&\begin{varwidth}[t]{\sphinxcolwidth{1}{5}}
\sphinxAtStartPar
当需要使用GMT绘图时
\sphinxbeforeendvarwidth
\end{varwidth}%
&\begin{varwidth}[t]{\sphinxcolwidth{1}{5}}
\sphinxAtStartPar
.dat文件（GMT格式）
\sphinxbeforeendvarwidth
\end{varwidth}%
&\begin{varwidth}[t]{\sphinxcolwidth{1}{5}}
\sphinxAtStartPar
时间序列数据
\sphinxbeforeendvarwidth
\end{varwidth}%
\\
\sphinxhline\begin{varwidth}[t]{\sphinxcolwidth{1}{5}}
\sphinxAtStartPar
导出HTML设置
\sphinxbeforeendvarwidth
\end{varwidth}%
&\begin{varwidth}[t]{\sphinxcolwidth{1}{5}}
\sphinxAtStartPar
导出HTML配置报告
\sphinxbeforeendvarwidth
\end{varwidth}%
&\begin{varwidth}[t]{\sphinxcolwidth{1}{5}}
\sphinxAtStartPar
当需要生成测点配置的网页报告时
\sphinxbeforeendvarwidth
\end{varwidth}%
&\begin{varwidth}[t]{\sphinxcolwidth{1}{5}}
\sphinxAtStartPar
.html文件
\sphinxbeforeendvarwidth
\end{varwidth}%
&\begin{varwidth}[t]{\sphinxcolwidth{1}{5}}
\sphinxAtStartPar
可在浏览器中查看
\sphinxbeforeendvarwidth
\end{varwidth}%
\\
\sphinxbottomrule
\end{tabulary}
\sphinxtableafterendhook\par
\sphinxattableend\end{savenotes}

\sphinxAtStartPar
\sphinxstylestrong{6. 数据显示}


\begin{savenotes}\sphinxattablestart
\sphinxthistablewithglobalstyle
\centering
\begin{tabulary}{\linewidth}[t]{TTTT}
\sphinxtoprule
\begin{varwidth}[t]{\sphinxcolwidth{1}{4}}
\sphinxstyletheadfamily \sphinxAtStartPar
功能
\sphinxbeforeendvarwidth
\end{varwidth}%
&\begin{varwidth}[t]{\sphinxcolwidth{1}{4}}
\sphinxstyletheadfamily \sphinxAtStartPar
说明
\sphinxbeforeendvarwidth
\end{varwidth}%
&\begin{varwidth}[t]{\sphinxcolwidth{1}{4}}
\sphinxstyletheadfamily \sphinxAtStartPar
使用场景
\sphinxbeforeendvarwidth
\end{varwidth}%
&\begin{varwidth}[t]{\sphinxcolwidth{1}{4}}
\sphinxstyletheadfamily \sphinxAtStartPar
操作步骤
\sphinxbeforeendvarwidth
\end{varwidth}%
\\
\sphinxmidrule
\sphinxtableatstartofbodyhook\begin{varwidth}[t]{\sphinxcolwidth{1}{4}}
\sphinxAtStartPar
添加时序显示
\sphinxbeforeendvarwidth
\end{varwidth}%
&\begin{varwidth}[t]{\sphinxcolwidth{1}{4}}
\sphinxAtStartPar
在时序窗口显示
\sphinxbeforeendvarwidth
\end{varwidth}%
&\begin{varwidth}[t]{\sphinxcolwidth{1}{4}}
\sphinxAtStartPar
当需要查看原始时间序列时
\sphinxbeforeendvarwidth
\end{varwidth}%
&\begin{varwidth}[t]{\sphinxcolwidth{1}{4}}
\sphinxAtStartPar
1. 右键测点名称2. 点击"添加时序显示"
\sphinxbeforeendvarwidth
\end{varwidth}%
\\
\sphinxhline\begin{varwidth}[t]{\sphinxcolwidth{1}{4}}
\sphinxAtStartPar
添加标定显示
\sphinxbeforeendvarwidth
\end{varwidth}%
&\begin{varwidth}[t]{\sphinxcolwidth{1}{4}}
\sphinxAtStartPar
在标定窗口显示
\sphinxbeforeendvarwidth
\end{varwidth}%
&\begin{varwidth}[t]{\sphinxcolwidth{1}{4}}
\sphinxAtStartPar
当需要查看标定曲线时
\sphinxbeforeendvarwidth
\end{varwidth}%
&\begin{varwidth}[t]{\sphinxcolwidth{1}{4}}
\sphinxAtStartPar
1. 右键测点名称2. 点击"添加标定显示"
\sphinxbeforeendvarwidth
\end{varwidth}%
\\
\sphinxbottomrule
\end{tabulary}
\sphinxtableafterendhook\par
\sphinxattableend\end{savenotes}


\subparagraph{右键菜单使用技巧}
\label{\detokenize{chapters/chapter1:id20}}\begin{enumerate}
\sphinxsetlistlabels{\arabic}{enumi}{enumii}{}{.}%
\item {} 
\sphinxAtStartPar
\sphinxstylestrong{批量操作优先使用分组级菜单}
\begin{itemize}
\item {} 
\sphinxAtStartPar
在分组级别进行批量处理比单点处理效率高得多

\item {} 
\sphinxAtStartPar
使用"处理全部"或"处理剩余"可自动处理所有测点

\end{itemize}

\item {} 
\sphinxAtStartPar
\sphinxstylestrong{合理使用数据导出功能}
\begin{itemize}
\item {} 
\sphinxAtStartPar
需要第三方软件绘图时，使用"导出绘图数据"功能

\item {} 
\sphinxAtStartPar
交付数据给反演团队时，使用"导出EDI"功能

\item {} 
\sphinxAtStartPar
需要GIS成图时，使用"导出KML/KMZ"功能

\end{itemize}

\item {} 
\sphinxAtStartPar
\sphinxstylestrong{及时标记数据质量}
\begin{itemize}
\item {} 
\sphinxAtStartPar
处理完成后立即标记数据质量

\item {} 
\sphinxAtStartPar
使用质量标记可快速筛选优质数据

\item {} 
\sphinxAtStartPar
批量导出时可按质量等级筛选

\end{itemize}

\item {} 
\sphinxAtStartPar
\sphinxstylestrong{使用复选框进行多选}
\begin{itemize}
\item {} 
\sphinxAtStartPar
在项目菜单中启用"复选框显示"

\item {} 
\sphinxAtStartPar
可多选测点进行批量操作

\item {} 
\sphinxAtStartPar
按住Ctrl键也可多选

\end{itemize}

\item {} 
\sphinxAtStartPar
\sphinxstylestrong{利用排序功能组织数据}
\begin{itemize}
\item {} 
\sphinxAtStartPar
按纬度排序可查看测线方向

\item {} 
\sphinxAtStartPar
按时间排序可追踪采集进度

\item {} 
\sphinxAtStartPar
按名称排序可快速定位测点

\end{itemize}

\item {} 
\sphinxAtStartPar
\sphinxstylestrong{注意操作顺序}
\begin{itemize}
\item {} 
\sphinxAtStartPar
先加载测点，再设置标定，最后处理

\item {} 
\sphinxAtStartPar
先处理，再筛选，最后导出

\item {} 
\sphinxAtStartPar
先标记质量，再批量导出

\end{itemize}

\end{enumerate}


\bigskip\hrule\bigskip



\subsection{1.4 工程文件结构}
\label{\detokenize{chapters/chapter1:id21}}

\subsubsection{1.4.1 目录组织}
\label{\detokenize{chapters/chapter1:id22}}
\begin{sphinxVerbatim}[commandchars=\\\{\}]
ProjectName/
├── ProjectName.MTDPE       \PYGZsh{} 工程主文件（加密）
├── Configurations/         \PYGZsh{} 配置文件
│   ├── MTDP.SETTING       \PYGZsh{} 软件设置
│   └── Lang/              \PYGZsh{} 语言文件
├── CLB/                    \PYGZsh{} 电场盒标定
├── CLC/                    \PYGZsh{} 磁传感器标定
└── temp/                   \PYGZsh{} 临时文件
\end{sphinxVerbatim}


\subsubsection{1.4.2 自动备份}
\label{\detokenize{chapters/chapter1:id23}}
\sphinxAtStartPar
MTDP自动维护工程备份：
\begin{itemize}
\item {} 
\sphinxAtStartPar
每次保存时自动创建备份

\item {} 
\sphinxAtStartPar
保留最近10个备份版本

\item {} 
\sphinxAtStartPar
备份文件位于工程目录

\end{itemize}


\bigskip\hrule\bigskip



\subsection{1.5 快速入门教程}
\label{\detokenize{chapters/chapter1:id24}}

\subsubsection{5分钟完成数据处理}
\label{\detokenize{chapters/chapter1:id25}}
\sphinxAtStartPar
\sphinxstylestrong{流程概览：}

\sphinxAtStartPar
\sphinxstylestrong{详细步骤：}

\sphinxAtStartPar
\sphinxstylestrong{步骤1：创建工程}
\begin{enumerate}
\sphinxsetlistlabels{\arabic}{enumi}{enumii}{}{.}%
\item {} 
\sphinxAtStartPar
选择 \sphinxcode{\sphinxupquote{文件 → 新建工程}}

\item {} 
\sphinxAtStartPar
设置工程名称和保存位置

\item {} 
\sphinxAtStartPar
配置工程基本信息

\end{enumerate}

\sphinxAtStartPar
\sphinxstylestrong{步骤2：创建工区和测段}
\begin{enumerate}
\sphinxsetlistlabels{\arabic}{enumi}{enumii}{}{.}%
\item {} 
\sphinxAtStartPar
右键工程 → \sphinxcode{\sphinxupquote{新建工区}}

\item {} 
\sphinxAtStartPar
右键工区 → \sphinxcode{\sphinxupquote{新建测段}}

\end{enumerate}

\sphinxAtStartPar
\sphinxstylestrong{步骤3：导入数据}
\begin{enumerate}
\sphinxsetlistlabels{\arabic}{enumi}{enumii}{}{.}%
\item {} 
\sphinxAtStartPar
右键测段 → \sphinxcode{\sphinxupquote{加载测点}}

\item {} 
\sphinxAtStartPar
选择数据目录

\item {} 
\sphinxAtStartPar
确认导入

\end{enumerate}

\sphinxAtStartPar
\sphinxstylestrong{步骤4：FFT处理}
\begin{enumerate}
\sphinxsetlistlabels{\arabic}{enumi}{enumii}{}{.}%
\item {} 
\sphinxAtStartPar
选择测段

\item {} 
\sphinxAtStartPar
右键 → \sphinxcode{\sphinxupquote{FFT处理}}

\item {} 
\sphinxAtStartPar
使用默认参数执行

\end{enumerate}

\sphinxAtStartPar
\sphinxstylestrong{步骤5：数据筛选}
\begin{enumerate}
\sphinxsetlistlabels{\arabic}{enumi}{enumii}{}{.}%
\item {} 
\sphinxAtStartPar
查看视电阻率/相位曲线

\item {} 
\sphinxAtStartPar
手动剔除异常点

\item {} 
\sphinxAtStartPar
应用Robust估计

\end{enumerate}

\sphinxAtStartPar
\sphinxstylestrong{步骤6：导出结果}
\begin{enumerate}
\sphinxsetlistlabels{\arabic}{enumi}{enumii}{}{.}%
\item {} 
\sphinxAtStartPar
右键测点 → \sphinxcode{\sphinxupquote{导出EDI}}

\item {} 
\sphinxAtStartPar
选择保存位置

\item {} 
\sphinxAtStartPar
选择导出内容（视电阻率/相位/阻抗）

\end{enumerate}


\bigskip\hrule\bigskip



\subsection{1.6 常见工作流程}
\label{\detokenize{chapters/chapter1:id26}}

\subsubsection{1.6.1 标准MT数据处理流程}
\label{\detokenize{chapters/chapter1:mt}}

\subsubsection{1.6.2 远参考处理流程}
\label{\detokenize{chapters/chapter1:id27}}
\sphinxAtStartPar
\sphinxstylestrong{适用场景}：有远参考站数据，需要消除本地噪声

\sphinxAtStartPar
\sphinxstylestrong{步骤：}
\begin{enumerate}
\sphinxsetlistlabels{\arabic}{enumi}{enumii}{}{.}%
\item {} 
\sphinxAtStartPar
\sphinxstylestrong{准备远参考数据}
\begin{itemize}
\item {} 
\sphinxAtStartPar
确保远参考站与测站时间同步

\item {} 
\sphinxAtStartPar
检查GPS时钟是否一致

\end{itemize}

\item {} 
\sphinxAtStartPar
\sphinxstylestrong{创建远参考测段}
\begin{itemize}
\item {} 
\sphinxAtStartPar
右键工区 → \sphinxcode{\sphinxupquote{新建测段}}

\item {} 
\sphinxAtStartPar
命名为 \sphinxcode{\sphinxupquote{RR\_Reference}}

\item {} 
\sphinxAtStartPar
导入远参考站数据

\end{itemize}

\item {} 
\sphinxAtStartPar
\sphinxstylestrong{关联远参考}
\begin{itemize}
\item {} 
\sphinxAtStartPar
右键本地测段 → \sphinxcode{\sphinxupquote{设置远参考}}

\item {} 
\sphinxAtStartPar
选择远参考测段

\end{itemize}

\item {} 
\sphinxAtStartPar
\sphinxstylestrong{执行处理}
\begin{itemize}
\item {} 
\sphinxAtStartPar
FFT处理时自动使用远参考

\end{itemize}

\end{enumerate}


\subsubsection{1.6.3 RMSMR强干扰处理流程}
\label{\detokenize{chapters/chapter1:rmsmr}}
\sphinxAtStartPar
\sphinxstylestrong{适用场景}：强干扰环境，常规方法效果不佳

\sphinxAtStartPar
\sphinxstylestrong{RMSMR参数设置：}


\begin{savenotes}\sphinxattablestart
\sphinxthistablewithglobalstyle
\centering
\begin{tabulary}{\linewidth}[t]{TTT}
\sphinxtoprule
\begin{varwidth}[t]{\sphinxcolwidth{1}{3}}
\sphinxstyletheadfamily \sphinxAtStartPar
参数
\sphinxbeforeendvarwidth
\end{varwidth}%
&\begin{varwidth}[t]{\sphinxcolwidth{1}{3}}
\sphinxstyletheadfamily \sphinxAtStartPar
推荐值
\sphinxbeforeendvarwidth
\end{varwidth}%
&\begin{varwidth}[t]{\sphinxcolwidth{1}{3}}
\sphinxstyletheadfamily \sphinxAtStartPar
说明
\sphinxbeforeendvarwidth
\end{varwidth}%
\\
\sphinxmidrule
\sphinxtableatstartofbodyhook\begin{varwidth}[t]{\sphinxcolwidth{1}{3}}
\sphinxAtStartPar
相干度阈值
\sphinxbeforeendvarwidth
\end{varwidth}%
&\begin{varwidth}[t]{\sphinxcolwidth{1}{3}}
\sphinxAtStartPar
0.7
\sphinxbeforeendvarwidth
\end{varwidth}%
&\begin{varwidth}[t]{\sphinxcolwidth{1}{3}}
\sphinxAtStartPar
低于此值的数据被剔除
\sphinxbeforeendvarwidth
\end{varwidth}%
\\
\sphinxhline\begin{varwidth}[t]{\sphinxcolwidth{1}{3}}
\sphinxAtStartPar
残差阈值
\sphinxbeforeendvarwidth
\end{varwidth}%
&\begin{varwidth}[t]{\sphinxcolwidth{1}{3}}
\sphinxAtStartPar
2.0σ
\sphinxbeforeendvarwidth
\end{varwidth}%
&\begin{varwidth}[t]{\sphinxcolwidth{1}{3}}
\sphinxAtStartPar
残差超过此值降权
\sphinxbeforeendvarwidth
\end{varwidth}%
\\
\sphinxhline\begin{varwidth}[t]{\sphinxcolwidth{1}{3}}
\sphinxAtStartPar
最大迭代次数
\sphinxbeforeendvarwidth
\end{varwidth}%
&\begin{varwidth}[t]{\sphinxcolwidth{1}{3}}
\sphinxAtStartPar
10
\sphinxbeforeendvarwidth
\end{varwidth}%
&\begin{varwidth}[t]{\sphinxcolwidth{1}{3}}
\sphinxAtStartPar
根据效果调整
\sphinxbeforeendvarwidth
\end{varwidth}%
\\
\sphinxbottomrule
\end{tabulary}
\sphinxtableafterendhook\par
\sphinxattableend\end{savenotes}


\subsubsection{1.6.4 批量处理流程}
\label{\detokenize{chapters/chapter1:id28}}
\sphinxAtStartPar
\sphinxstylestrong{适用场景}：多个测点需要相同处理

\sphinxAtStartPar
\sphinxstylestrong{步骤：}
\begin{enumerate}
\sphinxsetlistlabels{\arabic}{enumi}{enumii}{}{.}%
\item {} 
\sphinxAtStartPar
\sphinxstylestrong{设置处理模板}
\begin{itemize}
\item {} 
\sphinxAtStartPar
对单个测点进行完整处理

\item {} 
\sphinxAtStartPar
保存FFT参数为模板

\end{itemize}

\item {} 
\sphinxAtStartPar
\sphinxstylestrong{批量处理}
\begin{itemize}
\item {} 
\sphinxAtStartPar
右键测段 → \sphinxcode{\sphinxupquote{处理全部}}

\item {} 
\sphinxAtStartPar
选择处理方法

\item {} 
\sphinxAtStartPar
设置输出选项

\end{itemize}

\item {} 
\sphinxAtStartPar
\sphinxstylestrong{批量导出}
\begin{itemize}
\item {} 
\sphinxAtStartPar
右键工区 → \sphinxcode{\sphinxupquote{导出 → 批量导出EDI}}

\item {} 
\sphinxAtStartPar
设置导出目录和命名规则

\end{itemize}

\end{enumerate}


\subsubsection{1.6.5 数据质量评估流程}
\label{\detokenize{chapters/chapter1:id29}}
\sphinxAtStartPar
\sphinxstylestrong{检查清单：}


\subsubsection{1.6.6 典型应用场景}
\label{\detokenize{chapters/chapter1:id30}}

\begin{savenotes}\sphinxattablestart
\sphinxthistablewithglobalstyle
\centering
\begin{tabulary}{\linewidth}[t]{TTT}
\sphinxtoprule
\begin{varwidth}[t]{\sphinxcolwidth{1}{3}}
\sphinxstyletheadfamily \sphinxAtStartPar
场景
\sphinxbeforeendvarwidth
\end{varwidth}%
&\begin{varwidth}[t]{\sphinxcolwidth{1}{3}}
\sphinxstyletheadfamily \sphinxAtStartPar
推荐流程
\sphinxbeforeendvarwidth
\end{varwidth}%
&\begin{varwidth}[t]{\sphinxcolwidth{1}{3}}
\sphinxstyletheadfamily \sphinxAtStartPar
注意事项
\sphinxbeforeendvarwidth
\end{varwidth}%
\\
\sphinxmidrule
\sphinxtableatstartofbodyhook\begin{varwidth}[t]{\sphinxcolwidth{1}{3}}
\sphinxAtStartPar
\sphinxstylestrong{深部探测} (>10km)
\sphinxbeforeendvarwidth
\end{varwidth}%
&\begin{varwidth}[t]{\sphinxcolwidth{1}{3}}
\sphinxAtStartPar
长时间观测 + 远参考 + 低频筛选
\sphinxbeforeendvarwidth
\end{varwidth}%
&\begin{varwidth}[t]{\sphinxcolwidth{1}{3}}
\sphinxAtStartPar
确保低频数据充足
\sphinxbeforeendvarwidth
\end{varwidth}%
\\
\sphinxhline\begin{varwidth}[t]{\sphinxcolwidth{1}{3}}
\sphinxAtStartPar
\sphinxstylestrong{浅部探测} (<1km)
\sphinxbeforeendvarwidth
\end{varwidth}%
&\begin{varwidth}[t]{\sphinxcolwidth{1}{3}}
\sphinxAtStartPar
短时间观测 + 高频处理
\sphinxbeforeendvarwidth
\end{varwidth}%
&\begin{varwidth}[t]{\sphinxcolwidth{1}{3}}
\sphinxAtStartPar
注意文化噪声
\sphinxbeforeendvarwidth
\end{varwidth}%
\\
\sphinxhline\begin{varwidth}[t]{\sphinxcolwidth{1}{3}}
\sphinxAtStartPar
\sphinxstylestrong{城市环境}
\sphinxbeforeendvarwidth
\end{varwidth}%
&\begin{varwidth}[t]{\sphinxcolwidth{1}{3}}
\sphinxAtStartPar
RMSMR + 严格筛选
\sphinxbeforeendvarwidth
\end{varwidth}%
&\begin{varwidth}[t]{\sphinxcolwidth{1}{3}}
\sphinxAtStartPar
可能需要夜间观测
\sphinxbeforeendvarwidth
\end{varwidth}%
\\
\sphinxhline\begin{varwidth}[t]{\sphinxcolwidth{1}{3}}
\sphinxAtStartPar
\sphinxstylestrong{矿区勘探}
\sphinxbeforeendvarwidth
\end{varwidth}%
&\begin{varwidth}[t]{\sphinxcolwidth{1}{3}}
\sphinxAtStartPar
多次观测 + 远参考
\sphinxbeforeendvarwidth
\end{varwidth}%
&\begin{varwidth}[t]{\sphinxcolwidth{1}{3}}
\sphinxAtStartPar
注意矿山电磁干扰
\sphinxbeforeendvarwidth
\end{varwidth}%
\\
\sphinxhline\begin{varwidth}[t]{\sphinxcolwidth{1}{3}}
\sphinxAtStartPar
\sphinxstylestrong{地热勘探}
\sphinxbeforeendvarwidth
\end{varwidth}%
&\begin{varwidth}[t]{\sphinxcolwidth{1}{3}}
\sphinxAtStartPar
宽频带处理
\sphinxbeforeendvarwidth
\end{varwidth}%
&\begin{varwidth}[t]{\sphinxcolwidth{1}{3}}
\sphinxAtStartPar
关注低阻异常
\sphinxbeforeendvarwidth
\end{varwidth}%
\\
\sphinxbottomrule
\end{tabulary}
\sphinxtableafterendhook\par
\sphinxattableend\end{savenotes}


\bigskip\hrule\bigskip



\subsection{1.7 软件配置建议}
\label{\detokenize{chapters/chapter1:id31}}

\subsubsection{1.7.1 推荐设置}
\label{\detokenize{chapters/chapter1:id32}}

\begin{savenotes}\sphinxattablestart
\sphinxthistablewithglobalstyle
\centering
\begin{tabulary}{\linewidth}[t]{TTT}
\sphinxtoprule
\begin{varwidth}[t]{\sphinxcolwidth{1}{3}}
\sphinxstyletheadfamily \sphinxAtStartPar
设置项
\sphinxbeforeendvarwidth
\end{varwidth}%
&\begin{varwidth}[t]{\sphinxcolwidth{1}{3}}
\sphinxstyletheadfamily \sphinxAtStartPar
推荐值
\sphinxbeforeendvarwidth
\end{varwidth}%
&\begin{varwidth}[t]{\sphinxcolwidth{1}{3}}
\sphinxstyletheadfamily \sphinxAtStartPar
说明
\sphinxbeforeendvarwidth
\end{varwidth}%
\\
\sphinxmidrule
\sphinxtableatstartofbodyhook\begin{varwidth}[t]{\sphinxcolwidth{1}{3}}
\sphinxAtStartPar
\sphinxstylestrong{最大线程数}
\sphinxbeforeendvarwidth
\end{varwidth}%
&\begin{varwidth}[t]{\sphinxcolwidth{1}{3}}
\sphinxAtStartPar
CPU核心数
\sphinxbeforeendvarwidth
\end{varwidth}%
&\begin{varwidth}[t]{\sphinxcolwidth{1}{3}}
\sphinxAtStartPar
充分利用多核性能
\sphinxbeforeendvarwidth
\end{varwidth}%
\\
\sphinxhline\begin{varwidth}[t]{\sphinxcolwidth{1}{3}}
\sphinxAtStartPar
\sphinxstylestrong{FFT窗长度}
\sphinxbeforeendvarwidth
\end{varwidth}%
&\begin{varwidth}[t]{\sphinxcolwidth{1}{3}}
\sphinxAtStartPar
4096点
\sphinxbeforeendvarwidth
\end{varwidth}%
&\begin{varwidth}[t]{\sphinxcolwidth{1}{3}}
\sphinxAtStartPar
平衡频率分辨率和时间分辨率
\sphinxbeforeendvarwidth
\end{varwidth}%
\\
\sphinxhline\begin{varwidth}[t]{\sphinxcolwidth{1}{3}}
\sphinxAtStartPar
\sphinxstylestrong{标定路径}
\sphinxbeforeendvarwidth
\end{varwidth}%
&\begin{varwidth}[t]{\sphinxcolwidth{1}{3}}
\sphinxAtStartPar
工程目录下的CLB/CLC文件夹
\sphinxbeforeendvarwidth
\end{varwidth}%
&\begin{varwidth}[t]{\sphinxcolwidth{1}{3}}
\sphinxAtStartPar
便于管理
\sphinxbeforeendvarwidth
\end{varwidth}%
\\
\sphinxhline\begin{varwidth}[t]{\sphinxcolwidth{1}{3}}
\sphinxAtStartPar
\sphinxstylestrong{自动保存间隔}
\sphinxbeforeendvarwidth
\end{varwidth}%
&\begin{varwidth}[t]{\sphinxcolwidth{1}{3}}
\sphinxAtStartPar
10分钟
\sphinxbeforeendvarwidth
\end{varwidth}%
&\begin{varwidth}[t]{\sphinxcolwidth{1}{3}}
\sphinxAtStartPar
防止数据丢失
\sphinxbeforeendvarwidth
\end{varwidth}%
\\
\sphinxbottomrule
\end{tabulary}
\sphinxtableafterendhook\par
\sphinxattableend\end{savenotes}


\subsubsection{1.7.2 性能优化}
\label{\detokenize{chapters/chapter1:id33}}
\sphinxAtStartPar
\sphinxstylestrong{处理大量数据时：}
\begin{enumerate}
\sphinxsetlistlabels{\arabic}{enumi}{enumii}{}{.}%
\item {} 
\sphinxAtStartPar
关闭不必要的图表窗口

\item {} 
\sphinxAtStartPar
减少实时预览

\item {} 
\sphinxAtStartPar
使用批量处理代替单点处理

\item {} 
\sphinxAtStartPar
确保足够的磁盘空间

\end{enumerate}


\bigskip\hrule\bigskip



\subsection{1.8 获取帮助}
\label{\detokenize{chapters/chapter1:id34}}\begin{itemize}
\item {} 
\sphinxAtStartPar
菜单 \sphinxcode{\sphinxupquote{帮助 → 关于}} 查看版本信息

\item {} 
\sphinxAtStartPar
遇到问题请查看 {\hyperref[\detokenize{appendices/appendixD::doc}]{\sphinxcrossref{\DUrole{doc}{\DUrole{std}{\DUrole{std-doc}{常见问题FAQ}}}}}}

\item {} 
\sphinxAtStartPar
联系技术支持获取帮助

\end{itemize}

\sphinxstepscope


\section{📚 第2章 大地电磁数据处理教程}
\label{\detokenize{chapters/chapter2:id1}}\label{\detokenize{chapters/chapter2::doc}}\begin{quote}

\sphinxAtStartPar
\sphinxstylestrong{本章目标}：带你从零开始，一步步掌握大地电磁(MT)数据处理的完整流程，理解每个步骤背后的原理，学会使用MTDP软件获得高质量的处理结果。
\end{quote}


\bigskip\hrule\bigskip



\subsection{2.1 教程概述与学习路径}
\label{\detokenize{chapters/chapter2:id2}}

\subsubsection{2.1.1 MT数据处理全景图}
\label{\detokenize{chapters/chapter2:mt}}
\sphinxAtStartPar
大地电磁数据处理是一个从原始时间序列到最终电阻率模型的完整流程。下图展示了整个处理流程：


\subsubsection{2.1.2 本章学习路径}
\label{\detokenize{chapters/chapter2:id3}}

\begin{savenotes}\sphinxattablestart
\sphinxthistablewithglobalstyle
\centering
\begin{tabulary}{\linewidth}[t]{TTTT}
\sphinxtoprule
\begin{varwidth}[t]{\sphinxcolwidth{1}{4}}
\sphinxstyletheadfamily \sphinxAtStartPar
学习阶段
\sphinxbeforeendvarwidth
\end{varwidth}%
&\begin{varwidth}[t]{\sphinxcolwidth{1}{4}}
\sphinxstyletheadfamily \sphinxAtStartPar
章节
\sphinxbeforeendvarwidth
\end{varwidth}%
&\begin{varwidth}[t]{\sphinxcolwidth{1}{4}}
\sphinxstyletheadfamily \sphinxAtStartPar
内容
\sphinxbeforeendvarwidth
\end{varwidth}%
&\begin{varwidth}[t]{\sphinxcolwidth{1}{4}}
\sphinxstyletheadfamily \sphinxAtStartPar
预计时间
\sphinxbeforeendvarwidth
\end{varwidth}%
\\
\sphinxmidrule
\sphinxtableatstartofbodyhook\begin{varwidth}[t]{\sphinxcolwidth{1}{4}}
\sphinxAtStartPar
🎯 \sphinxstylestrong{入门}
\sphinxbeforeendvarwidth
\end{varwidth}%
&\begin{varwidth}[t]{\sphinxcolwidth{1}{4}}
\sphinxAtStartPar
2.2
\sphinxbeforeendvarwidth
\end{varwidth}%
&\begin{varwidth}[t]{\sphinxcolwidth{1}{4}}
\sphinxAtStartPar
MT基本原理：为什么电磁场能探测地下？
\sphinxbeforeendvarwidth
\end{varwidth}%
&\begin{varwidth}[t]{\sphinxcolwidth{1}{4}}
\sphinxAtStartPar
20分钟
\sphinxbeforeendvarwidth
\end{varwidth}%
\\
\sphinxhline\begin{varwidth}[t]{\sphinxcolwidth{1}{4}}
\sphinxAtStartPar
📋 \sphinxstylestrong{准备}
\sphinxbeforeendvarwidth
\end{varwidth}%
&\begin{varwidth}[t]{\sphinxcolwidth{1}{4}}
\sphinxAtStartPar
2.3
\sphinxbeforeendvarwidth
\end{varwidth}%
&\begin{varwidth}[t]{\sphinxcolwidth{1}{4}}
\sphinxAtStartPar
数据采集：野外工作需要注意什么？
\sphinxbeforeendvarwidth
\end{varwidth}%
&\begin{varwidth}[t]{\sphinxcolwidth{1}{4}}
\sphinxAtStartPar
15分钟
\sphinxbeforeendvarwidth
\end{varwidth}%
\\
\sphinxhline\begin{varwidth}[t]{\sphinxcolwidth{1}{4}}
\sphinxAtStartPar
🔧 \sphinxstylestrong{基础}
\sphinxbeforeendvarwidth
\end{varwidth}%
&\begin{varwidth}[t]{\sphinxcolwidth{1}{4}}
\sphinxAtStartPar
2.4\sphinxhyphen{}2.5
\sphinxbeforeendvarwidth
\end{varwidth}%
&\begin{varwidth}[t]{\sphinxcolwidth{1}{4}}
\sphinxAtStartPar
频谱分析与阻抗估计：核心处理步骤
\sphinxbeforeendvarwidth
\end{varwidth}%
&\begin{varwidth}[t]{\sphinxcolwidth{1}{4}}
\sphinxAtStartPar
30分钟
\sphinxbeforeendvarwidth
\end{varwidth}%
\\
\sphinxhline\begin{varwidth}[t]{\sphinxcolwidth{1}{4}}
\sphinxAtStartPar
✅ \sphinxstylestrong{质控}
\sphinxbeforeendvarwidth
\end{varwidth}%
&\begin{varwidth}[t]{\sphinxcolwidth{1}{4}}
\sphinxAtStartPar
2.6
\sphinxbeforeendvarwidth
\end{varwidth}%
&\begin{varwidth}[t]{\sphinxcolwidth{1}{4}}
\sphinxAtStartPar
数据质量评估：如何判断结果可靠？
\sphinxbeforeendvarwidth
\end{varwidth}%
&\begin{varwidth}[t]{\sphinxcolwidth{1}{4}}
\sphinxAtStartPar
20分钟
\sphinxbeforeendvarwidth
\end{varwidth}%
\\
\sphinxhline\begin{varwidth}[t]{\sphinxcolwidth{1}{4}}
\sphinxAtStartPar
🚀 \sphinxstylestrong{进阶}
\sphinxbeforeendvarwidth
\end{varwidth}%
&\begin{varwidth}[t]{\sphinxcolwidth{1}{4}}
\sphinxAtStartPar
2.7\sphinxhyphen{}2.8
\sphinxbeforeendvarwidth
\end{varwidth}%
&\begin{varwidth}[t]{\sphinxcolwidth{1}{4}}
\sphinxAtStartPar
高级工具与RMSMR：提升处理效果
\sphinxbeforeendvarwidth
\end{varwidth}%
&\begin{varwidth}[t]{\sphinxcolwidth{1}{4}}
\sphinxAtStartPar
25分钟
\sphinxbeforeendvarwidth
\end{varwidth}%
\\
\sphinxbottomrule
\end{tabulary}
\sphinxtableafterendhook\par
\sphinxattableend\end{savenotes}


\subsubsection{2.1.3 学习建议}
\label{\detokenize{chapters/chapter2:id4}}
\begin{sphinxuseclass}{tip}\begin{enumerate}
\sphinxsetlistlabels{\arabic}{enumi}{enumii}{}{.}%
\item {} 
\sphinxAtStartPar
\sphinxstylestrong{先理解原理}：不要急于操作软件，先搞清楚"为什么"

\item {} 
\sphinxAtStartPar
\sphinxstylestrong{动手实践}：用软件自带示例数据跟着教程操作

\item {} 
\sphinxAtStartPar
\sphinxstylestrong{对比学习}：处理同一数据，尝试不同方法，对比结果差异

\item {} 
\sphinxAtStartPar
\sphinxstylestrong{记录问题}：遇到不理解的地方记下来，后续章节可能有答案

\end{enumerate}

\end{sphinxuseclass}

\subsubsection{2.1.4 术语中英对照表}
\label{\detokenize{chapters/chapter2:id5}}

\begin{savenotes}\sphinxattablestart
\sphinxthistablewithglobalstyle
\centering
\begin{tabulary}{\linewidth}[t]{TTT}
\sphinxtoprule
\begin{varwidth}[t]{\sphinxcolwidth{1}{3}}
\sphinxstyletheadfamily \sphinxAtStartPar
中文术语
\sphinxbeforeendvarwidth
\end{varwidth}%
&\begin{varwidth}[t]{\sphinxcolwidth{1}{3}}
\sphinxstyletheadfamily \sphinxAtStartPar
英文术语
\sphinxbeforeendvarwidth
\end{varwidth}%
&\begin{varwidth}[t]{\sphinxcolwidth{1}{3}}
\sphinxstyletheadfamily \sphinxAtStartPar
缩写/符号
\sphinxbeforeendvarwidth
\end{varwidth}%
\\
\sphinxmidrule
\sphinxtableatstartofbodyhook\begin{varwidth}[t]{\sphinxcolwidth{1}{3}}
\sphinxAtStartPar
大地电磁
\sphinxbeforeendvarwidth
\end{varwidth}%
&\begin{varwidth}[t]{\sphinxcolwidth{1}{3}}
\sphinxAtStartPar
Magnetotelluric
\sphinxbeforeendvarwidth
\end{varwidth}%
&\begin{varwidth}[t]{\sphinxcolwidth{1}{3}}
\sphinxAtStartPar
MT
\sphinxbeforeendvarwidth
\end{varwidth}%
\\
\sphinxhline\begin{varwidth}[t]{\sphinxcolwidth{1}{3}}
\sphinxAtStartPar
阻抗张量
\sphinxbeforeendvarwidth
\end{varwidth}%
&\begin{varwidth}[t]{\sphinxcolwidth{1}{3}}
\sphinxAtStartPar
Impedance Tensor
\sphinxbeforeendvarwidth
\end{varwidth}%
&\begin{varwidth}[t]{\sphinxcolwidth{1}{3}}
\sphinxAtStartPar
\sphinxstylestrong{Z}
\sphinxbeforeendvarwidth
\end{varwidth}%
\\
\sphinxhline\begin{varwidth}[t]{\sphinxcolwidth{1}{3}}
\sphinxAtStartPar
视电阻率
\sphinxbeforeendvarwidth
\end{varwidth}%
&\begin{varwidth}[t]{\sphinxcolwidth{1}{3}}
\sphinxAtStartPar
Apparent Resistivity
\sphinxbeforeendvarwidth
\end{varwidth}%
&\begin{varwidth}[t]{\sphinxcolwidth{1}{3}}
\sphinxAtStartPar
ρₐ
\sphinxbeforeendvarwidth
\end{varwidth}%
\\
\sphinxhline\begin{varwidth}[t]{\sphinxcolwidth{1}{3}}
\sphinxAtStartPar
相位
\sphinxbeforeendvarwidth
\end{varwidth}%
&\begin{varwidth}[t]{\sphinxcolwidth{1}{3}}
\sphinxAtStartPar
Phase
\sphinxbeforeendvarwidth
\end{varwidth}%
&\begin{varwidth}[t]{\sphinxcolwidth{1}{3}}
\sphinxAtStartPar
φ
\sphinxbeforeendvarwidth
\end{varwidth}%
\\
\sphinxhline\begin{varwidth}[t]{\sphinxcolwidth{1}{3}}
\sphinxAtStartPar
倾子
\sphinxbeforeendvarwidth
\end{varwidth}%
&\begin{varwidth}[t]{\sphinxcolwidth{1}{3}}
\sphinxAtStartPar
Tipper
\sphinxbeforeendvarwidth
\end{varwidth}%
&\begin{varwidth}[t]{\sphinxcolwidth{1}{3}}
\sphinxAtStartPar
\sphinxstylestrong{T}
\sphinxbeforeendvarwidth
\end{varwidth}%
\\
\sphinxhline\begin{varwidth}[t]{\sphinxcolwidth{1}{3}}
\sphinxAtStartPar
趋肤深度
\sphinxbeforeendvarwidth
\end{varwidth}%
&\begin{varwidth}[t]{\sphinxcolwidth{1}{3}}
\sphinxAtStartPar
Skin Depth
\sphinxbeforeendvarwidth
\end{varwidth}%
&\begin{varwidth}[t]{\sphinxcolwidth{1}{3}}
\sphinxAtStartPar
δ
\sphinxbeforeendvarwidth
\end{varwidth}%
\\
\sphinxhline\begin{varwidth}[t]{\sphinxcolwidth{1}{3}}
\sphinxAtStartPar
功率谱密度
\sphinxbeforeendvarwidth
\end{varwidth}%
&\begin{varwidth}[t]{\sphinxcolwidth{1}{3}}
\sphinxAtStartPar
Power Spectral Density
\sphinxbeforeendvarwidth
\end{varwidth}%
&\begin{varwidth}[t]{\sphinxcolwidth{1}{3}}
\sphinxAtStartPar
PSD
\sphinxbeforeendvarwidth
\end{varwidth}%
\\
\sphinxhline\begin{varwidth}[t]{\sphinxcolwidth{1}{3}}
\sphinxAtStartPar
常相干度
\sphinxbeforeendvarwidth
\end{varwidth}%
&\begin{varwidth}[t]{\sphinxcolwidth{1}{3}}
\sphinxAtStartPar
Ordinary Coherence
\sphinxbeforeendvarwidth
\end{varwidth}%
&\begin{varwidth}[t]{\sphinxcolwidth{1}{3}}
\sphinxAtStartPar
Coh²
\sphinxbeforeendvarwidth
\end{varwidth}%
\\
\sphinxhline\begin{varwidth}[t]{\sphinxcolwidth{1}{3}}
\sphinxAtStartPar
重相干度
\sphinxbeforeendvarwidth
\end{varwidth}%
&\begin{varwidth}[t]{\sphinxcolwidth{1}{3}}
\sphinxAtStartPar
Multiple Coherence
\sphinxbeforeendvarwidth
\end{varwidth}%
&\begin{varwidth}[t]{\sphinxcolwidth{1}{3}}
\sphinxAtStartPar
Cohₘ²
\sphinxbeforeendvarwidth
\end{varwidth}%
\\
\sphinxhline\begin{varwidth}[t]{\sphinxcolwidth{1}{3}}
\sphinxAtStartPar
偏相干度
\sphinxbeforeendvarwidth
\end{varwidth}%
&\begin{varwidth}[t]{\sphinxcolwidth{1}{3}}
\sphinxAtStartPar
Partial Coherence
\sphinxbeforeendvarwidth
\end{varwidth}%
&\begin{varwidth}[t]{\sphinxcolwidth{1}{3}}
\sphinxAtStartPar
Cohₚ²
\sphinxbeforeendvarwidth
\end{varwidth}%
\\
\sphinxhline\begin{varwidth}[t]{\sphinxcolwidth{1}{3}}
\sphinxAtStartPar
远参考道
\sphinxbeforeendvarwidth
\end{varwidth}%
&\begin{varwidth}[t]{\sphinxcolwidth{1}{3}}
\sphinxAtStartPar
Remote Reference
\sphinxbeforeendvarwidth
\end{varwidth}%
&\begin{varwidth}[t]{\sphinxcolwidth{1}{3}}
\sphinxAtStartPar
RR
\sphinxbeforeendvarwidth
\end{varwidth}%
\\
\sphinxhline\begin{varwidth}[t]{\sphinxcolwidth{1}{3}}
\sphinxAtStartPar
稳健估计
\sphinxbeforeendvarwidth
\end{varwidth}%
&\begin{varwidth}[t]{\sphinxcolwidth{1}{3}}
\sphinxAtStartPar
Robust Estimation
\sphinxbeforeendvarwidth
\end{varwidth}%
&\begin{varwidth}[t]{\sphinxcolwidth{1}{3}}
\sphinxAtStartPar
\sphinxhyphen{}
\sphinxbeforeendvarwidth
\end{varwidth}%
\\
\sphinxhline\begin{varwidth}[t]{\sphinxcolwidth{1}{3}}
\sphinxAtStartPar
相位张量
\sphinxbeforeendvarwidth
\end{varwidth}%
&\begin{varwidth}[t]{\sphinxcolwidth{1}{3}}
\sphinxAtStartPar
Phase Tensor
\sphinxbeforeendvarwidth
\end{varwidth}%
&\begin{varwidth}[t]{\sphinxcolwidth{1}{3}}
\sphinxAtStartPar
\sphinxstylestrong{Φ}
\sphinxbeforeendvarwidth
\end{varwidth}%
\\
\sphinxbottomrule
\end{tabulary}
\sphinxtableafterendhook\par
\sphinxattableend\end{savenotes}


\bigskip\hrule\bigskip



\subsection{2.2 第一步：理解MT基本原理}
\label{\detokenize{chapters/chapter2:id6}}

\subsubsection{2.2.1 天然电磁场从哪里来？}
\label{\detokenize{chapters/chapter2:id7}}
\sphinxAtStartPar
\sphinxstylestrong{想象一下}：你站在野外，周围看不见任何电磁设备，但实际上空气中充满了看不见的电磁波——这就是MT探测的"光源"。


\paragraph{两种天然场源}
\label{\detokenize{chapters/chapter2:id8}}

\begin{savenotes}\sphinxattablestart
\sphinxthistablewithglobalstyle
\centering
\begin{tabulary}{\linewidth}[t]{TTTT}
\sphinxtoprule
\begin{varwidth}[t]{\sphinxcolwidth{1}{4}}
\sphinxstyletheadfamily \sphinxAtStartPar
场源类型
\sphinxbeforeendvarwidth
\end{varwidth}%
&\begin{varwidth}[t]{\sphinxcolwidth{1}{4}}
\sphinxstyletheadfamily \sphinxAtStartPar
频率范围
\sphinxbeforeendvarwidth
\end{varwidth}%
&\begin{varwidth}[t]{\sphinxcolwidth{1}{4}}
\sphinxstyletheadfamily \sphinxAtStartPar
来源
\sphinxbeforeendvarwidth
\end{varwidth}%
&\begin{varwidth}[t]{\sphinxcolwidth{1}{4}}
\sphinxstyletheadfamily \sphinxAtStartPar
特点
\sphinxbeforeendvarwidth
\end{varwidth}%
\\
\sphinxmidrule
\sphinxtableatstartofbodyhook\begin{varwidth}[t]{\sphinxcolwidth{1}{4}}
\sphinxAtStartPar
\sphinxstylestrong{电离层辐射}
\sphinxbeforeendvarwidth
\end{varwidth}%
&\begin{varwidth}[t]{\sphinxcolwidth{1}{4}}
\sphinxAtStartPar
< 1 Hz
\sphinxbeforeendvarwidth
\end{varwidth}%
&\begin{varwidth}[t]{\sphinxcolwidth{1}{4}}
\sphinxAtStartPar
太阳风与地球磁场相互作用
\sphinxbeforeendvarwidth
\end{varwidth}%
&\begin{varwidth}[t]{\sphinxcolwidth{1}{4}}
\sphinxAtStartPar
能量稳定，适合深部探测
\sphinxbeforeendvarwidth
\end{varwidth}%
\\
\sphinxhline\begin{varwidth}[t]{\sphinxcolwidth{1}{4}}
\sphinxAtStartPar
\sphinxstylestrong{雷电天电波}
\sphinxbeforeendvarwidth
\end{varwidth}%
&\begin{varwidth}[t]{\sphinxcolwidth{1}{4}}
\sphinxAtStartPar
> 1 Hz
\sphinxbeforeendvarwidth
\end{varwidth}%
&\begin{varwidth}[t]{\sphinxcolwidth{1}{4}}
\sphinxAtStartPar
全球雷电活动
\sphinxbeforeendvarwidth
\end{varwidth}%
&\begin{varwidth}[t]{\sphinxcolwidth{1}{4}}
\sphinxAtStartPar
能量较强，适合浅部探测
\sphinxbeforeendvarwidth
\end{varwidth}%
\\
\sphinxbottomrule
\end{tabulary}
\sphinxtableafterendhook\par
\sphinxattableend\end{savenotes}
\begin{quote}

\sphinxAtStartPar
🌍 \sphinxstylestrong{有趣的事实}：地球上每秒钟有约100次雷电，它们激发的电磁波在地表和电离层之间来回反射，形成了我们使用的"免费光源"。
\end{quote}


\paragraph{昼夜变化的影响}
\label{\detokenize{chapters/chapter2:id9}}
\sphinxAtStartPar
由于地球自转，天然场有明显的昼夜变化：

\begin{sphinxVerbatim}[commandchars=\\\{\}]
白天：大气电离度高 → 高频信号衰减快 → 信号较弱
夜间：大气电离度低 → 高频信号传播好 → 信号较强
\end{sphinxVerbatim}

\begin{sphinxuseclass}{warning}
\sphinxAtStartPar
对于音频MT（AMT，1Hz\sphinxhyphen{}10kHz），\sphinxstylestrong{夜间观测}效果通常更好，特别是在"死频带"（1\sphinxhyphen{}5kHz）区域。

\end{sphinxuseclass}

\subsubsection{2.2.2 电磁波如何"看见"地下？}
\label{\detokenize{chapters/chapter2:id10}}

\paragraph{核心概念：趋肤深度}
\label{\detokenize{chapters/chapter2:id11}}
\sphinxAtStartPar
\sphinxstylestrong{想象海浪}：
\begin{itemize}
\item {} 
\sphinxAtStartPar
长波海浪可以传播到很远的深海

\item {} 
\sphinxAtStartPar
短波涟漪只在水面上传播，很快消失

\end{itemize}

\sphinxAtStartPar
电磁波在地下也是如此！\sphinxstylestrong{频率越低，穿透越深}。

\sphinxAtStartPar
趋肤深度（电磁场衰减到1/e的深度）计算公式：
\begin{equation*}
\begin{split}\delta = 503\sqrt{\frac{\rho}{f}} \text{2.1}\end{split}
\end{equation*}
\sphinxAtStartPar
其中：
\begin{itemize}
\item {} 
\sphinxAtStartPar
\(\delta\) = 趋肤深度（米）

\item {} 
\sphinxAtStartPar
\(\rho\) = 地下电阻率（Ω·m）

\item {} 
\sphinxAtStartPar
\(f\) = 频率（Hz）

\end{itemize}


\paragraph{实际探测深度速查表}
\label{\detokenize{chapters/chapter2:id12}}

\begin{savenotes}\sphinxattablestart
\sphinxthistablewithglobalstyle
\centering
\begin{tabulary}{\linewidth}[t]{TTTTTT}
\sphinxtoprule
\begin{varwidth}[t]{\sphinxcolwidth{1}{6}}
\sphinxstyletheadfamily \sphinxAtStartPar
地下电阻率
\sphinxbeforeendvarwidth
\end{varwidth}%
&\begin{varwidth}[t]{\sphinxcolwidth{1}{6}}
\sphinxstyletheadfamily \sphinxAtStartPar
0.01 Hz
\sphinxbeforeendvarwidth
\end{varwidth}%
&\begin{varwidth}[t]{\sphinxcolwidth{1}{6}}
\sphinxstyletheadfamily \sphinxAtStartPar
0.1 Hz
\sphinxbeforeendvarwidth
\end{varwidth}%
&\begin{varwidth}[t]{\sphinxcolwidth{1}{6}}
\sphinxstyletheadfamily \sphinxAtStartPar
1 Hz
\sphinxbeforeendvarwidth
\end{varwidth}%
&\begin{varwidth}[t]{\sphinxcolwidth{1}{6}}
\sphinxstyletheadfamily \sphinxAtStartPar
10 Hz
\sphinxbeforeendvarwidth
\end{varwidth}%
&\begin{varwidth}[t]{\sphinxcolwidth{1}{6}}
\sphinxstyletheadfamily \sphinxAtStartPar
100 Hz
\sphinxbeforeendvarwidth
\end{varwidth}%
\\
\sphinxmidrule
\sphinxtableatstartofbodyhook\begin{varwidth}[t]{\sphinxcolwidth{1}{6}}
\sphinxAtStartPar
\sphinxstylestrong{10 Ω·m} (沉积岩)
\sphinxbeforeendvarwidth
\end{varwidth}%
&\begin{varwidth}[t]{\sphinxcolwidth{1}{6}}
\sphinxAtStartPar
50 km
\sphinxbeforeendvarwidth
\end{varwidth}%
&\begin{varwidth}[t]{\sphinxcolwidth{1}{6}}
\sphinxAtStartPar
16 km
\sphinxbeforeendvarwidth
\end{varwidth}%
&\begin{varwidth}[t]{\sphinxcolwidth{1}{6}}
\sphinxAtStartPar
5 km
\sphinxbeforeendvarwidth
\end{varwidth}%
&\begin{varwidth}[t]{\sphinxcolwidth{1}{6}}
\sphinxAtStartPar
1.6 km
\sphinxbeforeendvarwidth
\end{varwidth}%
&\begin{varwidth}[t]{\sphinxcolwidth{1}{6}}
\sphinxAtStartPar
500 m
\sphinxbeforeendvarwidth
\end{varwidth}%
\\
\sphinxhline\begin{varwidth}[t]{\sphinxcolwidth{1}{6}}
\sphinxAtStartPar
\sphinxstylestrong{100 Ω·m} (中等)
\sphinxbeforeendvarwidth
\end{varwidth}%
&\begin{varwidth}[t]{\sphinxcolwidth{1}{6}}
\sphinxAtStartPar
160 km
\sphinxbeforeendvarwidth
\end{varwidth}%
&\begin{varwidth}[t]{\sphinxcolwidth{1}{6}}
\sphinxAtStartPar
50 km
\sphinxbeforeendvarwidth
\end{varwidth}%
&\begin{varwidth}[t]{\sphinxcolwidth{1}{6}}
\sphinxAtStartPar
16 km
\sphinxbeforeendvarwidth
\end{varwidth}%
&\begin{varwidth}[t]{\sphinxcolwidth{1}{6}}
\sphinxAtStartPar
5 km
\sphinxbeforeendvarwidth
\end{varwidth}%
&\begin{varwidth}[t]{\sphinxcolwidth{1}{6}}
\sphinxAtStartPar
1.6 km
\sphinxbeforeendvarwidth
\end{varwidth}%
\\
\sphinxhline\begin{varwidth}[t]{\sphinxcolwidth{1}{6}}
\sphinxAtStartPar
\sphinxstylestrong{1000 Ω·m} (基岩)
\sphinxbeforeendvarwidth
\end{varwidth}%
&\begin{varwidth}[t]{\sphinxcolwidth{1}{6}}
\sphinxAtStartPar
500 km
\sphinxbeforeendvarwidth
\end{varwidth}%
&\begin{varwidth}[t]{\sphinxcolwidth{1}{6}}
\sphinxAtStartPar
160 km
\sphinxbeforeendvarwidth
\end{varwidth}%
&\begin{varwidth}[t]{\sphinxcolwidth{1}{6}}
\sphinxAtStartPar
50 km
\sphinxbeforeendvarwidth
\end{varwidth}%
&\begin{varwidth}[t]{\sphinxcolwidth{1}{6}}
\sphinxAtStartPar
16 km
\sphinxbeforeendvarwidth
\end{varwidth}%
&\begin{varwidth}[t]{\sphinxcolwidth{1}{6}}
\sphinxAtStartPar
5 km
\sphinxbeforeendvarwidth
\end{varwidth}%
\\
\sphinxbottomrule
\end{tabulary}
\sphinxtableafterendhook\par
\sphinxattableend\end{savenotes}

\begin{sphinxuseclass}{tip}\begin{itemize}
\item {} 
\sphinxAtStartPar
\sphinxstylestrong{低频 = 深部信息}：想探测更深，需要更低频率的信号

\item {} 
\sphinxAtStartPar
\sphinxstylestrong{高频 = 浅部信息}：浅部结构由高频信号反映

\item {} 
\sphinxAtStartPar
\sphinxstylestrong{电阻率影响}：同样频率，在高阻地区探测更深

\end{itemize}

\end{sphinxuseclass}

\subsubsection{2.2.3 我们测量的是什么？}
\label{\detokenize{chapters/chapter2:id13}}

\paragraph{阻抗：地下的"电阻"}
\label{\detokenize{chapters/chapter2:id14}}
\sphinxAtStartPar
\sphinxstylestrong{简单类比}：用万用表测电阻，我们测量电压和电流的比值。MT方法测量电场和磁场的比值，得到地下的"等效电阻"——这就是\sphinxstylestrong{阻抗}。
\begin{equation*}
\begin{split}\mathbf{E} = \mathbf{Z} \cdot \mathbf{H} \text{2.2}\end{split}
\end{equation*}
\sphinxAtStartPar
因为地下结构可能沿不同方向有不同电阻率（各向异性），所以阻抗是一个\sphinxstylestrong{2×2张量}：
\begin{equation*}
\begin{split}\begin{pmatrix} E_x \\ E_y \end{pmatrix} = \begin{pmatrix} Z_{xx} & Z_{xy} \\ Z_{yx} & Z_{yy} \end{pmatrix} \begin{pmatrix} H_x \\ H_y \end{pmatrix}\qquad\text{(2.3)}\end{split}
\end{equation*}

\paragraph{视电阻率与相位}
\label{\detokenize{chapters/chapter2:id15}}
\sphinxAtStartPar
从阻抗可以计算\sphinxstylestrong{视电阻率}和\sphinxstylestrong{相位}：
\begin{equation*}
\begin{split}\rho_a = \frac{|Z|^2}{\omega \mu_0} \text{2.4}\end{split}
\end{equation*}\begin{equation*}
\begin{split}\varphi = \arg(Z) \text{2.5}\end{split}
\end{equation*}
\sphinxAtStartPar
\sphinxstylestrong{相位的物理意义}：电场和磁场之间的时间延迟。
\begin{itemize}
\item {} 
\sphinxAtStartPar
均匀半空间：相位 = 45°

\item {} 
\sphinxAtStartPar
电阻率随深度增加：相位 > 45°

\item {} 
\sphinxAtStartPar
电阻率随深度减小：相位 < 45°

\end{itemize}


\paragraph{倾子：寻找侧向异常}
\label{\detokenize{chapters/chapter2:id16}}
\sphinxAtStartPar
如果地下是理想的水平层状结构，磁场只有水平分量。但如果旁边有电性异常体（如矿体、断裂带），磁场会"弯曲"，产生垂直分量\(H_z\)。

\sphinxAtStartPar
\sphinxstylestrong{倾子}描述这种关系：
\begin{equation*}
\begin{split}H_z = T_{zx}H_x + T_{zy}H_y \text{2.6}\end{split}
\end{equation*}\begin{quote}

\sphinxAtStartPar
🎯 \sphinxstylestrong{实际应用}：感应矢量（由倾子计算）会指向低阻体，是寻找矿体或断裂带的重要工具。
\end{quote}


\bigskip\hrule\bigskip


\sphinxAtStartPar
⬅️ {\hyperref[\detokenize{chapters/chapter2:id2}]{\sphinxcrossref{上一节：教程概述}}} | ➡️ {\hyperref[\detokenize{chapters/chapter2:id17}]{\sphinxcrossref{下一节：数据采集准备}}}


\bigskip\hrule\bigskip



\subsection{2.3 第二步：数据采集准备}
\label{\detokenize{chapters/chapter2:id17}}

\subsubsection{2.3.1 野外布置要点}
\label{\detokenize{chapters/chapter2:id18}}

\paragraph{测站组成}
\label{\detokenize{chapters/chapter2:id19}}
\sphinxAtStartPar
一个完整的MT测站需要记录5个分量：


\begin{savenotes}\sphinxattablestart
\sphinxthistablewithglobalstyle
\centering
\begin{tabulary}{\linewidth}[t]{TTT}
\sphinxtoprule
\begin{varwidth}[t]{\sphinxcolwidth{1}{3}}
\sphinxstyletheadfamily \sphinxAtStartPar
分量
\sphinxbeforeendvarwidth
\end{varwidth}%
&\begin{varwidth}[t]{\sphinxcolwidth{1}{3}}
\sphinxstyletheadfamily \sphinxAtStartPar
传感器
\sphinxbeforeendvarwidth
\end{varwidth}%
&\begin{varwidth}[t]{\sphinxcolwidth{1}{3}}
\sphinxstyletheadfamily \sphinxAtStartPar
测量方向
\sphinxbeforeendvarwidth
\end{varwidth}%
\\
\sphinxmidrule
\sphinxtableatstartofbodyhook\begin{varwidth}[t]{\sphinxcolwidth{1}{3}}
\sphinxAtStartPar
\(E_x\)
\sphinxbeforeendvarwidth
\end{varwidth}%
&\begin{varwidth}[t]{\sphinxcolwidth{1}{3}}
\sphinxAtStartPar
电极
\sphinxbeforeendvarwidth
\end{varwidth}%
&\begin{varwidth}[t]{\sphinxcolwidth{1}{3}}
\sphinxAtStartPar
南北方向
\sphinxbeforeendvarwidth
\end{varwidth}%
\\
\sphinxhline\begin{varwidth}[t]{\sphinxcolwidth{1}{3}}
\sphinxAtStartPar
\(E_y\)
\sphinxbeforeendvarwidth
\end{varwidth}%
&\begin{varwidth}[t]{\sphinxcolwidth{1}{3}}
\sphinxAtStartPar
电极
\sphinxbeforeendvarwidth
\end{varwidth}%
&\begin{varwidth}[t]{\sphinxcolwidth{1}{3}}
\sphinxAtStartPar
东西方向
\sphinxbeforeendvarwidth
\end{varwidth}%
\\
\sphinxhline\begin{varwidth}[t]{\sphinxcolwidth{1}{3}}
\sphinxAtStartPar
\(H_x\)
\sphinxbeforeendvarwidth
\end{varwidth}%
&\begin{varwidth}[t]{\sphinxcolwidth{1}{3}}
\sphinxAtStartPar
磁探头
\sphinxbeforeendvarwidth
\end{varwidth}%
&\begin{varwidth}[t]{\sphinxcolwidth{1}{3}}
\sphinxAtStartPar
南北方向
\sphinxbeforeendvarwidth
\end{varwidth}%
\\
\sphinxhline\begin{varwidth}[t]{\sphinxcolwidth{1}{3}}
\sphinxAtStartPar
\(H_y\)
\sphinxbeforeendvarwidth
\end{varwidth}%
&\begin{varwidth}[t]{\sphinxcolwidth{1}{3}}
\sphinxAtStartPar
磁探头
\sphinxbeforeendvarwidth
\end{varwidth}%
&\begin{varwidth}[t]{\sphinxcolwidth{1}{3}}
\sphinxAtStartPar
东西方向
\sphinxbeforeendvarwidth
\end{varwidth}%
\\
\sphinxhline\begin{varwidth}[t]{\sphinxcolwidth{1}{3}}
\sphinxAtStartPar
\(H_z\)
\sphinxbeforeendvarwidth
\end{varwidth}%
&\begin{varwidth}[t]{\sphinxcolwidth{1}{3}}
\sphinxAtStartPar
磁探头
\sphinxbeforeendvarwidth
\end{varwidth}%
&\begin{varwidth}[t]{\sphinxcolwidth{1}{3}}
\sphinxAtStartPar
垂直方向
\sphinxbeforeendvarwidth
\end{varwidth}%
\\
\sphinxbottomrule
\end{tabulary}
\sphinxtableafterendhook\par
\sphinxattableend\end{savenotes}


\paragraph{布置建议}
\label{\detokenize{chapters/chapter2:id20}}

\begin{savenotes}\sphinxattablestart
\sphinxthistablewithglobalstyle
\centering
\begin{tabulary}{\linewidth}[t]{TT}
\sphinxtoprule
\begin{varwidth}[t]{\sphinxcolwidth{1}{2}}
\sphinxstyletheadfamily \sphinxAtStartPar
注意事项
\sphinxbeforeendvarwidth
\end{varwidth}%
&\begin{varwidth}[t]{\sphinxcolwidth{1}{2}}
\sphinxstyletheadfamily \sphinxAtStartPar
具体要求
\sphinxbeforeendvarwidth
\end{varwidth}%
\\
\sphinxmidrule
\sphinxtableatstartofbodyhook\begin{varwidth}[t]{\sphinxcolwidth{1}{2}}
\sphinxAtStartPar
\sphinxstylestrong{避开干扰}
\sphinxbeforeendvarwidth
\end{varwidth}%
&\begin{varwidth}[t]{\sphinxcolwidth{1}{2}}
\sphinxAtStartPar
距离高压线>500m，工厂>1km，公路>200m
\sphinxbeforeendvarwidth
\end{varwidth}%
\\
\sphinxhline\begin{varwidth}[t]{\sphinxcolwidth{1}{2}}
\sphinxAtStartPar
\sphinxstylestrong{电极接触}
\sphinxbeforeendvarwidth
\end{varwidth}%
&\begin{varwidth}[t]{\sphinxcolwidth{1}{2}}
\sphinxAtStartPar
埋深>20cm，浇盐水降低接地电阻
\sphinxbeforeendvarwidth
\end{varwidth}%
\\
\sphinxhline\begin{varwidth}[t]{\sphinxcolwidth{1}{2}}
\sphinxAtStartPar
\sphinxstylestrong{磁探头}
\sphinxbeforeendvarwidth
\end{varwidth}%
&\begin{varwidth}[t]{\sphinxcolwidth{1}{2}}
\sphinxAtStartPar
水平放置，远离金属物体
\sphinxbeforeendvarwidth
\end{varwidth}%
\\
\sphinxhline\begin{varwidth}[t]{\sphinxcolwidth{1}{2}}
\sphinxAtStartPar
\sphinxstylestrong{电缆}
\sphinxbeforeendvarwidth
\end{varwidth}%
&\begin{varwidth}[t]{\sphinxcolwidth{1}{2}}
\sphinxAtStartPar
埋入土中或防水布覆盖，避免晃动
\sphinxbeforeendvarwidth
\end{varwidth}%
\\
\sphinxbottomrule
\end{tabulary}
\sphinxtableafterendhook\par
\sphinxattableend\end{savenotes}

\begin{sphinxuseclass}{danger}\begin{enumerate}
\sphinxsetlistlabels{\arabic}{enumi}{enumii}{}{.}%
\item {} 
\sphinxAtStartPar
\sphinxstylestrong{电极极化}：长时间观测后电极接触变差，需定期检查接地电阻（建议 < 1kΩ）

\item {} 
\sphinxAtStartPar
\sphinxstylestrong{磁探头倾斜}：水平磁探头倾斜会导致信号失真，使用水平仪校准

\item {} 
\sphinxAtStartPar
\sphinxstylestrong{电缆松动}：风吹动电缆会产生电磁干扰，务必固定好

\end{enumerate}

\end{sphinxuseclass}

\subsubsection{2.3.2 远参考站设置}
\label{\detokenize{chapters/chapter2:id21}}

\paragraph{为什么需要远参考站？}
\label{\detokenize{chapters/chapter2:id22}}
\sphinxAtStartPar
\sphinxstylestrong{问题}：本地测站可能受到附近噪声干扰（如工厂、电线），导致阻抗估计偏差。

\sphinxAtStartPar
\sphinxstylestrong{解决方案}：在远处设置参考站，利用"噪声不相关、信号相关"的特性消除偏差。


\paragraph{远参考站设置要求}
\label{\detokenize{chapters/chapter2:id23}}

\begin{savenotes}\sphinxattablestart
\sphinxthistablewithglobalstyle
\centering
\begin{tabulary}{\linewidth}[t]{TTT}
\sphinxtoprule
\begin{varwidth}[t]{\sphinxcolwidth{1}{3}}
\sphinxstyletheadfamily \sphinxAtStartPar
参数
\sphinxbeforeendvarwidth
\end{varwidth}%
&\begin{varwidth}[t]{\sphinxcolwidth{1}{3}}
\sphinxstyletheadfamily \sphinxAtStartPar
建议值
\sphinxbeforeendvarwidth
\end{varwidth}%
&\begin{varwidth}[t]{\sphinxcolwidth{1}{3}}
\sphinxstyletheadfamily \sphinxAtStartPar
原因
\sphinxbeforeendvarwidth
\end{varwidth}%
\\
\sphinxmidrule
\sphinxtableatstartofbodyhook\begin{varwidth}[t]{\sphinxcolwidth{1}{3}}
\sphinxAtStartPar
\sphinxstylestrong{距离}
\sphinxbeforeendvarwidth
\end{varwidth}%
&\begin{varwidth}[t]{\sphinxcolwidth{1}{3}}
\sphinxAtStartPar
几公里到几十公里
\sphinxbeforeendvarwidth
\end{varwidth}%
&\begin{varwidth}[t]{\sphinxcolwidth{1}{3}}
\sphinxAtStartPar
保证噪声不相关，信号仍相关
\sphinxbeforeendvarwidth
\end{varwidth}%
\\
\sphinxhline\begin{varwidth}[t]{\sphinxcolwidth{1}{3}}
\sphinxAtStartPar
\sphinxstylestrong{同步}
\sphinxbeforeendvarwidth
\end{varwidth}%
&\begin{varwidth}[t]{\sphinxcolwidth{1}{3}}
\sphinxAtStartPar
GPS时间同步
\sphinxbeforeendvarwidth
\end{varwidth}%
&\begin{varwidth}[t]{\sphinxcolwidth{1}{3}}
\sphinxAtStartPar
必须精确同步
\sphinxbeforeendvarwidth
\end{varwidth}%
\\
\sphinxhline\begin{varwidth}[t]{\sphinxcolwidth{1}{3}}
\sphinxAtStartPar
\sphinxstylestrong{环境}
\sphinxbeforeendvarwidth
\end{varwidth}%
&\begin{varwidth}[t]{\sphinxcolwidth{1}{3}}
\sphinxAtStartPar
低噪声区域
\sphinxbeforeendvarwidth
\end{varwidth}%
&\begin{varwidth}[t]{\sphinxcolwidth{1}{3}}
\sphinxAtStartPar
参考站质量影响处理效果
\sphinxbeforeendvarwidth
\end{varwidth}%
\\
\sphinxbottomrule
\end{tabulary}
\sphinxtableafterendhook\par
\sphinxattableend\end{savenotes}

\begin{sphinxuseclass}{danger}
\sphinxAtStartPar
远参考站与本地测站\sphinxstylestrong{不能}受同一噪声源影响。例如，两个站都靠近同一工厂就失去了远参考的意义。

\end{sphinxuseclass}

\subsubsection{2.3.3 观测时长建议}
\label{\detokenize{chapters/chapter2:id24}}

\paragraph{根据探测深度确定时长}
\label{\detokenize{chapters/chapter2:id25}}

\begin{savenotes}\sphinxattablestart
\sphinxthistablewithglobalstyle
\centering
\begin{tabulary}{\linewidth}[t]{TTT}
\sphinxtoprule
\begin{varwidth}[t]{\sphinxcolwidth{1}{3}}
\sphinxstyletheadfamily \sphinxAtStartPar
目标深度
\sphinxbeforeendvarwidth
\end{varwidth}%
&\begin{varwidth}[t]{\sphinxcolwidth{1}{3}}
\sphinxstyletheadfamily \sphinxAtStartPar
建议观测时长
\sphinxbeforeendvarwidth
\end{varwidth}%
&\begin{varwidth}[t]{\sphinxcolwidth{1}{3}}
\sphinxstyletheadfamily \sphinxAtStartPar
原因
\sphinxbeforeendvarwidth
\end{varwidth}%
\\
\sphinxmidrule
\sphinxtableatstartofbodyhook\begin{varwidth}[t]{\sphinxcolwidth{1}{3}}
\sphinxAtStartPar
< 1 km（浅部）
\sphinxbeforeendvarwidth
\end{varwidth}%
&\begin{varwidth}[t]{\sphinxcolwidth{1}{3}}
\sphinxAtStartPar
2\sphinxhyphen{}4 小时
\sphinxbeforeendvarwidth
\end{varwidth}%
&\begin{varwidth}[t]{\sphinxcolwidth{1}{3}}
\sphinxAtStartPar
高频信号充足
\sphinxbeforeendvarwidth
\end{varwidth}%
\\
\sphinxhline\begin{varwidth}[t]{\sphinxcolwidth{1}{3}}
\sphinxAtStartPar
1\sphinxhyphen{}10 km（中深部）
\sphinxbeforeendvarwidth
\end{varwidth}%
&\begin{varwidth}[t]{\sphinxcolwidth{1}{3}}
\sphinxAtStartPar
8\sphinxhyphen{}12 小时
\sphinxbeforeendvarwidth
\end{varwidth}%
&\begin{varwidth}[t]{\sphinxcolwidth{1}{3}}
\sphinxAtStartPar
需要足够的低频叠加
\sphinxbeforeendvarwidth
\end{varwidth}%
\\
\sphinxhline\begin{varwidth}[t]{\sphinxcolwidth{1}{3}}
\sphinxAtStartPar
> 10 km（深部）
\sphinxbeforeendvarwidth
\end{varwidth}%
&\begin{varwidth}[t]{\sphinxcolwidth{1}{3}}
\sphinxAtStartPar
24+ 小时
\sphinxbeforeendvarwidth
\end{varwidth}%
&\begin{varwidth}[t]{\sphinxcolwidth{1}{3}}
\sphinxAtStartPar
低频信号需要长时间积累
\sphinxbeforeendvarwidth
\end{varwidth}%
\\
\sphinxbottomrule
\end{tabulary}
\sphinxtableafterendhook\par
\sphinxattableend\end{savenotes}


\paragraph{提高数据质量的技巧}
\label{\detokenize{chapters/chapter2:id26}}\begin{enumerate}
\sphinxsetlistlabels{\arabic}{enumi}{enumii}{}{.}%
\item {} 
\sphinxAtStartPar
\sphinxstylestrong{跨越昼夜}：包含白天和夜间的数据，覆盖不同噪声环境

\item {} 
\sphinxAtStartPar
\sphinxstylestrong{多次叠加}：长时间观测可以增加叠加次数，提高信噪比

\item {} 
\sphinxAtStartPar
\sphinxstylestrong{避开干扰时段}：如附近的工厂工作时间

\end{enumerate}


\subsubsection{2.3.4 MTDP数据导入操作}
\label{\detokenize{chapters/chapter2:mtdp}}
\begin{sphinxuseclass}{details}
\sphinxAtStartPar
\sphinxstylestrong{导入时间序列数据：}
\begin{enumerate}
\sphinxsetlistlabels{\arabic}{enumi}{enumii}{}{.}%
\item {} 
\sphinxAtStartPar
菜单：\sphinxstylestrong{文件 → 导入 → 时间序列}（或快捷键 \sphinxcode{\sphinxupquote{Ctrl+I}}）

\item {} 
\sphinxAtStartPar
选择数据格式（Phoenix/Metronix/LEMI等）

\item {} 
\sphinxAtStartPar
选择数据文件夹

\item {} 
\sphinxAtStartPar
确认时间范围和采样率

\end{enumerate}

\sphinxAtStartPar
\sphinxstylestrong{导入远参考数据：}
\begin{enumerate}
\sphinxsetlistlabels{\arabic}{enumi}{enumii}{}{.}%
\item {} 
\sphinxAtStartPar
菜单：\sphinxstylestrong{文件 → 导入 → 远参考数据}

\item {} 
\sphinxAtStartPar
选择参考站数据文件

\item {} 
\sphinxAtStartPar
系统自动检测时间同步状态

\end{enumerate}

\end{sphinxuseclass}

\subsubsection{2.3.5 测点设置表单详解}
\label{\detokenize{chapters/chapter2:id27}}
\sphinxAtStartPar
在MTDP中，每个测点的详细参数通过测点设置表单进行管理。双击测点即可打开该表单。


\paragraph{2.3.5.1 基本信息}
\label{\detokenize{chapters/chapter2:id28}}

\begin{savenotes}\sphinxattablestart
\sphinxthistablewithglobalstyle
\centering
\begin{tabulary}{\linewidth}[t]{TTTT}
\sphinxtoprule
\begin{varwidth}[t]{\sphinxcolwidth{1}{4}}
\sphinxstyletheadfamily \sphinxAtStartPar
字段
\sphinxbeforeendvarwidth
\end{varwidth}%
&\begin{varwidth}[t]{\sphinxcolwidth{1}{4}}
\sphinxstyletheadfamily \sphinxAtStartPar
控件名称
\sphinxbeforeendvarwidth
\end{varwidth}%
&\begin{varwidth}[t]{\sphinxcolwidth{1}{4}}
\sphinxstyletheadfamily \sphinxAtStartPar
说明
\sphinxbeforeendvarwidth
\end{varwidth}%
&\begin{varwidth}[t]{\sphinxcolwidth{1}{4}}
\sphinxstyletheadfamily \sphinxAtStartPar
数据类型
\sphinxbeforeendvarwidth
\end{varwidth}%
\\
\sphinxmidrule
\sphinxtableatstartofbodyhook\begin{varwidth}[t]{\sphinxcolwidth{1}{4}}
\sphinxAtStartPar
测点名称
\sphinxbeforeendvarwidth
\end{varwidth}%
&\begin{varwidth}[t]{\sphinxcolwidth{1}{4}}
\sphinxAtStartPar
EditName
\sphinxbeforeendvarwidth
\end{varwidth}%
&\begin{varwidth}[t]{\sphinxcolwidth{1}{4}}
\sphinxAtStartPar
测点标识名称
\sphinxbeforeendvarwidth
\end{varwidth}%
&\begin{varwidth}[t]{\sphinxcolwidth{1}{4}}
\sphinxAtStartPar
字符串
\sphinxbeforeendvarwidth
\end{varwidth}%
\\
\sphinxhline\begin{varwidth}[t]{\sphinxcolwidth{1}{4}}
\sphinxAtStartPar
所有者
\sphinxbeforeendvarwidth
\end{varwidth}%
&\begin{varwidth}[t]{\sphinxcolwidth{1}{4}}
\sphinxAtStartPar
EditOwner
\sphinxbeforeendvarwidth
\end{varwidth}%
&\begin{varwidth}[t]{\sphinxcolwidth{1}{4}}
\sphinxAtStartPar
测点所有者/负责人
\sphinxbeforeendvarwidth
\end{varwidth}%
&\begin{varwidth}[t]{\sphinxcolwidth{1}{4}}
\sphinxAtStartPar
字符串
\sphinxbeforeendvarwidth
\end{varwidth}%
\\
\sphinxhline\begin{varwidth}[t]{\sphinxcolwidth{1}{4}}
\sphinxAtStartPar
描述
\sphinxbeforeendvarwidth
\end{varwidth}%
&\begin{varwidth}[t]{\sphinxcolwidth{1}{4}}
\sphinxAtStartPar
MemoDescription
\sphinxbeforeendvarwidth
\end{varwidth}%
&\begin{varwidth}[t]{\sphinxcolwidth{1}{4}}
\sphinxAtStartPar
测点详细描述信息
\sphinxbeforeendvarwidth
\end{varwidth}%
&\begin{varwidth}[t]{\sphinxcolwidth{1}{4}}
\sphinxAtStartPar
多行文本
\sphinxbeforeendvarwidth
\end{varwidth}%
\\
\sphinxhline\begin{varwidth}[t]{\sphinxcolwidth{1}{4}}
\sphinxAtStartPar
测量员
\sphinxbeforeendvarwidth
\end{varwidth}%
&\begin{varwidth}[t]{\sphinxcolwidth{1}{4}}
\sphinxAtStartPar
EditSurveyor
\sphinxbeforeendvarwidth
\end{varwidth}%
&\begin{varwidth}[t]{\sphinxcolwidth{1}{4}}
\sphinxAtStartPar
实施测量的技术人员
\sphinxbeforeendvarwidth
\end{varwidth}%
&\begin{varwidth}[t]{\sphinxcolwidth{1}{4}}
\sphinxAtStartPar
字符串
\sphinxbeforeendvarwidth
\end{varwidth}%
\\
\sphinxhline\begin{varwidth}[t]{\sphinxcolwidth{1}{4}}
\sphinxAtStartPar
数据采集者
\sphinxbeforeendvarwidth
\end{varwidth}%
&\begin{varwidth}[t]{\sphinxcolwidth{1}{4}}
\sphinxAtStartPar
EditDataCollector
\sphinxbeforeendvarwidth
\end{varwidth}%
&\begin{varwidth}[t]{\sphinxcolwidth{1}{4}}
\sphinxAtStartPar
采集设备操作人员
\sphinxbeforeendvarwidth
\end{varwidth}%
&\begin{varwidth}[t]{\sphinxcolwidth{1}{4}}
\sphinxAtStartPar
字符串
\sphinxbeforeendvarwidth
\end{varwidth}%
\\
\sphinxhline\begin{varwidth}[t]{\sphinxcolwidth{1}{4}}
\sphinxAtStartPar
数据处理者
\sphinxbeforeendvarwidth
\end{varwidth}%
&\begin{varwidth}[t]{\sphinxcolwidth{1}{4}}
\sphinxAtStartPar
EditDataProcessor
\sphinxbeforeendvarwidth
\end{varwidth}%
&\begin{varwidth}[t]{\sphinxcolwidth{1}{4}}
\sphinxAtStartPar
数据处理操作人员
\sphinxbeforeendvarwidth
\end{varwidth}%
&\begin{varwidth}[t]{\sphinxcolwidth{1}{4}}
\sphinxAtStartPar
字符串
\sphinxbeforeendvarwidth
\end{varwidth}%
\\
\sphinxhline\begin{varwidth}[t]{\sphinxcolwidth{1}{4}}
\sphinxAtStartPar
创建时间
\sphinxbeforeendvarwidth
\end{varwidth}%
&\begin{varwidth}[t]{\sphinxcolwidth{1}{4}}
\sphinxAtStartPar
EditCreationTime
\sphinxbeforeendvarwidth
\end{varwidth}%
&\begin{varwidth}[t]{\sphinxcolwidth{1}{4}}
\sphinxAtStartPar
测点创建时间戳
\sphinxbeforeendvarwidth
\end{varwidth}%
&\begin{varwidth}[t]{\sphinxcolwidth{1}{4}}
\sphinxAtStartPar
日期时间
\sphinxbeforeendvarwidth
\end{varwidth}%
\\
\sphinxhline\begin{varwidth}[t]{\sphinxcolwidth{1}{4}}
\sphinxAtStartPar
数据质量
\sphinxbeforeendvarwidth
\end{varwidth}%
&\begin{varwidth}[t]{\sphinxcolwidth{1}{4}}
\sphinxAtStartPar
ComboBoxDataQuality
\sphinxbeforeendvarwidth
\end{varwidth}%
&\begin{varwidth}[t]{\sphinxcolwidth{1}{4}}
\sphinxAtStartPar
数据质量评级(0\sphinxhyphen{}4级)
\sphinxbeforeendvarwidth
\end{varwidth}%
&\begin{varwidth}[t]{\sphinxcolwidth{1}{4}}
\sphinxAtStartPar
下拉选择
\sphinxbeforeendvarwidth
\end{varwidth}%
\\
\sphinxbottomrule
\end{tabulary}
\sphinxtableafterendhook\par
\sphinxattableend\end{savenotes}

\sphinxAtStartPar
\sphinxstylestrong{数据质量等级定义：}


\begin{savenotes}\sphinxattablestart
\sphinxthistablewithglobalstyle
\centering
\begin{tabulary}{\linewidth}[t]{TTT}
\sphinxtoprule
\begin{varwidth}[t]{\sphinxcolwidth{1}{3}}
\sphinxstyletheadfamily \sphinxAtStartPar
等级
\sphinxbeforeendvarwidth
\end{varwidth}%
&\begin{varwidth}[t]{\sphinxcolwidth{1}{3}}
\sphinxstyletheadfamily \sphinxAtStartPar
名称
\sphinxbeforeendvarwidth
\end{varwidth}%
&\begin{varwidth}[t]{\sphinxcolwidth{1}{3}}
\sphinxstyletheadfamily \sphinxAtStartPar
说明
\sphinxbeforeendvarwidth
\end{varwidth}%
\\
\sphinxmidrule
\sphinxtableatstartofbodyhook\begin{varwidth}[t]{\sphinxcolwidth{1}{3}}
\sphinxAtStartPar
0
\sphinxbeforeendvarwidth
\end{varwidth}%
&\begin{varwidth}[t]{\sphinxcolwidth{1}{3}}
\sphinxAtStartPar
未评估
\sphinxbeforeendvarwidth
\end{varwidth}%
&\begin{varwidth}[t]{\sphinxcolwidth{1}{3}}
\sphinxAtStartPar
尚未进行质量评估
\sphinxbeforeendvarwidth
\end{varwidth}%
\\
\sphinxhline\begin{varwidth}[t]{\sphinxcolwidth{1}{3}}
\sphinxAtStartPar
1
\sphinxbeforeendvarwidth
\end{varwidth}%
&\begin{varwidth}[t]{\sphinxcolwidth{1}{3}}
\sphinxAtStartPar
优秀
\sphinxbeforeendvarwidth
\end{varwidth}%
&\begin{varwidth}[t]{\sphinxcolwidth{1}{3}}
\sphinxAtStartPar
数据质量极佳可直接使用
\sphinxbeforeendvarwidth
\end{varwidth}%
\\
\sphinxhline\begin{varwidth}[t]{\sphinxcolwidth{1}{3}}
\sphinxAtStartPar
2
\sphinxbeforeendvarwidth
\end{varwidth}%
&\begin{varwidth}[t]{\sphinxcolwidth{1}{3}}
\sphinxAtStartPar
良好
\sphinxbeforeendvarwidth
\end{varwidth}%
&\begin{varwidth}[t]{\sphinxcolwidth{1}{3}}
\sphinxAtStartPar
数据质量正常可用
\sphinxbeforeendvarwidth
\end{varwidth}%
\\
\sphinxhline\begin{varwidth}[t]{\sphinxcolwidth{1}{3}}
\sphinxAtStartPar
3
\sphinxbeforeendvarwidth
\end{varwidth}%
&\begin{varwidth}[t]{\sphinxcolwidth{1}{3}}
\sphinxAtStartPar
一般
\sphinxbeforeendvarwidth
\end{varwidth}%
&\begin{varwidth}[t]{\sphinxcolwidth{1}{3}}
\sphinxAtStartPar
数据存在轻微问题需注意
\sphinxbeforeendvarwidth
\end{varwidth}%
\\
\sphinxhline\begin{varwidth}[t]{\sphinxcolwidth{1}{3}}
\sphinxAtStartPar
4
\sphinxbeforeendvarwidth
\end{varwidth}%
&\begin{varwidth}[t]{\sphinxcolwidth{1}{3}}
\sphinxAtStartPar
较差
\sphinxbeforeendvarwidth
\end{varwidth}%
&\begin{varwidth}[t]{\sphinxcolwidth{1}{3}}
\sphinxAtStartPar
数据质量问题较多
\sphinxbeforeendvarwidth
\end{varwidth}%
\\
\sphinxbottomrule
\end{tabulary}
\sphinxtableafterendhook\par
\sphinxattableend\end{savenotes}


\paragraph{2.3.5.2 坐标信息}
\label{\detokenize{chapters/chapter2:id29}}

\begin{savenotes}\sphinxattablestart
\sphinxthistablewithglobalstyle
\centering
\begin{tabulary}{\linewidth}[t]{TTTT}
\sphinxtoprule
\begin{varwidth}[t]{\sphinxcolwidth{1}{4}}
\sphinxstyletheadfamily \sphinxAtStartPar
字段
\sphinxbeforeendvarwidth
\end{varwidth}%
&\begin{varwidth}[t]{\sphinxcolwidth{1}{4}}
\sphinxstyletheadfamily \sphinxAtStartPar
控件名称
\sphinxbeforeendvarwidth
\end{varwidth}%
&\begin{varwidth}[t]{\sphinxcolwidth{1}{4}}
\sphinxstyletheadfamily \sphinxAtStartPar
说明
\sphinxbeforeendvarwidth
\end{varwidth}%
&\begin{varwidth}[t]{\sphinxcolwidth{1}{4}}
\sphinxstyletheadfamily \sphinxAtStartPar
单位
\sphinxbeforeendvarwidth
\end{varwidth}%
\\
\sphinxmidrule
\sphinxtableatstartofbodyhook\begin{varwidth}[t]{\sphinxcolwidth{1}{4}}
\sphinxAtStartPar
经度
\sphinxbeforeendvarwidth
\end{varwidth}%
&\begin{varwidth}[t]{\sphinxcolwidth{1}{4}}
\sphinxAtStartPar
EditLongitude
\sphinxbeforeendvarwidth
\end{varwidth}%
&\begin{varwidth}[t]{\sphinxcolwidth{1}{4}}
\sphinxAtStartPar
测点东向地理坐标
\sphinxbeforeendvarwidth
\end{varwidth}%
&\begin{varwidth}[t]{\sphinxcolwidth{1}{4}}
\sphinxAtStartPar
度
\sphinxbeforeendvarwidth
\end{varwidth}%
\\
\sphinxhline\begin{varwidth}[t]{\sphinxcolwidth{1}{4}}
\sphinxAtStartPar
纬度
\sphinxbeforeendvarwidth
\end{varwidth}%
&\begin{varwidth}[t]{\sphinxcolwidth{1}{4}}
\sphinxAtStartPar
EditLatitude
\sphinxbeforeendvarwidth
\end{varwidth}%
&\begin{varwidth}[t]{\sphinxcolwidth{1}{4}}
\sphinxAtStartPar
测点北向地理坐标
\sphinxbeforeendvarwidth
\end{varwidth}%
&\begin{varwidth}[t]{\sphinxcolwidth{1}{4}}
\sphinxAtStartPar
度
\sphinxbeforeendvarwidth
\end{varwidth}%
\\
\sphinxhline\begin{varwidth}[t]{\sphinxcolwidth{1}{4}}
\sphinxAtStartPar
高程
\sphinxbeforeendvarwidth
\end{varwidth}%
&\begin{varwidth}[t]{\sphinxcolwidth{1}{4}}
\sphinxAtStartPar
EditAltitude
\sphinxbeforeendvarwidth
\end{varwidth}%
&\begin{varwidth}[t]{\sphinxcolwidth{1}{4}}
\sphinxAtStartPar
测点海拔高度
\sphinxbeforeendvarwidth
\end{varwidth}%
&\begin{varwidth}[t]{\sphinxcolwidth{1}{4}}
\sphinxAtStartPar
米
\sphinxbeforeendvarwidth
\end{varwidth}%
\\
\sphinxbottomrule
\end{tabulary}
\sphinxtableafterendhook\par
\sphinxattableend\end{savenotes}


\paragraph{2.3.5.3 通道配置参数}
\label{\detokenize{chapters/chapter2:id30}}
\sphinxAtStartPar
每个测点包含5个电磁分量通道，各通道参数如下：


\begin{savenotes}\sphinxattablestart
\sphinxthistablewithglobalstyle
\centering
\begin{tabulary}{\linewidth}[t]{TTTTTT}
\sphinxtoprule
\begin{varwidth}[t]{\sphinxcolwidth{1}{6}}
\sphinxstyletheadfamily \sphinxAtStartPar
通道
\sphinxbeforeendvarwidth
\end{varwidth}%
&\begin{varwidth}[t]{\sphinxcolwidth{1}{6}}
\sphinxstyletheadfamily \sphinxAtStartPar
索引控件
\sphinxbeforeendvarwidth
\end{varwidth}%
&\begin{varwidth}[t]{\sphinxcolwidth{1}{6}}
\sphinxstyletheadfamily \sphinxAtStartPar
反向控件
\sphinxbeforeendvarwidth
\end{varwidth}%
&\begin{varwidth}[t]{\sphinxcolwidth{1}{6}}
\sphinxstyletheadfamily \sphinxAtStartPar
长度控件
\sphinxbeforeendvarwidth
\end{varwidth}%
&\begin{varwidth}[t]{\sphinxcolwidth{1}{6}}
\sphinxstyletheadfamily \sphinxAtStartPar
传感器ID
\sphinxbeforeendvarwidth
\end{varwidth}%
&\begin{varwidth}[t]{\sphinxcolwidth{1}{6}}
\sphinxstyletheadfamily \sphinxAtStartPar
旋转角控件
\sphinxbeforeendvarwidth
\end{varwidth}%
\\
\sphinxmidrule
\sphinxtableatstartofbodyhook\begin{varwidth}[t]{\sphinxcolwidth{1}{6}}
\sphinxAtStartPar
\sphinxstylestrong{Ex} (电场X)
\sphinxbeforeendvarwidth
\end{varwidth}%
&\begin{varwidth}[t]{\sphinxcolwidth{1}{6}}
\sphinxAtStartPar
NumberBoxExIndex
\sphinxbeforeendvarwidth
\end{varwidth}%
&\begin{varwidth}[t]{\sphinxcolwidth{1}{6}}
\sphinxAtStartPar
CheckBoxExReverse
\sphinxbeforeendvarwidth
\end{varwidth}%
&\begin{varwidth}[t]{\sphinxcolwidth{1}{6}}
\sphinxAtStartPar
NumberBoxExLength
\sphinxbeforeendvarwidth
\end{varwidth}%
&\begin{varwidth}[t]{\sphinxcolwidth{1}{6}}
\sphinxAtStartPar
\sphinxhyphen{}
\sphinxbeforeendvarwidth
\end{varwidth}%
&\begin{varwidth}[t]{\sphinxcolwidth{1}{6}}
\sphinxAtStartPar
NumberBoxExRotate
\sphinxbeforeendvarwidth
\end{varwidth}%
\\
\sphinxhline\begin{varwidth}[t]{\sphinxcolwidth{1}{6}}
\sphinxAtStartPar
\sphinxstylestrong{Ey} (电场Y)
\sphinxbeforeendvarwidth
\end{varwidth}%
&\begin{varwidth}[t]{\sphinxcolwidth{1}{6}}
\sphinxAtStartPar
NumberBoxEyIndex
\sphinxbeforeendvarwidth
\end{varwidth}%
&\begin{varwidth}[t]{\sphinxcolwidth{1}{6}}
\sphinxAtStartPar
CheckBoxEyReverse
\sphinxbeforeendvarwidth
\end{varwidth}%
&\begin{varwidth}[t]{\sphinxcolwidth{1}{6}}
\sphinxAtStartPar
NumberBoxEyLength
\sphinxbeforeendvarwidth
\end{varwidth}%
&\begin{varwidth}[t]{\sphinxcolwidth{1}{6}}
\sphinxAtStartPar
\sphinxhyphen{}
\sphinxbeforeendvarwidth
\end{varwidth}%
&\begin{varwidth}[t]{\sphinxcolwidth{1}{6}}
\sphinxAtStartPar
NumberBoxEyRotate
\sphinxbeforeendvarwidth
\end{varwidth}%
\\
\sphinxhline\begin{varwidth}[t]{\sphinxcolwidth{1}{6}}
\sphinxAtStartPar
\sphinxstylestrong{Hx} (磁场X)
\sphinxbeforeendvarwidth
\end{varwidth}%
&\begin{varwidth}[t]{\sphinxcolwidth{1}{6}}
\sphinxAtStartPar
NumberBoxHxIndex
\sphinxbeforeendvarwidth
\end{varwidth}%
&\begin{varwidth}[t]{\sphinxcolwidth{1}{6}}
\sphinxAtStartPar
CheckBoxHxReverse
\sphinxbeforeendvarwidth
\end{varwidth}%
&\begin{varwidth}[t]{\sphinxcolwidth{1}{6}}
\sphinxAtStartPar
NumberBoxHxLength
\sphinxbeforeendvarwidth
\end{varwidth}%
&\begin{varwidth}[t]{\sphinxcolwidth{1}{6}}
\sphinxAtStartPar
EditHxSensor
\sphinxbeforeendvarwidth
\end{varwidth}%
&\begin{varwidth}[t]{\sphinxcolwidth{1}{6}}
\sphinxAtStartPar
NumberBoxHxRotate
\sphinxbeforeendvarwidth
\end{varwidth}%
\\
\sphinxhline\begin{varwidth}[t]{\sphinxcolwidth{1}{6}}
\sphinxAtStartPar
\sphinxstylestrong{Hy} (磁场Y)
\sphinxbeforeendvarwidth
\end{varwidth}%
&\begin{varwidth}[t]{\sphinxcolwidth{1}{6}}
\sphinxAtStartPar
NumberBoxHyIndex
\sphinxbeforeendvarwidth
\end{varwidth}%
&\begin{varwidth}[t]{\sphinxcolwidth{1}{6}}
\sphinxAtStartPar
CheckBoxHyReverse
\sphinxbeforeendvarwidth
\end{varwidth}%
&\begin{varwidth}[t]{\sphinxcolwidth{1}{6}}
\sphinxAtStartPar
NumberBoxHyLength
\sphinxbeforeendvarwidth
\end{varwidth}%
&\begin{varwidth}[t]{\sphinxcolwidth{1}{6}}
\sphinxAtStartPar
EditHySensor
\sphinxbeforeendvarwidth
\end{varwidth}%
&\begin{varwidth}[t]{\sphinxcolwidth{1}{6}}
\sphinxAtStartPar
NumberBoxHyRotate
\sphinxbeforeendvarwidth
\end{varwidth}%
\\
\sphinxhline\begin{varwidth}[t]{\sphinxcolwidth{1}{6}}
\sphinxAtStartPar
\sphinxstylestrong{Hz} (磁场Z)
\sphinxbeforeendvarwidth
\end{varwidth}%
&\begin{varwidth}[t]{\sphinxcolwidth{1}{6}}
\sphinxAtStartPar
NumberBoxHzIndex
\sphinxbeforeendvarwidth
\end{varwidth}%
&\begin{varwidth}[t]{\sphinxcolwidth{1}{6}}
\sphinxAtStartPar
CheckBoxHzReverse
\sphinxbeforeendvarwidth
\end{varwidth}%
&\begin{varwidth}[t]{\sphinxcolwidth{1}{6}}
\sphinxAtStartPar
NumberBoxHzLength
\sphinxbeforeendvarwidth
\end{varwidth}%
&\begin{varwidth}[t]{\sphinxcolwidth{1}{6}}
\sphinxAtStartPar
EditHzSensor
\sphinxbeforeendvarwidth
\end{varwidth}%
&\begin{varwidth}[t]{\sphinxcolwidth{1}{6}}
\sphinxAtStartPar
NumberBoxHzRotate
\sphinxbeforeendvarwidth
\end{varwidth}%
\\
\sphinxbottomrule
\end{tabulary}
\sphinxtableafterendhook\par
\sphinxattableend\end{savenotes}

\sphinxAtStartPar
\sphinxstylestrong{通道参数说明：}


\begin{savenotes}\sphinxattablestart
\sphinxthistablewithglobalstyle
\centering
\begin{tabulary}{\linewidth}[t]{TTT}
\sphinxtoprule
\begin{varwidth}[t]{\sphinxcolwidth{1}{3}}
\sphinxstyletheadfamily \sphinxAtStartPar
参数
\sphinxbeforeendvarwidth
\end{varwidth}%
&\begin{varwidth}[t]{\sphinxcolwidth{1}{3}}
\sphinxstyletheadfamily \sphinxAtStartPar
说明
\sphinxbeforeendvarwidth
\end{varwidth}%
&\begin{varwidth}[t]{\sphinxcolwidth{1}{3}}
\sphinxstyletheadfamily \sphinxAtStartPar
注意事项
\sphinxbeforeendvarwidth
\end{varwidth}%
\\
\sphinxmidrule
\sphinxtableatstartofbodyhook\begin{varwidth}[t]{\sphinxcolwidth{1}{3}}
\sphinxAtStartPar
\sphinxstylestrong{索引}
\sphinxbeforeendvarwidth
\end{varwidth}%
&\begin{varwidth}[t]{\sphinxcolwidth{1}{3}}
\sphinxAtStartPar
通道在数据文件中的位置编号
\sphinxbeforeendvarwidth
\end{varwidth}%
&\begin{varwidth}[t]{\sphinxcolwidth{1}{3}}
\sphinxAtStartPar
通常自动识别
\sphinxbeforeendvarwidth
\end{varwidth}%
\\
\sphinxhline\begin{varwidth}[t]{\sphinxcolwidth{1}{3}}
\sphinxAtStartPar
\sphinxstylestrong{反向}
\sphinxbeforeendvarwidth
\end{varwidth}%
&\begin{varwidth}[t]{\sphinxcolwidth{1}{3}}
\sphinxAtStartPar
是否反转通道极性
\sphinxbeforeendvarwidth
\end{varwidth}%
&\begin{varwidth}[t]{\sphinxcolwidth{1}{3}}
\sphinxAtStartPar
用于修正电极接线错误
\sphinxbeforeendvarwidth
\end{varwidth}%
\\
\sphinxhline\begin{varwidth}[t]{\sphinxcolwidth{1}{3}}
\sphinxAtStartPar
\sphinxstylestrong{长度}
\sphinxbeforeendvarwidth
\end{varwidth}%
&\begin{varwidth}[t]{\sphinxcolwidth{1}{3}}
\sphinxAtStartPar
电极间距或传感器序列号
\sphinxbeforeendvarwidth
\end{varwidth}%
&\begin{varwidth}[t]{\sphinxcolwidth{1}{3}}
\sphinxAtStartPar
电场为电极距，磁场为传感器编号
\sphinxbeforeendvarwidth
\end{varwidth}%
\\
\sphinxhline\begin{varwidth}[t]{\sphinxcolwidth{1}{3}}
\sphinxAtStartPar
\sphinxstylestrong{传感器ID}
\sphinxbeforeendvarwidth
\end{varwidth}%
&\begin{varwidth}[t]{\sphinxcolwidth{1}{3}}
\sphinxAtStartPar
磁传感器的唯一标识
\sphinxbeforeendvarwidth
\end{varwidth}%
&\begin{varwidth}[t]{\sphinxcolwidth{1}{3}}
\sphinxAtStartPar
仅磁场通道有此参数
\sphinxbeforeendvarwidth
\end{varwidth}%
\\
\sphinxhline\begin{varwidth}[t]{\sphinxcolwidth{1}{3}}
\sphinxAtStartPar
\sphinxstylestrong{旋转角}
\sphinxbeforeendvarwidth
\end{varwidth}%
&\begin{varwidth}[t]{\sphinxcolwidth{1}{3}}
\sphinxAtStartPar
通道相对于真北的旋转角度
\sphinxbeforeendvarwidth
\end{varwidth}%
&\begin{varwidth}[t]{\sphinxcolwidth{1}{3}}
\sphinxAtStartPar
用于坐标旋转校正
\sphinxbeforeendvarwidth
\end{varwidth}%
\\
\sphinxbottomrule
\end{tabulary}
\sphinxtableafterendhook\par
\sphinxattableend\end{savenotes}


\paragraph{2.3.5.4 采集盒信息}
\label{\detokenize{chapters/chapter2:id31}}

\begin{savenotes}\sphinxattablestart
\sphinxthistablewithglobalstyle
\centering
\begin{tabulary}{\linewidth}[t]{TTTT}
\sphinxtoprule
\begin{varwidth}[t]{\sphinxcolwidth{1}{4}}
\sphinxstyletheadfamily \sphinxAtStartPar
字段
\sphinxbeforeendvarwidth
\end{varwidth}%
&\begin{varwidth}[t]{\sphinxcolwidth{1}{4}}
\sphinxstyletheadfamily \sphinxAtStartPar
控件名称
\sphinxbeforeendvarwidth
\end{varwidth}%
&\begin{varwidth}[t]{\sphinxcolwidth{1}{4}}
\sphinxstyletheadfamily \sphinxAtStartPar
说明
\sphinxbeforeendvarwidth
\end{varwidth}%
&\begin{varwidth}[t]{\sphinxcolwidth{1}{4}}
\sphinxstyletheadfamily \sphinxAtStartPar
单位
\sphinxbeforeendvarwidth
\end{varwidth}%
\\
\sphinxmidrule
\sphinxtableatstartofbodyhook\begin{varwidth}[t]{\sphinxcolwidth{1}{4}}
\sphinxAtStartPar
采集盒ID
\sphinxbeforeendvarwidth
\end{varwidth}%
&\begin{varwidth}[t]{\sphinxcolwidth{1}{4}}
\sphinxAtStartPar
EditBox
\sphinxbeforeendvarwidth
\end{varwidth}%
&\begin{varwidth}[t]{\sphinxcolwidth{1}{4}}
\sphinxAtStartPar
MT采集仪器的唯一序列号
\sphinxbeforeendvarwidth
\end{varwidth}%
&\begin{varwidth}[t]{\sphinxcolwidth{1}{4}}
\sphinxAtStartPar
\sphinxhyphen{}
\sphinxbeforeendvarwidth
\end{varwidth}%
\\
\sphinxhline\begin{varwidth}[t]{\sphinxcolwidth{1}{4}}
\sphinxAtStartPar
电场前置放大器
\sphinxbeforeendvarwidth
\end{varwidth}%
&\begin{varwidth}[t]{\sphinxcolwidth{1}{4}}
\sphinxAtStartPar
EditEPreamplifier
\sphinxbeforeendvarwidth
\end{varwidth}%
&\begin{varwidth}[t]{\sphinxcolwidth{1}{4}}
\sphinxAtStartPar
前置放大器型号/编号
\sphinxbeforeendvarwidth
\end{varwidth}%
&\begin{varwidth}[t]{\sphinxcolwidth{1}{4}}
\sphinxAtStartPar
\sphinxhyphen{}
\sphinxbeforeendvarwidth
\end{varwidth}%
\\
\sphinxhline\begin{varwidth}[t]{\sphinxcolwidth{1}{4}}
\sphinxAtStartPar
Ex接地电阻
\sphinxbeforeendvarwidth
\end{varwidth}%
&\begin{varwidth}[t]{\sphinxcolwidth{1}{4}}
\sphinxAtStartPar
NumberBoxGroundResEx
\sphinxbeforeendvarwidth
\end{varwidth}%
&\begin{varwidth}[t]{\sphinxcolwidth{1}{4}}
\sphinxAtStartPar
Ex通道电极接地电阻
\sphinxbeforeendvarwidth
\end{varwidth}%
&\begin{varwidth}[t]{\sphinxcolwidth{1}{4}}
\sphinxAtStartPar
Ω
\sphinxbeforeendvarwidth
\end{varwidth}%
\\
\sphinxhline\begin{varwidth}[t]{\sphinxcolwidth{1}{4}}
\sphinxAtStartPar
Ey接地电阻
\sphinxbeforeendvarwidth
\end{varwidth}%
&\begin{varwidth}[t]{\sphinxcolwidth{1}{4}}
\sphinxAtStartPar
NumberBoxGroundResEy
\sphinxbeforeendvarwidth
\end{varwidth}%
&\begin{varwidth}[t]{\sphinxcolwidth{1}{4}}
\sphinxAtStartPar
Ey通道电极接地电阻
\sphinxbeforeendvarwidth
\end{varwidth}%
&\begin{varwidth}[t]{\sphinxcolwidth{1}{4}}
\sphinxAtStartPar
Ω
\sphinxbeforeendvarwidth
\end{varwidth}%
\\
\sphinxbottomrule
\end{tabulary}
\sphinxtableafterendhook\par
\sphinxattableend\end{savenotes}

\sphinxAtStartPar
\sphinxstylestrong{接地电阻说明：}
\begin{itemize}
\item {} 
\sphinxAtStartPar
理想值：< 1kΩ

\item {} 
\sphinxAtStartPar
可接受：< 5kΩ

\item {} 
\sphinxAtStartPar
需改善：> 5kΩ（可能影响数据质量）

\end{itemize}


\paragraph{2.3.5.5 时间信息}
\label{\detokenize{chapters/chapter2:id32}}

\begin{savenotes}\sphinxattablestart
\sphinxthistablewithglobalstyle
\centering
\begin{tabulary}{\linewidth}[t]{TTT}
\sphinxtoprule
\begin{varwidth}[t]{\sphinxcolwidth{1}{3}}
\sphinxstyletheadfamily \sphinxAtStartPar
字段
\sphinxbeforeendvarwidth
\end{varwidth}%
&\begin{varwidth}[t]{\sphinxcolwidth{1}{3}}
\sphinxstyletheadfamily \sphinxAtStartPar
控件名称
\sphinxbeforeendvarwidth
\end{varwidth}%
&\begin{varwidth}[t]{\sphinxcolwidth{1}{3}}
\sphinxstyletheadfamily \sphinxAtStartPar
说明
\sphinxbeforeendvarwidth
\end{varwidth}%
\\
\sphinxmidrule
\sphinxtableatstartofbodyhook\begin{varwidth}[t]{\sphinxcolwidth{1}{3}}
\sphinxAtStartPar
采集开始时间
\sphinxbeforeendvarwidth
\end{varwidth}%
&\begin{varwidth}[t]{\sphinxcolwidth{1}{3}}
\sphinxAtStartPar
SiteBeginEnd.BeginDateTime
\sphinxbeforeendvarwidth
\end{varwidth}%
&\begin{varwidth}[t]{\sphinxcolwidth{1}{3}}
\sphinxAtStartPar
时间序列数据起始时刻
\sphinxbeforeendvarwidth
\end{varwidth}%
\\
\sphinxhline\begin{varwidth}[t]{\sphinxcolwidth{1}{3}}
\sphinxAtStartPar
采集结束时间
\sphinxbeforeendvarwidth
\end{varwidth}%
&\begin{varwidth}[t]{\sphinxcolwidth{1}{3}}
\sphinxAtStartPar
SiteBeginEnd.EndDateTime
\sphinxbeforeendvarwidth
\end{varwidth}%
&\begin{varwidth}[t]{\sphinxcolwidth{1}{3}}
\sphinxAtStartPar
时间序列数据结束时刻
\sphinxbeforeendvarwidth
\end{varwidth}%
\\
\sphinxhline\begin{varwidth}[t]{\sphinxcolwidth{1}{3}}
\sphinxAtStartPar
处理开始时间
\sphinxbeforeendvarwidth
\end{varwidth}%
&\begin{varwidth}[t]{\sphinxcolwidth{1}{3}}
\sphinxAtStartPar
ProcessBeginEnd.BeginDateTime
\sphinxbeforeendvarwidth
\end{varwidth}%
&\begin{varwidth}[t]{\sphinxcolwidth{1}{3}}
\sphinxAtStartPar
用于处理的时间窗口起始
\sphinxbeforeendvarwidth
\end{varwidth}%
\\
\sphinxhline\begin{varwidth}[t]{\sphinxcolwidth{1}{3}}
\sphinxAtStartPar
处理结束时间
\sphinxbeforeendvarwidth
\end{varwidth}%
&\begin{varwidth}[t]{\sphinxcolwidth{1}{3}}
\sphinxAtStartPar
ProcessBeginEnd.EndDateTime
\sphinxbeforeendvarwidth
\end{varwidth}%
&\begin{varwidth}[t]{\sphinxcolwidth{1}{3}}
\sphinxAtStartPar
用于处理的时间窗口结束
\sphinxbeforeendvarwidth
\end{varwidth}%
\\
\sphinxhline\begin{varwidth}[t]{\sphinxcolwidth{1}{3}}
\sphinxAtStartPar
采集时长
\sphinxbeforeendvarwidth
\end{varwidth}%
&\begin{varwidth}[t]{\sphinxcolwidth{1}{3}}
\sphinxAtStartPar
LabelTimeLengthH
\sphinxbeforeendvarwidth
\end{varwidth}%
&\begin{varwidth}[t]{\sphinxcolwidth{1}{3}}
\sphinxAtStartPar
自动计算的采集持续时间
\sphinxbeforeendvarwidth
\end{varwidth}%
\\
\sphinxbottomrule
\end{tabulary}
\sphinxtableafterendhook\par
\sphinxattableend\end{savenotes}


\paragraph{2.3.5.6 频率信息}
\label{\detokenize{chapters/chapter2:id33}}

\begin{savenotes}\sphinxattablestart
\sphinxthistablewithglobalstyle
\centering
\begin{tabulary}{\linewidth}[t]{TTTT}
\sphinxtoprule
\begin{varwidth}[t]{\sphinxcolwidth{1}{4}}
\sphinxstyletheadfamily \sphinxAtStartPar
字段
\sphinxbeforeendvarwidth
\end{varwidth}%
&\begin{varwidth}[t]{\sphinxcolwidth{1}{4}}
\sphinxstyletheadfamily \sphinxAtStartPar
控件名称
\sphinxbeforeendvarwidth
\end{varwidth}%
&\begin{varwidth}[t]{\sphinxcolwidth{1}{4}}
\sphinxstyletheadfamily \sphinxAtStartPar
说明
\sphinxbeforeendvarwidth
\end{varwidth}%
&\begin{varwidth}[t]{\sphinxcolwidth{1}{4}}
\sphinxstyletheadfamily \sphinxAtStartPar
单位
\sphinxbeforeendvarwidth
\end{varwidth}%
\\
\sphinxmidrule
\sphinxtableatstartofbodyhook\begin{varwidth}[t]{\sphinxcolwidth{1}{4}}
\sphinxAtStartPar
最大可用频率
\sphinxbeforeendvarwidth
\end{varwidth}%
&\begin{varwidth}[t]{\sphinxcolwidth{1}{4}}
\sphinxAtStartPar
EditMaxAvailableFrequency
\sphinxbeforeendvarwidth
\end{varwidth}%
&\begin{varwidth}[t]{\sphinxcolwidth{1}{4}}
\sphinxAtStartPar
数据中包含的最高频率
\sphinxbeforeendvarwidth
\end{varwidth}%
&\begin{varwidth}[t]{\sphinxcolwidth{1}{4}}
\sphinxAtStartPar
Hz
\sphinxbeforeendvarwidth
\end{varwidth}%
\\
\sphinxhline\begin{varwidth}[t]{\sphinxcolwidth{1}{4}}
\sphinxAtStartPar
最小可用频率
\sphinxbeforeendvarwidth
\end{varwidth}%
&\begin{varwidth}[t]{\sphinxcolwidth{1}{4}}
\sphinxAtStartPar
EditMinAvailableFrequency
\sphinxbeforeendvarwidth
\end{varwidth}%
&\begin{varwidth}[t]{\sphinxcolwidth{1}{4}}
\sphinxAtStartPar
数据中包含的最低频率
\sphinxbeforeendvarwidth
\end{varwidth}%
&\begin{varwidth}[t]{\sphinxcolwidth{1}{4}}
\sphinxAtStartPar
Hz
\sphinxbeforeendvarwidth
\end{varwidth}%
\\
\sphinxbottomrule
\end{tabulary}
\sphinxtableafterendhook\par
\sphinxattableend\end{savenotes}


\paragraph{2.3.5.7 特殊功能}
\label{\detokenize{chapters/chapter2:id34}}
\sphinxAtStartPar
\sphinxstylestrong{从TBL文件加载：}
\begin{itemize}
\item {} 
\sphinxAtStartPar
支持拖放 \sphinxcode{\sphinxupquote{.tbl}} 文件更新Phoenix测点信息

\item {} 
\sphinxAtStartPar
自动读取：坐标、时间、采集盒ID、传感器ID、电极长度、旋转角

\end{itemize}

\sphinxAtStartPar
\sphinxstylestrong{从LEMI文件加载：}
\begin{itemize}
\item {} 
\sphinxAtStartPar
支持拖放 \sphinxcode{\sphinxupquote{.lemi}} 文件

\item {} 
\sphinxAtStartPar
更新LEMI通道索引

\end{itemize}

\sphinxAtStartPar
\sphinxstylestrong{加载傅里叶系数：}
\begin{itemize}
\item {} 
\sphinxAtStartPar
支持拖放 \sphinxcode{\sphinxupquote{.stfc}} 文件

\end{itemize}

\sphinxAtStartPar
\sphinxstylestrong{更新Phoenix TBL：}
\begin{itemize}
\item {} 
\sphinxAtStartPar
修改坐标、采集盒、传感器信息后自动备份原TBL到 \sphinxcode{\sphinxupquote{.tbl.backup}}

\end{itemize}

\sphinxAtStartPar
\sphinxstylestrong{更新RMT JSON：}
\begin{itemize}
\item {} 
\sphinxAtStartPar
修改坐标、设备信息后更新JSON文件，自动备份原文件

\end{itemize}


\bigskip\hrule\bigskip


\sphinxAtStartPar
⬅️ {\hyperref[\detokenize{chapters/chapter2:id6}]{\sphinxcrossref{上一节：MT基本原理}}} | ➡️ {\hyperref[\detokenize{chapters/chapter2:id35}]{\sphinxcrossref{下一节：从时间序列到频谱}}}


\bigskip\hrule\bigskip



\subsection{2.4 第三步：从时间序列到频谱}
\label{\detokenize{chapters/chapter2:id35}}

\subsubsection{2.4.1 傅里叶变换基础}
\label{\detokenize{chapters/chapter2:id36}}

\paragraph{为什么需要FFT？}
\label{\detokenize{chapters/chapter2:fft}}
\sphinxAtStartPar
\sphinxstylestrong{原始数据}是时间域的电场和磁场信号，但MT分析需要在\sphinxstylestrong{频率域}进行——我们需要知道不同频率成分的振幅和相位。

\sphinxAtStartPar
\sphinxstylestrong{形象理解}：就像把一首音乐分解成不同音调（频率）的音符。


\paragraph{MTDP中的FFT参数}
\label{\detokenize{chapters/chapter2:mtdpfft}}

\begin{savenotes}\sphinxattablestart
\sphinxthistablewithglobalstyle
\centering
\begin{tabulary}{\linewidth}[t]{TTT}
\sphinxtoprule
\begin{varwidth}[t]{\sphinxcolwidth{1}{3}}
\sphinxstyletheadfamily \sphinxAtStartPar
参数
\sphinxbeforeendvarwidth
\end{varwidth}%
&\begin{varwidth}[t]{\sphinxcolwidth{1}{3}}
\sphinxstyletheadfamily \sphinxAtStartPar
推荐值
\sphinxbeforeendvarwidth
\end{varwidth}%
&\begin{varwidth}[t]{\sphinxcolwidth{1}{3}}
\sphinxstyletheadfamily \sphinxAtStartPar
说明
\sphinxbeforeendvarwidth
\end{varwidth}%
\\
\sphinxmidrule
\sphinxtableatstartofbodyhook\begin{varwidth}[t]{\sphinxcolwidth{1}{3}}
\sphinxAtStartPar
\sphinxstylestrong{窗长度}
\sphinxbeforeendvarwidth
\end{varwidth}%
&\begin{varwidth}[t]{\sphinxcolwidth{1}{3}}
\sphinxAtStartPar
2\textasciicircum{}n 点
\sphinxbeforeendvarwidth
\end{varwidth}%
&\begin{varwidth}[t]{\sphinxcolwidth{1}{3}}
\sphinxAtStartPar
常用 1024, 2048, 4096
\sphinxbeforeendvarwidth
\end{varwidth}%
\\
\sphinxhline\begin{varwidth}[t]{\sphinxcolwidth{1}{3}}
\sphinxAtStartPar
\sphinxstylestrong{窗函数}
\sphinxbeforeendvarwidth
\end{varwidth}%
&\begin{varwidth}[t]{\sphinxcolwidth{1}{3}}
\sphinxAtStartPar
汉宁窗
\sphinxbeforeendvarwidth
\end{varwidth}%
&\begin{varwidth}[t]{\sphinxcolwidth{1}{3}}
\sphinxAtStartPar
减少频谱泄漏
\sphinxbeforeendvarwidth
\end{varwidth}%
\\
\sphinxhline\begin{varwidth}[t]{\sphinxcolwidth{1}{3}}
\sphinxAtStartPar
\sphinxstylestrong{重叠率}
\sphinxbeforeendvarwidth
\end{varwidth}%
&\begin{varwidth}[t]{\sphinxcolwidth{1}{3}}
\sphinxAtStartPar
50\%
\sphinxbeforeendvarwidth
\end{varwidth}%
&\begin{varwidth}[t]{\sphinxcolwidth{1}{3}}
\sphinxAtStartPar
提高数据利用率
\sphinxbeforeendvarwidth
\end{varwidth}%
\\
\sphinxbottomrule
\end{tabulary}
\sphinxtableafterendhook\par
\sphinxattableend\end{savenotes}

\begin{sphinxuseclass}{danger}\begin{enumerate}
\sphinxsetlistlabels{\arabic}{enumi}{enumii}{}{.}%
\item {} 
\sphinxAtStartPar
\sphinxstylestrong{窗长度选择不当}：
\begin{itemize}
\item {} 
\sphinxAtStartPar
窗太短 → 频率分辨率低，低频精度差

\item {} 
\sphinxAtStartPar
窗太长 → 时间分辨率低，难以剔除短时噪声

\end{itemize}

\item {} 
\sphinxAtStartPar
\sphinxstylestrong{忽略去均值}：FFT前未去除直流分量会导致零频附近失真

\end{enumerate}

\end{sphinxuseclass}

\subsubsection{2.4.2 功率谱估计}
\label{\detokenize{chapters/chapter2:id37}}

\paragraph{自功率谱与互功率谱}
\label{\detokenize{chapters/chapter2:id38}}\begin{itemize}
\item {} 
\sphinxAtStartPar
\sphinxstylestrong{自功率谱}：\([XX^*]\) — 单个信号的能量

\item {} 
\sphinxAtStartPar
\sphinxstylestrong{互功率谱}：\([XY^*]\) — 两个信号的相关能量

\end{itemize}
\begin{equation*}
\begin{split}[XY^*] = \frac{1}{N}\sum_{i=1}^{N} X_i(f) \cdot Y_i^*(f) \text{2.7}\end{split}
\end{equation*}

\paragraph{叠加的作用}
\label{\detokenize{chapters/chapter2:id39}}
\sphinxAtStartPar
将时间序列分成多段，分别FFT后叠加平均：
\begin{equation*}
\begin{split}[XY^*]_{avg} = \frac{1}{N}\sum_{i=1}^{N} X_i Y_i^* \text{2.8}\end{split}
\end{equation*}
\sphinxAtStartPar
\sphinxstylestrong{叠加的好处}：
\begin{itemize}
\item {} 
\sphinxAtStartPar
随机噪声在平均中相互抵消

\item {} 
\sphinxAtStartPar
信号在平均中增强

\item {} 
\sphinxAtStartPar
信噪比提升 \(\sqrt{N}\) 倍

\end{itemize}

\begin{sphinxuseclass}{tip}
\sphinxAtStartPar
叠加次数 N = 100 时，信噪比可提升约 10 倍（\(\sqrt{100} = 10\)）。

\end{sphinxuseclass}

\subsubsection{2.4.3 数据分组与叠加}
\label{\detokenize{chapters/chapter2:id40}}

\paragraph{分组策略}
\label{\detokenize{chapters/chapter2:id41}}
\sphinxAtStartPar
MTDP支持将长时间序列分成多个**时间段（分组）**进行处理：


\begin{savenotes}\sphinxattablestart
\sphinxthistablewithglobalstyle
\centering
\begin{tabulary}{\linewidth}[t]{TT}
\sphinxtoprule
\begin{varwidth}[t]{\sphinxcolwidth{1}{2}}
\sphinxstyletheadfamily \sphinxAtStartPar
分组方式
\sphinxbeforeendvarwidth
\end{varwidth}%
&\begin{varwidth}[t]{\sphinxcolwidth{1}{2}}
\sphinxstyletheadfamily \sphinxAtStartPar
适用场景
\sphinxbeforeendvarwidth
\end{varwidth}%
\\
\sphinxmidrule
\sphinxtableatstartofbodyhook\begin{varwidth}[t]{\sphinxcolwidth{1}{2}}
\sphinxAtStartPar
\sphinxstylestrong{等长分组}
\sphinxbeforeendvarwidth
\end{varwidth}%
&\begin{varwidth}[t]{\sphinxcolwidth{1}{2}}
\sphinxAtStartPar
信号稳定的环境
\sphinxbeforeendvarwidth
\end{varwidth}%
\\
\sphinxhline\begin{varwidth}[t]{\sphinxcolwidth{1}{2}}
\sphinxAtStartPar
\sphinxstylestrong{按时间分段}
\sphinxbeforeendvarwidth
\end{varwidth}%
&\begin{varwidth}[t]{\sphinxcolwidth{1}{2}}
\sphinxAtStartPar
需要剔除特定时段（如白天干扰）
\sphinxbeforeendvarwidth
\end{varwidth}%
\\
\sphinxhline\begin{varwidth}[t]{\sphinxcolwidth{1}{2}}
\sphinxAtStartPar
\sphinxstylestrong{按质量筛选}
\sphinxbeforeendvarwidth
\end{varwidth}%
&\begin{varwidth}[t]{\sphinxcolwidth{1}{2}}
\sphinxAtStartPar
自动剔除低质量分组
\sphinxbeforeendvarwidth
\end{varwidth}%
\\
\sphinxbottomrule
\end{tabulary}
\sphinxtableafterendhook\par
\sphinxattableend\end{savenotes}


\subsubsection{2.4.4 MTDP的FFT操作}
\label{\detokenize{chapters/chapter2:id42}}
\begin{sphinxuseclass}{details}
\sphinxAtStartPar
\sphinxstylestrong{执行FFT处理：}
\begin{enumerate}
\sphinxsetlistlabels{\arabic}{enumi}{enumii}{}{.}%
\item {} 
\sphinxAtStartPar
菜单：\sphinxstylestrong{处理 → FFT设置}（或快捷键 \sphinxcode{\sphinxupquote{Ctrl+Shift+F}}）

\item {} 
\sphinxAtStartPar
设置参数：
\begin{itemize}
\item {} 
\sphinxAtStartPar
窗长度：4096点（推荐）

\item {} 
\sphinxAtStartPar
窗函数：Hann窗

\item {} 
\sphinxAtStartPar
重叠率：50\%

\end{itemize}

\item {} 
\sphinxAtStartPar
点击 \sphinxstylestrong{执行FFT}

\end{enumerate}

\sphinxAtStartPar
\sphinxstylestrong{查看功率谱：}
\begin{enumerate}
\sphinxsetlistlabels{\arabic}{enumi}{enumii}{}{.}%
\item {} 
\sphinxAtStartPar
切换到 \sphinxstylestrong{功率谱} 选项卡

\item {} 
\sphinxAtStartPar
选择要查看的通道（Ex, Ey, Hx, Hy, Hz）

\item {} 
\sphinxAtStartPar
可切换时间/频率视图

\end{enumerate}

\end{sphinxuseclass}

\bigskip\hrule\bigskip


\sphinxAtStartPar
⬅️ {\hyperref[\detokenize{chapters/chapter2:id17}]{\sphinxcrossref{上一节：数据采集准备}}} | ➡️ {\hyperref[\detokenize{chapters/chapter2:id43}]{\sphinxcrossref{下一节：阻抗估计}}}


\bigskip\hrule\bigskip



\subsection{2.5 第四步：阻抗估计}
\label{\detokenize{chapters/chapter2:id43}}

\subsubsection{2.5.1 方法对比总览}
\label{\detokenize{chapters/chapter2:id44}}

\begin{savenotes}\sphinxattablestart
\sphinxthistablewithglobalstyle
\centering
\begin{tabulary}{\linewidth}[t]{TTTT}
\sphinxtoprule
\begin{varwidth}[t]{\sphinxcolwidth{1}{4}}
\sphinxstyletheadfamily \sphinxAtStartPar
方法
\sphinxbeforeendvarwidth
\end{varwidth}%
&\begin{varwidth}[t]{\sphinxcolwidth{1}{4}}
\sphinxstyletheadfamily \sphinxAtStartPar
优点
\sphinxbeforeendvarwidth
\end{varwidth}%
&\begin{varwidth}[t]{\sphinxcolwidth{1}{4}}
\sphinxstyletheadfamily \sphinxAtStartPar
缺点
\sphinxbeforeendvarwidth
\end{varwidth}%
&\begin{varwidth}[t]{\sphinxcolwidth{1}{4}}
\sphinxstyletheadfamily \sphinxAtStartPar
适用场景
\sphinxbeforeendvarwidth
\end{varwidth}%
\\
\sphinxmidrule
\sphinxtableatstartofbodyhook\begin{varwidth}[t]{\sphinxcolwidth{1}{4}}
\sphinxAtStartPar
\sphinxstylestrong{最小二乘法}
\sphinxbeforeendvarwidth
\end{varwidth}%
&\begin{varwidth}[t]{\sphinxcolwidth{1}{4}}
\sphinxAtStartPar
简单快速
\sphinxbeforeendvarwidth
\end{varwidth}%
&\begin{varwidth}[t]{\sphinxcolwidth{1}{4}}
\sphinxAtStartPar
对噪声敏感
\sphinxbeforeendvarwidth
\end{varwidth}%
&\begin{varwidth}[t]{\sphinxcolwidth{1}{4}}
\sphinxAtStartPar
理想低噪声环境
\sphinxbeforeendvarwidth
\end{varwidth}%
\\
\sphinxhline\begin{varwidth}[t]{\sphinxcolwidth{1}{4}}
\sphinxAtStartPar
\sphinxstylestrong{远参考道法}
\sphinxbeforeendvarwidth
\end{varwidth}%
&\begin{varwidth}[t]{\sphinxcolwidth{1}{4}}
\sphinxAtStartPar
消除输入端偏差
\sphinxbeforeendvarwidth
\end{varwidth}%
&\begin{varwidth}[t]{\sphinxcolwidth{1}{4}}
\sphinxAtStartPar
需要参考站
\sphinxbeforeendvarwidth
\end{varwidth}%
&\begin{varwidth}[t]{\sphinxcolwidth{1}{4}}
\sphinxAtStartPar
\sphinxstylestrong{推荐首选}
\sphinxbeforeendvarwidth
\end{varwidth}%
\\
\sphinxhline\begin{varwidth}[t]{\sphinxcolwidth{1}{4}}
\sphinxAtStartPar
\sphinxstylestrong{稳健估计}
\sphinxbeforeendvarwidth
\end{varwidth}%
&\begin{varwidth}[t]{\sphinxcolwidth{1}{4}}
\sphinxAtStartPar
抗异常值
\sphinxbeforeendvarwidth
\end{varwidth}%
&\begin{varwidth}[t]{\sphinxcolwidth{1}{4}}
\sphinxAtStartPar
计算复杂
\sphinxbeforeendvarwidth
\end{varwidth}%
&\begin{varwidth}[t]{\sphinxcolwidth{1}{4}}
\sphinxAtStartPar
有少量强干扰
\sphinxbeforeendvarwidth
\end{varwidth}%
\\
\sphinxhline\begin{varwidth}[t]{\sphinxcolwidth{1}{4}}
\sphinxAtStartPar
\sphinxstylestrong{RMSMR}
\sphinxbeforeendvarwidth
\end{varwidth}%
&\begin{varwidth}[t]{\sphinxcolwidth{1}{4}}
\sphinxAtStartPar
综合处理能力强
\sphinxbeforeendvarwidth
\end{varwidth}%
&\begin{varwidth}[t]{\sphinxcolwidth{1}{4}}
\sphinxAtStartPar
需要迭代
\sphinxbeforeendvarwidth
\end{varwidth}%
&\begin{varwidth}[t]{\sphinxcolwidth{1}{4}}
\sphinxAtStartPar
\sphinxstylestrong{强干扰环境}
\sphinxbeforeendvarwidth
\end{varwidth}%
\\
\sphinxbottomrule
\end{tabulary}
\sphinxtableafterendhook\par
\sphinxattableend\end{savenotes}


\subsubsection{2.5.2 方法选择决策树}
\label{\detokenize{chapters/chapter2:id45}}

\subsubsection{2.5.3 最小二乘法（入门理解）}
\label{\detokenize{chapters/chapter2:id46}}

\paragraph{基本思想}
\label{\detokenize{chapters/chapter2:id47}}
\sphinxAtStartPar
找到阻抗 Z，使得预测误差最小：
\begin{equation*}
\begin{split}\min_Z \sum |E - ZH|^2 \text{2.9}\end{split}
\end{equation*}
\sphinxAtStartPar
解为：
\begin{equation*}
\begin{split}\mathbf{Z} = [\mathbf{E}\mathbf{H}^*][\mathbf{H}\mathbf{H}^*]^{-1} \text{2.10}\end{split}
\end{equation*}

\paragraph{为什么不推荐？}
\label{\detokenize{chapters/chapter2:id48}}
\sphinxAtStartPar
\sphinxstylestrong{致命弱点}：假设磁场 H 无噪声。但实际中：
\begin{itemize}
\item {} 
\sphinxAtStartPar
磁场同样受噪声污染

\item {} 
\sphinxAtStartPar
噪声会导致阻抗\sphinxstylestrong{系统偏差}

\item {} 
\sphinxAtStartPar
一个异常数据点就能破坏估计

\end{itemize}

\begin{sphinxuseclass}{warning}
\sphinxAtStartPar
最小二乘法主要用于理解原理，实际处理应使用远参考道法或稳健方法。

\end{sphinxuseclass}

\subsubsection{2.5.4 远参考道法（推荐）}
\label{\detokenize{chapters/chapter2:id49}}

\paragraph{核心思想}
\label{\detokenize{chapters/chapter2:id50}}
\sphinxAtStartPar
使用远参考站的磁场 R 替代本地磁场 H：
\begin{equation*}
\begin{split}\mathbf{Z} = [\mathbf{E}\mathbf{R}^*][\mathbf{H}\mathbf{R}^*]^{-1} \text{2.11}\end{split}
\end{equation*}

\paragraph{为什么有效？}
\label{\detokenize{chapters/chapter2:id51}}

\paragraph{使用条件}
\label{\detokenize{chapters/chapter2:id52}}\begin{enumerate}
\sphinxsetlistlabels{\arabic}{enumi}{enumii}{}{.}%
\item {} 
\sphinxAtStartPar
✅ 远参考站与测站\sphinxstylestrong{时间同步}

\item {} 
\sphinxAtStartPar
✅ 两站\sphinxstylestrong{噪声不相关}

\item {} 
\sphinxAtStartPar
✅ 两站距离适中（几公里到几十公里）

\end{enumerate}


\subsubsection{2.5.5 稳健估计（进阶）}
\label{\detokenize{chapters/chapter2:id53}}

\paragraph{回归M估计}
\label{\detokenize{chapters/chapter2:m}}
\sphinxAtStartPar
\sphinxstylestrong{核心思想}：给异常数据\sphinxstylestrong{降低权重}
\begin{equation*}
\begin{split}w_i = \begin{cases} 1, & |r_i| \leq r_0 \\ \frac{r_0}{|r_i|}, & |r_i| > r_0 \end{cases} \text{2.12}\end{split}
\end{equation*}
\sphinxAtStartPar
\sphinxstylestrong{形象理解}：
\begin{itemize}
\item {} 
\sphinxAtStartPar
正常数据：权重 = 1.0（完全参与）

\item {} 
\sphinxAtStartPar
轻微偏离：权重 = 0.8

\item {} 
\sphinxAtStartPar
严重偏离：权重 = 0.1（几乎忽略）

\end{itemize}


\paragraph{有界影响估计}
\label{\detokenize{chapters/chapter2:id54}}
\sphinxAtStartPar
在M估计基础上，进一步处理\sphinxstylestrong{输入端异常}（杠杆点）。这是BIRRP软件的核心算法。


\subsubsection{2.5.6 RMSMR方法（MTDP特色）}
\label{\detokenize{chapters/chapter2:rmsmr-mtdp}}

\paragraph{三维度处理策略}
\label{\detokenize{chapters/chapter2:id55}}
\sphinxAtStartPar
RMSMR = \sphinxstylestrong{R}obust + \sphinxstylestrong{M}ulti\sphinxhyphen{}parameter Screening + \sphinxstylestrong{M}ulti\sphinxhyphen{}angle \sphinxstylestrong{R}hoplus


\begin{savenotes}\sphinxattablestart
\sphinxthistablewithglobalstyle
\centering
\begin{tabulary}{\linewidth}[t]{TTT}
\sphinxtoprule
\begin{varwidth}[t]{\sphinxcolwidth{1}{3}}
\sphinxstyletheadfamily \sphinxAtStartPar
维度
\sphinxbeforeendvarwidth
\end{varwidth}%
&\begin{varwidth}[t]{\sphinxcolwidth{1}{3}}
\sphinxstyletheadfamily \sphinxAtStartPar
技术
\sphinxbeforeendvarwidth
\end{varwidth}%
&\begin{varwidth}[t]{\sphinxcolwidth{1}{3}}
\sphinxstyletheadfamily \sphinxAtStartPar
作用
\sphinxbeforeendvarwidth
\end{varwidth}%
\\
\sphinxmidrule
\sphinxtableatstartofbodyhook\begin{varwidth}[t]{\sphinxcolwidth{1}{3}}
\sphinxAtStartPar
\sphinxstylestrong{抑噪}
\sphinxbeforeendvarwidth
\end{varwidth}%
&\begin{varwidth}[t]{\sphinxcolwidth{1}{3}}
\sphinxAtStartPar
稳健估计
\sphinxbeforeendvarwidth
\end{varwidth}%
&\begin{varwidth}[t]{\sphinxcolwidth{1}{3}}
\sphinxAtStartPar
降低噪声数据的权重
\sphinxbeforeendvarwidth
\end{varwidth}%
\\
\sphinxhline\begin{varwidth}[t]{\sphinxcolwidth{1}{3}}
\sphinxAtStartPar
\sphinxstylestrong{去噪}
\sphinxbeforeendvarwidth
\end{varwidth}%
&\begin{varwidth}[t]{\sphinxcolwidth{1}{3}}
\sphinxAtStartPar
多参数筛选
\sphinxbeforeendvarwidth
\end{varwidth}%
&\begin{varwidth}[t]{\sphinxcolwidth{1}{3}}
\sphinxAtStartPar
剔除受污染的数据段
\sphinxbeforeendvarwidth
\end{varwidth}%
\\
\sphinxhline\begin{varwidth}[t]{\sphinxcolwidth{1}{3}}
\sphinxAtStartPar
\sphinxstylestrong{预测}
\sphinxbeforeendvarwidth
\end{varwidth}%
&\begin{varwidth}[t]{\sphinxcolwidth{1}{3}}
\sphinxAtStartPar
Rhoplus约束
\sphinxbeforeendvarwidth
\end{varwidth}%
&\begin{varwidth}[t]{\sphinxcolwidth{1}{3}}
\sphinxAtStartPar
用1D模型约束响应质量
\sphinxbeforeendvarwidth
\end{varwidth}%
\\
\sphinxbottomrule
\end{tabulary}
\sphinxtableafterendhook\par
\sphinxattableend\end{savenotes}
\begin{quote}

\sphinxAtStartPar
📖 \sphinxstylestrong{参考文献}：王培杰等. 基于稳健估计、数据筛选和Rhoplus约束的大地电磁数据处理方法. 地球物理学报, 2024.
\end{quote}


\subsubsection{2.5.7 MTDP阻抗估计操作}
\label{\detokenize{chapters/chapter2:id56}}
\begin{sphinxuseclass}{details}
\sphinxAtStartPar
\sphinxstylestrong{执行阻抗估计：}
\begin{enumerate}
\sphinxsetlistlabels{\arabic}{enumi}{enumii}{}{.}%
\item {} 
\sphinxAtStartPar
菜单：\sphinxstylestrong{处理 → 阻抗估计}（或快捷键 \sphinxcode{\sphinxupquote{Ctrl+Shift+Z}}）

\item {} 
\sphinxAtStartPar
选择处理方法：
\begin{itemize}
\item {} 
\sphinxAtStartPar
 最小二乘法

\item {} 
\sphinxAtStartPar
 远参考道法（推荐）

\item {} 
\sphinxAtStartPar
 稳健估计

\item {} 
\sphinxAtStartPar
 RMSMR迭代

\end{itemize}

\item {} 
\sphinxAtStartPar
点击 \sphinxstylestrong{执行}

\end{enumerate}

\sphinxAtStartPar
\sphinxstylestrong{查看结果：}
\begin{enumerate}
\sphinxsetlistlabels{\arabic}{enumi}{enumii}{}{.}%
\item {} 
\sphinxAtStartPar
切换到 \sphinxstylestrong{MT响应} 选项卡

\item {} 
\sphinxAtStartPar
选择查看：视电阻率 / 相位 / 阻抗张量

\item {} 
\sphinxAtStartPar
可切换 XY/YX 模式

\end{enumerate}

\end{sphinxuseclass}
\begin{sphinxuseclass}{danger}\begin{enumerate}
\sphinxsetlistlabels{\arabic}{enumi}{enumii}{}{.}%
\item {} 
\sphinxAtStartPar
\sphinxstylestrong{忘记设置远参考}：有参考数据但未在软件中关联

\item {} 
\sphinxAtStartPar
\sphinxstylestrong{方法不匹配}：高噪声数据却用最小二乘法

\item {} 
\sphinxAtStartPar
\sphinxstylestrong{忽略参数检查}：直接使用默认参数，不适配实际数据特点

\end{enumerate}

\end{sphinxuseclass}

\bigskip\hrule\bigskip


\sphinxAtStartPar
⬅️ {\hyperref[\detokenize{chapters/chapter2:id35}]{\sphinxcrossref{上一节：从时间序列到频谱}}} | ➡️ {\hyperref[\detokenize{chapters/chapter2:id57}]{\sphinxcrossref{下一节：数据质量评估}}}


\bigskip\hrule\bigskip



\subsection{2.6 第五步：数据质量评估}
\label{\detokenize{chapters/chapter2:id57}}

\subsubsection{2.6.1 质量评估指标总览}
\label{\detokenize{chapters/chapter2:id58}}

\begin{savenotes}\sphinxattablestart
\sphinxthistablewithglobalstyle
\centering
\begin{tabulary}{\linewidth}[t]{TTTT}
\sphinxtoprule
\begin{varwidth}[t]{\sphinxcolwidth{1}{4}}
\sphinxstyletheadfamily \sphinxAtStartPar
指标
\sphinxbeforeendvarwidth
\end{varwidth}%
&\begin{varwidth}[t]{\sphinxcolwidth{1}{4}}
\sphinxstyletheadfamily \sphinxAtStartPar
英文名称
\sphinxbeforeendvarwidth
\end{varwidth}%
&\begin{varwidth}[t]{\sphinxcolwidth{1}{4}}
\sphinxstyletheadfamily \sphinxAtStartPar
物理意义
\sphinxbeforeendvarwidth
\end{varwidth}%
&\begin{varwidth}[t]{\sphinxcolwidth{1}{4}}
\sphinxstyletheadfamily \sphinxAtStartPar
理想值
\sphinxbeforeendvarwidth
\end{varwidth}%
\\
\sphinxmidrule
\sphinxtableatstartofbodyhook\begin{varwidth}[t]{\sphinxcolwidth{1}{4}}
\sphinxAtStartPar
\sphinxstylestrong{常相干度}
\sphinxbeforeendvarwidth
\end{varwidth}%
&\begin{varwidth}[t]{\sphinxcolwidth{1}{4}}
\sphinxAtStartPar
Ordinary Coherence
\sphinxbeforeendvarwidth
\end{varwidth}%
&\begin{varwidth}[t]{\sphinxcolwidth{1}{4}}
\sphinxAtStartPar
两信号直接相关性
\sphinxbeforeendvarwidth
\end{varwidth}%
&\begin{varwidth}[t]{\sphinxcolwidth{1}{4}}
\sphinxAtStartPar
接近1
\sphinxbeforeendvarwidth
\end{varwidth}%
\\
\sphinxhline\begin{varwidth}[t]{\sphinxcolwidth{1}{4}}
\sphinxAtStartPar
\sphinxstylestrong{重相干度}
\sphinxbeforeendvarwidth
\end{varwidth}%
&\begin{varwidth}[t]{\sphinxcolwidth{1}{4}}
\sphinxAtStartPar
Multiple Coherence
\sphinxbeforeendvarwidth
\end{varwidth}%
&\begin{varwidth}[t]{\sphinxcolwidth{1}{4}}
\sphinxAtStartPar
输出与所有输入的综合相关性
\sphinxbeforeendvarwidth
\end{varwidth}%
&\begin{varwidth}[t]{\sphinxcolwidth{1}{4}}
\sphinxAtStartPar
接近1
\sphinxbeforeendvarwidth
\end{varwidth}%
\\
\sphinxhline\begin{varwidth}[t]{\sphinxcolwidth{1}{4}}
\sphinxAtStartPar
\sphinxstylestrong{偏相干度}
\sphinxbeforeendvarwidth
\end{varwidth}%
&\begin{varwidth}[t]{\sphinxcolwidth{1}{4}}
\sphinxAtStartPar
Partial Coherence
\sphinxbeforeendvarwidth
\end{varwidth}%
&\begin{varwidth}[t]{\sphinxcolwidth{1}{4}}
\sphinxAtStartPar
排除其他输入后某输入的独立贡献
\sphinxbeforeendvarwidth
\end{varwidth}%
&\begin{varwidth}[t]{\sphinxcolwidth{1}{4}}
\sphinxAtStartPar
>0.5
\sphinxbeforeendvarwidth
\end{varwidth}%
\\
\sphinxhline\begin{varwidth}[t]{\sphinxcolwidth{1}{4}}
\sphinxAtStartPar
\sphinxstylestrong{误差棒}
\sphinxbeforeendvarwidth
\end{varwidth}%
&\begin{varwidth}[t]{\sphinxcolwidth{1}{4}}
\sphinxAtStartPar
Error Bar
\sphinxbeforeendvarwidth
\end{varwidth}%
&\begin{varwidth}[t]{\sphinxcolwidth{1}{4}}
\sphinxAtStartPar
估计的不确定性
\sphinxbeforeendvarwidth
\end{varwidth}%
&\begin{varwidth}[t]{\sphinxcolwidth{1}{4}}
\sphinxAtStartPar
越小越好
\sphinxbeforeendvarwidth
\end{varwidth}%
\\
\sphinxbottomrule
\end{tabulary}
\sphinxtableafterendhook\par
\sphinxattableend\end{savenotes}


\subsubsection{2.6.2 三种相干度详解}
\label{\detokenize{chapters/chapter2:id59}}

\paragraph{常相干度 (Ordinary Coherence)}
\label{\detokenize{chapters/chapter2:ordinary-coherence}}
\sphinxAtStartPar
\sphinxstylestrong{定义}：描述两个信号之间的直接相关性（0\sphinxhyphen{}1）
\begin{equation*}
\begin{split}Coh^2(XY) = \frac{|[XY^*]|^2}{[XX^*][YY^*]} \text{2.13}\end{split}
\end{equation*}

\begin{savenotes}\sphinxattablestart
\sphinxthistablewithglobalstyle
\centering
\begin{tabulary}{\linewidth}[t]{TTT}
\sphinxtoprule
\begin{varwidth}[t]{\sphinxcolwidth{1}{3}}
\sphinxstyletheadfamily \sphinxAtStartPar
常相干度值
\sphinxbeforeendvarwidth
\end{varwidth}%
&\begin{varwidth}[t]{\sphinxcolwidth{1}{3}}
\sphinxstyletheadfamily \sphinxAtStartPar
数据质量
\sphinxbeforeendvarwidth
\end{varwidth}%
&\begin{varwidth}[t]{\sphinxcolwidth{1}{3}}
\sphinxstyletheadfamily \sphinxAtStartPar
建议
\sphinxbeforeendvarwidth
\end{varwidth}%
\\
\sphinxmidrule
\sphinxtableatstartofbodyhook\begin{varwidth}[t]{\sphinxcolwidth{1}{3}}
\sphinxAtStartPar
> 0.9
\sphinxbeforeendvarwidth
\end{varwidth}%
&\begin{varwidth}[t]{\sphinxcolwidth{1}{3}}
\sphinxAtStartPar
优秀
\sphinxbeforeendvarwidth
\end{varwidth}%
&\begin{varwidth}[t]{\sphinxcolwidth{1}{3}}
\sphinxAtStartPar
可直接使用
\sphinxbeforeendvarwidth
\end{varwidth}%
\\
\sphinxhline\begin{varwidth}[t]{\sphinxcolwidth{1}{3}}
\sphinxAtStartPar
0.7\sphinxhyphen{}0.9
\sphinxbeforeendvarwidth
\end{varwidth}%
&\begin{varwidth}[t]{\sphinxcolwidth{1}{3}}
\sphinxAtStartPar
良好
\sphinxbeforeendvarwidth
\end{varwidth}%
&\begin{varwidth}[t]{\sphinxcolwidth{1}{3}}
\sphinxAtStartPar
正常使用
\sphinxbeforeendvarwidth
\end{varwidth}%
\\
\sphinxhline\begin{varwidth}[t]{\sphinxcolwidth{1}{3}}
\sphinxAtStartPar
0.5\sphinxhyphen{}0.7
\sphinxbeforeendvarwidth
\end{varwidth}%
&\begin{varwidth}[t]{\sphinxcolwidth{1}{3}}
\sphinxAtStartPar
一般
\sphinxbeforeendvarwidth
\end{varwidth}%
&\begin{varwidth}[t]{\sphinxcolwidth{1}{3}}
\sphinxAtStartPar
考虑筛选
\sphinxbeforeendvarwidth
\end{varwidth}%
\\
\sphinxhline\begin{varwidth}[t]{\sphinxcolwidth{1}{3}}
\sphinxAtStartPar
< 0.5
\sphinxbeforeendvarwidth
\end{varwidth}%
&\begin{varwidth}[t]{\sphinxcolwidth{1}{3}}
\sphinxAtStartPar
较差
\sphinxbeforeendvarwidth
\end{varwidth}%
&\begin{varwidth}[t]{\sphinxcolwidth{1}{3}}
\sphinxAtStartPar
需要剔除或重新采集
\sphinxbeforeendvarwidth
\end{varwidth}%
\\
\sphinxbottomrule
\end{tabulary}
\sphinxtableafterendhook\par
\sphinxattableend\end{savenotes}


\paragraph{重相干度 (Multiple Coherence)}
\label{\detokenize{chapters/chapter2:multiple-coherence}}
\sphinxAtStartPar
\sphinxstylestrong{定义}：输出信号与所有输入信号的综合相关性
\begin{equation*}
\begin{split}Coh_m^2(E_x) = \frac{Z_{xy}[H_yE_x^*] + Z_{xx}[H_xE_x^*]}{[E_xE_x^*]} \text{2.14}\end{split}
\end{equation*}
\sphinxAtStartPar
\sphinxstylestrong{物理意义}：衡量两个水平磁场对电场的\sphinxstylestrong{综合解释能力}
\begin{itemize}
\item {} 
\sphinxAtStartPar
无噪声条件下，重相干度 = 1

\item {} 
\sphinxAtStartPar
噪声越强，重相干度越低

\item {} 
\sphinxAtStartPar
\sphinxstylestrong{同源噪声}也会导致高重相干度（需结合其他指标判断）

\end{itemize}


\paragraph{偏相干度 (Partial Coherence)}
\label{\detokenize{chapters/chapter2:partial-coherence}}
\sphinxAtStartPar
\sphinxstylestrong{定义}：排除其他输入影响后，某输入对输出的\sphinxstylestrong{独立贡献}
\begin{equation*}
\begin{split}Coh_p^2(E_xH_y) = \frac{Coh_m^2(E_x) - Coh^2(E_xH_x)}{1 - Coh^2(E_xH_x)} \text{2.15}\end{split}
\end{equation*}
\sphinxAtStartPar
\sphinxstylestrong{应用场景}：
\begin{itemize}
\item {} 
\sphinxAtStartPar
判断哪个磁场分量贡献更大

\item {} 
\sphinxAtStartPar
识别某个方向的信号是否被噪声污染

\item {} 
\sphinxAtStartPar
多输入系统中定位问题通道

\end{itemize}

\begin{sphinxuseclass}{tip}\begin{enumerate}
\sphinxsetlistlabels{\arabic}{enumi}{enumii}{}{.}%
\item {} 
\sphinxAtStartPar
\sphinxstylestrong{三种相干度结合看}：不能只看一个指标

\item {} 
\sphinxAtStartPar
\sphinxstylestrong{注意同源噪声}：高相干度 ≠ 数据好（可能是噪声相关）

\item {} 
\sphinxAtStartPar
\sphinxstylestrong{分段检查}：不同时段/频段的相干度可能差异很大

\end{enumerate}

\end{sphinxuseclass}

\subsubsection{2.6.3 其他质量指标}
\label{\detokenize{chapters/chapter2:id60}}

\paragraph{功率谱密度 (PSD)}
\label{\detokenize{chapters/chapter2:psd}}
\sphinxAtStartPar
\sphinxstylestrong{识别噪声}：功率谱密度异常高的时段可能包含噪声
\begin{equation*}
\begin{split}PSD = \frac{\sqrt{[AA^*]^2}}{\Delta T} \text{2.16}\end{split}
\end{equation*}
\sphinxAtStartPar
\sphinxstylestrong{判断标准}：
\begin{itemize}
\item {} 
\sphinxAtStartPar
正常天然场：PSD随频率缓慢变化

\item {} 
\sphinxAtStartPar
噪声干扰：PSD突然升高一个数量级

\end{itemize}


\paragraph{误差棒（方差估计）}
\label{\detokenize{chapters/chapter2:id61}}
\sphinxAtStartPar
\sphinxstylestrong{Gamble方差估计}：
\begin{equation*}
\begin{split}(\Delta Z_{ij})^2 = \frac{[|\eta_i|^2][|R_j|^2]}{N|D^R|^2} \text{2.17}\end{split}
\end{equation*}
\sphinxAtStartPar
\sphinxstylestrong{误差棒的意义}：
\begin{itemize}
\item {} 
\sphinxAtStartPar
误差小 → 数据可信度高

\item {} 
\sphinxAtStartPar
误差大 → 数据不确定性大

\item {} 
\sphinxAtStartPar
反演时：误差大的点权重小

\end{itemize}


\subsubsection{2.6.4 数据筛选方法}
\label{\detokenize{chapters/chapter2:id62}}

\paragraph{筛选参数一览}
\label{\detokenize{chapters/chapter2:id63}}

\begin{savenotes}\sphinxattablestart
\sphinxthistablewithglobalstyle
\centering
\begin{tabulary}{\linewidth}[t]{TTT}
\sphinxtoprule
\begin{varwidth}[t]{\sphinxcolwidth{1}{3}}
\sphinxstyletheadfamily \sphinxAtStartPar
参数
\sphinxbeforeendvarwidth
\end{varwidth}%
&\begin{varwidth}[t]{\sphinxcolwidth{1}{3}}
\sphinxstyletheadfamily \sphinxAtStartPar
物理意义
\sphinxbeforeendvarwidth
\end{varwidth}%
&\begin{varwidth}[t]{\sphinxcolwidth{1}{3}}
\sphinxstyletheadfamily \sphinxAtStartPar
筛选标准
\sphinxbeforeendvarwidth
\end{varwidth}%
\\
\sphinxmidrule
\sphinxtableatstartofbodyhook\begin{varwidth}[t]{\sphinxcolwidth{1}{3}}
\sphinxAtStartPar
\sphinxstylestrong{常相干度}
\sphinxbeforeendvarwidth
\end{varwidth}%
&\begin{varwidth}[t]{\sphinxcolwidth{1}{3}}
\sphinxAtStartPar
信号相关性
\sphinxbeforeendvarwidth
\end{varwidth}%
&\begin{varwidth}[t]{\sphinxcolwidth{1}{3}}
\sphinxAtStartPar
> 0.7
\sphinxbeforeendvarwidth
\end{varwidth}%
\\
\sphinxhline\begin{varwidth}[t]{\sphinxcolwidth{1}{3}}
\sphinxAtStartPar
\sphinxstylestrong{重相干度}
\sphinxbeforeendvarwidth
\end{varwidth}%
&\begin{varwidth}[t]{\sphinxcolwidth{1}{3}}
\sphinxAtStartPar
综合解释能力
\sphinxbeforeendvarwidth
\end{varwidth}%
&\begin{varwidth}[t]{\sphinxcolwidth{1}{3}}
\sphinxAtStartPar
> 0.7
\sphinxbeforeendvarwidth
\end{varwidth}%
\\
\sphinxhline\begin{varwidth}[t]{\sphinxcolwidth{1}{3}}
\sphinxAtStartPar
\sphinxstylestrong{功率谱密度}
\sphinxbeforeendvarwidth
\end{varwidth}%
&\begin{varwidth}[t]{\sphinxcolwidth{1}{3}}
\sphinxAtStartPar
信号能量
\sphinxbeforeendvarwidth
\end{varwidth}%
&\begin{varwidth}[t]{\sphinxcolwidth{1}{3}}
\sphinxAtStartPar
无异常峰值
\sphinxbeforeendvarwidth
\end{varwidth}%
\\
\sphinxhline\begin{varwidth}[t]{\sphinxcolwidth{1}{3}}
\sphinxAtStartPar
\sphinxstylestrong{极化方向}
\sphinxbeforeendvarwidth
\end{varwidth}%
&\begin{varwidth}[t]{\sphinxcolwidth{1}{3}}
\sphinxAtStartPar
信号方向
\sphinxbeforeendvarwidth
\end{varwidth}%
&\begin{varwidth}[t]{\sphinxcolwidth{1}{3}}
\sphinxAtStartPar
分布随机
\sphinxbeforeendvarwidth
\end{varwidth}%
\\
\sphinxhline\begin{varwidth}[t]{\sphinxcolwidth{1}{3}}
\sphinxAtStartPar
\sphinxstylestrong{残差}
\sphinxbeforeendvarwidth
\end{varwidth}%
&\begin{varwidth}[t]{\sphinxcolwidth{1}{3}}
\sphinxAtStartPar
拟合误差
\sphinxbeforeendvarwidth
\end{varwidth}%
&\begin{varwidth}[t]{\sphinxcolwidth{1}{3}}
\sphinxAtStartPar
< 2σ
\sphinxbeforeendvarwidth
\end{varwidth}%
\\
\sphinxbottomrule
\end{tabulary}
\sphinxtableafterendhook\par
\sphinxattableend\end{savenotes}


\paragraph{马氏距离筛选}
\label{\detokenize{chapters/chapter2:id64}}
\sphinxAtStartPar
\sphinxstylestrong{为什么用马氏距离}：综合考虑多个参数的相关性
\begin{equation*}
\begin{split}D_M = \sqrt{(\mathbf{x} - \boldsymbol{\mu})^T \mathbf{\Sigma}^{-1} (\mathbf{x} - \boldsymbol{\mu})} \text{2.18}\end{split}
\end{equation*}
\sphinxAtStartPar
\sphinxstylestrong{类比}：
\begin{itemize}
\item {} 
\sphinxAtStartPar
欧氏距离：只看"价格"一个维度

\item {} 
\sphinxAtStartPar
马氏距离：综合考虑"价格、品牌、重量"等多个维度的相关性

\end{itemize}


\subsubsection{2.6.5 好/坏数据对比示例}
\label{\detokenize{chapters/chapter2:id65}}

\paragraph{✅ 高质量数据特征}
\label{\detokenize{chapters/chapter2:id66}}
\begin{sphinxVerbatim}[commandchars=\\\{\}]
┌─────────────────────────────────────────────────────────┐
│                    高质量MT响应曲线                       │
├─────────────────────────────────────────────────────────┤
│                                                         │
│  视电阻率 (log scale)                                    │
│  ┌────────────────────────────────────────────────┐    │
│  │     ○ ○                                        │    │
│  │   ○     ○ ○                                    │    │
│  │  ○         ○ ○                                 │    │
│  │ ○             ○ ○ ○                            │    │
│  │○                  ○ ○ ○                        │    │
│  │                        ○ ○ ○                    │    │
│  └────────────────────────────────────────────────┘    │
│  特征：曲线光滑连续，误差棒小，相干度\PYGZgt{}0.8                  │
│                                                         │
│  相位 (度)                                              │
│  ┌────────────────────────────────────────────────┐    │
│  │            ○ ○ ○ ○ ○ ○ ○                       │    │
│  │         ○              ○                       │    │
│  │       ○                  ○                     │    │
│  │     ○                     ○                    │    │
│  │   ○                        ○                   │    │
│  └────────────────────────────────────────────────┘    │
│  特征：相位在30°\PYGZhy{}60°范围，曲线平滑                        │
└─────────────────────────────────────────────────────────┘
\end{sphinxVerbatim}


\paragraph{❌ 低质量数据特征}
\label{\detokenize{chapters/chapter2:id67}}
\begin{sphinxVerbatim}[commandchars=\\\{\}]
┌─────────────────────────────────────────────────────────┐
│                    低质量MT响应曲线                       │
├─────────────────────────────────────────────────────────┤
│                                                         │
│  视电阻率 (log scale)                                    │
│  ┌────────────────────────────────────────────────┐    │
│  │       ○                                        │    │
│  │     ○   ○     ○                                │    │
│  │   ○       ○     ○                              │    │
│  │ ○           ○ ○   ○  ○                         │    │
│  │                ○    ○                          │    │
│  │                     ○ ○                        │    │
│  └────────────────────────────────────────────────┘    │
│  问题：曲线脱节、跳动，误差棒大，相干度\PYGZlt{}0.5                 │
│        ↓ 死频带区域数据质量差                             │
│                                                         │
│  相位 (度)                                              │
│  ┌────────────────────────────────────────────────┐    │
│  │     ○           ○                              │    │
│  │   ○   ○       ○   ○                            │    │
│  │          ○ ○        ○                          │    │
│  │ ○                    ○                         │    │
│  │                        ○                       │    │
│  └────────────────────────────────────────────────┘    │
│  问题：相位异常（\PYGZgt{}90°或\PYGZlt{}0°），数据跳动                     │
└─────────────────────────────────────────────────────────┘
\end{sphinxVerbatim}


\paragraph{质量问题诊断表}
\label{\detokenize{chapters/chapter2:id68}}

\begin{savenotes}\sphinxattablestart
\sphinxthistablewithglobalstyle
\centering
\begin{tabulary}{\linewidth}[t]{TTTT}
\sphinxtoprule
\begin{varwidth}[t]{\sphinxcolwidth{1}{4}}
\sphinxstyletheadfamily \sphinxAtStartPar
问题现象
\sphinxbeforeendvarwidth
\end{varwidth}%
&\begin{varwidth}[t]{\sphinxcolwidth{1}{4}}
\sphinxstyletheadfamily \sphinxAtStartPar
可能原因
\sphinxbeforeendvarwidth
\end{varwidth}%
&\begin{varwidth}[t]{\sphinxcolwidth{1}{4}}
\sphinxstyletheadfamily \sphinxAtStartPar
诊断指标
\sphinxbeforeendvarwidth
\end{varwidth}%
&\begin{varwidth}[t]{\sphinxcolwidth{1}{4}}
\sphinxstyletheadfamily \sphinxAtStartPar
解决方案
\sphinxbeforeendvarwidth
\end{varwidth}%
\\
\sphinxmidrule
\sphinxtableatstartofbodyhook\begin{varwidth}[t]{\sphinxcolwidth{1}{4}}
\sphinxAtStartPar
\sphinxstylestrong{曲线脱节}
\sphinxbeforeendvarwidth
\end{varwidth}%
&\begin{varwidth}[t]{\sphinxcolwidth{1}{4}}
\sphinxAtStartPar
死频带噪声
\sphinxbeforeendvarwidth
\end{varwidth}%
&\begin{varwidth}[t]{\sphinxcolwidth{1}{4}}
\sphinxAtStartPar
相干度<0.5
\sphinxbeforeendvarwidth
\end{varwidth}%
&\begin{varwidth}[t]{\sphinxcolwidth{1}{4}}
\sphinxAtStartPar
Rhoplus校正
\sphinxbeforeendvarwidth
\end{varwidth}%
\\
\sphinxhline\begin{varwidth}[t]{\sphinxcolwidth{1}{4}}
\sphinxAtStartPar
\sphinxstylestrong{相位异常}
\sphinxbeforeendvarwidth
\end{varwidth}%
&\begin{varwidth}[t]{\sphinxcolwidth{1}{4}}
\sphinxAtStartPar
近场干扰
\sphinxbeforeendvarwidth
\end{varwidth}%
&\begin{varwidth}[t]{\sphinxcolwidth{1}{4}}
\sphinxAtStartPar
相位>90°
\sphinxbeforeendvarwidth
\end{varwidth}%
&\begin{varwidth}[t]{\sphinxcolwidth{1}{4}}
\sphinxAtStartPar
剔除相关频点
\sphinxbeforeendvarwidth
\end{varwidth}%
\\
\sphinxhline\begin{varwidth}[t]{\sphinxcolwidth{1}{4}}
\sphinxAtStartPar
\sphinxstylestrong{误差棒大}
\sphinxbeforeendvarwidth
\end{varwidth}%
&\begin{varwidth}[t]{\sphinxcolwidth{1}{4}}
\sphinxAtStartPar
叠加不足
\sphinxbeforeendvarwidth
\end{varwidth}%
&\begin{varwidth}[t]{\sphinxcolwidth{1}{4}}
\sphinxAtStartPar
误差>30\%
\sphinxbeforeendvarwidth
\end{varwidth}%
&\begin{varwidth}[t]{\sphinxcolwidth{1}{4}}
\sphinxAtStartPar
增加数据量
\sphinxbeforeendvarwidth
\end{varwidth}%
\\
\sphinxhline\begin{varwidth}[t]{\sphinxcolwidth{1}{4}}
\sphinxAtStartPar
\sphinxstylestrong{低相干度}
\sphinxbeforeendvarwidth
\end{varwidth}%
&\begin{varwidth}[t]{\sphinxcolwidth{1}{4}}
\sphinxAtStartPar
人文噪声
\sphinxbeforeendvarwidth
\end{varwidth}%
&\begin{varwidth}[t]{\sphinxcolwidth{1}{4}}
\sphinxAtStartPar
Coh<0.5
\sphinxbeforeendvarwidth
\end{varwidth}%
&\begin{varwidth}[t]{\sphinxcolwidth{1}{4}}
\sphinxAtStartPar
人工筛选/增加时长
\sphinxbeforeendvarwidth
\end{varwidth}%
\\
\sphinxhline\begin{varwidth}[t]{\sphinxcolwidth{1}{4}}
\sphinxAtStartPar
\sphinxstylestrong{曲线跳动}
\sphinxbeforeendvarwidth
\end{varwidth}%
&\begin{varwidth}[t]{\sphinxcolwidth{1}{4}}
\sphinxAtStartPar
强干扰
\sphinxbeforeendvarwidth
\end{varwidth}%
&\begin{varwidth}[t]{\sphinxcolwidth{1}{4}}
\sphinxAtStartPar
PSD异常
\sphinxbeforeendvarwidth
\end{varwidth}%
&\begin{varwidth}[t]{\sphinxcolwidth{1}{4}}
\sphinxAtStartPar
RMSMR处理
\sphinxbeforeendvarwidth
\end{varwidth}%
\\
\sphinxbottomrule
\end{tabulary}
\sphinxtableafterendhook\par
\sphinxattableend\end{savenotes}


\subsubsection{2.6.6 质量评估清单}
\label{\detokenize{chapters/chapter2:id69}}

\subsubsection{2.6.7 MTDP质量评估操作}
\label{\detokenize{chapters/chapter2:id70}}
\begin{sphinxuseclass}{details}
\sphinxAtStartPar
\sphinxstylestrong{查看相干度：}
\begin{enumerate}
\sphinxsetlistlabels{\arabic}{enumi}{enumii}{}{.}%
\item {} 
\sphinxAtStartPar
切换到 \sphinxstylestrong{数据筛选} 选项卡

\item {} 
\sphinxAtStartPar
勾选要显示的相干度类型：
\begin{itemize}
\item {} 
\sphinxAtStartPar
 常相干度 (Coh²)

\item {} 
\sphinxAtStartPar
 重相干度 (Cohₘ²)

\item {} 
\sphinxAtStartPar
 偏相干度 (Cohₚ²)

\end{itemize}

\item {} 
\sphinxAtStartPar
设置阈值线（推荐 0.7）

\end{enumerate}

\sphinxAtStartPar
\sphinxstylestrong{手动筛选数据：}
\begin{enumerate}
\sphinxsetlistlabels{\arabic}{enumi}{enumii}{}{.}%
\item {} 
\sphinxAtStartPar
在功率谱窗口框选异常区域

\item {} 
\sphinxAtStartPar
右键 → \sphinxstylestrong{剔除选中数据}

\item {} 
\sphinxAtStartPar
重新计算阻抗

\end{enumerate}

\sphinxAtStartPar
\sphinxstylestrong{自动筛选：}
\begin{enumerate}
\sphinxsetlistlabels{\arabic}{enumi}{enumii}{}{.}%
\item {} 
\sphinxAtStartPar
菜单：\sphinxstylestrong{处理 → 自动筛选}

\item {} 
\sphinxAtStartPar
设置筛选参数阈值

\item {} 
\sphinxAtStartPar
点击 \sphinxstylestrong{执行}

\end{enumerate}

\end{sphinxuseclass}

\bigskip\hrule\bigskip


\sphinxAtStartPar
⬅️ {\hyperref[\detokenize{chapters/chapter2:id43}]{\sphinxcrossref{上一节：阻抗估计}}} | ➡️ {\hyperref[\detokenize{chapters/chapter2:id71}]{\sphinxcrossref{下一节：高级分析工具}}}


\bigskip\hrule\bigskip



\subsection{2.7 第六步：高级分析工具}
\label{\detokenize{chapters/chapter2:id71}}

\subsubsection{2.7.1 相位张量分析}
\label{\detokenize{chapters/chapter2:id72}}

\paragraph{为什么用相位张量？}
\label{\detokenize{chapters/chapter2:id73}}
\sphinxAtStartPar
\sphinxstylestrong{传统方法的局限}：需要先假设地下是一维、二维还是三维。

\sphinxAtStartPar
\sphinxstylestrong{相位张量的优势}：\sphinxstylestrong{不需要任何维性假设}！


\paragraph{相位张量定义}
\label{\detokenize{chapters/chapter2:id74}}
\sphinxAtStartPar
将阻抗分解为实部和虚部：\(\mathbf{Z} = \mathbf{X} + i\mathbf{Y}\)

\sphinxAtStartPar
相位张量：
\begin{equation*}
\begin{split}\mathbf{\Phi} = \mathbf{X}^{-1}\mathbf{Y} \text{2.19}\end{split}
\end{equation*}

\paragraph{椭圆表示}
\label{\detokenize{chapters/chapter2:id75}}
\sphinxAtStartPar
相位张量可以用椭圆可视化：


\begin{savenotes}\sphinxattablestart
\sphinxthistablewithglobalstyle
\centering
\begin{tabulary}{\linewidth}[t]{TTT}
\sphinxtoprule
\begin{varwidth}[t]{\sphinxcolwidth{1}{3}}
\sphinxstyletheadfamily \sphinxAtStartPar
结构类型
\sphinxbeforeendvarwidth
\end{varwidth}%
&\begin{varwidth}[t]{\sphinxcolwidth{1}{3}}
\sphinxstyletheadfamily \sphinxAtStartPar
椭圆形状
\sphinxbeforeendvarwidth
\end{varwidth}%
&\begin{varwidth}[t]{\sphinxcolwidth{1}{3}}
\sphinxstyletheadfamily \sphinxAtStartPar
特征
\sphinxbeforeendvarwidth
\end{varwidth}%
\\
\sphinxmidrule
\sphinxtableatstartofbodyhook\begin{varwidth}[t]{\sphinxcolwidth{1}{3}}
\sphinxAtStartPar
\sphinxstylestrong{一维}
\sphinxbeforeendvarwidth
\end{varwidth}%
&\begin{varwidth}[t]{\sphinxcolwidth{1}{3}}
\sphinxAtStartPar
圆形
\sphinxbeforeendvarwidth
\end{varwidth}%
&\begin{varwidth}[t]{\sphinxcolwidth{1}{3}}
\sphinxAtStartPar
\(\Phi_{max} = \Phi_{min}\)
\sphinxbeforeendvarwidth
\end{varwidth}%
\\
\sphinxhline\begin{varwidth}[t]{\sphinxcolwidth{1}{3}}
\sphinxAtStartPar
\sphinxstylestrong{二维}
\sphinxbeforeendvarwidth
\end{varwidth}%
&\begin{varwidth}[t]{\sphinxcolwidth{1}{3}}
\sphinxAtStartPar
椭圆
\sphinxbeforeendvarwidth
\end{varwidth}%
&\begin{varwidth}[t]{\sphinxcolwidth{1}{3}}
\sphinxAtStartPar
扁平程度反映各向异性
\sphinxbeforeendvarwidth
\end{varwidth}%
\\
\sphinxhline\begin{varwidth}[t]{\sphinxcolwidth{1}{3}}
\sphinxAtStartPar
\sphinxstylestrong{三维}
\sphinxbeforeendvarwidth
\end{varwidth}%
&\begin{varwidth}[t]{\sphinxcolwidth{1}{3}}
\sphinxAtStartPar
不规则椭圆
\sphinxbeforeendvarwidth
\end{varwidth}%
&\begin{varwidth}[t]{\sphinxcolwidth{1}{3}}
\sphinxAtStartPar
椭圆倾斜
\sphinxbeforeendvarwidth
\end{varwidth}%
\\
\sphinxbottomrule
\end{tabulary}
\sphinxtableafterendhook\par
\sphinxattableend\end{savenotes}


\paragraph{关键参数}
\label{\detokenize{chapters/chapter2:id76}}

\begin{savenotes}\sphinxattablestart
\sphinxthistablewithglobalstyle
\centering
\begin{tabulary}{\linewidth}[t]{TTT}
\sphinxtoprule
\begin{varwidth}[t]{\sphinxcolwidth{1}{3}}
\sphinxstyletheadfamily \sphinxAtStartPar
参数
\sphinxbeforeendvarwidth
\end{varwidth}%
&\begin{varwidth}[t]{\sphinxcolwidth{1}{3}}
\sphinxstyletheadfamily \sphinxAtStartPar
定义
\sphinxbeforeendvarwidth
\end{varwidth}%
&\begin{varwidth}[t]{\sphinxcolwidth{1}{3}}
\sphinxstyletheadfamily \sphinxAtStartPar
判断标准
\sphinxbeforeendvarwidth
\end{varwidth}%
\\
\sphinxmidrule
\sphinxtableatstartofbodyhook\begin{varwidth}[t]{\sphinxcolwidth{1}{3}}
\sphinxAtStartPar
\sphinxstylestrong{主方向角 α}
\sphinxbeforeendvarwidth
\end{varwidth}%
&\begin{varwidth}[t]{\sphinxcolwidth{1}{3}}
\sphinxAtStartPar
电性主轴方向
\sphinxbeforeendvarwidth
\end{varwidth}%
&\begin{varwidth}[t]{\sphinxcolwidth{1}{3}}
\sphinxAtStartPar
椭圆长轴方向
\sphinxbeforeendvarwidth
\end{varwidth}%
\\
\sphinxhline\begin{varwidth}[t]{\sphinxcolwidth{1}{3}}
\sphinxAtStartPar
\sphinxstylestrong{二维偏离度 β}
\sphinxbeforeendvarwidth
\end{varwidth}%
&\begin{varwidth}[t]{\sphinxcolwidth{1}{3}}
\sphinxAtStartPar
偏离二维程度
\sphinxbeforeendvarwidth
\end{varwidth}%
&\begin{varwidth}[t]{\sphinxcolwidth{1}{3}}
\sphinxAtStartPar
\$
\sphinxbeforeendvarwidth
\end{varwidth}%
\\
\sphinxhline\begin{varwidth}[t]{\sphinxcolwidth{1}{3}}
\sphinxAtStartPar
\sphinxstylestrong{椭圆率 λ}
\sphinxbeforeendvarwidth
\end{varwidth}%
&\begin{varwidth}[t]{\sphinxcolwidth{1}{3}}
\sphinxAtStartPar
各向异性程度
\sphinxbeforeendvarwidth
\end{varwidth}%
&\begin{varwidth}[t]{\sphinxcolwidth{1}{3}}
\sphinxAtStartPar
λ=0为圆形
\sphinxbeforeendvarwidth
\end{varwidth}%
\\
\sphinxhline\begin{varwidth}[t]{\sphinxcolwidth{1}{3}}
\sphinxAtStartPar
\sphinxstylestrong{一维偏离度 κ}
\sphinxbeforeendvarwidth
\end{varwidth}%
&\begin{varwidth}[t]{\sphinxcolwidth{1}{3}}
\sphinxAtStartPar
偏离一维程度
\sphinxbeforeendvarwidth
\end{varwidth}%
&\begin{varwidth}[t]{\sphinxcolwidth{1}{3}}
\sphinxAtStartPar
κ=0为一维
\sphinxbeforeendvarwidth
\end{varwidth}%
\\
\sphinxbottomrule
\end{tabulary}
\sphinxtableafterendhook\par
\sphinxattableend\end{savenotes}


\subsubsection{2.7.2 Rhoplus分析}
\label{\detokenize{chapters/chapter2:rhoplus}}

\paragraph{核心思想}
\label{\detokenize{chapters/chapter2:id77}}
\sphinxAtStartPar
\sphinxstylestrong{奥卡姆剃刀原则}：寻找能解释观测数据的\sphinxstylestrong{最简单模型}

\sphinxAtStartPar
Rhoplus寻找与观测数据相容的\sphinxstylestrong{最少层数}一维模型。


\paragraph{目标函数}
\label{\detokenize{chapters/chapter2:id78}}\begin{equation*}
\begin{split}\min \sum_{i=1}^{N} \left[\frac{\rho_a^{obs}(\omega_i) - \rho_a^{cal}(\omega_i)}{\delta\rho_a(\omega_i)}\right]^2 \text{2.20}\end{split}
\end{equation*}

\paragraph{应用场景}
\label{\detokenize{chapters/chapter2:id79}}

\begin{savenotes}\sphinxattablestart
\sphinxthistablewithglobalstyle
\centering
\begin{tabulary}{\linewidth}[t]{TT}
\sphinxtoprule
\begin{varwidth}[t]{\sphinxcolwidth{1}{2}}
\sphinxstyletheadfamily \sphinxAtStartPar
应用
\sphinxbeforeendvarwidth
\end{varwidth}%
&\begin{varwidth}[t]{\sphinxcolwidth{1}{2}}
\sphinxstyletheadfamily \sphinxAtStartPar
说明
\sphinxbeforeendvarwidth
\end{varwidth}%
\\
\sphinxmidrule
\sphinxtableatstartofbodyhook\begin{varwidth}[t]{\sphinxcolwidth{1}{2}}
\sphinxAtStartPar
\sphinxstylestrong{数据质量检验}
\sphinxbeforeendvarwidth
\end{varwidth}%
&\begin{varwidth}[t]{\sphinxcolwidth{1}{2}}
\sphinxAtStartPar
数据能被简单模型拟合 → 数据可靠
\sphinxbeforeendvarwidth
\end{varwidth}%
\\
\sphinxhline\begin{varwidth}[t]{\sphinxcolwidth{1}{2}}
\sphinxAtStartPar
\sphinxstylestrong{维性分析}
\sphinxbeforeendvarwidth
\end{varwidth}%
&\begin{varwidth}[t]{\sphinxcolwidth{1}{2}}
\sphinxAtStartPar
不同方向拟合结果不同 → 结构复杂
\sphinxbeforeendvarwidth
\end{varwidth}%
\\
\sphinxhline\begin{varwidth}[t]{\sphinxcolwidth{1}{2}}
\sphinxAtStartPar
\sphinxstylestrong{静态位移检测}
\sphinxbeforeendvarwidth
\end{varwidth}%
&\begin{varwidth}[t]{\sphinxcolwidth{1}{2}}
\sphinxAtStartPar
模型预测与实测偏差
\sphinxbeforeendvarwidth
\end{varwidth}%
\\
\sphinxhline\begin{varwidth}[t]{\sphinxcolwidth{1}{2}}
\sphinxAtStartPar
\sphinxstylestrong{初始模型}
\sphinxbeforeendvarwidth
\end{varwidth}%
&\begin{varwidth}[t]{\sphinxcolwidth{1}{2}}
\sphinxAtStartPar
为复杂反演提供起点
\sphinxbeforeendvarwidth
\end{varwidth}%
\\
\sphinxbottomrule
\end{tabulary}
\sphinxtableafterendhook\par
\sphinxattableend\end{savenotes}


\paragraph{多角度Rhoplus}
\label{\detokenize{chapters/chapter2:id80}}
\sphinxAtStartPar
对于二维/三维结构，不同方向的视电阻率曲线不同。通过旋转计算不同角度的Rhoplus响应，可以判断维性特征。


\subsubsection{2.7.3 MTDP高级分析操作}
\label{\detokenize{chapters/chapter2:id81}}
\begin{sphinxuseclass}{details}
\sphinxAtStartPar
\sphinxstylestrong{相位张量分析：}
\begin{enumerate}
\sphinxsetlistlabels{\arabic}{enumi}{enumii}{}{.}%
\item {} 
\sphinxAtStartPar
切换到 \sphinxstylestrong{高级分析} 选项卡

\item {} 
\sphinxAtStartPar
选择 \sphinxstylestrong{相位张量} 子选项卡

\item {} 
\sphinxAtStartPar
查看椭圆图和参数曲线

\item {} 
\sphinxAtStartPar
可导出相位张量数据

\end{enumerate}

\sphinxAtStartPar
\sphinxstylestrong{Rhoplus分析：}
\begin{enumerate}
\sphinxsetlistlabels{\arabic}{enumi}{enumii}{}{.}%
\item {} 
\sphinxAtStartPar
切换到 \sphinxstylestrong{Rhoplus} 子选项卡

\item {} 
\sphinxAtStartPar
设置角度范围（默认 0°\sphinxhyphen{}180°，步长 10°）

\item {} 
\sphinxAtStartPar
点击 \sphinxstylestrong{计算}

\item {} 
\sphinxAtStartPar
查看不同角度的拟合结果

\end{enumerate}

\end{sphinxuseclass}

\bigskip\hrule\bigskip


\sphinxAtStartPar
⬅️ {\hyperref[\detokenize{chapters/chapter2:id57}]{\sphinxcrossref{上一节：数据质量评估}}} | ➡️ {\hyperref[\detokenize{chapters/chapter2:rmsmr}]{\sphinxcrossref{下一节：RMSMR完整流程}}}


\bigskip\hrule\bigskip



\subsection{2.8 第七步：RMSMR完整流程}
\label{\detokenize{chapters/chapter2:rmsmr}}

\subsubsection{2.8.1 RMSMR处理流程图}
\label{\detokenize{chapters/chapter2:id82}}

\subsubsection{2.8.2 参数设置指南}
\label{\detokenize{chapters/chapter2:id83}}

\paragraph{核心参数详解}
\label{\detokenize{chapters/chapter2:id84}}

\begin{savenotes}\sphinxattablestart
\sphinxthistablewithglobalstyle
\centering
\begin{tabulary}{\linewidth}[t]{TTTTT}
\sphinxtoprule
\begin{varwidth}[t]{\sphinxcolwidth{1}{5}}
\sphinxstyletheadfamily \sphinxAtStartPar
参数
\sphinxbeforeendvarwidth
\end{varwidth}%
&\begin{varwidth}[t]{\sphinxcolwidth{1}{5}}
\sphinxstyletheadfamily \sphinxAtStartPar
含义
\sphinxbeforeendvarwidth
\end{varwidth}%
&\begin{varwidth}[t]{\sphinxcolwidth{1}{5}}
\sphinxstyletheadfamily \sphinxAtStartPar
保守设置
\sphinxbeforeendvarwidth
\end{varwidth}%
&\begin{varwidth}[t]{\sphinxcolwidth{1}{5}}
\sphinxstyletheadfamily \sphinxAtStartPar
激进设置
\sphinxbeforeendvarwidth
\end{varwidth}%
&\begin{varwidth}[t]{\sphinxcolwidth{1}{5}}
\sphinxstyletheadfamily \sphinxAtStartPar
影响
\sphinxbeforeendvarwidth
\end{varwidth}%
\\
\sphinxmidrule
\sphinxtableatstartofbodyhook\begin{varwidth}[t]{\sphinxcolwidth{1}{5}}
\sphinxAtStartPar
\sphinxstylestrong{常相干度阈值}
\sphinxbeforeendvarwidth
\end{varwidth}%
&\begin{varwidth}[t]{\sphinxcolwidth{1}{5}}
\sphinxAtStartPar
相干度低于此值的数据被剔除
\sphinxbeforeendvarwidth
\end{varwidth}%
&\begin{varwidth}[t]{\sphinxcolwidth{1}{5}}
\sphinxAtStartPar
0.8
\sphinxbeforeendvarwidth
\end{varwidth}%
&\begin{varwidth}[t]{\sphinxcolwidth{1}{5}}
\sphinxAtStartPar
0.5
\sphinxbeforeendvarwidth
\end{varwidth}%
&\begin{varwidth}[t]{\sphinxcolwidth{1}{5}}
\sphinxAtStartPar
保守→数据少但质量高
\sphinxbeforeendvarwidth
\end{varwidth}%
\\
\sphinxhline\begin{varwidth}[t]{\sphinxcolwidth{1}{5}}
\sphinxAtStartPar
\sphinxstylestrong{重相干度阈值}
\sphinxbeforeendvarwidth
\end{varwidth}%
&\begin{varwidth}[t]{\sphinxcolwidth{1}{5}}
\sphinxAtStartPar
重相干度低于此值的数据被剔除
\sphinxbeforeendvarwidth
\end{varwidth}%
&\begin{varwidth}[t]{\sphinxcolwidth{1}{5}}
\sphinxAtStartPar
0.8
\sphinxbeforeendvarwidth
\end{varwidth}%
&\begin{varwidth}[t]{\sphinxcolwidth{1}{5}}
\sphinxAtStartPar
0.6
\sphinxbeforeendvarwidth
\end{varwidth}%
&\begin{varwidth}[t]{\sphinxcolwidth{1}{5}}
\sphinxAtStartPar
保守→更严格筛选
\sphinxbeforeendvarwidth
\end{varwidth}%
\\
\sphinxhline\begin{varwidth}[t]{\sphinxcolwidth{1}{5}}
\sphinxAtStartPar
\sphinxstylestrong{残差阈值}
\sphinxbeforeendvarwidth
\end{varwidth}%
&\begin{varwidth}[t]{\sphinxcolwidth{1}{5}}
\sphinxAtStartPar
残差超过此值的数据被降权/剔除
\sphinxbeforeendvarwidth
\end{varwidth}%
&\begin{varwidth}[t]{\sphinxcolwidth{1}{5}}
\sphinxAtStartPar
1.5σ
\sphinxbeforeendvarwidth
\end{varwidth}%
&\begin{varwidth}[t]{\sphinxcolwidth{1}{5}}
\sphinxAtStartPar
3.0σ
\sphinxbeforeendvarwidth
\end{varwidth}%
&\begin{varwidth}[t]{\sphinxcolwidth{1}{5}}
\sphinxAtStartPar
保守→剔除更多数据
\sphinxbeforeendvarwidth
\end{varwidth}%
\\
\sphinxhline\begin{varwidth}[t]{\sphinxcolwidth{1}{5}}
\sphinxAtStartPar
\sphinxstylestrong{马氏距离阈值}
\sphinxbeforeendvarwidth
\end{varwidth}%
&\begin{varwidth}[t]{\sphinxcolwidth{1}{5}}
\sphinxAtStartPar
马氏距离超过此值视为异常
\sphinxbeforeendvarwidth
\end{varwidth}%
&\begin{varwidth}[t]{\sphinxcolwidth{1}{5}}
\sphinxAtStartPar
χ²(0.90)
\sphinxbeforeendvarwidth
\end{varwidth}%
&\begin{varwidth}[t]{\sphinxcolwidth{1}{5}}
\sphinxAtStartPar
χ²(0.99)
\sphinxbeforeendvarwidth
\end{varwidth}%
&\begin{varwidth}[t]{\sphinxcolwidth{1}{5}}
\sphinxAtStartPar
保守→更严格异常检测
\sphinxbeforeendvarwidth
\end{varwidth}%
\\
\sphinxhline\begin{varwidth}[t]{\sphinxcolwidth{1}{5}}
\sphinxAtStartPar
\sphinxstylestrong{最大迭代次数}
\sphinxbeforeendvarwidth
\end{varwidth}%
&\begin{varwidth}[t]{\sphinxcolwidth{1}{5}}
\sphinxAtStartPar
迭代处理的最大次数
\sphinxbeforeendvarwidth
\end{varwidth}%
&\begin{varwidth}[t]{\sphinxcolwidth{1}{5}}
\sphinxAtStartPar
5
\sphinxbeforeendvarwidth
\end{varwidth}%
&\begin{varwidth}[t]{\sphinxcolwidth{1}{5}}
\sphinxAtStartPar
15
\sphinxbeforeendvarwidth
\end{varwidth}%
&\begin{varwidth}[t]{\sphinxcolwidth{1}{5}}
\sphinxAtStartPar
激进→可能过度筛选
\sphinxbeforeendvarwidth
\end{varwidth}%
\\
\sphinxbottomrule
\end{tabulary}
\sphinxtableafterendhook\par
\sphinxattableend\end{savenotes}


\paragraph{参数调优建议}
\label{\detokenize{chapters/chapter2:id85}}

\paragraph{不同场景的推荐参数}
\label{\detokenize{chapters/chapter2:id86}}

\begin{savenotes}\sphinxattablestart
\sphinxthistablewithglobalstyle
\centering
\begin{tabulary}{\linewidth}[t]{TTTTT}
\sphinxtoprule
\begin{varwidth}[t]{\sphinxcolwidth{1}{5}}
\sphinxstyletheadfamily \sphinxAtStartPar
场景
\sphinxbeforeendvarwidth
\end{varwidth}%
&\begin{varwidth}[t]{\sphinxcolwidth{1}{5}}
\sphinxstyletheadfamily \sphinxAtStartPar
常相干度
\sphinxbeforeendvarwidth
\end{varwidth}%
&\begin{varwidth}[t]{\sphinxcolwidth{1}{5}}
\sphinxstyletheadfamily \sphinxAtStartPar
重相干度
\sphinxbeforeendvarwidth
\end{varwidth}%
&\begin{varwidth}[t]{\sphinxcolwidth{1}{5}}
\sphinxstyletheadfamily \sphinxAtStartPar
残差阈值
\sphinxbeforeendvarwidth
\end{varwidth}%
&\begin{varwidth}[t]{\sphinxcolwidth{1}{5}}
\sphinxstyletheadfamily \sphinxAtStartPar
迭代次数
\sphinxbeforeendvarwidth
\end{varwidth}%
\\
\sphinxmidrule
\sphinxtableatstartofbodyhook\begin{varwidth}[t]{\sphinxcolwidth{1}{5}}
\sphinxAtStartPar
\sphinxstylestrong{城市郊区}
\sphinxbeforeendvarwidth
\end{varwidth}%
&\begin{varwidth}[t]{\sphinxcolwidth{1}{5}}
\sphinxAtStartPar
0.8
\sphinxbeforeendvarwidth
\end{varwidth}%
&\begin{varwidth}[t]{\sphinxcolwidth{1}{5}}
\sphinxAtStartPar
0.8
\sphinxbeforeendvarwidth
\end{varwidth}%
&\begin{varwidth}[t]{\sphinxcolwidth{1}{5}}
\sphinxAtStartPar
1.5σ
\sphinxbeforeendvarwidth
\end{varwidth}%
&\begin{varwidth}[t]{\sphinxcolwidth{1}{5}}
\sphinxAtStartPar
10
\sphinxbeforeendvarwidth
\end{varwidth}%
\\
\sphinxhline\begin{varwidth}[t]{\sphinxcolwidth{1}{5}}
\sphinxAtStartPar
\sphinxstylestrong{农村地区}
\sphinxbeforeendvarwidth
\end{varwidth}%
&\begin{varwidth}[t]{\sphinxcolwidth{1}{5}}
\sphinxAtStartPar
0.7
\sphinxbeforeendvarwidth
\end{varwidth}%
&\begin{varwidth}[t]{\sphinxcolwidth{1}{5}}
\sphinxAtStartPar
0.7
\sphinxbeforeendvarwidth
\end{varwidth}%
&\begin{varwidth}[t]{\sphinxcolwidth{1}{5}}
\sphinxAtStartPar
2.0σ
\sphinxbeforeendvarwidth
\end{varwidth}%
&\begin{varwidth}[t]{\sphinxcolwidth{1}{5}}
\sphinxAtStartPar
5
\sphinxbeforeendvarwidth
\end{varwidth}%
\\
\sphinxhline\begin{varwidth}[t]{\sphinxcolwidth{1}{5}}
\sphinxAtStartPar
\sphinxstylestrong{矿区附近}
\sphinxbeforeendvarwidth
\end{varwidth}%
&\begin{varwidth}[t]{\sphinxcolwidth{1}{5}}
\sphinxAtStartPar
0.8
\sphinxbeforeendvarwidth
\end{varwidth}%
&\begin{varwidth}[t]{\sphinxcolwidth{1}{5}}
\sphinxAtStartPar
0.8
\sphinxbeforeendvarwidth
\end{varwidth}%
&\begin{varwidth}[t]{\sphinxcolwidth{1}{5}}
\sphinxAtStartPar
1.5σ
\sphinxbeforeendvarwidth
\end{varwidth}%
&\begin{varwidth}[t]{\sphinxcolwidth{1}{5}}
\sphinxAtStartPar
15
\sphinxbeforeendvarwidth
\end{varwidth}%
\\
\sphinxhline\begin{varwidth}[t]{\sphinxcolwidth{1}{5}}
\sphinxAtStartPar
\sphinxstylestrong{远参考质量好}
\sphinxbeforeendvarwidth
\end{varwidth}%
&\begin{varwidth}[t]{\sphinxcolwidth{1}{5}}
\sphinxAtStartPar
0.6
\sphinxbeforeendvarwidth
\end{varwidth}%
&\begin{varwidth}[t]{\sphinxcolwidth{1}{5}}
\sphinxAtStartPar
0.6
\sphinxbeforeendvarwidth
\end{varwidth}%
&\begin{varwidth}[t]{\sphinxcolwidth{1}{5}}
\sphinxAtStartPar
2.5σ
\sphinxbeforeendvarwidth
\end{varwidth}%
&\begin{varwidth}[t]{\sphinxcolwidth{1}{5}}
\sphinxAtStartPar
5
\sphinxbeforeendvarwidth
\end{varwidth}%
\\
\sphinxbottomrule
\end{tabulary}
\sphinxtableafterendhook\par
\sphinxattableend\end{savenotes}


\subsubsection{2.8.3 在MTDP中执行RMSMR}
\label{\detokenize{chapters/chapter2:mtdprmsmr}}

\paragraph{步骤1：数据导入}
\label{\detokenize{chapters/chapter2:id87}}\begin{enumerate}
\sphinxsetlistlabels{\arabic}{enumi}{enumii}{}{.}%
\item {} 
\sphinxAtStartPar
导入本地测站时间序列

\item {} 
\sphinxAtStartPar
导入远参考站数据

\item {} 
\sphinxAtStartPar
确认时间同步

\end{enumerate}


\paragraph{步骤2：FFT设置}
\label{\detokenize{chapters/chapter2:id88}}

\begin{savenotes}\sphinxattablestart
\sphinxthistablewithglobalstyle
\centering
\begin{tabulary}{\linewidth}[t]{TT}
\sphinxtoprule
\begin{varwidth}[t]{\sphinxcolwidth{1}{2}}
\sphinxstyletheadfamily \sphinxAtStartPar
参数
\sphinxbeforeendvarwidth
\end{varwidth}%
&\begin{varwidth}[t]{\sphinxcolwidth{1}{2}}
\sphinxstyletheadfamily \sphinxAtStartPar
推荐设置
\sphinxbeforeendvarwidth
\end{varwidth}%
\\
\sphinxmidrule
\sphinxtableatstartofbodyhook\begin{varwidth}[t]{\sphinxcolwidth{1}{2}}
\sphinxAtStartPar
窗长度
\sphinxbeforeendvarwidth
\end{varwidth}%
&\begin{varwidth}[t]{\sphinxcolwidth{1}{2}}
\sphinxAtStartPar
4096点
\sphinxbeforeendvarwidth
\end{varwidth}%
\\
\sphinxhline\begin{varwidth}[t]{\sphinxcolwidth{1}{2}}
\sphinxAtStartPar
窗函数
\sphinxbeforeendvarwidth
\end{varwidth}%
&\begin{varwidth}[t]{\sphinxcolwidth{1}{2}}
\sphinxAtStartPar
Hann窗
\sphinxbeforeendvarwidth
\end{varwidth}%
\\
\sphinxhline\begin{varwidth}[t]{\sphinxcolwidth{1}{2}}
\sphinxAtStartPar
重叠率
\sphinxbeforeendvarwidth
\end{varwidth}%
&\begin{varwidth}[t]{\sphinxcolwidth{1}{2}}
\sphinxAtStartPar
50\%
\sphinxbeforeendvarwidth
\end{varwidth}%
\\
\sphinxbottomrule
\end{tabulary}
\sphinxtableafterendhook\par
\sphinxattableend\end{savenotes}


\paragraph{步骤3：选择处理方法}
\label{\detokenize{chapters/chapter2:id89}}
\sphinxAtStartPar
在MTDP中选择：
\begin{itemize}
\item {} 
\sphinxAtStartPar
✅ 远参考道处理

\item {} 
\sphinxAtStartPar
✅ 稳健估计

\item {} 
\sphinxAtStartPar
✅ RMSMR迭代

\end{itemize}

\begin{sphinxuseclass}{details}
\sphinxAtStartPar
\sphinxstylestrong{执行RMSMR处理：}
\begin{enumerate}
\sphinxsetlistlabels{\arabic}{enumi}{enumii}{}{.}%
\item {} 
\sphinxAtStartPar
菜单：\sphinxstylestrong{处理 → RMSMR处理}（或快捷键 \sphinxcode{\sphinxupquote{Ctrl+Shift+R}}）

\item {} 
\sphinxAtStartPar
在参数设置面板设置：

\begin{sphinxVerbatim}[commandchars=\\\{\}]
\PYG{p}{[}\PYG{n}{x}\PYG{p}{]} \PYG{n}{启用远参考道}
\PYG{p}{[}\PYG{n}{x}\PYG{p}{]} \PYG{n}{启用稳健估计}
\PYG{p}{[}\PYG{n}{x}\PYG{p}{]} \PYG{n}{启用多参数筛选}
\PYG{p}{[}\PYG{n}{x}\PYG{p}{]} \PYG{n}{启用Rhoplus约束}

\PYG{n}{相干度阈值}\PYG{p}{:} \PYG{l+m+mf}{0.7}
\PYG{n}{残差阈值}\PYG{p}{:} \PYG{l+m+mf}{2.0}
\PYG{n}{最大迭代次数}\PYG{p}{:} \PYG{l+m+mi}{10}
\end{sphinxVerbatim}

\item {} 
\sphinxAtStartPar
点击 \sphinxstylestrong{执行}

\item {} 
\sphinxAtStartPar
实时查看处理进度和中间结果

\end{enumerate}

\end{sphinxuseclass}

\subsubsection{2.8.4 常见问题处理（详细版）}
\label{\detokenize{chapters/chapter2:id90}}

\paragraph{问题诊断与解决步骤}
\label{\detokenize{chapters/chapter2:id91}}

\begin{savenotes}\sphinxattablestart
\sphinxthistablewithglobalstyle
\centering
\begin{tabulary}{\linewidth}[t]{TTTT}
\sphinxtoprule
\begin{varwidth}[t]{\sphinxcolwidth{1}{4}}
\sphinxstyletheadfamily \sphinxAtStartPar
问题现象
\sphinxbeforeendvarwidth
\end{varwidth}%
&\begin{varwidth}[t]{\sphinxcolwidth{1}{4}}
\sphinxstyletheadfamily \sphinxAtStartPar
可能原因
\sphinxbeforeendvarwidth
\end{varwidth}%
&\begin{varwidth}[t]{\sphinxcolwidth{1}{4}}
\sphinxstyletheadfamily \sphinxAtStartPar
诊断方法
\sphinxbeforeendvarwidth
\end{varwidth}%
&\begin{varwidth}[t]{\sphinxcolwidth{1}{4}}
\sphinxstyletheadfamily \sphinxAtStartPar
详细解决步骤
\sphinxbeforeendvarwidth
\end{varwidth}%
\\
\sphinxmidrule
\sphinxtableatstartofbodyhook\begin{varwidth}[t]{\sphinxcolwidth{1}{4}}
\sphinxAtStartPar
\sphinxstylestrong{曲线脱节}
\sphinxbeforeendvarwidth
\end{varwidth}%
&\begin{varwidth}[t]{\sphinxcolwidth{1}{4}}
\sphinxAtStartPar
死频带噪声(1\sphinxhyphen{}5kHz)
\sphinxbeforeendvarwidth
\end{varwidth}%
&\begin{varwidth}[t]{\sphinxcolwidth{1}{4}}
\sphinxAtStartPar
1. 检查1\sphinxhyphen{}5kHz频段相干度是否<0.52. 确认脱节是否发生在日间
\sphinxbeforeendvarwidth
\end{varwidth}%
&\begin{varwidth}[t]{\sphinxcolwidth{1}{4}}
\sphinxAtStartPar
\sphinxstylestrong{Step 1}: 切换到 \sphinxstylestrong{Rhoplus} 选项卡\sphinxstylestrong{Step 2}: 点击 \sphinxstylestrong{死频带校正} 按钮\sphinxstylestrong{Step 3}: 选择校正模式（自动/手动）\sphinxstylestrong{Step 4}: 预览校正效果\sphinxstylestrong{Step 5}: 点击 \sphinxstylestrong{应用校正}\sphinxstylestrong{Step 6}: 重新计算阻抗
\sphinxbeforeendvarwidth
\end{varwidth}%
\\
\sphinxhline\begin{varwidth}[t]{\sphinxcolwidth{1}{4}}
\sphinxAtStartPar
\sphinxstylestrong{相位异常}
\sphinxbeforeendvarwidth
\end{varwidth}%
&\begin{varwidth}[t]{\sphinxcolwidth{1}{4}}
\sphinxAtStartPar
近场干扰
\sphinxbeforeendvarwidth
\end{varwidth}%
&\begin{varwidth}[t]{\sphinxcolwidth{1}{4}}
\sphinxAtStartPar
1. 检查相位是否>90°或<0°2. 检查异常频点周围是否有强干扰源
\sphinxbeforeendvarwidth
\end{varwidth}%
&\begin{varwidth}[t]{\sphinxcolwidth{1}{4}}
\sphinxAtStartPar
\sphinxstylestrong{Step 1}: 在功率谱图中标记异常频点\sphinxstylestrong{Step 2}: 查看该频点的极化方向是否集中\sphinxstylestrong{Step 3}: 右键 → \sphinxstylestrong{剔除选中频点}\sphinxstylestrong{Step 4}: 重新计算阻抗\sphinxstylestrong{Step 5}: 验证相位是否恢复正常\sphinxstylestrong{Step 6}: 如果多个频点异常，考虑重新选址
\sphinxbeforeendvarwidth
\end{varwidth}%
\\
\sphinxhline\begin{varwidth}[t]{\sphinxcolwidth{1}{4}}
\sphinxAtStartPar
\sphinxstylestrong{误差棒大}
\sphinxbeforeendvarwidth
\end{varwidth}%
&\begin{varwidth}[t]{\sphinxcolwidth{1}{4}}
\sphinxAtStartPar
叠加次数不足
\sphinxbeforeendvarwidth
\end{varwidth}%
&\begin{varwidth}[t]{\sphinxcolwidth{1}{4}}
\sphinxAtStartPar
1. 检查相对误差是否>30\%2. 查看叠加次数是否<50次
\sphinxbeforeendvarwidth
\end{varwidth}%
&\begin{varwidth}[t]{\sphinxcolwidth{1}{4}}
\sphinxAtStartPar
\sphinxstylestrong{Step 1}: 检查当前叠加次数（建议>50次）\sphinxstylestrong{Step 2}: 如果叠加次数<50，尝试增加数据时长或合并相邻日期的数据\sphinxstylestrong{Step 3}: 在RMSMR参数中增加最大迭代次数\sphinxstylestrong{Step 4}: 重新执行处理\sphinxstylestrong{Step 5}: 如果误差棒仍大，检查是否有远参考数据缺失
\sphinxbeforeendvarwidth
\end{varwidth}%
\\
\sphinxhline\begin{varwidth}[t]{\sphinxcolwidth{1}{4}}
\sphinxAtStartPar
\sphinxstylestrong{低相干度}
\sphinxbeforeendvarwidth
\end{varwidth}%
&\begin{varwidth}[t]{\sphinxcolwidth{1}{4}}
\sphinxAtStartPar
人文噪声
\sphinxbeforeendvarwidth
\end{varwidth}%
&\begin{varwidth}[t]{\sphinxcolwidth{1}{4}}
\sphinxAtStartPar
1. 检查相干度曲线，找出Coh<0.5的频段/时段2. 检查对应时段的功率谱密度是否异常高
\sphinxbeforeendvarwidth
\end{varwidth}%
&\begin{varwidth}[t]{\sphinxcolwidth{1}{4}}
\sphinxAtStartPar
\sphinxstylestrong{Step 1}: 在功率谱窗口框选低相干度时段\sphinxstylestrong{Step 2}: 右键 → \sphinxstylestrong{剔除选中数据}\sphinxstylestrong{Step 3}: 重新计算阻抗\sphinxstylestrong{Step 4}: 如果相干度仍低，使用RMSMR方法迭代处理\sphinxstylestrong{Step 5}: 如果是系统性噪声，考虑增加观测时长或重新选址
\sphinxbeforeendvarwidth
\end{varwidth}%
\\
\sphinxhline\begin{varwidth}[t]{\sphinxcolwidth{1}{4}}
\sphinxAtStartPar
\sphinxstylestrong{曲线跳动}
\sphinxbeforeendvarwidth
\end{varwidth}%
&\begin{varwidth}[t]{\sphinxcolwidth{1}{4}}
\sphinxAtStartPar
强干扰
\sphinxbeforeendvarwidth
\end{varwidth}%
&\begin{varwidth}[t]{\sphinxcolwidth{1}{4}}
\sphinxAtStartPar
1. 检查功率谱密度是否有突然升高的时段2. 检查视电阻率曲线上是否有明显跳动的频点
\sphinxbeforeendvarwidth
\end{varwidth}%
&\begin{varwidth}[t]{\sphinxcolwidth{1}{4}}
\sphinxAtStartPar
\sphinxstylestrong{Step 1}: 在功率谱窗口框选异常高能量时段\sphinxstylestrong{Step 2}: 右键 → \sphinxstylestrong{剔除选中数据}\sphinxstylestrong{Step 3}: 重新计算阻抗\sphinxstylestrong{Step 4}: 如果仍跳动，使用RMSMR方法并适当提高残差阈值\sphinxstylestrong{Step 5}: 增加迭代次数至10\sphinxhyphen{}15次
\sphinxbeforeendvarwidth
\end{varwidth}%
\\
\sphinxhline\begin{varwidth}[t]{\sphinxcolwidth{1}{4}}
\sphinxAtStartPar
\sphinxstylestrong{静态位移}
\sphinxbeforeendvarwidth
\end{varwidth}%
&\begin{varwidth}[t]{\sphinxcolwidth{1}{4}}
\sphinxAtStartPar
近地表不均匀体
\sphinxbeforeendvarwidth
\end{varwidth}%
&\begin{varwidth}[t]{\sphinxcolwidth{1}{4}}
\sphinxAtStartPar
1. 对比同一测站不同方向的视电阻率曲线2. 检查是否存在整体平移（曲线形状相同但数值平移）
\sphinxbeforeendvarwidth
\end{varwidth}%
&\begin{varwidth}[t]{\sphinxcolwidth{1}{4}}
\sphinxAtStartPar
\sphinxstylestrong{Step 1}: 如果是单测站，使用Rhoplus约束校正\sphinxstylestrong{Step 2}: 如果是多测站，使用空间中值滤波校正\sphinxstylestrong{Step 3}: 检查近地表地质条件\sphinxstylestrong{Step 4}: 校正后重新评估数据质量
\sphinxbeforeendvarwidth
\end{varwidth}%
\\
\sphinxhline\begin{varwidth}[t]{\sphinxcolwidth{1}{4}}
\sphinxAtStartPar
\sphinxstylestrong{迭代不收敛}
\sphinxbeforeendvarwidth
\end{varwidth}%
&\begin{varwidth}[t]{\sphinxcolwidth{1}{4}}
\sphinxAtStartPar
参数过严或数据质量差
\sphinxbeforeendvarwidth
\end{varwidth}%
&\begin{varwidth}[t]{\sphinxcolwidth{1}{4}}
\sphinxAtStartPar
1. 检查每次迭代后数据量是否急剧减少2. 检查最终数据点是否<30\%原始量
\sphinxbeforeendvarwidth
\end{varwidth}%
&\begin{varwidth}[t]{\sphinxcolwidth{1}{4}}
\sphinxAtStartPar
\sphinxstylestrong{Step 1}: 放宽筛选参数（相干度0.7→0.6，残差2.0σ→3.0σ）\sphinxstylestrong{Step 2}: 增加最大迭代次数（10→20）\sphinxstylestrong{Step 3}: 重新执行RMSMR\sphinxstylestrong{Step 4}: 如果仍不收敛，检查原始数据质量\sphinxstylestrong{Step 5}: 如果数据本身质量差，可能需要重新采集
\sphinxbeforeendvarwidth
\end{varwidth}%
\\
\sphinxhline\begin{varwidth}[t]{\sphinxcolwidth{1}{4}}
\sphinxAtStartPar
\sphinxstylestrong{过度筛选}
\sphinxbeforeendvarwidth
\end{varwidth}%
&\begin{varwidth}[t]{\sphinxcolwidth{1}{4}}
\sphinxAtStartPar
数据量太少或阈值过严
\sphinxbeforeendvarwidth
\end{varwidth}%
&\begin{varwidth}[t]{\sphinxcolwidth{1}{4}}
\sphinxAtStartPar
1. 检查最终数据点数量是否<原始量的30\%2. 检查筛选参数设置是否过于严格
\sphinxbeforeendvarwidth
\end{varwidth}%
&\begin{varwidth}[t]{\sphinxcolwidth{1}{4}}
\sphinxAtStartPar
\sphinxstylestrong{Step 1}: 放宽筛选阈值（相干度0.8→0.6，残差1.5σ→2.5σ）\sphinxstylestrong{Step 2}: 减少迭代次数（10→5）\sphinxstylestrong{Step 3}: 检查远参考数据质量\sphinxstylestrong{Step 4}: 合并多天数据处理\sphinxstylestrong{Step 5}: 如果仍不够，需要增加观测时间重新采集
\sphinxbeforeendvarwidth
\end{varwidth}%
\\
\sphinxbottomrule
\end{tabulary}
\sphinxtableafterendhook\par
\sphinxattableend\end{savenotes}


\subsubsection{2.8.5 处理时间预估}
\label{\detokenize{chapters/chapter2:id92}}

\begin{savenotes}\sphinxattablestart
\sphinxthistablewithglobalstyle
\centering
\begin{tabulary}{\linewidth}[t]{TTTTTT}
\sphinxtoprule
\begin{varwidth}[t]{\sphinxcolwidth{1}{6}}
\sphinxstyletheadfamily \sphinxAtStartPar
数据量
\sphinxbeforeendvarwidth
\end{varwidth}%
&\begin{varwidth}[t]{\sphinxcolwidth{1}{6}}
\sphinxstyletheadfamily \sphinxAtStartPar
采样率
\sphinxbeforeendvarwidth
\end{varwidth}%
&\begin{varwidth}[t]{\sphinxcolwidth{1}{6}}
\sphinxstyletheadfamily \sphinxAtStartPar
FFT
\sphinxbeforeendvarwidth
\end{varwidth}%
&\begin{varwidth}[t]{\sphinxcolwidth{1}{6}}
\sphinxstyletheadfamily \sphinxAtStartPar
阻抗估计
\sphinxbeforeendvarwidth
\end{varwidth}%
&\begin{varwidth}[t]{\sphinxcolwidth{1}{6}}
\sphinxstyletheadfamily \sphinxAtStartPar
RMSMR(10次迭代)
\sphinxbeforeendvarwidth
\end{varwidth}%
&\begin{varwidth}[t]{\sphinxcolwidth{1}{6}}
\sphinxstyletheadfamily \sphinxAtStartPar
总计
\sphinxbeforeendvarwidth
\end{varwidth}%
\\
\sphinxmidrule
\sphinxtableatstartofbodyhook\begin{varwidth}[t]{\sphinxcolwidth{1}{6}}
\sphinxAtStartPar
\sphinxstylestrong{1小时}
\sphinxbeforeendvarwidth
\end{varwidth}%
&\begin{varwidth}[t]{\sphinxcolwidth{1}{6}}
\sphinxAtStartPar
512 Hz
\sphinxbeforeendvarwidth
\end{varwidth}%
&\begin{varwidth}[t]{\sphinxcolwidth{1}{6}}
\sphinxAtStartPar
\textasciitilde{}10秒
\sphinxbeforeendvarwidth
\end{varwidth}%
&\begin{varwidth}[t]{\sphinxcolwidth{1}{6}}
\sphinxAtStartPar
\textasciitilde{}5秒
\sphinxbeforeendvarwidth
\end{varwidth}%
&\begin{varwidth}[t]{\sphinxcolwidth{1}{6}}
\sphinxAtStartPar
\textasciitilde{}30秒
\sphinxbeforeendvarwidth
\end{varwidth}%
&\begin{varwidth}[t]{\sphinxcolwidth{1}{6}}
\sphinxAtStartPar
\textasciitilde{}1分钟
\sphinxbeforeendvarwidth
\end{varwidth}%
\\
\sphinxhline\begin{varwidth}[t]{\sphinxcolwidth{1}{6}}
\sphinxAtStartPar
\sphinxstylestrong{4小时}
\sphinxbeforeendvarwidth
\end{varwidth}%
&\begin{varwidth}[t]{\sphinxcolwidth{1}{6}}
\sphinxAtStartPar
512 Hz
\sphinxbeforeendvarwidth
\end{varwidth}%
&\begin{varwidth}[t]{\sphinxcolwidth{1}{6}}
\sphinxAtStartPar
\textasciitilde{}40秒
\sphinxbeforeendvarwidth
\end{varwidth}%
&\begin{varwidth}[t]{\sphinxcolwidth{1}{6}}
\sphinxAtStartPar
\textasciitilde{}15秒
\sphinxbeforeendvarwidth
\end{varwidth}%
&\begin{varwidth}[t]{\sphinxcolwidth{1}{6}}
\sphinxAtStartPar
\textasciitilde{}2分钟
\sphinxbeforeendvarwidth
\end{varwidth}%
&\begin{varwidth}[t]{\sphinxcolwidth{1}{6}}
\sphinxAtStartPar
\textasciitilde{}3分钟
\sphinxbeforeendvarwidth
\end{varwidth}%
\\
\sphinxhline\begin{varwidth}[t]{\sphinxcolwidth{1}{6}}
\sphinxAtStartPar
\sphinxstylestrong{12小时}
\sphinxbeforeendvarwidth
\end{varwidth}%
&\begin{varwidth}[t]{\sphinxcolwidth{1}{6}}
\sphinxAtStartPar
512 Hz
\sphinxbeforeendvarwidth
\end{varwidth}%
&\begin{varwidth}[t]{\sphinxcolwidth{1}{6}}
\sphinxAtStartPar
\textasciitilde{}2分钟
\sphinxbeforeendvarwidth
\end{varwidth}%
&\begin{varwidth}[t]{\sphinxcolwidth{1}{6}}
\sphinxAtStartPar
\textasciitilde{}30秒
\sphinxbeforeendvarwidth
\end{varwidth}%
&\begin{varwidth}[t]{\sphinxcolwidth{1}{6}}
\sphinxAtStartPar
\textasciitilde{}5分钟
\sphinxbeforeendvarwidth
\end{varwidth}%
&\begin{varwidth}[t]{\sphinxcolwidth{1}{6}}
\sphinxAtStartPar
\textasciitilde{}8分钟
\sphinxbeforeendvarwidth
\end{varwidth}%
\\
\sphinxhline\begin{varwidth}[t]{\sphinxcolwidth{1}{6}}
\sphinxAtStartPar
\sphinxstylestrong{24小时}
\sphinxbeforeendvarwidth
\end{varwidth}%
&\begin{varwidth}[t]{\sphinxcolwidth{1}{6}}
\sphinxAtStartPar
512 Hz
\sphinxbeforeendvarwidth
\end{varwidth}%
&\begin{varwidth}[t]{\sphinxcolwidth{1}{6}}
\sphinxAtStartPar
\textasciitilde{}4分钟
\sphinxbeforeendvarwidth
\end{varwidth}%
&\begin{varwidth}[t]{\sphinxcolwidth{1}{6}}
\sphinxAtStartPar
\textasciitilde{}1分钟
\sphinxbeforeendvarwidth
\end{varwidth}%
&\begin{varwidth}[t]{\sphinxcolwidth{1}{6}}
\sphinxAtStartPar
\textasciitilde{}10分钟
\sphinxbeforeendvarwidth
\end{varwidth}%
&\begin{varwidth}[t]{\sphinxcolwidth{1}{6}}
\sphinxAtStartPar
\textasciitilde{}15分钟
\sphinxbeforeendvarwidth
\end{varwidth}%
\\
\sphinxhline\begin{varwidth}[t]{\sphinxcolwidth{1}{6}}
\sphinxAtStartPar
\sphinxstylestrong{48小时}
\sphinxbeforeendvarwidth
\end{varwidth}%
&\begin{varwidth}[t]{\sphinxcolwidth{1}{6}}
\sphinxAtStartPar
512 Hz
\sphinxbeforeendvarwidth
\end{varwidth}%
&\begin{varwidth}[t]{\sphinxcolwidth{1}{6}}
\sphinxAtStartPar
\textasciitilde{}8分钟
\sphinxbeforeendvarwidth
\end{varwidth}%
&\begin{varwidth}[t]{\sphinxcolwidth{1}{6}}
\sphinxAtStartPar
\textasciitilde{}2分钟
\sphinxbeforeendvarwidth
\end{varwidth}%
&\begin{varwidth}[t]{\sphinxcolwidth{1}{6}}
\sphinxAtStartPar
\textasciitilde{}20分钟
\sphinxbeforeendvarwidth
\end{varwidth}%
&\begin{varwidth}[t]{\sphinxcolwidth{1}{6}}
\sphinxAtStartPar
\textasciitilde{}30分钟
\sphinxbeforeendvarwidth
\end{varwidth}%
\\
\sphinxbottomrule
\end{tabulary}
\sphinxtableafterendhook\par
\sphinxattableend\end{savenotes}
\begin{quote}

\sphinxAtStartPar
💡 \sphinxstylestrong{提示}：以上时间仅供参考，实际时间取决于计算机性能和数据复杂度。
\end{quote}


\subsubsection{2.8.5 处理结果验证}
\label{\detokenize{chapters/chapter2:id93}}

\paragraph{验证清单}
\label{\detokenize{chapters/chapter2:id94}}\begin{itemize}
\item {} 
\sphinxAtStartPar
 视电阻率曲线光滑连续

\item {} 
\sphinxAtStartPar
 相位在合理范围（30°\sphinxhyphen{}60°）

\item {} 
\sphinxAtStartPar
 相干度 > 0.7

\item {} 
\sphinxAtStartPar
 误差棒合理（<30\%）

\item {} 
\sphinxAtStartPar
 相位张量二维偏离度 |β| < 3°

\item {} 
\sphinxAtStartPar
 Rhoplus拟合残差小

\end{itemize}


\bigskip\hrule\bigskip


\sphinxAtStartPar
⬅️ {\hyperref[\detokenize{chapters/chapter2:id71}]{\sphinxcrossref{上一节：高级分析工具}}} | ➡️ {\hyperref[\detokenize{chapters/chapter2:id95}]{\sphinxcrossref{下一节：小结}}}


\bigskip\hrule\bigskip



\subsection{2.9 小结}
\label{\detokenize{chapters/chapter2:id95}}

\subsubsection{2.9.1 处理流程速查}
\label{\detokenize{chapters/chapter2:id96}}

\subsubsection{2.9.2 方法选择指南}
\label{\detokenize{chapters/chapter2:id97}}

\begin{savenotes}\sphinxattablestart
\sphinxthistablewithglobalstyle
\centering
\begin{tabulary}{\linewidth}[t]{TT}
\sphinxtoprule
\begin{varwidth}[t]{\sphinxcolwidth{1}{2}}
\sphinxstyletheadfamily \sphinxAtStartPar
数据条件
\sphinxbeforeendvarwidth
\end{varwidth}%
&\begin{varwidth}[t]{\sphinxcolwidth{1}{2}}
\sphinxstyletheadfamily \sphinxAtStartPar
推荐方法
\sphinxbeforeendvarwidth
\end{varwidth}%
\\
\sphinxmidrule
\sphinxtableatstartofbodyhook\begin{varwidth}[t]{\sphinxcolwidth{1}{2}}
\sphinxAtStartPar
优质数据 + 有远参考
\sphinxbeforeendvarwidth
\end{varwidth}%
&\begin{varwidth}[t]{\sphinxcolwidth{1}{2}}
\sphinxAtStartPar
远参考道法
\sphinxbeforeendvarwidth
\end{varwidth}%
\\
\sphinxhline\begin{varwidth}[t]{\sphinxcolwidth{1}{2}}
\sphinxAtStartPar
一般数据 + 有远参考
\sphinxbeforeendvarwidth
\end{varwidth}%
&\begin{varwidth}[t]{\sphinxcolwidth{1}{2}}
\sphinxAtStartPar
远参考 + 稳健估计
\sphinxbeforeendvarwidth
\end{varwidth}%
\\
\sphinxhline\begin{varwidth}[t]{\sphinxcolwidth{1}{2}}
\sphinxAtStartPar
强干扰环境
\sphinxbeforeendvarwidth
\end{varwidth}%
&\begin{varwidth}[t]{\sphinxcolwidth{1}{2}}
\sphinxAtStartPar
RMSMR迭代处理
\sphinxbeforeendvarwidth
\end{varwidth}%
\\
\sphinxhline\begin{varwidth}[t]{\sphinxcolwidth{1}{2}}
\sphinxAtStartPar
无远参考
\sphinxbeforeendvarwidth
\end{varwidth}%
&\begin{varwidth}[t]{\sphinxcolwidth{1}{2}}
\sphinxAtStartPar
稳健估计 + 严格筛选
\sphinxbeforeendvarwidth
\end{varwidth}%
\\
\sphinxbottomrule
\end{tabulary}
\sphinxtableafterendhook\par
\sphinxattableend\end{savenotes}


\subsubsection{2.9.3 质量标准}
\label{\detokenize{chapters/chapter2:id98}}

\begin{savenotes}\sphinxattablestart
\sphinxthistablewithglobalstyle
\centering
\begin{tabulary}{\linewidth}[t]{TT}
\sphinxtoprule
\begin{varwidth}[t]{\sphinxcolwidth{1}{2}}
\sphinxstyletheadfamily \sphinxAtStartPar
指标
\sphinxbeforeendvarwidth
\end{varwidth}%
&\begin{varwidth}[t]{\sphinxcolwidth{1}{2}}
\sphinxstyletheadfamily \sphinxAtStartPar
合格标准
\sphinxbeforeendvarwidth
\end{varwidth}%
\\
\sphinxmidrule
\sphinxtableatstartofbodyhook\begin{varwidth}[t]{\sphinxcolwidth{1}{2}}
\sphinxAtStartPar
常相干度
\sphinxbeforeendvarwidth
\end{varwidth}%
&\begin{varwidth}[t]{\sphinxcolwidth{1}{2}}
\sphinxAtStartPar
> 0.7
\sphinxbeforeendvarwidth
\end{varwidth}%
\\
\sphinxhline\begin{varwidth}[t]{\sphinxcolwidth{1}{2}}
\sphinxAtStartPar
重相干度
\sphinxbeforeendvarwidth
\end{varwidth}%
&\begin{varwidth}[t]{\sphinxcolwidth{1}{2}}
\sphinxAtStartPar
> 0.7
\sphinxbeforeendvarwidth
\end{varwidth}%
\\
\sphinxhline\begin{varwidth}[t]{\sphinxcolwidth{1}{2}}
\sphinxAtStartPar
相位范围
\sphinxbeforeendvarwidth
\end{varwidth}%
&\begin{varwidth}[t]{\sphinxcolwidth{1}{2}}
\sphinxAtStartPar
30° \sphinxhyphen{} 60°（TE模式）
\sphinxbeforeendvarwidth
\end{varwidth}%
\\
\sphinxhline\begin{varwidth}[t]{\sphinxcolwidth{1}{2}}
\sphinxAtStartPar
\sphinxstylestrong{二维偏离度}
\sphinxbeforeendvarwidth
\end{varwidth}%
&\begin{varwidth}[t]{\sphinxcolwidth{1}{2}}
\sphinxAtStartPar

\sphinxbeforeendvarwidth
\end{varwidth}%
\\
\sphinxhline\begin{varwidth}[t]{\sphinxcolwidth{1}{2}}
\sphinxAtStartPar
\sphinxstylestrong{曲线形态}
\sphinxbeforeendvarwidth
\end{varwidth}%
&\begin{varwidth}[t]{\sphinxcolwidth{1}{2}}
\sphinxAtStartPar
光滑连续
\sphinxbeforeendvarwidth
\end{varwidth}%
\\
\sphinxbottomrule
\end{tabulary}
\sphinxtableafterendhook\par
\sphinxattableend\end{savenotes}


\subsubsection{2.9.4 相关章节快速导航}
\label{\detokenize{chapters/chapter2:id99}}\begin{quote}

\sphinxAtStartPar
📖 \sphinxstylestrong{深入学习的章节关联}：


\begin{savenotes}\sphinxattablestart
\sphinxthistablewithglobalstyle
\centering
\begin{tabulary}{\linewidth}[t]{TTT}
\sphinxtoprule
\begin{varwidth}[t]{\sphinxcolwidth{1}{3}}
\sphinxstyletheadfamily \sphinxAtStartPar
想要深入学习...
\sphinxbeforeendvarwidth
\end{varwidth}%
&\begin{varwidth}[t]{\sphinxcolwidth{1}{3}}
\sphinxstyletheadfamily \sphinxAtStartPar
请参考章节
\sphinxbeforeendvarwidth
\end{varwidth}%
&\begin{varwidth}[t]{\sphinxcolwidth{1}{3}}
\sphinxstyletheadfamily \sphinxAtStartPar
位置
\sphinxbeforeendvarwidth
\end{varwidth}%
\\
\sphinxmidrule
\sphinxtableatstartofbodyhook\begin{varwidth}[t]{\sphinxcolwidth{1}{3}}
\sphinxAtStartPar
\sphinxstylestrong{FFT参数详细设置}
\sphinxbeforeendvarwidth
\end{varwidth}%
&\begin{varwidth}[t]{\sphinxcolwidth{1}{3}}
\sphinxAtStartPar
第4章 时间序列处理
\sphinxbeforeendvarwidth
\end{varwidth}%
&\begin{varwidth}[t]{\sphinxcolwidth{1}{3}}
\sphinxAtStartPar
4.2 FFT参数设置
\sphinxbeforeendvarwidth
\end{varwidth}%
\\
\sphinxhline\begin{varwidth}[t]{\sphinxcolwidth{1}{3}}
\sphinxAtStartPar
\sphinxstylestrong{RMSMR高级参数}
\sphinxbeforeendvarwidth
\end{varwidth}%
&\begin{varwidth}[t]{\sphinxcolwidth{1}{3}}
\sphinxAtStartPar
第5章 RMSMR处理
\sphinxbeforeendvarwidth
\end{varwidth}%
&\begin{varwidth}[t]{\sphinxcolwidth{1}{3}}
\sphinxAtStartPar
5.1\sphinxhyphen{}5.4
\sphinxbeforeendvarwidth
\end{varwidth}%
\\
\sphinxhline\begin{varwidth}[t]{\sphinxcolwidth{1}{3}}
\sphinxAtStartPar
\sphinxstylestrong{数据导入导出}
\sphinxbeforeendvarwidth
\end{varwidth}%
&\begin{varwidth}[t]{\sphinxcolwidth{1}{3}}
\sphinxAtStartPar
第3章 数据工程管理
\sphinxbeforeendvarwidth
\end{varwidth}%
&\begin{varwidth}[t]{\sphinxcolwidth{1}{3}}
\sphinxAtStartPar
3.2\sphinxhyphen{}3.4
\sphinxbeforeendvarwidth
\end{varwidth}%
\\
\sphinxhline\begin{varwidth}[t]{\sphinxcolwidth{1}{3}}
\sphinxAtStartPar
\sphinxstylestrong{远参考站设置}
\sphinxbeforeendvarwidth
\end{varwidth}%
&\begin{varwidth}[t]{\sphinxcolwidth{1}{3}}
\sphinxAtStartPar
第3章 数据工程管理
\sphinxbeforeendvarwidth
\end{varwidth}%
&\begin{varwidth}[t]{\sphinxcolwidth{1}{3}}
\sphinxAtStartPar
3.5 远参考管理
\sphinxbeforeendvarwidth
\end{varwidth}%
\\
\sphinxhline\begin{varwidth}[t]{\sphinxcolwidth{1}{3}}
\sphinxAtStartPar
\sphinxstylestrong{反演处理}
\sphinxbeforeendvarwidth
\end{varwidth}%
&\begin{varwidth}[t]{\sphinxcolwidth{1}{3}}
\sphinxAtStartPar
第6章 反演
\sphinxbeforeendvarwidth
\end{varwidth}%
&\begin{varwidth}[t]{\sphinxcolwidth{1}{3}}
\sphinxAtStartPar
6.1\sphinxhyphen{}6.5
\sphinxbeforeendvarwidth
\end{varwidth}%
\\
\sphinxhline\begin{varwidth}[t]{\sphinxcolwidth{1}{3}}
\sphinxAtStartPar
\sphinxstylestrong{结果可视化}
\sphinxbeforeendvarwidth
\end{varwidth}%
&\begin{varwidth}[t]{\sphinxcolwidth{1}{3}}
\sphinxAtStartPar
第7章 可视化
\sphinxbeforeendvarwidth
\end{varwidth}%
&\begin{varwidth}[t]{\sphinxcolwidth{1}{3}}
\sphinxAtStartPar
7.1\sphinxhyphen{}7.4
\sphinxbeforeendvarwidth
\end{varwidth}%
\\
\sphinxbottomrule
\end{tabulary}
\sphinxtableafterendhook\par
\sphinxattableend\end{savenotes}
\end{quote}


\bigskip\hrule\bigskip


\sphinxAtStartPar
⬅️ {\hyperref[\detokenize{chapters/chapter2:rmsmr}]{\sphinxcrossref{上一节：RMSMR完整流程}}} | ➡️ {\hyperref[\detokenize{chapters/chapter2:a}]{\sphinxcrossref{附录}}}


\bigskip\hrule\bigskip



\subsection{附录A：权威参考文献}
\label{\detokenize{chapters/chapter2:a}}

\subsubsection{经典教材}
\label{\detokenize{chapters/chapter2:id100}}

\begin{savenotes}\sphinxattablestart
\sphinxthistablewithglobalstyle
\centering
\begin{tabulary}{\linewidth}[t]{TTT}
\sphinxtoprule
\begin{varwidth}[t]{\sphinxcolwidth{1}{3}}
\sphinxstyletheadfamily \sphinxAtStartPar
文献
\sphinxbeforeendvarwidth
\end{varwidth}%
&\begin{varwidth}[t]{\sphinxcolwidth{1}{3}}
\sphinxstyletheadfamily \sphinxAtStartPar
年份
\sphinxbeforeendvarwidth
\end{varwidth}%
&\begin{varwidth}[t]{\sphinxcolwidth{1}{3}}
\sphinxstyletheadfamily \sphinxAtStartPar
说明
\sphinxbeforeendvarwidth
\end{varwidth}%
\\
\sphinxmidrule
\sphinxtableatstartofbodyhook\begin{varwidth}[t]{\sphinxcolwidth{1}{3}}
\sphinxAtStartPar
Chave \& Jones. \sphinxstylestrong{The Magnetotelluric Method: Theory and Practice}
\sphinxbeforeendvarwidth
\end{varwidth}%
&\begin{varwidth}[t]{\sphinxcolwidth{1}{3}}
\sphinxAtStartPar
2012
\sphinxbeforeendvarwidth
\end{varwidth}%
&\begin{varwidth}[t]{\sphinxcolwidth{1}{3}}
\sphinxAtStartPar
MT领域最权威教材
\sphinxbeforeendvarwidth
\end{varwidth}%
\\
\sphinxhline\begin{varwidth}[t]{\sphinxcolwidth{1}{3}}
\sphinxAtStartPar
Simpson \& Bahr. \sphinxstylestrong{Practical Magnetotellurics}
\sphinxbeforeendvarwidth
\end{varwidth}%
&\begin{varwidth}[t]{\sphinxcolwidth{1}{3}}
\sphinxAtStartPar
2005
\sphinxbeforeendvarwidth
\end{varwidth}%
&\begin{varwidth}[t]{\sphinxcolwidth{1}{3}}
\sphinxAtStartPar
实用指南
\sphinxbeforeendvarwidth
\end{varwidth}%
\\
\sphinxhline\begin{varwidth}[t]{\sphinxcolwidth{1}{3}}
\sphinxAtStartPar
Kaufman \& Keller. \sphinxstylestrong{The Magnetotelluric Sounding Method}
\sphinxbeforeendvarwidth
\end{varwidth}%
&\begin{varwidth}[t]{\sphinxcolwidth{1}{3}}
\sphinxAtStartPar
1981
\sphinxbeforeendvarwidth
\end{varwidth}%
&\begin{varwidth}[t]{\sphinxcolwidth{1}{3}}
\sphinxAtStartPar
经典理论
\sphinxbeforeendvarwidth
\end{varwidth}%
\\
\sphinxbottomrule
\end{tabulary}
\sphinxtableafterendhook\par
\sphinxattableend\end{savenotes}


\subsubsection{关键论文}
\label{\detokenize{chapters/chapter2:id101}}

\begin{savenotes}\sphinxattablestart
\sphinxthistablewithglobalstyle
\centering
\begin{tabulary}{\linewidth}[t]{TTT}
\sphinxtoprule
\begin{varwidth}[t]{\sphinxcolwidth{1}{3}}
\sphinxstyletheadfamily \sphinxAtStartPar
主题
\sphinxbeforeendvarwidth
\end{varwidth}%
&\begin{varwidth}[t]{\sphinxcolwidth{1}{3}}
\sphinxstyletheadfamily \sphinxAtStartPar
文献
\sphinxbeforeendvarwidth
\end{varwidth}%
&\begin{varwidth}[t]{\sphinxcolwidth{1}{3}}
\sphinxstyletheadfamily \sphinxAtStartPar
贡献
\sphinxbeforeendvarwidth
\end{varwidth}%
\\
\sphinxmidrule
\sphinxtableatstartofbodyhook\begin{varwidth}[t]{\sphinxcolwidth{1}{3}}
\sphinxAtStartPar
\sphinxstylestrong{远参考道法}
\sphinxbeforeendvarwidth
\end{varwidth}%
&\begin{varwidth}[t]{\sphinxcolwidth{1}{3}}
\sphinxAtStartPar
Gamble et al. (1979)
\sphinxbeforeendvarwidth
\end{varwidth}%
&\begin{varwidth}[t]{\sphinxcolwidth{1}{3}}
\sphinxAtStartPar
奠基性论文
\sphinxbeforeendvarwidth
\end{varwidth}%
\\
\sphinxhline\begin{varwidth}[t]{\sphinxcolwidth{1}{3}}
\sphinxAtStartPar
\sphinxstylestrong{稳健估计}
\sphinxbeforeendvarwidth
\end{varwidth}%
&\begin{varwidth}[t]{\sphinxcolwidth{1}{3}}
\sphinxAtStartPar
Chave et al. (1987, 2004)
\sphinxbeforeendvarwidth
\end{varwidth}%
&\begin{varwidth}[t]{\sphinxcolwidth{1}{3}}
\sphinxAtStartPar
BIRRP核心算法
\sphinxbeforeendvarwidth
\end{varwidth}%
\\
\sphinxhline\begin{varwidth}[t]{\sphinxcolwidth{1}{3}}
\sphinxAtStartPar
\sphinxstylestrong{相位张量}
\sphinxbeforeendvarwidth
\end{varwidth}%
&\begin{varwidth}[t]{\sphinxcolwidth{1}{3}}
\sphinxAtStartPar
Caldwell et al. (2004)
\sphinxbeforeendvarwidth
\end{varwidth}%
&\begin{varwidth}[t]{\sphinxcolwidth{1}{3}}
\sphinxAtStartPar
原创理论
\sphinxbeforeendvarwidth
\end{varwidth}%
\\
\sphinxhline\begin{varwidth}[t]{\sphinxcolwidth{1}{3}}
\sphinxAtStartPar
\sphinxstylestrong{Rhoplus}
\sphinxbeforeendvarwidth
\end{varwidth}%
&\begin{varwidth}[t]{\sphinxcolwidth{1}{3}}
\sphinxAtStartPar
Parker \& Booker (1996)
\sphinxbeforeendvarwidth
\end{varwidth}%
&\begin{varwidth}[t]{\sphinxcolwidth{1}{3}}
\sphinxAtStartPar
原创方法
\sphinxbeforeendvarwidth
\end{varwidth}%
\\
\sphinxhline\begin{varwidth}[t]{\sphinxcolwidth{1}{3}}
\sphinxAtStartPar
\sphinxstylestrong{RMSMR}
\sphinxbeforeendvarwidth
\end{varwidth}%
&\begin{varwidth}[t]{\sphinxcolwidth{1}{3}}
\sphinxAtStartPar
王培杰等 (2024)
\sphinxbeforeendvarwidth
\end{varwidth}%
&\begin{varwidth}[t]{\sphinxcolwidth{1}{3}}
\sphinxAtStartPar
MTDP采用的方法
\sphinxbeforeendvarwidth
\end{varwidth}%
\\
\sphinxbottomrule
\end{tabulary}
\sphinxtableafterendhook\par
\sphinxattableend\end{savenotes}


\subsubsection{开源工具}
\label{\detokenize{chapters/chapter2:id102}}

\begin{savenotes}\sphinxattablestart
\sphinxthistablewithglobalstyle
\centering
\begin{tabulary}{\linewidth}[t]{TTT}
\sphinxtoprule
\begin{varwidth}[t]{\sphinxcolwidth{1}{3}}
\sphinxstyletheadfamily \sphinxAtStartPar
项目
\sphinxbeforeendvarwidth
\end{varwidth}%
&\begin{varwidth}[t]{\sphinxcolwidth{1}{3}}
\sphinxstyletheadfamily \sphinxAtStartPar
网址
\sphinxbeforeendvarwidth
\end{varwidth}%
&\begin{varwidth}[t]{\sphinxcolwidth{1}{3}}
\sphinxstyletheadfamily \sphinxAtStartPar
说明
\sphinxbeforeendvarwidth
\end{varwidth}%
\\
\sphinxmidrule
\sphinxtableatstartofbodyhook\begin{varwidth}[t]{\sphinxcolwidth{1}{3}}
\sphinxAtStartPar
\sphinxstylestrong{MTpy}
\sphinxbeforeendvarwidth
\end{varwidth}%
&\begin{varwidth}[t]{\sphinxcolwidth{1}{3}}
\sphinxAtStartPar
\sphinxhref{http://github.com/MTgeophysics/mtpy}{github.com/MTgeophysics/mtpy} (http://github.com/MTgeophysics/mtpy)
\sphinxbeforeendvarwidth
\end{varwidth}%
&\begin{varwidth}[t]{\sphinxcolwidth{1}{3}}
\sphinxAtStartPar
Python MT库
\sphinxbeforeendvarwidth
\end{varwidth}%
\\
\sphinxhline\begin{varwidth}[t]{\sphinxcolwidth{1}{3}}
\sphinxAtStartPar
\sphinxstylestrong{BIRRP}
\sphinxbeforeendvarwidth
\end{varwidth}%
&\begin{varwidth}[t]{\sphinxcolwidth{1}{3}}
\sphinxAtStartPar
\sphinxhref{http://dias.ie}{dias.ie} (http://dias.ie)
\sphinxbeforeendvarwidth
\end{varwidth}%
&\begin{varwidth}[t]{\sphinxcolwidth{1}{3}}
\sphinxAtStartPar
稳健处理
\sphinxbeforeendvarwidth
\end{varwidth}%
\\
\sphinxhline\begin{varwidth}[t]{\sphinxcolwidth{1}{3}}
\sphinxAtStartPar
\sphinxstylestrong{EMTF}
\sphinxbeforeendvarwidth
\end{varwidth}%
&\begin{varwidth}[t]{\sphinxcolwidth{1}{3}}
\sphinxAtStartPar
\sphinxhref{http://iris.edu}{iris.edu} (http://iris.edu)
\sphinxbeforeendvarwidth
\end{varwidth}%
&\begin{varwidth}[t]{\sphinxcolwidth{1}{3}}
\sphinxAtStartPar
USGS工具
\sphinxbeforeendvarwidth
\end{varwidth}%
\\
\sphinxbottomrule
\end{tabulary}
\sphinxtableafterendhook\par
\sphinxattableend\end{savenotes}


\bigskip\hrule\bigskip



\subsection{附录B：公式速查表}
\label{\detokenize{chapters/chapter2:b}}

\subsubsection{基本公式}
\label{\detokenize{chapters/chapter2:id103}}

\begin{savenotes}\sphinxattablestart
\sphinxthistablewithglobalstyle
\centering
\begin{tabulary}{\linewidth}[t]{TTT}
\sphinxtoprule
\begin{varwidth}[t]{\sphinxcolwidth{1}{3}}
\sphinxstyletheadfamily \sphinxAtStartPar
名称
\sphinxbeforeendvarwidth
\end{varwidth}%
&\begin{varwidth}[t]{\sphinxcolwidth{1}{3}}
\sphinxstyletheadfamily \sphinxAtStartPar
公式
\sphinxbeforeendvarwidth
\end{varwidth}%
&\begin{varwidth}[t]{\sphinxcolwidth{1}{3}}
\sphinxstyletheadfamily \sphinxAtStartPar
编号
\sphinxbeforeendvarwidth
\end{varwidth}%
\\
\sphinxmidrule
\sphinxtableatstartofbodyhook\begin{varwidth}[t]{\sphinxcolwidth{1}{3}}
\sphinxAtStartPar
\sphinxstylestrong{趋肤深度}
\sphinxbeforeendvarwidth
\end{varwidth}%
&\begin{varwidth}[t]{\sphinxcolwidth{1}{3}}
\sphinxAtStartPar
\(\delta = 503\sqrt{\rho/f}\)
\sphinxbeforeendvarwidth
\end{varwidth}%
&\begin{varwidth}[t]{\sphinxcolwidth{1}{3}}
\sphinxAtStartPar
(2.1)
\sphinxbeforeendvarwidth
\end{varwidth}%
\\
\sphinxhline\begin{varwidth}[t]{\sphinxcolwidth{1}{3}}
\sphinxAtStartPar
\sphinxstylestrong{阻抗关系}
\sphinxbeforeendvarwidth
\end{varwidth}%
&\begin{varwidth}[t]{\sphinxcolwidth{1}{3}}
\sphinxAtStartPar
\(\mathbf{E} = \mathbf{Z}\mathbf{H}\)
\sphinxbeforeendvarwidth
\end{varwidth}%
&\begin{varwidth}[t]{\sphinxcolwidth{1}{3}}
\sphinxAtStartPar
(2.2)
\sphinxbeforeendvarwidth
\end{varwidth}%
\\
\sphinxhline\begin{varwidth}[t]{\sphinxcolwidth{1}{3}}
\sphinxAtStartPar
\sphinxstylestrong{视电阻率}
\sphinxbeforeendvarwidth
\end{varwidth}%
&\begin{varwidth}[t]{\sphinxcolwidth{1}{3}}
\sphinxAtStartPar
\(\rho_a = |Z|^2/(\omega\mu_0)\)
\sphinxbeforeendvarwidth
\end{varwidth}%
&\begin{varwidth}[t]{\sphinxcolwidth{1}{3}}
\sphinxAtStartPar
(2.4)
\sphinxbeforeendvarwidth
\end{varwidth}%
\\
\sphinxhline\begin{varwidth}[t]{\sphinxcolwidth{1}{3}}
\sphinxAtStartPar
\sphinxstylestrong{相位}
\sphinxbeforeendvarwidth
\end{varwidth}%
&\begin{varwidth}[t]{\sphinxcolwidth{1}{3}}
\sphinxAtStartPar
\(\varphi = \arg(Z)\)
\sphinxbeforeendvarwidth
\end{varwidth}%
&\begin{varwidth}[t]{\sphinxcolwidth{1}{3}}
\sphinxAtStartPar
(2.5)
\sphinxbeforeendvarwidth
\end{varwidth}%
\\
\sphinxhline\begin{varwidth}[t]{\sphinxcolwidth{1}{3}}
\sphinxAtStartPar
\sphinxstylestrong{倾子关系}
\sphinxbeforeendvarwidth
\end{varwidth}%
&\begin{varwidth}[t]{\sphinxcolwidth{1}{3}}
\sphinxAtStartPar
\(H_z = T_{zx}H_x + T_{zy}H_y\)
\sphinxbeforeendvarwidth
\end{varwidth}%
&\begin{varwidth}[t]{\sphinxcolwidth{1}{3}}
\sphinxAtStartPar
(2.6)
\sphinxbeforeendvarwidth
\end{varwidth}%
\\
\sphinxbottomrule
\end{tabulary}
\sphinxtableafterendhook\par
\sphinxattableend\end{savenotes}


\subsubsection{估计方法}
\label{\detokenize{chapters/chapter2:id104}}

\begin{savenotes}\sphinxattablestart
\sphinxthistablewithglobalstyle
\centering
\begin{tabulary}{\linewidth}[t]{TTT}
\sphinxtoprule
\begin{varwidth}[t]{\sphinxcolwidth{1}{3}}
\sphinxstyletheadfamily \sphinxAtStartPar
方法
\sphinxbeforeendvarwidth
\end{varwidth}%
&\begin{varwidth}[t]{\sphinxcolwidth{1}{3}}
\sphinxstyletheadfamily \sphinxAtStartPar
公式
\sphinxbeforeendvarwidth
\end{varwidth}%
&\begin{varwidth}[t]{\sphinxcolwidth{1}{3}}
\sphinxstyletheadfamily \sphinxAtStartPar
编号
\sphinxbeforeendvarwidth
\end{varwidth}%
\\
\sphinxmidrule
\sphinxtableatstartofbodyhook\begin{varwidth}[t]{\sphinxcolwidth{1}{3}}
\sphinxAtStartPar
\sphinxstylestrong{最小二乘}
\sphinxbeforeendvarwidth
\end{varwidth}%
&\begin{varwidth}[t]{\sphinxcolwidth{1}{3}}
\sphinxAtStartPar
\(\mathbf{Z} = [\mathbf{E}\mathbf{H}^*][\mathbf{H}\mathbf{H}^*]^{-1}\)
\sphinxbeforeendvarwidth
\end{varwidth}%
&\begin{varwidth}[t]{\sphinxcolwidth{1}{3}}
\sphinxAtStartPar
(2.10)
\sphinxbeforeendvarwidth
\end{varwidth}%
\\
\sphinxhline\begin{varwidth}[t]{\sphinxcolwidth{1}{3}}
\sphinxAtStartPar
\sphinxstylestrong{远参考道}
\sphinxbeforeendvarwidth
\end{varwidth}%
&\begin{varwidth}[t]{\sphinxcolwidth{1}{3}}
\sphinxAtStartPar
\(\mathbf{Z} = [\mathbf{E}\mathbf{R}^*][\mathbf{H}\mathbf{R}^*]^{-1}\)
\sphinxbeforeendvarwidth
\end{varwidth}%
&\begin{varwidth}[t]{\sphinxcolwidth{1}{3}}
\sphinxAtStartPar
(2.11)
\sphinxbeforeendvarwidth
\end{varwidth}%
\\
\sphinxbottomrule
\end{tabulary}
\sphinxtableafterendhook\par
\sphinxattableend\end{savenotes}


\subsubsection{相干度}
\label{\detokenize{chapters/chapter2:id105}}

\begin{savenotes}\sphinxattablestart
\sphinxthistablewithglobalstyle
\centering
\begin{tabulary}{\linewidth}[t]{TTT}
\sphinxtoprule
\begin{varwidth}[t]{\sphinxcolwidth{1}{3}}
\sphinxstyletheadfamily \sphinxAtStartPar
类型
\sphinxbeforeendvarwidth
\end{varwidth}%
&\begin{varwidth}[t]{\sphinxcolwidth{1}{3}}
\sphinxstyletheadfamily \sphinxAtStartPar
公式
\sphinxbeforeendvarwidth
\end{varwidth}%
&\begin{varwidth}[t]{\sphinxcolwidth{1}{3}}
\sphinxstyletheadfamily \sphinxAtStartPar
编号
\sphinxbeforeendvarwidth
\end{varwidth}%
\\
\sphinxmidrule
\sphinxtableatstartofbodyhook\begin{varwidth}[t]{\sphinxcolwidth{1}{3}}
\sphinxAtStartPar
\sphinxstylestrong{常相干度}
\sphinxbeforeendvarwidth
\end{varwidth}%
&\begin{varwidth}[t]{\sphinxcolwidth{1}{3}}
\sphinxAtStartPar
\(Coh^2(XY) = |[XY^*]|^2/([XX^*][YY^*])\)
\sphinxbeforeendvarwidth
\end{varwidth}%
&\begin{varwidth}[t]{\sphinxcolwidth{1}{3}}
\sphinxAtStartPar
(2.13)
\sphinxbeforeendvarwidth
\end{varwidth}%
\\
\sphinxhline\begin{varwidth}[t]{\sphinxcolwidth{1}{3}}
\sphinxAtStartPar
\sphinxstylestrong{重相干度}
\sphinxbeforeendvarwidth
\end{varwidth}%
&\begin{varwidth}[t]{\sphinxcolwidth{1}{3}}
\sphinxAtStartPar
\(Coh_m^2 = (Z_{xy}[H_yE_x^*] + Z_{xx}[H_xE_x^*])/[E_xE_x^*]\)
\sphinxbeforeendvarwidth
\end{varwidth}%
&\begin{varwidth}[t]{\sphinxcolwidth{1}{3}}
\sphinxAtStartPar
(2.14)
\sphinxbeforeendvarwidth
\end{varwidth}%
\\
\sphinxhline\begin{varwidth}[t]{\sphinxcolwidth{1}{3}}
\sphinxAtStartPar
\sphinxstylestrong{偏相干度}
\sphinxbeforeendvarwidth
\end{varwidth}%
&\begin{varwidth}[t]{\sphinxcolwidth{1}{3}}
\sphinxAtStartPar
\(Coh_p^2 = (Coh_m^2 - Coh^2)/(1 - Coh^2)\)
\sphinxbeforeendvarwidth
\end{varwidth}%
&\begin{varwidth}[t]{\sphinxcolwidth{1}{3}}
\sphinxAtStartPar
(2.15)
\sphinxbeforeendvarwidth
\end{varwidth}%
\\
\sphinxbottomrule
\end{tabulary}
\sphinxtableafterendhook\par
\sphinxattableend\end{savenotes}


\subsubsection{相位张量}
\label{\detokenize{chapters/chapter2:id106}}

\begin{savenotes}\sphinxattablestart
\sphinxthistablewithglobalstyle
\centering
\begin{tabulary}{\linewidth}[t]{TTT}
\sphinxtoprule
\begin{varwidth}[t]{\sphinxcolwidth{1}{3}}
\sphinxstyletheadfamily \sphinxAtStartPar
名称
\sphinxbeforeendvarwidth
\end{varwidth}%
&\begin{varwidth}[t]{\sphinxcolwidth{1}{3}}
\sphinxstyletheadfamily \sphinxAtStartPar
公式
\sphinxbeforeendvarwidth
\end{varwidth}%
&\begin{varwidth}[t]{\sphinxcolwidth{1}{3}}
\sphinxstyletheadfamily \sphinxAtStartPar
编号
\sphinxbeforeendvarwidth
\end{varwidth}%
\\
\sphinxmidrule
\sphinxtableatstartofbodyhook\begin{varwidth}[t]{\sphinxcolwidth{1}{3}}
\sphinxAtStartPar
\sphinxstylestrong{相位张量}
\sphinxbeforeendvarwidth
\end{varwidth}%
&\begin{varwidth}[t]{\sphinxcolwidth{1}{3}}
\sphinxAtStartPar
\(\mathbf{\Phi} = \mathbf{X}^{-1}\mathbf{Y}\)
\sphinxbeforeendvarwidth
\end{varwidth}%
&\begin{varwidth}[t]{\sphinxcolwidth{1}{3}}
\sphinxAtStartPar
(2.19)
\sphinxbeforeendvarwidth
\end{varwidth}%
\\
\sphinxhline\begin{varwidth}[t]{\sphinxcolwidth{1}{3}}
\sphinxAtStartPar
\sphinxstylestrong{主方向角}
\sphinxbeforeendvarwidth
\end{varwidth}%
&\begin{varwidth}[t]{\sphinxcolwidth{1}{3}}
\sphinxAtStartPar
\(\alpha = \frac{1}{2}\arctan(\frac{\Phi_{xy}+\Phi_{yx}}{\Phi_{xx}-\Phi_{yy}})\)
\sphinxbeforeendvarwidth
\end{varwidth}%
&\begin{varwidth}[t]{\sphinxcolwidth{1}{3}}
\sphinxAtStartPar
(2.21)
\sphinxbeforeendvarwidth
\end{varwidth}%
\\
\sphinxhline\begin{varwidth}[t]{\sphinxcolwidth{1}{3}}
\sphinxAtStartPar
\sphinxstylestrong{二维偏离度}
\sphinxbeforeendvarwidth
\end{varwidth}%
&\begin{varwidth}[t]{\sphinxcolwidth{1}{3}}
\sphinxAtStartPar
\(\beta = \frac{1}{2}\arctan(\frac{\Phi_{xy}-\Phi_{yx}}{\Phi_{xx}+\Phi_{yy}})\)
\sphinxbeforeendvarwidth
\end{varwidth}%
&\begin{varwidth}[t]{\sphinxcolwidth{1}{3}}
\sphinxAtStartPar
(2.22)
\sphinxbeforeendvarwidth
\end{varwidth}%
\\
\sphinxhline\begin{varwidth}[t]{\sphinxcolwidth{1}{3}}
\sphinxAtStartPar
\sphinxstylestrong{椭圆率}
\sphinxbeforeendvarwidth
\end{varwidth}%
&\begin{varwidth}[t]{\sphinxcolwidth{1}{3}}
\sphinxAtStartPar
\(\lambda = \frac{\Phi_{max} - \Phi_{min}}{\Phi_{max} + \Phi_{min}}\)
\sphinxbeforeendvarwidth
\end{varwidth}%
&\begin{varwidth}[t]{\sphinxcolwidth{1}{3}}
\sphinxAtStartPar
(2.23)
\sphinxbeforeendvarwidth
\end{varwidth}%
\\
\sphinxhline\begin{varwidth}[t]{\sphinxcolwidth{1}{3}}
\sphinxAtStartPar
\sphinxstylestrong{一维偏离度}
\sphinxbeforeendvarwidth
\end{varwidth}%
&\begin{varwidth}[t]{\sphinxcolwidth{1}{3}}
\sphinxAtStartPar
\(\kappa = \sqrt{\frac{(\Phi_{xx}-\Phi_{yy})^2 + 4\Phi_{xy}^2}{(\Phi_{xx}+\Phi_{yy})^2}}\)
\sphinxbeforeendvarwidth
\end{varwidth}%
&\begin{varwidth}[t]{\sphinxcolwidth{1}{3}}
\sphinxAtStartPar
(2.24)
\sphinxbeforeendvarwidth
\end{varwidth}%
\\
\sphinxbottomrule
\end{tabulary}
\sphinxtableafterendhook\par
\sphinxattableend\end{savenotes}


\subsubsection{其他公式}
\label{\detokenize{chapters/chapter2:id107}}

\begin{savenotes}\sphinxattablestart
\sphinxthistablewithglobalstyle
\centering
\begin{tabulary}{\linewidth}[t]{TTT}
\sphinxtoprule
\begin{varwidth}[t]{\sphinxcolwidth{1}{3}}
\sphinxstyletheadfamily \sphinxAtStartPar
名称
\sphinxbeforeendvarwidth
\end{varwidth}%
&\begin{varwidth}[t]{\sphinxcolwidth{1}{3}}
\sphinxstyletheadfamily \sphinxAtStartPar
公式
\sphinxbeforeendvarwidth
\end{varwidth}%
&\begin{varwidth}[t]{\sphinxcolwidth{1}{3}}
\sphinxstyletheadfamily \sphinxAtStartPar
编号
\sphinxbeforeendvarwidth
\end{varwidth}%
\\
\sphinxmidrule
\sphinxtableatstartofbodyhook\begin{varwidth}[t]{\sphinxcolwidth{1}{3}}
\sphinxAtStartPar
\sphinxstylestrong{马氏距离}
\sphinxbeforeendvarwidth
\end{varwidth}%
&\begin{varwidth}[t]{\sphinxcolwidth{1}{3}}
\sphinxAtStartPar
\(D_M = \sqrt{(\mathbf{x}-\boldsymbol{\mu})^T\mathbf{\Sigma}^{-1}(\mathbf{x}-\boldsymbol{\mu})}\)
\sphinxbeforeendvarwidth
\end{varwidth}%
&\begin{varwidth}[t]{\sphinxcolwidth{1}{3}}
\sphinxAtStartPar
(2.18)
\sphinxbeforeendvarwidth
\end{varwidth}%
\\
\sphinxhline\begin{varwidth}[t]{\sphinxcolwidth{1}{3}}
\sphinxAtStartPar
\sphinxstylestrong{Gamble方差}
\sphinxbeforeendvarwidth
\end{varwidth}%
&\begin{varwidth}[t]{\sphinxcolwidth{1}{3}}
\sphinxAtStartPar
\$(\textbackslash{}Delta Z\_\{ij\})\textasciicircum{}2 = \textbackslash{}frac\{{[}
\sphinxbeforeendvarwidth
\end{varwidth}%
&\begin{varwidth}[t]{\sphinxcolwidth{1}{3}}
\sphinxAtStartPar
\textbackslash{}eta\_i
\sphinxbeforeendvarwidth
\end{varwidth}%
\\
\sphinxbottomrule
\end{tabulary}
\sphinxtableafterendhook\par
\sphinxattableend\end{savenotes}


\bigskip\hrule\bigskip



\subsection{附录C：学习路径建议}
\label{\detokenize{chapters/chapter2:c}}

\subsubsection{初学者路径}
\label{\detokenize{chapters/chapter2:id108}}
\begin{sphinxVerbatim}[commandchars=\\\{\}]
1. 阅读 2.2 理解基本原理（20分钟）
2. 阅读 2.5 掌握阻抗估计方法（20分钟）
3. 用示例数据练习基本处理
4. 阅读 2.6 学习质量评估
\end{sphinxVerbatim}


\subsubsection{进阶用户路径}
\label{\detokenize{chapters/chapter2:id109}}
\begin{sphinxVerbatim}[commandchars=\\\{\}]
\PYG{l+m+mf}{1.} \PYG{n}{复习} \PYG{l+m+mf}{2.5} \PYG{n}{深入理解各种方法}
\PYG{l+m+mf}{2.} \PYG{n}{阅读} \PYG{l+m+mf}{2.7} \PYG{n}{掌握高级分析工具}
\PYG{l+m+mf}{3.} \PYG{n}{阅读} \PYG{l+m+mf}{2.8} \PYG{n}{学习RMSMR完整流程}
\PYG{l+m+mf}{4.} \PYG{n}{用实际数据练习高级处理}
\end{sphinxVerbatim}


\subsubsection{问题排查路径}
\label{\detokenize{chapters/chapter2:id110}}
\begin{sphinxVerbatim}[commandchars=\\\{\}]
遇到数据质量问题 → 
  查看 2.6 质量评估指标 → 
    查看 2.8.4 常见问题处理 → 
      调整参数重新处理
\end{sphinxVerbatim}


\bigskip\hrule\bigskip

\begin{quote}

\sphinxAtStartPar
\sphinxstylestrong{下一章预告}：第3章将介绍MTDP的数据工程管理功能，包括项目组织、数据导入导出等实用操作。
\end{quote}


\bigskip\hrule\bigskip



\subsection{附录D：术语索引}
\label{\detokenize{chapters/chapter2:d}}

\subsubsection{B}
\label{\detokenize{chapters/chapter2:id111}}\begin{itemize}
\item {} 
\sphinxAtStartPar
\sphinxstylestrong{BIRRP} .................... 2.5.5, 附录A

\item {} 
\sphinxAtStartPar
\sphinxstylestrong{半空间} .................... 2.2.3

\item {} 
\sphinxAtStartPar
\sphinxstylestrong{被影响估计} ................ 2.5.5

\end{itemize}


\subsubsection{C}
\label{\detokenize{chapters/chapter2:id112}}\begin{itemize}
\item {} 
\sphinxAtStartPar
\sphinxstylestrong{常相干度} .................. 2.6.2

\item {} 
\sphinxAtStartPar
\sphinxstylestrong{重相干度} .................. 2.6.2

\item {} 
\sphinxAtStartPar
\sphinxstylestrong{传感器} .................... 2.3.1

\item {} 
\sphinxAtStartPar
\sphinxstylestrong{传输函数} .................. 2.2

\end{itemize}


\subsubsection{D}
\label{\detokenize{chapters/chapter2:id113}}\begin{itemize}
\item {} 
\sphinxAtStartPar
\sphinxstylestrong{电极} ...................... 2.3.1

\item {} 
\sphinxAtStartPar
\sphinxstylestrong{电磁场} .................... 2.2.1

\item {} 
\sphinxAtStartPar
\sphinxstylestrong{叠加} ...................... 2.4.2

\item {} 
\sphinxAtStartPar
\sphinxstylestrong{对流层} .................... 2.2.1

\end{itemize}


\subsubsection{E}
\label{\detokenize{chapters/chapter2:e}}\begin{itemize}
\item {} 
\sphinxAtStartPar
\sphinxstylestrong{FFT} ....................... 2.4.1

\item {} 
\sphinxAtStartPar
\sphinxstylestrong{EMTF} ...................... 附录A

\end{itemize}


\subsubsection{F}
\label{\detokenize{chapters/chapter2:f}}\begin{itemize}
\item {} 
\sphinxAtStartPar
\sphinxstylestrong{方差估计} .................. 2.5

\item {} 
\sphinxAtStartPar
\sphinxstylestrong{傅里叶变换} ................ 2.4.1

\end{itemize}


\subsubsection{G}
\label{\detokenize{chapters/chapter2:g}}\begin{itemize}
\item {} 
\sphinxAtStartPar
\sphinxstylestrong{Gamble方差} ................ 2.5.1

\item {} 
\sphinxAtStartPar
\sphinxstylestrong{功率谱密度} ................ 2.4.2, 2.6.3

\end{itemize}


\subsubsection{H}
\label{\detokenize{chapters/chapter2:h}}\begin{itemize}
\item {} 
\sphinxAtStartPar
\sphinxstylestrong{Hann窗} .................... 2.4.1

\end{itemize}


\subsubsection{I}
\label{\detokenize{chapters/chapter2:i}}\begin{itemize}
\item {} 
\sphinxAtStartPar
\sphinxstylestrong{Impedance} ................. 见：阻抗

\end{itemize}


\subsubsection{J}
\label{\detokenize{chapters/chapter2:j}}\begin{itemize}
\item {} 
\sphinxAtStartPar
\sphinxstylestrong{极化方向} .................. 2.4.3

\item {} 
\sphinxAtStartPar
\sphinxstylestrong{静态位移} .................. 2.7.2, 2.8.4

\end{itemize}


\subsubsection{K}
\label{\detokenize{chapters/chapter2:k}}\begin{itemize}
\item {} 
\sphinxAtStartPar
\sphinxstylestrong{抗干扰} .................... 2.5.5

\end{itemize}


\subsubsection{L}
\label{\detokenize{chapters/chapter2:l}}\begin{itemize}
\item {} 
\sphinxAtStartPar
\sphinxstylestrong{离散傅里叶变换} ............ 2.4.1

\end{itemize}


\subsubsection{M}
\label{\detokenize{chapters/chapter2:id114}}\begin{itemize}
\item {} 
\sphinxAtStartPar
\sphinxstylestrong{MTpy} ...................... 附录A

\item {} 
\sphinxAtStartPar
\sphinxstylestrong{马氏距离} .................. 2.6.4, 2.8.2

\item {} 
\sphinxAtStartPar
\sphinxstylestrong{敏感度分析} ................ 2.8.5

\end{itemize}


\subsubsection{N}
\label{\detokenize{chapters/chapter2:n}}\begin{itemize}
\item {} 
\sphinxAtStartPar
\sphinxstylestrong{拟合} ...................... 2.7.2

\end{itemize}


\subsubsection{O}
\label{\detokenize{chapters/chapter2:o}}\begin{itemize}
\item {} 
\sphinxAtStartPar
\sphinxstylestrong{耦合} ...................... 2.2.3

\end{itemize}


\subsubsection{P}
\label{\detokenize{chapters/chapter2:p}}\begin{itemize}
\item {} 
\sphinxAtStartPar
\sphinxstylestrong{偏相干度} .................. 2.6.2

\item {} 
\sphinxAtStartPar
\sphinxstylestrong{频率} ...................... 2.2.2

\item {} 
\sphinxAtStartPar
\sphinxstylestrong{频谱} ...................... 2.4

\end{itemize}


\subsubsection{Q}
\label{\detokenize{chapters/chapter2:q}}\begin{itemize}
\item {} 
\sphinxAtStartPar
\sphinxstylestrong{倾子} ...................... 2.2.3

\item {} 
\sphinxAtStartPar
\sphinxstylestrong{趋肤深度} .................. 2.2.2

\end{itemize}


\subsubsection{R}
\label{\detokenize{chapters/chapter2:r}}\begin{itemize}
\item {} 
\sphinxAtStartPar
\sphinxstylestrong{Rhoplus} ................... 2.7.2

\item {} 
\sphinxAtStartPar
\sphinxstylestrong{RMSMR} ..................... 2.8

\item {} 
\sphinxAtStartPar
\sphinxstylestrong{人工筛选} .................. 2.6.7

\end{itemize}


\subsubsection{S}
\label{\detokenize{chapters/chapter2:s}}\begin{itemize}
\item {} 
\sphinxAtStartPar
\sphinxstylestrong{死频带} .................... 2.2.1, 2.8.4

\item {} 
\sphinxAtStartPar
\sphinxstylestrong{时间序列} .................. 2.4

\item {} 
\sphinxAtStartPar
\sphinxstylestrong{视电阻率} .................. 2.2.3

\item {} 
\sphinxAtStartPar
\sphinxstylestrong{相位} ...................... 2.2.3

\item {} 
\sphinxAtStartPar
\sphinxstylestrong{相位张量} .................. 2.7.1

\item {} 
\sphinxAtStartPar
\sphinxstylestrong{稳健估计} .................. 2.5.5

\end{itemize}


\subsubsection{T}
\label{\detokenize{chapters/chapter2:t}}\begin{itemize}
\item {} 
\sphinxAtStartPar
\sphinxstylestrong{Tipper} .................... 见：倾子

\item {} 
\sphinxAtStartPar
\sphinxstylestrong{探测深度} .................. 2.2.2

\end{itemize}


\subsubsection{W}
\label{\detokenize{chapters/chapter2:w}}\begin{itemize}
\item {} 
\sphinxAtStartPar
\sphinxstylestrong{误差棒} .................... 2.6.3

\end{itemize}


\subsubsection{Y}
\label{\detokenize{chapters/chapter2:y}}\begin{itemize}
\item {} 
\sphinxAtStartPar
\sphinxstylestrong{远参考道} .................. 2.3.2, 2.5.4

\item {} 
\sphinxAtStartPar
\sphinxstylestrong{阈值} ...................... 2.8.2

\end{itemize}


\subsubsection{Z}
\label{\detokenize{chapters/chapter2:z}}\begin{itemize}
\item {} 
\sphinxAtStartPar
\sphinxstylestrong{噪声} ...................... 2.3, 2.6

\item {} 
\sphinxAtStartPar
\sphinxstylestrong{张量} ...................... 2.2.1

\item {} 
\sphinxAtStartPar
\sphinxstylestrong{阻抗} ...................... 2.2.3, 2.5

\item {} 
\sphinxAtStartPar
\sphinxstylestrong{阻抗估计} .................. 2.5

\end{itemize}


\bigskip\hrule\bigskip



\subsection{附录E：常见问题 FAQ}
\label{\detokenize{chapters/chapter2:e-faq}}

\subsubsection{Q1: 为什么我的视电阻率曲线在某个频段突然跳动？}
\label{\detokenize{chapters/chapter2:q1}}
\sphinxAtStartPar
\sphinxstylestrong{A}: 可能是死频带(1\sphinxhyphen{}5kHz)噪声干扰。建议按以下步骤处理：
\begin{enumerate}
\sphinxsetlistlabels{\arabic}{enumi}{enumii}{}{.}%
\item {} 
\sphinxAtStartPar
检查该频段相干度是否<0.5

\item {} 
\sphinxAtStartPar
如果是日间数据，死频带问题更常见

\item {} 
\sphinxAtStartPar
使用 \sphinxstylestrong{Rhoplus死频带校正}功能

\item {} 
\sphinxAtStartPar
如果校正效果不好，考虑夜间重新采集

\end{enumerate}


\bigskip\hrule\bigskip



\subsubsection{Q2: 远参考站距离多远合适？}
\label{\detokenize{chapters/chapter2:q2}}
\sphinxAtStartPar
\sphinxstylestrong{A}: 一般建议 \sphinxstylestrong{5\sphinxhyphen{}50km}。具体参考：


\begin{savenotes}\sphinxattablestart
\sphinxthistablewithglobalstyle
\centering
\begin{tabulary}{\linewidth}[t]{TTT}
\sphinxtoprule
\begin{varwidth}[t]{\sphinxcolwidth{1}{3}}
\sphinxstyletheadfamily \sphinxAtStartPar
距离
\sphinxbeforeendvarwidth
\end{varwidth}%
&\begin{varwidth}[t]{\sphinxcolwidth{1}{3}}
\sphinxstyletheadfamily \sphinxAtStartPar
优点
\sphinxbeforeendvarwidth
\end{varwidth}%
&\begin{varwidth}[t]{\sphinxcolwidth{1}{3}}
\sphinxstyletheadfamily \sphinxAtStartPar
缺点
\sphinxbeforeendvarwidth
\end{varwidth}%
\\
\sphinxmidrule
\sphinxtableatstartofbodyhook\begin{varwidth}[t]{\sphinxcolwidth{1}{3}}
\sphinxAtStartPar
< 3km
\sphinxbeforeendvarwidth
\end{varwidth}%
&\begin{varwidth}[t]{\sphinxcolwidth{1}{3}}
\sphinxAtStartPar
信号相关性强
\sphinxbeforeendvarwidth
\end{varwidth}%
&\begin{varwidth}[t]{\sphinxcolwidth{1}{3}}
\sphinxAtStartPar
噪声可能相关
\sphinxbeforeendvarwidth
\end{varwidth}%
\\
\sphinxhline\begin{varwidth}[t]{\sphinxcolwidth{1}{3}}
\sphinxAtStartPar
3\sphinxhyphen{}10km
\sphinxbeforeendvarwidth
\end{varwidth}%
&\begin{varwidth}[t]{\sphinxcolwidth{1}{3}}
\sphinxAtStartPar
平衡
\sphinxbeforeendvarwidth
\end{varwidth}%
&\begin{varwidth}[t]{\sphinxcolwidth{1}{3}}
\sphinxAtStartPar
一般推荐
\sphinxbeforeendvarwidth
\end{varwidth}%
\\
\sphinxhline\begin{varwidth}[t]{\sphinxcolwidth{1}{3}}
\sphinxAtStartPar
10\sphinxhyphen{}50km
\sphinxbeforeendvarwidth
\end{varwidth}%
&\begin{varwidth}[t]{\sphinxcolwidth{1}{3}}
\sphinxAtStartPar
噪声不相关
\sphinxbeforeendvarwidth
\end{varwidth}%
&\begin{varwidth}[t]{\sphinxcolwidth{1}{3}}
\sphinxAtStartPar
\sphinxstylestrong{最佳范围}
\sphinxbeforeendvarwidth
\end{varwidth}%
\\
\sphinxhline\begin{varwidth}[t]{\sphinxcolwidth{1}{3}}
\sphinxAtStartPar
> 100km
\sphinxbeforeendvarwidth
\end{varwidth}%
&\begin{varwidth}[t]{\sphinxcolwidth{1}{3}}
\sphinxAtStartPar
噪声完全不相关
\sphinxbeforeendvarwidth
\end{varwidth}%
&\begin{varwidth}[t]{\sphinxcolwidth{1}{3}}
\sphinxAtStartPar
信号相关性降低
\sphinxbeforeendvarwidth
\end{varwidth}%
\\
\sphinxbottomrule
\end{tabulary}
\sphinxtableafterendhook\par
\sphinxattableend\end{savenotes}


\bigskip\hrule\bigskip



\subsubsection{Q3: RMSMR迭代多少次合适？}
\label{\detokenize{chapters/chapter2:q3-rmsmr}}
\sphinxAtStartPar
\sphinxstylestrong{A}: 通常 \sphinxstylestrong{5\sphinxhyphen{}10次}。判断标准：
\begin{itemize}
\item {} 
\sphinxAtStartPar
\sphinxstylestrong{5次后质量达标} → 数据质量好，处理完成

\item {} 
\sphinxAtStartPar
\sphinxstylestrong{10次后仍不达标} → 检查原始数据质量

\item {} 
\sphinxAtStartPar
\sphinxstylestrong{每次迭代数据量减少>20\%} → 参数过严，需要放宽阈值

\end{itemize}


\bigskip\hrule\bigskip



\subsubsection{Q4: 相干度一直很低怎么办？}
\label{\detokenize{chapters/chapter2:q4}}
\sphinxAtStartPar
\sphinxstylestrong{A}: 按以下优先级排查：
\begin{enumerate}
\sphinxsetlistlabels{\arabic}{enumi}{enumii}{}{.}%
\item {} 
\sphinxAtStartPar
\sphinxstylestrong{检查原始数据}：是否有明显的干扰时段？

\item {} 
\sphinxAtStartPar
\sphinxstylestrong{检查远参考}：远参考数据是否正常？

\item {} 
\sphinxAtStartPar
\sphinxstylestrong{调整阈值}：尝试降低相干度阈值（0.7→0.6）

\item {} 
\sphinxAtStartPar
\sphinxstylestrong{增加时长}：是否观测时间太短？

\item {} 
\sphinxAtStartPar
\sphinxstylestrong{重新选址}：如果周围有持续性干扰源，可能需要更换测站位置

\end{enumerate}


\bigskip\hrule\bigskip



\subsubsection{Q5: 如何判断数据是否可以用于反演？}
\label{\detokenize{chapters/chapter2:q5}}
\sphinxAtStartPar
\sphinxstylestrong{A}: 使用以下检查清单：
\begin{itemize}
\item {} 
\sphinxAtStartPar
 视电阻率曲线光滑连续，无脱节

\item {} 
\sphinxAtStartPar
 相位在30°\sphinxhyphen{}60°范围（TE模式）

\item {} 
\sphinxAtStartPar
 常相干度 > 0.7（大部分频点）

\item {} 
\sphinxAtStartPar
 误差棒 < 30\%（大部分频点）

\item {} 
\sphinxAtStartPar
 相位张量二维偏离度 |β| < 5°

\item {} 
\sphinxAtStartPar
 没有明显的静态位移

\end{itemize}

\sphinxAtStartPar
如果以上6项中5项以上满足，数据通常可以用于反演。


\bigskip\hrule\bigskip



\subsubsection{Q6: 不同仪器格式的数据如何导入？}
\label{\detokenize{chapters/chapter2:q6}}
\sphinxAtStartPar
\sphinxstylestrong{A}: MTDP支持多种仪器格式：


\begin{savenotes}\sphinxattablestart
\sphinxthistablewithglobalstyle
\centering
\begin{tabulary}{\linewidth}[t]{TTT}
\sphinxtoprule
\begin{varwidth}[t]{\sphinxcolwidth{1}{3}}
\sphinxstyletheadfamily \sphinxAtStartPar
仪器品牌
\sphinxbeforeendvarwidth
\end{varwidth}%
&\begin{varwidth}[t]{\sphinxcolwidth{1}{3}}
\sphinxstyletheadfamily \sphinxAtStartPar
导入方式
\sphinxbeforeendvarwidth
\end{varwidth}%
&\begin{varwidth}[t]{\sphinxcolwidth{1}{3}}
\sphinxstyletheadfamily \sphinxAtStartPar
注意事项
\sphinxbeforeendvarwidth
\end{varwidth}%
\\
\sphinxmidrule
\sphinxtableatstartofbodyhook\begin{varwidth}[t]{\sphinxcolwidth{1}{3}}
\sphinxAtStartPar
\sphinxstylestrong{Phoenix}
\sphinxbeforeendvarwidth
\end{varwidth}%
&\begin{varwidth}[t]{\sphinxcolwidth{1}{3}}
\sphinxAtStartPar
自动识别
\sphinxbeforeendvarwidth
\end{varwidth}%
&\begin{varwidth}[t]{\sphinxcolwidth{1}{3}}
\sphinxAtStartPar
需要完整的标定文件
\sphinxbeforeendvarwidth
\end{varwidth}%
\\
\sphinxhline\begin{varwidth}[t]{\sphinxcolwidth{1}{3}}
\sphinxAtStartPar
\sphinxstylestrong{Metronix}
\sphinxbeforeendvarwidth
\end{varwidth}%
&\begin{varwidth}[t]{\sphinxcolwidth{1}{3}}
\sphinxAtStartPar
自动识别
\sphinxbeforeendvarwidth
\end{varwidth}%
&\begin{varwidth}[t]{\sphinxcolwidth{1}{3}}
\sphinxAtStartPar
检查采样率设置
\sphinxbeforeendvarwidth
\end{varwidth}%
\\
\sphinxhline\begin{varwidth}[t]{\sphinxcolwidth{1}{3}}
\sphinxAtStartPar
\sphinxstylestrong{LEMI}
\sphinxbeforeendvarwidth
\end{varwidth}%
&\begin{varwidth}[t]{\sphinxcolwidth{1}{3}}
\sphinxAtStartPar
自动识别
\sphinxbeforeendvarwidth
\end{varwidth}%
&\begin{varwidth}[t]{\sphinxcolwidth{1}{3}}
\sphinxAtStartPar
注意时间格式
\sphinxbeforeendvarwidth
\end{varwidth}%
\\
\sphinxhline\begin{varwidth}[t]{\sphinxcolwidth{1}{3}}
\sphinxAtStartPar
\sphinxstylestrong{其他}
\sphinxbeforeendvarwidth
\end{varwidth}%
&\begin{varwidth}[t]{\sphinxcolwidth{1}{3}}
\sphinxAtStartPar
通用格式导入
\sphinxbeforeendvarwidth
\end{varwidth}%
&\begin{varwidth}[t]{\sphinxcolwidth{1}{3}}
\sphinxAtStartPar
需要ASCII/CSV格式
\sphinxbeforeendvarwidth
\end{varwidth}%
\\
\sphinxbottomrule
\end{tabulary}
\sphinxtableafterendhook\par
\sphinxattableend\end{savenotes}


\bigskip\hrule\bigskip



\subsubsection{Q7: 处理结果导出什么格式？}
\label{\detokenize{chapters/chapter2:q7}}
\sphinxAtStartPar
\sphinxstylestrong{A}: MTDP支持多种导出格式：


\begin{savenotes}\sphinxattablestart
\sphinxthistablewithglobalstyle
\centering
\begin{tabulary}{\linewidth}[t]{TTT}
\sphinxtoprule
\begin{varwidth}[t]{\sphinxcolwidth{1}{3}}
\sphinxstyletheadfamily \sphinxAtStartPar
格式
\sphinxbeforeendvarwidth
\end{varwidth}%
&\begin{varwidth}[t]{\sphinxcolwidth{1}{3}}
\sphinxstyletheadfamily \sphinxAtStartPar
用途
\sphinxbeforeendvarwidth
\end{varwidth}%
&\begin{varwidth}[t]{\sphinxcolwidth{1}{3}}
\sphinxstyletheadfamily \sphinxAtStartPar
说明
\sphinxbeforeendvarwidth
\end{varwidth}%
\\
\sphinxmidrule
\sphinxtableatstartofbodyhook\begin{varwidth}[t]{\sphinxcolwidth{1}{3}}
\sphinxAtStartPar
\sphinxstylestrong{EDI}
\sphinxbeforeendvarwidth
\end{varwidth}%
&\begin{varwidth}[t]{\sphinxcolwidth{1}{3}}
\sphinxAtStartPar
标准交换格式
\sphinxbeforeendvarwidth
\end{varwidth}%
&\begin{varwidth}[t]{\sphinxcolwidth{1}{3}}
\sphinxAtStartPar
推荐用于反演软件
\sphinxbeforeendvarwidth
\end{varwidth}%
\\
\sphinxhline\begin{varwidth}[t]{\sphinxcolwidth{1}{3}}
\sphinxAtStartPar
\sphinxstylestrong{J\sphinxhyphen{}file}
\sphinxbeforeendvarwidth
\end{varwidth}%
&\begin{varwidth}[t]{\sphinxcolwidth{1}{3}}
\sphinxAtStartPar
BIRRP格式
\sphinxbeforeendvarwidth
\end{varwidth}%
&\begin{varwidth}[t]{\sphinxcolwidth{1}{3}}
\sphinxAtStartPar
用于进一步处理
\sphinxbeforeendvarwidth
\end{varwidth}%
\\
\sphinxhline\begin{varwidth}[t]{\sphinxcolwidth{1}{3}}
\sphinxAtStartPar
\sphinxstylestrong{CSV}
\sphinxbeforeendvarwidth
\end{varwidth}%
&\begin{varwidth}[t]{\sphinxcolwidth{1}{3}}
\sphinxAtStartPar
通用格式
\sphinxbeforeendvarwidth
\end{varwidth}%
&\begin{varwidth}[t]{\sphinxcolwidth{1}{3}}
\sphinxAtStartPar
用于自定义分析
\sphinxbeforeendvarwidth
\end{varwidth}%
\\
\sphinxhline\begin{varwidth}[t]{\sphinxcolwidth{1}{3}}
\sphinxAtStartPar
\sphinxstylestrong{PNG/PDF}
\sphinxbeforeendvarwidth
\end{varwidth}%
&\begin{varwidth}[t]{\sphinxcolwidth{1}{3}}
\sphinxAtStartPar
图形导出
\sphinxbeforeendvarwidth
\end{varwidth}%
&\begin{varwidth}[t]{\sphinxcolwidth{1}{3}}
\sphinxAtStartPar
用于报告
\sphinxbeforeendvarwidth
\end{varwidth}%
\\
\sphinxbottomrule
\end{tabulary}
\sphinxtableafterendhook\par
\sphinxattableend\end{savenotes}


\bigskip\hrule\bigskip

\begin{quote}

\sphinxAtStartPar
\sphinxstylestrong{下一章预告}：第3章将介绍MTDP的数据工程管理功能，包括项目组织、数据导入导出等实用操作。
\end{quote}

\sphinxstepscope


\section{📂 第3章 数据工程管理}
\label{\detokenize{chapters/chapter3:id1}}\label{\detokenize{chapters/chapter3::doc}}

\subsection{3.1 工程结构概述}
\label{\detokenize{chapters/chapter3:id2}}
\sphinxAtStartPar
MTDP采用四级层次结构管理数据：

\begin{sphinxVerbatim}[commandchars=\\\{\}]
📦 工程 (Project) → 🗺️ 工区 (WorkArea) → 📋 测段 (Group) → 📍 测点 (Site)
\end{sphinxVerbatim}


\subsubsection{3.1.1 各层级说明}
\label{\detokenize{chapters/chapter3:id3}}

\begin{savenotes}\sphinxattablestart
\sphinxthistablewithglobalstyle
\centering
\begin{tabulary}{\linewidth}[t]{TTT}
\sphinxtoprule
\begin{varwidth}[t]{\sphinxcolwidth{1}{3}}
\sphinxstyletheadfamily \sphinxAtStartPar
层级
\sphinxbeforeendvarwidth
\end{varwidth}%
&\begin{varwidth}[t]{\sphinxcolwidth{1}{3}}
\sphinxstyletheadfamily \sphinxAtStartPar
说明
\sphinxbeforeendvarwidth
\end{varwidth}%
&\begin{varwidth}[t]{\sphinxcolwidth{1}{3}}
\sphinxstyletheadfamily \sphinxAtStartPar
典型用途
\sphinxbeforeendvarwidth
\end{varwidth}%
\\
\sphinxmidrule
\sphinxtableatstartofbodyhook\begin{varwidth}[t]{\sphinxcolwidth{1}{3}}
\sphinxAtStartPar
工程
\sphinxbeforeendvarwidth
\end{varwidth}%
&\begin{varwidth}[t]{\sphinxcolwidth{1}{3}}
\sphinxAtStartPar
最顶层容器
\sphinxbeforeendvarwidth
\end{varwidth}%
&\begin{varwidth}[t]{\sphinxcolwidth{1}{3}}
\sphinxAtStartPar
整个勘探项目
\sphinxbeforeendvarwidth
\end{varwidth}%
\\
\sphinxhline\begin{varwidth}[t]{\sphinxcolwidth{1}{3}}
\sphinxAtStartPar
工区
\sphinxbeforeendvarwidth
\end{varwidth}%
&\begin{varwidth}[t]{\sphinxcolwidth{1}{3}}
\sphinxAtStartPar
按区域划分
\sphinxbeforeendvarwidth
\end{varwidth}%
&\begin{varwidth}[t]{\sphinxcolwidth{1}{3}}
\sphinxAtStartPar
同一区域的测点集合
\sphinxbeforeendvarwidth
\end{varwidth}%
\\
\sphinxhline\begin{varwidth}[t]{\sphinxcolwidth{1}{3}}
\sphinxAtStartPar
测段
\sphinxbeforeendvarwidth
\end{varwidth}%
&\begin{varwidth}[t]{\sphinxcolwidth{1}{3}}
\sphinxAtStartPar
按时间/任务划分
\sphinxbeforeendvarwidth
\end{varwidth}%
&\begin{varwidth}[t]{\sphinxcolwidth{1}{3}}
\sphinxAtStartPar
同一批处理的测点
\sphinxbeforeendvarwidth
\end{varwidth}%
\\
\sphinxhline\begin{varwidth}[t]{\sphinxcolwidth{1}{3}}
\sphinxAtStartPar
测点
\sphinxbeforeendvarwidth
\end{varwidth}%
&\begin{varwidth}[t]{\sphinxcolwidth{1}{3}}
\sphinxAtStartPar
最小数据单元
\sphinxbeforeendvarwidth
\end{varwidth}%
&\begin{varwidth}[t]{\sphinxcolwidth{1}{3}}
\sphinxAtStartPar
单个观测点数据
\sphinxbeforeendvarwidth
\end{varwidth}%
\\
\sphinxbottomrule
\end{tabulary}
\sphinxtableafterendhook\par
\sphinxattableend\end{savenotes}


\subsubsection{3.1.2 测点类型}
\label{\detokenize{chapters/chapter3:id4}}
\sphinxAtStartPar
MTDP支持多种测点类型：


\begin{savenotes}\sphinxattablestart
\sphinxthistablewithglobalstyle
\centering
\begin{tabulary}{\linewidth}[t]{TTT}
\sphinxtoprule
\begin{varwidth}[t]{\sphinxcolwidth{1}{3}}
\sphinxstyletheadfamily \sphinxAtStartPar
类型
\sphinxbeforeendvarwidth
\end{varwidth}%
&\begin{varwidth}[t]{\sphinxcolwidth{1}{3}}
\sphinxstyletheadfamily \sphinxAtStartPar
📻 仪器
\sphinxbeforeendvarwidth
\end{varwidth}%
&\begin{varwidth}[t]{\sphinxcolwidth{1}{3}}
\sphinxstyletheadfamily \sphinxAtStartPar
📄 格式
\sphinxbeforeendvarwidth
\end{varwidth}%
\\
\sphinxmidrule
\sphinxtableatstartofbodyhook\begin{varwidth}[t]{\sphinxcolwidth{1}{3}}
\sphinxAtStartPar
PhoenixSite
\sphinxbeforeendvarwidth
\end{varwidth}%
&\begin{varwidth}[t]{\sphinxcolwidth{1}{3}}
\sphinxAtStartPar
Phoenix MTU\sphinxhyphen{}5/5A
\sphinxbeforeendvarwidth
\end{varwidth}%
&\begin{varwidth}[t]{\sphinxcolwidth{1}{3}}
\sphinxAtStartPar
TBL格式数据
\sphinxbeforeendvarwidth
\end{varwidth}%
\\
\sphinxhline\begin{varwidth}[t]{\sphinxcolwidth{1}{3}}
\sphinxAtStartPar
MTUSite
\sphinxbeforeendvarwidth
\end{varwidth}%
&\begin{varwidth}[t]{\sphinxcolwidth{1}{3}}
\sphinxAtStartPar
Phoenix MTU\sphinxhyphen{}5C/8A
\sphinxbeforeendvarwidth
\end{varwidth}%
&\begin{varwidth}[t]{\sphinxcolwidth{1}{3}}
\sphinxAtStartPar
JSON格式数据
\sphinxbeforeendvarwidth
\end{varwidth}%
\\
\sphinxhline\begin{varwidth}[t]{\sphinxcolwidth{1}{3}}
\sphinxAtStartPar
MetronixSite
\sphinxbeforeendvarwidth
\end{varwidth}%
&\begin{varwidth}[t]{\sphinxcolwidth{1}{3}}
\sphinxAtStartPar
Metronix仪器
\sphinxbeforeendvarwidth
\end{varwidth}%
&\begin{varwidth}[t]{\sphinxcolwidth{1}{3}}
\sphinxAtStartPar
ATM格式数据
\sphinxbeforeendvarwidth
\end{varwidth}%
\\
\sphinxhline\begin{varwidth}[t]{\sphinxcolwidth{1}{3}}
\sphinxAtStartPar
LEMISite
\sphinxbeforeendvarwidth
\end{varwidth}%
&\begin{varwidth}[t]{\sphinxcolwidth{1}{3}}
\sphinxAtStartPar
LEMI长周期
\sphinxbeforeendvarwidth
\end{varwidth}%
&\begin{varwidth}[t]{\sphinxcolwidth{1}{3}}
\sphinxAtStartPar
长周期数据
\sphinxbeforeendvarwidth
\end{varwidth}%
\\
\sphinxhline\begin{varwidth}[t]{\sphinxcolwidth{1}{3}}
\sphinxAtStartPar
RMTSite
\sphinxbeforeendvarwidth
\end{varwidth}%
&\begin{varwidth}[t]{\sphinxcolwidth{1}{3}}
\sphinxAtStartPar
RMT格式
\sphinxbeforeendvarwidth
\end{varwidth}%
&\begin{varwidth}[t]{\sphinxcolwidth{1}{3}}
\sphinxAtStartPar
无线电MT
\sphinxbeforeendvarwidth
\end{varwidth}%
\\
\sphinxhline\begin{varwidth}[t]{\sphinxcolwidth{1}{3}}
\sphinxAtStartPar
SBFSite
\sphinxbeforeendvarwidth
\end{varwidth}%
&\begin{varwidth}[t]{\sphinxcolwidth{1}{3}}
\sphinxAtStartPar
SBF格式
\sphinxbeforeendvarwidth
\end{varwidth}%
&\begin{varwidth}[t]{\sphinxcolwidth{1}{3}}
\sphinxAtStartPar
RMT数据
\sphinxbeforeendvarwidth
\end{varwidth}%
\\
\sphinxhline\begin{varwidth}[t]{\sphinxcolwidth{1}{3}}
\sphinxAtStartPar
CASRMTSite
\sphinxbeforeendvarwidth
\end{varwidth}%
&\begin{varwidth}[t]{\sphinxcolwidth{1}{3}}
\sphinxAtStartPar
CASRMT格式
\sphinxbeforeendvarwidth
\end{varwidth}%
&\begin{varwidth}[t]{\sphinxcolwidth{1}{3}}
\sphinxAtStartPar
可控源AMT
\sphinxbeforeendvarwidth
\end{varwidth}%
\\
\sphinxhline\begin{varwidth}[t]{\sphinxcolwidth{1}{3}}
\sphinxAtStartPar
ATTSSite
\sphinxbeforeendvarwidth
\end{varwidth}%
&\begin{varwidth}[t]{\sphinxcolwidth{1}{3}}
\sphinxAtStartPar
Aether仪器
\sphinxbeforeendvarwidth
\end{varwidth}%
&\begin{varwidth}[t]{\sphinxcolwidth{1}{3}}
\sphinxAtStartPar
ATTS格式
\sphinxbeforeendvarwidth
\end{varwidth}%
\\
\sphinxhline\begin{varwidth}[t]{\sphinxcolwidth{1}{3}}
\sphinxAtStartPar
SyntheticSite
\sphinxbeforeendvarwidth
\end{varwidth}%
&\begin{varwidth}[t]{\sphinxcolwidth{1}{3}}
\sphinxAtStartPar
合成数据
\sphinxbeforeendvarwidth
\end{varwidth}%
&\begin{varwidth}[t]{\sphinxcolwidth{1}{3}}
\sphinxAtStartPar
正演模拟
\sphinxbeforeendvarwidth
\end{varwidth}%
\\
\sphinxhline\begin{varwidth}[t]{\sphinxcolwidth{1}{3}}
\sphinxAtStartPar
EDISite
\sphinxbeforeendvarwidth
\end{varwidth}%
&\begin{varwidth}[t]{\sphinxcolwidth{1}{3}}
\sphinxAtStartPar
EDI格式
\sphinxbeforeendvarwidth
\end{varwidth}%
&\begin{varwidth}[t]{\sphinxcolwidth{1}{3}}
\sphinxAtStartPar
已处理数据
\sphinxbeforeendvarwidth
\end{varwidth}%
\\
\sphinxbottomrule
\end{tabulary}
\sphinxtableafterendhook\par
\sphinxattableend\end{savenotes}


\bigskip\hrule\bigskip



\subsection{3.2 创建与管理工程}
\label{\detokenize{chapters/chapter3:id5}}

\subsubsection{3.2.1 新建工程}
\label{\detokenize{chapters/chapter3:id6}}
\sphinxAtStartPar
1️⃣ 选择菜单 \sphinxcode{\sphinxupquote{文件 → 新建工程}}
2️⃣ 在弹出的对话框中设置：
\begin{itemize}
\item {} 
\sphinxAtStartPar
📝 工程名称

\item {} 
\sphinxAtStartPar
📁 保存位置

\item {} 
\sphinxAtStartPar
📖 工程描述（可选）
3️⃣ 点击"确定"创建

\end{itemize}


\subsubsection{3.2.2 打开已有工程}
\label{\detokenize{chapters/chapter3:id7}}
\sphinxAtStartPar
\sphinxstylestrong{📄 方式1：直接打开}
\begin{itemize}
\item {} 
\sphinxAtStartPar
选择 \sphinxcode{\sphinxupquote{文件 → 打开工程}}

\item {} 
\sphinxAtStartPar
选择 .MTDPE 文件

\end{itemize}

\sphinxAtStartPar
\sphinxstylestrong{🕐 方式2：最近文件}
\begin{itemize}
\item {} 
\sphinxAtStartPar
选择 \sphinxcode{\sphinxupquote{文件 → 重新打开}}

\item {} 
\sphinxAtStartPar
从列表中选择

\end{itemize}


\subsubsection{3.2.3 工程属性设置 (ProjectSetForm)}
\label{\detokenize{chapters/chapter3:projectsetform}}
\sphinxAtStartPar
打开工程设置窗口可以配置以下属性：

\sphinxAtStartPar
\sphinxstylestrong{📋 基本信息：}


\begin{savenotes}\sphinxattablestart
\sphinxthistablewithglobalstyle
\centering
\begin{tabulary}{\linewidth}[t]{TT}
\sphinxtoprule
\begin{varwidth}[t]{\sphinxcolwidth{1}{2}}
\sphinxstyletheadfamily \sphinxAtStartPar
字段
\sphinxbeforeendvarwidth
\end{varwidth}%
&\begin{varwidth}[t]{\sphinxcolwidth{1}{2}}
\sphinxstyletheadfamily \sphinxAtStartPar
说明
\sphinxbeforeendvarwidth
\end{varwidth}%
\\
\sphinxmidrule
\sphinxtableatstartofbodyhook\begin{varwidth}[t]{\sphinxcolwidth{1}{2}}
\sphinxAtStartPar
Name
\sphinxbeforeendvarwidth
\end{varwidth}%
&\begin{varwidth}[t]{\sphinxcolwidth{1}{2}}
\sphinxAtStartPar
工程名称
\sphinxbeforeendvarwidth
\end{varwidth}%
\\
\sphinxhline\begin{varwidth}[t]{\sphinxcolwidth{1}{2}}
\sphinxAtStartPar
Owner
\sphinxbeforeendvarwidth
\end{varwidth}%
&\begin{varwidth}[t]{\sphinxcolwidth{1}{2}}
\sphinxAtStartPar
所有者
\sphinxbeforeendvarwidth
\end{varwidth}%
\\
\sphinxhline\begin{varwidth}[t]{\sphinxcolwidth{1}{2}}
\sphinxAtStartPar
Creator
\sphinxbeforeendvarwidth
\end{varwidth}%
&\begin{varwidth}[t]{\sphinxcolwidth{1}{2}}
\sphinxAtStartPar
创建者
\sphinxbeforeendvarwidth
\end{varwidth}%
\\
\sphinxhline\begin{varwidth}[t]{\sphinxcolwidth{1}{2}}
\sphinxAtStartPar
Editor
\sphinxbeforeendvarwidth
\end{varwidth}%
&\begin{varwidth}[t]{\sphinxcolwidth{1}{2}}
\sphinxAtStartPar
编辑者
\sphinxbeforeendvarwidth
\end{varwidth}%
\\
\sphinxhline\begin{varwidth}[t]{\sphinxcolwidth{1}{2}}
\sphinxAtStartPar
Purpose
\sphinxbeforeendvarwidth
\end{varwidth}%
&\begin{varwidth}[t]{\sphinxcolwidth{1}{2}}
\sphinxAtStartPar
工程目的
\sphinxbeforeendvarwidth
\end{varwidth}%
\\
\sphinxhline\begin{varwidth}[t]{\sphinxcolwidth{1}{2}}
\sphinxAtStartPar
Description
\sphinxbeforeendvarwidth
\end{varwidth}%
&\begin{varwidth}[t]{\sphinxcolwidth{1}{2}}
\sphinxAtStartPar
描述信息
\sphinxbeforeendvarwidth
\end{varwidth}%
\\
\sphinxhline\begin{varwidth}[t]{\sphinxcolwidth{1}{2}}
\sphinxAtStartPar
Version
\sphinxbeforeendvarwidth
\end{varwidth}%
&\begin{varwidth}[t]{\sphinxcolwidth{1}{2}}
\sphinxAtStartPar
版本号
\sphinxbeforeendvarwidth
\end{varwidth}%
\\
\sphinxbottomrule
\end{tabulary}
\sphinxtableafterendhook\par
\sphinxattableend\end{savenotes}

\sphinxAtStartPar
\sphinxstylestrong{📁 目录配置：}


\begin{savenotes}\sphinxattablestart
\sphinxthistablewithglobalstyle
\centering
\begin{tabulary}{\linewidth}[t]{TT}
\sphinxtoprule
\begin{varwidth}[t]{\sphinxcolwidth{1}{2}}
\sphinxstyletheadfamily \sphinxAtStartPar
字段
\sphinxbeforeendvarwidth
\end{varwidth}%
&\begin{varwidth}[t]{\sphinxcolwidth{1}{2}}
\sphinxstyletheadfamily \sphinxAtStartPar
说明
\sphinxbeforeendvarwidth
\end{varwidth}%
\\
\sphinxmidrule
\sphinxtableatstartofbodyhook\begin{varwidth}[t]{\sphinxcolwidth{1}{2}}
\sphinxAtStartPar
Project Directory
\sphinxbeforeendvarwidth
\end{varwidth}%
&\begin{varwidth}[t]{\sphinxcolwidth{1}{2}}
\sphinxAtStartPar
工程目录
\sphinxbeforeendvarwidth
\end{varwidth}%
\\
\sphinxhline\begin{varwidth}[t]{\sphinxcolwidth{1}{2}}
\sphinxAtStartPar
Data Directory
\sphinxbeforeendvarwidth
\end{varwidth}%
&\begin{varwidth}[t]{\sphinxcolwidth{1}{2}}
\sphinxAtStartPar
数据目录
\sphinxbeforeendvarwidth
\end{varwidth}%
\\
\sphinxbottomrule
\end{tabulary}
\sphinxtableafterendhook\par
\sphinxattableend\end{savenotes}

\sphinxAtStartPar
\sphinxstylestrong{⏰ 时间信息（自动记录）：}
\begin{itemize}
\item {} 
\sphinxAtStartPar
创建时间

\item {} 
\sphinxAtStartPar
最后编辑时间

\end{itemize}
\begin{quote}

\sphinxAtStartPar
💡 \sphinxstylestrong{操作方法}：右键工程 → \sphinxcode{\sphinxupquote{工程设置}}
\end{quote}


\subsubsection{3.2.4 人员信息管理 (PersonSetForm)}
\label{\detokenize{chapters/chapter3:personsetform}}
\sphinxAtStartPar
管理项目中参与人员的信息，便于数据溯源和协作。

\sphinxAtStartPar
\sphinxstylestrong{👤 人员字段：}


\begin{savenotes}\sphinxattablestart
\sphinxthistablewithglobalstyle
\centering
\begin{tabulary}{\linewidth}[t]{TT}
\sphinxtoprule
\begin{varwidth}[t]{\sphinxcolwidth{1}{2}}
\sphinxstyletheadfamily \sphinxAtStartPar
字段
\sphinxbeforeendvarwidth
\end{varwidth}%
&\begin{varwidth}[t]{\sphinxcolwidth{1}{2}}
\sphinxstyletheadfamily \sphinxAtStartPar
说明
\sphinxbeforeendvarwidth
\end{varwidth}%
\\
\sphinxmidrule
\sphinxtableatstartofbodyhook\begin{varwidth}[t]{\sphinxcolwidth{1}{2}}
\sphinxAtStartPar
Name
\sphinxbeforeendvarwidth
\end{varwidth}%
&\begin{varwidth}[t]{\sphinxcolwidth{1}{2}}
\sphinxAtStartPar
姓名
\sphinxbeforeendvarwidth
\end{varwidth}%
\\
\sphinxhline\begin{varwidth}[t]{\sphinxcolwidth{1}{2}}
\sphinxAtStartPar
Phone
\sphinxbeforeendvarwidth
\end{varwidth}%
&\begin{varwidth}[t]{\sphinxcolwidth{1}{2}}
\sphinxAtStartPar
联系电话
\sphinxbeforeendvarwidth
\end{varwidth}%
\\
\sphinxhline\begin{varwidth}[t]{\sphinxcolwidth{1}{2}}
\sphinxAtStartPar
Email
\sphinxbeforeendvarwidth
\end{varwidth}%
&\begin{varwidth}[t]{\sphinxcolwidth{1}{2}}
\sphinxAtStartPar
电子邮箱
\sphinxbeforeendvarwidth
\end{varwidth}%
\\
\sphinxhline\begin{varwidth}[t]{\sphinxcolwidth{1}{2}}
\sphinxAtStartPar
Institution
\sphinxbeforeendvarwidth
\end{varwidth}%
&\begin{varwidth}[t]{\sphinxcolwidth{1}{2}}
\sphinxAtStartPar
所属机构
\sphinxbeforeendvarwidth
\end{varwidth}%
\\
\sphinxhline\begin{varwidth}[t]{\sphinxcolwidth{1}{2}}
\sphinxAtStartPar
Position
\sphinxbeforeendvarwidth
\end{varwidth}%
&\begin{varwidth}[t]{\sphinxcolwidth{1}{2}}
\sphinxAtStartPar
职位
\sphinxbeforeendvarwidth
\end{varwidth}%
\\
\sphinxhline\begin{varwidth}[t]{\sphinxcolwidth{1}{2}}
\sphinxAtStartPar
Description
\sphinxbeforeendvarwidth
\end{varwidth}%
&\begin{varwidth}[t]{\sphinxcolwidth{1}{2}}
\sphinxAtStartPar
备注
\sphinxbeforeendvarwidth
\end{varwidth}%
\\
\sphinxbottomrule
\end{tabulary}
\sphinxtableafterendhook\par
\sphinxattableend\end{savenotes}
\begin{quote}

\sphinxAtStartPar
\sphinxstylestrong{操作方法}：在工程设置或工区设置中点击人员字段进行编辑
\end{quote}


\bigskip\hrule\bigskip



\subsection{3.3 工区管理}
\label{\detokenize{chapters/chapter3:id8}}

\subsubsection{3.3.1 新建工区}
\label{\detokenize{chapters/chapter3:id9}}\begin{enumerate}
\sphinxsetlistlabels{\arabic}{enumi}{enumii}{}{.}%
\item {} 
\sphinxAtStartPar
右键工程 → \sphinxcode{\sphinxupquote{新建工区}}

\item {} 
\sphinxAtStartPar
设置工区属性：
\begin{itemize}
\item {} 
\sphinxAtStartPar
名称

\item {} 
\sphinxAtStartPar
描述

\item {} 
\sphinxAtStartPar
坐标范围（可选）

\end{itemize}

\end{enumerate}


\subsubsection{3.3.2 工区属性设置 (WorkAreaSetForm)}
\label{\detokenize{chapters/chapter3:workareasetform}}
\sphinxAtStartPar
打开工区设置窗口可以配置以下属性：

\sphinxAtStartPar
\sphinxstylestrong{📋 基本信息：}


\begin{savenotes}\sphinxattablestart
\sphinxthistablewithglobalstyle
\centering
\begin{tabulary}{\linewidth}[t]{TT}
\sphinxtoprule
\begin{varwidth}[t]{\sphinxcolwidth{1}{2}}
\sphinxstyletheadfamily \sphinxAtStartPar
字段
\sphinxbeforeendvarwidth
\end{varwidth}%
&\begin{varwidth}[t]{\sphinxcolwidth{1}{2}}
\sphinxstyletheadfamily \sphinxAtStartPar
说明
\sphinxbeforeendvarwidth
\end{varwidth}%
\\
\sphinxmidrule
\sphinxtableatstartofbodyhook\begin{varwidth}[t]{\sphinxcolwidth{1}{2}}
\sphinxAtStartPar
Name
\sphinxbeforeendvarwidth
\end{varwidth}%
&\begin{varwidth}[t]{\sphinxcolwidth{1}{2}}
\sphinxAtStartPar
工区名称
\sphinxbeforeendvarwidth
\end{varwidth}%
\\
\sphinxhline\begin{varwidth}[t]{\sphinxcolwidth{1}{2}}
\sphinxAtStartPar
Description
\sphinxbeforeendvarwidth
\end{varwidth}%
&\begin{varwidth}[t]{\sphinxcolwidth{1}{2}}
\sphinxAtStartPar
描述信息
\sphinxbeforeendvarwidth
\end{varwidth}%
\\
\sphinxhline\begin{varwidth}[t]{\sphinxcolwidth{1}{2}}
\sphinxAtStartPar
Owner
\sphinxbeforeendvarwidth
\end{varwidth}%
&\begin{varwidth}[t]{\sphinxcolwidth{1}{2}}
\sphinxAtStartPar
所有者
\sphinxbeforeendvarwidth
\end{varwidth}%
\\
\sphinxhline\begin{varwidth}[t]{\sphinxcolwidth{1}{2}}
\sphinxAtStartPar
Surveyor
\sphinxbeforeendvarwidth
\end{varwidth}%
&\begin{varwidth}[t]{\sphinxcolwidth{1}{2}}
\sphinxAtStartPar
测量员
\sphinxbeforeendvarwidth
\end{varwidth}%
\\
\sphinxhline\begin{varwidth}[t]{\sphinxcolwidth{1}{2}}
\sphinxAtStartPar
DataCollector
\sphinxbeforeendvarwidth
\end{varwidth}%
&\begin{varwidth}[t]{\sphinxcolwidth{1}{2}}
\sphinxAtStartPar
数据采集者
\sphinxbeforeendvarwidth
\end{varwidth}%
\\
\sphinxhline\begin{varwidth}[t]{\sphinxcolwidth{1}{2}}
\sphinxAtStartPar
DataProcessor
\sphinxbeforeendvarwidth
\end{varwidth}%
&\begin{varwidth}[t]{\sphinxcolwidth{1}{2}}
\sphinxAtStartPar
数据处理者
\sphinxbeforeendvarwidth
\end{varwidth}%
\\
\sphinxbottomrule
\end{tabulary}
\sphinxtableafterendhook\par
\sphinxattableend\end{savenotes}

\sphinxAtStartPar
\sphinxstylestrong{🌍 坐标范围（可选）：}


\begin{savenotes}\sphinxattablestart
\sphinxthistablewithglobalstyle
\centering
\begin{tabulary}{\linewidth}[t]{TT}
\sphinxtoprule
\begin{varwidth}[t]{\sphinxcolwidth{1}{2}}
\sphinxstyletheadfamily \sphinxAtStartPar
字段
\sphinxbeforeendvarwidth
\end{varwidth}%
&\begin{varwidth}[t]{\sphinxcolwidth{1}{2}}
\sphinxstyletheadfamily \sphinxAtStartPar
说明
\sphinxbeforeendvarwidth
\end{varwidth}%
\\
\sphinxmidrule
\sphinxtableatstartofbodyhook\begin{varwidth}[t]{\sphinxcolwidth{1}{2}}
\sphinxAtStartPar
MinLongitude
\sphinxbeforeendvarwidth
\end{varwidth}%
&\begin{varwidth}[t]{\sphinxcolwidth{1}{2}}
\sphinxAtStartPar
最小经度
\sphinxbeforeendvarwidth
\end{varwidth}%
\\
\sphinxhline\begin{varwidth}[t]{\sphinxcolwidth{1}{2}}
\sphinxAtStartPar
MaxLongitude
\sphinxbeforeendvarwidth
\end{varwidth}%
&\begin{varwidth}[t]{\sphinxcolwidth{1}{2}}
\sphinxAtStartPar
最大经度
\sphinxbeforeendvarwidth
\end{varwidth}%
\\
\sphinxhline\begin{varwidth}[t]{\sphinxcolwidth{1}{2}}
\sphinxAtStartPar
MinLatitude
\sphinxbeforeendvarwidth
\end{varwidth}%
&\begin{varwidth}[t]{\sphinxcolwidth{1}{2}}
\sphinxAtStartPar
最小纬度
\sphinxbeforeendvarwidth
\end{varwidth}%
\\
\sphinxhline\begin{varwidth}[t]{\sphinxcolwidth{1}{2}}
\sphinxAtStartPar
MaxLatitude
\sphinxbeforeendvarwidth
\end{varwidth}%
&\begin{varwidth}[t]{\sphinxcolwidth{1}{2}}
\sphinxAtStartPar
最大纬度
\sphinxbeforeendvarwidth
\end{varwidth}%
\\
\sphinxbottomrule
\end{tabulary}
\sphinxtableafterendhook\par
\sphinxattableend\end{savenotes}
\begin{quote}

\sphinxAtStartPar
💡 \sphinxstylestrong{提示}：坐标范围可用于在地图上显示工区边界，方便数据管理和可视化。
\end{quote}

\sphinxAtStartPar
\sphinxstylestrong{📊 工区统计信息：}


\begin{savenotes}\sphinxattablestart
\sphinxthistablewithglobalstyle
\centering
\begin{tabulary}{\linewidth}[t]{TT}
\sphinxtoprule
\begin{varwidth}[t]{\sphinxcolwidth{1}{2}}
\sphinxstyletheadfamily \sphinxAtStartPar
属性
\sphinxbeforeendvarwidth
\end{varwidth}%
&\begin{varwidth}[t]{\sphinxcolwidth{1}{2}}
\sphinxstyletheadfamily \sphinxAtStartPar
说明
\sphinxbeforeendvarwidth
\end{varwidth}%
\\
\sphinxmidrule
\sphinxtableatstartofbodyhook\begin{varwidth}[t]{\sphinxcolwidth{1}{2}}
\sphinxAtStartPar
CreateTime
\sphinxbeforeendvarwidth
\end{varwidth}%
&\begin{varwidth}[t]{\sphinxcolwidth{1}{2}}
\sphinxAtStartPar
创建时间
\sphinxbeforeendvarwidth
\end{varwidth}%
\\
\sphinxhline\begin{varwidth}[t]{\sphinxcolwidth{1}{2}}
\sphinxAtStartPar
Groups
\sphinxbeforeendvarwidth
\end{varwidth}%
&\begin{varwidth}[t]{\sphinxcolwidth{1}{2}}
\sphinxAtStartPar
测段列表
\sphinxbeforeendvarwidth
\end{varwidth}%
\\
\sphinxhline\begin{varwidth}[t]{\sphinxcolwidth{1}{2}}
\sphinxAtStartPar
SiteCount
\sphinxbeforeendvarwidth
\end{varwidth}%
&\begin{varwidth}[t]{\sphinxcolwidth{1}{2}}
\sphinxAtStartPar
测点总数（自动计算）
\sphinxbeforeendvarwidth
\end{varwidth}%
\\
\sphinxbottomrule
\end{tabulary}
\sphinxtableafterendhook\par
\sphinxattableend\end{savenotes}


\subsubsection{3.3.3 工区操作}
\label{\detokenize{chapters/chapter3:id10}}\begin{itemize}
\item {} 
\sphinxAtStartPar
\sphinxstylestrong{编辑}：右键 → 工区设置

\item {} 
\sphinxAtStartPar
\sphinxstylestrong{复制}：右键 → 复制工区

\item {} 
\sphinxAtStartPar
\sphinxstylestrong{删除}：右键 → 删除工区

\item {} 
\sphinxAtStartPar
\sphinxstylestrong{导出坐标}：导出工区内所有测点坐标

\end{itemize}


\bigskip\hrule\bigskip



\subsection{3.4 测段管理}
\label{\detokenize{chapters/chapter3:id11}}

\subsubsection{3.4.1 新建测段}
\label{\detokenize{chapters/chapter3:id12}}\begin{enumerate}
\sphinxsetlistlabels{\arabic}{enumi}{enumii}{}{.}%
\item {} 
\sphinxAtStartPar
右键工区 → \sphinxcode{\sphinxupquote{新建测段}}

\item {} 
\sphinxAtStartPar
设置测段属性

\end{enumerate}


\subsubsection{3.4.2 测段属性设置 (GroupSetForm)}
\label{\detokenize{chapters/chapter3:groupsetform}}
\sphinxAtStartPar
打开测段设置窗口可以配置以下属性：

\sphinxAtStartPar
\sphinxstylestrong{📋 基本信息：}


\begin{savenotes}\sphinxattablestart
\sphinxthistablewithglobalstyle
\centering
\begin{tabulary}{\linewidth}[t]{TT}
\sphinxtoprule
\begin{varwidth}[t]{\sphinxcolwidth{1}{2}}
\sphinxstyletheadfamily \sphinxAtStartPar
字段
\sphinxbeforeendvarwidth
\end{varwidth}%
&\begin{varwidth}[t]{\sphinxcolwidth{1}{2}}
\sphinxstyletheadfamily \sphinxAtStartPar
说明
\sphinxbeforeendvarwidth
\end{varwidth}%
\\
\sphinxmidrule
\sphinxtableatstartofbodyhook\begin{varwidth}[t]{\sphinxcolwidth{1}{2}}
\sphinxAtStartPar
Name
\sphinxbeforeendvarwidth
\end{varwidth}%
&\begin{varwidth}[t]{\sphinxcolwidth{1}{2}}
\sphinxAtStartPar
测段名称
\sphinxbeforeendvarwidth
\end{varwidth}%
\\
\sphinxhline\begin{varwidth}[t]{\sphinxcolwidth{1}{2}}
\sphinxAtStartPar
Description
\sphinxbeforeendvarwidth
\end{varwidth}%
&\begin{varwidth}[t]{\sphinxcolwidth{1}{2}}
\sphinxAtStartPar
描述信息
\sphinxbeforeendvarwidth
\end{varwidth}%
\\
\sphinxbottomrule
\end{tabulary}
\sphinxtableafterendhook\par
\sphinxattableend\end{savenotes}

\sphinxAtStartPar
\sphinxstylestrong{⚙️ 处理方案：}

\sphinxAtStartPar
测段可以关联FFT处理方案，用于后续数据处理：


\begin{savenotes}\sphinxattablestart
\sphinxthistablewithglobalstyle
\centering
\begin{tabulary}{\linewidth}[t]{TT}
\sphinxtoprule
\begin{varwidth}[t]{\sphinxcolwidth{1}{2}}
\sphinxstyletheadfamily \sphinxAtStartPar
选项
\sphinxbeforeendvarwidth
\end{varwidth}%
&\begin{varwidth}[t]{\sphinxcolwidth{1}{2}}
\sphinxstyletheadfamily \sphinxAtStartPar
说明
\sphinxbeforeendvarwidth
\end{varwidth}%
\\
\sphinxmidrule
\sphinxtableatstartofbodyhook\begin{varwidth}[t]{\sphinxcolwidth{1}{2}}
\sphinxAtStartPar
选择已有方案
\sphinxbeforeendvarwidth
\end{varwidth}%
&\begin{varwidth}[t]{\sphinxcolwidth{1}{2}}
\sphinxAtStartPar
从下拉列表中选择已创建的FFT处理方案
\sphinxbeforeendvarwidth
\end{varwidth}%
\\
\sphinxhline\begin{varwidth}[t]{\sphinxcolwidth{1}{2}}
\sphinxAtStartPar
使用默认方案
\sphinxbeforeendvarwidth
\end{varwidth}%
&\begin{varwidth}[t]{\sphinxcolwidth{1}{2}}
\sphinxAtStartPar
不指定方案时使用系统默认参数
\sphinxbeforeendvarwidth
\end{varwidth}%
\\
\sphinxbottomrule
\end{tabulary}
\sphinxtableafterendhook\par
\sphinxattableend\end{savenotes}
\begin{quote}

\sphinxAtStartPar
💡 \sphinxstylestrong{提示}：处理方案可在后续处理阶段修改，不影响数据导入。
\end{quote}

\sphinxAtStartPar
\sphinxstylestrong{📊 测段统计信息：}


\begin{savenotes}\sphinxattablestart
\sphinxthistablewithglobalstyle
\centering
\begin{tabulary}{\linewidth}[t]{TT}
\sphinxtoprule
\begin{varwidth}[t]{\sphinxcolwidth{1}{2}}
\sphinxstyletheadfamily \sphinxAtStartPar
属性
\sphinxbeforeendvarwidth
\end{varwidth}%
&\begin{varwidth}[t]{\sphinxcolwidth{1}{2}}
\sphinxstyletheadfamily \sphinxAtStartPar
说明
\sphinxbeforeendvarwidth
\end{varwidth}%
\\
\sphinxmidrule
\sphinxtableatstartofbodyhook\begin{varwidth}[t]{\sphinxcolwidth{1}{2}}
\sphinxAtStartPar
Sites
\sphinxbeforeendvarwidth
\end{varwidth}%
&\begin{varwidth}[t]{\sphinxcolwidth{1}{2}}
\sphinxAtStartPar
测点列表
\sphinxbeforeendvarwidth
\end{varwidth}%
\\
\sphinxhline\begin{varwidth}[t]{\sphinxcolwidth{1}{2}}
\sphinxAtStartPar
SiteCount
\sphinxbeforeendvarwidth
\end{varwidth}%
&\begin{varwidth}[t]{\sphinxcolwidth{1}{2}}
\sphinxAtStartPar
测点数量（自动计算）
\sphinxbeforeendvarwidth
\end{varwidth}%
\\
\sphinxhline\begin{varwidth}[t]{\sphinxcolwidth{1}{2}}
\sphinxAtStartPar
ProcessSchema
\sphinxbeforeendvarwidth
\end{varwidth}%
&\begin{varwidth}[t]{\sphinxcolwidth{1}{2}}
\sphinxAtStartPar
关联的处理方案
\sphinxbeforeendvarwidth
\end{varwidth}%
\\
\sphinxbottomrule
\end{tabulary}
\sphinxtableafterendhook\par
\sphinxattableend\end{savenotes}


\subsubsection{3.4.3 批量加载测点}
\label{\detokenize{chapters/chapter3:id13}}
\sphinxAtStartPar
\sphinxstylestrong{从目录加载：}
\begin{enumerate}
\sphinxsetlistlabels{\arabic}{enumi}{enumii}{}{.}%
\item {} 
\sphinxAtStartPar
右键测段 → \sphinxcode{\sphinxupquote{加载测点 → 从目录}}

\item {} 
\sphinxAtStartPar
选择数据目录

\item {} 
\sphinxAtStartPar
系统自动识别数据类型

\end{enumerate}

\sphinxAtStartPar
\sphinxstylestrong{从EDI文件加载：}
\begin{enumerate}
\sphinxsetlistlabels{\arabic}{enumi}{enumii}{}{.}%
\item {} 
\sphinxAtStartPar
右键测段 → \sphinxcode{\sphinxupquote{加载测点 → 从EDI文件}}

\item {} 
\sphinxAtStartPar
选择EDI文件（支持多选）

\end{enumerate}


\bigskip\hrule\bigskip



\subsection{3.5 测点管理}
\label{\detokenize{chapters/chapter3:id14}}

\subsubsection{3.5.1 测点属性}
\label{\detokenize{chapters/chapter3:id15}}

\begin{savenotes}\sphinxattablestart
\sphinxthistablewithglobalstyle
\centering
\begin{tabulary}{\linewidth}[t]{TT}
\sphinxtoprule
\begin{varwidth}[t]{\sphinxcolwidth{1}{2}}
\sphinxstyletheadfamily \sphinxAtStartPar
属性
\sphinxbeforeendvarwidth
\end{varwidth}%
&\begin{varwidth}[t]{\sphinxcolwidth{1}{2}}
\sphinxstyletheadfamily \sphinxAtStartPar
说明
\sphinxbeforeendvarwidth
\end{varwidth}%
\\
\sphinxmidrule
\sphinxtableatstartofbodyhook\begin{varwidth}[t]{\sphinxcolwidth{1}{2}}
\sphinxAtStartPar
\sphinxstylestrong{📋 基本信息}
\sphinxbeforeendvarwidth
\end{varwidth}%
&\begin{varwidth}[t]{\sphinxcolwidth{1}{2}}
\sphinxAtStartPar

\sphinxbeforeendvarwidth
\end{varwidth}%
\\
\sphinxhline\begin{varwidth}[t]{\sphinxcolwidth{1}{2}}
\sphinxAtStartPar
Name
\sphinxbeforeendvarwidth
\end{varwidth}%
&\begin{varwidth}[t]{\sphinxcolwidth{1}{2}}
\sphinxAtStartPar
测点名称
\sphinxbeforeendvarwidth
\end{varwidth}%
\\
\sphinxhline\begin{varwidth}[t]{\sphinxcolwidth{1}{2}}
\sphinxAtStartPar
Owner
\sphinxbeforeendvarwidth
\end{varwidth}%
&\begin{varwidth}[t]{\sphinxcolwidth{1}{2}}
\sphinxAtStartPar
所有者
\sphinxbeforeendvarwidth
\end{varwidth}%
\\
\sphinxhline\begin{varwidth}[t]{\sphinxcolwidth{1}{2}}
\sphinxAtStartPar
Description
\sphinxbeforeendvarwidth
\end{varwidth}%
&\begin{varwidth}[t]{\sphinxcolwidth{1}{2}}
\sphinxAtStartPar
描述
\sphinxbeforeendvarwidth
\end{varwidth}%
\\
\sphinxhline\begin{varwidth}[t]{\sphinxcolwidth{1}{2}}
\sphinxAtStartPar
Surveyor
\sphinxbeforeendvarwidth
\end{varwidth}%
&\begin{varwidth}[t]{\sphinxcolwidth{1}{2}}
\sphinxAtStartPar
测量员
\sphinxbeforeendvarwidth
\end{varwidth}%
\\
\sphinxhline\begin{varwidth}[t]{\sphinxcolwidth{1}{2}}
\sphinxAtStartPar
DataCollector
\sphinxbeforeendvarwidth
\end{varwidth}%
&\begin{varwidth}[t]{\sphinxcolwidth{1}{2}}
\sphinxAtStartPar
数据采集者
\sphinxbeforeendvarwidth
\end{varwidth}%
\\
\sphinxhline\begin{varwidth}[t]{\sphinxcolwidth{1}{2}}
\sphinxAtStartPar
DataProcessor
\sphinxbeforeendvarwidth
\end{varwidth}%
&\begin{varwidth}[t]{\sphinxcolwidth{1}{2}}
\sphinxAtStartPar
数据处理者
\sphinxbeforeendvarwidth
\end{varwidth}%
\\
\sphinxhline\begin{varwidth}[t]{\sphinxcolwidth{1}{2}}
\sphinxAtStartPar
CreationTime
\sphinxbeforeendvarwidth
\end{varwidth}%
&\begin{varwidth}[t]{\sphinxcolwidth{1}{2}}
\sphinxAtStartPar
创建时间
\sphinxbeforeendvarwidth
\end{varwidth}%
\\
\sphinxhline\begin{varwidth}[t]{\sphinxcolwidth{1}{2}}
\sphinxAtStartPar
DataQuality
\sphinxbeforeendvarwidth
\end{varwidth}%
&\begin{varwidth}[t]{\sphinxcolwidth{1}{2}}
\sphinxAtStartPar
数据质量（0\sphinxhyphen{}4级）
\sphinxbeforeendvarwidth
\end{varwidth}%
\\
\sphinxhline\begin{varwidth}[t]{\sphinxcolwidth{1}{2}}
\sphinxAtStartPar
\sphinxstylestrong{🌍 坐标信息}
\sphinxbeforeendvarwidth
\end{varwidth}%
&\begin{varwidth}[t]{\sphinxcolwidth{1}{2}}
\sphinxAtStartPar

\sphinxbeforeendvarwidth
\end{varwidth}%
\\
\sphinxhline\begin{varwidth}[t]{\sphinxcolwidth{1}{2}}
\sphinxAtStartPar
Longitude
\sphinxbeforeendvarwidth
\end{varwidth}%
&\begin{varwidth}[t]{\sphinxcolwidth{1}{2}}
\sphinxAtStartPar
经度
\sphinxbeforeendvarwidth
\end{varwidth}%
\\
\sphinxhline\begin{varwidth}[t]{\sphinxcolwidth{1}{2}}
\sphinxAtStartPar
Latitude
\sphinxbeforeendvarwidth
\end{varwidth}%
&\begin{varwidth}[t]{\sphinxcolwidth{1}{2}}
\sphinxAtStartPar
纬度
\sphinxbeforeendvarwidth
\end{varwidth}%
\\
\sphinxhline\begin{varwidth}[t]{\sphinxcolwidth{1}{2}}
\sphinxAtStartPar
Altitude
\sphinxbeforeendvarwidth
\end{varwidth}%
&\begin{varwidth}[t]{\sphinxcolwidth{1}{2}}
\sphinxAtStartPar
高程
\sphinxbeforeendvarwidth
\end{varwidth}%
\\
\sphinxhline\begin{varwidth}[t]{\sphinxcolwidth{1}{2}}
\sphinxAtStartPar
\sphinxstylestrong{📦 采集盒信息}
\sphinxbeforeendvarwidth
\end{varwidth}%
&\begin{varwidth}[t]{\sphinxcolwidth{1}{2}}
\sphinxAtStartPar

\sphinxbeforeendvarwidth
\end{varwidth}%
\\
\sphinxhline\begin{varwidth}[t]{\sphinxcolwidth{1}{2}}
\sphinxAtStartPar
BoxID
\sphinxbeforeendvarwidth
\end{varwidth}%
&\begin{varwidth}[t]{\sphinxcolwidth{1}{2}}
\sphinxAtStartPar
采集盒编号
\sphinxbeforeendvarwidth
\end{varwidth}%
\\
\sphinxhline\begin{varwidth}[t]{\sphinxcolwidth{1}{2}}
\sphinxAtStartPar
EPreamplifier
\sphinxbeforeendvarwidth
\end{varwidth}%
&\begin{varwidth}[t]{\sphinxcolwidth{1}{2}}
\sphinxAtStartPar
电场前置放大器
\sphinxbeforeendvarwidth
\end{varwidth}%
\\
\sphinxhline\begin{varwidth}[t]{\sphinxcolwidth{1}{2}}
\sphinxAtStartPar
ExGroundRes
\sphinxbeforeendvarwidth
\end{varwidth}%
&\begin{varwidth}[t]{\sphinxcolwidth{1}{2}}
\sphinxAtStartPar
Ex接地电阻
\sphinxbeforeendvarwidth
\end{varwidth}%
\\
\sphinxhline\begin{varwidth}[t]{\sphinxcolwidth{1}{2}}
\sphinxAtStartPar
EyGroundRes
\sphinxbeforeendvarwidth
\end{varwidth}%
&\begin{varwidth}[t]{\sphinxcolwidth{1}{2}}
\sphinxAtStartPar
Ey接地电阻
\sphinxbeforeendvarwidth
\end{varwidth}%
\\
\sphinxhline\begin{varwidth}[t]{\sphinxcolwidth{1}{2}}
\sphinxAtStartPar
\sphinxstylestrong{📊 频率信息}
\sphinxbeforeendvarwidth
\end{varwidth}%
&\begin{varwidth}[t]{\sphinxcolwidth{1}{2}}
\sphinxAtStartPar

\sphinxbeforeendvarwidth
\end{varwidth}%
\\
\sphinxhline\begin{varwidth}[t]{\sphinxcolwidth{1}{2}}
\sphinxAtStartPar
MaxAvailableFrequency
\sphinxbeforeendvarwidth
\end{varwidth}%
&\begin{varwidth}[t]{\sphinxcolwidth{1}{2}}
\sphinxAtStartPar
最大可用频率
\sphinxbeforeendvarwidth
\end{varwidth}%
\\
\sphinxhline\begin{varwidth}[t]{\sphinxcolwidth{1}{2}}
\sphinxAtStartPar
MinAvailableFrequency
\sphinxbeforeendvarwidth
\end{varwidth}%
&\begin{varwidth}[t]{\sphinxcolwidth{1}{2}}
\sphinxAtStartPar
最小可用频率
\sphinxbeforeendvarwidth
\end{varwidth}%
\\
\sphinxbottomrule
\end{tabulary}
\sphinxtableafterendhook\par
\sphinxattableend\end{savenotes}


\subsubsection{3.5.2 测点操作}
\label{\detokenize{chapters/chapter3:id16}}\begin{itemize}
\item {} 
\sphinxAtStartPar
\sphinxstylestrong{编辑属性}：右键 → 测点设置

\item {} 
\sphinxAtStartPar
\sphinxstylestrong{设置电场类型}：右键 → 电场类型

\item {} 
\sphinxAtStartPar
\sphinxstylestrong{设置数据质量}：右键 → 数据质量

\item {} 
\sphinxAtStartPar
\sphinxstylestrong{导出EDI}：右键 → 导出EDI

\end{itemize}


\subsubsection{3.5.3 数据质量等级}
\label{\detokenize{chapters/chapter3:id17}}

\begin{savenotes}\sphinxattablestart
\sphinxthistablewithglobalstyle
\centering
\begin{tabulary}{\linewidth}[t]{TT}
\sphinxtoprule
\begin{varwidth}[t]{\sphinxcolwidth{1}{2}}
\sphinxstyletheadfamily \sphinxAtStartPar
等级
\sphinxbeforeendvarwidth
\end{varwidth}%
&\begin{varwidth}[t]{\sphinxcolwidth{1}{2}}
\sphinxstyletheadfamily \sphinxAtStartPar
说明
\sphinxbeforeendvarwidth
\end{varwidth}%
\\
\sphinxmidrule
\sphinxtableatstartofbodyhook\begin{varwidth}[t]{\sphinxcolwidth{1}{2}}
\sphinxAtStartPar
0
\sphinxbeforeendvarwidth
\end{varwidth}%
&\begin{varwidth}[t]{\sphinxcolwidth{1}{2}}
\sphinxAtStartPar
未评估
\sphinxbeforeendvarwidth
\end{varwidth}%
\\
\sphinxhline\begin{varwidth}[t]{\sphinxcolwidth{1}{2}}
\sphinxAtStartPar
1
\sphinxbeforeendvarwidth
\end{varwidth}%
&\begin{varwidth}[t]{\sphinxcolwidth{1}{2}}
\sphinxAtStartPar
优秀
\sphinxbeforeendvarwidth
\end{varwidth}%
\\
\sphinxhline\begin{varwidth}[t]{\sphinxcolwidth{1}{2}}
\sphinxAtStartPar
2
\sphinxbeforeendvarwidth
\end{varwidth}%
&\begin{varwidth}[t]{\sphinxcolwidth{1}{2}}
\sphinxAtStartPar
良好
\sphinxbeforeendvarwidth
\end{varwidth}%
\\
\sphinxhline\begin{varwidth}[t]{\sphinxcolwidth{1}{2}}
\sphinxAtStartPar
3
\sphinxbeforeendvarwidth
\end{varwidth}%
&\begin{varwidth}[t]{\sphinxcolwidth{1}{2}}
\sphinxAtStartPar
一般
\sphinxbeforeendvarwidth
\end{varwidth}%
\\
\sphinxhline\begin{varwidth}[t]{\sphinxcolwidth{1}{2}}
\sphinxAtStartPar
4
\sphinxbeforeendvarwidth
\end{varwidth}%
&\begin{varwidth}[t]{\sphinxcolwidth{1}{2}}
\sphinxAtStartPar
较差
\sphinxbeforeendvarwidth
\end{varwidth}%
\\
\sphinxbottomrule
\end{tabulary}
\sphinxtableafterendhook\par
\sphinxattableend\end{savenotes}


\subsubsection{3.5.4 测点排序}
\label{\detokenize{chapters/chapter3:id18}}
\sphinxAtStartPar
可按以下方式对测点排序：


\begin{savenotes}\sphinxattablestart
\sphinxthistablewithglobalstyle
\centering
\begin{tabulary}{\linewidth}[t]{TT}
\sphinxtoprule
\begin{varwidth}[t]{\sphinxcolwidth{1}{2}}
\sphinxstyletheadfamily \sphinxAtStartPar
排序方式
\sphinxbeforeendvarwidth
\end{varwidth}%
&\begin{varwidth}[t]{\sphinxcolwidth{1}{2}}
\sphinxstyletheadfamily \sphinxAtStartPar
说明
\sphinxbeforeendvarwidth
\end{varwidth}%
\\
\sphinxmidrule
\sphinxtableatstartofbodyhook\begin{varwidth}[t]{\sphinxcolwidth{1}{2}}
\sphinxAtStartPar
按纬度
\sphinxbeforeendvarwidth
\end{varwidth}%
&\begin{varwidth}[t]{\sphinxcolwidth{1}{2}}
\sphinxAtStartPar
从南到北或从北到南
\sphinxbeforeendvarwidth
\end{varwidth}%
\\
\sphinxhline\begin{varwidth}[t]{\sphinxcolwidth{1}{2}}
\sphinxAtStartPar
按经度
\sphinxbeforeendvarwidth
\end{varwidth}%
&\begin{varwidth}[t]{\sphinxcolwidth{1}{2}}
\sphinxAtStartPar
从西到东或从东到西
\sphinxbeforeendvarwidth
\end{varwidth}%
\\
\sphinxhline\begin{varwidth}[t]{\sphinxcolwidth{1}{2}}
\sphinxAtStartPar
按名称
\sphinxbeforeendvarwidth
\end{varwidth}%
&\begin{varwidth}[t]{\sphinxcolwidth{1}{2}}
\sphinxAtStartPar
按字母顺序
\sphinxbeforeendvarwidth
\end{varwidth}%
\\
\sphinxhline\begin{varwidth}[t]{\sphinxcolwidth{1}{2}}
\sphinxAtStartPar
按开始时间
\sphinxbeforeendvarwidth
\end{varwidth}%
&\begin{varwidth}[t]{\sphinxcolwidth{1}{2}}
\sphinxAtStartPar
按采集开始时间
\sphinxbeforeendvarwidth
\end{varwidth}%
\\
\sphinxhline\begin{varwidth}[t]{\sphinxcolwidth{1}{2}}
\sphinxAtStartPar
按结束时间
\sphinxbeforeendvarwidth
\end{varwidth}%
&\begin{varwidth}[t]{\sphinxcolwidth{1}{2}}
\sphinxAtStartPar
按采集结束时间
\sphinxbeforeendvarwidth
\end{varwidth}%
\\
\sphinxbottomrule
\end{tabulary}
\sphinxtableafterendhook\par
\sphinxattableend\end{savenotes}

\sphinxAtStartPar
操作方法：右键测段 → 排序测点 → 选择排序方式


\subsubsection{3.5.5 通道配置}
\label{\detokenize{chapters/chapter3:id19}}
\sphinxAtStartPar
每个测点包含5个电磁场通道，可独立配置：


\begin{savenotes}\sphinxattablestart
\sphinxthistablewithglobalstyle
\centering
\begin{tabulary}{\linewidth}[t]{TTTTTT}
\sphinxtoprule
\begin{varwidth}[t]{\sphinxcolwidth{1}{6}}
\sphinxstyletheadfamily \sphinxAtStartPar
通道
\sphinxbeforeendvarwidth
\end{varwidth}%
&\begin{varwidth}[t]{\sphinxcolwidth{1}{6}}
\sphinxstyletheadfamily \sphinxAtStartPar
索引
\sphinxbeforeendvarwidth
\end{varwidth}%
&\begin{varwidth}[t]{\sphinxcolwidth{1}{6}}
\sphinxstyletheadfamily \sphinxAtStartPar
反向
\sphinxbeforeendvarwidth
\end{varwidth}%
&\begin{varwidth}[t]{\sphinxcolwidth{1}{6}}
\sphinxstyletheadfamily \sphinxAtStartPar
长度(m)
\sphinxbeforeendvarwidth
\end{varwidth}%
&\begin{varwidth}[t]{\sphinxcolwidth{1}{6}}
\sphinxstyletheadfamily \sphinxAtStartPar
传感器
\sphinxbeforeendvarwidth
\end{varwidth}%
&\begin{varwidth}[t]{\sphinxcolwidth{1}{6}}
\sphinxstyletheadfamily \sphinxAtStartPar
旋转角(°)
\sphinxbeforeendvarwidth
\end{varwidth}%
\\
\sphinxmidrule
\sphinxtableatstartofbodyhook\begin{varwidth}[t]{\sphinxcolwidth{1}{6}}
\sphinxAtStartPar
Ex
\sphinxbeforeendvarwidth
\end{varwidth}%
&\begin{varwidth}[t]{\sphinxcolwidth{1}{6}}
\sphinxAtStartPar
\(\checkmark\)
\sphinxbeforeendvarwidth
\end{varwidth}%
&\begin{varwidth}[t]{\sphinxcolwidth{1}{6}}
\sphinxAtStartPar
\(\checkmark\)
\sphinxbeforeendvarwidth
\end{varwidth}%
&\begin{varwidth}[t]{\sphinxcolwidth{1}{6}}
\sphinxAtStartPar
\(\checkmark\)
\sphinxbeforeendvarwidth
\end{varwidth}%
&\begin{varwidth}[t]{\sphinxcolwidth{1}{6}}
\sphinxAtStartPar
\sphinxhyphen{}
\sphinxbeforeendvarwidth
\end{varwidth}%
&\begin{varwidth}[t]{\sphinxcolwidth{1}{6}}
\sphinxAtStartPar
\(\checkmark\)
\sphinxbeforeendvarwidth
\end{varwidth}%
\\
\sphinxhline\begin{varwidth}[t]{\sphinxcolwidth{1}{6}}
\sphinxAtStartPar
Ey
\sphinxbeforeendvarwidth
\end{varwidth}%
&\begin{varwidth}[t]{\sphinxcolwidth{1}{6}}
\sphinxAtStartPar
\(\checkmark\)
\sphinxbeforeendvarwidth
\end{varwidth}%
&\begin{varwidth}[t]{\sphinxcolwidth{1}{6}}
\sphinxAtStartPar
\(\checkmark\)
\sphinxbeforeendvarwidth
\end{varwidth}%
&\begin{varwidth}[t]{\sphinxcolwidth{1}{6}}
\sphinxAtStartPar
\(\checkmark\)
\sphinxbeforeendvarwidth
\end{varwidth}%
&\begin{varwidth}[t]{\sphinxcolwidth{1}{6}}
\sphinxAtStartPar
\sphinxhyphen{}
\sphinxbeforeendvarwidth
\end{varwidth}%
&\begin{varwidth}[t]{\sphinxcolwidth{1}{6}}
\sphinxAtStartPar
\(\checkmark\)
\sphinxbeforeendvarwidth
\end{varwidth}%
\\
\sphinxhline\begin{varwidth}[t]{\sphinxcolwidth{1}{6}}
\sphinxAtStartPar
Hx
\sphinxbeforeendvarwidth
\end{varwidth}%
&\begin{varwidth}[t]{\sphinxcolwidth{1}{6}}
\sphinxAtStartPar
\(\checkmark\)
\sphinxbeforeendvarwidth
\end{varwidth}%
&\begin{varwidth}[t]{\sphinxcolwidth{1}{6}}
\sphinxAtStartPar
\(\checkmark\)
\sphinxbeforeendvarwidth
\end{varwidth}%
&\begin{varwidth}[t]{\sphinxcolwidth{1}{6}}
\sphinxAtStartPar
\(\checkmark\)
\sphinxbeforeendvarwidth
\end{varwidth}%
&\begin{varwidth}[t]{\sphinxcolwidth{1}{6}}
\sphinxAtStartPar
\(\checkmark\)
\sphinxbeforeendvarwidth
\end{varwidth}%
&\begin{varwidth}[t]{\sphinxcolwidth{1}{6}}
\sphinxAtStartPar
\(\checkmark\)
\sphinxbeforeendvarwidth
\end{varwidth}%
\\
\sphinxhline\begin{varwidth}[t]{\sphinxcolwidth{1}{6}}
\sphinxAtStartPar
Hy
\sphinxbeforeendvarwidth
\end{varwidth}%
&\begin{varwidth}[t]{\sphinxcolwidth{1}{6}}
\sphinxAtStartPar
\(\checkmark\)
\sphinxbeforeendvarwidth
\end{varwidth}%
&\begin{varwidth}[t]{\sphinxcolwidth{1}{6}}
\sphinxAtStartPar
\(\checkmark\)
\sphinxbeforeendvarwidth
\end{varwidth}%
&\begin{varwidth}[t]{\sphinxcolwidth{1}{6}}
\sphinxAtStartPar
\(\checkmark\)
\sphinxbeforeendvarwidth
\end{varwidth}%
&\begin{varwidth}[t]{\sphinxcolwidth{1}{6}}
\sphinxAtStartPar
\(\checkmark\)
\sphinxbeforeendvarwidth
\end{varwidth}%
&\begin{varwidth}[t]{\sphinxcolwidth{1}{6}}
\sphinxAtStartPar
\(\checkmark\)
\sphinxbeforeendvarwidth
\end{varwidth}%
\\
\sphinxhline\begin{varwidth}[t]{\sphinxcolwidth{1}{6}}
\sphinxAtStartPar
Hz
\sphinxbeforeendvarwidth
\end{varwidth}%
&\begin{varwidth}[t]{\sphinxcolwidth{1}{6}}
\sphinxAtStartPar
\(\checkmark\)
\sphinxbeforeendvarwidth
\end{varwidth}%
&\begin{varwidth}[t]{\sphinxcolwidth{1}{6}}
\sphinxAtStartPar
\(\checkmark\)
\sphinxbeforeendvarwidth
\end{varwidth}%
&\begin{varwidth}[t]{\sphinxcolwidth{1}{6}}
\sphinxAtStartPar
\(\checkmark\)
\sphinxbeforeendvarwidth
\end{varwidth}%
&\begin{varwidth}[t]{\sphinxcolwidth{1}{6}}
\sphinxAtStartPar
\(\checkmark\)
\sphinxbeforeendvarwidth
\end{varwidth}%
&\begin{varwidth}[t]{\sphinxcolwidth{1}{6}}
\sphinxAtStartPar
\(\checkmark\)
\sphinxbeforeendvarwidth
\end{varwidth}%
\\
\sphinxbottomrule
\end{tabulary}
\sphinxtableafterendhook\par
\sphinxattableend\end{savenotes}

\sphinxAtStartPar
\sphinxstylestrong{配置项说明：}


\begin{savenotes}\sphinxattablestart
\sphinxthistablewithglobalstyle
\centering
\begin{tabulary}{\linewidth}[t]{TT}
\sphinxtoprule
\begin{varwidth}[t]{\sphinxcolwidth{1}{2}}
\sphinxstyletheadfamily \sphinxAtStartPar
配置项
\sphinxbeforeendvarwidth
\end{varwidth}%
&\begin{varwidth}[t]{\sphinxcolwidth{1}{2}}
\sphinxstyletheadfamily \sphinxAtStartPar
说明
\sphinxbeforeendvarwidth
\end{varwidth}%
\\
\sphinxmidrule
\sphinxtableatstartofbodyhook\begin{varwidth}[t]{\sphinxcolwidth{1}{2}}
\sphinxAtStartPar
索引
\sphinxbeforeendvarwidth
\end{varwidth}%
&\begin{varwidth}[t]{\sphinxcolwidth{1}{2}}
\sphinxAtStartPar
通道在数据文件中的顺序号
\sphinxbeforeendvarwidth
\end{varwidth}%
\\
\sphinxhline\begin{varwidth}[t]{\sphinxcolwidth{1}{2}}
\sphinxAtStartPar
反向
\sphinxbeforeendvarwidth
\end{varwidth}%
&\begin{varwidth}[t]{\sphinxcolwidth{1}{2}}
\sphinxAtStartPar
是否翻转信号方向
\sphinxbeforeendvarwidth
\end{varwidth}%
\\
\sphinxhline\begin{varwidth}[t]{\sphinxcolwidth{1}{2}}
\sphinxAtStartPar
长度
\sphinxbeforeendvarwidth
\end{varwidth}%
&\begin{varwidth}[t]{\sphinxcolwidth{1}{2}}
\sphinxAtStartPar
电场电极距离（仅Ex、Ey）
\sphinxbeforeendvarwidth
\end{varwidth}%
\\
\sphinxhline\begin{varwidth}[t]{\sphinxcolwidth{1}{2}}
\sphinxAtStartPar
传感器
\sphinxbeforeendvarwidth
\end{varwidth}%
&\begin{varwidth}[t]{\sphinxcolwidth{1}{2}}
\sphinxAtStartPar
磁传感器编号（仅Hx、Hy、Hz）
\sphinxbeforeendvarwidth
\end{varwidth}%
\\
\sphinxhline\begin{varwidth}[t]{\sphinxcolwidth{1}{2}}
\sphinxAtStartPar
旋转角
\sphinxbeforeendvarwidth
\end{varwidth}%
&\begin{varwidth}[t]{\sphinxcolwidth{1}{2}}
\sphinxAtStartPar
通道相对于地理北的旋转角度
\sphinxbeforeendvarwidth
\end{varwidth}%
\\
\sphinxbottomrule
\end{tabulary}
\sphinxtableafterendhook\par
\sphinxattableend\end{savenotes}


\subsubsection{3.5.6 时间信息}
\label{\detokenize{chapters/chapter3:id20}}

\begin{savenotes}\sphinxattablestart
\sphinxthistablewithglobalstyle
\centering
\begin{tabulary}{\linewidth}[t]{TT}
\sphinxtoprule
\begin{varwidth}[t]{\sphinxcolwidth{1}{2}}
\sphinxstyletheadfamily \sphinxAtStartPar
属性
\sphinxbeforeendvarwidth
\end{varwidth}%
&\begin{varwidth}[t]{\sphinxcolwidth{1}{2}}
\sphinxstyletheadfamily \sphinxAtStartPar
说明
\sphinxbeforeendvarwidth
\end{varwidth}%
\\
\sphinxmidrule
\sphinxtableatstartofbodyhook\begin{varwidth}[t]{\sphinxcolwidth{1}{2}}
\sphinxAtStartPar
BeginTime
\sphinxbeforeendvarwidth
\end{varwidth}%
&\begin{varwidth}[t]{\sphinxcolwidth{1}{2}}
\sphinxAtStartPar
采集开始时间
\sphinxbeforeendvarwidth
\end{varwidth}%
\\
\sphinxhline\begin{varwidth}[t]{\sphinxcolwidth{1}{2}}
\sphinxAtStartPar
EndTime
\sphinxbeforeendvarwidth
\end{varwidth}%
&\begin{varwidth}[t]{\sphinxcolwidth{1}{2}}
\sphinxAtStartPar
采集结束时间
\sphinxbeforeendvarwidth
\end{varwidth}%
\\
\sphinxhline\begin{varwidth}[t]{\sphinxcolwidth{1}{2}}
\sphinxAtStartPar
ProcessBeginTime
\sphinxbeforeendvarwidth
\end{varwidth}%
&\begin{varwidth}[t]{\sphinxcolwidth{1}{2}}
\sphinxAtStartPar
处理开始时间（可选）
\sphinxbeforeendvarwidth
\end{varwidth}%
\\
\sphinxhline\begin{varwidth}[t]{\sphinxcolwidth{1}{2}}
\sphinxAtStartPar
ProcessEndTime
\sphinxbeforeendvarwidth
\end{varwidth}%
&\begin{varwidth}[t]{\sphinxcolwidth{1}{2}}
\sphinxAtStartPar
处理结束时间（可选）
\sphinxbeforeendvarwidth
\end{varwidth}%
\\
\sphinxhline\begin{varwidth}[t]{\sphinxcolwidth{1}{2}}
\sphinxAtStartPar
TimeLength
\sphinxbeforeendvarwidth
\end{varwidth}%
&\begin{varwidth}[t]{\sphinxcolwidth{1}{2}}
\sphinxAtStartPar
采集时长（小时，自动计算）
\sphinxbeforeendvarwidth
\end{varwidth}%
\\
\sphinxbottomrule
\end{tabulary}
\sphinxtableafterendhook\par
\sphinxattableend\end{savenotes}


\subsubsection{3.5.7 特殊功能}
\label{\detokenize{chapters/chapter3:id21}}
\sphinxAtStartPar
\sphinxstylestrong{🔄 拖放文件更新}

\sphinxAtStartPar
支持通过拖放文件自动更新测点信息：


\begin{savenotes}\sphinxattablestart
\sphinxthistablewithglobalstyle
\centering
\begin{tabulary}{\linewidth}[t]{TT}
\sphinxtoprule
\begin{varwidth}[t]{\sphinxcolwidth{1}{2}}
\sphinxstyletheadfamily \sphinxAtStartPar
文件类型
\sphinxbeforeendvarwidth
\end{varwidth}%
&\begin{varwidth}[t]{\sphinxcolwidth{1}{2}}
\sphinxstyletheadfamily \sphinxAtStartPar
更新内容
\sphinxbeforeendvarwidth
\end{varwidth}%
\\
\sphinxmidrule
\sphinxtableatstartofbodyhook\begin{varwidth}[t]{\sphinxcolwidth{1}{2}}
\sphinxAtStartPar
.tbl
\sphinxbeforeendvarwidth
\end{varwidth}%
&\begin{varwidth}[t]{\sphinxcolwidth{1}{2}}
\sphinxAtStartPar
自动更新Phoenix测点坐标、时间、采集盒ID、传感器ID、电极长度、旋转角
\sphinxbeforeendvarwidth
\end{varwidth}%
\\
\sphinxhline\begin{varwidth}[t]{\sphinxcolwidth{1}{2}}
\sphinxAtStartPar
.lemi
\sphinxbeforeendvarwidth
\end{varwidth}%
&\begin{varwidth}[t]{\sphinxcolwidth{1}{2}}
\sphinxAtStartPar
更新LEMI通道索引
\sphinxbeforeendvarwidth
\end{varwidth}%
\\
\sphinxhline\begin{varwidth}[t]{\sphinxcolwidth{1}{2}}
\sphinxAtStartPar
.stfc
\sphinxbeforeendvarwidth
\end{varwidth}%
&\begin{varwidth}[t]{\sphinxcolwidth{1}{2}}
\sphinxAtStartPar
加载傅里叶系数
\sphinxbeforeendvarwidth
\end{varwidth}%
\\
\sphinxbottomrule
\end{tabulary}
\sphinxtableafterendhook\par
\sphinxattableend\end{savenotes}

\sphinxAtStartPar
\sphinxstylestrong{💾 自动备份}


\begin{savenotes}\sphinxattablestart
\sphinxthistablewithglobalstyle
\centering
\begin{tabulary}{\linewidth}[t]{TT}
\sphinxtoprule
\begin{varwidth}[t]{\sphinxcolwidth{1}{2}}
\sphinxstyletheadfamily \sphinxAtStartPar
操作
\sphinxbeforeendvarwidth
\end{varwidth}%
&\begin{varwidth}[t]{\sphinxcolwidth{1}{2}}
\sphinxstyletheadfamily \sphinxAtStartPar
备份行为
\sphinxbeforeendvarwidth
\end{varwidth}%
\\
\sphinxmidrule
\sphinxtableatstartofbodyhook\begin{varwidth}[t]{\sphinxcolwidth{1}{2}}
\sphinxAtStartPar
修改Phoenix TBL
\sphinxbeforeendvarwidth
\end{varwidth}%
&\begin{varwidth}[t]{\sphinxcolwidth{1}{2}}
\sphinxAtStartPar
自动备份为 .tbl.backup
\sphinxbeforeendvarwidth
\end{varwidth}%
\\
\sphinxhline\begin{varwidth}[t]{\sphinxcolwidth{1}{2}}
\sphinxAtStartPar
修改RMT JSON
\sphinxbeforeendvarwidth
\end{varwidth}%
&\begin{varwidth}[t]{\sphinxcolwidth{1}{2}}
\sphinxAtStartPar
自动备份原文件
\sphinxbeforeendvarwidth
\end{varwidth}%
\\
\sphinxbottomrule
\end{tabulary}
\sphinxtableafterendhook\par
\sphinxattableend\end{savenotes}


\bigskip\hrule\bigskip



\subsection{3.6 远参考站管理}
\label{\detokenize{chapters/chapter3:id22}}

\subsubsection{3.6.1 远参考站计算}
\label{\detokenize{chapters/chapter3:id23}}
\sphinxAtStartPar
当测段内有多个测点同时采集时，可自动计算远参考关系：

\sphinxAtStartPar
1️⃣ 右键测段 → \sphinxcode{\sphinxupquote{计算远参考}}
2️⃣ 系统自动计算所有测点间的傅里叶系数关系
3️⃣ 生成 .RSTFC 文件（远参考傅里叶系数）
\begin{quote}

\sphinxAtStartPar
💡 \sphinxstylestrong{提示}：远参考处理需要测点间有时间重叠的采集数据。
\end{quote}


\subsubsection{3.6.2 设置公共磁场站}
\label{\detokenize{chapters/chapter3:id24}}
\sphinxAtStartPar
将一个测点作为公共磁场站：

\sphinxAtStartPar
1️⃣ 选择作为磁场站的测点
2️⃣ 右键测段 → \sphinxcode{\sphinxupquote{设为公共磁场站}}
3️⃣ 所有其他测点将使用该站的磁场数据


\subsubsection{3.6.3 设置远参考站}
\label{\detokenize{chapters/chapter3:id25}}
\sphinxAtStartPar
将一个测点作为远参考站：

\sphinxAtStartPar
1️⃣ 选择作为远参考的测点
2️⃣ 右键测段 → \sphinxcode{\sphinxupquote{设为远参考站}}
3️⃣ 所有其他测点将使用该站进行远参考处理


\bigskip\hrule\bigskip



\subsection{3.7 坐标管理}
\label{\detokenize{chapters/chapter3:id26}}

\subsubsection{3.7.1 从KML导入坐标}
\label{\detokenize{chapters/chapter3:kml}}
\sphinxAtStartPar
从KML文件批量更新测点坐标：
\begin{enumerate}
\sphinxsetlistlabels{\arabic}{enumi}{enumii}{}{.}%
\item {} 
\sphinxAtStartPar
准备包含测点位置信息的KML文件

\item {} 
\sphinxAtStartPar
右键测段 → \sphinxcode{\sphinxupquote{从KML读取坐标}}

\item {} 
\sphinxAtStartPar
选择KML文件

\item {} 
\sphinxAtStartPar
系统自动匹配测点名称并更新坐标

\end{enumerate}


\subsubsection{3.7.2 导出测点坐标}
\label{\detokenize{chapters/chapter3:id27}}\begin{enumerate}
\sphinxsetlistlabels{\arabic}{enumi}{enumii}{}{.}%
\item {} 
\sphinxAtStartPar
右键测段 → \sphinxcode{\sphinxupquote{导出坐标}}

\item {} 
\sphinxAtStartPar
选择保存位置

\item {} 
\sphinxAtStartPar
生成坐标文件（包含名称、经度、纬度、高程）

\end{enumerate}


\subsubsection{3.7.3 坐标格式}
\label{\detokenize{chapters/chapter3:id28}}
\sphinxAtStartPar
MTDP支持以下坐标格式：
\begin{itemize}
\item {} 
\sphinxAtStartPar
十进制度（推荐）

\item {} 
\sphinxAtStartPar
度分秒（自动转换）

\end{itemize}


\bigskip\hrule\bigskip



\subsection{3.8 工程文件管理}
\label{\detokenize{chapters/chapter3:id29}}

\subsubsection{3.6.1 保存与备份}
\label{\detokenize{chapters/chapter3:id30}}
\sphinxAtStartPar
\sphinxstylestrong{手动保存：}
\begin{itemize}
\item {} 
\sphinxAtStartPar
\sphinxcode{\sphinxupquote{文件 → 保存}} (Ctrl+S)

\item {} 
\sphinxAtStartPar
\sphinxcode{\sphinxupquote{文件 → 另存为}}

\end{itemize}

\sphinxAtStartPar
\sphinxstylestrong{自动备份：}
\begin{itemize}
\item {} 
\sphinxAtStartPar
每次保存时自动创建备份

\item {} 
\sphinxAtStartPar
保留最近10个版本

\item {} 
\sphinxAtStartPar
备份文件在工程目录下

\end{itemize}


\subsubsection{3.6.2 工程信息存储}
\label{\detokenize{chapters/chapter3:id31}}
\sphinxAtStartPar
工程信息保存在以下位置：
\begin{itemize}
\item {} 
\sphinxAtStartPar
\sphinxstylestrong{工程文件列表}: \sphinxcode{\sphinxupquote{Configurations/ProjectInfos.XML}}

\item {} 
\sphinxAtStartPar
\sphinxstylestrong{语言文件}: \sphinxcode{\sphinxupquote{Configurations/Lang/}} 目录

\end{itemize}


\subsubsection{3.6.3 导入导出}
\label{\detokenize{chapters/chapter3:id32}}
\sphinxAtStartPar
\sphinxstylestrong{导出版本：}
\begin{itemize}
\item {} 
\sphinxAtStartPar
\sphinxcode{\sphinxupquote{文件 → 导出版本}}

\item {} 
\sphinxAtStartPar
创建完整工程副本

\end{itemize}

\sphinxAtStartPar
\sphinxstylestrong{导出选中项：}
\begin{itemize}
\item {} 
\sphinxAtStartPar
\sphinxcode{\sphinxupquote{文件 → 导出选中项}}

\item {} 
\sphinxAtStartPar
仅导出选中的数据

\end{itemize}


\bigskip\hrule\bigskip



\subsection{3.9 多语言支持}
\label{\detokenize{chapters/chapter3:id33}}
\sphinxAtStartPar
MTDP支持多语言界面：
\begin{enumerate}
\sphinxsetlistlabels{\arabic}{enumi}{enumii}{}{.}%
\item {} 
\sphinxAtStartPar
选择 \sphinxcode{\sphinxupquote{设置 → 语言}}

\item {} 
\sphinxAtStartPar
选择语言：
\begin{itemize}
\item {} 
\sphinxAtStartPar
中文

\item {} 
\sphinxAtStartPar
English

\end{itemize}

\item {} 
\sphinxAtStartPar
重启软件生效

\end{enumerate}

\sphinxAtStartPar
语言文件存储在 \sphinxcode{\sphinxupquote{Configurations/Lang/}} 目录。


\bigskip\hrule\bigskip



\subsection{3.10 常用操作快捷方式}
\label{\detokenize{chapters/chapter3:id34}}

\begin{savenotes}\sphinxattablestart
\sphinxthistablewithglobalstyle
\centering
\begin{tabulary}{\linewidth}[t]{TT}
\sphinxtoprule
\begin{varwidth}[t]{\sphinxcolwidth{1}{2}}
\sphinxstyletheadfamily \sphinxAtStartPar
操作
\sphinxbeforeendvarwidth
\end{varwidth}%
&\begin{varwidth}[t]{\sphinxcolwidth{1}{2}}
\sphinxstyletheadfamily \sphinxAtStartPar
方法
\sphinxbeforeendvarwidth
\end{varwidth}%
\\
\sphinxmidrule
\sphinxtableatstartofbodyhook\begin{varwidth}[t]{\sphinxcolwidth{1}{2}}
\sphinxAtStartPar
复制测点
\sphinxbeforeendvarwidth
\end{varwidth}%
&\begin{varwidth}[t]{\sphinxcolwidth{1}{2}}
\sphinxAtStartPar
右键 → 复制 / Ctrl+C
\sphinxbeforeendvarwidth
\end{varwidth}%
\\
\sphinxhline\begin{varwidth}[t]{\sphinxcolwidth{1}{2}}
\sphinxAtStartPar
粘贴测点
\sphinxbeforeendvarwidth
\end{varwidth}%
&\begin{varwidth}[t]{\sphinxcolwidth{1}{2}}
\sphinxAtStartPar
右键 → 粘贴 / Ctrl+V
\sphinxbeforeendvarwidth
\end{varwidth}%
\\
\sphinxhline\begin{varwidth}[t]{\sphinxcolwidth{1}{2}}
\sphinxAtStartPar
删除
\sphinxbeforeendvarwidth
\end{varwidth}%
&\begin{varwidth}[t]{\sphinxcolwidth{1}{2}}
\sphinxAtStartPar
右键 → 删除 / Delete
\sphinxbeforeendvarwidth
\end{varwidth}%
\\
\sphinxhline\begin{varwidth}[t]{\sphinxcolwidth{1}{2}}
\sphinxAtStartPar
重命名
\sphinxbeforeendvarwidth
\end{varwidth}%
&\begin{varwidth}[t]{\sphinxcolwidth{1}{2}}
\sphinxAtStartPar
右键 → 重命名 / F2
\sphinxbeforeendvarwidth
\end{varwidth}%
\\
\sphinxhline\begin{varwidth}[t]{\sphinxcolwidth{1}{2}}
\sphinxAtStartPar
批量选择
\sphinxbeforeendvarwidth
\end{varwidth}%
&\begin{varwidth}[t]{\sphinxcolwidth{1}{2}}
\sphinxAtStartPar
按住Ctrl多选
\sphinxbeforeendvarwidth
\end{varwidth}%
\\
\sphinxbottomrule
\end{tabulary}
\sphinxtableafterendhook\par
\sphinxattableend\end{savenotes}

\sphinxstepscope


\section{🧪 第4章 仪器数据支持}
\label{\detokenize{chapters/chapter4:id1}}\label{\detokenize{chapters/chapter4::doc}}
\sphinxAtStartPar
本章介绍MTDP支持的各类MT仪器数据格式及其导入处理流程。


\bigskip\hrule\bigskip



\subsection{4.1 Phoenix MTU\sphinxhyphen{}5/5A (TBL格式)}
\label{\detokenize{chapters/chapter4:phoenix-mtu-5-5a-tbl}}

\subsubsection{4.1.1 支持的仪器型号}
\label{\detokenize{chapters/chapter4:id2}}\begin{itemize}
\item {} 
\sphinxAtStartPar
\sphinxstylestrong{MTU\sphinxhyphen{}5}：经典款MT采集系统

\item {} 
\sphinxAtStartPar
\sphinxstylestrong{MTU\sphinxhyphen{}5A}：升级版本，支持更高采样率

\end{itemize}


\subsubsection{4.1.2 数据文件}
\label{\detokenize{chapters/chapter4:id3}}

\begin{savenotes}\sphinxattablestart
\sphinxthistablewithglobalstyle
\centering
\begin{tabulary}{\linewidth}[t]{TT}
\sphinxtoprule
\begin{varwidth}[t]{\sphinxcolwidth{1}{2}}
\sphinxstyletheadfamily \sphinxAtStartPar
文件
\sphinxbeforeendvarwidth
\end{varwidth}%
&\begin{varwidth}[t]{\sphinxcolwidth{1}{2}}
\sphinxstyletheadfamily \sphinxAtStartPar
说明
\sphinxbeforeendvarwidth
\end{varwidth}%
\\
\sphinxmidrule
\sphinxtableatstartofbodyhook\begin{varwidth}[t]{\sphinxcolwidth{1}{2}}
\sphinxAtStartPar
.tbl
\sphinxbeforeendvarwidth
\end{varwidth}%
&\begin{varwidth}[t]{\sphinxcolwidth{1}{2}}
\sphinxAtStartPar
头文件，记录采集参数
\sphinxbeforeendvarwidth
\end{varwidth}%
\\
\sphinxhline\begin{varwidth}[t]{\sphinxcolwidth{1}{2}}
\sphinxAtStartPar
.ts
\sphinxbeforeendvarwidth
\end{varwidth}%
&\begin{varwidth}[t]{\sphinxcolwidth{1}{2}}
\sphinxAtStartPar
时间序列数据
\sphinxbeforeendvarwidth
\end{varwidth}%
\\
\sphinxhline\begin{varwidth}[t]{\sphinxcolwidth{1}{2}}
\sphinxAtStartPar
.mt
\sphinxbeforeendvarwidth
\end{varwidth}%
&\begin{varwidth}[t]{\sphinxcolwidth{1}{2}}
\sphinxAtStartPar
连续时间序列数据
\sphinxbeforeendvarwidth
\end{varwidth}%
\\
\sphinxbottomrule
\end{tabulary}
\sphinxtableafterendhook\par
\sphinxattableend\end{savenotes}


\subsubsection{4.1.3 Phoenix MTU\sphinxhyphen{}5/5A 数据导入（详细步骤）}
\label{\detokenize{chapters/chapter4:phoenix-mtu-5-5a}}

\paragraph{4.1.3.1 准备工作}
\label{\detokenize{chapters/chapter4:id4}}
\sphinxAtStartPar
\sphinxstylestrong{步骤1：检查数据文件}
\begin{itemize}
\item {} 
\sphinxAtStartPar
 确认数据目录存在

\item {} 
\sphinxAtStartPar
 检查.tbl头文件完整性

\item {} 
\sphinxAtStartPar
 验证标定文件（CLB/CLC）是否齐全

\item {} 
\sphinxAtStartPar
 确认.ts时间序列文件完整

\end{itemize}

\sphinxAtStartPar
\sphinxstylestrong{步骤2：在MTDP中准备工程结构}
\begin{itemize}
\item {} 
\sphinxAtStartPar
 创建或选择目标工区

\item {} 
\sphinxAtStartPar
 在工区下创建测段

\item {} 
\sphinxAtStartPar
 确认工程路径设置正确

\end{itemize}


\paragraph{4.1.3.2 导入流程}
\label{\detokenize{chapters/chapter4:id5}}
\sphinxAtStartPar
\sphinxstylestrong{步骤3：打开导入对话框}
\begin{enumerate}
\sphinxsetlistlabels{\arabic}{enumi}{enumii}{}{.}%
\item {} 
\sphinxAtStartPar
在工程树中选择目标测段

\item {} 
\sphinxAtStartPar
右键选择 \sphinxcode{\sphinxupquote{加载测点 → 从目录加载Phoenix测点}}

\item {} 
\sphinxAtStartPar
系统打开Phoenix测点导入对话框

\end{enumerate}

\sphinxAtStartPar
\sphinxstylestrong{步骤4：配置导入参数}

\sphinxAtStartPar
在导入对话框中：


\begin{savenotes}\sphinxattablestart
\sphinxthistablewithglobalstyle
\centering
\begin{tabulary}{\linewidth}[t]{TTT}
\sphinxtoprule
\begin{varwidth}[t]{\sphinxcolwidth{1}{3}}
\sphinxstyletheadfamily \sphinxAtStartPar
参数
\sphinxbeforeendvarwidth
\end{varwidth}%
&\begin{varwidth}[t]{\sphinxcolwidth{1}{3}}
\sphinxstyletheadfamily \sphinxAtStartPar
说明
\sphinxbeforeendvarwidth
\end{varwidth}%
&\begin{varwidth}[t]{\sphinxcolwidth{1}{3}}
\sphinxstyletheadfamily \sphinxAtStartPar
默认值
\sphinxbeforeendvarwidth
\end{varwidth}%
\\
\sphinxmidrule
\sphinxtableatstartofbodyhook\begin{varwidth}[t]{\sphinxcolwidth{1}{3}}
\sphinxAtStartPar
数据目录
\sphinxbeforeendvarwidth
\end{varwidth}%
&\begin{varwidth}[t]{\sphinxcolwidth{1}{3}}
\sphinxAtStartPar
选择Phoenix数据所在文件夹
\sphinxbeforeendvarwidth
\end{varwidth}%
&\begin{varwidth}[t]{\sphinxcolwidth{1}{3}}
\sphinxAtStartPar
系统记住上次位置
\sphinxbeforeendvarwidth
\end{varwidth}%
\\
\sphinxhline\begin{varwidth}[t]{\sphinxcolwidth{1}{3}}
\sphinxAtStartPar
递归搜索
\sphinxbeforeendvarwidth
\end{varwidth}%
&\begin{varwidth}[t]{\sphinxcolwidth{1}{3}}
\sphinxAtStartPar
勾选是否搜索子目录
\sphinxbeforeendvarwidth
\end{varwidth}%
&\begin{varwidth}[t]{\sphinxcolwidth{1}{3}}
\sphinxAtStartPar
勾选
\sphinxbeforeendvarwidth
\end{varwidth}%
\\
\sphinxhline\begin{varwidth}[t]{\sphinxcolwidth{1}{3}}
\sphinxAtStartPar
文件格式
\sphinxbeforeendvarwidth
\end{varwidth}%
&\begin{varwidth}[t]{\sphinxcolwidth{1}{3}}
\sphinxAtStartPar
自动识别
\sphinxbeforeendvarwidth
\end{varwidth}%
&\begin{varwidth}[t]{\sphinxcolwidth{1}{3}}
\sphinxAtStartPar
自动
\sphinxbeforeendvarwidth
\end{varwidth}%
\\
\sphinxhline\begin{varwidth}[t]{\sphinxcolwidth{1}{3}}
\sphinxAtStartPar
测点命名规则
\sphinxbeforeendvarwidth
\end{varwidth}%
&\begin{varwidth}[t]{\sphinxcolwidth{1}{3}}
\sphinxAtStartPar
从文件名提取
\sphinxbeforeendvarwidth
\end{varwidth}%
&\begin{varwidth}[t]{\sphinxcolwidth{1}{3}}
\sphinxAtStartPar
从文件名提取
\sphinxbeforeendvarwidth
\end{varwidth}%
\\
\sphinxhline\begin{varwidth}[t]{\sphinxcolwidth{1}{3}}
\sphinxAtStartPar
标定文件路径
\sphinxbeforeendvarwidth
\end{varwidth}%
&\begin{varwidth}[t]{\sphinxcolwidth{1}{3}}
\sphinxAtStartPar
自动搜索
\sphinxbeforeendvarwidth
\end{varwidth}%
&\begin{varwidth}[t]{\sphinxcolwidth{1}{3}}
\sphinxAtStartPar
系统搜索路径
\sphinxbeforeendvarwidth
\end{varwidth}%
\\
\sphinxbottomrule
\end{tabulary}
\sphinxtableafterendhook\par
\sphinxattableend\end{savenotes}

\sphinxAtStartPar
\sphinxstylestrong{步骤5：执行导入}

\sphinxAtStartPar
点击"确定"后，系统将：
\begin{enumerate}
\sphinxsetlistlabels{\arabic}{enumi}{enumii}{}{.}%
\item {} 
\sphinxAtStartPar
扫描目录中的.tbl文件

\item {} 
\sphinxAtStartPar
解析每个.tbl头文件获取测点信息

\item {} 
\sphinxAtStartPar
提取时间序列文件路径

\item {} 
\sphinxAtStartPar
创建测点结构并添加到工程树

\item {} 
\sphinxAtStartPar
显示导入进度

\end{enumerate}


\paragraph{4.1.3.3 导入后验证}
\label{\detokenize{chapters/chapter4:id6}}
\sphinxAtStartPar
\sphinxstylestrong{检查项目：}


\begin{savenotes}\sphinxattablestart
\sphinxthistablewithglobalstyle
\centering
\begin{tabulary}{\linewidth}[t]{TT}
\sphinxtoprule
\begin{varwidth}[t]{\sphinxcolwidth{1}{2}}
\sphinxstyletheadfamily \sphinxAtStartPar
检查项
\sphinxbeforeendvarwidth
\end{varwidth}%
&\begin{varwidth}[t]{\sphinxcolwidth{1}{2}}
\sphinxstyletheadfamily \sphinxAtStartPar
验证方法
\sphinxbeforeendvarwidth
\end{varwidth}%
\\
\sphinxmidrule
\sphinxtableatstartofbodyhook\begin{varwidth}[t]{\sphinxcolwidth{1}{2}}
\sphinxAtStartPar
测点数量
\sphinxbeforeendvarwidth
\end{varwidth}%
&\begin{varwidth}[t]{\sphinxcolwidth{1}{2}}
\sphinxAtStartPar
与源文件数量对比
\sphinxbeforeendvarwidth
\end{varwidth}%
\\
\sphinxhline\begin{varwidth}[t]{\sphinxcolwidth{1}{2}}
\sphinxAtStartPar
采样率匹配
\sphinxbeforeendvarwidth
\end{varwidth}%
&\begin{varwidth}[t]{\sphinxcolwidth{1}{2}}
\sphinxAtStartPar
查看TS2/TS3/TS4/TS5分配
\sphinxbeforeendvarwidth
\end{varwidth}%
\\
\sphinxhline\begin{varwidth}[t]{\sphinxcolwidth{1}{2}}
\sphinxAtStartPar
标定文件关联
\sphinxbeforeendvarwidth
\end{varwidth}%
&\begin{varwidth}[t]{\sphinxcolwidth{1}{2}}
\sphinxAtStartPar
打开测点设置查看标定文件
\sphinxbeforeendvarwidth
\end{varwidth}%
\\
\sphinxhline\begin{varwidth}[t]{\sphinxcolwidth{1}{2}}
\sphinxAtStartPar
坐标信息
\sphinxbeforeendvarwidth
\end{varwidth}%
&\begin{varwidth}[t]{\sphinxcolwidth{1}{2}}
\sphinxAtStartPar
确认经纬度正确
\sphinxbeforeendvarwidth
\end{varwidth}%
\\
\sphinxbottomrule
\end{tabulary}
\sphinxtableafterendhook\par
\sphinxattableend\end{savenotes}


\paragraph{4.1.3.4 常见问题与解决}
\label{\detokenize{chapters/chapter4:id7}}
\sphinxAtStartPar
\sphinxstylestrong{问题1：提示"文件格式不识别"}
\begin{itemize}
\item {} 
\sphinxAtStartPar
原因：不是标准的Phoenix格式

\item {} 
\sphinxAtStartPar
解决：
\begin{enumerate}
\sphinxsetlistlabels{\arabic}{enumi}{enumii}{}{.}%
\item {} 
\sphinxAtStartPar
检查.tbl文件是否为空或损坏

\item {} 
\sphinxAtStartPar
确认.tbl头文件中必需字段（BoxID、采样率等）完整

\item {} 
\sphinxAtStartPar
尝试使用Phoenix官方工具重新导出数据

\end{enumerate}

\end{itemize}

\sphinxAtStartPar
\sphinxstylestrong{问题2：导入后时间序列显示异常}
\begin{itemize}
\item {} 
\sphinxAtStartPar
原因：采样率设置不匹配或文件损坏

\item {} 
\sphinxAtStartPar
解决：
\begin{enumerate}
\sphinxsetlistlabels{\arabic}{enumi}{enumii}{}{.}%
\item {} 
\sphinxAtStartPar
检查.ts文件大小是否合理

\item {} 
\sphinxAtStartPar
验证采样率与.tbl头文件一致

\item {} 
\sphinxAtStartPar
尝试使用\sphinxcode{\sphinxupquote{工具 → 时间序列查看器}}单独查看数据

\end{enumerate}

\end{itemize}

\sphinxAtStartPar
\sphinxstylestrong{问题3：标定文件未找到}
\begin{itemize}
\item {} 
\sphinxAtStartPar
原因：标定文件不在标准搜索路径

\item {} 
\sphinxAtStartPar
解决：
\begin{enumerate}
\sphinxsetlistlabels{\arabic}{enumi}{enumii}{}{.}%
\item {} 
\sphinxAtStartPar
检查采集盒序列号与标定文件名匹配

\item {} 
\sphinxAtStartPar
将标定文件复制到工程目录

\item {} 
\sphinxAtStartPar
使用\sphinxcode{\sphinxupquote{设置 → 管理标定路径}}添加标定文件目录

\end{enumerate}

\end{itemize}

\sphinxAtStartPar
\sphinxstylestrong{问题4：导入速度缓慢}
\begin{itemize}
\item {} 
\sphinxAtStartPar
原因：数据量大或磁盘IO性能问题

\item {} 
\sphinxAtStartPar
解决：
\begin{enumerate}
\sphinxsetlistlabels{\arabic}{enumi}{enumii}{}{.}%
\item {} 
\sphinxAtStartPar
关闭其他程序减少系统负载

\item {} 
\sphinxAtStartPar
分批导入而非一次性导入全部测点

\item {} 
\sphinxAtStartPar
使用SSD存储工程文件

\end{enumerate}

\end{itemize}


\paragraph{4.1.3.5 最佳实践}
\label{\detokenize{chapters/chapter4:id8}}
\sphinxAtStartPar
\sphinxstylestrong{文件组织建议：}
\begin{itemize}
\item {} 
\sphinxAtStartPar
将同一测区的所有Phoenix数据放在同一文件夹

\item {} 
\sphinxAtStartPar
使用有意义的文件夹命名（如工区名称\_测线）

\item {} 
\sphinxAtStartPar
备份原始数据文件到单独目录

\item {} 
\sphinxAtStartPar
定期清理不再使用的临时文件

\end{itemize}

\sphinxAtStartPar
\sphinxstylestrong{导入工作流建议：}
\begin{enumerate}
\sphinxsetlistlabels{\arabic}{enumi}{enumii}{}{.}%
\item {} 
\sphinxAtStartPar
先导入少量测试数据验证流程

\item {} 
\sphinxAtStartPar
确认导入结果后再批量导入剩余数据

\item {} 
\sphinxAtStartPar
导入后立即检查测点数量和质量

\item {} 
\sphinxAtStartPar
保存工程备份以便回溯

\end{enumerate}


\subsubsection{4.1.4 标定文件}
\label{\detokenize{chapters/chapter4:id9}}

\begin{savenotes}\sphinxattablestart
\sphinxthistablewithglobalstyle
\centering
\begin{tabulary}{\linewidth}[t]{TT}
\sphinxtoprule
\begin{varwidth}[t]{\sphinxcolwidth{1}{2}}
\sphinxstyletheadfamily \sphinxAtStartPar
文件
\sphinxbeforeendvarwidth
\end{varwidth}%
&\begin{varwidth}[t]{\sphinxcolwidth{1}{2}}
\sphinxstyletheadfamily \sphinxAtStartPar
用途
\sphinxbeforeendvarwidth
\end{varwidth}%
\\
\sphinxmidrule
\sphinxtableatstartofbodyhook\begin{varwidth}[t]{\sphinxcolwidth{1}{2}}
\sphinxAtStartPar
CLB
\sphinxbeforeendvarwidth
\end{varwidth}%
&\begin{varwidth}[t]{\sphinxcolwidth{1}{2}}
\sphinxAtStartPar
电场盒标定
\sphinxbeforeendvarwidth
\end{varwidth}%
\\
\sphinxhline\begin{varwidth}[t]{\sphinxcolwidth{1}{2}}
\sphinxAtStartPar
CLC
\sphinxbeforeendvarwidth
\end{varwidth}%
&\begin{varwidth}[t]{\sphinxcolwidth{1}{2}}
\sphinxAtStartPar
磁传感器标定
\sphinxbeforeendvarwidth
\end{varwidth}%
\\
\sphinxbottomrule
\end{tabulary}
\sphinxtableafterendhook\par
\sphinxattableend\end{savenotes}
\begin{quote}

\sphinxAtStartPar
💡 \sphinxstylestrong{提示}：确保标定文件与采集盒序列号匹配。
\end{quote}


\subsubsection{4.1.5 专用工具}
\label{\detokenize{chapters/chapter4:id10}}
\sphinxAtStartPar
Phoenix仪器提供多种专用处理工具：


\paragraph{Phoenix数据处理 (PhoenixDataProcessForm)}
\label{\detokenize{chapters/chapter4:phoenix-phoenixdataprocessform}}
\sphinxAtStartPar
\sphinxstylestrong{功能：}
\begin{itemize}
\item {} 
\sphinxAtStartPar
TBL文件管理和编辑

\item {} 
\sphinxAtStartPar
标定文件关联（CLB/CLC）

\item {} 
\sphinxAtStartPar
图表查看和分析

\end{itemize}

\sphinxAtStartPar
\sphinxstylestrong{操作步骤：}
\begin{enumerate}
\sphinxsetlistlabels{\arabic}{enumi}{enumii}{}{.}%
\item {} 
\sphinxAtStartPar
选择菜单 \sphinxcode{\sphinxupquote{工具 → Phoenix工具 → 数据处理}}

\item {} 
\sphinxAtStartPar
加载TBL文件

\item {} 
\sphinxAtStartPar
配置标定文件路径

\item {} 
\sphinxAtStartPar
查看时间序列和频谱

\end{enumerate}


\paragraph{SSMT2000处理 (PhoenixNormalDataProcessForm)}
\label{\detokenize{chapters/chapter4:ssmt2000-phoenixnormaldataprocessform}}
\sphinxAtStartPar
\sphinxstylestrong{支持格式：}


\begin{savenotes}\sphinxattablestart
\sphinxthistablewithglobalstyle
\centering
\begin{tabulary}{\linewidth}[t]{TT}
\sphinxtoprule
\begin{varwidth}[t]{\sphinxcolwidth{1}{2}}
\sphinxstyletheadfamily \sphinxAtStartPar
格式
\sphinxbeforeendvarwidth
\end{varwidth}%
&\begin{varwidth}[t]{\sphinxcolwidth{1}{2}}
\sphinxstyletheadfamily \sphinxAtStartPar
说明
\sphinxbeforeendvarwidth
\end{varwidth}%
\\
\sphinxmidrule
\sphinxtableatstartofbodyhook\begin{varwidth}[t]{\sphinxcolwidth{1}{2}}
\sphinxAtStartPar
CLB
\sphinxbeforeendvarwidth
\end{varwidth}%
&\begin{varwidth}[t]{\sphinxcolwidth{1}{2}}
\sphinxAtStartPar
电场盒标定文件
\sphinxbeforeendvarwidth
\end{varwidth}%
\\
\sphinxhline\begin{varwidth}[t]{\sphinxcolwidth{1}{2}}
\sphinxAtStartPar
CLC
\sphinxbeforeendvarwidth
\end{varwidth}%
&\begin{varwidth}[t]{\sphinxcolwidth{1}{2}}
\sphinxAtStartPar
磁传感器标定文件
\sphinxbeforeendvarwidth
\end{varwidth}%
\\
\sphinxhline\begin{varwidth}[t]{\sphinxcolwidth{1}{2}}
\sphinxAtStartPar
PFT
\sphinxbeforeendvarwidth
\end{varwidth}%
&\begin{varwidth}[t]{\sphinxcolwidth{1}{2}}
\sphinxAtStartPar
功率谱文件
\sphinxbeforeendvarwidth
\end{varwidth}%
\\
\sphinxhline\begin{varwidth}[t]{\sphinxcolwidth{1}{2}}
\sphinxAtStartPar
PRM
\sphinxbeforeendvarwidth
\end{varwidth}%
&\begin{varwidth}[t]{\sphinxcolwidth{1}{2}}
\sphinxAtStartPar
参数文件
\sphinxbeforeendvarwidth
\end{varwidth}%
\\
\sphinxhline\begin{varwidth}[t]{\sphinxcolwidth{1}{2}}
\sphinxAtStartPar
FC
\sphinxbeforeendvarwidth
\end{varwidth}%
&\begin{varwidth}[t]{\sphinxcolwidth{1}{2}}
\sphinxAtStartPar
傅里叶系数文件
\sphinxbeforeendvarwidth
\end{varwidth}%
\\
\sphinxhline\begin{varwidth}[t]{\sphinxcolwidth{1}{2}}
\sphinxAtStartPar
MT
\sphinxbeforeendvarwidth
\end{varwidth}%
&\begin{varwidth}[t]{\sphinxcolwidth{1}{2}}
\sphinxAtStartPar
MT响应文件
\sphinxbeforeendvarwidth
\end{varwidth}%
\\
\sphinxhline\begin{varwidth}[t]{\sphinxcolwidth{1}{2}}
\sphinxAtStartPar
EDI
\sphinxbeforeendvarwidth
\end{varwidth}%
&\begin{varwidth}[t]{\sphinxcolwidth{1}{2}}
\sphinxAtStartPar
EDI格式文件
\sphinxbeforeendvarwidth
\end{varwidth}%
\\
\sphinxbottomrule
\end{tabulary}
\sphinxtableafterendhook\par
\sphinxattableend\end{savenotes}


\paragraph{数据合并 (PhoenixDataMergeForm)}
\label{\detokenize{chapters/chapter4:phoenixdatamergeform}}
\sphinxAtStartPar
\sphinxstylestrong{功能：}
\begin{itemize}
\item {} 
\sphinxAtStartPar
MTU\sphinxhyphen{}5A数据合并

\item {} 
\sphinxAtStartPar
多文件自动合并

\item {} 
\sphinxAtStartPar
时间对齐

\end{itemize}

\sphinxAtStartPar
\sphinxstylestrong{操作步骤：}
\begin{enumerate}
\sphinxsetlistlabels{\arabic}{enumi}{enumii}{}{.}%
\item {} 
\sphinxAtStartPar
选择菜单 \sphinxcode{\sphinxupquote{工具 → Phoenix TS合并}}

\item {} 
\sphinxAtStartPar
选择要合并的文件

\item {} 
\sphinxAtStartPar
设置合并参数

\item {} 
\sphinxAtStartPar
执行合并

\end{enumerate}


\paragraph{时间序列分割 (PhoenixTSSplit)}
\label{\detokenize{chapters/chapter4:phoenixtssplit}}
\sphinxAtStartPar
\sphinxstylestrong{功能：}
\begin{itemize}
\item {} 
\sphinxAtStartPar
分割长时间连续数据

\item {} 
\sphinxAtStartPar
按时间段分割

\item {} 
\sphinxAtStartPar
段管理

\end{itemize}


\paragraph{通道重排 (PhoenixChannelReArrengeForm)}
\label{\detokenize{chapters/chapter4:phoenixchannelrearrengeform}}
\sphinxAtStartPar
\sphinxstylestrong{通道控制：}


\begin{savenotes}\sphinxattablestart
\sphinxthistablewithglobalstyle
\centering
\begin{tabulary}{\linewidth}[t]{TT}
\sphinxtoprule
\begin{varwidth}[t]{\sphinxcolwidth{1}{2}}
\sphinxstyletheadfamily \sphinxAtStartPar
通道
\sphinxbeforeendvarwidth
\end{varwidth}%
&\begin{varwidth}[t]{\sphinxcolwidth{1}{2}}
\sphinxstyletheadfamily \sphinxAtStartPar
功能
\sphinxbeforeendvarwidth
\end{varwidth}%
\\
\sphinxmidrule
\sphinxtableatstartofbodyhook\begin{varwidth}[t]{\sphinxcolwidth{1}{2}}
\sphinxAtStartPar
Ex
\sphinxbeforeendvarwidth
\end{varwidth}%
&\begin{varwidth}[t]{\sphinxcolwidth{1}{2}}
\sphinxAtStartPar
X方向电场
\sphinxbeforeendvarwidth
\end{varwidth}%
\\
\sphinxhline\begin{varwidth}[t]{\sphinxcolwidth{1}{2}}
\sphinxAtStartPar
Ey
\sphinxbeforeendvarwidth
\end{varwidth}%
&\begin{varwidth}[t]{\sphinxcolwidth{1}{2}}
\sphinxAtStartPar
Y方向电场
\sphinxbeforeendvarwidth
\end{varwidth}%
\\
\sphinxhline\begin{varwidth}[t]{\sphinxcolwidth{1}{2}}
\sphinxAtStartPar
Hx
\sphinxbeforeendvarwidth
\end{varwidth}%
&\begin{varwidth}[t]{\sphinxcolwidth{1}{2}}
\sphinxAtStartPar
X方向磁场
\sphinxbeforeendvarwidth
\end{varwidth}%
\\
\sphinxhline\begin{varwidth}[t]{\sphinxcolwidth{1}{2}}
\sphinxAtStartPar
Hy
\sphinxbeforeendvarwidth
\end{varwidth}%
&\begin{varwidth}[t]{\sphinxcolwidth{1}{2}}
\sphinxAtStartPar
Y方向磁场
\sphinxbeforeendvarwidth
\end{varwidth}%
\\
\sphinxhline\begin{varwidth}[t]{\sphinxcolwidth{1}{2}}
\sphinxAtStartPar
Hz
\sphinxbeforeendvarwidth
\end{varwidth}%
&\begin{varwidth}[t]{\sphinxcolwidth{1}{2}}
\sphinxAtStartPar
Z方向磁场
\sphinxbeforeendvarwidth
\end{varwidth}%
\\
\sphinxbottomrule
\end{tabulary}
\sphinxtableafterendhook\par
\sphinxattableend\end{savenotes}

\sphinxAtStartPar
**操作：**拖动通道滑块调整顺序


\paragraph{时间序列对齐修复 (PhoenixTSMisalignmentRepairForm)}
\label{\detokenize{chapters/chapter4:phoenixtsmisalignmentrepairform}}
\sphinxAtStartPar
**功能：**修复时钟同步问题导致的时间序列错位


\paragraph{TBL参数替换 (TBLParameterReplaceForm)}
\label{\detokenize{chapters/chapter4:tbl-tblparameterreplaceform}}
\sphinxAtStartPar
**功能：**批量替换TBL文件中的参数值


\paragraph{TS→FC转换 (PhoenixTsToFtMenu)}
\label{\detokenize{chapters/chapter4:tsfc-phoenixtstoftmenu}}
\sphinxAtStartPar
**功能：**将时间序列转换为傅里叶系数文件


\bigskip\hrule\bigskip



\subsection{4.2 Phoenix MTU\sphinxhyphen{}5C/8A (JSON格式)}
\label{\detokenize{chapters/chapter4:phoenix-mtu-5c-8a-json}}

\subsubsection{4.2.1 支持的仪器型号}
\label{\detokenize{chapters/chapter4:id11}}\begin{itemize}
\item {} 
\sphinxAtStartPar
\sphinxstylestrong{MTU\sphinxhyphen{}5C}：新一代宽带MT采集系统

\item {} 
\sphinxAtStartPar
\sphinxstylestrong{MTU\sphinxhyphen{}8A}：多通道MT采集系统

\end{itemize}


\subsubsection{4.2.2 数据导入}
\label{\detokenize{chapters/chapter4:id12}}\begin{enumerate}
\sphinxsetlistlabels{\arabic}{enumi}{enumii}{}{.}%
\item {} 
\sphinxAtStartPar
选择目标测段

\item {} 
\sphinxAtStartPar
右键选择 \sphinxcode{\sphinxupquote{加载测点 → 从目录加载MTU测点}}

\item {} 
\sphinxAtStartPar
选择JSON数据目录

\item {} 
\sphinxAtStartPar
确认导入

\end{enumerate}


\subsubsection{4.2.3 降采样处理}
\label{\detokenize{chapters/chapter4:id13}}
\sphinxAtStartPar
高采样率数据可降采样处理：
\begin{enumerate}
\sphinxsetlistlabels{\arabic}{enumi}{enumii}{}{.}%
\item {} 
\sphinxAtStartPar
选择菜单 \sphinxcode{\sphinxupquote{工具 → MTU工具 → 降采样}}

\item {} 
\sphinxAtStartPar
设置目标采样率

\item {} 
\sphinxAtStartPar
执行降采样

\end{enumerate}


\bigskip\hrule\bigskip



\subsection{4.3 Metronix仪器}
\label{\detokenize{chapters/chapter4:metronix}}

\subsubsection{4.3.1 支持的仪器型号}
\label{\detokenize{chapters/chapter4:id14}}\begin{itemize}
\item {} 
\sphinxAtStartPar
ADU\sphinxhyphen{}06

\item {} 
\sphinxAtStartPar
ADU\sphinxhyphen{}07

\item {} 
\sphinxAtStartPar
MMS系列

\end{itemize}


\subsubsection{4.3.2 数据文件}
\label{\detokenize{chapters/chapter4:id15}}

\begin{savenotes}\sphinxattablestart
\sphinxthistablewithglobalstyle
\centering
\begin{tabulary}{\linewidth}[t]{TT}
\sphinxtoprule
\begin{varwidth}[t]{\sphinxcolwidth{1}{2}}
\sphinxstyletheadfamily \sphinxAtStartPar
文件
\sphinxbeforeendvarwidth
\end{varwidth}%
&\begin{varwidth}[t]{\sphinxcolwidth{1}{2}}
\sphinxstyletheadfamily \sphinxAtStartPar
说明
\sphinxbeforeendvarwidth
\end{varwidth}%
\\
\sphinxmidrule
\sphinxtableatstartofbodyhook\begin{varwidth}[t]{\sphinxcolwidth{1}{2}}
\sphinxAtStartPar
.ats
\sphinxbeforeendvarwidth
\end{varwidth}%
&\begin{varwidth}[t]{\sphinxcolwidth{1}{2}}
\sphinxAtStartPar
时间序列数据
\sphinxbeforeendvarwidth
\end{varwidth}%
\\
\sphinxhline\begin{varwidth}[t]{\sphinxcolwidth{1}{2}}
\sphinxAtStartPar
.atm
\sphinxbeforeendvarwidth
\end{varwidth}%
&\begin{varwidth}[t]{\sphinxcolwidth{1}{2}}
\sphinxAtStartPar
ATM格式数据
\sphinxbeforeendvarwidth
\end{varwidth}%
\\
\sphinxhline\begin{varwidth}[t]{\sphinxcolwidth{1}{2}}
\sphinxAtStartPar
.cal
\sphinxbeforeendvarwidth
\end{varwidth}%
&\begin{varwidth}[t]{\sphinxcolwidth{1}{2}}
\sphinxAtStartPar
标定文件
\sphinxbeforeendvarwidth
\end{varwidth}%
\\
\sphinxbottomrule
\end{tabulary}
\sphinxtableafterendhook\par
\sphinxattableend\end{savenotes}


\subsubsection{4.3.3 数据导入}
\label{\detokenize{chapters/chapter4:id16}}\begin{enumerate}
\sphinxsetlistlabels{\arabic}{enumi}{enumii}{}{.}%
\item {} 
\sphinxAtStartPar
选择目标测段

\item {} 
\sphinxAtStartPar
右键选择 \sphinxcode{\sphinxupquote{加载测点 → 加载Metronix测点}}

\item {} 
\sphinxAtStartPar
选择数据目录

\item {} 
\sphinxAtStartPar
配置标定文件路径

\item {} 
\sphinxAtStartPar
执行导入

\end{enumerate}


\subsubsection{4.3.4 WEM数据处理器}
\label{\detokenize{chapters/chapter4:wem}}\begin{enumerate}
\sphinxsetlistlabels{\arabic}{enumi}{enumii}{}{.}%
\item {} 
\sphinxAtStartPar
选择菜单 \sphinxcode{\sphinxupquote{工具 → Metronix工具 → WEM处理器}}

\item {} 
\sphinxAtStartPar
加载WEM格式数据

\item {} 
\sphinxAtStartPar
执行处理

\end{enumerate}


\subsubsection{4.3.5 Metronix专用工具}
\label{\detokenize{chapters/chapter4:id17}}

\paragraph{WEM处理器 (WEMProcessorForm)}
\label{\detokenize{chapters/chapter4:wem-wemprocessorform}}
\sphinxAtStartPar
**功能：**处理Western Electromagnetic数据

\sphinxAtStartPar
\sphinxstylestrong{操作步骤：}
\begin{enumerate}
\sphinxsetlistlabels{\arabic}{enumi}{enumii}{}{.}%
\item {} 
\sphinxAtStartPar
选择菜单 \sphinxcode{\sphinxupquote{工具 → Metronix工具 → WEM处理器}}

\item {} 
\sphinxAtStartPar
加载WEM格式数据

\item {} 
\sphinxAtStartPar
执行处理

\item {} 
\sphinxAtStartPar
导出结果

\end{enumerate}


\paragraph{ATS→DAT转换}
\label{\detokenize{chapters/chapter4:atsdat}}
\sphinxAtStartPar
**功能：**将ATS格式转换为DAT格式

\sphinxAtStartPar
**操作：**选择菜单 \sphinxcode{\sphinxupquote{工具 → Metronix工具 → ATS→DAT}}


\bigskip\hrule\bigskip



\subsection{4.4 LEMI长周期仪器}
\label{\detokenize{chapters/chapter4:lemi}}

\subsubsection{4.4.1 数据特点}
\label{\detokenize{chapters/chapter4:id18}}
\sphinxAtStartPar
LEMI仪器适合长周期MT观测，记录周期可达数万秒。


\subsubsection{4.4.2 数据导入}
\label{\detokenize{chapters/chapter4:id19}}\begin{enumerate}
\sphinxsetlistlabels{\arabic}{enumi}{enumii}{}{.}%
\item {} 
\sphinxAtStartPar
选择目标测段

\item {} 
\sphinxAtStartPar
右键选择 \sphinxcode{\sphinxupquote{加载测点 → 加载LEMI测点}}

\item {} 
\sphinxAtStartPar
选择数据目录

\item {} 
\sphinxAtStartPar
确认通道映射

\item {} 
\sphinxAtStartPar
执行导入

\end{enumerate}


\subsubsection{4.4.3 与高频数据联合处理}
\label{\detokenize{chapters/chapter4:id20}}
\begin{sphinxVerbatim}[commandchars=\\\{\}]

1. 分别导入高频和长周期数据
2. 分别处理两个频段
3. 导出全频段结果
\PYGZsh{}\PYGZsh{}\PYGZsh{} 4.4.4 LEMI专用工具

\PYGZsh{}\PYGZsh{}\PYGZsh{}\PYGZsh{} 测点合并 (LEMISiteMergeForm)

**功能：**合并多个LEMI测点数据

**操作步骤：**
1. 选择菜单 `工具 → LEMI测点合并`
2. 选择要合并的文件
3. 执行合并

\PYGZsh{}\PYGZsh{}\PYGZsh{}\PYGZsh{} H自动恢复 (LEMISiteHAutoRestoreMenu)

**功能：**自动恢复LEMI磁场数据

\PYGZhy{}\PYGZhy{}\PYGZhy{}

\PYGZsh{}\PYGZsh{} 4.5 RMT/CASRMT格式

\PYGZsh{}\PYGZsh{}\PYGZsh{} 4.5.1 RMT与标准MT处理流程的区别

RMT（Radio Magnetotellurics）数据处理流程与标准MT数据有本质区别：

```mermaid
graph TB
    subgraph 标准MT处理流程
        A1[原始时间序列] \PYGZhy{}\PYGZhy{}\PYGZgt{} A2[FFT频谱分析]
        A2 \PYGZhy{}\PYGZhy{}\PYGZgt{} A3[校准处理]
        A3 \PYGZhy{}\PYGZhy{}\PYGZgt{} A4[传递函数估计]
        A4 \PYGZhy{}\PYGZhy{}\PYGZgt{} A5[阻抗/相位计算]
        A5 \PYGZhy{}\PYGZhy{}\PYGZgt{} A6[视电阻率计算]
    end
    
    subgraph RMT/CASRMT处理流程
        B1[已处理数据文件] \PYGZhy{}\PYGZhy{}\PYGZgt{} B2[导入数据]
        B2 \PYGZhy{}\PYGZhy{}\PYGZgt{} B3[格式解析]
        B3 \PYGZhy{}\PYGZhy{}\PYGZgt{} B4[系统响应校正]
        B4 \PYGZhy{}\PYGZhy{}\PYGZgt{} B5[质量检查]
        B5 \PYGZhy{}\PYGZhy{}\PYGZgt{} B6[导出结果]
    end
    
    style A1 fill:\PYGZsh{}e3f2fd
    style A6 fill:\PYGZsh{}c8e6c9
    style B1 fill:\PYGZsh{}fff8e1
    style B6 fill:\PYGZsh{}c8e6c9
\end{sphinxVerbatim}

\sphinxAtStartPar
\sphinxstylestrong{关键区别：}


\begin{savenotes}\sphinxattablestart
\sphinxthistablewithglobalstyle
\centering
\begin{tabulary}{\linewidth}[t]{TTT}
\sphinxtoprule
\begin{varwidth}[t]{\sphinxcolwidth{1}{3}}
\sphinxstyletheadfamily \sphinxAtStartPar
特征
\sphinxbeforeendvarwidth
\end{varwidth}%
&\begin{varwidth}[t]{\sphinxcolwidth{1}{3}}
\sphinxstyletheadfamily \sphinxAtStartPar
标准MT
\sphinxbeforeendvarwidth
\end{varwidth}%
&\begin{varwidth}[t]{\sphinxcolwidth{1}{3}}
\sphinxstyletheadfamily \sphinxAtStartPar
RMT/CASRMT
\sphinxbeforeendvarwidth
\end{varwidth}%
\\
\sphinxmidrule
\sphinxtableatstartofbodyhook\begin{varwidth}[t]{\sphinxcolwidth{1}{3}}
\sphinxAtStartPar
数据类型
\sphinxbeforeendvarwidth
\end{varwidth}%
&\begin{varwidth}[t]{\sphinxcolwidth{1}{3}}
\sphinxAtStartPar
原始时间序列
\sphinxbeforeendvarwidth
\end{varwidth}%
&\begin{varwidth}[t]{\sphinxcolwidth{1}{3}}
\sphinxAtStartPar
已处理阻抗数据
\sphinxbeforeendvarwidth
\end{varwidth}%
\\
\sphinxhline\begin{varwidth}[t]{\sphinxcolwidth{1}{3}}
\sphinxAtStartPar
处理步骤
\sphinxbeforeendvarwidth
\end{varwidth}%
&\begin{varwidth}[t]{\sphinxcolwidth{1}{3}}
\sphinxAtStartPar
需要FFT/校准/传递函数估计
\sphinxbeforeendvarwidth
\end{varwidth}%
&\begin{varwidth}[t]{\sphinxcolwidth{1}{3}}
\sphinxAtStartPar
仅需格式解析和校正
\sphinxbeforeendvarwidth
\end{varwidth}%
\\
\sphinxhline\begin{varwidth}[t]{\sphinxcolwidth{1}{3}}
\sphinxAtStartPar
校准方式
\sphinxbeforeendvarwidth
\end{varwidth}%
&\begin{varwidth}[t]{\sphinxcolwidth{1}{3}}
\sphinxAtStartPar
多种仪器校准文件
\sphinxbeforeendvarwidth
\end{varwidth}%
&\begin{varwidth}[t]{\sphinxcolwidth{1}{3}}
\sphinxAtStartPar
内置系统响应文件
\sphinxbeforeendvarwidth
\end{varwidth}%
\\
\sphinxhline\begin{varwidth}[t]{\sphinxcolwidth{1}{3}}
\sphinxAtStartPar
数据格式
\sphinxbeforeendvarwidth
\end{varwidth}%
&\begin{varwidth}[t]{\sphinxcolwidth{1}{3}}
\sphinxAtStartPar
.ts/.mt/.ats等
\sphinxbeforeendvarwidth
\end{varwidth}%
&\begin{varwidth}[t]{\sphinxcolwidth{1}{3}}
\sphinxAtStartPar
.sbf/.tr1/.tr2/.json
\sphinxbeforeendvarwidth
\end{varwidth}%
\\
\sphinxhline\begin{varwidth}[t]{\sphinxcolwidth{1}{3}}
\sphinxAtStartPar
频率范围
\sphinxbeforeendvarwidth
\end{varwidth}%
&\begin{varwidth}[t]{\sphinxcolwidth{1}{3}}
\sphinxAtStartPar
全频段可调
\sphinxbeforeendvarwidth
\end{varwidth}%
&\begin{varwidth}[t]{\sphinxcolwidth{1}{3}}
\sphinxAtStartPar
固定频率（10Hz\sphinxhyphen{}20kHz）
\sphinxbeforeendvarwidth
\end{varwidth}%
\\
\sphinxhline\begin{varwidth}[t]{\sphinxcolwidth{1}{3}}
\sphinxAtStartPar
处理时间
\sphinxbeforeendvarwidth
\end{varwidth}%
&\begin{varwidth}[t]{\sphinxcolwidth{1}{3}}
\sphinxAtStartPar
较长
\sphinxbeforeendvarwidth
\end{varwidth}%
&\begin{varwidth}[t]{\sphinxcolwidth{1}{3}}
\sphinxAtStartPar
极短（仅导入导出）
\sphinxbeforeendvarwidth
\end{varwidth}%
\\
\sphinxbottomrule
\end{tabulary}
\sphinxtableafterendhook\par
\sphinxattableend\end{savenotes}


\subsubsection{4.5.2 数据特点}
\label{\detokenize{chapters/chapter4:id21}}
\sphinxAtStartPar
\sphinxstylestrong{RMT数据包含的信息：}

\sphinxAtStartPar
RMT数据文件已包含完整的处理结果，无需额外处理：


\begin{savenotes}\sphinxattablestart
\sphinxthistablewithglobalstyle
\centering
\begin{tabulary}{\linewidth}[t]{TTT}
\sphinxtoprule
\begin{varwidth}[t]{\sphinxcolwidth{1}{3}}
\sphinxstyletheadfamily \sphinxAtStartPar
数据类型
\sphinxbeforeendvarwidth
\end{varwidth}%
&\begin{varwidth}[t]{\sphinxcolwidth{1}{3}}
\sphinxstyletheadfamily \sphinxAtStartPar
说明
\sphinxbeforeendvarwidth
\end{varwidth}%
&\begin{varwidth}[t]{\sphinxcolwidth{1}{3}}
\sphinxstyletheadfamily \sphinxAtStartPar
用途
\sphinxbeforeendvarwidth
\end{varwidth}%
\\
\sphinxmidrule
\sphinxtableatstartofbodyhook\begin{varwidth}[t]{\sphinxcolwidth{1}{3}}
\sphinxAtStartPar
频率点
\sphinxbeforeendvarwidth
\end{varwidth}%
&\begin{varwidth}[t]{\sphinxcolwidth{1}{3}}
\sphinxAtStartPar
固定对数间隔
\sphinxbeforeendvarwidth
\end{varwidth}%
&\begin{varwidth}[t]{\sphinxcolwidth{1}{3}}
\sphinxAtStartPar
10 Hz \sphinxhyphen{} 20 kHz
\sphinxbeforeendvarwidth
\end{varwidth}%
\\
\sphinxhline\begin{varwidth}[t]{\sphinxcolwidth{1}{3}}
\sphinxAtStartPar
阻抗张量
\sphinxbeforeendvarwidth
\end{varwidth}%
&\begin{varwidth}[t]{\sphinxcolwidth{1}{3}}
\sphinxAtStartPar
Zxx, Zxy, Zyx, Zyy
\sphinxbeforeendvarwidth
\end{varwidth}%
&\begin{varwidth}[t]{\sphinxcolwidth{1}{3}}
\sphinxAtStartPar
MT响应函数
\sphinxbeforeendvarwidth
\end{varwidth}%
\\
\sphinxhline\begin{varwidth}[t]{\sphinxcolwidth{1}{3}}
\sphinxAtStartPar
视电阻率
\sphinxbeforeendvarwidth
\end{varwidth}%
&\begin{varwidth}[t]{\sphinxcolwidth{1}{3}}
\sphinxAtStartPar
ρxy, ρyx
\sphinxbeforeendvarwidth
\end{varwidth}%
&\begin{varwidth}[t]{\sphinxcolwidth{1}{3}}
\sphinxAtStartPar
电阻率解释
\sphinxbeforeendvarwidth
\end{varwidth}%
\\
\sphinxhline\begin{varwidth}[t]{\sphinxcolwidth{1}{3}}
\sphinxAtStartPar
相位
\sphinxbeforeendvarwidth
\end{varwidth}%
&\begin{varwidth}[t]{\sphinxcolwidth{1}{3}}
\sphinxAtStartPar
φxy, φyx
\sphinxbeforeendvarwidth
\end{varwidth}%
&\begin{varwidth}[t]{\sphinxcolwidth{1}{3}}
\sphinxAtStartPar
相位分析
\sphinxbeforeendvarwidth
\end{varwidth}%
\\
\sphinxhline\begin{varwidth}[t]{\sphinxcolwidth{1}{3}}
\sphinxAtStartPar
误差
\sphinxbeforeendvarwidth
\end{varwidth}%
&\begin{varwidth}[t]{\sphinxcolwidth{1}{3}}
\sphinxAtStartPar
各分量方差
\sphinxbeforeendvarwidth
\end{varwidth}%
&\begin{varwidth}[t]{\sphinxcolwidth{1}{3}}
\sphinxAtStartPar
质量评估
\sphinxbeforeendvarwidth
\end{varwidth}%
\\
\sphinxhline\begin{varwidth}[t]{\sphinxcolwidth{1}{3}}
\sphinxAtStartPar
系统响应
\sphinxbeforeendvarwidth
\end{varwidth}%
&\begin{varwidth}[t]{\sphinxcolwidth{1}{3}}
\sphinxAtStartPar
仪器特定校准
\sphinxbeforeendvarwidth
\end{varwidth}%
&\begin{varwidth}[t]{\sphinxcolwidth{1}{3}}
\sphinxAtStartPar
数据校正
\sphinxbeforeendvarwidth
\end{varwidth}%
\\
\sphinxbottomrule
\end{tabulary}
\sphinxtableafterendhook\par
\sphinxattableend\end{savenotes}


\subsubsection{4.5.3 数据导入}
\label{\detokenize{chapters/chapter4:id22}}
\sphinxAtStartPar
\sphinxstylestrong{RMT格式（SBF）：}
\begin{enumerate}
\sphinxsetlistlabels{\arabic}{enumi}{enumii}{}{.}%
\item {} 
\sphinxAtStartPar
选择工区级别

\item {} 
\sphinxAtStartPar
右键选择 \sphinxcode{\sphinxupquote{加载SBF测段}}

\item {} 
\sphinxAtStartPar
选择.sbf文件

\item {} 
\sphinxAtStartPar
系统自动解析并应用系统响应

\end{enumerate}

\sphinxAtStartPar
\sphinxstylestrong{CASRMT格式：}
\begin{enumerate}
\sphinxsetlistlabels{\arabic}{enumi}{enumii}{}{.}%
\item {} 
\sphinxAtStartPar
选择目标测段

\item {} 
\sphinxAtStartPar
右键选择 \sphinxcode{\sphinxupquote{加载测点 → 加载CASRMT测点}}

\item {} 
\sphinxAtStartPar
选择数据目录

\item {} 
\sphinxAtStartPar
系统自动识别JSON格式

\end{enumerate}


\subsubsection{4.5.4 RMT数据文件格式}
\label{\detokenize{chapters/chapter4:rmt}}

\paragraph{SBF格式 (.sbf)}
\label{\detokenize{chapters/chapter4:sbf-sbf}}
\sphinxAtStartPar
SBF（Station Binary Format）是一种紧凑的二进制格式：


\begin{savenotes}\sphinxattablestart
\sphinxthistablewithglobalstyle
\centering
\begin{tabulary}{\linewidth}[t]{TT}
\sphinxtoprule
\begin{varwidth}[t]{\sphinxcolwidth{1}{2}}
\sphinxstyletheadfamily \sphinxAtStartPar
段类型
\sphinxbeforeendvarwidth
\end{varwidth}%
&\begin{varwidth}[t]{\sphinxcolwidth{1}{2}}
\sphinxstyletheadfamily \sphinxAtStartPar
说明
\sphinxbeforeendvarwidth
\end{varwidth}%
\\
\sphinxmidrule
\sphinxtableatstartofbodyhook\begin{varwidth}[t]{\sphinxcolwidth{1}{2}}
\sphinxAtStartPar
TextSection
\sphinxbeforeendvarwidth
\end{varwidth}%
&\begin{varwidth}[t]{\sphinxcolwidth{1}{2}}
\sphinxAtStartPar
文本信息头
\sphinxbeforeendvarwidth
\end{varwidth}%
\\
\sphinxhline\begin{varwidth}[t]{\sphinxcolwidth{1}{2}}
\sphinxAtStartPar
IniSection
\sphinxbeforeendvarwidth
\end{varwidth}%
&\begin{varwidth}[t]{\sphinxcolwidth{1}{2}}
\sphinxAtStartPar
配置参数
\sphinxbeforeendvarwidth
\end{varwidth}%
\\
\sphinxhline\begin{varwidth}[t]{\sphinxcolwidth{1}{2}}
\sphinxAtStartPar
TrassSection
\sphinxbeforeendvarwidth
\end{varwidth}%
&\begin{varwidth}[t]{\sphinxcolwidth{1}{2}}
\sphinxAtStartPar
时间序列数据（原始）
\sphinxbeforeendvarwidth
\end{varwidth}%
\\
\sphinxhline\begin{varwidth}[t]{\sphinxcolwidth{1}{2}}
\sphinxAtStartPar
SpectrumSection
\sphinxbeforeendvarwidth
\end{varwidth}%
&\begin{varwidth}[t]{\sphinxcolwidth{1}{2}}
\sphinxAtStartPar
频谱数据
\sphinxbeforeendvarwidth
\end{varwidth}%
\\
\sphinxhline\begin{varwidth}[t]{\sphinxcolwidth{1}{2}}
\sphinxAtStartPar
SpectraFrequencySection
\sphinxbeforeendvarwidth
\end{varwidth}%
&\begin{varwidth}[t]{\sphinxcolwidth{1}{2}}
\sphinxAtStartPar
频率点定义
\sphinxbeforeendvarwidth
\end{varwidth}%
\\
\sphinxhline\begin{varwidth}[t]{\sphinxcolwidth{1}{2}}
\sphinxAtStartPar
DeviceParameterSection
\sphinxbeforeendvarwidth
\end{varwidth}%
&\begin{varwidth}[t]{\sphinxcolwidth{1}{2}}
\sphinxAtStartPar
设备参数
\sphinxbeforeendvarwidth
\end{varwidth}%
\\
\sphinxhline\begin{varwidth}[t]{\sphinxcolwidth{1}{2}}
\sphinxAtStartPar
RegistrarParameterSection
\sphinxbeforeendvarwidth
\end{varwidth}%
&\begin{varwidth}[t]{\sphinxcolwidth{1}{2}}
\sphinxAtStartPar
注册参数
\sphinxbeforeendvarwidth
\end{varwidth}%
\\
\sphinxbottomrule
\end{tabulary}
\sphinxtableafterendhook\par
\sphinxattableend\end{savenotes}


\paragraph{CASRMT格式 (.json/.rmtjson)}
\label{\detokenize{chapters/chapter4:casrmt-json-rmtjson}}
\sphinxAtStartPar
JSON格式便于解析和共享：

\begin{sphinxVerbatim}[commandchars=\\\{\}]
\PYGZob{}
  \PYGZdq{}station\PYGZus{}id\PYGZdq{}: \PYGZdq{}CAS001\PYGZdq{},
  \PYGZdq{}frequencies\PYGZdq{}: [10000, 8000, 6000, ...],
  \PYGZdq{}impedance\PYGZdq{}: \PYGZob{}
    \PYGZdq{}Zxx\PYGZdq{}: [\PYGZob{}\PYGZdq{}r\PYGZdq{}: 0.123, \PYGZdq{}i\PYGZdq{}: \PYGZhy{}0.045\PYGZcb{}, ...],
    \PYGZdq{}Zxy\PYGZdq{}: [\PYGZob{}\PYGZdq{}r\PYGZdq{}: 1.234, \PYGZdq{}i\PYGZdq{}: \PYGZhy{}0.567\PYGZcb{}, ...],
    \PYGZdq{}Zyx\PYGZdq{}: [\PYGZob{}\PYGZdq{}r\PYGZdq{}: \PYGZhy{}1.456, \PYGZdq{}i\PYGZdq{}: 0.789\PYGZcb{}, ...],
    \PYGZdq{}Zyy\PYGZdq{}: [\PYGZob{}\PYGZdq{}r\PYGZdq{}: 0.098, \PYGZdq{}i\PYGZdq{}: \PYGZhy{}0.032\PYGZcb{}, ...]
  \PYGZcb{},
  \PYGZdq{}apparent\PYGZus{}resistivity\PYGZdq{}: \PYGZob{}
    \PYGZdq{}rho\PYGZus{}xy\PYGZdq{}: [15.2, 18.4, 22.1, ...],
    \PYGZdq{}rho\PYGZus{}yx\PYGZdq{}: [18.5, 21.2, 25.6, ...]
  \PYGZcb{},
  \PYGZdq{}phase\PYGZdq{}: \PYGZob{}
    \PYGZdq{}phi\PYGZus{}xy\PYGZdq{}: [42.5, 45.2, 48.1, ...],
    \PYGZdq{}phi\PYGZus{}yx\PYGZdq{}: [47.3, 50.1, 52.8, ...]
  \PYGZcb{}
\PYGZcb{}
\end{sphinxVerbatim}


\subsubsection{4.5.5 系统响应校准}
\label{\detokenize{chapters/chapter4:id23}}
\sphinxAtStartPar
RMT数据需要应用仪器系统响应进行校准：

\sphinxAtStartPar
\sphinxstylestrong{校准参数：}


\begin{savenotes}\sphinxattablestart
\sphinxthistablewithglobalstyle
\centering
\begin{tabulary}{\linewidth}[t]{TTT}
\sphinxtoprule
\begin{varwidth}[t]{\sphinxcolwidth{1}{3}}
\sphinxstyletheadfamily \sphinxAtStartPar
参数
\sphinxbeforeendvarwidth
\end{varwidth}%
&\begin{varwidth}[t]{\sphinxcolwidth{1}{3}}
\sphinxstyletheadfamily \sphinxAtStartPar
说明
\sphinxbeforeendvarwidth
\end{varwidth}%
&\begin{varwidth}[t]{\sphinxcolwidth{1}{3}}
\sphinxstyletheadfamily \sphinxAtStartPar
来源
\sphinxbeforeendvarwidth
\end{varwidth}%
\\
\sphinxmidrule
\sphinxtableatstartofbodyhook\begin{varwidth}[t]{\sphinxcolwidth{1}{3}}
\sphinxAtStartPar
传感器类型
\sphinxbeforeendvarwidth
\end{varwidth}%
&\begin{varwidth}[t]{\sphinxcolwidth{1}{3}}
\sphinxAtStartPar
Hx/Hy/Hz传感器型号
\sphinxbeforeendvarwidth
\end{varwidth}%
&\begin{varwidth}[t]{\sphinxcolwidth{1}{3}}
\sphinxAtStartPar
SBF设备参数段
\sphinxbeforeendvarwidth
\end{varwidth}%
\\
\sphinxhline\begin{varwidth}[t]{\sphinxcolwidth{1}{3}}
\sphinxAtStartPar
校准日期
\sphinxbeforeendvarwidth
\end{varwidth}%
&\begin{varwidth}[t]{\sphinxcolwidth{1}{3}}
\sphinxAtStartPar
上次校准时间
\sphinxbeforeendvarwidth
\end{varwidth}%
&\begin{varwidth}[t]{\sphinxcolwidth{1}{3}}
\sphinxAtStartPar
SBF文本段
\sphinxbeforeendvarwidth
\end{varwidth}%
\\
\sphinxhline\begin{varwidth}[t]{\sphinxcolwidth{1}{3}}
\sphinxAtStartPar
幅值响应
\sphinxbeforeendvarwidth
\end{varwidth}%
&\begin{varwidth}[t]{\sphinxcolwidth{1}{3}}
\sphinxAtStartPar
频率\sphinxhyphen{}幅度曲线
\sphinxbeforeendvarwidth
\end{varwidth}%
&\begin{varwidth}[t]{\sphinxcolwidth{1}{3}}
\sphinxAtStartPar
系统响应文件
\sphinxbeforeendvarwidth
\end{varwidth}%
\\
\sphinxhline\begin{varwidth}[t]{\sphinxcolwidth{1}{3}}
\sphinxAtStartPar
相位响应
\sphinxbeforeendvarwidth
\end{varwidth}%
&\begin{varwidth}[t]{\sphinxcolwidth{1}{3}}
\sphinxAtStartPar
频率\sphinxhyphen{}相位曲线
\sphinxbeforeendvarwidth
\end{varwidth}%
&\begin{varwidth}[t]{\sphinxcolwidth{1}{3}}
\sphinxAtStartPar
系统响应文件
\sphinxbeforeendvarwidth
\end{varwidth}%
\\
\sphinxbottomrule
\end{tabulary}
\sphinxtableafterendhook\par
\sphinxattableend\end{savenotes}


\subsubsection{4.5.6 RMT/CASRMT专用工具}
\label{\detokenize{chapters/chapter4:rmt-casrmt}}

\paragraph{RMT站点管理 (RMTSiteForm)}
\label{\detokenize{chapters/chapter4:rmt-rmtsiteform}}
\sphinxAtStartPar
**功能：**管理RMT格式测点配置

\sphinxAtStartPar
\sphinxstylestrong{设置项：}


\begin{savenotes}\sphinxattablestart
\sphinxthistablewithglobalstyle
\centering
\begin{tabulary}{\linewidth}[t]{TT}
\sphinxtoprule
\begin{varwidth}[t]{\sphinxcolwidth{1}{2}}
\sphinxstyletheadfamily \sphinxAtStartPar
字段
\sphinxbeforeendvarwidth
\end{varwidth}%
&\begin{varwidth}[t]{\sphinxcolwidth{1}{2}}
\sphinxstyletheadfamily \sphinxAtStartPar
说明
\sphinxbeforeendvarwidth
\end{varwidth}%
\\
\sphinxmidrule
\sphinxtableatstartofbodyhook\begin{varwidth}[t]{\sphinxcolwidth{1}{2}}
\sphinxAtStartPar
站点名称
\sphinxbeforeendvarwidth
\end{varwidth}%
&\begin{varwidth}[t]{\sphinxcolwidth{1}{2}}
\sphinxAtStartPar
测点标识
\sphinxbeforeendvarwidth
\end{varwidth}%
\\
\sphinxhline\begin{varwidth}[t]{\sphinxcolwidth{1}{2}}
\sphinxAtStartPar
坐标信息
\sphinxbeforeendvarwidth
\end{varwidth}%
&\begin{varwidth}[t]{\sphinxcolwidth{1}{2}}
\sphinxAtStartPar
经纬度高程
\sphinxbeforeendvarwidth
\end{varwidth}%
\\
\sphinxhline\begin{varwidth}[t]{\sphinxcolwidth{1}{2}}
\sphinxAtStartPar
采集参数
\sphinxbeforeendvarwidth
\end{varwidth}%
&\begin{varwidth}[t]{\sphinxcolwidth{1}{2}}
\sphinxAtStartPar
采样率、频率范围
\sphinxbeforeendvarwidth
\end{varwidth}%
\\
\sphinxhline\begin{varwidth}[t]{\sphinxcolwidth{1}{2}}
\sphinxAtStartPar
系统响应
\sphinxbeforeendvarwidth
\end{varwidth}%
&\begin{varwidth}[t]{\sphinxcolwidth{1}{2}}
\sphinxAtStartPar
校准文件路径
\sphinxbeforeendvarwidth
\end{varwidth}%
\\
\sphinxbottomrule
\end{tabulary}
\sphinxtableafterendhook\par
\sphinxattableend\end{savenotes}


\bigskip\hrule\bigskip



\subsection{4.6 Aether仪器 (ATTS格式)}
\label{\detokenize{chapters/chapter4:aether-atts}}

\subsubsection{4.6.1 数据导入}
\label{\detokenize{chapters/chapter4:id24}}\begin{enumerate}
\sphinxsetlistlabels{\arabic}{enumi}{enumii}{}{.}%
\item {} 
\sphinxAtStartPar
选择目标测段

\item {} 
\sphinxAtStartPar
右键选择 \sphinxcode{\sphinxupquote{加载测点 → 加载ATTS测点}}

\item {} 
\sphinxAtStartPar
选择数据目录

\end{enumerate}


\subsubsection{4.6.2 降采样处理}
\label{\detokenize{chapters/chapter4:id25}}\begin{enumerate}
\sphinxsetlistlabels{\arabic}{enumi}{enumii}{}{.}%
\item {} 
\sphinxAtStartPar
选择菜单 \sphinxcode{\sphinxupquote{时间序列处理 → ATTS降采样}}

\item {} 
\sphinxAtStartPar
设置目标采样率

\item {} 
\sphinxAtStartPar
执行降采样

\end{enumerate}


\bigskip\hrule\bigskip



\subsection{4.7 通用EDI格式}
\label{\detokenize{chapters/chapter4:edi}}

\subsubsection{4.7.1 导入已有EDI数据}
\label{\detokenize{chapters/chapter4:id26}}\begin{enumerate}
\sphinxsetlistlabels{\arabic}{enumi}{enumii}{}{.}%
\item {} 
\sphinxAtStartPar
选择目标测段

\item {} 
\sphinxAtStartPar
右键选择 \sphinxcode{\sphinxupquote{加载测点 → 从EDI文件加载}}

\item {} 
\sphinxAtStartPar
选择EDI文件

\end{enumerate}


\subsubsection{4.7.2 数据编辑}
\label{\detokenize{chapters/chapter4:id27}}
\sphinxAtStartPar
导入EDI后可进行：
\begin{itemize}
\item {} 
\sphinxAtStartPar
修改测点元数据

\item {} 
\sphinxAtStartPar
删除异常频点

\item {} 
\sphinxAtStartPar
数据修复

\end{itemize}


\bigskip\hrule\bigskip



\subsection{4.8 合成数据 (Synthetic)}
\label{\detokenize{chapters/chapter4:synthetic}}

\subsubsection{4.8.1 数据格式}
\label{\detokenize{chapters/chapter4:id28}}
\sphinxAtStartPar
合成数据用于正演模拟测试，支持以下格式：


\begin{savenotes}\sphinxattablestart
\sphinxthistablewithglobalstyle
\centering
\begin{tabulary}{\linewidth}[t]{TT}
\sphinxtoprule
\begin{varwidth}[t]{\sphinxcolwidth{1}{2}}
\sphinxstyletheadfamily \sphinxAtStartPar
扩展名
\sphinxbeforeendvarwidth
\end{varwidth}%
&\begin{varwidth}[t]{\sphinxcolwidth{1}{2}}
\sphinxstyletheadfamily \sphinxAtStartPar
说明
\sphinxbeforeendvarwidth
\end{varwidth}%
\\
\sphinxmidrule
\sphinxtableatstartofbodyhook\begin{varwidth}[t]{\sphinxcolwidth{1}{2}}
\sphinxAtStartPar
.timeseries
\sphinxbeforeendvarwidth
\end{varwidth}%
&\begin{varwidth}[t]{\sphinxcolwidth{1}{2}}
\sphinxAtStartPar
合成时间序列数据
\sphinxbeforeendvarwidth
\end{varwidth}%
\\
\sphinxhline\begin{varwidth}[t]{\sphinxcolwidth{1}{2}}
\sphinxAtStartPar
.timeseriesb
\sphinxbeforeendvarwidth
\end{varwidth}%
&\begin{varwidth}[t]{\sphinxcolwidth{1}{2}}
\sphinxAtStartPar
二进制格式合成数据
\sphinxbeforeendvarwidth
\end{varwidth}%
\\
\sphinxbottomrule
\end{tabulary}
\sphinxtableafterendhook\par
\sphinxattableend\end{savenotes}


\subsubsection{4.8.2 数据导入}
\label{\detokenize{chapters/chapter4:id29}}\begin{enumerate}
\sphinxsetlistlabels{\arabic}{enumi}{enumii}{}{.}%
\item {} 
\sphinxAtStartPar
选择目标测段

\item {} 
\sphinxAtStartPar
右键选择 \sphinxcode{\sphinxupquote{加载测点 → 加载合成数据}}

\item {} 
\sphinxAtStartPar
选择 .timeseries 文件

\end{enumerate}


\bigskip\hrule\bigskip



\subsection{4.9 AGEXXL格式}
\label{\detokenize{chapters/chapter4:agexxl}}

\subsubsection{4.9.1 数据导入}
\label{\detokenize{chapters/chapter4:id30}}\begin{enumerate}
\sphinxsetlistlabels{\arabic}{enumi}{enumii}{}{.}%
\item {} 
\sphinxAtStartPar
选择目标测段

\item {} 
\sphinxAtStartPar
右键选择 \sphinxcode{\sphinxupquote{加载测点 → 加载AGEXXL测点}}

\item {} 
\sphinxAtStartPar
选择 .dat 数据文件

\end{enumerate}


\bigskip\hrule\bigskip



\subsection{4.10 数据格式汇总}
\label{\detokenize{chapters/chapter4:id31}}

\begin{savenotes}\sphinxattablestart
\sphinxthistablewithglobalstyle
\centering
\begin{tabulary}{\linewidth}[t]{TTTT}
\sphinxtoprule
\begin{varwidth}[t]{\sphinxcolwidth{1}{4}}
\sphinxstyletheadfamily \sphinxAtStartPar
仪器/格式
\sphinxbeforeendvarwidth
\end{varwidth}%
&\begin{varwidth}[t]{\sphinxcolwidth{1}{4}}
\sphinxstyletheadfamily \sphinxAtStartPar
数据文件
\sphinxbeforeendvarwidth
\end{varwidth}%
&\begin{varwidth}[t]{\sphinxcolwidth{1}{4}}
\sphinxstyletheadfamily \sphinxAtStartPar
标定文件
\sphinxbeforeendvarwidth
\end{varwidth}%
&\begin{varwidth}[t]{\sphinxcolwidth{1}{4}}
\sphinxstyletheadfamily \sphinxAtStartPar
自动识别扩展名
\sphinxbeforeendvarwidth
\end{varwidth}%
\\
\sphinxmidrule
\sphinxtableatstartofbodyhook\begin{varwidth}[t]{\sphinxcolwidth{1}{4}}
\sphinxAtStartPar
Phoenix MTU\sphinxhyphen{}5/5A
\sphinxbeforeendvarwidth
\end{varwidth}%
&\begin{varwidth}[t]{\sphinxcolwidth{1}{4}}
\sphinxAtStartPar
.ts, .mt
\sphinxbeforeendvarwidth
\end{varwidth}%
&\begin{varwidth}[t]{\sphinxcolwidth{1}{4}}
\sphinxAtStartPar
.clb, .clc
\sphinxbeforeendvarwidth
\end{varwidth}%
&\begin{varwidth}[t]{\sphinxcolwidth{1}{4}}
\sphinxAtStartPar
.tbl
\sphinxbeforeendvarwidth
\end{varwidth}%
\\
\sphinxhline\begin{varwidth}[t]{\sphinxcolwidth{1}{4}}
\sphinxAtStartPar
Phoenix MTU\sphinxhyphen{}5C/8A
\sphinxbeforeendvarwidth
\end{varwidth}%
&\begin{varwidth}[t]{\sphinxcolwidth{1}{4}}
\sphinxAtStartPar
.ts
\sphinxbeforeendvarwidth
\end{varwidth}%
&\begin{varwidth}[t]{\sphinxcolwidth{1}{4}}
\sphinxAtStartPar
.json
\sphinxbeforeendvarwidth
\end{varwidth}%
&\begin{varwidth}[t]{\sphinxcolwidth{1}{4}}
\sphinxAtStartPar
recmeta.json
\sphinxbeforeendvarwidth
\end{varwidth}%
\\
\sphinxhline\begin{varwidth}[t]{\sphinxcolwidth{1}{4}}
\sphinxAtStartPar
Metronix
\sphinxbeforeendvarwidth
\end{varwidth}%
&\begin{varwidth}[t]{\sphinxcolwidth{1}{4}}
\sphinxAtStartPar
.ats, .atm
\sphinxbeforeendvarwidth
\end{varwidth}%
&\begin{varwidth}[t]{\sphinxcolwidth{1}{4}}
\sphinxAtStartPar
.cal
\sphinxbeforeendvarwidth
\end{varwidth}%
&\begin{varwidth}[t]{\sphinxcolwidth{1}{4}}
\sphinxAtStartPar
.mxsite
\sphinxbeforeendvarwidth
\end{varwidth}%
\\
\sphinxhline\begin{varwidth}[t]{\sphinxcolwidth{1}{4}}
\sphinxAtStartPar
LEMI
\sphinxbeforeendvarwidth
\end{varwidth}%
&\begin{varwidth}[t]{\sphinxcolwidth{1}{4}}
\sphinxAtStartPar
二进制
\sphinxbeforeendvarwidth
\end{varwidth}%
&\begin{varwidth}[t]{\sphinxcolwidth{1}{4}}
\sphinxAtStartPar
.txt
\sphinxbeforeendvarwidth
\end{varwidth}%
&\begin{varwidth}[t]{\sphinxcolwidth{1}{4}}
\sphinxAtStartPar
.lemi, .lemijson
\sphinxbeforeendvarwidth
\end{varwidth}%
\\
\sphinxhline\begin{varwidth}[t]{\sphinxcolwidth{1}{4}}
\sphinxAtStartPar
SBF
\sphinxbeforeendvarwidth
\end{varwidth}%
&\begin{varwidth}[t]{\sphinxcolwidth{1}{4}}
\sphinxAtStartPar
.sbf
\sphinxbeforeendvarwidth
\end{varwidth}%
&\begin{varwidth}[t]{\sphinxcolwidth{1}{4}}
\sphinxAtStartPar
内置
\sphinxbeforeendvarwidth
\end{varwidth}%
&\begin{varwidth}[t]{\sphinxcolwidth{1}{4}}
\sphinxAtStartPar
.sbf
\sphinxbeforeendvarwidth
\end{varwidth}%
\\
\sphinxhline\begin{varwidth}[t]{\sphinxcolwidth{1}{4}}
\sphinxAtStartPar
RMT
\sphinxbeforeendvarwidth
\end{varwidth}%
&\begin{varwidth}[t]{\sphinxcolwidth{1}{4}}
\sphinxAtStartPar
.tr1/.tr2/.tr3
\sphinxbeforeendvarwidth
\end{varwidth}%
&\begin{varwidth}[t]{\sphinxcolwidth{1}{4}}
\sphinxAtStartPar
内置
\sphinxbeforeendvarwidth
\end{varwidth}%
&\begin{varwidth}[t]{\sphinxcolwidth{1}{4}}
\sphinxAtStartPar
.tr1, .tr2, .tr3
\sphinxbeforeendvarwidth
\end{varwidth}%
\\
\sphinxhline\begin{varwidth}[t]{\sphinxcolwidth{1}{4}}
\sphinxAtStartPar
CASRMT
\sphinxbeforeendvarwidth
\end{varwidth}%
&\begin{varwidth}[t]{\sphinxcolwidth{1}{4}}
\sphinxAtStartPar
.json
\sphinxbeforeendvarwidth
\end{varwidth}%
&\begin{varwidth}[t]{\sphinxcolwidth{1}{4}}
\sphinxAtStartPar
.csv
\sphinxbeforeendvarwidth
\end{varwidth}%
&\begin{varwidth}[t]{\sphinxcolwidth{1}{4}}
\sphinxAtStartPar
.rmtjson, .json
\sphinxbeforeendvarwidth
\end{varwidth}%
\\
\sphinxhline\begin{varwidth}[t]{\sphinxcolwidth{1}{4}}
\sphinxAtStartPar
ATTS
\sphinxbeforeendvarwidth
\end{varwidth}%
&\begin{varwidth}[t]{\sphinxcolwidth{1}{4}}
\sphinxAtStartPar
二进制
\sphinxbeforeendvarwidth
\end{varwidth}%
&\begin{varwidth}[t]{\sphinxcolwidth{1}{4}}
\sphinxAtStartPar
.csv
\sphinxbeforeendvarwidth
\end{varwidth}%
&\begin{varwidth}[t]{\sphinxcolwidth{1}{4}}
\sphinxAtStartPar
.atts, .atinfo
\sphinxbeforeendvarwidth
\end{varwidth}%
\\
\sphinxhline\begin{varwidth}[t]{\sphinxcolwidth{1}{4}}
\sphinxAtStartPar
AGEXXL
\sphinxbeforeendvarwidth
\end{varwidth}%
&\begin{varwidth}[t]{\sphinxcolwidth{1}{4}}
\sphinxAtStartPar
.dat
\sphinxbeforeendvarwidth
\end{varwidth}%
&\begin{varwidth}[t]{\sphinxcolwidth{1}{4}}
\sphinxAtStartPar
.csv
\sphinxbeforeendvarwidth
\end{varwidth}%
&\begin{varwidth}[t]{\sphinxcolwidth{1}{4}}
\sphinxAtStartPar
.dat
\sphinxbeforeendvarwidth
\end{varwidth}%
\\
\sphinxhline\begin{varwidth}[t]{\sphinxcolwidth{1}{4}}
\sphinxAtStartPar
Synthetic
\sphinxbeforeendvarwidth
\end{varwidth}%
&\begin{varwidth}[t]{\sphinxcolwidth{1}{4}}
\sphinxAtStartPar
.timeseries
\sphinxbeforeendvarwidth
\end{varwidth}%
&\begin{varwidth}[t]{\sphinxcolwidth{1}{4}}
\sphinxAtStartPar
\sphinxhyphen{}
\sphinxbeforeendvarwidth
\end{varwidth}%
&\begin{varwidth}[t]{\sphinxcolwidth{1}{4}}
\sphinxAtStartPar
.timeseries, .timeseriesb
\sphinxbeforeendvarwidth
\end{varwidth}%
\\
\sphinxhline\begin{varwidth}[t]{\sphinxcolwidth{1}{4}}
\sphinxAtStartPar
EDI
\sphinxbeforeendvarwidth
\end{varwidth}%
&\begin{varwidth}[t]{\sphinxcolwidth{1}{4}}
\sphinxAtStartPar
.edi
\sphinxbeforeendvarwidth
\end{varwidth}%
&\begin{varwidth}[t]{\sphinxcolwidth{1}{4}}
\sphinxAtStartPar
\sphinxhyphen{}
\sphinxbeforeendvarwidth
\end{varwidth}%
&\begin{varwidth}[t]{\sphinxcolwidth{1}{4}}
\sphinxAtStartPar
.edi
\sphinxbeforeendvarwidth
\end{varwidth}%
\\
\sphinxbottomrule
\end{tabulary}
\sphinxtableafterendhook\par
\sphinxattableend\end{savenotes}


\bigskip\hrule\bigskip



\subsection{4.11 标定文件查找机制}
\label{\detokenize{chapters/chapter4:id32}}
\sphinxAtStartPar
MTDP采用多级路径搜索标定文件，确保标定文件能够被自动找到：


\subsubsection{4.11.1 搜索顺序}
\label{\detokenize{chapters/chapter4:id33}}

\begin{savenotes}\sphinxattablestart
\sphinxthistablewithglobalstyle
\centering
\begin{tabulary}{\linewidth}[t]{TTT}
\sphinxtoprule
\begin{varwidth}[t]{\sphinxcolwidth{1}{3}}
\sphinxstyletheadfamily \sphinxAtStartPar
优先级
\sphinxbeforeendvarwidth
\end{varwidth}%
&\begin{varwidth}[t]{\sphinxcolwidth{1}{3}}
\sphinxstyletheadfamily \sphinxAtStartPar
搜索位置
\sphinxbeforeendvarwidth
\end{varwidth}%
&\begin{varwidth}[t]{\sphinxcolwidth{1}{3}}
\sphinxstyletheadfamily \sphinxAtStartPar
说明
\sphinxbeforeendvarwidth
\end{varwidth}%
\\
\sphinxmidrule
\sphinxtableatstartofbodyhook\begin{varwidth}[t]{\sphinxcolwidth{1}{3}}
\sphinxAtStartPar
1
\sphinxbeforeendvarwidth
\end{varwidth}%
&\begin{varwidth}[t]{\sphinxcolwidth{1}{3}}
\sphinxAtStartPar
数据文件所在目录
\sphinxbeforeendvarwidth
\end{varwidth}%
&\begin{varwidth}[t]{\sphinxcolwidth{1}{3}}
\sphinxAtStartPar
TBL/数据文件同目录
\sphinxbeforeendvarwidth
\end{varwidth}%
\\
\sphinxhline\begin{varwidth}[t]{\sphinxcolwidth{1}{3}}
\sphinxAtStartPar
2
\sphinxbeforeendvarwidth
\end{varwidth}%
&\begin{varwidth}[t]{\sphinxcolwidth{1}{3}}
\sphinxAtStartPar
用户指定路径
\sphinxbeforeendvarwidth
\end{varwidth}%
&\begin{varwidth}[t]{\sphinxcolwidth{1}{3}}
\sphinxAtStartPar
设置中配置的额外路径
\sphinxbeforeendvarwidth
\end{varwidth}%
\\
\sphinxhline\begin{varwidth}[t]{\sphinxcolwidth{1}{3}}
\sphinxAtStartPar
3
\sphinxbeforeendvarwidth
\end{varwidth}%
&\begin{varwidth}[t]{\sphinxcolwidth{1}{3}}
\sphinxAtStartPar
全局配置路径
\sphinxbeforeendvarwidth
\end{varwidth}%
&\begin{varwidth}[t]{\sphinxcolwidth{1}{3}}
\sphinxAtStartPar
默认标定目录
\sphinxbeforeendvarwidth
\end{varwidth}%
\\
\sphinxhline\begin{varwidth}[t]{\sphinxcolwidth{1}{3}}
\sphinxAtStartPar
4
\sphinxbeforeendvarwidth
\end{varwidth}%
&\begin{varwidth}[t]{\sphinxcolwidth{1}{3}}
\sphinxAtStartPar
默认路径
\sphinxbeforeendvarwidth
\end{varwidth}%
&\begin{varwidth}[t]{\sphinxcolwidth{1}{3}}
\sphinxAtStartPar
CalibrationFiles/Text 或 CalibrationFiles/CSV
\sphinxbeforeendvarwidth
\end{varwidth}%
\\
\sphinxbottomrule
\end{tabulary}
\sphinxtableafterendhook\par
\sphinxattableend\end{savenotes}


\subsubsection{4.11.2 标定文件格式}
\label{\detokenize{chapters/chapter4:id34}}

\begin{savenotes}\sphinxattablestart
\sphinxthistablewithglobalstyle
\centering
\begin{tabulary}{\linewidth}[t]{TTT}
\sphinxtoprule
\begin{varwidth}[t]{\sphinxcolwidth{1}{3}}
\sphinxstyletheadfamily \sphinxAtStartPar
格式
\sphinxbeforeendvarwidth
\end{varwidth}%
&\begin{varwidth}[t]{\sphinxcolwidth{1}{3}}
\sphinxstyletheadfamily \sphinxAtStartPar
扩展名
\sphinxbeforeendvarwidth
\end{varwidth}%
&\begin{varwidth}[t]{\sphinxcolwidth{1}{3}}
\sphinxstyletheadfamily \sphinxAtStartPar
内容
\sphinxbeforeendvarwidth
\end{varwidth}%
\\
\sphinxmidrule
\sphinxtableatstartofbodyhook\begin{varwidth}[t]{\sphinxcolwidth{1}{3}}
\sphinxAtStartPar
幅值相位格式
\sphinxbeforeendvarwidth
\end{varwidth}%
&\begin{varwidth}[t]{\sphinxcolwidth{1}{3}}
\sphinxAtStartPar
.txt
\sphinxbeforeendvarwidth
\end{varwidth}%
&\begin{varwidth}[t]{\sphinxcolwidth{1}{3}}
\sphinxAtStartPar
频率、幅度、相位
\sphinxbeforeendvarwidth
\end{varwidth}%
\\
\sphinxhline\begin{varwidth}[t]{\sphinxcolwidth{1}{3}}
\sphinxAtStartPar
CSV格式
\sphinxbeforeendvarwidth
\end{varwidth}%
&\begin{varwidth}[t]{\sphinxcolwidth{1}{3}}
\sphinxAtStartPar
\_CLB.csv, \_CLC.csv
\sphinxbeforeendvarwidth
\end{varwidth}%
&\begin{varwidth}[t]{\sphinxcolwidth{1}{3}}
\sphinxAtStartPar
频率、复数响应
\sphinxbeforeendvarwidth
\end{varwidth}%
\\
\sphinxhline\begin{varwidth}[t]{\sphinxcolwidth{1}{3}}
\sphinxAtStartPar
实部虚部格式
\sphinxbeforeendvarwidth
\end{varwidth}%
&\begin{varwidth}[t]{\sphinxcolwidth{1}{3}}
\sphinxAtStartPar
.txt
\sphinxbeforeendvarwidth
\end{varwidth}%
&\begin{varwidth}[t]{\sphinxcolwidth{1}{3}}
\sphinxAtStartPar
频率、实部、虚部
\sphinxbeforeendvarwidth
\end{varwidth}%
\\
\sphinxbottomrule
\end{tabulary}
\sphinxtableafterendhook\par
\sphinxattableend\end{savenotes}


\subsubsection{4.11.3 配置标定路径}
\label{\detokenize{chapters/chapter4:id35}}\begin{enumerate}
\sphinxsetlistlabels{\arabic}{enumi}{enumii}{}{.}%
\item {} 
\sphinxAtStartPar
选择 \sphinxcode{\sphinxupquote{设置 → 管理标定路径}}

\item {} 
\sphinxAtStartPar
添加或删除标定文件搜索路径

\item {} 
\sphinxAtStartPar
系统自动在配置的路径中查找标定文件

\end{enumerate}


\bigskip\hrule\bigskip



\subsection{4.12 多文件加载方式}
\label{\detokenize{chapters/chapter4:id36}}

\subsubsection{4.12.1 单文件加载}
\label{\detokenize{chapters/chapter4:id37}}
\sphinxAtStartPar
标准方式，从单个数据文件创建测点。


\subsubsection{4.12.2 双源文件加载（X/Y源）}
\label{\detokenize{chapters/chapter4:x-y}}
\sphinxAtStartPar
用于可控源MT数据，分别加载X方向和Y方向源数据：
\begin{enumerate}
\sphinxsetlistlabels{\arabic}{enumi}{enumii}{}{.}%
\item {} 
\sphinxAtStartPar
右键测段 → \sphinxcode{\sphinxupquote{加载测点 → 双源文件}}

\item {} 
\sphinxAtStartPar
选择X方向数据文件

\item {} 
\sphinxAtStartPar
选择Y方向数据文件

\end{enumerate}


\subsubsection{4.12.3 多文件合并加载}
\label{\detokenize{chapters/chapter4:id38}}
\sphinxAtStartPar
将多个数据文件合并为一个测点：
\begin{enumerate}
\sphinxsetlistlabels{\arabic}{enumi}{enumii}{}{.}%
\item {} 
\sphinxAtStartPar
右键测段 → \sphinxcode{\sphinxupquote{加载测点 → 多文件合并}}

\item {} 
\sphinxAtStartPar
选择多个数据文件

\item {} 
\sphinxAtStartPar
系统自动合并

\end{enumerate}


\bigskip\hrule\bigskip



\subsection{4.13 仪器类型自动识别}
\label{\detokenize{chapters/chapter4:id39}}
\sphinxAtStartPar
MTDP根据文件扩展名自动识别仪器类型：


\begin{savenotes}\sphinxattablestart
\sphinxthistablewithglobalstyle
\centering
\begin{tabulary}{\linewidth}[t]{TT}
\sphinxtoprule
\begin{varwidth}[t]{\sphinxcolwidth{1}{2}}
\sphinxstyletheadfamily \sphinxAtStartPar
文件特征
\sphinxbeforeendvarwidth
\end{varwidth}%
&\begin{varwidth}[t]{\sphinxcolwidth{1}{2}}
\sphinxstyletheadfamily \sphinxAtStartPar
识别为
\sphinxbeforeendvarwidth
\end{varwidth}%
\\
\sphinxmidrule
\sphinxtableatstartofbodyhook\begin{varwidth}[t]{\sphinxcolwidth{1}{2}}
\sphinxAtStartPar
.tbl
\sphinxbeforeendvarwidth
\end{varwidth}%
&\begin{varwidth}[t]{\sphinxcolwidth{1}{2}}
\sphinxAtStartPar
Phoenix
\sphinxbeforeendvarwidth
\end{varwidth}%
\\
\sphinxhline\begin{varwidth}[t]{\sphinxcolwidth{1}{2}}
\sphinxAtStartPar
recmeta.json
\sphinxbeforeendvarwidth
\end{varwidth}%
&\begin{varwidth}[t]{\sphinxcolwidth{1}{2}}
\sphinxAtStartPar
MTU
\sphinxbeforeendvarwidth
\end{varwidth}%
\\
\sphinxhline\begin{varwidth}[t]{\sphinxcolwidth{1}{2}}
\sphinxAtStartPar
.mxsite
\sphinxbeforeendvarwidth
\end{varwidth}%
&\begin{varwidth}[t]{\sphinxcolwidth{1}{2}}
\sphinxAtStartPar
Metronix
\sphinxbeforeendvarwidth
\end{varwidth}%
\\
\sphinxhline\begin{varwidth}[t]{\sphinxcolwidth{1}{2}}
\sphinxAtStartPar
.lemi, .lemijson
\sphinxbeforeendvarwidth
\end{varwidth}%
&\begin{varwidth}[t]{\sphinxcolwidth{1}{2}}
\sphinxAtStartPar
LEMI
\sphinxbeforeendvarwidth
\end{varwidth}%
\\
\sphinxhline\begin{varwidth}[t]{\sphinxcolwidth{1}{2}}
\sphinxAtStartPar
.sbf
\sphinxbeforeendvarwidth
\end{varwidth}%
&\begin{varwidth}[t]{\sphinxcolwidth{1}{2}}
\sphinxAtStartPar
SBF
\sphinxbeforeendvarwidth
\end{varwidth}%
\\
\sphinxhline\begin{varwidth}[t]{\sphinxcolwidth{1}{2}}
\sphinxAtStartPar
.tr1, .tr2, .tr3
\sphinxbeforeendvarwidth
\end{varwidth}%
&\begin{varwidth}[t]{\sphinxcolwidth{1}{2}}
\sphinxAtStartPar
RMT
\sphinxbeforeendvarwidth
\end{varwidth}%
\\
\sphinxhline\begin{varwidth}[t]{\sphinxcolwidth{1}{2}}
\sphinxAtStartPar
.rmtjson, .json (非recmeta)
\sphinxbeforeendvarwidth
\end{varwidth}%
&\begin{varwidth}[t]{\sphinxcolwidth{1}{2}}
\sphinxAtStartPar
CASRMT
\sphinxbeforeendvarwidth
\end{varwidth}%
\\
\sphinxhline\begin{varwidth}[t]{\sphinxcolwidth{1}{2}}
\sphinxAtStartPar
.atts, .atinfo
\sphinxbeforeendvarwidth
\end{varwidth}%
&\begin{varwidth}[t]{\sphinxcolwidth{1}{2}}
\sphinxAtStartPar
ATTS
\sphinxbeforeendvarwidth
\end{varwidth}%
\\
\sphinxhline\begin{varwidth}[t]{\sphinxcolwidth{1}{2}}
\sphinxAtStartPar
.dat
\sphinxbeforeendvarwidth
\end{varwidth}%
&\begin{varwidth}[t]{\sphinxcolwidth{1}{2}}
\sphinxAtStartPar
AGEXXL
\sphinxbeforeendvarwidth
\end{varwidth}%
\\
\sphinxhline\begin{varwidth}[t]{\sphinxcolwidth{1}{2}}
\sphinxAtStartPar
.timeseries, .timeseriesb
\sphinxbeforeendvarwidth
\end{varwidth}%
&\begin{varwidth}[t]{\sphinxcolwidth{1}{2}}
\sphinxAtStartPar
Synthetic
\sphinxbeforeendvarwidth
\end{varwidth}%
\\
\sphinxhline\begin{varwidth}[t]{\sphinxcolwidth{1}{2}}
\sphinxAtStartPar
.edi
\sphinxbeforeendvarwidth
\end{varwidth}%
&\begin{varwidth}[t]{\sphinxcolwidth{1}{2}}
\sphinxAtStartPar
EDI
\sphinxbeforeendvarwidth
\end{varwidth}%
\\
\sphinxbottomrule
\end{tabulary}
\sphinxtableafterendhook\par
\sphinxattableend\end{savenotes}


\bigskip\hrule\bigskip



\subsection{4.14 数据导入建议}
\label{\detokenize{chapters/chapter4:id40}}

\subsubsection{4.14.1 导入前检查}
\label{\detokenize{chapters/chapter4:id41}}\begin{itemize}
\item {} 
\sphinxAtStartPar
确认数据文件完整

\item {} 
\sphinxAtStartPar
检查标定文件匹配

\item {} 
\sphinxAtStartPar
验证通道配置

\end{itemize}


\subsubsection{4.14.2 数据组织}
\label{\detokenize{chapters/chapter4:id42}}\begin{itemize}
\item {} 
\sphinxAtStartPar
同一仪器数据放在同一目录

\item {} 
\sphinxAtStartPar
保持文件命名规范

\item {} 
\sphinxAtStartPar
备份原始数据

\end{itemize}


\bigskip\hrule\bigskip



\subsection{4.15 批量添加频点}
\label{\detokenize{chapters/chapter4:id43}}
\sphinxAtStartPar
在MTDP中处理数据时，可能需要批量添加或修改频点。MTDP提供了强大的频点管理功能。


\subsubsection{4.15.1 批量添加频点表单 (BatchAddFreqForm)}
\label{\detokenize{chapters/chapter4:batchaddfreqform}}
\sphinxAtStartPar
\sphinxstylestrong{位置}: \sphinxcode{\sphinxupquote{处理 → 批量添加频点}}

\sphinxAtStartPar
\sphinxstylestrong{功能}: 一次性生成多个频点，支持三种分布方式


\subsubsection{4.15.2 频率生成方式}
\label{\detokenize{chapters/chapter4:id44}}
\sphinxAtStartPar
MTDP支持三种频率生成方式：


\begin{savenotes}\sphinxattablestart
\sphinxthistablewithglobalstyle
\centering
\begin{tabulary}{\linewidth}[t]{TTTT}
\sphinxtoprule
\begin{varwidth}[t]{\sphinxcolwidth{1}{4}}
\sphinxstyletheadfamily \sphinxAtStartPar
方式
\sphinxbeforeendvarwidth
\end{varwidth}%
&\begin{varwidth}[t]{\sphinxcolwidth{1}{4}}
\sphinxstyletheadfamily \sphinxAtStartPar
索引
\sphinxbeforeendvarwidth
\end{varwidth}%
&\begin{varwidth}[t]{\sphinxcolwidth{1}{4}}
\sphinxstyletheadfamily \sphinxAtStartPar
说明
\sphinxbeforeendvarwidth
\end{varwidth}%
&\begin{varwidth}[t]{\sphinxcolwidth{1}{4}}
\sphinxstyletheadfamily \sphinxAtStartPar
适用场景
\sphinxbeforeendvarwidth
\end{varwidth}%
\\
\sphinxmidrule
\sphinxtableatstartofbodyhook\begin{varwidth}[t]{\sphinxcolwidth{1}{4}}
\sphinxAtStartPar
\sphinxstylestrong{等比分布（对数）}
\sphinxbeforeendvarwidth
\end{varwidth}%
&\begin{varwidth}[t]{\sphinxcolwidth{1}{4}}
\sphinxAtStartPar
0
\sphinxbeforeendvarwidth
\end{varwidth}%
&\begin{varwidth}[t]{\sphinxcolwidth{1}{4}}
\sphinxAtStartPar
按对数等间隔生成
\sphinxbeforeendvarwidth
\end{varwidth}%
&\begin{varwidth}[t]{\sphinxcolwidth{1}{4}}
\sphinxAtStartPar
宽频带探测，覆盖全频段
\sphinxbeforeendvarwidth
\end{varwidth}%
\\
\sphinxhline\begin{varwidth}[t]{\sphinxcolwidth{1}{4}}
\sphinxAtStartPar
\sphinxstylestrong{等差分布（线性）}
\sphinxbeforeendvarwidth
\end{varwidth}%
&\begin{varwidth}[t]{\sphinxcolwidth{1}{4}}
\sphinxAtStartPar
1
\sphinxbeforeendvarwidth
\end{varwidth}%
&\begin{varwidth}[t]{\sphinxcolwidth{1}{4}}
\sphinxAtStartPar
按线性等间隔生成
\sphinxbeforeendvarwidth
\end{varwidth}%
&\begin{varwidth}[t]{\sphinxcolwidth{1}{4}}
\sphinxAtStartPar
固定频率间隔研究
\sphinxbeforeendvarwidth
\end{varwidth}%
\\
\sphinxhline\begin{varwidth}[t]{\sphinxcolwidth{1}{4}}
\sphinxAtStartPar
\sphinxstylestrong{自定义间距}
\sphinxbeforeendvarwidth
\end{varwidth}%
&\begin{varwidth}[t]{\sphinxcolwidth{1}{4}}
\sphinxAtStartPar
2
\sphinxbeforeendvarwidth
\end{varwidth}%
&\begin{varwidth}[t]{\sphinxcolwidth{1}{4}}
\sphinxAtStartPar
按自定义倍数递增
\sphinxbeforeendvarwidth
\end{varwidth}%
&\begin{varwidth}[t]{\sphinxcolwidth{1}{4}}
\sphinxAtStartPar
非均匀采样研究
\sphinxbeforeendvarwidth
\end{varwidth}%
\\
\sphinxbottomrule
\end{tabulary}
\sphinxtableafterendhook\par
\sphinxattableend\end{savenotes}


\subsubsection{4.15.3 参数说明}
\label{\detokenize{chapters/chapter4:id45}}

\begin{savenotes}\sphinxattablestart
\sphinxthistablewithglobalstyle
\centering
\begin{tabulary}{\linewidth}[t]{TTTT}
\sphinxtoprule
\begin{varwidth}[t]{\sphinxcolwidth{1}{4}}
\sphinxstyletheadfamily \sphinxAtStartPar
参数
\sphinxbeforeendvarwidth
\end{varwidth}%
&\begin{varwidth}[t]{\sphinxcolwidth{1}{4}}
\sphinxstyletheadfamily \sphinxAtStartPar
控件名称
\sphinxbeforeendvarwidth
\end{varwidth}%
&\begin{varwidth}[t]{\sphinxcolwidth{1}{4}}
\sphinxstyletheadfamily \sphinxAtStartPar
说明
\sphinxbeforeendvarwidth
\end{varwidth}%
&\begin{varwidth}[t]{\sphinxcolwidth{1}{4}}
\sphinxstyletheadfamily \sphinxAtStartPar
单位
\sphinxbeforeendvarwidth
\end{varwidth}%
\\
\sphinxmidrule
\sphinxtableatstartofbodyhook\begin{varwidth}[t]{\sphinxcolwidth{1}{4}}
\sphinxAtStartPar
起始频率
\sphinxbeforeendvarwidth
\end{varwidth}%
&\begin{varwidth}[t]{\sphinxcolwidth{1}{4}}
\sphinxAtStartPar
NumberBoxStartFreq
\sphinxbeforeendvarwidth
\end{varwidth}%
&\begin{varwidth}[t]{\sphinxcolwidth{1}{4}}
\sphinxAtStartPar
频带低端频率
\sphinxbeforeendvarwidth
\end{varwidth}%
&\begin{varwidth}[t]{\sphinxcolwidth{1}{4}}
\sphinxAtStartPar
Hz
\sphinxbeforeendvarwidth
\end{varwidth}%
\\
\sphinxhline\begin{varwidth}[t]{\sphinxcolwidth{1}{4}}
\sphinxAtStartPar
终止频率
\sphinxbeforeendvarwidth
\end{varwidth}%
&\begin{varwidth}[t]{\sphinxcolwidth{1}{4}}
\sphinxAtStartPar
NumberBoxEndFreq
\sphinxbeforeendvarwidth
\end{varwidth}%
&\begin{varwidth}[t]{\sphinxcolwidth{1}{4}}
\sphinxAtStartPar
频带高端频率
\sphinxbeforeendvarwidth
\end{varwidth}%
&\begin{varwidth}[t]{\sphinxcolwidth{1}{4}}
\sphinxAtStartPar
Hz
\sphinxbeforeendvarwidth
\end{varwidth}%
\\
\sphinxhline\begin{varwidth}[t]{\sphinxcolwidth{1}{4}}
\sphinxAtStartPar
频点数量
\sphinxbeforeendvarwidth
\end{varwidth}%
&\begin{varwidth}[t]{\sphinxcolwidth{1}{4}}
\sphinxAtStartPar
NumberBoxFreqCount
\sphinxbeforeendvarwidth
\end{varwidth}%
&\begin{varwidth}[t]{\sphinxcolwidth{1}{4}}
\sphinxAtStartPar
要生成的频点总数
\sphinxbeforeendvarwidth
\end{varwidth}%
&\begin{varwidth}[t]{\sphinxcolwidth{1}{4}}
\sphinxAtStartPar
\sphinxhyphen{}
\sphinxbeforeendvarwidth
\end{varwidth}%
\\
\sphinxhline\begin{varwidth}[t]{\sphinxcolwidth{1}{4}}
\sphinxAtStartPar
频率间距
\sphinxbeforeendvarwidth
\end{varwidth}%
&\begin{varwidth}[t]{\sphinxcolwidth{1}{4}}
\sphinxAtStartPar
NumberBoxSpacing
\sphinxbeforeendvarwidth
\end{varwidth}%
&\begin{varwidth}[t]{\sphinxcolwidth{1}{4}}
\sphinxAtStartPar
线性间隔值或对数倍数
\sphinxbeforeendvarwidth
\end{varwidth}%
&\begin{varwidth}[t]{\sphinxcolwidth{1}{4}}
\sphinxAtStartPar
Hz或比例
\sphinxbeforeendvarwidth
\end{varwidth}%
\\
\sphinxbottomrule
\end{tabulary}
\sphinxtableafterendhook\par
\sphinxattableend\end{savenotes}


\subsubsection{4.15.4 生成算法}
\label{\detokenize{chapters/chapter4:id46}}

\paragraph{等比分布（对数间隔）}
\label{\detokenize{chapters/chapter4:id47}}
\sphinxAtStartPar
频点按对数等间隔分布，公式为：

\begin{sphinxVerbatim}[commandchars=\\\{\}]
Freq[i] = StartFreq × (EndFreq / StartFreq)\PYGZca{}(i / (Count \PYGZhy{} 1))
\end{sphinxVerbatim}

\sphinxAtStartPar
\sphinxstylestrong{特点}:
\begin{itemize}
\item {} 
\sphinxAtStartPar
频点在对数尺度上均匀分布

\item {} 
\sphinxAtStartPar
低频端频点密集，高频端稀疏

\item {} 
\sphinxAtStartPar
适合宽频带MT探测

\end{itemize}

\sphinxAtStartPar
\sphinxstylestrong{示例}:

\begin{sphinxVerbatim}[commandchars=\\\{\}]
\PYG{n}{起始频率}\PYG{p}{:} \PYG{l+m+mf}{0.001} \PYG{n}{Hz}
\PYG{n}{终止频率}\PYG{p}{:} \PYG{l+m+mi}{1000} \PYG{n}{Hz}
\PYG{n}{频点数量}\PYG{p}{:} \PYG{l+m+mi}{32}

\PYG{n}{结果}\PYG{p}{:} \PYG{l+m+mf}{0.001}\PYG{p}{,} \PYG{l+m+mf}{0.002}\PYG{p}{,} \PYG{l+m+mf}{0.004}\PYG{p}{,} \PYG{l+m+mf}{0.008}\PYG{p}{,} \PYG{o}{.}\PYG{o}{.}\PYG{o}{.} \PYG{p}{(}\PYG{n}{对数等比}\PYG{p}{)}
\end{sphinxVerbatim}


\paragraph{等差分布（线性间隔）}
\label{\detokenize{chapters/chapter4:id48}}
\sphinxAtStartPar
频点按固定频率间距均匀分布，公式为：

\begin{sphinxVerbatim}[commandchars=\\\{\}]
\PYG{n}{Freq}\PYG{p}{[}\PYG{l+m+mi}{0}\PYG{p}{]} \PYG{o}{=} \PYG{n}{StartFreq}
\PYG{n}{Freq}\PYG{p}{[}\PYG{n}{i}\PYG{o}{+}\PYG{l+m+mi}{1}\PYG{p}{]} \PYG{o}{=} \PYG{n}{Freq}\PYG{p}{[}\PYG{n}{i}\PYG{p}{]} \PYG{o}{+} \PYG{n}{Spacing}
\PYG{n}{直到} \PYG{n}{Freq}\PYG{p}{[}\PYG{n}{i}\PYG{p}{]} \PYG{o}{\PYGZgt{}} \PYG{n}{EndFreq}
\end{sphinxVerbatim}

\sphinxAtStartPar
\sphinxstylestrong{特点}:
\begin{itemize}
\item {} 
\sphinxAtStartPar
频点在线性尺度上均匀分布

\item {} 
\sphinxAtStartPar
每个频段间隔相同

\item {} 
\sphinxAtStartPar
适合固定频率研究

\end{itemize}

\sphinxAtStartPar
\sphinxstylestrong{示例}:

\begin{sphinxVerbatim}[commandchars=\\\{\}]
\PYG{n}{起始频率}\PYG{p}{:} \PYG{l+m+mi}{10} \PYG{n}{Hz}
\PYG{n}{终止频率}\PYG{p}{:} \PYG{l+m+mi}{100} \PYG{n}{Hz}
\PYG{n}{频率间距}\PYG{p}{:} \PYG{l+m+mi}{5} \PYG{n}{Hz}

\PYG{n}{结果}\PYG{p}{:} \PYG{l+m+mi}{10}\PYG{p}{,} \PYG{l+m+mi}{15}\PYG{p}{,} \PYG{l+m+mi}{20}\PYG{p}{,} \PYG{l+m+mi}{25}\PYG{p}{,} \PYG{l+m+mi}{30}\PYG{p}{,} \PYG{o}{.}\PYG{o}{.}\PYG{o}{.} \PYG{l+m+mi}{95}\PYG{p}{,} \PYG{l+m+mi}{100}
\end{sphinxVerbatim}


\paragraph{自定义间距（乘法递增）}
\label{\detokenize{chapters/chapter4:id49}}
\sphinxAtStartPar
频点按倍数因子递增，公式为：

\begin{sphinxVerbatim}[commandchars=\\\{\}]
Freq[0] = StartFreq
Freq[i+1] = Freq[i] × Spacing
直到 Freq[i] \PYGZgt{} EndFreq
\end{sphinxVerbatim}

\sphinxAtStartPar
\sphinxstylestrong{特点}:
\begin{itemize}
\item {} 
\sphinxAtStartPar
频点间隔按倍数增长

\item {} 
\sphinxAtStartPar
频率间隔逐渐增大

\item {} 
\sphinxAtStartPar
适合对数尺度研究

\end{itemize}

\sphinxAtStartPar
\sphinxstylestrong{示例}:

\begin{sphinxVerbatim}[commandchars=\\\{\}]
\PYG{n}{起始频率}\PYG{p}{:} \PYG{l+m+mi}{1} \PYG{n}{Hz}
\PYG{n}{终止频率}\PYG{p}{:} \PYG{l+m+mi}{1000} \PYG{n}{Hz}
\PYG{n}{间距因子}\PYG{p}{:} \PYG{l+m+mf}{1.5}

\PYG{n}{结果}\PYG{p}{:} \PYG{l+m+mi}{1}\PYG{p}{,} \PYG{l+m+mf}{1.5}\PYG{p}{,} \PYG{l+m+mf}{2.25}\PYG{p}{,} \PYG{l+m+mf}{3.375}\PYG{p}{,} \PYG{o}{.}\PYG{o}{.}\PYG{o}{.} \PYG{p}{(}\PYG{n}{倍增}\PYG{p}{)}
\end{sphinxVerbatim}


\subsubsection{4.15.5 频段编辑功能 (EditBandForm)}
\label{\detokenize{chapters/chapter4:editbandform}}
\sphinxAtStartPar
\sphinxstylestrong{位置}: \sphinxcode{\sphinxupquote{处理 → 编辑频段}}

\sphinxAtStartPar
\sphinxstylestrong{功能}: 添加新频段或编辑已有频段


\paragraph{界面参数}
\label{\detokenize{chapters/chapter4:id50}}

\begin{savenotes}\sphinxattablestart
\sphinxthistablewithglobalstyle
\centering
\begin{tabulary}{\linewidth}[t]{TTT}
\sphinxtoprule
\begin{varwidth}[t]{\sphinxcolwidth{1}{3}}
\sphinxstyletheadfamily \sphinxAtStartPar
参数
\sphinxbeforeendvarwidth
\end{varwidth}%
&\begin{varwidth}[t]{\sphinxcolwidth{1}{3}}
\sphinxstyletheadfamily \sphinxAtStartPar
控件
\sphinxbeforeendvarwidth
\end{varwidth}%
&\begin{varwidth}[t]{\sphinxcolwidth{1}{3}}
\sphinxstyletheadfamily \sphinxAtStartPar
说明
\sphinxbeforeendvarwidth
\end{varwidth}%
\\
\sphinxmidrule
\sphinxtableatstartofbodyhook\begin{varwidth}[t]{\sphinxcolwidth{1}{3}}
\sphinxAtStartPar
频段名称
\sphinxbeforeendvarwidth
\end{varwidth}%
&\begin{varwidth}[t]{\sphinxcolwidth{1}{3}}
\sphinxAtStartPar
Edit
\sphinxbeforeendvarwidth
\end{varwidth}%
&\begin{varwidth}[t]{\sphinxcolwidth{1}{3}}
\sphinxAtStartPar
频段标识名称
\sphinxbeforeendvarwidth
\end{varwidth}%
\\
\sphinxhline\begin{varwidth}[t]{\sphinxcolwidth{1}{3}}
\sphinxAtStartPar
采样率
\sphinxbeforeendvarwidth
\end{varwidth}%
&\begin{varwidth}[t]{\sphinxcolwidth{1}{3}}
\sphinxAtStartPar
Edit
\sphinxbeforeendvarwidth
\end{varwidth}%
&\begin{varwidth}[t]{\sphinxcolwidth{1}{3}}
\sphinxAtStartPar
采样率设置
\sphinxbeforeendvarwidth
\end{varwidth}%
\\
\sphinxbottomrule
\end{tabulary}
\sphinxtableafterendhook\par
\sphinxattableend\end{savenotes}


\paragraph{采样率输入格式}
\label{\detokenize{chapters/chapter4:id51}}
\sphinxAtStartPar
支持多种分隔符格式：


\begin{savenotes}\sphinxattablestart
\sphinxthistablewithglobalstyle
\centering
\begin{tabulary}{\linewidth}[t]{TTT}
\sphinxtoprule
\begin{varwidth}[t]{\sphinxcolwidth{1}{3}}
\sphinxstyletheadfamily \sphinxAtStartPar
格式
\sphinxbeforeendvarwidth
\end{varwidth}%
&\begin{varwidth}[t]{\sphinxcolwidth{1}{3}}
\sphinxstyletheadfamily \sphinxAtStartPar
示例输入
\sphinxbeforeendvarwidth
\end{varwidth}%
&\begin{varwidth}[t]{\sphinxcolwidth{1}{3}}
\sphinxstyletheadfamily \sphinxAtStartPar
解析结果
\sphinxbeforeendvarwidth
\end{varwidth}%
\\
\sphinxmidrule
\sphinxtableatstartofbodyhook\begin{varwidth}[t]{\sphinxcolwidth{1}{3}}
\sphinxAtStartPar
单值
\sphinxbeforeendvarwidth
\end{varwidth}%
&\begin{varwidth}[t]{\sphinxcolwidth{1}{3}}
\sphinxAtStartPar
\sphinxcode{\sphinxupquote{2400}}
\sphinxbeforeendvarwidth
\end{varwidth}%
&\begin{varwidth}[t]{\sphinxcolwidth{1}{3}}
\sphinxAtStartPar
2400 Hz
\sphinxbeforeendvarwidth
\end{varwidth}%
\\
\sphinxhline\begin{varwidth}[t]{\sphinxcolwidth{1}{3}}
\sphinxAtStartPar
逗号分隔
\sphinxbeforeendvarwidth
\end{varwidth}%
&\begin{varwidth}[t]{\sphinxcolwidth{1}{3}}
\sphinxAtStartPar
\sphinxcode{\sphinxupquote{2400, 150, 15}}
\sphinxbeforeendvarwidth
\end{varwidth}%
&\begin{varwidth}[t]{\sphinxcolwidth{1}{3}}
\sphinxAtStartPar
2400, 150, 15 Hz
\sphinxbeforeendvarwidth
\end{varwidth}%
\\
\sphinxhline\begin{varwidth}[t]{\sphinxcolwidth{1}{3}}
\sphinxAtStartPar
空格分隔
\sphinxbeforeendvarwidth
\end{varwidth}%
&\begin{varwidth}[t]{\sphinxcolwidth{1}{3}}
\sphinxAtStartPar
\sphinxcode{\sphinxupquote{2400 150 15}}
\sphinxbeforeendvarwidth
\end{varwidth}%
&\begin{varwidth}[t]{\sphinxcolwidth{1}{3}}
\sphinxAtStartPar
2400, 150, 15 Hz
\sphinxbeforeendvarwidth
\end{varwidth}%
\\
\sphinxhline\begin{varwidth}[t]{\sphinxcolwidth{1}{3}}
\sphinxAtStartPar
分号分隔
\sphinxbeforeendvarwidth
\end{varwidth}%
&\begin{varwidth}[t]{\sphinxcolwidth{1}{3}}
\sphinxAtStartPar
\sphinxcode{\sphinxupquote{2400;150;15}}
\sphinxbeforeendvarwidth
\end{varwidth}%
&\begin{varwidth}[t]{\sphinxcolwidth{1}{3}}
\sphinxAtStartPar
2400, 150, 15 Hz
\sphinxbeforeendvarwidth
\end{varwidth}%
\\
\sphinxhline\begin{varwidth}[t]{\sphinxcolwidth{1}{3}}
\sphinxAtStartPar
换行分隔
\sphinxbeforeendvarwidth
\end{varwidth}%
&\begin{varwidth}[t]{\sphinxcolwidth{1}{3}}
\sphinxAtStartPar
\sphinxcode{\sphinxupquote{2400}} + \sphinxcode{\sphinxupquote{150}} + \sphinxcode{\sphinxupquote{15}}
\sphinxbeforeendvarwidth
\end{varwidth}%
&\begin{varwidth}[t]{\sphinxcolwidth{1}{3}}
\sphinxAtStartPar
2400, 150, 15 Hz
\sphinxbeforeendvarwidth
\end{varwidth}%
\\
\sphinxbottomrule
\end{tabulary}
\sphinxtableafterendhook\par
\sphinxattableend\end{savenotes}


\paragraph{验证规则}
\label{\detokenize{chapters/chapter4:id52}}\begin{itemize}
\item {} 
\sphinxAtStartPar
频段名称不能为空

\item {} 
\sphinxAtStartPar
采样率必须大于0

\item {} 
\sphinxAtStartPar
至少输入一个采样率

\end{itemize}


\bigskip\hrule\bigskip



\subsection{4.16 分组类型设置 (GroupTypeForm)}
\label{\detokenize{chapters/chapter4:grouptypeform}}
\sphinxAtStartPar
在MTDP中处理数据时，需要设置XPR（交叉功率比）分组方式。

\sphinxAtStartPar
\sphinxstylestrong{位置}: \sphinxcode{\sphinxupquote{处理 → 分组类型}}


\subsubsection{分组类型说明}
\label{\detokenize{chapters/chapter4:id53}}

\begin{savenotes}\sphinxattablestart
\sphinxthistablewithglobalstyle
\centering
\begin{tabulary}{\linewidth}[t]{TTT}
\sphinxtoprule
\begin{varwidth}[t]{\sphinxcolwidth{1}{3}}
\sphinxstyletheadfamily \sphinxAtStartPar
类型
\sphinxbeforeendvarwidth
\end{varwidth}%
&\begin{varwidth}[t]{\sphinxcolwidth{1}{3}}
\sphinxstyletheadfamily \sphinxAtStartPar
值
\sphinxbeforeendvarwidth
\end{varwidth}%
&\begin{varwidth}[t]{\sphinxcolwidth{1}{3}}
\sphinxstyletheadfamily \sphinxAtStartPar
说明
\sphinxbeforeendvarwidth
\end{varwidth}%
\\
\sphinxmidrule
\sphinxtableatstartofbodyhook\begin{varwidth}[t]{\sphinxcolwidth{1}{3}}
\sphinxAtStartPar
平均
\sphinxbeforeendvarwidth
\end{varwidth}%
&\begin{varwidth}[t]{\sphinxcolwidth{1}{3}}
\sphinxAtStartPar
0
\sphinxbeforeendvarwidth
\end{varwidth}%
&\begin{varwidth}[t]{\sphinxcolwidth{1}{3}}
\sphinxAtStartPar
所有数据平均分组
\sphinxbeforeendvarwidth
\end{varwidth}%
\\
\sphinxhline\begin{varwidth}[t]{\sphinxcolwidth{1}{3}}
\sphinxAtStartPar
排序
\sphinxbeforeendvarwidth
\end{varwidth}%
&\begin{varwidth}[t]{\sphinxcolwidth{1}{3}}
\sphinxAtStartPar
1
\sphinxbeforeendvarwidth
\end{varwidth}%
&\begin{varwidth}[t]{\sphinxcolwidth{1}{3}}
\sphinxAtStartPar
按相干度排序后分组
\sphinxbeforeendvarwidth
\end{varwidth}%
\\
\sphinxhline\begin{varwidth}[t]{\sphinxcolwidth{1}{3}}
\sphinxAtStartPar
随机
\sphinxbeforeendvarwidth
\end{varwidth}%
&\begin{varwidth}[t]{\sphinxcolwidth{1}{3}}
\sphinxAtStartPar
2
\sphinxbeforeendvarwidth
\end{varwidth}%
&\begin{varwidth}[t]{\sphinxcolwidth{1}{3}}
\sphinxAtStartPar
随机方式分组
\sphinxbeforeendvarwidth
\end{varwidth}%
\\
\sphinxbottomrule
\end{tabulary}
\sphinxtableafterendhook\par
\sphinxattableend\end{savenotes}


\subsubsection{MaxXPR参数}
\label{\detokenize{chapters/chapter4:maxxpr}}
\sphinxAtStartPar
MaxXPR（最大交叉功率比）用于设置分组阈值：
\begin{itemize}
\item {} 
\sphinxAtStartPar
值越大：接受更多数据，但可能包含噪声

\item {} 
\sphinxAtStartPar
值越小：剔除更多噪声，但可能丢失有效数据

\end{itemize}

\sphinxAtStartPar
\sphinxstylestrong{建议值}:
\begin{itemize}
\item {} 
\sphinxAtStartPar
保守设置: 1.0

\item {} 
\sphinxAtStartPar
标准设置: 1.5

\item {} 
\sphinxAtStartPar
激进设置: 2.0

\end{itemize}

\sphinxstepscope


\section{📊 第5章 时间序列处理}
\label{\detokenize{chapters/chapter5:id1}}\label{\detokenize{chapters/chapter5::doc}}
\sphinxAtStartPar
本章介绍时间序列数据的处理方法，包括FFT参数设置、频段管理、标定应用和滤波处理。


\bigskip\hrule\bigskip



\subsection{5.1 ⚙️ FFT参数设置概述}
\label{\detokenize{chapters/chapter5:fft}}
\sphinxAtStartPar
FFT参数设置对话框用于配置时间序列处理的各项参数。可通过以下方式打开：

\sphinxAtStartPar
1️⃣ 右键测段 → \sphinxcode{\sphinxupquote{FFT参数设置}}
2️⃣ 菜单 \sphinxcode{\sphinxupquote{处理 → FFT参数设置}}


\subsubsection{4.1.1 参数设置对话框结构}
\label{\detokenize{chapters/chapter5:id2}}
\sphinxAtStartPar
对话框分为以下几个主要区域：


\begin{savenotes}\sphinxattablestart
\sphinxthistablewithglobalstyle
\centering
\begin{tabulary}{\linewidth}[t]{TT}
\sphinxtoprule
\begin{varwidth}[t]{\sphinxcolwidth{1}{2}}
\sphinxstyletheadfamily \sphinxAtStartPar
区域
\sphinxbeforeendvarwidth
\end{varwidth}%
&\begin{varwidth}[t]{\sphinxcolwidth{1}{2}}
\sphinxstyletheadfamily \sphinxAtStartPar
内容
\sphinxbeforeendvarwidth
\end{varwidth}%
\\
\sphinxmidrule
\sphinxtableatstartofbodyhook\begin{varwidth}[t]{\sphinxcolwidth{1}{2}}
\sphinxAtStartPar
🌳 左侧树形视图
\sphinxbeforeendvarwidth
\end{varwidth}%
&\begin{varwidth}[t]{\sphinxcolwidth{1}{2}}
\sphinxAtStartPar
频段和子Schema管理
\sphinxbeforeendvarwidth
\end{varwidth}%
\\
\sphinxhline\begin{varwidth}[t]{\sphinxcolwidth{1}{2}}
\sphinxAtStartPar
📋 右侧频率列表
\sphinxbeforeendvarwidth
\end{varwidth}%
&\begin{varwidth}[t]{\sphinxcolwidth{1}{2}}
\sphinxAtStartPar
显示当前选中项的频率
\sphinxbeforeendvarwidth
\end{varwidth}%
\\
\sphinxhline\begin{varwidth}[t]{\sphinxcolwidth{1}{2}}
\sphinxAtStartPar
📊 上部参数区
\sphinxbeforeendvarwidth
\end{varwidth}%
&\begin{varwidth}[t]{\sphinxcolwidth{1}{2}}
\sphinxAtStartPar
FFT/估计/自动筛选参数
\sphinxbeforeendvarwidth
\end{varwidth}%
\\
\sphinxhline\begin{varwidth}[t]{\sphinxcolwidth{1}{2}}
\sphinxAtStartPar
🔘 下部按钮区
\sphinxbeforeendvarwidth
\end{varwidth}%
&\begin{varwidth}[t]{\sphinxcolwidth{1}{2}}
\sphinxAtStartPar
打开/保存/确认/取消
\sphinxbeforeendvarwidth
\end{varwidth}%
\\
\sphinxbottomrule
\end{tabulary}
\sphinxtableafterendhook\par
\sphinxattableend\end{savenotes}


\bigskip\hrule\bigskip



\subsection{5.2 📂 频段管理（MainSchema）}
\label{\detokenize{chapters/chapter5:mainschema}}

\subsubsection{5.2.1 频段结构}
\label{\detokenize{chapters/chapter5:id3}}
\sphinxAtStartPar
MTDP采用三级层次结构组织FFT参数：

\begin{sphinxVerbatim}[commandchars=\\\{\}]
📦 处理方案 (TTSPS)
└── 📊 频段 (TTSMainSchema)
    └── 📋 子Schema (TTSSubSchema)
        └── 📍 频率点
\end{sphinxVerbatim}


\subsubsection{5.2.2 频段属性}
\label{\detokenize{chapters/chapter5:id4}}

\begin{savenotes}\sphinxattablestart
\sphinxthistablewithglobalstyle
\centering
\begin{tabulary}{\linewidth}[t]{TTT}
\sphinxtoprule
\begin{varwidth}[t]{\sphinxcolwidth{1}{3}}
\sphinxstyletheadfamily \sphinxAtStartPar
属性
\sphinxbeforeendvarwidth
\end{varwidth}%
&\begin{varwidth}[t]{\sphinxcolwidth{1}{3}}
\sphinxstyletheadfamily \sphinxAtStartPar
类型
\sphinxbeforeendvarwidth
\end{varwidth}%
&\begin{varwidth}[t]{\sphinxcolwidth{1}{3}}
\sphinxstyletheadfamily \sphinxAtStartPar
说明
\sphinxbeforeendvarwidth
\end{varwidth}%
\\
\sphinxmidrule
\sphinxtableatstartofbodyhook\begin{varwidth}[t]{\sphinxcolwidth{1}{3}}
\sphinxAtStartPar
Name
\sphinxbeforeendvarwidth
\end{varwidth}%
&\begin{varwidth}[t]{\sphinxcolwidth{1}{3}}
\sphinxAtStartPar
String
\sphinxbeforeendvarwidth
\end{varwidth}%
&\begin{varwidth}[t]{\sphinxcolwidth{1}{3}}
\sphinxAtStartPar
频段名称（如TS2、TS3、TS4、TS5）
\sphinxbeforeendvarwidth
\end{varwidth}%
\\
\sphinxhline\begin{varwidth}[t]{\sphinxcolwidth{1}{3}}
\sphinxAtStartPar
SampleRate
\sphinxbeforeendvarwidth
\end{varwidth}%
&\begin{varwidth}[t]{\sphinxcolwidth{1}{3}}
\sphinxAtStartPar
Double1D
\sphinxbeforeendvarwidth
\end{varwidth}%
&\begin{varwidth}[t]{\sphinxcolwidth{1}{3}}
\sphinxAtStartPar
采样率数组（支持多个采样率）
\sphinxbeforeendvarwidth
\end{varwidth}%
\\
\sphinxhline\begin{varwidth}[t]{\sphinxcolwidth{1}{3}}
\sphinxAtStartPar
SubSchema
\sphinxbeforeendvarwidth
\end{varwidth}%
&\begin{varwidth}[t]{\sphinxcolwidth{1}{3}}
\sphinxAtStartPar
List
\sphinxbeforeendvarwidth
\end{varwidth}%
&\begin{varwidth}[t]{\sphinxcolwidth{1}{3}}
\sphinxAtStartPar
子Schema列表
\sphinxbeforeendvarwidth
\end{varwidth}%
\\
\sphinxbottomrule
\end{tabulary}
\sphinxtableafterendhook\par
\sphinxattableend\end{savenotes}


\subsubsection{5.2.3 频段操作}
\label{\detokenize{chapters/chapter5:id5}}

\begin{savenotes}\sphinxattablestart
\sphinxthistablewithglobalstyle
\centering
\begin{tabulary}{\linewidth}[t]{TT}
\sphinxtoprule
\begin{varwidth}[t]{\sphinxcolwidth{1}{2}}
\sphinxstyletheadfamily \sphinxAtStartPar
操作
\sphinxbeforeendvarwidth
\end{varwidth}%
&\begin{varwidth}[t]{\sphinxcolwidth{1}{2}}
\sphinxstyletheadfamily \sphinxAtStartPar
方法
\sphinxbeforeendvarwidth
\end{varwidth}%
\\
\sphinxmidrule
\sphinxtableatstartofbodyhook\begin{varwidth}[t]{\sphinxcolwidth{1}{2}}
\sphinxAtStartPar
添加频段
\sphinxbeforeendvarwidth
\end{varwidth}%
&\begin{varwidth}[t]{\sphinxcolwidth{1}{2}}
\sphinxAtStartPar
点击"添加频段"按钮或右键菜单
\sphinxbeforeendvarwidth
\end{varwidth}%
\\
\sphinxhline\begin{varwidth}[t]{\sphinxcolwidth{1}{2}}
\sphinxAtStartPar
编辑频段
\sphinxbeforeendvarwidth
\end{varwidth}%
&\begin{varwidth}[t]{\sphinxcolwidth{1}{2}}
\sphinxAtStartPar
双击频段节点或点击"编辑频段"
\sphinxbeforeendvarwidth
\end{varwidth}%
\\
\sphinxhline\begin{varwidth}[t]{\sphinxcolwidth{1}{2}}
\sphinxAtStartPar
删除频段
\sphinxbeforeendvarwidth
\end{varwidth}%
&\begin{varwidth}[t]{\sphinxcolwidth{1}{2}}
\sphinxAtStartPar
选中频段后点击"删除频段"
\sphinxbeforeendvarwidth
\end{varwidth}%
\\
\sphinxhline\begin{varwidth}[t]{\sphinxcolwidth{1}{2}}
\sphinxAtStartPar
复制频段
\sphinxbeforeendvarwidth
\end{varwidth}%
&\begin{varwidth}[t]{\sphinxcolwidth{1}{2}}
\sphinxAtStartPar
右键菜单 → 复制频段
\sphinxbeforeendvarwidth
\end{varwidth}%
\\
\sphinxbottomrule
\end{tabulary}
\sphinxtableafterendhook\par
\sphinxattableend\end{savenotes}


\subsubsection{5.2.4 预设频段配置}
\label{\detokenize{chapters/chapter5:id6}}
\sphinxAtStartPar
系统提供预设频段配置，位于 \sphinxcode{\sphinxupquote{Configurations\textbackslash{}}} 目录：


\begin{savenotes}\sphinxattablestart
\sphinxthistablewithglobalstyle
\centering
\begin{tabulary}{\linewidth}[t]{TT}
\sphinxtoprule
\begin{varwidth}[t]{\sphinxcolwidth{1}{2}}
\sphinxstyletheadfamily \sphinxAtStartPar
配置文件
\sphinxbeforeendvarwidth
\end{varwidth}%
&\begin{varwidth}[t]{\sphinxcolwidth{1}{2}}
\sphinxstyletheadfamily \sphinxAtStartPar
说明
\sphinxbeforeendvarwidth
\end{varwidth}%
\\
\sphinxmidrule
\sphinxtableatstartofbodyhook\begin{varwidth}[t]{\sphinxcolwidth{1}{2}}
\sphinxAtStartPar
MT2Octave.XMLTSPS
\sphinxbeforeendvarwidth
\end{varwidth}%
&\begin{varwidth}[t]{\sphinxcolwidth{1}{2}}
\sphinxAtStartPar
MT标准二倍频配置
\sphinxbeforeendvarwidth
\end{varwidth}%
\\
\sphinxhline\begin{varwidth}[t]{\sphinxcolwidth{1}{2}}
\sphinxAtStartPar
其他.XMLTSPS
\sphinxbeforeendvarwidth
\end{varwidth}%
&\begin{varwidth}[t]{\sphinxcolwidth{1}{2}}
\sphinxAtStartPar
用户自定义配置
\sphinxbeforeendvarwidth
\end{varwidth}%
\\
\sphinxbottomrule
\end{tabulary}
\sphinxtableafterendhook\par
\sphinxattableend\end{savenotes}


\bigskip\hrule\bigskip



\subsection{5.3 子Schema管理（SubSchema）}
\label{\detokenize{chapters/chapter5:schema-subschema}}

\subsubsection{4.3.1 子Schema属性}
\label{\detokenize{chapters/chapter5:schema}}

\begin{savenotes}\sphinxattablestart
\sphinxthistablewithglobalstyle
\centering
\begin{tabulary}{\linewidth}[t]{TTTT}
\sphinxtoprule
\begin{varwidth}[t]{\sphinxcolwidth{1}{4}}
\sphinxstyletheadfamily \sphinxAtStartPar
属性
\sphinxbeforeendvarwidth
\end{varwidth}%
&\begin{varwidth}[t]{\sphinxcolwidth{1}{4}}
\sphinxstyletheadfamily \sphinxAtStartPar
类型
\sphinxbeforeendvarwidth
\end{varwidth}%
&\begin{varwidth}[t]{\sphinxcolwidth{1}{4}}
\sphinxstyletheadfamily \sphinxAtStartPar
默认值
\sphinxbeforeendvarwidth
\end{varwidth}%
&\begin{varwidth}[t]{\sphinxcolwidth{1}{4}}
\sphinxstyletheadfamily \sphinxAtStartPar
说明
\sphinxbeforeendvarwidth
\end{varwidth}%
\\
\sphinxmidrule
\sphinxtableatstartofbodyhook\begin{varwidth}[t]{\sphinxcolwidth{1}{4}}
\sphinxAtStartPar
SampleLength
\sphinxbeforeendvarwidth
\end{varwidth}%
&\begin{varwidth}[t]{\sphinxcolwidth{1}{4}}
\sphinxAtStartPar
Double
\sphinxbeforeendvarwidth
\end{varwidth}%
&\begin{varwidth}[t]{\sphinxcolwidth{1}{4}}
\sphinxAtStartPar
4096
\sphinxbeforeendvarwidth
\end{varwidth}%
&\begin{varwidth}[t]{\sphinxcolwidth{1}{4}}
\sphinxAtStartPar
FFT长度（样本数）
\sphinxbeforeendvarwidth
\end{varwidth}%
\\
\sphinxhline\begin{varwidth}[t]{\sphinxcolwidth{1}{4}}
\sphinxAtStartPar
Overlap
\sphinxbeforeendvarwidth
\end{varwidth}%
&\begin{varwidth}[t]{\sphinxcolwidth{1}{4}}
\sphinxAtStartPar
Double
\sphinxbeforeendvarwidth
\end{varwidth}%
&\begin{varwidth}[t]{\sphinxcolwidth{1}{4}}
\sphinxAtStartPar
0.5
\sphinxbeforeendvarwidth
\end{varwidth}%
&\begin{varwidth}[t]{\sphinxcolwidth{1}{4}}
\sphinxAtStartPar
重叠率（0\sphinxhyphen{}0.75）
\sphinxbeforeendvarwidth
\end{varwidth}%
\\
\sphinxhline\begin{varwidth}[t]{\sphinxcolwidth{1}{4}}
\sphinxAtStartPar
Frequency
\sphinxbeforeendvarwidth
\end{varwidth}%
&\begin{varwidth}[t]{\sphinxcolwidth{1}{4}}
\sphinxAtStartPar
Double1D
\sphinxbeforeendvarwidth
\end{varwidth}%
&\begin{varwidth}[t]{\sphinxcolwidth{1}{4}}
\sphinxAtStartPar
\sphinxhyphen{}
\sphinxbeforeendvarwidth
\end{varwidth}%
&\begin{varwidth}[t]{\sphinxcolwidth{1}{4}}
\sphinxAtStartPar
频率数组
\sphinxbeforeendvarwidth
\end{varwidth}%
\\
\sphinxhline\begin{varwidth}[t]{\sphinxcolwidth{1}{4}}
\sphinxAtStartPar
MaxXPR
\sphinxbeforeendvarwidth
\end{varwidth}%
&\begin{varwidth}[t]{\sphinxcolwidth{1}{4}}
\sphinxAtStartPar
Integer
\sphinxbeforeendvarwidth
\end{varwidth}%
&\begin{varwidth}[t]{\sphinxcolwidth{1}{4}}
\sphinxAtStartPar
100
\sphinxbeforeendvarwidth
\end{varwidth}%
&\begin{varwidth}[t]{\sphinxcolwidth{1}{4}}
\sphinxAtStartPar
最大XPR值
\sphinxbeforeendvarwidth
\end{varwidth}%
\\
\sphinxhline\begin{varwidth}[t]{\sphinxcolwidth{1}{4}}
\sphinxAtStartPar
GroupingType
\sphinxbeforeendvarwidth
\end{varwidth}%
&\begin{varwidth}[t]{\sphinxcolwidth{1}{4}}
\sphinxAtStartPar
Integer
\sphinxbeforeendvarwidth
\end{varwidth}%
&\begin{varwidth}[t]{\sphinxcolwidth{1}{4}}
\sphinxAtStartPar
0
\sphinxbeforeendvarwidth
\end{varwidth}%
&\begin{varwidth}[t]{\sphinxcolwidth{1}{4}}
\sphinxAtStartPar
分组类型
\sphinxbeforeendvarwidth
\end{varwidth}%
\\
\sphinxbottomrule
\end{tabulary}
\sphinxtableafterendhook\par
\sphinxattableend\end{savenotes}


\subsubsection{4.3.2 FFT长度设置}
\label{\detokenize{chapters/chapter5:id7}}
\sphinxAtStartPar
FFT长度影响频率分辨率和计算时间：


\begin{savenotes}\sphinxattablestart
\sphinxthistablewithglobalstyle
\centering
\begin{tabulary}{\linewidth}[t]{TTT}
\sphinxtoprule
\begin{varwidth}[t]{\sphinxcolwidth{1}{3}}
\sphinxstyletheadfamily \sphinxAtStartPar
采样率
\sphinxbeforeendvarwidth
\end{varwidth}%
&\begin{varwidth}[t]{\sphinxcolwidth{1}{3}}
\sphinxstyletheadfamily \sphinxAtStartPar
推荐FFT长度
\sphinxbeforeendvarwidth
\end{varwidth}%
&\begin{varwidth}[t]{\sphinxcolwidth{1}{3}}
\sphinxstyletheadfamily \sphinxAtStartPar
频率分辨率
\sphinxbeforeendvarwidth
\end{varwidth}%
\\
\sphinxmidrule
\sphinxtableatstartofbodyhook\begin{varwidth}[t]{\sphinxcolwidth{1}{3}}
\sphinxAtStartPar
24000 Hz
\sphinxbeforeendvarwidth
\end{varwidth}%
&\begin{varwidth}[t]{\sphinxcolwidth{1}{3}}
\sphinxAtStartPar
65536
\sphinxbeforeendvarwidth
\end{varwidth}%
&\begin{varwidth}[t]{\sphinxcolwidth{1}{3}}
\sphinxAtStartPar
0.37 Hz
\sphinxbeforeendvarwidth
\end{varwidth}%
\\
\sphinxhline\begin{varwidth}[t]{\sphinxcolwidth{1}{3}}
\sphinxAtStartPar
2400 Hz
\sphinxbeforeendvarwidth
\end{varwidth}%
&\begin{varwidth}[t]{\sphinxcolwidth{1}{3}}
\sphinxAtStartPar
8192
\sphinxbeforeendvarwidth
\end{varwidth}%
&\begin{varwidth}[t]{\sphinxcolwidth{1}{3}}
\sphinxAtStartPar
0.29 Hz
\sphinxbeforeendvarwidth
\end{varwidth}%
\\
\sphinxhline\begin{varwidth}[t]{\sphinxcolwidth{1}{3}}
\sphinxAtStartPar
150 Hz
\sphinxbeforeendvarwidth
\end{varwidth}%
&\begin{varwidth}[t]{\sphinxcolwidth{1}{3}}
\sphinxAtStartPar
4096
\sphinxbeforeendvarwidth
\end{varwidth}%
&\begin{varwidth}[t]{\sphinxcolwidth{1}{3}}
\sphinxAtStartPar
0.04 Hz
\sphinxbeforeendvarwidth
\end{varwidth}%
\\
\sphinxhline\begin{varwidth}[t]{\sphinxcolwidth{1}{3}}
\sphinxAtStartPar
15 Hz
\sphinxbeforeendvarwidth
\end{varwidth}%
&\begin{varwidth}[t]{\sphinxcolwidth{1}{3}}
\sphinxAtStartPar
4096
\sphinxbeforeendvarwidth
\end{varwidth}%
&\begin{varwidth}[t]{\sphinxcolwidth{1}{3}}
\sphinxAtStartPar
0.004 Hz
\sphinxbeforeendvarwidth
\end{varwidth}%
\\
\sphinxbottomrule
\end{tabulary}
\sphinxtableafterendhook\par
\sphinxattableend\end{savenotes}


\subsubsection{4.3.3 重叠率设置}
\label{\detokenize{chapters/chapter5:id8}}

\begin{savenotes}\sphinxattablestart
\sphinxthistablewithglobalstyle
\centering
\begin{tabulary}{\linewidth}[t]{TT}
\sphinxtoprule
\begin{varwidth}[t]{\sphinxcolwidth{1}{2}}
\sphinxstyletheadfamily \sphinxAtStartPar
重叠率
\sphinxbeforeendvarwidth
\end{varwidth}%
&\begin{varwidth}[t]{\sphinxcolwidth{1}{2}}
\sphinxstyletheadfamily \sphinxAtStartPar
适用场景
\sphinxbeforeendvarwidth
\end{varwidth}%
\\
\sphinxmidrule
\sphinxtableatstartofbodyhook\begin{varwidth}[t]{\sphinxcolwidth{1}{2}}
\sphinxAtStartPar
50\%
\sphinxbeforeendvarwidth
\end{varwidth}%
&\begin{varwidth}[t]{\sphinxcolwidth{1}{2}}
\sphinxAtStartPar
最常用，平衡效率和精度
\sphinxbeforeendvarwidth
\end{varwidth}%
\\
\sphinxhline\begin{varwidth}[t]{\sphinxcolwidth{1}{2}}
\sphinxAtStartPar
75\%
\sphinxbeforeendvarwidth
\end{varwidth}%
&\begin{varwidth}[t]{\sphinxcolwidth{1}{2}}
\sphinxAtStartPar
更高数据利用率，更多FFT窗口
\sphinxbeforeendvarwidth
\end{varwidth}%
\\
\sphinxhline\begin{varwidth}[t]{\sphinxcolwidth{1}{2}}
\sphinxAtStartPar
0\%
\sphinxbeforeendvarwidth
\end{varwidth}%
&\begin{varwidth}[t]{\sphinxcolwidth{1}{2}}
\sphinxAtStartPar
计算最快，但窗口数最少
\sphinxbeforeendvarwidth
\end{varwidth}%
\\
\sphinxbottomrule
\end{tabulary}
\sphinxtableafterendhook\par
\sphinxattableend\end{savenotes}


\subsubsection{4.3.4 子Schema操作}
\label{\detokenize{chapters/chapter5:id9}}

\begin{savenotes}\sphinxattablestart
\sphinxthistablewithglobalstyle
\centering
\begin{tabulary}{\linewidth}[t]{TT}
\sphinxtoprule
\begin{varwidth}[t]{\sphinxcolwidth{1}{2}}
\sphinxstyletheadfamily \sphinxAtStartPar
操作
\sphinxbeforeendvarwidth
\end{varwidth}%
&\begin{varwidth}[t]{\sphinxcolwidth{1}{2}}
\sphinxstyletheadfamily \sphinxAtStartPar
方法
\sphinxbeforeendvarwidth
\end{varwidth}%
\\
\sphinxmidrule
\sphinxtableatstartofbodyhook\begin{varwidth}[t]{\sphinxcolwidth{1}{2}}
\sphinxAtStartPar
添加子Schema
\sphinxbeforeendvarwidth
\end{varwidth}%
&\begin{varwidth}[t]{\sphinxcolwidth{1}{2}}
\sphinxAtStartPar
右键频段 → 添加子Schema
\sphinxbeforeendvarwidth
\end{varwidth}%
\\
\sphinxhline\begin{varwidth}[t]{\sphinxcolwidth{1}{2}}
\sphinxAtStartPar
删除子Schema
\sphinxbeforeendvarwidth
\end{varwidth}%
&\begin{varwidth}[t]{\sphinxcolwidth{1}{2}}
\sphinxAtStartPar
选中子Schema → 右键删除
\sphinxbeforeendvarwidth
\end{varwidth}%
\\
\sphinxbottomrule
\end{tabulary}
\sphinxtableafterendhook\par
\sphinxattableend\end{savenotes}


\bigskip\hrule\bigskip



\subsection{5.4 频率管理}
\label{\detokenize{chapters/chapter5:id10}}

\subsubsection{4.4.1 频率列表操作}
\label{\detokenize{chapters/chapter5:id11}}

\begin{savenotes}\sphinxattablestart
\sphinxthistablewithglobalstyle
\centering
\begin{tabulary}{\linewidth}[t]{TT}
\sphinxtoprule
\begin{varwidth}[t]{\sphinxcolwidth{1}{2}}
\sphinxstyletheadfamily \sphinxAtStartPar
操作
\sphinxbeforeendvarwidth
\end{varwidth}%
&\begin{varwidth}[t]{\sphinxcolwidth{1}{2}}
\sphinxstyletheadfamily \sphinxAtStartPar
说明
\sphinxbeforeendvarwidth
\end{varwidth}%
\\
\sphinxmidrule
\sphinxtableatstartofbodyhook\begin{varwidth}[t]{\sphinxcolwidth{1}{2}}
\sphinxAtStartPar
添加频率
\sphinxbeforeendvarwidth
\end{varwidth}%
&\begin{varwidth}[t]{\sphinxcolwidth{1}{2}}
\sphinxAtStartPar
手动输入单个频率值
\sphinxbeforeendvarwidth
\end{varwidth}%
\\
\sphinxhline\begin{varwidth}[t]{\sphinxcolwidth{1}{2}}
\sphinxAtStartPar
批量添加
\sphinxbeforeendvarwidth
\end{varwidth}%
&\begin{varwidth}[t]{\sphinxcolwidth{1}{2}}
\sphinxAtStartPar
打开批量添加对话框，生成多个频率
\sphinxbeforeendvarwidth
\end{varwidth}%
\\
\sphinxhline\begin{varwidth}[t]{\sphinxcolwidth{1}{2}}
\sphinxAtStartPar
编辑频率
\sphinxbeforeendvarwidth
\end{varwidth}%
&\begin{varwidth}[t]{\sphinxcolwidth{1}{2}}
\sphinxAtStartPar
修改选中频率的值
\sphinxbeforeendvarwidth
\end{varwidth}%
\\
\sphinxhline\begin{varwidth}[t]{\sphinxcolwidth{1}{2}}
\sphinxAtStartPar
删除频率
\sphinxbeforeendvarwidth
\end{varwidth}%
&\begin{varwidth}[t]{\sphinxcolwidth{1}{2}}
\sphinxAtStartPar
删除选中的频率
\sphinxbeforeendvarwidth
\end{varwidth}%
\\
\sphinxhline\begin{varwidth}[t]{\sphinxcolwidth{1}{2}}
\sphinxAtStartPar
清空频率
\sphinxbeforeendvarwidth
\end{varwidth}%
&\begin{varwidth}[t]{\sphinxcolwidth{1}{2}}
\sphinxAtStartPar
清空当前子Schema的所有频率
\sphinxbeforeendvarwidth
\end{varwidth}%
\\
\sphinxhline\begin{varwidth}[t]{\sphinxcolwidth{1}{2}}
\sphinxAtStartPar
升序排序
\sphinxbeforeendvarwidth
\end{varwidth}%
&\begin{varwidth}[t]{\sphinxcolwidth{1}{2}}
\sphinxAtStartPar
按频率升序排列
\sphinxbeforeendvarwidth
\end{varwidth}%
\\
\sphinxhline\begin{varwidth}[t]{\sphinxcolwidth{1}{2}}
\sphinxAtStartPar
降序排序
\sphinxbeforeendvarwidth
\end{varwidth}%
&\begin{varwidth}[t]{\sphinxcolwidth{1}{2}}
\sphinxAtStartPar
按频率降序排列
\sphinxbeforeendvarwidth
\end{varwidth}%
\\
\sphinxbottomrule
\end{tabulary}
\sphinxtableafterendhook\par
\sphinxattableend\end{savenotes}


\subsubsection{4.4.2 批量生成频率}
\label{\detokenize{chapters/chapter5:id12}}
\sphinxAtStartPar
批量添加频点对话框支持三种频率生成方式：


\paragraph{📊 方式1：等比分布（Logarithmic）}
\label{\detokenize{chapters/chapter5:logarithmic}}
\sphinxAtStartPar
按对数等间隔生成频率，频率比为常数。

\sphinxAtStartPar
\sphinxstylestrong{参数：}
\begin{itemize}
\item {} 
\sphinxAtStartPar
最低频率（Hz）

\item {} 
\sphinxAtStartPar
最高频率（Hz）

\item {} 
\sphinxAtStartPar
频点数量

\end{itemize}

\sphinxAtStartPar
\sphinxstylestrong{计算公式：}

\begin{sphinxVerbatim}[commandchars=\\\{\}]
\PYG{n}{ratio} \PYG{o}{=} \PYG{p}{(}\PYG{n}{f\PYGZus{}max} \PYG{o}{/} \PYG{n}{f\PYGZus{}min}\PYG{p}{)} \PYG{o}{\PYGZca{}} \PYG{p}{(}\PYG{l+m+mi}{1} \PYG{o}{/} \PYG{p}{(}\PYG{n}{n}\PYG{o}{\PYGZhy{}}\PYG{l+m+mi}{1}\PYG{p}{)}\PYG{p}{)}
\PYG{n}{f}\PYG{p}{[}\PYG{n}{i}\PYG{p}{]} \PYG{o}{=} \PYG{n}{f\PYGZus{}min} \PYG{o}{*} \PYG{n}{ratio}\PYG{o}{\PYGZca{}}\PYG{n}{i}
\end{sphinxVerbatim}

\sphinxAtStartPar
\sphinxstylestrong{适用场景：}
\begin{itemize}
\item {} 
\sphinxAtStartPar
MT数据处理的标准方式

\item {} 
\sphinxAtStartPar
符合电磁测深的频率分布规律

\item {} 
\sphinxAtStartPar
低频段和高频段都有合理的覆盖

\end{itemize}


\paragraph{📏 方式2：等差分布（Linear）}
\label{\detokenize{chapters/chapter5:linear}}
\sphinxAtStartPar
按线性等间隔生成频率，频率差为常数。

\sphinxAtStartPar
\sphinxstylestrong{参数：}
\begin{itemize}
\item {} 
\sphinxAtStartPar
最低频率（Hz）

\item {} 
\sphinxAtStartPar
最高频率（Hz）

\item {} 
\sphinxAtStartPar
频点数量

\end{itemize}

\sphinxAtStartPar
\sphinxstylestrong{计算公式：}

\begin{sphinxVerbatim}[commandchars=\\\{\}]
\PYG{n}{step} \PYG{o}{=} \PYG{p}{(}\PYG{n}{f\PYGZus{}max} \PYG{o}{\PYGZhy{}} \PYG{n}{f\PYGZus{}min}\PYG{p}{)} \PYG{o}{/} \PYG{p}{(}\PYG{n}{n}\PYG{o}{\PYGZhy{}}\PYG{l+m+mi}{1}\PYG{p}{)}
\PYG{n}{f}\PYG{p}{[}\PYG{n}{i}\PYG{p}{]} \PYG{o}{=} \PYG{n}{f\PYGZus{}min} \PYG{o}{+} \PYG{n}{step} \PYG{o}{*} \PYG{n}{i}
\end{sphinxVerbatim}

\sphinxAtStartPar
\sphinxstylestrong{适用场景：}
\begin{itemize}
\item {} 
\sphinxAtStartPar
特定频段的精细分析

\item {} 
\sphinxAtStartPar
需要均匀频率分辨率

\item {} 
\sphinxAtStartPar
CSAMT等人工源数据处理

\end{itemize}


\paragraph{🔧 方式3：自定义间距（Custom）}
\label{\detokenize{chapters/chapter5:custom}}
\sphinxAtStartPar
按用户指定的频率比值生成频率。

\sphinxAtStartPar
\sphinxstylestrong{参数：}
\begin{itemize}
\item {} 
\sphinxAtStartPar
起始频率（Hz）

\item {} 
\sphinxAtStartPar
频率比值（ratio）

\item {} 
\sphinxAtStartPar
频点数量

\end{itemize}

\sphinxAtStartPar
\sphinxstylestrong{计算公式：}

\begin{sphinxVerbatim}[commandchars=\\\{\}]
\PYG{n}{f}\PYG{p}{[}\PYG{n}{i}\PYG{p}{]} \PYG{o}{=} \PYG{n}{f\PYGZus{}start} \PYG{o}{*} \PYG{n}{ratio}\PYG{o}{\PYGZca{}}\PYG{n}{i}
\end{sphinxVerbatim}

\sphinxAtStartPar
\sphinxstylestrong{适用场景：}
\begin{itemize}
\item {} 
\sphinxAtStartPar
特殊频率分布需求

\item {} 
\sphinxAtStartPar
针对特定频段的优化

\item {} 
\sphinxAtStartPar
与其他软件频率点匹配

\end{itemize}


\paragraph{📌 Phoenix多频段模式}
\label{\detokenize{chapters/chapter5:phoenix}}
\sphinxAtStartPar
专门为Phoenix仪器设计的自动分配模式：
\begin{itemize}
\item {} 
\sphinxAtStartPar
自动将频率分配到TS2/TS3/TS4/TS5频段

\item {} 
\sphinxAtStartPar
根据采样率自动选择合适的频段

\item {} 
\sphinxAtStartPar
支持多采样率混合处理

\end{itemize}

\sphinxAtStartPar
\sphinxstylestrong{附加参数（所有方式通用）：}
\begin{itemize}
\item {} 
\sphinxAtStartPar
FFT长度：影响频率分辨率

\item {} 
\sphinxAtStartPar
重叠率：0\%\sphinxhyphen{}75\%

\item {} 
\sphinxAtStartPar
MaxXPR：最大XPR值限制

\end{itemize}


\subsubsection{4.4.3 频率显示格式}
\label{\detokenize{chapters/chapter5:id13}}
\sphinxAtStartPar
系统根据频率大小自动调整显示精度：


\begin{savenotes}\sphinxattablestart
\sphinxthistablewithglobalstyle
\centering
\begin{tabulary}{\linewidth}[t]{TT}
\sphinxtoprule
\begin{varwidth}[t]{\sphinxcolwidth{1}{2}}
\sphinxstyletheadfamily \sphinxAtStartPar
频率范围
\sphinxbeforeendvarwidth
\end{varwidth}%
&\begin{varwidth}[t]{\sphinxcolwidth{1}{2}}
\sphinxstyletheadfamily \sphinxAtStartPar
小数位数
\sphinxbeforeendvarwidth
\end{varwidth}%
\\
\sphinxmidrule
\sphinxtableatstartofbodyhook\begin{varwidth}[t]{\sphinxcolwidth{1}{2}}
\sphinxAtStartPar
>= 100 kHz
\sphinxbeforeendvarwidth
\end{varwidth}%
&\begin{varwidth}[t]{\sphinxcolwidth{1}{2}}
\sphinxAtStartPar
0位
\sphinxbeforeendvarwidth
\end{varwidth}%
\\
\sphinxhline\begin{varwidth}[t]{\sphinxcolwidth{1}{2}}
\sphinxAtStartPar
>= 10 kHz
\sphinxbeforeendvarwidth
\end{varwidth}%
&\begin{varwidth}[t]{\sphinxcolwidth{1}{2}}
\sphinxAtStartPar
1位
\sphinxbeforeendvarwidth
\end{varwidth}%
\\
\sphinxhline\begin{varwidth}[t]{\sphinxcolwidth{1}{2}}
\sphinxAtStartPar
>= 1 kHz
\sphinxbeforeendvarwidth
\end{varwidth}%
&\begin{varwidth}[t]{\sphinxcolwidth{1}{2}}
\sphinxAtStartPar
2位
\sphinxbeforeendvarwidth
\end{varwidth}%
\\
\sphinxhline\begin{varwidth}[t]{\sphinxcolwidth{1}{2}}
\sphinxAtStartPar
>= 100 Hz
\sphinxbeforeendvarwidth
\end{varwidth}%
&\begin{varwidth}[t]{\sphinxcolwidth{1}{2}}
\sphinxAtStartPar
3位
\sphinxbeforeendvarwidth
\end{varwidth}%
\\
\sphinxhline\begin{varwidth}[t]{\sphinxcolwidth{1}{2}}
\sphinxAtStartPar
>= 10 Hz
\sphinxbeforeendvarwidth
\end{varwidth}%
&\begin{varwidth}[t]{\sphinxcolwidth{1}{2}}
\sphinxAtStartPar
4位
\sphinxbeforeendvarwidth
\end{varwidth}%
\\
\sphinxhline\begin{varwidth}[t]{\sphinxcolwidth{1}{2}}
\sphinxAtStartPar
>= 1 Hz
\sphinxbeforeendvarwidth
\end{varwidth}%
&\begin{varwidth}[t]{\sphinxcolwidth{1}{2}}
\sphinxAtStartPar
5位
\sphinxbeforeendvarwidth
\end{varwidth}%
\\
\sphinxhline\begin{varwidth}[t]{\sphinxcolwidth{1}{2}}
\sphinxAtStartPar
< 1 Hz
\sphinxbeforeendvarwidth
\end{varwidth}%
&\begin{varwidth}[t]{\sphinxcolwidth{1}{2}}
\sphinxAtStartPar
6位
\sphinxbeforeendvarwidth
\end{varwidth}%
\\
\sphinxbottomrule
\end{tabulary}
\sphinxtableafterendhook\par
\sphinxattableend\end{savenotes}


\bigskip\hrule\bigskip



\subsection{5.5 全局FFT参数}
\label{\detokenize{chapters/chapter5:id14}}

\subsubsection{4.5.1 窗函数}
\label{\detokenize{chapters/chapter5:id15}}
\sphinxAtStartPar
FFT窗函数影响频谱分析的频率分辨率和泄漏特性。MTDP支持多种窗函数：


\begin{savenotes}\sphinxattablestart
\sphinxthistablewithglobalstyle
\centering
\begin{tabulary}{\linewidth}[t]{TTTTTT}
\sphinxtoprule
\begin{varwidth}[t]{\sphinxcolwidth{1}{6}}
\sphinxstyletheadfamily \sphinxAtStartPar
值
\sphinxbeforeendvarwidth
\end{varwidth}%
&\begin{varwidth}[t]{\sphinxcolwidth{1}{6}}
\sphinxstyletheadfamily \sphinxAtStartPar
窗函数
\sphinxbeforeendvarwidth
\end{varwidth}%
&\begin{varwidth}[t]{\sphinxcolwidth{1}{6}}
\sphinxstyletheadfamily \sphinxAtStartPar
特性
\sphinxbeforeendvarwidth
\end{varwidth}%
&\begin{varwidth}[t]{\sphinxcolwidth{1}{6}}
\sphinxstyletheadfamily \sphinxAtStartPar
适用场景
\sphinxbeforeendvarwidth
\end{varwidth}%
&\begin{varwidth}[t]{\sphinxcolwidth{1}{6}}
\sphinxstyletheadfamily \sphinxAtStartPar
频率泄漏
\sphinxbeforeendvarwidth
\end{varwidth}%
&\begin{varwidth}[t]{\sphinxcolwidth{1}{6}}
\sphinxstyletheadfamily \sphinxAtStartPar
主瓣宽
\sphinxbeforeendvarwidth
\end{varwidth}%
\\
\sphinxmidrule
\sphinxtableatstartofbodyhook\begin{varwidth}[t]{\sphinxcolwidth{1}{6}}
\sphinxAtStartPar
0
\sphinxbeforeendvarwidth
\end{varwidth}%
&\begin{varwidth}[t]{\sphinxcolwidth{1}{6}}
\sphinxAtStartPar
矩形窗
\sphinxbeforeendvarwidth
\end{varwidth}%
&\begin{varwidth}[t]{\sphinxcolwidth{1}{6}}
\sphinxAtStartPar
矩形截断
\sphinxbeforeendvarwidth
\end{varwidth}%
&\begin{varwidth}[t]{\sphinxcolwidth{1}{6}}
\sphinxAtStartPar
瞬态信号分析
\sphinxbeforeendvarwidth
\end{varwidth}%
&\begin{varwidth}[t]{\sphinxcolwidth{1}{6}}
\sphinxAtStartPar
严重
\sphinxbeforeendvarwidth
\end{varwidth}%
&\begin{varwidth}[t]{\sphinxcolwidth{1}{6}}
\sphinxAtStartPar
4π/N
\sphinxbeforeendvarwidth
\end{varwidth}%
\\
\sphinxhline\begin{varwidth}[t]{\sphinxcolwidth{1}{6}}
\sphinxAtStartPar
1
\sphinxbeforeendvarwidth
\end{varwidth}%
&\begin{varwidth}[t]{\sphinxcolwidth{1}{6}}
\sphinxAtStartPar
汉明窗
\sphinxbeforeendvarwidth
\end{varwidth}%
&\begin{varwidth}[t]{\sphinxcolwidth{1}{6}}
\sphinxAtStartPar
起始和结束逐渐衰减
\sphinxbeforeendvarwidth
\end{varwidth}%
&\begin{varwidth}[t]{\sphinxcolwidth{1}{6}}
\sphinxAtStartPar
一般信号
\sphinxbeforeendvarwidth
\end{varwidth}%
&\begin{varwidth}[t]{\sphinxcolwidth{1}{6}}
\sphinxAtStartPar
中等
\sphinxbeforeendvarwidth
\end{varwidth}%
&\begin{varwidth}[t]{\sphinxcolwidth{1}{6}}
\sphinxAtStartPar
8π/N
\sphinxbeforeendvarwidth
\end{varwidth}%
\\
\sphinxhline\begin{varwidth}[t]{\sphinxcolwidth{1}{6}}
\sphinxAtStartPar
2
\sphinxbeforeendvarwidth
\end{varwidth}%
&\begin{varwidth}[t]{\sphinxcolwidth{1}{6}}
\sphinxAtStartPar
汉宁窗
\sphinxbeforeendvarwidth
\end{varwidth}%
&\begin{varwidth}[t]{\sphinxcolwidth{1}{6}}
\sphinxAtStartPar
主瓣宽较窄
\sphinxbeforeendvarwidth
\end{varwidth}%
&\begin{varwidth}[t]{\sphinxcolwidth{1}{6}}
\sphinxAtStartPar
较平稳信号
\sphinxbeforeendvarwidth
\end{varwidth}%
&\begin{varwidth}[t]{\sphinxcolwidth{1}{6}}
\sphinxAtStartPar
较小
\sphinxbeforeendvarwidth
\end{varwidth}%
&\begin{varwidth}[t]{\sphinxcolwidth{1}{6}}
\sphinxAtStartPar
8π/N
\sphinxbeforeendvarwidth
\end{varwidth}%
\\
\sphinxhline\begin{varwidth}[t]{\sphinxcolwidth{1}{6}}
\sphinxAtStartPar
3
\sphinxbeforeendvarwidth
\end{varwidth}%
&\begin{varwidth}[t]{\sphinxcolwidth{1}{6}}
\sphinxAtStartPar
布莱克曼窗
\sphinxbeforeendvarwidth
\end{varwidth}%
&\begin{varwidth}[t]{\sphinxcolwidth{1}{6}}
\sphinxAtStartPar
主瓣宽度最小
\sphinxbeforeendvarwidth
\end{varwidth}%
&\begin{varwidth}[t]{\sphinxcolwidth{1}{6}}
\sphinxAtStartPar
短暂态信号
\sphinxbeforeendvarwidth
\end{varwidth}%
&\begin{varwidth}[t]{\sphinxcolwidth{1}{6}}
\sphinxAtStartPar
极小
\sphinxbeforeendvarwidth
\end{varwidth}%
&\begin{varwidth}[t]{\sphinxcolwidth{1}{6}}
\sphinxAtStartPar
12π/N
\sphinxbeforeendvarwidth
\end{varwidth}%
\\
\sphinxhline\begin{varwidth}[t]{\sphinxcolwidth{1}{6}}
\sphinxAtStartPar
4
\sphinxbeforeendvarwidth
\end{varwidth}%
&\begin{varwidth}[t]{\sphinxcolwidth{1}{6}}
\sphinxAtStartPar
巴特利特窗
\sphinxbeforeendvarwidth
\end{varwidth}%
&\begin{varwidth}[t]{\sphinxcolwidth{1}{6}}
\sphinxAtStartPar
多个主瓣叠加
\sphinxbeforeendvarwidth
\end{varwidth}%
&\begin{varwidth}[t]{\sphinxcolwidth{1}{6}}
\sphinxAtStartPar
调制信号
\sphinxbeforeendvarwidth
\end{varwidth}%
&\begin{varwidth}[t]{\sphinxcolwidth{1}{6}}
\sphinxAtStartPar
较小
\sphinxbeforeendvarwidth
\end{varwidth}%
&\begin{varwidth}[t]{\sphinxcolwidth{1}{6}}
\sphinxAtStartPar
8π/N
\sphinxbeforeendvarwidth
\end{varwidth}%
\\
\sphinxhline\begin{varwidth}[t]{\sphinxcolwidth{1}{6}}
\sphinxAtStartPar
5
\sphinxbeforeendvarwidth
\end{varwidth}%
&\begin{varwidth}[t]{\sphinxcolwidth{1}{6}}
\sphinxAtStartPar
平顶窗
\sphinxbeforeendvarwidth
\end{varwidth}%
&\begin{varwidth}[t]{\sphinxcolwidth{1}{6}}
\sphinxAtStartPar
旁瓣抑制
\sphinxbeforeendvarwidth
\end{varwidth}%
&\begin{varwidth}[t]{\sphinxcolwidth{1}{6}}
\sphinxAtStartPar
稳态/谐波
\sphinxbeforeendvarwidth
\end{varwidth}%
&\begin{varwidth}[t]{\sphinxcolwidth{1}{6}}
\sphinxAtStartPar
极小
\sphinxbeforeendvarwidth
\end{varwidth}%
&\begin{varwidth}[t]{\sphinxcolwidth{1}{6}}
\sphinxAtStartPar
12π/N
\sphinxbeforeendvarwidth
\end{varwidth}%
\\
\sphinxhline\begin{varwidth}[t]{\sphinxcolwidth{1}{6}}
\sphinxAtStartPar
6
\sphinxbeforeendvarwidth
\end{varwidth}%
&\begin{varwidth}[t]{\sphinxcolwidth{1}{6}}
\sphinxAtStartPar
高斯窗
\sphinxbeforeendvarwidth
\end{varwidth}%
&\begin{varwidth}[t]{\sphinxcolwidth{1}{6}}
\sphinxAtStartPar
指数级衰减
\sphinxbeforeendvarwidth
\end{varwidth}%
&\begin{varwidth}[t]{\sphinxcolwidth{1}{6}}
\sphinxAtStartPar
瞬态/平滑
\sphinxbeforeendvarwidth
\end{varwidth}%
&\begin{varwidth}[t]{\sphinxcolwidth{1}{6}}
\sphinxAtStartPar
小
\sphinxbeforeendvarwidth
\end{varwidth}%
&\begin{varwidth}[t]{\sphinxcolwidth{1}{6}}
\sphinxAtStartPar
8π/N
\sphinxbeforeendvarwidth
\end{varwidth}%
\\
\sphinxbottomrule
\end{tabulary}
\sphinxtableafterendhook\par
\sphinxattableend\end{savenotes}


\subsubsection{4.5.2 窗函数选择指南}
\label{\detokenize{chapters/chapter5:id16}}
\sphinxAtStartPar
\sphinxstylestrong{选择原则：}
\begin{enumerate}
\sphinxsetlistlabels{\arabic}{enumi}{enumii}{}{.}%
\item {} 
\sphinxAtStartPar
\sphinxstylestrong{根据信号特征选择}
\begin{itemize}
\item {} 
\sphinxAtStartPar
瞬态脉冲 → 矩形窗

\item {} 
\sphinxAtStartPar
连续平稳信号 → 汉宁窗

\item {} 
\sphinxAtStartPar
强谐波分量 → 布莱克曼窗

\item {} 
\sphinxAtStartPar
未知信号特征 → 汉明窗

\end{itemize}

\item {} 
\sphinxAtStartPar
\sphinxstylestrong{根据频率范围选择}
\begin{itemize}
\item {} 
\sphinxAtStartPar
高频段（>1000Hz）→ 矩形窗、汉宁窗

\item {} 
\sphinxAtStartPar
中频段（1\sphinxhyphen{}1000Hz）→ 汉宁窗、布莱克曼窗

\item {} 
\sphinxAtStartPar
低频段（<10Hz）→ 高斯窗

\end{itemize}

\item {} 
\sphinxAtStartPar
\sphinxstylestrong{泄漏考虑}
\begin{itemize}
\item {} 
\sphinxAtStartPar
频率分辨率要求高 → 矩形窗、高斯窗

\item {} 
\sphinxAtStartPar
需要良好幅度响应 → 布莱克曼窗

\item {} 
\sphinxAtStartPar
调制频谱分析 → 汉明窗、布莱克曼窗

\end{itemize}

\end{enumerate}


\subsubsection{4.5.3 实用设置建议}
\label{\detokenize{chapters/chapter5:id17}}
\sphinxAtStartPar
\sphinxstylestrong{典型场景配置：}


\begin{savenotes}\sphinxattablestart
\sphinxthistablewithglobalstyle
\centering
\begin{tabulary}{\linewidth}[t]{TTTTT}
\sphinxtoprule
\begin{varwidth}[t]{\sphinxcolwidth{1}{5}}
\sphinxstyletheadfamily \sphinxAtStartPar
应用场景
\sphinxbeforeendvarwidth
\end{varwidth}%
&\begin{varwidth}[t]{\sphinxcolwidth{1}{5}}
\sphinxstyletheadfamily \sphinxAtStartPar
推荐窗函数
\sphinxbeforeendvarwidth
\end{varwidth}%
&\begin{varwidth}[t]{\sphinxcolwidth{1}{5}}
\sphinxstyletheadfamily \sphinxAtStartPar
FFT长度
\sphinxbeforeendvarwidth
\end{varwidth}%
&\begin{varwidth}[t]{\sphinxcolwidth{1}{5}}
\sphinxstyletheadfamily \sphinxAtStartPar
重叠率
\sphinxbeforeendvarwidth
\end{varwidth}%
&\begin{varwidth}[t]{\sphinxcolwidth{1}{5}}
\sphinxstyletheadfamily \sphinxAtStartPar
说明
\sphinxbeforeendvarwidth
\end{varwidth}%
\\
\sphinxmidrule
\sphinxtableatstartofbodyhook\begin{varwidth}[t]{\sphinxcolwidth{1}{5}}
\sphinxAtStartPar
Phoenix高频（TS2）
\sphinxbeforeendvarwidth
\end{varwidth}%
&\begin{varwidth}[t]{\sphinxcolwidth{1}{5}}
\sphinxAtStartPar
矩形窗
\sphinxbeforeendvarwidth
\end{varwidth}%
&\begin{varwidth}[t]{\sphinxcolwidth{1}{5}}
\sphinxAtStartPar
24000\sphinxhyphen{}48000
\sphinxbeforeendvarwidth
\end{varwidth}%
&\begin{varwidth}[t]{\sphinxcolwidth{1}{5}}
\sphinxAtStartPar
50\%
\sphinxbeforeendvarwidth
\end{varwidth}%
&\begin{varwidth}[t]{\sphinxcolwidth{1}{5}}
\sphinxAtStartPar
高频瞬态分析
\sphinxbeforeendvarwidth
\end{varwidth}%
\\
\sphinxhline\begin{varwidth}[t]{\sphinxcolwidth{1}{5}}
\sphinxAtStartPar
Phoenix中频（TS3/TS4）
\sphinxbeforeendvarwidth
\end{varwidth}%
&\begin{varwidth}[t]{\sphinxcolwidth{1}{5}}
\sphinxAtStartPar
汉宁窗
\sphinxbeforeendvarwidth
\end{varwidth}%
&\begin{varwidth}[t]{\sphinxcolwidth{1}{5}}
\sphinxAtStartPar
12000\sphinxhyphen{}24000
\sphinxbeforeendvarwidth
\end{varwidth}%
&\begin{varwidth}[t]{\sphinxcolwidth{1}{5}}
\sphinxAtStartPar
50\%
\sphinxbeforeendvarwidth
\end{varwidth}%
&\begin{varwidth}[t]{\sphinxcolwidth{1}{5}}
\sphinxAtStartPar
平衡时频分辨率
\sphinxbeforeendvarwidth
\end{varwidth}%
\\
\sphinxhline\begin{varwidth}[t]{\sphinxcolwidth{1}{5}}
\sphinxAtStartPar
Phoenix低频（TS5）
\sphinxbeforeendvarwidth
\end{varwidth}%
&\begin{varwidth}[t]{\sphinxcolwidth{1}{5}}
\sphinxAtStartPar
布莱克曼窗
\sphinxbeforeendvarwidth
\end{varwidth}%
&\begin{varwidth}[t]{\sphinxcolwidth{1}{5}}
\sphinxAtStartPar
4800\sphinxhyphen{}24000
\sphinxbeforeendvarwidth
\end{varwidth}%
&\begin{varwidth}[t]{\sphinxcolwidth{1}{5}}
\sphinxAtStartPar
75\%
\sphinxbeforeendvarwidth
\end{varwidth}%
&\begin{varwidth}[t]{\sphinxcolwidth{1}{5}}
\sphinxAtStartPar
低频宽频响应
\sphinxbeforeendvarwidth
\end{varwidth}%
\\
\sphinxhline\begin{varwidth}[t]{\sphinxcolwidth{1}{5}}
\sphinxAtStartPar
MTU\sphinxhyphen{}5A宽频
\sphinxbeforeendvarwidth
\end{varwidth}%
&\begin{varwidth}[t]{\sphinxcolwidth{1}{5}}
\sphinxAtStartPar
高斯窗
\sphinxbeforeendvarwidth
\end{varwidth}%
&\begin{varwidth}[t]{\sphinxcolwidth{1}{5}}
\sphinxAtStartPar
24000\sphinxhyphen{}48000
\sphinxbeforeendvarwidth
\end{varwidth}%
&\begin{varwidth}[t]{\sphinxcolwidth{1}{5}}
\sphinxAtStartPar
50\%
\sphinxbeforeendvarwidth
\end{varwidth}%
&\begin{varwidth}[t]{\sphinxcolwidth{1}{5}}
\sphinxAtStartPar
瞬态信号平滑
\sphinxbeforeendvarwidth
\end{varwidth}%
\\
\sphinxhline\begin{varwidth}[t]{\sphinxcolwidth{1}{5}}
\sphinxAtStartPar
远参考分析
\sphinxbeforeendvarwidth
\end{varwidth}%
&\begin{varwidth}[t]{\sphinxcolwidth{1}{5}}
\sphinxAtStartPar
汉宁窗
\sphinxbeforeendvarwidth
\end{varwidth}%
&\begin{varwidth}[t]{\sphinxcolwidth{1}{5}}
\sphinxAtStartPar
9600\sphinxhyphen{}19200
\sphinxbeforeendvarwidth
\end{varwidth}%
&\begin{varwidth}[t]{\sphinxcolwidth{1}{5}}
\sphinxAtStartPar
75\%
\sphinxbeforeendvarwidth
\end{varwidth}%
&\begin{varwidth}[t]{\sphinxcolwidth{1}{5}}
\sphinxAtStartPar
相干性分析优化
\sphinxbeforeendvarwidth
\end{varwidth}%
\\
\sphinxbottomrule
\end{tabulary}
\sphinxtableafterendhook\par
\sphinxattableend\end{savenotes}


\subsubsection{4.5.4 窗函数对比示例}
\label{\detokenize{chapters/chapter5:id18}}
\sphinxAtStartPar
\sphinxstylestrong{示例：对比矩形窗和汉宁窗}

\begin{sphinxVerbatim}[commandchars=\\\{\}]
信号：1Hz余弦波，采样率2400Hz

矩形窗（截断）：
\PYGZhy{} 频谱：出现明显旁瓣（虚假频率）
\PYGZhy{} 主瓣分辨率：4π/N
\PYGZhy{} 适用：快速变化信号分析

汉宁窗（渐变）：
\PYGZhy{} 频谱：无旁瓣，更干净
\PYGZhy{} 主瓣分辨率：8π/N（较宽）
\PYGZhy{} 适用：一般MT数据分析
\end{sphinxVerbatim}


\subsubsection{4.5.5 单频计算方式}
\label{\detokenize{chapters/chapter5:id19}}

\begin{savenotes}\sphinxattablestart
\sphinxthistablewithglobalstyle
\centering
\begin{tabulary}{\linewidth}[t]{TTT}
\sphinxtoprule
\begin{varwidth}[t]{\sphinxcolwidth{1}{3}}
\sphinxstyletheadfamily \sphinxAtStartPar
值
\sphinxbeforeendvarwidth
\end{varwidth}%
&\begin{varwidth}[t]{\sphinxcolwidth{1}{3}}
\sphinxstyletheadfamily \sphinxAtStartPar
方式
\sphinxbeforeendvarwidth
\end{varwidth}%
&\begin{varwidth}[t]{\sphinxcolwidth{1}{3}}
\sphinxstyletheadfamily \sphinxAtStartPar
说明
\sphinxbeforeendvarwidth
\end{varwidth}%
\\
\sphinxmidrule
\sphinxtableatstartofbodyhook\begin{varwidth}[t]{\sphinxcolwidth{1}{3}}
\sphinxAtStartPar
0
\sphinxbeforeendvarwidth
\end{varwidth}%
&\begin{varwidth}[t]{\sphinxcolwidth{1}{3}}
\sphinxAtStartPar
指定频率
\sphinxbeforeendvarwidth
\end{varwidth}%
&\begin{varwidth}[t]{\sphinxcolwidth{1}{3}}
\sphinxAtStartPar
使用指定的中心频率
\sphinxbeforeendvarwidth
\end{varwidth}%
\\
\sphinxhline\begin{varwidth}[t]{\sphinxcolwidth{1}{3}}
\sphinxAtStartPar
1
\sphinxbeforeendvarwidth
\end{varwidth}%
&\begin{varwidth}[t]{\sphinxcolwidth{1}{3}}
\sphinxAtStartPar
频带平均
\sphinxbeforeendvarwidth
\end{varwidth}%
&\begin{varwidth}[t]{\sphinxcolwidth{1}{3}}
\sphinxAtStartPar
在指定频带内平均
\sphinxbeforeendvarwidth
\end{varwidth}%
\\
\sphinxhline\begin{varwidth}[t]{\sphinxcolwidth{1}{3}}
\sphinxAtStartPar
2
\sphinxbeforeendvarwidth
\end{varwidth}%
&\begin{varwidth}[t]{\sphinxcolwidth{1}{3}}
\sphinxAtStartPar
峰值搜索
\sphinxbeforeendvarwidth
\end{varwidth}%
&\begin{varwidth}[t]{\sphinxcolwidth{1}{3}}
\sphinxAtStartPar
搜索频带内峰值
\sphinxbeforeendvarwidth
\end{varwidth}%
\\
\sphinxbottomrule
\end{tabulary}
\sphinxtableafterendhook\par
\sphinxattableend\end{savenotes}

\sphinxAtStartPar
\sphinxstylestrong{单频计算范围：} 设置频带宽度（对数单位）


\subsubsection{4.5.6 多锥谱分析（Multi\sphinxhyphen{}Taper）}
\label{\detokenize{chapters/chapter5:multi-taper}}

\begin{savenotes}\sphinxattablestart
\sphinxthistablewithglobalstyle
\centering
\begin{tabulary}{\linewidth}[t]{TT}
\sphinxtoprule
\begin{varwidth}[t]{\sphinxcolwidth{1}{2}}
\sphinxstyletheadfamily \sphinxAtStartPar
参数
\sphinxbeforeendvarwidth
\end{varwidth}%
&\begin{varwidth}[t]{\sphinxcolwidth{1}{2}}
\sphinxstyletheadfamily \sphinxAtStartPar
说明
\sphinxbeforeendvarwidth
\end{varwidth}%
\\
\sphinxmidrule
\sphinxtableatstartofbodyhook\begin{varwidth}[t]{\sphinxcolwidth{1}{2}}
\sphinxAtStartPar
MultiTaper
\sphinxbeforeendvarwidth
\end{varwidth}%
&\begin{varwidth}[t]{\sphinxcolwidth{1}{2}}
\sphinxAtStartPar
锥数量，0表示禁用
\sphinxbeforeendvarwidth
\end{varwidth}%
\\
\sphinxbottomrule
\end{tabulary}
\sphinxtableafterendhook\par
\sphinxattableend\end{savenotes}

\sphinxAtStartPar
多锥谱分析可减少频谱泄漏，提高估计精度。


\subsubsection{4.5.7 零填充（Zero Padding）}
\label{\detokenize{chapters/chapter5:zero-padding}}

\begin{savenotes}\sphinxattablestart
\sphinxthistablewithglobalstyle
\centering
\begin{tabulary}{\linewidth}[t]{TT}
\sphinxtoprule
\begin{varwidth}[t]{\sphinxcolwidth{1}{2}}
\sphinxstyletheadfamily \sphinxAtStartPar
参数
\sphinxbeforeendvarwidth
\end{varwidth}%
&\begin{varwidth}[t]{\sphinxcolwidth{1}{2}}
\sphinxstyletheadfamily \sphinxAtStartPar
说明
\sphinxbeforeendvarwidth
\end{varwidth}%
\\
\sphinxmidrule
\sphinxtableatstartofbodyhook\begin{varwidth}[t]{\sphinxcolwidth{1}{2}}
\sphinxAtStartPar
ZeroPadding
\sphinxbeforeendvarwidth
\end{varwidth}%
&\begin{varwidth}[t]{\sphinxcolwidth{1}{2}}
\sphinxAtStartPar
零填充因子
\sphinxbeforeendvarwidth
\end{varwidth}%
\\
\sphinxbottomrule
\end{tabulary}
\sphinxtableafterendhook\par
\sphinxattableend\end{savenotes}

\sphinxAtStartPar
零填充可提高频率分辨率显示，但不会增加实际信息。


\subsubsection{4.5.8 搜索带宽（Search Band）}
\label{\detokenize{chapters/chapter5:search-band}}

\begin{savenotes}\sphinxattablestart
\sphinxthistablewithglobalstyle
\centering
\begin{tabulary}{\linewidth}[t]{TT}
\sphinxtoprule
\begin{varwidth}[t]{\sphinxcolwidth{1}{2}}
\sphinxstyletheadfamily \sphinxAtStartPar
参数
\sphinxbeforeendvarwidth
\end{varwidth}%
&\begin{varwidth}[t]{\sphinxcolwidth{1}{2}}
\sphinxstyletheadfamily \sphinxAtStartPar
说明
\sphinxbeforeendvarwidth
\end{varwidth}%
\\
\sphinxmidrule
\sphinxtableatstartofbodyhook\begin{varwidth}[t]{\sphinxcolwidth{1}{2}}
\sphinxAtStartPar
SearchBand
\sphinxbeforeendvarwidth
\end{varwidth}%
&\begin{varwidth}[t]{\sphinxcolwidth{1}{2}}
\sphinxAtStartPar
频率搜索带宽
\sphinxbeforeendvarwidth
\end{varwidth}%
\\
\sphinxbottomrule
\end{tabulary}
\sphinxtableafterendhook\par
\sphinxattableend\end{savenotes}


\bigskip\hrule\bigskip



\subsection{5.6 🎯 估计方法设置}
\label{\detokenize{chapters/chapter5:id20}}

\subsubsection{4.6.1 参考道估计方法}
\label{\detokenize{chapters/chapter5:id21}}

\begin{savenotes}\sphinxattablestart
\sphinxthistablewithglobalstyle
\centering
\begin{tabulary}{\linewidth}[t]{TTTT}
\sphinxtoprule
\begin{varwidth}[t]{\sphinxcolwidth{1}{4}}
\sphinxstyletheadfamily \sphinxAtStartPar
值
\sphinxbeforeendvarwidth
\end{varwidth}%
&\begin{varwidth}[t]{\sphinxcolwidth{1}{4}}
\sphinxstyletheadfamily \sphinxAtStartPar
代码
\sphinxbeforeendvarwidth
\end{varwidth}%
&\begin{varwidth}[t]{\sphinxcolwidth{1}{4}}
\sphinxstyletheadfamily \sphinxAtStartPar
方法
\sphinxbeforeendvarwidth
\end{varwidth}%
&\begin{varwidth}[t]{\sphinxcolwidth{1}{4}}
\sphinxstyletheadfamily \sphinxAtStartPar
说明
\sphinxbeforeendvarwidth
\end{varwidth}%
\\
\sphinxmidrule
\sphinxtableatstartofbodyhook\begin{varwidth}[t]{\sphinxcolwidth{1}{4}}
\sphinxAtStartPar
0
\sphinxbeforeendvarwidth
\end{varwidth}%
&\begin{varwidth}[t]{\sphinxcolwidth{1}{4}}
\sphinxAtStartPar
LE
\sphinxbeforeendvarwidth
\end{varwidth}%
&\begin{varwidth}[t]{\sphinxcolwidth{1}{4}}
\sphinxAtStartPar
Local E
\sphinxbeforeendvarwidth
\end{varwidth}%
&\begin{varwidth}[t]{\sphinxcolwidth{1}{4}}
\sphinxAtStartPar
本地电场参考
\sphinxbeforeendvarwidth
\end{varwidth}%
\\
\sphinxhline\begin{varwidth}[t]{\sphinxcolwidth{1}{4}}
\sphinxAtStartPar
1
\sphinxbeforeendvarwidth
\end{varwidth}%
&\begin{varwidth}[t]{\sphinxcolwidth{1}{4}}
\sphinxAtStartPar
LH
\sphinxbeforeendvarwidth
\end{varwidth}%
&\begin{varwidth}[t]{\sphinxcolwidth{1}{4}}
\sphinxAtStartPar
Local H
\sphinxbeforeendvarwidth
\end{varwidth}%
&\begin{varwidth}[t]{\sphinxcolwidth{1}{4}}
\sphinxAtStartPar
本地磁场参考
\sphinxbeforeendvarwidth
\end{varwidth}%
\\
\sphinxhline\begin{varwidth}[t]{\sphinxcolwidth{1}{4}}
\sphinxAtStartPar
2
\sphinxbeforeendvarwidth
\end{varwidth}%
&\begin{varwidth}[t]{\sphinxcolwidth{1}{4}}
\sphinxAtStartPar
LEH
\sphinxbeforeendvarwidth
\end{varwidth}%
&\begin{varwidth}[t]{\sphinxcolwidth{1}{4}}
\sphinxAtStartPar
Local E/H
\sphinxbeforeendvarwidth
\end{varwidth}%
&\begin{varwidth}[t]{\sphinxcolwidth{1}{4}}
\sphinxAtStartPar
本地电磁场参考
\sphinxbeforeendvarwidth
\end{varwidth}%
\\
\sphinxhline\begin{varwidth}[t]{\sphinxcolwidth{1}{4}}
\sphinxAtStartPar
3
\sphinxbeforeendvarwidth
\end{varwidth}%
&\begin{varwidth}[t]{\sphinxcolwidth{1}{4}}
\sphinxAtStartPar
RE
\sphinxbeforeendvarwidth
\end{varwidth}%
&\begin{varwidth}[t]{\sphinxcolwidth{1}{4}}
\sphinxAtStartPar
Remote E
\sphinxbeforeendvarwidth
\end{varwidth}%
&\begin{varwidth}[t]{\sphinxcolwidth{1}{4}}
\sphinxAtStartPar
远参考电场
\sphinxbeforeendvarwidth
\end{varwidth}%
\\
\sphinxhline\begin{varwidth}[t]{\sphinxcolwidth{1}{4}}
\sphinxAtStartPar
4
\sphinxbeforeendvarwidth
\end{varwidth}%
&\begin{varwidth}[t]{\sphinxcolwidth{1}{4}}
\sphinxAtStartPar
RH
\sphinxbeforeendvarwidth
\end{varwidth}%
&\begin{varwidth}[t]{\sphinxcolwidth{1}{4}}
\sphinxAtStartPar
Remote H
\sphinxbeforeendvarwidth
\end{varwidth}%
&\begin{varwidth}[t]{\sphinxcolwidth{1}{4}}
\sphinxAtStartPar
远参考磁场
\sphinxbeforeendvarwidth
\end{varwidth}%
\\
\sphinxhline\begin{varwidth}[t]{\sphinxcolwidth{1}{4}}
\sphinxAtStartPar
5
\sphinxbeforeendvarwidth
\end{varwidth}%
&\begin{varwidth}[t]{\sphinxcolwidth{1}{4}}
\sphinxAtStartPar
REH
\sphinxbeforeendvarwidth
\end{varwidth}%
&\begin{varwidth}[t]{\sphinxcolwidth{1}{4}}
\sphinxAtStartPar
Remote E/H
\sphinxbeforeendvarwidth
\end{varwidth}%
&\begin{varwidth}[t]{\sphinxcolwidth{1}{4}}
\sphinxAtStartPar
远参考电磁场
\sphinxbeforeendvarwidth
\end{varwidth}%
\\
\sphinxhline\begin{varwidth}[t]{\sphinxcolwidth{1}{4}}
\sphinxAtStartPar
6
\sphinxbeforeendvarwidth
\end{varwidth}%
&\begin{varwidth}[t]{\sphinxcolwidth{1}{4}}
\sphinxAtStartPar
RELH
\sphinxbeforeendvarwidth
\end{varwidth}%
&\begin{varwidth}[t]{\sphinxcolwidth{1}{4}}
\sphinxAtStartPar
Remote E/Local H
\sphinxbeforeendvarwidth
\end{varwidth}%
&\begin{varwidth}[t]{\sphinxcolwidth{1}{4}}
\sphinxAtStartPar
远参考E/本地H
\sphinxbeforeendvarwidth
\end{varwidth}%
\\
\sphinxbottomrule
\end{tabulary}
\sphinxtableafterendhook\par
\sphinxattableend\end{savenotes}
\begin{quote}

\sphinxAtStartPar
💡 \sphinxstylestrong{提示}：使用远参考道可有效消除本地噪声影响，提高数据质量。
\end{quote}


\subsubsection{4.6.2 稳健估计方法}
\label{\detokenize{chapters/chapter5:id22}}
\sphinxAtStartPar
\sphinxstylestrong{简化流程：} 1. \sphinxstylestrong{📥 导入数据} → 2. \sphinxstylestrong{🔍 检查数据质量} → 3. \sphinxstylestrong{⚙️ 设置FFT参数} → 4. \sphinxstylestrong{▶️ 执行FFT} → 5. \sphinxstylestrong{📐 应用标定}


\subsubsection{4.15.1 数据检查要点}
\label{\detokenize{chapters/chapter5:id23}}\begin{itemize}
\item {} 
\sphinxAtStartPar
查看时间序列是否完整

\item {} 
\sphinxAtStartPar
检查是否有明显噪声

\item {} 
\sphinxAtStartPar
确认采样率正确

\end{itemize}


\subsubsection{4.15.2 强干扰环境建议}
\label{\detokenize{chapters/chapter5:id24}}\begin{itemize}
\item {} 
\sphinxAtStartPar
启用工频陷波滤波

\item {} 
\sphinxAtStartPar
使用远参考道（RE/RH/REH）

\item {} 
\sphinxAtStartPar
采用稳健估计方法（Regression\sphinxhyphen{}M或Repeated Median）

\item {} 
\sphinxAtStartPar
启用自动筛选

\end{itemize}


\bigskip\hrule\bigskip



\subsection{5.16 🔬 测点精细处理界面（频谱编辑窗口）}
\label{\detokenize{chapters/chapter5:id25}}\begin{quote}

\sphinxAtStartPar
\sphinxstylestrong{📖 核心参考文献}

\sphinxAtStartPar
王培杰, 陈小斌, 韩鹏, 张赟昀. 基于稳健估计、数据筛选和Rhoplus约束的大地电磁数据处理方法. 地球物理学报, 2024, 67(11): 4325\sphinxhyphen{}4342.

\sphinxAtStartPar
\sphinxstyleemphasis{Strong interference magnetotelluric data processing method based on robust estimation, data screening and Rhoplus constraint. Chinese Journal of Geophysics, 2024.}
\end{quote}

\sphinxAtStartPar
测点精细处理界面是MTDP中进行单点数据精细处理的核心工具，通过右键测点 → \sphinxcode{\sphinxupquote{频谱编辑}} 打开。该界面完整实现了\sphinxstylestrong{RMSMR方法}——即\sphinxstylestrong{稳健估计(Robust) + 多参数筛选(Multi\sphinxhyphen{}parameter Screening) + 多角度Rhoplus约束(Multi\sphinxhyphen{}angle Rhoplus)}，专门用于处理强干扰环境下的MT数据。


\bigskip\hrule\bigskip



\subsubsection{5.16.0 RMSMR理论基础（为什么需要精细处理？）}
\label{\detokenize{chapters/chapter5:rmsmr}}

\paragraph{为什么MT数据需要精细处理？}
\label{\detokenize{chapters/chapter5:mt}}
\sphinxAtStartPar
天然电磁场信号非常微弱，容易受到各种干扰：


\begin{savenotes}\sphinxattablestart
\sphinxthistablewithglobalstyle
\centering
\begin{tabulary}{\linewidth}[t]{TTT}
\sphinxtoprule
\begin{varwidth}[t]{\sphinxcolwidth{1}{3}}
\sphinxstyletheadfamily \sphinxAtStartPar
干扰类型
\sphinxbeforeendvarwidth
\end{varwidth}%
&\begin{varwidth}[t]{\sphinxcolwidth{1}{3}}
\sphinxstyletheadfamily \sphinxAtStartPar
表现特征
\sphinxbeforeendvarwidth
\end{varwidth}%
&\begin{varwidth}[t]{\sphinxcolwidth{1}{3}}
\sphinxstyletheadfamily \sphinxAtStartPar
影响频段
\sphinxbeforeendvarwidth
\end{varwidth}%
\\
\sphinxmidrule
\sphinxtableatstartofbodyhook\begin{varwidth}[t]{\sphinxcolwidth{1}{3}}
\sphinxAtStartPar
\sphinxstylestrong{人文噪声}
\sphinxbeforeendvarwidth
\end{varwidth}%
&\begin{varwidth}[t]{\sphinxcolwidth{1}{3}}
\sphinxAtStartPar
工频干扰、电气化铁路
\sphinxbeforeendvarwidth
\end{varwidth}%
&\begin{varwidth}[t]{\sphinxcolwidth{1}{3}}
\sphinxAtStartPar
50Hz及其谐波
\sphinxbeforeendvarwidth
\end{varwidth}%
\\
\sphinxhline\begin{varwidth}[t]{\sphinxcolwidth{1}{3}}
\sphinxAtStartPar
\sphinxstylestrong{死频带问题}
\sphinxbeforeendvarwidth
\end{varwidth}%
&\begin{varwidth}[t]{\sphinxcolwidth{1}{3}}
\sphinxAtStartPar
信号天然较弱
\sphinxbeforeendvarwidth
\end{varwidth}%
&\begin{varwidth}[t]{\sphinxcolwidth{1}{3}}
\sphinxAtStartPar
1\sphinxhyphen{}5 kHz (AMT死频带)
\sphinxbeforeendvarwidth
\end{varwidth}%
\\
\sphinxhline\begin{varwidth}[t]{\sphinxcolwidth{1}{3}}
\sphinxAtStartPar
\sphinxstylestrong{相干噪声}
\sphinxbeforeendvarwidth
\end{varwidth}%
&\begin{varwidth}[t]{\sphinxcolwidth{1}{3}}
\sphinxAtStartPar
多频率同时受影响
\sphinxbeforeendvarwidth
\end{varwidth}%
&\begin{varwidth}[t]{\sphinxcolwidth{1}{3}}
\sphinxAtStartPar
宽频带
\sphinxbeforeendvarwidth
\end{varwidth}%
\\
\sphinxhline\begin{varwidth}[t]{\sphinxcolwidth{1}{3}}
\sphinxAtStartPar
\sphinxstylestrong{近场效应}
\sphinxbeforeendvarwidth
\end{varwidth}%
&\begin{varwidth}[t]{\sphinxcolwidth{1}{3}}
\sphinxAtStartPar
阻抗失真、相位异常
\sphinxbeforeendvarwidth
\end{varwidth}%
&\begin{varwidth}[t]{\sphinxcolwidth{1}{3}}
\sphinxAtStartPar
低频
\sphinxbeforeendvarwidth
\end{varwidth}%
\\
\sphinxbottomrule
\end{tabulary}
\sphinxtableafterendhook\par
\sphinxattableend\end{savenotes}


\paragraph{RMSMR方法的三重保障}
\label{\detokenize{chapters/chapter5:id26}}
\sphinxAtStartPar
\sphinxstylestrong{三层保护的物理意义：}


\begin{savenotes}\sphinxattablestart
\sphinxthistablewithglobalstyle
\centering
\begin{tabulary}{\linewidth}[t]{TTTT}
\sphinxtoprule
\begin{varwidth}[t]{\sphinxcolwidth{1}{4}}
\sphinxstyletheadfamily \sphinxAtStartPar
层次
\sphinxbeforeendvarwidth
\end{varwidth}%
&\begin{varwidth}[t]{\sphinxcolwidth{1}{4}}
\sphinxstyletheadfamily \sphinxAtStartPar
方法
\sphinxbeforeendvarwidth
\end{varwidth}%
&\begin{varwidth}[t]{\sphinxcolwidth{1}{4}}
\sphinxstyletheadfamily \sphinxAtStartPar
解决的问题
\sphinxbeforeendvarwidth
\end{varwidth}%
&\begin{varwidth}[t]{\sphinxcolwidth{1}{4}}
\sphinxstyletheadfamily \sphinxAtStartPar
界面位置
\sphinxbeforeendvarwidth
\end{varwidth}%
\\
\sphinxmidrule
\sphinxtableatstartofbodyhook\begin{varwidth}[t]{\sphinxcolwidth{1}{4}}
\sphinxAtStartPar
\sphinxstylestrong{抑噪层}
\sphinxbeforeendvarwidth
\end{varwidth}%
&\begin{varwidth}[t]{\sphinxcolwidth{1}{4}}
\sphinxAtStartPar
稳健估计(Robust)
\sphinxbeforeendvarwidth
\end{varwidth}%
&\begin{varwidth}[t]{\sphinxcolwidth{1}{4}}
\sphinxAtStartPar
降低异常数据权重，不让少数坏点扭曲结果
\sphinxbeforeendvarwidth
\end{varwidth}%
&\begin{varwidth}[t]{\sphinxcolwidth{1}{4}}
\sphinxAtStartPar
Robust方案下拉框
\sphinxbeforeendvarwidth
\end{varwidth}%
\\
\sphinxhline\begin{varwidth}[t]{\sphinxcolwidth{1}{4}}
\sphinxAtStartPar
\sphinxstylestrong{去噪层}
\sphinxbeforeendvarwidth
\end{varwidth}%
&\begin{varwidth}[t]{\sphinxcolwidth{1}{4}}
\sphinxAtStartPar
多参数筛选(Screening)
\sphinxbeforeendvarwidth
\end{varwidth}%
&\begin{varwidth}[t]{\sphinxcolwidth{1}{4}}
\sphinxAtStartPar
识别并剔除受污染的数据段
\sphinxbeforeendvarwidth
\end{varwidth}%
&\begin{varwidth}[t]{\sphinxcolwidth{1}{4}}
\sphinxAtStartPar
自动筛选选项卡
\sphinxbeforeendvarwidth
\end{varwidth}%
\\
\sphinxhline\begin{varwidth}[t]{\sphinxcolwidth{1}{4}}
\sphinxAtStartPar
\sphinxstylestrong{约束层}
\sphinxbeforeendvarwidth
\end{varwidth}%
&\begin{varwidth}[t]{\sphinxcolwidth{1}{4}}
\sphinxAtStartPar
多角度Rhoplus(Rhoplus)
\sphinxbeforeendvarwidth
\end{varwidth}%
&\begin{varwidth}[t]{\sphinxcolwidth{1}{4}}
\sphinxAtStartPar
用高质量频段约束低质量频段
\sphinxbeforeendvarwidth
\end{varwidth}%
&\begin{varwidth}[t]{\sphinxcolwidth{1}{4}}
\sphinxAtStartPar
RhoPlus GA/MOGA按钮
\sphinxbeforeendvarwidth
\end{varwidth}%
\\
\sphinxbottomrule
\end{tabulary}
\sphinxtableafterendhook\par
\sphinxattableend\end{savenotes}


\paragraph{从理论到实践的映射}
\label{\detokenize{chapters/chapter5:id27}}

\begin{savenotes}\sphinxattablestart
\sphinxthistablewithglobalstyle
\centering
\begin{tabulary}{\linewidth}[t]{TTT}
\sphinxtoprule
\begin{varwidth}[t]{\sphinxcolwidth{1}{3}}
\sphinxstyletheadfamily \sphinxAtStartPar
RMSMR理论步骤
\sphinxbeforeendvarwidth
\end{varwidth}%
&\begin{varwidth}[t]{\sphinxcolwidth{1}{3}}
\sphinxstyletheadfamily \sphinxAtStartPar
界面操作
\sphinxbeforeendvarwidth
\end{varwidth}%
&\begin{varwidth}[t]{\sphinxcolwidth{1}{3}}
\sphinxstyletheadfamily \sphinxAtStartPar
参数建议
\sphinxbeforeendvarwidth
\end{varwidth}%
\\
\sphinxmidrule
\sphinxtableatstartofbodyhook\begin{varwidth}[t]{\sphinxcolwidth{1}{3}}
\sphinxAtStartPar
1. 初步筛选
\sphinxbeforeendvarwidth
\end{varwidth}%
&\begin{varwidth}[t]{\sphinxcolwidth{1}{3}}
\sphinxAtStartPar
查看相干度选项卡，框选剔除低相干度数据
\sphinxbeforeendvarwidth
\end{varwidth}%
&\begin{varwidth}[t]{\sphinxcolwidth{1}{3}}
\sphinxAtStartPar
相干度阈值 > 0.7
\sphinxbeforeendvarwidth
\end{varwidth}%
\\
\sphinxhline\begin{varwidth}[t]{\sphinxcolwidth{1}{3}}
\sphinxAtStartPar
2. 稳健估计
\sphinxbeforeendvarwidth
\end{varwidth}%
&\begin{varwidth}[t]{\sphinxcolwidth{1}{3}}
\sphinxAtStartPar
选择Robust方法为"M估计"
\sphinxbeforeendvarwidth
\end{varwidth}%
&\begin{varwidth}[t]{\sphinxcolwidth{1}{3}}
\sphinxAtStartPar
推荐Regression\sphinxhyphen{}M
\sphinxbeforeendvarwidth
\end{varwidth}%
\\
\sphinxhline\begin{varwidth}[t]{\sphinxcolwidth{1}{3}}
\sphinxAtStartPar
3. 残差分析
\sphinxbeforeendvarwidth
\end{varwidth}%
&\begin{varwidth}[t]{\sphinxcolwidth{1}{3}}
\sphinxAtStartPar
设置残差剔除参数
\sphinxbeforeendvarwidth
\end{varwidth}%
&\begin{varwidth}[t]{\sphinxcolwidth{1}{3}}
\sphinxAtStartPar
残差阈值 1.5\sphinxhyphen{}2.5
\sphinxbeforeendvarwidth
\end{varwidth}%
\\
\sphinxhline\begin{varwidth}[t]{\sphinxcolwidth{1}{3}}
\sphinxAtStartPar
4. 多参数筛选
\sphinxbeforeendvarwidth
\end{varwidth}%
&\begin{varwidth}[t]{\sphinxcolwidth{1}{3}}
\sphinxAtStartPar
使用马氏距离自动筛选
\sphinxbeforeendvarwidth
\end{varwidth}%
&\begin{varwidth}[t]{\sphinxcolwidth{1}{3}}
\sphinxAtStartPar
阈值 2.0\sphinxhyphen{}3.0
\sphinxbeforeendvarwidth
\end{varwidth}%
\\
\sphinxhline\begin{varwidth}[t]{\sphinxcolwidth{1}{3}}
\sphinxAtStartPar
5. Rhoplus约束
\sphinxbeforeendvarwidth
\end{varwidth}%
&\begin{varwidth}[t]{\sphinxcolwidth{1}{3}}
\sphinxAtStartPar
执行RhoPlus MOGA筛选
\sphinxbeforeendvarwidth
\end{varwidth}%
&\begin{varwidth}[t]{\sphinxcolwidth{1}{3}}
\sphinxAtStartPar
保留比例 30\%\sphinxhyphen{}50\%
\sphinxbeforeendvarwidth
\end{varwidth}%
\\
\sphinxhline\begin{varwidth}[t]{\sphinxcolwidth{1}{3}}
\sphinxAtStartPar
6. 迭代优化
\sphinxbeforeendvarwidth
\end{varwidth}%
&\begin{varwidth}[t]{\sphinxcolwidth{1}{3}}
\sphinxAtStartPar
重复2\sphinxhyphen{}5直到结果稳定
\sphinxbeforeendvarwidth
\end{varwidth}%
&\begin{varwidth}[t]{\sphinxcolwidth{1}{3}}
\sphinxAtStartPar
迭代3\sphinxhyphen{}10次
\sphinxbeforeendvarwidth
\end{varwidth}%
\\
\sphinxbottomrule
\end{tabulary}
\sphinxtableafterendhook\par
\sphinxattableend\end{savenotes}


\bigskip\hrule\bigskip



\subsubsection{5.16.1 界面布局概述}
\label{\detokenize{chapters/chapter5:id28}}
\sphinxAtStartPar
测点精细处理界面是MTDP中进行单点数据精细处理的核心工具，通过右键测点 → \sphinxcode{\sphinxupquote{频谱编辑}} 打开。该界面提供完整的频谱数据查看、筛选、阻抗估计和数据导出功能。


\bigskip\hrule\bigskip



\subsubsection{5.16.1 界面布局概述}
\label{\detokenize{chapters/chapter5:id29}}
\sphinxAtStartPar
\sphinxstylestrong{主要区域说明：}


\begin{savenotes}\sphinxattablestart
\sphinxthistablewithglobalstyle
\centering
\begin{tabulary}{\linewidth}[t]{TTT}
\sphinxtoprule
\begin{varwidth}[t]{\sphinxcolwidth{1}{3}}
\sphinxstyletheadfamily \sphinxAtStartPar
区域
\sphinxbeforeendvarwidth
\end{varwidth}%
&\begin{varwidth}[t]{\sphinxcolwidth{1}{3}}
\sphinxstyletheadfamily \sphinxAtStartPar
内容
\sphinxbeforeendvarwidth
\end{varwidth}%
&\begin{varwidth}[t]{\sphinxcolwidth{1}{3}}
\sphinxstyletheadfamily \sphinxAtStartPar
功能
\sphinxbeforeendvarwidth
\end{varwidth}%
\\
\sphinxmidrule
\sphinxtableatstartofbodyhook\begin{varwidth}[t]{\sphinxcolwidth{1}{3}}
\sphinxAtStartPar
\sphinxstylestrong{工具栏}
\sphinxbeforeendvarwidth
\end{varwidth}%
&\begin{varwidth}[t]{\sphinxcolwidth{1}{3}}
\sphinxAtStartPar
保存、导出、处理按钮
\sphinxbeforeendvarwidth
\end{varwidth}%
&\begin{varwidth}[t]{\sphinxcolwidth{1}{3}}
\sphinxAtStartPar
快捷操作入口
\sphinxbeforeendvarwidth
\end{varwidth}%
\\
\sphinxhline\begin{varwidth}[t]{\sphinxcolwidth{1}{3}}
\sphinxAtStartPar
\sphinxstylestrong{FC列表}
\sphinxbeforeendvarwidth
\end{varwidth}%
&\begin{varwidth}[t]{\sphinxcolwidth{1}{3}}
\sphinxAtStartPar
傅里叶系数文件列表
\sphinxbeforeendvarwidth
\end{varwidth}%
&\begin{varwidth}[t]{\sphinxcolwidth{1}{3}}
\sphinxAtStartPar
管理和选择FC文件
\sphinxbeforeendvarwidth
\end{varwidth}%
\\
\sphinxhline\begin{varwidth}[t]{\sphinxcolwidth{1}{3}}
\sphinxAtStartPar
\sphinxstylestrong{编辑FC列表}
\sphinxbeforeendvarwidth
\end{varwidth}%
&\begin{varwidth}[t]{\sphinxcolwidth{1}{3}}
\sphinxAtStartPar
已编辑的FC版本列表
\sphinxbeforeendvarwidth
\end{varwidth}%
&\begin{varwidth}[t]{\sphinxcolwidth{1}{3}}
\sphinxAtStartPar
管理多个处理版本
\sphinxbeforeendvarwidth
\end{varwidth}%
\\
\sphinxhline\begin{varwidth}[t]{\sphinxcolwidth{1}{3}}
\sphinxAtStartPar
\sphinxstylestrong{远参考站列表}
\sphinxbeforeendvarwidth
\end{varwidth}%
&\begin{varwidth}[t]{\sphinxcolwidth{1}{3}}
\sphinxAtStartPar
可用的远参考站
\sphinxbeforeendvarwidth
\end{varwidth}%
&\begin{varwidth}[t]{\sphinxcolwidth{1}{3}}
\sphinxAtStartPar
选择参考站进行远参考处理
\sphinxbeforeendvarwidth
\end{varwidth}%
\\
\sphinxhline\begin{varwidth}[t]{\sphinxcolwidth{1}{3}}
\sphinxAtStartPar
\sphinxstylestrong{反向参考站列表}
\sphinxbeforeendvarwidth
\end{varwidth}%
&\begin{varwidth}[t]{\sphinxcolwidth{1}{3}}
\sphinxAtStartPar
反向参考站
\sphinxbeforeendvarwidth
\end{varwidth}%
&\begin{varwidth}[t]{\sphinxcolwidth{1}{3}}
\sphinxAtStartPar
用于反向参考估计
\sphinxbeforeendvarwidth
\end{varwidth}%
\\
\sphinxhline\begin{varwidth}[t]{\sphinxcolwidth{1}{3}}
\sphinxAtStartPar
\sphinxstylestrong{频谱组编辑框架}
\sphinxbeforeendvarwidth
\end{varwidth}%
&\begin{varwidth}[t]{\sphinxcolwidth{1}{3}}
\sphinxAtStartPar
多选项卡编辑区
\sphinxbeforeendvarwidth
\end{varwidth}%
&\begin{varwidth}[t]{\sphinxcolwidth{1}{3}}
\sphinxAtStartPar
数据查看和筛选
\sphinxbeforeendvarwidth
\end{varwidth}%
\\
\sphinxbottomrule
\end{tabulary}
\sphinxtableafterendhook\par
\sphinxattableend\end{savenotes}


\bigskip\hrule\bigskip



\subsubsection{5.16.2 FC文件管理}
\label{\detokenize{chapters/chapter5:fc}}

\paragraph{傅里叶系数（FC）文件说明}
\label{\detokenize{chapters/chapter5:id30}}
\sphinxAtStartPar
FC文件存储FFT处理后得到的频谱数据，是后续所有处理的基础。


\begin{savenotes}\sphinxattablestart
\sphinxthistablewithglobalstyle
\centering
\begin{tabulary}{\linewidth}[t]{TTT}
\sphinxtoprule
\begin{varwidth}[t]{\sphinxcolwidth{1}{3}}
\sphinxstyletheadfamily \sphinxAtStartPar
文件类型
\sphinxbeforeendvarwidth
\end{varwidth}%
&\begin{varwidth}[t]{\sphinxcolwidth{1}{3}}
\sphinxstyletheadfamily \sphinxAtStartPar
扩展名
\sphinxbeforeendvarwidth
\end{varwidth}%
&\begin{varwidth}[t]{\sphinxcolwidth{1}{3}}
\sphinxstyletheadfamily \sphinxAtStartPar
说明
\sphinxbeforeendvarwidth
\end{varwidth}%
\\
\sphinxmidrule
\sphinxtableatstartofbodyhook\begin{varwidth}[t]{\sphinxcolwidth{1}{3}}
\sphinxAtStartPar
\sphinxstylestrong{原始FC}
\sphinxbeforeendvarwidth
\end{varwidth}%
&\begin{varwidth}[t]{\sphinxcolwidth{1}{3}}
\sphinxAtStartPar
.STFC
\sphinxbeforeendvarwidth
\end{varwidth}%
&\begin{varwidth}[t]{\sphinxcolwidth{1}{3}}
\sphinxAtStartPar
FFT处理后的原始频谱数据
\sphinxbeforeendvarwidth
\end{varwidth}%
\\
\sphinxhline\begin{varwidth}[t]{\sphinxcolwidth{1}{3}}
\sphinxAtStartPar
\sphinxstylestrong{编辑FC}
\sphinxbeforeendvarwidth
\end{varwidth}%
&\begin{varwidth}[t]{\sphinxcolwidth{1}{3}}
\sphinxAtStartPar
.STFCG
\sphinxbeforeendvarwidth
\end{varwidth}%
&\begin{varwidth}[t]{\sphinxcolwidth{1}{3}}
\sphinxAtStartPar
经过筛选和处理的FC数据
\sphinxbeforeendvarwidth
\end{varwidth}%
\\
\sphinxbottomrule
\end{tabulary}
\sphinxtableafterendhook\par
\sphinxattableend\end{savenotes}


\paragraph{FC列表操作}
\label{\detokenize{chapters/chapter5:id31}}
\sphinxAtStartPar
右键FC列表项可执行以下操作：


\begin{savenotes}\sphinxattablestart
\sphinxthistablewithglobalstyle
\centering
\begin{tabulary}{\linewidth}[t]{TTT}
\sphinxtoprule
\begin{varwidth}[t]{\sphinxcolwidth{1}{3}}
\sphinxstyletheadfamily \sphinxAtStartPar
操作
\sphinxbeforeendvarwidth
\end{varwidth}%
&\begin{varwidth}[t]{\sphinxcolwidth{1}{3}}
\sphinxstyletheadfamily \sphinxAtStartPar
功能
\sphinxbeforeendvarwidth
\end{varwidth}%
&\begin{varwidth}[t]{\sphinxcolwidth{1}{3}}
\sphinxstyletheadfamily \sphinxAtStartPar
说明
\sphinxbeforeendvarwidth
\end{varwidth}%
\\
\sphinxmidrule
\sphinxtableatstartofbodyhook\begin{varwidth}[t]{\sphinxcolwidth{1}{3}}
\sphinxAtStartPar
\sphinxstylestrong{设为当前}
\sphinxbeforeendvarwidth
\end{varwidth}%
&\begin{varwidth}[t]{\sphinxcolwidth{1}{3}}
\sphinxAtStartPar
激活选中的FC
\sphinxbeforeendvarwidth
\end{varwidth}%
&\begin{varwidth}[t]{\sphinxcolwidth{1}{3}}
\sphinxAtStartPar
将选中的FC设为当前处理对象
\sphinxbeforeendvarwidth
\end{varwidth}%
\\
\sphinxhline\begin{varwidth}[t]{\sphinxcolwidth{1}{3}}
\sphinxAtStartPar
\sphinxstylestrong{删除}
\sphinxbeforeendvarwidth
\end{varwidth}%
&\begin{varwidth}[t]{\sphinxcolwidth{1}{3}}
\sphinxAtStartPar
删除FC文件
\sphinxbeforeendvarwidth
\end{varwidth}%
&\begin{varwidth}[t]{\sphinxcolwidth{1}{3}}
\sphinxAtStartPar
永久删除FC文件（不可恢复）
\sphinxbeforeendvarwidth
\end{varwidth}%
\\
\sphinxhline\begin{varwidth}[t]{\sphinxcolwidth{1}{3}}
\sphinxAtStartPar
\sphinxstylestrong{合并}
\sphinxbeforeendvarwidth
\end{varwidth}%
&\begin{varwidth}[t]{\sphinxcolwidth{1}{3}}
\sphinxAtStartPar
合并多个FC
\sphinxbeforeendvarwidth
\end{varwidth}%
&\begin{varwidth}[t]{\sphinxcolwidth{1}{3}}
\sphinxAtStartPar
将多个FC合并为一个
\sphinxbeforeendvarwidth
\end{varwidth}%
\\
\sphinxhline\begin{varwidth}[t]{\sphinxcolwidth{1}{3}}
\sphinxAtStartPar
\sphinxstylestrong{导出JSON}
\sphinxbeforeendvarwidth
\end{varwidth}%
&\begin{varwidth}[t]{\sphinxcolwidth{1}{3}}
\sphinxAtStartPar
导出为JSON格式
\sphinxbeforeendvarwidth
\end{varwidth}%
&\begin{varwidth}[t]{\sphinxcolwidth{1}{3}}
\sphinxAtStartPar
便于程序处理和数据交换
\sphinxbeforeendvarwidth
\end{varwidth}%
\\
\sphinxhline\begin{varwidth}[t]{\sphinxcolwidth{1}{3}}
\sphinxAtStartPar
\sphinxstylestrong{导出振幅谱}
\sphinxbeforeendvarwidth
\end{varwidth}%
&\begin{varwidth}[t]{\sphinxcolwidth{1}{3}}
\sphinxAtStartPar
导出振幅谱数据
\sphinxbeforeendvarwidth
\end{varwidth}%
&\begin{varwidth}[t]{\sphinxcolwidth{1}{3}}
\sphinxAtStartPar
导出各通道振幅谱
\sphinxbeforeendvarwidth
\end{varwidth}%
\\
\sphinxbottomrule
\end{tabulary}
\sphinxtableafterendhook\par
\sphinxattableend\end{savenotes}


\paragraph{编辑FC版本管理}
\label{\detokenize{chapters/chapter5:id32}}
\sphinxAtStartPar
编辑FC列表支持创建多个处理版本，便于对比不同处理方法的效果。

\sphinxAtStartPar
\sphinxstylestrong{操作菜单：}


\begin{savenotes}\sphinxattablestart
\sphinxthistablewithglobalstyle
\centering
\begin{tabulary}{\linewidth}[t]{TTT}
\sphinxtoprule
\begin{varwidth}[t]{\sphinxcolwidth{1}{3}}
\sphinxstyletheadfamily \sphinxAtStartPar
操作
\sphinxbeforeendvarwidth
\end{varwidth}%
&\begin{varwidth}[t]{\sphinxcolwidth{1}{3}}
\sphinxstyletheadfamily \sphinxAtStartPar
功能
\sphinxbeforeendvarwidth
\end{varwidth}%
&\begin{varwidth}[t]{\sphinxcolwidth{1}{3}}
\sphinxstyletheadfamily \sphinxAtStartPar
使用场景
\sphinxbeforeendvarwidth
\end{varwidth}%
\\
\sphinxmidrule
\sphinxtableatstartofbodyhook\begin{varwidth}[t]{\sphinxcolwidth{1}{3}}
\sphinxAtStartPar
\sphinxstylestrong{复制}
\sphinxbeforeendvarwidth
\end{varwidth}%
&\begin{varwidth}[t]{\sphinxcolwidth{1}{3}}
\sphinxAtStartPar
创建FC副本
\sphinxbeforeendvarwidth
\end{varwidth}%
&\begin{varwidth}[t]{\sphinxcolwidth{1}{3}}
\sphinxAtStartPar
保留原始版本，在新版本上处理
\sphinxbeforeendvarwidth
\end{varwidth}%
\\
\sphinxhline\begin{varwidth}[t]{\sphinxcolwidth{1}{3}}
\sphinxAtStartPar
\sphinxstylestrong{重命名}
\sphinxbeforeendvarwidth
\end{varwidth}%
&\begin{varwidth}[t]{\sphinxcolwidth{1}{3}}
\sphinxAtStartPar
修改FC名称
\sphinxbeforeendvarwidth
\end{varwidth}%
&\begin{varwidth}[t]{\sphinxcolwidth{1}{3}}
\sphinxAtStartPar
便于区分不同处理版本
\sphinxbeforeendvarwidth
\end{varwidth}%
\\
\sphinxhline\begin{varwidth}[t]{\sphinxcolwidth{1}{3}}
\sphinxAtStartPar
\sphinxstylestrong{删除}
\sphinxbeforeendvarwidth
\end{varwidth}%
&\begin{varwidth}[t]{\sphinxcolwidth{1}{3}}
\sphinxAtStartPar
删除编辑FC
\sphinxbeforeendvarwidth
\end{varwidth}%
&\begin{varwidth}[t]{\sphinxcolwidth{1}{3}}
\sphinxAtStartPar
清理不需要的版本
\sphinxbeforeendvarwidth
\end{varwidth}%
\\
\sphinxhline\begin{varwidth}[t]{\sphinxcolwidth{1}{3}}
\sphinxAtStartPar
\sphinxstylestrong{设为测点FC}
\sphinxbeforeendvarwidth
\end{varwidth}%
&\begin{varwidth}[t]{\sphinxcolwidth{1}{3}}
\sphinxAtStartPar
应用到测点
\sphinxbeforeendvarwidth
\end{varwidth}%
&\begin{varwidth}[t]{\sphinxcolwidth{1}{3}}
\sphinxAtStartPar
将编辑结果应用到测点
\sphinxbeforeendvarwidth
\end{varwidth}%
\\
\sphinxhline\begin{varwidth}[t]{\sphinxcolwidth{1}{3}}
\sphinxAtStartPar
\sphinxstylestrong{合并选中}
\sphinxbeforeendvarwidth
\end{varwidth}%
&\begin{varwidth}[t]{\sphinxcolwidth{1}{3}}
\sphinxAtStartPar
合并多个FC
\sphinxbeforeendvarwidth
\end{varwidth}%
&\begin{varwidth}[t]{\sphinxcolwidth{1}{3}}
\sphinxAtStartPar
合并不同时段的处理结果
\sphinxbeforeendvarwidth
\end{varwidth}%
\\
\sphinxhline\begin{varwidth}[t]{\sphinxcolwidth{1}{3}}
\sphinxAtStartPar
\sphinxstylestrong{XY合并}
\sphinxbeforeendvarwidth
\end{varwidth}%
&\begin{varwidth}[t]{\sphinxcolwidth{1}{3}}
\sphinxAtStartPar
X/Y站数据合并
\sphinxbeforeendvarwidth
\end{varwidth}%
&\begin{varwidth}[t]{\sphinxcolwidth{1}{3}}
\sphinxAtStartPar
合并不同方向观测的数据
\sphinxbeforeendvarwidth
\end{varwidth}%
\\
\sphinxhline\begin{varwidth}[t]{\sphinxcolwidth{1}{3}}
\sphinxAtStartPar
\sphinxstylestrong{复制剔除}
\sphinxbeforeendvarwidth
\end{varwidth}%
&\begin{varwidth}[t]{\sphinxcolwidth{1}{3}}
\sphinxAtStartPar
复制剔除标记
\sphinxbeforeendvarwidth
\end{varwidth}%
&\begin{varwidth}[t]{\sphinxcolwidth{1}{3}}
\sphinxAtStartPar
在FC间共享剔除信息
\sphinxbeforeendvarwidth
\end{varwidth}%
\\
\sphinxhline\begin{varwidth}[t]{\sphinxcolwidth{1}{3}}
\sphinxAtStartPar
\sphinxstylestrong{粘贴剔除}
\sphinxbeforeendvarwidth
\end{varwidth}%
&\begin{varwidth}[t]{\sphinxcolwidth{1}{3}}
\sphinxAtStartPar
粘贴剔除标记
\sphinxbeforeendvarwidth
\end{varwidth}%
&\begin{varwidth}[t]{\sphinxcolwidth{1}{3}}
\sphinxAtStartPar
应用之前复制的剔除信息
\sphinxbeforeendvarwidth
\end{varwidth}%
\\
\sphinxhline\begin{varwidth}[t]{\sphinxcolwidth{1}{3}}
\sphinxAtStartPar
\sphinxstylestrong{删除剔除数据}
\sphinxbeforeendvarwidth
\end{varwidth}%
&\begin{varwidth}[t]{\sphinxcolwidth{1}{3}}
\sphinxAtStartPar
创建无剔除FC
\sphinxbeforeendvarwidth
\end{varwidth}%
&\begin{varwidth}[t]{\sphinxcolwidth{1}{3}}
\sphinxAtStartPar
创建不含剔除数据的新FC
\sphinxbeforeendvarwidth
\end{varwidth}%
\\
\sphinxbottomrule
\end{tabulary}
\sphinxtableafterendhook\par
\sphinxattableend\end{savenotes}


\bigskip\hrule\bigskip



\subsubsection{5.16.3 远参考站配置}
\label{\detokenize{chapters/chapter5:id33}}

\paragraph{远参考站列表}
\label{\detokenize{chapters/chapter5:id34}}
\sphinxAtStartPar
远参考处理是提高数据质量的重要手段。在频谱编辑窗口中可以管理和选择远参考站。

\sphinxAtStartPar
\sphinxstylestrong{远参考站列表操作：}


\begin{savenotes}\sphinxattablestart
\sphinxthistablewithglobalstyle
\centering
\begin{tabulary}{\linewidth}[t]{TT}
\sphinxtoprule
\begin{varwidth}[t]{\sphinxcolwidth{1}{2}}
\sphinxstyletheadfamily \sphinxAtStartPar
操作
\sphinxbeforeendvarwidth
\end{varwidth}%
&\begin{varwidth}[t]{\sphinxcolwidth{1}{2}}
\sphinxstyletheadfamily \sphinxAtStartPar
功能
\sphinxbeforeendvarwidth
\end{varwidth}%
\\
\sphinxmidrule
\sphinxtableatstartofbodyhook\begin{varwidth}[t]{\sphinxcolwidth{1}{2}}
\sphinxAtStartPar
\sphinxstylestrong{勾选}
\sphinxbeforeendvarwidth
\end{varwidth}%
&\begin{varwidth}[t]{\sphinxcolwidth{1}{2}}
\sphinxAtStartPar
选择参与处理的参考站
\sphinxbeforeendvarwidth
\end{varwidth}%
\\
\sphinxhline\begin{varwidth}[t]{\sphinxcolwidth{1}{2}}
\sphinxAtStartPar
\sphinxstylestrong{删除选中}
\sphinxbeforeendvarwidth
\end{varwidth}%
&\begin{varwidth}[t]{\sphinxcolwidth{1}{2}}
\sphinxAtStartPar
删除已勾选的参考站
\sphinxbeforeendvarwidth
\end{varwidth}%
\\
\sphinxbottomrule
\end{tabulary}
\sphinxtableafterendhook\par
\sphinxattableend\end{savenotes}


\paragraph{远参考方案选择}
\label{\detokenize{chapters/chapter5:id35}}

\begin{savenotes}\sphinxattablestart
\sphinxthistablewithglobalstyle
\centering
\begin{tabulary}{\linewidth}[t]{TTT}
\sphinxtoprule
\begin{varwidth}[t]{\sphinxcolwidth{1}{3}}
\sphinxstyletheadfamily \sphinxAtStartPar
方案
\sphinxbeforeendvarwidth
\end{varwidth}%
&\begin{varwidth}[t]{\sphinxcolwidth{1}{3}}
\sphinxstyletheadfamily \sphinxAtStartPar
代码
\sphinxbeforeendvarwidth
\end{varwidth}%
&\begin{varwidth}[t]{\sphinxcolwidth{1}{3}}
\sphinxstyletheadfamily \sphinxAtStartPar
说明
\sphinxbeforeendvarwidth
\end{varwidth}%
\\
\sphinxmidrule
\sphinxtableatstartofbodyhook\begin{varwidth}[t]{\sphinxcolwidth{1}{3}}
\sphinxAtStartPar
\sphinxstylestrong{Local E}
\sphinxbeforeendvarwidth
\end{varwidth}%
&\begin{varwidth}[t]{\sphinxcolwidth{1}{3}}
\sphinxAtStartPar
LE
\sphinxbeforeendvarwidth
\end{varwidth}%
&\begin{varwidth}[t]{\sphinxcolwidth{1}{3}}
\sphinxAtStartPar
本地电场参考
\sphinxbeforeendvarwidth
\end{varwidth}%
\\
\sphinxhline\begin{varwidth}[t]{\sphinxcolwidth{1}{3}}
\sphinxAtStartPar
\sphinxstylestrong{Local H}
\sphinxbeforeendvarwidth
\end{varwidth}%
&\begin{varwidth}[t]{\sphinxcolwidth{1}{3}}
\sphinxAtStartPar
LH
\sphinxbeforeendvarwidth
\end{varwidth}%
&\begin{varwidth}[t]{\sphinxcolwidth{1}{3}}
\sphinxAtStartPar
本地磁场参考
\sphinxbeforeendvarwidth
\end{varwidth}%
\\
\sphinxhline\begin{varwidth}[t]{\sphinxcolwidth{1}{3}}
\sphinxAtStartPar
\sphinxstylestrong{Local E/H}
\sphinxbeforeendvarwidth
\end{varwidth}%
&\begin{varwidth}[t]{\sphinxcolwidth{1}{3}}
\sphinxAtStartPar
LEH
\sphinxbeforeendvarwidth
\end{varwidth}%
&\begin{varwidth}[t]{\sphinxcolwidth{1}{3}}
\sphinxAtStartPar
本地电磁场参考
\sphinxbeforeendvarwidth
\end{varwidth}%
\\
\sphinxhline\begin{varwidth}[t]{\sphinxcolwidth{1}{3}}
\sphinxAtStartPar
\sphinxstylestrong{Remote E}
\sphinxbeforeendvarwidth
\end{varwidth}%
&\begin{varwidth}[t]{\sphinxcolwidth{1}{3}}
\sphinxAtStartPar
RE
\sphinxbeforeendvarwidth
\end{varwidth}%
&\begin{varwidth}[t]{\sphinxcolwidth{1}{3}}
\sphinxAtStartPar
远参考电场 ⭐
\sphinxbeforeendvarwidth
\end{varwidth}%
\\
\sphinxhline\begin{varwidth}[t]{\sphinxcolwidth{1}{3}}
\sphinxAtStartPar
\sphinxstylestrong{Remote H}
\sphinxbeforeendvarwidth
\end{varwidth}%
&\begin{varwidth}[t]{\sphinxcolwidth{1}{3}}
\sphinxAtStartPar
RH
\sphinxbeforeendvarwidth
\end{varwidth}%
&\begin{varwidth}[t]{\sphinxcolwidth{1}{3}}
\sphinxAtStartPar
远参考磁场 ⭐
\sphinxbeforeendvarwidth
\end{varwidth}%
\\
\sphinxhline\begin{varwidth}[t]{\sphinxcolwidth{1}{3}}
\sphinxAtStartPar
\sphinxstylestrong{Remote E/H}
\sphinxbeforeendvarwidth
\end{varwidth}%
&\begin{varwidth}[t]{\sphinxcolwidth{1}{3}}
\sphinxAtStartPar
REH
\sphinxbeforeendvarwidth
\end{varwidth}%
&\begin{varwidth}[t]{\sphinxcolwidth{1}{3}}
\sphinxAtStartPar
远参考电磁场 ⭐
\sphinxbeforeendvarwidth
\end{varwidth}%
\\
\sphinxhline\begin{varwidth}[t]{\sphinxcolwidth{1}{3}}
\sphinxAtStartPar
\sphinxstylestrong{Remote E/Local H}
\sphinxbeforeendvarwidth
\end{varwidth}%
&\begin{varwidth}[t]{\sphinxcolwidth{1}{3}}
\sphinxAtStartPar
RELH
\sphinxbeforeendvarwidth
\end{varwidth}%
&\begin{varwidth}[t]{\sphinxcolwidth{1}{3}}
\sphinxAtStartPar
远E/本地H
\sphinxbeforeendvarwidth
\end{varwidth}%
\\
\sphinxbottomrule
\end{tabulary}
\sphinxtableafterendhook\par
\sphinxattableend\end{savenotes}
\begin{quote}

\sphinxAtStartPar
💡 \sphinxstylestrong{推荐}：使用 RE（远参考电场）或 RH（远参考磁场）可获得最佳效果。
\end{quote}


\paragraph{反向参考站配置}
\label{\detokenize{chapters/chapter5:id36}}
\sphinxAtStartPar
反向参考用于特殊处理场景，可将本地测站作为参考站使用。


\begin{savenotes}\sphinxattablestart
\sphinxthistablewithglobalstyle
\centering
\begin{tabulary}{\linewidth}[t]{TT}
\sphinxtoprule
\begin{varwidth}[t]{\sphinxcolwidth{1}{2}}
\sphinxstyletheadfamily \sphinxAtStartPar
方案
\sphinxbeforeendvarwidth
\end{varwidth}%
&\begin{varwidth}[t]{\sphinxcolwidth{1}{2}}
\sphinxstyletheadfamily \sphinxAtStartPar
说明
\sphinxbeforeendvarwidth
\end{varwidth}%
\\
\sphinxmidrule
\sphinxtableatstartofbodyhook\begin{varwidth}[t]{\sphinxcolwidth{1}{2}}
\sphinxAtStartPar
\sphinxstylestrong{None}
\sphinxbeforeendvarwidth
\end{varwidth}%
&\begin{varwidth}[t]{\sphinxcolwidth{1}{2}}
\sphinxAtStartPar
不使用反向参考
\sphinxbeforeendvarwidth
\end{varwidth}%
\\
\sphinxhline\begin{varwidth}[t]{\sphinxcolwidth{1}{2}}
\sphinxAtStartPar
\sphinxstylestrong{LE/LH/LEH}
\sphinxbeforeendvarwidth
\end{varwidth}%
&\begin{varwidth}[t]{\sphinxcolwidth{1}{2}}
\sphinxAtStartPar
本地参考方案
\sphinxbeforeendvarwidth
\end{varwidth}%
\\
\sphinxhline\begin{varwidth}[t]{\sphinxcolwidth{1}{2}}
\sphinxAtStartPar
\sphinxstylestrong{RE/RH/REH}
\sphinxbeforeendvarwidth
\end{varwidth}%
&\begin{varwidth}[t]{\sphinxcolwidth{1}{2}}
\sphinxAtStartPar
远参考方案
\sphinxbeforeendvarwidth
\end{varwidth}%
\\
\sphinxbottomrule
\end{tabulary}
\sphinxtableafterendhook\par
\sphinxattableend\end{savenotes}


\paragraph{旋转角度设置}
\label{\detokenize{chapters/chapter5:id37}}

\begin{savenotes}\sphinxattablestart
\sphinxthistablewithglobalstyle
\centering
\begin{tabulary}{\linewidth}[t]{TT}
\sphinxtoprule
\begin{varwidth}[t]{\sphinxcolwidth{1}{2}}
\sphinxstyletheadfamily \sphinxAtStartPar
参数
\sphinxbeforeendvarwidth
\end{varwidth}%
&\begin{varwidth}[t]{\sphinxcolwidth{1}{2}}
\sphinxstyletheadfamily \sphinxAtStartPar
说明
\sphinxbeforeendvarwidth
\end{varwidth}%
\\
\sphinxmidrule
\sphinxtableatstartofbodyhook\begin{varwidth}[t]{\sphinxcolwidth{1}{2}}
\sphinxAtStartPar
\sphinxstylestrong{测站旋转角}
\sphinxbeforeendvarwidth
\end{varwidth}%
&\begin{varwidth}[t]{\sphinxcolwidth{1}{2}}
\sphinxAtStartPar
本地测站的旋转角度（度）
\sphinxbeforeendvarwidth
\end{varwidth}%
\\
\sphinxhline\begin{varwidth}[t]{\sphinxcolwidth{1}{2}}
\sphinxAtStartPar
\sphinxstylestrong{参考站旋转角}
\sphinxbeforeendvarwidth
\end{varwidth}%
&\begin{varwidth}[t]{\sphinxcolwidth{1}{2}}
\sphinxAtStartPar
远参考站的旋转角度（度）
\sphinxbeforeendvarwidth
\end{varwidth}%
\\
\sphinxbottomrule
\end{tabulary}
\sphinxtableafterendhook\par
\sphinxattableend\end{savenotes}


\bigskip\hrule\bigskip



\subsubsection{5.16.4 估计方法与残差剔除}
\label{\detokenize{chapters/chapter5:id38}}

\paragraph{Robust估计方法选择}
\label{\detokenize{chapters/chapter5:robust}}

\begin{savenotes}\sphinxattablestart
\sphinxthistablewithglobalstyle
\centering
\begin{tabulary}{\linewidth}[t]{TTTT}
\sphinxtoprule
\begin{varwidth}[t]{\sphinxcolwidth{1}{4}}
\sphinxstyletheadfamily \sphinxAtStartPar
方法
\sphinxbeforeendvarwidth
\end{varwidth}%
&\begin{varwidth}[t]{\sphinxcolwidth{1}{4}}
\sphinxstyletheadfamily \sphinxAtStartPar
代码
\sphinxbeforeendvarwidth
\end{varwidth}%
&\begin{varwidth}[t]{\sphinxcolwidth{1}{4}}
\sphinxstyletheadfamily \sphinxAtStartPar
说明
\sphinxbeforeendvarwidth
\end{varwidth}%
&\begin{varwidth}[t]{\sphinxcolwidth{1}{4}}
\sphinxstyletheadfamily \sphinxAtStartPar
推荐场景
\sphinxbeforeendvarwidth
\end{varwidth}%
\\
\sphinxmidrule
\sphinxtableatstartofbodyhook\begin{varwidth}[t]{\sphinxcolwidth{1}{4}}
\sphinxAtStartPar
\sphinxstylestrong{最小二乘}
\sphinxbeforeendvarwidth
\end{varwidth}%
&\begin{varwidth}[t]{\sphinxcolwidth{1}{4}}
\sphinxAtStartPar
LS
\sphinxbeforeendvarwidth
\end{varwidth}%
&\begin{varwidth}[t]{\sphinxcolwidth{1}{4}}
\sphinxAtStartPar
标准最小二乘法
\sphinxbeforeendvarwidth
\end{varwidth}%
&\begin{varwidth}[t]{\sphinxcolwidth{1}{4}}
\sphinxAtStartPar
低噪声数据
\sphinxbeforeendvarwidth
\end{varwidth}%
\\
\sphinxhline\begin{varwidth}[t]{\sphinxcolwidth{1}{4}}
\sphinxAtStartPar
\sphinxstylestrong{M估计}
\sphinxbeforeendvarwidth
\end{varwidth}%
&\begin{varwidth}[t]{\sphinxcolwidth{1}{4}}
\sphinxAtStartPar
ME
\sphinxbeforeendvarwidth
\end{varwidth}%
&\begin{varwidth}[t]{\sphinxcolwidth{1}{4}}
\sphinxAtStartPar
回归M估计 ⭐
\sphinxbeforeendvarwidth
\end{varwidth}%
&\begin{varwidth}[t]{\sphinxcolwidth{1}{4}}
\sphinxAtStartPar
一般MT数据
\sphinxbeforeendvarwidth
\end{varwidth}%
\\
\sphinxhline\begin{varwidth}[t]{\sphinxcolwidth{1}{4}}
\sphinxAtStartPar
\sphinxstylestrong{重复中位数}
\sphinxbeforeendvarwidth
\end{varwidth}%
&\begin{varwidth}[t]{\sphinxcolwidth{1}{4}}
\sphinxAtStartPar
RM
\sphinxbeforeendvarwidth
\end{varwidth}%
&\begin{varwidth}[t]{\sphinxcolwidth{1}{4}}
\sphinxAtStartPar
重复中位数法
\sphinxbeforeendvarwidth
\end{varwidth}%
&\begin{varwidth}[t]{\sphinxcolwidth{1}{4}}
\sphinxAtStartPar
强干扰环境
\sphinxbeforeendvarwidth
\end{varwidth}%
\\
\sphinxhline\begin{varwidth}[t]{\sphinxcolwidth{1}{4}}
\sphinxAtStartPar
\sphinxstylestrong{有界影响}
\sphinxbeforeendvarwidth
\end{varwidth}%
&\begin{varwidth}[t]{\sphinxcolwidth{1}{4}}
\sphinxAtStartPar
BI
\sphinxbeforeendvarwidth
\end{varwidth}%
&\begin{varwidth}[t]{\sphinxcolwidth{1}{4}}
\sphinxAtStartPar
有界影响估计
\sphinxbeforeendvarwidth
\end{varwidth}%
&\begin{varwidth}[t]{\sphinxcolwidth{1}{4}}
\sphinxAtStartPar
杠杆点处理
\sphinxbeforeendvarwidth
\end{varwidth}%
\\
\sphinxhline\begin{varwidth}[t]{\sphinxcolwidth{1}{4}}
\sphinxAtStartPar
\sphinxstylestrong{Huber预加权}
\sphinxbeforeendvarwidth
\end{varwidth}%
&\begin{varwidth}[t]{\sphinxcolwidth{1}{4}}
\sphinxAtStartPar
HPW
\sphinxbeforeendvarwidth
\end{varwidth}%
&\begin{varwidth}[t]{\sphinxcolwidth{1}{4}}
\sphinxAtStartPar
Huber预加权
\sphinxbeforeendvarwidth
\end{varwidth}%
&\begin{varwidth}[t]{\sphinxcolwidth{1}{4}}
\sphinxAtStartPar
稳健处理
\sphinxbeforeendvarwidth
\end{varwidth}%
\\
\sphinxhline\begin{varwidth}[t]{\sphinxcolwidth{1}{4}}
\sphinxAtStartPar
\sphinxstylestrong{Thomson预加权}
\sphinxbeforeendvarwidth
\end{varwidth}%
&\begin{varwidth}[t]{\sphinxcolwidth{1}{4}}
\sphinxAtStartPar
TPW
\sphinxbeforeendvarwidth
\end{varwidth}%
&\begin{varwidth}[t]{\sphinxcolwidth{1}{4}}
\sphinxAtStartPar
Thomson预加权
\sphinxbeforeendvarwidth
\end{varwidth}%
&\begin{varwidth}[t]{\sphinxcolwidth{1}{4}}
\sphinxAtStartPar
稳健处理
\sphinxbeforeendvarwidth
\end{varwidth}%
\\
\sphinxhline\begin{varwidth}[t]{\sphinxcolwidth{1}{4}}
\sphinxAtStartPar
\sphinxstylestrong{Robust+AI}
\sphinxbeforeendvarwidth
\end{varwidth}%
&\begin{varwidth}[t]{\sphinxcolwidth{1}{4}}
\sphinxAtStartPar
AI
\sphinxbeforeendvarwidth
\end{varwidth}%
&\begin{varwidth}[t]{\sphinxcolwidth{1}{4}}
\sphinxAtStartPar
稳健估计+AI预测
\sphinxbeforeendvarwidth
\end{varwidth}%
&\begin{varwidth}[t]{\sphinxcolwidth{1}{4}}
\sphinxAtStartPar
复杂噪声
\sphinxbeforeendvarwidth
\end{varwidth}%
\\
\sphinxbottomrule
\end{tabulary}
\sphinxtableafterendhook\par
\sphinxattableend\end{savenotes}


\paragraph{残差剔除参数}
\label{\detokenize{chapters/chapter5:id39}}

\begin{savenotes}\sphinxattablestart
\sphinxthistablewithglobalstyle
\centering
\begin{tabulary}{\linewidth}[t]{TTT}
\sphinxtoprule
\begin{varwidth}[t]{\sphinxcolwidth{1}{3}}
\sphinxstyletheadfamily \sphinxAtStartPar
参数
\sphinxbeforeendvarwidth
\end{varwidth}%
&\begin{varwidth}[t]{\sphinxcolwidth{1}{3}}
\sphinxstyletheadfamily \sphinxAtStartPar
说明
\sphinxbeforeendvarwidth
\end{varwidth}%
&\begin{varwidth}[t]{\sphinxcolwidth{1}{3}}
\sphinxstyletheadfamily \sphinxAtStartPar
推荐值
\sphinxbeforeendvarwidth
\end{varwidth}%
\\
\sphinxmidrule
\sphinxtableatstartofbodyhook\begin{varwidth}[t]{\sphinxcolwidth{1}{3}}
\sphinxAtStartPar
\sphinxstylestrong{残差阈值}
\sphinxbeforeendvarwidth
\end{varwidth}%
&\begin{varwidth}[t]{\sphinxcolwidth{1}{3}}
\sphinxAtStartPar
残差超过此值的数据被降权
\sphinxbeforeendvarwidth
\end{varwidth}%
&\begin{varwidth}[t]{\sphinxcolwidth{1}{3}}
\sphinxAtStartPar
1.5\sphinxhyphen{}2.5
\sphinxbeforeendvarwidth
\end{varwidth}%
\\
\sphinxhline\begin{varwidth}[t]{\sphinxcolwidth{1}{3}}
\sphinxAtStartPar
\sphinxstylestrong{剩余数据比例}
\sphinxbeforeendvarwidth
\end{varwidth}%
&\begin{varwidth}[t]{\sphinxcolwidth{1}{3}}
\sphinxAtStartPar
保留数据的最小比例
\sphinxbeforeendvarwidth
\end{varwidth}%
&\begin{varwidth}[t]{\sphinxcolwidth{1}{3}}
\sphinxAtStartPar
0.3\sphinxhyphen{}0.5
\sphinxbeforeendvarwidth
\end{varwidth}%
\\
\sphinxbottomrule
\end{tabulary}
\sphinxtableafterendhook\par
\sphinxattableend\end{savenotes}

\sphinxAtStartPar
\sphinxstylestrong{添加剔除功能：}


\begin{savenotes}\sphinxattablestart
\sphinxthistablewithglobalstyle
\centering
\begin{tabulary}{\linewidth}[t]{TT}
\sphinxtoprule
\begin{varwidth}[t]{\sphinxcolwidth{1}{2}}
\sphinxstyletheadfamily \sphinxAtStartPar
按钮
\sphinxbeforeendvarwidth
\end{varwidth}%
&\begin{varwidth}[t]{\sphinxcolwidth{1}{2}}
\sphinxstyletheadfamily \sphinxAtStartPar
功能
\sphinxbeforeendvarwidth
\end{varwidth}%
\\
\sphinxmidrule
\sphinxtableatstartofbodyhook\begin{varwidth}[t]{\sphinxcolwidth{1}{2}}
\sphinxAtStartPar
\sphinxstylestrong{添加剔除异常频点}
\sphinxbeforeendvarwidth
\end{varwidth}%
&\begin{varwidth}[t]{\sphinxcolwidth{1}{2}}
\sphinxAtStartPar
将当前频点的异常数据加入剔除列表
\sphinxbeforeendvarwidth
\end{varwidth}%
\\
\sphinxhline\begin{varwidth}[t]{\sphinxcolwidth{1}{2}}
\sphinxAtStartPar
\sphinxstylestrong{根据传输方案剔除}
\sphinxbeforeendvarwidth
\end{varwidth}%
&\begin{varwidth}[t]{\sphinxcolwidth{1}{2}}
\sphinxAtStartPar
按预设传输方案剔除数据
\sphinxbeforeendvarwidth
\end{varwidth}%
\\
\sphinxbottomrule
\end{tabulary}
\sphinxtableafterendhook\par
\sphinxattableend\end{savenotes}


\bigskip\hrule\bigskip



\subsubsection{5.16.5 频谱组编辑框架（选项卡详解）}
\label{\detokenize{chapters/chapter5:id40}}
\sphinxAtStartPar
频谱组编辑框架包含多个选项卡，用于查看和处理不同类型的数据。


\paragraph{选项卡1：视电阻率/相位（RhoPhs）}
\label{\detokenize{chapters/chapter5:rhophs}}
\sphinxAtStartPar
\sphinxstylestrong{显示内容：}
\begin{itemize}
\item {} 
\sphinxAtStartPar
ρxy、ρyx 视电阻率曲线

\item {} 
\sphinxAtStartPar
φxy、φyx 相位曲线

\item {} 
\sphinxAtStartPar
误差棒

\end{itemize}

\sphinxAtStartPar
\sphinxstylestrong{交互操作：}
\begin{itemize}
\item {} 
\sphinxAtStartPar
左键点击：选择频点

\item {} 
\sphinxAtStartPar
框选：选择多个数据点

\item {} 
\sphinxAtStartPar
右键：剔除/恢复选中数据

\end{itemize}


\paragraph{选项卡2：倾子（Tipper）}
\label{\detokenize{chapters/chapter5:tipper}}
\sphinxAtStartPar
\sphinxstylestrong{显示内容：}
\begin{itemize}
\item {} 
\sphinxAtStartPar
Tzx、Tzy 倾子分量（振幅和相位）

\item {} 
\sphinxAtStartPar
感应矢量指示

\end{itemize}

\sphinxAtStartPar
\sphinxstylestrong{倾子应用：}
\begin{itemize}
\item {} 
\sphinxAtStartPar
识别侧向电性异常

\item {} 
\sphinxAtStartPar
判断二维/三维结构

\end{itemize}


\paragraph{选项卡3：阻抗张量（Z）}
\label{\detokenize{chapters/chapter5:z}}
\sphinxAtStartPar
\sphinxstylestrong{显示内容：}
\begin{itemize}
\item {} 
\sphinxAtStartPar
Zxx、Zxy、Zyx、Zyy 阻抗张量元素

\item {} 
\sphinxAtStartPar
实部和虚部分别显示

\end{itemize}


\paragraph{选项卡4：相位张量（PhaseTensor）}
\label{\detokenize{chapters/chapter5:phasetensor}}
\sphinxAtStartPar
\sphinxstylestrong{显示内容：}
\begin{itemize}
\item {} 
\sphinxAtStartPar
α（主方向角）

\item {} 
\sphinxAtStartPar
β（二维偏离度）

\item {} 
\sphinxAtStartPar
λ（椭圆率）

\item {} 
\sphinxAtStartPar
Φmax、Φmin（相位张量主值）

\item {} 
\sphinxAtStartPar
Skew1D、Skew2D（偏斜度）

\end{itemize}


\paragraph{选项卡5：相干度（Coherence）}
\label{\detokenize{chapters/chapter5:coherence}}
\sphinxAtStartPar
\sphinxstylestrong{显示内容：}
\begin{itemize}
\item {} 
\sphinxAtStartPar
常相干度（Coh²）

\item {} 
\sphinxAtStartPar
重相干度（Cohₘ²）

\item {} 
\sphinxAtStartPar
偏相干度（Cohₚ²）

\item {} 
\sphinxAtStartPar
Ex\sphinxhyphen{}Rx、Ex\sphinxhyphen{}Hy、Ey\sphinxhyphen{}Rx、Ey\sphinxhyphen{}Ry、Ey\sphinxhyphen{}Hx、Hx\sphinxhyphen{}Rx、Hx\sphinxhyphen{}Ry、Hy\sphinxhyphen{}Rx、Hy\sphinxhyphen{}Ry

\end{itemize}

\sphinxAtStartPar
\sphinxstylestrong{质量判断：}


\begin{savenotes}\sphinxattablestart
\sphinxthistablewithglobalstyle
\centering
\begin{tabulary}{\linewidth}[t]{TTT}
\sphinxtoprule
\begin{varwidth}[t]{\sphinxcolwidth{1}{3}}
\sphinxstyletheadfamily \sphinxAtStartPar
相干度值
\sphinxbeforeendvarwidth
\end{varwidth}%
&\begin{varwidth}[t]{\sphinxcolwidth{1}{3}}
\sphinxstyletheadfamily \sphinxAtStartPar
数据质量
\sphinxbeforeendvarwidth
\end{varwidth}%
&\begin{varwidth}[t]{\sphinxcolwidth{1}{3}}
\sphinxstyletheadfamily \sphinxAtStartPar
建议
\sphinxbeforeendvarwidth
\end{varwidth}%
\\
\sphinxmidrule
\sphinxtableatstartofbodyhook\begin{varwidth}[t]{\sphinxcolwidth{1}{3}}
\sphinxAtStartPar
> 0.9
\sphinxbeforeendvarwidth
\end{varwidth}%
&\begin{varwidth}[t]{\sphinxcolwidth{1}{3}}
\sphinxAtStartPar
优秀
\sphinxbeforeendvarwidth
\end{varwidth}%
&\begin{varwidth}[t]{\sphinxcolwidth{1}{3}}
\sphinxAtStartPar
可直接使用
\sphinxbeforeendvarwidth
\end{varwidth}%
\\
\sphinxhline\begin{varwidth}[t]{\sphinxcolwidth{1}{3}}
\sphinxAtStartPar
0.7\sphinxhyphen{}0.9
\sphinxbeforeendvarwidth
\end{varwidth}%
&\begin{varwidth}[t]{\sphinxcolwidth{1}{3}}
\sphinxAtStartPar
良好
\sphinxbeforeendvarwidth
\end{varwidth}%
&\begin{varwidth}[t]{\sphinxcolwidth{1}{3}}
\sphinxAtStartPar
正常使用
\sphinxbeforeendvarwidth
\end{varwidth}%
\\
\sphinxhline\begin{varwidth}[t]{\sphinxcolwidth{1}{3}}
\sphinxAtStartPar
0.5\sphinxhyphen{}0.7
\sphinxbeforeendvarwidth
\end{varwidth}%
&\begin{varwidth}[t]{\sphinxcolwidth{1}{3}}
\sphinxAtStartPar
一般
\sphinxbeforeendvarwidth
\end{varwidth}%
&\begin{varwidth}[t]{\sphinxcolwidth{1}{3}}
\sphinxAtStartPar
考虑筛选
\sphinxbeforeendvarwidth
\end{varwidth}%
\\
\sphinxhline\begin{varwidth}[t]{\sphinxcolwidth{1}{3}}
\sphinxAtStartPar
< 0.5
\sphinxbeforeendvarwidth
\end{varwidth}%
&\begin{varwidth}[t]{\sphinxcolwidth{1}{3}}
\sphinxAtStartPar
较差
\sphinxbeforeendvarwidth
\end{varwidth}%
&\begin{varwidth}[t]{\sphinxcolwidth{1}{3}}
\sphinxAtStartPar
建议剔除
\sphinxbeforeendvarwidth
\end{varwidth}%
\\
\sphinxbottomrule
\end{tabulary}
\sphinxtableafterendhook\par
\sphinxattableend\end{savenotes}


\paragraph{选项卡6：CCZ阻抗}
\label{\detokenize{chapters/chapter5:ccz}}
\sphinxAtStartPar
\sphinxstylestrong{显示内容：}
\begin{itemize}
\item {} 
\sphinxAtStartPar
CCZ阻抗张量

\item {} 
\sphinxAtStartPar
θ（方位角）

\item {} 
\sphinxAtStartPar
Skew1D、Skew2D

\end{itemize}


\paragraph{选项卡7：极化（Polar）}
\label{\detokenize{chapters/chapter5:polar}}
\sphinxAtStartPar
\sphinxstylestrong{显示内容：}
\begin{itemize}
\item {} 
\sphinxAtStartPar
电场极化方向

\item {} 
\sphinxAtStartPar
磁场极化方向

\item {} 
\sphinxAtStartPar
极化椭圆参数

\end{itemize}


\paragraph{选项卡8：自动筛选参数（AutoSelect）}
\label{\detokenize{chapters/chapter5:autoselect}}
\sphinxAtStartPar
\sphinxstylestrong{自动筛选方案配置：}


\begin{savenotes}\sphinxattablestart
\sphinxthistablewithglobalstyle
\centering
\begin{tabulary}{\linewidth}[t]{TT}
\sphinxtoprule
\begin{varwidth}[t]{\sphinxcolwidth{1}{2}}
\sphinxstyletheadfamily \sphinxAtStartPar
参数
\sphinxbeforeendvarwidth
\end{varwidth}%
&\begin{varwidth}[t]{\sphinxcolwidth{1}{2}}
\sphinxstyletheadfamily \sphinxAtStartPar
说明
\sphinxbeforeendvarwidth
\end{varwidth}%
\\
\sphinxmidrule
\sphinxtableatstartofbodyhook\begin{varwidth}[t]{\sphinxcolwidth{1}{2}}
\sphinxAtStartPar
\sphinxstylestrong{筛选方案}
\sphinxbeforeendvarwidth
\end{varwidth}%
&\begin{varwidth}[t]{\sphinxcolwidth{1}{2}}
\sphinxAtStartPar
选择预设的筛选方案
\sphinxbeforeendvarwidth
\end{varwidth}%
\\
\sphinxhline\begin{varwidth}[t]{\sphinxcolwidth{1}{2}}
\sphinxAtStartPar
\sphinxstylestrong{旋转角度列表}
\sphinxbeforeendvarwidth
\end{varwidth}%
&\begin{varwidth}[t]{\sphinxcolwidth{1}{2}}
\sphinxAtStartPar
用于多角度筛选的旋转角度
\sphinxbeforeendvarwidth
\end{varwidth}%
\\
\sphinxhline\begin{varwidth}[t]{\sphinxcolwidth{1}{2}}
\sphinxAtStartPar
\sphinxstylestrong{筛选参数}
\sphinxbeforeendvarwidth
\end{varwidth}%
&\begin{varwidth}[t]{\sphinxcolwidth{1}{2}}
\sphinxAtStartPar
参与筛选的MT参数类型
\sphinxbeforeendvarwidth
\end{varwidth}%
\\
\sphinxhline\begin{varwidth}[t]{\sphinxcolwidth{1}{2}}
\sphinxAtStartPar
\sphinxstylestrong{排序方式}
\sphinxbeforeendvarwidth
\end{varwidth}%
&\begin{varwidth}[t]{\sphinxcolwidth{1}{2}}
\sphinxAtStartPar
升序/降序
\sphinxbeforeendvarwidth
\end{varwidth}%
\\
\sphinxhline\begin{varwidth}[t]{\sphinxcolwidth{1}{2}}
\sphinxAtStartPar
\sphinxstylestrong{权重}
\sphinxbeforeendvarwidth
\end{varwidth}%
&\begin{varwidth}[t]{\sphinxcolwidth{1}{2}}
\sphinxAtStartPar
各参数的权重
\sphinxbeforeendvarwidth
\end{varwidth}%
\\
\sphinxhline\begin{varwidth}[t]{\sphinxcolwidth{1}{2}}
\sphinxAtStartPar
\sphinxstylestrong{保留比例}
\sphinxbeforeendvarwidth
\end{varwidth}%
&\begin{varwidth}[t]{\sphinxcolwidth{1}{2}}
\sphinxAtStartPar
筛选后保留的数据比例
\sphinxbeforeendvarwidth
\end{varwidth}%
\\
\sphinxhline\begin{varwidth}[t]{\sphinxcolwidth{1}{2}}
\sphinxAtStartPar
\sphinxstylestrong{退出误差}
\sphinxbeforeendvarwidth
\end{varwidth}%
&\begin{varwidth}[t]{\sphinxcolwidth{1}{2}}
\sphinxAtStartPar
迭代退出条件
\sphinxbeforeendvarwidth
\end{varwidth}%
\\
\sphinxbottomrule
\end{tabulary}
\sphinxtableafterendhook\par
\sphinxattableend\end{savenotes}

\sphinxAtStartPar
\sphinxstylestrong{可用筛选参数：}


\begin{savenotes}\sphinxattablestart
\sphinxthistablewithglobalstyle
\centering
\begin{tabulary}{\linewidth}[t]{TT}
\sphinxtoprule
\begin{varwidth}[t]{\sphinxcolwidth{1}{2}}
\sphinxstyletheadfamily \sphinxAtStartPar
参数类型
\sphinxbeforeendvarwidth
\end{varwidth}%
&\begin{varwidth}[t]{\sphinxcolwidth{1}{2}}
\sphinxstyletheadfamily \sphinxAtStartPar
说明
\sphinxbeforeendvarwidth
\end{varwidth}%
\\
\sphinxmidrule
\sphinxtableatstartofbodyhook\begin{varwidth}[t]{\sphinxcolwidth{1}{2}}
\sphinxAtStartPar
视电阻率
\sphinxbeforeendvarwidth
\end{varwidth}%
&\begin{varwidth}[t]{\sphinxcolwidth{1}{2}}
\sphinxAtStartPar
ρxy、ρyx
\sphinxbeforeendvarwidth
\end{varwidth}%
\\
\sphinxhline\begin{varwidth}[t]{\sphinxcolwidth{1}{2}}
\sphinxAtStartPar
相位
\sphinxbeforeendvarwidth
\end{varwidth}%
&\begin{varwidth}[t]{\sphinxcolwidth{1}{2}}
\sphinxAtStartPar
φxy、φyx
\sphinxbeforeendvarwidth
\end{varwidth}%
\\
\sphinxhline\begin{varwidth}[t]{\sphinxcolwidth{1}{2}}
\sphinxAtStartPar
相干度
\sphinxbeforeendvarwidth
\end{varwidth}%
&\begin{varwidth}[t]{\sphinxcolwidth{1}{2}}
\sphinxAtStartPar
Ex\sphinxhyphen{}Hx、Ex\sphinxhyphen{}Hy等
\sphinxbeforeendvarwidth
\end{varwidth}%
\\
\sphinxhline\begin{varwidth}[t]{\sphinxcolwidth{1}{2}}
\sphinxAtStartPar
倾子
\sphinxbeforeendvarwidth
\end{varwidth}%
&\begin{varwidth}[t]{\sphinxcolwidth{1}{2}}
\sphinxAtStartPar
Tzx、Tzy
\sphinxbeforeendvarwidth
\end{varwidth}%
\\
\sphinxhline\begin{varwidth}[t]{\sphinxcolwidth{1}{2}}
\sphinxAtStartPar
相位张量
\sphinxbeforeendvarwidth
\end{varwidth}%
&\begin{varwidth}[t]{\sphinxcolwidth{1}{2}}
\sphinxAtStartPar
α、β、λ等
\sphinxbeforeendvarwidth
\end{varwidth}%
\\
\sphinxbottomrule
\end{tabulary}
\sphinxtableafterendhook\par
\sphinxattableend\end{savenotes}


\paragraph{选项卡9：传输方案（Transmit）}
\label{\detokenize{chapters/chapter5:transmit}}
\sphinxAtStartPar
传输方案用于管理频谱数据的传输和剔除规则。


\bigskip\hrule\bigskip



\subsubsection{5.16.6 自动筛选算法详解（RMSMR核心）}
\label{\detokenize{chapters/chapter5:id41}}\begin{quote}

\sphinxAtStartPar
\sphinxstylestrong{📚 理论基础}：自动筛选是RMSMR方法的核心环节。通过智能算法从大量频谱数据中识别并剔除受污染的数据段，保留高质量信号。
\end{quote}


\paragraph{为什么需要自动筛选？}
\label{\detokenize{chapters/chapter5:id42}}
\sphinxAtStartPar
MT数据处理面临的根本问题：\sphinxstylestrong{如何从大量数据中识别哪些是"好数据"？}

\sphinxAtStartPar
\sphinxstylestrong{传统方法的局限：}
\begin{itemize}
\item {} 
\sphinxAtStartPar
只看单一指标（如相干度）

\item {} 
\sphinxAtStartPar
阈值选择主观

\item {} 
\sphinxAtStartPar
难以处理复杂干扰模式

\end{itemize}

\sphinxAtStartPar
\sphinxstylestrong{MTDP的解决方案：多算法融合}
\begin{itemize}
\item {} 
\sphinxAtStartPar
多目标遗传算法(MOGA)：同时优化多个质量指标

\item {} 
\sphinxAtStartPar
Rhoplus约束：用物理模型约束数据质量

\item {} 
\sphinxAtStartPar
马氏距离：统计异常检测

\end{itemize}


\bigskip\hrule\bigskip



\paragraph{筛选类型总览}
\label{\detokenize{chapters/chapter5:id43}}

\bigskip\hrule\bigskip



\paragraph{各算法详解}
\label{\detokenize{chapters/chapter5:id44}}

\subparagraph{1. RhoPlus SA（模拟退火）}
\label{\detokenize{chapters/chapter5:rhoplus-sa}}
\sphinxAtStartPar
\sphinxstylestrong{原理：} 模拟金属退火过程，在解空间中随机搜索，逐步降低"温度"收敛到最优解。

\sphinxAtStartPar
\sphinxstylestrong{物理意义：}
\begin{itemize}
\item {} 
\sphinxAtStartPar
在不同旋转角度下计算视电阻率

\item {} 
\sphinxAtStartPar
寻找使模型最简单的角度组合

\item {} 
\sphinxAtStartPar
用高质量频段的1D响应约束低质量频段

\end{itemize}

\sphinxAtStartPar
\sphinxstylestrong{适用场景：} 快速筛选，作为初次处理


\begin{savenotes}\sphinxattablestart
\sphinxthistablewithglobalstyle
\centering
\begin{tabulary}{\linewidth}[t]{TTT}
\sphinxtoprule
\begin{varwidth}[t]{\sphinxcolwidth{1}{3}}
\sphinxstyletheadfamily \sphinxAtStartPar
参数
\sphinxbeforeendvarwidth
\end{varwidth}%
&\begin{varwidth}[t]{\sphinxcolwidth{1}{3}}
\sphinxstyletheadfamily \sphinxAtStartPar
说明
\sphinxbeforeendvarwidth
\end{varwidth}%
&\begin{varwidth}[t]{\sphinxcolwidth{1}{3}}
\sphinxstyletheadfamily \sphinxAtStartPar
推荐值
\sphinxbeforeendvarwidth
\end{varwidth}%
\\
\sphinxmidrule
\sphinxtableatstartofbodyhook\begin{varwidth}[t]{\sphinxcolwidth{1}{3}}
\sphinxAtStartPar
最大迭代次数
\sphinxbeforeendvarwidth
\end{varwidth}%
&\begin{varwidth}[t]{\sphinxcolwidth{1}{3}}
\sphinxAtStartPar
搜索的最大步数
\sphinxbeforeendvarwidth
\end{varwidth}%
&\begin{varwidth}[t]{\sphinxcolwidth{1}{3}}
\sphinxAtStartPar
10000\sphinxhyphen{}100000
\sphinxbeforeendvarwidth
\end{varwidth}%
\\
\sphinxhline\begin{varwidth}[t]{\sphinxcolwidth{1}{3}}
\sphinxAtStartPar
最小RMS
\sphinxbeforeendvarwidth
\end{varwidth}%
&\begin{varwidth}[t]{\sphinxcolwidth{1}{3}}
\sphinxAtStartPar
低于此值不筛选
\sphinxbeforeendvarwidth
\end{varwidth}%
&\begin{varwidth}[t]{\sphinxcolwidth{1}{3}}
\sphinxAtStartPar
0.5
\sphinxbeforeendvarwidth
\end{varwidth}%
\\
\sphinxbottomrule
\end{tabulary}
\sphinxtableafterendhook\par
\sphinxattableend\end{savenotes}


\subparagraph{2. RhoPlus GA（遗传算法）}
\label{\detokenize{chapters/chapter5:rhoplus-ga}}
\sphinxAtStartPar
\sphinxstylestrong{原理：} 模拟生物进化过程，通过选择、交叉、变异操作逐代优化。

\sphinxAtStartPar
\sphinxstylestrong{与SA的区别：}
\begin{itemize}
\item {} 
\sphinxAtStartPar
维护一个种群而非单个解

\item {} 
\sphinxAtStartPar
更容易跳出局部最优

\item {} 
\sphinxAtStartPar
计算量更大但结果更稳定

\end{itemize}

\sphinxAtStartPar
\sphinxstylestrong{适用场景：} SA效果不佳时的升级选择


\subparagraph{3. RhoPlus MOGA（多目标遗传算法）⭐ 推荐}
\label{\detokenize{chapters/chapter5:rhoplus-moga}}
\sphinxAtStartPar
\sphinxstylestrong{核心创新：} 同时优化多个目标函数，而非单一目标

\sphinxAtStartPar
\sphinxstylestrong{多目标的物理意义：}


\begin{savenotes}\sphinxattablestart
\sphinxthistablewithglobalstyle
\centering
\begin{tabulary}{\linewidth}[t]{TTT}
\sphinxtoprule
\begin{varwidth}[t]{\sphinxcolwidth{1}{3}}
\sphinxstyletheadfamily \sphinxAtStartPar
目标函数
\sphinxbeforeendvarwidth
\end{varwidth}%
&\begin{varwidth}[t]{\sphinxcolwidth{1}{3}}
\sphinxstyletheadfamily \sphinxAtStartPar
物理意义
\sphinxbeforeendvarwidth
\end{varwidth}%
&\begin{varwidth}[t]{\sphinxcolwidth{1}{3}}
\sphinxstyletheadfamily \sphinxAtStartPar
为什么要优化？
\sphinxbeforeendvarwidth
\end{varwidth}%
\\
\sphinxmidrule
\sphinxtableatstartofbodyhook\begin{varwidth}[t]{\sphinxcolwidth{1}{3}}
\sphinxAtStartPar
Zxy实部RMS
\sphinxbeforeendvarwidth
\end{varwidth}%
&\begin{varwidth}[t]{\sphinxcolwidth{1}{3}}
\sphinxAtStartPar
阻抗实部与预测值的偏差
\sphinxbeforeendvarwidth
\end{varwidth}%
&\begin{varwidth}[t]{\sphinxcolwidth{1}{3}}
\sphinxAtStartPar
确保振幅准确
\sphinxbeforeendvarwidth
\end{varwidth}%
\\
\sphinxhline\begin{varwidth}[t]{\sphinxcolwidth{1}{3}}
\sphinxAtStartPar
Zxy虚部RMS
\sphinxbeforeendvarwidth
\end{varwidth}%
&\begin{varwidth}[t]{\sphinxcolwidth{1}{3}}
\sphinxAtStartPar
阻抗虚部与预测值的偏差
\sphinxbeforeendvarwidth
\end{varwidth}%
&\begin{varwidth}[t]{\sphinxcolwidth{1}{3}}
\sphinxAtStartPar
确保相位准确
\sphinxbeforeendvarwidth
\end{varwidth}%
\\
\sphinxhline\begin{varwidth}[t]{\sphinxcolwidth{1}{3}}
\sphinxAtStartPar
Zyx实部RMS
\sphinxbeforeendvarwidth
\end{varwidth}%
&\begin{varwidth}[t]{\sphinxcolwidth{1}{3}}
\sphinxAtStartPar
正交方向阻抗实部
\sphinxbeforeendvarwidth
\end{varwidth}%
&\begin{varwidth}[t]{\sphinxcolwidth{1}{3}}
\sphinxAtStartPar
确保张量完整
\sphinxbeforeendvarwidth
\end{varwidth}%
\\
\sphinxhline\begin{varwidth}[t]{\sphinxcolwidth{1}{3}}
\sphinxAtStartPar
Zyx虚部RMS
\sphinxbeforeendvarwidth
\end{varwidth}%
&\begin{varwidth}[t]{\sphinxcolwidth{1}{3}}
\sphinxAtStartPar
正交方向阻抗虚部
\sphinxbeforeendvarwidth
\end{varwidth}%
&\begin{varwidth}[t]{\sphinxcolwidth{1}{3}}
\sphinxAtStartPar
确保张量完整
\sphinxbeforeendvarwidth
\end{varwidth}%
\\
\sphinxbottomrule
\end{tabulary}
\sphinxtableafterendhook\par
\sphinxattableend\end{savenotes}

\sphinxAtStartPar
\sphinxstylestrong{为什么MOGA比GA更好？}

\sphinxAtStartPar
\sphinxstylestrong{适用场景：} 强干扰环境，作为最终处理步骤


\subparagraph{4. Full Z MOGA（全阻抗MOGA）⭐ 高精度}
\label{\detokenize{chapters/chapter5:full-z-moga-moga}}
\sphinxAtStartPar
\sphinxstylestrong{与RhoPlus MOGA的区别：}


\begin{savenotes}\sphinxattablestart
\sphinxthistablewithglobalstyle
\centering
\begin{tabulary}{\linewidth}[t]{TTT}
\sphinxtoprule
\begin{varwidth}[t]{\sphinxcolwidth{1}{3}}
\sphinxstyletheadfamily \sphinxAtStartPar
对比项
\sphinxbeforeendvarwidth
\end{varwidth}%
&\begin{varwidth}[t]{\sphinxcolwidth{1}{3}}
\sphinxstyletheadfamily \sphinxAtStartPar
RhoPlus MOGA
\sphinxbeforeendvarwidth
\end{varwidth}%
&\begin{varwidth}[t]{\sphinxcolwidth{1}{3}}
\sphinxstyletheadfamily \sphinxAtStartPar
Full Z MOGA
\sphinxbeforeendvarwidth
\end{varwidth}%
\\
\sphinxmidrule
\sphinxtableatstartofbodyhook\begin{varwidth}[t]{\sphinxcolwidth{1}{3}}
\sphinxAtStartPar
约束数据
\sphinxbeforeendvarwidth
\end{varwidth}%
&\begin{varwidth}[t]{\sphinxcolwidth{1}{3}}
\sphinxAtStartPar
Zxy, Zyx
\sphinxbeforeendvarwidth
\end{varwidth}%
&\begin{varwidth}[t]{\sphinxcolwidth{1}{3}}
\sphinxAtStartPar
Zxx, Zxy, Zyx, Zyy
\sphinxbeforeendvarwidth
\end{varwidth}%
\\
\sphinxhline\begin{varwidth}[t]{\sphinxcolwidth{1}{3}}
\sphinxAtStartPar
约束来源
\sphinxbeforeendvarwidth
\end{varwidth}%
&\begin{varwidth}[t]{\sphinxcolwidth{1}{3}}
\sphinxAtStartPar
多角度旋转
\sphinxbeforeendvarwidth
\end{varwidth}%
&\begin{varwidth}[t]{\sphinxcolwidth{1}{3}}
\sphinxAtStartPar
理论模型预测
\sphinxbeforeendvarwidth
\end{varwidth}%
\\
\sphinxhline\begin{varwidth}[t]{\sphinxcolwidth{1}{3}}
\sphinxAtStartPar
精度
\sphinxbeforeendvarwidth
\end{varwidth}%
&\begin{varwidth}[t]{\sphinxcolwidth{1}{3}}
\sphinxAtStartPar
高
\sphinxbeforeendvarwidth
\end{varwidth}%
&\begin{varwidth}[t]{\sphinxcolwidth{1}{3}}
\sphinxAtStartPar
最高
\sphinxbeforeendvarwidth
\end{varwidth}%
\\
\sphinxhline\begin{varwidth}[t]{\sphinxcolwidth{1}{3}}
\sphinxAtStartPar
计算量
\sphinxbeforeendvarwidth
\end{varwidth}%
&\begin{varwidth}[t]{\sphinxcolwidth{1}{3}}
\sphinxAtStartPar
中等
\sphinxbeforeendvarwidth
\end{varwidth}%
&\begin{varwidth}[t]{\sphinxcolwidth{1}{3}}
\sphinxAtStartPar
大
\sphinxbeforeendvarwidth
\end{varwidth}%
\\
\sphinxbottomrule
\end{tabulary}
\sphinxtableafterendhook\par
\sphinxattableend\end{savenotes}

\sphinxAtStartPar
\sphinxstylestrong{物理意义：}
\begin{itemize}
\item {} 
\sphinxAtStartPar
不仅约束主阻抗(Zxy, Zyx)

\item {} 
\sphinxAtStartPar
同时约束对角元素(Zxx, Zyy)

\item {} 
\sphinxAtStartPar
对于三维结构更准确

\end{itemize}

\sphinxAtStartPar
\sphinxstylestrong{适用场景：} 需要最高精度的研究项目


\bigskip\hrule\bigskip



\paragraph{马氏距离筛选（预处理步骤）}
\label{\detokenize{chapters/chapter5:id45}}
\sphinxAtStartPar
\sphinxstylestrong{理论基础：} 马氏距离是考虑变量间相关性的广义距离。
\begin{equation*}
\begin{split}D_M = \sqrt{(\mathbf{x} - \boldsymbol{\mu})^T \mathbf{\Sigma}^{-1} (\mathbf{x} - \boldsymbol{\mu})}\end{split}
\end{equation*}
\sphinxAtStartPar
\sphinxstylestrong{与欧氏距离的区别：}

\sphinxAtStartPar
\sphinxstylestrong{实际应用：}

\sphinxAtStartPar
当多个参数同时出现轻微异常时，单独看每个参数都不明显，但马氏距离能识别出这种"组合异常"。

\sphinxAtStartPar
\sphinxstylestrong{参数设置：}


\begin{savenotes}\sphinxattablestart
\sphinxthistablewithglobalstyle
\centering
\begin{tabulary}{\linewidth}[t]{TTTT}
\sphinxtoprule
\begin{varwidth}[t]{\sphinxcolwidth{1}{4}}
\sphinxstyletheadfamily \sphinxAtStartPar
参数
\sphinxbeforeendvarwidth
\end{varwidth}%
&\begin{varwidth}[t]{\sphinxcolwidth{1}{4}}
\sphinxstyletheadfamily \sphinxAtStartPar
说明
\sphinxbeforeendvarwidth
\end{varwidth}%
&\begin{varwidth}[t]{\sphinxcolwidth{1}{4}}
\sphinxstyletheadfamily \sphinxAtStartPar
推荐值
\sphinxbeforeendvarwidth
\end{varwidth}%
&\begin{varwidth}[t]{\sphinxcolwidth{1}{4}}
\sphinxstyletheadfamily \sphinxAtStartPar
物理意义
\sphinxbeforeendvarwidth
\end{varwidth}%
\\
\sphinxmidrule
\sphinxtableatstartofbodyhook\begin{varwidth}[t]{\sphinxcolwidth{1}{4}}
\sphinxAtStartPar
阈值
\sphinxbeforeendvarwidth
\end{varwidth}%
&\begin{varwidth}[t]{\sphinxcolwidth{1}{4}}
\sphinxAtStartPar
马氏距离阈值
\sphinxbeforeendvarwidth
\end{varwidth}%
&\begin{varwidth}[t]{\sphinxcolwidth{1}{4}}
\sphinxAtStartPar
2.0\sphinxhyphen{}3.0
\sphinxbeforeendvarwidth
\end{varwidth}%
&\begin{varwidth}[t]{\sphinxcolwidth{1}{4}}
\sphinxAtStartPar
对应95\%\sphinxhyphen{}99\%置信区间
\sphinxbeforeendvarwidth
\end{varwidth}%
\\
\sphinxhline\begin{varwidth}[t]{\sphinxcolwidth{1}{4}}
\sphinxAtStartPar
最大迭代次数
\sphinxbeforeendvarwidth
\end{varwidth}%
&\begin{varwidth}[t]{\sphinxcolwidth{1}{4}}
\sphinxAtStartPar
迭代筛选次数
\sphinxbeforeendvarwidth
\end{varwidth}%
&\begin{varwidth}[t]{\sphinxcolwidth{1}{4}}
\sphinxAtStartPar
10\sphinxhyphen{}50
\sphinxbeforeendvarwidth
\end{varwidth}%
&\begin{varwidth}[t]{\sphinxcolwidth{1}{4}}
\sphinxAtStartPar
迭代越多越严格
\sphinxbeforeendvarwidth
\end{varwidth}%
\\
\sphinxhline\begin{varwidth}[t]{\sphinxcolwidth{1}{4}}
\sphinxAtStartPar
参数集
\sphinxbeforeendvarwidth
\end{varwidth}%
&\begin{varwidth}[t]{\sphinxcolwidth{1}{4}}
\sphinxAtStartPar
参与计算的MT参数
\sphinxbeforeendvarwidth
\end{varwidth}%
&\begin{varwidth}[t]{\sphinxcolwidth{1}{4}}
\sphinxAtStartPar
选择主要参数
\sphinxbeforeendvarwidth
\end{varwidth}%
&\begin{varwidth}[t]{\sphinxcolwidth{1}{4}}
\sphinxAtStartPar
视电阻率+相位+相干度
\sphinxbeforeendvarwidth
\end{varwidth}%
\\
\sphinxbottomrule
\end{tabulary}
\sphinxtableafterendhook\par
\sphinxattableend\end{savenotes}

\sphinxAtStartPar
\sphinxstylestrong{操作按钮：}


\begin{savenotes}\sphinxattablestart
\sphinxthistablewithglobalstyle
\centering
\begin{tabulary}{\linewidth}[t]{TTT}
\sphinxtoprule
\begin{varwidth}[t]{\sphinxcolwidth{1}{3}}
\sphinxstyletheadfamily \sphinxAtStartPar
按钮
\sphinxbeforeendvarwidth
\end{varwidth}%
&\begin{varwidth}[t]{\sphinxcolwidth{1}{3}}
\sphinxstyletheadfamily \sphinxAtStartPar
功能
\sphinxbeforeendvarwidth
\end{varwidth}%
&\begin{varwidth}[t]{\sphinxcolwidth{1}{3}}
\sphinxstyletheadfamily \sphinxAtStartPar
使用时机
\sphinxbeforeendvarwidth
\end{varwidth}%
\\
\sphinxmidrule
\sphinxtableatstartofbodyhook\begin{varwidth}[t]{\sphinxcolwidth{1}{3}}
\sphinxAtStartPar
\sphinxstylestrong{MD自动筛选}
\sphinxbeforeendvarwidth
\end{varwidth}%
&\begin{varwidth}[t]{\sphinxcolwidth{1}{3}}
\sphinxAtStartPar
对当前频率执行
\sphinxbeforeendvarwidth
\end{varwidth}%
&\begin{varwidth}[t]{\sphinxcolwidth{1}{3}}
\sphinxAtStartPar
单频点问题诊断
\sphinxbeforeendvarwidth
\end{varwidth}%
\\
\sphinxhline\begin{varwidth}[t]{\sphinxcolwidth{1}{3}}
\sphinxAtStartPar
\sphinxstylestrong{MD自动筛选（固定H）}
\sphinxbeforeendvarwidth
\end{varwidth}%
&\begin{varwidth}[t]{\sphinxcolwidth{1}{3}}
\sphinxAtStartPar
固定H方向的筛选
\sphinxbeforeendvarwidth
\end{varwidth}%
&\begin{varwidth}[t]{\sphinxcolwidth{1}{3}}
\sphinxAtStartPar
H场数据质量好时
\sphinxbeforeendvarwidth
\end{varwidth}%
\\
\sphinxhline\begin{varwidth}[t]{\sphinxcolwidth{1}{3}}
\sphinxAtStartPar
\sphinxstylestrong{全部MD筛选}
\sphinxbeforeendvarwidth
\end{varwidth}%
&\begin{varwidth}[t]{\sphinxcolwidth{1}{3}}
\sphinxAtStartPar
对所有频率执行
\sphinxbeforeendvarwidth
\end{varwidth}%
&\begin{varwidth}[t]{\sphinxcolwidth{1}{3}}
\sphinxAtStartPar
整体质量提升
\sphinxbeforeendvarwidth
\end{varwidth}%
\\
\sphinxhline\begin{varwidth}[t]{\sphinxcolwidth{1}{3}}
\sphinxAtStartPar
\sphinxstylestrong{全部MD筛选（固定H）}
\sphinxbeforeendvarwidth
\end{varwidth}%
&\begin{varwidth}[t]{\sphinxcolwidth{1}{3}}
\sphinxAtStartPar
对所有频率执行
\sphinxbeforeendvarwidth
\end{varwidth}%
&\begin{varwidth}[t]{\sphinxcolwidth{1}{3}}
\sphinxAtStartPar
H场数据质量好时
\sphinxbeforeendvarwidth
\end{varwidth}%
\\
\sphinxbottomrule
\end{tabulary}
\sphinxtableafterendhook\par
\sphinxattableend\end{savenotes}


\bigskip\hrule\bigskip



\paragraph{筛选算法选择指南}
\label{\detokenize{chapters/chapter5:id46}}
\sphinxAtStartPar
\sphinxstylestrong{推荐流程：}
\begin{enumerate}
\sphinxsetlistlabels{\arabic}{enumi}{enumii}{}{.}%
\item {} 
\sphinxAtStartPar
\sphinxstylestrong{预处理}：马氏距离筛选 → 剔除明显异常

\item {} 
\sphinxAtStartPar
\sphinxstylestrong{主处理}：RhoPlus MOGA → 精细筛选

\item {} 
\sphinxAtStartPar
\sphinxstylestrong{可选}：Full Z MOGA → 最终优化（如需要）

\end{enumerate}

\sphinxAtStartPar
\sphinxstylestrong{保留比例建议：}


\begin{savenotes}\sphinxattablestart
\sphinxthistablewithglobalstyle
\centering
\begin{tabulary}{\linewidth}[t]{TTT}
\sphinxtoprule
\begin{varwidth}[t]{\sphinxcolwidth{1}{3}}
\sphinxstyletheadfamily \sphinxAtStartPar
环境类型
\sphinxbeforeendvarwidth
\end{varwidth}%
&\begin{varwidth}[t]{\sphinxcolwidth{1}{3}}
\sphinxstyletheadfamily \sphinxAtStartPar
保留比例
\sphinxbeforeendvarwidth
\end{varwidth}%
&\begin{varwidth}[t]{\sphinxcolwidth{1}{3}}
\sphinxstyletheadfamily \sphinxAtStartPar
说明
\sphinxbeforeendvarwidth
\end{varwidth}%
\\
\sphinxmidrule
\sphinxtableatstartofbodyhook\begin{varwidth}[t]{\sphinxcolwidth{1}{3}}
\sphinxAtStartPar
低噪声环境
\sphinxbeforeendvarwidth
\end{varwidth}%
&\begin{varwidth}[t]{\sphinxcolwidth{1}{3}}
\sphinxAtStartPar
50\%\sphinxhyphen{}70\%
\sphinxbeforeendvarwidth
\end{varwidth}%
&\begin{varwidth}[t]{\sphinxcolwidth{1}{3}}
\sphinxAtStartPar
可以保留更多数据
\sphinxbeforeendvarwidth
\end{varwidth}%
\\
\sphinxhline\begin{varwidth}[t]{\sphinxcolwidth{1}{3}}
\sphinxAtStartPar
中等噪声
\sphinxbeforeendvarwidth
\end{varwidth}%
&\begin{varwidth}[t]{\sphinxcolwidth{1}{3}}
\sphinxAtStartPar
30\%\sphinxhyphen{}50\%
\sphinxbeforeendvarwidth
\end{varwidth}%
&\begin{varwidth}[t]{\sphinxcolwidth{1}{3}}
\sphinxAtStartPar
平衡数据量和质量
\sphinxbeforeendvarwidth
\end{varwidth}%
\\
\sphinxhline\begin{varwidth}[t]{\sphinxcolwidth{1}{3}}
\sphinxAtStartPar
强干扰环境
\sphinxbeforeendvarwidth
\end{varwidth}%
&\begin{varwidth}[t]{\sphinxcolwidth{1}{3}}
\sphinxAtStartPar
20\%\sphinxhyphen{}30\%
\sphinxbeforeendvarwidth
\end{varwidth}%
&\begin{varwidth}[t]{\sphinxcolwidth{1}{3}}
\sphinxAtStartPar
严格筛选，宁可少而精
\sphinxbeforeendvarwidth
\end{varwidth}%
\\
\sphinxbottomrule
\end{tabulary}
\sphinxtableafterendhook\par
\sphinxattableend\end{savenotes}


\bigskip\hrule\bigskip



\subsubsection{5.16.7 数据导出功能}
\label{\detokenize{chapters/chapter5:id47}}
\sphinxAtStartPar
频谱编辑窗口支持多种数据导出格式。


\paragraph{导出菜单概览}
\label{\detokenize{chapters/chapter5:id48}}

\paragraph{EDI格式导出}
\label{\detokenize{chapters/chapter5:edi}}

\begin{savenotes}\sphinxattablestart
\sphinxthistablewithglobalstyle
\centering
\begin{tabulary}{\linewidth}[t]{TTT}
\sphinxtoprule
\begin{varwidth}[t]{\sphinxcolwidth{1}{3}}
\sphinxstyletheadfamily \sphinxAtStartPar
选项
\sphinxbeforeendvarwidth
\end{varwidth}%
&\begin{varwidth}[t]{\sphinxcolwidth{1}{3}}
\sphinxstyletheadfamily \sphinxAtStartPar
说明
\sphinxbeforeendvarwidth
\end{varwidth}%
&\begin{varwidth}[t]{\sphinxcolwidth{1}{3}}
\sphinxstyletheadfamily \sphinxAtStartPar
内容
\sphinxbeforeendvarwidth
\end{varwidth}%
\\
\sphinxmidrule
\sphinxtableatstartofbodyhook\begin{varwidth}[t]{\sphinxcolwidth{1}{3}}
\sphinxAtStartPar
\sphinxstylestrong{导出阻抗EDI}
\sphinxbeforeendvarwidth
\end{varwidth}%
&\begin{varwidth}[t]{\sphinxcolwidth{1}{3}}
\sphinxAtStartPar
导出阻抗张量EDI
\sphinxbeforeendvarwidth
\end{varwidth}%
&\begin{varwidth}[t]{\sphinxcolwidth{1}{3}}
\sphinxAtStartPar
Zxx, Zxy, Zyx, Zyy
\sphinxbeforeendvarwidth
\end{varwidth}%
\\
\sphinxhline\begin{varwidth}[t]{\sphinxcolwidth{1}{3}}
\sphinxAtStartPar
\sphinxstylestrong{导出频谱EDI}
\sphinxbeforeendvarwidth
\end{varwidth}%
&\begin{varwidth}[t]{\sphinxcolwidth{1}{3}}
\sphinxAtStartPar
导出频谱数据EDI
\sphinxbeforeendvarwidth
\end{varwidth}%
&\begin{varwidth}[t]{\sphinxcolwidth{1}{3}}
\sphinxAtStartPar
功率谱、互功率谱
\sphinxbeforeendvarwidth
\end{varwidth}%
\\
\sphinxbottomrule
\end{tabulary}
\sphinxtableafterendhook\par
\sphinxattableend\end{savenotes}


\paragraph{JSON格式导出}
\label{\detokenize{chapters/chapter5:json}}\begin{itemize}
\item {} 
\sphinxAtStartPar
将FC数据导出为JSON格式

\item {} 
\sphinxAtStartPar
便于程序处理和数据交换

\item {} 
\sphinxAtStartPar
文件路径自动复制到剪贴板

\end{itemize}


\paragraph{参数导出}
\label{\detokenize{chapters/chapter5:id49}}\begin{itemize}
\item {} 
\sphinxAtStartPar
导出所有MT参数到指定目录

\item {} 
\sphinxAtStartPar
包括视电阻率、相位、阻抗、倾子等

\item {} 
\sphinxAtStartPar
每个参数保存为单独文件

\end{itemize}


\paragraph{残差权重导出}
\label{\detokenize{chapters/chapter5:id50}}
\sphinxAtStartPar
导出残差和权重数据，用于分析处理效果。

\sphinxAtStartPar
\sphinxstylestrong{输出格式：}

\begin{sphinxVerbatim}[commandchars=\\\{\}]
\PYG{c+c1}{\PYGZsh{}Index ExRes EyRes HzRes ExHuberW EyHuberW HzHuberW ExThomsonW EyThomsonW HzThomsonW Time}
\PYG{l+m+mi}{000001}  \PYG{l+m+mf}{1.234e\PYGZhy{}03} \PYG{l+m+mf}{2.345e\PYGZhy{}03} \PYG{l+m+mf}{3.456e\PYGZhy{}04} \PYG{l+m+mf}{1.000} \PYG{l+m+mf}{0.950} \PYG{l+m+mf}{1.000} \PYG{l+m+mf}{0.980} \PYG{l+m+mf}{0.930} \PYG{l+m+mf}{0.990} \PYG{l+m+mi}{2024}\PYG{o}{\PYGZhy{}}\PYG{l+m+mi}{01}\PYG{o}{\PYGZhy{}}\PYG{l+m+mi}{15}\PYG{n}{T10}\PYG{p}{:}\PYG{l+m+mi}{30}\PYG{p}{:}\PYG{l+m+mf}{00.000}
\PYG{o}{.}\PYG{o}{.}\PYG{o}{.}
\end{sphinxVerbatim}

\sphinxAtStartPar
\sphinxstylestrong{字段说明：}


\begin{savenotes}\sphinxattablestart
\sphinxthistablewithglobalstyle
\centering
\begin{tabulary}{\linewidth}[t]{TT}
\sphinxtoprule
\begin{varwidth}[t]{\sphinxcolwidth{1}{2}}
\sphinxstyletheadfamily \sphinxAtStartPar
字段
\sphinxbeforeendvarwidth
\end{varwidth}%
&\begin{varwidth}[t]{\sphinxcolwidth{1}{2}}
\sphinxstyletheadfamily \sphinxAtStartPar
说明
\sphinxbeforeendvarwidth
\end{varwidth}%
\\
\sphinxmidrule
\sphinxtableatstartofbodyhook\begin{varwidth}[t]{\sphinxcolwidth{1}{2}}
\sphinxAtStartPar
Index
\sphinxbeforeendvarwidth
\end{varwidth}%
&\begin{varwidth}[t]{\sphinxcolwidth{1}{2}}
\sphinxAtStartPar
数据点索引
\sphinxbeforeendvarwidth
\end{varwidth}%
\\
\sphinxhline\begin{varwidth}[t]{\sphinxcolwidth{1}{2}}
\sphinxAtStartPar
ExRes/EyRes/HzRes
\sphinxbeforeendvarwidth
\end{varwidth}%
&\begin{varwidth}[t]{\sphinxcolwidth{1}{2}}
\sphinxAtStartPar
Ex、Ey、Hz通道的残差
\sphinxbeforeendvarwidth
\end{varwidth}%
\\
\sphinxhline\begin{varwidth}[t]{\sphinxcolwidth{1}{2}}
\sphinxAtStartPar
ExHuberW/EyHuberW/HzHuberW
\sphinxbeforeendvarwidth
\end{varwidth}%
&\begin{varwidth}[t]{\sphinxcolwidth{1}{2}}
\sphinxAtStartPar
Huber权重
\sphinxbeforeendvarwidth
\end{varwidth}%
\\
\sphinxhline\begin{varwidth}[t]{\sphinxcolwidth{1}{2}}
\sphinxAtStartPar
ExThomsonW/EyThomsonW/HzThomsonW
\sphinxbeforeendvarwidth
\end{varwidth}%
&\begin{varwidth}[t]{\sphinxcolwidth{1}{2}}
\sphinxAtStartPar
Thomson权重
\sphinxbeforeendvarwidth
\end{varwidth}%
\\
\sphinxhline\begin{varwidth}[t]{\sphinxcolwidth{1}{2}}
\sphinxAtStartPar
Time
\sphinxbeforeendvarwidth
\end{varwidth}%
&\begin{varwidth}[t]{\sphinxcolwidth{1}{2}}
\sphinxAtStartPar
数据时间戳（GMT格式）
\sphinxbeforeendvarwidth
\end{varwidth}%
\\
\sphinxbottomrule
\end{tabulary}
\sphinxtableafterendhook\par
\sphinxattableend\end{savenotes}


\bigskip\hrule\bigskip



\subsubsection{5.16.8 多参考站处理}
\label{\detokenize{chapters/chapter5:id51}}
\sphinxAtStartPar
MTDP支持多参考站处理，可以利用多个远参考站提高数据质量。

\sphinxAtStartPar
\sphinxstylestrong{操作步骤：}
\begin{enumerate}
\sphinxsetlistlabels{\arabic}{enumi}{enumii}{}{.}%
\item {} 
\sphinxAtStartPar
在远参考站列表中勾选要使用的参考站

\item {} 
\sphinxAtStartPar
选择远参考方案（推荐REH）

\item {} 
\sphinxAtStartPar
点击"添加多参考站处理"按钮

\item {} 
\sphinxAtStartPar
系统自动创建新的编辑FC

\end{enumerate}


\bigskip\hrule\bigskip



\subsubsection{5.16.9 理论模型对比}
\label{\detokenize{chapters/chapter5:id52}}
\sphinxAtStartPar
频谱编辑窗口支持加载理论模型进行对比分析。

\sphinxAtStartPar
\sphinxstylestrong{功能：}


\begin{savenotes}\sphinxattablestart
\sphinxthistablewithglobalstyle
\centering
\begin{tabulary}{\linewidth}[t]{TT}
\sphinxtoprule
\begin{varwidth}[t]{\sphinxcolwidth{1}{2}}
\sphinxstyletheadfamily \sphinxAtStartPar
功能
\sphinxbeforeendvarwidth
\end{varwidth}%
&\begin{varwidth}[t]{\sphinxcolwidth{1}{2}}
\sphinxstyletheadfamily \sphinxAtStartPar
说明
\sphinxbeforeendvarwidth
\end{varwidth}%
\\
\sphinxmidrule
\sphinxtableatstartofbodyhook\begin{varwidth}[t]{\sphinxcolwidth{1}{2}}
\sphinxAtStartPar
\sphinxstylestrong{加载预测值}
\sphinxbeforeendvarwidth
\end{varwidth}%
&\begin{varwidth}[t]{\sphinxcolwidth{1}{2}}
\sphinxAtStartPar
从文件加载理论模型预测值
\sphinxbeforeendvarwidth
\end{varwidth}%
\\
\sphinxhline\begin{varwidth}[t]{\sphinxcolwidth{1}{2}}
\sphinxAtStartPar
\sphinxstylestrong{保存预测值}
\sphinxbeforeendvarwidth
\end{varwidth}%
&\begin{varwidth}[t]{\sphinxcolwidth{1}{2}}
\sphinxAtStartPar
保存当前预测值到文件
\sphinxbeforeendvarwidth
\end{varwidth}%
\\
\sphinxhline\begin{varwidth}[t]{\sphinxcolwidth{1}{2}}
\sphinxAtStartPar
\sphinxstylestrong{理论旋转}
\sphinxbeforeendvarwidth
\end{varwidth}%
&\begin{varwidth}[t]{\sphinxcolwidth{1}{2}}
\sphinxAtStartPar
旋转理论模型
\sphinxbeforeendvarwidth
\end{varwidth}%
\\
\sphinxbottomrule
\end{tabulary}
\sphinxtableafterendhook\par
\sphinxattableend\end{savenotes}

\sphinxAtStartPar
\sphinxstylestrong{预测值文件格式：}

\begin{sphinxVerbatim}[commandchars=\\\{\}]
\PYG{n}{fre} \PYG{n}{zxxr} \PYG{n}{zxxi} \PYG{n}{zxyr} \PYG{n}{zxyi} \PYG{n}{zyxr} \PYG{n}{zyxi} \PYG{n}{zyyr} \PYG{n}{zyyi}
\PYG{l+m+mf}{100.0} \PYG{l+m+mf}{0.1234} \PYG{o}{\PYGZhy{}}\PYG{l+m+mf}{0.0567} \PYG{l+m+mf}{1.2345} \PYG{o}{\PYGZhy{}}\PYG{l+m+mf}{0.7890} \PYG{o}{\PYGZhy{}}\PYG{l+m+mf}{1.4567} \PYG{l+m+mf}{0.8901} \PYG{l+m+mf}{0.0987} \PYG{o}{\PYGZhy{}}\PYG{l+m+mf}{0.0432}
\PYG{o}{.}\PYG{o}{.}\PYG{o}{.}
\end{sphinxVerbatim}


\bigskip\hrule\bigskip



\subsubsection{5.16.10 常用操作流程（RMSMR实践）}
\label{\detokenize{chapters/chapter5:id53}}\begin{quote}

\sphinxAtStartPar
\sphinxstylestrong{💡 RMSMR方法的核心思想：} 迭代优化，逐步提升数据质量。每一轮处理后检查结果，不满意则调整参数重试。
\end{quote}


\paragraph{流程1：快速质量检查（1分钟）}
\label{\detokenize{chapters/chapter5:id54}}
\sphinxAtStartPar
\sphinxstylestrong{目的：} 快速评估数据质量，决定是否需要精细处理


\bigskip\hrule\bigskip



\paragraph{流程2：远参考处理（5分钟）}
\label{\detokenize{chapters/chapter5:id55}}
\sphinxAtStartPar
\sphinxstylestrong{理论基础：} 远参考利用"噪声不相关、信号相关"的特性消除本地噪声偏差。


\bigskip\hrule\bigskip



\paragraph{流程3：自动筛选（10分钟）}
\label{\detokenize{chapters/chapter5:id56}}
\sphinxAtStartPar
\sphinxstylestrong{理论基础：} 使用多目标遗传算法(MOGA)，同时优化多个质量指标。


\bigskip\hrule\bigskip



\paragraph{流程4：完整RMSMR处理流程（20\sphinxhyphen{}30分钟）}
\label{\detokenize{chapters/chapter5:rmsmr-20-30}}
\sphinxAtStartPar
\sphinxstylestrong{适用于强干扰环境或高精度要求项目}


\bigskip\hrule\bigskip



\subsubsection{5.16.11 快捷键参考}
\label{\detokenize{chapters/chapter5:id57}}

\begin{savenotes}\sphinxattablestart
\sphinxthistablewithglobalstyle
\centering
\begin{tabulary}{\linewidth}[t]{TT}
\sphinxtoprule
\begin{varwidth}[t]{\sphinxcolwidth{1}{2}}
\sphinxstyletheadfamily \sphinxAtStartPar
快捷键
\sphinxbeforeendvarwidth
\end{varwidth}%
&\begin{varwidth}[t]{\sphinxcolwidth{1}{2}}
\sphinxstyletheadfamily \sphinxAtStartPar
功能
\sphinxbeforeendvarwidth
\end{varwidth}%
\\
\sphinxmidrule
\sphinxtableatstartofbodyhook\begin{varwidth}[t]{\sphinxcolwidth{1}{2}}
\sphinxAtStartPar
←
\sphinxbeforeendvarwidth
\end{varwidth}%
&\begin{varwidth}[t]{\sphinxcolwidth{1}{2}}
\sphinxAtStartPar
切换到上一个频率
\sphinxbeforeendvarwidth
\end{varwidth}%
\\
\sphinxhline\begin{varwidth}[t]{\sphinxcolwidth{1}{2}}
\sphinxAtStartPar
→
\sphinxbeforeendvarwidth
\end{varwidth}%
&\begin{varwidth}[t]{\sphinxcolwidth{1}{2}}
\sphinxAtStartPar
切换到下一个频率
\sphinxbeforeendvarwidth
\end{varwidth}%
\\
\sphinxhline\begin{varwidth}[t]{\sphinxcolwidth{1}{2}}
\sphinxAtStartPar
Ctrl+S
\sphinxbeforeendvarwidth
\end{varwidth}%
&\begin{varwidth}[t]{\sphinxcolwidth{1}{2}}
\sphinxAtStartPar
保存当前FC
\sphinxbeforeendvarwidth
\end{varwidth}%
\\
\sphinxhline\begin{varwidth}[t]{\sphinxcolwidth{1}{2}}
\sphinxAtStartPar
Ctrl+Shift+S
\sphinxbeforeendvarwidth
\end{varwidth}%
&\begin{varwidth}[t]{\sphinxcolwidth{1}{2}}
\sphinxAtStartPar
保存所有FC
\sphinxbeforeendvarwidth
\end{varwidth}%
\\
\sphinxbottomrule
\end{tabulary}
\sphinxtableafterendhook\par
\sphinxattableend\end{savenotes}


\bigskip\hrule\bigskip



\subsubsection{5.16.12 常见问题与解决}
\label{\detokenize{chapters/chapter5:id58}}

\begin{savenotes}\sphinxattablestart
\sphinxthistablewithglobalstyle
\centering
\begin{tabulary}{\linewidth}[t]{TTT}
\sphinxtoprule
\begin{varwidth}[t]{\sphinxcolwidth{1}{3}}
\sphinxstyletheadfamily \sphinxAtStartPar
问题
\sphinxbeforeendvarwidth
\end{varwidth}%
&\begin{varwidth}[t]{\sphinxcolwidth{1}{3}}
\sphinxstyletheadfamily \sphinxAtStartPar
可能原因
\sphinxbeforeendvarwidth
\end{varwidth}%
&\begin{varwidth}[t]{\sphinxcolwidth{1}{3}}
\sphinxstyletheadfamily \sphinxAtStartPar
解决方案
\sphinxbeforeendvarwidth
\end{varwidth}%
\\
\sphinxmidrule
\sphinxtableatstartofbodyhook\begin{varwidth}[t]{\sphinxcolwidth{1}{3}}
\sphinxAtStartPar
\sphinxstylestrong{曲线脱节}
\sphinxbeforeendvarwidth
\end{varwidth}%
&\begin{varwidth}[t]{\sphinxcolwidth{1}{3}}
\sphinxAtStartPar
死频带噪声
\sphinxbeforeendvarwidth
\end{varwidth}%
&\begin{varwidth}[t]{\sphinxcolwidth{1}{3}}
\sphinxAtStartPar
使用RhoPlus MOGA筛选
\sphinxbeforeendvarwidth
\end{varwidth}%
\\
\sphinxhline\begin{varwidth}[t]{\sphinxcolwidth{1}{3}}
\sphinxAtStartPar
\sphinxstylestrong{相位异常}
\sphinxbeforeendvarwidth
\end{varwidth}%
&\begin{varwidth}[t]{\sphinxcolwidth{1}{3}}
\sphinxAtStartPar
近场干扰
\sphinxbeforeendvarwidth
\end{varwidth}%
&\begin{varwidth}[t]{\sphinxcolwidth{1}{3}}
\sphinxAtStartPar
手动剔除异常频点
\sphinxbeforeendvarwidth
\end{varwidth}%
\\
\sphinxhline\begin{varwidth}[t]{\sphinxcolwidth{1}{3}}
\sphinxAtStartPar
\sphinxstylestrong{误差棒过大}
\sphinxbeforeendvarwidth
\end{varwidth}%
&\begin{varwidth}[t]{\sphinxcolwidth{1}{3}}
\sphinxAtStartPar
叠加次数不足
\sphinxbeforeendvarwidth
\end{varwidth}%
&\begin{varwidth}[t]{\sphinxcolwidth{1}{3}}
\sphinxAtStartPar
增加数据时长或合并FC
\sphinxbeforeendvarwidth
\end{varwidth}%
\\
\sphinxhline\begin{varwidth}[t]{\sphinxcolwidth{1}{3}}
\sphinxAtStartPar
\sphinxstylestrong{相干度低}
\sphinxbeforeendvarwidth
\end{varwidth}%
&\begin{varwidth}[t]{\sphinxcolwidth{1}{3}}
\sphinxAtStartPar
人文噪声
\sphinxbeforeendvarwidth
\end{varwidth}%
&\begin{varwidth}[t]{\sphinxcolwidth{1}{3}}
\sphinxAtStartPar
远参考+Robust估计
\sphinxbeforeendvarwidth
\end{varwidth}%
\\
\sphinxhline\begin{varwidth}[t]{\sphinxcolwidth{1}{3}}
\sphinxAtStartPar
\sphinxstylestrong{自动筛选效果差}
\sphinxbeforeendvarwidth
\end{varwidth}%
&\begin{varwidth}[t]{\sphinxcolwidth{1}{3}}
\sphinxAtStartPar
参数设置不当
\sphinxbeforeendvarwidth
\end{varwidth}%
&\begin{varwidth}[t]{\sphinxcolwidth{1}{3}}
\sphinxAtStartPar
降低保留比例或换算法
\sphinxbeforeendvarwidth
\end{varwidth}%
\\
\sphinxhline\begin{varwidth}[t]{\sphinxcolwidth{1}{3}}
\sphinxAtStartPar
\sphinxstylestrong{远参考无效果}
\sphinxbeforeendvarwidth
\end{varwidth}%
&\begin{varwidth}[t]{\sphinxcolwidth{1}{3}}
\sphinxAtStartPar
参考站质量差
\sphinxbeforeendvarwidth
\end{varwidth}%
&\begin{varwidth}[t]{\sphinxcolwidth{1}{3}}
\sphinxAtStartPar
更换参考站或检查时间同步
\sphinxbeforeendvarwidth
\end{varwidth}%
\\
\sphinxbottomrule
\end{tabulary}
\sphinxtableafterendhook\par
\sphinxattableend\end{savenotes}


\bigskip\hrule\bigskip



\subsubsection{5.16.13 参考文献}
\label{\detokenize{chapters/chapter5:id59}}\begin{quote}

\sphinxAtStartPar
{[}1{]} 王培杰, 陈小斌, 韩鹏, 张赟昀. 基于稳健估计、数据筛选和Rhoplus约束的大地电磁数据处理方法. 地球物理学报, 2024, 67(11): 4325\sphinxhyphen{}4342.

\sphinxAtStartPar
{[}2{]} Zhou, C., et al. Application of the Rhoplus method to audio magnetotelluric dead band distortion data. Chinese J. Geophys., 2015.

\sphinxAtStartPar
{[}3{]} Platz, A., \& Weckmann, U. An automated new pre\sphinxhyphen{}selection tool for noisy MT data using the Mahalanobis distance. GJI, 2019.
\end{quote}
\begin{quote}

\sphinxAtStartPar
💡 \sphinxstylestrong{提示}：处理完成后，建议使用"设为测点FC"将结果应用到测点，然后导出EDI文件。
\end{quote}

\sphinxstepscope


\section{🔍 第6章 数据筛选与估计}
\label{\detokenize{chapters/chapter6:id1}}\label{\detokenize{chapters/chapter6::doc}}
\sphinxAtStartPar
本章详细介绍MTDP的数据筛选和阻抗估计功能。


\bigskip\hrule\bigskip



\subsection{6.1 🖥️ 数据筛选界面}
\label{\detokenize{chapters/chapter6:id2}}

\subsubsection{6.1.1 界面布局}
\label{\detokenize{chapters/chapter6:id3}}
\sphinxAtStartPar
数据筛选界面分为多个选项卡：


\begin{savenotes}\sphinxattablestart
\sphinxthistablewithglobalstyle
\centering
\begin{tabulary}{\linewidth}[t]{TT}
\sphinxtoprule
\begin{varwidth}[t]{\sphinxcolwidth{1}{2}}
\sphinxstyletheadfamily \sphinxAtStartPar
选项卡
\sphinxbeforeendvarwidth
\end{varwidth}%
&\begin{varwidth}[t]{\sphinxcolwidth{1}{2}}
\sphinxstyletheadfamily \sphinxAtStartPar
功能
\sphinxbeforeendvarwidth
\end{varwidth}%
\\
\sphinxmidrule
\sphinxtableatstartofbodyhook\begin{varwidth}[t]{\sphinxcolwidth{1}{2}}
\sphinxAtStartPar
视电阻率/相位
\sphinxbeforeendvarwidth
\end{varwidth}%
&\begin{varwidth}[t]{\sphinxcolwidth{1}{2}}
\sphinxAtStartPar
编辑Rho和Phase数据
\sphinxbeforeendvarwidth
\end{varwidth}%
\\
\sphinxhline\begin{varwidth}[t]{\sphinxcolwidth{1}{2}}
\sphinxAtStartPar
阻抗
\sphinxbeforeendvarwidth
\end{varwidth}%
&\begin{varwidth}[t]{\sphinxcolwidth{1}{2}}
\sphinxAtStartPar
编辑阻抗张量
\sphinxbeforeendvarwidth
\end{varwidth}%
\\
\sphinxhline\begin{varwidth}[t]{\sphinxcolwidth{1}{2}}
\sphinxAtStartPar
倾子
\sphinxbeforeendvarwidth
\end{varwidth}%
&\begin{varwidth}[t]{\sphinxcolwidth{1}{2}}
\sphinxAtStartPar
编辑倾子向量
\sphinxbeforeendvarwidth
\end{varwidth}%
\\
\sphinxhline\begin{varwidth}[t]{\sphinxcolwidth{1}{2}}
\sphinxAtStartPar
相位张量
\sphinxbeforeendvarwidth
\end{varwidth}%
&\begin{varwidth}[t]{\sphinxcolwidth{1}{2}}
\sphinxAtStartPar
查看相位张量参数
\sphinxbeforeendvarwidth
\end{varwidth}%
\\
\sphinxhline\begin{varwidth}[t]{\sphinxcolwidth{1}{2}}
\sphinxAtStartPar
相干度
\sphinxbeforeendvarwidth
\end{varwidth}%
&\begin{varwidth}[t]{\sphinxcolwidth{1}{2}}
\sphinxAtStartPar
查看相干度曲线
\sphinxbeforeendvarwidth
\end{varwidth}%
\\
\sphinxhline\begin{varwidth}[t]{\sphinxcolwidth{1}{2}}
\sphinxAtStartPar
信噪比
\sphinxbeforeendvarwidth
\end{varwidth}%
&\begin{varwidth}[t]{\sphinxcolwidth{1}{2}}
\sphinxAtStartPar
查看信噪比曲线
\sphinxbeforeendvarwidth
\end{varwidth}%
\\
\sphinxhline\begin{varwidth}[t]{\sphinxcolwidth{1}{2}}
\sphinxAtStartPar
频谱
\sphinxbeforeendvarwidth
\end{varwidth}%
&\begin{varwidth}[t]{\sphinxcolwidth{1}{2}}
\sphinxAtStartPar
查看功率谱密度
\sphinxbeforeendvarwidth
\end{varwidth}%
\\
\sphinxhline\begin{varwidth}[t]{\sphinxcolwidth{1}{2}}
\sphinxAtStartPar
系统响应
\sphinxbeforeendvarwidth
\end{varwidth}%
&\begin{varwidth}[t]{\sphinxcolwidth{1}{2}}
\sphinxAtStartPar
查看标定曲线
\sphinxbeforeendvarwidth
\end{varwidth}%
\\
\sphinxhline\begin{varwidth}[t]{\sphinxcolwidth{1}{2}}
\sphinxAtStartPar
双参数编辑
\sphinxbeforeendvarwidth
\end{varwidth}%
&\begin{varwidth}[t]{\sphinxcolwidth{1}{2}}
\sphinxAtStartPar
双参数散点图筛选
\sphinxbeforeendvarwidth
\end{varwidth}%
\\
\sphinxhline\begin{varwidth}[t]{\sphinxcolwidth{1}{2}}
\sphinxAtStartPar
自动筛选
\sphinxbeforeendvarwidth
\end{varwidth}%
&\begin{varwidth}[t]{\sphinxcolwidth{1}{2}}
\sphinxAtStartPar
设置自动筛选参数
\sphinxbeforeendvarwidth
\end{varwidth}%
\\
\sphinxbottomrule
\end{tabulary}
\sphinxtableafterendhook\par
\sphinxattableend\end{savenotes}


\subsubsection{5.1.2 工具栏}
\label{\detokenize{chapters/chapter6:id4}}

\begin{savenotes}\sphinxattablestart
\sphinxthistablewithglobalstyle
\centering
\begin{tabulary}{\linewidth}[t]{TT}
\sphinxtoprule
\begin{varwidth}[t]{\sphinxcolwidth{1}{2}}
\sphinxstyletheadfamily \sphinxAtStartPar
按钮
\sphinxbeforeendvarwidth
\end{varwidth}%
&\begin{varwidth}[t]{\sphinxcolwidth{1}{2}}
\sphinxstyletheadfamily \sphinxAtStartPar
功能
\sphinxbeforeendvarwidth
\end{varwidth}%
\\
\sphinxmidrule
\sphinxtableatstartofbodyhook\begin{varwidth}[t]{\sphinxcolwidth{1}{2}}
\sphinxAtStartPar
保存FC
\sphinxbeforeendvarwidth
\end{varwidth}%
&\begin{varwidth}[t]{\sphinxcolwidth{1}{2}}
\sphinxAtStartPar
保存傅里叶系数
\sphinxbeforeendvarwidth
\end{varwidth}%
\\
\sphinxhline\begin{varwidth}[t]{\sphinxcolwidth{1}{2}}
\sphinxAtStartPar
全选
\sphinxbeforeendvarwidth
\end{varwidth}%
&\begin{varwidth}[t]{\sphinxcolwidth{1}{2}}
\sphinxAtStartPar
选中所有数据
\sphinxbeforeendvarwidth
\end{varwidth}%
\\
\sphinxhline\begin{varwidth}[t]{\sphinxcolwidth{1}{2}}
\sphinxAtStartPar
当前频率全选
\sphinxbeforeendvarwidth
\end{varwidth}%
&\begin{varwidth}[t]{\sphinxcolwidth{1}{2}}
\sphinxAtStartPar
选中当前频率的所有数据
\sphinxbeforeendvarwidth
\end{varwidth}%
\\
\sphinxhline\begin{varwidth}[t]{\sphinxcolwidth{1}{2}}
\sphinxAtStartPar
显示误差棒
\sphinxbeforeendvarwidth
\end{varwidth}%
&\begin{varwidth}[t]{\sphinxcolwidth{1}{2}}
\sphinxAtStartPar
显示/隐藏误差棒
\sphinxbeforeendvarwidth
\end{varwidth}%
\\
\sphinxhline\begin{varwidth}[t]{\sphinxcolwidth{1}{2}}
\sphinxAtStartPar
旋转角度
\sphinxbeforeendvarwidth
\end{varwidth}%
&\begin{varwidth}[t]{\sphinxcolwidth{1}{2}}
\sphinxAtStartPar
设置坐标旋转角度
\sphinxbeforeendvarwidth
\end{varwidth}%
\\
\sphinxbottomrule
\end{tabulary}
\sphinxtableafterendhook\par
\sphinxattableend\end{savenotes}


\bigskip\hrule\bigskip



\subsection{6.2 📈 图表类型}
\label{\detokenize{chapters/chapter6:id5}}

\subsubsection{5.2.1 视电阻率/相位}
\label{\detokenize{chapters/chapter6:id6}}

\begin{savenotes}\sphinxattablestart
\sphinxthistablewithglobalstyle
\centering
\begin{tabulary}{\linewidth}[t]{TT}
\sphinxtoprule
\begin{varwidth}[t]{\sphinxcolwidth{1}{2}}
\sphinxstyletheadfamily \sphinxAtStartPar
曲线
\sphinxbeforeendvarwidth
\end{varwidth}%
&\begin{varwidth}[t]{\sphinxcolwidth{1}{2}}
\sphinxstyletheadfamily \sphinxAtStartPar
说明
\sphinxbeforeendvarwidth
\end{varwidth}%
\\
\sphinxmidrule
\sphinxtableatstartofbodyhook\begin{varwidth}[t]{\sphinxcolwidth{1}{2}}
\sphinxAtStartPar
Rxx/Ryy
\sphinxbeforeendvarwidth
\end{varwidth}%
&\begin{varwidth}[t]{\sphinxcolwidth{1}{2}}
\sphinxAtStartPar
对角元素视电阻率
\sphinxbeforeendvarwidth
\end{varwidth}%
\\
\sphinxhline\begin{varwidth}[t]{\sphinxcolwidth{1}{2}}
\sphinxAtStartPar
Rxy/Ryx
\sphinxbeforeendvarwidth
\end{varwidth}%
&\begin{varwidth}[t]{\sphinxcolwidth{1}{2}}
\sphinxAtStartPar
非对角元素视电阻率
\sphinxbeforeendvarwidth
\end{varwidth}%
\\
\sphinxhline\begin{varwidth}[t]{\sphinxcolwidth{1}{2}}
\sphinxAtStartPar
Pxx/Pyy
\sphinxbeforeendvarwidth
\end{varwidth}%
&\begin{varwidth}[t]{\sphinxcolwidth{1}{2}}
\sphinxAtStartPar
对角元素相位
\sphinxbeforeendvarwidth
\end{varwidth}%
\\
\sphinxhline\begin{varwidth}[t]{\sphinxcolwidth{1}{2}}
\sphinxAtStartPar
Pxy/Pyx
\sphinxbeforeendvarwidth
\end{varwidth}%
&\begin{varwidth}[t]{\sphinxcolwidth{1}{2}}
\sphinxAtStartPar
非对角元素相位
\sphinxbeforeendvarwidth
\end{varwidth}%
\\
\sphinxbottomrule
\end{tabulary}
\sphinxtableafterendhook\par
\sphinxattableend\end{savenotes}


\subsubsection{5.2.2 阻抗张量}
\label{\detokenize{chapters/chapter6:id7}}

\begin{savenotes}\sphinxattablestart
\sphinxthistablewithglobalstyle
\centering
\begin{tabulary}{\linewidth}[t]{TT}
\sphinxtoprule
\begin{varwidth}[t]{\sphinxcolwidth{1}{2}}
\sphinxstyletheadfamily \sphinxAtStartPar
曲线
\sphinxbeforeendvarwidth
\end{varwidth}%
&\begin{varwidth}[t]{\sphinxcolwidth{1}{2}}
\sphinxstyletheadfamily \sphinxAtStartPar
说明
\sphinxbeforeendvarwidth
\end{varwidth}%
\\
\sphinxmidrule
\sphinxtableatstartofbodyhook\begin{varwidth}[t]{\sphinxcolwidth{1}{2}}
\sphinxAtStartPar
Z振幅
\sphinxbeforeendvarwidth
\end{varwidth}%
&\begin{varwidth}[t]{\sphinxcolwidth{1}{2}}
\sphinxAtStartPar
阻抗振幅
\sphinxbeforeendvarwidth
\end{varwidth}%
\\
\sphinxhline\begin{varwidth}[t]{\sphinxcolwidth{1}{2}}
\sphinxAtStartPar
Z相位
\sphinxbeforeendvarwidth
\end{varwidth}%
&\begin{varwidth}[t]{\sphinxcolwidth{1}{2}}
\sphinxAtStartPar
阻抗相位
\sphinxbeforeendvarwidth
\end{varwidth}%
\\
\sphinxhline\begin{varwidth}[t]{\sphinxcolwidth{1}{2}}
\sphinxAtStartPar
Z实部
\sphinxbeforeendvarwidth
\end{varwidth}%
&\begin{varwidth}[t]{\sphinxcolwidth{1}{2}}
\sphinxAtStartPar
阻抗实部
\sphinxbeforeendvarwidth
\end{varwidth}%
\\
\sphinxhline\begin{varwidth}[t]{\sphinxcolwidth{1}{2}}
\sphinxAtStartPar
Z虚部
\sphinxbeforeendvarwidth
\end{varwidth}%
&\begin{varwidth}[t]{\sphinxcolwidth{1}{2}}
\sphinxAtStartPar
阻抗虚部
\sphinxbeforeendvarwidth
\end{varwidth}%
\\
\sphinxbottomrule
\end{tabulary}
\sphinxtableafterendhook\par
\sphinxattableend\end{savenotes}


\subsubsection{5.2.3 倾子向量}
\label{\detokenize{chapters/chapter6:id8}}

\begin{savenotes}\sphinxattablestart
\sphinxthistablewithglobalstyle
\centering
\begin{tabulary}{\linewidth}[t]{TT}
\sphinxtoprule
\begin{varwidth}[t]{\sphinxcolwidth{1}{2}}
\sphinxstyletheadfamily \sphinxAtStartPar
曲线
\sphinxbeforeendvarwidth
\end{varwidth}%
&\begin{varwidth}[t]{\sphinxcolwidth{1}{2}}
\sphinxstyletheadfamily \sphinxAtStartPar
说明
\sphinxbeforeendvarwidth
\end{varwidth}%
\\
\sphinxmidrule
\sphinxtableatstartofbodyhook\begin{varwidth}[t]{\sphinxcolwidth{1}{2}}
\sphinxAtStartPar
T振幅
\sphinxbeforeendvarwidth
\end{varwidth}%
&\begin{varwidth}[t]{\sphinxcolwidth{1}{2}}
\sphinxAtStartPar
倾子振幅
\sphinxbeforeendvarwidth
\end{varwidth}%
\\
\sphinxhline\begin{varwidth}[t]{\sphinxcolwidth{1}{2}}
\sphinxAtStartPar
T相位
\sphinxbeforeendvarwidth
\end{varwidth}%
&\begin{varwidth}[t]{\sphinxcolwidth{1}{2}}
\sphinxAtStartPar
倾子相位
\sphinxbeforeendvarwidth
\end{varwidth}%
\\
\sphinxhline\begin{varwidth}[t]{\sphinxcolwidth{1}{2}}
\sphinxAtStartPar
T实部
\sphinxbeforeendvarwidth
\end{varwidth}%
&\begin{varwidth}[t]{\sphinxcolwidth{1}{2}}
\sphinxAtStartPar
倾子实部
\sphinxbeforeendvarwidth
\end{varwidth}%
\\
\sphinxhline\begin{varwidth}[t]{\sphinxcolwidth{1}{2}}
\sphinxAtStartPar
T虚部
\sphinxbeforeendvarwidth
\end{varwidth}%
&\begin{varwidth}[t]{\sphinxcolwidth{1}{2}}
\sphinxAtStartPar
倾子虚部
\sphinxbeforeendvarwidth
\end{varwidth}%
\\
\sphinxbottomrule
\end{tabulary}
\sphinxtableafterendhook\par
\sphinxattableend\end{savenotes}


\subsubsection{5.2.4 相位张量}
\label{\detokenize{chapters/chapter6:id9}}

\begin{savenotes}\sphinxattablestart
\sphinxthistablewithglobalstyle
\centering
\begin{tabulary}{\linewidth}[t]{TT}
\sphinxtoprule
\begin{varwidth}[t]{\sphinxcolwidth{1}{2}}
\sphinxstyletheadfamily \sphinxAtStartPar
曲线
\sphinxbeforeendvarwidth
\end{varwidth}%
&\begin{varwidth}[t]{\sphinxcolwidth{1}{2}}
\sphinxstyletheadfamily \sphinxAtStartPar
说明
\sphinxbeforeendvarwidth
\end{varwidth}%
\\
\sphinxmidrule
\sphinxtableatstartofbodyhook\begin{varwidth}[t]{\sphinxcolwidth{1}{2}}
\sphinxAtStartPar
PT元素
\sphinxbeforeendvarwidth
\end{varwidth}%
&\begin{varwidth}[t]{\sphinxcolwidth{1}{2}}
\sphinxAtStartPar
相位张量四个分量
\sphinxbeforeendvarwidth
\end{varwidth}%
\\
\sphinxhline\begin{varwidth}[t]{\sphinxcolwidth{1}{2}}
\sphinxAtStartPar
α/β
\sphinxbeforeendvarwidth
\end{varwidth}%
&\begin{varwidth}[t]{\sphinxcolwidth{1}{2}}
\sphinxAtStartPar
相位张量主方向角
\sphinxbeforeendvarwidth
\end{varwidth}%
\\
\sphinxhline\begin{varwidth}[t]{\sphinxcolwidth{1}{2}}
\sphinxAtStartPar
Max/Min
\sphinxbeforeendvarwidth
\end{varwidth}%
&\begin{varwidth}[t]{\sphinxcolwidth{1}{2}}
\sphinxAtStartPar
相位张量最大/最小值
\sphinxbeforeendvarwidth
\end{varwidth}%
\\
\sphinxhline\begin{varwidth}[t]{\sphinxcolwidth{1}{2}}
\sphinxAtStartPar
Skew1D/2D
\sphinxbeforeendvarwidth
\end{varwidth}%
&\begin{varwidth}[t]{\sphinxcolwidth{1}{2}}
\sphinxAtStartPar
相位张量偏斜度
\sphinxbeforeendvarwidth
\end{varwidth}%
\\
\sphinxbottomrule
\end{tabulary}
\sphinxtableafterendhook\par
\sphinxattableend\end{savenotes}
\begin{quote}

\sphinxAtStartPar
\sphinxstylestrong{📖 理论背景：相位张量（Phase Tensor）}

\sphinxAtStartPar
相位张量是由Caldwell等（2004）提出的一种分析MT数据的新方法，\sphinxstylestrong{不依赖于地下电性结构的维性假设}。

\sphinxAtStartPar
\sphinxstylestrong{定义}：将阻抗张量分解为实部和虚部 \(\mathbf{Z} = \mathbf{X} + i\mathbf{Y}\)，相位张量定义为：
\begin{equation*}
\begin{split}\mathbf{\Phi} = \mathbf{X}^{-1}\mathbf{Y} = \begin{bmatrix} \Phi_{xx} & \Phi_{xy} \\ \Phi_{yx} & \Phi_{yy} \end{bmatrix}\end{split}
\end{equation*}
\sphinxAtStartPar
\sphinxstylestrong{主要参数：}


\begin{savenotes}\sphinxattablestart
\sphinxthistablewithglobalstyle
\centering
\begin{tabulary}{\linewidth}[t]{TTT}
\sphinxtoprule
\begin{varwidth}[t]{\sphinxcolwidth{1}{3}}
\sphinxstyletheadfamily \sphinxAtStartPar
参数
\sphinxbeforeendvarwidth
\end{varwidth}%
&\begin{varwidth}[t]{\sphinxcolwidth{1}{3}}
\sphinxstyletheadfamily \sphinxAtStartPar
公式
\sphinxbeforeendvarwidth
\end{varwidth}%
&\begin{varwidth}[t]{\sphinxcolwidth{1}{3}}
\sphinxstyletheadfamily \sphinxAtStartPar
物理意义
\sphinxbeforeendvarwidth
\end{varwidth}%
\\
\sphinxmidrule
\sphinxtableatstartofbodyhook\begin{varwidth}[t]{\sphinxcolwidth{1}{3}}
\sphinxAtStartPar
\sphinxstylestrong{主方向角 α}
\sphinxbeforeendvarwidth
\end{varwidth}%
&\begin{varwidth}[t]{\sphinxcolwidth{1}{3}}
\sphinxAtStartPar
\(\alpha = \frac{1}{2}\arctan\left(\frac{\Phi_{xy}+\Phi_{yx}}{\Phi_{xx}-\Phi_{yy}}\right)\)
\sphinxbeforeendvarwidth
\end{varwidth}%
&\begin{varwidth}[t]{\sphinxcolwidth{1}{3}}
\sphinxAtStartPar
电性结构主走向方向
\sphinxbeforeendvarwidth
\end{varwidth}%
\\
\sphinxhline\begin{varwidth}[t]{\sphinxcolwidth{1}{3}}
\sphinxAtStartPar
\sphinxstylestrong{二维偏离度 β}
\sphinxbeforeendvarwidth
\end{varwidth}%
&\begin{varwidth}[t]{\sphinxcolwidth{1}{3}}
\sphinxAtStartPar
\(\beta = \frac{1}{2}\arctan\left(\frac{\Phi_{xy}-\Phi_{yx}}{\Phi_{xx}+\Phi_{yy}}\right)\)
\sphinxbeforeendvarwidth
\end{varwidth}%
&\begin{varwidth}[t]{\sphinxcolwidth{1}{3}}
\sphinxAtStartPar
偏离二维程度，\$
\sphinxbeforeendvarwidth
\end{varwidth}%
\\
\sphinxhline\begin{varwidth}[t]{\sphinxcolwidth{1}{3}}
\sphinxAtStartPar
\sphinxstylestrong{椭圆率 λ}
\sphinxbeforeendvarwidth
\end{varwidth}%
&\begin{varwidth}[t]{\sphinxcolwidth{1}{3}}
\sphinxAtStartPar
\(\lambda = \frac{\Phi_{max} - \Phi_{min}}{\Phi_{max} + \Phi_{min}}\)
\sphinxbeforeendvarwidth
\end{varwidth}%
&\begin{varwidth}[t]{\sphinxcolwidth{1}{3}}
\sphinxAtStartPar
各向异性程度
\sphinxbeforeendvarwidth
\end{varwidth}%
\\
\sphinxhline\begin{varwidth}[t]{\sphinxcolwidth{1}{3}}
\sphinxAtStartPar
\sphinxstylestrong{一维偏离度 κ}
\sphinxbeforeendvarwidth
\end{varwidth}%
&\begin{varwidth}[t]{\sphinxcolwidth{1}{3}}
\sphinxAtStartPar
\(\kappa = \sqrt{\frac{(\Phi_{xx}-\Phi_{yy})^2 + 4\Phi_{xy}^2}{(\Phi_{xx}+\Phi_{yy})^2}}\)
\sphinxbeforeendvarwidth
\end{varwidth}%
&\begin{varwidth}[t]{\sphinxcolwidth{1}{3}}
\sphinxAtStartPar
偏离一维程度
\sphinxbeforeendvarwidth
\end{varwidth}%
\\
\sphinxbottomrule
\end{tabulary}
\sphinxtableafterendhook\par
\sphinxattableend\end{savenotes}

\sphinxAtStartPar
\sphinxstylestrong{椭圆表示}：相位张量可用椭圆可视化，长轴=\(\Phi_{max}\)，短轴=\(\Phi_{min}\)，方向=主方向角\(\alpha\)。对于一维层状介质，相位张量为圆形。
\end{quote}


\subsubsection{5.2.5 相干度}
\label{\detokenize{chapters/chapter6:id10}}
\sphinxAtStartPar
\sphinxstylestrong{常相干度（Ordinary Coherency）：}


\begin{savenotes}\sphinxattablestart
\sphinxthistablewithglobalstyle
\centering
\begin{tabulary}{\linewidth}[t]{TT}
\sphinxtoprule
\begin{varwidth}[t]{\sphinxcolwidth{1}{2}}
\sphinxstyletheadfamily \sphinxAtStartPar
曲线
\sphinxbeforeendvarwidth
\end{varwidth}%
&\begin{varwidth}[t]{\sphinxcolwidth{1}{2}}
\sphinxstyletheadfamily \sphinxAtStartPar
说明
\sphinxbeforeendvarwidth
\end{varwidth}%
\\
\sphinxmidrule
\sphinxtableatstartofbodyhook\begin{varwidth}[t]{\sphinxcolwidth{1}{2}}
\sphinxAtStartPar
CohEx
\sphinxbeforeendvarwidth
\end{varwidth}%
&\begin{varwidth}[t]{\sphinxcolwidth{1}{2}}
\sphinxAtStartPar
Ex道常相干度
\sphinxbeforeendvarwidth
\end{varwidth}%
\\
\sphinxhline\begin{varwidth}[t]{\sphinxcolwidth{1}{2}}
\sphinxAtStartPar
CohEy
\sphinxbeforeendvarwidth
\end{varwidth}%
&\begin{varwidth}[t]{\sphinxcolwidth{1}{2}}
\sphinxAtStartPar
Ey道常相干度
\sphinxbeforeendvarwidth
\end{varwidth}%
\\
\sphinxhline\begin{varwidth}[t]{\sphinxcolwidth{1}{2}}
\sphinxAtStartPar
CohHx
\sphinxbeforeendvarwidth
\end{varwidth}%
&\begin{varwidth}[t]{\sphinxcolwidth{1}{2}}
\sphinxAtStartPar
Hx道常相干度
\sphinxbeforeendvarwidth
\end{varwidth}%
\\
\sphinxhline\begin{varwidth}[t]{\sphinxcolwidth{1}{2}}
\sphinxAtStartPar
CohHy
\sphinxbeforeendvarwidth
\end{varwidth}%
&\begin{varwidth}[t]{\sphinxcolwidth{1}{2}}
\sphinxAtStartPar
Hy道常相干度
\sphinxbeforeendvarwidth
\end{varwidth}%
\\
\sphinxhline\begin{varwidth}[t]{\sphinxcolwidth{1}{2}}
\sphinxAtStartPar
CohHz
\sphinxbeforeendvarwidth
\end{varwidth}%
&\begin{varwidth}[t]{\sphinxcolwidth{1}{2}}
\sphinxAtStartPar
Hz道常相干度
\sphinxbeforeendvarwidth
\end{varwidth}%
\\
\sphinxbottomrule
\end{tabulary}
\sphinxtableafterendhook\par
\sphinxattableend\end{savenotes}

\sphinxAtStartPar
\sphinxstylestrong{多道相干度（Multiple Channel Coherency）：}


\begin{savenotes}\sphinxattablestart
\sphinxthistablewithglobalstyle
\centering
\begin{tabulary}{\linewidth}[t]{TT}
\sphinxtoprule
\begin{varwidth}[t]{\sphinxcolwidth{1}{2}}
\sphinxstyletheadfamily \sphinxAtStartPar
曲线
\sphinxbeforeendvarwidth
\end{varwidth}%
&\begin{varwidth}[t]{\sphinxcolwidth{1}{2}}
\sphinxstyletheadfamily \sphinxAtStartPar
说明
\sphinxbeforeendvarwidth
\end{varwidth}%
\\
\sphinxmidrule
\sphinxtableatstartofbodyhook\begin{varwidth}[t]{\sphinxcolwidth{1}{2}}
\sphinxAtStartPar
CohMEx
\sphinxbeforeendvarwidth
\end{varwidth}%
&\begin{varwidth}[t]{\sphinxcolwidth{1}{2}}
\sphinxAtStartPar
Ex与磁场多道相干
\sphinxbeforeendvarwidth
\end{varwidth}%
\\
\sphinxhline\begin{varwidth}[t]{\sphinxcolwidth{1}{2}}
\sphinxAtStartPar
CohMEy
\sphinxbeforeendvarwidth
\end{varwidth}%
&\begin{varwidth}[t]{\sphinxcolwidth{1}{2}}
\sphinxAtStartPar
Ey与磁场多道相干
\sphinxbeforeendvarwidth
\end{varwidth}%
\\
\sphinxhline\begin{varwidth}[t]{\sphinxcolwidth{1}{2}}
\sphinxAtStartPar
CohMHx
\sphinxbeforeendvarwidth
\end{varwidth}%
&\begin{varwidth}[t]{\sphinxcolwidth{1}{2}}
\sphinxAtStartPar
Hx与磁场多道相干
\sphinxbeforeendvarwidth
\end{varwidth}%
\\
\sphinxhline\begin{varwidth}[t]{\sphinxcolwidth{1}{2}}
\sphinxAtStartPar
CohMHy
\sphinxbeforeendvarwidth
\end{varwidth}%
&\begin{varwidth}[t]{\sphinxcolwidth{1}{2}}
\sphinxAtStartPar
Hy与磁场多道相干
\sphinxbeforeendvarwidth
\end{varwidth}%
\\
\sphinxhline\begin{varwidth}[t]{\sphinxcolwidth{1}{2}}
\sphinxAtStartPar
CohMHz
\sphinxbeforeendvarwidth
\end{varwidth}%
&\begin{varwidth}[t]{\sphinxcolwidth{1}{2}}
\sphinxAtStartPar
Hz与磁场多道相干
\sphinxbeforeendvarwidth
\end{varwidth}%
\\
\sphinxbottomrule
\end{tabulary}
\sphinxtableafterendhook\par
\sphinxattableend\end{savenotes}

\sphinxAtStartPar
\sphinxstylestrong{偏相干度（Partial Coherency）：}


\begin{savenotes}\sphinxattablestart
\sphinxthistablewithglobalstyle
\centering
\begin{tabulary}{\linewidth}[t]{TT}
\sphinxtoprule
\begin{varwidth}[t]{\sphinxcolwidth{1}{2}}
\sphinxstyletheadfamily \sphinxAtStartPar
曲线
\sphinxbeforeendvarwidth
\end{varwidth}%
&\begin{varwidth}[t]{\sphinxcolwidth{1}{2}}
\sphinxstyletheadfamily \sphinxAtStartPar
说明
\sphinxbeforeendvarwidth
\end{varwidth}%
\\
\sphinxmidrule
\sphinxtableatstartofbodyhook\begin{varwidth}[t]{\sphinxcolwidth{1}{2}}
\sphinxAtStartPar
PCohExHx
\sphinxbeforeendvarwidth
\end{varwidth}%
&\begin{varwidth}[t]{\sphinxcolwidth{1}{2}}
\sphinxAtStartPar
Ex\sphinxhyphen{}Hx偏相干
\sphinxbeforeendvarwidth
\end{varwidth}%
\\
\sphinxhline\begin{varwidth}[t]{\sphinxcolwidth{1}{2}}
\sphinxAtStartPar
PCohExHy
\sphinxbeforeendvarwidth
\end{varwidth}%
&\begin{varwidth}[t]{\sphinxcolwidth{1}{2}}
\sphinxAtStartPar
Ex\sphinxhyphen{}Hy偏相干
\sphinxbeforeendvarwidth
\end{varwidth}%
\\
\sphinxhline\begin{varwidth}[t]{\sphinxcolwidth{1}{2}}
\sphinxAtStartPar
PCohEyHx
\sphinxbeforeendvarwidth
\end{varwidth}%
&\begin{varwidth}[t]{\sphinxcolwidth{1}{2}}
\sphinxAtStartPar
Ey\sphinxhyphen{}Hx偏相干
\sphinxbeforeendvarwidth
\end{varwidth}%
\\
\sphinxhline\begin{varwidth}[t]{\sphinxcolwidth{1}{2}}
\sphinxAtStartPar
PCohEyHy
\sphinxbeforeendvarwidth
\end{varwidth}%
&\begin{varwidth}[t]{\sphinxcolwidth{1}{2}}
\sphinxAtStartPar
Ey\sphinxhyphen{}Hy偏相干
\sphinxbeforeendvarwidth
\end{varwidth}%
\\
\sphinxbottomrule
\end{tabulary}
\sphinxtableafterendhook\par
\sphinxattableend\end{savenotes}

\sphinxAtStartPar
\sphinxstylestrong{重相干度（Bi\sphinxhyphen{}coherency）：}


\begin{savenotes}\sphinxattablestart
\sphinxthistablewithglobalstyle
\centering
\begin{tabulary}{\linewidth}[t]{TT}
\sphinxtoprule
\begin{varwidth}[t]{\sphinxcolwidth{1}{2}}
\sphinxstyletheadfamily \sphinxAtStartPar
曲线
\sphinxbeforeendvarwidth
\end{varwidth}%
&\begin{varwidth}[t]{\sphinxcolwidth{1}{2}}
\sphinxstyletheadfamily \sphinxAtStartPar
说明
\sphinxbeforeendvarwidth
\end{varwidth}%
\\
\sphinxmidrule
\sphinxtableatstartofbodyhook\begin{varwidth}[t]{\sphinxcolwidth{1}{2}}
\sphinxAtStartPar
BiCohEx
\sphinxbeforeendvarwidth
\end{varwidth}%
&\begin{varwidth}[t]{\sphinxcolwidth{1}{2}}
\sphinxAtStartPar
Ex重相干度
\sphinxbeforeendvarwidth
\end{varwidth}%
\\
\sphinxhline\begin{varwidth}[t]{\sphinxcolwidth{1}{2}}
\sphinxAtStartPar
BiCohEy
\sphinxbeforeendvarwidth
\end{varwidth}%
&\begin{varwidth}[t]{\sphinxcolwidth{1}{2}}
\sphinxAtStartPar
Ey重相干度
\sphinxbeforeendvarwidth
\end{varwidth}%
\\
\sphinxhline\begin{varwidth}[t]{\sphinxcolwidth{1}{2}}
\sphinxAtStartPar
BiCohHx
\sphinxbeforeendvarwidth
\end{varwidth}%
&\begin{varwidth}[t]{\sphinxcolwidth{1}{2}}
\sphinxAtStartPar
Hx重相干度
\sphinxbeforeendvarwidth
\end{varwidth}%
\\
\sphinxhline\begin{varwidth}[t]{\sphinxcolwidth{1}{2}}
\sphinxAtStartPar
BiCohHy
\sphinxbeforeendvarwidth
\end{varwidth}%
&\begin{varwidth}[t]{\sphinxcolwidth{1}{2}}
\sphinxAtStartPar
Hy重相干度
\sphinxbeforeendvarwidth
\end{varwidth}%
\\
\sphinxbottomrule
\end{tabulary}
\sphinxtableafterendhook\par
\sphinxattableend\end{savenotes}


\subsubsection{5.2.6 其他参数}
\label{\detokenize{chapters/chapter6:id11}}

\begin{savenotes}\sphinxattablestart
\sphinxthistablewithglobalstyle
\centering
\begin{tabulary}{\linewidth}[t]{TT}
\sphinxtoprule
\begin{varwidth}[t]{\sphinxcolwidth{1}{2}}
\sphinxstyletheadfamily \sphinxAtStartPar
曲线
\sphinxbeforeendvarwidth
\end{varwidth}%
&\begin{varwidth}[t]{\sphinxcolwidth{1}{2}}
\sphinxstyletheadfamily \sphinxAtStartPar
说明
\sphinxbeforeendvarwidth
\end{varwidth}%
\\
\sphinxmidrule
\sphinxtableatstartofbodyhook\begin{varwidth}[t]{\sphinxcolwidth{1}{2}}
\sphinxAtStartPar
CCZ Skew
\sphinxbeforeendvarwidth
\end{varwidth}%
&\begin{varwidth}[t]{\sphinxcolwidth{1}{2}}
\sphinxAtStartPar
共轭阻抗偏斜度
\sphinxbeforeendvarwidth
\end{varwidth}%
\\
\sphinxhline\begin{varwidth}[t]{\sphinxcolwidth{1}{2}}
\sphinxAtStartPar
CCZ Theta
\sphinxbeforeendvarwidth
\end{varwidth}%
&\begin{varwidth}[t]{\sphinxcolwidth{1}{2}}
\sphinxAtStartPar
共轭阻抗角度
\sphinxbeforeendvarwidth
\end{varwidth}%
\\
\sphinxhline\begin{varwidth}[t]{\sphinxcolwidth{1}{2}}
\sphinxAtStartPar
SNR
\sphinxbeforeendvarwidth
\end{varwidth}%
&\begin{varwidth}[t]{\sphinxcolwidth{1}{2}}
\sphinxAtStartPar
信噪比
\sphinxbeforeendvarwidth
\end{varwidth}%
\\
\sphinxhline\begin{varwidth}[t]{\sphinxcolwidth{1}{2}}
\sphinxAtStartPar
功率谱
\sphinxbeforeendvarwidth
\end{varwidth}%
&\begin{varwidth}[t]{\sphinxcolwidth{1}{2}}
\sphinxAtStartPar
各道功率谱密度
\sphinxbeforeendvarwidth
\end{varwidth}%
\\
\sphinxhline\begin{varwidth}[t]{\sphinxcolwidth{1}{2}}
\sphinxAtStartPar
系统响应
\sphinxbeforeendvarwidth
\end{varwidth}%
&\begin{varwidth}[t]{\sphinxcolwidth{1}{2}}
\sphinxAtStartPar
标定曲线
\sphinxbeforeendvarwidth
\end{varwidth}%
\\
\sphinxhline\begin{varwidth}[t]{\sphinxcolwidth{1}{2}}
\sphinxAtStartPar
极化方向
\sphinxbeforeendvarwidth
\end{varwidth}%
&\begin{varwidth}[t]{\sphinxcolwidth{1}{2}}
\sphinxAtStartPar
电/磁场极化
\sphinxbeforeendvarwidth
\end{varwidth}%
\\
\sphinxbottomrule
\end{tabulary}
\sphinxtableafterendhook\par
\sphinxattableend\end{savenotes}


\bigskip\hrule\bigskip



\subsection{6.3 筛选操作}
\label{\detokenize{chapters/chapter6:id12}}

\subsubsection{5.3.1 单参数筛选}
\label{\detokenize{chapters/chapter6:id13}}\begin{enumerate}
\sphinxsetlistlabels{\arabic}{enumi}{enumii}{}{.}%
\item {} 
\sphinxAtStartPar
切换到要编辑的参数选项卡

\item {} 
\sphinxAtStartPar
在图表上进行筛选：
\begin{itemize}
\item {} 
\sphinxAtStartPar
点选：点击数据点

\item {} 
\sphinxAtStartPar
矩形选：拖动绘制矩形

\item {} 
\sphinxAtStartPar
多边形选：依次点击顶点

\end{itemize}

\item {} 
\sphinxAtStartPar
按 Enter 确认选择

\end{enumerate}


\subsubsection{5.3.2 双参数筛选}
\label{\detokenize{chapters/chapter6:id14}}\begin{enumerate}
\sphinxsetlistlabels{\arabic}{enumi}{enumii}{}{.}%
\item {} 
\sphinxAtStartPar
切换到"双参数编辑"选项卡

\item {} 
\sphinxAtStartPar
选择X轴和Y轴参数

\item {} 
\sphinxAtStartPar
点击"添加"创建新的筛选图表

\item {} 
\sphinxAtStartPar
在散点图上筛选

\end{enumerate}


\subsubsection{5.3.3 撤销与重做}
\label{\detokenize{chapters/chapter6:id15}}

\begin{savenotes}\sphinxattablestart
\sphinxthistablewithglobalstyle
\centering
\begin{tabulary}{\linewidth}[t]{TT}
\sphinxtoprule
\begin{varwidth}[t]{\sphinxcolwidth{1}{2}}
\sphinxstyletheadfamily \sphinxAtStartPar
快捷键
\sphinxbeforeendvarwidth
\end{varwidth}%
&\begin{varwidth}[t]{\sphinxcolwidth{1}{2}}
\sphinxstyletheadfamily \sphinxAtStartPar
功能
\sphinxbeforeendvarwidth
\end{varwidth}%
\\
\sphinxmidrule
\sphinxtableatstartofbodyhook\begin{varwidth}[t]{\sphinxcolwidth{1}{2}}
\sphinxAtStartPar
Ctrl + Z
\sphinxbeforeendvarwidth
\end{varwidth}%
&\begin{varwidth}[t]{\sphinxcolwidth{1}{2}}
\sphinxAtStartPar
撤销
\sphinxbeforeendvarwidth
\end{varwidth}%
\\
\sphinxhline\begin{varwidth}[t]{\sphinxcolwidth{1}{2}}
\sphinxAtStartPar
Ctrl + Y
\sphinxbeforeendvarwidth
\end{varwidth}%
&\begin{varwidth}[t]{\sphinxcolwidth{1}{2}}
\sphinxAtStartPar
重做
\sphinxbeforeendvarwidth
\end{varwidth}%
\\
\sphinxbottomrule
\end{tabulary}
\sphinxtableafterendhook\par
\sphinxattableend\end{savenotes}


\bigskip\hrule\bigskip



\subsection{6.4 旋转角度}
\label{\detokenize{chapters/chapter6:id16}}

\subsubsection{5.4.1 功能说明}
\label{\detokenize{chapters/chapter6:id17}}
\sphinxAtStartPar
设置坐标旋转角度，用于：
\begin{itemize}
\item {} 
\sphinxAtStartPar
旋转到构造主方向

\item {} 
\sphinxAtStartPar
获取TE/TM模式

\item {} 
\sphinxAtStartPar
最佳主轴角度

\end{itemize}


\subsubsection{5.4.2 操作方法}
\label{\detokenize{chapters/chapter6:id18}}\begin{enumerate}
\sphinxsetlistlabels{\arabic}{enumi}{enumii}{}{.}%
\item {} 
\sphinxAtStartPar
在工具栏中设置旋转角度（度）

\item {} 
\sphinxAtStartPar
或点击"最佳旋转角度"自动计算

\item {} 
\sphinxAtStartPar
曲线自动更新

\end{enumerate}


\bigskip\hrule\bigskip



\subsection{6.5 📏 马氏距离筛选}
\label{\detokenize{chapters/chapter6:id19}}

\subsubsection{5.5.1 功能说明}
\label{\detokenize{chapters/chapter6:id20}}
\sphinxAtStartPar
使用马氏距离自动识别异常数据点。


\subsubsection{5.5.2 操作方法}
\label{\detokenize{chapters/chapter6:id21}}
\sphinxAtStartPar
1️⃣ 设置阈值（默认4.0）
2️⃣ 设置迭代次数（默认5）
3️⃣ 点击"马氏距离筛选"
4️⃣ 查看筛选结果
\begin{quote}

\sphinxAtStartPar
\sphinxstylestrong{📖 理论背景：马氏距离（Mahalanobis Distance）}

\sphinxAtStartPar
马氏距离考虑了变量之间的相关性和各变量的方差，能够更准确地衡量样本的异常程度。

\sphinxAtStartPar
\sphinxstylestrong{定义}：对于p维随机向量\(\mathbf{x}\)，马氏距离定义为：
\begin{equation*}
\begin{split}D_M(\mathbf{x}) = \sqrt{(\mathbf{x} - \boldsymbol{\mu})^T \mathbf{\Sigma}^{-1} (\mathbf{x} - \boldsymbol{\mu})}\end{split}
\end{equation*}
\sphinxAtStartPar
式中\(\boldsymbol{\mu}\)为均值向量，\(\mathbf{\Sigma}\)为协方差矩阵。

\sphinxAtStartPar
\sphinxstylestrong{在MT数据筛选中的应用}：将每个频点的数据视为多维向量，例如：
\begin{equation*}
\begin{split}\mathbf{x}_i = [\log\rho_{xy}, \log\rho_{yx}, \varphi_{xy}, \varphi_{yx}]^T\end{split}
\end{equation*}
\sphinxAtStartPar
计算每个频点的马氏距离，距离较大的频点被认为是离群点。

\sphinxAtStartPar
\sphinxstylestrong{阈值确定方法：}


\begin{savenotes}\sphinxattablestart
\sphinxthistablewithglobalstyle
\centering
\begin{tabulary}{\linewidth}[t]{TT}
\sphinxtoprule
\begin{varwidth}[t]{\sphinxcolwidth{1}{2}}
\sphinxstyletheadfamily \sphinxAtStartPar
方法
\sphinxbeforeendvarwidth
\end{varwidth}%
&\begin{varwidth}[t]{\sphinxcolwidth{1}{2}}
\sphinxstyletheadfamily \sphinxAtStartPar
说明
\sphinxbeforeendvarwidth
\end{varwidth}%
\\
\sphinxmidrule
\sphinxtableatstartofbodyhook\begin{varwidth}[t]{\sphinxcolwidth{1}{2}}
\sphinxAtStartPar
\sphinxstylestrong{卡方分布法}
\sphinxbeforeendvarwidth
\end{varwidth}%
&\begin{varwidth}[t]{\sphinxcolwidth{1}{2}}
\sphinxAtStartPar
\(D_M^2\)服从自由度为p的卡方分布，取\(\chi_{p,1-\alpha}^2\)为阈值
\sphinxbeforeendvarwidth
\end{varwidth}%
\\
\sphinxhline\begin{varwidth}[t]{\sphinxcolwidth{1}{2}}
\sphinxAtStartPar
\sphinxstylestrong{百分位数法}
\sphinxbeforeendvarwidth
\end{varwidth}%
&\begin{varwidth}[t]{\sphinxcolwidth{1}{2}}
\sphinxAtStartPar
选取某个百分位数（如95\%或99\%）作为阈值
\sphinxbeforeendvarwidth
\end{varwidth}%
\\
\sphinxhline\begin{varwidth}[t]{\sphinxcolwidth{1}{2}}
\sphinxAtStartPar
\sphinxstylestrong{迭代法}
\sphinxbeforeendvarwidth
\end{varwidth}%
&\begin{varwidth}[t]{\sphinxcolwidth{1}{2}}
\sphinxAtStartPar
每次剔除最大马氏距离点，直到分布稳定
\sphinxbeforeendvarwidth
\end{varwidth}%
\\
\sphinxbottomrule
\end{tabulary}
\sphinxtableafterendhook\par
\sphinxattableend\end{savenotes}

\sphinxAtStartPar
\sphinxstylestrong{优点}：综合考虑多个参数的相关性，对变量尺度不敏感，比单一参数筛选更加稳健。
\end{quote}


\bigskip\hrule\bigskip



\subsection{6.6 🤖 自动筛选}
\label{\detokenize{chapters/chapter6:id22}}

\subsubsection{5.6.1 筛选参数}
\label{\detokenize{chapters/chapter6:id23}}

\begin{savenotes}\sphinxattablestart
\sphinxthistablewithglobalstyle
\centering
\begin{tabulary}{\linewidth}[t]{TT}
\sphinxtoprule
\begin{varwidth}[t]{\sphinxcolwidth{1}{2}}
\sphinxstyletheadfamily \sphinxAtStartPar
参数
\sphinxbeforeendvarwidth
\end{varwidth}%
&\begin{varwidth}[t]{\sphinxcolwidth{1}{2}}
\sphinxstyletheadfamily \sphinxAtStartPar
说明
\sphinxbeforeendvarwidth
\end{varwidth}%
\\
\sphinxmidrule
\sphinxtableatstartofbodyhook\begin{varwidth}[t]{\sphinxcolwidth{1}{2}}
\sphinxAtStartPar
最大保留比例
\sphinxbeforeendvarwidth
\end{varwidth}%
&\begin{varwidth}[t]{\sphinxcolwidth{1}{2}}
\sphinxAtStartPar
保留数据的最大比例
\sphinxbeforeendvarwidth
\end{varwidth}%
\\
\sphinxhline\begin{varwidth}[t]{\sphinxcolwidth{1}{2}}
\sphinxAtStartPar
最小保留比例
\sphinxbeforeendvarwidth
\end{varwidth}%
&\begin{varwidth}[t]{\sphinxcolwidth{1}{2}}
\sphinxAtStartPar
保留数据的最小比例
\sphinxbeforeendvarwidth
\end{varwidth}%
\\
\sphinxhline\begin{varwidth}[t]{\sphinxcolwidth{1}{2}}
\sphinxAtStartPar
退出误差
\sphinxbeforeendvarwidth
\end{varwidth}%
&\begin{varwidth}[t]{\sphinxcolwidth{1}{2}}
\sphinxAtStartPar
迭代退出的误差阈值
\sphinxbeforeendvarwidth
\end{varwidth}%
\\
\sphinxhline\begin{varwidth}[t]{\sphinxcolwidth{1}{2}}
\sphinxAtStartPar
自动百分比
\sphinxbeforeendvarwidth
\end{varwidth}%
&\begin{varwidth}[t]{\sphinxcolwidth{1}{2}}
\sphinxAtStartPar
自动筛选的百分比
\sphinxbeforeendvarwidth
\end{varwidth}%
\\
\sphinxbottomrule
\end{tabulary}
\sphinxtableafterendhook\par
\sphinxattableend\end{savenotes}


\subsubsection{5.6.2 自动旋转角度}
\label{\detokenize{chapters/chapter6:id24}}
\sphinxAtStartPar
可设置多个旋转角度进行自动筛选：
\begin{enumerate}
\sphinxsetlistlabels{\arabic}{enumi}{enumii}{}{.}%
\item {} 
\sphinxAtStartPar
在"自动旋转角度"列表中添加角度

\item {} 
\sphinxAtStartPar
系统在各角度下分别筛选

\item {} 
\sphinxAtStartPar
选择最优结果

\end{enumerate}


\subsubsection{5.6.3 筛选类型}
\label{\detokenize{chapters/chapter6:id25}}

\begin{savenotes}\sphinxattablestart
\sphinxthistablewithglobalstyle
\centering
\begin{tabulary}{\linewidth}[t]{TT}
\sphinxtoprule
\begin{varwidth}[t]{\sphinxcolwidth{1}{2}}
\sphinxstyletheadfamily \sphinxAtStartPar
类型
\sphinxbeforeendvarwidth
\end{varwidth}%
&\begin{varwidth}[t]{\sphinxcolwidth{1}{2}}
\sphinxstyletheadfamily \sphinxAtStartPar
说明
\sphinxbeforeendvarwidth
\end{varwidth}%
\\
\sphinxmidrule
\sphinxtableatstartofbodyhook\begin{varwidth}[t]{\sphinxcolwidth{1}{2}}
\sphinxAtStartPar
Rhoplus SA
\sphinxbeforeendvarwidth
\end{varwidth}%
&\begin{varwidth}[t]{\sphinxcolwidth{1}{2}}
\sphinxAtStartPar
Rhoplus谱分析筛选
\sphinxbeforeendvarwidth
\end{varwidth}%
\\
\sphinxhline\begin{varwidth}[t]{\sphinxcolwidth{1}{2}}
\sphinxAtStartPar
Rhoplus GA
\sphinxbeforeendvarwidth
\end{varwidth}%
&\begin{varwidth}[t]{\sphinxcolwidth{1}{2}}
\sphinxAtStartPar
Rhoplus遗传算法筛选
\sphinxbeforeendvarwidth
\end{varwidth}%
\\
\sphinxhline\begin{varwidth}[t]{\sphinxcolwidth{1}{2}}
\sphinxAtStartPar
Rhoplus MOGA
\sphinxbeforeendvarwidth
\end{varwidth}%
&\begin{varwidth}[t]{\sphinxcolwidth{1}{2}}
\sphinxAtStartPar
Rhoplus多目标遗传算法
\sphinxbeforeendvarwidth
\end{varwidth}%
\\
\sphinxhline\begin{varwidth}[t]{\sphinxcolwidth{1}{2}}
\sphinxAtStartPar
Full Z MOGA
\sphinxbeforeendvarwidth
\end{varwidth}%
&\begin{varwidth}[t]{\sphinxcolwidth{1}{2}}
\sphinxAtStartPar
全阻抗多目标遗传算法
\sphinxbeforeendvarwidth
\end{varwidth}%
\\
\sphinxbottomrule
\end{tabulary}
\sphinxtableafterendhook\par
\sphinxattableend\end{savenotes}


\bigskip\hrule\bigskip



\subsection{6.7 📐 Rhoplus参考曲线}
\label{\detokenize{chapters/chapter6:rhoplus}}

\subsubsection{5.7.1 功能说明}
\label{\detokenize{chapters/chapter6:id26}}
\sphinxAtStartPar
Rhoplus是一种基于1D层状模型的反演方法，用于：
\begin{itemize}
\item {} 
\sphinxAtStartPar
数据质量评估

\item {} 
\sphinxAtStartPar
生成平滑的参考曲线

\item {} 
\sphinxAtStartPar
识别异常数据点

\end{itemize}
\begin{quote}

\sphinxAtStartPar
\sphinxstylestrong{📖 理论背景：Rhoplus方法}

\sphinxAtStartPar
Rhoplus方法由Parker和Booker（1996）提出，核心思想是\sphinxstylestrong{寻找与观测数据相容的最简单地下电阻率结构}。

\sphinxAtStartPar
\sphinxstylestrong{基本原理}：如果存在一个一维层状模型能够完全解释观测的视电阻率数据，那么该模型就是与数据相容的最简单模型。Rhoplus方法寻找的是使拟合误差最小的一维模型。

\sphinxAtStartPar
\sphinxstylestrong{目标函数}：
\begin{equation*}
\begin{split}\min_{\boldsymbol{\rho}, \boldsymbol{h}} \sum_{i=1}^{N} w_i \left[\frac{\rho_a^{obs}(\omega_i) - \rho_a^{cal}(\omega_i, \boldsymbol{\rho}, \boldsymbol{h})}{\delta\rho_a(\omega_i)}\right]^2\end{split}
\end{equation*}
\sphinxAtStartPar
式中：\(\rho_a^{obs}\)为观测视电阻率，\(\rho_a^{cal}\)为计算视电阻率，\(\delta\rho_a\)为观测误差，\(\boldsymbol{\rho}\)为各层电阻率，\(\boldsymbol{h}\)为各层厚度。

\sphinxAtStartPar
\sphinxstylestrong{多角度Rhoplus}：对于二维或三维结构，不同方向的视电阻率曲线不同。通过旋转坐标系计算不同角度的Rhoplus响应：
\begin{equation*}
\begin{split}\rho_a(\theta, \omega) = \frac{|Z_{xy}\cos^2\theta + (Z_{yy}-Z_{xx})\sin\theta\cos\theta - Z_{yx}\sin^2\theta|^2}{\omega\mu_0}\end{split}
\end{equation*}
\sphinxAtStartPar
\sphinxstylestrong{主要应用：}


\begin{savenotes}\sphinxattablestart
\sphinxthistablewithglobalstyle
\centering
\begin{tabulary}{\linewidth}[t]{TT}
\sphinxtoprule
\begin{varwidth}[t]{\sphinxcolwidth{1}{2}}
\sphinxstyletheadfamily \sphinxAtStartPar
应用
\sphinxbeforeendvarwidth
\end{varwidth}%
&\begin{varwidth}[t]{\sphinxcolwidth{1}{2}}
\sphinxstyletheadfamily \sphinxAtStartPar
说明
\sphinxbeforeendvarwidth
\end{varwidth}%
\\
\sphinxmidrule
\sphinxtableatstartofbodyhook\begin{varwidth}[t]{\sphinxcolwidth{1}{2}}
\sphinxAtStartPar
\sphinxstylestrong{数据质量评估}
\sphinxbeforeendvarwidth
\end{varwidth}%
&\begin{varwidth}[t]{\sphinxcolwidth{1}{2}}
\sphinxAtStartPar
如果数据能被一维Rhoplus很好拟合，说明数据质量高
\sphinxbeforeendvarwidth
\end{varwidth}%
\\
\sphinxhline\begin{varwidth}[t]{\sphinxcolwidth{1}{2}}
\sphinxAtStartPar
\sphinxstylestrong{维性分析}
\sphinxbeforeendvarwidth
\end{varwidth}%
&\begin{varwidth}[t]{\sphinxcolwidth{1}{2}}
\sphinxAtStartPar
比较不同方向Rhoplus响应差异，判断维性特征
\sphinxbeforeendvarwidth
\end{varwidth}%
\\
\sphinxhline\begin{varwidth}[t]{\sphinxcolwidth{1}{2}}
\sphinxAtStartPar
\sphinxstylestrong{静态位移校正}
\sphinxbeforeendvarwidth
\end{varwidth}%
&\begin{varwidth}[t]{\sphinxcolwidth{1}{2}}
\sphinxAtStartPar
检测近地表不均匀体引起的静态位移效应
\sphinxbeforeendvarwidth
\end{varwidth}%
\\
\sphinxhline\begin{varwidth}[t]{\sphinxcolwidth{1}{2}}
\sphinxAtStartPar
\sphinxstylestrong{初始模型构建}
\sphinxbeforeendvarwidth
\end{varwidth}%
&\begin{varwidth}[t]{\sphinxcolwidth{1}{2}}
\sphinxAtStartPar
Rhoplus反演结果可作为2D/3D反演的初始模型
\sphinxbeforeendvarwidth
\end{varwidth}%
\\
\sphinxbottomrule
\end{tabulary}
\sphinxtableafterendhook\par
\sphinxattableend\end{savenotes}
\end{quote}


\subsubsection{5.7.2 Rhoplus计算方式}
\label{\detokenize{chapters/chapter6:id27}}

\begin{savenotes}\sphinxattablestart
\sphinxthistablewithglobalstyle
\centering
\begin{tabulary}{\linewidth}[t]{TT}
\sphinxtoprule
\begin{varwidth}[t]{\sphinxcolwidth{1}{2}}
\sphinxstyletheadfamily \sphinxAtStartPar
方法
\sphinxbeforeendvarwidth
\end{varwidth}%
&\begin{varwidth}[t]{\sphinxcolwidth{1}{2}}
\sphinxstyletheadfamily \sphinxAtStartPar
说明
\sphinxbeforeendvarwidth
\end{varwidth}%
\\
\sphinxmidrule
\sphinxtableatstartofbodyhook\begin{varwidth}[t]{\sphinxcolwidth{1}{2}}
\sphinxAtStartPar
CalRhoplus
\sphinxbeforeendvarwidth
\end{varwidth}%
&\begin{varwidth}[t]{\sphinxcolwidth{1}{2}}
\sphinxAtStartPar
XY和YX同时反演
\sphinxbeforeendvarwidth
\end{varwidth}%
\\
\sphinxhline\begin{varwidth}[t]{\sphinxcolwidth{1}{2}}
\sphinxAtStartPar
CalRhoplusXY
\sphinxbeforeendvarwidth
\end{varwidth}%
&\begin{varwidth}[t]{\sphinxcolwidth{1}{2}}
\sphinxAtStartPar
仅XY模式反演
\sphinxbeforeendvarwidth
\end{varwidth}%
\\
\sphinxhline\begin{varwidth}[t]{\sphinxcolwidth{1}{2}}
\sphinxAtStartPar
CalRhoplusYX
\sphinxbeforeendvarwidth
\end{varwidth}%
&\begin{varwidth}[t]{\sphinxcolwidth{1}{2}}
\sphinxAtStartPar
仅YX模式反演
\sphinxbeforeendvarwidth
\end{varwidth}%
\\
\sphinxbottomrule
\end{tabulary}
\sphinxtableafterendhook\par
\sphinxattableend\end{savenotes}


\subsubsection{5.7.3 操作方法}
\label{\detokenize{chapters/chapter6:id28}}\begin{enumerate}
\sphinxsetlistlabels{\arabic}{enumi}{enumii}{}{.}%
\item {} 
\sphinxAtStartPar
勾选"自动更新Rhoplus"

\item {} 
\sphinxAtStartPar
系统自动计算参考曲线

\item {} 
\sphinxAtStartPar
曲线叠加显示在图表上

\end{enumerate}


\subsubsection{5.7.4 输出结果}
\label{\detokenize{chapters/chapter6:id29}}
\sphinxAtStartPar
Rhoplus输出包括：
\begin{itemize}
\item {} 
\sphinxAtStartPar
平滑后的视电阻率曲线

\item {} 
\sphinxAtStartPar
平滑后的相位曲线

\item {} 
\sphinxAtStartPar
1D层状模型参数（层数、层厚、电阻率）

\end{itemize}


\subsubsection{5.7.5 添加预测曲线}
\label{\detokenize{chapters/chapter6:id30}}\begin{enumerate}
\sphinxsetlistlabels{\arabic}{enumi}{enumii}{}{.}%
\item {} 
\sphinxAtStartPar
点击"添加Rhoplus角度"

\item {} 
\sphinxAtStartPar
选择预测模型

\item {} 
\sphinxAtStartPar
预测曲线显示在图表上

\end{enumerate}


\bigskip\hrule\bigskip



\subsection{6.8 🧠 AI模型预测}
\label{\detokenize{chapters/chapter6:ai}}

\subsubsection{5.8.1 功能说明}
\label{\detokenize{chapters/chapter6:id31}}
\sphinxAtStartPar
使用深度学习模型预测完整的阻抗张量，用于：
\begin{itemize}
\item {} 
\sphinxAtStartPar
填补缺失数据

\item {} 
\sphinxAtStartPar
识别异常频点

\item {} 
\sphinxAtStartPar
提供参考曲线

\end{itemize}


\subsubsection{5.8.2 预测方法}
\label{\detokenize{chapters/chapter6:id32}}

\begin{savenotes}\sphinxattablestart
\sphinxthistablewithglobalstyle
\centering
\begin{tabulary}{\linewidth}[t]{TTT}
\sphinxtoprule
\begin{varwidth}[t]{\sphinxcolwidth{1}{3}}
\sphinxstyletheadfamily \sphinxAtStartPar
方法
\sphinxbeforeendvarwidth
\end{varwidth}%
&\begin{varwidth}[t]{\sphinxcolwidth{1}{3}}
\sphinxstyletheadfamily \sphinxAtStartPar
算法
\sphinxbeforeendvarwidth
\end{varwidth}%
&\begin{varwidth}[t]{\sphinxcolwidth{1}{3}}
\sphinxstyletheadfamily \sphinxAtStartPar
适用场景
\sphinxbeforeendvarwidth
\end{varwidth}%
\\
\sphinxmidrule
\sphinxtableatstartofbodyhook\begin{varwidth}[t]{\sphinxcolwidth{1}{3}}
\sphinxAtStartPar
ModelPredict
\sphinxbeforeendvarwidth
\end{varwidth}%
&\begin{varwidth}[t]{\sphinxcolwidth{1}{3}}
\sphinxAtStartPar
ZPredict深度学习模型
\sphinxbeforeendvarwidth
\end{varwidth}%
&\begin{varwidth}[t]{\sphinxcolwidth{1}{3}}
\sphinxAtStartPar
一般数据
\sphinxbeforeendvarwidth
\end{varwidth}%
\\
\sphinxhline\begin{varwidth}[t]{\sphinxcolwidth{1}{3}}
\sphinxAtStartPar
ModelPredict1
\sphinxbeforeendvarwidth
\end{varwidth}%
&\begin{varwidth}[t]{\sphinxcolwidth{1}{3}}
\sphinxAtStartPar
中值滤波方法
\sphinxbeforeendvarwidth
\end{varwidth}%
&\begin{varwidth}[t]{\sphinxcolwidth{1}{3}}
\sphinxAtStartPar
简单平滑
\sphinxbeforeendvarwidth
\end{varwidth}%
\\
\sphinxhline\begin{varwidth}[t]{\sphinxcolwidth{1}{3}}
\sphinxAtStartPar
ModelPredict2
\sphinxbeforeendvarwidth
\end{varwidth}%
&\begin{varwidth}[t]{\sphinxcolwidth{1}{3}}
\sphinxAtStartPar
迭代优化方法
\sphinxbeforeendvarwidth
\end{varwidth}%
&\begin{varwidth}[t]{\sphinxcolwidth{1}{3}}
\sphinxAtStartPar
精细预测
\sphinxbeforeendvarwidth
\end{varwidth}%
\\
\sphinxbottomrule
\end{tabulary}
\sphinxtableafterendhook\par
\sphinxattableend\end{savenotes}


\subsubsection{5.8.3 操作方法}
\label{\detokenize{chapters/chapter6:id33}}\begin{enumerate}
\sphinxsetlistlabels{\arabic}{enumi}{enumii}{}{.}%
\item {} 
\sphinxAtStartPar
点击"AI预测"按钮

\item {} 
\sphinxAtStartPar
选择预测方法

\item {} 
\sphinxAtStartPar
系统计算预测阻抗

\item {} 
\sphinxAtStartPar
预测曲线叠加显示

\end{enumerate}


\subsubsection{5.8.4 使用建议}
\label{\detokenize{chapters/chapter6:id34}}\begin{itemize}
\item {} 
\sphinxAtStartPar
先进行基本筛选再使用AI预测

\item {} 
\sphinxAtStartPar
对比预测结果与实际数据

\item {} 
\sphinxAtStartPar
异常频点可参考预测值进行判断

\end{itemize}


\bigskip\hrule\bigskip



\subsection{6.9 理论曲线预测}
\label{\detokenize{chapters/chapter6:id35}}

\subsubsection{5.9.1 功能说明}
\label{\detokenize{chapters/chapter6:id36}}
\sphinxAtStartPar
根据正演模型计算理论MT响应曲线。


\subsubsection{5.9.2 操作方法}
\label{\detokenize{chapters/chapter6:id37}}\begin{enumerate}
\sphinxsetlistlabels{\arabic}{enumi}{enumii}{}{.}%
\item {} 
\sphinxAtStartPar
点击"曲线预测"按钮

\item {} 
\sphinxAtStartPar
设置地电模型参数

\item {} 
\sphinxAtStartPar
系统计算理论曲线

\item {} 
\sphinxAtStartPar
理论曲线叠加显示

\end{enumerate}


\subsubsection{5.9.3 导入/导出预测}
\label{\detokenize{chapters/chapter6:id38}}\begin{itemize}
\item {} 
\sphinxAtStartPar
点击"加载预测"导入预测模型

\item {} 
\sphinxAtStartPar
点击"保存预测"保存当前预测

\end{itemize}


\bigskip\hrule\bigskip



\subsection{6.10 🎯 Robust稳健估计}
\label{\detokenize{chapters/chapter6:robust}}

\subsubsection{5.9.1 可用方法}
\label{\detokenize{chapters/chapter6:id39}}

\begin{savenotes}\sphinxattablestart
\sphinxthistablewithglobalstyle
\centering
\begin{tabulary}{\linewidth}[t]{TT}
\sphinxtoprule
\begin{varwidth}[t]{\sphinxcolwidth{1}{2}}
\sphinxstyletheadfamily \sphinxAtStartPar
方法
\sphinxbeforeendvarwidth
\end{varwidth}%
&\begin{varwidth}[t]{\sphinxcolwidth{1}{2}}
\sphinxstyletheadfamily \sphinxAtStartPar
适用场景
\sphinxbeforeendvarwidth
\end{varwidth}%
\\
\sphinxmidrule
\sphinxtableatstartofbodyhook\begin{varwidth}[t]{\sphinxcolwidth{1}{2}}
\sphinxAtStartPar
最小二乘法
\sphinxbeforeendvarwidth
\end{varwidth}%
&\begin{varwidth}[t]{\sphinxcolwidth{1}{2}}
\sphinxAtStartPar
高质量数据
\sphinxbeforeendvarwidth
\end{varwidth}%
\\
\sphinxhline\begin{varwidth}[t]{\sphinxcolwidth{1}{2}}
\sphinxAtStartPar
Regression\sphinxhyphen{}M
\sphinxbeforeendvarwidth
\end{varwidth}%
&\begin{varwidth}[t]{\sphinxcolwidth{1}{2}}
\sphinxAtStartPar
一般数据（推荐）
\sphinxbeforeendvarwidth
\end{varwidth}%
\\
\sphinxhline\begin{varwidth}[t]{\sphinxcolwidth{1}{2}}
\sphinxAtStartPar
重复中位数法
\sphinxbeforeendvarwidth
\end{varwidth}%
&\begin{varwidth}[t]{\sphinxcolwidth{1}{2}}
\sphinxAtStartPar
强干扰环境
\sphinxbeforeendvarwidth
\end{varwidth}%
\\
\sphinxhline\begin{varwidth}[t]{\sphinxcolwidth{1}{2}}
\sphinxAtStartPar
Robust+AI
\sphinxbeforeendvarwidth
\end{varwidth}%
&\begin{varwidth}[t]{\sphinxcolwidth{1}{2}}
\sphinxAtStartPar
复杂噪声
\sphinxbeforeendvarwidth
\end{varwidth}%
\\
\sphinxbottomrule
\end{tabulary}
\sphinxtableafterendhook\par
\sphinxattableend\end{savenotes}


\subsubsection{5.9.2 参数设置}
\label{\detokenize{chapters/chapter6:id40}}\begin{enumerate}
\sphinxsetlistlabels{\arabic}{enumi}{enumii}{}{.}%
\item {} 
\sphinxAtStartPar
点击"MD参数"按钮

\item {} 
\sphinxAtStartPar
选择估计方法

\item {} 
\sphinxAtStartPar
调整参数（通常使用默认值）

\item {} 
\sphinxAtStartPar
应用设置

\end{enumerate}


\bigskip\hrule\bigskip



\subsection{6.11 时轴显示}
\label{\detokenize{chapters/chapter6:id41}}

\subsubsection{5.10.1 功能说明}
\label{\detokenize{chapters/chapter6:id42}}
\sphinxAtStartPar
按采集时间显示数据，用于识别时间相关的噪声。


\subsubsection{5.10.2 操作方法}
\label{\detokenize{chapters/chapter6:id43}}\begin{enumerate}
\sphinxsetlistlabels{\arabic}{enumi}{enumii}{}{.}%
\item {} 
\sphinxAtStartPar
勾选"时轴"选项

\item {} 
\sphinxAtStartPar
设置时区偏移

\item {} 
\sphinxAtStartPar
数据按时间顺序显示

\end{enumerate}


\bigskip\hrule\bigskip



\subsection{6.12 传递调度}
\label{\detokenize{chapters/chapter6:id44}}

\subsubsection{5.11.1 功能说明}
\label{\detokenize{chapters/chapter6:id45}}
\sphinxAtStartPar
查看和管理测点的传递函数估计结果。


\subsubsection{5.11.2 操作方法}
\label{\detokenize{chapters/chapter6:id46}}\begin{itemize}
\item {} 
\sphinxAtStartPar
选择要保留的数据

\item {} 
\sphinxAtStartPar
选择要剔除的数据

\end{itemize}


\bigskip\hrule\bigskip



\subsection{6.13 🚀 实用处理流程}
\label{\detokenize{chapters/chapter6:id47}}

\subsubsection{5.12.1 👶 新手推荐}
\label{\detokenize{chapters/chapter6:id48}}
\sphinxAtStartPar
1️⃣ 查看视电阻率/相位曲线
2️⃣ 应用Robust估计
3️⃣ 手动剔除明显异常点
4️⃣ 保存结果


\subsubsection{5.12.2 ⭐ 高质量数据}
\label{\detokenize{chapters/chapter6:id49}}
\sphinxAtStartPar
1️⃣ 查看相干度曲线
2️⃣ 设置旋转角度
3️⃣ 启用Rhoplus参考曲线
4️⃣ 精细筛选
5️⃣ 保存结果


\subsubsection{5.12.3 ⚠️ 强干扰环境}
\label{\detokenize{chapters/chapter6:id50}}
\sphinxAtStartPar
1️⃣ 查看相干度和SNR
2️⃣ 使用马氏距离筛选
3️⃣ 应用自动筛选
4️⃣ 多角度对比
5️⃣ 人工审核
6️⃣ 保存结果


\subsection{6.14 XPR分组设置}
\label{\detokenize{chapters/chapter6:xpr}}
\sphinxAtStartPar
XPR（交叉功率比）分组用于确定FFT窗口如何组合计算传递函数。


\subsubsection{5.13.1 分组类型}
\label{\detokenize{chapters/chapter6:id51}}

\begin{savenotes}\sphinxattablestart
\sphinxthistablewithglobalstyle
\centering
\begin{tabulary}{\linewidth}[t]{TTT}
\sphinxtoprule
\begin{varwidth}[t]{\sphinxcolwidth{1}{3}}
\sphinxstyletheadfamily \sphinxAtStartPar
类型
\sphinxbeforeendvarwidth
\end{varwidth}%
&\begin{varwidth}[t]{\sphinxcolwidth{1}{3}}
\sphinxstyletheadfamily \sphinxAtStartPar
名称
\sphinxbeforeendvarwidth
\end{varwidth}%
&\begin{varwidth}[t]{\sphinxcolwidth{1}{3}}
\sphinxstyletheadfamily \sphinxAtStartPar
说明
\sphinxbeforeendvarwidth
\end{varwidth}%
\\
\sphinxmidrule
\sphinxtableatstartofbodyhook\begin{varwidth}[t]{\sphinxcolwidth{1}{3}}
\sphinxAtStartPar
0
\sphinxbeforeendvarwidth
\end{varwidth}%
&\begin{varwidth}[t]{\sphinxcolwidth{1}{3}}
\sphinxAtStartPar
平均
\sphinxbeforeendvarwidth
\end{varwidth}%
&\begin{varwidth}[t]{\sphinxcolwidth{1}{3}}
\sphinxAtStartPar
所有窗口等权平均计算
\sphinxbeforeendvarwidth
\end{varwidth}%
\\
\sphinxhline\begin{varwidth}[t]{\sphinxcolwidth{1}{3}}
\sphinxAtStartPar
1
\sphinxbeforeendvarwidth
\end{varwidth}%
&\begin{varwidth}[t]{\sphinxcolwidth{1}{3}}
\sphinxAtStartPar
排序
\sphinxbeforeendvarwidth
\end{varwidth}%
&\begin{varwidth}[t]{\sphinxcolwidth{1}{3}}
\sphinxAtStartPar
按XPR值排序后加权计算
\sphinxbeforeendvarwidth
\end{varwidth}%
\\
\sphinxhline\begin{varwidth}[t]{\sphinxcolwidth{1}{3}}
\sphinxAtStartPar
2
\sphinxbeforeendvarwidth
\end{varwidth}%
&\begin{varwidth}[t]{\sphinxcolwidth{1}{3}}
\sphinxAtStartPar
随机
\sphinxbeforeendvarwidth
\end{varwidth}%
&\begin{varwidth}[t]{\sphinxcolwidth{1}{3}}
\sphinxAtStartPar
随机选取窗口组合
\sphinxbeforeendvarwidth
\end{varwidth}%
\\
\sphinxbottomrule
\end{tabulary}
\sphinxtableafterendhook\par
\sphinxattableend\end{savenotes}


\subsubsection{5.13.2 MaxXPR参数}
\label{\detokenize{chapters/chapter6:maxxpr}}\begin{itemize}
\item {} 
\sphinxAtStartPar
\sphinxstylestrong{说明}: 最大允许的XPR值

\item {} 
\sphinxAtStartPar
\sphinxstylestrong{默认值}: 100

\item {} 
\sphinxAtStartPar
\sphinxstylestrong{作用}: 限制参与计算的窗口XPR范围，剔除XPR过大的异常窗口

\end{itemize}


\subsubsection{5.13.3 使用建议}
\label{\detokenize{chapters/chapter6:id52}}\begin{itemize}
\item {} 
\sphinxAtStartPar
\sphinxstylestrong{平均模式}: 适用于高质量数据，计算速度快

\item {} 
\sphinxAtStartPar
\sphinxstylestrong{排序模式}: 适用于中等质量数据，可提高估计精度

\item {} 
\sphinxAtStartPar
\sphinxstylestrong{随机模式}: 用于统计检验和误差估计

\end{itemize}


\bigskip\hrule\bigskip



\subsection{6.15 🌐 远参考站管理}
\label{\detokenize{chapters/chapter6:id53}}

\subsubsection{5.14.1 功能说明}
\label{\detokenize{chapters/chapter6:id54}}
\sphinxAtStartPar
远参考（Remote Reference）是一种提高MT数据质量的技术，通过使用远离测点的参考站数据来降低噪声影响。


\subsubsection{5.14.2 适用条件}
\label{\detokenize{chapters/chapter6:id55}}

\begin{savenotes}\sphinxattablestart
\sphinxthistablewithglobalstyle
\centering
\begin{tabulary}{\linewidth}[t]{TT}
\sphinxtoprule
\begin{varwidth}[t]{\sphinxcolwidth{1}{2}}
\sphinxstyletheadfamily \sphinxAtStartPar
条件
\sphinxbeforeendvarwidth
\end{varwidth}%
&\begin{varwidth}[t]{\sphinxcolwidth{1}{2}}
\sphinxstyletheadfamily \sphinxAtStartPar
要求
\sphinxbeforeendvarwidth
\end{varwidth}%
\\
\sphinxmidrule
\sphinxtableatstartofbodyhook\begin{varwidth}[t]{\sphinxcolwidth{1}{2}}
\sphinxAtStartPar
参考站距离
\sphinxbeforeendvarwidth
\end{varwidth}%
&\begin{varwidth}[t]{\sphinxcolwidth{1}{2}}
\sphinxAtStartPar
通常 > 10km
\sphinxbeforeendvarwidth
\end{varwidth}%
\\
\sphinxhline\begin{varwidth}[t]{\sphinxcolwidth{1}{2}}
\sphinxAtStartPar
时间同步
\sphinxbeforeendvarwidth
\end{varwidth}%
&\begin{varwidth}[t]{\sphinxcolwidth{1}{2}}
\sphinxAtStartPar
与本地站同时采集
\sphinxbeforeendvarwidth
\end{varwidth}%
\\
\sphinxhline\begin{varwidth}[t]{\sphinxcolwidth{1}{2}}
\sphinxAtStartPar
电磁环境
\sphinxbeforeendvarwidth
\end{varwidth}%
&\begin{varwidth}[t]{\sphinxcolwidth{1}{2}}
\sphinxAtStartPar
参考站环境安静
\sphinxbeforeendvarwidth
\end{varwidth}%
\\
\sphinxhline\begin{varwidth}[t]{\sphinxcolwidth{1}{2}}
\sphinxAtStartPar
相干性
\sphinxbeforeendvarwidth
\end{varwidth}%
&\begin{varwidth}[t]{\sphinxcolwidth{1}{2}}
\sphinxAtStartPar
与本地站相干性 > 0.7
\sphinxbeforeendvarwidth
\end{varwidth}%
\\
\sphinxbottomrule
\end{tabulary}
\sphinxtableafterendhook\par
\sphinxattableend\end{savenotes}


\subsubsection{5.14.3 远参考数据编辑 (RemoteReferenceEditForm)}
\label{\detokenize{chapters/chapter6:remotereferenceeditform}}
\sphinxAtStartPar
\sphinxstylestrong{功能说明：}

\sphinxAtStartPar
编辑和配置远参考站的傅里叶系数数据。

\sphinxAtStartPar
\sphinxstylestrong{编辑内容：}


\begin{savenotes}\sphinxattablestart
\sphinxthistablewithglobalstyle
\centering
\begin{tabulary}{\linewidth}[t]{TT}
\sphinxtoprule
\begin{varwidth}[t]{\sphinxcolwidth{1}{2}}
\sphinxstyletheadfamily \sphinxAtStartPar
项目
\sphinxbeforeendvarwidth
\end{varwidth}%
&\begin{varwidth}[t]{\sphinxcolwidth{1}{2}}
\sphinxstyletheadfamily \sphinxAtStartPar
说明
\sphinxbeforeendvarwidth
\end{varwidth}%
\\
\sphinxmidrule
\sphinxtableatstartofbodyhook\begin{varwidth}[t]{\sphinxcolwidth{1}{2}}
\sphinxAtStartPar
参考站名称
\sphinxbeforeendvarwidth
\end{varwidth}%
&\begin{varwidth}[t]{\sphinxcolwidth{1}{2}}
\sphinxAtStartPar
远参考站标识
\sphinxbeforeendvarwidth
\end{varwidth}%
\\
\sphinxhline\begin{varwidth}[t]{\sphinxcolwidth{1}{2}}
\sphinxAtStartPar
时间范围
\sphinxbeforeendvarwidth
\end{varwidth}%
&\begin{varwidth}[t]{\sphinxcolwidth{1}{2}}
\sphinxAtStartPar
数据采集时间
\sphinxbeforeendvarwidth
\end{varwidth}%
\\
\sphinxhline\begin{varwidth}[t]{\sphinxcolwidth{1}{2}}
\sphinxAtStartPar
频率范围
\sphinxbeforeendvarwidth
\end{varwidth}%
&\begin{varwidth}[t]{\sphinxcolwidth{1}{2}}
\sphinxAtStartPar
可用频率范围
\sphinxbeforeendvarwidth
\end{varwidth}%
\\
\sphinxhline\begin{varwidth}[t]{\sphinxcolwidth{1}{2}}
\sphinxAtStartPar
相干性
\sphinxbeforeendvarwidth
\end{varwidth}%
&\begin{varwidth}[t]{\sphinxcolwidth{1}{2}}
\sphinxAtStartPar
与本地站的相干性
\sphinxbeforeendvarwidth
\end{varwidth}%
\\
\sphinxbottomrule
\end{tabulary}
\sphinxtableafterendhook\par
\sphinxattableend\end{savenotes}

\sphinxAtStartPar
\sphinxstylestrong{操作步骤：}
\begin{enumerate}
\sphinxsetlistlabels{\arabic}{enumi}{enumii}{}{.}%
\item {} 
\sphinxAtStartPar
右键测点 → \sphinxcode{\sphinxupquote{远参考设置}}

\item {} 
\sphinxAtStartPar
在对话框中配置远参考站

\item {} 
\sphinxAtStartPar
查看相干性曲线

\item {} 
\sphinxAtStartPar
保存设置

\end{enumerate}


\subsubsection{5.14.4 远参考处理建议}
\label{\detokenize{chapters/chapter6:id56}}
\sphinxAtStartPar
\sphinxstylestrong{高质量远参考条件：}

\sphinxAtStartPar
\sphinxstylestrong{常见问题：}


\begin{savenotes}\sphinxattablestart
\sphinxthistablewithglobalstyle
\centering
\begin{tabulary}{\linewidth}[t]{TT}
\sphinxtoprule
\begin{varwidth}[t]{\sphinxcolwidth{1}{2}}
\sphinxstyletheadfamily \sphinxAtStartPar
问题
\sphinxbeforeendvarwidth
\end{varwidth}%
&\begin{varwidth}[t]{\sphinxcolwidth{1}{2}}
\sphinxstyletheadfamily \sphinxAtStartPar
解决方案
\sphinxbeforeendvarwidth
\end{varwidth}%
\\
\sphinxmidrule
\sphinxtableatstartofbodyhook\begin{varwidth}[t]{\sphinxcolwidth{1}{2}}
\sphinxAtStartPar
相干性低
\sphinxbeforeendvarwidth
\end{varwidth}%
&\begin{varwidth}[t]{\sphinxcolwidth{1}{2}}
\sphinxAtStartPar
检查时间同步
\sphinxbeforeendvarwidth
\end{varwidth}%
\\
\sphinxhline\begin{varwidth}[t]{\sphinxcolwidth{1}{2}}
\sphinxAtStartPar
数据缺失
\sphinxbeforeendvarwidth
\end{varwidth}%
&\begin{varwidth}[t]{\sphinxcolwidth{1}{2}}
\sphinxAtStartPar
确认参考站数据完整
\sphinxbeforeendvarwidth
\end{varwidth}%
\\
\sphinxhline\begin{varwidth}[t]{\sphinxcolwidth{1}{2}}
\sphinxAtStartPar
噪声干扰
\sphinxbeforeendvarwidth
\end{varwidth}%
&\begin{varwidth}[t]{\sphinxcolwidth{1}{2}}
\sphinxAtStartPar
选择更安静的参考站
\sphinxbeforeendvarwidth
\end{varwidth}%
\\
\sphinxbottomrule
\end{tabulary}
\sphinxtableafterendhook\par
\sphinxattableend\end{savenotes}


\bigskip\hrule\bigskip



\subsection{6.16 🔄 通道旋转恢复}
\label{\detokenize{chapters/chapter6:id57}}

\subsubsection{5.15.1 功能说明}
\label{\detokenize{chapters/chapter6:id58}}
\sphinxAtStartPar
通道旋转恢复功能用于校正传感器方向偏差,通过旋转电磁场分量实现坐标系变换。


\subsubsection{5.15.2 旋转角度设置}
\label{\detokenize{chapters/chapter6:id59}}

\begin{savenotes}\sphinxattablestart
\sphinxthistablewithglobalstyle
\centering
\begin{tabulary}{\linewidth}[t]{TTT}
\sphinxtoprule
\begin{varwidth}[t]{\sphinxcolwidth{1}{3}}
\sphinxstyletheadfamily \sphinxAtStartPar
通道
\sphinxbeforeendvarwidth
\end{varwidth}%
&\begin{varwidth}[t]{\sphinxcolwidth{1}{3}}
\sphinxstyletheadfamily \sphinxAtStartPar
范围
\sphinxbeforeendvarwidth
\end{varwidth}%
&\begin{varwidth}[t]{\sphinxcolwidth{1}{3}}
\sphinxstyletheadfamily \sphinxAtStartPar
说明
\sphinxbeforeendvarwidth
\end{varwidth}%
\\
\sphinxmidrule
\sphinxtableatstartofbodyhook\begin{varwidth}[t]{\sphinxcolwidth{1}{3}}
\sphinxAtStartPar
Ex
\sphinxbeforeendvarwidth
\end{varwidth}%
&\begin{varwidth}[t]{\sphinxcolwidth{1}{3}}
\sphinxAtStartPar
\sphinxhyphen{}360° \textasciitilde{} +360°
\sphinxbeforeendvarwidth
\end{varwidth}%
&\begin{varwidth}[t]{\sphinxcolwidth{1}{3}}
\sphinxAtStartPar
X方向电场旋转角
\sphinxbeforeendvarwidth
\end{varwidth}%
\\
\sphinxhline\begin{varwidth}[t]{\sphinxcolwidth{1}{3}}
\sphinxAtStartPar
Ey
\sphinxbeforeendvarwidth
\end{varwidth}%
&\begin{varwidth}[t]{\sphinxcolwidth{1}{3}}
\sphinxAtStartPar
\sphinxhyphen{}360° \textasciitilde{} +360°
\sphinxbeforeendvarwidth
\end{varwidth}%
&\begin{varwidth}[t]{\sphinxcolwidth{1}{3}}
\sphinxAtStartPar
Y方向电场旋转角
\sphinxbeforeendvarwidth
\end{varwidth}%
\\
\sphinxhline\begin{varwidth}[t]{\sphinxcolwidth{1}{3}}
\sphinxAtStartPar
Hx
\sphinxbeforeendvarwidth
\end{varwidth}%
&\begin{varwidth}[t]{\sphinxcolwidth{1}{3}}
\sphinxAtStartPar
\sphinxhyphen{}360° \textasciitilde{} +360°
\sphinxbeforeendvarwidth
\end{varwidth}%
&\begin{varwidth}[t]{\sphinxcolwidth{1}{3}}
\sphinxAtStartPar
X方向磁场旋转角
\sphinxbeforeendvarwidth
\end{varwidth}%
\\
\sphinxhline\begin{varwidth}[t]{\sphinxcolwidth{1}{3}}
\sphinxAtStartPar
Hy
\sphinxbeforeendvarwidth
\end{varwidth}%
&\begin{varwidth}[t]{\sphinxcolwidth{1}{3}}
\sphinxAtStartPar
\sphinxhyphen{}360° \textasciitilde{} +360°
\sphinxbeforeendvarwidth
\end{varwidth}%
&\begin{varwidth}[t]{\sphinxcolwidth{1}{3}}
\sphinxAtStartPar
Y方向磁场旋转角
\sphinxbeforeendvarwidth
\end{varwidth}%
\\
\sphinxbottomrule
\end{tabulary}
\sphinxtableafterendhook\par
\sphinxattableend\end{savenotes}


\subsubsection{5.15.3 使用场景}
\label{\detokenize{chapters/chapter6:id60}}\begin{itemize}
\item {} 
\sphinxAtStartPar
传感器安装方向与地理北不一致

\item {} 
\sphinxAtStartPar
需要将数据转换到统一坐标系

\item {} 
\sphinxAtStartPar
远参考数据处理时的坐标对齐

\item {} 
\sphinxAtStartPar
多测点数据对比时的坐标标准化

\end{itemize}


\subsubsection{5.15.4 操作步骤}
\label{\detokenize{chapters/chapter6:id61}}\begin{enumerate}
\sphinxsetlistlabels{\arabic}{enumi}{enumii}{}{.}%
\item {} 
\sphinxAtStartPar
在工程树中选择测点

\item {} 
\sphinxAtStartPar
右键选择 \sphinxcode{\sphinxupquote{编辑 → 通道旋转恢复}}

\item {} 
\sphinxAtStartPar
设置各通道的旋转角度

\item {} 
\sphinxAtStartPar
点击确定应用旋转

\item {} 
\sphinxAtStartPar
系统自动更新傅里叶系数数据

\end{enumerate}


\subsubsection{5.15.5 注意事项}
\label{\detokenize{chapters/chapter6:id62}}\begin{itemize}
\item {} 
\sphinxAtStartPar
旋转操作会修改原始傅里叶系数

\item {} 
\sphinxAtStartPar
建议在旋转前备份数据

\item {} 
\sphinxAtStartPar
旋转角度以逆时针为正方向

\item {} 
\sphinxAtStartPar
Hz通道不受旋转影响

\end{itemize}


\bigskip\hrule\bigskip



\subsection{6.17 🧬 智能筛选算法}
\label{\detokenize{chapters/chapter6:id63}}

\subsubsection{5.16.1 遗传算法概述}
\label{\detokenize{chapters/chapter6:id64}}
\sphinxAtStartPar
MTDP使用遗传算法进行MT数据智能筛选，自动优化数据质量。遗传算法模拟自然进化过程，通过选择、交叉和变异操作，逐步优化数据筛选方案。


\subsubsection{5.16.2 算法参数}
\label{\detokenize{chapters/chapter6:id65}}

\begin{savenotes}\sphinxattablestart
\sphinxthistablewithglobalstyle
\centering
\begin{tabulary}{\linewidth}[t]{TTT}
\sphinxtoprule
\begin{varwidth}[t]{\sphinxcolwidth{1}{3}}
\sphinxstyletheadfamily \sphinxAtStartPar
参数
\sphinxbeforeendvarwidth
\end{varwidth}%
&\begin{varwidth}[t]{\sphinxcolwidth{1}{3}}
\sphinxstyletheadfamily \sphinxAtStartPar
说明
\sphinxbeforeendvarwidth
\end{varwidth}%
&\begin{varwidth}[t]{\sphinxcolwidth{1}{3}}
\sphinxstyletheadfamily \sphinxAtStartPar
典型值
\sphinxbeforeendvarwidth
\end{varwidth}%
\\
\sphinxmidrule
\sphinxtableatstartofbodyhook\begin{varwidth}[t]{\sphinxcolwidth{1}{3}}
\sphinxAtStartPar
种群大小
\sphinxbeforeendvarwidth
\end{varwidth}%
&\begin{varwidth}[t]{\sphinxcolwidth{1}{3}}
\sphinxAtStartPar
每代个体数量
\sphinxbeforeendvarwidth
\end{varwidth}%
&\begin{varwidth}[t]{\sphinxcolwidth{1}{3}}
\sphinxAtStartPar
50\sphinxhyphen{}200
\sphinxbeforeendvarwidth
\end{varwidth}%
\\
\sphinxhline\begin{varwidth}[t]{\sphinxcolwidth{1}{3}}
\sphinxAtStartPar
迭代次数
\sphinxbeforeendvarwidth
\end{varwidth}%
&\begin{varwidth}[t]{\sphinxcolwidth{1}{3}}
\sphinxAtStartPar
进化代数
\sphinxbeforeendvarwidth
\end{varwidth}%
&\begin{varwidth}[t]{\sphinxcolwidth{1}{3}}
\sphinxAtStartPar
100\sphinxhyphen{}500
\sphinxbeforeendvarwidth
\end{varwidth}%
\\
\sphinxhline\begin{varwidth}[t]{\sphinxcolwidth{1}{3}}
\sphinxAtStartPar
交叉率
\sphinxbeforeendvarwidth
\end{varwidth}%
&\begin{varwidth}[t]{\sphinxcolwidth{1}{3}}
\sphinxAtStartPar
个体交叉概率
\sphinxbeforeendvarwidth
\end{varwidth}%
&\begin{varwidth}[t]{\sphinxcolwidth{1}{3}}
\sphinxAtStartPar
0.7\sphinxhyphen{}0.9
\sphinxbeforeendvarwidth
\end{varwidth}%
\\
\sphinxhline\begin{varwidth}[t]{\sphinxcolwidth{1}{3}}
\sphinxAtStartPar
变异率
\sphinxbeforeendvarwidth
\end{varwidth}%
&\begin{varwidth}[t]{\sphinxcolwidth{1}{3}}
\sphinxAtStartPar
基因变异概率
\sphinxbeforeendvarwidth
\end{varwidth}%
&\begin{varwidth}[t]{\sphinxcolwidth{1}{3}}
\sphinxAtStartPar
0.01\sphinxhyphen{}0.1
\sphinxbeforeendvarwidth
\end{varwidth}%
\\
\sphinxhline\begin{varwidth}[t]{\sphinxcolwidth{1}{3}}
\sphinxAtStartPar
精英保留
\sphinxbeforeendvarwidth
\end{varwidth}%
&\begin{varwidth}[t]{\sphinxcolwidth{1}{3}}
\sphinxAtStartPar
保留最优个体数
\sphinxbeforeendvarwidth
\end{varwidth}%
&\begin{varwidth}[t]{\sphinxcolwidth{1}{3}}
\sphinxAtStartPar
1\sphinxhyphen{}5
\sphinxbeforeendvarwidth
\end{varwidth}%
\\
\sphinxbottomrule
\end{tabulary}
\sphinxtableafterendhook\par
\sphinxattableend\end{savenotes}


\subsubsection{5.16.3 优化目标}
\label{\detokenize{chapters/chapter6:id66}}\begin{itemize}
\item {} 
\sphinxAtStartPar
\sphinxstylestrong{RMS最小化}：最小化数据与模型的残差

\item {} 
\sphinxAtStartPar
\sphinxstylestrong{相干性最大化}：最大化通道间相干性

\item {} 
\sphinxAtStartPar
\sphinxstylestrong{数据利用率}：最大化有效数据点数

\end{itemize}


\subsubsection{5.16.4 多目标优化}
\label{\detokenize{chapters/chapter6:id67}}
\sphinxAtStartPar
支持NSGA\sphinxhyphen{}II多目标优化算法：
\begin{itemize}
\item {} 
\sphinxAtStartPar
Pareto最优解集

\item {} 
\sphinxAtStartPar
拥挤距离排序

\item {} 
\sphinxAtStartPar
外部档案管理

\end{itemize}


\subsubsection{5.16.5 统计质量控制}
\label{\detokenize{chapters/chapter6:id68}}
\sphinxAtStartPar
\sphinxstylestrong{马氏距离过滤：}
\begin{itemize}
\item {} 
\sphinxAtStartPar
协方差矩阵分析

\item {} 
\sphinxAtStartPar
Huber加权稳健统计

\item {} 
\sphinxAtStartPar
迭代阈值优化

\item {} 
\sphinxAtStartPar
异常值自动检测

\end{itemize}


\subsubsection{5.16.6 使用建议}
\label{\detokenize{chapters/chapter6:id69}}\begin{enumerate}
\sphinxsetlistlabels{\arabic}{enumi}{enumii}{}{.}%
\item {} 
\sphinxAtStartPar
首先使用默认参数进行初步筛选

\item {} 
\sphinxAtStartPar
根据结果调整参数进行精细优化

\item {} 
\sphinxAtStartPar
对比多次运行结果确保稳定性

\item {} 
\sphinxAtStartPar
结合人工筛选进行最终确认

\end{enumerate}


\subsubsection{5.17 MT参数完整列表}
\label{\detokenize{chapters/chapter6:mt}}
\sphinxAtStartPar
MTDP支持75+种MT参数用于数据分析和筛选。以下按类别详细介绍各参数的定义、理论解释和计算公式。

\sphinxAtStartPar
MTDP支持75+种MT参数用于数据分析和筛选。以下按类别详细介绍各参数的定义、理论解释和计算公式。


\paragraph{5.17.1 基本参数}
\label{\detokenize{chapters/chapter6:id70}}

\begin{savenotes}\sphinxattablestart
\sphinxthistablewithglobalstyle
\centering
\begin{tabulary}{\linewidth}[t]{TTTT}
\sphinxtoprule
\begin{varwidth}[t]{\sphinxcolwidth{1}{4}}
\sphinxstyletheadfamily \sphinxAtStartPar
参数名
\sphinxbeforeendvarwidth
\end{varwidth}%
&\begin{varwidth}[t]{\sphinxcolwidth{1}{4}}
\sphinxstyletheadfamily \sphinxAtStartPar
内部名称
\sphinxbeforeendvarwidth
\end{varwidth}%
&\begin{varwidth}[t]{\sphinxcolwidth{1}{4}}
\sphinxstyletheadfamily \sphinxAtStartPar
说明
\sphinxbeforeendvarwidth
\end{varwidth}%
&\begin{varwidth}[t]{\sphinxcolwidth{1}{4}}
\sphinxstyletheadfamily \sphinxAtStartPar
单位
\sphinxbeforeendvarwidth
\end{varwidth}%
\\
\sphinxmidrule
\sphinxtableatstartofbodyhook\begin{varwidth}[t]{\sphinxcolwidth{1}{4}}
\sphinxAtStartPar
fre
\sphinxbeforeendvarwidth
\end{varwidth}%
&\begin{varwidth}[t]{\sphinxcolwidth{1}{4}}
\sphinxAtStartPar
MT\_fre
\sphinxbeforeendvarwidth
\end{varwidth}%
&\begin{varwidth}[t]{\sphinxcolwidth{1}{4}}
\sphinxAtStartPar
频率
\sphinxbeforeendvarwidth
\end{varwidth}%
&\begin{varwidth}[t]{\sphinxcolwidth{1}{4}}
\sphinxAtStartPar
Hz
\sphinxbeforeendvarwidth
\end{varwidth}%
\\
\sphinxhline\begin{varwidth}[t]{\sphinxcolwidth{1}{4}}
\sphinxAtStartPar
lgfre
\sphinxbeforeendvarwidth
\end{varwidth}%
&\begin{varwidth}[t]{\sphinxcolwidth{1}{4}}
\sphinxAtStartPar
MT\_lgfre
\sphinxbeforeendvarwidth
\end{varwidth}%
&\begin{varwidth}[t]{\sphinxcolwidth{1}{4}}
\sphinxAtStartPar
对数频率（log10）
\sphinxbeforeendvarwidth
\end{varwidth}%
&\begin{varwidth}[t]{\sphinxcolwidth{1}{4}}
\sphinxAtStartPar
log(Hz)
\sphinxbeforeendvarwidth
\end{varwidth}%
\\
\sphinxhline\begin{varwidth}[t]{\sphinxcolwidth{1}{4}}
\sphinxAtStartPar
Time
\sphinxbeforeendvarwidth
\end{varwidth}%
&\begin{varwidth}[t]{\sphinxcolwidth{1}{4}}
\sphinxAtStartPar
MT\_Time
\sphinxbeforeendvarwidth
\end{varwidth}%
&\begin{varwidth}[t]{\sphinxcolwidth{1}{4}}
\sphinxAtStartPar
时间
\sphinxbeforeendvarwidth
\end{varwidth}%
&\begin{varwidth}[t]{\sphinxcolwidth{1}{4}}
\sphinxAtStartPar
s
\sphinxbeforeendvarwidth
\end{varwidth}%
\\
\sphinxbottomrule
\end{tabulary}
\sphinxtableafterendhook\par
\sphinxattableend\end{savenotes}


\paragraph{5.17.2 阻抗张量实部与虚部（单位：Ω）}
\label{\detokenize{chapters/chapter6:id71}}
\sphinxAtStartPar
阻抗张量Z描述电场与磁场之间的关系：

\sphinxAtStartPar
\sphinxstylestrong{理论公式：}

\begin{sphinxVerbatim}[commandchars=\\\{\}]
\PYG{p}{[}\PYG{n}{Ex}\PYG{p}{]}   \PYG{p}{[}\PYG{n}{Zxx}  \PYG{n}{Zxy}\PYG{p}{]} \PYG{p}{[}\PYG{n}{Hx}\PYG{p}{]}
\PYG{p}{[}\PYG{n}{Ey}\PYG{p}{]} \PYG{o}{=} \PYG{p}{[}\PYG{n}{Zyx}  \PYG{n}{Zyy}\PYG{p}{]} \PYG{p}{[}\PYG{n}{Hy}\PYG{p}{]}
\end{sphinxVerbatim}


\begin{savenotes}\sphinxattablestart
\sphinxthistablewithglobalstyle
\centering
\begin{tabulary}{\linewidth}[t]{TTTT}
\sphinxtoprule
\begin{varwidth}[t]{\sphinxcolwidth{1}{4}}
\sphinxstyletheadfamily \sphinxAtStartPar
参数名
\sphinxbeforeendvarwidth
\end{varwidth}%
&\begin{varwidth}[t]{\sphinxcolwidth{1}{4}}
\sphinxstyletheadfamily \sphinxAtStartPar
内部名称
\sphinxbeforeendvarwidth
\end{varwidth}%
&\begin{varwidth}[t]{\sphinxcolwidth{1}{4}}
\sphinxstyletheadfamily \sphinxAtStartPar
说明
\sphinxbeforeendvarwidth
\end{varwidth}%
&\begin{varwidth}[t]{\sphinxcolwidth{1}{4}}
\sphinxstyletheadfamily \sphinxAtStartPar
计算公式
\sphinxbeforeendvarwidth
\end{varwidth}%
\\
\sphinxmidrule
\sphinxtableatstartofbodyhook\begin{varwidth}[t]{\sphinxcolwidth{1}{4}}
\sphinxAtStartPar
zxxr
\sphinxbeforeendvarwidth
\end{varwidth}%
&\begin{varwidth}[t]{\sphinxcolwidth{1}{4}}
\sphinxAtStartPar
MT\_zxxr
\sphinxbeforeendvarwidth
\end{varwidth}%
&\begin{varwidth}[t]{\sphinxcolwidth{1}{4}}
\sphinxAtStartPar
Zxx实部
\sphinxbeforeendvarwidth
\end{varwidth}%
&\begin{varwidth}[t]{\sphinxcolwidth{1}{4}}
\sphinxAtStartPar
Re(Zxx)
\sphinxbeforeendvarwidth
\end{varwidth}%
\\
\sphinxhline\begin{varwidth}[t]{\sphinxcolwidth{1}{4}}
\sphinxAtStartPar
zxxi
\sphinxbeforeendvarwidth
\end{varwidth}%
&\begin{varwidth}[t]{\sphinxcolwidth{1}{4}}
\sphinxAtStartPar
MT\_zxxi
\sphinxbeforeendvarwidth
\end{varwidth}%
&\begin{varwidth}[t]{\sphinxcolwidth{1}{4}}
\sphinxAtStartPar
Zxx虚部
\sphinxbeforeendvarwidth
\end{varwidth}%
&\begin{varwidth}[t]{\sphinxcolwidth{1}{4}}
\sphinxAtStartPar
Im(Zxx)
\sphinxbeforeendvarwidth
\end{varwidth}%
\\
\sphinxhline\begin{varwidth}[t]{\sphinxcolwidth{1}{4}}
\sphinxAtStartPar
zxyr
\sphinxbeforeendvarwidth
\end{varwidth}%
&\begin{varwidth}[t]{\sphinxcolwidth{1}{4}}
\sphinxAtStartPar
MT\_zxyr
\sphinxbeforeendvarwidth
\end{varwidth}%
&\begin{varwidth}[t]{\sphinxcolwidth{1}{4}}
\sphinxAtStartPar
Zxy实部
\sphinxbeforeendvarwidth
\end{varwidth}%
&\begin{varwidth}[t]{\sphinxcolwidth{1}{4}}
\sphinxAtStartPar
Re(Zxy)
\sphinxbeforeendvarwidth
\end{varwidth}%
\\
\sphinxhline\begin{varwidth}[t]{\sphinxcolwidth{1}{4}}
\sphinxAtStartPar
zxyi
\sphinxbeforeendvarwidth
\end{varwidth}%
&\begin{varwidth}[t]{\sphinxcolwidth{1}{4}}
\sphinxAtStartPar
MT\_zxyi
\sphinxbeforeendvarwidth
\end{varwidth}%
&\begin{varwidth}[t]{\sphinxcolwidth{1}{4}}
\sphinxAtStartPar
Zxy虚部
\sphinxbeforeendvarwidth
\end{varwidth}%
&\begin{varwidth}[t]{\sphinxcolwidth{1}{4}}
\sphinxAtStartPar
Im(Zxy)
\sphinxbeforeendvarwidth
\end{varwidth}%
\\
\sphinxhline\begin{varwidth}[t]{\sphinxcolwidth{1}{4}}
\sphinxAtStartPar
zyxr
\sphinxbeforeendvarwidth
\end{varwidth}%
&\begin{varwidth}[t]{\sphinxcolwidth{1}{4}}
\sphinxAtStartPar
MT\_zyxr
\sphinxbeforeendvarwidth
\end{varwidth}%
&\begin{varwidth}[t]{\sphinxcolwidth{1}{4}}
\sphinxAtStartPar
Zyx实部
\sphinxbeforeendvarwidth
\end{varwidth}%
&\begin{varwidth}[t]{\sphinxcolwidth{1}{4}}
\sphinxAtStartPar
Re(Zyx)
\sphinxbeforeendvarwidth
\end{varwidth}%
\\
\sphinxhline\begin{varwidth}[t]{\sphinxcolwidth{1}{4}}
\sphinxAtStartPar
zyxi
\sphinxbeforeendvarwidth
\end{varwidth}%
&\begin{varwidth}[t]{\sphinxcolwidth{1}{4}}
\sphinxAtStartPar
MT\_zyxi
\sphinxbeforeendvarwidth
\end{varwidth}%
&\begin{varwidth}[t]{\sphinxcolwidth{1}{4}}
\sphinxAtStartPar
Zyx虚部
\sphinxbeforeendvarwidth
\end{varwidth}%
&\begin{varwidth}[t]{\sphinxcolwidth{1}{4}}
\sphinxAtStartPar
Im(Zyx)
\sphinxbeforeendvarwidth
\end{varwidth}%
\\
\sphinxhline\begin{varwidth}[t]{\sphinxcolwidth{1}{4}}
\sphinxAtStartPar
zyyr
\sphinxbeforeendvarwidth
\end{varwidth}%
&\begin{varwidth}[t]{\sphinxcolwidth{1}{4}}
\sphinxAtStartPar
MT\_zyyr
\sphinxbeforeendvarwidth
\end{varwidth}%
&\begin{varwidth}[t]{\sphinxcolwidth{1}{4}}
\sphinxAtStartPar
Zyy实部
\sphinxbeforeendvarwidth
\end{varwidth}%
&\begin{varwidth}[t]{\sphinxcolwidth{1}{4}}
\sphinxAtStartPar
Re(Zyy)
\sphinxbeforeendvarwidth
\end{varwidth}%
\\
\sphinxhline\begin{varwidth}[t]{\sphinxcolwidth{1}{4}}
\sphinxAtStartPar
zyyi
\sphinxbeforeendvarwidth
\end{varwidth}%
&\begin{varwidth}[t]{\sphinxcolwidth{1}{4}}
\sphinxAtStartPar
MT\_zyyi
\sphinxbeforeendvarwidth
\end{varwidth}%
&\begin{varwidth}[t]{\sphinxcolwidth{1}{4}}
\sphinxAtStartPar
Zyy虚部
\sphinxbeforeendvarwidth
\end{varwidth}%
&\begin{varwidth}[t]{\sphinxcolwidth{1}{4}}
\sphinxAtStartPar
Im(Zyy)
\sphinxbeforeendvarwidth
\end{varwidth}%
\\
\sphinxbottomrule
\end{tabulary}
\sphinxtableafterendhook\par
\sphinxattableend\end{savenotes}

\sphinxAtStartPar
\sphinxstylestrong{Gamble估计公式：}

\begin{sphinxVerbatim}[commandchars=\\\{\}]
Z = \PYGZlt{}EH*\PYGZgt{} / \PYGZlt{}HH*\PYGZgt{}
其中：E为电场，H为磁场，*表示共轭，\PYGZlt{}\PYGZgt{}表示平均
\end{sphinxVerbatim}


\paragraph{5.17.3 阻抗张量幅值与相位}
\label{\detokenize{chapters/chapter6:id72}}

\begin{savenotes}\sphinxattablestart
\sphinxthistablewithglobalstyle
\centering
\begin{tabulary}{\linewidth}[t]{TTTT}
\sphinxtoprule
\begin{varwidth}[t]{\sphinxcolwidth{1}{4}}
\sphinxstyletheadfamily \sphinxAtStartPar
参数名
\sphinxbeforeendvarwidth
\end{varwidth}%
&\begin{varwidth}[t]{\sphinxcolwidth{1}{4}}
\sphinxstyletheadfamily \sphinxAtStartPar
内部名称
\sphinxbeforeendvarwidth
\end{varwidth}%
&\begin{varwidth}[t]{\sphinxcolwidth{1}{4}}
\sphinxstyletheadfamily \sphinxAtStartPar
说明
\sphinxbeforeendvarwidth
\end{varwidth}%
&\begin{varwidth}[t]{\sphinxcolwidth{1}{4}}
\sphinxstyletheadfamily \sphinxAtStartPar
计算公式
\sphinxbeforeendvarwidth
\end{varwidth}%
\\
\sphinxmidrule
\sphinxtableatstartofbodyhook\begin{varwidth}[t]{\sphinxcolwidth{1}{4}}
\sphinxAtStartPar
zxxa
\sphinxbeforeendvarwidth
\end{varwidth}%
&\begin{varwidth}[t]{\sphinxcolwidth{1}{4}}
\sphinxAtStartPar
MT\_zxxa
\sphinxbeforeendvarwidth
\end{varwidth}%
&\begin{varwidth}[t]{\sphinxcolwidth{1}{4}}
\sphinxAtStartPar
Zxx幅值
\sphinxbeforeendvarwidth
\end{varwidth}%
&\begin{varwidth}[t]{\sphinxcolwidth{1}{4}}
\sphinxAtStartPar
|Zxx| = √(zxxr² + zxxi²)
\sphinxbeforeendvarwidth
\end{varwidth}%
\\
\sphinxhline\begin{varwidth}[t]{\sphinxcolwidth{1}{4}}
\sphinxAtStartPar
zxxp
\sphinxbeforeendvarwidth
\end{varwidth}%
&\begin{varwidth}[t]{\sphinxcolwidth{1}{4}}
\sphinxAtStartPar
MT\_zxxp
\sphinxbeforeendvarwidth
\end{varwidth}%
&\begin{varwidth}[t]{\sphinxcolwidth{1}{4}}
\sphinxAtStartPar
Zxx相位
\sphinxbeforeendvarwidth
\end{varwidth}%
&\begin{varwidth}[t]{\sphinxcolwidth{1}{4}}
\sphinxAtStartPar
atan2(zxxi, zxxr) × 180/π
\sphinxbeforeendvarwidth
\end{varwidth}%
\\
\sphinxhline\begin{varwidth}[t]{\sphinxcolwidth{1}{4}}
\sphinxAtStartPar
zxya
\sphinxbeforeendvarwidth
\end{varwidth}%
&\begin{varwidth}[t]{\sphinxcolwidth{1}{4}}
\sphinxAtStartPar
MT\_zxya
\sphinxbeforeendvarwidth
\end{varwidth}%
&\begin{varwidth}[t]{\sphinxcolwidth{1}{4}}
\sphinxAtStartPar
Zxy幅值
\sphinxbeforeendvarwidth
\end{varwidth}%
&\begin{varwidth}[t]{\sphinxcolwidth{1}{4}}
\sphinxAtStartPar
|Zxy| = √(zxyr² + zxyi²)
\sphinxbeforeendvarwidth
\end{varwidth}%
\\
\sphinxhline\begin{varwidth}[t]{\sphinxcolwidth{1}{4}}
\sphinxAtStartPar
zxyp
\sphinxbeforeendvarwidth
\end{varwidth}%
&\begin{varwidth}[t]{\sphinxcolwidth{1}{4}}
\sphinxAtStartPar
MT\_zxyp
\sphinxbeforeendvarwidth
\end{varwidth}%
&\begin{varwidth}[t]{\sphinxcolwidth{1}{4}}
\sphinxAtStartPar
Zxy相位
\sphinxbeforeendvarwidth
\end{varwidth}%
&\begin{varwidth}[t]{\sphinxcolwidth{1}{4}}
\sphinxAtStartPar
atan2(zxyi, zxyr) × 180/π
\sphinxbeforeendvarwidth
\end{varwidth}%
\\
\sphinxhline\begin{varwidth}[t]{\sphinxcolwidth{1}{4}}
\sphinxAtStartPar
zyxa
\sphinxbeforeendvarwidth
\end{varwidth}%
&\begin{varwidth}[t]{\sphinxcolwidth{1}{4}}
\sphinxAtStartPar
MT\_zyxa
\sphinxbeforeendvarwidth
\end{varwidth}%
&\begin{varwidth}[t]{\sphinxcolwidth{1}{4}}
\sphinxAtStartPar
Zyx幅值
\sphinxbeforeendvarwidth
\end{varwidth}%
&\begin{varwidth}[t]{\sphinxcolwidth{1}{4}}
\sphinxAtStartPar
|Zyx| = √(zyxr² + zyxi²)
\sphinxbeforeendvarwidth
\end{varwidth}%
\\
\sphinxhline\begin{varwidth}[t]{\sphinxcolwidth{1}{4}}
\sphinxAtStartPar
zyxp
\sphinxbeforeendvarwidth
\end{varwidth}%
&\begin{varwidth}[t]{\sphinxcolwidth{1}{4}}
\sphinxAtStartPar
MT\_zyxp
\sphinxbeforeendvarwidth
\end{varwidth}%
&\begin{varwidth}[t]{\sphinxcolwidth{1}{4}}
\sphinxAtStartPar
Zyx相位
\sphinxbeforeendvarwidth
\end{varwidth}%
&\begin{varwidth}[t]{\sphinxcolwidth{1}{4}}
\sphinxAtStartPar
atan2(zyxi, zyxr) × 180/π
\sphinxbeforeendvarwidth
\end{varwidth}%
\\
\sphinxhline\begin{varwidth}[t]{\sphinxcolwidth{1}{4}}
\sphinxAtStartPar
zyya
\sphinxbeforeendvarwidth
\end{varwidth}%
&\begin{varwidth}[t]{\sphinxcolwidth{1}{4}}
\sphinxAtStartPar
MT\_zyya
\sphinxbeforeendvarwidth
\end{varwidth}%
&\begin{varwidth}[t]{\sphinxcolwidth{1}{4}}
\sphinxAtStartPar
Zyy幅值
\sphinxbeforeendvarwidth
\end{varwidth}%
&\begin{varwidth}[t]{\sphinxcolwidth{1}{4}}
\sphinxAtStartPar
|Zyy| = √(zyyr² + zyyi²)
\sphinxbeforeendvarwidth
\end{varwidth}%
\\
\sphinxhline\begin{varwidth}[t]{\sphinxcolwidth{1}{4}}
\sphinxAtStartPar
zyyp
\sphinxbeforeendvarwidth
\end{varwidth}%
&\begin{varwidth}[t]{\sphinxcolwidth{1}{4}}
\sphinxAtStartPar
MT\_zyyp
\sphinxbeforeendvarwidth
\end{varwidth}%
&\begin{varwidth}[t]{\sphinxcolwidth{1}{4}}
\sphinxAtStartPar
Zyy相位
\sphinxbeforeendvarwidth
\end{varwidth}%
&\begin{varwidth}[t]{\sphinxcolwidth{1}{4}}
\sphinxAtStartPar
atan2(zyyi, zyyr) × 180/π
\sphinxbeforeendvarwidth
\end{varwidth}%
\\
\sphinxbottomrule
\end{tabulary}
\sphinxtableafterendhook\par
\sphinxattableend\end{savenotes}
\begin{quote}

\sphinxAtStartPar
\sphinxstylestrong{注意：} zxxp1/zxyp1/zyxp1/zyyp1为第一象限相位（0\sphinxhyphen{}90°）
\end{quote}


\paragraph{5.17.4 倾子张量（无量纲）}
\label{\detokenize{chapters/chapter6:id73}}
\sphinxAtStartPar
倾子T描述垂直磁场与水平磁场的关系：

\sphinxAtStartPar
\sphinxstylestrong{理论公式：}

\begin{sphinxVerbatim}[commandchars=\\\{\}]
\PYG{n}{Hz} \PYG{o}{=} \PYG{n}{Tzx·Hx} \PYG{o}{+} \PYG{n}{Tzy·Hy}
\end{sphinxVerbatim}


\begin{savenotes}\sphinxattablestart
\sphinxthistablewithglobalstyle
\centering
\begin{tabulary}{\linewidth}[t]{TTTT}
\sphinxtoprule
\begin{varwidth}[t]{\sphinxcolwidth{1}{4}}
\sphinxstyletheadfamily \sphinxAtStartPar
参数名
\sphinxbeforeendvarwidth
\end{varwidth}%
&\begin{varwidth}[t]{\sphinxcolwidth{1}{4}}
\sphinxstyletheadfamily \sphinxAtStartPar
内部名称
\sphinxbeforeendvarwidth
\end{varwidth}%
&\begin{varwidth}[t]{\sphinxcolwidth{1}{4}}
\sphinxstyletheadfamily \sphinxAtStartPar
说明
\sphinxbeforeendvarwidth
\end{varwidth}%
&\begin{varwidth}[t]{\sphinxcolwidth{1}{4}}
\sphinxstyletheadfamily \sphinxAtStartPar
计算公式
\sphinxbeforeendvarwidth
\end{varwidth}%
\\
\sphinxmidrule
\sphinxtableatstartofbodyhook\begin{varwidth}[t]{\sphinxcolwidth{1}{4}}
\sphinxAtStartPar
tzxr
\sphinxbeforeendvarwidth
\end{varwidth}%
&\begin{varwidth}[t]{\sphinxcolwidth{1}{4}}
\sphinxAtStartPar
MT\_tzxr
\sphinxbeforeendvarwidth
\end{varwidth}%
&\begin{varwidth}[t]{\sphinxcolwidth{1}{4}}
\sphinxAtStartPar
Tzx实部
\sphinxbeforeendvarwidth
\end{varwidth}%
&\begin{varwidth}[t]{\sphinxcolwidth{1}{4}}
\sphinxAtStartPar
Re(Tzx)
\sphinxbeforeendvarwidth
\end{varwidth}%
\\
\sphinxhline\begin{varwidth}[t]{\sphinxcolwidth{1}{4}}
\sphinxAtStartPar
tzxi
\sphinxbeforeendvarwidth
\end{varwidth}%
&\begin{varwidth}[t]{\sphinxcolwidth{1}{4}}
\sphinxAtStartPar
MT\_tzxi
\sphinxbeforeendvarwidth
\end{varwidth}%
&\begin{varwidth}[t]{\sphinxcolwidth{1}{4}}
\sphinxAtStartPar
Tzx虚部
\sphinxbeforeendvarwidth
\end{varwidth}%
&\begin{varwidth}[t]{\sphinxcolwidth{1}{4}}
\sphinxAtStartPar
Im(Tzx)
\sphinxbeforeendvarwidth
\end{varwidth}%
\\
\sphinxhline\begin{varwidth}[t]{\sphinxcolwidth{1}{4}}
\sphinxAtStartPar
tzyr
\sphinxbeforeendvarwidth
\end{varwidth}%
&\begin{varwidth}[t]{\sphinxcolwidth{1}{4}}
\sphinxAtStartPar
MT\_tzyr
\sphinxbeforeendvarwidth
\end{varwidth}%
&\begin{varwidth}[t]{\sphinxcolwidth{1}{4}}
\sphinxAtStartPar
Tzy实部
\sphinxbeforeendvarwidth
\end{varwidth}%
&\begin{varwidth}[t]{\sphinxcolwidth{1}{4}}
\sphinxAtStartPar
Re(Tzy)
\sphinxbeforeendvarwidth
\end{varwidth}%
\\
\sphinxhline\begin{varwidth}[t]{\sphinxcolwidth{1}{4}}
\sphinxAtStartPar
tzyi
\sphinxbeforeendvarwidth
\end{varwidth}%
&\begin{varwidth}[t]{\sphinxcolwidth{1}{4}}
\sphinxAtStartPar
MT\_tzyi
\sphinxbeforeendvarwidth
\end{varwidth}%
&\begin{varwidth}[t]{\sphinxcolwidth{1}{4}}
\sphinxAtStartPar
Tzy虚部
\sphinxbeforeendvarwidth
\end{varwidth}%
&\begin{varwidth}[t]{\sphinxcolwidth{1}{4}}
\sphinxAtStartPar
Im(Tzy)
\sphinxbeforeendvarwidth
\end{varwidth}%
\\
\sphinxhline\begin{varwidth}[t]{\sphinxcolwidth{1}{4}}
\sphinxAtStartPar
tzxa
\sphinxbeforeendvarwidth
\end{varwidth}%
&\begin{varwidth}[t]{\sphinxcolwidth{1}{4}}
\sphinxAtStartPar
MT\_tzxa
\sphinxbeforeendvarwidth
\end{varwidth}%
&\begin{varwidth}[t]{\sphinxcolwidth{1}{4}}
\sphinxAtStartPar
Tzx幅值
\sphinxbeforeendvarwidth
\end{varwidth}%
&\begin{varwidth}[t]{\sphinxcolwidth{1}{4}}
\sphinxAtStartPar
|Tzx|
\sphinxbeforeendvarwidth
\end{varwidth}%
\\
\sphinxhline\begin{varwidth}[t]{\sphinxcolwidth{1}{4}}
\sphinxAtStartPar
tzxp
\sphinxbeforeendvarwidth
\end{varwidth}%
&\begin{varwidth}[t]{\sphinxcolwidth{1}{4}}
\sphinxAtStartPar
MT\_tzxp
\sphinxbeforeendvarwidth
\end{varwidth}%
&\begin{varwidth}[t]{\sphinxcolwidth{1}{4}}
\sphinxAtStartPar
Tzx相位
\sphinxbeforeendvarwidth
\end{varwidth}%
&\begin{varwidth}[t]{\sphinxcolwidth{1}{4}}
\sphinxAtStartPar
atan2(tzxi, tzxr) × 180/π
\sphinxbeforeendvarwidth
\end{varwidth}%
\\
\sphinxhline\begin{varwidth}[t]{\sphinxcolwidth{1}{4}}
\sphinxAtStartPar
tzya
\sphinxbeforeendvarwidth
\end{varwidth}%
&\begin{varwidth}[t]{\sphinxcolwidth{1}{4}}
\sphinxAtStartPar
MT\_tzya
\sphinxbeforeendvarwidth
\end{varwidth}%
&\begin{varwidth}[t]{\sphinxcolwidth{1}{4}}
\sphinxAtStartPar
Tzy幅值
\sphinxbeforeendvarwidth
\end{varwidth}%
&\begin{varwidth}[t]{\sphinxcolwidth{1}{4}}
\sphinxAtStartPar
|Tzy|
\sphinxbeforeendvarwidth
\end{varwidth}%
\\
\sphinxhline\begin{varwidth}[t]{\sphinxcolwidth{1}{4}}
\sphinxAtStartPar
tzyp
\sphinxbeforeendvarwidth
\end{varwidth}%
&\begin{varwidth}[t]{\sphinxcolwidth{1}{4}}
\sphinxAtStartPar
MT\_tzyp
\sphinxbeforeendvarwidth
\end{varwidth}%
&\begin{varwidth}[t]{\sphinxcolwidth{1}{4}}
\sphinxAtStartPar
Tzy相位
\sphinxbeforeendvarwidth
\end{varwidth}%
&\begin{varwidth}[t]{\sphinxcolwidth{1}{4}}
\sphinxAtStartPar
atan2(tzyi, tzyr) × 180/π
\sphinxbeforeendvarwidth
\end{varwidth}%
\\
\sphinxbottomrule
\end{tabulary}
\sphinxtableafterendhook\par
\sphinxattableend\end{savenotes}


\paragraph{5.17.5 视电阻率（单位：Ω·m）}
\label{\detokenize{chapters/chapter6:m}}
\sphinxAtStartPar
\sphinxstylestrong{理论公式：}

\begin{sphinxVerbatim}[commandchars=\\\{\}]
ρij = (1/ωμ₀) × |Zij|²
其中：ω = 2πf，μ₀ = 4π×10⁻⁷ H/m
\end{sphinxVerbatim}


\begin{savenotes}\sphinxattablestart
\sphinxthistablewithglobalstyle
\centering
\begin{tabulary}{\linewidth}[t]{TTTT}
\sphinxtoprule
\begin{varwidth}[t]{\sphinxcolwidth{1}{4}}
\sphinxstyletheadfamily \sphinxAtStartPar
参数名
\sphinxbeforeendvarwidth
\end{varwidth}%
&\begin{varwidth}[t]{\sphinxcolwidth{1}{4}}
\sphinxstyletheadfamily \sphinxAtStartPar
内部名称
\sphinxbeforeendvarwidth
\end{varwidth}%
&\begin{varwidth}[t]{\sphinxcolwidth{1}{4}}
\sphinxstyletheadfamily \sphinxAtStartPar
说明
\sphinxbeforeendvarwidth
\end{varwidth}%
&\begin{varwidth}[t]{\sphinxcolwidth{1}{4}}
\sphinxstyletheadfamily \sphinxAtStartPar
计算公式
\sphinxbeforeendvarwidth
\end{varwidth}%
\\
\sphinxmidrule
\sphinxtableatstartofbodyhook\begin{varwidth}[t]{\sphinxcolwidth{1}{4}}
\sphinxAtStartPar
rxx
\sphinxbeforeendvarwidth
\end{varwidth}%
&\begin{varwidth}[t]{\sphinxcolwidth{1}{4}}
\sphinxAtStartPar
MT\_rxx
\sphinxbeforeendvarwidth
\end{varwidth}%
&\begin{varwidth}[t]{\sphinxcolwidth{1}{4}}
\sphinxAtStartPar
ρxx视电阻率
\sphinxbeforeendvarwidth
\end{varwidth}%
&\begin{varwidth}[t]{\sphinxcolwidth{1}{4}}
\sphinxAtStartPar
(1/ωμ₀) × |Zxx|²
\sphinxbeforeendvarwidth
\end{varwidth}%
\\
\sphinxhline\begin{varwidth}[t]{\sphinxcolwidth{1}{4}}
\sphinxAtStartPar
rxy
\sphinxbeforeendvarwidth
\end{varwidth}%
&\begin{varwidth}[t]{\sphinxcolwidth{1}{4}}
\sphinxAtStartPar
MT\_rxy
\sphinxbeforeendvarwidth
\end{varwidth}%
&\begin{varwidth}[t]{\sphinxcolwidth{1}{4}}
\sphinxAtStartPar
ρxy视电阻率
\sphinxbeforeendvarwidth
\end{varwidth}%
&\begin{varwidth}[t]{\sphinxcolwidth{1}{4}}
\sphinxAtStartPar
(1/ωμ₀) × |Zxy|²
\sphinxbeforeendvarwidth
\end{varwidth}%
\\
\sphinxhline\begin{varwidth}[t]{\sphinxcolwidth{1}{4}}
\sphinxAtStartPar
ryx
\sphinxbeforeendvarwidth
\end{varwidth}%
&\begin{varwidth}[t]{\sphinxcolwidth{1}{4}}
\sphinxAtStartPar
MT\_ryx
\sphinxbeforeendvarwidth
\end{varwidth}%
&\begin{varwidth}[t]{\sphinxcolwidth{1}{4}}
\sphinxAtStartPar
ρyx视电阻率
\sphinxbeforeendvarwidth
\end{varwidth}%
&\begin{varwidth}[t]{\sphinxcolwidth{1}{4}}
\sphinxAtStartPar
(1/ωμ₀) × |Zyx|²
\sphinxbeforeendvarwidth
\end{varwidth}%
\\
\sphinxhline\begin{varwidth}[t]{\sphinxcolwidth{1}{4}}
\sphinxAtStartPar
ryy
\sphinxbeforeendvarwidth
\end{varwidth}%
&\begin{varwidth}[t]{\sphinxcolwidth{1}{4}}
\sphinxAtStartPar
MT\_ryy
\sphinxbeforeendvarwidth
\end{varwidth}%
&\begin{varwidth}[t]{\sphinxcolwidth{1}{4}}
\sphinxAtStartPar
ρyy视电阻率
\sphinxbeforeendvarwidth
\end{varwidth}%
&\begin{varwidth}[t]{\sphinxcolwidth{1}{4}}
\sphinxAtStartPar
(1/ωμ₀) × |Zyy|²
\sphinxbeforeendvarwidth
\end{varwidth}%
\\
\sphinxhline\begin{varwidth}[t]{\sphinxcolwidth{1}{4}}
\sphinxAtStartPar
lgrxx
\sphinxbeforeendvarwidth
\end{varwidth}%
&\begin{varwidth}[t]{\sphinxcolwidth{1}{4}}
\sphinxAtStartPar
MT\_lgrxx
\sphinxbeforeendvarwidth
\end{varwidth}%
&\begin{varwidth}[t]{\sphinxcolwidth{1}{4}}
\sphinxAtStartPar
log₁₀(ρxx)
\sphinxbeforeendvarwidth
\end{varwidth}%
&\begin{varwidth}[t]{\sphinxcolwidth{1}{4}}
\sphinxAtStartPar
log₁₀(rxx)
\sphinxbeforeendvarwidth
\end{varwidth}%
\\
\sphinxhline\begin{varwidth}[t]{\sphinxcolwidth{1}{4}}
\sphinxAtStartPar
lgrxy
\sphinxbeforeendvarwidth
\end{varwidth}%
&\begin{varwidth}[t]{\sphinxcolwidth{1}{4}}
\sphinxAtStartPar
MT\_lgrxy
\sphinxbeforeendvarwidth
\end{varwidth}%
&\begin{varwidth}[t]{\sphinxcolwidth{1}{4}}
\sphinxAtStartPar
log₁₀(ρxy)
\sphinxbeforeendvarwidth
\end{varwidth}%
&\begin{varwidth}[t]{\sphinxcolwidth{1}{4}}
\sphinxAtStartPar
log₁₀(rxy)
\sphinxbeforeendvarwidth
\end{varwidth}%
\\
\sphinxhline\begin{varwidth}[t]{\sphinxcolwidth{1}{4}}
\sphinxAtStartPar
lgryx
\sphinxbeforeendvarwidth
\end{varwidth}%
&\begin{varwidth}[t]{\sphinxcolwidth{1}{4}}
\sphinxAtStartPar
MT\_lgryx
\sphinxbeforeendvarwidth
\end{varwidth}%
&\begin{varwidth}[t]{\sphinxcolwidth{1}{4}}
\sphinxAtStartPar
log₁₀(ρyx)
\sphinxbeforeendvarwidth
\end{varwidth}%
&\begin{varwidth}[t]{\sphinxcolwidth{1}{4}}
\sphinxAtStartPar
log₁₀(ryx)
\sphinxbeforeendvarwidth
\end{varwidth}%
\\
\sphinxhline\begin{varwidth}[t]{\sphinxcolwidth{1}{4}}
\sphinxAtStartPar
lgryy
\sphinxbeforeendvarwidth
\end{varwidth}%
&\begin{varwidth}[t]{\sphinxcolwidth{1}{4}}
\sphinxAtStartPar
MT\_lgryy
\sphinxbeforeendvarwidth
\end{varwidth}%
&\begin{varwidth}[t]{\sphinxcolwidth{1}{4}}
\sphinxAtStartPar
log₁₀(ρyy)
\sphinxbeforeendvarwidth
\end{varwidth}%
&\begin{varwidth}[t]{\sphinxcolwidth{1}{4}}
\sphinxAtStartPar
log₁₀(ryy)
\sphinxbeforeendvarwidth
\end{varwidth}%
\\
\sphinxbottomrule
\end{tabulary}
\sphinxtableafterendhook\par
\sphinxattableend\end{savenotes}
\begin{quote}

\sphinxAtStartPar
\sphinxstylestrong{常用：} ρxy和ρyx是MT解释的主要参数
\end{quote}


\paragraph{5.17.6 阻抗方差}
\label{\detokenize{chapters/chapter6:id74}}

\begin{savenotes}\sphinxattablestart
\sphinxthistablewithglobalstyle
\centering
\begin{tabulary}{\linewidth}[t]{TTTT}
\sphinxtoprule
\begin{varwidth}[t]{\sphinxcolwidth{1}{4}}
\sphinxstyletheadfamily \sphinxAtStartPar
参数名
\sphinxbeforeendvarwidth
\end{varwidth}%
&\begin{varwidth}[t]{\sphinxcolwidth{1}{4}}
\sphinxstyletheadfamily \sphinxAtStartPar
内部名称
\sphinxbeforeendvarwidth
\end{varwidth}%
&\begin{varwidth}[t]{\sphinxcolwidth{1}{4}}
\sphinxstyletheadfamily \sphinxAtStartPar
说明
\sphinxbeforeendvarwidth
\end{varwidth}%
&\begin{varwidth}[t]{\sphinxcolwidth{1}{4}}
\sphinxstyletheadfamily \sphinxAtStartPar
用途
\sphinxbeforeendvarwidth
\end{varwidth}%
\\
\sphinxmidrule
\sphinxtableatstartofbodyhook\begin{varwidth}[t]{\sphinxcolwidth{1}{4}}
\sphinxAtStartPar
ZxxVar
\sphinxbeforeendvarwidth
\end{varwidth}%
&\begin{varwidth}[t]{\sphinxcolwidth{1}{4}}
\sphinxAtStartPar
MT\_ZxxVar
\sphinxbeforeendvarwidth
\end{varwidth}%
&\begin{varwidth}[t]{\sphinxcolwidth{1}{4}}
\sphinxAtStartPar
Zxx方差
\sphinxbeforeendvarwidth
\end{varwidth}%
&\begin{varwidth}[t]{\sphinxcolwidth{1}{4}}
\sphinxAtStartPar
数据质量评估
\sphinxbeforeendvarwidth
\end{varwidth}%
\\
\sphinxhline\begin{varwidth}[t]{\sphinxcolwidth{1}{4}}
\sphinxAtStartPar
ZxyVar
\sphinxbeforeendvarwidth
\end{varwidth}%
&\begin{varwidth}[t]{\sphinxcolwidth{1}{4}}
\sphinxAtStartPar
MT\_ZxyVar
\sphinxbeforeendvarwidth
\end{varwidth}%
&\begin{varwidth}[t]{\sphinxcolwidth{1}{4}}
\sphinxAtStartPar
Zxy方差
\sphinxbeforeendvarwidth
\end{varwidth}%
&\begin{varwidth}[t]{\sphinxcolwidth{1}{4}}
\sphinxAtStartPar
数据质量评估
\sphinxbeforeendvarwidth
\end{varwidth}%
\\
\sphinxhline\begin{varwidth}[t]{\sphinxcolwidth{1}{4}}
\sphinxAtStartPar
ZyxVar
\sphinxbeforeendvarwidth
\end{varwidth}%
&\begin{varwidth}[t]{\sphinxcolwidth{1}{4}}
\sphinxAtStartPar
MT\_ZyxVar
\sphinxbeforeendvarwidth
\end{varwidth}%
&\begin{varwidth}[t]{\sphinxcolwidth{1}{4}}
\sphinxAtStartPar
Zyx方差
\sphinxbeforeendvarwidth
\end{varwidth}%
&\begin{varwidth}[t]{\sphinxcolwidth{1}{4}}
\sphinxAtStartPar
数据质量评估
\sphinxbeforeendvarwidth
\end{varwidth}%
\\
\sphinxhline\begin{varwidth}[t]{\sphinxcolwidth{1}{4}}
\sphinxAtStartPar
ZyyVar
\sphinxbeforeendvarwidth
\end{varwidth}%
&\begin{varwidth}[t]{\sphinxcolwidth{1}{4}}
\sphinxAtStartPar
MT\_ZyyVar
\sphinxbeforeendvarwidth
\end{varwidth}%
&\begin{varwidth}[t]{\sphinxcolwidth{1}{4}}
\sphinxAtStartPar
Zyy方差
\sphinxbeforeendvarwidth
\end{varwidth}%
&\begin{varwidth}[t]{\sphinxcolwidth{1}{4}}
\sphinxAtStartPar
数据质量评估
\sphinxbeforeendvarwidth
\end{varwidth}%
\\
\sphinxhline\begin{varwidth}[t]{\sphinxcolwidth{1}{4}}
\sphinxAtStartPar
TzxVar
\sphinxbeforeendvarwidth
\end{varwidth}%
&\begin{varwidth}[t]{\sphinxcolwidth{1}{4}}
\sphinxAtStartPar
MT\_TzxVar
\sphinxbeforeendvarwidth
\end{varwidth}%
&\begin{varwidth}[t]{\sphinxcolwidth{1}{4}}
\sphinxAtStartPar
Tzx方差
\sphinxbeforeendvarwidth
\end{varwidth}%
&\begin{varwidth}[t]{\sphinxcolwidth{1}{4}}
\sphinxAtStartPar
数据质量评估
\sphinxbeforeendvarwidth
\end{varwidth}%
\\
\sphinxhline\begin{varwidth}[t]{\sphinxcolwidth{1}{4}}
\sphinxAtStartPar
TzyVar
\sphinxbeforeendvarwidth
\end{varwidth}%
&\begin{varwidth}[t]{\sphinxcolwidth{1}{4}}
\sphinxAtStartPar
MT\_TzyVar
\sphinxbeforeendvarwidth
\end{varwidth}%
&\begin{varwidth}[t]{\sphinxcolwidth{1}{4}}
\sphinxAtStartPar
Tzy方差
\sphinxbeforeendvarwidth
\end{varwidth}%
&\begin{varwidth}[t]{\sphinxcolwidth{1}{4}}
\sphinxAtStartPar
数据质量评估
\sphinxbeforeendvarwidth
\end{varwidth}%
\\
\sphinxbottomrule
\end{tabulary}
\sphinxtableafterendhook\par
\sphinxattableend\end{savenotes}


\paragraph{5.17.7 视电阻率与相位方差}
\label{\detokenize{chapters/chapter6:id75}}

\begin{savenotes}\sphinxattablestart
\sphinxthistablewithglobalstyle
\centering
\begin{tabulary}{\linewidth}[t]{TTTT}
\sphinxtoprule
\begin{varwidth}[t]{\sphinxcolwidth{1}{4}}
\sphinxstyletheadfamily \sphinxAtStartPar
参数名
\sphinxbeforeendvarwidth
\end{varwidth}%
&\begin{varwidth}[t]{\sphinxcolwidth{1}{4}}
\sphinxstyletheadfamily \sphinxAtStartPar
内部名称
\sphinxbeforeendvarwidth
\end{varwidth}%
&\begin{varwidth}[t]{\sphinxcolwidth{1}{4}}
\sphinxstyletheadfamily \sphinxAtStartPar
说明
\sphinxbeforeendvarwidth
\end{varwidth}%
&\begin{varwidth}[t]{\sphinxcolwidth{1}{4}}
\sphinxstyletheadfamily \sphinxAtStartPar
计算公式
\sphinxbeforeendvarwidth
\end{varwidth}%
\\
\sphinxmidrule
\sphinxtableatstartofbodyhook\begin{varwidth}[t]{\sphinxcolwidth{1}{4}}
\sphinxAtStartPar
rxxvar
\sphinxbeforeendvarwidth
\end{varwidth}%
&\begin{varwidth}[t]{\sphinxcolwidth{1}{4}}
\sphinxAtStartPar
MT\_rxxvar
\sphinxbeforeendvarwidth
\end{varwidth}%
&\begin{varwidth}[t]{\sphinxcolwidth{1}{4}}
\sphinxAtStartPar
ρxx方差
\sphinxbeforeendvarwidth
\end{varwidth}%
&\begin{varwidth}[t]{\sphinxcolwidth{1}{4}}
\sphinxAtStartPar
基于Zxx方差传播
\sphinxbeforeendvarwidth
\end{varwidth}%
\\
\sphinxhline\begin{varwidth}[t]{\sphinxcolwidth{1}{4}}
\sphinxAtStartPar
rxyvar
\sphinxbeforeendvarwidth
\end{varwidth}%
&\begin{varwidth}[t]{\sphinxcolwidth{1}{4}}
\sphinxAtStartPar
MT\_rxyvar
\sphinxbeforeendvarwidth
\end{varwidth}%
&\begin{varwidth}[t]{\sphinxcolwidth{1}{4}}
\sphinxAtStartPar
ρxy方差
\sphinxbeforeendvarwidth
\end{varwidth}%
&\begin{varwidth}[t]{\sphinxcolwidth{1}{4}}
\sphinxAtStartPar
基于Zxy方差传播
\sphinxbeforeendvarwidth
\end{varwidth}%
\\
\sphinxhline\begin{varwidth}[t]{\sphinxcolwidth{1}{4}}
\sphinxAtStartPar
ryxvar
\sphinxbeforeendvarwidth
\end{varwidth}%
&\begin{varwidth}[t]{\sphinxcolwidth{1}{4}}
\sphinxAtStartPar
MT\_ryxvar
\sphinxbeforeendvarwidth
\end{varwidth}%
&\begin{varwidth}[t]{\sphinxcolwidth{1}{4}}
\sphinxAtStartPar
ρyx方差
\sphinxbeforeendvarwidth
\end{varwidth}%
&\begin{varwidth}[t]{\sphinxcolwidth{1}{4}}
\sphinxAtStartPar
基于Zyx方差传播
\sphinxbeforeendvarwidth
\end{varwidth}%
\\
\sphinxhline\begin{varwidth}[t]{\sphinxcolwidth{1}{4}}
\sphinxAtStartPar
ryyvar
\sphinxbeforeendvarwidth
\end{varwidth}%
&\begin{varwidth}[t]{\sphinxcolwidth{1}{4}}
\sphinxAtStartPar
MT\_ryyvar
\sphinxbeforeendvarwidth
\end{varwidth}%
&\begin{varwidth}[t]{\sphinxcolwidth{1}{4}}
\sphinxAtStartPar
ρyy方差
\sphinxbeforeendvarwidth
\end{varwidth}%
&\begin{varwidth}[t]{\sphinxcolwidth{1}{4}}
\sphinxAtStartPar
基于Zyy方差传播
\sphinxbeforeendvarwidth
\end{varwidth}%
\\
\sphinxhline\begin{varwidth}[t]{\sphinxcolwidth{1}{4}}
\sphinxAtStartPar
lgrxxvar
\sphinxbeforeendvarwidth
\end{varwidth}%
&\begin{varwidth}[t]{\sphinxcolwidth{1}{4}}
\sphinxAtStartPar
MT\_lgrxxvar
\sphinxbeforeendvarwidth
\end{varwidth}%
&\begin{varwidth}[t]{\sphinxcolwidth{1}{4}}
\sphinxAtStartPar
log(ρxx)方差
\sphinxbeforeendvarwidth
\end{varwidth}%
&\begin{varwidth}[t]{\sphinxcolwidth{1}{4}}
\sphinxAtStartPar
数据筛选
\sphinxbeforeendvarwidth
\end{varwidth}%
\\
\sphinxhline\begin{varwidth}[t]{\sphinxcolwidth{1}{4}}
\sphinxAtStartPar
lgrxyvar
\sphinxbeforeendvarwidth
\end{varwidth}%
&\begin{varwidth}[t]{\sphinxcolwidth{1}{4}}
\sphinxAtStartPar
MT\_lgrxyvar
\sphinxbeforeendvarwidth
\end{varwidth}%
&\begin{varwidth}[t]{\sphinxcolwidth{1}{4}}
\sphinxAtStartPar
log(ρxy)方差
\sphinxbeforeendvarwidth
\end{varwidth}%
&\begin{varwidth}[t]{\sphinxcolwidth{1}{4}}
\sphinxAtStartPar
数据筛选
\sphinxbeforeendvarwidth
\end{varwidth}%
\\
\sphinxhline\begin{varwidth}[t]{\sphinxcolwidth{1}{4}}
\sphinxAtStartPar
lgryxvar
\sphinxbeforeendvarwidth
\end{varwidth}%
&\begin{varwidth}[t]{\sphinxcolwidth{1}{4}}
\sphinxAtStartPar
MT\_lgryxvar
\sphinxbeforeendvarwidth
\end{varwidth}%
&\begin{varwidth}[t]{\sphinxcolwidth{1}{4}}
\sphinxAtStartPar
log(ρyx)方差
\sphinxbeforeendvarwidth
\end{varwidth}%
&\begin{varwidth}[t]{\sphinxcolwidth{1}{4}}
\sphinxAtStartPar
数据筛选
\sphinxbeforeendvarwidth
\end{varwidth}%
\\
\sphinxhline\begin{varwidth}[t]{\sphinxcolwidth{1}{4}}
\sphinxAtStartPar
lgryyvar
\sphinxbeforeendvarwidth
\end{varwidth}%
&\begin{varwidth}[t]{\sphinxcolwidth{1}{4}}
\sphinxAtStartPar
MT\_lgryyvar
\sphinxbeforeendvarwidth
\end{varwidth}%
&\begin{varwidth}[t]{\sphinxcolwidth{1}{4}}
\sphinxAtStartPar
log(ρyy)方差
\sphinxbeforeendvarwidth
\end{varwidth}%
&\begin{varwidth}[t]{\sphinxcolwidth{1}{4}}
\sphinxAtStartPar
数据筛选
\sphinxbeforeendvarwidth
\end{varwidth}%
\\
\sphinxhline\begin{varwidth}[t]{\sphinxcolwidth{1}{4}}
\sphinxAtStartPar
pxxvar
\sphinxbeforeendvarwidth
\end{varwidth}%
&\begin{varwidth}[t]{\sphinxcolwidth{1}{4}}
\sphinxAtStartPar
MT\_pxxvar
\sphinxbeforeendvarwidth
\end{varwidth}%
&\begin{varwidth}[t]{\sphinxcolwidth{1}{4}}
\sphinxAtStartPar
φxx方差
\sphinxbeforeendvarwidth
\end{varwidth}%
&\begin{varwidth}[t]{\sphinxcolwidth{1}{4}}
\sphinxAtStartPar
相位误差
\sphinxbeforeendvarwidth
\end{varwidth}%
\\
\sphinxhline\begin{varwidth}[t]{\sphinxcolwidth{1}{4}}
\sphinxAtStartPar
pxyvar
\sphinxbeforeendvarwidth
\end{varwidth}%
&\begin{varwidth}[t]{\sphinxcolwidth{1}{4}}
\sphinxAtStartPar
MT\_pxyvar
\sphinxbeforeendvarwidth
\end{varwidth}%
&\begin{varwidth}[t]{\sphinxcolwidth{1}{4}}
\sphinxAtStartPar
φxy方差
\sphinxbeforeendvarwidth
\end{varwidth}%
&\begin{varwidth}[t]{\sphinxcolwidth{1}{4}}
\sphinxAtStartPar
相位误差
\sphinxbeforeendvarwidth
\end{varwidth}%
\\
\sphinxhline\begin{varwidth}[t]{\sphinxcolwidth{1}{4}}
\sphinxAtStartPar
pyxvar
\sphinxbeforeendvarwidth
\end{varwidth}%
&\begin{varwidth}[t]{\sphinxcolwidth{1}{4}}
\sphinxAtStartPar
MT\_pyxvar
\sphinxbeforeendvarwidth
\end{varwidth}%
&\begin{varwidth}[t]{\sphinxcolwidth{1}{4}}
\sphinxAtStartPar
φyx方差
\sphinxbeforeendvarwidth
\end{varwidth}%
&\begin{varwidth}[t]{\sphinxcolwidth{1}{4}}
\sphinxAtStartPar
相位误差
\sphinxbeforeendvarwidth
\end{varwidth}%
\\
\sphinxhline\begin{varwidth}[t]{\sphinxcolwidth{1}{4}}
\sphinxAtStartPar
pyyvar
\sphinxbeforeendvarwidth
\end{varwidth}%
&\begin{varwidth}[t]{\sphinxcolwidth{1}{4}}
\sphinxAtStartPar
MT\_pyyvar
\sphinxbeforeendvarwidth
\end{varwidth}%
&\begin{varwidth}[t]{\sphinxcolwidth{1}{4}}
\sphinxAtStartPar
φyy方差
\sphinxbeforeendvarwidth
\end{varwidth}%
&\begin{varwidth}[t]{\sphinxcolwidth{1}{4}}
\sphinxAtStartPar
相位误差
\sphinxbeforeendvarwidth
\end{varwidth}%
\\
\sphinxbottomrule
\end{tabulary}
\sphinxtableafterendhook\par
\sphinxattableend\end{savenotes}


\paragraph{5.17.8 相位张量}
\label{\detokenize{chapters/chapter6:id76}}
\sphinxAtStartPar
相位张量Φ是归一化阻抗张量的实部，用于分析地下构造：

\sphinxAtStartPar
\sphinxstylestrong{理论公式：}

\begin{sphinxVerbatim}[commandchars=\\\{\}]
Φ = Re(Z) / |ωμ₀| = [Φxx Φxy]
                        [Φyx Φyy]
\end{sphinxVerbatim}


\begin{savenotes}\sphinxattablestart
\sphinxthistablewithglobalstyle
\centering
\begin{tabulary}{\linewidth}[t]{TTTT}
\sphinxtoprule
\begin{varwidth}[t]{\sphinxcolwidth{1}{4}}
\sphinxstyletheadfamily \sphinxAtStartPar
参数名
\sphinxbeforeendvarwidth
\end{varwidth}%
&\begin{varwidth}[t]{\sphinxcolwidth{1}{4}}
\sphinxstyletheadfamily \sphinxAtStartPar
内部名称
\sphinxbeforeendvarwidth
\end{varwidth}%
&\begin{varwidth}[t]{\sphinxcolwidth{1}{4}}
\sphinxstyletheadfamily \sphinxAtStartPar
说明
\sphinxbeforeendvarwidth
\end{varwidth}%
&\begin{varwidth}[t]{\sphinxcolwidth{1}{4}}
\sphinxstyletheadfamily \sphinxAtStartPar
用途
\sphinxbeforeendvarwidth
\end{varwidth}%
\\
\sphinxmidrule
\sphinxtableatstartofbodyhook\begin{varwidth}[t]{\sphinxcolwidth{1}{4}}
\sphinxAtStartPar
ptxx
\sphinxbeforeendvarwidth
\end{varwidth}%
&\begin{varwidth}[t]{\sphinxcolwidth{1}{4}}
\sphinxAtStartPar
MT\_ptxx
\sphinxbeforeendvarwidth
\end{varwidth}%
&\begin{varwidth}[t]{\sphinxcolwidth{1}{4}}
\sphinxAtStartPar
Φxx
\sphinxbeforeendvarwidth
\end{varwidth}%
&\begin{varwidth}[t]{\sphinxcolwidth{1}{4}}
\sphinxAtStartPar
相位张量分量
\sphinxbeforeendvarwidth
\end{varwidth}%
\\
\sphinxhline\begin{varwidth}[t]{\sphinxcolwidth{1}{4}}
\sphinxAtStartPar
ptxy
\sphinxbeforeendvarwidth
\end{varwidth}%
&\begin{varwidth}[t]{\sphinxcolwidth{1}{4}}
\sphinxAtStartPar
MT\_ptxy
\sphinxbeforeendvarwidth
\end{varwidth}%
&\begin{varwidth}[t]{\sphinxcolwidth{1}{4}}
\sphinxAtStartPar
Φxy
\sphinxbeforeendvarwidth
\end{varwidth}%
&\begin{varwidth}[t]{\sphinxcolwidth{1}{4}}
\sphinxAtStartPar
相位张量分量
\sphinxbeforeendvarwidth
\end{varwidth}%
\\
\sphinxhline\begin{varwidth}[t]{\sphinxcolwidth{1}{4}}
\sphinxAtStartPar
ptyx
\sphinxbeforeendvarwidth
\end{varwidth}%
&\begin{varwidth}[t]{\sphinxcolwidth{1}{4}}
\sphinxAtStartPar
MT\_ptyx
\sphinxbeforeendvarwidth
\end{varwidth}%
&\begin{varwidth}[t]{\sphinxcolwidth{1}{4}}
\sphinxAtStartPar
Φyx
\sphinxbeforeendvarwidth
\end{varwidth}%
&\begin{varwidth}[t]{\sphinxcolwidth{1}{4}}
\sphinxAtStartPar
相位张量分量
\sphinxbeforeendvarwidth
\end{varwidth}%
\\
\sphinxhline\begin{varwidth}[t]{\sphinxcolwidth{1}{4}}
\sphinxAtStartPar
ptyy
\sphinxbeforeendvarwidth
\end{varwidth}%
&\begin{varwidth}[t]{\sphinxcolwidth{1}{4}}
\sphinxAtStartPar
MT\_ptyy
\sphinxbeforeendvarwidth
\end{varwidth}%
&\begin{varwidth}[t]{\sphinxcolwidth{1}{4}}
\sphinxAtStartPar
Φyy
\sphinxbeforeendvarwidth
\end{varwidth}%
&\begin{varwidth}[t]{\sphinxcolwidth{1}{4}}
\sphinxAtStartPar
相位张量分量
\sphinxbeforeendvarwidth
\end{varwidth}%
\\
\sphinxhline\begin{varwidth}[t]{\sphinxcolwidth{1}{4}}
\sphinxAtStartPar
alpha
\sphinxbeforeendvarwidth
\end{varwidth}%
&\begin{varwidth}[t]{\sphinxcolwidth{1}{4}}
\sphinxAtStartPar
MT\_alpha
\sphinxbeforeendvarwidth
\end{varwidth}%
&\begin{varwidth}[t]{\sphinxcolwidth{1}{4}}
\sphinxAtStartPar
α角
\sphinxbeforeendvarwidth
\end{varwidth}%
&\begin{varwidth}[t]{\sphinxcolwidth{1}{4}}
\sphinxAtStartPar
主轴方向角
\sphinxbeforeendvarwidth
\end{varwidth}%
\\
\sphinxhline\begin{varwidth}[t]{\sphinxcolwidth{1}{4}}
\sphinxAtStartPar
beta
\sphinxbeforeendvarwidth
\end{varwidth}%
&\begin{varwidth}[t]{\sphinxcolwidth{1}{4}}
\sphinxAtStartPar
MT\_beta
\sphinxbeforeendvarwidth
\end{varwidth}%
&\begin{varwidth}[t]{\sphinxcolwidth{1}{4}}
\sphinxAtStartPar
β角
\sphinxbeforeendvarwidth
\end{varwidth}%
&\begin{varwidth}[t]{\sphinxcolwidth{1}{4}}
\sphinxAtStartPar
构造三维性指示
\sphinxbeforeendvarwidth
\end{varwidth}%
\\
\sphinxhline\begin{varwidth}[t]{\sphinxcolwidth{1}{4}}
\sphinxAtStartPar
pmax
\sphinxbeforeendvarwidth
\end{varwidth}%
&\begin{varwidth}[t]{\sphinxcolwidth{1}{4}}
\sphinxAtStartPar
MT\_pmax
\sphinxbeforeendvarwidth
\end{varwidth}%
&\begin{varwidth}[t]{\sphinxcolwidth{1}{4}}
\sphinxAtStartPar
Φmax
\sphinxbeforeendvarwidth
\end{varwidth}%
&\begin{varwidth}[t]{\sphinxcolwidth{1}{4}}
\sphinxAtStartPar
最大主值
\sphinxbeforeendvarwidth
\end{varwidth}%
\\
\sphinxhline\begin{varwidth}[t]{\sphinxcolwidth{1}{4}}
\sphinxAtStartPar
pmin
\sphinxbeforeendvarwidth
\end{varwidth}%
&\begin{varwidth}[t]{\sphinxcolwidth{1}{4}}
\sphinxAtStartPar
MT\_pmin
\sphinxbeforeendvarwidth
\end{varwidth}%
&\begin{varwidth}[t]{\sphinxcolwidth{1}{4}}
\sphinxAtStartPar
Φmin
\sphinxbeforeendvarwidth
\end{varwidth}%
&\begin{varwidth}[t]{\sphinxcolwidth{1}{4}}
\sphinxAtStartPar
最小主值
\sphinxbeforeendvarwidth
\end{varwidth}%
\\
\sphinxhline\begin{varwidth}[t]{\sphinxcolwidth{1}{4}}
\sphinxAtStartPar
ptskew1d
\sphinxbeforeendvarwidth
\end{varwidth}%
&\begin{varwidth}[t]{\sphinxcolwidth{1}{4}}
\sphinxAtStartPar
MT\_ptskew1d
\sphinxbeforeendvarwidth
\end{varwidth}%
&\begin{varwidth}[t]{\sphinxcolwidth{1}{4}}
\sphinxAtStartPar
1D偏斜度
\sphinxbeforeendvarwidth
\end{varwidth}%
&\begin{varwidth}[t]{\sphinxcolwidth{1}{4}}
\sphinxAtStartPar
一维偏离度
\sphinxbeforeendvarwidth
\end{varwidth}%
\\
\sphinxhline\begin{varwidth}[t]{\sphinxcolwidth{1}{4}}
\sphinxAtStartPar
ptskew2d
\sphinxbeforeendvarwidth
\end{varwidth}%
&\begin{varwidth}[t]{\sphinxcolwidth{1}{4}}
\sphinxAtStartPar
MT\_ptskew2d
\sphinxbeforeendvarwidth
\end{varwidth}%
&\begin{varwidth}[t]{\sphinxcolwidth{1}{4}}
\sphinxAtStartPar
2D偏斜度
\sphinxbeforeendvarwidth
\end{varwidth}%
&\begin{varwidth}[t]{\sphinxcolwidth{1}{4}}
\sphinxAtStartPar
二维偏离度
\sphinxbeforeendvarwidth
\end{varwidth}%
\\
\sphinxbottomrule
\end{tabulary}
\sphinxtableafterendhook\par
\sphinxattableend\end{savenotes}

\sphinxAtStartPar
\sphinxstylestrong{解释：}
\begin{itemize}
\item {} 
\sphinxAtStartPar
\sphinxstylestrong{β ≈ 0°}：二维构造

\item {} 
\sphinxAtStartPar
\sphinxstylestrong{β > 0°}：三维构造

\item {} 
\sphinxAtStartPar
\sphinxstylestrong{skew较小}：近二维构造

\end{itemize}


\paragraph{5.17.9 共轭阻抗变换}
\label{\detokenize{chapters/chapter6:id77}}

\begin{savenotes}\sphinxattablestart
\sphinxthistablewithglobalstyle
\centering
\begin{tabulary}{\linewidth}[t]{TTTT}
\sphinxtoprule
\begin{varwidth}[t]{\sphinxcolwidth{1}{4}}
\sphinxstyletheadfamily \sphinxAtStartPar
参数名
\sphinxbeforeendvarwidth
\end{varwidth}%
&\begin{varwidth}[t]{\sphinxcolwidth{1}{4}}
\sphinxstyletheadfamily \sphinxAtStartPar
内部名称
\sphinxbeforeendvarwidth
\end{varwidth}%
&\begin{varwidth}[t]{\sphinxcolwidth{1}{4}}
\sphinxstyletheadfamily \sphinxAtStartPar
说明
\sphinxbeforeendvarwidth
\end{varwidth}%
&\begin{varwidth}[t]{\sphinxcolwidth{1}{4}}
\sphinxstyletheadfamily \sphinxAtStartPar
用途
\sphinxbeforeendvarwidth
\end{varwidth}%
\\
\sphinxmidrule
\sphinxtableatstartofbodyhook\begin{varwidth}[t]{\sphinxcolwidth{1}{4}}
\sphinxAtStartPar
cczxxr
\sphinxbeforeendvarwidth
\end{varwidth}%
&\begin{varwidth}[t]{\sphinxcolwidth{1}{4}}
\sphinxAtStartPar
MT\_cczxxr
\sphinxbeforeendvarwidth
\end{varwidth}%
&\begin{varwidth}[t]{\sphinxcolwidth{1}{4}}
\sphinxAtStartPar
共轭Zxx实部
\sphinxbeforeendvarwidth
\end{varwidth}%
&\begin{varwidth}[t]{\sphinxcolwidth{1}{4}}
\sphinxAtStartPar
构造分析
\sphinxbeforeendvarwidth
\end{varwidth}%
\\
\sphinxhline\begin{varwidth}[t]{\sphinxcolwidth{1}{4}}
\sphinxAtStartPar
cczxxi
\sphinxbeforeendvarwidth
\end{varwidth}%
&\begin{varwidth}[t]{\sphinxcolwidth{1}{4}}
\sphinxAtStartPar
MT\_cczxxi
\sphinxbeforeendvarwidth
\end{varwidth}%
&\begin{varwidth}[t]{\sphinxcolwidth{1}{4}}
\sphinxAtStartPar
共轭Zxx虚部
\sphinxbeforeendvarwidth
\end{varwidth}%
&\begin{varwidth}[t]{\sphinxcolwidth{1}{4}}
\sphinxAtStartPar
构造分析
\sphinxbeforeendvarwidth
\end{varwidth}%
\\
\sphinxhline\begin{varwidth}[t]{\sphinxcolwidth{1}{4}}
\sphinxAtStartPar
cczxyr
\sphinxbeforeendvarwidth
\end{varwidth}%
&\begin{varwidth}[t]{\sphinxcolwidth{1}{4}}
\sphinxAtStartPar
MT\_cczxyr
\sphinxbeforeendvarwidth
\end{varwidth}%
&\begin{varwidth}[t]{\sphinxcolwidth{1}{4}}
\sphinxAtStartPar
共轭Zxy实部
\sphinxbeforeendvarwidth
\end{varwidth}%
&\begin{varwidth}[t]{\sphinxcolwidth{1}{4}}
\sphinxAtStartPar
构造分析
\sphinxbeforeendvarwidth
\end{varwidth}%
\\
\sphinxhline\begin{varwidth}[t]{\sphinxcolwidth{1}{4}}
\sphinxAtStartPar
cczxyi
\sphinxbeforeendvarwidth
\end{varwidth}%
&\begin{varwidth}[t]{\sphinxcolwidth{1}{4}}
\sphinxAtStartPar
MT\_cczxyi
\sphinxbeforeendvarwidth
\end{varwidth}%
&\begin{varwidth}[t]{\sphinxcolwidth{1}{4}}
\sphinxAtStartPar
共轭Zxy虚部
\sphinxbeforeendvarwidth
\end{varwidth}%
&\begin{varwidth}[t]{\sphinxcolwidth{1}{4}}
\sphinxAtStartPar
构造分析
\sphinxbeforeendvarwidth
\end{varwidth}%
\\
\sphinxhline\begin{varwidth}[t]{\sphinxcolwidth{1}{4}}
\sphinxAtStartPar
cczyxr
\sphinxbeforeendvarwidth
\end{varwidth}%
&\begin{varwidth}[t]{\sphinxcolwidth{1}{4}}
\sphinxAtStartPar
MT\_cczyxr
\sphinxbeforeendvarwidth
\end{varwidth}%
&\begin{varwidth}[t]{\sphinxcolwidth{1}{4}}
\sphinxAtStartPar
共轭Zyx实部
\sphinxbeforeendvarwidth
\end{varwidth}%
&\begin{varwidth}[t]{\sphinxcolwidth{1}{4}}
\sphinxAtStartPar
构造分析
\sphinxbeforeendvarwidth
\end{varwidth}%
\\
\sphinxhline\begin{varwidth}[t]{\sphinxcolwidth{1}{4}}
\sphinxAtStartPar
cczyxi
\sphinxbeforeendvarwidth
\end{varwidth}%
&\begin{varwidth}[t]{\sphinxcolwidth{1}{4}}
\sphinxAtStartPar
MT\_cczyxi
\sphinxbeforeendvarwidth
\end{varwidth}%
&\begin{varwidth}[t]{\sphinxcolwidth{1}{4}}
\sphinxAtStartPar
共轭Zyx虚部
\sphinxbeforeendvarwidth
\end{varwidth}%
&\begin{varwidth}[t]{\sphinxcolwidth{1}{4}}
\sphinxAtStartPar
构造分析
\sphinxbeforeendvarwidth
\end{varwidth}%
\\
\sphinxhline\begin{varwidth}[t]{\sphinxcolwidth{1}{4}}
\sphinxAtStartPar
cczyyr
\sphinxbeforeendvarwidth
\end{varwidth}%
&\begin{varwidth}[t]{\sphinxcolwidth{1}{4}}
\sphinxAtStartPar
MT\_cczyyr
\sphinxbeforeendvarwidth
\end{varwidth}%
&\begin{varwidth}[t]{\sphinxcolwidth{1}{4}}
\sphinxAtStartPar
共轭Zyy实部
\sphinxbeforeendvarwidth
\end{varwidth}%
&\begin{varwidth}[t]{\sphinxcolwidth{1}{4}}
\sphinxAtStartPar
构造分析
\sphinxbeforeendvarwidth
\end{varwidth}%
\\
\sphinxhline\begin{varwidth}[t]{\sphinxcolwidth{1}{4}}
\sphinxAtStartPar
cczyyi
\sphinxbeforeendvarwidth
\end{varwidth}%
&\begin{varwidth}[t]{\sphinxcolwidth{1}{4}}
\sphinxAtStartPar
MT\_cczyyi
\sphinxbeforeendvarwidth
\end{varwidth}%
&\begin{varwidth}[t]{\sphinxcolwidth{1}{4}}
\sphinxAtStartPar
共轭Zyy虚部
\sphinxbeforeendvarwidth
\end{varwidth}%
&\begin{varwidth}[t]{\sphinxcolwidth{1}{4}}
\sphinxAtStartPar
构造分析
\sphinxbeforeendvarwidth
\end{varwidth}%
\\
\sphinxhline\begin{varwidth}[t]{\sphinxcolwidth{1}{4}}
\sphinxAtStartPar
ccztheta
\sphinxbeforeendvarwidth
\end{varwidth}%
&\begin{varwidth}[t]{\sphinxcolwidth{1}{4}}
\sphinxAtStartPar
MT\_ccztheta
\sphinxbeforeendvarwidth
\end{varwidth}%
&\begin{varwidth}[t]{\sphinxcolwidth{1}{4}}
\sphinxAtStartPar
共轭旋转角
\sphinxbeforeendvarwidth
\end{varwidth}%
&\begin{varwidth}[t]{\sphinxcolwidth{1}{4}}
\sphinxAtStartPar
主轴方向
\sphinxbeforeendvarwidth
\end{varwidth}%
\\
\sphinxhline\begin{varwidth}[t]{\sphinxcolwidth{1}{4}}
\sphinxAtStartPar
cczskew1d
\sphinxbeforeendvarwidth
\end{varwidth}%
&\begin{varwidth}[t]{\sphinxcolwidth{1}{4}}
\sphinxAtStartPar
MT\_cczskew1d
\sphinxbeforeendvarwidth
\end{varwidth}%
&\begin{varwidth}[t]{\sphinxcolwidth{1}{4}}
\sphinxAtStartPar
共轭1D偏斜
\sphinxbeforeendvarwidth
\end{varwidth}%
&\begin{varwidth}[t]{\sphinxcolwidth{1}{4}}
\sphinxAtStartPar
构造维数
\sphinxbeforeendvarwidth
\end{varwidth}%
\\
\sphinxhline\begin{varwidth}[t]{\sphinxcolwidth{1}{4}}
\sphinxAtStartPar
cczskew2d
\sphinxbeforeendvarwidth
\end{varwidth}%
&\begin{varwidth}[t]{\sphinxcolwidth{1}{4}}
\sphinxAtStartPar
MT\_cczskew2d
\sphinxbeforeendvarwidth
\end{varwidth}%
&\begin{varwidth}[t]{\sphinxcolwidth{1}{4}}
\sphinxAtStartPar
共轭2D偏斜
\sphinxbeforeendvarwidth
\end{varwidth}%
&\begin{varwidth}[t]{\sphinxcolwidth{1}{4}}
\sphinxAtStartPar
构造维数
\sphinxbeforeendvarwidth
\end{varwidth}%
\\
\sphinxbottomrule
\end{tabulary}
\sphinxtableafterendhook\par
\sphinxattableend\end{savenotes}


\paragraph{5.17.10 相干度}
\label{\detokenize{chapters/chapter6:id78}}
\sphinxAtStartPar
相干度衡量两个信号之间的线性相关程度：

\sphinxAtStartPar
\sphinxstylestrong{理论公式：}

\begin{sphinxVerbatim}[commandchars=\\\{\}]
γ² = |\PYGZlt{}XY*\PYGZgt{}|² / (\PYGZlt{}XX*\PYGZgt{}\PYGZlt{}YY*\PYGZgt{})
范围：0 ≤ γ² ≤ 1
\end{sphinxVerbatim}

\sphinxAtStartPar
\sphinxstylestrong{通道间相干度：}


\begin{savenotes}\sphinxattablestart
\sphinxthistablewithglobalstyle
\centering
\begin{tabulary}{\linewidth}[t]{TTTT}
\sphinxtoprule
\begin{varwidth}[t]{\sphinxcolwidth{1}{4}}
\sphinxstyletheadfamily \sphinxAtStartPar
参数名
\sphinxbeforeendvarwidth
\end{varwidth}%
&\begin{varwidth}[t]{\sphinxcolwidth{1}{4}}
\sphinxstyletheadfamily \sphinxAtStartPar
内部名称
\sphinxbeforeendvarwidth
\end{varwidth}%
&\begin{varwidth}[t]{\sphinxcolwidth{1}{4}}
\sphinxstyletheadfamily \sphinxAtStartPar
说明
\sphinxbeforeendvarwidth
\end{varwidth}%
&\begin{varwidth}[t]{\sphinxcolwidth{1}{4}}
\sphinxstyletheadfamily \sphinxAtStartPar
典型范围
\sphinxbeforeendvarwidth
\end{varwidth}%
\\
\sphinxmidrule
\sphinxtableatstartofbodyhook\begin{varwidth}[t]{\sphinxcolwidth{1}{4}}
\sphinxAtStartPar
CohExHy
\sphinxbeforeendvarwidth
\end{varwidth}%
&\begin{varwidth}[t]{\sphinxcolwidth{1}{4}}
\sphinxAtStartPar
MT\_CohExHy
\sphinxbeforeendvarwidth
\end{varwidth}%
&\begin{varwidth}[t]{\sphinxcolwidth{1}{4}}
\sphinxAtStartPar
Ex\sphinxhyphen{}Hy相干度
\sphinxbeforeendvarwidth
\end{varwidth}%
&\begin{varwidth}[t]{\sphinxcolwidth{1}{4}}
\sphinxAtStartPar
0.5\sphinxhyphen{}0.95
\sphinxbeforeendvarwidth
\end{varwidth}%
\\
\sphinxhline\begin{varwidth}[t]{\sphinxcolwidth{1}{4}}
\sphinxAtStartPar
CohExHx
\sphinxbeforeendvarwidth
\end{varwidth}%
&\begin{varwidth}[t]{\sphinxcolwidth{1}{4}}
\sphinxAtStartPar
MT\_CohExHx
\sphinxbeforeendvarwidth
\end{varwidth}%
&\begin{varwidth}[t]{\sphinxcolwidth{1}{4}}
\sphinxAtStartPar
Ex\sphinxhyphen{}Hx相干度
\sphinxbeforeendvarwidth
\end{varwidth}%
&\begin{varwidth}[t]{\sphinxcolwidth{1}{4}}
\sphinxAtStartPar
0.3\sphinxhyphen{}0.8
\sphinxbeforeendvarwidth
\end{varwidth}%
\\
\sphinxhline\begin{varwidth}[t]{\sphinxcolwidth{1}{4}}
\sphinxAtStartPar
CohEyHx
\sphinxbeforeendvarwidth
\end{varwidth}%
&\begin{varwidth}[t]{\sphinxcolwidth{1}{4}}
\sphinxAtStartPar
MT\_CohEyHx
\sphinxbeforeendvarwidth
\end{varwidth}%
&\begin{varwidth}[t]{\sphinxcolwidth{1}{4}}
\sphinxAtStartPar
Ey\sphinxhyphen{}Hx相干度
\sphinxbeforeendvarwidth
\end{varwidth}%
&\begin{varwidth}[t]{\sphinxcolwidth{1}{4}}
\sphinxAtStartPar
0.5\sphinxhyphen{}0.95
\sphinxbeforeendvarwidth
\end{varwidth}%
\\
\sphinxhline\begin{varwidth}[t]{\sphinxcolwidth{1}{4}}
\sphinxAtStartPar
CohEyHy
\sphinxbeforeendvarwidth
\end{varwidth}%
&\begin{varwidth}[t]{\sphinxcolwidth{1}{4}}
\sphinxAtStartPar
MT\_CohEyHy
\sphinxbeforeendvarwidth
\end{varwidth}%
&\begin{varwidth}[t]{\sphinxcolwidth{1}{4}}
\sphinxAtStartPar
Ey\sphinxhyphen{}Hy相干度
\sphinxbeforeendvarwidth
\end{varwidth}%
&\begin{varwidth}[t]{\sphinxcolwidth{1}{4}}
\sphinxAtStartPar
0.3\sphinxhyphen{}0.8
\sphinxbeforeendvarwidth
\end{varwidth}%
\\
\sphinxbottomrule
\end{tabulary}
\sphinxtableafterendhook\par
\sphinxattableend\end{savenotes}

\sphinxAtStartPar
\sphinxstylestrong{远参考相干度（带Rx/Ry）：}


\begin{savenotes}\sphinxattablestart
\sphinxthistablewithglobalstyle
\centering
\begin{tabulary}{\linewidth}[t]{TTT}
\sphinxtoprule
\begin{varwidth}[t]{\sphinxcolwidth{1}{3}}
\sphinxstyletheadfamily \sphinxAtStartPar
参数名
\sphinxbeforeendvarwidth
\end{varwidth}%
&\begin{varwidth}[t]{\sphinxcolwidth{1}{3}}
\sphinxstyletheadfamily \sphinxAtStartPar
内部名称
\sphinxbeforeendvarwidth
\end{varwidth}%
&\begin{varwidth}[t]{\sphinxcolwidth{1}{3}}
\sphinxstyletheadfamily \sphinxAtStartPar
说明
\sphinxbeforeendvarwidth
\end{varwidth}%
\\
\sphinxmidrule
\sphinxtableatstartofbodyhook\begin{varwidth}[t]{\sphinxcolwidth{1}{3}}
\sphinxAtStartPar
CohExRx
\sphinxbeforeendvarwidth
\end{varwidth}%
&\begin{varwidth}[t]{\sphinxcolwidth{1}{3}}
\sphinxAtStartPar
MT\_CohExRx
\sphinxbeforeendvarwidth
\end{varwidth}%
&\begin{varwidth}[t]{\sphinxcolwidth{1}{3}}
\sphinxAtStartPar
Ex\sphinxhyphen{}远参考X相干度
\sphinxbeforeendvarwidth
\end{varwidth}%
\\
\sphinxhline\begin{varwidth}[t]{\sphinxcolwidth{1}{3}}
\sphinxAtStartPar
CohExRy
\sphinxbeforeendvarwidth
\end{varwidth}%
&\begin{varwidth}[t]{\sphinxcolwidth{1}{3}}
\sphinxAtStartPar
MT\_CohExRy
\sphinxbeforeendvarwidth
\end{varwidth}%
&\begin{varwidth}[t]{\sphinxcolwidth{1}{3}}
\sphinxAtStartPar
Ex\sphinxhyphen{}远参考Y相干度
\sphinxbeforeendvarwidth
\end{varwidth}%
\\
\sphinxhline\begin{varwidth}[t]{\sphinxcolwidth{1}{3}}
\sphinxAtStartPar
CohEyRx
\sphinxbeforeendvarwidth
\end{varwidth}%
&\begin{varwidth}[t]{\sphinxcolwidth{1}{3}}
\sphinxAtStartPar
MT\_CohEyRx
\sphinxbeforeendvarwidth
\end{varwidth}%
&\begin{varwidth}[t]{\sphinxcolwidth{1}{3}}
\sphinxAtStartPar
Ey\sphinxhyphen{}远参考X相干度
\sphinxbeforeendvarwidth
\end{varwidth}%
\\
\sphinxhline\begin{varwidth}[t]{\sphinxcolwidth{1}{3}}
\sphinxAtStartPar
CohEyRy
\sphinxbeforeendvarwidth
\end{varwidth}%
&\begin{varwidth}[t]{\sphinxcolwidth{1}{3}}
\sphinxAtStartPar
MT\_CohEyRy
\sphinxbeforeendvarwidth
\end{varwidth}%
&\begin{varwidth}[t]{\sphinxcolwidth{1}{3}}
\sphinxAtStartPar
Ey\sphinxhyphen{}远参考Y相干度
\sphinxbeforeendvarwidth
\end{varwidth}%
\\
\sphinxhline\begin{varwidth}[t]{\sphinxcolwidth{1}{3}}
\sphinxAtStartPar
CohHxRx
\sphinxbeforeendvarwidth
\end{varwidth}%
&\begin{varwidth}[t]{\sphinxcolwidth{1}{3}}
\sphinxAtStartPar
MT\_CohHxRx
\sphinxbeforeendvarwidth
\end{varwidth}%
&\begin{varwidth}[t]{\sphinxcolwidth{1}{3}}
\sphinxAtStartPar
Hx\sphinxhyphen{}远参考X相干度
\sphinxbeforeendvarwidth
\end{varwidth}%
\\
\sphinxhline\begin{varwidth}[t]{\sphinxcolwidth{1}{3}}
\sphinxAtStartPar
CohHxRy
\sphinxbeforeendvarwidth
\end{varwidth}%
&\begin{varwidth}[t]{\sphinxcolwidth{1}{3}}
\sphinxAtStartPar
MT\_CohHxRy
\sphinxbeforeendvarwidth
\end{varwidth}%
&\begin{varwidth}[t]{\sphinxcolwidth{1}{3}}
\sphinxAtStartPar
Hx\sphinxhyphen{}远参考Y相干度
\sphinxbeforeendvarwidth
\end{varwidth}%
\\
\sphinxhline\begin{varwidth}[t]{\sphinxcolwidth{1}{3}}
\sphinxAtStartPar
CohHyRx
\sphinxbeforeendvarwidth
\end{varwidth}%
&\begin{varwidth}[t]{\sphinxcolwidth{1}{3}}
\sphinxAtStartPar
MT\_CohHyRx
\sphinxbeforeendvarwidth
\end{varwidth}%
&\begin{varwidth}[t]{\sphinxcolwidth{1}{3}}
\sphinxAtStartPar
Hy\sphinxhyphen{}远参考X相干度
\sphinxbeforeendvarwidth
\end{varwidth}%
\\
\sphinxhline\begin{varwidth}[t]{\sphinxcolwidth{1}{3}}
\sphinxAtStartPar
CohHyRy
\sphinxbeforeendvarwidth
\end{varwidth}%
&\begin{varwidth}[t]{\sphinxcolwidth{1}{3}}
\sphinxAtStartPar
MT\_CohHyRy
\sphinxbeforeendvarwidth
\end{varwidth}%
&\begin{varwidth}[t]{\sphinxcolwidth{1}{3}}
\sphinxAtStartPar
Hy\sphinxhyphen{}远参考Y相干度
\sphinxbeforeendvarwidth
\end{varwidth}%
\\
\sphinxbottomrule
\end{tabulary}
\sphinxtableafterendhook\par
\sphinxattableend\end{savenotes}

\sphinxAtStartPar
\sphinxstylestrong{偏相干度：}


\begin{savenotes}\sphinxattablestart
\sphinxthistablewithglobalstyle
\centering
\begin{tabulary}{\linewidth}[t]{TTTT}
\sphinxtoprule
\begin{varwidth}[t]{\sphinxcolwidth{1}{4}}
\sphinxstyletheadfamily \sphinxAtStartPar
参数名
\sphinxbeforeendvarwidth
\end{varwidth}%
&\begin{varwidth}[t]{\sphinxcolwidth{1}{4}}
\sphinxstyletheadfamily \sphinxAtStartPar
内部名称
\sphinxbeforeendvarwidth
\end{varwidth}%
&\begin{varwidth}[t]{\sphinxcolwidth{1}{4}}
\sphinxstyletheadfamily \sphinxAtStartPar
说明
\sphinxbeforeendvarwidth
\end{varwidth}%
&\begin{varwidth}[t]{\sphinxcolwidth{1}{4}}
\sphinxstyletheadfamily \sphinxAtStartPar
计算公式
\sphinxbeforeendvarwidth
\end{varwidth}%
\\
\sphinxmidrule
\sphinxtableatstartofbodyhook\begin{varwidth}[t]{\sphinxcolwidth{1}{4}}
\sphinxAtStartPar
CohpExHx
\sphinxbeforeendvarwidth
\end{varwidth}%
&\begin{varwidth}[t]{\sphinxcolwidth{1}{4}}
\sphinxAtStartPar
MT\_CohpExHx
\sphinxbeforeendvarwidth
\end{varwidth}%
&\begin{varwidth}[t]{\sphinxcolwidth{1}{4}}
\sphinxAtStartPar
Ex\sphinxhyphen{}Hx偏相干
\sphinxbeforeendvarwidth
\end{varwidth}%
&\begin{varwidth}[t]{\sphinxcolwidth{1}{4}}
\sphinxAtStartPar
(γb \sphinxhyphen{} γ1)/(1\sphinxhyphen{}γ1)
\sphinxbeforeendvarwidth
\end{varwidth}%
\\
\sphinxhline\begin{varwidth}[t]{\sphinxcolwidth{1}{4}}
\sphinxAtStartPar
CohpExHy
\sphinxbeforeendvarwidth
\end{varwidth}%
&\begin{varwidth}[t]{\sphinxcolwidth{1}{4}}
\sphinxAtStartPar
MT\_CohpExHy
\sphinxbeforeendvarwidth
\end{varwidth}%
&\begin{varwidth}[t]{\sphinxcolwidth{1}{4}}
\sphinxAtStartPar
Ex\sphinxhyphen{}Hy偏相干
\sphinxbeforeendvarwidth
\end{varwidth}%
&\begin{varwidth}[t]{\sphinxcolwidth{1}{4}}
\sphinxAtStartPar
(γb \sphinxhyphen{} γ1)/(1\sphinxhyphen{}γ1)
\sphinxbeforeendvarwidth
\end{varwidth}%
\\
\sphinxhline\begin{varwidth}[t]{\sphinxcolwidth{1}{4}}
\sphinxAtStartPar
CohpEyHx
\sphinxbeforeendvarwidth
\end{varwidth}%
&\begin{varwidth}[t]{\sphinxcolwidth{1}{4}}
\sphinxAtStartPar
MT\_CohpEyHx
\sphinxbeforeendvarwidth
\end{varwidth}%
&\begin{varwidth}[t]{\sphinxcolwidth{1}{4}}
\sphinxAtStartPar
Ey\sphinxhyphen{}Hx偏相干
\sphinxbeforeendvarwidth
\end{varwidth}%
&\begin{varwidth}[t]{\sphinxcolwidth{1}{4}}
\sphinxAtStartPar
(γb \sphinxhyphen{} γ1)/(1\sphinxhyphen{}γ1)
\sphinxbeforeendvarwidth
\end{varwidth}%
\\
\sphinxhline\begin{varwidth}[t]{\sphinxcolwidth{1}{4}}
\sphinxAtStartPar
CohpEyHy
\sphinxbeforeendvarwidth
\end{varwidth}%
&\begin{varwidth}[t]{\sphinxcolwidth{1}{4}}
\sphinxAtStartPar
MT\_CohpEyHy
\sphinxbeforeendvarwidth
\end{varwidth}%
&\begin{varwidth}[t]{\sphinxcolwidth{1}{4}}
\sphinxAtStartPar
Ey\sphinxhyphen{}Hy偏相干
\sphinxbeforeendvarwidth
\end{varwidth}%
&\begin{varwidth}[t]{\sphinxcolwidth{1}{4}}
\sphinxAtStartPar
(γb \sphinxhyphen{} γ1)/(1\sphinxhyphen{}γ1)
\sphinxbeforeendvarwidth
\end{varwidth}%
\\
\sphinxhline\begin{varwidth}[t]{\sphinxcolwidth{1}{4}}
\sphinxAtStartPar
CohpHxEx
\sphinxbeforeendvarwidth
\end{varwidth}%
&\begin{varwidth}[t]{\sphinxcolwidth{1}{4}}
\sphinxAtStartPar
MT\_CohpHxEx
\sphinxbeforeendvarwidth
\end{varwidth}%
&\begin{varwidth}[t]{\sphinxcolwidth{1}{4}}
\sphinxAtStartPar
Hx\sphinxhyphen{}Ex偏相干
\sphinxbeforeendvarwidth
\end{varwidth}%
&\begin{varwidth}[t]{\sphinxcolwidth{1}{4}}
\sphinxAtStartPar
(γb \sphinxhyphen{} γ1)/(1\sphinxhyphen{}γ1)
\sphinxbeforeendvarwidth
\end{varwidth}%
\\
\sphinxhline\begin{varwidth}[t]{\sphinxcolwidth{1}{4}}
\sphinxAtStartPar
CohpHxEy
\sphinxbeforeendvarwidth
\end{varwidth}%
&\begin{varwidth}[t]{\sphinxcolwidth{1}{4}}
\sphinxAtStartPar
MT\_CohpHxEy
\sphinxbeforeendvarwidth
\end{varwidth}%
&\begin{varwidth}[t]{\sphinxcolwidth{1}{4}}
\sphinxAtStartPar
Hx\sphinxhyphen{}Ey偏相干
\sphinxbeforeendvarwidth
\end{varwidth}%
&\begin{varwidth}[t]{\sphinxcolwidth{1}{4}}
\sphinxAtStartPar
(γb \sphinxhyphen{} γ1)/(1\sphinxhyphen{}γ1)
\sphinxbeforeendvarwidth
\end{varwidth}%
\\
\sphinxhline\begin{varwidth}[t]{\sphinxcolwidth{1}{4}}
\sphinxAtStartPar
CohpHyEx
\sphinxbeforeendvarwidth
\end{varwidth}%
&\begin{varwidth}[t]{\sphinxcolwidth{1}{4}}
\sphinxAtStartPar
MT\_CohpHyEx
\sphinxbeforeendvarwidth
\end{varwidth}%
&\begin{varwidth}[t]{\sphinxcolwidth{1}{4}}
\sphinxAtStartPar
Hy\sphinxhyphen{}Ex偏相干
\sphinxbeforeendvarwidth
\end{varwidth}%
&\begin{varwidth}[t]{\sphinxcolwidth{1}{4}}
\sphinxAtStartPar
(γb \sphinxhyphen{} γ1)/(1\sphinxhyphen{}γ1)
\sphinxbeforeendvarwidth
\end{varwidth}%
\\
\sphinxhline\begin{varwidth}[t]{\sphinxcolwidth{1}{4}}
\sphinxAtStartPar
CohpHyEy
\sphinxbeforeendvarwidth
\end{varwidth}%
&\begin{varwidth}[t]{\sphinxcolwidth{1}{4}}
\sphinxAtStartPar
MT\_CohpHyEy
\sphinxbeforeendvarwidth
\end{varwidth}%
&\begin{varwidth}[t]{\sphinxcolwidth{1}{4}}
\sphinxAtStartPar
Hy\sphinxhyphen{}Ey偏相干
\sphinxbeforeendvarwidth
\end{varwidth}%
&\begin{varwidth}[t]{\sphinxcolwidth{1}{4}}
\sphinxAtStartPar
(γb \sphinxhyphen{} γ1)/(1\sphinxhyphen{}γ1)
\sphinxbeforeendvarwidth
\end{varwidth}%
\\
\sphinxhline\begin{varwidth}[t]{\sphinxcolwidth{1}{4}}
\sphinxAtStartPar
CohpHzHx
\sphinxbeforeendvarwidth
\end{varwidth}%
&\begin{varwidth}[t]{\sphinxcolwidth{1}{4}}
\sphinxAtStartPar
MT\_CohpHzHx
\sphinxbeforeendvarwidth
\end{varwidth}%
&\begin{varwidth}[t]{\sphinxcolwidth{1}{4}}
\sphinxAtStartPar
Hz\sphinxhyphen{}Hx偏相干
\sphinxbeforeendvarwidth
\end{varwidth}%
&\begin{varwidth}[t]{\sphinxcolwidth{1}{4}}
\sphinxAtStartPar
(γb \sphinxhyphen{} γ1)/(1\sphinxhyphen{}γ1)
\sphinxbeforeendvarwidth
\end{varwidth}%
\\
\sphinxhline\begin{varwidth}[t]{\sphinxcolwidth{1}{4}}
\sphinxAtStartPar
CohpHzHy
\sphinxbeforeendvarwidth
\end{varwidth}%
&\begin{varwidth}[t]{\sphinxcolwidth{1}{4}}
\sphinxAtStartPar
MT\_CohpHzHy
\sphinxbeforeendvarwidth
\end{varwidth}%
&\begin{varwidth}[t]{\sphinxcolwidth{1}{4}}
\sphinxAtStartPar
Hz\sphinxhyphen{}Hy偏相干
\sphinxbeforeendvarwidth
\end{varwidth}%
&\begin{varwidth}[t]{\sphinxcolwidth{1}{4}}
\sphinxAtStartPar
(γb \sphinxhyphen{} γ1)/(1\sphinxhyphen{}γ1)
\sphinxbeforeendvarwidth
\end{varwidth}%
\\
\sphinxbottomrule
\end{tabulary}
\sphinxtableafterendhook\par
\sphinxattableend\end{savenotes}

\sphinxAtStartPar
\sphinxstylestrong{多重相干度：}


\begin{savenotes}\sphinxattablestart
\sphinxthistablewithglobalstyle
\centering
\begin{tabulary}{\linewidth}[t]{TTT}
\sphinxtoprule
\begin{varwidth}[t]{\sphinxcolwidth{1}{3}}
\sphinxstyletheadfamily \sphinxAtStartPar
参数名
\sphinxbeforeendvarwidth
\end{varwidth}%
&\begin{varwidth}[t]{\sphinxcolwidth{1}{3}}
\sphinxstyletheadfamily \sphinxAtStartPar
内部名称
\sphinxbeforeendvarwidth
\end{varwidth}%
&\begin{varwidth}[t]{\sphinxcolwidth{1}{3}}
\sphinxstyletheadfamily \sphinxAtStartPar
说明
\sphinxbeforeendvarwidth
\end{varwidth}%
\\
\sphinxmidrule
\sphinxtableatstartofbodyhook\begin{varwidth}[t]{\sphinxcolwidth{1}{3}}
\sphinxAtStartPar
CohbEx
\sphinxbeforeendvarwidth
\end{varwidth}%
&\begin{varwidth}[t]{\sphinxcolwidth{1}{3}}
\sphinxAtStartPar
MT\_CohbEx
\sphinxbeforeendvarwidth
\end{varwidth}%
&\begin{varwidth}[t]{\sphinxcolwidth{1}{3}}
\sphinxAtStartPar
Ex多重相干度
\sphinxbeforeendvarwidth
\end{varwidth}%
\\
\sphinxhline\begin{varwidth}[t]{\sphinxcolwidth{1}{3}}
\sphinxAtStartPar
CohbEy
\sphinxbeforeendvarwidth
\end{varwidth}%
&\begin{varwidth}[t]{\sphinxcolwidth{1}{3}}
\sphinxAtStartPar
MT\_CohbEy
\sphinxbeforeendvarwidth
\end{varwidth}%
&\begin{varwidth}[t]{\sphinxcolwidth{1}{3}}
\sphinxAtStartPar
Ey多重相干度
\sphinxbeforeendvarwidth
\end{varwidth}%
\\
\sphinxhline\begin{varwidth}[t]{\sphinxcolwidth{1}{3}}
\sphinxAtStartPar
CohbHx
\sphinxbeforeendvarwidth
\end{varwidth}%
&\begin{varwidth}[t]{\sphinxcolwidth{1}{3}}
\sphinxAtStartPar
MT\_CohbHx
\sphinxbeforeendvarwidth
\end{varwidth}%
&\begin{varwidth}[t]{\sphinxcolwidth{1}{3}}
\sphinxAtStartPar
Hx多重相干度
\sphinxbeforeendvarwidth
\end{varwidth}%
\\
\sphinxhline\begin{varwidth}[t]{\sphinxcolwidth{1}{3}}
\sphinxAtStartPar
CohbHy
\sphinxbeforeendvarwidth
\end{varwidth}%
&\begin{varwidth}[t]{\sphinxcolwidth{1}{3}}
\sphinxAtStartPar
MT\_CohbHy
\sphinxbeforeendvarwidth
\end{varwidth}%
&\begin{varwidth}[t]{\sphinxcolwidth{1}{3}}
\sphinxAtStartPar
Hy多重相干度
\sphinxbeforeendvarwidth
\end{varwidth}%
\\
\sphinxhline\begin{varwidth}[t]{\sphinxcolwidth{1}{3}}
\sphinxAtStartPar
CohbHz
\sphinxbeforeendvarwidth
\end{varwidth}%
&\begin{varwidth}[t]{\sphinxcolwidth{1}{3}}
\sphinxAtStartPar
MT\_CohbHz
\sphinxbeforeendvarwidth
\end{varwidth}%
&\begin{varwidth}[t]{\sphinxcolwidth{1}{3}}
\sphinxAtStartPar
Hz多重相干度
\sphinxbeforeendvarwidth
\end{varwidth}%
\\
\sphinxbottomrule
\end{tabulary}
\sphinxtableafterendhook\par
\sphinxattableend\end{savenotes}


\paragraph{5.17.11 信号与噪声}
\label{\detokenize{chapters/chapter6:id79}}

\begin{savenotes}\sphinxattablestart
\sphinxthistablewithglobalstyle
\centering
\begin{tabulary}{\linewidth}[t]{TTTT}
\sphinxtoprule
\begin{varwidth}[t]{\sphinxcolwidth{1}{4}}
\sphinxstyletheadfamily \sphinxAtStartPar
参数名
\sphinxbeforeendvarwidth
\end{varwidth}%
&\begin{varwidth}[t]{\sphinxcolwidth{1}{4}}
\sphinxstyletheadfamily \sphinxAtStartPar
内部名称
\sphinxbeforeendvarwidth
\end{varwidth}%
&\begin{varwidth}[t]{\sphinxcolwidth{1}{4}}
\sphinxstyletheadfamily \sphinxAtStartPar
说明
\sphinxbeforeendvarwidth
\end{varwidth}%
&\begin{varwidth}[t]{\sphinxcolwidth{1}{4}}
\sphinxstyletheadfamily \sphinxAtStartPar
单位
\sphinxbeforeendvarwidth
\end{varwidth}%
\\
\sphinxmidrule
\sphinxtableatstartofbodyhook\begin{varwidth}[t]{\sphinxcolwidth{1}{4}}
\sphinxAtStartPar
exsignal
\sphinxbeforeendvarwidth
\end{varwidth}%
&\begin{varwidth}[t]{\sphinxcolwidth{1}{4}}
\sphinxAtStartPar
MT\_exsignal
\sphinxbeforeendvarwidth
\end{varwidth}%
&\begin{varwidth}[t]{\sphinxcolwidth{1}{4}}
\sphinxAtStartPar
Ex信号功率
\sphinxbeforeendvarwidth
\end{varwidth}%
&\begin{varwidth}[t]{\sphinxcolwidth{1}{4}}
\sphinxAtStartPar
V²/m²
\sphinxbeforeendvarwidth
\end{varwidth}%
\\
\sphinxhline\begin{varwidth}[t]{\sphinxcolwidth{1}{4}}
\sphinxAtStartPar
eysignal
\sphinxbeforeendvarwidth
\end{varwidth}%
&\begin{varwidth}[t]{\sphinxcolwidth{1}{4}}
\sphinxAtStartPar
MT\_eysignal
\sphinxbeforeendvarwidth
\end{varwidth}%
&\begin{varwidth}[t]{\sphinxcolwidth{1}{4}}
\sphinxAtStartPar
Ey信号功率
\sphinxbeforeendvarwidth
\end{varwidth}%
&\begin{varwidth}[t]{\sphinxcolwidth{1}{4}}
\sphinxAtStartPar
V²/m²
\sphinxbeforeendvarwidth
\end{varwidth}%
\\
\sphinxhline\begin{varwidth}[t]{\sphinxcolwidth{1}{4}}
\sphinxAtStartPar
hxsignal
\sphinxbeforeendvarwidth
\end{varwidth}%
&\begin{varwidth}[t]{\sphinxcolwidth{1}{4}}
\sphinxAtStartPar
MT\_hxsignal
\sphinxbeforeendvarwidth
\end{varwidth}%
&\begin{varwidth}[t]{\sphinxcolwidth{1}{4}}
\sphinxAtStartPar
Hx信号功率
\sphinxbeforeendvarwidth
\end{varwidth}%
&\begin{varwidth}[t]{\sphinxcolwidth{1}{4}}
\sphinxAtStartPar
nT²
\sphinxbeforeendvarwidth
\end{varwidth}%
\\
\sphinxhline\begin{varwidth}[t]{\sphinxcolwidth{1}{4}}
\sphinxAtStartPar
hysignal
\sphinxbeforeendvarwidth
\end{varwidth}%
&\begin{varwidth}[t]{\sphinxcolwidth{1}{4}}
\sphinxAtStartPar
MT\_hysignal
\sphinxbeforeendvarwidth
\end{varwidth}%
&\begin{varwidth}[t]{\sphinxcolwidth{1}{4}}
\sphinxAtStartPar
Hy信号功率
\sphinxbeforeendvarwidth
\end{varwidth}%
&\begin{varwidth}[t]{\sphinxcolwidth{1}{4}}
\sphinxAtStartPar
nT²
\sphinxbeforeendvarwidth
\end{varwidth}%
\\
\sphinxhline\begin{varwidth}[t]{\sphinxcolwidth{1}{4}}
\sphinxAtStartPar
rxsignal
\sphinxbeforeendvarwidth
\end{varwidth}%
&\begin{varwidth}[t]{\sphinxcolwidth{1}{4}}
\sphinxAtStartPar
MT\_rxsignal
\sphinxbeforeendvarwidth
\end{varwidth}%
&\begin{varwidth}[t]{\sphinxcolwidth{1}{4}}
\sphinxAtStartPar
Rx信号功率
\sphinxbeforeendvarwidth
\end{varwidth}%
&\begin{varwidth}[t]{\sphinxcolwidth{1}{4}}
\sphinxAtStartPar
nT²
\sphinxbeforeendvarwidth
\end{varwidth}%
\\
\sphinxhline\begin{varwidth}[t]{\sphinxcolwidth{1}{4}}
\sphinxAtStartPar
rysignal
\sphinxbeforeendvarwidth
\end{varwidth}%
&\begin{varwidth}[t]{\sphinxcolwidth{1}{4}}
\sphinxAtStartPar
MT\_rysignal
\sphinxbeforeendvarwidth
\end{varwidth}%
&\begin{varwidth}[t]{\sphinxcolwidth{1}{4}}
\sphinxAtStartPar
Ry信号功率
\sphinxbeforeendvarwidth
\end{varwidth}%
&\begin{varwidth}[t]{\sphinxcolwidth{1}{4}}
\sphinxAtStartPar
nT²
\sphinxbeforeendvarwidth
\end{varwidth}%
\\
\sphinxhline\begin{varwidth}[t]{\sphinxcolwidth{1}{4}}
\sphinxAtStartPar
exnoise
\sphinxbeforeendvarwidth
\end{varwidth}%
&\begin{varwidth}[t]{\sphinxcolwidth{1}{4}}
\sphinxAtStartPar
MT\_exnoise
\sphinxbeforeendvarwidth
\end{varwidth}%
&\begin{varwidth}[t]{\sphinxcolwidth{1}{4}}
\sphinxAtStartPar
Ex噪声功率
\sphinxbeforeendvarwidth
\end{varwidth}%
&\begin{varwidth}[t]{\sphinxcolwidth{1}{4}}
\sphinxAtStartPar
V²/m²
\sphinxbeforeendvarwidth
\end{varwidth}%
\\
\sphinxhline\begin{varwidth}[t]{\sphinxcolwidth{1}{4}}
\sphinxAtStartPar
eynoise
\sphinxbeforeendvarwidth
\end{varwidth}%
&\begin{varwidth}[t]{\sphinxcolwidth{1}{4}}
\sphinxAtStartPar
MT\_eynoise
\sphinxbeforeendvarwidth
\end{varwidth}%
&\begin{varwidth}[t]{\sphinxcolwidth{1}{4}}
\sphinxAtStartPar
Ey噪声功率
\sphinxbeforeendvarwidth
\end{varwidth}%
&\begin{varwidth}[t]{\sphinxcolwidth{1}{4}}
\sphinxAtStartPar
V²/m²
\sphinxbeforeendvarwidth
\end{varwidth}%
\\
\sphinxhline\begin{varwidth}[t]{\sphinxcolwidth{1}{4}}
\sphinxAtStartPar
hxnoise
\sphinxbeforeendvarwidth
\end{varwidth}%
&\begin{varwidth}[t]{\sphinxcolwidth{1}{4}}
\sphinxAtStartPar
MT\_hxnoise
\sphinxbeforeendvarwidth
\end{varwidth}%
&\begin{varwidth}[t]{\sphinxcolwidth{1}{4}}
\sphinxAtStartPar
Hx噪声功率
\sphinxbeforeendvarwidth
\end{varwidth}%
&\begin{varwidth}[t]{\sphinxcolwidth{1}{4}}
\sphinxAtStartPar
nT²
\sphinxbeforeendvarwidth
\end{varwidth}%
\\
\sphinxhline\begin{varwidth}[t]{\sphinxcolwidth{1}{4}}
\sphinxAtStartPar
hynoise
\sphinxbeforeendvarwidth
\end{varwidth}%
&\begin{varwidth}[t]{\sphinxcolwidth{1}{4}}
\sphinxAtStartPar
MT\_hynoise
\sphinxbeforeendvarwidth
\end{varwidth}%
&\begin{varwidth}[t]{\sphinxcolwidth{1}{4}}
\sphinxAtStartPar
Hy噪声功率
\sphinxbeforeendvarwidth
\end{varwidth}%
&\begin{varwidth}[t]{\sphinxcolwidth{1}{4}}
\sphinxAtStartPar
nT²
\sphinxbeforeendvarwidth
\end{varwidth}%
\\
\sphinxhline\begin{varwidth}[t]{\sphinxcolwidth{1}{4}}
\sphinxAtStartPar
rxnoise
\sphinxbeforeendvarwidth
\end{varwidth}%
&\begin{varwidth}[t]{\sphinxcolwidth{1}{4}}
\sphinxAtStartPar
MT\_rxnoise
\sphinxbeforeendvarwidth
\end{varwidth}%
&\begin{varwidth}[t]{\sphinxcolwidth{1}{4}}
\sphinxAtStartPar
Rx噪声功率
\sphinxbeforeendvarwidth
\end{varwidth}%
&\begin{varwidth}[t]{\sphinxcolwidth{1}{4}}
\sphinxAtStartPar
nT²
\sphinxbeforeendvarwidth
\end{varwidth}%
\\
\sphinxhline\begin{varwidth}[t]{\sphinxcolwidth{1}{4}}
\sphinxAtStartPar
rynoise
\sphinxbeforeendvarwidth
\end{varwidth}%
&\begin{varwidth}[t]{\sphinxcolwidth{1}{4}}
\sphinxAtStartPar
MT\_rynoise
\sphinxbeforeendvarwidth
\end{varwidth}%
&\begin{varwidth}[t]{\sphinxcolwidth{1}{4}}
\sphinxAtStartPar
Ry噪声功率
\sphinxbeforeendvarwidth
\end{varwidth}%
&\begin{varwidth}[t]{\sphinxcolwidth{1}{4}}
\sphinxAtStartPar
nT²
\sphinxbeforeendvarwidth
\end{varwidth}%
\\
\sphinxhline\begin{varwidth}[t]{\sphinxcolwidth{1}{4}}
\sphinxAtStartPar
\sphinxstylestrong{信噪比（SNR）：}
\sphinxbeforeendvarwidth
\end{varwidth}%
&\begin{varwidth}[t]{\sphinxcolwidth{1}{4}}
\sphinxAtStartPar

\sphinxbeforeendvarwidth
\end{varwidth}%
&\begin{varwidth}[t]{\sphinxcolwidth{1}{4}}
\sphinxAtStartPar

\sphinxbeforeendvarwidth
\end{varwidth}%
&\begin{varwidth}[t]{\sphinxcolwidth{1}{4}}
\sphinxAtStartPar

\sphinxbeforeendvarwidth
\end{varwidth}%
\\
\sphinxbottomrule
\end{tabulary}
\sphinxtableafterendhook\par
\sphinxattableend\end{savenotes}


\begin{savenotes}\sphinxattablestart
\sphinxthistablewithglobalstyle
\centering
\begin{tabulary}{\linewidth}[t]{TTTT}
\sphinxtoprule
\begin{varwidth}[t]{\sphinxcolwidth{1}{4}}
\sphinxstyletheadfamily \sphinxAtStartPar
参数名
\sphinxbeforeendvarwidth
\end{varwidth}%
&\begin{varwidth}[t]{\sphinxcolwidth{1}{4}}
\sphinxstyletheadfamily \sphinxAtStartPar
内部名称
\sphinxbeforeendvarwidth
\end{varwidth}%
&\begin{varwidth}[t]{\sphinxcolwidth{1}{4}}
\sphinxstyletheadfamily \sphinxAtStartPar
说明
\sphinxbeforeendvarwidth
\end{varwidth}%
&\begin{varwidth}[t]{\sphinxcolwidth{1}{4}}
\sphinxstyletheadfamily \sphinxAtStartPar
计算公式
\sphinxbeforeendvarwidth
\end{varwidth}%
\\
\sphinxmidrule
\sphinxtableatstartofbodyhook\begin{varwidth}[t]{\sphinxcolwidth{1}{4}}
\sphinxAtStartPar
exsnr
\sphinxbeforeendvarwidth
\end{varwidth}%
&\begin{varwidth}[t]{\sphinxcolwidth{1}{4}}
\sphinxAtStartPar
MT\_exsnr
\sphinxbeforeendvarwidth
\end{varwidth}%
&\begin{varwidth}[t]{\sphinxcolwidth{1}{4}}
\sphinxAtStartPar
Ex信噪比
\sphinxbeforeendvarwidth
\end{varwidth}%
&\begin{varwidth}[t]{\sphinxcolwidth{1}{4}}
\sphinxAtStartPar
signal/noise
\sphinxbeforeendvarwidth
\end{varwidth}%
\\
\sphinxhline\begin{varwidth}[t]{\sphinxcolwidth{1}{4}}
\sphinxAtStartPar
eysnr
\sphinxbeforeendvarwidth
\end{varwidth}%
&\begin{varwidth}[t]{\sphinxcolwidth{1}{4}}
\sphinxAtStartPar
MT\_eysnr
\sphinxbeforeendvarwidth
\end{varwidth}%
&\begin{varwidth}[t]{\sphinxcolwidth{1}{4}}
\sphinxAtStartPar
Ey信噪比
\sphinxbeforeendvarwidth
\end{varwidth}%
&\begin{varwidth}[t]{\sphinxcolwidth{1}{4}}
\sphinxAtStartPar
signal/noise
\sphinxbeforeendvarwidth
\end{varwidth}%
\\
\sphinxhline\begin{varwidth}[t]{\sphinxcolwidth{1}{4}}
\sphinxAtStartPar
hxsnr
\sphinxbeforeendvarwidth
\end{varwidth}%
&\begin{varwidth}[t]{\sphinxcolwidth{1}{4}}
\sphinxAtStartPar
MT\_hxsnr
\sphinxbeforeendvarwidth
\end{varwidth}%
&\begin{varwidth}[t]{\sphinxcolwidth{1}{4}}
\sphinxAtStartPar
Hx信噪比
\sphinxbeforeendvarwidth
\end{varwidth}%
&\begin{varwidth}[t]{\sphinxcolwidth{1}{4}}
\sphinxAtStartPar
signal/noise
\sphinxbeforeendvarwidth
\end{varwidth}%
\\
\sphinxhline\begin{varwidth}[t]{\sphinxcolwidth{1}{4}}
\sphinxAtStartPar
hysnr
\sphinxbeforeendvarwidth
\end{varwidth}%
&\begin{varwidth}[t]{\sphinxcolwidth{1}{4}}
\sphinxAtStartPar
MT\_hysnr
\sphinxbeforeendvarwidth
\end{varwidth}%
&\begin{varwidth}[t]{\sphinxcolwidth{1}{4}}
\sphinxAtStartPar
Hy信噪比
\sphinxbeforeendvarwidth
\end{varwidth}%
&\begin{varwidth}[t]{\sphinxcolwidth{1}{4}}
\sphinxAtStartPar
signal/noise
\sphinxbeforeendvarwidth
\end{varwidth}%
\\
\sphinxhline\begin{varwidth}[t]{\sphinxcolwidth{1}{4}}
\sphinxAtStartPar
rxsnr
\sphinxbeforeendvarwidth
\end{varwidth}%
&\begin{varwidth}[t]{\sphinxcolwidth{1}{4}}
\sphinxAtStartPar
MT\_rxsnr
\sphinxbeforeendvarwidth
\end{varwidth}%
&\begin{varwidth}[t]{\sphinxcolwidth{1}{4}}
\sphinxAtStartPar
Rx信噪比
\sphinxbeforeendvarwidth
\end{varwidth}%
&\begin{varwidth}[t]{\sphinxcolwidth{1}{4}}
\sphinxAtStartPar
signal/noise
\sphinxbeforeendvarwidth
\end{varwidth}%
\\
\sphinxhline\begin{varwidth}[t]{\sphinxcolwidth{1}{4}}
\sphinxAtStartPar
rysnr
\sphinxbeforeendvarwidth
\end{varwidth}%
&\begin{varwidth}[t]{\sphinxcolwidth{1}{4}}
\sphinxAtStartPar
MT\_rysnr
\sphinxbeforeendvarwidth
\end{varwidth}%
&\begin{varwidth}[t]{\sphinxcolwidth{1}{4}}
\sphinxAtStartPar
Ry信噪比
\sphinxbeforeendvarwidth
\end{varwidth}%
&\begin{varwidth}[t]{\sphinxcolwidth{1}{4}}
\sphinxAtStartPar
signal/noise
\sphinxbeforeendvarwidth
\end{varwidth}%
\\
\sphinxbottomrule
\end{tabulary}
\sphinxtableafterendhook\par
\sphinxattableend\end{savenotes}


\paragraph{5.17.12 功率谱密度}
\label{\detokenize{chapters/chapter6:id80}}

\begin{savenotes}\sphinxattablestart
\sphinxthistablewithglobalstyle
\centering
\begin{tabulary}{\linewidth}[t]{TTTT}
\sphinxtoprule
\begin{varwidth}[t]{\sphinxcolwidth{1}{4}}
\sphinxstyletheadfamily \sphinxAtStartPar
参数名
\sphinxbeforeendvarwidth
\end{varwidth}%
&\begin{varwidth}[t]{\sphinxcolwidth{1}{4}}
\sphinxstyletheadfamily \sphinxAtStartPar
内部名称
\sphinxbeforeendvarwidth
\end{varwidth}%
&\begin{varwidth}[t]{\sphinxcolwidth{1}{4}}
\sphinxstyletheadfamily \sphinxAtStartPar
说明
\sphinxbeforeendvarwidth
\end{varwidth}%
&\begin{varwidth}[t]{\sphinxcolwidth{1}{4}}
\sphinxstyletheadfamily \sphinxAtStartPar
计算公式
\sphinxbeforeendvarwidth
\end{varwidth}%
\\
\sphinxmidrule
\sphinxtableatstartofbodyhook\begin{varwidth}[t]{\sphinxcolwidth{1}{4}}
\sphinxAtStartPar
ExPSD
\sphinxbeforeendvarwidth
\end{varwidth}%
&\begin{varwidth}[t]{\sphinxcolwidth{1}{4}}
\sphinxAtStartPar
MT\_ExPSD
\sphinxbeforeendvarwidth
\end{varwidth}%
&\begin{varwidth}[t]{\sphinxcolwidth{1}{4}}
\sphinxAtStartPar
Ex功率谱密度
\sphinxbeforeendvarwidth
\end{varwidth}%
&\begin{varwidth}[t]{\sphinxcolwidth{1}{4}}
\sphinxAtStartPar
<Ex·Ex*>
\sphinxbeforeendvarwidth
\end{varwidth}%
\\
\sphinxhline\begin{varwidth}[t]{\sphinxcolwidth{1}{4}}
\sphinxAtStartPar
EyPSD
\sphinxbeforeendvarwidth
\end{varwidth}%
&\begin{varwidth}[t]{\sphinxcolwidth{1}{4}}
\sphinxAtStartPar
MT\_EyPSD
\sphinxbeforeendvarwidth
\end{varwidth}%
&\begin{varwidth}[t]{\sphinxcolwidth{1}{4}}
\sphinxAtStartPar
Ey功率谱密度
\sphinxbeforeendvarwidth
\end{varwidth}%
&\begin{varwidth}[t]{\sphinxcolwidth{1}{4}}
\sphinxAtStartPar
<Ey·Ey*>
\sphinxbeforeendvarwidth
\end{varwidth}%
\\
\sphinxhline\begin{varwidth}[t]{\sphinxcolwidth{1}{4}}
\sphinxAtStartPar
HxPSD
\sphinxbeforeendvarwidth
\end{varwidth}%
&\begin{varwidth}[t]{\sphinxcolwidth{1}{4}}
\sphinxAtStartPar
MT\_HxPSD
\sphinxbeforeendvarwidth
\end{varwidth}%
&\begin{varwidth}[t]{\sphinxcolwidth{1}{4}}
\sphinxAtStartPar
Hx功率谱密度
\sphinxbeforeendvarwidth
\end{varwidth}%
&\begin{varwidth}[t]{\sphinxcolwidth{1}{4}}
\sphinxAtStartPar
<Hx·Hx*>
\sphinxbeforeendvarwidth
\end{varwidth}%
\\
\sphinxhline\begin{varwidth}[t]{\sphinxcolwidth{1}{4}}
\sphinxAtStartPar
HyPSD
\sphinxbeforeendvarwidth
\end{varwidth}%
&\begin{varwidth}[t]{\sphinxcolwidth{1}{4}}
\sphinxAtStartPar
MT\_HyPSD
\sphinxbeforeendvarwidth
\end{varwidth}%
&\begin{varwidth}[t]{\sphinxcolwidth{1}{4}}
\sphinxAtStartPar
Hy功率谱密度
\sphinxbeforeendvarwidth
\end{varwidth}%
&\begin{varwidth}[t]{\sphinxcolwidth{1}{4}}
\sphinxAtStartPar
<Hy·Hy*>
\sphinxbeforeendvarwidth
\end{varwidth}%
\\
\sphinxhline\begin{varwidth}[t]{\sphinxcolwidth{1}{4}}
\sphinxAtStartPar
HzPSD
\sphinxbeforeendvarwidth
\end{varwidth}%
&\begin{varwidth}[t]{\sphinxcolwidth{1}{4}}
\sphinxAtStartPar
MT\_HzPSD
\sphinxbeforeendvarwidth
\end{varwidth}%
&\begin{varwidth}[t]{\sphinxcolwidth{1}{4}}
\sphinxAtStartPar
Hz功率谱密度
\sphinxbeforeendvarwidth
\end{varwidth}%
&\begin{varwidth}[t]{\sphinxcolwidth{1}{4}}
\sphinxAtStartPar
<Hz·Hz*>
\sphinxbeforeendvarwidth
\end{varwidth}%
\\
\sphinxhline\begin{varwidth}[t]{\sphinxcolwidth{1}{4}}
\sphinxAtStartPar
lgExPSD
\sphinxbeforeendvarwidth
\end{varwidth}%
&\begin{varwidth}[t]{\sphinxcolwidth{1}{4}}
\sphinxAtStartPar
MT\_lgExPSD
\sphinxbeforeendvarwidth
\end{varwidth}%
&\begin{varwidth}[t]{\sphinxcolwidth{1}{4}}
\sphinxAtStartPar
log(Ex PSD)
\sphinxbeforeendvarwidth
\end{varwidth}%
&\begin{varwidth}[t]{\sphinxcolwidth{1}{4}}
\sphinxAtStartPar
log₁₀(ExPSD)
\sphinxbeforeendvarwidth
\end{varwidth}%
\\
\sphinxhline\begin{varwidth}[t]{\sphinxcolwidth{1}{4}}
\sphinxAtStartPar
lgEyPSD
\sphinxbeforeendvarwidth
\end{varwidth}%
&\begin{varwidth}[t]{\sphinxcolwidth{1}{4}}
\sphinxAtStartPar
MT\_lgEyPSD
\sphinxbeforeendvarwidth
\end{varwidth}%
&\begin{varwidth}[t]{\sphinxcolwidth{1}{4}}
\sphinxAtStartPar
log(Ey PSD)
\sphinxbeforeendvarwidth
\end{varwidth}%
&\begin{varwidth}[t]{\sphinxcolwidth{1}{4}}
\sphinxAtStartPar
log₁₀(EyPSD)
\sphinxbeforeendvarwidth
\end{varwidth}%
\\
\sphinxhline\begin{varwidth}[t]{\sphinxcolwidth{1}{4}}
\sphinxAtStartPar
lgHxPSD
\sphinxbeforeendvarwidth
\end{varwidth}%
&\begin{varwidth}[t]{\sphinxcolwidth{1}{4}}
\sphinxAtStartPar
MT\_lgHxPSD
\sphinxbeforeendvarwidth
\end{varwidth}%
&\begin{varwidth}[t]{\sphinxcolwidth{1}{4}}
\sphinxAtStartPar
log(Hx PSD)
\sphinxbeforeendvarwidth
\end{varwidth}%
&\begin{varwidth}[t]{\sphinxcolwidth{1}{4}}
\sphinxAtStartPar
log₁₀(HxPSD)
\sphinxbeforeendvarwidth
\end{varwidth}%
\\
\sphinxhline\begin{varwidth}[t]{\sphinxcolwidth{1}{4}}
\sphinxAtStartPar
lgHyPSD
\sphinxbeforeendvarwidth
\end{varwidth}%
&\begin{varwidth}[t]{\sphinxcolwidth{1}{4}}
\sphinxAtStartPar
MT\_lgHyPSD
\sphinxbeforeendvarwidth
\end{varwidth}%
&\begin{varwidth}[t]{\sphinxcolwidth{1}{4}}
\sphinxAtStartPar
log(Hy PSD)
\sphinxbeforeendvarwidth
\end{varwidth}%
&\begin{varwidth}[t]{\sphinxcolwidth{1}{4}}
\sphinxAtStartPar
log₁₀(HyPSD)
\sphinxbeforeendvarwidth
\end{varwidth}%
\\
\sphinxhline\begin{varwidth}[t]{\sphinxcolwidth{1}{4}}
\sphinxAtStartPar
lgHzPSD
\sphinxbeforeendvarwidth
\end{varwidth}%
&\begin{varwidth}[t]{\sphinxcolwidth{1}{4}}
\sphinxAtStartPar
MT\_lgHzPSD
\sphinxbeforeendvarwidth
\end{varwidth}%
&\begin{varwidth}[t]{\sphinxcolwidth{1}{4}}
\sphinxAtStartPar
log(Hz PSD)
\sphinxbeforeendvarwidth
\end{varwidth}%
&\begin{varwidth}[t]{\sphinxcolwidth{1}{4}}
\sphinxAtStartPar
log₁₀(HzPSD)
\sphinxbeforeendvarwidth
\end{varwidth}%
\\
\sphinxbottomrule
\end{tabulary}
\sphinxtableafterendhook\par
\sphinxattableend\end{savenotes}


\paragraph{5.17.13 频谱幅值}
\label{\detokenize{chapters/chapter6:id81}}

\begin{savenotes}\sphinxattablestart
\sphinxthistablewithglobalstyle
\centering
\begin{tabulary}{\linewidth}[t]{TTTT}
\sphinxtoprule
\begin{varwidth}[t]{\sphinxcolwidth{1}{4}}
\sphinxstyletheadfamily \sphinxAtStartPar
参数名
\sphinxbeforeendvarwidth
\end{varwidth}%
&\begin{varwidth}[t]{\sphinxcolwidth{1}{4}}
\sphinxstyletheadfamily \sphinxAtStartPar
内部名称
\sphinxbeforeendvarwidth
\end{varwidth}%
&\begin{varwidth}[t]{\sphinxcolwidth{1}{4}}
\sphinxstyletheadfamily \sphinxAtStartPar
说明
\sphinxbeforeendvarwidth
\end{varwidth}%
&\begin{varwidth}[t]{\sphinxcolwidth{1}{4}}
\sphinxstyletheadfamily \sphinxAtStartPar
计算公式
\sphinxbeforeendvarwidth
\end{varwidth}%
\\
\sphinxmidrule
\sphinxtableatstartofbodyhook\begin{varwidth}[t]{\sphinxcolwidth{1}{4}}
\sphinxAtStartPar
ExSpectra
\sphinxbeforeendvarwidth
\end{varwidth}%
&\begin{varwidth}[t]{\sphinxcolwidth{1}{4}}
\sphinxAtStartPar
MT\_ExSpectra
\sphinxbeforeendvarwidth
\end{varwidth}%
&\begin{varwidth}[t]{\sphinxcolwidth{1}{4}}
\sphinxAtStartPar
Ex频谱幅值
\sphinxbeforeendvarwidth
\end{varwidth}%
&\begin{varwidth}[t]{\sphinxcolwidth{1}{4}}
\sphinxAtStartPar
√(ExPSD)
\sphinxbeforeendvarwidth
\end{varwidth}%
\\
\sphinxhline\begin{varwidth}[t]{\sphinxcolwidth{1}{4}}
\sphinxAtStartPar
EySpectra
\sphinxbeforeendvarwidth
\end{varwidth}%
&\begin{varwidth}[t]{\sphinxcolwidth{1}{4}}
\sphinxAtStartPar
MT\_EySpectra
\sphinxbeforeendvarwidth
\end{varwidth}%
&\begin{varwidth}[t]{\sphinxcolwidth{1}{4}}
\sphinxAtStartPar
Ey频谱幅值
\sphinxbeforeendvarwidth
\end{varwidth}%
&\begin{varwidth}[t]{\sphinxcolwidth{1}{4}}
\sphinxAtStartPar
√(EyPSD)
\sphinxbeforeendvarwidth
\end{varwidth}%
\\
\sphinxhline\begin{varwidth}[t]{\sphinxcolwidth{1}{4}}
\sphinxAtStartPar
HxSpectra
\sphinxbeforeendvarwidth
\end{varwidth}%
&\begin{varwidth}[t]{\sphinxcolwidth{1}{4}}
\sphinxAtStartPar
MT\_HxSpectra
\sphinxbeforeendvarwidth
\end{varwidth}%
&\begin{varwidth}[t]{\sphinxcolwidth{1}{4}}
\sphinxAtStartPar
Hx频谱幅值
\sphinxbeforeendvarwidth
\end{varwidth}%
&\begin{varwidth}[t]{\sphinxcolwidth{1}{4}}
\sphinxAtStartPar
√(HxPSD)
\sphinxbeforeendvarwidth
\end{varwidth}%
\\
\sphinxhline\begin{varwidth}[t]{\sphinxcolwidth{1}{4}}
\sphinxAtStartPar
HySpectra
\sphinxbeforeendvarwidth
\end{varwidth}%
&\begin{varwidth}[t]{\sphinxcolwidth{1}{4}}
\sphinxAtStartPar
MT\_HySpectra
\sphinxbeforeendvarwidth
\end{varwidth}%
&\begin{varwidth}[t]{\sphinxcolwidth{1}{4}}
\sphinxAtStartPar
Hy频谱幅值
\sphinxbeforeendvarwidth
\end{varwidth}%
&\begin{varwidth}[t]{\sphinxcolwidth{1}{4}}
\sphinxAtStartPar
√(HyPSD)
\sphinxbeforeendvarwidth
\end{varwidth}%
\\
\sphinxhline\begin{varwidth}[t]{\sphinxcolwidth{1}{4}}
\sphinxAtStartPar
HzSpectra
\sphinxbeforeendvarwidth
\end{varwidth}%
&\begin{varwidth}[t]{\sphinxcolwidth{1}{4}}
\sphinxAtStartPar
MT\_HzSpectra
\sphinxbeforeendvarwidth
\end{varwidth}%
&\begin{varwidth}[t]{\sphinxcolwidth{1}{4}}
\sphinxAtStartPar
Hz频谱幅值
\sphinxbeforeendvarwidth
\end{varwidth}%
&\begin{varwidth}[t]{\sphinxcolwidth{1}{4}}
\sphinxAtStartPar
√(HzPSD)
\sphinxbeforeendvarwidth
\end{varwidth}%
\\
\sphinxhline\begin{varwidth}[t]{\sphinxcolwidth{1}{4}}
\sphinxAtStartPar
lgExSpectra
\sphinxbeforeendvarwidth
\end{varwidth}%
&\begin{varwidth}[t]{\sphinxcolwidth{1}{4}}
\sphinxAtStartPar
MT\_lgExSpectra
\sphinxbeforeendvarwidth
\end{varwidth}%
&\begin{varwidth}[t]{\sphinxcolwidth{1}{4}}
\sphinxAtStartPar
log(Ex Spectra)
\sphinxbeforeendvarwidth
\end{varwidth}%
&\begin{varwidth}[t]{\sphinxcolwidth{1}{4}}
\sphinxAtStartPar
log₁₀(ExSpectra)
\sphinxbeforeendvarwidth
\end{varwidth}%
\\
\sphinxhline\begin{varwidth}[t]{\sphinxcolwidth{1}{4}}
\sphinxAtStartPar
lgEySpectra
\sphinxbeforeendvarwidth
\end{varwidth}%
&\begin{varwidth}[t]{\sphinxcolwidth{1}{4}}
\sphinxAtStartPar
MT\_lgEySpectra
\sphinxbeforeendvarwidth
\end{varwidth}%
&\begin{varwidth}[t]{\sphinxcolwidth{1}{4}}
\sphinxAtStartPar
log(Ey Spectra)
\sphinxbeforeendvarwidth
\end{varwidth}%
&\begin{varwidth}[t]{\sphinxcolwidth{1}{4}}
\sphinxAtStartPar
log₁₀(EySpectra)
\sphinxbeforeendvarwidth
\end{varwidth}%
\\
\sphinxhline\begin{varwidth}[t]{\sphinxcolwidth{1}{4}}
\sphinxAtStartPar
lgHxSpectra
\sphinxbeforeendvarwidth
\end{varwidth}%
&\begin{varwidth}[t]{\sphinxcolwidth{1}{4}}
\sphinxAtStartPar
MT\_lgHxSpectra
\sphinxbeforeendvarwidth
\end{varwidth}%
&\begin{varwidth}[t]{\sphinxcolwidth{1}{4}}
\sphinxAtStartPar
log(Hx Spectra)
\sphinxbeforeendvarwidth
\end{varwidth}%
&\begin{varwidth}[t]{\sphinxcolwidth{1}{4}}
\sphinxAtStartPar
log₁₀(HxSpectra)
\sphinxbeforeendvarwidth
\end{varwidth}%
\\
\sphinxhline\begin{varwidth}[t]{\sphinxcolwidth{1}{4}}
\sphinxAtStartPar
lgHySpectra
\sphinxbeforeendvarwidth
\end{varwidth}%
&\begin{varwidth}[t]{\sphinxcolwidth{1}{4}}
\sphinxAtStartPar
MT\_lgHySpectra
\sphinxbeforeendvarwidth
\end{varwidth}%
&\begin{varwidth}[t]{\sphinxcolwidth{1}{4}}
\sphinxAtStartPar
log(Hy Spectra)
\sphinxbeforeendvarwidth
\end{varwidth}%
&\begin{varwidth}[t]{\sphinxcolwidth{1}{4}}
\sphinxAtStartPar
log₁₀(HySpectra)
\sphinxbeforeendvarwidth
\end{varwidth}%
\\
\sphinxhline\begin{varwidth}[t]{\sphinxcolwidth{1}{4}}
\sphinxAtStartPar
lgHzSpectra
\sphinxbeforeendvarwidth
\end{varwidth}%
&\begin{varwidth}[t]{\sphinxcolwidth{1}{4}}
\sphinxAtStartPar
MT\_lgHzSpectra
\sphinxbeforeendvarwidth
\end{varwidth}%
&\begin{varwidth}[t]{\sphinxcolwidth{1}{4}}
\sphinxAtStartPar
log(Hz Spectra)
\sphinxbeforeendvarwidth
\end{varwidth}%
&\begin{varwidth}[t]{\sphinxcolwidth{1}{4}}
\sphinxAtStartPar
log₁₀(HzSpectra)
\sphinxbeforeendvarwidth
\end{varwidth}%
\\
\sphinxbottomrule
\end{tabulary}
\sphinxtableafterendhook\par
\sphinxattableend\end{savenotes}


\paragraph{5.17.14 极化参数}
\label{\detokenize{chapters/chapter6:id82}}

\begin{savenotes}\sphinxattablestart
\sphinxthistablewithglobalstyle
\centering
\begin{tabulary}{\linewidth}[t]{TTTT}
\sphinxtoprule
\begin{varwidth}[t]{\sphinxcolwidth{1}{4}}
\sphinxstyletheadfamily \sphinxAtStartPar
参数名
\sphinxbeforeendvarwidth
\end{varwidth}%
&\begin{varwidth}[t]{\sphinxcolwidth{1}{4}}
\sphinxstyletheadfamily \sphinxAtStartPar
内部名称
\sphinxbeforeendvarwidth
\end{varwidth}%
&\begin{varwidth}[t]{\sphinxcolwidth{1}{4}}
\sphinxstyletheadfamily \sphinxAtStartPar
说明
\sphinxbeforeendvarwidth
\end{varwidth}%
&\begin{varwidth}[t]{\sphinxcolwidth{1}{4}}
\sphinxstyletheadfamily \sphinxAtStartPar
计算公式
\sphinxbeforeendvarwidth
\end{varwidth}%
\\
\sphinxmidrule
\sphinxtableatstartofbodyhook\begin{varwidth}[t]{\sphinxcolwidth{1}{4}}
\sphinxAtStartPar
EPolar
\sphinxbeforeendvarwidth
\end{varwidth}%
&\begin{varwidth}[t]{\sphinxcolwidth{1}{4}}
\sphinxAtStartPar
MT\_EPolar
\sphinxbeforeendvarwidth
\end{varwidth}%
&\begin{varwidth}[t]{\sphinxcolwidth{1}{4}}
\sphinxAtStartPar
电场极化角
\sphinxbeforeendvarwidth
\end{varwidth}%
&\begin{varwidth}[t]{\sphinxcolwidth{1}{4}}
\sphinxAtStartPar
arctan(2×Exy/(Exx\sphinxhyphen{}Eyy))×180/π
\sphinxbeforeendvarwidth
\end{varwidth}%
\\
\sphinxhline\begin{varwidth}[t]{\sphinxcolwidth{1}{4}}
\sphinxAtStartPar
HPolar
\sphinxbeforeendvarwidth
\end{varwidth}%
&\begin{varwidth}[t]{\sphinxcolwidth{1}{4}}
\sphinxAtStartPar
MT\_HPolar
\sphinxbeforeendvarwidth
\end{varwidth}%
&\begin{varwidth}[t]{\sphinxcolwidth{1}{4}}
\sphinxAtStartPar
磁场极化角
\sphinxbeforeendvarwidth
\end{varwidth}%
&\begin{varwidth}[t]{\sphinxcolwidth{1}{4}}
\sphinxAtStartPar
arctan(2×Hxy/(Hxx\sphinxhyphen{}Hyy))×180/π
\sphinxbeforeendvarwidth
\end{varwidth}%
\\
\sphinxhline\begin{varwidth}[t]{\sphinxcolwidth{1}{4}}
\sphinxAtStartPar
HRPolar
\sphinxbeforeendvarwidth
\end{varwidth}%
&\begin{varwidth}[t]{\sphinxcolwidth{1}{4}}
\sphinxAtStartPar
MT\_HRPolar
\sphinxbeforeendvarwidth
\end{varwidth}%
&\begin{varwidth}[t]{\sphinxcolwidth{1}{4}}
\sphinxAtStartPar
远参考极化差
\sphinxbeforeendvarwidth
\end{varwidth}%
&\begin{varwidth}[t]{\sphinxcolwidth{1}{4}}
\sphinxAtStartPar
HPolar \sphinxhyphen{} RRPolar
\sphinxbeforeendvarwidth
\end{varwidth}%
\\
\sphinxbottomrule
\end{tabulary}
\sphinxtableafterendhook\par
\sphinxattableend\end{savenotes}


\paragraph{5.17.15 阻抗行列式}
\label{\detokenize{chapters/chapter6:id83}}

\begin{savenotes}\sphinxattablestart
\sphinxthistablewithglobalstyle
\centering
\begin{tabulary}{\linewidth}[t]{TTTT}
\sphinxtoprule
\begin{varwidth}[t]{\sphinxcolwidth{1}{4}}
\sphinxstyletheadfamily \sphinxAtStartPar
参数名
\sphinxbeforeendvarwidth
\end{varwidth}%
&\begin{varwidth}[t]{\sphinxcolwidth{1}{4}}
\sphinxstyletheadfamily \sphinxAtStartPar
内部名称
\sphinxbeforeendvarwidth
\end{varwidth}%
&\begin{varwidth}[t]{\sphinxcolwidth{1}{4}}
\sphinxstyletheadfamily \sphinxAtStartPar
说明
\sphinxbeforeendvarwidth
\end{varwidth}%
&\begin{varwidth}[t]{\sphinxcolwidth{1}{4}}
\sphinxstyletheadfamily \sphinxAtStartPar
计算公式
\sphinxbeforeendvarwidth
\end{varwidth}%
\\
\sphinxmidrule
\sphinxtableatstartofbodyhook\begin{varwidth}[t]{\sphinxcolwidth{1}{4}}
\sphinxAtStartPar
DetZr
\sphinxbeforeendvarwidth
\end{varwidth}%
&\begin{varwidth}[t]{\sphinxcolwidth{1}{4}}
\sphinxAtStartPar
MT\_DetZr
\sphinxbeforeendvarwidth
\end{varwidth}%
&\begin{varwidth}[t]{\sphinxcolwidth{1}{4}}
\sphinxAtStartPar
det(Z)实部
\sphinxbeforeendvarwidth
\end{varwidth}%
&\begin{varwidth}[t]{\sphinxcolwidth{1}{4}}
\sphinxAtStartPar
Re(Zxx·Zyy \sphinxhyphen{} Zxy·Zyx)
\sphinxbeforeendvarwidth
\end{varwidth}%
\\
\sphinxhline\begin{varwidth}[t]{\sphinxcolwidth{1}{4}}
\sphinxAtStartPar
DetZi
\sphinxbeforeendvarwidth
\end{varwidth}%
&\begin{varwidth}[t]{\sphinxcolwidth{1}{4}}
\sphinxAtStartPar
MT\_DetZi
\sphinxbeforeendvarwidth
\end{varwidth}%
&\begin{varwidth}[t]{\sphinxcolwidth{1}{4}}
\sphinxAtStartPar
det(Z)虚部
\sphinxbeforeendvarwidth
\end{varwidth}%
&\begin{varwidth}[t]{\sphinxcolwidth{1}{4}}
\sphinxAtStartPar
Im(Zxx·Zyy \sphinxhyphen{} Zxy·Zyx)
\sphinxbeforeendvarwidth
\end{varwidth}%
\\
\sphinxhline\begin{varwidth}[t]{\sphinxcolwidth{1}{4}}
\sphinxAtStartPar
DetZa
\sphinxbeforeendvarwidth
\end{varwidth}%
&\begin{varwidth}[t]{\sphinxcolwidth{1}{4}}
\sphinxAtStartPar
MT\_DetZa
\sphinxbeforeendvarwidth
\end{varwidth}%
&\begin{varwidth}[t]{\sphinxcolwidth{1}{4}}
\sphinxAtStartPar
det(Z)幅值
\sphinxbeforeendvarwidth
\end{varwidth}%
&\begin{varwidth}[t]{\sphinxcolwidth{1}{4}}
\sphinxAtStartPar
|Zxx·Zyy \sphinxhyphen{} Zxy·Zyx|
\sphinxbeforeendvarwidth
\end{varwidth}%
\\
\sphinxhline\begin{varwidth}[t]{\sphinxcolwidth{1}{4}}
\sphinxAtStartPar
DetZp
\sphinxbeforeendvarwidth
\end{varwidth}%
&\begin{varwidth}[t]{\sphinxcolwidth{1}{4}}
\sphinxAtStartPar
MT\_DetZp
\sphinxbeforeendvarwidth
\end{varwidth}%
&\begin{varwidth}[t]{\sphinxcolwidth{1}{4}}
\sphinxAtStartPar
det(Z)相位
\sphinxbeforeendvarwidth
\end{varwidth}%
&\begin{varwidth}[t]{\sphinxcolwidth{1}{4}}
\sphinxAtStartPar
arg(Zxx·Zyy \sphinxhyphen{} Zxy·Zyx)
\sphinxbeforeendvarwidth
\end{varwidth}%
\\
\sphinxbottomrule
\end{tabulary}
\sphinxtableafterendhook\par
\sphinxattableend\end{savenotes}


\paragraph{5.17.16 数据筛选常用参数推荐}
\label{\detokenize{chapters/chapter6:id84}}
\sphinxAtStartPar
\sphinxstylestrong{推荐用于数据质量筛选的参数组合：}


\begin{savenotes}\sphinxattablestart
\sphinxthistablewithglobalstyle
\centering
\begin{tabulary}{\linewidth}[t]{TTT}
\sphinxtoprule
\begin{varwidth}[t]{\sphinxcolwidth{1}{3}}
\sphinxstyletheadfamily \sphinxAtStartPar
筛选类型
\sphinxbeforeendvarwidth
\end{varwidth}%
&\begin{varwidth}[t]{\sphinxcolwidth{1}{3}}
\sphinxstyletheadfamily \sphinxAtStartPar
推荐参数
\sphinxbeforeendvarwidth
\end{varwidth}%
&\begin{varwidth}[t]{\sphinxcolwidth{1}{3}}
\sphinxstyletheadfamily \sphinxAtStartPar
筛选标准
\sphinxbeforeendvarwidth
\end{varwidth}%
\\
\sphinxmidrule
\sphinxtableatstartofbodyhook\begin{varwidth}[t]{\sphinxcolwidth{1}{3}}
\sphinxAtStartPar
相干度筛选
\sphinxbeforeendvarwidth
\end{varwidth}%
&\begin{varwidth}[t]{\sphinxcolwidth{1}{3}}
\sphinxAtStartPar
CohExHy, CohEyHx
\sphinxbeforeendvarwidth
\end{varwidth}%
&\begin{varwidth}[t]{\sphinxcolwidth{1}{3}}
\sphinxAtStartPar
> 0.7 (优秀), > 0.5 (合格)
\sphinxbeforeendvarwidth
\end{varwidth}%
\\
\sphinxhline\begin{varwidth}[t]{\sphinxcolwidth{1}{3}}
\sphinxAtStartPar
信噪比筛选
\sphinxbeforeendvarwidth
\end{varwidth}%
&\begin{varwidth}[t]{\sphinxcolwidth{1}{3}}
\sphinxAtStartPar
exsnr, eysnr, hxsnr, hysnr
\sphinxbeforeendvarwidth
\end{varwidth}%
&\begin{varwidth}[t]{\sphinxcolwidth{1}{3}}
\sphinxAtStartPar
> 3 (推荐)
\sphinxbeforeendvarwidth
\end{varwidth}%
\\
\sphinxhline\begin{varwidth}[t]{\sphinxcolwidth{1}{3}}
\sphinxAtStartPar
误差筛选
\sphinxbeforeendvarwidth
\end{varwidth}%
&\begin{varwidth}[t]{\sphinxcolwidth{1}{3}}
\sphinxAtStartPar
rxyvar, ryxvar, pxyvar, pyxvar
\sphinxbeforeendvarwidth
\end{varwidth}%
&\begin{varwidth}[t]{\sphinxcolwidth{1}{3}}
\sphinxAtStartPar
越小越好
\sphinxbeforeendvarwidth
\end{varwidth}%
\\
\sphinxhline\begin{varwidth}[t]{\sphinxcolwidth{1}{3}}
\sphinxAtStartPar
三维性筛选
\sphinxbeforeendvarwidth
\end{varwidth}%
&\begin{varwidth}[t]{\sphinxcolwidth{1}{3}}
\sphinxAtStartPar
beta, ptskew2d
\sphinxbeforeendvarwidth
\end{varwidth}%
&\begin{varwidth}[t]{\sphinxcolwidth{1}{3}}
\sphinxAtStartPar
< 3° (近2D), < 6° (可接受)
\sphinxbeforeendvarwidth
\end{varwidth}%
\\
\sphinxbottomrule
\end{tabulary}
\sphinxtableafterendhook\par
\sphinxattableend\end{savenotes}


\paragraph{5.17.17 参数使用建议}
\label{\detokenize{chapters/chapter6:id85}}\begin{enumerate}
\sphinxsetlistlabels{\arabic}{enumi}{enumii}{}{.}%
\item {} 
\sphinxAtStartPar
\sphinxstylestrong{初学者}：优先关注 rxy、ryx（视电阻率）和 zxyp、zyxp（相位）

\item {} 
\sphinxAtStartPar
\sphinxstylestrong{质量控制}：使用相干度（CohExHy、CohEyHx）和信噪比（SNR）

\item {} 
\sphinxAtStartPar
\sphinxstylestrong{构造分析}：关注相位张量参数（alpha、beta、skew）

\item {} 
\sphinxAtStartPar
\sphinxstylestrong{噪声诊断}：检查功率谱密度和各通道噪声水平

\end{enumerate}


\paragraph{5.18 EDI导出格式}
\label{\detokenize{chapters/chapter6:edi}}
\sphinxAtStartPar
MTDataPro支持多种EDI格式导出，方便与其他MT软件进行数据交换。

\sphinxAtStartPar
\sphinxstylestrong{各EDI导出格式说明：}


\begin{savenotes}\sphinxattablestart
\sphinxthistablewithglobalstyle
\centering
\begin{tabulary}{\linewidth}[t]{TTTT}
\sphinxtoprule
\begin{varwidth}[t]{\sphinxcolwidth{1}{4}}
\sphinxstyletheadfamily \sphinxAtStartPar
格式名称
\sphinxbeforeendvarwidth
\end{varwidth}%
&\begin{varwidth}[t]{\sphinxcolwidth{1}{4}}
\sphinxstyletheadfamily \sphinxAtStartPar
菜单路径
\sphinxbeforeendvarwidth
\end{varwidth}%
&\begin{varwidth}[t]{\sphinxcolwidth{1}{4}}
\sphinxstyletheadfamily \sphinxAtStartPar
说明
\sphinxbeforeendvarwidth
\end{varwidth}%
&\begin{varwidth}[t]{\sphinxcolwidth{1}{4}}
\sphinxstyletheadfamily \sphinxAtStartPar
适用场景
\sphinxbeforeendvarwidth
\end{varwidth}%
\\
\sphinxmidrule
\sphinxtableatstartofbodyhook\begin{varwidth}[t]{\sphinxcolwidth{1}{4}}
\sphinxAtStartPar
SpeEDI
\sphinxbeforeendvarwidth
\end{varwidth}%
&\begin{varwidth}[t]{\sphinxcolwidth{1}{4}}
\sphinxAtStartPar
组导出 → SpeEDI
\sphinxbeforeendvarwidth
\end{varwidth}%
&\begin{varwidth}[t]{\sphinxcolwidth{1}{4}}
\sphinxAtStartPar
光谱EDI格式
\sphinxbeforeendvarwidth
\end{varwidth}%
&\begin{varwidth}[t]{\sphinxcolwidth{1}{4}}
\sphinxAtStartPar
Phoenix MT系统数据兼容
\sphinxbeforeendvarwidth
\end{varwidth}%
\\
\sphinxhline\begin{varwidth}[t]{\sphinxcolwidth{1}{4}}
\sphinxAtStartPar
ZTEDI
\sphinxbeforeendvarwidth
\end{varwidth}%
&\begin{varwidth}[t]{\sphinxcolwidth{1}{4}}
\sphinxAtStartPar
组导出 → ZTEDI
\sphinxbeforeendvarwidth
\end{varwidth}%
&\begin{varwidth}[t]{\sphinxcolwidth{1}{4}}
\sphinxAtStartPar
Z\sphinxhyphen{}Tensor EDI格式
\sphinxbeforeendvarwidth
\end{varwidth}%
&\begin{varwidth}[t]{\sphinxcolwidth{1}{4}}
\sphinxAtStartPar
垂直分量数据导出
\sphinxbeforeendvarwidth
\end{varwidth}%
\\
\sphinxhline\begin{varwidth}[t]{\sphinxcolwidth{1}{4}}
\sphinxAtStartPar
MTpyEDI
\sphinxbeforeendvarwidth
\end{varwidth}%
&\begin{varwidth}[t]{\sphinxcolwidth{1}{4}}
\sphinxAtStartPar
组导出 → MTpyEDI
\sphinxbeforeendvarwidth
\end{varwidth}%
&\begin{varwidth}[t]{\sphinxcolwidth{1}{4}}
\sphinxAtStartPar
MTpy库兼容格式
\sphinxbeforeendvarwidth
\end{varwidth}%
&\begin{varwidth}[t]{\sphinxcolwidth{1}{4}}
\sphinxAtStartPar
Python MTpy数据分析
\sphinxbeforeendvarwidth
\end{varwidth}%
\\
\sphinxhline\begin{varwidth}[t]{\sphinxcolwidth{1}{4}}
\sphinxAtStartPar
PLT
\sphinxbeforeendvarwidth
\end{varwidth}%
&\begin{varwidth}[t]{\sphinxcolwidth{1}{4}}
\sphinxAtStartPar
组导出 → PLT
\sphinxbeforeendvarwidth
\end{varwidth}%
&\begin{varwidth}[t]{\sphinxcolwidth{1}{4}}
\sphinxAtStartPar
绘图格式
\sphinxbeforeendvarwidth
\end{varwidth}%
&\begin{varwidth}[t]{\sphinxcolwidth{1}{4}}
\sphinxAtStartPar
数据可视化与绘图
\sphinxbeforeendvarwidth
\end{varwidth}%
\\
\sphinxbottomrule
\end{tabulary}
\sphinxtableafterendhook\par
\sphinxattableend\end{savenotes}

\sphinxAtStartPar
\sphinxstylestrong{SpeEDI格式}
\begin{itemize}
\item {} 
\sphinxAtStartPar
基于标准EDI格式的光谱版本

\item {} 
\sphinxAtStartPar
保留完整频率信息

\item {} 
\sphinxAtStartPar
适用于Phoenix系列仪器数据

\end{itemize}

\sphinxAtStartPar
\sphinxstylestrong{ZTEDI格式}
\begin{itemize}
\item {} 
\sphinxAtStartPar
专门用于垂直电场(Z)分量

\item {} 
\sphinxAtStartPar
包含Z数据和张量信息

\item {} 
\sphinxAtStartPar
适用于大地电磁测深研究

\end{itemize}

\sphinxAtStartPar
\sphinxstylestrong{MTpyEDI格式}
\begin{itemize}
\item {} 
\sphinxAtStartPar
兼容Python的MTpy库

\item {} 
\sphinxAtStartPar
优化了数据类型和结构

\item {} 
\sphinxAtStartPar
便于在Python环境中进行数据分析

\end{itemize}

\sphinxAtStartPar
\sphinxstylestrong{PLT格式}
\begin{itemize}
\item {} 
\sphinxAtStartPar
用于数据绘图和可视化

\item {} 
\sphinxAtStartPar
包含绘图所需的坐标数据

\item {} 
\sphinxAtStartPar
适用于制作publication\sphinxhyphen{}quality图形

\end{itemize}

\sphinxstepscope


\section{📤 第7章 结果输出与可视化}
\label{\detokenize{chapters/chapter7:id1}}\label{\detokenize{chapters/chapter7::doc}}
\sphinxAtStartPar
本章介绍MTDP中数据处理结果的查看、导出和可视化功能。


\bigskip\hrule\bigskip



\subsection{7.1 数据图表查看}
\label{\detokenize{chapters/chapter7:id2}}

\subsubsection{7.1.1 视电阻率/相位曲线}
\label{\detokenize{chapters/chapter7:id3}}\begin{enumerate}
\sphinxsetlistlabels{\arabic}{enumi}{enumii}{}{.}%
\item {} 
\sphinxAtStartPar
在工程树中选择测点

\item {} 
\sphinxAtStartPar
双击或右键选择"查看曲线"

\item {} 
\sphinxAtStartPar
窗口显示视电阻率（ρa）和相位（φ）曲线

\end{enumerate}

\sphinxAtStartPar
\sphinxstylestrong{曲线显示选项：}
\begin{itemize}
\item {} 
\sphinxAtStartPar
\sphinxstylestrong{XY/YX模式}：显示 ρxy/φxy 和 ρyx/φyx

\item {} 
\sphinxAtStartPar
\sphinxstylestrong{误差棒}：显示/隐藏误差棒

\item {} 
\sphinxAtStartPar
\sphinxstylestrong{数据点标记}：标记数据点位置

\item {} 
\sphinxAtStartPar
\sphinxstylestrong{对数/线性坐标}：切换坐标轴类型

\end{itemize}


\subsubsection{7.1.2 多测点对比}
\label{\detokenize{chapters/chapter7:id4}}\begin{enumerate}
\sphinxsetlistlabels{\arabic}{enumi}{enumii}{}{.}%
\item {} 
\sphinxAtStartPar
在工程树中选择多个测点（按住Ctrl多选）

\item {} 
\sphinxAtStartPar
右键选择"对比曲线"

\item {} 
\sphinxAtStartPar
所有选中的测点曲线将显示在同一图表中

\end{enumerate}


\subsubsection{7.1.3 图表交互操作}
\label{\detokenize{chapters/chapter7:id5}}

\begin{savenotes}\sphinxattablestart
\sphinxthistablewithglobalstyle
\centering
\begin{tabulary}{\linewidth}[t]{TT}
\sphinxtoprule
\begin{varwidth}[t]{\sphinxcolwidth{1}{2}}
\sphinxstyletheadfamily \sphinxAtStartPar
操作
\sphinxbeforeendvarwidth
\end{varwidth}%
&\begin{varwidth}[t]{\sphinxcolwidth{1}{2}}
\sphinxstyletheadfamily \sphinxAtStartPar
说明
\sphinxbeforeendvarwidth
\end{varwidth}%
\\
\sphinxmidrule
\sphinxtableatstartofbodyhook\begin{varwidth}[t]{\sphinxcolwidth{1}{2}}
\sphinxAtStartPar
鼠标滚轮
\sphinxbeforeendvarwidth
\end{varwidth}%
&\begin{varwidth}[t]{\sphinxcolwidth{1}{2}}
\sphinxAtStartPar
缩放图表
\sphinxbeforeendvarwidth
\end{varwidth}%
\\
\sphinxhline\begin{varwidth}[t]{\sphinxcolwidth{1}{2}}
\sphinxAtStartPar
左键拖动
\sphinxbeforeendvarwidth
\end{varwidth}%
&\begin{varwidth}[t]{\sphinxcolwidth{1}{2}}
\sphinxAtStartPar
平移图表
\sphinxbeforeendvarwidth
\end{varwidth}%
\\
\sphinxhline\begin{varwidth}[t]{\sphinxcolwidth{1}{2}}
\sphinxAtStartPar
双击
\sphinxbeforeendvarwidth
\end{varwidth}%
&\begin{varwidth}[t]{\sphinxcolwidth{1}{2}}
\sphinxAtStartPar
恢复原始视图
\sphinxbeforeendvarwidth
\end{varwidth}%
\\
\sphinxhline\begin{varwidth}[t]{\sphinxcolwidth{1}{2}}
\sphinxAtStartPar
右键菜单
\sphinxbeforeendvarwidth
\end{varwidth}%
&\begin{varwidth}[t]{\sphinxcolwidth{1}{2}}
\sphinxAtStartPar
导出、设置等选项
\sphinxbeforeendvarwidth
\end{varwidth}%
\\
\sphinxbottomrule
\end{tabulary}
\sphinxtableafterendhook\par
\sphinxattableend\end{savenotes}


\bigskip\hrule\bigskip



\subsection{7.2 EDI格式导出}
\label{\detokenize{chapters/chapter7:edi}}

\subsubsection{6.2.1 批量导出（详细流程）}
\label{\detokenize{chapters/chapter7:id6}}

\paragraph{7.2.1.1 批量导出概述}
\label{\detokenize{chapters/chapter7:id7}}
\sphinxAtStartPar
批量导出功能支持一次性导出多个测点的处理结果，大幅提高工作效率。

\sphinxAtStartPar
\sphinxstylestrong{支持格式：} SpeEDI | ZTEDI | MTpyEDI | PLT | KML/KMZ


\paragraph{7.2.1.2 导入前准备}
\label{\detokenize{chapters/chapter7:id8}}
\sphinxAtStartPar
\sphinxstylestrong{检查清单：}
\begin{itemize}
\item {} 
\sphinxAtStartPar
 确认所有测点已完成FFT计算

\item {} 
\sphinxAtStartPar
 检查数据质量评估（相干性、残差等）

\item {} 
\sphinxAtStartPar
 确认需要导出的格式

\item {} 
\sphinxAtStartPar
 准备输出目录和命名规则

\item {} 
\sphinxAtStartPar
 检查磁盘空间充足

\end{itemize}

\sphinxAtStartPar
\sphinxstylestrong{质量评估标准：}


\begin{savenotes}\sphinxattablestart
\sphinxthistablewithglobalstyle
\centering
\begin{tabulary}{\linewidth}[t]{TTT}
\sphinxtoprule
\begin{varwidth}[t]{\sphinxcolwidth{1}{3}}
\sphinxstyletheadfamily \sphinxAtStartPar
指标
\sphinxbeforeendvarwidth
\end{varwidth}%
&\begin{varwidth}[t]{\sphinxcolwidth{1}{3}}
\sphinxstyletheadfamily \sphinxAtStartPar
合格范围
\sphinxbeforeendvarwidth
\end{varwidth}%
&\begin{varwidth}[t]{\sphinxcolwidth{1}{3}}
\sphinxstyletheadfamily \sphinxAtStartPar
说明
\sphinxbeforeendvarwidth
\end{varwidth}%
\\
\sphinxmidrule
\sphinxtableatstartofbodyhook\begin{varwidth}[t]{\sphinxcolwidth{1}{3}}
\sphinxAtStartPar
相干性
\sphinxbeforeendvarwidth
\end{varwidth}%
&\begin{varwidth}[t]{\sphinxcolwidth{1}{3}}
\sphinxAtStartPar
>0.7
\sphinxbeforeendvarwidth
\end{varwidth}%
&\begin{varwidth}[t]{\sphinxcolwidth{1}{3}}
\sphinxAtStartPar
优秀，可导出
\sphinxbeforeendvarwidth
\end{varwidth}%
\\
\sphinxhline\begin{varwidth}[t]{\sphinxcolwidth{1}{3}}
\sphinxAtStartPar
相干性
\sphinxbeforeendvarwidth
\end{varwidth}%
&\begin{varwidth}[t]{\sphinxcolwidth{1}{3}}
\sphinxAtStartPar
0.5\sphinxhyphen{}0.7
\sphinxbeforeendvarwidth
\end{varwidth}%
&\begin{varwidth}[t]{\sphinxcolwidth{1}{3}}
\sphinxAtStartPar
良好，谨慎使用
\sphinxbeforeendvarwidth
\end{varwidth}%
\\
\sphinxhline\begin{varwidth}[t]{\sphinxcolwidth{1}{3}}
\sphinxAtStartPar
相干性
\sphinxbeforeendvarwidth
\end{varwidth}%
&\begin{varwidth}[t]{\sphinxcolwidth{1}{3}}
\sphinxAtStartPar
<0.5
\sphinxbeforeendvarwidth
\end{varwidth}%
&\begin{varwidth}[t]{\sphinxcolwidth{1}{3}}
\sphinxAtStartPar
数据质量差，建议重新处理
\sphinxbeforeendvarwidth
\end{varwidth}%
\\
\sphinxhline\begin{varwidth}[t]{\sphinxcolwidth{1}{3}}
\sphinxAtStartPar
残差
\sphinxbeforeendvarwidth
\end{varwidth}%
&\begin{varwidth}[t]{\sphinxcolwidth{1}{3}}
\sphinxAtStartPar
<10°
\sphinxbeforeendvarwidth
\end{varwidth}%
&\begin{varwidth}[t]{\sphinxcolwidth{1}{3}}
\sphinxAtStartPar
优秀
\sphinxbeforeendvarwidth
\end{varwidth}%
\\
\sphinxhline\begin{varwidth}[t]{\sphinxcolwidth{1}{3}}
\sphinxAtStartPar
残差
\sphinxbeforeendvarwidth
\end{varwidth}%
&\begin{varwidth}[t]{\sphinxcolwidth{1}{3}}
\sphinxAtStartPar
10°\sphinxhyphen{}20°
\sphinxbeforeendvarwidth
\end{varwidth}%
&\begin{varwidth}[t]{\sphinxcolwidth{1}{3}}
\sphinxAtStartPar
良好
\sphinxbeforeendvarwidth
\end{varwidth}%
\\
\sphinxhline\begin{varwidth}[t]{\sphinxcolwidth{1}{3}}
\sphinxAtStartPar
残差
\sphinxbeforeendvarwidth
\end{varwidth}%
&\begin{varwidth}[t]{\sphinxcolwidth{1}{3}}
\sphinxAtStartPar
>20°
\sphinxbeforeendvarwidth
\end{varwidth}%
&\begin{varwidth}[t]{\sphinxcolwidth{1}{3}}
\sphinxAtStartPar
可疑，检查数据
\sphinxbeforeendvarwidth
\end{varwidth}%
\\
\sphinxbottomrule
\end{tabulary}
\sphinxtableafterendhook\par
\sphinxattableend\end{savenotes}


\paragraph{7.2.1.3 批量导出操作步骤}
\label{\detokenize{chapters/chapter7:id9}}
\sphinxAtStartPar
\sphinxstylestrong{步骤1：启动批量导出}
\begin{enumerate}
\sphinxsetlistlabels{\arabic}{enumi}{enumii}{}{.}%
\item {} 
\sphinxAtStartPar
在工程树中选择工区或测段

\item {} 
\sphinxAtStartPar
右键选择 \sphinxcode{\sphinxupquote{导出 → 批量导出}}

\item {} 
\sphinxAtStartPar
系统打开批量导出对话框

\end{enumerate}

\sphinxAtStartPar
\sphinxstylestrong{步骤2：配置导出参数}

\sphinxAtStartPar
在批量导出对话框中设置：


\begin{savenotes}\sphinxattablestart
\sphinxthistablewithglobalstyle
\centering
\begin{tabulary}{\linewidth}[t]{TTT}
\sphinxtoprule
\begin{varwidth}[t]{\sphinxcolwidth{1}{3}}
\sphinxstyletheadfamily \sphinxAtStartPar
参数
\sphinxbeforeendvarwidth
\end{varwidth}%
&\begin{varwidth}[t]{\sphinxcolwidth{1}{3}}
\sphinxstyletheadfamily \sphinxAtStartPar
说明
\sphinxbeforeendvarwidth
\end{varwidth}%
&\begin{varwidth}[t]{\sphinxcolwidth{1}{3}}
\sphinxstyletheadfamily \sphinxAtStartPar
推荐设置
\sphinxbeforeendvarwidth
\end{varwidth}%
\\
\sphinxmidrule
\sphinxtableatstartofbodyhook\begin{varwidth}[t]{\sphinxcolwidth{1}{3}}
\sphinxAtStartPar
导出格式
\sphinxbeforeendvarwidth
\end{varwidth}%
&\begin{varwidth}[t]{\sphinxcolwidth{1}{3}}
\sphinxAtStartPar
根据接收软件选择
\sphinxbeforeendvarwidth
\end{varwidth}%
&\begin{varwidth}[t]{\sphinxcolwidth{1}{3}}
\sphinxAtStartPar
SpeEDI/ZTEDI/MTpyEDI
\sphinxbeforeendvarwidth
\end{varwidth}%
\\
\sphinxhline\begin{varwidth}[t]{\sphinxcolwidth{1}{3}}
\sphinxAtStartPar
文件命名规则
\sphinxbeforeendvarwidth
\end{varwidth}%
&\begin{varwidth}[t]{\sphinxcolwidth{1}{3}}
\sphinxAtStartPar
\{测点名\}.\{格式扩展名\}
\sphinxbeforeendvarwidth
\end{varwidth}%
&\begin{varwidth}[t]{\sphinxcolwidth{1}{3}}
\sphinxAtStartPar
\{测点\}\_EDI.edi
\sphinxbeforeendvarwidth
\end{varwidth}%
\\
\sphinxhline\begin{varwidth}[t]{\sphinxcolwidth{1}{3}}
\sphinxAtStartPar
覆盖已存在
\sphinxbeforeendvarwidth
\end{varwidth}%
&\begin{varwidth}[t]{\sphinxcolwidth{1}{3}}
\sphinxAtStartPar
不覆盖（安全）
\sphinxbeforeendvarwidth
\end{varwidth}%
&\begin{varwidth}[t]{\sphinxcolwidth{1}{3}}
\sphinxAtStartPar
不覆盖
\sphinxbeforeendvarwidth
\end{varwidth}%
\\
\sphinxhline\begin{varwidth}[t]{\sphinxcolwidth{1}{3}}
\sphinxAtStartPar
导出范围
\sphinxbeforeendvarwidth
\end{varwidth}%
&\begin{varwidth}[t]{\sphinxcolwidth{1}{3}}
\sphinxAtStartPar
全部测点/选中测点
\sphinxbeforeendvarwidth
\end{varwidth}%
&\begin{varwidth}[t]{\sphinxcolwidth{1}{3}}
\sphinxAtStartPar
全部测点
\sphinxbeforeendvarwidth
\end{varwidth}%
\\
\sphinxhline\begin{varwidth}[t]{\sphinxcolwidth{1}{3}}
\sphinxAtStartPar
质量过滤
\sphinxbeforeendvarwidth
\end{varwidth}%
&\begin{varwidth}[t]{\sphinxcolwidth{1}{3}}
\sphinxAtStartPar
仅导出相干性>X的测点
\sphinxbeforeendvarwidth
\end{varwidth}%
&\begin{varwidth}[t]{\sphinxcolwidth{1}{3}}
\sphinxAtStartPar
不限
\sphinxbeforeendvarwidth
\end{varwidth}%
\\
\sphinxhline\begin{varwidth}[t]{\sphinxcolwidth{1}{3}}
\sphinxAtStartPar
包含误差棒
\sphinxbeforeendvarwidth
\end{varwidth}%
&\begin{varwidth}[t]{\sphinxcolwidth{1}{3}}
\sphinxAtStartPar
是/否
\sphinxbeforeendvarwidth
\end{varwidth}%
&\begin{varwidth}[t]{\sphinxcolwidth{1}{3}}
\sphinxAtStartPar
是
\sphinxbeforeendvarwidth
\end{varwidth}%
\\
\sphinxbottomrule
\end{tabulary}
\sphinxtableafterendhook\par
\sphinxattableend\end{savenotes}

\sphinxAtStartPar
\sphinxstylestrong{步骤3：选择导出项目}

\sphinxAtStartPar
可选择导出的数据项：
\begin{itemize}
\item {} 
\sphinxAtStartPar
 阻抗张量（Zxx, Zxy, Zyx, Zyy）

\item {} 
\sphinxAtStartPar
 相位张量（Pxy, Pyx）

\item {} 
\sphinxAtStartPar
 倾子张量（Tzx, Tzy）

\item {} 
\sphinxAtStartPar
 视电阻率张量（Rhoxx, Rhoxy, Rhoyx, Rhoyy）

\item {} 
\sphinxAtStartPar
 相干度曲线

\item {} 
\sphinxAtStartPar
 数据质量信息

\end{itemize}

\sphinxAtStartPar
\sphinxstylestrong{步骤4：执行导出}

\sphinxAtStartPar
点击"开始导出"后：
\begin{enumerate}
\sphinxsetlistlabels{\arabic}{enumi}{enumii}{}{.}%
\item {} 
\sphinxAtStartPar
系统显示进度条

\item {} 
\sphinxAtStartPar
逐个处理测点数据

\item {} 
\sphinxAtStartPar
生成EDI文件并保存

\item {} 
\sphinxAtStartPar
完成后显示导出摘要

\end{enumerate}


\paragraph{7.2.1.4 导出后验证}
\label{\detokenize{chapters/chapter7:id10}}
\sphinxAtStartPar
\sphinxstylestrong{检查项目：}


\begin{savenotes}\sphinxattablestart
\sphinxthistablewithglobalstyle
\centering
\begin{tabulary}{\linewidth}[t]{TT}
\sphinxtoprule
\begin{varwidth}[t]{\sphinxcolwidth{1}{2}}
\sphinxstyletheadfamily \sphinxAtStartPar
检查项
\sphinxbeforeendvarwidth
\end{varwidth}%
&\begin{varwidth}[t]{\sphinxcolwidth{1}{2}}
\sphinxstyletheadfamily \sphinxAtStartPar
验证方法
\sphinxbeforeendvarwidth
\end{varwidth}%
\\
\sphinxmidrule
\sphinxtableatstartofbodyhook\begin{varwidth}[t]{\sphinxcolwidth{1}{2}}
\sphinxAtStartPar
文件数量
\sphinxbeforeendvarwidth
\end{varwidth}%
&\begin{varwidth}[t]{\sphinxcolwidth{1}{2}}
\sphinxAtStartPar
对比导出数量与源测点数
\sphinxbeforeendvarwidth
\end{varwidth}%
\\
\sphinxhline\begin{varwidth}[t]{\sphinxcolwidth{1}{2}}
\sphinxAtStartPar
文件大小
\sphinxbeforeendvarwidth
\end{varwidth}%
&\begin{varwidth}[t]{\sphinxcolwidth{1}{2}}
\sphinxAtStartPar
确认EDI文件大小合理
\sphinxbeforeendvarwidth
\end{varwidth}%
\\
\sphinxhline\begin{varwidth}[t]{\sphinxcolwidth{1}{2}}
\sphinxAtStartPar
打开测试
\sphinxbeforeendvarwidth
\end{varwidth}%
&\begin{varwidth}[t]{\sphinxcolwidth{1}{2}}
\sphinxAtStartPar
用专业软件打开验证
\sphinxbeforeendvarwidth
\end{varwidth}%
\\
\sphinxhline\begin{varwidth}[t]{\sphinxcolwidth{1}{2}}
\sphinxAtStartPar
数据完整性
\sphinxbeforeendvarwidth
\end{varwidth}%
&\begin{varwidth}[t]{\sphinxcolwidth{1}{2}}
\sphinxAtStartPar
检查关键频点是否导出
\sphinxbeforeendvarwidth
\end{varwidth}%
\\
\sphinxbottomrule
\end{tabulary}
\sphinxtableafterendhook\par
\sphinxattableend\end{savenotes}


\paragraph{7.2.1.5 常见问题}
\label{\detokenize{chapters/chapter7:id11}}
\sphinxAtStartPar
\sphinxstylestrong{问题1：部分测点导出失败}
\begin{itemize}
\item {} 
\sphinxAtStartPar
原因：数据质量不足或计算未完成

\item {} 
\sphinxAtStartPar
解决：
\begin{enumerate}
\sphinxsetlistlabels{\arabic}{enumi}{enumii}{}{.}%
\item {} 
\sphinxAtStartPar
检查失败测点的FFT状态

\item {} 
\sphinxAtStartPar
查看"处理"选项卡中的计算日志

\item {} 
\sphinxAtStartPar
重新完成FFT计算后再导出

\end{enumerate}

\end{itemize}

\sphinxAtStartPar
\sphinxstylestrong{问题2：导出格式不兼容}
\begin{itemize}
\item {} 
\sphinxAtStartPar
原因：EDI版本不符合接收软件要求

\item {} 
\sphinxAtStartPar
解决：
\begin{enumerate}
\sphinxsetlistlabels{\arabic}{enumi}{enumii}{}{.}%
\item {} 
\sphinxAtStartPar
确认目标软件的EDI版本要求

\item {} 
\sphinxAtStartPar
使用批量导出的格式选择功能

\item {} 
\sphinxAtStartPar
重新导出兼容版本

\end{enumerate}

\end{itemize}

\sphinxAtStartPar
\sphinxstylestrong{问题3：文件名冲突}
\begin{itemize}
\item {} 
\sphinxAtStartPar
原因：导出目录中已存在同名文件

\item {} 
\sphinxAtStartPar
解决：
\begin{enumerate}
\sphinxsetlistlabels{\arabic}{enumi}{enumii}{}{.}%
\item {} 
\sphinxAtStartPar
使用"覆盖已存在"选项（如需要）

\item {} 
\sphinxAtStartPar
使用不同的命名规则

\item {} 
\sphinxAtStartPar
手动重命名冲突文件

\end{enumerate}

\end{itemize}

\sphinxAtStartPar
\sphinxstylestrong{问题4：磁盘空间不足}
\begin{itemize}
\item {} 
\sphinxAtStartPar
原因：批量导出数据量大

\item {} 
\sphinxAtStartPar
解决：
\begin{enumerate}
\sphinxsetlistlabels{\arabic}{enumi}{enumii}{}{.}%
\item {} 
\sphinxAtStartPar
评估所需磁盘空间（约1MB/测点）

\item {} 
\sphinxAtStartPar
分批导出而非一次性全部导出

\item {} 
\sphinxAtStartPar
使用压缩或分盘存储

\end{enumerate}

\end{itemize}


\paragraph{7.2.1.6 批量导出最佳实践}
\label{\detokenize{chapters/chapter7:id12}}
\sphinxAtStartPar
\sphinxstylestrong{文件组织建议：}
\begin{itemize}
\item {} 
\sphinxAtStartPar
为每个测线创建独立文件夹

\item {} 
\sphinxAtStartPar
使用统一的命名约定（如\{测线\}\sphinxstyleemphasis{\{日期\}}\{测点\}.edi）

\item {} 
\sphinxAtStartPar
按时间戳或质量级别分批存储

\item {} 
\sphinxAtStartPar
保留原始工程文件作为备份

\end{itemize}

\sphinxAtStartPar
\sphinxstylestrong{导出工作流建议：}
\begin{enumerate}
\sphinxsetlistlabels{\arabic}{enumi}{enumii}{}{.}%
\item {} 
\sphinxAtStartPar
先用少量测点测试导出流程

\item {} 
\sphinxAtStartPar
验证导出文件可以被目标软件正确读取

\item {} 
\sphinxAtStartPar
确认所有必需元数据都包含在导出文件中

\item {} 
\sphinxAtStartPar
建立导出质量检查清单，确保数据完整性

\end{enumerate}


\subsubsection{7.2.2 导出操作}
\label{\detokenize{chapters/chapter7:id13}}
\sphinxAtStartPar
\sphinxstylestrong{单测点导出：}
\begin{enumerate}
\sphinxsetlistlabels{\arabic}{enumi}{enumii}{}{.}%
\item {} 
\sphinxAtStartPar
在工程树中选择测点

\item {} 
\sphinxAtStartPar
右键选择相应的导出选项

\item {} 
\sphinxAtStartPar
选择保存位置

\end{enumerate}

\sphinxAtStartPar
\sphinxstylestrong{批量导出：}
\begin{enumerate}
\sphinxsetlistlabels{\arabic}{enumi}{enumii}{}{.}%
\item {} 
\sphinxAtStartPar
在工程树中选择测段（Group）

\item {} 
\sphinxAtStartPar
右键选择批量导出选项

\item {} 
\sphinxAtStartPar
选择输出目录

\item {} 
\sphinxAtStartPar
系统自动为每个测点生成EDI文件

\end{enumerate}

\sphinxAtStartPar
\sphinxstylestrong{工区级导出：}
\begin{enumerate}
\sphinxsetlistlabels{\arabic}{enumi}{enumii}{}{.}%
\item {} 
\sphinxAtStartPar
选择工区

\item {} 
\sphinxAtStartPar
右键选择 \sphinxcode{\sphinxupquote{导出EDI}}

\item {} 
\sphinxAtStartPar
选择输出目录

\end{enumerate}

\sphinxAtStartPar
\sphinxstylestrong{HEAD段详细信息：}
\begin{itemize}
\item {} 
\sphinxAtStartPar
DataID：数据标识

\item {} 
\sphinxAtStartPar
AcquiredBy：采集者

\item {} 
\sphinxAtStartPar
AcquiredDate/EndDate：采集日期

\item {} 
\sphinxAtStartPar
Name：测点名称

\item {} 
\sphinxAtStartPar
Latitude/Longitude/Elevation：坐标信息

\end{itemize}


\subsubsection{7.2.4 EDI格式详细规范}
\label{\detokenize{chapters/chapter7:id14}}

\paragraph{EDI文件完整结构}
\label{\detokenize{chapters/chapter7:id15}}
\begin{sphinxVerbatim}[commandchars=\\\{\}]
\PYGZgt{}HEAD
    DATAID=MTSurvey2024
    PROJECT=Area1
    SURVEY=Line1
    YEAR=2024
    MONTH=01
    DAY=15
    PROJLAT=30.5
    PROJLONG=120.5
    PROJELEV=150.5
    ACQBY=MTTeam
    FILEBY=MTDP v1.9.4

\PYGZgt{}INFO
    MAXSITES=1
    UNITS=M
    ELEVUNIT=METERS
    COORDINATE=Geographic

\PYGZgt{}DEFINEMEAS
    MAXCHAN=5
    MAXFREQ=30
    UNITS=M
    REFTIM=2024\PYGZhy{}01\PYGZhy{}15 00:00:00

\PYGZgt{}SITE
    BIRPTIM=24
    LAT=30:15:30
    LONG=120:34:12
    ELEV=150.5
    ANGLE=0

\PYGZgt{}FREQ
    100.0000
    80.0000
    ...

\PYGZgt{}MTEDATA
    FREQ=100.0
    ZXXR=0.1234
    ZXXI=\PYGZhy{}0.0567
    ZXYR=1.2345
    ZXYI=\PYGZhy{}0.7890
    ZYXR=\PYGZhy{}1.4567
    ZYXI=0.8901
    ZYYR=0.0987
    ZYYI=\PYGZhy{}0.0432
    ...

\PYGZgt{}END
\end{sphinxVerbatim}


\paragraph{EDI数据段详解}
\label{\detokenize{chapters/chapter7:id16}}

\begin{savenotes}\sphinxattablestart
\sphinxthistablewithglobalstyle
\centering
\begin{tabulary}{\linewidth}[t]{TTTT}
\sphinxtoprule
\begin{varwidth}[t]{\sphinxcolwidth{1}{4}}
\sphinxstyletheadfamily \sphinxAtStartPar
段名
\sphinxbeforeendvarwidth
\end{varwidth}%
&\begin{varwidth}[t]{\sphinxcolwidth{1}{4}}
\sphinxstyletheadfamily \sphinxAtStartPar
必需
\sphinxbeforeendvarwidth
\end{varwidth}%
&\begin{varwidth}[t]{\sphinxcolwidth{1}{4}}
\sphinxstyletheadfamily \sphinxAtStartPar
内容
\sphinxbeforeendvarwidth
\end{varwidth}%
&\begin{varwidth}[t]{\sphinxcolwidth{1}{4}}
\sphinxstyletheadfamily \sphinxAtStartPar
说明
\sphinxbeforeendvarwidth
\end{varwidth}%
\\
\sphinxmidrule
\sphinxtableatstartofbodyhook\begin{varwidth}[t]{\sphinxcolwidth{1}{4}}
\sphinxAtStartPar
\sphinxstylestrong{HEAD}
\sphinxbeforeendvarwidth
\end{varwidth}%
&\begin{varwidth}[t]{\sphinxcolwidth{1}{4}}
\sphinxAtStartPar
是
\sphinxbeforeendvarwidth
\end{varwidth}%
&\begin{varwidth}[t]{\sphinxcolwidth{1}{4}}
\sphinxAtStartPar
头信息
\sphinxbeforeendvarwidth
\end{varwidth}%
&\begin{varwidth}[t]{\sphinxcolwidth{1}{4}}
\sphinxAtStartPar
包含项目、日期、坐标等基本信息
\sphinxbeforeendvarwidth
\end{varwidth}%
\\
\sphinxhline\begin{varwidth}[t]{\sphinxcolwidth{1}{4}}
\sphinxAtStartPar
\sphinxstylestrong{INFO}
\sphinxbeforeendvarwidth
\end{varwidth}%
&\begin{varwidth}[t]{\sphinxcolwidth{1}{4}}
\sphinxAtStartPar
是
\sphinxbeforeendvarwidth
\end{varwidth}%
&\begin{varwidth}[t]{\sphinxcolwidth{1}{4}}
\sphinxAtStartPar
数据说明
\sphinxbeforeendvarwidth
\end{varwidth}%
&\begin{varwidth}[t]{\sphinxcolwidth{1}{4}}
\sphinxAtStartPar
包含数据单位、坐标系统等
\sphinxbeforeendvarwidth
\end{varwidth}%
\\
\sphinxhline\begin{varwidth}[t]{\sphinxcolwidth{1}{4}}
\sphinxAtStartPar
\sphinxstylestrong{DEFINEMEAS}
\sphinxbeforeendvarwidth
\end{varwidth}%
&\begin{varwidth}[t]{\sphinxcolwidth{1}{4}}
\sphinxAtStartPar
是
\sphinxbeforeendvarwidth
\end{varwidth}%
&\begin{varwidth}[t]{\sphinxcolwidth{1}{4}}
\sphinxAtStartPar
测量定义
\sphinxbeforeendvarwidth
\end{varwidth}%
&\begin{varwidth}[t]{\sphinxcolwidth{1}{4}}
\sphinxAtStartPar
定义通道数、频率数等
\sphinxbeforeendvarwidth
\end{varwidth}%
\\
\sphinxhline\begin{varwidth}[t]{\sphinxcolwidth{1}{4}}
\sphinxAtStartPar
\sphinxstylestrong{SITE}
\sphinxbeforeendvarwidth
\end{varwidth}%
&\begin{varwidth}[t]{\sphinxcolwidth{1}{4}}
\sphinxAtStartPar
是
\sphinxbeforeendvarwidth
\end{varwidth}%
&\begin{varwidth}[t]{\sphinxcolwidth{1}{4}}
\sphinxAtStartPar
测点信息
\sphinxbeforeendvarwidth
\end{varwidth}%
&\begin{varwidth}[t]{\sphinxcolwidth{1}{4}}
\sphinxAtStartPar
测点坐标、角度等
\sphinxbeforeendvarwidth
\end{varwidth}%
\\
\sphinxhline\begin{varwidth}[t]{\sphinxcolwidth{1}{4}}
\sphinxAtStartPar
\sphinxstylestrong{FREQ}
\sphinxbeforeendvarwidth
\end{varwidth}%
&\begin{varwidth}[t]{\sphinxcolwidth{1}{4}}
\sphinxAtStartPar
是
\sphinxbeforeendvarwidth
\end{varwidth}%
&\begin{varwidth}[t]{\sphinxcolwidth{1}{4}}
\sphinxAtStartPar
频率列表
\sphinxbeforeendvarwidth
\end{varwidth}%
&\begin{varwidth}[t]{\sphinxcolwidth{1}{4}}
\sphinxAtStartPar
所有频率点（Hz）
\sphinxbeforeendvarwidth
\end{varwidth}%
\\
\sphinxhline\begin{varwidth}[t]{\sphinxcolwidth{1}{4}}
\sphinxAtStartPar
\sphinxstylestrong{MTEDATA}
\sphinxbeforeendvarwidth
\end{varwidth}%
&\begin{varwidth}[t]{\sphinxcolwidth{1}{4}}
\sphinxAtStartPar
是
\sphinxbeforeendvarwidth
\end{varwidth}%
&\begin{varwidth}[t]{\sphinxcolwidth{1}{4}}
\sphinxAtStartPar
MT数据
\sphinxbeforeendvarwidth
\end{varwidth}%
&\begin{varwidth}[t]{\sphinxcolwidth{1}{4}}
\sphinxAtStartPar
阻抗、倾子、误差等
\sphinxbeforeendvarwidth
\end{varwidth}%
\\
\sphinxhline\begin{varwidth}[t]{\sphinxcolwidth{1}{4}}
\sphinxAtStartPar
\sphinxstylestrong{END}
\sphinxbeforeendvarwidth
\end{varwidth}%
&\begin{varwidth}[t]{\sphinxcolwidth{1}{4}}
\sphinxAtStartPar
是
\sphinxbeforeendvarwidth
\end{varwidth}%
&\begin{varwidth}[t]{\sphinxcolwidth{1}{4}}
\sphinxAtStartPar
结束标记
\sphinxbeforeendvarwidth
\end{varwidth}%
&\begin{varwidth}[t]{\sphinxcolwidth{1}{4}}
\sphinxAtStartPar
文件结束
\sphinxbeforeendvarwidth
\end{varwidth}%
\\
\sphinxbottomrule
\end{tabulary}
\sphinxtableafterendhook\par
\sphinxattableend\end{savenotes}


\paragraph{MTEDATA数据字段}
\label{\detokenize{chapters/chapter7:mtedata}}

\begin{savenotes}\sphinxattablestart
\sphinxthistablewithglobalstyle
\centering
\begin{tabulary}{\linewidth}[t]{TTT}
\sphinxtoprule
\begin{varwidth}[t]{\sphinxcolwidth{1}{3}}
\sphinxstyletheadfamily \sphinxAtStartPar
字段
\sphinxbeforeendvarwidth
\end{varwidth}%
&\begin{varwidth}[t]{\sphinxcolwidth{1}{3}}
\sphinxstyletheadfamily \sphinxAtStartPar
单位
\sphinxbeforeendvarwidth
\end{varwidth}%
&\begin{varwidth}[t]{\sphinxcolwidth{1}{3}}
\sphinxstyletheadfamily \sphinxAtStartPar
说明
\sphinxbeforeendvarwidth
\end{varwidth}%
\\
\sphinxmidrule
\sphinxtableatstartofbodyhook\begin{varwidth}[t]{\sphinxcolwidth{1}{3}}
\sphinxAtStartPar
\sphinxstylestrong{FREQ}
\sphinxbeforeendvarwidth
\end{varwidth}%
&\begin{varwidth}[t]{\sphinxcolwidth{1}{3}}
\sphinxAtStartPar
Hz
\sphinxbeforeendvarwidth
\end{varwidth}%
&\begin{varwidth}[t]{\sphinxcolwidth{1}{3}}
\sphinxAtStartPar
频率
\sphinxbeforeendvarwidth
\end{varwidth}%
\\
\sphinxhline\begin{varwidth}[t]{\sphinxcolwidth{1}{3}}
\sphinxAtStartPar
\sphinxstylestrong{ZXXR/ZXXI}
\sphinxbeforeendvarwidth
\end{varwidth}%
&\begin{varwidth}[t]{\sphinxcolwidth{1}{3}}
\sphinxAtStartPar
Ω/km
\sphinxbeforeendvarwidth
\end{varwidth}%
&\begin{varwidth}[t]{\sphinxcolwidth{1}{3}}
\sphinxAtStartPar
阻抗Zxx实部/虚部
\sphinxbeforeendvarwidth
\end{varwidth}%
\\
\sphinxhline\begin{varwidth}[t]{\sphinxcolwidth{1}{3}}
\sphinxAtStartPar
\sphinxstylestrong{ZXYR/ZXYI}
\sphinxbeforeendvarwidth
\end{varwidth}%
&\begin{varwidth}[t]{\sphinxcolwidth{1}{3}}
\sphinxAtStartPar
Ω/km
\sphinxbeforeendvarwidth
\end{varwidth}%
&\begin{varwidth}[t]{\sphinxcolwidth{1}{3}}
\sphinxAtStartPar
阻抗Zxy实部/虚部
\sphinxbeforeendvarwidth
\end{varwidth}%
\\
\sphinxhline\begin{varwidth}[t]{\sphinxcolwidth{1}{3}}
\sphinxAtStartPar
\sphinxstylestrong{ZYXR/ZYXI}
\sphinxbeforeendvarwidth
\end{varwidth}%
&\begin{varwidth}[t]{\sphinxcolwidth{1}{3}}
\sphinxAtStartPar
Ω/km
\sphinxbeforeendvarwidth
\end{varwidth}%
&\begin{varwidth}[t]{\sphinxcolwidth{1}{3}}
\sphinxAtStartPar
阻抗Zyx实部/虚部
\sphinxbeforeendvarwidth
\end{varwidth}%
\\
\sphinxhline\begin{varwidth}[t]{\sphinxcolwidth{1}{3}}
\sphinxAtStartPar
\sphinxstylestrong{ZYYR/ZYYI}
\sphinxbeforeendvarwidth
\end{varwidth}%
&\begin{varwidth}[t]{\sphinxcolwidth{1}{3}}
\sphinxAtStartPar
Ω/km
\sphinxbeforeendvarwidth
\end{varwidth}%
&\begin{varwidth}[t]{\sphinxcolwidth{1}{3}}
\sphinxAtStartPar
阻抗Zyy实部/虚部
\sphinxbeforeendvarwidth
\end{varwidth}%
\\
\sphinxhline\begin{varwidth}[t]{\sphinxcolwidth{1}{3}}
\sphinxAtStartPar
\sphinxstylestrong{TXR/TXI}
\sphinxbeforeendvarwidth
\end{varwidth}%
&\begin{varwidth}[t]{\sphinxcolwidth{1}{3}}
\sphinxAtStartPar
无量纲
\sphinxbeforeendvarwidth
\end{varwidth}%
&\begin{varwidth}[t]{\sphinxcolwidth{1}{3}}
\sphinxAtStartPar
倾子Tx实部/虚部
\sphinxbeforeendvarwidth
\end{varwidth}%
\\
\sphinxhline\begin{varwidth}[t]{\sphinxcolwidth{1}{3}}
\sphinxAtStartPar
\sphinxstylestrong{TYR/TYI}
\sphinxbeforeendvarwidth
\end{varwidth}%
&\begin{varwidth}[t]{\sphinxcolwidth{1}{3}}
\sphinxAtStartPar
无量纲
\sphinxbeforeendvarwidth
\end{varwidth}%
&\begin{varwidth}[t]{\sphinxcolwidth{1}{3}}
\sphinxAtStartPar
倾子Ty实部/虚部
\sphinxbeforeendvarwidth
\end{varwidth}%
\\
\sphinxhline\begin{varwidth}[t]{\sphinxcolwidth{1}{3}}
\sphinxAtStartPar
\sphinxstylestrong{ZXX.VAR}
\sphinxbeforeendvarwidth
\end{varwidth}%
&\begin{varwidth}[t]{\sphinxcolwidth{1}{3}}
\sphinxAtStartPar
\sphinxhyphen{}
\sphinxbeforeendvarwidth
\end{varwidth}%
&\begin{varwidth}[t]{\sphinxcolwidth{1}{3}}
\sphinxAtStartPar
Zxx方差
\sphinxbeforeendvarwidth
\end{varwidth}%
\\
\sphinxhline\begin{varwidth}[t]{\sphinxcolwidth{1}{3}}
\sphinxAtStartPar
\sphinxstylestrong{ZXY.VAR}
\sphinxbeforeendvarwidth
\end{varwidth}%
&\begin{varwidth}[t]{\sphinxcolwidth{1}{3}}
\sphinxAtStartPar
\sphinxhyphen{}
\sphinxbeforeendvarwidth
\end{varwidth}%
&\begin{varwidth}[t]{\sphinxcolwidth{1}{3}}
\sphinxAtStartPar
Zxy方差
\sphinxbeforeendvarwidth
\end{varwidth}%
\\
\sphinxhline\begin{varwidth}[t]{\sphinxcolwidth{1}{3}}
\sphinxAtStartPar
...
\sphinxbeforeendvarwidth
\end{varwidth}%
&\begin{varwidth}[t]{\sphinxcolwidth{1}{3}}
\sphinxAtStartPar
...
\sphinxbeforeendvarwidth
\end{varwidth}%
&\begin{varwidth}[t]{\sphinxcolwidth{1}{3}}
\sphinxAtStartPar
...
\sphinxbeforeendvarwidth
\end{varwidth}%
\\
\sphinxbottomrule
\end{tabulary}
\sphinxtableafterendhook\par
\sphinxattableend\end{savenotes}


\paragraph{EDI版本兼容性}
\label{\detokenize{chapters/chapter7:id17}}

\begin{savenotes}\sphinxattablestart
\sphinxthistablewithglobalstyle
\centering
\begin{tabulary}{\linewidth}[t]{TTT}
\sphinxtoprule
\begin{varwidth}[t]{\sphinxcolwidth{1}{3}}
\sphinxstyletheadfamily \sphinxAtStartPar
版本
\sphinxbeforeendvarwidth
\end{varwidth}%
&\begin{varwidth}[t]{\sphinxcolwidth{1}{3}}
\sphinxstyletheadfamily \sphinxAtStartPar
特点
\sphinxbeforeendvarwidth
\end{varwidth}%
&\begin{varwidth}[t]{\sphinxcolwidth{1}{3}}
\sphinxstyletheadfamily \sphinxAtStartPar
兼容软件
\sphinxbeforeendvarwidth
\end{varwidth}%
\\
\sphinxmidrule
\sphinxtableatstartofbodyhook\begin{varwidth}[t]{\sphinxcolwidth{1}{3}}
\sphinxAtStartPar
\sphinxstylestrong{SpeEDI}
\sphinxbeforeendvarwidth
\end{varwidth}%
&\begin{varwidth}[t]{\sphinxcolwidth{1}{3}}
\sphinxAtStartPar
标准格式，广泛支持
\sphinxbeforeendvarwidth
\end{varwidth}%
&\begin{varwidth}[t]{\sphinxcolwidth{1}{3}}
\sphinxAtStartPar
MT\sphinxhyphen{}Pioneer, WinGLink
\sphinxbeforeendvarwidth
\end{varwidth}%
\\
\sphinxhline\begin{varwidth}[t]{\sphinxcolwidth{1}{3}}
\sphinxAtStartPar
\sphinxstylestrong{ZTEDI}
\sphinxbeforeendvarwidth
\end{varwidth}%
&\begin{varwidth}[t]{\sphinxcolwidth{1}{3}}
\sphinxAtStartPar
支持张量数据
\sphinxbeforeendvarwidth
\end{varwidth}%
&\begin{varwidth}[t]{\sphinxcolwidth{1}{3}}
\sphinxAtStartPar
Zonge软件
\sphinxbeforeendvarwidth
\end{varwidth}%
\\
\sphinxhline\begin{varwidth}[t]{\sphinxcolwidth{1}{3}}
\sphinxAtStartPar
\sphinxstylestrong{MTpyEDI}
\sphinxbeforeendvarwidth
\end{varwidth}%
&\begin{varwidth}[t]{\sphinxcolwidth{1}{3}}
\sphinxAtStartPar
Python兼容格式
\sphinxbeforeendvarwidth
\end{varwidth}%
&\begin{varwidth}[t]{\sphinxcolwidth{1}{3}}
\sphinxAtStartPar
MTpy, SimPEG
\sphinxbeforeendvarwidth
\end{varwidth}%
\\
\sphinxhline\begin{varwidth}[t]{\sphinxcolwidth{1}{3}}
\sphinxAtStartPar
\sphinxstylestrong{SEG\sphinxhyphen{}EDI}
\sphinxbeforeendvarwidth
\end{varwidth}%
&\begin{varwidth}[t]{\sphinxcolwidth{1}{3}}
\sphinxAtStartPar
SEG标准格式
\sphinxbeforeendvarwidth
\end{varwidth}%
&\begin{varwidth}[t]{\sphinxcolwidth{1}{3}}
\sphinxAtStartPar
大多数商业软件
\sphinxbeforeendvarwidth
\end{varwidth}%
\\
\sphinxbottomrule
\end{tabulary}
\sphinxtableafterendhook\par
\sphinxattableend\end{savenotes}


\subsubsection{7.2.5 EDI格式验证}
\label{\detokenize{chapters/chapter7:id18}}
\sphinxAtStartPar
\sphinxstylestrong{验证清单：}

\sphinxAtStartPar
\sphinxstylestrong{数据合理性检查：}


\begin{savenotes}\sphinxattablestart
\sphinxthistablewithglobalstyle
\centering
\begin{tabulary}{\linewidth}[t]{TTT}
\sphinxtoprule
\begin{varwidth}[t]{\sphinxcolwidth{1}{3}}
\sphinxstyletheadfamily \sphinxAtStartPar
检查项
\sphinxbeforeendvarwidth
\end{varwidth}%
&\begin{varwidth}[t]{\sphinxcolwidth{1}{3}}
\sphinxstyletheadfamily \sphinxAtStartPar
合理范围
\sphinxbeforeendvarwidth
\end{varwidth}%
&\begin{varwidth}[t]{\sphinxcolwidth{1}{3}}
\sphinxstyletheadfamily \sphinxAtStartPar
异常处理
\sphinxbeforeendvarwidth
\end{varwidth}%
\\
\sphinxmidrule
\sphinxtableatstartofbodyhook\begin{varwidth}[t]{\sphinxcolwidth{1}{3}}
\sphinxAtStartPar
\sphinxstylestrong{视电阻率}
\sphinxbeforeendvarwidth
\end{varwidth}%
&\begin{varwidth}[t]{\sphinxcolwidth{1}{3}}
\sphinxAtStartPar
0.01 \sphinxhyphen{} 100000 Ω·m
\sphinxbeforeendvarwidth
\end{varwidth}%
&\begin{varwidth}[t]{\sphinxcolwidth{1}{3}}
\sphinxAtStartPar
检查标定
\sphinxbeforeendvarwidth
\end{varwidth}%
\\
\sphinxhline\begin{varwidth}[t]{\sphinxcolwidth{1}{3}}
\sphinxAtStartPar
\sphinxstylestrong{相位}
\sphinxbeforeendvarwidth
\end{varwidth}%
&\begin{varwidth}[t]{\sphinxcolwidth{1}{3}}
\sphinxAtStartPar
0° \sphinxhyphen{} 90°
\sphinxbeforeendvarwidth
\end{varwidth}%
&\begin{varwidth}[t]{\sphinxcolwidth{1}{3}}
\sphinxAtStartPar
检查象限
\sphinxbeforeendvarwidth
\end{varwidth}%
\\
\sphinxhline\begin{varwidth}[t]{\sphinxcolwidth{1}{3}}
\sphinxAtStartPar
\sphinxstylestrong{方差}
\sphinxbeforeendvarwidth
\end{varwidth}%
&\begin{varwidth}[t]{\sphinxcolwidth{1}{3}}
\sphinxAtStartPar
< 50\%
\sphinxbeforeendvarwidth
\end{varwidth}%
&\begin{varwidth}[t]{\sphinxcolwidth{1}{3}}
\sphinxAtStartPar
增加叠加
\sphinxbeforeendvarwidth
\end{varwidth}%
\\
\sphinxhline\begin{varwidth}[t]{\sphinxcolwidth{1}{3}}
\sphinxAtStartPar
\sphinxstylestrong{频率间隔}
\sphinxbeforeendvarwidth
\end{varwidth}%
&\begin{varwidth}[t]{\sphinxcolwidth{1}{3}}
\sphinxAtStartPar
对数均匀
\sphinxbeforeendvarwidth
\end{varwidth}%
&\begin{varwidth}[t]{\sphinxcolwidth{1}{3}}
\sphinxAtStartPar
检查FFT设置
\sphinxbeforeendvarwidth
\end{varwidth}%
\\
\sphinxbottomrule
\end{tabulary}
\sphinxtableafterendhook\par
\sphinxattableend\end{savenotes}


\subsubsection{7.2.6 其他导出格式}
\label{\detokenize{chapters/chapter7:id19}}

\paragraph{PLT格式}
\label{\detokenize{chapters/chapter7:plt}}
\sphinxAtStartPar
PLT格式用于绘图软件。

\sphinxAtStartPar
\sphinxstylestrong{内容：}
\begin{itemize}
\item {} 
\sphinxAtStartPar
视电阻率曲线数据

\item {} 
\sphinxAtStartPar
相位曲线数据

\item {} 
\sphinxAtStartPar
坐标信息

\end{itemize}

\sphinxAtStartPar
\sphinxstylestrong{导出步骤：}
\begin{enumerate}
\sphinxsetlistlabels{\arabic}{enumi}{enumii}{}{.}%
\item {} 
\sphinxAtStartPar
右键测点 → \sphinxcode{\sphinxupquote{导出PLT}}

\item {} 
\sphinxAtStartPar
选择保存位置

\end{enumerate}

\sphinxAtStartPar
\sphinxstylestrong{PLT文件示例：}

\begin{sphinxVerbatim}[commandchars=\\\{\}]
\PYGZsh{} MTDP PLT File
\PYGZsh{} Site: SITE001
\PYGZsh{} Coordinate: 30.2567, 120.5678
\PYGZsh{} Frequency(Hz) Rhoxy(Ohm\PYGZhy{}m) Rhoyx(Ohm\PYGZhy{}m) Phasexy(deg) Phaseyx(deg)
100.0  15.23  18.45  42.5  47.3
80.0   16.78  19.23  43.2  46.8
...
\end{sphinxVerbatim}


\paragraph{CSV格式}
\label{\detokenize{chapters/chapter7:csv}}
\sphinxAtStartPar
CSV格式便于数据处理和分析。

\sphinxAtStartPar
\sphinxstylestrong{导出步骤：}
\begin{enumerate}
\sphinxsetlistlabels{\arabic}{enumi}{enumii}{}{.}%
\item {} 
\sphinxAtStartPar
在曲线窗口右键 → \sphinxcode{\sphinxupquote{导出数据}}

\item {} 
\sphinxAtStartPar
选择 CSV 格式

\item {} 
\sphinxAtStartPar
设置分隔符（逗号/制表符）

\end{enumerate}

\sphinxAtStartPar
\sphinxstylestrong{CSV文件结构：}

\begin{sphinxVerbatim}[commandchars=\\\{\}]
Frequency,ExEy,HyHx,Rhoxy,Rhoyx,Phasexy,Phaseyx,Error\PYGZus{}xy,Error\PYGZus{}yx
100.0,0.1234,0.0567,15.23,18.45,42.5,47.3,0.05,0.06
80.0,0.1456,0.0678,16.78,19.23,43.2,46.8,0.04,0.05
...
\end{sphinxVerbatim}


\bigskip\hrule\bigskip



\subsection{7.3 地理信息导出}
\label{\detokenize{chapters/chapter7:id20}}

\subsubsection{7.3.1 KML格式导出}
\label{\detokenize{chapters/chapter7:kml}}
\sphinxAtStartPar
\sphinxstylestrong{导出级别：}


\begin{savenotes}\sphinxattablestart
\sphinxthistablewithglobalstyle
\centering
\begin{tabulary}{\linewidth}[t]{TT}
\sphinxtoprule
\begin{varwidth}[t]{\sphinxcolwidth{1}{2}}
\sphinxstyletheadfamily \sphinxAtStartPar
级别
\sphinxbeforeendvarwidth
\end{varwidth}%
&\begin{varwidth}[t]{\sphinxcolwidth{1}{2}}
\sphinxstyletheadfamily \sphinxAtStartPar
菜单选项
\sphinxbeforeendvarwidth
\end{varwidth}%
\\
\sphinxmidrule
\sphinxtableatstartofbodyhook\begin{varwidth}[t]{\sphinxcolwidth{1}{2}}
\sphinxAtStartPar
工程
\sphinxbeforeendvarwidth
\end{varwidth}%
&\begin{varwidth}[t]{\sphinxcolwidth{1}{2}}
\sphinxAtStartPar
\sphinxcode{\sphinxupquote{工程 → 导出KML}}
\sphinxbeforeendvarwidth
\end{varwidth}%
\\
\sphinxhline\begin{varwidth}[t]{\sphinxcolwidth{1}{2}}
\sphinxAtStartPar
工区
\sphinxbeforeendvarwidth
\end{varwidth}%
&\begin{varwidth}[t]{\sphinxcolwidth{1}{2}}
\sphinxAtStartPar
\sphinxcode{\sphinxupquote{工区 → 导出普通KML}}
\sphinxbeforeendvarwidth
\end{varwidth}%
\\
\sphinxhline\begin{varwidth}[t]{\sphinxcolwidth{1}{2}}
\sphinxAtStartPar
测段
\sphinxbeforeendvarwidth
\end{varwidth}%
&\begin{varwidth}[t]{\sphinxcolwidth{1}{2}}
\sphinxAtStartPar
\sphinxcode{\sphinxupquote{测段 → 导出普通KML}}
\sphinxbeforeendvarwidth
\end{varwidth}%
\\
\sphinxbottomrule
\end{tabulary}
\sphinxtableafterendhook\par
\sphinxattableend\end{savenotes}

\sphinxAtStartPar
\sphinxstylestrong{导出步骤：}
\begin{enumerate}
\sphinxsetlistlabels{\arabic}{enumi}{enumii}{}{.}%
\item {} 
\sphinxAtStartPar
选择要导出的工程/工区/测段

\item {} 
\sphinxAtStartPar
右键选择相应的KML导出选项

\item {} 
\sphinxAtStartPar
设置保存位置

\item {} 
\sphinxAtStartPar
用Google Earth打开查看

\end{enumerate}


\subsubsection{7.3.2 KMZ格式导出}
\label{\detokenize{chapters/chapter7:kmz}}
\sphinxAtStartPar
KMZ是KML的压缩格式：
\begin{enumerate}
\sphinxsetlistlabels{\arabic}{enumi}{enumii}{}{.}%
\item {} 
\sphinxAtStartPar
选择工区或测段

\item {} 
\sphinxAtStartPar
右键选择 \sphinxcode{\sphinxupquote{导出KMZ}}

\item {} 
\sphinxAtStartPar
选择保存位置

\end{enumerate}


\subsubsection{7.3.3 KML文件结构}
\label{\detokenize{chapters/chapter7:id21}}\begin{itemize}
\item {} 
\sphinxAtStartPar
\sphinxstylestrong{Document}：文档容器

\item {} 
\sphinxAtStartPar
\sphinxstylestrong{Folder}：文件夹组织

\item {} 
\sphinxAtStartPar
\sphinxstylestrong{Placemark}：地标（测点）

\item {} 
\sphinxAtStartPar
\sphinxstylestrong{Point}：点坐标

\item {} 
\sphinxAtStartPar
\sphinxstylestrong{Style}：样式定义（IconStyle、LineStyle、BalloonStyle）

\item {} 
\sphinxAtStartPar
\sphinxstylestrong{LookAt}：视角设置

\item {} 
\sphinxAtStartPar
\sphinxstylestrong{ExtendedData}：扩展数据

\end{itemize}


\bigskip\hrule\bigskip



\subsection{7.4 坐标数据导出}
\label{\detokenize{chapters/chapter7:id22}}\begin{enumerate}
\sphinxsetlistlabels{\arabic}{enumi}{enumii}{}{.}%
\item {} 
\sphinxAtStartPar
选择工区

\item {} 
\sphinxAtStartPar
右键选择 \sphinxcode{\sphinxupquote{导出测点坐标}}

\item {} 
\sphinxAtStartPar
选择导出格式（TXT/CSV）

\end{enumerate}

\sphinxAtStartPar
导出内容：
\begin{itemize}
\item {} 
\sphinxAtStartPar
测点名称

\item {} 
\sphinxAtStartPar
经度

\item {} 
\sphinxAtStartPar
纬度

\item {} 
\sphinxAtStartPar
高程

\end{itemize}


\bigskip\hrule\bigskip



\subsection{7.5 批量导出功能}
\label{\detokenize{chapters/chapter7:id23}}

\subsubsection{7.5.1 导出测段参数}
\label{\detokenize{chapters/chapter7:id24}}
\sphinxAtStartPar
将测段内所有测点的参数导出到单个文件：
\begin{enumerate}
\sphinxsetlistlabels{\arabic}{enumi}{enumii}{}{.}%
\item {} 
\sphinxAtStartPar
选择测段

\item {} 
\sphinxAtStartPar
右键选择 \sphinxcode{\sphinxupquote{导出参数}}

\item {} 
\sphinxAtStartPar
选择导出内容

\end{enumerate}

\sphinxAtStartPar
\sphinxstylestrong{导出文件格式：}
\begin{itemize}
\item {} 
\sphinxAtStartPar
文件名：\sphinxcode{\sphinxupquote{测段名\sphinxhyphen{}Paint.dat}}

\item {} 
\sphinxAtStartPar
内容：测点名、坐标、频率、所有MT参数

\end{itemize}


\subsubsection{7.5.2 导出相位张量}
\label{\detokenize{chapters/chapter7:id25}}
\sphinxAtStartPar
按频率导出所有测点的相位张量参数：
\begin{enumerate}
\sphinxsetlistlabels{\arabic}{enumi}{enumii}{}{.}%
\item {} 
\sphinxAtStartPar
选择测段

\item {} 
\sphinxAtStartPar
右键选择 \sphinxcode{\sphinxupquote{导出相位张量}}

\item {} 
\sphinxAtStartPar
系统为每个频率生成一个文件

\end{enumerate}

\sphinxAtStartPar
\sphinxstylestrong{导出文件格式：}
\begin{itemize}
\item {} 
\sphinxAtStartPar
文件名：\sphinxcode{\sphinxupquote{频率\sphinxhyphen{}PhaseTensor.dat}}

\item {} 
\sphinxAtStartPar
内容：测点名、坐标、alpha、beta、PtMax、PtMin、Skew1D、Skew2D等

\end{itemize}


\subsubsection{7.5.3 导出数据质量}
\label{\detokenize{chapters/chapter7:id26}}\begin{enumerate}
\sphinxsetlistlabels{\arabic}{enumi}{enumii}{}{.}%
\item {} 
\sphinxAtStartPar
选择测段

\item {} 
\sphinxAtStartPar
右键选择 \sphinxcode{\sphinxupquote{导出数据质量}}

\item {} 
\sphinxAtStartPar
生成质量等级文件

\end{enumerate}

\sphinxAtStartPar
\sphinxstylestrong{导出文件格式：}
\begin{itemize}
\item {} 
\sphinxAtStartPar
文件名：\sphinxcode{\sphinxupquote{测段名\sphinxhyphen{}DataQuality.dat}}

\item {} 
\sphinxAtStartPar
内容：测点名、质量等级

\end{itemize}


\subsubsection{7.5.4 导出测点坐标}
\label{\detokenize{chapters/chapter7:id27}}\begin{enumerate}
\sphinxsetlistlabels{\arabic}{enumi}{enumii}{}{.}%
\item {} 
\sphinxAtStartPar
选择测段

\item {} 
\sphinxAtStartPar
右键选择 \sphinxcode{\sphinxupquote{导出坐标}}

\item {} 
\sphinxAtStartPar
生成坐标文件

\end{enumerate}

\sphinxAtStartPar
\sphinxstylestrong{导出文件格式：}
\begin{itemize}
\item {} 
\sphinxAtStartPar
文件名：\sphinxcode{\sphinxupquote{测段名\sphinxhyphen{}Coor.dat}}

\item {} 
\sphinxAtStartPar
内容：测点名、经度、纬度、高程

\end{itemize}


\bigskip\hrule\bigskip



\subsection{7.6 其他导出格式}
\label{\detokenize{chapters/chapter7:id28}}

\subsubsection{7.6.1 视电阻率/相位数据导出}
\label{\detokenize{chapters/chapter7:id29}}\begin{enumerate}
\sphinxsetlistlabels{\arabic}{enumi}{enumii}{}{.}%
\item {} 
\sphinxAtStartPar
选择测点

\item {} 
\sphinxAtStartPar
在曲线窗口中右键选择"导出数据"

\item {} 
\sphinxAtStartPar
选择格式（CSV、TXT）

\end{enumerate}


\subsubsection{7.6.2 图表导出}
\label{\detokenize{chapters/chapter7:id30}}

\begin{savenotes}\sphinxattablestart
\sphinxthistablewithglobalstyle
\centering
\begin{tabulary}{\linewidth}[t]{TT}
\sphinxtoprule
\begin{varwidth}[t]{\sphinxcolwidth{1}{2}}
\sphinxstyletheadfamily \sphinxAtStartPar
格式
\sphinxbeforeendvarwidth
\end{varwidth}%
&\begin{varwidth}[t]{\sphinxcolwidth{1}{2}}
\sphinxstyletheadfamily \sphinxAtStartPar
说明
\sphinxbeforeendvarwidth
\end{varwidth}%
\\
\sphinxmidrule
\sphinxtableatstartofbodyhook\begin{varwidth}[t]{\sphinxcolwidth{1}{2}}
\sphinxAtStartPar
EMF/WMF
\sphinxbeforeendvarwidth
\end{varwidth}%
&\begin{varwidth}[t]{\sphinxcolwidth{1}{2}}
\sphinxAtStartPar
矢量图（Windows）
\sphinxbeforeendvarwidth
\end{varwidth}%
\\
\sphinxhline\begin{varwidth}[t]{\sphinxcolwidth{1}{2}}
\sphinxAtStartPar
PNG
\sphinxbeforeendvarwidth
\end{varwidth}%
&\begin{varwidth}[t]{\sphinxcolwidth{1}{2}}
\sphinxAtStartPar
位图，透明背景
\sphinxbeforeendvarwidth
\end{varwidth}%
\\
\sphinxhline\begin{varwidth}[t]{\sphinxcolwidth{1}{2}}
\sphinxAtStartPar
BMP
\sphinxbeforeendvarwidth
\end{varwidth}%
&\begin{varwidth}[t]{\sphinxcolwidth{1}{2}}
\sphinxAtStartPar
位图，无压缩
\sphinxbeforeendvarwidth
\end{varwidth}%
\\
\sphinxhline\begin{varwidth}[t]{\sphinxcolwidth{1}{2}}
\sphinxAtStartPar
JPG
\sphinxbeforeendvarwidth
\end{varwidth}%
&\begin{varwidth}[t]{\sphinxcolwidth{1}{2}}
\sphinxAtStartPar
位图，有损压缩
\sphinxbeforeendvarwidth
\end{varwidth}%
\\
\sphinxbottomrule
\end{tabulary}
\sphinxtableafterendhook\par
\sphinxattableend\end{savenotes}


\bigskip\hrule\bigskip



\subsection{7.7 数据质量控制}
\label{\detokenize{chapters/chapter7:id31}}

\subsubsection{7.7.1 数据质量标记}
\label{\detokenize{chapters/chapter7:id32}}

\begin{savenotes}\sphinxattablestart
\sphinxthistablewithglobalstyle
\centering
\begin{tabulary}{\linewidth}[t]{TT}
\sphinxtoprule
\begin{varwidth}[t]{\sphinxcolwidth{1}{2}}
\sphinxstyletheadfamily \sphinxAtStartPar
等级
\sphinxbeforeendvarwidth
\end{varwidth}%
&\begin{varwidth}[t]{\sphinxcolwidth{1}{2}}
\sphinxstyletheadfamily \sphinxAtStartPar
说明
\sphinxbeforeendvarwidth
\end{varwidth}%
\\
\sphinxmidrule
\sphinxtableatstartofbodyhook\begin{varwidth}[t]{\sphinxcolwidth{1}{2}}
\sphinxAtStartPar
0
\sphinxbeforeendvarwidth
\end{varwidth}%
&\begin{varwidth}[t]{\sphinxcolwidth{1}{2}}
\sphinxAtStartPar
未评估
\sphinxbeforeendvarwidth
\end{varwidth}%
\\
\sphinxhline\begin{varwidth}[t]{\sphinxcolwidth{1}{2}}
\sphinxAtStartPar
1
\sphinxbeforeendvarwidth
\end{varwidth}%
&\begin{varwidth}[t]{\sphinxcolwidth{1}{2}}
\sphinxAtStartPar
优秀
\sphinxbeforeendvarwidth
\end{varwidth}%
\\
\sphinxhline\begin{varwidth}[t]{\sphinxcolwidth{1}{2}}
\sphinxAtStartPar
2
\sphinxbeforeendvarwidth
\end{varwidth}%
&\begin{varwidth}[t]{\sphinxcolwidth{1}{2}}
\sphinxAtStartPar
良好
\sphinxbeforeendvarwidth
\end{varwidth}%
\\
\sphinxhline\begin{varwidth}[t]{\sphinxcolwidth{1}{2}}
\sphinxAtStartPar
3
\sphinxbeforeendvarwidth
\end{varwidth}%
&\begin{varwidth}[t]{\sphinxcolwidth{1}{2}}
\sphinxAtStartPar
一般
\sphinxbeforeendvarwidth
\end{varwidth}%
\\
\sphinxhline\begin{varwidth}[t]{\sphinxcolwidth{1}{2}}
\sphinxAtStartPar
4
\sphinxbeforeendvarwidth
\end{varwidth}%
&\begin{varwidth}[t]{\sphinxcolwidth{1}{2}}
\sphinxAtStartPar
较差
\sphinxbeforeendvarwidth
\end{varwidth}%
\\
\sphinxbottomrule
\end{tabulary}
\sphinxtableafterendhook\par
\sphinxattableend\end{savenotes}


\subsubsection{7.7.2 批量质量导出}
\label{\detokenize{chapters/chapter7:id33}}\begin{enumerate}
\sphinxsetlistlabels{\arabic}{enumi}{enumii}{}{.}%
\item {} 
\sphinxAtStartPar
选择测段

\item {} 
\sphinxAtStartPar
右键选择 \sphinxcode{\sphinxupquote{导出数据质量}}

\item {} 
\sphinxAtStartPar
导出所有测点的质量信息

\end{enumerate}


\bigskip\hrule\bigskip



\subsection{7.8 与GIS软件联动}
\label{\detokenize{chapters/chapter7:gis}}
\sphinxAtStartPar
MTDP导出的地理数据可与以下软件联动：
\begin{itemize}
\item {} 
\sphinxAtStartPar
\sphinxstylestrong{Google Earth}：查看测点空间分布

\item {} 
\sphinxAtStartPar
\sphinxstylestrong{QGIS}：进行空间分析和制图

\item {} 
\sphinxAtStartPar
\sphinxstylestrong{ArcGIS}：专业GIS处理

\item {} 
\sphinxAtStartPar
\sphinxstylestrong{Global Mapper}：多格式数据集成

\end{itemize}
\begin{quote}

\sphinxAtStartPar
\sphinxstylestrong{提示}：导出KML/KMZ前，请确保测点的经纬度坐标已正确设置。
\end{quote}


\bigskip\hrule\bigskip



\subsection{7.9 工程文件备份}
\label{\detokenize{chapters/chapter7:id34}}

\subsubsection{7.9.1 导出版本}
\label{\detokenize{chapters/chapter7:id35}}\begin{itemize}
\item {} 
\sphinxAtStartPar
选择 \sphinxcode{\sphinxupquote{文件 → 导出版本}}

\item {} 
\sphinxAtStartPar
创建工程的完整副本

\end{itemize}


\subsubsection{7.9.2 自动备份}
\label{\detokenize{chapters/chapter7:id36}}\begin{itemize}
\item {} 
\sphinxAtStartPar
每次保存时自动创建备份

\item {} 
\sphinxAtStartPar
保留最近10个备份版本

\item {} 
\sphinxAtStartPar
备份文件存储在工程目录下

\end{itemize}


\bigskip\hrule\bigskip



\subsection{7.10 时间序列导出}
\label{\detokenize{chapters/chapter7:id37}}

\subsubsection{7.10.1 GMT格式导出}
\label{\detokenize{chapters/chapter7:gmt}}\begin{enumerate}
\sphinxsetlistlabels{\arabic}{enumi}{enumii}{}{.}%
\item {} 
\sphinxAtStartPar
选择测点

\item {} 
\sphinxAtStartPar
选择 \sphinxcode{\sphinxupquote{导出 → GMT时间序列}}

\item {} 
\sphinxAtStartPar
时间序列导出为GMT兼容格式

\end{enumerate}


\subsubsection{7.10.2 ATGMT格式导出}
\label{\detokenize{chapters/chapter7:atgmt}}\begin{enumerate}
\sphinxsetlistlabels{\arabic}{enumi}{enumii}{}{.}%
\item {} 
\sphinxAtStartPar
选择测点

\item {} 
\sphinxAtStartPar
选择 \sphinxcode{\sphinxupquote{工具 → 时间序列 → ATGMT格式}}

\item {} 
\sphinxAtStartPar
配置导出参数

\end{enumerate}


\bigskip\hrule\bigskip



\subsection{7.11 HTML报告生成}
\label{\detokenize{chapters/chapter7:html}}\begin{enumerate}
\sphinxsetlistlabels{\arabic}{enumi}{enumii}{}{.}%
\item {} 
\sphinxAtStartPar
选择测点

\item {} 
\sphinxAtStartPar
选择 \sphinxcode{\sphinxupquote{测点 → 导出HTML设置}}

\item {} 
\sphinxAtStartPar
生成测点配置的HTML报告

\end{enumerate}

\sphinxstepscope


\section{🧮 第8章 工具集}
\label{\detokenize{chapters/chapter8:id1}}\label{\detokenize{chapters/chapter8::doc}}
\sphinxAtStartPar
本章介绍MTDP提供的各种实用工具。


\bigskip\hrule\bigskip



\subsection{8.1 🛠️ 工具菜单完整列表}
\label{\detokenize{chapters/chapter8:id2}}
\sphinxAtStartPar
MTDataPro提供丰富的工具菜单，支持时间序列处理、格式转换、数据修复等功能。


\subsubsection{8.1.1 标定与查看工具}
\label{\detokenize{chapters/chapter8:id3}}

\begin{savenotes}\sphinxattablestart
\sphinxthistablewithglobalstyle
\centering
\begin{tabulary}{\linewidth}[t]{TT}
\sphinxtoprule
\begin{varwidth}[t]{\sphinxcolwidth{1}{2}}
\sphinxstyletheadfamily \sphinxAtStartPar
菜单项
\sphinxbeforeendvarwidth
\end{varwidth}%
&\begin{varwidth}[t]{\sphinxcolwidth{1}{2}}
\sphinxstyletheadfamily \sphinxAtStartPar
功能说明
\sphinxbeforeendvarwidth
\end{varwidth}%
\\
\sphinxmidrule
\sphinxtableatstartofbodyhook\begin{varwidth}[t]{\sphinxcolwidth{1}{2}}
\sphinxAtStartPar
标定查看 (CalViewMenu)
\sphinxbeforeendvarwidth
\end{varwidth}%
&\begin{varwidth}[t]{\sphinxcolwidth{1}{2}}
\sphinxAtStartPar
查看和校准仪器标定数据
\sphinxbeforeendvarwidth
\end{varwidth}%
\\
\sphinxhline\begin{varwidth}[t]{\sphinxcolwidth{1}{2}}
\sphinxAtStartPar
TS查看 (TSViewMenu)
\sphinxbeforeendvarwidth
\end{varwidth}%
&\begin{varwidth}[t]{\sphinxcolwidth{1}{2}}
\sphinxAtStartPar
时间序列数据查看器
\sphinxbeforeendvarwidth
\end{varwidth}%
\\
\sphinxbottomrule
\end{tabulary}
\sphinxtableafterendhook\par
\sphinxattableend\end{savenotes}


\subsubsection{8.1.2 Phoenix时间序列工具}
\label{\detokenize{chapters/chapter8:phoenix}}

\begin{savenotes}\sphinxattablestart
\sphinxthistablewithglobalstyle
\centering
\begin{tabulary}{\linewidth}[t]{TT}
\sphinxtoprule
\begin{varwidth}[t]{\sphinxcolwidth{1}{2}}
\sphinxstyletheadfamily \sphinxAtStartPar
菜单项
\sphinxbeforeendvarwidth
\end{varwidth}%
&\begin{varwidth}[t]{\sphinxcolwidth{1}{2}}
\sphinxstyletheadfamily \sphinxAtStartPar
功能说明
\sphinxbeforeendvarwidth
\end{varwidth}%
\\
\sphinxmidrule
\sphinxtableatstartofbodyhook\begin{varwidth}[t]{\sphinxcolwidth{1}{2}}
\sphinxAtStartPar
Phoenix TS分割 (SplitPhoenixTSMenu)
\sphinxbeforeendvarwidth
\end{varwidth}%
&\begin{varwidth}[t]{\sphinxcolwidth{1}{2}}
\sphinxAtStartPar
分割Phoenix格式时间序列
\sphinxbeforeendvarwidth
\end{varwidth}%
\\
\sphinxhline\begin{varwidth}[t]{\sphinxcolwidth{1}{2}}
\sphinxAtStartPar
Phoenix TS修复 (PhoneixTSRepairMenu)
\sphinxbeforeendvarwidth
\end{varwidth}%
&\begin{varwidth}[t]{\sphinxcolwidth{1}{2}}
\sphinxAtStartPar
修复损坏的Phoenix TS文件
\sphinxbeforeendvarwidth
\end{varwidth}%
\\
\sphinxhline\begin{varwidth}[t]{\sphinxcolwidth{1}{2}}
\sphinxAtStartPar
Phoenix TS合并 (MergePhoneixTSMenu)
\sphinxbeforeendvarwidth
\end{varwidth}%
&\begin{varwidth}[t]{\sphinxcolwidth{1}{2}}
\sphinxAtStartPar
合并多个Phoenix TS文件
\sphinxbeforeendvarwidth
\end{varwidth}%
\\
\sphinxhline\begin{varwidth}[t]{\sphinxcolwidth{1}{2}}
\sphinxAtStartPar
Phoenix通道重排 (PhoenixChannelRearrenageMenu)
\sphinxbeforeendvarwidth
\end{varwidth}%
&\begin{varwidth}[t]{\sphinxcolwidth{1}{2}}
\sphinxAtStartPar
重新排列Phoenix通道顺序
\sphinxbeforeendvarwidth
\end{varwidth}%
\\
\sphinxhline\begin{varwidth}[t]{\sphinxcolwidth{1}{2}}
\sphinxAtStartPar
TBL坐标修复 (TBLCoorFixMenu)
\sphinxbeforeendvarwidth
\end{varwidth}%
&\begin{varwidth}[t]{\sphinxcolwidth{1}{2}}
\sphinxAtStartPar
修复TBL格式坐标数据
\sphinxbeforeendvarwidth
\end{varwidth}%
\\
\sphinxhline\begin{varwidth}[t]{\sphinxcolwidth{1}{2}}
\sphinxAtStartPar
TBL参数替换 (TBLParameterReplaceMenu)
\sphinxbeforeendvarwidth
\end{varwidth}%
&\begin{varwidth}[t]{\sphinxcolwidth{1}{2}}
\sphinxAtStartPar
批量替换TBL参数
\sphinxbeforeendvarwidth
\end{varwidth}%
\\
\sphinxbottomrule
\end{tabulary}
\sphinxtableafterendhook\par
\sphinxattableend\end{savenotes}


\subsubsection{8.1.3 时间序列转换工具 (TSConvertMenu)}
\label{\detokenize{chapters/chapter8:tsconvertmenu}}

\begin{savenotes}\sphinxattablestart
\sphinxthistablewithglobalstyle
\centering
\begin{tabulary}{\linewidth}[t]{TT}
\sphinxtoprule
\begin{varwidth}[t]{\sphinxcolwidth{1}{2}}
\sphinxstyletheadfamily \sphinxAtStartPar
子菜单项
\sphinxbeforeendvarwidth
\end{varwidth}%
&\begin{varwidth}[t]{\sphinxcolwidth{1}{2}}
\sphinxstyletheadfamily \sphinxAtStartPar
功能说明
\sphinxbeforeendvarwidth
\end{varwidth}%
\\
\sphinxmidrule
\sphinxtableatstartofbodyhook\begin{varwidth}[t]{\sphinxcolwidth{1}{2}}
\sphinxAtStartPar
TXT→TSN (PhoneixTxt2TSNMenu)
\sphinxbeforeendvarwidth
\end{varwidth}%
&\begin{varwidth}[t]{\sphinxcolwidth{1}{2}}
\sphinxAtStartPar
Phoenix文本格式转换为TSN
\sphinxbeforeendvarwidth
\end{varwidth}%
\\
\sphinxhline\begin{varwidth}[t]{\sphinxcolwidth{1}{2}}
\sphinxAtStartPar
TSN→DAT (PhoenixTSN2DatMenu)
\sphinxbeforeendvarwidth
\end{varwidth}%
&\begin{varwidth}[t]{\sphinxcolwidth{1}{2}}
\sphinxAtStartPar
TSN格式转换为DAT
\sphinxbeforeendvarwidth
\end{varwidth}%
\\
\sphinxhline\begin{varwidth}[t]{\sphinxcolwidth{1}{2}}
\sphinxAtStartPar
DAT→TSN (PhoenixDat2TSNMenu)
\sphinxbeforeendvarwidth
\end{varwidth}%
&\begin{varwidth}[t]{\sphinxcolwidth{1}{2}}
\sphinxAtStartPar
DAT格式转换为TSN
\sphinxbeforeendvarwidth
\end{varwidth}%
\\
\sphinxhline\begin{varwidth}[t]{\sphinxcolwidth{1}{2}}
\sphinxAtStartPar
TSN→GMT (PhoneixTSN2GMTMenu)
\sphinxbeforeendvarwidth
\end{varwidth}%
&\begin{varwidth}[t]{\sphinxcolwidth{1}{2}}
\sphinxAtStartPar
TSN格式转换为GMT
\sphinxbeforeendvarwidth
\end{varwidth}%
\\
\sphinxhline\begin{varwidth}[t]{\sphinxcolwidth{1}{2}}
\sphinxAtStartPar
时间序列→GMT (TimeSeries2GMTMenu)
\sphinxbeforeendvarwidth
\end{varwidth}%
&\begin{varwidth}[t]{\sphinxcolwidth{1}{2}}
\sphinxAtStartPar
通用时间序列转GMT
\sphinxbeforeendvarwidth
\end{varwidth}%
\\
\sphinxhline\begin{varwidth}[t]{\sphinxcolwidth{1}{2}}
\sphinxAtStartPar
时间序列→ATGMT (TimeSeries2ATGMTMenu)
\sphinxbeforeendvarwidth
\end{varwidth}%
&\begin{varwidth}[t]{\sphinxcolwidth{1}{2}}
\sphinxAtStartPar
时间序列转ATGMT格式
\sphinxbeforeendvarwidth
\end{varwidth}%
\\
\sphinxbottomrule
\end{tabulary}
\sphinxtableafterendhook\par
\sphinxattableend\end{savenotes}


\subsubsection{8.1.4 LEMI工具}
\label{\detokenize{chapters/chapter8:lemi}}

\begin{savenotes}\sphinxattablestart
\sphinxthistablewithglobalstyle
\centering
\begin{tabulary}{\linewidth}[t]{TT}
\sphinxtoprule
\begin{varwidth}[t]{\sphinxcolwidth{1}{2}}
\sphinxstyletheadfamily \sphinxAtStartPar
菜单项
\sphinxbeforeendvarwidth
\end{varwidth}%
&\begin{varwidth}[t]{\sphinxcolwidth{1}{2}}
\sphinxstyletheadfamily \sphinxAtStartPar
功能说明
\sphinxbeforeendvarwidth
\end{varwidth}%
\\
\sphinxmidrule
\sphinxtableatstartofbodyhook\begin{varwidth}[t]{\sphinxcolwidth{1}{2}}
\sphinxAtStartPar
LEMI测点合并 (LEMISiteMergeMenu)
\sphinxbeforeendvarwidth
\end{varwidth}%
&\begin{varwidth}[t]{\sphinxcolwidth{1}{2}}
\sphinxAtStartPar
合并LEMI格式测点数据
\sphinxbeforeendvarwidth
\end{varwidth}%
\\
\sphinxhline\begin{varwidth}[t]{\sphinxcolwidth{1}{2}}
\sphinxAtStartPar
LEMI H自动恢复 (LEMISiteHAutoRestoreMenu)
\sphinxbeforeendvarwidth
\end{varwidth}%
&\begin{varwidth}[t]{\sphinxcolwidth{1}{2}}
\sphinxAtStartPar
自动恢复LEMI H通道数据
\sphinxbeforeendvarwidth
\end{varwidth}%
\\
\sphinxbottomrule
\end{tabulary}
\sphinxtableafterendhook\par
\sphinxattableend\end{savenotes}


\subsubsection{8.1.5 Metronix工具}
\label{\detokenize{chapters/chapter8:metronix}}

\begin{savenotes}\sphinxattablestart
\sphinxthistablewithglobalstyle
\centering
\begin{tabulary}{\linewidth}[t]{TT}
\sphinxtoprule
\begin{varwidth}[t]{\sphinxcolwidth{1}{2}}
\sphinxstyletheadfamily \sphinxAtStartPar
菜单项
\sphinxbeforeendvarwidth
\end{varwidth}%
&\begin{varwidth}[t]{\sphinxcolwidth{1}{2}}
\sphinxstyletheadfamily \sphinxAtStartPar
功能说明
\sphinxbeforeendvarwidth
\end{varwidth}%
\\
\sphinxmidrule
\sphinxtableatstartofbodyhook\begin{varwidth}[t]{\sphinxcolwidth{1}{2}}
\sphinxAtStartPar
Metronix测点合并 (MetronixSiteMergeMen)
\sphinxbeforeendvarwidth
\end{varwidth}%
&\begin{varwidth}[t]{\sphinxcolwidth{1}{2}}
\sphinxAtStartPar
合并Metronix格式测点数据
\sphinxbeforeendvarwidth
\end{varwidth}%
\\
\sphinxhline\begin{varwidth}[t]{\sphinxcolwidth{1}{2}}
\sphinxAtStartPar
Metronix ATS→DAT (MetronixAtsToDatMenu)
\sphinxbeforeendvarwidth
\end{varwidth}%
&\begin{varwidth}[t]{\sphinxcolwidth{1}{2}}
\sphinxAtStartPar
Metronix ATS格式转换为DAT
\sphinxbeforeendvarwidth
\end{varwidth}%
\\
\sphinxbottomrule
\end{tabulary}
\sphinxtableafterendhook\par
\sphinxattableend\end{savenotes}


\subsubsection{8.1.6 WEM数据处理}
\label{\detokenize{chapters/chapter8:wem}}

\begin{savenotes}\sphinxattablestart
\sphinxthistablewithglobalstyle
\centering
\begin{tabulary}{\linewidth}[t]{TT}
\sphinxtoprule
\begin{varwidth}[t]{\sphinxcolwidth{1}{2}}
\sphinxstyletheadfamily \sphinxAtStartPar
菜单项
\sphinxbeforeendvarwidth
\end{varwidth}%
&\begin{varwidth}[t]{\sphinxcolwidth{1}{2}}
\sphinxstyletheadfamily \sphinxAtStartPar
功能说明
\sphinxbeforeendvarwidth
\end{varwidth}%
\\
\sphinxmidrule
\sphinxtableatstartofbodyhook\begin{varwidth}[t]{\sphinxcolwidth{1}{2}}
\sphinxAtStartPar
WEM数据处理 (WEMDataProcessMenu)
\sphinxbeforeendvarwidth
\end{varwidth}%
&\begin{varwidth}[t]{\sphinxcolwidth{1}{2}}
\sphinxAtStartPar
WEM格式数据处理工具
\sphinxbeforeendvarwidth
\end{varwidth}%
\\
\sphinxbottomrule
\end{tabulary}
\sphinxtableafterendhook\par
\sphinxattableend\end{savenotes}


\subsubsection{8.1.7 EDI处理工具}
\label{\detokenize{chapters/chapter8:edi}}

\begin{savenotes}\sphinxattablestart
\sphinxthistablewithglobalstyle
\centering
\begin{tabulary}{\linewidth}[t]{TT}
\sphinxtoprule
\begin{varwidth}[t]{\sphinxcolwidth{1}{2}}
\sphinxstyletheadfamily \sphinxAtStartPar
菜单项
\sphinxbeforeendvarwidth
\end{varwidth}%
&\begin{varwidth}[t]{\sphinxcolwidth{1}{2}}
\sphinxstyletheadfamily \sphinxAtStartPar
功能说明
\sphinxbeforeendvarwidth
\end{varwidth}%
\\
\sphinxmidrule
\sphinxtableatstartofbodyhook\begin{varwidth}[t]{\sphinxcolwidth{1}{2}}
\sphinxAtStartPar
EDI修复 (EDIRepairMenu)
\sphinxbeforeendvarwidth
\end{varwidth}%
&\begin{varwidth}[t]{\sphinxcolwidth{1}{2}}
\sphinxAtStartPar
修复损坏的EDI文件
\sphinxbeforeendvarwidth
\end{varwidth}%
\\
\sphinxhline\begin{varwidth}[t]{\sphinxcolwidth{1}{2}}
\sphinxAtStartPar
EDI转换相干度 (EDIConvertCOHMenu)
\sphinxbeforeendvarwidth
\end{varwidth}%
&\begin{varwidth}[t]{\sphinxcolwidth{1}{2}}
\sphinxAtStartPar
转换EDI相干度数据
\sphinxbeforeendvarwidth
\end{varwidth}%
\\
\sphinxhline\begin{varwidth}[t]{\sphinxcolwidth{1}{2}}
\sphinxAtStartPar
旧版MT→EDI (OldMT2EdiMenu)
\sphinxbeforeendvarwidth
\end{varwidth}%
&\begin{varwidth}[t]{\sphinxcolwidth{1}{2}}
\sphinxAtStartPar
旧版MT格式转换为EDI
\sphinxbeforeendvarwidth
\end{varwidth}%
\\
\sphinxbottomrule
\end{tabulary}
\sphinxtableafterendhook\par
\sphinxattableend\end{savenotes}


\subsubsection{8.1.8 Phoenix格式转换工具 (PhoenixConvertMenu)}
\label{\detokenize{chapters/chapter8:phoenix-phoenixconvertmenu}}

\begin{savenotes}\sphinxattablestart
\sphinxthistablewithglobalstyle
\centering
\begin{tabulary}{\linewidth}[t]{TT}
\sphinxtoprule
\begin{varwidth}[t]{\sphinxcolwidth{1}{2}}
\sphinxstyletheadfamily \sphinxAtStartPar
子菜单项
\sphinxbeforeendvarwidth
\end{varwidth}%
&\begin{varwidth}[t]{\sphinxcolwidth{1}{2}}
\sphinxstyletheadfamily \sphinxAtStartPar
功能说明
\sphinxbeforeendvarwidth
\end{varwidth}%
\\
\sphinxmidrule
\sphinxtableatstartofbodyhook\begin{varwidth}[t]{\sphinxcolwidth{1}{2}}
\sphinxAtStartPar
CLB→JSON (CLB2JSONMenu)
\sphinxbeforeendvarwidth
\end{varwidth}%
&\begin{varwidth}[t]{\sphinxcolwidth{1}{2}}
\sphinxAtStartPar
CLB标定文件转JSON
\sphinxbeforeendvarwidth
\end{varwidth}%
\\
\sphinxhline\begin{varwidth}[t]{\sphinxcolwidth{1}{2}}
\sphinxAtStartPar
CLC→JSON (CLC2JSONMenu)
\sphinxbeforeendvarwidth
\end{varwidth}%
&\begin{varwidth}[t]{\sphinxcolwidth{1}{2}}
\sphinxAtStartPar
CLC标定文件转JSON
\sphinxbeforeendvarwidth
\end{varwidth}%
\\
\sphinxhline\begin{varwidth}[t]{\sphinxcolwidth{1}{2}}
\sphinxAtStartPar
JSON→CLB (JSON2CLBMenu)
\sphinxbeforeendvarwidth
\end{varwidth}%
&\begin{varwidth}[t]{\sphinxcolwidth{1}{2}}
\sphinxAtStartPar
JSON转CLB标定文件
\sphinxbeforeendvarwidth
\end{varwidth}%
\\
\sphinxhline\begin{varwidth}[t]{\sphinxcolwidth{1}{2}}
\sphinxAtStartPar
JSON→CLC (JSON2CLCMenu)
\sphinxbeforeendvarwidth
\end{varwidth}%
&\begin{varwidth}[t]{\sphinxcolwidth{1}{2}}
\sphinxAtStartPar
JSON转CLC标定文件
\sphinxbeforeendvarwidth
\end{varwidth}%
\\
\sphinxhline\begin{varwidth}[t]{\sphinxcolwidth{1}{2}}
\sphinxAtStartPar
TBL→JSON (TBL2JSONMenu)
\sphinxbeforeendvarwidth
\end{varwidth}%
&\begin{varwidth}[t]{\sphinxcolwidth{1}{2}}
\sphinxAtStartPar
TBL参数文件转JSON
\sphinxbeforeendvarwidth
\end{varwidth}%
\\
\sphinxhline\begin{varwidth}[t]{\sphinxcolwidth{1}{2}}
\sphinxAtStartPar
JSON→TBL (JSON2TBLMenu)
\sphinxbeforeendvarwidth
\end{varwidth}%
&\begin{varwidth}[t]{\sphinxcolwidth{1}{2}}
\sphinxAtStartPar
JSON转TBL参数文件
\sphinxbeforeendvarwidth
\end{varwidth}%
\\
\sphinxbottomrule
\end{tabulary}
\sphinxtableafterendhook\par
\sphinxattableend\end{savenotes}


\bigskip\hrule\bigskip



\subsection{8.2 💡 使用建议}
\label{\detokenize{chapters/chapter8:id4}}

\subsubsection{8.2.1 工具选择建议}
\label{\detokenize{chapters/chapter8:id5}}\begin{itemize}
\item {} 
\sphinxAtStartPar
根据数据格式选择合适的转换工具

\item {} 
\sphinxAtStartPar
批量处理时建议先备份原始数据

\item {} 
\sphinxAtStartPar
定期清理临时文件以节省磁盘空间

\end{itemize}


\subsubsection{8.2.2 常见问题处理}
\label{\detokenize{chapters/chapter8:id6}}

\begin{savenotes}\sphinxattablestart
\sphinxthistablewithglobalstyle
\centering
\begin{tabulary}{\linewidth}[t]{TT}
\sphinxtoprule
\begin{varwidth}[t]{\sphinxcolwidth{1}{2}}
\sphinxstyletheadfamily \sphinxAtStartPar
问题
\sphinxbeforeendvarwidth
\end{varwidth}%
&\begin{varwidth}[t]{\sphinxcolwidth{1}{2}}
\sphinxstyletheadfamily \sphinxAtStartPar
解决方案
\sphinxbeforeendvarwidth
\end{varwidth}%
\\
\sphinxmidrule
\sphinxtableatstartofbodyhook\begin{varwidth}[t]{\sphinxcolwidth{1}{2}}
\sphinxAtStartPar
转换失败
\sphinxbeforeendvarwidth
\end{varwidth}%
&\begin{varwidth}[t]{\sphinxcolwidth{1}{2}}
\sphinxAtStartPar
检查原始文件格式是否正确
\sphinxbeforeendvarwidth
\end{varwidth}%
\\
\sphinxhline\begin{varwidth}[t]{\sphinxcolwidth{1}{2}}
\sphinxAtStartPar
路径不存在
\sphinxbeforeendvarwidth
\end{varwidth}%
&\begin{varwidth}[t]{\sphinxcolwidth{1}{2}}
\sphinxAtStartPar
确认标定文件搜索路径设置正确
\sphinxbeforeendvarwidth
\end{varwidth}%
\\
\sphinxhline\begin{varwidth}[t]{\sphinxcolwidth{1}{2}}
\sphinxAtStartPar
数据丢失
\sphinxbeforeendvarwidth
\end{varwidth}%
&\begin{varwidth}[t]{\sphinxcolwidth{1}{2}}
\sphinxAtStartPar
使用备份文件重新处理
\sphinxbeforeendvarwidth
\end{varwidth}%
\\
\sphinxbottomrule
\end{tabulary}
\sphinxtableafterendhook\par
\sphinxattableend\end{savenotes}

\sphinxstepscope


\section{👁️ 第9章 视图功能}
\label{\detokenize{chapters/chapter9:id1}}\label{\detokenize{chapters/chapter9::doc}}
\sphinxAtStartPar
本章介绍MTDP的各种视图功能，包括时间序列查看、FFT实时计算、地图视图等。


\bigskip\hrule\bigskip



\subsection{9.1 🖥️ 界面布局}
\label{\detokenize{chapters/chapter9:id2}}

\subsubsection{9.1.1 主要选项卡}
\label{\detokenize{chapters/chapter9:id3}}

\begin{savenotes}\sphinxattablestart
\sphinxthistablewithglobalstyle
\centering
\begin{tabulary}{\linewidth}[t]{TT}
\sphinxtoprule
\begin{varwidth}[t]{\sphinxcolwidth{1}{2}}
\sphinxstyletheadfamily \sphinxAtStartPar
选项卡
\sphinxbeforeendvarwidth
\end{varwidth}%
&\begin{varwidth}[t]{\sphinxcolwidth{1}{2}}
\sphinxstyletheadfamily \sphinxAtStartPar
功能
\sphinxbeforeendvarwidth
\end{varwidth}%
\\
\sphinxmidrule
\sphinxtableatstartofbodyhook\begin{varwidth}[t]{\sphinxcolwidth{1}{2}}
\sphinxAtStartPar
工程
\sphinxbeforeendvarwidth
\end{varwidth}%
&\begin{varwidth}[t]{\sphinxcolwidth{1}{2}}
\sphinxAtStartPar
工程信息和结构
\sphinxbeforeendvarwidth
\end{varwidth}%
\\
\sphinxhline\begin{varwidth}[t]{\sphinxcolwidth{1}{2}}
\sphinxAtStartPar
视图
\sphinxbeforeendvarwidth
\end{varwidth}%
&\begin{varwidth}[t]{\sphinxcolwidth{1}{2}}
\sphinxAtStartPar
数据可视化
\sphinxbeforeendvarwidth
\end{varwidth}%
\\
\sphinxhline\begin{varwidth}[t]{\sphinxcolwidth{1}{2}}
\sphinxAtStartPar
标定
\sphinxbeforeendvarwidth
\end{varwidth}%
&\begin{varwidth}[t]{\sphinxcolwidth{1}{2}}
\sphinxAtStartPar
标定曲线查看
\sphinxbeforeendvarwidth
\end{varwidth}%
\\
\sphinxhline\begin{varwidth}[t]{\sphinxcolwidth{1}{2}}
\sphinxAtStartPar
线程
\sphinxbeforeendvarwidth
\end{varwidth}%
&\begin{varwidth}[t]{\sphinxcolwidth{1}{2}}
\sphinxAtStartPar
多线程状态监控
\sphinxbeforeendvarwidth
\end{varwidth}%
\\
\sphinxbottomrule
\end{tabulary}
\sphinxtableafterendhook\par
\sphinxattableend\end{savenotes}


\subsubsection{9.1.2 视图子页面}
\label{\detokenize{chapters/chapter9:id4}}

\begin{savenotes}\sphinxattablestart
\sphinxthistablewithglobalstyle
\centering
\begin{tabulary}{\linewidth}[t]{TT}
\sphinxtoprule
\begin{varwidth}[t]{\sphinxcolwidth{1}{2}}
\sphinxstyletheadfamily \sphinxAtStartPar
子页面
\sphinxbeforeendvarwidth
\end{varwidth}%
&\begin{varwidth}[t]{\sphinxcolwidth{1}{2}}
\sphinxstyletheadfamily \sphinxAtStartPar
功能
\sphinxbeforeendvarwidth
\end{varwidth}%
\\
\sphinxmidrule
\sphinxtableatstartofbodyhook\begin{varwidth}[t]{\sphinxcolwidth{1}{2}}
\sphinxAtStartPar
地图
\sphinxbeforeendvarwidth
\end{varwidth}%
&\begin{varwidth}[t]{\sphinxcolwidth{1}{2}}
\sphinxAtStartPar
测点空间分布
\sphinxbeforeendvarwidth
\end{varwidth}%
\\
\sphinxhline\begin{varwidth}[t]{\sphinxcolwidth{1}{2}}
\sphinxAtStartPar
时间序列
\sphinxbeforeendvarwidth
\end{varwidth}%
&\begin{varwidth}[t]{\sphinxcolwidth{1}{2}}
\sphinxAtStartPar
时间序列查看
\sphinxbeforeendvarwidth
\end{varwidth}%
\\
\sphinxhline\begin{varwidth}[t]{\sphinxcolwidth{1}{2}}
\sphinxAtStartPar
曲线
\sphinxbeforeendvarwidth
\end{varwidth}%
&\begin{varwidth}[t]{\sphinxcolwidth{1}{2}}
\sphinxAtStartPar
视电阻率/相位
\sphinxbeforeendvarwidth
\end{varwidth}%
\\
\sphinxbottomrule
\end{tabulary}
\sphinxtableafterendhook\par
\sphinxattableend\end{savenotes}


\bigskip\hrule\bigskip



\subsection{9.2 📊 时间序列查看器}
\label{\detokenize{chapters/chapter9:id5}}

\subsubsection{9.2.1 打开时间序列查看器}
\label{\detokenize{chapters/chapter9:id6}}
\sphinxAtStartPar
\sphinxstylestrong{方式1：从工程树添加}
\begin{enumerate}
\sphinxsetlistlabels{\arabic}{enumi}{enumii}{}{.}%
\item {} 
\sphinxAtStartPar
在工程树中选择测点

\item {} 
\sphinxAtStartPar
右键选择 \sphinxcode{\sphinxupquote{添加到时间序列显示}}

\item {} 
\sphinxAtStartPar
时间序列显示在视图窗口

\end{enumerate}

\sphinxAtStartPar
\sphinxstylestrong{方式2：打开时间序列文件}
\begin{enumerate}
\sphinxsetlistlabels{\arabic}{enumi}{enumii}{}{.}%
\item {} 
\sphinxAtStartPar
选择菜单 \sphinxcode{\sphinxupquote{工具 → 时间序列查看器}}

\item {} 
\sphinxAtStartPar
点击 \sphinxcode{\sphinxupquote{打开}} 按钮

\item {} 
\sphinxAtStartPar
选择时间序列文件（.TS、.tbl、.timeseries等）

\item {} 
\sphinxAtStartPar
系统自动加载并显示

\end{enumerate}


\subsubsection{9.2.2 界面结构}
\label{\detokenize{chapters/chapter9:id7}}
\sphinxAtStartPar
时间序列查看器分为以下区域：


\begin{savenotes}\sphinxattablestart
\sphinxthistablewithglobalstyle
\centering
\begin{tabulary}{\linewidth}[t]{TTT}
\sphinxtoprule
\begin{varwidth}[t]{\sphinxcolwidth{1}{3}}
\sphinxstyletheadfamily \sphinxAtStartPar
区域
\sphinxbeforeendvarwidth
\end{varwidth}%
&\begin{varwidth}[t]{\sphinxcolwidth{1}{3}}
\sphinxstyletheadfamily \sphinxAtStartPar
位置
\sphinxbeforeendvarwidth
\end{varwidth}%
&\begin{varwidth}[t]{\sphinxcolwidth{1}{3}}
\sphinxstyletheadfamily \sphinxAtStartPar
功能
\sphinxbeforeendvarwidth
\end{varwidth}%
\\
\sphinxmidrule
\sphinxtableatstartofbodyhook\begin{varwidth}[t]{\sphinxcolwidth{1}{3}}
\sphinxAtStartPar
工具栏
\sphinxbeforeendvarwidth
\end{varwidth}%
&\begin{varwidth}[t]{\sphinxcolwidth{1}{3}}
\sphinxAtStartPar
顶部
\sphinxbeforeendvarwidth
\end{varwidth}%
&\begin{varwidth}[t]{\sphinxcolwidth{1}{3}}
\sphinxAtStartPar
控制按钮和参数设置
\sphinxbeforeendvarwidth
\end{varwidth}%
\\
\sphinxhline\begin{varwidth}[t]{\sphinxcolwidth{1}{3}}
\sphinxAtStartPar
通道列表
\sphinxbeforeendvarwidth
\end{varwidth}%
&\begin{varwidth}[t]{\sphinxcolwidth{1}{3}}
\sphinxAtStartPar
左侧
\sphinxbeforeendvarwidth
\end{varwidth}%
&\begin{varwidth}[t]{\sphinxcolwidth{1}{3}}
\sphinxAtStartPar
通道显示控制
\sphinxbeforeendvarwidth
\end{varwidth}%
\\
\sphinxhline\begin{varwidth}[t]{\sphinxcolwidth{1}{3}}
\sphinxAtStartPar
测点列表
\sphinxbeforeendvarwidth
\end{varwidth}%
&\begin{varwidth}[t]{\sphinxcolwidth{1}{3}}
\sphinxAtStartPar
左侧
\sphinxbeforeendvarwidth
\end{varwidth}%
&\begin{varwidth}[t]{\sphinxcolwidth{1}{3}}
\sphinxAtStartPar
测点显示控制
\sphinxbeforeendvarwidth
\end{varwidth}%
\\
\sphinxhline\begin{varwidth}[t]{\sphinxcolwidth{1}{3}}
\sphinxAtStartPar
图表区
\sphinxbeforeendvarwidth
\end{varwidth}%
&\begin{varwidth}[t]{\sphinxcolwidth{1}{3}}
\sphinxAtStartPar
中央
\sphinxbeforeendvarwidth
\end{varwidth}%
&\begin{varwidth}[t]{\sphinxcolwidth{1}{3}}
\sphinxAtStartPar
时间序列波形显示
\sphinxbeforeendvarwidth
\end{varwidth}%
\\
\sphinxhline\begin{varwidth}[t]{\sphinxcolwidth{1}{3}}
\sphinxAtStartPar
时间轴
\sphinxbeforeendvarwidth
\end{varwidth}%
&\begin{varwidth}[t]{\sphinxcolwidth{1}{3}}
\sphinxAtStartPar
底部
\sphinxbeforeendvarwidth
\end{varwidth}%
&\begin{varwidth}[t]{\sphinxcolwidth{1}{3}}
\sphinxAtStartPar
时间范围控制
\sphinxbeforeendvarwidth
\end{varwidth}%
\\
\sphinxbottomrule
\end{tabulary}
\sphinxtableafterendhook\par
\sphinxattableend\end{savenotes}


\subsubsection{9.2.3 通道显示控制}
\label{\detokenize{chapters/chapter9:id8}}
\sphinxAtStartPar
\sphinxstylestrong{通道列表功能：}


\begin{savenotes}\sphinxattablestart
\sphinxthistablewithglobalstyle
\centering
\begin{tabulary}{\linewidth}[t]{TT}
\sphinxtoprule
\begin{varwidth}[t]{\sphinxcolwidth{1}{2}}
\sphinxstyletheadfamily \sphinxAtStartPar
操作
\sphinxbeforeendvarwidth
\end{varwidth}%
&\begin{varwidth}[t]{\sphinxcolwidth{1}{2}}
\sphinxstyletheadfamily \sphinxAtStartPar
说明
\sphinxbeforeendvarwidth
\end{varwidth}%
\\
\sphinxmidrule
\sphinxtableatstartofbodyhook\begin{varwidth}[t]{\sphinxcolwidth{1}{2}}
\sphinxAtStartPar
勾选通道
\sphinxbeforeendvarwidth
\end{varwidth}%
&\begin{varwidth}[t]{\sphinxcolwidth{1}{2}}
\sphinxAtStartPar
显示该通道波形
\sphinxbeforeendvarwidth
\end{varwidth}%
\\
\sphinxhline\begin{varwidth}[t]{\sphinxcolwidth{1}{2}}
\sphinxAtStartPar
取消勾选
\sphinxbeforeendvarwidth
\end{varwidth}%
&\begin{varwidth}[t]{\sphinxcolwidth{1}{2}}
\sphinxAtStartPar
隐藏该通道波形
\sphinxbeforeendvarwidth
\end{varwidth}%
\\
\sphinxhline\begin{varwidth}[t]{\sphinxcolwidth{1}{2}}
\sphinxAtStartPar
全选
\sphinxbeforeendvarwidth
\end{varwidth}%
&\begin{varwidth}[t]{\sphinxcolwidth{1}{2}}
\sphinxAtStartPar
显示所有通道
\sphinxbeforeendvarwidth
\end{varwidth}%
\\
\sphinxhline\begin{varwidth}[t]{\sphinxcolwidth{1}{2}}
\sphinxAtStartPar
全不选
\sphinxbeforeendvarwidth
\end{varwidth}%
&\begin{varwidth}[t]{\sphinxcolwidth{1}{2}}
\sphinxAtStartPar
隐藏所有通道
\sphinxbeforeendvarwidth
\end{varwidth}%
\\
\sphinxbottomrule
\end{tabulary}
\sphinxtableafterendhook\par
\sphinxattableend\end{savenotes}

\sphinxAtStartPar
\sphinxstylestrong{标准5通道：}
\begin{itemize}
\item {} 
\sphinxAtStartPar
Ex \sphinxhyphen{} X方向电场

\item {} 
\sphinxAtStartPar
Ey \sphinxhyphen{} Y方向电场

\item {} 
\sphinxAtStartPar
Hx \sphinxhyphen{} X方向磁场

\item {} 
\sphinxAtStartPar
Hy \sphinxhyphen{} Y方向磁场

\item {} 
\sphinxAtStartPar
Hz \sphinxhyphen{} Z方向磁场

\end{itemize}


\subsubsection{9.2.4 测点显示控制}
\label{\detokenize{chapters/chapter9:id9}}
\sphinxAtStartPar
\sphinxstylestrong{测点列表功能：}


\begin{savenotes}\sphinxattablestart
\sphinxthistablewithglobalstyle
\centering
\begin{tabulary}{\linewidth}[t]{TT}
\sphinxtoprule
\begin{varwidth}[t]{\sphinxcolwidth{1}{2}}
\sphinxstyletheadfamily \sphinxAtStartPar
操作
\sphinxbeforeendvarwidth
\end{varwidth}%
&\begin{varwidth}[t]{\sphinxcolwidth{1}{2}}
\sphinxstyletheadfamily \sphinxAtStartPar
说明
\sphinxbeforeendvarwidth
\end{varwidth}%
\\
\sphinxmidrule
\sphinxtableatstartofbodyhook\begin{varwidth}[t]{\sphinxcolwidth{1}{2}}
\sphinxAtStartPar
勾选测点
\sphinxbeforeendvarwidth
\end{varwidth}%
&\begin{varwidth}[t]{\sphinxcolwidth{1}{2}}
\sphinxAtStartPar
显示该测点数据
\sphinxbeforeendvarwidth
\end{varwidth}%
\\
\sphinxhline\begin{varwidth}[t]{\sphinxcolwidth{1}{2}}
\sphinxAtStartPar
取消勾选
\sphinxbeforeendvarwidth
\end{varwidth}%
&\begin{varwidth}[t]{\sphinxcolwidth{1}{2}}
\sphinxAtStartPar
隐藏该测点数据
\sphinxbeforeendvarwidth
\end{varwidth}%
\\
\sphinxhline\begin{varwidth}[t]{\sphinxcolwidth{1}{2}}
\sphinxAtStartPar
点击测点
\sphinxbeforeendvarwidth
\end{varwidth}%
&\begin{varwidth}[t]{\sphinxcolwidth{1}{2}}
\sphinxAtStartPar
选中测点，可修改颜色
\sphinxbeforeendvarwidth
\end{varwidth}%
\\
\sphinxhline\begin{varwidth}[t]{\sphinxcolwidth{1}{2}}
\sphinxAtStartPar
颜色面板
\sphinxbeforeendvarwidth
\end{varwidth}%
&\begin{varwidth}[t]{\sphinxcolwidth{1}{2}}
\sphinxAtStartPar
设置选中测点的颜色
\sphinxbeforeendvarwidth
\end{varwidth}%
\\
\sphinxbottomrule
\end{tabulary}
\sphinxtableafterendhook\par
\sphinxattableend\end{savenotes}

\sphinxAtStartPar
\sphinxstylestrong{颜色设置：}
\begin{enumerate}
\sphinxsetlistlabels{\arabic}{enumi}{enumii}{}{.}%
\item {} 
\sphinxAtStartPar
在测点列表中点击要修改的测点

\item {} 
\sphinxAtStartPar
在颜色面板中选择新颜色

\item {} 
\sphinxAtStartPar
该测点的所有通道曲线颜色更新

\end{enumerate}


\subsubsection{9.2.5 显示参数设置}
\label{\detokenize{chapters/chapter9:id10}}

\begin{savenotes}\sphinxattablestart
\sphinxthistablewithglobalstyle
\centering
\begin{tabulary}{\linewidth}[t]{TTT}
\sphinxtoprule
\begin{varwidth}[t]{\sphinxcolwidth{1}{3}}
\sphinxstyletheadfamily \sphinxAtStartPar
参数
\sphinxbeforeendvarwidth
\end{varwidth}%
&\begin{varwidth}[t]{\sphinxcolwidth{1}{3}}
\sphinxstyletheadfamily \sphinxAtStartPar
说明
\sphinxbeforeendvarwidth
\end{varwidth}%
&\begin{varwidth}[t]{\sphinxcolwidth{1}{3}}
\sphinxstyletheadfamily \sphinxAtStartPar
默认值
\sphinxbeforeendvarwidth
\end{varwidth}%
\\
\sphinxmidrule
\sphinxtableatstartofbodyhook\begin{varwidth}[t]{\sphinxcolwidth{1}{3}}
\sphinxAtStartPar
显示长度
\sphinxbeforeendvarwidth
\end{varwidth}%
&\begin{varwidth}[t]{\sphinxcolwidth{1}{3}}
\sphinxAtStartPar
单屏显示的数据点数
\sphinxbeforeendvarwidth
\end{varwidth}%
&\begin{varwidth}[t]{\sphinxcolwidth{1}{3}}
\sphinxAtStartPar
24000
\sphinxbeforeendvarwidth
\end{varwidth}%
\\
\sphinxhline\begin{varwidth}[t]{\sphinxcolwidth{1}{3}}
\sphinxAtStartPar
起始位置
\sphinxbeforeendvarwidth
\end{varwidth}%
&\begin{varwidth}[t]{\sphinxcolwidth{1}{3}}
\sphinxAtStartPar
显示的起始索引
\sphinxbeforeendvarwidth
\end{varwidth}%
&\begin{varwidth}[t]{\sphinxcolwidth{1}{3}}
\sphinxAtStartPar
0
\sphinxbeforeendvarwidth
\end{varwidth}%
\\
\sphinxhline\begin{varwidth}[t]{\sphinxcolwidth{1}{3}}
\sphinxAtStartPar
采样间隔
\sphinxbeforeendvarwidth
\end{varwidth}%
&\begin{varwidth}[t]{\sphinxcolwidth{1}{3}}
\sphinxAtStartPar
数据采样间隔
\sphinxbeforeendvarwidth
\end{varwidth}%
&\begin{varwidth}[t]{\sphinxcolwidth{1}{3}}
\sphinxAtStartPar
1
\sphinxbeforeendvarwidth
\end{varwidth}%
\\
\sphinxhline\begin{varwidth}[t]{\sphinxcolwidth{1}{3}}
\sphinxAtStartPar
时区偏移
\sphinxbeforeendvarwidth
\end{varwidth}%
&\begin{varwidth}[t]{\sphinxcolwidth{1}{3}}
\sphinxAtStartPar
时间显示的时区偏移（小时）
\sphinxbeforeendvarwidth
\end{varwidth}%
&\begin{varwidth}[t]{\sphinxcolwidth{1}{3}}
\sphinxAtStartPar
0
\sphinxbeforeendvarwidth
\end{varwidth}%
\\
\sphinxbottomrule
\end{tabulary}
\sphinxtableafterendhook\par
\sphinxattableend\end{savenotes}

\sphinxAtStartPar
\sphinxstylestrong{显示长度说明：}
\begin{itemize}
\item {} 
\sphinxAtStartPar
设置过小：显示数据少，细节清晰

\item {} 
\sphinxAtStartPar
设置过大：显示数据多，细节模糊

\item {} 
\sphinxAtStartPar
推荐值：采样率 × 10秒（如2400Hz采样率设为24000）

\end{itemize}


\subsubsection{9.2.6 时间轴控制}
\label{\detokenize{chapters/chapter9:id11}}
\sphinxAtStartPar
\sphinxstylestrong{时间轴滑块：}
\begin{itemize}
\item {} 
\sphinxAtStartPar
拖动滑块浏览不同时间段的数据

\item {} 
\sphinxAtStartPar
滑块范围 = 总数据量 / 显示长度

\end{itemize}

\sphinxAtStartPar
\sphinxstylestrong{起始位置输入：}
\begin{itemize}
\item {} 
\sphinxAtStartPar
直接输入起始索引值精确定位

\end{itemize}

\sphinxAtStartPar
\sphinxstylestrong{时区设置：}
\begin{itemize}
\item {} 
\sphinxAtStartPar
调整时间显示的时区偏移

\item {} 
\sphinxAtStartPar
用于统一不同时区采集的数据

\end{itemize}


\subsubsection{9.2.7 显示选项}
\label{\detokenize{chapters/chapter9:id12}}

\begin{savenotes}\sphinxattablestart
\sphinxthistablewithglobalstyle
\centering
\begin{tabulary}{\linewidth}[t]{TT}
\sphinxtoprule
\begin{varwidth}[t]{\sphinxcolwidth{1}{2}}
\sphinxstyletheadfamily \sphinxAtStartPar
选项
\sphinxbeforeendvarwidth
\end{varwidth}%
&\begin{varwidth}[t]{\sphinxcolwidth{1}{2}}
\sphinxstyletheadfamily \sphinxAtStartPar
说明
\sphinxbeforeendvarwidth
\end{varwidth}%
\\
\sphinxmidrule
\sphinxtableatstartofbodyhook\begin{varwidth}[t]{\sphinxcolwidth{1}{2}}
\sphinxAtStartPar
自动X轴
\sphinxbeforeendvarwidth
\end{varwidth}%
&\begin{varwidth}[t]{\sphinxcolwidth{1}{2}}
\sphinxAtStartPar
X轴自动调整范围
\sphinxbeforeendvarwidth
\end{varwidth}%
\\
\sphinxhline\begin{varwidth}[t]{\sphinxcolwidth{1}{2}}
\sphinxAtStartPar
自动Y偏移
\sphinxbeforeendvarwidth
\end{varwidth}%
&\begin{varwidth}[t]{\sphinxcolwidth{1}{2}}
\sphinxAtStartPar
多曲线自动Y偏移，避免重叠
\sphinxbeforeendvarwidth
\end{varwidth}%
\\
\sphinxhline\begin{varwidth}[t]{\sphinxcolwidth{1}{2}}
\sphinxAtStartPar
差分显示
\sphinxbeforeendvarwidth
\end{varwidth}%
&\begin{varwidth}[t]{\sphinxcolwidth{1}{2}}
\sphinxAtStartPar
显示相邻点差分值
\sphinxbeforeendvarwidth
\end{varwidth}%
\\
\sphinxbottomrule
\end{tabulary}
\sphinxtableafterendhook\par
\sphinxattableend\end{savenotes}

\sphinxAtStartPar
\sphinxstylestrong{自动Y偏移：}
\begin{itemize}
\item {} 
\sphinxAtStartPar
开启时，多个测点的曲线自动在Y方向错开

\item {} 
\sphinxAtStartPar
便于对比不同测点的波形特征

\end{itemize}

\sphinxAtStartPar
\sphinxstylestrong{差分显示：}
\begin{itemize}
\item {} 
\sphinxAtStartPar
显示相邻数据点的差值

\item {} 
\sphinxAtStartPar
用于检测突变和跳变

\end{itemize}


\subsubsection{9.2.8 多采样率支持}
\label{\detokenize{chapters/chapter9:id13}}
\sphinxAtStartPar
对于多采样率数据（如Phoenix MTU）：
\begin{enumerate}
\sphinxsetlistlabels{\arabic}{enumi}{enumii}{}{.}%
\item {} 
\sphinxAtStartPar
使用采样率下拉框选择要查看的频段

\item {} 
\sphinxAtStartPar
不同采样率的数据分开显示

\item {} 
\sphinxAtStartPar
常见采样率：
\begin{itemize}
\item {} 
\sphinxAtStartPar
TS2: 24000 Hz

\item {} 
\sphinxAtStartPar
TS3: 2400 Hz

\item {} 
\sphinxAtStartPar
TS4: 150 Hz

\item {} 
\sphinxAtStartPar
TS5: 15 Hz

\end{itemize}

\end{enumerate}


\bigskip\hrule\bigskip



\subsection{9.3 📈 FFT实时计算}
\label{\detokenize{chapters/chapter9:fft}}

\subsubsection{9.3.1 功能说明}
\label{\detokenize{chapters/chapter9:id14}}
\sphinxAtStartPar
时间序列查看器集成了FFT实时计算功能，可在浏览时间序列的同时查看频谱分析结果。


\subsubsection{9.3.2 FFT选项卡}
\label{\detokenize{chapters/chapter9:id15}}
\sphinxAtStartPar
切换到 \sphinxcode{\sphinxupquote{FFT}} 选项卡查看频谱分析结果：


\begin{savenotes}\sphinxattablestart
\sphinxthistablewithglobalstyle
\centering
\begin{tabulary}{\linewidth}[t]{TT}
\sphinxtoprule
\begin{varwidth}[t]{\sphinxcolwidth{1}{2}}
\sphinxstyletheadfamily \sphinxAtStartPar
子选项卡
\sphinxbeforeendvarwidth
\end{varwidth}%
&\begin{varwidth}[t]{\sphinxcolwidth{1}{2}}
\sphinxstyletheadfamily \sphinxAtStartPar
内容
\sphinxbeforeendvarwidth
\end{varwidth}%
\\
\sphinxmidrule
\sphinxtableatstartofbodyhook\begin{varwidth}[t]{\sphinxcolwidth{1}{2}}
\sphinxAtStartPar
频谱
\sphinxbeforeendvarwidth
\end{varwidth}%
&\begin{varwidth}[t]{\sphinxcolwidth{1}{2}}
\sphinxAtStartPar
各通道功率谱密度
\sphinxbeforeendvarwidth
\end{varwidth}%
\\
\sphinxhline\begin{varwidth}[t]{\sphinxcolwidth{1}{2}}
\sphinxAtStartPar
相干性
\sphinxbeforeendvarwidth
\end{varwidth}%
&\begin{varwidth}[t]{\sphinxcolwidth{1}{2}}
\sphinxAtStartPar
通道间相干性曲线
\sphinxbeforeendvarwidth
\end{varwidth}%
\\
\sphinxbottomrule
\end{tabulary}
\sphinxtableafterendhook\par
\sphinxattableend\end{savenotes}


\subsubsection{9.3.3 FFT参数设置}
\label{\detokenize{chapters/chapter9:id16}}

\begin{savenotes}\sphinxattablestart
\sphinxthistablewithglobalstyle
\centering
\begin{tabulary}{\linewidth}[t]{TT}
\sphinxtoprule
\begin{varwidth}[t]{\sphinxcolwidth{1}{2}}
\sphinxstyletheadfamily \sphinxAtStartPar
参数
\sphinxbeforeendvarwidth
\end{varwidth}%
&\begin{varwidth}[t]{\sphinxcolwidth{1}{2}}
\sphinxstyletheadfamily \sphinxAtStartPar
说明
\sphinxbeforeendvarwidth
\end{varwidth}%
\\
\sphinxmidrule
\sphinxtableatstartofbodyhook\begin{varwidth}[t]{\sphinxcolwidth{1}{2}}
\sphinxAtStartPar
窗函数
\sphinxbeforeendvarwidth
\end{varwidth}%
&\begin{varwidth}[t]{\sphinxcolwidth{1}{2}}
\sphinxAtStartPar
FFT窗函数类型
\sphinxbeforeendvarwidth
\end{varwidth}%
\\
\sphinxhline\begin{varwidth}[t]{\sphinxcolwidth{1}{2}}
\sphinxAtStartPar
FFT长度
\sphinxbeforeendvarwidth
\end{varwidth}%
&\begin{varwidth}[t]{\sphinxcolwidth{1}{2}}
\sphinxAtStartPar
FFT变换长度
\sphinxbeforeendvarwidth
\end{varwidth}%
\\
\sphinxhline\begin{varwidth}[t]{\sphinxcolwidth{1}{2}}
\sphinxAtStartPar
重叠率
\sphinxbeforeendvarwidth
\end{varwidth}%
&\begin{varwidth}[t]{\sphinxcolwidth{1}{2}}
\sphinxAtStartPar
相邻FFT窗口的重叠比例
\sphinxbeforeendvarwidth
\end{varwidth}%
\\
\sphinxhline\begin{varwidth}[t]{\sphinxcolwidth{1}{2}}
\sphinxAtStartPar
多锥谱数
\sphinxbeforeendvarwidth
\end{varwidth}%
&\begin{varwidth}[t]{\sphinxcolwidth{1}{2}}
\sphinxAtStartPar
多锥谱分析的锥数
\sphinxbeforeendvarwidth
\end{varwidth}%
\\
\sphinxhline\begin{varwidth}[t]{\sphinxcolwidth{1}{2}}
\sphinxAtStartPar
零填充
\sphinxbeforeendvarwidth
\end{varwidth}%
&\begin{varwidth}[t]{\sphinxcolwidth{1}{2}}
\sphinxAtStartPar
零填充因子
\sphinxbeforeendvarwidth
\end{varwidth}%
\\
\sphinxbottomrule
\end{tabulary}
\sphinxtableafterendhook\par
\sphinxattableend\end{savenotes}

\sphinxAtStartPar
\sphinxstylestrong{窗函数选项：}


\begin{savenotes}\sphinxattablestart
\sphinxthistablewithglobalstyle
\centering
\begin{tabulary}{\linewidth}[t]{TT}
\sphinxtoprule
\begin{varwidth}[t]{\sphinxcolwidth{1}{2}}
\sphinxstyletheadfamily \sphinxAtStartPar
窗函数
\sphinxbeforeendvarwidth
\end{varwidth}%
&\begin{varwidth}[t]{\sphinxcolwidth{1}{2}}
\sphinxstyletheadfamily \sphinxAtStartPar
适用场景
\sphinxbeforeendvarwidth
\end{varwidth}%
\\
\sphinxmidrule
\sphinxtableatstartofbodyhook\begin{varwidth}[t]{\sphinxcolwidth{1}{2}}
\sphinxAtStartPar
矩形窗
\sphinxbeforeendvarwidth
\end{varwidth}%
&\begin{varwidth}[t]{\sphinxcolwidth{1}{2}}
\sphinxAtStartPar
瞬态信号
\sphinxbeforeendvarwidth
\end{varwidth}%
\\
\sphinxhline\begin{varwidth}[t]{\sphinxcolwidth{1}{2}}
\sphinxAtStartPar
汉明窗
\sphinxbeforeendvarwidth
\end{varwidth}%
&\begin{varwidth}[t]{\sphinxcolwidth{1}{2}}
\sphinxAtStartPar
一般数据
\sphinxbeforeendvarwidth
\end{varwidth}%
\\
\sphinxhline\begin{varwidth}[t]{\sphinxcolwidth{1}{2}}
\sphinxAtStartPar
汉宁窗
\sphinxbeforeendvarwidth
\end{varwidth}%
&\begin{varwidth}[t]{\sphinxcolwidth{1}{2}}
\sphinxAtStartPar
一般数据（推荐）
\sphinxbeforeendvarwidth
\end{varwidth}%
\\
\sphinxhline\begin{varwidth}[t]{\sphinxcolwidth{1}{2}}
\sphinxAtStartPar
布莱克曼窗
\sphinxbeforeendvarwidth
\end{varwidth}%
&\begin{varwidth}[t]{\sphinxcolwidth{1}{2}}
\sphinxAtStartPar
高精度频谱
\sphinxbeforeendvarwidth
\end{varwidth}%
\\
\sphinxhline\begin{varwidth}[t]{\sphinxcolwidth{1}{2}}
\sphinxAtStartPar
平顶窗
\sphinxbeforeendvarwidth
\end{varwidth}%
&\begin{varwidth}[t]{\sphinxcolwidth{1}{2}}
\sphinxAtStartPar
幅度测量
\sphinxbeforeendvarwidth
\end{varwidth}%
\\
\sphinxhline\begin{varwidth}[t]{\sphinxcolwidth{1}{2}}
\sphinxAtStartPar
高斯窗
\sphinxbeforeendvarwidth
\end{varwidth}%
&\begin{varwidth}[t]{\sphinxcolwidth{1}{2}}
\sphinxAtStartPar
瞬态分析
\sphinxbeforeendvarwidth
\end{varwidth}%
\\
\sphinxhline\begin{varwidth}[t]{\sphinxcolwidth{1}{2}}
\sphinxAtStartPar
巴特利特窗
\sphinxbeforeendvarwidth
\end{varwidth}%
&\begin{varwidth}[t]{\sphinxcolwidth{1}{2}}
\sphinxAtStartPar
频谱估计
\sphinxbeforeendvarwidth
\end{varwidth}%
\\
\sphinxhline\begin{varwidth}[t]{\sphinxcolwidth{1}{2}}
\sphinxAtStartPar
多锥正弦窗
\sphinxbeforeendvarwidth
\end{varwidth}%
&\begin{varwidth}[t]{\sphinxcolwidth{1}{2}}
\sphinxAtStartPar
高精度谱估计
\sphinxbeforeendvarwidth
\end{varwidth}%
\\
\sphinxhline\begin{varwidth}[t]{\sphinxcolwidth{1}{2}}
\sphinxAtStartPar
多锥Slepian窗
\sphinxbeforeendvarwidth
\end{varwidth}%
&\begin{varwidth}[t]{\sphinxcolwidth{1}{2}}
\sphinxAtStartPar
最优谱估计
\sphinxbeforeendvarwidth
\end{varwidth}%
\\
\sphinxbottomrule
\end{tabulary}
\sphinxtableafterendhook\par
\sphinxattableend\end{savenotes}


\subsubsection{9.3.4 FFT计算方式}
\label{\detokenize{chapters/chapter9:id17}}

\begin{savenotes}\sphinxattablestart
\sphinxthistablewithglobalstyle
\centering
\begin{tabulary}{\linewidth}[t]{TT}
\sphinxtoprule
\begin{varwidth}[t]{\sphinxcolwidth{1}{2}}
\sphinxstyletheadfamily \sphinxAtStartPar
选项
\sphinxbeforeendvarwidth
\end{varwidth}%
&\begin{varwidth}[t]{\sphinxcolwidth{1}{2}}
\sphinxstyletheadfamily \sphinxAtStartPar
说明
\sphinxbeforeendvarwidth
\end{varwidth}%
\\
\sphinxmidrule
\sphinxtableatstartofbodyhook\begin{varwidth}[t]{\sphinxcolwidth{1}{2}}
\sphinxAtStartPar
自动FFT
\sphinxbeforeendvarwidth
\end{varwidth}%
&\begin{varwidth}[t]{\sphinxcolwidth{1}{2}}
\sphinxAtStartPar
滚动时间轴时自动更新FFT
\sphinxbeforeendvarwidth
\end{varwidth}%
\\
\sphinxhline\begin{varwidth}[t]{\sphinxcolwidth{1}{2}}
\sphinxAtStartPar
全数据FFT
\sphinxbeforeendvarwidth
\end{varwidth}%
&\begin{varwidth}[t]{\sphinxcolwidth{1}{2}}
\sphinxAtStartPar
对全部显示数据计算FFT
\sphinxbeforeendvarwidth
\end{varwidth}%
\\
\sphinxbottomrule
\end{tabulary}
\sphinxtableafterendhook\par
\sphinxattableend\end{savenotes}

\sphinxAtStartPar
\sphinxstylestrong{自动FFT模式：}
\begin{itemize}
\item {} 
\sphinxAtStartPar
拖动时间轴时，FFT结果自动更新

\item {} 
\sphinxAtStartPar
便于快速查看不同时间段的频谱特征

\end{itemize}

\sphinxAtStartPar
\sphinxstylestrong{全数据FFT模式：}
\begin{itemize}
\item {} 
\sphinxAtStartPar
对当前显示的全部数据计算FFT

\item {} 
\sphinxAtStartPar
频率分辨率更高，但更新较慢

\end{itemize}


\subsubsection{9.3.5 频谱显示}
\label{\detokenize{chapters/chapter9:id18}}
\sphinxAtStartPar
频谱图显示各通道的功率谱密度：
\begin{itemize}
\item {} 
\sphinxAtStartPar
X轴：频率（Hz）

\item {} 
\sphinxAtStartPar
Y轴：功率谱密度（对数坐标）

\item {} 
\sphinxAtStartPar
不同颜色表示不同通道

\end{itemize}

\sphinxAtStartPar
\sphinxstylestrong{通道选择：}
\begin{itemize}
\item {} 
\sphinxAtStartPar
在左侧通道列表中勾选要显示的通道

\end{itemize}


\subsubsection{9.3.6 相干性显示}
\label{\detokenize{chapters/chapter9:id19}}
\sphinxAtStartPar
相干性图显示通道间的相干性：
\begin{enumerate}
\sphinxsetlistlabels{\arabic}{enumi}{enumii}{}{.}%
\item {} 
\sphinxAtStartPar
选择参考通道（从下拉框选择）

\item {} 
\sphinxAtStartPar
显示该通道与其他所有通道的相干性曲线

\item {} 
\sphinxAtStartPar
相干性值范围：0\sphinxhyphen{}1

\item {} 
\sphinxAtStartPar
高相干性（接近1）表示两通道信号高度相关

\end{enumerate}

\sphinxAtStartPar
\sphinxstylestrong{相干性应用：}
\begin{itemize}
\item {} 
\sphinxAtStartPar
检查信号质量

\item {} 
\sphinxAtStartPar
识别相关噪声

\item {} 
\sphinxAtStartPar
评估远参考效果

\end{itemize}


\bigskip\hrule\bigskip



\subsection{9.4 📻 Phoenix时间序列查看器}
\label{\detokenize{chapters/chapter9:phoenix}}

\subsubsection{9.4.1 专用功能}
\label{\detokenize{chapters/chapter9:id20}}
\sphinxAtStartPar
Phoenix时间序列查看器针对Phoenix仪器数据提供专用功能：
\begin{itemize}
\item {} 
\sphinxAtStartPar
直接打开 .tbl 文件

\item {} 
\sphinxAtStartPar
自动识别多采样率数据

\item {} 
\sphinxAtStartPar
支持 .timeseries 格式

\end{itemize}


\subsubsection{9.4.2 从TBL文件加载}
\label{\detokenize{chapters/chapter9:tbl}}\begin{enumerate}
\sphinxsetlistlabels{\arabic}{enumi}{enumii}{}{.}%
\item {} 
\sphinxAtStartPar
点击 \sphinxcode{\sphinxupquote{打开}} 按钮

\item {} 
\sphinxAtStartPar
选择 .tbl 文件

\item {} 
\sphinxAtStartPar
系统自动：
\begin{itemize}
\item {} 
\sphinxAtStartPar
读取头文件信息

\item {} 
\sphinxAtStartPar
加载关联的时间序列文件

\item {} 
\sphinxAtStartPar
识别采样率配置

\end{itemize}

\end{enumerate}


\subsubsection{9.4.3 数据类型支持}
\label{\detokenize{chapters/chapter9:id21}}

\begin{savenotes}\sphinxattablestart
\sphinxthistablewithglobalstyle
\centering
\begin{tabulary}{\linewidth}[t]{TT}
\sphinxtoprule
\begin{varwidth}[t]{\sphinxcolwidth{1}{2}}
\sphinxstyletheadfamily \sphinxAtStartPar
数据类型
\sphinxbeforeendvarwidth
\end{varwidth}%
&\begin{varwidth}[t]{\sphinxcolwidth{1}{2}}
\sphinxstyletheadfamily \sphinxAtStartPar
说明
\sphinxbeforeendvarwidth
\end{varwidth}%
\\
\sphinxmidrule
\sphinxtableatstartofbodyhook\begin{varwidth}[t]{\sphinxcolwidth{1}{2}}
\sphinxAtStartPar
连续时间序列
\sphinxbeforeendvarwidth
\end{varwidth}%
&\begin{varwidth}[t]{\sphinxcolwidth{1}{2}}
\sphinxAtStartPar
单一连续采集段
\sphinxbeforeendvarwidth
\end{varwidth}%
\\
\sphinxhline\begin{varwidth}[t]{\sphinxcolwidth{1}{2}}
\sphinxAtStartPar
多段时间序列
\sphinxbeforeendvarwidth
\end{varwidth}%
&\begin{varwidth}[t]{\sphinxcolwidth{1}{2}}
\sphinxAtStartPar
多个采集段，支持时间对齐
\sphinxbeforeendvarwidth
\end{varwidth}%
\\
\sphinxbottomrule
\end{tabulary}
\sphinxtableafterendhook\par
\sphinxattableend\end{savenotes}

\sphinxAtStartPar
\sphinxstylestrong{多段时间序列处理：}
\begin{itemize}
\item {} 
\sphinxAtStartPar
自动按时间顺序排列

\item {} 
\sphinxAtStartPar
支持多测点时间对齐

\item {} 
\sphinxAtStartPar
便于远参考数据处理

\end{itemize}


\bigskip\hrule\bigskip



\subsection{9.5 🗺️ 地图视图}
\label{\detokenize{chapters/chapter9:id22}}

\subsubsection{9.5.1 功能说明}
\label{\detokenize{chapters/chapter9:id23}}
\sphinxAtStartPar
地图视图显示测点的空间分布，支持与GIS软件联动。


\subsubsection{9.5.2 使用方法}
\label{\detokenize{chapters/chapter9:id24}}\begin{enumerate}
\sphinxsetlistlabels{\arabic}{enumi}{enumii}{}{.}%
\item {} 
\sphinxAtStartPar
切换到 \sphinxcode{\sphinxupquote{视图 → 地图}} 选项卡

\item {} 
\sphinxAtStartPar
选择工区或测段

\item {} 
\sphinxAtStartPar
查看测点空间分布

\item {} 
\sphinxAtStartPar
点击测点查看详情

\end{enumerate}


\subsubsection{9.5.3 地图操作}
\label{\detokenize{chapters/chapter9:id25}}

\begin{savenotes}\sphinxattablestart
\sphinxthistablewithglobalstyle
\centering
\begin{tabulary}{\linewidth}[t]{TT}
\sphinxtoprule
\begin{varwidth}[t]{\sphinxcolwidth{1}{2}}
\sphinxstyletheadfamily \sphinxAtStartPar
操作
\sphinxbeforeendvarwidth
\end{varwidth}%
&\begin{varwidth}[t]{\sphinxcolwidth{1}{2}}
\sphinxstyletheadfamily \sphinxAtStartPar
功能
\sphinxbeforeendvarwidth
\end{varwidth}%
\\
\sphinxmidrule
\sphinxtableatstartofbodyhook\begin{varwidth}[t]{\sphinxcolwidth{1}{2}}
\sphinxAtStartPar
左键点击
\sphinxbeforeendvarwidth
\end{varwidth}%
&\begin{varwidth}[t]{\sphinxcolwidth{1}{2}}
\sphinxAtStartPar
选中测点
\sphinxbeforeendvarwidth
\end{varwidth}%
\\
\sphinxhline\begin{varwidth}[t]{\sphinxcolwidth{1}{2}}
\sphinxAtStartPar
左键拖动
\sphinxbeforeendvarwidth
\end{varwidth}%
&\begin{varwidth}[t]{\sphinxcolwidth{1}{2}}
\sphinxAtStartPar
平移地图
\sphinxbeforeendvarwidth
\end{varwidth}%
\\
\sphinxhline\begin{varwidth}[t]{\sphinxcolwidth{1}{2}}
\sphinxAtStartPar
滚轮
\sphinxbeforeendvarwidth
\end{varwidth}%
&\begin{varwidth}[t]{\sphinxcolwidth{1}{2}}
\sphinxAtStartPar
缩放地图
\sphinxbeforeendvarwidth
\end{varwidth}%
\\
\sphinxhline\begin{varwidth}[t]{\sphinxcolwidth{1}{2}}
\sphinxAtStartPar
双击
\sphinxbeforeendvarwidth
\end{varwidth}%
&\begin{varwidth}[t]{\sphinxcolwidth{1}{2}}
\sphinxAtStartPar
定位到测点
\sphinxbeforeendvarwidth
\end{varwidth}%
\\
\sphinxbottomrule
\end{tabulary}
\sphinxtableafterendhook\par
\sphinxattableend\end{savenotes}


\bigskip\hrule\bigskip



\subsection{9.6 📉 视电阻率/相位曲线}
\label{\detokenize{chapters/chapter9:id26}}

\subsubsection{9.6.1 打开曲线视图}
\label{\detokenize{chapters/chapter9:id27}}\begin{enumerate}
\sphinxsetlistlabels{\arabic}{enumi}{enumii}{}{.}%
\item {} 
\sphinxAtStartPar
双击测点

\item {} 
\sphinxAtStartPar
或切换到 \sphinxcode{\sphinxupquote{视图 → 曲线}} 选项卡

\end{enumerate}


\subsubsection{9.6.2 显示选项}
\label{\detokenize{chapters/chapter9:id28}}

\begin{savenotes}\sphinxattablestart
\sphinxthistablewithglobalstyle
\centering
\begin{tabulary}{\linewidth}[t]{TT}
\sphinxtoprule
\begin{varwidth}[t]{\sphinxcolwidth{1}{2}}
\sphinxstyletheadfamily \sphinxAtStartPar
选项
\sphinxbeforeendvarwidth
\end{varwidth}%
&\begin{varwidth}[t]{\sphinxcolwidth{1}{2}}
\sphinxstyletheadfamily \sphinxAtStartPar
说明
\sphinxbeforeendvarwidth
\end{varwidth}%
\\
\sphinxmidrule
\sphinxtableatstartofbodyhook\begin{varwidth}[t]{\sphinxcolwidth{1}{2}}
\sphinxAtStartPar
XY/YX模式
\sphinxbeforeendvarwidth
\end{varwidth}%
&\begin{varwidth}[t]{\sphinxcolwidth{1}{2}}
\sphinxAtStartPar
切换显示模式
\sphinxbeforeendvarwidth
\end{varwidth}%
\\
\sphinxhline\begin{varwidth}[t]{\sphinxcolwidth{1}{2}}
\sphinxAtStartPar
误差棒
\sphinxbeforeendvarwidth
\end{varwidth}%
&\begin{varwidth}[t]{\sphinxcolwidth{1}{2}}
\sphinxAtStartPar
显示/隐藏误差棒
\sphinxbeforeendvarwidth
\end{varwidth}%
\\
\sphinxhline\begin{varwidth}[t]{\sphinxcolwidth{1}{2}}
\sphinxAtStartPar
多测点叠加
\sphinxbeforeendvarwidth
\end{varwidth}%
&\begin{varwidth}[t]{\sphinxcolwidth{1}{2}}
\sphinxAtStartPar
多个测点曲线叠加显示
\sphinxbeforeendvarwidth
\end{varwidth}%
\\
\sphinxbottomrule
\end{tabulary}
\sphinxtableafterendhook\par
\sphinxattableend\end{savenotes}


\bigskip\hrule\bigskip



\subsection{9.7 📡 标定曲线查看}
\label{\detokenize{chapters/chapter9:id29}}

\subsubsection{9.7.1 标定查看器类型}
\label{\detokenize{chapters/chapter9:id30}}
\sphinxAtStartPar
MTDP提供两种标定查看器，满足不同用户的需求：


\begin{savenotes}\sphinxattablestart
\sphinxthistablewithglobalstyle
\centering
\begin{tabulary}{\linewidth}[t]{TTT}
\sphinxtoprule
\begin{varwidth}[t]{\sphinxcolwidth{1}{3}}
\sphinxstyletheadfamily \sphinxAtStartPar
查看器类型
\sphinxbeforeendvarwidth
\end{varwidth}%
&\begin{varwidth}[t]{\sphinxcolwidth{1}{3}}
\sphinxstyletheadfamily \sphinxAtStartPar
用途
\sphinxbeforeendvarwidth
\end{varwidth}%
&\begin{varwidth}[t]{\sphinxcolwidth{1}{3}}
\sphinxstyletheadfamily \sphinxAtStartPar
特点
\sphinxbeforeendvarwidth
\end{varwidth}%
\\
\sphinxmidrule
\sphinxtableatstartofbodyhook\begin{varwidth}[t]{\sphinxcolwidth{1}{3}}
\sphinxAtStartPar
通用标定查看器
\sphinxbeforeendvarwidth
\end{varwidth}%
&\begin{varwidth}[t]{\sphinxcolwidth{1}{3}}
\sphinxAtStartPar
多种格式的标定数据
\sphinxbeforeendvarwidth
\end{varwidth}%
&\begin{varwidth}[t]{\sphinxcolwidth{1}{3}}
\sphinxAtStartPar
支持多种仪器格式，灵活通用
\sphinxbeforeendvarwidth
\end{varwidth}%
\\
\sphinxhline\begin{varwidth}[t]{\sphinxcolwidth{1}{3}}
\sphinxAtStartPar
Phoenix标定查看器
\sphinxbeforeendvarwidth
\end{varwidth}%
&\begin{varwidth}[t]{\sphinxcolwidth{1}{3}}
\sphinxAtStartPar
Phoenix仪器专用
\sphinxbeforeendvarwidth
\end{varwidth}%
&\begin{varwidth}[t]{\sphinxcolwidth{1}{3}}
\sphinxAtStartPar
专为CLB/CLC文件优化，功能完整
\sphinxbeforeendvarwidth
\end{varwidth}%
\\
\sphinxbottomrule
\end{tabulary}
\sphinxtableafterendhook\par
\sphinxattableend\end{savenotes}


\subsubsection{9.7.2 通用标定查看器 (CalCalibrationForm)}
\label{\detokenize{chapters/chapter9:calcalibrationform}}
\sphinxAtStartPar
通用标定查看器集成在主界面中，支持多种格式的标定数据查看和分析。

\sphinxAtStartPar
\sphinxstylestrong{选项卡结构：}


\begin{savenotes}\sphinxattablestart
\sphinxthistablewithglobalstyle
\centering
\begin{tabulary}{\linewidth}[t]{TT}
\sphinxtoprule
\begin{varwidth}[t]{\sphinxcolwidth{1}{2}}
\sphinxstyletheadfamily \sphinxAtStartPar
选项卡
\sphinxbeforeendvarwidth
\end{varwidth}%
&\begin{varwidth}[t]{\sphinxcolwidth{1}{2}}
\sphinxstyletheadfamily \sphinxAtStartPar
功能
\sphinxbeforeendvarwidth
\end{varwidth}%
\\
\sphinxmidrule
\sphinxtableatstartofbodyhook\begin{varwidth}[t]{\sphinxcolwidth{1}{2}}
\sphinxAtStartPar
通道
\sphinxbeforeendvarwidth
\end{varwidth}%
&\begin{varwidth}[t]{\sphinxcolwidth{1}{2}}
\sphinxAtStartPar
通道级标定查看，分析采集盒各通道的响应
\sphinxbeforeendvarwidth
\end{varwidth}%
\\
\sphinxhline\begin{varwidth}[t]{\sphinxcolwidth{1}{2}}
\sphinxAtStartPar
传感器
\sphinxbeforeendvarwidth
\end{varwidth}%
&\begin{varwidth}[t]{\sphinxcolwidth{1}{2}}
\sphinxAtStartPar
传感器级标定查看，分析磁传感器等的响应
\sphinxbeforeendvarwidth
\end{varwidth}%
\\
\sphinxbottomrule
\end{tabulary}
\sphinxtableafterendhook\par
\sphinxattableend\end{savenotes}

\sphinxAtStartPar
\sphinxstylestrong{主要功能：}


\begin{savenotes}\sphinxattablestart
\sphinxthistablewithglobalstyle
\centering
\begin{tabulary}{\linewidth}[t]{TT}
\sphinxtoprule
\begin{varwidth}[t]{\sphinxcolwidth{1}{2}}
\sphinxstyletheadfamily \sphinxAtStartPar
功能
\sphinxbeforeendvarwidth
\end{varwidth}%
&\begin{varwidth}[t]{\sphinxcolwidth{1}{2}}
\sphinxstyletheadfamily \sphinxAtStartPar
说明
\sphinxbeforeendvarwidth
\end{varwidth}%
\\
\sphinxmidrule
\sphinxtableatstartofbodyhook\begin{varwidth}[t]{\sphinxcolwidth{1}{2}}
\sphinxAtStartPar
参考站点选择
\sphinxbeforeendvarwidth
\end{varwidth}%
&\begin{varwidth}[t]{\sphinxcolwidth{1}{2}}
\sphinxAtStartPar
选择参考测点用于标定对比
\sphinxbeforeendvarwidth
\end{varwidth}%
\\
\sphinxhline\begin{varwidth}[t]{\sphinxcolwidth{1}{2}}
\sphinxAtStartPar
幅度响应曲线
\sphinxbeforeendvarwidth
\end{varwidth}%
&\begin{varwidth}[t]{\sphinxcolwidth{1}{2}}
\sphinxAtStartPar
显示幅度随频率的变化
\sphinxbeforeendvarwidth
\end{varwidth}%
\\
\sphinxhline\begin{varwidth}[t]{\sphinxcolwidth{1}{2}}
\sphinxAtStartPar
相位响应曲线
\sphinxbeforeendvarwidth
\end{varwidth}%
&\begin{varwidth}[t]{\sphinxcolwidth{1}{2}}
\sphinxAtStartPar
显示相位随频率的变化
\sphinxbeforeendvarwidth
\end{varwidth}%
\\
\sphinxhline\begin{varwidth}[t]{\sphinxcolwidth{1}{2}}
\sphinxAtStartPar
计算并保存
\sphinxbeforeendvarwidth
\end{varwidth}%
&\begin{varwidth}[t]{\sphinxcolwidth{1}{2}}
\sphinxAtStartPar
计算标定结果并保存到项目
\sphinxbeforeendvarwidth
\end{varwidth}%
\\
\sphinxbottomrule
\end{tabulary}
\sphinxtableafterendhook\par
\sphinxattableend\end{savenotes}

\sphinxAtStartPar
\sphinxstylestrong{操作步骤：}
\begin{enumerate}
\sphinxsetlistlabels{\arabic}{enumi}{enumii}{}{.}%
\item {} 
\sphinxAtStartPar
在主界面切换到 \sphinxcode{\sphinxupquote{标定}} 选项卡

\item {} 
\sphinxAtStartPar
根据需要选择 \sphinxcode{\sphinxupquote{采集盒}} 或 \sphinxcode{\sphinxupquote{传感器}} 子选项卡

\item {} 
\sphinxAtStartPar
右键点击空白区域 → \sphinxcode{\sphinxupquote{添加标定曲线}}

\item {} 
\sphinxAtStartPar
在文件对话框中选择标定文件

\item {} 
\sphinxAtStartPar
系统自动加载并显示幅度/相位曲线

\item {} 
\sphinxAtStartPar
可继续添加多条曲线进行对比

\end{enumerate}

\sphinxAtStartPar
\sphinxstylestrong{右键菜单选项：}


\begin{savenotes}\sphinxattablestart
\sphinxthistablewithglobalstyle
\centering
\begin{tabulary}{\linewidth}[t]{TT}
\sphinxtoprule
\begin{varwidth}[t]{\sphinxcolwidth{1}{2}}
\sphinxstyletheadfamily \sphinxAtStartPar
菜单项
\sphinxbeforeendvarwidth
\end{varwidth}%
&\begin{varwidth}[t]{\sphinxcolwidth{1}{2}}
\sphinxstyletheadfamily \sphinxAtStartPar
功能
\sphinxbeforeendvarwidth
\end{varwidth}%
\\
\sphinxmidrule
\sphinxtableatstartofbodyhook\begin{varwidth}[t]{\sphinxcolwidth{1}{2}}
\sphinxAtStartPar
添加标定曲线
\sphinxbeforeendvarwidth
\end{varwidth}%
&\begin{varwidth}[t]{\sphinxcolwidth{1}{2}}
\sphinxAtStartPar
加载新的标定文件
\sphinxbeforeendvarwidth
\end{varwidth}%
\\
\sphinxhline\begin{varwidth}[t]{\sphinxcolwidth{1}{2}}
\sphinxAtStartPar
清除选中
\sphinxbeforeendvarwidth
\end{varwidth}%
&\begin{varwidth}[t]{\sphinxcolwidth{1}{2}}
\sphinxAtStartPar
清除当前选中的曲线
\sphinxbeforeendvarwidth
\end{varwidth}%
\\
\sphinxhline\begin{varwidth}[t]{\sphinxcolwidth{1}{2}}
\sphinxAtStartPar
清除全部
\sphinxbeforeendvarwidth
\end{varwidth}%
&\begin{varwidth}[t]{\sphinxcolwidth{1}{2}}
\sphinxAtStartPar
清除所有已加载的曲线
\sphinxbeforeendvarwidth
\end{varwidth}%
\\
\sphinxhline\begin{varwidth}[t]{\sphinxcolwidth{1}{2}}
\sphinxAtStartPar
导出图像
\sphinxbeforeendvarwidth
\end{varwidth}%
&\begin{varwidth}[t]{\sphinxcolwidth{1}{2}}
\sphinxAtStartPar
将当前图表导出为图片
\sphinxbeforeendvarwidth
\end{varwidth}%
\\
\sphinxhline\begin{varwidth}[t]{\sphinxcolwidth{1}{2}}
\sphinxAtStartPar
复制数据
\sphinxbeforeendvarwidth
\end{varwidth}%
&\begin{varwidth}[t]{\sphinxcolwidth{1}{2}}
\sphinxAtStartPar
复制标定数据到剪贴板
\sphinxbeforeendvarwidth
\end{varwidth}%
\\
\sphinxbottomrule
\end{tabulary}
\sphinxtableafterendhook\par
\sphinxattableend\end{savenotes}


\subsubsection{9.7.3 Phoenix标定查看器 (CalibrationViewForm)}
\label{\detokenize{chapters/chapter9:phoenix-calibrationviewform}}
\sphinxAtStartPar
Phoenix标定查看器是专门为Phoenix仪器设计的独立工具，提供更专业的标定数据分析功能。

\sphinxAtStartPar
\sphinxstylestrong{支持的文件格式：}


\begin{savenotes}\sphinxattablestart
\sphinxthistablewithglobalstyle
\centering
\begin{tabulary}{\linewidth}[t]{TTT}
\sphinxtoprule
\begin{varwidth}[t]{\sphinxcolwidth{1}{3}}
\sphinxstyletheadfamily \sphinxAtStartPar
格式
\sphinxbeforeendvarwidth
\end{varwidth}%
&\begin{varwidth}[t]{\sphinxcolwidth{1}{3}}
\sphinxstyletheadfamily \sphinxAtStartPar
说明
\sphinxbeforeendvarwidth
\end{varwidth}%
&\begin{varwidth}[t]{\sphinxcolwidth{1}{3}}
\sphinxstyletheadfamily \sphinxAtStartPar
内容
\sphinxbeforeendvarwidth
\end{varwidth}%
\\
\sphinxmidrule
\sphinxtableatstartofbodyhook\begin{varwidth}[t]{\sphinxcolwidth{1}{3}}
\sphinxAtStartPar
CLB
\sphinxbeforeendvarwidth
\end{varwidth}%
&\begin{varwidth}[t]{\sphinxcolwidth{1}{3}}
\sphinxAtStartPar
电场盒标定文件
\sphinxbeforeendvarwidth
\end{varwidth}%
&\begin{varwidth}[t]{\sphinxcolwidth{1}{3}}
\sphinxAtStartPar
电场通道的幅度和相位响应
\sphinxbeforeendvarwidth
\end{varwidth}%
\\
\sphinxhline\begin{varwidth}[t]{\sphinxcolwidth{1}{3}}
\sphinxAtStartPar
CLC
\sphinxbeforeendvarwidth
\end{varwidth}%
&\begin{varwidth}[t]{\sphinxcolwidth{1}{3}}
\sphinxAtStartPar
磁传感器标定文件
\sphinxbeforeendvarwidth
\end{varwidth}%
&\begin{varwidth}[t]{\sphinxcolwidth{1}{3}}
\sphinxAtStartPar
磁场通道的幅度和相位响应
\sphinxbeforeendvarwidth
\end{varwidth}%
\\
\sphinxhline\begin{varwidth}[t]{\sphinxcolwidth{1}{3}}
\sphinxAtStartPar
Site
\sphinxbeforeendvarwidth
\end{varwidth}%
&\begin{varwidth}[t]{\sphinxcolwidth{1}{3}}
\sphinxAtStartPar
测点标定数据
\sphinxbeforeendvarwidth
\end{varwidth}%
&\begin{varwidth}[t]{\sphinxcolwidth{1}{3}}
\sphinxAtStartPar
测点级别的综合标定信息
\sphinxbeforeendvarwidth
\end{varwidth}%
\\
\sphinxbottomrule
\end{tabulary}
\sphinxtableafterendhook\par
\sphinxattableend\end{savenotes}

\sphinxAtStartPar
\sphinxstylestrong{打开Phoenix标定查看器：}
\begin{enumerate}
\sphinxsetlistlabels{\arabic}{enumi}{enumii}{}{.}%
\item {} 
\sphinxAtStartPar
选择菜单 \sphinxcode{\sphinxupquote{工具 → 标定查看}}

\item {} 
\sphinxAtStartPar
或使用快捷键（如果已配置）

\end{enumerate}

\sphinxAtStartPar
\sphinxstylestrong{操作步骤：}
\begin{enumerate}
\sphinxsetlistlabels{\arabic}{enumi}{enumii}{}{.}%
\item {} 
\sphinxAtStartPar
打开Phoenix标定查看器

\item {} 
\sphinxAtStartPar
点击 \sphinxcode{\sphinxupquote{打开}} 按钮或使用 \sphinxcode{\sphinxupquote{文件 → 打开}}

\item {} 
\sphinxAtStartPar
在文件对话框中选择标定文件（.clb、.clc等）

\item {} 
\sphinxAtStartPar
系统自动解析并显示标定曲线

\item {} 
\sphinxAtStartPar
查看幅度和相位响应

\end{enumerate}

\sphinxAtStartPar
\sphinxstylestrong{右键菜单操作：}


\begin{savenotes}\sphinxattablestart
\sphinxthistablewithglobalstyle
\centering
\begin{tabulary}{\linewidth}[t]{TT}
\sphinxtoprule
\begin{varwidth}[t]{\sphinxcolwidth{1}{2}}
\sphinxstyletheadfamily \sphinxAtStartPar
操作
\sphinxbeforeendvarwidth
\end{varwidth}%
&\begin{varwidth}[t]{\sphinxcolwidth{1}{2}}
\sphinxstyletheadfamily \sphinxAtStartPar
说明
\sphinxbeforeendvarwidth
\end{varwidth}%
\\
\sphinxmidrule
\sphinxtableatstartofbodyhook\begin{varwidth}[t]{\sphinxcolwidth{1}{2}}
\sphinxAtStartPar
加载其他标定
\sphinxbeforeendvarwidth
\end{varwidth}%
&\begin{varwidth}[t]{\sphinxcolwidth{1}{2}}
\sphinxAtStartPar
在当前视图中叠加其他标定文件
\sphinxbeforeendvarwidth
\end{varwidth}%
\\
\sphinxhline\begin{varwidth}[t]{\sphinxcolwidth{1}{2}}
\sphinxAtStartPar
清除当前标定
\sphinxbeforeendvarwidth
\end{varwidth}%
&\begin{varwidth}[t]{\sphinxcolwidth{1}{2}}
\sphinxAtStartPar
移除当前选中的标定曲线
\sphinxbeforeendvarwidth
\end{varwidth}%
\\
\sphinxhline\begin{varwidth}[t]{\sphinxcolwidth{1}{2}}
\sphinxAtStartPar
清除所有标定
\sphinxbeforeendvarwidth
\end{varwidth}%
&\begin{varwidth}[t]{\sphinxcolwidth{1}{2}}
\sphinxAtStartPar
移除所有已加载的标定曲线
\sphinxbeforeendvarwidth
\end{varwidth}%
\\
\sphinxhline\begin{varwidth}[t]{\sphinxcolwidth{1}{2}}
\sphinxAtStartPar
导出标定数据
\sphinxbeforeendvarwidth
\end{varwidth}%
&\begin{varwidth}[t]{\sphinxcolwidth{1}{2}}
\sphinxAtStartPar
将标定数据导出为CSV或其他格式
\sphinxbeforeendvarwidth
\end{varwidth}%
\\
\sphinxhline\begin{varwidth}[t]{\sphinxcolwidth{1}{2}}
\sphinxAtStartPar
打印
\sphinxbeforeendvarwidth
\end{varwidth}%
&\begin{varwidth}[t]{\sphinxcolwidth{1}{2}}
\sphinxAtStartPar
打印当前标定曲线图
\sphinxbeforeendvarwidth
\end{varwidth}%
\\
\sphinxbottomrule
\end{tabulary}
\sphinxtableafterendhook\par
\sphinxattableend\end{savenotes}

\sphinxAtStartPar
\sphinxstylestrong{显示选项：}


\begin{savenotes}\sphinxattablestart
\sphinxthistablewithglobalstyle
\centering
\begin{tabulary}{\linewidth}[t]{TT}
\sphinxtoprule
\begin{varwidth}[t]{\sphinxcolwidth{1}{2}}
\sphinxstyletheadfamily \sphinxAtStartPar
选项
\sphinxbeforeendvarwidth
\end{varwidth}%
&\begin{varwidth}[t]{\sphinxcolwidth{1}{2}}
\sphinxstyletheadfamily \sphinxAtStartPar
说明
\sphinxbeforeendvarwidth
\end{varwidth}%
\\
\sphinxmidrule
\sphinxtableatstartofbodyhook\begin{varwidth}[t]{\sphinxcolwidth{1}{2}}
\sphinxAtStartPar
显示幅度
\sphinxbeforeendvarwidth
\end{varwidth}%
&\begin{varwidth}[t]{\sphinxcolwidth{1}{2}}
\sphinxAtStartPar
显示/隐藏幅度曲线
\sphinxbeforeendvarwidth
\end{varwidth}%
\\
\sphinxhline\begin{varwidth}[t]{\sphinxcolwidth{1}{2}}
\sphinxAtStartPar
显示相位
\sphinxbeforeendvarwidth
\end{varwidth}%
&\begin{varwidth}[t]{\sphinxcolwidth{1}{2}}
\sphinxAtStartPar
显示/隐藏相位曲线
\sphinxbeforeendvarwidth
\end{varwidth}%
\\
\sphinxhline\begin{varwidth}[t]{\sphinxcolwidth{1}{2}}
\sphinxAtStartPar
对数坐标
\sphinxbeforeendvarwidth
\end{varwidth}%
&\begin{varwidth}[t]{\sphinxcolwidth{1}{2}}
\sphinxAtStartPar
X轴使用对数坐标
\sphinxbeforeendvarwidth
\end{varwidth}%
\\
\sphinxhline\begin{varwidth}[t]{\sphinxcolwidth{1}{2}}
\sphinxAtStartPar
网格线
\sphinxbeforeendvarwidth
\end{varwidth}%
&\begin{varwidth}[t]{\sphinxcolwidth{1}{2}}
\sphinxAtStartPar
显示/隐藏网格
\sphinxbeforeendvarwidth
\end{varwidth}%
\\
\sphinxhline\begin{varwidth}[t]{\sphinxcolwidth{1}{2}}
\sphinxAtStartPar
图例
\sphinxbeforeendvarwidth
\end{varwidth}%
&\begin{varwidth}[t]{\sphinxcolwidth{1}{2}}
\sphinxAtStartPar
显示/隐藏图例
\sphinxbeforeendvarwidth
\end{varwidth}%
\\
\sphinxbottomrule
\end{tabulary}
\sphinxtableafterendhook\par
\sphinxattableend\end{savenotes}


\subsubsection{9.7.4 标定曲线解读}
\label{\detokenize{chapters/chapter9:id31}}
\sphinxAtStartPar
正确解读标定曲线对于评估仪器性能和数据质量至关重要。


\paragraph{幅度曲线解读}
\label{\detokenize{chapters/chapter9:id32}}
\sphinxAtStartPar
\sphinxstylestrong{坐标轴说明：}
\begin{itemize}
\item {} 
\sphinxAtStartPar
X轴：频率（对数坐标）

\item {} 
\sphinxAtStartPar
Y轴：幅度（dB或线性单位）

\end{itemize}

\sphinxAtStartPar
\sphinxstylestrong{曲线特征分析：}


\begin{savenotes}\sphinxattablestart
\sphinxthistablewithglobalstyle
\centering
\begin{tabulary}{\linewidth}[t]{TTT}
\sphinxtoprule
\begin{varwidth}[t]{\sphinxcolwidth{1}{3}}
\sphinxstyletheadfamily \sphinxAtStartPar
曲线特征
\sphinxbeforeendvarwidth
\end{varwidth}%
&\begin{varwidth}[t]{\sphinxcolwidth{1}{3}}
\sphinxstyletheadfamily \sphinxAtStartPar
含义
\sphinxbeforeendvarwidth
\end{varwidth}%
&\begin{varwidth}[t]{\sphinxcolwidth{1}{3}}
\sphinxstyletheadfamily \sphinxAtStartPar
说明
\sphinxbeforeendvarwidth
\end{varwidth}%
\\
\sphinxmidrule
\sphinxtableatstartofbodyhook\begin{varwidth}[t]{\sphinxcolwidth{1}{3}}
\sphinxAtStartPar
平坦区域
\sphinxbeforeendvarwidth
\end{varwidth}%
&\begin{varwidth}[t]{\sphinxcolwidth{1}{3}}
\sphinxAtStartPar
频率响应稳定
\sphinxbeforeendvarwidth
\end{varwidth}%
&\begin{varwidth}[t]{\sphinxcolwidth{1}{3}}
\sphinxAtStartPar
该频段内仪器响应均匀
\sphinxbeforeendvarwidth
\end{varwidth}%
\\
\sphinxhline\begin{varwidth}[t]{\sphinxcolwidth{1}{3}}
\sphinxAtStartPar
下降区域
\sphinxbeforeendvarwidth
\end{varwidth}%
&\begin{varwidth}[t]{\sphinxcolwidth{1}{3}}
\sphinxAtStartPar
高频衰减
\sphinxbeforeendvarwidth
\end{varwidth}%
&\begin{varwidth}[t]{\sphinxcolwidth{1}{3}}
\sphinxAtStartPar
传感器在高频端灵敏度下降
\sphinxbeforeendvarwidth
\end{varwidth}%
\\
\sphinxhline\begin{varwidth}[t]{\sphinxcolwidth{1}{3}}
\sphinxAtStartPar
上升区域
\sphinxbeforeendvarwidth
\end{varwidth}%
&\begin{varwidth}[t]{\sphinxcolwidth{1}{3}}
\sphinxAtStartPar
低频提升
\sphinxbeforeendvarwidth
\end{varwidth}%
&\begin{varwidth}[t]{\sphinxcolwidth{1}{3}}
\sphinxAtStartPar
可能存在低频补偿
\sphinxbeforeendvarwidth
\end{varwidth}%
\\
\sphinxhline\begin{varwidth}[t]{\sphinxcolwidth{1}{3}}
\sphinxAtStartPar
波动
\sphinxbeforeendvarwidth
\end{varwidth}%
&\begin{varwidth}[t]{\sphinxcolwidth{1}{3}}
\sphinxAtStartPar
共振或不稳定
\sphinxbeforeendvarwidth
\end{varwidth}%
&\begin{varwidth}[t]{\sphinxcolwidth{1}{3}}
\sphinxAtStartPar
可能存在机械共振或电路问题
\sphinxbeforeendvarwidth
\end{varwidth}%
\\
\sphinxbottomrule
\end{tabulary}
\sphinxtableafterendhook\par
\sphinxattableend\end{savenotes}

\sphinxAtStartPar
\sphinxstylestrong{理想曲线特征：}
\begin{itemize}
\item {} 
\sphinxAtStartPar
在工作频段内幅度平坦

\item {} 
\sphinxAtStartPar
过渡频段平滑下降

\item {} 
\sphinxAtStartPar
无明显波动或毛刺

\end{itemize}


\paragraph{相位曲线解读}
\label{\detokenize{chapters/chapter9:id33}}
\sphinxAtStartPar
\sphinxstylestrong{坐标轴说明：}
\begin{itemize}
\item {} 
\sphinxAtStartPar
X轴：频率（对数坐标）

\item {} 
\sphinxAtStartPar
Y轴：相位（度）

\end{itemize}

\sphinxAtStartPar
\sphinxstylestrong{曲线特征分析：}


\begin{savenotes}\sphinxattablestart
\sphinxthistablewithglobalstyle
\centering
\begin{tabulary}{\linewidth}[t]{TTT}
\sphinxtoprule
\begin{varwidth}[t]{\sphinxcolwidth{1}{3}}
\sphinxstyletheadfamily \sphinxAtStartPar
曲线特征
\sphinxbeforeendvarwidth
\end{varwidth}%
&\begin{varwidth}[t]{\sphinxcolwidth{1}{3}}
\sphinxstyletheadfamily \sphinxAtStartPar
含义
\sphinxbeforeendvarwidth
\end{varwidth}%
&\begin{varwidth}[t]{\sphinxcolwidth{1}{3}}
\sphinxstyletheadfamily \sphinxAtStartPar
说明
\sphinxbeforeendvarwidth
\end{varwidth}%
\\
\sphinxmidrule
\sphinxtableatstartofbodyhook\begin{varwidth}[t]{\sphinxcolwidth{1}{3}}
\sphinxAtStartPar
线性区域
\sphinxbeforeendvarwidth
\end{varwidth}%
&\begin{varwidth}[t]{\sphinxcolwidth{1}{3}}
\sphinxAtStartPar
相位响应良好
\sphinxbeforeendvarwidth
\end{varwidth}%
&\begin{varwidth}[t]{\sphinxcolwidth{1}{3}}
\sphinxAtStartPar
相位随频率线性变化
\sphinxbeforeendvarwidth
\end{varwidth}%
\\
\sphinxhline\begin{varwidth}[t]{\sphinxcolwidth{1}{3}}
\sphinxAtStartPar
阶跃变化
\sphinxbeforeendvarwidth
\end{varwidth}%
&\begin{varwidth}[t]{\sphinxcolwidth{1}{3}}
\sphinxAtStartPar
相位跳变
\sphinxbeforeendvarwidth
\end{varwidth}%
&\begin{varwidth}[t]{\sphinxcolwidth{1}{3}}
\sphinxAtStartPar
可能存在滤波器响应
\sphinxbeforeendvarwidth
\end{varwidth}%
\\
\sphinxhline\begin{varwidth}[t]{\sphinxcolwidth{1}{3}}
\sphinxAtStartPar
波动
\sphinxbeforeendvarwidth
\end{varwidth}%
&\begin{varwidth}[t]{\sphinxcolwidth{1}{3}}
\sphinxAtStartPar
非线性响应
\sphinxbeforeendvarwidth
\end{varwidth}%
&\begin{varwidth}[t]{\sphinxcolwidth{1}{3}}
\sphinxAtStartPar
可能影响信号重构
\sphinxbeforeendvarwidth
\end{varwidth}%
\\
\sphinxhline\begin{varwidth}[t]{\sphinxcolwidth{1}{3}}
\sphinxAtStartPar
趋于0°
\sphinxbeforeendvarwidth
\end{varwidth}%
&\begin{varwidth}[t]{\sphinxcolwidth{1}{3}}
\sphinxAtStartPar
最小相位系统
\sphinxbeforeendvarwidth
\end{varwidth}%
&\begin{varwidth}[t]{\sphinxcolwidth{1}{3}}
\sphinxAtStartPar
理想响应特征
\sphinxbeforeendvarwidth
\end{varwidth}%
\\
\sphinxbottomrule
\end{tabulary}
\sphinxtableafterendhook\par
\sphinxattableend\end{savenotes}

\sphinxAtStartPar
\sphinxstylestrong{注意事项：}
\begin{itemize}
\item {} 
\sphinxAtStartPar
相位曲线应在工作频段内连续

\item {} 
\sphinxAtStartPar
避免出现大的相位跳变

\item {} 
\sphinxAtStartPar
线性度越好，信号还原越准确

\end{itemize}


\paragraph{多曲线对比}
\label{\detokenize{chapters/chapter9:id34}}
\sphinxAtStartPar
当同时加载多条标定曲线时：
\begin{enumerate}
\sphinxsetlistlabels{\arabic}{enumi}{enumii}{}{.}%
\item {} 
\sphinxAtStartPar
\sphinxstylestrong{同类型对比}：比较同一类型不同传感器的响应一致性

\item {} 
\sphinxAtStartPar
\sphinxstylestrong{时序对比}：比较同一传感器不同时间的标定，评估稳定性

\item {} 
\sphinxAtStartPar
\sphinxstylestrong{频率范围评估}：确定各传感器的有效工作频段

\end{enumerate}

\sphinxAtStartPar
\sphinxstylestrong{对比分析要点：}
\begin{itemize}
\item {} 
\sphinxAtStartPar
曲线重合度：一致性好的曲线应高度重合

\item {} 
\sphinxAtStartPar
偏差分析：记录偏差较大的频段

\item {} 
\sphinxAtStartPar
长期稳定性：定期标定对比，监控仪器漂移

\end{itemize}


\subsection{9.8 🧵 多线程监控}
\label{\detokenize{chapters/chapter9:id35}}

\subsubsection{9.8.1 功能说明}
\label{\detokenize{chapters/chapter9:id36}}
\sphinxAtStartPar
多线程监控选项卡显示后台处理任务的状态。


\subsubsection{9.8.2 监控内容}
\label{\detokenize{chapters/chapter9:id37}}

\begin{savenotes}\sphinxattablestart
\sphinxthistablewithglobalstyle
\centering
\begin{tabulary}{\linewidth}[t]{TT}
\sphinxtoprule
\begin{varwidth}[t]{\sphinxcolwidth{1}{2}}
\sphinxstyletheadfamily \sphinxAtStartPar
信息
\sphinxbeforeendvarwidth
\end{varwidth}%
&\begin{varwidth}[t]{\sphinxcolwidth{1}{2}}
\sphinxstyletheadfamily \sphinxAtStartPar
说明
\sphinxbeforeendvarwidth
\end{varwidth}%
\\
\sphinxmidrule
\sphinxtableatstartofbodyhook\begin{varwidth}[t]{\sphinxcolwidth{1}{2}}
\sphinxAtStartPar
运行线程数
\sphinxbeforeendvarwidth
\end{varwidth}%
&\begin{varwidth}[t]{\sphinxcolwidth{1}{2}}
\sphinxAtStartPar
当前活动的处理线程数
\sphinxbeforeendvarwidth
\end{varwidth}%
\\
\sphinxhline\begin{varwidth}[t]{\sphinxcolwidth{1}{2}}
\sphinxAtStartPar
处理进度
\sphinxbeforeendvarwidth
\end{varwidth}%
&\begin{varwidth}[t]{\sphinxcolwidth{1}{2}}
\sphinxAtStartPar
当前任务的完成百分比
\sphinxbeforeendvarwidth
\end{varwidth}%
\\
\sphinxhline\begin{varwidth}[t]{\sphinxcolwidth{1}{2}}
\sphinxAtStartPar
当前任务
\sphinxbeforeendvarwidth
\end{varwidth}%
&\begin{varwidth}[t]{\sphinxcolwidth{1}{2}}
\sphinxAtStartPar
正在执行的操作
\sphinxbeforeendvarwidth
\end{varwidth}%
\\
\sphinxhline\begin{varwidth}[t]{\sphinxcolwidth{1}{2}}
\sphinxAtStartPar
状态信息
\sphinxbeforeendvarwidth
\end{varwidth}%
&\begin{varwidth}[t]{\sphinxcolwidth{1}{2}}
\sphinxAtStartPar
任务执行状态
\sphinxbeforeendvarwidth
\end{varwidth}%
\\
\sphinxbottomrule
\end{tabulary}
\sphinxtableafterendhook\par
\sphinxattableend\end{savenotes}


\subsubsection{9.8.3 任务管理}
\label{\detokenize{chapters/chapter9:id38}}\begin{itemize}
\item {} 
\sphinxAtStartPar
查看队列中的待处理任务

\item {} 
\sphinxAtStartPar
监控长时间任务进度

\item {} 
\sphinxAtStartPar
识别阻塞任务

\end{itemize}


\bigskip\hrule\bigskip



\subsection{9.9 工程树操作}
\label{\detokenize{chapters/chapter9:id39}}

\subsubsection{9.9.1 快捷操作}
\label{\detokenize{chapters/chapter9:id40}}

\begin{savenotes}\sphinxattablestart
\sphinxthistablewithglobalstyle
\centering
\begin{tabulary}{\linewidth}[t]{TT}
\sphinxtoprule
\begin{varwidth}[t]{\sphinxcolwidth{1}{2}}
\sphinxstyletheadfamily \sphinxAtStartPar
操作
\sphinxbeforeendvarwidth
\end{varwidth}%
&\begin{varwidth}[t]{\sphinxcolwidth{1}{2}}
\sphinxstyletheadfamily \sphinxAtStartPar
方法
\sphinxbeforeendvarwidth
\end{varwidth}%
\\
\sphinxmidrule
\sphinxtableatstartofbodyhook\begin{varwidth}[t]{\sphinxcolwidth{1}{2}}
\sphinxAtStartPar
展开/折叠
\sphinxbeforeendvarwidth
\end{varwidth}%
&\begin{varwidth}[t]{\sphinxcolwidth{1}{2}}
\sphinxAtStartPar
双击或点击+/\sphinxhyphen{}
\sphinxbeforeendvarwidth
\end{varwidth}%
\\
\sphinxhline\begin{varwidth}[t]{\sphinxcolwidth{1}{2}}
\sphinxAtStartPar
多选
\sphinxbeforeendvarwidth
\end{varwidth}%
&\begin{varwidth}[t]{\sphinxcolwidth{1}{2}}
\sphinxAtStartPar
Ctrl+点击
\sphinxbeforeendvarwidth
\end{varwidth}%
\\
\sphinxhline\begin{varwidth}[t]{\sphinxcolwidth{1}{2}}
\sphinxAtStartPar
范围选择
\sphinxbeforeendvarwidth
\end{varwidth}%
&\begin{varwidth}[t]{\sphinxcolwidth{1}{2}}
\sphinxAtStartPar
Shift+点击
\sphinxbeforeendvarwidth
\end{varwidth}%
\\
\sphinxhline\begin{varwidth}[t]{\sphinxcolwidth{1}{2}}
\sphinxAtStartPar
右键菜单
\sphinxbeforeendvarwidth
\end{varwidth}%
&\begin{varwidth}[t]{\sphinxcolwidth{1}{2}}
\sphinxAtStartPar
右键点击
\sphinxbeforeendvarwidth
\end{varwidth}%
\\
\sphinxbottomrule
\end{tabulary}
\sphinxtableafterendhook\par
\sphinxattableend\end{savenotes}


\subsubsection{9.9.2 层次结构}
\label{\detokenize{chapters/chapter9:id41}}
\begin{sphinxVerbatim}[commandchars=\\\{\}]
工程 → 工区 → 测段 → 测点
\end{sphinxVerbatim}


\subsubsection{9.9.3 拖拽操作}
\label{\detokenize{chapters/chapter9:id42}}\begin{itemize}
\item {} 
\sphinxAtStartPar
拖动测点调整位置

\item {} 
\sphinxAtStartPar
拖动测点到其他测段

\end{itemize}


\bigskip\hrule\bigskip



\subsection{9.10 图表通用操作}
\label{\detokenize{chapters/chapter9:id43}}

\subsubsection{9.10.1 鼠标操作}
\label{\detokenize{chapters/chapter9:id44}}

\begin{savenotes}\sphinxattablestart
\sphinxthistablewithglobalstyle
\centering
\begin{tabulary}{\linewidth}[t]{TT}
\sphinxtoprule
\begin{varwidth}[t]{\sphinxcolwidth{1}{2}}
\sphinxstyletheadfamily \sphinxAtStartPar
操作
\sphinxbeforeendvarwidth
\end{varwidth}%
&\begin{varwidth}[t]{\sphinxcolwidth{1}{2}}
\sphinxstyletheadfamily \sphinxAtStartPar
功能
\sphinxbeforeendvarwidth
\end{varwidth}%
\\
\sphinxmidrule
\sphinxtableatstartofbodyhook\begin{varwidth}[t]{\sphinxcolwidth{1}{2}}
\sphinxAtStartPar
左键拖动
\sphinxbeforeendvarwidth
\end{varwidth}%
&\begin{varwidth}[t]{\sphinxcolwidth{1}{2}}
\sphinxAtStartPar
平移图表
\sphinxbeforeendvarwidth
\end{varwidth}%
\\
\sphinxhline\begin{varwidth}[t]{\sphinxcolwidth{1}{2}}
\sphinxAtStartPar
滚轮向上
\sphinxbeforeendvarwidth
\end{varwidth}%
&\begin{varwidth}[t]{\sphinxcolwidth{1}{2}}
\sphinxAtStartPar
放大
\sphinxbeforeendvarwidth
\end{varwidth}%
\\
\sphinxhline\begin{varwidth}[t]{\sphinxcolwidth{1}{2}}
\sphinxAtStartPar
滚轮向下
\sphinxbeforeendvarwidth
\end{varwidth}%
&\begin{varwidth}[t]{\sphinxcolwidth{1}{2}}
\sphinxAtStartPar
缩小
\sphinxbeforeendvarwidth
\end{varwidth}%
\\
\sphinxhline\begin{varwidth}[t]{\sphinxcolwidth{1}{2}}
\sphinxAtStartPar
双击
\sphinxbeforeendvarwidth
\end{varwidth}%
&\begin{varwidth}[t]{\sphinxcolwidth{1}{2}}
\sphinxAtStartPar
恢复原始视图
\sphinxbeforeendvarwidth
\end{varwidth}%
\\
\sphinxhline\begin{varwidth}[t]{\sphinxcolwidth{1}{2}}
\sphinxAtStartPar
右键
\sphinxbeforeendvarwidth
\end{varwidth}%
&\begin{varwidth}[t]{\sphinxcolwidth{1}{2}}
\sphinxAtStartPar
打开上下文菜单
\sphinxbeforeendvarwidth
\end{varwidth}%
\\
\sphinxbottomrule
\end{tabulary}
\sphinxtableafterendhook\par
\sphinxattableend\end{savenotes}


\subsubsection{9.10.2 键盘快捷键}
\label{\detokenize{chapters/chapter9:id45}}

\begin{savenotes}\sphinxattablestart
\sphinxthistablewithglobalstyle
\centering
\begin{tabulary}{\linewidth}[t]{TT}
\sphinxtoprule
\begin{varwidth}[t]{\sphinxcolwidth{1}{2}}
\sphinxstyletheadfamily \sphinxAtStartPar
快捷键
\sphinxbeforeendvarwidth
\end{varwidth}%
&\begin{varwidth}[t]{\sphinxcolwidth{1}{2}}
\sphinxstyletheadfamily \sphinxAtStartPar
功能
\sphinxbeforeendvarwidth
\end{varwidth}%
\\
\sphinxmidrule
\sphinxtableatstartofbodyhook\begin{varwidth}[t]{\sphinxcolwidth{1}{2}}
\sphinxAtStartPar
Ctrl + +
\sphinxbeforeendvarwidth
\end{varwidth}%
&\begin{varwidth}[t]{\sphinxcolwidth{1}{2}}
\sphinxAtStartPar
放大
\sphinxbeforeendvarwidth
\end{varwidth}%
\\
\sphinxhline\begin{varwidth}[t]{\sphinxcolwidth{1}{2}}
\sphinxAtStartPar
Ctrl + \sphinxhyphen{}
\sphinxbeforeendvarwidth
\end{varwidth}%
&\begin{varwidth}[t]{\sphinxcolwidth{1}{2}}
\sphinxAtStartPar
缩小
\sphinxbeforeendvarwidth
\end{varwidth}%
\\
\sphinxhline\begin{varwidth}[t]{\sphinxcolwidth{1}{2}}
\sphinxAtStartPar
Ctrl + 0
\sphinxbeforeendvarwidth
\end{varwidth}%
&\begin{varwidth}[t]{\sphinxcolwidth{1}{2}}
\sphinxAtStartPar
恢复原始大小
\sphinxbeforeendvarwidth
\end{varwidth}%
\\
\sphinxhline\begin{varwidth}[t]{\sphinxcolwidth{1}{2}}
\sphinxAtStartPar
Home
\sphinxbeforeendvarwidth
\end{varwidth}%
&\begin{varwidth}[t]{\sphinxcolwidth{1}{2}}
\sphinxAtStartPar
跳转到开始
\sphinxbeforeendvarwidth
\end{varwidth}%
\\
\sphinxhline\begin{varwidth}[t]{\sphinxcolwidth{1}{2}}
\sphinxAtStartPar
End
\sphinxbeforeendvarwidth
\end{varwidth}%
&\begin{varwidth}[t]{\sphinxcolwidth{1}{2}}
\sphinxAtStartPar
跳转到结束
\sphinxbeforeendvarwidth
\end{varwidth}%
\\
\sphinxbottomrule
\end{tabulary}
\sphinxtableafterendhook\par
\sphinxattableend\end{savenotes}


\bigskip\hrule\bigskip



\subsection{9.11 🔬 SSA奇异谱分析}
\label{\detokenize{chapters/chapter9:ssa}}

\subsubsection{9.11.1 功能说明}
\label{\detokenize{chapters/chapter9:id46}}
\sphinxAtStartPar
奇异谱分析（Singular Spectrum Analysis）是一种强大的时间序列分析方法，用于：
\begin{itemize}
\item {} 
\sphinxAtStartPar
提取时间序列中的趋势成分

\item {} 
\sphinxAtStartPar
识别周期性信号

\item {} 
\sphinxAtStartPar
去除噪声

\item {} 
\sphinxAtStartPar
数据重构

\end{itemize}


\subsubsection{9.11.2 打开SSA分析}
\label{\detokenize{chapters/chapter9:id47}}\begin{enumerate}
\sphinxsetlistlabels{\arabic}{enumi}{enumii}{}{.}%
\item {} 
\sphinxAtStartPar
选择菜单 \sphinxcode{\sphinxupquote{时间序列处理 → SSA/MSSA}}

\item {} 
\sphinxAtStartPar
选择要分析的时间序列

\item {} 
\sphinxAtStartPar
设置SSA参数

\end{enumerate}


\subsubsection{9.11.3 SSA参数}
\label{\detokenize{chapters/chapter9:id48}}

\begin{savenotes}\sphinxattablestart
\sphinxthistablewithglobalstyle
\centering
\begin{tabulary}{\linewidth}[t]{TT}
\sphinxtoprule
\begin{varwidth}[t]{\sphinxcolwidth{1}{2}}
\sphinxstyletheadfamily \sphinxAtStartPar
参数
\sphinxbeforeendvarwidth
\end{varwidth}%
&\begin{varwidth}[t]{\sphinxcolwidth{1}{2}}
\sphinxstyletheadfamily \sphinxAtStartPar
说明
\sphinxbeforeendvarwidth
\end{varwidth}%
\\
\sphinxmidrule
\sphinxtableatstartofbodyhook\begin{varwidth}[t]{\sphinxcolwidth{1}{2}}
\sphinxAtStartPar
窗口长度
\sphinxbeforeendvarwidth
\end{varwidth}%
&\begin{varwidth}[t]{\sphinxcolwidth{1}{2}}
\sphinxAtStartPar
嵌入维度，影响分解精度
\sphinxbeforeendvarwidth
\end{varwidth}%
\\
\sphinxhline\begin{varwidth}[t]{\sphinxcolwidth{1}{2}}
\sphinxAtStartPar
主成分数
\sphinxbeforeendvarwidth
\end{varwidth}%
&\begin{varwidth}[t]{\sphinxcolwidth{1}{2}}
\sphinxAtStartPar
保留的主成分数量
\sphinxbeforeendvarwidth
\end{varwidth}%
\\
\sphinxhline\begin{varwidth}[t]{\sphinxcolwidth{1}{2}}
\sphinxAtStartPar
重构组
\sphinxbeforeendvarwidth
\end{varwidth}%
&\begin{varwidth}[t]{\sphinxcolwidth{1}{2}}
\sphinxAtStartPar
用于重构的成分组
\sphinxbeforeendvarwidth
\end{varwidth}%
\\
\sphinxbottomrule
\end{tabulary}
\sphinxtableafterendhook\par
\sphinxattableend\end{savenotes}


\subsubsection{9.11.4 分析步骤}
\label{\detokenize{chapters/chapter9:id49}}\begin{enumerate}
\sphinxsetlistlabels{\arabic}{enumi}{enumii}{}{.}%
\item {} 
\sphinxAtStartPar
\sphinxstylestrong{加载时间序列}：选择要分析的数据文件

\item {} 
\sphinxAtStartPar
\sphinxstylestrong{设置参数}：配置窗口长度和成分数

\item {} 
\sphinxAtStartPar
\sphinxstylestrong{执行分解}：计算特征值和特征向量

\item {} 
\sphinxAtStartPar
\sphinxstylestrong{选择成分}：查看各成分的贡献率

\item {} 
\sphinxAtStartPar
\sphinxstylestrong{重构信号}：选择成分进行信号重构

\item {} 
\sphinxAtStartPar
\sphinxstylestrong{导出结果}：保存分解和重构结果

\end{enumerate}


\subsubsection{9.11.5 MSSA多通道分析}
\label{\detokenize{chapters/chapter9:mssa}}
\sphinxAtStartPar
多通道奇异谱分析（MSSA）扩展了SSA，可同时分析多个通道：
\begin{itemize}
\item {} 
\sphinxAtStartPar
分析通道间的共同模式

\item {} 
\sphinxAtStartPar
识别相干信号成分

\item {} 
\sphinxAtStartPar
多通道联合去噪

\end{itemize}

\sphinxAtStartPar
\sphinxstylestrong{支持的数据格式：}


\begin{savenotes}\sphinxattablestart
\sphinxthistablewithglobalstyle
\centering
\begin{tabulary}{\linewidth}[t]{TT}
\sphinxtoprule
\begin{varwidth}[t]{\sphinxcolwidth{1}{2}}
\sphinxstyletheadfamily \sphinxAtStartPar
格式
\sphinxbeforeendvarwidth
\end{varwidth}%
&\begin{varwidth}[t]{\sphinxcolwidth{1}{2}}
\sphinxstyletheadfamily \sphinxAtStartPar
说明
\sphinxbeforeendvarwidth
\end{varwidth}%
\\
\sphinxmidrule
\sphinxtableatstartofbodyhook\begin{varwidth}[t]{\sphinxcolwidth{1}{2}}
\sphinxAtStartPar
Phoenix TS2
\sphinxbeforeendvarwidth
\end{varwidth}%
&\begin{varwidth}[t]{\sphinxcolwidth{1}{2}}
\sphinxAtStartPar
24000 Hz 采样率
\sphinxbeforeendvarwidth
\end{varwidth}%
\\
\sphinxhline\begin{varwidth}[t]{\sphinxcolwidth{1}{2}}
\sphinxAtStartPar
Phoenix TS3
\sphinxbeforeendvarwidth
\end{varwidth}%
&\begin{varwidth}[t]{\sphinxcolwidth{1}{2}}
\sphinxAtStartPar
2400 Hz 采样率
\sphinxbeforeendvarwidth
\end{varwidth}%
\\
\sphinxhline\begin{varwidth}[t]{\sphinxcolwidth{1}{2}}
\sphinxAtStartPar
Phoenix TS4
\sphinxbeforeendvarwidth
\end{varwidth}%
&\begin{varwidth}[t]{\sphinxcolwidth{1}{2}}
\sphinxAtStartPar
150 Hz 采样率
\sphinxbeforeendvarwidth
\end{varwidth}%
\\
\sphinxhline\begin{varwidth}[t]{\sphinxcolwidth{1}{2}}
\sphinxAtStartPar
Phoenix TS5
\sphinxbeforeendvarwidth
\end{varwidth}%
&\begin{varwidth}[t]{\sphinxcolwidth{1}{2}}
\sphinxAtStartPar
15 Hz 采样率
\sphinxbeforeendvarwidth
\end{varwidth}%
\\
\sphinxhline\begin{varwidth}[t]{\sphinxcolwidth{1}{2}}
\sphinxAtStartPar
LEMI
\sphinxbeforeendvarwidth
\end{varwidth}%
&\begin{varwidth}[t]{\sphinxcolwidth{1}{2}}
\sphinxAtStartPar
LEMI仪器数据
\sphinxbeforeendvarwidth
\end{varwidth}%
\\
\sphinxbottomrule
\end{tabulary}
\sphinxtableafterendhook\par
\sphinxattableend\end{savenotes}

\sphinxAtStartPar
\sphinxstylestrong{通道选择：}
\begin{itemize}
\item {} 
\sphinxAtStartPar
Ex：X方向电场

\item {} 
\sphinxAtStartPar
Ey：Y方向电场

\item {} 
\sphinxAtStartPar
Hx：X方向磁场

\item {} 
\sphinxAtStartPar
Hy：Y方向磁场

\item {} 
\sphinxAtStartPar
Hz：Z方向磁场

\end{itemize}

\sphinxAtStartPar
\sphinxstylestrong{时间对齐：}

\sphinxAtStartPar
当加载多个测点的时间序列时，可以选择启用时间对齐功能：
\begin{itemize}
\item {} 
\sphinxAtStartPar
自动对齐不同测点的起始时间

\item {} 
\sphinxAtStartPar
使用相同时间段的数据进行分析

\end{itemize}


\subsubsection{9.11.6 结果解读}
\label{\detokenize{chapters/chapter9:id50}}

\begin{savenotes}\sphinxattablestart
\sphinxthistablewithglobalstyle
\centering
\begin{tabulary}{\linewidth}[t]{TT}
\sphinxtoprule
\begin{varwidth}[t]{\sphinxcolwidth{1}{2}}
\sphinxstyletheadfamily \sphinxAtStartPar
成分类型
\sphinxbeforeendvarwidth
\end{varwidth}%
&\begin{varwidth}[t]{\sphinxcolwidth{1}{2}}
\sphinxstyletheadfamily \sphinxAtStartPar
特征
\sphinxbeforeendvarwidth
\end{varwidth}%
\\
\sphinxmidrule
\sphinxtableatstartofbodyhook\begin{varwidth}[t]{\sphinxcolwidth{1}{2}}
\sphinxAtStartPar
趋势成分
\sphinxbeforeendvarwidth
\end{varwidth}%
&\begin{varwidth}[t]{\sphinxcolwidth{1}{2}}
\sphinxAtStartPar
缓慢变化，低频
\sphinxbeforeendvarwidth
\end{varwidth}%
\\
\sphinxhline\begin{varwidth}[t]{\sphinxcolwidth{1}{2}}
\sphinxAtStartPar
周期成分
\sphinxbeforeendvarwidth
\end{varwidth}%
&\begin{varwidth}[t]{\sphinxcolwidth{1}{2}}
\sphinxAtStartPar
特定频率的振荡
\sphinxbeforeendvarwidth
\end{varwidth}%
\\
\sphinxhline\begin{varwidth}[t]{\sphinxcolwidth{1}{2}}
\sphinxAtStartPar
噪声成分
\sphinxbeforeendvarwidth
\end{varwidth}%
&\begin{varwidth}[t]{\sphinxcolwidth{1}{2}}
\sphinxAtStartPar
高频，无规律
\sphinxbeforeendvarwidth
\end{varwidth}%
\\
\sphinxbottomrule
\end{tabulary}
\sphinxtableafterendhook\par
\sphinxattableend\end{savenotes}


\subsubsection{9.11.7 应用场景}
\label{\detokenize{chapters/chapter9:id51}}\begin{itemize}
\item {} 
\sphinxAtStartPar
\sphinxstylestrong{去噪处理}：保留主要成分，去除噪声成分

\item {} 
\sphinxAtStartPar
\sphinxstylestrong{信号分离}：分离不同频率的信号

\item {} 
\sphinxAtStartPar
\sphinxstylestrong{缺失数据插值}：利用SSA重构填补缺失

\item {} 
\sphinxAtStartPar
\sphinxstylestrong{异常检测}：识别异常成分

\end{itemize}


\subsubsection{9.11.8 导出结果}
\label{\detokenize{chapters/chapter9:id52}}
\sphinxAtStartPar
MSSA分析完成后，可以将重构后的数据导出：
\begin{enumerate}
\sphinxsetlistlabels{\arabic}{enumi}{enumii}{}{.}%
\item {} 
\sphinxAtStartPar
点击 \sphinxcode{\sphinxupquote{导出结果}} 按钮

\item {} 
\sphinxAtStartPar
选择保存目录

\item {} 
\sphinxAtStartPar
系统自动将重构后的时间序列保存为原始格式

\item {} 
\sphinxAtStartPar
支持导出为Phoenix TS格式或LEMI格式

\end{enumerate}

\sphinxstepscope


\section{实例展示}
\label{\detokenize{gallery/index:id1}}\label{\detokenize{gallery/index::doc}}
\sphinxAtStartPar
本章节展示 MTDataPro 的典型应用实例和处理效果。


\subsection{图件展示}
\label{\detokenize{gallery/index:id2}}

\subsubsection{视电阻率曲线}
\label{\detokenize{gallery/index:id3}}
\sphinxAtStartPar
典型的 MT 视电阻率和相位曲线：

\begin{sphinxVerbatim}[commandchars=\\\{\}]
ρ\PYGZus{}xy (Ω·m)          φ\PYGZus{}xy (°)
    |                  |
100 ┤    ○           45┤      ○
    |   /               |     /
 10 ┤  ○              0┤    ○
    | /                 |   /
  1 ┤○               \PYGZhy{}45┤  ○
    |\PYGZus{}\PYGZus{}\PYGZus{}\PYGZus{}\PYGZus{}\PYGZus{}\PYGZus{}\PYGZus{}\PYGZus{}\PYGZus{}\PYGZus{}\PYGZus{}\PYGZus{}\PYGZus{}\PYGZus{}\PYGZus{}\PYGZus{}   |\PYGZus{}\PYGZus{}\PYGZus{}\PYGZus{}\PYGZus{}\PYGZus{}\PYGZus{}\PYGZus{}\PYGZus{}\PYGZus{}\PYGZus{}\PYGZus{}\PYGZus{}\PYGZus{}\PYGZus{}\PYGZus{}\PYGZus{}
      10    100  1000      10    100  1000
           T (s)                  T (s)
\end{sphinxVerbatim}


\subsubsection{地图视图}
\label{\detokenize{gallery/index:id4}}
\sphinxAtStartPar
MTDataPro 支持在地图上显示测点位置：
\begin{itemize}
\item {} 
\sphinxAtStartPar
支持导入 KML/GPX 等地理数据

\item {} 
\sphinxAtStartPar
可叠加地形图、地质图

\item {} 
\sphinxAtStartPar
支持测点编辑和导航

\end{itemize}


\subsection{实例列表}
\label{\detokenize{gallery/index:id5}}
\sphinxAtStartPar
更多实例正在整理中，包括：
\begin{itemize}
\item {} 
\sphinxAtStartPar
Phoenix 数据处理流程

\item {} 
\sphinxAtStartPar
Metronix 数据处理

\item {} 
\sphinxAtStartPar
参考道处理技术

\item {} 
\sphinxAtStartPar
批量处理方法

\end{itemize}


\subsection{贡献实例}
\label{\detokenize{gallery/index:id6}}
\sphinxAtStartPar
欢迎贡献您的处理实例！请参考 {\hyperref[\detokenize{appendix/contributing::doc}]{\sphinxcrossref{\DUrole{std}{\DUrole{std-doc}{贡献指南}}}}}。

\sphinxstepscope


\section{MTDataPro 数据导出与 Python 绘图教程}
\label{\detokenize{tutorial/plotting:mtdatapro-python}}\label{\detokenize{tutorial/plotting::doc}}
\sphinxAtStartPar
本章节介绍如何使用 MTDataPro 导出的数据配合 Python 脚本进行高质量科学绘图。


\subsection{数据导出}
\label{\detokenize{tutorial/plotting:id1}}
\sphinxAtStartPar
在使用绘图脚本之前，需要先从 MTDataPro 导出绘图所需的数据文件。


\subsubsection{导出步骤}
\label{\detokenize{tutorial/plotting:id2}}\begin{enumerate}
\sphinxsetlistlabels{\arabic}{enumi}{enumii}{}{.}%
\item {} 
\sphinxAtStartPar
在工程树中选择测段（Group）

\item {} 
\sphinxAtStartPar
右键选择 \sphinxcode{\sphinxupquote{导出绘图数据}}

\item {} 
\sphinxAtStartPar
选择导出内容：
\begin{itemize}
\item {} 
\sphinxAtStartPar
视电阻率数据（ρxy, ρyx）

\item {} 
\sphinxAtStartPar
相位数据（φxy, φyx）

\item {} 
\sphinxAtStartPar
误差数据

\item {} 
\sphinxAtStartPar
倾子数据（Tx, Ty）

\item {} 
\sphinxAtStartPar
相位张量参数（α, β, φmax, φmin）

\end{itemize}

\end{enumerate}


\subsubsection{导出文件命名规则}
\label{\detokenize{tutorial/plotting:id3}}
\sphinxAtStartPar
导出的数据文件遵循统一的命名规则 \sphinxcode{\sphinxupquote{\{测点名\}\sphinxhyphen{}\{参数名\}.dat}}：


\begin{savenotes}\sphinxattablestart
\sphinxthistablewithglobalstyle
\centering
\begin{tabulary}{\linewidth}[t]{TTT}
\sphinxtoprule
\begin{varwidth}[t]{\sphinxcolwidth{1}{3}}
\sphinxstyletheadfamily \sphinxAtStartPar
文件名
\sphinxbeforeendvarwidth
\end{varwidth}%
&\begin{varwidth}[t]{\sphinxcolwidth{1}{3}}
\sphinxstyletheadfamily \sphinxAtStartPar
内容
\sphinxbeforeendvarwidth
\end{varwidth}%
&\begin{varwidth}[t]{\sphinxcolwidth{1}{3}}
\sphinxstyletheadfamily \sphinxAtStartPar
说明
\sphinxbeforeendvarwidth
\end{varwidth}%
\\
\sphinxmidrule
\sphinxtableatstartofbodyhook\begin{varwidth}[t]{\sphinxcolwidth{1}{3}}
\sphinxAtStartPar
\sphinxcode{\sphinxupquote{\{name\}\sphinxhyphen{}rxy.dat}}
\sphinxbeforeendvarwidth
\end{varwidth}%
&\begin{varwidth}[t]{\sphinxcolwidth{1}{3}}
\sphinxAtStartPar
视电阻率 ρxy
\sphinxbeforeendvarwidth
\end{varwidth}%
&\begin{varwidth}[t]{\sphinxcolwidth{1}{3}}
\sphinxAtStartPar
第1列：频率，第2列：数值
\sphinxbeforeendvarwidth
\end{varwidth}%
\\
\sphinxhline\begin{varwidth}[t]{\sphinxcolwidth{1}{3}}
\sphinxAtStartPar
\sphinxcode{\sphinxupquote{\{name\}\sphinxhyphen{}ryx.dat}}
\sphinxbeforeendvarwidth
\end{varwidth}%
&\begin{varwidth}[t]{\sphinxcolwidth{1}{3}}
\sphinxAtStartPar
视电阻率 ρyx
\sphinxbeforeendvarwidth
\end{varwidth}%
&\begin{varwidth}[t]{\sphinxcolwidth{1}{3}}
\sphinxAtStartPar
第1列：频率，第2列：数值
\sphinxbeforeendvarwidth
\end{varwidth}%
\\
\sphinxhline\begin{varwidth}[t]{\sphinxcolwidth{1}{3}}
\sphinxAtStartPar
\sphinxcode{\sphinxupquote{\{name\}\sphinxhyphen{}zxyp.dat}}
\sphinxbeforeendvarwidth
\end{varwidth}%
&\begin{varwidth}[t]{\sphinxcolwidth{1}{3}}
\sphinxAtStartPar
相位 φxy
\sphinxbeforeendvarwidth
\end{varwidth}%
&\begin{varwidth}[t]{\sphinxcolwidth{1}{3}}
\sphinxAtStartPar
第1列：频率，第2列：数值（度）
\sphinxbeforeendvarwidth
\end{varwidth}%
\\
\sphinxhline\begin{varwidth}[t]{\sphinxcolwidth{1}{3}}
\sphinxAtStartPar
\sphinxcode{\sphinxupquote{\{name\}\sphinxhyphen{}zyxp.dat}}
\sphinxbeforeendvarwidth
\end{varwidth}%
&\begin{varwidth}[t]{\sphinxcolwidth{1}{3}}
\sphinxAtStartPar
相位 φyx
\sphinxbeforeendvarwidth
\end{varwidth}%
&\begin{varwidth}[t]{\sphinxcolwidth{1}{3}}
\sphinxAtStartPar
第1列：频率，第2列：数值（度）
\sphinxbeforeendvarwidth
\end{varwidth}%
\\
\sphinxhline\begin{varwidth}[t]{\sphinxcolwidth{1}{3}}
\sphinxAtStartPar
\sphinxcode{\sphinxupquote{\{name\}\sphinxhyphen{}rxyvar.dat}}
\sphinxbeforeendvarwidth
\end{varwidth}%
&\begin{varwidth}[t]{\sphinxcolwidth{1}{3}}
\sphinxAtStartPar
ρxy 方差
\sphinxbeforeendvarwidth
\end{varwidth}%
&\begin{varwidth}[t]{\sphinxcolwidth{1}{3}}
\sphinxAtStartPar
误差数据
\sphinxbeforeendvarwidth
\end{varwidth}%
\\
\sphinxhline\begin{varwidth}[t]{\sphinxcolwidth{1}{3}}
\sphinxAtStartPar
\sphinxcode{\sphinxupquote{\{name\}\sphinxhyphen{}ryxvar.dat}}
\sphinxbeforeendvarwidth
\end{varwidth}%
&\begin{varwidth}[t]{\sphinxcolwidth{1}{3}}
\sphinxAtStartPar
ρyx 方差
\sphinxbeforeendvarwidth
\end{varwidth}%
&\begin{varwidth}[t]{\sphinxcolwidth{1}{3}}
\sphinxAtStartPar
误差数据
\sphinxbeforeendvarwidth
\end{varwidth}%
\\
\sphinxhline\begin{varwidth}[t]{\sphinxcolwidth{1}{3}}
\sphinxAtStartPar
\sphinxcode{\sphinxupquote{\{name\}\sphinxhyphen{}tzxr.dat}}
\sphinxbeforeendvarwidth
\end{varwidth}%
&\begin{varwidth}[t]{\sphinxcolwidth{1}{3}}
\sphinxAtStartPar
倾子 Tx 实部
\sphinxbeforeendvarwidth
\end{varwidth}%
&\begin{varwidth}[t]{\sphinxcolwidth{1}{3}}
\sphinxAtStartPar
复数形式
\sphinxbeforeendvarwidth
\end{varwidth}%
\\
\sphinxhline\begin{varwidth}[t]{\sphinxcolwidth{1}{3}}
\sphinxAtStartPar
\sphinxcode{\sphinxupquote{\{name\}\sphinxhyphen{}tzxi.dat}}
\sphinxbeforeendvarwidth
\end{varwidth}%
&\begin{varwidth}[t]{\sphinxcolwidth{1}{3}}
\sphinxAtStartPar
倾子 Tx 虚部
\sphinxbeforeendvarwidth
\end{varwidth}%
&\begin{varwidth}[t]{\sphinxcolwidth{1}{3}}
\sphinxAtStartPar
复数形式
\sphinxbeforeendvarwidth
\end{varwidth}%
\\
\sphinxhline\begin{varwidth}[t]{\sphinxcolwidth{1}{3}}
\sphinxAtStartPar
\sphinxcode{\sphinxupquote{\{name\}\sphinxhyphen{}tzyr.dat}}
\sphinxbeforeendvarwidth
\end{varwidth}%
&\begin{varwidth}[t]{\sphinxcolwidth{1}{3}}
\sphinxAtStartPar
倾子 Ty 实部
\sphinxbeforeendvarwidth
\end{varwidth}%
&\begin{varwidth}[t]{\sphinxcolwidth{1}{3}}
\sphinxAtStartPar
复数形式
\sphinxbeforeendvarwidth
\end{varwidth}%
\\
\sphinxhline\begin{varwidth}[t]{\sphinxcolwidth{1}{3}}
\sphinxAtStartPar
\sphinxcode{\sphinxupquote{\{name\}\sphinxhyphen{}tzyi.dat}}
\sphinxbeforeendvarwidth
\end{varwidth}%
&\begin{varwidth}[t]{\sphinxcolwidth{1}{3}}
\sphinxAtStartPar
倾子 Ty 虚部
\sphinxbeforeendvarwidth
\end{varwidth}%
&\begin{varwidth}[t]{\sphinxcolwidth{1}{3}}
\sphinxAtStartPar
复数形式
\sphinxbeforeendvarwidth
\end{varwidth}%
\\
\sphinxhline\begin{varwidth}[t]{\sphinxcolwidth{1}{3}}
\sphinxAtStartPar
\sphinxcode{\sphinxupquote{\{name\}\sphinxhyphen{}alpha.dat}}
\sphinxbeforeendvarwidth
\end{varwidth}%
&\begin{varwidth}[t]{\sphinxcolwidth{1}{3}}
\sphinxAtStartPar
相位张量 α
\sphinxbeforeendvarwidth
\end{varwidth}%
&\begin{varwidth}[t]{\sphinxcolwidth{1}{3}}
\sphinxAtStartPar
旋转角度
\sphinxbeforeendvarwidth
\end{varwidth}%
\\
\sphinxhline\begin{varwidth}[t]{\sphinxcolwidth{1}{3}}
\sphinxAtStartPar
\sphinxcode{\sphinxupquote{\{name\}\sphinxhyphen{}beta.dat}}
\sphinxbeforeendvarwidth
\end{varwidth}%
&\begin{varwidth}[t]{\sphinxcolwidth{1}{3}}
\sphinxAtStartPar
相位张量 β
\sphinxbeforeendvarwidth
\end{varwidth}%
&\begin{varwidth}[t]{\sphinxcolwidth{1}{3}}
\sphinxAtStartPar
偏离角度
\sphinxbeforeendvarwidth
\end{varwidth}%
\\
\sphinxhline\begin{varwidth}[t]{\sphinxcolwidth{1}{3}}
\sphinxAtStartPar
\sphinxcode{\sphinxupquote{\{name\}\sphinxhyphen{}pmin.dat}}
\sphinxbeforeendvarwidth
\end{varwidth}%
&\begin{varwidth}[t]{\sphinxcolwidth{1}{3}}
\sphinxAtStartPar
相位张量 φmin
\sphinxbeforeendvarwidth
\end{varwidth}%
&\begin{varwidth}[t]{\sphinxcolwidth{1}{3}}
\sphinxAtStartPar
最小主轴
\sphinxbeforeendvarwidth
\end{varwidth}%
\\
\sphinxhline\begin{varwidth}[t]{\sphinxcolwidth{1}{3}}
\sphinxAtStartPar
\sphinxcode{\sphinxupquote{\{name\}\sphinxhyphen{}pmax.dat}}
\sphinxbeforeendvarwidth
\end{varwidth}%
&\begin{varwidth}[t]{\sphinxcolwidth{1}{3}}
\sphinxAtStartPar
相位张量 φmax
\sphinxbeforeendvarwidth
\end{varwidth}%
&\begin{varwidth}[t]{\sphinxcolwidth{1}{3}}
\sphinxAtStartPar
最大主轴
\sphinxbeforeendvarwidth
\end{varwidth}%
\\
\sphinxbottomrule
\end{tabulary}
\sphinxtableafterendhook\par
\sphinxattableend\end{savenotes}


\bigskip\hrule\bigskip



\subsection{Python 绘图环境配置}
\label{\detokenize{tutorial/plotting:python}}

\subsubsection{依赖库}
\label{\detokenize{tutorial/plotting:id4}}
\begin{sphinxVerbatim}[commandchars=\\\{\}]
pip\PYG{+w}{ }install\PYG{+w}{ }numpy\PYG{+w}{ }matplotlib\PYG{+w}{ }pdf2image\PYG{+w}{ }PyPDF2\PYG{+w}{ }shapely\PYG{+w}{ }pyproj
\end{sphinxVerbatim}


\subsubsection{matplotlib 学术绘图配置}
\label{\detokenize{tutorial/plotting:matplotlib}}
\sphinxAtStartPar
建议在脚本开头添加以下配置以符合学术规范：

\begin{sphinxVerbatim}[commandchars=\\\{\}]
\PYG{k+kn}{import}\PYG{+w}{ }\PYG{n+nn}{matplotlib}\PYG{n+nn}{.}\PYG{n+nn}{pyplot}\PYG{+w}{ }\PYG{k}{as}\PYG{+w}{ }\PYG{n+nn}{plt}
\PYG{k+kn}{import}\PYG{+w}{ }\PYG{n+nn}{numpy}\PYG{+w}{ }\PYG{k}{as}\PYG{+w}{ }\PYG{n+nn}{np}

\PYG{c+c1}{\PYGZsh{} 设置 matplotlib 字体样式以符合学术规范}
\PYG{n}{plt}\PYG{o}{.}\PYG{n}{rcParams}\PYG{o}{.}\PYG{n}{update}\PYG{p}{(}\PYG{p}{\PYGZob{}}
    \PYG{l+s+s1}{\PYGZsq{}}\PYG{l+s+s1}{font.size}\PYG{l+s+s1}{\PYGZsq{}}\PYG{p}{:} \PYG{l+m+mi}{12}\PYG{p}{,}
    \PYG{l+s+s1}{\PYGZsq{}}\PYG{l+s+s1}{font.family}\PYG{l+s+s1}{\PYGZsq{}}\PYG{p}{:} \PYG{l+s+s1}{\PYGZsq{}}\PYG{l+s+s1}{serif}\PYG{l+s+s1}{\PYGZsq{}}\PYG{p}{,}
    \PYG{l+s+s1}{\PYGZsq{}}\PYG{l+s+s1}{font.serif}\PYG{l+s+s1}{\PYGZsq{}}\PYG{p}{:} \PYG{l+s+s1}{\PYGZsq{}}\PYG{l+s+s1}{Times New Roman}\PYG{l+s+s1}{\PYGZsq{}}\PYG{p}{,}
    \PYG{l+s+s1}{\PYGZsq{}}\PYG{l+s+s1}{axes.titlesize}\PYG{l+s+s1}{\PYGZsq{}}\PYG{p}{:} \PYG{l+m+mi}{18}\PYG{p}{,}   \PYG{c+c1}{\PYGZsh{} 标题字体大小}
    \PYG{l+s+s1}{\PYGZsq{}}\PYG{l+s+s1}{axes.labelsize}\PYG{l+s+s1}{\PYGZsq{}}\PYG{p}{:} \PYG{l+m+mi}{15}\PYG{p}{,}    \PYG{c+c1}{\PYGZsh{} 轴标签字体大小}
    \PYG{l+s+s1}{\PYGZsq{}}\PYG{l+s+s1}{xtick.labelsize}\PYG{l+s+s1}{\PYGZsq{}}\PYG{p}{:} \PYG{l+m+mi}{12}\PYG{p}{,}  \PYG{c+c1}{\PYGZsh{} x轴刻度标签字体大小}
    \PYG{l+s+s1}{\PYGZsq{}}\PYG{l+s+s1}{ytick.labelsize}\PYG{l+s+s1}{\PYGZsq{}}\PYG{p}{:} \PYG{l+m+mi}{12}\PYG{p}{,}   \PYG{c+c1}{\PYGZsh{} y轴刻度标签字体大小}
    \PYG{l+s+s1}{\PYGZsq{}}\PYG{l+s+s1}{legend.fontsize}\PYG{l+s+s1}{\PYGZsq{}}\PYG{p}{:} \PYG{l+m+mi}{12}    \PYG{c+c1}{\PYGZsh{} 图例字体大小}
\PYG{p}{\PYGZcb{}}\PYG{p}{)}
\end{sphinxVerbatim}


\bigskip\hrule\bigskip



\subsection{单测点视电阻率/相位曲线}
\label{\detokenize{tutorial/plotting:id5}}
\sphinxAtStartPar
绘制单个测点的视电阻率和相位曲线，支持误差棒显示。


\subsubsection{脚本：\sphinxhref{http://PlotFullZ-Single.py}{PlotFullZ\sphinxhyphen{}Single.py} (http://PlotFullZ\sphinxhyphen{}Single.py)}
\label{\detokenize{tutorial/plotting:plotfullz-single-py}}
\begin{sphinxVerbatim}[commandchars=\\\{\}]
\PYG{k+kn}{import}\PYG{+w}{ }\PYG{n+nn}{os}
\PYG{k+kn}{import}\PYG{+w}{ }\PYG{n+nn}{numpy}\PYG{+w}{ }\PYG{k}{as}\PYG{+w}{ }\PYG{n+nn}{np}
\PYG{k+kn}{import}\PYG{+w}{ }\PYG{n+nn}{matplotlib}\PYG{n+nn}{.}\PYG{n+nn}{pyplot}\PYG{+w}{ }\PYG{k}{as}\PYG{+w}{ }\PYG{n+nn}{plt}
\PYG{k+kn}{from}\PYG{+w}{ }\PYG{n+nn}{matplotlib}\PYG{n+nn}{.}\PYG{n+nn}{ticker}\PYG{+w}{ }\PYG{k+kn}{import} \PYG{n}{ScalarFormatter}

\PYG{c+c1}{\PYGZsh{} ===== 配置参数 =====}
\PYG{n}{path} \PYG{o}{=} \PYG{l+s+sa}{r}\PYG{l+s+s1}{\PYGZsq{}}\PYG{l+s+s1}{E:}\PYG{l+s+s1}{\PYGZbs{}}\PYG{l+s+s1}{MTData}\PYG{l+s+s1}{\PYGZbs{}}\PYG{l+s+s1}{Project}\PYG{l+s+s1}{\PYGZbs{}}\PYG{l+s+s1}{SiteData}\PYG{l+s+s1}{\PYGZsq{}}  \PYG{c+c1}{\PYGZsh{} 数据文件夹路径}
\PYG{n}{ymax} \PYG{o}{=} \PYG{l+m+mf}{1e4}   \PYG{c+c1}{\PYGZsh{} 视电阻率 y 轴最大值}
\PYG{n}{ymin} \PYG{o}{=} \PYG{l+m+mf}{1e0}   \PYG{c+c1}{\PYGZsh{} 视电阻率 y 轴最小值}
\PYG{n}{xmax} \PYG{o}{=} \PYG{l+m+mf}{1e3}   \PYG{c+c1}{\PYGZsh{} 频率 x 轴最大值}
\PYG{n}{xmin} \PYG{o}{=} \PYG{l+m+mf}{1e\PYGZhy{}5}  \PYG{c+c1}{\PYGZsh{} 频率 x 轴最小值}

\PYG{c+c1}{\PYGZsh{} 获取所有测点文件夹}
\PYG{n}{names} \PYG{o}{=} \PYG{p}{[}\PYG{n}{f} \PYG{k}{for} \PYG{n}{f} \PYG{o+ow}{in} \PYG{n}{os}\PYG{o}{.}\PYG{n}{listdir}\PYG{p}{(}\PYG{n}{path}\PYG{p}{)} \PYG{k}{if} \PYG{n}{os}\PYG{o}{.}\PYG{n}{path}\PYG{o}{.}\PYG{n}{isdir}\PYG{p}{(}\PYG{n}{os}\PYG{o}{.}\PYG{n}{path}\PYG{o}{.}\PYG{n}{join}\PYG{p}{(}\PYG{n}{path}\PYG{p}{,} \PYG{n}{f}\PYG{p}{)}\PYG{p}{)}\PYG{p}{]}

\PYG{k}{for} \PYG{n}{name} \PYG{o+ow}{in} \PYG{n}{names}\PYG{p}{:}
    \PYG{n+nb}{print}\PYG{p}{(}\PYG{l+s+sa}{f}\PYG{l+s+s1}{\PYGZsq{}}\PYG{l+s+s1}{Processing: }\PYG{l+s+si}{\PYGZob{}}\PYG{n}{name}\PYG{l+s+si}{\PYGZcb{}}\PYG{l+s+s1}{\PYGZsq{}}\PYG{p}{)}
    
    \PYG{c+c1}{\PYGZsh{} 读取数据}
    \PYG{n}{fre} \PYG{o}{=} \PYG{n}{np}\PYG{o}{.}\PYG{n}{loadtxt}\PYG{p}{(}\PYG{l+s+sa}{f}\PYG{l+s+s1}{\PYGZsq{}}\PYG{l+s+si}{\PYGZob{}}\PYG{n}{path}\PYG{l+s+si}{\PYGZcb{}}\PYG{l+s+se}{\PYGZbs{}\PYGZbs{}}\PYG{l+s+si}{\PYGZob{}}\PYG{n}{name}\PYG{l+s+si}{\PYGZcb{}}\PYG{l+s+se}{\PYGZbs{}\PYGZbs{}}\PYG{l+s+si}{\PYGZob{}}\PYG{n}{name}\PYG{l+s+si}{\PYGZcb{}}\PYG{l+s+s1}{\PYGZhy{}rxx.dat}\PYG{l+s+s1}{\PYGZsq{}}\PYG{p}{)}\PYG{p}{[}\PYG{p}{:}\PYG{p}{,} \PYG{l+m+mi}{0}\PYG{p}{]}
    \PYG{n}{rxy} \PYG{o}{=} \PYG{n}{np}\PYG{o}{.}\PYG{n}{loadtxt}\PYG{p}{(}\PYG{l+s+sa}{f}\PYG{l+s+s1}{\PYGZsq{}}\PYG{l+s+si}{\PYGZob{}}\PYG{n}{path}\PYG{l+s+si}{\PYGZcb{}}\PYG{l+s+se}{\PYGZbs{}\PYGZbs{}}\PYG{l+s+si}{\PYGZob{}}\PYG{n}{name}\PYG{l+s+si}{\PYGZcb{}}\PYG{l+s+se}{\PYGZbs{}\PYGZbs{}}\PYG{l+s+si}{\PYGZob{}}\PYG{n}{name}\PYG{l+s+si}{\PYGZcb{}}\PYG{l+s+s1}{\PYGZhy{}rxy.dat}\PYG{l+s+s1}{\PYGZsq{}}\PYG{p}{)}\PYG{p}{[}\PYG{p}{:}\PYG{p}{,} \PYG{l+m+mi}{1}\PYG{p}{]}
    \PYG{n}{ryx} \PYG{o}{=} \PYG{n}{np}\PYG{o}{.}\PYG{n}{loadtxt}\PYG{p}{(}\PYG{l+s+sa}{f}\PYG{l+s+s1}{\PYGZsq{}}\PYG{l+s+si}{\PYGZob{}}\PYG{n}{path}\PYG{l+s+si}{\PYGZcb{}}\PYG{l+s+se}{\PYGZbs{}\PYGZbs{}}\PYG{l+s+si}{\PYGZob{}}\PYG{n}{name}\PYG{l+s+si}{\PYGZcb{}}\PYG{l+s+se}{\PYGZbs{}\PYGZbs{}}\PYG{l+s+si}{\PYGZob{}}\PYG{n}{name}\PYG{l+s+si}{\PYGZcb{}}\PYG{l+s+s1}{\PYGZhy{}ryx.dat}\PYG{l+s+s1}{\PYGZsq{}}\PYG{p}{)}\PYG{p}{[}\PYG{p}{:}\PYG{p}{,} \PYG{l+m+mi}{1}\PYG{p}{]}
    \PYG{n}{pxy} \PYG{o}{=} \PYG{n}{np}\PYG{o}{.}\PYG{n}{loadtxt}\PYG{p}{(}\PYG{l+s+sa}{f}\PYG{l+s+s1}{\PYGZsq{}}\PYG{l+s+si}{\PYGZob{}}\PYG{n}{path}\PYG{l+s+si}{\PYGZcb{}}\PYG{l+s+se}{\PYGZbs{}\PYGZbs{}}\PYG{l+s+si}{\PYGZob{}}\PYG{n}{name}\PYG{l+s+si}{\PYGZcb{}}\PYG{l+s+se}{\PYGZbs{}\PYGZbs{}}\PYG{l+s+si}{\PYGZob{}}\PYG{n}{name}\PYG{l+s+si}{\PYGZcb{}}\PYG{l+s+s1}{\PYGZhy{}zxyp.dat}\PYG{l+s+s1}{\PYGZsq{}}\PYG{p}{)}\PYG{p}{[}\PYG{p}{:}\PYG{p}{,} \PYG{l+m+mi}{1}\PYG{p}{]}
    \PYG{n}{pyx} \PYG{o}{=} \PYG{n}{np}\PYG{o}{.}\PYG{n}{loadtxt}\PYG{p}{(}\PYG{l+s+sa}{f}\PYG{l+s+s1}{\PYGZsq{}}\PYG{l+s+si}{\PYGZob{}}\PYG{n}{path}\PYG{l+s+si}{\PYGZcb{}}\PYG{l+s+se}{\PYGZbs{}\PYGZbs{}}\PYG{l+s+si}{\PYGZob{}}\PYG{n}{name}\PYG{l+s+si}{\PYGZcb{}}\PYG{l+s+se}{\PYGZbs{}\PYGZbs{}}\PYG{l+s+si}{\PYGZob{}}\PYG{n}{name}\PYG{l+s+si}{\PYGZcb{}}\PYG{l+s+s1}{\PYGZhy{}zyxp.dat}\PYG{l+s+s1}{\PYGZsq{}}\PYG{p}{)}\PYG{p}{[}\PYG{p}{:}\PYG{p}{,} \PYG{l+m+mi}{1}\PYG{p}{]}
    
    \PYG{c+c1}{\PYGZsh{} 读取误差数据}
    \PYG{n}{rxy\PYGZus{}var} \PYG{o}{=} \PYG{n}{np}\PYG{o}{.}\PYG{n}{loadtxt}\PYG{p}{(}\PYG{l+s+sa}{f}\PYG{l+s+s1}{\PYGZsq{}}\PYG{l+s+si}{\PYGZob{}}\PYG{n}{path}\PYG{l+s+si}{\PYGZcb{}}\PYG{l+s+se}{\PYGZbs{}\PYGZbs{}}\PYG{l+s+si}{\PYGZob{}}\PYG{n}{name}\PYG{l+s+si}{\PYGZcb{}}\PYG{l+s+se}{\PYGZbs{}\PYGZbs{}}\PYG{l+s+si}{\PYGZob{}}\PYG{n}{name}\PYG{l+s+si}{\PYGZcb{}}\PYG{l+s+s1}{\PYGZhy{}rxyvar.dat}\PYG{l+s+s1}{\PYGZsq{}}\PYG{p}{)}\PYG{p}{[}\PYG{p}{:}\PYG{p}{,} \PYG{l+m+mi}{1}\PYG{p}{]}
    \PYG{n}{ryx\PYGZus{}var} \PYG{o}{=} \PYG{n}{np}\PYG{o}{.}\PYG{n}{loadtxt}\PYG{p}{(}\PYG{l+s+sa}{f}\PYG{l+s+s1}{\PYGZsq{}}\PYG{l+s+si}{\PYGZob{}}\PYG{n}{path}\PYG{l+s+si}{\PYGZcb{}}\PYG{l+s+se}{\PYGZbs{}\PYGZbs{}}\PYG{l+s+si}{\PYGZob{}}\PYG{n}{name}\PYG{l+s+si}{\PYGZcb{}}\PYG{l+s+se}{\PYGZbs{}\PYGZbs{}}\PYG{l+s+si}{\PYGZob{}}\PYG{n}{name}\PYG{l+s+si}{\PYGZcb{}}\PYG{l+s+s1}{\PYGZhy{}ryxvar.dat}\PYG{l+s+s1}{\PYGZsq{}}\PYG{p}{)}\PYG{p}{[}\PYG{p}{:}\PYG{p}{,} \PYG{l+m+mi}{1}\PYG{p}{]}
    \PYG{n}{pxy\PYGZus{}var} \PYG{o}{=} \PYG{n}{np}\PYG{o}{.}\PYG{n}{loadtxt}\PYG{p}{(}\PYG{l+s+sa}{f}\PYG{l+s+s1}{\PYGZsq{}}\PYG{l+s+si}{\PYGZob{}}\PYG{n}{path}\PYG{l+s+si}{\PYGZcb{}}\PYG{l+s+se}{\PYGZbs{}\PYGZbs{}}\PYG{l+s+si}{\PYGZob{}}\PYG{n}{name}\PYG{l+s+si}{\PYGZcb{}}\PYG{l+s+se}{\PYGZbs{}\PYGZbs{}}\PYG{l+s+si}{\PYGZob{}}\PYG{n}{name}\PYG{l+s+si}{\PYGZcb{}}\PYG{l+s+s1}{\PYGZhy{}pxyvar.dat}\PYG{l+s+s1}{\PYGZsq{}}\PYG{p}{)}\PYG{p}{[}\PYG{p}{:}\PYG{p}{,} \PYG{l+m+mi}{1}\PYG{p}{]}
    \PYG{n}{pyx\PYGZus{}var} \PYG{o}{=} \PYG{n}{np}\PYG{o}{.}\PYG{n}{loadtxt}\PYG{p}{(}\PYG{l+s+sa}{f}\PYG{l+s+s1}{\PYGZsq{}}\PYG{l+s+si}{\PYGZob{}}\PYG{n}{path}\PYG{l+s+si}{\PYGZcb{}}\PYG{l+s+se}{\PYGZbs{}\PYGZbs{}}\PYG{l+s+si}{\PYGZob{}}\PYG{n}{name}\PYG{l+s+si}{\PYGZcb{}}\PYG{l+s+se}{\PYGZbs{}\PYGZbs{}}\PYG{l+s+si}{\PYGZob{}}\PYG{n}{name}\PYG{l+s+si}{\PYGZcb{}}\PYG{l+s+s1}{\PYGZhy{}pyxvar.dat}\PYG{l+s+s1}{\PYGZsq{}}\PYG{p}{)}\PYG{p}{[}\PYG{p}{:}\PYG{p}{,} \PYG{l+m+mi}{1}\PYG{p}{]}
    
    \PYG{c+c1}{\PYGZsh{} 创建图形和子图}
    \PYG{n}{fig}\PYG{p}{,} \PYG{n}{axs} \PYG{o}{=} \PYG{n}{plt}\PYG{o}{.}\PYG{n}{subplots}\PYG{p}{(}\PYG{l+m+mi}{2}\PYG{p}{,} \PYG{l+m+mi}{1}\PYG{p}{,} \PYG{n}{figsize}\PYG{o}{=}\PYG{p}{(}\PYG{l+m+mi}{8}\PYG{p}{,} \PYG{l+m+mi}{7}\PYG{p}{)}\PYG{p}{,} \PYG{n}{sharex}\PYG{o}{=}\PYG{l+s+s1}{\PYGZsq{}}\PYG{l+s+s1}{all}\PYG{l+s+s1}{\PYGZsq{}}\PYG{p}{,} \PYG{n}{constrained\PYGZus{}layout}\PYG{o}{=}\PYG{k+kc}{True}\PYG{p}{)}
    
    \PYG{c+c1}{\PYGZsh{} 视电阻率子图}
    \PYG{n}{ax} \PYG{o}{=} \PYG{n}{axs}\PYG{p}{[}\PYG{l+m+mi}{0}\PYG{p}{]}
    \PYG{n}{ax2} \PYG{o}{=} \PYG{n}{ax}\PYG{o}{.}\PYG{n}{twinx}\PYG{p}{(}\PYG{p}{)}
    
    \PYG{c+c1}{\PYGZsh{} 误差转换到对数坐标}
    \PYG{n}{log\PYGZus{}rxy} \PYG{o}{=} \PYG{n}{np}\PYG{o}{.}\PYG{n}{log10}\PYG{p}{(}\PYG{n}{rxy}\PYG{p}{)}
    \PYG{n}{log\PYGZus{}ryx} \PYG{o}{=} \PYG{n}{np}\PYG{o}{.}\PYG{n}{log10}\PYG{p}{(}\PYG{n}{ryx}\PYG{p}{)}
    \PYG{n}{log\PYGZus{}rxy\PYGZus{}var} \PYG{o}{=} \PYG{n}{rxy\PYGZus{}var} \PYG{o}{/} \PYG{p}{(}\PYG{n}{rxy} \PYG{o}{*} \PYG{n}{np}\PYG{o}{.}\PYG{n}{log}\PYG{p}{(}\PYG{l+m+mi}{10}\PYG{p}{)}\PYG{p}{)}
    \PYG{n}{log\PYGZus{}ryx\PYGZus{}var} \PYG{o}{=} \PYG{n}{ryx\PYGZus{}var} \PYG{o}{/} \PYG{p}{(}\PYG{n}{ryx} \PYG{o}{*} \PYG{n}{np}\PYG{o}{.}\PYG{n}{log}\PYG{p}{(}\PYG{l+m+mi}{10}\PYG{p}{)}\PYG{p}{)}
    
    \PYG{n}{ax2}\PYG{o}{.}\PYG{n}{errorbar}\PYG{p}{(}\PYG{n}{fre}\PYG{p}{,} \PYG{n}{log\PYGZus{}rxy}\PYG{p}{,} \PYG{n}{yerr}\PYG{o}{=}\PYG{n}{log\PYGZus{}rxy\PYGZus{}var}\PYG{p}{,} \PYG{n}{fmt}\PYG{o}{=}\PYG{l+s+s1}{\PYGZsq{}}\PYG{l+s+s1}{ro}\PYG{l+s+s1}{\PYGZsq{}}\PYG{p}{,} \PYG{n}{label}\PYG{o}{=}\PYG{l+s+s1}{\PYGZsq{}}\PYG{l+s+s1}{XY}\PYG{l+s+s1}{\PYGZsq{}}\PYG{p}{,} \PYG{n}{markerfacecolor}\PYG{o}{=}\PYG{l+s+s1}{\PYGZsq{}}\PYG{l+s+s1}{none}\PYG{l+s+s1}{\PYGZsq{}}\PYG{p}{)}
    \PYG{n}{ax2}\PYG{o}{.}\PYG{n}{errorbar}\PYG{p}{(}\PYG{n}{fre}\PYG{p}{,} \PYG{n}{log\PYGZus{}ryx}\PYG{p}{,} \PYG{n}{yerr}\PYG{o}{=}\PYG{n}{log\PYGZus{}ryx\PYGZus{}var}\PYG{p}{,} \PYG{n}{fmt}\PYG{o}{=}\PYG{l+s+s1}{\PYGZsq{}}\PYG{l+s+s1}{bs}\PYG{l+s+s1}{\PYGZsq{}}\PYG{p}{,} \PYG{n}{label}\PYG{o}{=}\PYG{l+s+s1}{\PYGZsq{}}\PYG{l+s+s1}{YX}\PYG{l+s+s1}{\PYGZsq{}}\PYG{p}{,} \PYG{n}{markerfacecolor}\PYG{o}{=}\PYG{l+s+s1}{\PYGZsq{}}\PYG{l+s+s1}{none}\PYG{l+s+s1}{\PYGZsq{}}\PYG{p}{)}
    \PYG{n}{ax2}\PYG{o}{.}\PYG{n}{set\PYGZus{}ylim}\PYG{p}{(}\PYG{n}{np}\PYG{o}{.}\PYG{n}{log10}\PYG{p}{(}\PYG{n}{ymin}\PYG{p}{)}\PYG{p}{,} \PYG{n}{np}\PYG{o}{.}\PYG{n}{log10}\PYG{p}{(}\PYG{n}{ymax}\PYG{p}{)}\PYG{p}{)}
    \PYG{n}{ax2}\PYG{o}{.}\PYG{n}{invert\PYGZus{}xaxis}\PYG{p}{(}\PYG{p}{)}
    \PYG{n}{ax2}\PYG{o}{.}\PYG{n}{legend}\PYG{p}{(}\PYG{n}{loc}\PYG{o}{=}\PYG{l+s+s1}{\PYGZsq{}}\PYG{l+s+s1}{upper right}\PYG{l+s+s1}{\PYGZsq{}}\PYG{p}{,} \PYG{n}{bbox\PYGZus{}to\PYGZus{}anchor}\PYG{o}{=}\PYG{p}{(}\PYG{l+m+mf}{1.2}\PYG{p}{,} \PYG{l+m+mf}{0.8}\PYG{p}{)}\PYG{p}{)}
    
    \PYG{n}{ax}\PYG{o}{.}\PYG{n}{set\PYGZus{}ylim}\PYG{p}{(}\PYG{n}{ymin}\PYG{p}{,} \PYG{n}{ymax}\PYG{p}{)}
    \PYG{n}{ax}\PYG{o}{.}\PYG{n}{set\PYGZus{}yscale}\PYG{p}{(}\PYG{l+s+s1}{\PYGZsq{}}\PYG{l+s+s1}{log}\PYG{l+s+s1}{\PYGZsq{}}\PYG{p}{)}
    \PYG{n}{ax}\PYG{o}{.}\PYG{n}{set\PYGZus{}xscale}\PYG{p}{(}\PYG{l+s+s1}{\PYGZsq{}}\PYG{l+s+s1}{log}\PYG{l+s+s1}{\PYGZsq{}}\PYG{p}{)}
    \PYG{n}{ax}\PYG{o}{.}\PYG{n}{set\PYGZus{}xlim}\PYG{p}{(}\PYG{n}{xmin}\PYG{p}{,} \PYG{n}{xmax}\PYG{p}{)}
    \PYG{n}{ax}\PYG{o}{.}\PYG{n}{invert\PYGZus{}xaxis}\PYG{p}{(}\PYG{p}{)}
    \PYG{n}{ax}\PYG{o}{.}\PYG{n}{grid}\PYG{p}{(}\PYG{k+kc}{True}\PYG{p}{,} \PYG{n}{which}\PYG{o}{=}\PYG{l+s+s1}{\PYGZsq{}}\PYG{l+s+s1}{major}\PYG{l+s+s1}{\PYGZsq{}}\PYG{p}{,} \PYG{n}{linestyle}\PYG{o}{=}\PYG{l+s+s1}{\PYGZsq{}}\PYG{l+s+s1}{\PYGZhy{}\PYGZhy{}}\PYG{l+s+s1}{\PYGZsq{}}\PYG{p}{,} \PYG{n}{linewidth}\PYG{o}{=}\PYG{l+m+mf}{0.5}\PYG{p}{)}
    \PYG{n}{ax}\PYG{o}{.}\PYG{n}{set\PYGZus{}ylabel}\PYG{p}{(}\PYG{l+s+sa}{r}\PYG{l+s+s1}{\PYGZsq{}}\PYG{l+s+s1}{\PYGZdl{}}\PYG{l+s+s1}{\PYGZbs{}}\PYG{l+s+s1}{rho}\PYG{l+s+s1}{\PYGZbs{}}\PYG{l+s+s1}{;(}\PYG{l+s+s1}{\PYGZbs{}}\PYG{l+s+s1}{Omega }\PYG{l+s+s1}{\PYGZbs{}}\PYG{l+s+s1}{cdot m)\PYGZdl{}}\PYG{l+s+s1}{\PYGZsq{}}\PYG{p}{)}
    
    \PYG{c+c1}{\PYGZsh{} 相位子图}
    \PYG{n}{ax} \PYG{o}{=} \PYG{n}{axs}\PYG{p}{[}\PYG{l+m+mi}{1}\PYG{p}{]}
    \PYG{n}{ax}\PYG{o}{.}\PYG{n}{errorbar}\PYG{p}{(}\PYG{n}{fre}\PYG{p}{,} \PYG{n}{pxy}\PYG{p}{,} \PYG{n}{yerr}\PYG{o}{=}\PYG{n}{pxy\PYGZus{}var}\PYG{p}{,} \PYG{n}{fmt}\PYG{o}{=}\PYG{l+s+s1}{\PYGZsq{}}\PYG{l+s+s1}{ro}\PYG{l+s+s1}{\PYGZsq{}}\PYG{p}{,} \PYG{n}{label}\PYG{o}{=}\PYG{l+s+s1}{\PYGZsq{}}\PYG{l+s+s1}{XY}\PYG{l+s+s1}{\PYGZsq{}}\PYG{p}{,} \PYG{n}{markerfacecolor}\PYG{o}{=}\PYG{l+s+s1}{\PYGZsq{}}\PYG{l+s+s1}{none}\PYG{l+s+s1}{\PYGZsq{}}\PYG{p}{)}
    \PYG{n}{ax}\PYG{o}{.}\PYG{n}{errorbar}\PYG{p}{(}\PYG{n}{fre}\PYG{p}{,} \PYG{n}{pyx}\PYG{p}{,} \PYG{n}{yerr}\PYG{o}{=}\PYG{n}{pyx\PYGZus{}var}\PYG{p}{,} \PYG{n}{fmt}\PYG{o}{=}\PYG{l+s+s1}{\PYGZsq{}}\PYG{l+s+s1}{bs}\PYG{l+s+s1}{\PYGZsq{}}\PYG{p}{,} \PYG{n}{label}\PYG{o}{=}\PYG{l+s+s1}{\PYGZsq{}}\PYG{l+s+s1}{YX}\PYG{l+s+s1}{\PYGZsq{}}\PYG{p}{,} \PYG{n}{markerfacecolor}\PYG{o}{=}\PYG{l+s+s1}{\PYGZsq{}}\PYG{l+s+s1}{none}\PYG{l+s+s1}{\PYGZsq{}}\PYG{p}{)}
    \PYG{n}{ax}\PYG{o}{.}\PYG{n}{set\PYGZus{}ylim}\PYG{p}{(}\PYG{o}{\PYGZhy{}}\PYG{l+m+mi}{180}\PYG{p}{,} \PYG{l+m+mi}{180}\PYG{p}{)}
    \PYG{n}{ax}\PYG{o}{.}\PYG{n}{set\PYGZus{}yticks}\PYG{p}{(}\PYG{n}{np}\PYG{o}{.}\PYG{n}{arange}\PYG{p}{(}\PYG{o}{\PYGZhy{}}\PYG{l+m+mi}{180}\PYG{p}{,} \PYG{l+m+mi}{181}\PYG{p}{,} \PYG{l+m+mi}{90}\PYG{p}{)}\PYG{p}{)}
    \PYG{n}{ax}\PYG{o}{.}\PYG{n}{set\PYGZus{}xscale}\PYG{p}{(}\PYG{l+s+s1}{\PYGZsq{}}\PYG{l+s+s1}{log}\PYG{l+s+s1}{\PYGZsq{}}\PYG{p}{)}
    \PYG{n}{ax}\PYG{o}{.}\PYG{n}{set\PYGZus{}xlabel}\PYG{p}{(}\PYG{l+s+s1}{\PYGZsq{}}\PYG{l+s+s1}{Frequency (Hz)}\PYG{l+s+s1}{\PYGZsq{}}\PYG{p}{)}
    \PYG{n}{ax}\PYG{o}{.}\PYG{n}{set\PYGZus{}ylabel}\PYG{p}{(}\PYG{l+s+sa}{r}\PYG{l+s+s1}{\PYGZsq{}}\PYG{l+s+s1}{\PYGZdl{}}\PYG{l+s+s1}{\PYGZbs{}}\PYG{l+s+s1}{varphi}\PYG{l+s+s1}{\PYGZbs{}}\PYG{l+s+s1}{;(}\PYG{l+s+s1}{\PYGZbs{}}\PYG{l+s+s1}{degree)\PYGZdl{}}\PYG{l+s+s1}{\PYGZsq{}}\PYG{p}{)}
    \PYG{n}{ax}\PYG{o}{.}\PYG{n}{grid}\PYG{p}{(}\PYG{k+kc}{True}\PYG{p}{,} \PYG{n}{which}\PYG{o}{=}\PYG{l+s+s1}{\PYGZsq{}}\PYG{l+s+s1}{major}\PYG{l+s+s1}{\PYGZsq{}}\PYG{p}{,} \PYG{n}{linestyle}\PYG{o}{=}\PYG{l+s+s1}{\PYGZsq{}}\PYG{l+s+s1}{\PYGZhy{}\PYGZhy{}}\PYG{l+s+s1}{\PYGZsq{}}\PYG{p}{,} \PYG{n}{linewidth}\PYG{o}{=}\PYG{l+m+mf}{0.5}\PYG{p}{)}
    
    \PYG{c+c1}{\PYGZsh{} 保存图形}
    \PYG{n}{plt}\PYG{o}{.}\PYG{n}{savefig}\PYG{p}{(}\PYG{l+s+sa}{f}\PYG{l+s+s1}{\PYGZsq{}}\PYG{l+s+si}{\PYGZob{}}\PYG{n}{path}\PYG{l+s+si}{\PYGZcb{}}\PYG{l+s+se}{\PYGZbs{}\PYGZbs{}}\PYG{l+s+si}{\PYGZob{}}\PYG{n}{name}\PYG{l+s+si}{\PYGZcb{}}\PYG{l+s+s1}{\PYGZus{}Site.png}\PYG{l+s+s1}{\PYGZsq{}}\PYG{p}{,} \PYG{n}{dpi}\PYG{o}{=}\PYG{l+m+mi}{300}\PYG{p}{,} \PYG{n}{bbox\PYGZus{}inches}\PYG{o}{=}\PYG{l+s+s1}{\PYGZsq{}}\PYG{l+s+s1}{tight}\PYG{l+s+s1}{\PYGZsq{}}\PYG{p}{)}
    \PYG{n}{plt}\PYG{o}{.}\PYG{n}{savefig}\PYG{p}{(}\PYG{l+s+sa}{f}\PYG{l+s+s1}{\PYGZsq{}}\PYG{l+s+si}{\PYGZob{}}\PYG{n}{path}\PYG{l+s+si}{\PYGZcb{}}\PYG{l+s+se}{\PYGZbs{}\PYGZbs{}}\PYG{l+s+si}{\PYGZob{}}\PYG{n}{name}\PYG{l+s+si}{\PYGZcb{}}\PYG{l+s+s1}{\PYGZus{}Site.pdf}\PYG{l+s+s1}{\PYGZsq{}}\PYG{p}{,} \PYG{n+nb}{format}\PYG{o}{=}\PYG{l+s+s1}{\PYGZsq{}}\PYG{l+s+s1}{pdf}\PYG{l+s+s1}{\PYGZsq{}}\PYG{p}{,} \PYG{n}{dpi}\PYG{o}{=}\PYG{l+m+mi}{300}\PYG{p}{,} \PYG{n}{bbox\PYGZus{}inches}\PYG{o}{=}\PYG{l+s+s1}{\PYGZsq{}}\PYG{l+s+s1}{tight}\PYG{l+s+s1}{\PYGZsq{}}\PYG{p}{)}
    \PYG{n}{plt}\PYG{o}{.}\PYG{n}{close}\PYG{p}{(}\PYG{p}{)}

\PYG{n+nb}{print}\PYG{p}{(}\PYG{l+s+s1}{\PYGZsq{}}\PYG{l+s+s1}{Finished!}\PYG{l+s+s1}{\PYGZsq{}}\PYG{p}{)}
\end{sphinxVerbatim}


\subsubsection{输出效果}
\label{\detokenize{tutorial/plotting:id6}}
\sphinxAtStartPar
生成的图像包含两个子图：
\begin{itemize}
\item {} 
\sphinxAtStartPar
上图：视电阻率 ρxy（红色）和 ρyx（蓝色）随频率变化

\item {} 
\sphinxAtStartPar
下图：对应的相位 φxy 和 φyx

\end{itemize}


\bigskip\hrule\bigskip



\subsection{多测点联合对比图}
\label{\detokenize{tutorial/plotting:id7}}
\sphinxAtStartPar
将多个测点的曲线绘制在同一张图中，便于对比分析。


\subsubsection{脚本：\sphinxhref{http://PlotFullZ-All.py}{PlotFullZ\sphinxhyphen{}All.py} (http://PlotFullZ\sphinxhyphen{}All.py)}
\label{\detokenize{tutorial/plotting:plotfullz-all-py}}
\begin{sphinxVerbatim}[commandchars=\\\{\}]
\PYG{k+kn}{import}\PYG{+w}{ }\PYG{n+nn}{os}
\PYG{k+kn}{import}\PYG{+w}{ }\PYG{n+nn}{numpy}\PYG{+w}{ }\PYG{k}{as}\PYG{+w}{ }\PYG{n+nn}{np}
\PYG{k+kn}{import}\PYG{+w}{ }\PYG{n+nn}{matplotlib}\PYG{n+nn}{.}\PYG{n+nn}{pyplot}\PYG{+w}{ }\PYG{k}{as}\PYG{+w}{ }\PYG{n+nn}{plt}
\PYG{k+kn}{from}\PYG{+w}{ }\PYG{n+nn}{matplotlib}\PYG{n+nn}{.}\PYG{n+nn}{ticker}\PYG{+w}{ }\PYG{k+kn}{import} \PYG{n}{ScalarFormatter}

\PYG{c+c1}{\PYGZsh{} ===== 配置参数 =====}
\PYG{n}{path} \PYG{o}{=} \PYG{l+s+sa}{r}\PYG{l+s+s1}{\PYGZsq{}}\PYG{l+s+s1}{F:}\PYG{l+s+s1}{\PYGZbs{}}\PYG{l+s+s1}{ZhangYeAMT}\PYG{l+s+s1}{\PYGZbs{}}\PYG{l+s+s1}{MTDPProject}\PYG{l+s+s1}{\PYGZbs{}}\PYG{l+s+s1}{ZhangYe}\PYG{l+s+s1}{\PYGZbs{}}\PYG{l+s+s1}{0901}\PYG{l+s+s1}{\PYGZsq{}}
\PYG{n}{nrows} \PYG{o}{=} \PYG{l+m+mi}{8}   \PYG{c+c1}{\PYGZsh{} 网格行数（每个测点占2行：电阻率+相位）}
\PYG{n}{ncols} \PYG{o}{=} \PYG{l+m+mi}{5}    \PYG{c+c1}{\PYGZsh{} 网格列数}
\PYG{n}{ymax} \PYG{o}{=} \PYG{l+m+mf}{1e4}
\PYG{n}{ymin} \PYG{o}{=} \PYG{l+m+mf}{1e\PYGZhy{}1}
\PYG{n}{xmax} \PYG{o}{=} \PYG{l+m+mf}{1e3}
\PYG{n}{xmin} \PYG{o}{=} \PYG{l+m+mf}{1e\PYGZhy{}4}

\PYG{c+c1}{\PYGZsh{} 获取所有测点文件夹}
\PYG{n}{names} \PYG{o}{=} \PYG{p}{[}\PYG{n}{f} \PYG{k}{for} \PYG{n}{f} \PYG{o+ow}{in} \PYG{n}{os}\PYG{o}{.}\PYG{n}{listdir}\PYG{p}{(}\PYG{n}{path}\PYG{p}{)} \PYG{k}{if} \PYG{n}{os}\PYG{o}{.}\PYG{n}{path}\PYG{o}{.}\PYG{n}{isdir}\PYG{p}{(}\PYG{n}{os}\PYG{o}{.}\PYG{n}{path}\PYG{o}{.}\PYG{n}{join}\PYG{p}{(}\PYG{n}{path}\PYG{p}{,} \PYG{n}{f}\PYG{p}{)}\PYG{p}{)}\PYG{p}{]}

\PYG{c+c1}{\PYGZsh{} 创建大图形}
\PYG{n}{fig}\PYG{p}{,} \PYG{n}{axs} \PYG{o}{=} \PYG{n}{plt}\PYG{o}{.}\PYG{n}{subplots}\PYG{p}{(}\PYG{n}{nrows}\PYG{o}{=}\PYG{n}{nrows}\PYG{o}{*}\PYG{l+m+mi}{2}\PYG{p}{,} \PYG{n}{ncols}\PYG{o}{=}\PYG{n}{ncols}\PYG{p}{,} \PYG{n}{figsize}\PYG{o}{=}\PYG{p}{(}\PYG{l+m+mi}{25}\PYG{p}{,} \PYG{l+m+mi}{40}\PYG{p}{)}\PYG{p}{,} \PYG{n}{sharex}\PYG{o}{=}\PYG{l+s+s1}{\PYGZsq{}}\PYG{l+s+s1}{all}\PYG{l+s+s1}{\PYGZsq{}}\PYG{p}{,} \PYG{n}{constrained\PYGZus{}layout}\PYG{o}{=}\PYG{k+kc}{True}\PYG{p}{)}

\PYG{k}{for} \PYG{n}{i}\PYG{p}{,} \PYG{n}{name} \PYG{o+ow}{in} \PYG{n+nb}{enumerate}\PYG{p}{(}\PYG{n}{names}\PYG{p}{)}\PYG{p}{:}
    \PYG{n+nb}{print}\PYG{p}{(}\PYG{l+s+sa}{f}\PYG{l+s+s1}{\PYGZsq{}}\PYG{l+s+se}{\PYGZbs{}r}\PYG{l+s+s1}{Processing }\PYG{l+s+si}{\PYGZob{}}\PYG{n}{name}\PYG{l+s+si}{\PYGZcb{}}\PYG{l+s+s1}{\PYGZsq{}}\PYG{p}{,} \PYG{n}{end}\PYG{o}{=}\PYG{l+s+s1}{\PYGZsq{}}\PYG{l+s+s1}{\PYGZsq{}}\PYG{p}{)}
    
    \PYG{c+c1}{\PYGZsh{} 读取数据}
    \PYG{n}{fre} \PYG{o}{=} \PYG{n}{np}\PYG{o}{.}\PYG{n}{loadtxt}\PYG{p}{(}\PYG{l+s+sa}{f}\PYG{l+s+s1}{\PYGZsq{}}\PYG{l+s+si}{\PYGZob{}}\PYG{n}{path}\PYG{l+s+si}{\PYGZcb{}}\PYG{l+s+se}{\PYGZbs{}\PYGZbs{}}\PYG{l+s+si}{\PYGZob{}}\PYG{n}{name}\PYG{l+s+si}{\PYGZcb{}}\PYG{l+s+se}{\PYGZbs{}\PYGZbs{}}\PYG{l+s+si}{\PYGZob{}}\PYG{n}{name}\PYG{l+s+si}{\PYGZcb{}}\PYG{l+s+s1}{\PYGZhy{}rxx.dat}\PYG{l+s+s1}{\PYGZsq{}}\PYG{p}{)}\PYG{p}{[}\PYG{p}{:}\PYG{p}{,} \PYG{l+m+mi}{0}\PYG{p}{]}
    \PYG{n}{rxy} \PYG{o}{=} \PYG{n}{np}\PYG{o}{.}\PYG{n}{loadtxt}\PYG{p}{(}\PYG{l+s+sa}{f}\PYG{l+s+s1}{\PYGZsq{}}\PYG{l+s+si}{\PYGZob{}}\PYG{n}{path}\PYG{l+s+si}{\PYGZcb{}}\PYG{l+s+se}{\PYGZbs{}\PYGZbs{}}\PYG{l+s+si}{\PYGZob{}}\PYG{n}{name}\PYG{l+s+si}{\PYGZcb{}}\PYG{l+s+se}{\PYGZbs{}\PYGZbs{}}\PYG{l+s+si}{\PYGZob{}}\PYG{n}{name}\PYG{l+s+si}{\PYGZcb{}}\PYG{l+s+s1}{\PYGZhy{}rxy.dat}\PYG{l+s+s1}{\PYGZsq{}}\PYG{p}{)}\PYG{p}{[}\PYG{p}{:}\PYG{p}{,} \PYG{l+m+mi}{1}\PYG{p}{]}
    \PYG{n}{ryx} \PYG{o}{=} \PYG{n}{np}\PYG{o}{.}\PYG{n}{loadtxt}\PYG{p}{(}\PYG{l+s+sa}{f}\PYG{l+s+s1}{\PYGZsq{}}\PYG{l+s+si}{\PYGZob{}}\PYG{n}{path}\PYG{l+s+si}{\PYGZcb{}}\PYG{l+s+se}{\PYGZbs{}\PYGZbs{}}\PYG{l+s+si}{\PYGZob{}}\PYG{n}{name}\PYG{l+s+si}{\PYGZcb{}}\PYG{l+s+se}{\PYGZbs{}\PYGZbs{}}\PYG{l+s+si}{\PYGZob{}}\PYG{n}{name}\PYG{l+s+si}{\PYGZcb{}}\PYG{l+s+s1}{\PYGZhy{}ryx.dat}\PYG{l+s+s1}{\PYGZsq{}}\PYG{p}{)}\PYG{p}{[}\PYG{p}{:}\PYG{p}{,} \PYG{l+m+mi}{1}\PYG{p}{]}
    \PYG{n}{pxy} \PYG{o}{=} \PYG{n}{np}\PYG{o}{.}\PYG{n}{loadtxt}\PYG{p}{(}\PYG{l+s+sa}{f}\PYG{l+s+s1}{\PYGZsq{}}\PYG{l+s+si}{\PYGZob{}}\PYG{n}{path}\PYG{l+s+si}{\PYGZcb{}}\PYG{l+s+se}{\PYGZbs{}\PYGZbs{}}\PYG{l+s+si}{\PYGZob{}}\PYG{n}{name}\PYG{l+s+si}{\PYGZcb{}}\PYG{l+s+se}{\PYGZbs{}\PYGZbs{}}\PYG{l+s+si}{\PYGZob{}}\PYG{n}{name}\PYG{l+s+si}{\PYGZcb{}}\PYG{l+s+s1}{\PYGZhy{}zxyp.dat}\PYG{l+s+s1}{\PYGZsq{}}\PYG{p}{)}\PYG{p}{[}\PYG{p}{:}\PYG{p}{,} \PYG{l+m+mi}{1}\PYG{p}{]}
    \PYG{n}{pyx} \PYG{o}{=} \PYG{n}{np}\PYG{o}{.}\PYG{n}{loadtxt}\PYG{p}{(}\PYG{l+s+sa}{f}\PYG{l+s+s1}{\PYGZsq{}}\PYG{l+s+si}{\PYGZob{}}\PYG{n}{path}\PYG{l+s+si}{\PYGZcb{}}\PYG{l+s+se}{\PYGZbs{}\PYGZbs{}}\PYG{l+s+si}{\PYGZob{}}\PYG{n}{name}\PYG{l+s+si}{\PYGZcb{}}\PYG{l+s+se}{\PYGZbs{}\PYGZbs{}}\PYG{l+s+si}{\PYGZob{}}\PYG{n}{name}\PYG{l+s+si}{\PYGZcb{}}\PYG{l+s+s1}{\PYGZhy{}zyxp.dat}\PYG{l+s+s1}{\PYGZsq{}}\PYG{p}{)}\PYG{p}{[}\PYG{p}{:}\PYG{p}{,} \PYG{l+m+mi}{1}\PYG{p}{]}
    
    \PYG{c+c1}{\PYGZsh{} 计算子图位置}
    \PYG{n}{row} \PYG{o}{=} \PYG{p}{(}\PYG{n}{i} \PYG{o}{/}\PYG{o}{/} \PYG{n}{ncols}\PYG{p}{)} \PYG{o}{*} \PYG{l+m+mi}{2}
    \PYG{n}{col} \PYG{o}{=} \PYG{n}{i} \PYG{o}{\PYGZpc{}} \PYG{n}{ncols}
    
    \PYG{c+c1}{\PYGZsh{} 视电阻率子图}
    \PYG{n}{ax} \PYG{o}{=} \PYG{n}{axs}\PYG{p}{[}\PYG{n}{row}\PYG{p}{,} \PYG{n}{col}\PYG{p}{]}
    \PYG{n}{ax2} \PYG{o}{=} \PYG{n}{ax}\PYG{o}{.}\PYG{n}{twinx}\PYG{p}{(}\PYG{p}{)}
    
    \PYG{n}{log\PYGZus{}rxy} \PYG{o}{=} \PYG{n}{np}\PYG{o}{.}\PYG{n}{log10}\PYG{p}{(}\PYG{n}{rxy}\PYG{p}{)}
    \PYG{n}{log\PYGZus{}ryx} \PYG{o}{=} \PYG{n}{np}\PYG{o}{.}\PYG{n}{log10}\PYG{p}{(}\PYG{n}{ryx}\PYG{p}{)}
    \PYG{n}{log\PYGZus{}rxy\PYGZus{}var} \PYG{o}{=} \PYG{n}{rxy\PYGZus{}var} \PYG{o}{/} \PYG{p}{(}\PYG{n}{rxy} \PYG{o}{*} \PYG{n}{np}\PYG{o}{.}\PYG{n}{log}\PYG{p}{(}\PYG{l+m+mi}{10}\PYG{p}{)}\PYG{p}{)}
    \PYG{n}{log\PYGZus{}ryx\PYGZus{}var} \PYG{o}{=} \PYG{n}{ryx\PYGZus{}var} \PYG{o}{/} \PYG{p}{(}\PYG{n}{ryx} \PYG{o}{*} \PYG{n}{np}\PYG{o}{.}\PYG{n}{log}\PYG{p}{(}\PYG{l+m+mi}{10}\PYG{p}{)}\PYG{p}{)}
    
    \PYG{n}{ax2}\PYG{o}{.}\PYG{n}{errorbar}\PYG{p}{(}\PYG{n}{fre}\PYG{p}{,} \PYG{n}{log\PYGZus{}rxy}\PYG{p}{,} \PYG{n}{yerr}\PYG{o}{=}\PYG{n}{log\PYGZus{}rxy\PYGZus{}var}\PYG{p}{,} \PYG{n}{fmt}\PYG{o}{=}\PYG{l+s+s1}{\PYGZsq{}}\PYG{l+s+s1}{ro}\PYG{l+s+s1}{\PYGZsq{}}\PYG{p}{,} \PYG{n}{label}\PYG{o}{=}\PYG{l+s+s1}{\PYGZsq{}}\PYG{l+s+s1}{XY}\PYG{l+s+s1}{\PYGZsq{}}\PYG{p}{,} \PYG{n}{markerfacecolor}\PYG{o}{=}\PYG{l+s+s1}{\PYGZsq{}}\PYG{l+s+s1}{none}\PYG{l+s+s1}{\PYGZsq{}}\PYG{p}{)}
    \PYG{n}{ax2}\PYG{o}{.}\PYG{n}{errorbar}\PYG{p}{(}\PYG{n}{fre}\PYG{p}{,} \PYG{n}{log\PYGZus{}ryx}\PYG{p}{,} \PYG{n}{yerr}\PYG{o}{=}\PYG{n}{log\PYGZus{}ryx\PYGZus{}var}\PYG{p}{,} \PYG{n}{fmt}\PYG{o}{=}\PYG{l+s+s1}{\PYGZsq{}}\PYG{l+s+s1}{bs}\PYG{l+s+s1}{\PYGZsq{}}\PYG{p}{,} \PYG{n}{label}\PYG{o}{=}\PYG{l+s+s1}{\PYGZsq{}}\PYG{l+s+s1}{YX}\PYG{l+s+s1}{\PYGZsq{}}\PYG{p}{,} \PYG{n}{markerfacecolor}\PYG{o}{=}\PYG{l+s+s1}{\PYGZsq{}}\PYG{l+s+s1}{none}\PYG{l+s+s1}{\PYGZsq{}}\PYG{p}{)}
    \PYG{n}{ax2}\PYG{o}{.}\PYG{n}{set\PYGZus{}ylim}\PYG{p}{(}\PYG{n}{np}\PYG{o}{.}\PYG{n}{log10}\PYG{p}{(}\PYG{n}{ymin}\PYG{p}{)}\PYG{p}{,} \PYG{n}{np}\PYG{o}{.}\PYG{n}{log10}\PYG{p}{(}\PYG{n}{ymax}\PYG{p}{)}\PYG{p}{)}
    \PYG{n}{ax2}\PYG{o}{.}\PYG{n}{invert\PYGZus{}xaxis}\PYG{p}{(}\PYG{p}{)}
    \PYG{n}{ax2}\PYG{o}{.}\PYG{n}{set\PYGZus{}yticklabels}\PYG{p}{(}\PYG{p}{[}\PYG{p}{]}\PYG{p}{)}
    
    \PYG{n}{ax}\PYG{o}{.}\PYG{n}{set\PYGZus{}ylim}\PYG{p}{(}\PYG{n}{ymin}\PYG{p}{,} \PYG{n}{ymax}\PYG{p}{)}
    \PYG{n}{ax}\PYG{o}{.}\PYG{n}{set\PYGZus{}yscale}\PYG{p}{(}\PYG{l+s+s1}{\PYGZsq{}}\PYG{l+s+s1}{log}\PYG{l+s+s1}{\PYGZsq{}}\PYG{p}{)}
    \PYG{n}{ax}\PYG{o}{.}\PYG{n}{set\PYGZus{}xscale}\PYG{p}{(}\PYG{l+s+s1}{\PYGZsq{}}\PYG{l+s+s1}{log}\PYG{l+s+s1}{\PYGZsq{}}\PYG{p}{)}
    \PYG{n}{ax}\PYG{o}{.}\PYG{n}{set\PYGZus{}xlim}\PYG{p}{(}\PYG{n}{xmin}\PYG{p}{,} \PYG{n}{xmax}\PYG{p}{)}
    \PYG{n}{ax}\PYG{o}{.}\PYG{n}{invert\PYGZus{}xaxis}\PYG{p}{(}\PYG{p}{)}
    \PYG{n}{ax}\PYG{o}{.}\PYG{n}{grid}\PYG{p}{(}\PYG{k+kc}{True}\PYG{p}{,} \PYG{n}{which}\PYG{o}{=}\PYG{l+s+s1}{\PYGZsq{}}\PYG{l+s+s1}{major}\PYG{l+s+s1}{\PYGZsq{}}\PYG{p}{,} \PYG{n}{linestyle}\PYG{o}{=}\PYG{l+s+s1}{\PYGZsq{}}\PYG{l+s+s1}{\PYGZhy{}\PYGZhy{}}\PYG{l+s+s1}{\PYGZsq{}}\PYG{p}{,} \PYG{n}{linewidth}\PYG{o}{=}\PYG{l+m+mf}{0.5}\PYG{p}{)}
    \PYG{n}{ax}\PYG{o}{.}\PYG{n}{text}\PYG{p}{(}\PYG{l+m+mf}{0.5}\PYG{p}{,} \PYG{l+m+mf}{0.88}\PYG{p}{,} \PYG{n}{name}\PYG{p}{,} \PYG{n}{ha}\PYG{o}{=}\PYG{l+s+s1}{\PYGZsq{}}\PYG{l+s+s1}{center}\PYG{l+s+s1}{\PYGZsq{}}\PYG{p}{,} \PYG{n}{va}\PYG{o}{=}\PYG{l+s+s1}{\PYGZsq{}}\PYG{l+s+s1}{bottom}\PYG{l+s+s1}{\PYGZsq{}}\PYG{p}{,} \PYG{n}{transform}\PYG{o}{=}\PYG{n}{ax}\PYG{o}{.}\PYG{n}{transAxes}\PYG{p}{,} 
            \PYG{n}{fontsize}\PYG{o}{=}\PYG{l+m+mi}{12}\PYG{p}{,} \PYG{n}{fontweight}\PYG{o}{=}\PYG{l+s+s1}{\PYGZsq{}}\PYG{l+s+s1}{bold}\PYG{l+s+s1}{\PYGZsq{}}\PYG{p}{)}
    
    \PYG{k}{if} \PYG{n}{col} \PYG{o}{==} \PYG{l+m+mi}{0}\PYG{p}{:}
        \PYG{n}{ax}\PYG{o}{.}\PYG{n}{set\PYGZus{}ylabel}\PYG{p}{(}\PYG{l+s+sa}{r}\PYG{l+s+s1}{\PYGZsq{}}\PYG{l+s+s1}{\PYGZdl{}}\PYG{l+s+s1}{\PYGZbs{}}\PYG{l+s+s1}{rho}\PYG{l+s+s1}{\PYGZbs{}}\PYG{l+s+s1}{;(}\PYG{l+s+s1}{\PYGZbs{}}\PYG{l+s+s1}{Omega }\PYG{l+s+s1}{\PYGZbs{}}\PYG{l+s+s1}{cdot m)\PYGZdl{}}\PYG{l+s+s1}{\PYGZsq{}}\PYG{p}{,} \PYG{n}{fontsize}\PYG{o}{=}\PYG{l+m+mi}{8}\PYG{p}{)}
    
    \PYG{c+c1}{\PYGZsh{} 相位子图}
    \PYG{n}{ax} \PYG{o}{=} \PYG{n}{axs}\PYG{p}{[}\PYG{n}{row} \PYG{o}{+} \PYG{l+m+mi}{1}\PYG{p}{,} \PYG{n}{col}\PYG{p}{]}
    \PYG{n}{ax}\PYG{o}{.}\PYG{n}{errorbar}\PYG{p}{(}\PYG{n}{fre}\PYG{p}{,} \PYG{n}{pxy}\PYG{p}{,} \PYG{n}{fmt}\PYG{o}{=}\PYG{l+s+s1}{\PYGZsq{}}\PYG{l+s+s1}{ro}\PYG{l+s+s1}{\PYGZsq{}}\PYG{p}{,} \PYG{n}{label}\PYG{o}{=}\PYG{l+s+s1}{\PYGZsq{}}\PYG{l+s+s1}{XY}\PYG{l+s+s1}{\PYGZsq{}}\PYG{p}{,} \PYG{n}{markerfacecolor}\PYG{o}{=}\PYG{l+s+s1}{\PYGZsq{}}\PYG{l+s+s1}{none}\PYG{l+s+s1}{\PYGZsq{}}\PYG{p}{)}
    \PYG{n}{ax}\PYG{o}{.}\PYG{n}{errorbar}\PYG{p}{(}\PYG{n}{fre}\PYG{p}{,} \PYG{n}{pyx}\PYG{p}{,} \PYG{n}{fmt}\PYG{o}{=}\PYG{l+s+s1}{\PYGZsq{}}\PYG{l+s+s1}{bs}\PYG{l+s+s1}{\PYGZsq{}}\PYG{p}{,} \PYG{n}{label}\PYG{o}{=}\PYG{l+s+s1}{\PYGZsq{}}\PYG{l+s+s1}{YX}\PYG{l+s+s1}{\PYGZsq{}}\PYG{p}{,} \PYG{n}{markerfacecolor}\PYG{o}{=}\PYG{l+s+s1}{\PYGZsq{}}\PYG{l+s+s1}{none}\PYG{l+s+s1}{\PYGZsq{}}\PYG{p}{)}
    \PYG{n}{ax}\PYG{o}{.}\PYG{n}{set\PYGZus{}ylim}\PYG{p}{(}\PYG{o}{\PYGZhy{}}\PYG{l+m+mi}{180}\PYG{p}{,} \PYG{l+m+mi}{180}\PYG{p}{)}
    \PYG{n}{ax}\PYG{o}{.}\PYG{n}{set\PYGZus{}yticks}\PYG{p}{(}\PYG{n}{np}\PYG{o}{.}\PYG{n}{arange}\PYG{p}{(}\PYG{o}{\PYGZhy{}}\PYG{l+m+mi}{180}\PYG{p}{,} \PYG{l+m+mi}{181}\PYG{p}{,} \PYG{l+m+mi}{90}\PYG{p}{)}\PYG{p}{)}
    \PYG{n}{ax}\PYG{o}{.}\PYG{n}{set\PYGZus{}xscale}\PYG{p}{(}\PYG{l+s+s1}{\PYGZsq{}}\PYG{l+s+s1}{log}\PYG{l+s+s1}{\PYGZsq{}}\PYG{p}{)}
    \PYG{n}{ax}\PYG{o}{.}\PYG{n}{grid}\PYG{p}{(}\PYG{k+kc}{True}\PYG{p}{,} \PYG{n}{which}\PYG{o}{=}\PYG{l+s+s1}{\PYGZsq{}}\PYG{l+s+s1}{major}\PYG{l+s+s1}{\PYGZsq{}}\PYG{p}{,} \PYG{n}{linestyle}\PYG{o}{=}\PYG{l+s+s1}{\PYGZsq{}}\PYG{l+s+s1}{\PYGZhy{}\PYGZhy{}}\PYG{l+s+s1}{\PYGZsq{}}\PYG{p}{,} \PYG{n}{linewidth}\PYG{o}{=}\PYG{l+m+mf}{0.5}\PYG{p}{)}
    
    \PYG{k}{if} \PYG{n}{row} \PYG{o}{+} \PYG{l+m+mi}{1} \PYG{o}{==} \PYG{n}{nrows} \PYG{o}{*} \PYG{l+m+mi}{2} \PYG{o}{\PYGZhy{}} \PYG{l+m+mi}{1}\PYG{p}{:}
        \PYG{n}{ax}\PYG{o}{.}\PYG{n}{set\PYGZus{}xlabel}\PYG{p}{(}\PYG{l+s+s1}{\PYGZsq{}}\PYG{l+s+s1}{Frequency (Hz)}\PYG{l+s+s1}{\PYGZsq{}}\PYG{p}{,} \PYG{n}{fontsize}\PYG{o}{=}\PYG{l+m+mi}{8}\PYG{p}{)}
    \PYG{k}{if} \PYG{n}{col} \PYG{o}{==} \PYG{l+m+mi}{0}\PYG{p}{:}
        \PYG{n}{ax}\PYG{o}{.}\PYG{n}{set\PYGZus{}ylabel}\PYG{p}{(}\PYG{l+s+sa}{r}\PYG{l+s+s1}{\PYGZsq{}}\PYG{l+s+s1}{\PYGZdl{}}\PYG{l+s+s1}{\PYGZbs{}}\PYG{l+s+s1}{varphi}\PYG{l+s+s1}{\PYGZbs{}}\PYG{l+s+s1}{;(}\PYG{l+s+s1}{\PYGZbs{}}\PYG{l+s+s1}{degree)\PYGZdl{}}\PYG{l+s+s1}{\PYGZsq{}}\PYG{p}{,} \PYG{n}{fontsize}\PYG{o}{=}\PYG{l+m+mi}{8}\PYG{p}{)}

\PYG{n}{fig}\PYG{o}{.}\PYG{n}{subplots\PYGZus{}adjust}\PYG{p}{(}\PYG{n}{hspace}\PYG{o}{=}\PYG{l+m+mf}{0.4}\PYG{p}{,} \PYG{n}{wspace}\PYG{o}{=}\PYG{l+m+mf}{0.4}\PYG{p}{)}
\PYG{n}{plt}\PYG{o}{.}\PYG{n}{savefig}\PYG{p}{(}\PYG{l+s+sa}{f}\PYG{l+s+s1}{\PYGZsq{}}\PYG{l+s+si}{\PYGZob{}}\PYG{n}{path}\PYG{l+s+si}{\PYGZcb{}}\PYG{l+s+se}{\PYGZbs{}\PYGZbs{}}\PYG{l+s+s1}{All\PYGZus{}Sites.png}\PYG{l+s+s1}{\PYGZsq{}}\PYG{p}{,} \PYG{n}{dpi}\PYG{o}{=}\PYG{l+m+mi}{300}\PYG{p}{,} \PYG{n}{bbox\PYGZus{}inches}\PYG{o}{=}\PYG{l+s+s1}{\PYGZsq{}}\PYG{l+s+s1}{tight}\PYG{l+s+s1}{\PYGZsq{}}\PYG{p}{)}
\PYG{n}{plt}\PYG{o}{.}\PYG{n}{savefig}\PYG{p}{(}\PYG{l+s+sa}{f}\PYG{l+s+s1}{\PYGZsq{}}\PYG{l+s+si}{\PYGZob{}}\PYG{n}{path}\PYG{l+s+si}{\PYGZcb{}}\PYG{l+s+se}{\PYGZbs{}\PYGZbs{}}\PYG{l+s+s1}{All\PYGZus{}Sites.pdf}\PYG{l+s+s1}{\PYGZsq{}}\PYG{p}{,} \PYG{n}{dpi}\PYG{o}{=}\PYG{l+m+mi}{300}\PYG{p}{,} \PYG{n}{bbox\PYGZus{}inches}\PYG{o}{=}\PYG{l+s+s1}{\PYGZsq{}}\PYG{l+s+s1}{tight}\PYG{l+s+s1}{\PYGZsq{}}\PYG{p}{)}
\PYG{n}{plt}\PYG{o}{.}\PYG{n}{savefig}\PYG{p}{(}\PYG{l+s+sa}{f}\PYG{l+s+s1}{\PYGZsq{}}\PYG{l+s+si}{\PYGZob{}}\PYG{n}{path}\PYG{l+s+si}{\PYGZcb{}}\PYG{l+s+se}{\PYGZbs{}\PYGZbs{}}\PYG{l+s+s1}{All\PYGZus{}Sites.svg}\PYG{l+s+s1}{\PYGZsq{}}\PYG{p}{,} \PYG{n}{dpi}\PYG{o}{=}\PYG{l+m+mi}{300}\PYG{p}{,} \PYG{n}{bbox\PYGZus{}inches}\PYG{o}{=}\PYG{l+s+s1}{\PYGZsq{}}\PYG{l+s+s1}{tight}\PYG{l+s+s1}{\PYGZsq{}}\PYG{p}{)}

\PYG{n+nb}{print}\PYG{p}{(}\PYG{l+s+s1}{\PYGZsq{}}\PYG{l+s+se}{\PYGZbs{}n}\PYG{l+s+s1}{Finished!}\PYG{l+s+s1}{\PYGZsq{}}\PYG{p}{)}
\end{sphinxVerbatim}


\bigskip\hrule\bigskip



\subsection{测点对比分析}
\label{\detokenize{tutorial/plotting:id8}}
\sphinxAtStartPar
比较两个测点或两组数据的差异，计算 RMS 误差。


\subsubsection{脚本：\sphinxhref{http://PlotFullZ-Compare.py}{PlotFullZ\sphinxhyphen{}Compare.py} (http://PlotFullZ\sphinxhyphen{}Compare.py)}
\label{\detokenize{tutorial/plotting:plotfullz-compare-py}}
\begin{sphinxVerbatim}[commandchars=\\\{\}]
\PYG{k+kn}{import}\PYG{+w}{ }\PYG{n+nn}{os}
\PYG{k+kn}{import}\PYG{+w}{ }\PYG{n+nn}{numpy}\PYG{+w}{ }\PYG{k}{as}\PYG{+w}{ }\PYG{n+nn}{np}
\PYG{k+kn}{import}\PYG{+w}{ }\PYG{n+nn}{matplotlib}\PYG{n+nn}{.}\PYG{n+nn}{pyplot}\PYG{+w}{ }\PYG{k}{as}\PYG{+w}{ }\PYG{n+nn}{plt}

\PYG{k}{def}\PYG{+w}{ }\PYG{n+nf}{plotMTDP}\PYG{p}{(}\PYG{n}{path}\PYG{p}{,} \PYG{n}{name1}\PYG{p}{,} \PYG{n}{name2}\PYG{p}{)}\PYG{p}{:}
\PYG{+w}{    }\PYG{l+s+sd}{\PYGZdq{}\PYGZdq{}\PYGZdq{}绘制两个测点的对比图\PYGZdq{}\PYGZdq{}\PYGZdq{}}
    
    \PYG{c+c1}{\PYGZsh{} 配置参数}
    \PYG{n}{ymax} \PYG{o}{=} \PYG{l+m+mf}{1e3}
    \PYG{n}{ymin} \PYG{o}{=} \PYG{l+m+mf}{1e0}
    \PYG{n}{xmax} \PYG{o}{=} \PYG{l+m+mf}{1e3}
    \PYG{n}{xmin} \PYG{o}{=} \PYG{l+m+mf}{5e\PYGZhy{}4}
    
    \PYG{c+c1}{\PYGZsh{} 创建输出目录}
    \PYG{n}{out\PYGZus{}dir} \PYG{o}{=} \PYG{n}{os}\PYG{o}{.}\PYG{n}{path}\PYG{o}{.}\PYG{n}{join}\PYG{p}{(}\PYG{n}{path}\PYG{p}{,} \PYG{l+s+s1}{\PYGZsq{}}\PYG{l+s+s1}{Compare}\PYG{l+s+s1}{\PYGZsq{}}\PYG{p}{)}
    \PYG{k}{if} \PYG{o+ow}{not} \PYG{n}{os}\PYG{o}{.}\PYG{n}{path}\PYG{o}{.}\PYG{n}{exists}\PYG{p}{(}\PYG{n}{out\PYGZus{}dir}\PYG{p}{)}\PYG{p}{:}
        \PYG{n}{os}\PYG{o}{.}\PYG{n}{makedirs}\PYG{p}{(}\PYG{n}{out\PYGZus{}dir}\PYG{p}{)}
    
    \PYG{c+c1}{\PYGZsh{} 读取测点1数据}
    \PYG{n}{fre1} \PYG{o}{=} \PYG{n}{np}\PYG{o}{.}\PYG{n}{loadtxt}\PYG{p}{(}\PYG{l+s+sa}{f}\PYG{l+s+s1}{\PYGZsq{}}\PYG{l+s+si}{\PYGZob{}}\PYG{n}{path}\PYG{l+s+si}{\PYGZcb{}}\PYG{l+s+se}{\PYGZbs{}\PYGZbs{}}\PYG{l+s+si}{\PYGZob{}}\PYG{n}{name1}\PYG{l+s+si}{\PYGZcb{}}\PYG{l+s+se}{\PYGZbs{}\PYGZbs{}}\PYG{l+s+si}{\PYGZob{}}\PYG{n}{name1}\PYG{l+s+si}{\PYGZcb{}}\PYG{l+s+s1}{\PYGZhy{}rxy.dat}\PYG{l+s+s1}{\PYGZsq{}}\PYG{p}{)}\PYG{p}{[}\PYG{p}{:}\PYG{p}{,} \PYG{l+m+mi}{0}\PYG{p}{]}
    \PYG{n}{rxy1} \PYG{o}{=} \PYG{n}{np}\PYG{o}{.}\PYG{n}{loadtxt}\PYG{p}{(}\PYG{l+s+sa}{f}\PYG{l+s+s1}{\PYGZsq{}}\PYG{l+s+si}{\PYGZob{}}\PYG{n}{path}\PYG{l+s+si}{\PYGZcb{}}\PYG{l+s+se}{\PYGZbs{}\PYGZbs{}}\PYG{l+s+si}{\PYGZob{}}\PYG{n}{name1}\PYG{l+s+si}{\PYGZcb{}}\PYG{l+s+se}{\PYGZbs{}\PYGZbs{}}\PYG{l+s+si}{\PYGZob{}}\PYG{n}{name1}\PYG{l+s+si}{\PYGZcb{}}\PYG{l+s+s1}{\PYGZhy{}rxy.dat}\PYG{l+s+s1}{\PYGZsq{}}\PYG{p}{)}\PYG{p}{[}\PYG{p}{:}\PYG{p}{,} \PYG{l+m+mi}{1}\PYG{p}{]}
    \PYG{n}{ryx1} \PYG{o}{=} \PYG{n}{np}\PYG{o}{.}\PYG{n}{loadtxt}\PYG{p}{(}\PYG{l+s+sa}{f}\PYG{l+s+s1}{\PYGZsq{}}\PYG{l+s+si}{\PYGZob{}}\PYG{n}{path}\PYG{l+s+si}{\PYGZcb{}}\PYG{l+s+se}{\PYGZbs{}\PYGZbs{}}\PYG{l+s+si}{\PYGZob{}}\PYG{n}{name1}\PYG{l+s+si}{\PYGZcb{}}\PYG{l+s+se}{\PYGZbs{}\PYGZbs{}}\PYG{l+s+si}{\PYGZob{}}\PYG{n}{name1}\PYG{l+s+si}{\PYGZcb{}}\PYG{l+s+s1}{\PYGZhy{}ryx.dat}\PYG{l+s+s1}{\PYGZsq{}}\PYG{p}{)}\PYG{p}{[}\PYG{p}{:}\PYG{p}{,} \PYG{l+m+mi}{1}\PYG{p}{]}
    \PYG{n}{pxy1} \PYG{o}{=} \PYG{n}{np}\PYG{o}{.}\PYG{n}{loadtxt}\PYG{p}{(}\PYG{l+s+sa}{f}\PYG{l+s+s1}{\PYGZsq{}}\PYG{l+s+si}{\PYGZob{}}\PYG{n}{path}\PYG{l+s+si}{\PYGZcb{}}\PYG{l+s+se}{\PYGZbs{}\PYGZbs{}}\PYG{l+s+si}{\PYGZob{}}\PYG{n}{name1}\PYG{l+s+si}{\PYGZcb{}}\PYG{l+s+se}{\PYGZbs{}\PYGZbs{}}\PYG{l+s+si}{\PYGZob{}}\PYG{n}{name1}\PYG{l+s+si}{\PYGZcb{}}\PYG{l+s+s1}{\PYGZhy{}zxyp.dat}\PYG{l+s+s1}{\PYGZsq{}}\PYG{p}{)}\PYG{p}{[}\PYG{p}{:}\PYG{p}{,} \PYG{l+m+mi}{1}\PYG{p}{]}
    \PYG{n}{pyx1} \PYG{o}{=} \PYG{n}{np}\PYG{o}{.}\PYG{n}{loadtxt}\PYG{p}{(}\PYG{l+s+sa}{f}\PYG{l+s+s1}{\PYGZsq{}}\PYG{l+s+si}{\PYGZob{}}\PYG{n}{path}\PYG{l+s+si}{\PYGZcb{}}\PYG{l+s+se}{\PYGZbs{}\PYGZbs{}}\PYG{l+s+si}{\PYGZob{}}\PYG{n}{name1}\PYG{l+s+si}{\PYGZcb{}}\PYG{l+s+se}{\PYGZbs{}\PYGZbs{}}\PYG{l+s+si}{\PYGZob{}}\PYG{n}{name1}\PYG{l+s+si}{\PYGZcb{}}\PYG{l+s+s1}{\PYGZhy{}zyxp.dat}\PYG{l+s+s1}{\PYGZsq{}}\PYG{p}{)}\PYG{p}{[}\PYG{p}{:}\PYG{p}{,} \PYG{l+m+mi}{1}\PYG{p}{]}
    
    \PYG{c+c1}{\PYGZsh{} 读取测点2数据}
    \PYG{n}{fre2} \PYG{o}{=} \PYG{n}{np}\PYG{o}{.}\PYG{n}{loadtxt}\PYG{p}{(}\PYG{l+s+sa}{f}\PYG{l+s+s1}{\PYGZsq{}}\PYG{l+s+si}{\PYGZob{}}\PYG{n}{path}\PYG{l+s+si}{\PYGZcb{}}\PYG{l+s+se}{\PYGZbs{}\PYGZbs{}}\PYG{l+s+si}{\PYGZob{}}\PYG{n}{name2}\PYG{l+s+si}{\PYGZcb{}}\PYG{l+s+se}{\PYGZbs{}\PYGZbs{}}\PYG{l+s+si}{\PYGZob{}}\PYG{n}{name2}\PYG{l+s+si}{\PYGZcb{}}\PYG{l+s+s1}{\PYGZhy{}rxy.dat}\PYG{l+s+s1}{\PYGZsq{}}\PYG{p}{)}\PYG{p}{[}\PYG{p}{:}\PYG{p}{,} \PYG{l+m+mi}{0}\PYG{p}{]}
    \PYG{n}{rxy2} \PYG{o}{=} \PYG{n}{np}\PYG{o}{.}\PYG{n}{loadtxt}\PYG{p}{(}\PYG{l+s+sa}{f}\PYG{l+s+s1}{\PYGZsq{}}\PYG{l+s+si}{\PYGZob{}}\PYG{n}{path}\PYG{l+s+si}{\PYGZcb{}}\PYG{l+s+se}{\PYGZbs{}\PYGZbs{}}\PYG{l+s+si}{\PYGZob{}}\PYG{n}{name2}\PYG{l+s+si}{\PYGZcb{}}\PYG{l+s+se}{\PYGZbs{}\PYGZbs{}}\PYG{l+s+si}{\PYGZob{}}\PYG{n}{name2}\PYG{l+s+si}{\PYGZcb{}}\PYG{l+s+s1}{\PYGZhy{}rxy.dat}\PYG{l+s+s1}{\PYGZsq{}}\PYG{p}{)}\PYG{p}{[}\PYG{p}{:}\PYG{p}{,} \PYG{l+m+mi}{1}\PYG{p}{]}
    \PYG{n}{ryx2} \PYG{o}{=} \PYG{n}{np}\PYG{o}{.}\PYG{n}{loadtxt}\PYG{p}{(}\PYG{l+s+sa}{f}\PYG{l+s+s1}{\PYGZsq{}}\PYG{l+s+si}{\PYGZob{}}\PYG{n}{path}\PYG{l+s+si}{\PYGZcb{}}\PYG{l+s+se}{\PYGZbs{}\PYGZbs{}}\PYG{l+s+si}{\PYGZob{}}\PYG{n}{name2}\PYG{l+s+si}{\PYGZcb{}}\PYG{l+s+se}{\PYGZbs{}\PYGZbs{}}\PYG{l+s+si}{\PYGZob{}}\PYG{n}{name2}\PYG{l+s+si}{\PYGZcb{}}\PYG{l+s+s1}{\PYGZhy{}ryx.dat}\PYG{l+s+s1}{\PYGZsq{}}\PYG{p}{)}\PYG{p}{[}\PYG{p}{:}\PYG{p}{,} \PYG{l+m+mi}{1}\PYG{p}{]}
    \PYG{n}{pxy2} \PYG{o}{=} \PYG{n}{np}\PYG{o}{.}\PYG{n}{loadtxt}\PYG{p}{(}\PYG{l+s+sa}{f}\PYG{l+s+s1}{\PYGZsq{}}\PYG{l+s+si}{\PYGZob{}}\PYG{n}{path}\PYG{l+s+si}{\PYGZcb{}}\PYG{l+s+se}{\PYGZbs{}\PYGZbs{}}\PYG{l+s+si}{\PYGZob{}}\PYG{n}{name2}\PYG{l+s+si}{\PYGZcb{}}\PYG{l+s+se}{\PYGZbs{}\PYGZbs{}}\PYG{l+s+si}{\PYGZob{}}\PYG{n}{name2}\PYG{l+s+si}{\PYGZcb{}}\PYG{l+s+s1}{\PYGZhy{}zxyp.dat}\PYG{l+s+s1}{\PYGZsq{}}\PYG{p}{)}\PYG{p}{[}\PYG{p}{:}\PYG{p}{,} \PYG{l+m+mi}{1}\PYG{p}{]}
    \PYG{n}{pyx2} \PYG{o}{=} \PYG{n}{np}\PYG{o}{.}\PYG{n}{loadtxt}\PYG{p}{(}\PYG{l+s+sa}{f}\PYG{l+s+s1}{\PYGZsq{}}\PYG{l+s+si}{\PYGZob{}}\PYG{n}{path}\PYG{l+s+si}{\PYGZcb{}}\PYG{l+s+se}{\PYGZbs{}\PYGZbs{}}\PYG{l+s+si}{\PYGZob{}}\PYG{n}{name2}\PYG{l+s+si}{\PYGZcb{}}\PYG{l+s+se}{\PYGZbs{}\PYGZbs{}}\PYG{l+s+si}{\PYGZob{}}\PYG{n}{name2}\PYG{l+s+si}{\PYGZcb{}}\PYG{l+s+s1}{\PYGZhy{}zyxp.dat}\PYG{l+s+s1}{\PYGZsq{}}\PYG{p}{)}\PYG{p}{[}\PYG{p}{:}\PYG{p}{,} \PYG{l+m+mi}{1}\PYG{p}{]}
    
    \PYG{c+c1}{\PYGZsh{} 计算 RMS 误差}
    \PYG{n}{endindex} \PYG{o}{=} \PYG{l+m+mi}{39}
    \PYG{n}{rxyerror} \PYG{o}{=} \PYG{p}{(}\PYG{n}{rxy1}\PYG{p}{[}\PYG{l+m+mi}{0}\PYG{p}{:}\PYG{n}{endindex}\PYG{p}{]} \PYG{o}{\PYGZhy{}} \PYG{n}{rxy2}\PYG{p}{[}\PYG{l+m+mi}{0}\PYG{p}{:}\PYG{n}{endindex}\PYG{p}{]}\PYG{p}{)} \PYG{o}{/} \PYG{n}{rxy1}\PYG{p}{[}\PYG{l+m+mi}{0}\PYG{p}{:}\PYG{n}{endindex}\PYG{p}{]}
    \PYG{n}{rxyrms} \PYG{o}{=} \PYG{n}{np}\PYG{o}{.}\PYG{n}{sqrt}\PYG{p}{(}\PYG{n}{np}\PYG{o}{.}\PYG{n}{mean}\PYG{p}{(}\PYG{n}{rxyerror}\PYG{o}{*}\PYG{o}{*}\PYG{l+m+mi}{2}\PYG{p}{)} \PYG{o}{/} \PYG{l+m+mi}{2}\PYG{p}{)}
    
    \PYG{n}{ryxerror} \PYG{o}{=} \PYG{p}{(}\PYG{n}{ryx1}\PYG{p}{[}\PYG{l+m+mi}{0}\PYG{p}{:}\PYG{n}{endindex}\PYG{p}{]} \PYG{o}{\PYGZhy{}} \PYG{n}{ryx2}\PYG{p}{[}\PYG{l+m+mi}{0}\PYG{p}{:}\PYG{n}{endindex}\PYG{p}{]}\PYG{p}{)} \PYG{o}{/} \PYG{n}{ryx1}\PYG{p}{[}\PYG{l+m+mi}{0}\PYG{p}{:}\PYG{n}{endindex}\PYG{p}{]}
    \PYG{n}{ryxrms} \PYG{o}{=} \PYG{n}{np}\PYG{o}{.}\PYG{n}{sqrt}\PYG{p}{(}\PYG{n}{np}\PYG{o}{.}\PYG{n}{mean}\PYG{p}{(}\PYG{n}{ryxerror}\PYG{o}{*}\PYG{o}{*}\PYG{l+m+mi}{2}\PYG{p}{)} \PYG{o}{/} \PYG{l+m+mi}{2}\PYG{p}{)}
    
    \PYG{n}{pxyerror} \PYG{o}{=} \PYG{p}{(}\PYG{n}{pxy1}\PYG{p}{[}\PYG{l+m+mi}{0}\PYG{p}{:}\PYG{n}{endindex}\PYG{p}{]} \PYG{o}{\PYGZhy{}} \PYG{n}{pxy2}\PYG{p}{[}\PYG{l+m+mi}{0}\PYG{p}{:}\PYG{n}{endindex}\PYG{p}{]}\PYG{p}{)} \PYG{o}{/} \PYG{n}{pxy1}\PYG{p}{[}\PYG{l+m+mi}{0}\PYG{p}{:}\PYG{n}{endindex}\PYG{p}{]}
    \PYG{n}{pxyrms} \PYG{o}{=} \PYG{n}{np}\PYG{o}{.}\PYG{n}{sqrt}\PYG{p}{(}\PYG{n}{np}\PYG{o}{.}\PYG{n}{mean}\PYG{p}{(}\PYG{n}{pxyerror}\PYG{o}{*}\PYG{o}{*}\PYG{l+m+mi}{2}\PYG{p}{)} \PYG{o}{/} \PYG{l+m+mi}{2}\PYG{p}{)}
    
    \PYG{n}{pyxerror} \PYG{o}{=} \PYG{p}{(}\PYG{n}{pyx1}\PYG{p}{[}\PYG{l+m+mi}{0}\PYG{p}{:}\PYG{n}{endindex}\PYG{p}{]} \PYG{o}{\PYGZhy{}} \PYG{n}{pyx2}\PYG{p}{[}\PYG{l+m+mi}{0}\PYG{p}{:}\PYG{n}{endindex}\PYG{p}{]}\PYG{p}{)} \PYG{o}{/} \PYG{n}{pyx1}\PYG{p}{[}\PYG{l+m+mi}{0}\PYG{p}{:}\PYG{n}{endindex}\PYG{p}{]}
    \PYG{n}{pyxrms} \PYG{o}{=} \PYG{n}{np}\PYG{o}{.}\PYG{n}{sqrt}\PYG{p}{(}\PYG{n}{np}\PYG{o}{.}\PYG{n}{mean}\PYG{p}{(}\PYG{n}{pyxerror}\PYG{o}{*}\PYG{o}{*}\PYG{l+m+mi}{2}\PYG{p}{)} \PYG{o}{/} \PYG{l+m+mi}{2}\PYG{p}{)}
    
    \PYG{n+nb}{print}\PYG{p}{(}\PYG{l+s+sa}{f}\PYG{l+s+s1}{\PYGZsq{}}\PYG{l+s+s1}{RMS Error Rxy: }\PYG{l+s+si}{\PYGZob{}}\PYG{n}{rxyrms}\PYG{o}{*}\PYG{l+m+mi}{100}\PYG{l+s+si}{:}\PYG{l+s+s1}{.1f}\PYG{l+s+si}{\PYGZcb{}}\PYG{l+s+s1}{\PYGZpc{}, Ryx: }\PYG{l+s+si}{\PYGZob{}}\PYG{n}{ryxrms}\PYG{o}{*}\PYG{l+m+mi}{100}\PYG{l+s+si}{:}\PYG{l+s+s1}{.1f}\PYG{l+s+si}{\PYGZcb{}}\PYG{l+s+s1}{\PYGZpc{}}\PYG{l+s+s1}{\PYGZsq{}}\PYG{p}{)}
    \PYG{n+nb}{print}\PYG{p}{(}\PYG{l+s+sa}{f}\PYG{l+s+s1}{\PYGZsq{}}\PYG{l+s+s1}{RMS Error Pxy: }\PYG{l+s+si}{\PYGZob{}}\PYG{n}{pxyrms}\PYG{o}{*}\PYG{l+m+mi}{100}\PYG{l+s+si}{:}\PYG{l+s+s1}{.1f}\PYG{l+s+si}{\PYGZcb{}}\PYG{l+s+s1}{\PYGZpc{}, Pyx: }\PYG{l+s+si}{\PYGZob{}}\PYG{n}{pyxrms}\PYG{o}{*}\PYG{l+m+mi}{100}\PYG{l+s+si}{:}\PYG{l+s+s1}{.1f}\PYG{l+s+si}{\PYGZcb{}}\PYG{l+s+s1}{\PYGZpc{}}\PYG{l+s+s1}{\PYGZsq{}}\PYG{p}{)}
    
    \PYG{c+c1}{\PYGZsh{} 创建 2x2 子图}
    \PYG{n}{fig}\PYG{p}{,} \PYG{n}{axs} \PYG{o}{=} \PYG{n}{plt}\PYG{o}{.}\PYG{n}{subplots}\PYG{p}{(}\PYG{l+m+mi}{2}\PYG{p}{,} \PYG{l+m+mi}{2}\PYG{p}{,} \PYG{n}{figsize}\PYG{o}{=}\PYG{p}{(}\PYG{l+m+mi}{10}\PYG{p}{,} \PYG{l+m+mi}{6}\PYG{p}{)}\PYG{p}{,} \PYG{n}{sharex}\PYG{o}{=}\PYG{l+s+s1}{\PYGZsq{}}\PYG{l+s+s1}{all}\PYG{l+s+s1}{\PYGZsq{}}\PYG{p}{,} \PYG{n}{constrained\PYGZus{}layout}\PYG{o}{=}\PYG{k+kc}{True}\PYG{p}{)}
    
    \PYG{c+c1}{\PYGZsh{} ρxy}
    \PYG{n}{axs}\PYG{p}{[}\PYG{l+m+mi}{0}\PYG{p}{,}\PYG{l+m+mi}{0}\PYG{p}{]}\PYG{o}{.}\PYG{n}{loglog}\PYG{p}{(}\PYG{n}{fre1}\PYG{p}{,} \PYG{n}{rxy1}\PYG{p}{,} \PYG{l+s+s1}{\PYGZsq{}}\PYG{l+s+s1}{r}\PYG{l+s+s1}{\PYGZsq{}}\PYG{p}{,} \PYG{n}{label}\PYG{o}{=}\PYG{n}{name1}\PYG{p}{)}
    \PYG{n}{axs}\PYG{p}{[}\PYG{l+m+mi}{0}\PYG{p}{,}\PYG{l+m+mi}{0}\PYG{p}{]}\PYG{o}{.}\PYG{n}{loglog}\PYG{p}{(}\PYG{n}{fre2}\PYG{p}{,} \PYG{n}{rxy2}\PYG{p}{,} \PYG{l+s+s1}{\PYGZsq{}}\PYG{l+s+s1}{b}\PYG{l+s+s1}{\PYGZsq{}}\PYG{p}{,} \PYG{n}{label}\PYG{o}{=}\PYG{n}{name2}\PYG{p}{)}
    \PYG{n}{axs}\PYG{p}{[}\PYG{l+m+mi}{0}\PYG{p}{,}\PYG{l+m+mi}{0}\PYG{p}{]}\PYG{o}{.}\PYG{n}{set\PYGZus{}ylim}\PYG{p}{(}\PYG{n}{ymin}\PYG{p}{,} \PYG{n}{ymax}\PYG{p}{)}
    \PYG{n}{axs}\PYG{p}{[}\PYG{l+m+mi}{0}\PYG{p}{,}\PYG{l+m+mi}{0}\PYG{p}{]}\PYG{o}{.}\PYG{n}{set\PYGZus{}xlim}\PYG{p}{(}\PYG{n}{xmin}\PYG{p}{,} \PYG{n}{xmax}\PYG{p}{)}
    \PYG{n}{axs}\PYG{p}{[}\PYG{l+m+mi}{0}\PYG{p}{,}\PYG{l+m+mi}{0}\PYG{p}{]}\PYG{o}{.}\PYG{n}{invert\PYGZus{}xaxis}\PYG{p}{(}\PYG{p}{)}
    \PYG{n}{axs}\PYG{p}{[}\PYG{l+m+mi}{0}\PYG{p}{,}\PYG{l+m+mi}{0}\PYG{p}{]}\PYG{o}{.}\PYG{n}{set\PYGZus{}ylabel}\PYG{p}{(}\PYG{l+s+sa}{r}\PYG{l+s+s1}{\PYGZsq{}}\PYG{l+s+s1}{\PYGZdl{}}\PYG{l+s+s1}{\PYGZbs{}}\PYG{l+s+s1}{rho}\PYG{l+s+s1}{\PYGZbs{}}\PYG{l+s+s1}{;(}\PYG{l+s+s1}{\PYGZbs{}}\PYG{l+s+s1}{Omega }\PYG{l+s+s1}{\PYGZbs{}}\PYG{l+s+s1}{cdot m)\PYGZdl{}}\PYG{l+s+s1}{\PYGZsq{}}\PYG{p}{)}
    \PYG{n}{axs}\PYG{p}{[}\PYG{l+m+mi}{0}\PYG{p}{,}\PYG{l+m+mi}{0}\PYG{p}{]}\PYG{o}{.}\PYG{n}{set\PYGZus{}title}\PYG{p}{(}\PYG{l+s+sa}{r}\PYG{l+s+s1}{\PYGZsq{}}\PYG{l+s+s1}{\PYGZdl{}}\PYG{l+s+s1}{\PYGZbs{}}\PYG{l+s+s1}{rho\PYGZus{}}\PYG{l+s+si}{\PYGZob{}xy\PYGZcb{}}\PYG{l+s+s1}{\PYGZdl{} RMS = }\PYG{l+s+s1}{\PYGZsq{}} \PYG{o}{+} \PYG{l+s+sa}{f}\PYG{l+s+s1}{\PYGZsq{}}\PYG{l+s+si}{\PYGZob{}}\PYG{n}{rxyrms}\PYG{o}{*}\PYG{l+m+mi}{100}\PYG{l+s+si}{:}\PYG{l+s+s1}{.1f}\PYG{l+s+si}{\PYGZcb{}}\PYG{l+s+s1}{\PYGZpc{}}\PYG{l+s+s1}{\PYGZsq{}}\PYG{p}{)}
    \PYG{n}{axs}\PYG{p}{[}\PYG{l+m+mi}{0}\PYG{p}{,}\PYG{l+m+mi}{0}\PYG{p}{]}\PYG{o}{.}\PYG{n}{grid}\PYG{p}{(}\PYG{k+kc}{True}\PYG{p}{,} \PYG{n}{which}\PYG{o}{=}\PYG{l+s+s1}{\PYGZsq{}}\PYG{l+s+s1}{major}\PYG{l+s+s1}{\PYGZsq{}}\PYG{p}{,} \PYG{n}{linestyle}\PYG{o}{=}\PYG{l+s+s1}{\PYGZsq{}}\PYG{l+s+s1}{\PYGZhy{}\PYGZhy{}}\PYG{l+s+s1}{\PYGZsq{}}\PYG{p}{,} \PYG{n}{linewidth}\PYG{o}{=}\PYG{l+m+mf}{0.5}\PYG{p}{)}
    \PYG{n}{axs}\PYG{p}{[}\PYG{l+m+mi}{0}\PYG{p}{,}\PYG{l+m+mi}{0}\PYG{p}{]}\PYG{o}{.}\PYG{n}{legend}\PYG{p}{(}\PYG{p}{)}
    
    \PYG{c+c1}{\PYGZsh{} ρyx}
    \PYG{n}{axs}\PYG{p}{[}\PYG{l+m+mi}{0}\PYG{p}{,}\PYG{l+m+mi}{1}\PYG{p}{]}\PYG{o}{.}\PYG{n}{loglog}\PYG{p}{(}\PYG{n}{fre1}\PYG{p}{,} \PYG{n}{ryx1}\PYG{p}{,} \PYG{l+s+s1}{\PYGZsq{}}\PYG{l+s+s1}{r}\PYG{l+s+s1}{\PYGZsq{}}\PYG{p}{,} \PYG{n}{label}\PYG{o}{=}\PYG{n}{name1}\PYG{p}{)}
    \PYG{n}{axs}\PYG{p}{[}\PYG{l+m+mi}{0}\PYG{p}{,}\PYG{l+m+mi}{1}\PYG{p}{]}\PYG{o}{.}\PYG{n}{loglog}\PYG{p}{(}\PYG{n}{fre2}\PYG{p}{,} \PYG{n}{ryx2}\PYG{p}{,} \PYG{l+s+s1}{\PYGZsq{}}\PYG{l+s+s1}{b}\PYG{l+s+s1}{\PYGZsq{}}\PYG{p}{,} \PYG{n}{label}\PYG{o}{=}\PYG{n}{name2}\PYG{p}{)}
    \PYG{n}{axs}\PYG{p}{[}\PYG{l+m+mi}{0}\PYG{p}{,}\PYG{l+m+mi}{1}\PYG{p}{]}\PYG{o}{.}\PYG{n}{set\PYGZus{}ylim}\PYG{p}{(}\PYG{n}{ymin}\PYG{p}{,} \PYG{n}{ymax}\PYG{p}{)}
    \PYG{n}{axs}\PYG{p}{[}\PYG{l+m+mi}{0}\PYG{p}{,}\PYG{l+m+mi}{1}\PYG{p}{]}\PYG{o}{.}\PYG{n}{set\PYGZus{}xlim}\PYG{p}{(}\PYG{n}{xmin}\PYG{p}{,} \PYG{n}{xmax}\PYG{p}{)}
    \PYG{n}{axs}\PYG{p}{[}\PYG{l+m+mi}{0}\PYG{p}{,}\PYG{l+m+mi}{1}\PYG{p}{]}\PYG{o}{.}\PYG{n}{invert\PYGZus{}xaxis}\PYG{p}{(}\PYG{p}{)}
    \PYG{n}{axs}\PYG{p}{[}\PYG{l+m+mi}{0}\PYG{p}{,}\PYG{l+m+mi}{1}\PYG{p}{]}\PYG{o}{.}\PYG{n}{set\PYGZus{}title}\PYG{p}{(}\PYG{l+s+sa}{r}\PYG{l+s+s1}{\PYGZsq{}}\PYG{l+s+s1}{\PYGZdl{}}\PYG{l+s+s1}{\PYGZbs{}}\PYG{l+s+s1}{rho\PYGZus{}}\PYG{l+s+si}{\PYGZob{}yx\PYGZcb{}}\PYG{l+s+s1}{\PYGZdl{} RMS = }\PYG{l+s+s1}{\PYGZsq{}} \PYG{o}{+} \PYG{l+s+sa}{f}\PYG{l+s+s1}{\PYGZsq{}}\PYG{l+s+si}{\PYGZob{}}\PYG{n}{ryxrms}\PYG{o}{*}\PYG{l+m+mi}{100}\PYG{l+s+si}{:}\PYG{l+s+s1}{.1f}\PYG{l+s+si}{\PYGZcb{}}\PYG{l+s+s1}{\PYGZpc{}}\PYG{l+s+s1}{\PYGZsq{}}\PYG{p}{)}
    \PYG{n}{axs}\PYG{p}{[}\PYG{l+m+mi}{0}\PYG{p}{,}\PYG{l+m+mi}{1}\PYG{p}{]}\PYG{o}{.}\PYG{n}{grid}\PYG{p}{(}\PYG{k+kc}{True}\PYG{p}{,} \PYG{n}{which}\PYG{o}{=}\PYG{l+s+s1}{\PYGZsq{}}\PYG{l+s+s1}{major}\PYG{l+s+s1}{\PYGZsq{}}\PYG{p}{,} \PYG{n}{linestyle}\PYG{o}{=}\PYG{l+s+s1}{\PYGZsq{}}\PYG{l+s+s1}{\PYGZhy{}\PYGZhy{}}\PYG{l+s+s1}{\PYGZsq{}}\PYG{p}{,} \PYG{n}{linewidth}\PYG{o}{=}\PYG{l+m+mf}{0.5}\PYG{p}{)}
    \PYG{n}{axs}\PYG{p}{[}\PYG{l+m+mi}{0}\PYG{p}{,}\PYG{l+m+mi}{1}\PYG{p}{]}\PYG{o}{.}\PYG{n}{legend}\PYG{p}{(}\PYG{p}{)}
    
    \PYG{c+c1}{\PYGZsh{} φxy}
    \PYG{n}{axs}\PYG{p}{[}\PYG{l+m+mi}{1}\PYG{p}{,}\PYG{l+m+mi}{0}\PYG{p}{]}\PYG{o}{.}\PYG{n}{semilogx}\PYG{p}{(}\PYG{n}{fre1}\PYG{p}{,} \PYG{n}{pxy1}\PYG{p}{,} \PYG{l+s+s1}{\PYGZsq{}}\PYG{l+s+s1}{r}\PYG{l+s+s1}{\PYGZsq{}}\PYG{p}{,} \PYG{n}{label}\PYG{o}{=}\PYG{n}{name1}\PYG{p}{)}
    \PYG{n}{axs}\PYG{p}{[}\PYG{l+m+mi}{1}\PYG{p}{,}\PYG{l+m+mi}{0}\PYG{p}{]}\PYG{o}{.}\PYG{n}{semilogx}\PYG{p}{(}\PYG{n}{fre2}\PYG{p}{,} \PYG{n}{pxy2}\PYG{p}{,} \PYG{l+s+s1}{\PYGZsq{}}\PYG{l+s+s1}{b}\PYG{l+s+s1}{\PYGZsq{}}\PYG{p}{,} \PYG{n}{label}\PYG{o}{=}\PYG{n}{name2}\PYG{p}{)}
    \PYG{n}{axs}\PYG{p}{[}\PYG{l+m+mi}{1}\PYG{p}{,}\PYG{l+m+mi}{0}\PYG{p}{]}\PYG{o}{.}\PYG{n}{set\PYGZus{}ylim}\PYG{p}{(}\PYG{o}{\PYGZhy{}}\PYG{l+m+mi}{180}\PYG{p}{,} \PYG{l+m+mi}{180}\PYG{p}{)}
    \PYG{n}{axs}\PYG{p}{[}\PYG{l+m+mi}{1}\PYG{p}{,}\PYG{l+m+mi}{0}\PYG{p}{]}\PYG{o}{.}\PYG{n}{set\PYGZus{}yticks}\PYG{p}{(}\PYG{n}{np}\PYG{o}{.}\PYG{n}{arange}\PYG{p}{(}\PYG{o}{\PYGZhy{}}\PYG{l+m+mi}{180}\PYG{p}{,} \PYG{l+m+mi}{181}\PYG{p}{,} \PYG{l+m+mi}{90}\PYG{p}{)}\PYG{p}{)}
    \PYG{n}{axs}\PYG{p}{[}\PYG{l+m+mi}{1}\PYG{p}{,}\PYG{l+m+mi}{0}\PYG{p}{]}\PYG{o}{.}\PYG{n}{set\PYGZus{}xlim}\PYG{p}{(}\PYG{n}{xmin}\PYG{p}{,} \PYG{n}{xmax}\PYG{p}{)}
    \PYG{n}{axs}\PYG{p}{[}\PYG{l+m+mi}{1}\PYG{p}{,}\PYG{l+m+mi}{0}\PYG{p}{]}\PYG{o}{.}\PYG{n}{invert\PYGZus{}xaxis}\PYG{p}{(}\PYG{p}{)}
    \PYG{n}{axs}\PYG{p}{[}\PYG{l+m+mi}{1}\PYG{p}{,}\PYG{l+m+mi}{0}\PYG{p}{]}\PYG{o}{.}\PYG{n}{set\PYGZus{}xlabel}\PYG{p}{(}\PYG{l+s+s1}{\PYGZsq{}}\PYG{l+s+s1}{Frequency (Hz)}\PYG{l+s+s1}{\PYGZsq{}}\PYG{p}{)}
    \PYG{n}{axs}\PYG{p}{[}\PYG{l+m+mi}{1}\PYG{p}{,}\PYG{l+m+mi}{0}\PYG{p}{]}\PYG{o}{.}\PYG{n}{set\PYGZus{}ylabel}\PYG{p}{(}\PYG{l+s+sa}{r}\PYG{l+s+s1}{\PYGZsq{}}\PYG{l+s+s1}{\PYGZdl{}}\PYG{l+s+s1}{\PYGZbs{}}\PYG{l+s+s1}{varphi}\PYG{l+s+s1}{\PYGZbs{}}\PYG{l+s+s1}{;(}\PYG{l+s+s1}{\PYGZbs{}}\PYG{l+s+s1}{degree)\PYGZdl{}}\PYG{l+s+s1}{\PYGZsq{}}\PYG{p}{)}
    \PYG{n}{axs}\PYG{p}{[}\PYG{l+m+mi}{1}\PYG{p}{,}\PYG{l+m+mi}{0}\PYG{p}{]}\PYG{o}{.}\PYG{n}{set\PYGZus{}title}\PYG{p}{(}\PYG{l+s+sa}{r}\PYG{l+s+s1}{\PYGZsq{}}\PYG{l+s+s1}{\PYGZdl{}}\PYG{l+s+s1}{\PYGZbs{}}\PYG{l+s+s1}{varphi\PYGZus{}}\PYG{l+s+si}{\PYGZob{}xy\PYGZcb{}}\PYG{l+s+s1}{\PYGZdl{} RMS = }\PYG{l+s+s1}{\PYGZsq{}} \PYG{o}{+} \PYG{l+s+sa}{f}\PYG{l+s+s1}{\PYGZsq{}}\PYG{l+s+si}{\PYGZob{}}\PYG{n}{pxyrms}\PYG{o}{*}\PYG{l+m+mi}{100}\PYG{l+s+si}{:}\PYG{l+s+s1}{.1f}\PYG{l+s+si}{\PYGZcb{}}\PYG{l+s+s1}{\PYGZpc{}}\PYG{l+s+s1}{\PYGZsq{}}\PYG{p}{)}
    \PYG{n}{axs}\PYG{p}{[}\PYG{l+m+mi}{1}\PYG{p}{,}\PYG{l+m+mi}{0}\PYG{p}{]}\PYG{o}{.}\PYG{n}{grid}\PYG{p}{(}\PYG{k+kc}{True}\PYG{p}{,} \PYG{n}{which}\PYG{o}{=}\PYG{l+s+s1}{\PYGZsq{}}\PYG{l+s+s1}{major}\PYG{l+s+s1}{\PYGZsq{}}\PYG{p}{,} \PYG{n}{linestyle}\PYG{o}{=}\PYG{l+s+s1}{\PYGZsq{}}\PYG{l+s+s1}{\PYGZhy{}\PYGZhy{}}\PYG{l+s+s1}{\PYGZsq{}}\PYG{p}{,} \PYG{n}{linewidth}\PYG{o}{=}\PYG{l+m+mf}{0.5}\PYG{p}{)}
    \PYG{n}{axs}\PYG{p}{[}\PYG{l+m+mi}{1}\PYG{p}{,}\PYG{l+m+mi}{0}\PYG{p}{]}\PYG{o}{.}\PYG{n}{legend}\PYG{p}{(}\PYG{p}{)}
    
    \PYG{c+c1}{\PYGZsh{} φyx}
    \PYG{n}{axs}\PYG{p}{[}\PYG{l+m+mi}{1}\PYG{p}{,}\PYG{l+m+mi}{1}\PYG{p}{]}\PYG{o}{.}\PYG{n}{semilogx}\PYG{p}{(}\PYG{n}{fre1}\PYG{p}{,} \PYG{n}{pyx1}\PYG{p}{,} \PYG{l+s+s1}{\PYGZsq{}}\PYG{l+s+s1}{r}\PYG{l+s+s1}{\PYGZsq{}}\PYG{p}{,} \PYG{n}{label}\PYG{o}{=}\PYG{n}{name1}\PYG{p}{)}
    \PYG{n}{axs}\PYG{p}{[}\PYG{l+m+mi}{1}\PYG{p}{,}\PYG{l+m+mi}{1}\PYG{p}{]}\PYG{o}{.}\PYG{n}{semilogx}\PYG{p}{(}\PYG{n}{fre2}\PYG{p}{,} \PYG{n}{pyx2}\PYG{p}{,} \PYG{l+s+s1}{\PYGZsq{}}\PYG{l+s+s1}{b}\PYG{l+s+s1}{\PYGZsq{}}\PYG{p}{,} \PYG{n}{label}\PYG{o}{=}\PYG{n}{name2}\PYG{p}{)}
    \PYG{n}{axs}\PYG{p}{[}\PYG{l+m+mi}{1}\PYG{p}{,}\PYG{l+m+mi}{1}\PYG{p}{]}\PYG{o}{.}\PYG{n}{set\PYGZus{}ylim}\PYG{p}{(}\PYG{o}{\PYGZhy{}}\PYG{l+m+mi}{180}\PYG{p}{,} \PYG{l+m+mi}{180}\PYG{p}{)}
    \PYG{n}{axs}\PYG{p}{[}\PYG{l+m+mi}{1}\PYG{p}{,}\PYG{l+m+mi}{1}\PYG{p}{]}\PYG{o}{.}\PYG{n}{set\PYGZus{}yticks}\PYG{p}{(}\PYG{n}{np}\PYG{o}{.}\PYG{n}{arange}\PYG{p}{(}\PYG{o}{\PYGZhy{}}\PYG{l+m+mi}{180}\PYG{p}{,} \PYG{l+m+mi}{181}\PYG{p}{,} \PYG{l+m+mi}{90}\PYG{p}{)}\PYG{p}{)}
    \PYG{n}{axs}\PYG{p}{[}\PYG{l+m+mi}{1}\PYG{p}{,}\PYG{l+m+mi}{1}\PYG{p}{]}\PYG{o}{.}\PYG{n}{set\PYGZus{}xlim}\PYG{p}{(}\PYG{n}{xmin}\PYG{p}{,} \PYG{n}{xmax}\PYG{p}{)}
    \PYG{n}{axs}\PYG{p}{[}\PYG{l+m+mi}{1}\PYG{p}{,}\PYG{l+m+mi}{1}\PYG{p}{]}\PYG{o}{.}\PYG{n}{invert\PYGZus{}xaxis}\PYG{p}{(}\PYG{p}{)}
    \PYG{n}{axs}\PYG{p}{[}\PYG{l+m+mi}{1}\PYG{p}{,}\PYG{l+m+mi}{1}\PYG{p}{]}\PYG{o}{.}\PYG{n}{set\PYGZus{}xlabel}\PYG{p}{(}\PYG{l+s+s1}{\PYGZsq{}}\PYG{l+s+s1}{Frequency (Hz)}\PYG{l+s+s1}{\PYGZsq{}}\PYG{p}{)}
    \PYG{n}{axs}\PYG{p}{[}\PYG{l+m+mi}{1}\PYG{p}{,}\PYG{l+m+mi}{1}\PYG{p}{]}\PYG{o}{.}\PYG{n}{set\PYGZus{}title}\PYG{p}{(}\PYG{l+s+sa}{r}\PYG{l+s+s1}{\PYGZsq{}}\PYG{l+s+s1}{\PYGZdl{}}\PYG{l+s+s1}{\PYGZbs{}}\PYG{l+s+s1}{varphi\PYGZus{}}\PYG{l+s+si}{\PYGZob{}yx\PYGZcb{}}\PYG{l+s+s1}{\PYGZdl{} RMS = }\PYG{l+s+s1}{\PYGZsq{}} \PYG{o}{+} \PYG{l+s+sa}{f}\PYG{l+s+s1}{\PYGZsq{}}\PYG{l+s+si}{\PYGZob{}}\PYG{n}{pyxrms}\PYG{o}{*}\PYG{l+m+mi}{100}\PYG{l+s+si}{:}\PYG{l+s+s1}{.1f}\PYG{l+s+si}{\PYGZcb{}}\PYG{l+s+s1}{\PYGZpc{}}\PYG{l+s+s1}{\PYGZsq{}}\PYG{p}{)}
    \PYG{n}{axs}\PYG{p}{[}\PYG{l+m+mi}{1}\PYG{p}{,}\PYG{l+m+mi}{1}\PYG{p}{]}\PYG{o}{.}\PYG{n}{grid}\PYG{p}{(}\PYG{k+kc}{True}\PYG{p}{,} \PYG{n}{which}\PYG{o}{=}\PYG{l+s+s1}{\PYGZsq{}}\PYG{l+s+s1}{major}\PYG{l+s+s1}{\PYGZsq{}}\PYG{p}{,} \PYG{n}{linestyle}\PYG{o}{=}\PYG{l+s+s1}{\PYGZsq{}}\PYG{l+s+s1}{\PYGZhy{}\PYGZhy{}}\PYG{l+s+s1}{\PYGZsq{}}\PYG{p}{,} \PYG{n}{linewidth}\PYG{o}{=}\PYG{l+m+mf}{0.5}\PYG{p}{)}
    \PYG{n}{axs}\PYG{p}{[}\PYG{l+m+mi}{1}\PYG{p}{,}\PYG{l+m+mi}{1}\PYG{p}{]}\PYG{o}{.}\PYG{n}{legend}\PYG{p}{(}\PYG{p}{)}
    
    \PYG{c+c1}{\PYGZsh{} 保存}
    \PYG{n}{fig}\PYG{o}{.}\PYG{n}{savefig}\PYG{p}{(}\PYG{n}{os}\PYG{o}{.}\PYG{n}{path}\PYG{o}{.}\PYG{n}{join}\PYG{p}{(}\PYG{n}{out\PYGZus{}dir}\PYG{p}{,} \PYG{l+s+sa}{f}\PYG{l+s+s1}{\PYGZsq{}}\PYG{l+s+si}{\PYGZob{}}\PYG{n}{name1}\PYG{l+s+si}{\PYGZcb{}}\PYG{l+s+s1}{\PYGZhy{}}\PYG{l+s+si}{\PYGZob{}}\PYG{n}{name2}\PYG{l+s+si}{\PYGZcb{}}\PYG{l+s+s1}{.png}\PYG{l+s+s1}{\PYGZsq{}}\PYG{p}{)}\PYG{p}{,} \PYG{n}{dpi}\PYG{o}{=}\PYG{l+m+mi}{300}\PYG{p}{,} \PYG{n}{bbox\PYGZus{}inches}\PYG{o}{=}\PYG{l+s+s1}{\PYGZsq{}}\PYG{l+s+s1}{tight}\PYG{l+s+s1}{\PYGZsq{}}\PYG{p}{)}
    \PYG{n}{fig}\PYG{o}{.}\PYG{n}{savefig}\PYG{p}{(}\PYG{n}{os}\PYG{o}{.}\PYG{n}{path}\PYG{o}{.}\PYG{n}{join}\PYG{p}{(}\PYG{n}{out\PYGZus{}dir}\PYG{p}{,} \PYG{l+s+sa}{f}\PYG{l+s+s1}{\PYGZsq{}}\PYG{l+s+si}{\PYGZob{}}\PYG{n}{name1}\PYG{l+s+si}{\PYGZcb{}}\PYG{l+s+s1}{\PYGZhy{}}\PYG{l+s+si}{\PYGZob{}}\PYG{n}{name2}\PYG{l+s+si}{\PYGZcb{}}\PYG{l+s+s1}{.pdf}\PYG{l+s+s1}{\PYGZsq{}}\PYG{p}{)}\PYG{p}{,} \PYG{n}{dpi}\PYG{o}{=}\PYG{l+m+mi}{300}\PYG{p}{,} \PYG{n}{bbox\PYGZus{}inches}\PYG{o}{=}\PYG{l+s+s1}{\PYGZsq{}}\PYG{l+s+s1}{tight}\PYG{l+s+s1}{\PYGZsq{}}\PYG{p}{)}
    \PYG{n}{fig}\PYG{o}{.}\PYG{n}{savefig}\PYG{p}{(}\PYG{n}{os}\PYG{o}{.}\PYG{n}{path}\PYG{o}{.}\PYG{n}{join}\PYG{p}{(}\PYG{n}{out\PYGZus{}dir}\PYG{p}{,} \PYG{l+s+sa}{f}\PYG{l+s+s1}{\PYGZsq{}}\PYG{l+s+si}{\PYGZob{}}\PYG{n}{name1}\PYG{l+s+si}{\PYGZcb{}}\PYG{l+s+s1}{\PYGZhy{}}\PYG{l+s+si}{\PYGZob{}}\PYG{n}{name2}\PYG{l+s+si}{\PYGZcb{}}\PYG{l+s+s1}{.svg}\PYG{l+s+s1}{\PYGZsq{}}\PYG{p}{)}\PYG{p}{,} \PYG{n}{dpi}\PYG{o}{=}\PYG{l+m+mi}{300}\PYG{p}{,} \PYG{n}{bbox\PYGZus{}inches}\PYG{o}{=}\PYG{l+s+s1}{\PYGZsq{}}\PYG{l+s+s1}{tight}\PYG{l+s+s1}{\PYGZsq{}}\PYG{p}{)}
    \PYG{n}{plt}\PYG{o}{.}\PYG{n}{close}\PYG{p}{(}\PYG{p}{)}


\PYG{c+c1}{\PYGZsh{} ===== 使用示例 =====}
\PYG{n}{path} \PYG{o}{=} \PYG{l+s+sa}{r}\PYG{l+s+s1}{\PYGZsq{}}\PYG{l+s+s1}{D:}\PYG{l+s+s1}{\PYGZbs{}}\PYG{l+s+s1}{DataProcess}\PYG{l+s+s1}{\PYGZbs{}}\PYG{l+s+s1}{JianghanPlain}\PYG{l+s+s1}{\PYGZbs{}}\PYG{l+s+s1}{MTDPProject}\PYG{l+s+s1}{\PYGZbs{}}\PYG{l+s+s1}{Consistency}\PYG{l+s+s1}{\PYGZsq{}}
\PYG{n}{comparison\PYGZus{}pairs} \PYG{o}{=} \PYG{p}{[}
    \PYG{p}{[}\PYG{l+s+s1}{\PYGZsq{}}\PYG{l+s+s1}{JHMT41SA}\PYG{l+s+s1}{\PYGZsq{}}\PYG{p}{,} \PYG{l+s+s1}{\PYGZsq{}}\PYG{l+s+s1}{JHMT41SB}\PYG{l+s+s1}{\PYGZsq{}}\PYG{p}{]}\PYG{p}{,}
    \PYG{p}{[}\PYG{l+s+s1}{\PYGZsq{}}\PYG{l+s+s1}{JHMT45SA}\PYG{l+s+s1}{\PYGZsq{}}\PYG{p}{,} \PYG{l+s+s1}{\PYGZsq{}}\PYG{l+s+s1}{JHMT45SB}\PYG{l+s+s1}{\PYGZsq{}}\PYG{p}{]}\PYG{p}{,}
    \PYG{p}{[}\PYG{l+s+s1}{\PYGZsq{}}\PYG{l+s+s1}{JHMT46SA}\PYG{l+s+s1}{\PYGZsq{}}\PYG{p}{,} \PYG{l+s+s1}{\PYGZsq{}}\PYG{l+s+s1}{JHMT46SJ}\PYG{l+s+s1}{\PYGZsq{}}\PYG{p}{]}\PYG{p}{,}
\PYG{p}{]}

\PYG{k}{for} \PYG{n}{name1}\PYG{p}{,} \PYG{n}{name2} \PYG{o+ow}{in} \PYG{n}{comparison\PYGZus{}pairs}\PYG{p}{:}
    \PYG{n+nb}{print}\PYG{p}{(}\PYG{l+s+sa}{f}\PYG{l+s+s1}{\PYGZsq{}}\PYG{l+s+s1}{Processing }\PYG{l+s+si}{\PYGZob{}}\PYG{n}{name1}\PYG{l+s+si}{\PYGZcb{}}\PYG{l+s+s1}{ vs }\PYG{l+s+si}{\PYGZob{}}\PYG{n}{name2}\PYG{l+s+si}{\PYGZcb{}}\PYG{l+s+s1}{\PYGZsq{}}\PYG{p}{)}
    \PYG{n}{plotMTDP}\PYG{p}{(}\PYG{n}{path}\PYG{p}{,} \PYG{n}{name1}\PYG{p}{,} \PYG{n}{name2}\PYG{p}{)}

\PYG{n+nb}{print}\PYG{p}{(}\PYG{l+s+s1}{\PYGZsq{}}\PYG{l+s+s1}{Done!}\PYG{l+s+s1}{\PYGZsq{}}\PYG{p}{)}
\end{sphinxVerbatim}


\bigskip\hrule\bigskip



\subsection{全阻抗与相位张量图}
\label{\detokenize{tutorial/plotting:id9}}
\sphinxAtStartPar
绘制包含完整阻抗张量、倾子和相位张量的综合图件。


\subsubsection{脚本：\sphinxhref{http://PlotFullZ-Site.py}{PlotFullZ\sphinxhyphen{}Site.py} (http://PlotFullZ\sphinxhyphen{}Site.py)}
\label{\detokenize{tutorial/plotting:plotfullz-site-py}}
\begin{sphinxVerbatim}[commandchars=\\\{\}]
\PYG{k+kn}{import}\PYG{+w}{ }\PYG{n+nn}{os}
\PYG{k+kn}{import}\PYG{+w}{ }\PYG{n+nn}{numpy}\PYG{+w}{ }\PYG{k}{as}\PYG{+w}{ }\PYG{n+nn}{np}
\PYG{k+kn}{import}\PYG{+w}{ }\PYG{n+nn}{matplotlib}\PYG{n+nn}{.}\PYG{n+nn}{pyplot}\PYG{+w}{ }\PYG{k}{as}\PYG{+w}{ }\PYG{n+nn}{plt}
\PYG{k+kn}{from}\PYG{+w}{ }\PYG{n+nn}{matplotlib}\PYG{n+nn}{.}\PYG{n+nn}{ticker}\PYG{+w}{ }\PYG{k+kn}{import} \PYG{n}{ScalarFormatter}
\PYG{k+kn}{from}\PYG{+w}{ }\PYG{n+nn}{matplotlib}\PYG{n+nn}{.}\PYG{n+nn}{collections}\PYG{+w}{ }\PYG{k+kn}{import} \PYG{n}{EllipseCollection}
\PYG{k+kn}{from}\PYG{+w}{ }\PYG{n+nn}{matplotlib}\PYG{n+nn}{.}\PYG{n+nn}{colors}\PYG{+w}{ }\PYG{k+kn}{import} \PYG{n}{Normalize}
\PYG{k+kn}{from}\PYG{+w}{ }\PYG{n+nn}{matplotlib}\PYG{n+nn}{.}\PYG{n+nn}{cm}\PYG{+w}{ }\PYG{k+kn}{import} \PYG{n}{ScalarMappable}
\PYG{k+kn}{from}\PYG{+w}{ }\PYG{n+nn}{mpl\PYGZus{}toolkits}\PYG{n+nn}{.}\PYG{n+nn}{axes\PYGZus{}grid1}\PYG{n+nn}{.}\PYG{n+nn}{inset\PYGZus{}locator}\PYG{+w}{ }\PYG{k+kn}{import} \PYG{n}{inset\PYGZus{}axes}

\PYG{c+c1}{\PYGZsh{} ===== 配置参数 =====}
\PYG{n}{path} \PYG{o}{=} \PYG{l+s+sa}{r}\PYG{l+s+s1}{\PYGZsq{}}\PYG{l+s+s1}{D:}\PYG{l+s+s1}{\PYGZbs{}}\PYG{l+s+s1}{Thesis}\PYG{l+s+s1}{\PYGZbs{}}\PYG{l+s+s1}{Chapter5}\PYG{l+s+s1}{\PYGZbs{}}\PYG{l+s+s1}{MTDPProject}\PYG{l+s+s1}{\PYGZbs{}}\PYG{l+s+s1}{LModel}\PYG{l+s+s1}{\PYGZbs{}}\PYG{l+s+s1}{A}\PYG{l+s+s1}{\PYGZsq{}}
\PYG{n}{name} \PYG{o}{=} \PYG{l+s+s1}{\PYGZsq{}}\PYG{l+s+s1}{A}\PYG{l+s+s1}{\PYGZsq{}}

\PYG{c+c1}{\PYGZsh{} ===== 读取数据 =====}
\PYG{n}{fre} \PYG{o}{=} \PYG{n}{np}\PYG{o}{.}\PYG{n}{loadtxt}\PYG{p}{(}\PYG{l+s+sa}{f}\PYG{l+s+s1}{\PYGZsq{}}\PYG{l+s+si}{\PYGZob{}}\PYG{n}{path}\PYG{l+s+si}{\PYGZcb{}}\PYG{l+s+s1}{/}\PYG{l+s+si}{\PYGZob{}}\PYG{n}{name}\PYG{l+s+si}{\PYGZcb{}}\PYG{l+s+s1}{\PYGZhy{}rxx.dat}\PYG{l+s+s1}{\PYGZsq{}}\PYG{p}{)}\PYG{p}{[}\PYG{p}{:}\PYG{p}{,} \PYG{l+m+mi}{0}\PYG{p}{]}
\PYG{n}{rxx} \PYG{o}{=} \PYG{n}{np}\PYG{o}{.}\PYG{n}{loadtxt}\PYG{p}{(}\PYG{l+s+sa}{f}\PYG{l+s+s1}{\PYGZsq{}}\PYG{l+s+si}{\PYGZob{}}\PYG{n}{path}\PYG{l+s+si}{\PYGZcb{}}\PYG{l+s+s1}{/}\PYG{l+s+si}{\PYGZob{}}\PYG{n}{name}\PYG{l+s+si}{\PYGZcb{}}\PYG{l+s+s1}{\PYGZhy{}rxx.dat}\PYG{l+s+s1}{\PYGZsq{}}\PYG{p}{)}\PYG{p}{[}\PYG{p}{:}\PYG{p}{,} \PYG{l+m+mi}{1}\PYG{p}{]}
\PYG{n}{rxy} \PYG{o}{=} \PYG{n}{np}\PYG{o}{.}\PYG{n}{loadtxt}\PYG{p}{(}\PYG{l+s+sa}{f}\PYG{l+s+s1}{\PYGZsq{}}\PYG{l+s+si}{\PYGZob{}}\PYG{n}{path}\PYG{l+s+si}{\PYGZcb{}}\PYG{l+s+s1}{/}\PYG{l+s+si}{\PYGZob{}}\PYG{n}{name}\PYG{l+s+si}{\PYGZcb{}}\PYG{l+s+s1}{\PYGZhy{}rxy.dat}\PYG{l+s+s1}{\PYGZsq{}}\PYG{p}{)}\PYG{p}{[}\PYG{p}{:}\PYG{p}{,} \PYG{l+m+mi}{1}\PYG{p}{]}
\PYG{n}{ryx} \PYG{o}{=} \PYG{n}{np}\PYG{o}{.}\PYG{n}{loadtxt}\PYG{p}{(}\PYG{l+s+sa}{f}\PYG{l+s+s1}{\PYGZsq{}}\PYG{l+s+si}{\PYGZob{}}\PYG{n}{path}\PYG{l+s+si}{\PYGZcb{}}\PYG{l+s+s1}{/}\PYG{l+s+si}{\PYGZob{}}\PYG{n}{name}\PYG{l+s+si}{\PYGZcb{}}\PYG{l+s+s1}{\PYGZhy{}ryx.dat}\PYG{l+s+s1}{\PYGZsq{}}\PYG{p}{)}\PYG{p}{[}\PYG{p}{:}\PYG{p}{,} \PYG{l+m+mi}{1}\PYG{p}{]}
\PYG{n}{ryy} \PYG{o}{=} \PYG{n}{np}\PYG{o}{.}\PYG{n}{loadtxt}\PYG{p}{(}\PYG{l+s+sa}{f}\PYG{l+s+s1}{\PYGZsq{}}\PYG{l+s+si}{\PYGZob{}}\PYG{n}{path}\PYG{l+s+si}{\PYGZcb{}}\PYG{l+s+s1}{/}\PYG{l+s+si}{\PYGZob{}}\PYG{n}{name}\PYG{l+s+si}{\PYGZcb{}}\PYG{l+s+s1}{\PYGZhy{}rxy.dat}\PYG{l+s+s1}{\PYGZsq{}}\PYG{p}{)}\PYG{p}{[}\PYG{p}{:}\PYG{p}{,} \PYG{l+m+mi}{1}\PYG{p}{]}

\PYG{n}{pxx} \PYG{o}{=} \PYG{n}{np}\PYG{o}{.}\PYG{n}{loadtxt}\PYG{p}{(}\PYG{l+s+sa}{f}\PYG{l+s+s1}{\PYGZsq{}}\PYG{l+s+si}{\PYGZob{}}\PYG{n}{path}\PYG{l+s+si}{\PYGZcb{}}\PYG{l+s+s1}{/}\PYG{l+s+si}{\PYGZob{}}\PYG{n}{name}\PYG{l+s+si}{\PYGZcb{}}\PYG{l+s+s1}{\PYGZhy{}zxxp.dat}\PYG{l+s+s1}{\PYGZsq{}}\PYG{p}{)}\PYG{p}{[}\PYG{p}{:}\PYG{p}{,} \PYG{l+m+mi}{1}\PYG{p}{]}
\PYG{n}{pxy} \PYG{o}{=} \PYG{n}{np}\PYG{o}{.}\PYG{n}{loadtxt}\PYG{p}{(}\PYG{l+s+sa}{f}\PYG{l+s+s1}{\PYGZsq{}}\PYG{l+s+si}{\PYGZob{}}\PYG{n}{path}\PYG{l+s+si}{\PYGZcb{}}\PYG{l+s+s1}{/}\PYG{l+s+si}{\PYGZob{}}\PYG{n}{name}\PYG{l+s+si}{\PYGZcb{}}\PYG{l+s+s1}{\PYGZhy{}zxyp.dat}\PYG{l+s+s1}{\PYGZsq{}}\PYG{p}{)}\PYG{p}{[}\PYG{p}{:}\PYG{p}{,} \PYG{l+m+mi}{1}\PYG{p}{]}
\PYG{n}{pyx} \PYG{o}{=} \PYG{n}{np}\PYG{o}{.}\PYG{n}{loadtxt}\PYG{p}{(}\PYG{l+s+sa}{f}\PYG{l+s+s1}{\PYGZsq{}}\PYG{l+s+si}{\PYGZob{}}\PYG{n}{path}\PYG{l+s+si}{\PYGZcb{}}\PYG{l+s+s1}{/}\PYG{l+s+si}{\PYGZob{}}\PYG{n}{name}\PYG{l+s+si}{\PYGZcb{}}\PYG{l+s+s1}{\PYGZhy{}zyxp.dat}\PYG{l+s+s1}{\PYGZsq{}}\PYG{p}{)}\PYG{p}{[}\PYG{p}{:}\PYG{p}{,} \PYG{l+m+mi}{1}\PYG{p}{]}
\PYG{n}{pyy} \PYG{o}{=} \PYG{n}{np}\PYG{o}{.}\PYG{n}{loadtxt}\PYG{p}{(}\PYG{l+s+sa}{f}\PYG{l+s+s1}{\PYGZsq{}}\PYG{l+s+si}{\PYGZob{}}\PYG{n}{path}\PYG{l+s+si}{\PYGZcb{}}\PYG{l+s+s1}{/}\PYG{l+s+si}{\PYGZob{}}\PYG{n}{name}\PYG{l+s+si}{\PYGZcb{}}\PYG{l+s+s1}{\PYGZhy{}zyyp.dat}\PYG{l+s+s1}{\PYGZsq{}}\PYG{p}{)}\PYG{p}{[}\PYG{p}{:}\PYG{p}{,} \PYG{l+m+mi}{1}\PYG{p}{]}

\PYG{c+c1}{\PYGZsh{} 倾子数据}
\PYG{n}{tzxr} \PYG{o}{=} \PYG{n}{np}\PYG{o}{.}\PYG{n}{loadtxt}\PYG{p}{(}\PYG{l+s+sa}{f}\PYG{l+s+s1}{\PYGZsq{}}\PYG{l+s+si}{\PYGZob{}}\PYG{n}{path}\PYG{l+s+si}{\PYGZcb{}}\PYG{l+s+s1}{/}\PYG{l+s+si}{\PYGZob{}}\PYG{n}{name}\PYG{l+s+si}{\PYGZcb{}}\PYG{l+s+s1}{\PYGZhy{}tzxr.dat}\PYG{l+s+s1}{\PYGZsq{}}\PYG{p}{)}\PYG{p}{[}\PYG{p}{:}\PYG{p}{,} \PYG{l+m+mi}{1}\PYG{p}{]}
\PYG{n}{tzxi} \PYG{o}{=} \PYG{n}{np}\PYG{o}{.}\PYG{n}{loadtxt}\PYG{p}{(}\PYG{l+s+sa}{f}\PYG{l+s+s1}{\PYGZsq{}}\PYG{l+s+si}{\PYGZob{}}\PYG{n}{path}\PYG{l+s+si}{\PYGZcb{}}\PYG{l+s+s1}{/}\PYG{l+s+si}{\PYGZob{}}\PYG{n}{name}\PYG{l+s+si}{\PYGZcb{}}\PYG{l+s+s1}{\PYGZhy{}tzxi.dat}\PYG{l+s+s1}{\PYGZsq{}}\PYG{p}{)}\PYG{p}{[}\PYG{p}{:}\PYG{p}{,} \PYG{l+m+mi}{1}\PYG{p}{]}
\PYG{n}{tzyr} \PYG{o}{=} \PYG{n}{np}\PYG{o}{.}\PYG{n}{loadtxt}\PYG{p}{(}\PYG{l+s+sa}{f}\PYG{l+s+s1}{\PYGZsq{}}\PYG{l+s+si}{\PYGZob{}}\PYG{n}{path}\PYG{l+s+si}{\PYGZcb{}}\PYG{l+s+s1}{/}\PYG{l+s+si}{\PYGZob{}}\PYG{n}{name}\PYG{l+s+si}{\PYGZcb{}}\PYG{l+s+s1}{\PYGZhy{}tzyr.dat}\PYG{l+s+s1}{\PYGZsq{}}\PYG{p}{)}\PYG{p}{[}\PYG{p}{:}\PYG{p}{,} \PYG{l+m+mi}{1}\PYG{p}{]}
\PYG{n}{tzyi} \PYG{o}{=} \PYG{n}{np}\PYG{o}{.}\PYG{n}{loadtxt}\PYG{p}{(}\PYG{l+s+sa}{f}\PYG{l+s+s1}{\PYGZsq{}}\PYG{l+s+si}{\PYGZob{}}\PYG{n}{path}\PYG{l+s+si}{\PYGZcb{}}\PYG{l+s+s1}{/}\PYG{l+s+si}{\PYGZob{}}\PYG{n}{name}\PYG{l+s+si}{\PYGZcb{}}\PYG{l+s+s1}{\PYGZhy{}tzyi.dat}\PYG{l+s+s1}{\PYGZsq{}}\PYG{p}{)}\PYG{p}{[}\PYG{p}{:}\PYG{p}{,} \PYG{l+m+mi}{1}\PYG{p}{]}

\PYG{c+c1}{\PYGZsh{} ===== 配置坐标范围 =====}
\PYG{n}{yfmin}\PYG{p}{,} \PYG{n}{yfmax} \PYG{o}{=} \PYG{l+m+mi}{2}\PYG{p}{,} \PYG{l+m+mi}{35}
\PYG{n}{y\PYGZus{}min} \PYG{o}{=} \PYG{l+m+mf}{1e\PYGZhy{}1}
\PYG{n}{y\PYGZus{}max} \PYG{o}{=} \PYG{l+m+mf}{1e5}
\PYG{n}{x\PYGZus{}min} \PYG{o}{=} \PYG{n}{np}\PYG{o}{.}\PYG{n}{min}\PYG{p}{(}\PYG{n}{fre}\PYG{p}{)}
\PYG{n}{x\PYGZus{}max} \PYG{o}{=} \PYG{n}{np}\PYG{o}{.}\PYG{n}{max}\PYG{p}{(}\PYG{n}{fre}\PYG{p}{)}
\PYG{n}{xdelta} \PYG{o}{=} \PYG{p}{(}\PYG{n}{np}\PYG{o}{.}\PYG{n}{log}\PYG{p}{(}\PYG{n}{x\PYGZus{}max}\PYG{p}{)} \PYG{o}{\PYGZhy{}} \PYG{n}{np}\PYG{o}{.}\PYG{n}{log}\PYG{p}{(}\PYG{n}{x\PYGZus{}min}\PYG{p}{)}\PYG{p}{)} \PYG{o}{/} \PYG{l+m+mi}{10}
\PYG{n}{x\PYGZus{}min} \PYG{o}{=} \PYG{n}{np}\PYG{o}{.}\PYG{n}{exp}\PYG{p}{(}\PYG{n}{np}\PYG{o}{.}\PYG{n}{log}\PYG{p}{(}\PYG{n}{x\PYGZus{}min}\PYG{p}{)} \PYG{o}{\PYGZhy{}} \PYG{n}{xdelta}\PYG{p}{)}
\PYG{n}{x\PYGZus{}max} \PYG{o}{=} \PYG{n}{np}\PYG{o}{.}\PYG{n}{exp}\PYG{p}{(}\PYG{n}{np}\PYG{o}{.}\PYG{n}{log}\PYG{p}{(}\PYG{n}{x\PYGZus{}max}\PYG{p}{)} \PYG{o}{+} \PYG{n}{xdelta}\PYG{p}{)}

\PYG{c+c1}{\PYGZsh{} ===== 绘图 =====}
\PYG{n}{fig}\PYG{p}{,} \PYG{n}{axs} \PYG{o}{=} \PYG{n}{plt}\PYG{o}{.}\PYG{n}{subplots}\PYG{p}{(}\PYG{l+m+mi}{4}\PYG{p}{,} \PYG{l+m+mi}{1}\PYG{p}{,} \PYG{n}{figsize}\PYG{o}{=}\PYG{p}{(}\PYG{l+m+mi}{8}\PYG{p}{,} \PYG{l+m+mi}{10}\PYG{p}{)}\PYG{p}{,} \PYG{n}{sharex}\PYG{o}{=}\PYG{l+s+s1}{\PYGZsq{}}\PYG{l+s+s1}{all}\PYG{l+s+s1}{\PYGZsq{}}\PYG{p}{,}
                        \PYG{n}{gridspec\PYGZus{}kw}\PYG{o}{=}\PYG{p}{\PYGZob{}}\PYG{l+s+s1}{\PYGZsq{}}\PYG{l+s+s1}{height\PYGZus{}ratios}\PYG{l+s+s1}{\PYGZsq{}}\PYG{p}{:} \PYG{p}{[}\PYG{l+m+mi}{4}\PYG{p}{,} \PYG{l+m+mi}{4}\PYG{p}{,} \PYG{l+m+mi}{2}\PYG{p}{,} \PYG{l+m+mi}{2}\PYG{p}{]}\PYG{p}{\PYGZcb{}}\PYG{p}{,}
                        \PYG{n}{constrained\PYGZus{}layout}\PYG{o}{=}\PYG{k+kc}{True}\PYG{p}{)}

\PYG{c+c1}{\PYGZsh{} 1. 视电阻率}
\PYG{n}{ax} \PYG{o}{=} \PYG{n}{axs}\PYG{p}{[}\PYG{l+m+mi}{0}\PYG{p}{]}
\PYG{n}{ax}\PYG{o}{.}\PYG{n}{loglog}\PYG{p}{(}\PYG{n}{fre}\PYG{p}{,} \PYG{n}{rxx}\PYG{p}{,} \PYG{l+s+s1}{\PYGZsq{}}\PYG{l+s+s1}{mx}\PYG{l+s+s1}{\PYGZsq{}}\PYG{p}{,} \PYG{n}{label}\PYG{o}{=}\PYG{l+s+s1}{\PYGZsq{}}\PYG{l+s+s1}{XX}\PYG{l+s+s1}{\PYGZsq{}}\PYG{p}{,} \PYG{n}{markerfacecolor}\PYG{o}{=}\PYG{l+s+s1}{\PYGZsq{}}\PYG{l+s+s1}{none}\PYG{l+s+s1}{\PYGZsq{}}\PYG{p}{)}
\PYG{n}{ax}\PYG{o}{.}\PYG{n}{loglog}\PYG{p}{(}\PYG{n}{fre}\PYG{p}{,} \PYG{n}{rxy}\PYG{p}{,} \PYG{l+s+s1}{\PYGZsq{}}\PYG{l+s+s1}{ro}\PYG{l+s+s1}{\PYGZsq{}}\PYG{p}{,} \PYG{n}{label}\PYG{o}{=}\PYG{l+s+s1}{\PYGZsq{}}\PYG{l+s+s1}{XY}\PYG{l+s+s1}{\PYGZsq{}}\PYG{p}{,} \PYG{n}{markerfacecolor}\PYG{o}{=}\PYG{l+s+s1}{\PYGZsq{}}\PYG{l+s+s1}{none}\PYG{l+s+s1}{\PYGZsq{}}\PYG{p}{)}
\PYG{n}{ax}\PYG{o}{.}\PYG{n}{loglog}\PYG{p}{(}\PYG{n}{fre}\PYG{p}{,} \PYG{n}{ryx}\PYG{p}{,} \PYG{l+s+s1}{\PYGZsq{}}\PYG{l+s+s1}{bs}\PYG{l+s+s1}{\PYGZsq{}}\PYG{p}{,} \PYG{n}{label}\PYG{o}{=}\PYG{l+s+s1}{\PYGZsq{}}\PYG{l+s+s1}{YX}\PYG{l+s+s1}{\PYGZsq{}}\PYG{p}{,} \PYG{n}{markerfacecolor}\PYG{o}{=}\PYG{l+s+s1}{\PYGZsq{}}\PYG{l+s+s1}{none}\PYG{l+s+s1}{\PYGZsq{}}\PYG{p}{)}
\PYG{n}{ax}\PYG{o}{.}\PYG{n}{loglog}\PYG{p}{(}\PYG{n}{fre}\PYG{p}{,} \PYG{n}{ryy}\PYG{p}{,} \PYG{l+s+s1}{\PYGZsq{}}\PYG{l+s+s1}{c+}\PYG{l+s+s1}{\PYGZsq{}}\PYG{p}{,} \PYG{n}{label}\PYG{o}{=}\PYG{l+s+s1}{\PYGZsq{}}\PYG{l+s+s1}{YY}\PYG{l+s+s1}{\PYGZsq{}}\PYG{p}{,} \PYG{n}{markerfacecolor}\PYG{o}{=}\PYG{l+s+s1}{\PYGZsq{}}\PYG{l+s+s1}{none}\PYG{l+s+s1}{\PYGZsq{}}\PYG{p}{)}
\PYG{n}{ax}\PYG{o}{.}\PYG{n}{set\PYGZus{}ylim}\PYG{p}{(}\PYG{n}{y\PYGZus{}min}\PYG{p}{,} \PYG{n}{y\PYGZus{}max}\PYG{p}{)}
\PYG{n}{ax}\PYG{o}{.}\PYG{n}{set\PYGZus{}yscale}\PYG{p}{(}\PYG{l+s+s1}{\PYGZsq{}}\PYG{l+s+s1}{log}\PYG{l+s+s1}{\PYGZsq{}}\PYG{p}{)}
\PYG{n}{ax}\PYG{o}{.}\PYG{n}{set\PYGZus{}xlim}\PYG{p}{(}\PYG{n}{x\PYGZus{}min}\PYG{p}{,} \PYG{n}{x\PYGZus{}max}\PYG{p}{)}
\PYG{n}{ax}\PYG{o}{.}\PYG{n}{invert\PYGZus{}xaxis}\PYG{p}{(}\PYG{p}{)}
\PYG{n}{ax}\PYG{o}{.}\PYG{n}{set\PYGZus{}ylabel}\PYG{p}{(}\PYG{l+s+sa}{r}\PYG{l+s+s1}{\PYGZsq{}}\PYG{l+s+s1}{\PYGZdl{}}\PYG{l+s+s1}{\PYGZbs{}}\PYG{l+s+s1}{rho}\PYG{l+s+s1}{\PYGZbs{}}\PYG{l+s+s1}{;(}\PYG{l+s+s1}{\PYGZbs{}}\PYG{l+s+s1}{Omega }\PYG{l+s+s1}{\PYGZbs{}}\PYG{l+s+s1}{cdot m)\PYGZdl{}}\PYG{l+s+s1}{\PYGZsq{}}\PYG{p}{)}
\PYG{n}{ax}\PYG{o}{.}\PYG{n}{grid}\PYG{p}{(}\PYG{k+kc}{True}\PYG{p}{,} \PYG{n}{which}\PYG{o}{=}\PYG{l+s+s1}{\PYGZsq{}}\PYG{l+s+s1}{major}\PYG{l+s+s1}{\PYGZsq{}}\PYG{p}{,} \PYG{n}{linestyle}\PYG{o}{=}\PYG{l+s+s1}{\PYGZsq{}}\PYG{l+s+s1}{\PYGZhy{}\PYGZhy{}}\PYG{l+s+s1}{\PYGZsq{}}\PYG{p}{,} \PYG{n}{linewidth}\PYG{o}{=}\PYG{l+m+mf}{0.5}\PYG{p}{)}
\PYG{n}{ax}\PYG{o}{.}\PYG{n}{legend}\PYG{p}{(}\PYG{n}{loc}\PYG{o}{=}\PYG{l+s+s1}{\PYGZsq{}}\PYG{l+s+s1}{upper right}\PYG{l+s+s1}{\PYGZsq{}}\PYG{p}{,} \PYG{n}{bbox\PYGZus{}to\PYGZus{}anchor}\PYG{o}{=}\PYG{p}{(}\PYG{l+m+mf}{1.15}\PYG{p}{,} \PYG{l+m+mf}{1.0}\PYG{p}{)}\PYG{p}{)}

\PYG{c+c1}{\PYGZsh{} 2. 相位}
\PYG{n}{ax} \PYG{o}{=} \PYG{n}{axs}\PYG{p}{[}\PYG{l+m+mi}{1}\PYG{p}{]}
\PYG{n}{ax}\PYG{o}{.}\PYG{n}{semilogx}\PYG{p}{(}\PYG{n}{fre}\PYG{p}{,} \PYG{n}{pxx}\PYG{p}{,} \PYG{l+s+s1}{\PYGZsq{}}\PYG{l+s+s1}{mx}\PYG{l+s+s1}{\PYGZsq{}}\PYG{p}{,} \PYG{n}{label}\PYG{o}{=}\PYG{l+s+s1}{\PYGZsq{}}\PYG{l+s+s1}{XX}\PYG{l+s+s1}{\PYGZsq{}}\PYG{p}{,} \PYG{n}{markerfacecolor}\PYG{o}{=}\PYG{l+s+s1}{\PYGZsq{}}\PYG{l+s+s1}{none}\PYG{l+s+s1}{\PYGZsq{}}\PYG{p}{)}
\PYG{n}{ax}\PYG{o}{.}\PYG{n}{semilogx}\PYG{p}{(}\PYG{n}{fre}\PYG{p}{,} \PYG{n}{pxy}\PYG{p}{,} \PYG{l+s+s1}{\PYGZsq{}}\PYG{l+s+s1}{ro}\PYG{l+s+s1}{\PYGZsq{}}\PYG{p}{,} \PYG{n}{label}\PYG{o}{=}\PYG{l+s+s1}{\PYGZsq{}}\PYG{l+s+s1}{XY}\PYG{l+s+s1}{\PYGZsq{}}\PYG{p}{,} \PYG{n}{markerfacecolor}\PYG{o}{=}\PYG{l+s+s1}{\PYGZsq{}}\PYG{l+s+s1}{none}\PYG{l+s+s1}{\PYGZsq{}}\PYG{p}{)}
\PYG{n}{ax}\PYG{o}{.}\PYG{n}{semilogx}\PYG{p}{(}\PYG{n}{fre}\PYG{p}{,} \PYG{n}{pyx}\PYG{p}{,} \PYG{l+s+s1}{\PYGZsq{}}\PYG{l+s+s1}{bs}\PYG{l+s+s1}{\PYGZsq{}}\PYG{p}{,} \PYG{n}{label}\PYG{o}{=}\PYG{l+s+s1}{\PYGZsq{}}\PYG{l+s+s1}{YX}\PYG{l+s+s1}{\PYGZsq{}}\PYG{p}{,} \PYG{n}{markerfacecolor}\PYG{o}{=}\PYG{l+s+s1}{\PYGZsq{}}\PYG{l+s+s1}{none}\PYG{l+s+s1}{\PYGZsq{}}\PYG{p}{)}
\PYG{n}{ax}\PYG{o}{.}\PYG{n}{semilogx}\PYG{p}{(}\PYG{n}{fre}\PYG{p}{,} \PYG{n}{pyy}\PYG{p}{,} \PYG{l+s+s1}{\PYGZsq{}}\PYG{l+s+s1}{c+}\PYG{l+s+s1}{\PYGZsq{}}\PYG{p}{,} \PYG{n}{label}\PYG{o}{=}\PYG{l+s+s1}{\PYGZsq{}}\PYG{l+s+s1}{YY}\PYG{l+s+s1}{\PYGZsq{}}\PYG{p}{,} \PYG{n}{markerfacecolor}\PYG{o}{=}\PYG{l+s+s1}{\PYGZsq{}}\PYG{l+s+s1}{none}\PYG{l+s+s1}{\PYGZsq{}}\PYG{p}{)}
\PYG{n}{ax}\PYG{o}{.}\PYG{n}{set\PYGZus{}ylim}\PYG{p}{(}\PYG{o}{\PYGZhy{}}\PYG{l+m+mi}{180}\PYG{p}{,} \PYG{l+m+mi}{180}\PYG{p}{)}
\PYG{n}{ax}\PYG{o}{.}\PYG{n}{set\PYGZus{}yticks}\PYG{p}{(}\PYG{n}{np}\PYG{o}{.}\PYG{n}{arange}\PYG{p}{(}\PYG{o}{\PYGZhy{}}\PYG{l+m+mi}{180}\PYG{p}{,} \PYG{l+m+mi}{181}\PYG{p}{,} \PYG{l+m+mi}{90}\PYG{p}{)}\PYG{p}{)}
\PYG{n}{ax}\PYG{o}{.}\PYG{n}{set\PYGZus{}ylabel}\PYG{p}{(}\PYG{l+s+sa}{r}\PYG{l+s+s1}{\PYGZsq{}}\PYG{l+s+s1}{\PYGZdl{}}\PYG{l+s+s1}{\PYGZbs{}}\PYG{l+s+s1}{varphi}\PYG{l+s+s1}{\PYGZbs{}}\PYG{l+s+s1}{;(}\PYG{l+s+s1}{\PYGZbs{}}\PYG{l+s+s1}{degree)\PYGZdl{}}\PYG{l+s+s1}{\PYGZsq{}}\PYG{p}{)}
\PYG{n}{ax}\PYG{o}{.}\PYG{n}{grid}\PYG{p}{(}\PYG{k+kc}{True}\PYG{p}{,} \PYG{n}{which}\PYG{o}{=}\PYG{l+s+s1}{\PYGZsq{}}\PYG{l+s+s1}{major}\PYG{l+s+s1}{\PYGZsq{}}\PYG{p}{,} \PYG{n}{linestyle}\PYG{o}{=}\PYG{l+s+s1}{\PYGZsq{}}\PYG{l+s+s1}{\PYGZhy{}\PYGZhy{}}\PYG{l+s+s1}{\PYGZsq{}}\PYG{p}{,} \PYG{n}{linewidth}\PYG{o}{=}\PYG{l+m+mf}{0.5}\PYG{p}{)}
\PYG{n}{ax}\PYG{o}{.}\PYG{n}{legend}\PYG{p}{(}\PYG{n}{loc}\PYG{o}{=}\PYG{l+s+s1}{\PYGZsq{}}\PYG{l+s+s1}{upper right}\PYG{l+s+s1}{\PYGZsq{}}\PYG{p}{,} \PYG{n}{bbox\PYGZus{}to\PYGZus{}anchor}\PYG{o}{=}\PYG{p}{(}\PYG{l+m+mf}{1.15}\PYG{p}{,} \PYG{l+m+mf}{1.0}\PYG{p}{)}\PYG{p}{)}

\PYG{c+c1}{\PYGZsh{} 3. 倾子}
\PYG{n}{ax} \PYG{o}{=} \PYG{n}{axs}\PYG{p}{[}\PYG{l+m+mi}{2}\PYG{p}{]}
\PYG{n}{ax}\PYG{o}{.}\PYG{n}{semilogx}\PYG{p}{(}\PYG{n}{fre}\PYG{p}{,} \PYG{n}{tzxr}\PYG{p}{,} \PYG{l+s+s1}{\PYGZsq{}}\PYG{l+s+s1}{rv}\PYG{l+s+s1}{\PYGZsq{}}\PYG{p}{,} \PYG{n}{label}\PYG{o}{=}\PYG{l+s+s1}{\PYGZsq{}}\PYG{l+s+s1}{Xr}\PYG{l+s+s1}{\PYGZsq{}}\PYG{p}{,} \PYG{n}{markerfacecolor}\PYG{o}{=}\PYG{l+s+s1}{\PYGZsq{}}\PYG{l+s+s1}{none}\PYG{l+s+s1}{\PYGZsq{}}\PYG{p}{)}
\PYG{n}{ax}\PYG{o}{.}\PYG{n}{semilogx}\PYG{p}{(}\PYG{n}{fre}\PYG{p}{,} \PYG{n}{tzxi}\PYG{p}{,} \PYG{l+s+s1}{\PYGZsq{}}\PYG{l+s+s1}{r\PYGZca{}}\PYG{l+s+s1}{\PYGZsq{}}\PYG{p}{,} \PYG{n}{label}\PYG{o}{=}\PYG{l+s+s1}{\PYGZsq{}}\PYG{l+s+s1}{Xi}\PYG{l+s+s1}{\PYGZsq{}}\PYG{p}{,} \PYG{n}{markerfacecolor}\PYG{o}{=}\PYG{l+s+s1}{\PYGZsq{}}\PYG{l+s+s1}{none}\PYG{l+s+s1}{\PYGZsq{}}\PYG{p}{)}
\PYG{n}{ax}\PYG{o}{.}\PYG{n}{semilogx}\PYG{p}{(}\PYG{n}{fre}\PYG{p}{,} \PYG{n}{tzyr}\PYG{p}{,} \PYG{l+s+s1}{\PYGZsq{}}\PYG{l+s+s1}{bo}\PYG{l+s+s1}{\PYGZsq{}}\PYG{p}{,} \PYG{n}{label}\PYG{o}{=}\PYG{l+s+s1}{\PYGZsq{}}\PYG{l+s+s1}{Yr}\PYG{l+s+s1}{\PYGZsq{}}\PYG{p}{,} \PYG{n}{markerfacecolor}\PYG{o}{=}\PYG{l+s+s1}{\PYGZsq{}}\PYG{l+s+s1}{none}\PYG{l+s+s1}{\PYGZsq{}}\PYG{p}{)}
\PYG{n}{ax}\PYG{o}{.}\PYG{n}{semilogx}\PYG{p}{(}\PYG{n}{fre}\PYG{p}{,} \PYG{n}{tzyi}\PYG{p}{,} \PYG{l+s+s1}{\PYGZsq{}}\PYG{l+s+s1}{bs}\PYG{l+s+s1}{\PYGZsq{}}\PYG{p}{,} \PYG{n}{label}\PYG{o}{=}\PYG{l+s+s1}{\PYGZsq{}}\PYG{l+s+s1}{Yi}\PYG{l+s+s1}{\PYGZsq{}}\PYG{p}{,} \PYG{n}{markerfacecolor}\PYG{o}{=}\PYG{l+s+s1}{\PYGZsq{}}\PYG{l+s+s1}{none}\PYG{l+s+s1}{\PYGZsq{}}\PYG{p}{)}
\PYG{n}{ax}\PYG{o}{.}\PYG{n}{set\PYGZus{}ylim}\PYG{p}{(}\PYG{o}{\PYGZhy{}}\PYG{l+m+mf}{0.5}\PYG{p}{,} \PYG{l+m+mf}{0.5}\PYG{p}{)}
\PYG{n}{ax}\PYG{o}{.}\PYG{n}{set\PYGZus{}ylabel}\PYG{p}{(}\PYG{l+s+s1}{\PYGZsq{}}\PYG{l+s+s1}{Tipper}\PYG{l+s+s1}{\PYGZsq{}}\PYG{p}{)}
\PYG{n}{ax}\PYG{o}{.}\PYG{n}{grid}\PYG{p}{(}\PYG{k+kc}{True}\PYG{p}{,} \PYG{n}{which}\PYG{o}{=}\PYG{l+s+s1}{\PYGZsq{}}\PYG{l+s+s1}{major}\PYG{l+s+s1}{\PYGZsq{}}\PYG{p}{,} \PYG{n}{linestyle}\PYG{o}{=}\PYG{l+s+s1}{\PYGZsq{}}\PYG{l+s+s1}{\PYGZhy{}\PYGZhy{}}\PYG{l+s+s1}{\PYGZsq{}}\PYG{p}{,} \PYG{n}{linewidth}\PYG{o}{=}\PYG{l+m+mf}{0.5}\PYG{p}{)}
\PYG{n}{ax}\PYG{o}{.}\PYG{n}{legend}\PYG{p}{(}\PYG{n}{loc}\PYG{o}{=}\PYG{l+s+s1}{\PYGZsq{}}\PYG{l+s+s1}{upper right}\PYG{l+s+s1}{\PYGZsq{}}\PYG{p}{,} \PYG{n}{bbox\PYGZus{}to\PYGZus{}anchor}\PYG{o}{=}\PYG{p}{(}\PYG{l+m+mf}{1.15}\PYG{p}{,} \PYG{l+m+mf}{1.0}\PYG{p}{)}\PYG{p}{)}

\PYG{c+c1}{\PYGZsh{} 4. 相位张量椭圆}
\PYG{n}{ax} \PYG{o}{=} \PYG{n}{axs}\PYG{p}{[}\PYG{l+m+mi}{3}\PYG{p}{]}
\PYG{n}{axLinear} \PYG{o}{=} \PYG{n}{ax}\PYG{o}{.}\PYG{n}{twiny}\PYG{p}{(}\PYG{p}{)}
\PYG{n}{axLinear}\PYG{o}{.}\PYG{n}{set\PYGZus{}xscale}\PYG{p}{(}\PYG{l+s+s1}{\PYGZsq{}}\PYG{l+s+s1}{linear}\PYG{l+s+s1}{\PYGZsq{}}\PYG{p}{)}
\PYG{n}{axLinear}\PYG{o}{.}\PYG{n}{set\PYGZus{}xlim}\PYG{p}{(}\PYG{n}{np}\PYG{o}{.}\PYG{n}{log10}\PYG{p}{(}\PYG{n}{x\PYGZus{}min}\PYG{p}{)}\PYG{p}{,} \PYG{n}{np}\PYG{o}{.}\PYG{n}{log10}\PYG{p}{(}\PYG{n}{x\PYGZus{}max}\PYG{p}{)}\PYG{p}{)}
\PYG{n}{axLinear}\PYG{o}{.}\PYG{n}{invert\PYGZus{}xaxis}\PYG{p}{(}\PYG{p}{)}

\PYG{c+c1}{\PYGZsh{} 相位张量参数}
\PYG{n}{beta} \PYG{o}{=} \PYG{n}{np}\PYG{o}{.}\PYG{n}{loadtxt}\PYG{p}{(}\PYG{l+s+sa}{f}\PYG{l+s+s1}{\PYGZsq{}}\PYG{l+s+si}{\PYGZob{}}\PYG{n}{path}\PYG{l+s+si}{\PYGZcb{}}\PYG{l+s+s1}{/}\PYG{l+s+si}{\PYGZob{}}\PYG{n}{name}\PYG{l+s+si}{\PYGZcb{}}\PYG{l+s+s1}{\PYGZhy{}beta.dat}\PYG{l+s+s1}{\PYGZsq{}}\PYG{p}{)}\PYG{p}{[}\PYG{p}{:}\PYG{p}{,} \PYG{l+m+mi}{1}\PYG{p}{]}
\PYG{n}{alpha} \PYG{o}{=} \PYG{n}{np}\PYG{o}{.}\PYG{n}{loadtxt}\PYG{p}{(}\PYG{l+s+sa}{f}\PYG{l+s+s1}{\PYGZsq{}}\PYG{l+s+si}{\PYGZob{}}\PYG{n}{path}\PYG{l+s+si}{\PYGZcb{}}\PYG{l+s+s1}{/}\PYG{l+s+si}{\PYGZob{}}\PYG{n}{name}\PYG{l+s+si}{\PYGZcb{}}\PYG{l+s+s1}{\PYGZhy{}alpha.dat}\PYG{l+s+s1}{\PYGZsq{}}\PYG{p}{)}\PYG{p}{[}\PYG{p}{:}\PYG{p}{,} \PYG{l+m+mi}{1}\PYG{p}{]}
\PYG{n}{pmin} \PYG{o}{=} \PYG{n}{np}\PYG{o}{.}\PYG{n}{abs}\PYG{p}{(}\PYG{n}{np}\PYG{o}{.}\PYG{n}{loadtxt}\PYG{p}{(}\PYG{l+s+sa}{f}\PYG{l+s+s1}{\PYGZsq{}}\PYG{l+s+si}{\PYGZob{}}\PYG{n}{path}\PYG{l+s+si}{\PYGZcb{}}\PYG{l+s+s1}{/}\PYG{l+s+si}{\PYGZob{}}\PYG{n}{name}\PYG{l+s+si}{\PYGZcb{}}\PYG{l+s+s1}{\PYGZhy{}pmin.dat}\PYG{l+s+s1}{\PYGZsq{}}\PYG{p}{)}\PYG{p}{[}\PYG{p}{:}\PYG{p}{,} \PYG{l+m+mi}{1}\PYG{p}{]}\PYG{p}{)}
\PYG{n}{pmax} \PYG{o}{=} \PYG{n}{np}\PYG{o}{.}\PYG{n}{abs}\PYG{p}{(}\PYG{n}{np}\PYG{o}{.}\PYG{n}{loadtxt}\PYG{p}{(}\PYG{l+s+sa}{f}\PYG{l+s+s1}{\PYGZsq{}}\PYG{l+s+si}{\PYGZob{}}\PYG{n}{path}\PYG{l+s+si}{\PYGZcb{}}\PYG{l+s+s1}{/}\PYG{l+s+si}{\PYGZob{}}\PYG{n}{name}\PYG{l+s+si}{\PYGZcb{}}\PYG{l+s+s1}{\PYGZhy{}pmax.dat}\PYG{l+s+s1}{\PYGZsq{}}\PYG{p}{)}\PYG{p}{[}\PYG{p}{:}\PYG{p}{,} \PYG{l+m+mi}{1}\PYG{p}{]}\PYG{p}{)}

\PYG{c+c1}{\PYGZsh{} 归一化椭圆尺寸}
\PYG{n}{midp} \PYG{o}{=} \PYG{l+m+mi}{1}
\PYG{n}{ptscale} \PYG{o}{=} \PYG{l+m+mi}{1} \PYG{o}{/} \PYG{p}{(}\PYG{n}{np}\PYG{o}{.}\PYG{n}{log10}\PYG{p}{(}\PYG{n}{x\PYGZus{}max}\PYG{p}{)} \PYG{o}{\PYGZhy{}} \PYG{n}{np}\PYG{o}{.}\PYG{n}{log10}\PYG{p}{(}\PYG{n}{x\PYGZus{}min}\PYG{p}{)}\PYG{p}{)}
\PYG{n}{pmax} \PYG{o}{=} \PYG{n}{pmax} \PYG{o}{/} \PYG{n}{midp} \PYG{o}{*} \PYG{n}{ptscale}
\PYG{n}{pmin} \PYG{o}{=} \PYG{n}{pmin} \PYG{o}{/} \PYG{n}{midp} \PYG{o}{*} \PYG{n}{ptscale}
\PYG{n}{pmax}\PYG{p}{[}\PYG{n}{pmax} \PYG{o}{\PYGZgt{}} \PYG{l+m+mi}{3} \PYG{o}{*} \PYG{n}{ptscale}\PYG{p}{]} \PYG{o}{=} \PYG{l+m+mi}{3} \PYG{o}{*} \PYG{n}{ptscale}

\PYG{c+c1}{\PYGZsh{} 绘制相位张量椭圆}
\PYG{n}{cmap} \PYG{o}{=} \PYG{n}{plt}\PYG{o}{.}\PYG{n}{get\PYGZus{}cmap}\PYG{p}{(}\PYG{l+s+s1}{\PYGZsq{}}\PYG{l+s+s1}{rainbow}\PYG{l+s+s1}{\PYGZsq{}}\PYG{p}{)}
\PYG{n}{tmpbeta} \PYG{o}{=} \PYG{n}{np}\PYG{o}{.}\PYG{n}{abs}\PYG{p}{(}\PYG{n}{beta}\PYG{o}{.}\PYG{n}{copy}\PYG{p}{(}\PYG{p}{)}\PYG{p}{)}
\PYG{n}{tmpbeta}\PYG{p}{[}\PYG{n}{tmpbeta} \PYG{o}{\PYGZgt{}} \PYG{l+m+mi}{10}\PYG{p}{]} \PYG{o}{=} \PYG{l+m+mi}{10}
\PYG{n}{tmpbeta} \PYG{o}{=} \PYG{n}{tmpbeta} \PYG{o}{/} \PYG{l+m+mi}{10}

\PYG{n}{ellipses} \PYG{o}{=} \PYG{n}{EllipseCollection}\PYG{p}{(}
    \PYG{n}{widths}\PYG{o}{=}\PYG{n}{pmax}\PYG{p}{,} \PYG{n}{heights}\PYG{o}{=}\PYG{n}{pmin}\PYG{p}{,} \PYG{n}{angles}\PYG{o}{=}\PYG{n}{alpha} \PYG{o}{\PYGZhy{}} \PYG{n}{beta}\PYG{p}{,}
    \PYG{n}{units}\PYG{o}{=}\PYG{l+s+s1}{\PYGZsq{}}\PYG{l+s+s1}{xy}\PYG{l+s+s1}{\PYGZsq{}}\PYG{p}{,} \PYG{n}{offsets}\PYG{o}{=}\PYG{n}{np}\PYG{o}{.}\PYG{n}{column\PYGZus{}stack}\PYG{p}{(}\PYG{p}{(}\PYG{n}{np}\PYG{o}{.}\PYG{n}{log10}\PYG{p}{(}\PYG{n}{fre}\PYG{p}{)}\PYG{p}{,} \PYG{n}{np}\PYG{o}{.}\PYG{n}{ones\PYGZus{}like}\PYG{p}{(}\PYG{n}{fre}\PYG{p}{)} \PYG{o}{*} \PYG{l+m+mf}{0.5}\PYG{p}{)}\PYG{p}{)}\PYG{p}{,}
    \PYG{n}{transOffset}\PYG{o}{=}\PYG{n}{axLinear}\PYG{o}{.}\PYG{n}{transData}\PYG{p}{,} \PYG{n}{facecolor}\PYG{o}{=}\PYG{n}{cmap}\PYG{p}{(}\PYG{n}{tmpbeta}\PYG{p}{)}\PYG{p}{,} \PYG{n}{lw}\PYG{o}{=}\PYG{l+m+mi}{0}
\PYG{p}{)}
\PYG{n}{axLinear}\PYG{o}{.}\PYG{n}{add\PYGZus{}collection}\PYG{p}{(}\PYG{n}{ellipses}\PYG{p}{)}
\PYG{n}{axLinear}\PYG{o}{.}\PYG{n}{get\PYGZus{}xaxis}\PYG{p}{(}\PYG{p}{)}\PYG{o}{.}\PYG{n}{set\PYGZus{}visible}\PYG{p}{(}\PYG{k+kc}{False}\PYG{p}{)}

\PYG{c+c1}{\PYGZsh{} 添加色标}
\PYG{n}{norm} \PYG{o}{=} \PYG{n}{Normalize}\PYG{p}{(}\PYG{n}{vmin}\PYG{o}{=}\PYG{l+m+mi}{0}\PYG{p}{,} \PYG{n}{vmax}\PYG{o}{=}\PYG{l+m+mi}{10}\PYG{p}{)}
\PYG{n}{cax} \PYG{o}{=} \PYG{n}{inset\PYGZus{}axes}\PYG{p}{(}\PYG{n}{axs}\PYG{p}{[}\PYG{l+m+mi}{3}\PYG{p}{]}\PYG{p}{,} \PYG{n}{width}\PYG{o}{=}\PYG{l+s+s2}{\PYGZdq{}}\PYG{l+s+s2}{5}\PYG{l+s+s2}{\PYGZpc{}}\PYG{l+s+s2}{\PYGZdq{}}\PYG{p}{,} \PYG{n}{height}\PYG{o}{=}\PYG{l+s+s2}{\PYGZdq{}}\PYG{l+s+s2}{100}\PYG{l+s+s2}{\PYGZpc{}}\PYG{l+s+s2}{\PYGZdq{}}\PYG{p}{,} \PYG{n}{loc}\PYG{o}{=}\PYG{l+s+s1}{\PYGZsq{}}\PYG{l+s+s1}{upper right}\PYG{l+s+s1}{\PYGZsq{}}\PYG{p}{,}
                 \PYG{n}{bbox\PYGZus{}to\PYGZus{}anchor}\PYG{o}{=}\PYG{p}{(}\PYG{l+m+mf}{0.1}\PYG{p}{,} \PYG{l+m+mf}{0.}\PYG{p}{,} \PYG{l+m+mi}{1}\PYG{p}{,} \PYG{l+m+mi}{1}\PYG{p}{)}\PYG{p}{,} \PYG{n}{bbox\PYGZus{}transform}\PYG{o}{=}\PYG{n}{axs}\PYG{p}{[}\PYG{l+m+mi}{3}\PYG{p}{]}\PYG{o}{.}\PYG{n}{transAxes}\PYG{p}{)}
\PYG{n}{cb} \PYG{o}{=} \PYG{n}{plt}\PYG{o}{.}\PYG{n}{colorbar}\PYG{p}{(}\PYG{n}{ScalarMappable}\PYG{p}{(}\PYG{n}{cmap}\PYG{o}{=}\PYG{n}{cmap}\PYG{p}{,} \PYG{n}{norm}\PYG{o}{=}\PYG{n}{norm}\PYG{p}{)}\PYG{p}{,} \PYG{n}{cax}\PYG{o}{=}\PYG{n}{cax}\PYG{p}{,} \PYG{n}{orientation}\PYG{o}{=}\PYG{l+s+s1}{\PYGZsq{}}\PYG{l+s+s1}{vertical}\PYG{l+s+s1}{\PYGZsq{}}\PYG{p}{)}
\PYG{n}{cb}\PYG{o}{.}\PYG{n}{set\PYGZus{}label}\PYG{p}{(}\PYG{l+s+sa}{r}\PYG{l+s+s1}{\PYGZsq{}}\PYG{l+s+s1}{\PYGZdl{}|}\PYG{l+s+s1}{\PYGZbs{}}\PYG{l+s+s1}{beta|\PYGZdl{}}\PYG{l+s+s1}{\PYGZsq{}}\PYG{p}{)}

\PYG{n}{ax}\PYG{o}{.}\PYG{n}{set\PYGZus{}xlabel}\PYG{p}{(}\PYG{l+s+s1}{\PYGZsq{}}\PYG{l+s+s1}{Frequency (Hz)}\PYG{l+s+s1}{\PYGZsq{}}\PYG{p}{)}
\PYG{n}{ax}\PYG{o}{.}\PYG{n}{set\PYGZus{}ylabel}\PYG{p}{(}\PYG{l+s+s1}{\PYGZsq{}}\PYG{l+s+s1}{Phase Tensor}\PYG{l+s+s1}{\PYGZsq{}}\PYG{p}{)}

\PYG{c+c1}{\PYGZsh{} 保存}
\PYG{n}{plt}\PYG{o}{.}\PYG{n}{savefig}\PYG{p}{(}\PYG{l+s+sa}{f}\PYG{l+s+s1}{\PYGZsq{}}\PYG{l+s+si}{\PYGZob{}}\PYG{n}{path}\PYG{l+s+si}{\PYGZcb{}}\PYG{l+s+s1}{/}\PYG{l+s+si}{\PYGZob{}}\PYG{n}{name}\PYG{l+s+si}{\PYGZcb{}}\PYG{l+s+s1}{\PYGZhy{}Z.png}\PYG{l+s+s1}{\PYGZsq{}}\PYG{p}{,} \PYG{n}{dpi}\PYG{o}{=}\PYG{l+m+mi}{300}\PYG{p}{,} \PYG{n}{bbox\PYGZus{}inches}\PYG{o}{=}\PYG{l+s+s1}{\PYGZsq{}}\PYG{l+s+s1}{tight}\PYG{l+s+s1}{\PYGZsq{}}\PYG{p}{)}
\PYG{n}{plt}\PYG{o}{.}\PYG{n}{savefig}\PYG{p}{(}\PYG{l+s+sa}{f}\PYG{l+s+s1}{\PYGZsq{}}\PYG{l+s+si}{\PYGZob{}}\PYG{n}{path}\PYG{l+s+si}{\PYGZcb{}}\PYG{l+s+s1}{/}\PYG{l+s+si}{\PYGZob{}}\PYG{n}{name}\PYG{l+s+si}{\PYGZcb{}}\PYG{l+s+s1}{\PYGZhy{}Z.pdf}\PYG{l+s+s1}{\PYGZsq{}}\PYG{p}{,} \PYG{n+nb}{format}\PYG{o}{=}\PYG{l+s+s1}{\PYGZsq{}}\PYG{l+s+s1}{pdf}\PYG{l+s+s1}{\PYGZsq{}}\PYG{p}{,} \PYG{n}{dpi}\PYG{o}{=}\PYG{l+m+mi}{300}\PYG{p}{,} \PYG{n}{bbox\PYGZus{}inches}\PYG{o}{=}\PYG{l+s+s1}{\PYGZsq{}}\PYG{l+s+s1}{tight}\PYG{l+s+s1}{\PYGZsq{}}\PYG{p}{)}
\PYG{n}{plt}\PYG{o}{.}\PYG{n}{show}\PYG{p}{(}\PYG{p}{)}
\end{sphinxVerbatim}


\bigskip\hrule\bigskip



\subsection{PDF 合并工具}
\label{\detokenize{tutorial/plotting:pdf}}
\sphinxAtStartPar
将多个 PDF 文件合并为一个。


\subsubsection{脚本：\sphinxhref{http://MergerAllZPDF.py}{MergerAllZPDF.py} (http://MergerAllZPDF.py)}
\label{\detokenize{tutorial/plotting:mergerallzpdf-py}}
\begin{sphinxVerbatim}[commandchars=\\\{\}]
\PYG{k+kn}{import}\PYG{+w}{ }\PYG{n+nn}{os}
\PYG{k+kn}{from}\PYG{+w}{ }\PYG{n+nn}{PyPDF2}\PYG{+w}{ }\PYG{k+kn}{import} \PYG{n}{PdfMerger}

\PYG{c+c1}{\PYGZsh{} 配置参数}
\PYG{n}{path} \PYG{o}{=} \PYG{l+s+sa}{r}\PYG{l+s+s1}{\PYGZsq{}}\PYG{l+s+s1}{D:}\PYG{l+s+s1}{\PYGZbs{}}\PYG{l+s+s1}{Research}\PYG{l+s+s1}{\PYGZbs{}}\PYG{l+s+s1}{MT3D}\PYG{l+s+s1}{\PYGZbs{}}\PYG{l+s+s1}{Figures}\PYG{l+s+s1}{\PYGZsq{}}
\PYG{n}{output\PYGZus{}name} \PYG{o}{=} \PYG{l+s+s1}{\PYGZsq{}}\PYG{l+s+s1}{ALL}\PYG{l+s+s1}{\PYGZsq{}}

\PYG{c+c1}{\PYGZsh{} 获取所有 PDF 文件}
\PYG{n}{pdf\PYGZus{}files} \PYG{o}{=} \PYG{p}{[}\PYG{n}{os}\PYG{o}{.}\PYG{n}{path}\PYG{o}{.}\PYG{n}{join}\PYG{p}{(}\PYG{n}{path}\PYG{p}{,} \PYG{n}{f}\PYG{p}{)} \PYG{k}{for} \PYG{n}{f} \PYG{o+ow}{in} \PYG{n}{os}\PYG{o}{.}\PYG{n}{listdir}\PYG{p}{(}\PYG{n}{path}\PYG{p}{)} \PYG{k}{if} \PYG{n}{f}\PYG{o}{.}\PYG{n}{endswith}\PYG{p}{(}\PYG{l+s+s1}{\PYGZsq{}}\PYG{l+s+s1}{.pdf}\PYG{l+s+s1}{\PYGZsq{}}\PYG{p}{)}\PYG{p}{]}

\PYG{c+c1}{\PYGZsh{} 合并 PDF}
\PYG{n}{merger} \PYG{o}{=} \PYG{n}{PdfMerger}\PYG{p}{(}\PYG{p}{)}
\PYG{k}{for} \PYG{n}{pdf\PYGZus{}file} \PYG{o+ow}{in} \PYG{n}{pdf\PYGZus{}files}\PYG{p}{:}
    \PYG{n}{merger}\PYG{o}{.}\PYG{n}{append}\PYG{p}{(}\PYG{n}{pdf\PYGZus{}file}\PYG{p}{)}

\PYG{c+c1}{\PYGZsh{} 保存}
\PYG{n}{output\PYGZus{}path} \PYG{o}{=} \PYG{n}{os}\PYG{o}{.}\PYG{n}{path}\PYG{o}{.}\PYG{n}{join}\PYG{p}{(}\PYG{n}{path}\PYG{p}{,} \PYG{l+s+sa}{f}\PYG{l+s+s1}{\PYGZsq{}}\PYG{l+s+si}{\PYGZob{}}\PYG{n}{output\PYGZus{}name}\PYG{l+s+si}{\PYGZcb{}}\PYG{l+s+s1}{.pdf}\PYG{l+s+s1}{\PYGZsq{}}\PYG{p}{)}
\PYG{n}{merger}\PYG{o}{.}\PYG{n}{write}\PYG{p}{(}\PYG{n}{output\PYGZus{}path}\PYG{p}{)}
\PYG{n}{merger}\PYG{o}{.}\PYG{n}{close}\PYG{p}{(}\PYG{p}{)}

\PYG{n+nb}{print}\PYG{p}{(}\PYG{l+s+sa}{f}\PYG{l+s+s1}{\PYGZsq{}}\PYG{l+s+s1}{Combined PDF saved: }\PYG{l+s+si}{\PYGZob{}}\PYG{n}{output\PYGZus{}path}\PYG{l+s+si}{\PYGZcb{}}\PYG{l+s+s1}{\PYGZsq{}}\PYG{p}{)}
\end{sphinxVerbatim}


\bigskip\hrule\bigskip



\subsection{坐标转换工具}
\label{\detokenize{tutorial/plotting:id10}}
\sphinxAtStartPar
在 GIS 应用中，可能需要将 MTDataPro 导出的坐标转换为不同的坐标系。


\subsubsection{脚本：\sphinxhref{http://coorConvert.py}{coorConvert.py} (http://coorConvert.py)}
\label{\detokenize{tutorial/plotting:coorconvert-py}}
\begin{sphinxVerbatim}[commandchars=\\\{\}]
\PYG{k+kn}{from}\PYG{+w}{ }\PYG{n+nn}{shapely}\PYG{n+nn}{.}\PYG{n+nn}{geometry}\PYG{+w}{ }\PYG{k+kn}{import} \PYG{n}{Point}
\PYG{k+kn}{from}\PYG{+w}{ }\PYG{n+nn}{pyproj}\PYG{+w}{ }\PYG{k+kn}{import} \PYG{n}{CRS}\PYG{p}{,} \PYG{n}{Transformer}

\PYG{c+c1}{\PYGZsh{} WGS84 坐标}
\PYG{n}{wgs84\PYGZus{}point} \PYG{o}{=} \PYG{n}{Point}\PYG{p}{(}\PYG{l+m+mf}{120.5}\PYG{p}{,} \PYG{l+m+mf}{30.25}\PYG{p}{)}  \PYG{c+c1}{\PYGZsh{} (经度, 纬度)}

\PYG{c+c1}{\PYGZsh{} 转换到目标坐标系（如 Web Mercator）}
\PYG{n}{target\PYGZus{}crs} \PYG{o}{=} \PYG{n}{CRS}\PYG{p}{(}\PYG{l+s+s2}{\PYGZdq{}}\PYG{l+s+s2}{EPSG:3857}\PYG{l+s+s2}{\PYGZdq{}}\PYG{p}{)}
\PYG{n}{transformer} \PYG{o}{=} \PYG{n}{Transformer}\PYG{o}{.}\PYG{n}{from\PYGZus{}crs}\PYG{p}{(}\PYG{n}{CRS}\PYG{p}{(}\PYG{l+s+s2}{\PYGZdq{}}\PYG{l+s+s2}{EPSG:4326}\PYG{l+s+s2}{\PYGZdq{}}\PYG{p}{)}\PYG{p}{,} \PYG{n}{target\PYGZus{}crs}\PYG{p}{)}

\PYG{c+c1}{\PYGZsh{} 执行转换}
\PYG{n}{x}\PYG{p}{,} \PYG{n}{y} \PYG{o}{=} \PYG{n}{transformer}\PYG{o}{.}\PYG{n}{transform}\PYG{p}{(}\PYG{n}{wgs84\PYGZus{}point}\PYG{o}{.}\PYG{n}{x}\PYG{p}{,} \PYG{n}{wgs84\PYGZus{}point}\PYG{o}{.}\PYG{n}{y}\PYG{p}{)}
\PYG{n+nb}{print}\PYG{p}{(}\PYG{l+s+sa}{f}\PYG{l+s+s1}{\PYGZsq{}}\PYG{l+s+s1}{转换后坐标: }\PYG{l+s+si}{\PYGZob{}}\PYG{n}{x}\PYG{l+s+si}{\PYGZcb{}}\PYG{l+s+s1}{, }\PYG{l+s+si}{\PYGZob{}}\PYG{n}{y}\PYG{l+s+si}{\PYGZcb{}}\PYG{l+s+s1}{\PYGZsq{}}\PYG{p}{)}

\PYG{c+c1}{\PYGZsh{} 支持的坐标系}
\PYG{c+c1}{\PYGZsh{} EPSG:4326 \PYGZhy{} WGS84}
\PYG{c+c1}{\PYGZsh{} EPSG:3857 \PYGZhy{} Web Mercator}
\PYG{c+c1}{\PYGZsh{} EPSG:4490 \PYGZhy{} China Geodetic Coordinate System 2000}
\PYG{c+c1}{\PYGZsh{} EPSG:4549 \PYGZhy{} China TM\PYGZhy{}3 117E}
\end{sphinxVerbatim}


\bigskip\hrule\bigskip



\subsection{常见问题}
\label{\detokenize{tutorial/plotting:id11}}

\subsubsection{Q: 绘图的频率范围如何设置？}
\label{\detokenize{tutorial/plotting:q}}
\sphinxAtStartPar
A: 根据您的数据频率范围调整 \sphinxcode{\sphinxupquote{xmin}} 和 \sphinxcode{\sphinxupquote{xmax}}：
\begin{itemize}
\item {} 
\sphinxAtStartPar
AMT 数据：通常 10 Hz \textasciitilde{} 10 kHz

\item {} 
\sphinxAtStartPar
MT 数据：通常 0.001 Hz \textasciitilde{} 100 Hz

\item {} 
\sphinxAtStartPar
远参考 MT：可低至 0.0001 Hz

\end{itemize}


\subsubsection{Q: 误差棒不显示？}
\label{\detokenize{tutorial/plotting:id12}}
\sphinxAtStartPar
A: 确保从 MTDataPro 导出了误差数据文件（\sphinxcode{\sphinxupquote{*\sphinxhyphen{}rxyvar.dat}} 等）。误差在对数坐标下需要转换：

\begin{sphinxVerbatim}[commandchars=\\\{\}]
\PYG{n}{log\PYGZus{}error} \PYG{o}{=} \PYG{n}{error} \PYG{o}{/} \PYG{p}{(}\PYG{n}{value} \PYG{o}{*} \PYG{n}{np}\PYG{o}{.}\PYG{n}{log}\PYG{p}{(}\PYG{l+m+mi}{10}\PYG{p}{)}\PYG{p}{)}
\end{sphinxVerbatim}


\subsubsection{Q: 如何调整图形尺寸？}
\label{\detokenize{tutorial/plotting:id13}}
\sphinxAtStartPar
A: 修改 \sphinxcode{\sphinxupquote{figsize}} 参数：

\begin{sphinxVerbatim}[commandchars=\\\{\}]
\PYG{n}{fig}\PYG{p}{,} \PYG{n}{axs} \PYG{o}{=} \PYG{n}{plt}\PYG{o}{.}\PYG{n}{subplots}\PYG{p}{(}\PYG{l+m+mi}{2}\PYG{p}{,} \PYG{l+m+mi}{1}\PYG{p}{,} \PYG{n}{figsize}\PYG{o}{=}\PYG{p}{(}\PYG{l+m+mi}{10}\PYG{p}{,} \PYG{l+m+mi}{8}\PYG{p}{)}\PYG{p}{)}  \PYG{c+c1}{\PYGZsh{} 宽10英寸，高8英寸}
\end{sphinxVerbatim}


\subsubsection{Q: 中文标签显示为方块？}
\label{\detokenize{tutorial/plotting:id14}}
\sphinxAtStartPar
A: 使用系统支持的中文字体：

\begin{sphinxVerbatim}[commandchars=\\\{\}]
\PYG{n}{plt}\PYG{o}{.}\PYG{n}{rcParams}\PYG{p}{[}\PYG{l+s+s1}{\PYGZsq{}}\PYG{l+s+s1}{font.sans\PYGZhy{}serif}\PYG{l+s+s1}{\PYGZsq{}}\PYG{p}{]} \PYG{o}{=} \PYG{p}{[}\PYG{l+s+s1}{\PYGZsq{}}\PYG{l+s+s1}{SimHei}\PYG{l+s+s1}{\PYGZsq{}}\PYG{p}{,} \PYG{l+s+s1}{\PYGZsq{}}\PYG{l+s+s1}{DejaVu Sans}\PYG{l+s+s1}{\PYGZsq{}}\PYG{p}{]}
\PYG{n}{plt}\PYG{o}{.}\PYG{n}{rcParams}\PYG{p}{[}\PYG{l+s+s1}{\PYGZsq{}}\PYG{l+s+s1}{axes.unicode\PYGZus{}minus}\PYG{l+s+s1}{\PYGZsq{}}\PYG{p}{]} \PYG{o}{=} \PYG{k+kc}{False}
\end{sphinxVerbatim}

\sphinxAtStartPar
建议使用英文标签以确保跨平台兼容性。

\sphinxstepscope


\section{附录A: 快捷键参考}
\label{\detokenize{appendices/appendixA:a}}\label{\detokenize{appendices/appendixA::doc}}

\subsection{常用快捷键}
\label{\detokenize{appendices/appendixA:id1}}

\begin{savenotes}\sphinxattablestart
\sphinxthistablewithglobalstyle
\centering
\begin{tabulary}{\linewidth}[t]{TT}
\sphinxtoprule
\begin{varwidth}[t]{\sphinxcolwidth{1}{2}}
\sphinxstyletheadfamily \sphinxAtStartPar
快捷键
\sphinxbeforeendvarwidth
\end{varwidth}%
&\begin{varwidth}[t]{\sphinxcolwidth{1}{2}}
\sphinxstyletheadfamily \sphinxAtStartPar
功能
\sphinxbeforeendvarwidth
\end{varwidth}%
\\
\sphinxmidrule
\sphinxtableatstartofbodyhook\begin{varwidth}[t]{\sphinxcolwidth{1}{2}}
\sphinxAtStartPar
Ctrl+N
\sphinxbeforeendvarwidth
\end{varwidth}%
&\begin{varwidth}[t]{\sphinxcolwidth{1}{2}}
\sphinxAtStartPar
新建工程
\sphinxbeforeendvarwidth
\end{varwidth}%
\\
\sphinxhline\begin{varwidth}[t]{\sphinxcolwidth{1}{2}}
\sphinxAtStartPar
Ctrl+O
\sphinxbeforeendvarwidth
\end{varwidth}%
&\begin{varwidth}[t]{\sphinxcolwidth{1}{2}}
\sphinxAtStartPar
打开工程
\sphinxbeforeendvarwidth
\end{varwidth}%
\\
\sphinxhline\begin{varwidth}[t]{\sphinxcolwidth{1}{2}}
\sphinxAtStartPar
Ctrl+S
\sphinxbeforeendvarwidth
\end{varwidth}%
&\begin{varwidth}[t]{\sphinxcolwidth{1}{2}}
\sphinxAtStartPar
保存工程
\sphinxbeforeendvarwidth
\end{varwidth}%
\\
\sphinxhline\begin{varwidth}[t]{\sphinxcolwidth{1}{2}}
\sphinxAtStartPar
Ctrl+Z
\sphinxbeforeendvarwidth
\end{varwidth}%
&\begin{varwidth}[t]{\sphinxcolwidth{1}{2}}
\sphinxAtStartPar
撤销
\sphinxbeforeendvarwidth
\end{varwidth}%
\\
\sphinxhline\begin{varwidth}[t]{\sphinxcolwidth{1}{2}}
\sphinxAtStartPar
Ctrl+Y
\sphinxbeforeendvarwidth
\end{varwidth}%
&\begin{varwidth}[t]{\sphinxcolwidth{1}{2}}
\sphinxAtStartPar
重做
\sphinxbeforeendvarwidth
\end{varwidth}%
\\
\sphinxhline\begin{varwidth}[t]{\sphinxcolwidth{1}{2}}
\sphinxAtStartPar
Ctrl+A
\sphinxbeforeendvarwidth
\end{varwidth}%
&\begin{varwidth}[t]{\sphinxcolwidth{1}{2}}
\sphinxAtStartPar
全选
\sphinxbeforeendvarwidth
\end{varwidth}%
\\
\sphinxhline\begin{varwidth}[t]{\sphinxcolwidth{1}{2}}
\sphinxAtStartPar
Ctrl+C
\sphinxbeforeendvarwidth
\end{varwidth}%
&\begin{varwidth}[t]{\sphinxcolwidth{1}{2}}
\sphinxAtStartPar
复制
\sphinxbeforeendvarwidth
\end{varwidth}%
\\
\sphinxhline\begin{varwidth}[t]{\sphinxcolwidth{1}{2}}
\sphinxAtStartPar
Ctrl+V
\sphinxbeforeendvarwidth
\end{varwidth}%
&\begin{varwidth}[t]{\sphinxcolwidth{1}{2}}
\sphinxAtStartPar
粘贴
\sphinxbeforeendvarwidth
\end{varwidth}%
\\
\sphinxhline\begin{varwidth}[t]{\sphinxcolwidth{1}{2}}
\sphinxAtStartPar
Ctrl+X
\sphinxbeforeendvarwidth
\end{varwidth}%
&\begin{varwidth}[t]{\sphinxcolwidth{1}{2}}
\sphinxAtStartPar
剪切
\sphinxbeforeendvarwidth
\end{varwidth}%
\\
\sphinxhline\begin{varwidth}[t]{\sphinxcolwidth{1}{2}}
\sphinxAtStartPar
Delete
\sphinxbeforeendvarwidth
\end{varwidth}%
&\begin{varwidth}[t]{\sphinxcolwidth{1}{2}}
\sphinxAtStartPar
删除选中项
\sphinxbeforeendvarwidth
\end{varwidth}%
\\
\sphinxhline\begin{varwidth}[t]{\sphinxcolwidth{1}{2}}
\sphinxAtStartPar
F5
\sphinxbeforeendvarwidth
\end{varwidth}%
&\begin{varwidth}[t]{\sphinxcolwidth{1}{2}}
\sphinxAtStartPar
刷新
\sphinxbeforeendvarwidth
\end{varwidth}%
\\
\sphinxhline\begin{varwidth}[t]{\sphinxcolwidth{1}{2}}
\sphinxAtStartPar
F11
\sphinxbeforeendvarwidth
\end{varwidth}%
&\begin{varwidth}[t]{\sphinxcolwidth{1}{2}}
\sphinxAtStartPar
全屏模式
\sphinxbeforeendvarwidth
\end{varwidth}%
\\
\sphinxhline\begin{varwidth}[t]{\sphinxcolwidth{1}{2}}
\sphinxAtStartPar
Ctrl+Q
\sphinxbeforeendvarwidth
\end{varwidth}%
&\begin{varwidth}[t]{\sphinxcolwidth{1}{2}}
\sphinxAtStartPar
退出程序
\sphinxbeforeendvarwidth
\end{varwidth}%
\\
\sphinxbottomrule
\end{tabulary}
\sphinxtableafterendhook\par
\sphinxattableend\end{savenotes}


\bigskip\hrule\bigskip



\subsection{处理相关快捷键}
\label{\detokenize{appendices/appendixA:id2}}

\begin{savenotes}\sphinxattablestart
\sphinxthistablewithglobalstyle
\centering
\begin{tabulary}{\linewidth}[t]{TT}
\sphinxtoprule
\begin{varwidth}[t]{\sphinxcolwidth{1}{2}}
\sphinxstyletheadfamily \sphinxAtStartPar
快捷键
\sphinxbeforeendvarwidth
\end{varwidth}%
&\begin{varwidth}[t]{\sphinxcolwidth{1}{2}}
\sphinxstyletheadfamily \sphinxAtStartPar
功能
\sphinxbeforeendvarwidth
\end{varwidth}%
\\
\sphinxmidrule
\sphinxtableatstartofbodyhook\begin{varwidth}[t]{\sphinxcolwidth{1}{2}}
\sphinxAtStartPar
F9
\sphinxbeforeendvarwidth
\end{varwidth}%
&\begin{varwidth}[t]{\sphinxcolwidth{1}{2}}
\sphinxAtStartPar
执行FFT计算
\sphinxbeforeendvarwidth
\end{varwidth}%
\\
\sphinxhline\begin{varwidth}[t]{\sphinxcolwidth{1}{2}}
\sphinxAtStartPar
Ctrl+F9
\sphinxbeforeendvarwidth
\end{varwidth}%
&\begin{varwidth}[t]{\sphinxcolwidth{1}{2}}
\sphinxAtStartPar
批量FFT计算
\sphinxbeforeendvarwidth
\end{varwidth}%
\\
\sphinxhline\begin{varwidth}[t]{\sphinxcolwidth{1}{2}}
\sphinxAtStartPar
F10
\sphinxbeforeendvarwidth
\end{varwidth}%
&\begin{varwidth}[t]{\sphinxcolwidth{1}{2}}
\sphinxAtStartPar
数据筛选
\sphinxbeforeendvarwidth
\end{varwidth}%
\\
\sphinxhline\begin{varwidth}[t]{\sphinxcolwidth{1}{2}}
\sphinxAtStartPar
Ctrl+R
\sphinxbeforeendvarwidth
\end{varwidth}%
&\begin{varwidth}[t]{\sphinxcolwidth{1}{2}}
\sphinxAtStartPar
Robust估计
\sphinxbeforeendvarwidth
\end{varwidth}%
\\
\sphinxbottomrule
\end{tabulary}
\sphinxtableafterendhook\par
\sphinxattableend\end{savenotes}


\bigskip\hrule\bigskip



\subsection{视图相关快捷键}
\label{\detokenize{appendices/appendixA:id3}}

\begin{savenotes}\sphinxattablestart
\sphinxthistablewithglobalstyle
\centering
\begin{tabulary}{\linewidth}[t]{TT}
\sphinxtoprule
\begin{varwidth}[t]{\sphinxcolwidth{1}{2}}
\sphinxstyletheadfamily \sphinxAtStartPar
快捷键
\sphinxbeforeendvarwidth
\end{varwidth}%
&\begin{varwidth}[t]{\sphinxcolwidth{1}{2}}
\sphinxstyletheadfamily \sphinxAtStartPar
功能
\sphinxbeforeendvarwidth
\end{varwidth}%
\\
\sphinxmidrule
\sphinxtableatstartofbodyhook\begin{varwidth}[t]{\sphinxcolwidth{1}{2}}
\sphinxAtStartPar
Ctrl+1
\sphinxbeforeendvarwidth
\end{varwidth}%
&\begin{varwidth}[t]{\sphinxcolwidth{1}{2}}
\sphinxAtStartPar
时间序列视图
\sphinxbeforeendvarwidth
\end{varwidth}%
\\
\sphinxhline\begin{varwidth}[t]{\sphinxcolwidth{1}{2}}
\sphinxAtStartPar
Ctrl+2
\sphinxbeforeendvarwidth
\end{varwidth}%
&\begin{varwidth}[t]{\sphinxcolwidth{1}{2}}
\sphinxAtStartPar
傅里叶系数视图
\sphinxbeforeendvarwidth
\end{varwidth}%
\\
\sphinxhline\begin{varwidth}[t]{\sphinxcolwidth{1}{2}}
\sphinxAtStartPar
Ctrl+3
\sphinxbeforeendvarwidth
\end{varwidth}%
&\begin{varwidth}[t]{\sphinxcolwidth{1}{2}}
\sphinxAtStartPar
频谱图视图
\sphinxbeforeendvarwidth
\end{varwidth}%
\\
\sphinxhline\begin{varwidth}[t]{\sphinxcolwidth{1}{2}}
\sphinxAtStartPar
Ctrl+4
\sphinxbeforeendvarwidth
\end{varwidth}%
&\begin{varwidth}[t]{\sphinxcolwidth{1}{2}}
\sphinxAtStartPar
响应函数视图
\sphinxbeforeendvarwidth
\end{varwidth}%
\\
\sphinxhline\begin{varwidth}[t]{\sphinxcolwidth{1}{2}}
\sphinxAtStartPar
Ctrl+W
\sphinxbeforeendvarwidth
\end{varwidth}%
&\begin{varwidth}[t]{\sphinxcolwidth{1}{2}}
\sphinxAtStartPar
关闭当前标签页
\sphinxbeforeendvarwidth
\end{varwidth}%
\\
\sphinxhline\begin{varwidth}[t]{\sphinxcolwidth{1}{2}}
\sphinxAtStartPar
Ctrl+Tab
\sphinxbeforeendvarwidth
\end{varwidth}%
&\begin{varwidth}[t]{\sphinxcolwidth{1}{2}}
\sphinxAtStartPar
切换到下一个标签页
\sphinxbeforeendvarwidth
\end{varwidth}%
\\
\sphinxhline\begin{varwidth}[t]{\sphinxcolwidth{1}{2}}
\sphinxAtStartPar
Ctrl+Shift+Tab
\sphinxbeforeendvarwidth
\end{varwidth}%
&\begin{varwidth}[t]{\sphinxcolwidth{1}{2}}
\sphinxAtStartPar
切换到上一个标签页
\sphinxbeforeendvarwidth
\end{varwidth}%
\\
\sphinxbottomrule
\end{tabulary}
\sphinxtableafterendhook\par
\sphinxattableend\end{savenotes}


\bigskip\hrule\bigskip



\subsection{工程管理快捷键}
\label{\detokenize{appendices/appendixA:id4}}

\begin{savenotes}\sphinxattablestart
\sphinxthistablewithglobalstyle
\centering
\begin{tabulary}{\linewidth}[t]{TT}
\sphinxtoprule
\begin{varwidth}[t]{\sphinxcolwidth{1}{2}}
\sphinxstyletheadfamily \sphinxAtStartPar
快捷键
\sphinxbeforeendvarwidth
\end{varwidth}%
&\begin{varwidth}[t]{\sphinxcolwidth{1}{2}}
\sphinxstyletheadfamily \sphinxAtStartPar
功能
\sphinxbeforeendvarwidth
\end{varwidth}%
\\
\sphinxmidrule
\sphinxtableatstartofbodyhook\begin{varwidth}[t]{\sphinxcolwidth{1}{2}}
\sphinxAtStartPar
Ctrl+Shift+N
\sphinxbeforeendvarwidth
\end{varwidth}%
&\begin{varwidth}[t]{\sphinxcolwidth{1}{2}}
\sphinxAtStartPar
新建工区
\sphinxbeforeendvarwidth
\end{varwidth}%
\\
\sphinxhline\begin{varwidth}[t]{\sphinxcolwidth{1}{2}}
\sphinxAtStartPar
Ctrl+Shift+S
\sphinxbeforeendvarwidth
\end{varwidth}%
&\begin{varwidth}[t]{\sphinxcolwidth{1}{2}}
\sphinxAtStartPar
另存为
\sphinxbeforeendvarwidth
\end{varwidth}%
\\
\sphinxhline\begin{varwidth}[t]{\sphinxcolwidth{1}{2}}
\sphinxAtStartPar
Ctrl+E
\sphinxbeforeendvarwidth
\end{varwidth}%
&\begin{varwidth}[t]{\sphinxcolwidth{1}{2}}
\sphinxAtStartPar
导出数据
\sphinxbeforeendvarwidth
\end{varwidth}%
\\
\sphinxhline\begin{varwidth}[t]{\sphinxcolwidth{1}{2}}
\sphinxAtStartPar
Ctrl+I
\sphinxbeforeendvarwidth
\end{varwidth}%
&\begin{varwidth}[t]{\sphinxcolwidth{1}{2}}
\sphinxAtStartPar
导入数据
\sphinxbeforeendvarwidth
\end{varwidth}%
\\
\sphinxbottomrule
\end{tabulary}
\sphinxtableafterendhook\par
\sphinxattableend\end{savenotes}

\sphinxstepscope


\section{附录B: 参数说明表}
\label{\detokenize{appendices/appendixB:b}}\label{\detokenize{appendices/appendixB::doc}}

\subsection{处理方案参数}
\label{\detokenize{appendices/appendixB:id1}}

\begin{savenotes}\sphinxattablestart
\sphinxthistablewithglobalstyle
\centering
\begin{tabulary}{\linewidth}[t]{TTTT}
\sphinxtoprule
\begin{varwidth}[t]{\sphinxcolwidth{1}{4}}
\sphinxstyletheadfamily \sphinxAtStartPar
参数名
\sphinxbeforeendvarwidth
\end{varwidth}%
&\begin{varwidth}[t]{\sphinxcolwidth{1}{4}}
\sphinxstyletheadfamily \sphinxAtStartPar
类型
\sphinxbeforeendvarwidth
\end{varwidth}%
&\begin{varwidth}[t]{\sphinxcolwidth{1}{4}}
\sphinxstyletheadfamily \sphinxAtStartPar
默认值
\sphinxbeforeendvarwidth
\end{varwidth}%
&\begin{varwidth}[t]{\sphinxcolwidth{1}{4}}
\sphinxstyletheadfamily \sphinxAtStartPar
说明
\sphinxbeforeendvarwidth
\end{varwidth}%
\\
\sphinxmidrule
\sphinxtableatstartofbodyhook\begin{varwidth}[t]{\sphinxcolwidth{1}{4}}
\sphinxAtStartPar
Window
\sphinxbeforeendvarwidth
\end{varwidth}%
&\begin{varwidth}[t]{\sphinxcolwidth{1}{4}}
\sphinxAtStartPar
Integer
\sphinxbeforeendvarwidth
\end{varwidth}%
&\begin{varwidth}[t]{\sphinxcolwidth{1}{4}}
\sphinxAtStartPar
2
\sphinxbeforeendvarwidth
\end{varwidth}%
&\begin{varwidth}[t]{\sphinxcolwidth{1}{4}}
\sphinxAtStartPar
窗函数类型（0=矩形,1=汉明,2=汉宁）
\sphinxbeforeendvarwidth
\end{varwidth}%
\\
\sphinxhline\begin{varwidth}[t]{\sphinxcolwidth{1}{4}}
\sphinxAtStartPar
FFTDFT
\sphinxbeforeendvarwidth
\end{varwidth}%
&\begin{varwidth}[t]{\sphinxcolwidth{1}{4}}
\sphinxAtStartPar
Boolean
\sphinxbeforeendvarwidth
\end{varwidth}%
&\begin{varwidth}[t]{\sphinxcolwidth{1}{4}}
\sphinxAtStartPar
False
\sphinxbeforeendvarwidth
\end{varwidth}%
&\begin{varwidth}[t]{\sphinxcolwidth{1}{4}}
\sphinxAtStartPar
FFT/DFT选择
\sphinxbeforeendvarwidth
\end{varwidth}%
\\
\sphinxhline\begin{varwidth}[t]{\sphinxcolwidth{1}{4}}
\sphinxAtStartPar
SingleFrequencyCalculationWay
\sphinxbeforeendvarwidth
\end{varwidth}%
&\begin{varwidth}[t]{\sphinxcolwidth{1}{4}}
\sphinxAtStartPar
Integer
\sphinxbeforeendvarwidth
\end{varwidth}%
&\begin{varwidth}[t]{\sphinxcolwidth{1}{4}}
\sphinxAtStartPar
0
\sphinxbeforeendvarwidth
\end{varwidth}%
&\begin{varwidth}[t]{\sphinxcolwidth{1}{4}}
\sphinxAtStartPar
单频计算方式
\sphinxbeforeendvarwidth
\end{varwidth}%
\\
\sphinxhline\begin{varwidth}[t]{\sphinxcolwidth{1}{4}}
\sphinxAtStartPar
ZeroPadding
\sphinxbeforeendvarwidth
\end{varwidth}%
&\begin{varwidth}[t]{\sphinxcolwidth{1}{4}}
\sphinxAtStartPar
Double
\sphinxbeforeendvarwidth
\end{varwidth}%
&\begin{varwidth}[t]{\sphinxcolwidth{1}{4}}
\sphinxAtStartPar
0
\sphinxbeforeendvarwidth
\end{varwidth}%
&\begin{varwidth}[t]{\sphinxcolwidth{1}{4}}
\sphinxAtStartPar
零填充因子
\sphinxbeforeendvarwidth
\end{varwidth}%
\\
\sphinxhline\begin{varwidth}[t]{\sphinxcolwidth{1}{4}}
\sphinxAtStartPar
MultiTaper
\sphinxbeforeendvarwidth
\end{varwidth}%
&\begin{varwidth}[t]{\sphinxcolwidth{1}{4}}
\sphinxAtStartPar
Integer
\sphinxbeforeendvarwidth
\end{varwidth}%
&\begin{varwidth}[t]{\sphinxcolwidth{1}{4}}
\sphinxAtStartPar
0
\sphinxbeforeendvarwidth
\end{varwidth}%
&\begin{varwidth}[t]{\sphinxcolwidth{1}{4}}
\sphinxAtStartPar
多锥谱分析锥数
\sphinxbeforeendvarwidth
\end{varwidth}%
\\
\sphinxhline\begin{varwidth}[t]{\sphinxcolwidth{1}{4}}
\sphinxAtStartPar
SearchBand
\sphinxbeforeendvarwidth
\end{varwidth}%
&\begin{varwidth}[t]{\sphinxcolwidth{1}{4}}
\sphinxAtStartPar
Double
\sphinxbeforeendvarwidth
\end{varwidth}%
&\begin{varwidth}[t]{\sphinxcolwidth{1}{4}}
\sphinxAtStartPar
0.1
\sphinxbeforeendvarwidth
\end{varwidth}%
&\begin{varwidth}[t]{\sphinxcolwidth{1}{4}}
\sphinxAtStartPar
搜索带宽
\sphinxbeforeendvarwidth
\end{varwidth}%
\\
\sphinxhline\begin{varwidth}[t]{\sphinxcolwidth{1}{4}}
\sphinxAtStartPar
ReferenceEstimator
\sphinxbeforeendvarwidth
\end{varwidth}%
&\begin{varwidth}[t]{\sphinxcolwidth{1}{4}}
\sphinxAtStartPar
Integer
\sphinxbeforeendvarwidth
\end{varwidth}%
&\begin{varwidth}[t]{\sphinxcolwidth{1}{4}}
\sphinxAtStartPar
0
\sphinxbeforeendvarwidth
\end{varwidth}%
&\begin{varwidth}[t]{\sphinxcolwidth{1}{4}}
\sphinxAtStartPar
参考道估计方法
\sphinxbeforeendvarwidth
\end{varwidth}%
\\
\sphinxhline\begin{varwidth}[t]{\sphinxcolwidth{1}{4}}
\sphinxAtStartPar
RobustEstimator
\sphinxbeforeendvarwidth
\end{varwidth}%
&\begin{varwidth}[t]{\sphinxcolwidth{1}{4}}
\sphinxAtStartPar
Integer
\sphinxbeforeendvarwidth
\end{varwidth}%
&\begin{varwidth}[t]{\sphinxcolwidth{1}{4}}
\sphinxAtStartPar
1
\sphinxbeforeendvarwidth
\end{varwidth}%
&\begin{varwidth}[t]{\sphinxcolwidth{1}{4}}
\sphinxAtStartPar
稳健估计方法
\sphinxbeforeendvarwidth
\end{varwidth}%
\\
\sphinxbottomrule
\end{tabulary}
\sphinxtableafterendhook\par
\sphinxattableend\end{savenotes}


\bigskip\hrule\bigskip



\subsection{FFT参数}
\label{\detokenize{appendices/appendixB:fft}}

\begin{savenotes}\sphinxattablestart
\sphinxthistablewithglobalstyle
\centering
\begin{tabulary}{\linewidth}[t]{TTTT}
\sphinxtoprule
\begin{varwidth}[t]{\sphinxcolwidth{1}{4}}
\sphinxstyletheadfamily \sphinxAtStartPar
参数名
\sphinxbeforeendvarwidth
\end{varwidth}%
&\begin{varwidth}[t]{\sphinxcolwidth{1}{4}}
\sphinxstyletheadfamily \sphinxAtStartPar
类型
\sphinxbeforeendvarwidth
\end{varwidth}%
&\begin{varwidth}[t]{\sphinxcolwidth{1}{4}}
\sphinxstyletheadfamily \sphinxAtStartPar
默认值
\sphinxbeforeendvarwidth
\end{varwidth}%
&\begin{varwidth}[t]{\sphinxcolwidth{1}{4}}
\sphinxstyletheadfamily \sphinxAtStartPar
说明
\sphinxbeforeendvarwidth
\end{varwidth}%
\\
\sphinxmidrule
\sphinxtableatstartofbodyhook\begin{varwidth}[t]{\sphinxcolwidth{1}{4}}
\sphinxAtStartPar
SampleLength
\sphinxbeforeendvarwidth
\end{varwidth}%
&\begin{varwidth}[t]{\sphinxcolwidth{1}{4}}
\sphinxAtStartPar
Double
\sphinxbeforeendvarwidth
\end{varwidth}%
&\begin{varwidth}[t]{\sphinxcolwidth{1}{4}}
\sphinxAtStartPar
4096
\sphinxbeforeendvarwidth
\end{varwidth}%
&\begin{varwidth}[t]{\sphinxcolwidth{1}{4}}
\sphinxAtStartPar
样本长度（FFT长度）
\sphinxbeforeendvarwidth
\end{varwidth}%
\\
\sphinxhline\begin{varwidth}[t]{\sphinxcolwidth{1}{4}}
\sphinxAtStartPar
Overlap
\sphinxbeforeendvarwidth
\end{varwidth}%
&\begin{varwidth}[t]{\sphinxcolwidth{1}{4}}
\sphinxAtStartPar
Double
\sphinxbeforeendvarwidth
\end{varwidth}%
&\begin{varwidth}[t]{\sphinxcolwidth{1}{4}}
\sphinxAtStartPar
0.5
\sphinxbeforeendvarwidth
\end{varwidth}%
&\begin{varwidth}[t]{\sphinxcolwidth{1}{4}}
\sphinxAtStartPar
重叠率（0\sphinxhyphen{}0.75）
\sphinxbeforeendvarwidth
\end{varwidth}%
\\
\sphinxhline\begin{varwidth}[t]{\sphinxcolwidth{1}{4}}
\sphinxAtStartPar
MaxXPR
\sphinxbeforeendvarwidth
\end{varwidth}%
&\begin{varwidth}[t]{\sphinxcolwidth{1}{4}}
\sphinxAtStartPar
Integer
\sphinxbeforeendvarwidth
\end{varwidth}%
&\begin{varwidth}[t]{\sphinxcolwidth{1}{4}}
\sphinxAtStartPar
100
\sphinxbeforeendvarwidth
\end{varwidth}%
&\begin{varwidth}[t]{\sphinxcolwidth{1}{4}}
\sphinxAtStartPar
最大XPR值
\sphinxbeforeendvarwidth
\end{varwidth}%
\\
\sphinxhline\begin{varwidth}[t]{\sphinxcolwidth{1}{4}}
\sphinxAtStartPar
GroupingType
\sphinxbeforeendvarwidth
\end{varwidth}%
&\begin{varwidth}[t]{\sphinxcolwidth{1}{4}}
\sphinxAtStartPar
Integer
\sphinxbeforeendvarwidth
\end{varwidth}%
&\begin{varwidth}[t]{\sphinxcolwidth{1}{4}}
\sphinxAtStartPar
0
\sphinxbeforeendvarwidth
\end{varwidth}%
&\begin{varwidth}[t]{\sphinxcolwidth{1}{4}}
\sphinxAtStartPar
分组类型
\sphinxbeforeendvarwidth
\end{varwidth}%
\\
\sphinxbottomrule
\end{tabulary}
\sphinxtableafterendhook\par
\sphinxattableend\end{savenotes}


\bigskip\hrule\bigskip



\subsection{稳健估计参数}
\label{\detokenize{appendices/appendixB:id2}}

\begin{savenotes}\sphinxattablestart
\sphinxthistablewithglobalstyle
\centering
\begin{tabulary}{\linewidth}[t]{TTTT}
\sphinxtoprule
\begin{varwidth}[t]{\sphinxcolwidth{1}{4}}
\sphinxstyletheadfamily \sphinxAtStartPar
参数名
\sphinxbeforeendvarwidth
\end{varwidth}%
&\begin{varwidth}[t]{\sphinxcolwidth{1}{4}}
\sphinxstyletheadfamily \sphinxAtStartPar
类型
\sphinxbeforeendvarwidth
\end{varwidth}%
&\begin{varwidth}[t]{\sphinxcolwidth{1}{4}}
\sphinxstyletheadfamily \sphinxAtStartPar
默认值
\sphinxbeforeendvarwidth
\end{varwidth}%
&\begin{varwidth}[t]{\sphinxcolwidth{1}{4}}
\sphinxstyletheadfamily \sphinxAtStartPar
说明
\sphinxbeforeendvarwidth
\end{varwidth}%
\\
\sphinxmidrule
\sphinxtableatstartofbodyhook\begin{varwidth}[t]{\sphinxcolwidth{1}{4}}
\sphinxAtStartPar
maxiters
\sphinxbeforeendvarwidth
\end{varwidth}%
&\begin{varwidth}[t]{\sphinxcolwidth{1}{4}}
\sphinxAtStartPar
Integer
\sphinxbeforeendvarwidth
\end{varwidth}%
&\begin{varwidth}[t]{\sphinxcolwidth{1}{4}}
\sphinxAtStartPar
100
\sphinxbeforeendvarwidth
\end{varwidth}%
&\begin{varwidth}[t]{\sphinxcolwidth{1}{4}}
\sphinxAtStartPar
最大迭代次数
\sphinxbeforeendvarwidth
\end{varwidth}%
\\
\sphinxhline\begin{varwidth}[t]{\sphinxcolwidth{1}{4}}
\sphinxAtStartPar
eps
\sphinxbeforeendvarwidth
\end{varwidth}%
&\begin{varwidth}[t]{\sphinxcolwidth{1}{4}}
\sphinxAtStartPar
Double
\sphinxbeforeendvarwidth
\end{varwidth}%
&\begin{varwidth}[t]{\sphinxcolwidth{1}{4}}
\sphinxAtStartPar
0.01
\sphinxbeforeendvarwidth
\end{varwidth}%
&\begin{varwidth}[t]{\sphinxcolwidth{1}{4}}
\sphinxAtStartPar
迭代步长容差
\sphinxbeforeendvarwidth
\end{varwidth}%
\\
\sphinxhline\begin{varwidth}[t]{\sphinxcolwidth{1}{4}}
\sphinxAtStartPar
epslast
\sphinxbeforeendvarwidth
\end{varwidth}%
&\begin{varwidth}[t]{\sphinxcolwidth{1}{4}}
\sphinxAtStartPar
Double
\sphinxbeforeendvarwidth
\end{varwidth}%
&\begin{varwidth}[t]{\sphinxcolwidth{1}{4}}
\sphinxAtStartPar
0.0001
\sphinxbeforeendvarwidth
\end{varwidth}%
&\begin{varwidth}[t]{\sphinxcolwidth{1}{4}}
\sphinxAtStartPar
最终收敛容差
\sphinxbeforeendvarwidth
\end{varwidth}%
\\
\sphinxhline\begin{varwidth}[t]{\sphinxcolwidth{1}{4}}
\sphinxAtStartPar
r0
\sphinxbeforeendvarwidth
\end{varwidth}%
&\begin{varwidth}[t]{\sphinxcolwidth{1}{4}}
\sphinxAtStartPar
Double
\sphinxbeforeendvarwidth
\end{varwidth}%
&\begin{varwidth}[t]{\sphinxcolwidth{1}{4}}
\sphinxAtStartPar
1.5
\sphinxbeforeendvarwidth
\end{varwidth}%
&\begin{varwidth}[t]{\sphinxcolwidth{1}{4}}
\sphinxAtStartPar
调谐参数（λ）
\sphinxbeforeendvarwidth
\end{varwidth}%
\\
\sphinxhline\begin{varwidth}[t]{\sphinxcolwidth{1}{4}}
\sphinxAtStartPar
ainuin
\sphinxbeforeendvarwidth
\end{varwidth}%
&\begin{varwidth}[t]{\sphinxcolwidth{1}{4}}
\sphinxAtStartPar
Double
\sphinxbeforeendvarwidth
\end{varwidth}%
&\begin{varwidth}[t]{\sphinxcolwidth{1}{4}}
\sphinxAtStartPar
0.9999
\sphinxbeforeendvarwidth
\end{varwidth}%
&\begin{varwidth}[t]{\sphinxcolwidth{1}{4}}
\sphinxAtStartPar
有界影响上界
\sphinxbeforeendvarwidth
\end{varwidth}%
\\
\sphinxhline\begin{varwidth}[t]{\sphinxcolwidth{1}{4}}
\sphinxAtStartPar
ainlin
\sphinxbeforeendvarwidth
\end{varwidth}%
&\begin{varwidth}[t]{\sphinxcolwidth{1}{4}}
\sphinxAtStartPar
Double
\sphinxbeforeendvarwidth
\end{varwidth}%
&\begin{varwidth}[t]{\sphinxcolwidth{1}{4}}
\sphinxAtStartPar
0.0001
\sphinxbeforeendvarwidth
\end{varwidth}%
&\begin{varwidth}[t]{\sphinxcolwidth{1}{4}}
\sphinxAtStartPar
有界影响下界
\sphinxbeforeendvarwidth
\end{varwidth}%
\\
\sphinxbottomrule
\end{tabulary}
\sphinxtableafterendhook\par
\sphinxattableend\end{savenotes}


\bigskip\hrule\bigskip



\subsection{参考道估计方法}
\label{\detokenize{appendices/appendixB:id3}}

\begin{savenotes}\sphinxattablestart
\sphinxthistablewithglobalstyle
\centering
\begin{tabulary}{\linewidth}[t]{TTT}
\sphinxtoprule
\begin{varwidth}[t]{\sphinxcolwidth{1}{3}}
\sphinxstyletheadfamily \sphinxAtStartPar
值
\sphinxbeforeendvarwidth
\end{varwidth}%
&\begin{varwidth}[t]{\sphinxcolwidth{1}{3}}
\sphinxstyletheadfamily \sphinxAtStartPar
方法
\sphinxbeforeendvarwidth
\end{varwidth}%
&\begin{varwidth}[t]{\sphinxcolwidth{1}{3}}
\sphinxstyletheadfamily \sphinxAtStartPar
代码
\sphinxbeforeendvarwidth
\end{varwidth}%
\\
\sphinxmidrule
\sphinxtableatstartofbodyhook\begin{varwidth}[t]{\sphinxcolwidth{1}{3}}
\sphinxAtStartPar
0
\sphinxbeforeendvarwidth
\end{varwidth}%
&\begin{varwidth}[t]{\sphinxcolwidth{1}{3}}
\sphinxAtStartPar
Local E
\sphinxbeforeendvarwidth
\end{varwidth}%
&\begin{varwidth}[t]{\sphinxcolwidth{1}{3}}
\sphinxAtStartPar
LE
\sphinxbeforeendvarwidth
\end{varwidth}%
\\
\sphinxhline\begin{varwidth}[t]{\sphinxcolwidth{1}{3}}
\sphinxAtStartPar
1
\sphinxbeforeendvarwidth
\end{varwidth}%
&\begin{varwidth}[t]{\sphinxcolwidth{1}{3}}
\sphinxAtStartPar
Local H
\sphinxbeforeendvarwidth
\end{varwidth}%
&\begin{varwidth}[t]{\sphinxcolwidth{1}{3}}
\sphinxAtStartPar
LH
\sphinxbeforeendvarwidth
\end{varwidth}%
\\
\sphinxhline\begin{varwidth}[t]{\sphinxcolwidth{1}{3}}
\sphinxAtStartPar
2
\sphinxbeforeendvarwidth
\end{varwidth}%
&\begin{varwidth}[t]{\sphinxcolwidth{1}{3}}
\sphinxAtStartPar
Local E/H
\sphinxbeforeendvarwidth
\end{varwidth}%
&\begin{varwidth}[t]{\sphinxcolwidth{1}{3}}
\sphinxAtStartPar
LEH
\sphinxbeforeendvarwidth
\end{varwidth}%
\\
\sphinxhline\begin{varwidth}[t]{\sphinxcolwidth{1}{3}}
\sphinxAtStartPar
3
\sphinxbeforeendvarwidth
\end{varwidth}%
&\begin{varwidth}[t]{\sphinxcolwidth{1}{3}}
\sphinxAtStartPar
Remote E
\sphinxbeforeendvarwidth
\end{varwidth}%
&\begin{varwidth}[t]{\sphinxcolwidth{1}{3}}
\sphinxAtStartPar
RE
\sphinxbeforeendvarwidth
\end{varwidth}%
\\
\sphinxhline\begin{varwidth}[t]{\sphinxcolwidth{1}{3}}
\sphinxAtStartPar
4
\sphinxbeforeendvarwidth
\end{varwidth}%
&\begin{varwidth}[t]{\sphinxcolwidth{1}{3}}
\sphinxAtStartPar
Remote H
\sphinxbeforeendvarwidth
\end{varwidth}%
&\begin{varwidth}[t]{\sphinxcolwidth{1}{3}}
\sphinxAtStartPar
RH
\sphinxbeforeendvarwidth
\end{varwidth}%
\\
\sphinxhline\begin{varwidth}[t]{\sphinxcolwidth{1}{3}}
\sphinxAtStartPar
5
\sphinxbeforeendvarwidth
\end{varwidth}%
&\begin{varwidth}[t]{\sphinxcolwidth{1}{3}}
\sphinxAtStartPar
Remote E/H
\sphinxbeforeendvarwidth
\end{varwidth}%
&\begin{varwidth}[t]{\sphinxcolwidth{1}{3}}
\sphinxAtStartPar
REH
\sphinxbeforeendvarwidth
\end{varwidth}%
\\
\sphinxhline\begin{varwidth}[t]{\sphinxcolwidth{1}{3}}
\sphinxAtStartPar
6
\sphinxbeforeendvarwidth
\end{varwidth}%
&\begin{varwidth}[t]{\sphinxcolwidth{1}{3}}
\sphinxAtStartPar
Remote E/Local H
\sphinxbeforeendvarwidth
\end{varwidth}%
&\begin{varwidth}[t]{\sphinxcolwidth{1}{3}}
\sphinxAtStartPar
RELH
\sphinxbeforeendvarwidth
\end{varwidth}%
\\
\sphinxbottomrule
\end{tabulary}
\sphinxtableafterendhook\par
\sphinxattableend\end{savenotes}


\bigskip\hrule\bigskip



\subsection{稳健估计方法}
\label{\detokenize{appendices/appendixB:id4}}

\begin{savenotes}\sphinxattablestart
\sphinxthistablewithglobalstyle
\centering
\begin{tabulary}{\linewidth}[t]{TTT}
\sphinxtoprule
\begin{varwidth}[t]{\sphinxcolwidth{1}{3}}
\sphinxstyletheadfamily \sphinxAtStartPar
值
\sphinxbeforeendvarwidth
\end{varwidth}%
&\begin{varwidth}[t]{\sphinxcolwidth{1}{3}}
\sphinxstyletheadfamily \sphinxAtStartPar
方法
\sphinxbeforeendvarwidth
\end{varwidth}%
&\begin{varwidth}[t]{\sphinxcolwidth{1}{3}}
\sphinxstyletheadfamily \sphinxAtStartPar
说明
\sphinxbeforeendvarwidth
\end{varwidth}%
\\
\sphinxmidrule
\sphinxtableatstartofbodyhook\begin{varwidth}[t]{\sphinxcolwidth{1}{3}}
\sphinxAtStartPar
0
\sphinxbeforeendvarwidth
\end{varwidth}%
&\begin{varwidth}[t]{\sphinxcolwidth{1}{3}}
\sphinxAtStartPar
LeastSquares
\sphinxbeforeendvarwidth
\end{varwidth}%
&\begin{varwidth}[t]{\sphinxcolwidth{1}{3}}
\sphinxAtStartPar
标准最小二乘法
\sphinxbeforeendvarwidth
\end{varwidth}%
\\
\sphinxhline\begin{varwidth}[t]{\sphinxcolwidth{1}{3}}
\sphinxAtStartPar
1
\sphinxbeforeendvarwidth
\end{varwidth}%
&\begin{varwidth}[t]{\sphinxcolwidth{1}{3}}
\sphinxAtStartPar
Regression\sphinxhyphen{}M
\sphinxbeforeendvarwidth
\end{varwidth}%
&\begin{varwidth}[t]{\sphinxcolwidth{1}{3}}
\sphinxAtStartPar
M估计回归法
\sphinxbeforeendvarwidth
\end{varwidth}%
\\
\sphinxhline\begin{varwidth}[t]{\sphinxcolwidth{1}{3}}
\sphinxAtStartPar
2
\sphinxbeforeendvarwidth
\end{varwidth}%
&\begin{varwidth}[t]{\sphinxcolwidth{1}{3}}
\sphinxAtStartPar
Repeated Median
\sphinxbeforeendvarwidth
\end{varwidth}%
&\begin{varwidth}[t]{\sphinxcolwidth{1}{3}}
\sphinxAtStartPar
重复中位数法（高抗噪性）
\sphinxbeforeendvarwidth
\end{varwidth}%
\\
\sphinxhline\begin{varwidth}[t]{\sphinxcolwidth{1}{3}}
\sphinxAtStartPar
3
\sphinxbeforeendvarwidth
\end{varwidth}%
&\begin{varwidth}[t]{\sphinxcolwidth{1}{3}}
\sphinxAtStartPar
Bounded Influence
\sphinxbeforeendvarwidth
\end{varwidth}%
&\begin{varwidth}[t]{\sphinxcolwidth{1}{3}}
\sphinxAtStartPar
有界影响估计
\sphinxbeforeendvarwidth
\end{varwidth}%
\\
\sphinxhline\begin{varwidth}[t]{\sphinxcolwidth{1}{3}}
\sphinxAtStartPar
4
\sphinxbeforeendvarwidth
\end{varwidth}%
&\begin{varwidth}[t]{\sphinxcolwidth{1}{3}}
\sphinxAtStartPar
Huber Pre\sphinxhyphen{}Weighted
\sphinxbeforeendvarwidth
\end{varwidth}%
&\begin{varwidth}[t]{\sphinxcolwidth{1}{3}}
\sphinxAtStartPar
Huber预加权估计
\sphinxbeforeendvarwidth
\end{varwidth}%
\\
\sphinxhline\begin{varwidth}[t]{\sphinxcolwidth{1}{3}}
\sphinxAtStartPar
5
\sphinxbeforeendvarwidth
\end{varwidth}%
&\begin{varwidth}[t]{\sphinxcolwidth{1}{3}}
\sphinxAtStartPar
Thomson Pre\sphinxhyphen{}Weighted
\sphinxbeforeendvarwidth
\end{varwidth}%
&\begin{varwidth}[t]{\sphinxcolwidth{1}{3}}
\sphinxAtStartPar
Thomson预加权估计
\sphinxbeforeendvarwidth
\end{varwidth}%
\\
\sphinxhline\begin{varwidth}[t]{\sphinxcolwidth{1}{3}}
\sphinxAtStartPar
6
\sphinxbeforeendvarwidth
\end{varwidth}%
&\begin{varwidth}[t]{\sphinxcolwidth{1}{3}}
\sphinxAtStartPar
Robust+AI
\sphinxbeforeendvarwidth
\end{varwidth}%
&\begin{varwidth}[t]{\sphinxcolwidth{1}{3}}
\sphinxAtStartPar
稳健估计结合AI预测
\sphinxbeforeendvarwidth
\end{varwidth}%
\\
\sphinxbottomrule
\end{tabulary}
\sphinxtableafterendhook\par
\sphinxattableend\end{savenotes}


\bigskip\hrule\bigskip



\subsection{自动筛选参数}
\label{\detokenize{appendices/appendixB:id5}}

\begin{savenotes}\sphinxattablestart
\sphinxthistablewithglobalstyle
\centering
\begin{tabulary}{\linewidth}[t]{TTTT}
\sphinxtoprule
\begin{varwidth}[t]{\sphinxcolwidth{1}{4}}
\sphinxstyletheadfamily \sphinxAtStartPar
参数名
\sphinxbeforeendvarwidth
\end{varwidth}%
&\begin{varwidth}[t]{\sphinxcolwidth{1}{4}}
\sphinxstyletheadfamily \sphinxAtStartPar
类型
\sphinxbeforeendvarwidth
\end{varwidth}%
&\begin{varwidth}[t]{\sphinxcolwidth{1}{4}}
\sphinxstyletheadfamily \sphinxAtStartPar
默认值
\sphinxbeforeendvarwidth
\end{varwidth}%
&\begin{varwidth}[t]{\sphinxcolwidth{1}{4}}
\sphinxstyletheadfamily \sphinxAtStartPar
说明
\sphinxbeforeendvarwidth
\end{varwidth}%
\\
\sphinxmidrule
\sphinxtableatstartofbodyhook\begin{varwidth}[t]{\sphinxcolwidth{1}{4}}
\sphinxAtStartPar
MaxProportion
\sphinxbeforeendvarwidth
\end{varwidth}%
&\begin{varwidth}[t]{\sphinxcolwidth{1}{4}}
\sphinxAtStartPar
Double
\sphinxbeforeendvarwidth
\end{varwidth}%
&\begin{varwidth}[t]{\sphinxcolwidth{1}{4}}
\sphinxAtStartPar
0.9
\sphinxbeforeendvarwidth
\end{varwidth}%
&\begin{varwidth}[t]{\sphinxcolwidth{1}{4}}
\sphinxAtStartPar
最大保留比例
\sphinxbeforeendvarwidth
\end{varwidth}%
\\
\sphinxhline\begin{varwidth}[t]{\sphinxcolwidth{1}{4}}
\sphinxAtStartPar
MinProportion
\sphinxbeforeendvarwidth
\end{varwidth}%
&\begin{varwidth}[t]{\sphinxcolwidth{1}{4}}
\sphinxAtStartPar
Double
\sphinxbeforeendvarwidth
\end{varwidth}%
&\begin{varwidth}[t]{\sphinxcolwidth{1}{4}}
\sphinxAtStartPar
0.1
\sphinxbeforeendvarwidth
\end{varwidth}%
&\begin{varwidth}[t]{\sphinxcolwidth{1}{4}}
\sphinxAtStartPar
最小保留比例
\sphinxbeforeendvarwidth
\end{varwidth}%
\\
\sphinxhline\begin{varwidth}[t]{\sphinxcolwidth{1}{4}}
\sphinxAtStartPar
ExitError
\sphinxbeforeendvarwidth
\end{varwidth}%
&\begin{varwidth}[t]{\sphinxcolwidth{1}{4}}
\sphinxAtStartPar
Double
\sphinxbeforeendvarwidth
\end{varwidth}%
&\begin{varwidth}[t]{\sphinxcolwidth{1}{4}}
\sphinxAtStartPar
0.01
\sphinxbeforeendvarwidth
\end{varwidth}%
&\begin{varwidth}[t]{\sphinxcolwidth{1}{4}}
\sphinxAtStartPar
退出误差阈值
\sphinxbeforeendvarwidth
\end{varwidth}%
\\
\sphinxhline\begin{varwidth}[t]{\sphinxcolwidth{1}{4}}
\sphinxAtStartPar
AutoPrecent
\sphinxbeforeendvarwidth
\end{varwidth}%
&\begin{varwidth}[t]{\sphinxcolwidth{1}{4}}
\sphinxAtStartPar
Double
\sphinxbeforeendvarwidth
\end{varwidth}%
&\begin{varwidth}[t]{\sphinxcolwidth{1}{4}}
\sphinxAtStartPar
0.5
\sphinxbeforeendvarwidth
\end{varwidth}%
&\begin{varwidth}[t]{\sphinxcolwidth{1}{4}}
\sphinxAtStartPar
自动选择百分比
\sphinxbeforeendvarwidth
\end{varwidth}%
\\
\sphinxbottomrule
\end{tabulary}
\sphinxtableafterendhook\par
\sphinxattableend\end{savenotes}


\bigskip\hrule\bigskip



\subsection{马氏距离参数}
\label{\detokenize{appendices/appendixB:id6}}

\begin{savenotes}\sphinxattablestart
\sphinxthistablewithglobalstyle
\centering
\begin{tabulary}{\linewidth}[t]{TTTT}
\sphinxtoprule
\begin{varwidth}[t]{\sphinxcolwidth{1}{4}}
\sphinxstyletheadfamily \sphinxAtStartPar
参数名
\sphinxbeforeendvarwidth
\end{varwidth}%
&\begin{varwidth}[t]{\sphinxcolwidth{1}{4}}
\sphinxstyletheadfamily \sphinxAtStartPar
类型
\sphinxbeforeendvarwidth
\end{varwidth}%
&\begin{varwidth}[t]{\sphinxcolwidth{1}{4}}
\sphinxstyletheadfamily \sphinxAtStartPar
默认值
\sphinxbeforeendvarwidth
\end{varwidth}%
&\begin{varwidth}[t]{\sphinxcolwidth{1}{4}}
\sphinxstyletheadfamily \sphinxAtStartPar
说明
\sphinxbeforeendvarwidth
\end{varwidth}%
\\
\sphinxmidrule
\sphinxtableatstartofbodyhook\begin{varwidth}[t]{\sphinxcolwidth{1}{4}}
\sphinxAtStartPar
RobustC
\sphinxbeforeendvarwidth
\end{varwidth}%
&\begin{varwidth}[t]{\sphinxcolwidth{1}{4}}
\sphinxAtStartPar
Double
\sphinxbeforeendvarwidth
\end{varwidth}%
&\begin{varwidth}[t]{\sphinxcolwidth{1}{4}}
\sphinxAtStartPar
4.0
\sphinxbeforeendvarwidth
\end{varwidth}%
&\begin{varwidth}[t]{\sphinxcolwidth{1}{4}}
\sphinxAtStartPar
Huber权函数阈值
\sphinxbeforeendvarwidth
\end{varwidth}%
\\
\sphinxhline\begin{varwidth}[t]{\sphinxcolwidth{1}{4}}
\sphinxAtStartPar
MaxIterations
\sphinxbeforeendvarwidth
\end{varwidth}%
&\begin{varwidth}[t]{\sphinxcolwidth{1}{4}}
\sphinxAtStartPar
Integer
\sphinxbeforeendvarwidth
\end{varwidth}%
&\begin{varwidth}[t]{\sphinxcolwidth{1}{4}}
\sphinxAtStartPar
5
\sphinxbeforeendvarwidth
\end{varwidth}%
&\begin{varwidth}[t]{\sphinxcolwidth{1}{4}}
\sphinxAtStartPar
最大迭代次数
\sphinxbeforeendvarwidth
\end{varwidth}%
\\
\sphinxbottomrule
\end{tabulary}
\sphinxtableafterendhook\par
\sphinxattableend\end{savenotes}


\bigskip\hrule\bigskip



\subsection{遗传算法参数}
\label{\detokenize{appendices/appendixB:id7}}

\begin{savenotes}\sphinxattablestart
\sphinxthistablewithglobalstyle
\centering
\begin{tabulary}{\linewidth}[t]{TTTT}
\sphinxtoprule
\begin{varwidth}[t]{\sphinxcolwidth{1}{4}}
\sphinxstyletheadfamily \sphinxAtStartPar
参数名
\sphinxbeforeendvarwidth
\end{varwidth}%
&\begin{varwidth}[t]{\sphinxcolwidth{1}{4}}
\sphinxstyletheadfamily \sphinxAtStartPar
类型
\sphinxbeforeendvarwidth
\end{varwidth}%
&\begin{varwidth}[t]{\sphinxcolwidth{1}{4}}
\sphinxstyletheadfamily \sphinxAtStartPar
默认值
\sphinxbeforeendvarwidth
\end{varwidth}%
&\begin{varwidth}[t]{\sphinxcolwidth{1}{4}}
\sphinxstyletheadfamily \sphinxAtStartPar
说明
\sphinxbeforeendvarwidth
\end{varwidth}%
\\
\sphinxmidrule
\sphinxtableatstartofbodyhook\begin{varwidth}[t]{\sphinxcolwidth{1}{4}}
\sphinxAtStartPar
PopulationSize
\sphinxbeforeendvarwidth
\end{varwidth}%
&\begin{varwidth}[t]{\sphinxcolwidth{1}{4}}
\sphinxAtStartPar
Integer
\sphinxbeforeendvarwidth
\end{varwidth}%
&\begin{varwidth}[t]{\sphinxcolwidth{1}{4}}
\sphinxAtStartPar
100
\sphinxbeforeendvarwidth
\end{varwidth}%
&\begin{varwidth}[t]{\sphinxcolwidth{1}{4}}
\sphinxAtStartPar
种群大小
\sphinxbeforeendvarwidth
\end{varwidth}%
\\
\sphinxhline\begin{varwidth}[t]{\sphinxcolwidth{1}{4}}
\sphinxAtStartPar
Generations
\sphinxbeforeendvarwidth
\end{varwidth}%
&\begin{varwidth}[t]{\sphinxcolwidth{1}{4}}
\sphinxAtStartPar
Integer
\sphinxbeforeendvarwidth
\end{varwidth}%
&\begin{varwidth}[t]{\sphinxcolwidth{1}{4}}
\sphinxAtStartPar
200
\sphinxbeforeendvarwidth
\end{varwidth}%
&\begin{varwidth}[t]{\sphinxcolwidth{1}{4}}
\sphinxAtStartPar
最大进化代数
\sphinxbeforeendvarwidth
\end{varwidth}%
\\
\sphinxhline\begin{varwidth}[t]{\sphinxcolwidth{1}{4}}
\sphinxAtStartPar
CrossoverRate
\sphinxbeforeendvarwidth
\end{varwidth}%
&\begin{varwidth}[t]{\sphinxcolwidth{1}{4}}
\sphinxAtStartPar
Double
\sphinxbeforeendvarwidth
\end{varwidth}%
&\begin{varwidth}[t]{\sphinxcolwidth{1}{4}}
\sphinxAtStartPar
0.8
\sphinxbeforeendvarwidth
\end{varwidth}%
&\begin{varwidth}[t]{\sphinxcolwidth{1}{4}}
\sphinxAtStartPar
交叉概率
\sphinxbeforeendvarwidth
\end{varwidth}%
\\
\sphinxhline\begin{varwidth}[t]{\sphinxcolwidth{1}{4}}
\sphinxAtStartPar
MutationRate
\sphinxbeforeendvarwidth
\end{varwidth}%
&\begin{varwidth}[t]{\sphinxcolwidth{1}{4}}
\sphinxAtStartPar
Double
\sphinxbeforeendvarwidth
\end{varwidth}%
&\begin{varwidth}[t]{\sphinxcolwidth{1}{4}}
\sphinxAtStartPar
0.1
\sphinxbeforeendvarwidth
\end{varwidth}%
&\begin{varwidth}[t]{\sphinxcolwidth{1}{4}}
\sphinxAtStartPar
变异概率
\sphinxbeforeendvarwidth
\end{varwidth}%
\\
\sphinxbottomrule
\end{tabulary}
\sphinxtableafterendhook\par
\sphinxattableend\end{savenotes}


\bigskip\hrule\bigskip



\subsection{降采样参数}
\label{\detokenize{appendices/appendixB:id8}}

\begin{savenotes}\sphinxattablestart
\sphinxthistablewithglobalstyle
\centering
\begin{tabulary}{\linewidth}[t]{TTTT}
\sphinxtoprule
\begin{varwidth}[t]{\sphinxcolwidth{1}{4}}
\sphinxstyletheadfamily \sphinxAtStartPar
参数名
\sphinxbeforeendvarwidth
\end{varwidth}%
&\begin{varwidth}[t]{\sphinxcolwidth{1}{4}}
\sphinxstyletheadfamily \sphinxAtStartPar
类型
\sphinxbeforeendvarwidth
\end{varwidth}%
&\begin{varwidth}[t]{\sphinxcolwidth{1}{4}}
\sphinxstyletheadfamily \sphinxAtStartPar
可选值
\sphinxbeforeendvarwidth
\end{varwidth}%
&\begin{varwidth}[t]{\sphinxcolwidth{1}{4}}
\sphinxstyletheadfamily \sphinxAtStartPar
说明
\sphinxbeforeendvarwidth
\end{varwidth}%
\\
\sphinxmidrule
\sphinxtableatstartofbodyhook\begin{varwidth}[t]{\sphinxcolwidth{1}{4}}
\sphinxAtStartPar
Factor
\sphinxbeforeendvarwidth
\end{varwidth}%
&\begin{varwidth}[t]{\sphinxcolwidth{1}{4}}
\sphinxAtStartPar
Enum
\sphinxbeforeendvarwidth
\end{varwidth}%
&\begin{varwidth}[t]{\sphinxcolwidth{1}{4}}
\sphinxAtStartPar
2x\sphinxhyphen{}4800x
\sphinxbeforeendvarwidth
\end{varwidth}%
&\begin{varwidth}[t]{\sphinxcolwidth{1}{4}}
\sphinxAtStartPar
降采样因子
\sphinxbeforeendvarwidth
\end{varwidth}%
\\
\sphinxhline\begin{varwidth}[t]{\sphinxcolwidth{1}{4}}
\sphinxAtStartPar
FilterType
\sphinxbeforeendvarwidth
\end{varwidth}%
&\begin{varwidth}[t]{\sphinxcolwidth{1}{4}}
\sphinxAtStartPar
Enum
\sphinxbeforeendvarwidth
\end{varwidth}%
&\begin{varwidth}[t]{\sphinxcolwidth{1}{4}}
\sphinxAtStartPar
FIR/CIC
\sphinxbeforeendvarwidth
\end{varwidth}%
&\begin{varwidth}[t]{\sphinxcolwidth{1}{4}}
\sphinxAtStartPar
抗混叠滤波器类型
\sphinxbeforeendvarwidth
\end{varwidth}%
\\
\sphinxhline\begin{varwidth}[t]{\sphinxcolwidth{1}{4}}
\sphinxAtStartPar
BoundaryMethod
\sphinxbeforeendvarwidth
\end{varwidth}%
&\begin{varwidth}[t]{\sphinxcolwidth{1}{4}}
\sphinxAtStartPar
Enum
\sphinxbeforeendvarwidth
\end{varwidth}%
&\begin{varwidth}[t]{\sphinxcolwidth{1}{4}}
\sphinxAtStartPar
Zero/Mirror/Constant
\sphinxbeforeendvarwidth
\end{varwidth}%
&\begin{varwidth}[t]{\sphinxcolwidth{1}{4}}
\sphinxAtStartPar
边界处理方法
\sphinxbeforeendvarwidth
\end{varwidth}%
\\
\sphinxhline\begin{varwidth}[t]{\sphinxcolwidth{1}{4}}
\sphinxAtStartPar
PowerLineFilter
\sphinxbeforeendvarwidth
\end{varwidth}%
&\begin{varwidth}[t]{\sphinxcolwidth{1}{4}}
\sphinxAtStartPar
Boolean
\sphinxbeforeendvarwidth
\end{varwidth}%
&\begin{varwidth}[t]{\sphinxcolwidth{1}{4}}
\sphinxAtStartPar
True/False
\sphinxbeforeendvarwidth
\end{varwidth}%
&\begin{varwidth}[t]{\sphinxcolwidth{1}{4}}
\sphinxAtStartPar
是否启用工频滤波
\sphinxbeforeendvarwidth
\end{varwidth}%
\\
\sphinxhline\begin{varwidth}[t]{\sphinxcolwidth{1}{4}}
\sphinxAtStartPar
PowerLineFreq
\sphinxbeforeendvarwidth
\end{varwidth}%
&\begin{varwidth}[t]{\sphinxcolwidth{1}{4}}
\sphinxAtStartPar
Integer
\sphinxbeforeendvarwidth
\end{varwidth}%
&\begin{varwidth}[t]{\sphinxcolwidth{1}{4}}
\sphinxAtStartPar
50/60
\sphinxbeforeendvarwidth
\end{varwidth}%
&\begin{varwidth}[t]{\sphinxcolwidth{1}{4}}
\sphinxAtStartPar
工频频率
\sphinxbeforeendvarwidth
\end{varwidth}%
\\
\sphinxhline\begin{varwidth}[t]{\sphinxcolwidth{1}{4}}
\sphinxAtStartPar
PowerLineHarmonics
\sphinxbeforeendvarwidth
\end{varwidth}%
&\begin{varwidth}[t]{\sphinxcolwidth{1}{4}}
\sphinxAtStartPar
Integer
\sphinxbeforeendvarwidth
\end{varwidth}%
&\begin{varwidth}[t]{\sphinxcolwidth{1}{4}}
\sphinxAtStartPar
1\sphinxhyphen{}9
\sphinxbeforeendvarwidth
\end{varwidth}%
&\begin{varwidth}[t]{\sphinxcolwidth{1}{4}}
\sphinxAtStartPar
谐波数量
\sphinxbeforeendvarwidth
\end{varwidth}%
\\
\sphinxbottomrule
\end{tabulary}
\sphinxtableafterendhook\par
\sphinxattableend\end{savenotes}


\bigskip\hrule\bigskip



\subsection{工频滤波参数}
\label{\detokenize{appendices/appendixB:id9}}

\begin{savenotes}\sphinxattablestart
\sphinxthistablewithglobalstyle
\centering
\begin{tabulary}{\linewidth}[t]{TTTT}
\sphinxtoprule
\begin{varwidth}[t]{\sphinxcolwidth{1}{4}}
\sphinxstyletheadfamily \sphinxAtStartPar
参数名
\sphinxbeforeendvarwidth
\end{varwidth}%
&\begin{varwidth}[t]{\sphinxcolwidth{1}{4}}
\sphinxstyletheadfamily \sphinxAtStartPar
类型
\sphinxbeforeendvarwidth
\end{varwidth}%
&\begin{varwidth}[t]{\sphinxcolwidth{1}{4}}
\sphinxstyletheadfamily \sphinxAtStartPar
默认值
\sphinxbeforeendvarwidth
\end{varwidth}%
&\begin{varwidth}[t]{\sphinxcolwidth{1}{4}}
\sphinxstyletheadfamily \sphinxAtStartPar
说明
\sphinxbeforeendvarwidth
\end{varwidth}%
\\
\sphinxmidrule
\sphinxtableatstartofbodyhook\begin{varwidth}[t]{\sphinxcolwidth{1}{4}}
\sphinxAtStartPar
SampleRate
\sphinxbeforeendvarwidth
\end{varwidth}%
&\begin{varwidth}[t]{\sphinxcolwidth{1}{4}}
\sphinxAtStartPar
Double
\sphinxbeforeendvarwidth
\end{varwidth}%
&\begin{varwidth}[t]{\sphinxcolwidth{1}{4}}
\sphinxAtStartPar
\sphinxhyphen{}\sphinxhyphen{}
\sphinxbeforeendvarwidth
\end{varwidth}%
&\begin{varwidth}[t]{\sphinxcolwidth{1}{4}}
\sphinxAtStartPar
采样率（Hz）
\sphinxbeforeendvarwidth
\end{varwidth}%
\\
\sphinxhline\begin{varwidth}[t]{\sphinxcolwidth{1}{4}}
\sphinxAtStartPar
BaseFreq
\sphinxbeforeendvarwidth
\end{varwidth}%
&\begin{varwidth}[t]{\sphinxcolwidth{1}{4}}
\sphinxAtStartPar
Integer
\sphinxbeforeendvarwidth
\end{varwidth}%
&\begin{varwidth}[t]{\sphinxcolwidth{1}{4}}
\sphinxAtStartPar
50
\sphinxbeforeendvarwidth
\end{varwidth}%
&\begin{varwidth}[t]{\sphinxcolwidth{1}{4}}
\sphinxAtStartPar
基频（50或60Hz）
\sphinxbeforeendvarwidth
\end{varwidth}%
\\
\sphinxhline\begin{varwidth}[t]{\sphinxcolwidth{1}{4}}
\sphinxAtStartPar
HarmonicCount
\sphinxbeforeendvarwidth
\end{varwidth}%
&\begin{varwidth}[t]{\sphinxcolwidth{1}{4}}
\sphinxAtStartPar
Integer
\sphinxbeforeendvarwidth
\end{varwidth}%
&\begin{varwidth}[t]{\sphinxcolwidth{1}{4}}
\sphinxAtStartPar
9
\sphinxbeforeendvarwidth
\end{varwidth}%
&\begin{varwidth}[t]{\sphinxcolwidth{1}{4}}
\sphinxAtStartPar
谐波数量（含基频）
\sphinxbeforeendvarwidth
\end{varwidth}%
\\
\sphinxbottomrule
\end{tabulary}
\sphinxtableafterendhook\par
\sphinxattableend\end{savenotes}


\bigskip\hrule\bigskip



\subsection{Rhoplus参数}
\label{\detokenize{appendices/appendixB:rhoplus}}

\begin{savenotes}\sphinxattablestart
\sphinxthistablewithglobalstyle
\centering
\begin{tabulary}{\linewidth}[t]{TTTT}
\sphinxtoprule
\begin{varwidth}[t]{\sphinxcolwidth{1}{4}}
\sphinxstyletheadfamily \sphinxAtStartPar
参数名
\sphinxbeforeendvarwidth
\end{varwidth}%
&\begin{varwidth}[t]{\sphinxcolwidth{1}{4}}
\sphinxstyletheadfamily \sphinxAtStartPar
类型
\sphinxbeforeendvarwidth
\end{varwidth}%
&\begin{varwidth}[t]{\sphinxcolwidth{1}{4}}
\sphinxstyletheadfamily \sphinxAtStartPar
默认值
\sphinxbeforeendvarwidth
\end{varwidth}%
&\begin{varwidth}[t]{\sphinxcolwidth{1}{4}}
\sphinxstyletheadfamily \sphinxAtStartPar
说明
\sphinxbeforeendvarwidth
\end{varwidth}%
\\
\sphinxmidrule
\sphinxtableatstartofbodyhook\begin{varwidth}[t]{\sphinxcolwidth{1}{4}}
\sphinxAtStartPar
DefaultLayerCount
\sphinxbeforeendvarwidth
\end{varwidth}%
&\begin{varwidth}[t]{\sphinxcolwidth{1}{4}}
\sphinxAtStartPar
Integer
\sphinxbeforeendvarwidth
\end{varwidth}%
&\begin{varwidth}[t]{\sphinxcolwidth{1}{4}}
\sphinxAtStartPar
20
\sphinxbeforeendvarwidth
\end{varwidth}%
&\begin{varwidth}[t]{\sphinxcolwidth{1}{4}}
\sphinxAtStartPar
默认层数
\sphinxbeforeendvarwidth
\end{varwidth}%
\\
\sphinxhline\begin{varwidth}[t]{\sphinxcolwidth{1}{4}}
\sphinxAtStartPar
nfre
\sphinxbeforeendvarwidth
\end{varwidth}%
&\begin{varwidth}[t]{\sphinxcolwidth{1}{4}}
\sphinxAtStartPar
LongInt
\sphinxbeforeendvarwidth
\end{varwidth}%
&\begin{varwidth}[t]{\sphinxcolwidth{1}{4}}
\sphinxAtStartPar
\sphinxhyphen{}\sphinxhyphen{}
\sphinxbeforeendvarwidth
\end{varwidth}%
&\begin{varwidth}[t]{\sphinxcolwidth{1}{4}}
\sphinxAtStartPar
频率点数
\sphinxbeforeendvarwidth
\end{varwidth}%
\\
\sphinxhline\begin{varwidth}[t]{\sphinxcolwidth{1}{4}}
\sphinxAtStartPar
lcount
\sphinxbeforeendvarwidth
\end{varwidth}%
&\begin{varwidth}[t]{\sphinxcolwidth{1}{4}}
\sphinxAtStartPar
Integer
\sphinxbeforeendvarwidth
\end{varwidth}%
&\begin{varwidth}[t]{\sphinxcolwidth{1}{4}}
\sphinxAtStartPar
20
\sphinxbeforeendvarwidth
\end{varwidth}%
&\begin{varwidth}[t]{\sphinxcolwidth{1}{4}}
\sphinxAtStartPar
输出层数
\sphinxbeforeendvarwidth
\end{varwidth}%
\\
\sphinxbottomrule
\end{tabulary}
\sphinxtableafterendhook\par
\sphinxattableend\end{savenotes}


\bigskip\hrule\bigskip



\subsection{AI预测参数}
\label{\detokenize{appendices/appendixB:ai}}

\begin{savenotes}\sphinxattablestart
\sphinxthistablewithglobalstyle
\centering
\begin{tabulary}{\linewidth}[t]{TTT}
\sphinxtoprule
\begin{varwidth}[t]{\sphinxcolwidth{1}{3}}
\sphinxstyletheadfamily \sphinxAtStartPar
方法
\sphinxbeforeendvarwidth
\end{varwidth}%
&\begin{varwidth}[t]{\sphinxcolwidth{1}{3}}
\sphinxstyletheadfamily \sphinxAtStartPar
说明
\sphinxbeforeendvarwidth
\end{varwidth}%
&\begin{varwidth}[t]{\sphinxcolwidth{1}{3}}
\sphinxstyletheadfamily \sphinxAtStartPar
适用场景
\sphinxbeforeendvarwidth
\end{varwidth}%
\\
\sphinxmidrule
\sphinxtableatstartofbodyhook\begin{varwidth}[t]{\sphinxcolwidth{1}{3}}
\sphinxAtStartPar
ModelPredict
\sphinxbeforeendvarwidth
\end{varwidth}%
&\begin{varwidth}[t]{\sphinxcolwidth{1}{3}}
\sphinxAtStartPar
ZPredict深度学习
\sphinxbeforeendvarwidth
\end{varwidth}%
&\begin{varwidth}[t]{\sphinxcolwidth{1}{3}}
\sphinxAtStartPar
一般数据预测
\sphinxbeforeendvarwidth
\end{varwidth}%
\\
\sphinxhline\begin{varwidth}[t]{\sphinxcolwidth{1}{3}}
\sphinxAtStartPar
ModelPredict1
\sphinxbeforeendvarwidth
\end{varwidth}%
&\begin{varwidth}[t]{\sphinxcolwidth{1}{3}}
\sphinxAtStartPar
中值滤波（窗口3）
\sphinxbeforeendvarwidth
\end{varwidth}%
&\begin{varwidth}[t]{\sphinxcolwidth{1}{3}}
\sphinxAtStartPar
简单平滑
\sphinxbeforeendvarwidth
\end{varwidth}%
\\
\sphinxhline\begin{varwidth}[t]{\sphinxcolwidth{1}{3}}
\sphinxAtStartPar
ModelPredict2
\sphinxbeforeendvarwidth
\end{varwidth}%
&\begin{varwidth}[t]{\sphinxcolwidth{1}{3}}
\sphinxAtStartPar
迭代优化（最多100次）
\sphinxbeforeendvarwidth
\end{varwidth}%
&\begin{varwidth}[t]{\sphinxcolwidth{1}{3}}
\sphinxAtStartPar
精细预测
\sphinxbeforeendvarwidth
\end{varwidth}%
\\
\sphinxbottomrule
\end{tabulary}
\sphinxtableafterendhook\par
\sphinxattableend\end{savenotes}


\bigskip\hrule\bigskip



\subsection{测点排序方式}
\label{\detokenize{appendices/appendixB:id10}}

\begin{savenotes}\sphinxattablestart
\sphinxthistablewithglobalstyle
\centering
\begin{tabulary}{\linewidth}[t]{TTT}
\sphinxtoprule
\begin{varwidth}[t]{\sphinxcolwidth{1}{3}}
\sphinxstyletheadfamily \sphinxAtStartPar
值
\sphinxbeforeendvarwidth
\end{varwidth}%
&\begin{varwidth}[t]{\sphinxcolwidth{1}{3}}
\sphinxstyletheadfamily \sphinxAtStartPar
排序方式
\sphinxbeforeendvarwidth
\end{varwidth}%
&\begin{varwidth}[t]{\sphinxcolwidth{1}{3}}
\sphinxstyletheadfamily \sphinxAtStartPar
说明
\sphinxbeforeendvarwidth
\end{varwidth}%
\\
\sphinxmidrule
\sphinxtableatstartofbodyhook\begin{varwidth}[t]{\sphinxcolwidth{1}{3}}
\sphinxAtStartPar
0
\sphinxbeforeendvarwidth
\end{varwidth}%
&\begin{varwidth}[t]{\sphinxcolwidth{1}{3}}
\sphinxAtStartPar
按纬度
\sphinxbeforeendvarwidth
\end{varwidth}%
&\begin{varwidth}[t]{\sphinxcolwidth{1}{3}}
\sphinxAtStartPar
从南到北
\sphinxbeforeendvarwidth
\end{varwidth}%
\\
\sphinxhline\begin{varwidth}[t]{\sphinxcolwidth{1}{3}}
\sphinxAtStartPar
1
\sphinxbeforeendvarwidth
\end{varwidth}%
&\begin{varwidth}[t]{\sphinxcolwidth{1}{3}}
\sphinxAtStartPar
按经度
\sphinxbeforeendvarwidth
\end{varwidth}%
&\begin{varwidth}[t]{\sphinxcolwidth{1}{3}}
\sphinxAtStartPar
从西到东
\sphinxbeforeendvarwidth
\end{varwidth}%
\\
\sphinxhline\begin{varwidth}[t]{\sphinxcolwidth{1}{3}}
\sphinxAtStartPar
2
\sphinxbeforeendvarwidth
\end{varwidth}%
&\begin{varwidth}[t]{\sphinxcolwidth{1}{3}}
\sphinxAtStartPar
按名称
\sphinxbeforeendvarwidth
\end{varwidth}%
&\begin{varwidth}[t]{\sphinxcolwidth{1}{3}}
\sphinxAtStartPar
字母顺序
\sphinxbeforeendvarwidth
\end{varwidth}%
\\
\sphinxhline\begin{varwidth}[t]{\sphinxcolwidth{1}{3}}
\sphinxAtStartPar
3
\sphinxbeforeendvarwidth
\end{varwidth}%
&\begin{varwidth}[t]{\sphinxcolwidth{1}{3}}
\sphinxAtStartPar
按开始时间
\sphinxbeforeendvarwidth
\end{varwidth}%
&\begin{varwidth}[t]{\sphinxcolwidth{1}{3}}
\sphinxAtStartPar
时间先后
\sphinxbeforeendvarwidth
\end{varwidth}%
\\
\sphinxhline\begin{varwidth}[t]{\sphinxcolwidth{1}{3}}
\sphinxAtStartPar
4
\sphinxbeforeendvarwidth
\end{varwidth}%
&\begin{varwidth}[t]{\sphinxcolwidth{1}{3}}
\sphinxAtStartPar
按结束时间
\sphinxbeforeendvarwidth
\end{varwidth}%
&\begin{varwidth}[t]{\sphinxcolwidth{1}{3}}
\sphinxAtStartPar
时间先后
\sphinxbeforeendvarwidth
\end{varwidth}%
\\
\sphinxbottomrule
\end{tabulary}
\sphinxtableafterendhook\par
\sphinxattableend\end{savenotes}


\bigskip\hrule\bigskip



\subsection{仪器类型代码}
\label{\detokenize{appendices/appendixB:id11}}

\begin{savenotes}\sphinxattablestart
\sphinxthistablewithglobalstyle
\centering
\begin{tabulary}{\linewidth}[t]{TTT}
\sphinxtoprule
\begin{varwidth}[t]{\sphinxcolwidth{1}{3}}
\sphinxstyletheadfamily \sphinxAtStartPar
代码
\sphinxbeforeendvarwidth
\end{varwidth}%
&\begin{varwidth}[t]{\sphinxcolwidth{1}{3}}
\sphinxstyletheadfamily \sphinxAtStartPar
类型
\sphinxbeforeendvarwidth
\end{varwidth}%
&\begin{varwidth}[t]{\sphinxcolwidth{1}{3}}
\sphinxstyletheadfamily \sphinxAtStartPar
文件特征
\sphinxbeforeendvarwidth
\end{varwidth}%
\\
\sphinxmidrule
\sphinxtableatstartofbodyhook\begin{varwidth}[t]{\sphinxcolwidth{1}{3}}
\sphinxAtStartPar
0
\sphinxbeforeendvarwidth
\end{varwidth}%
&\begin{varwidth}[t]{\sphinxcolwidth{1}{3}}
\sphinxAtStartPar
Phoenix Site
\sphinxbeforeendvarwidth
\end{varwidth}%
&\begin{varwidth}[t]{\sphinxcolwidth{1}{3}}
\sphinxAtStartPar
.tbl
\sphinxbeforeendvarwidth
\end{varwidth}%
\\
\sphinxhline\begin{varwidth}[t]{\sphinxcolwidth{1}{3}}
\sphinxAtStartPar
1
\sphinxbeforeendvarwidth
\end{varwidth}%
&\begin{varwidth}[t]{\sphinxcolwidth{1}{3}}
\sphinxAtStartPar
LEMI Site
\sphinxbeforeendvarwidth
\end{varwidth}%
&\begin{varwidth}[t]{\sphinxcolwidth{1}{3}}
\sphinxAtStartPar
.lemi
\sphinxbeforeendvarwidth
\end{varwidth}%
\\
\sphinxhline\begin{varwidth}[t]{\sphinxcolwidth{1}{3}}
\sphinxAtStartPar
2
\sphinxbeforeendvarwidth
\end{varwidth}%
&\begin{varwidth}[t]{\sphinxcolwidth{1}{3}}
\sphinxAtStartPar
Metronix Site
\sphinxbeforeendvarwidth
\end{varwidth}%
&\begin{varwidth}[t]{\sphinxcolwidth{1}{3}}
\sphinxAtStartPar
.mxsite
\sphinxbeforeendvarwidth
\end{varwidth}%
\\
\sphinxhline\begin{varwidth}[t]{\sphinxcolwidth{1}{3}}
\sphinxAtStartPar
3
\sphinxbeforeendvarwidth
\end{varwidth}%
&\begin{varwidth}[t]{\sphinxcolwidth{1}{3}}
\sphinxAtStartPar
SBF Site
\sphinxbeforeendvarwidth
\end{varwidth}%
&\begin{varwidth}[t]{\sphinxcolwidth{1}{3}}
\sphinxAtStartPar
.sbf
\sphinxbeforeendvarwidth
\end{varwidth}%
\\
\sphinxhline\begin{varwidth}[t]{\sphinxcolwidth{1}{3}}
\sphinxAtStartPar
4
\sphinxbeforeendvarwidth
\end{varwidth}%
&\begin{varwidth}[t]{\sphinxcolwidth{1}{3}}
\sphinxAtStartPar
CASRMT Site
\sphinxbeforeendvarwidth
\end{varwidth}%
&\begin{varwidth}[t]{\sphinxcolwidth{1}{3}}
\sphinxAtStartPar
.rmtjson, .json
\sphinxbeforeendvarwidth
\end{varwidth}%
\\
\sphinxhline\begin{varwidth}[t]{\sphinxcolwidth{1}{3}}
\sphinxAtStartPar
5
\sphinxbeforeendvarwidth
\end{varwidth}%
&\begin{varwidth}[t]{\sphinxcolwidth{1}{3}}
\sphinxAtStartPar
MTU Site
\sphinxbeforeendvarwidth
\end{varwidth}%
&\begin{varwidth}[t]{\sphinxcolwidth{1}{3}}
\sphinxAtStartPar
recmeta.json
\sphinxbeforeendvarwidth
\end{varwidth}%
\\
\sphinxhline\begin{varwidth}[t]{\sphinxcolwidth{1}{3}}
\sphinxAtStartPar
6
\sphinxbeforeendvarwidth
\end{varwidth}%
&\begin{varwidth}[t]{\sphinxcolwidth{1}{3}}
\sphinxAtStartPar
RMT Site
\sphinxbeforeendvarwidth
\end{varwidth}%
&\begin{varwidth}[t]{\sphinxcolwidth{1}{3}}
\sphinxAtStartPar
.tr1, .tr2, .tr3
\sphinxbeforeendvarwidth
\end{varwidth}%
\\
\sphinxhline\begin{varwidth}[t]{\sphinxcolwidth{1}{3}}
\sphinxAtStartPar
7
\sphinxbeforeendvarwidth
\end{varwidth}%
&\begin{varwidth}[t]{\sphinxcolwidth{1}{3}}
\sphinxAtStartPar
ATTS Site
\sphinxbeforeendvarwidth
\end{varwidth}%
&\begin{varwidth}[t]{\sphinxcolwidth{1}{3}}
\sphinxAtStartPar
.atts, .atinfo
\sphinxbeforeendvarwidth
\end{varwidth}%
\\
\sphinxhline\begin{varwidth}[t]{\sphinxcolwidth{1}{3}}
\sphinxAtStartPar
8
\sphinxbeforeendvarwidth
\end{varwidth}%
&\begin{varwidth}[t]{\sphinxcolwidth{1}{3}}
\sphinxAtStartPar
AGEXXL Site
\sphinxbeforeendvarwidth
\end{varwidth}%
&\begin{varwidth}[t]{\sphinxcolwidth{1}{3}}
\sphinxAtStartPar
.dat
\sphinxbeforeendvarwidth
\end{varwidth}%
\\
\sphinxhline\begin{varwidth}[t]{\sphinxcolwidth{1}{3}}
\sphinxAtStartPar
\sphinxhyphen{}1
\sphinxbeforeendvarwidth
\end{varwidth}%
&\begin{varwidth}[t]{\sphinxcolwidth{1}{3}}
\sphinxAtStartPar
Synthetic Site
\sphinxbeforeendvarwidth
\end{varwidth}%
&\begin{varwidth}[t]{\sphinxcolwidth{1}{3}}
\sphinxAtStartPar
.timeseries
\sphinxbeforeendvarwidth
\end{varwidth}%
\\
\sphinxhline\begin{varwidth}[t]{\sphinxcolwidth{1}{3}}
\sphinxAtStartPar
100+
\sphinxbeforeendvarwidth
\end{varwidth}%
&\begin{varwidth}[t]{\sphinxcolwidth{1}{3}}
\sphinxAtStartPar
Source Site
\sphinxbeforeendvarwidth
\end{varwidth}%
&\begin{varwidth}[t]{\sphinxcolwidth{1}{3}}
\sphinxAtStartPar
源站点
\sphinxbeforeendvarwidth
\end{varwidth}%
\\
\sphinxhline\begin{varwidth}[t]{\sphinxcolwidth{1}{3}}
\sphinxAtStartPar
200+
\sphinxbeforeendvarwidth
\end{varwidth}%
&\begin{varwidth}[t]{\sphinxcolwidth{1}{3}}
\sphinxAtStartPar
Multiple Site
\sphinxbeforeendvarwidth
\end{varwidth}%
&\begin{varwidth}[t]{\sphinxcolwidth{1}{3}}
\sphinxAtStartPar
多文件合并站点
\sphinxbeforeendvarwidth
\end{varwidth}%
\\
\sphinxbottomrule
\end{tabulary}
\sphinxtableafterendhook\par
\sphinxattableend\end{savenotes}


\bigskip\hrule\bigskip



\subsection{数据质量等级}
\label{\detokenize{appendices/appendixB:id12}}

\begin{savenotes}\sphinxattablestart
\sphinxthistablewithglobalstyle
\centering
\begin{tabulary}{\linewidth}[t]{TTT}
\sphinxtoprule
\begin{varwidth}[t]{\sphinxcolwidth{1}{3}}
\sphinxstyletheadfamily \sphinxAtStartPar
等级
\sphinxbeforeendvarwidth
\end{varwidth}%
&\begin{varwidth}[t]{\sphinxcolwidth{1}{3}}
\sphinxstyletheadfamily \sphinxAtStartPar
说明
\sphinxbeforeendvarwidth
\end{varwidth}%
&\begin{varwidth}[t]{\sphinxcolwidth{1}{3}}
\sphinxstyletheadfamily \sphinxAtStartPar
建议处理
\sphinxbeforeendvarwidth
\end{varwidth}%
\\
\sphinxmidrule
\sphinxtableatstartofbodyhook\begin{varwidth}[t]{\sphinxcolwidth{1}{3}}
\sphinxAtStartPar
0
\sphinxbeforeendvarwidth
\end{varwidth}%
&\begin{varwidth}[t]{\sphinxcolwidth{1}{3}}
\sphinxAtStartPar
未评估
\sphinxbeforeendvarwidth
\end{varwidth}%
&\begin{varwidth}[t]{\sphinxcolwidth{1}{3}}
\sphinxAtStartPar
需要评估
\sphinxbeforeendvarwidth
\end{varwidth}%
\\
\sphinxhline\begin{varwidth}[t]{\sphinxcolwidth{1}{3}}
\sphinxAtStartPar
1
\sphinxbeforeendvarwidth
\end{varwidth}%
&\begin{varwidth}[t]{\sphinxcolwidth{1}{3}}
\sphinxAtStartPar
优秀
\sphinxbeforeendvarwidth
\end{varwidth}%
&\begin{varwidth}[t]{\sphinxcolwidth{1}{3}}
\sphinxAtStartPar
可直接使用
\sphinxbeforeendvarwidth
\end{varwidth}%
\\
\sphinxhline\begin{varwidth}[t]{\sphinxcolwidth{1}{3}}
\sphinxAtStartPar
2
\sphinxbeforeendvarwidth
\end{varwidth}%
&\begin{varwidth}[t]{\sphinxcolwidth{1}{3}}
\sphinxAtStartPar
良好
\sphinxbeforeendvarwidth
\end{varwidth}%
&\begin{varwidth}[t]{\sphinxcolwidth{1}{3}}
\sphinxAtStartPar
轻度筛选
\sphinxbeforeendvarwidth
\end{varwidth}%
\\
\sphinxhline\begin{varwidth}[t]{\sphinxcolwidth{1}{3}}
\sphinxAtStartPar
3
\sphinxbeforeendvarwidth
\end{varwidth}%
&\begin{varwidth}[t]{\sphinxcolwidth{1}{3}}
\sphinxAtStartPar
一般
\sphinxbeforeendvarwidth
\end{varwidth}%
&\begin{varwidth}[t]{\sphinxcolwidth{1}{3}}
\sphinxAtStartPar
需要筛选
\sphinxbeforeendvarwidth
\end{varwidth}%
\\
\sphinxhline\begin{varwidth}[t]{\sphinxcolwidth{1}{3}}
\sphinxAtStartPar
4
\sphinxbeforeendvarwidth
\end{varwidth}%
&\begin{varwidth}[t]{\sphinxcolwidth{1}{3}}
\sphinxAtStartPar
较差
\sphinxbeforeendvarwidth
\end{varwidth}%
&\begin{varwidth}[t]{\sphinxcolwidth{1}{3}}
\sphinxAtStartPar
建议重采或舍弃
\sphinxbeforeendvarwidth
\end{varwidth}%
\\
\sphinxbottomrule
\end{tabulary}
\sphinxtableafterendhook\par
\sphinxattableend\end{savenotes}


\bigskip\hrule\bigskip



\subsection{MT参数类型}
\label{\detokenize{appendices/appendixB:mt}}

\begin{savenotes}\sphinxattablestart
\sphinxthistablewithglobalstyle
\centering
\begin{tabulary}{\linewidth}[t]{TT}
\sphinxtoprule
\begin{varwidth}[t]{\sphinxcolwidth{1}{2}}
\sphinxstyletheadfamily \sphinxAtStartPar
参数
\sphinxbeforeendvarwidth
\end{varwidth}%
&\begin{varwidth}[t]{\sphinxcolwidth{1}{2}}
\sphinxstyletheadfamily \sphinxAtStartPar
说明
\sphinxbeforeendvarwidth
\end{varwidth}%
\\
\sphinxmidrule
\sphinxtableatstartofbodyhook\begin{varwidth}[t]{\sphinxcolwidth{1}{2}}
\sphinxAtStartPar
MT\_fre
\sphinxbeforeendvarwidth
\end{varwidth}%
&\begin{varwidth}[t]{\sphinxcolwidth{1}{2}}
\sphinxAtStartPar
频率
\sphinxbeforeendvarwidth
\end{varwidth}%
\\
\sphinxhline\begin{varwidth}[t]{\sphinxcolwidth{1}{2}}
\sphinxAtStartPar
MT\_Zxxr/Zxxi
\sphinxbeforeendvarwidth
\end{varwidth}%
&\begin{varwidth}[t]{\sphinxcolwidth{1}{2}}
\sphinxAtStartPar
Zxx实部/虚部
\sphinxbeforeendvarwidth
\end{varwidth}%
\\
\sphinxhline\begin{varwidth}[t]{\sphinxcolwidth{1}{2}}
\sphinxAtStartPar
MT\_Zxyr/Zxyi
\sphinxbeforeendvarwidth
\end{varwidth}%
&\begin{varwidth}[t]{\sphinxcolwidth{1}{2}}
\sphinxAtStartPar
Zxy实部/虚部
\sphinxbeforeendvarwidth
\end{varwidth}%
\\
\sphinxhline\begin{varwidth}[t]{\sphinxcolwidth{1}{2}}
\sphinxAtStartPar
MT\_Zyxr/Zyxi
\sphinxbeforeendvarwidth
\end{varwidth}%
&\begin{varwidth}[t]{\sphinxcolwidth{1}{2}}
\sphinxAtStartPar
Zyx实部/虚部
\sphinxbeforeendvarwidth
\end{varwidth}%
\\
\sphinxhline\begin{varwidth}[t]{\sphinxcolwidth{1}{2}}
\sphinxAtStartPar
MT\_Zyyr/Zyyi
\sphinxbeforeendvarwidth
\end{varwidth}%
&\begin{varwidth}[t]{\sphinxcolwidth{1}{2}}
\sphinxAtStartPar
Zyy实部/虚部
\sphinxbeforeendvarwidth
\end{varwidth}%
\\
\sphinxhline\begin{varwidth}[t]{\sphinxcolwidth{1}{2}}
\sphinxAtStartPar
MT\_Tzxr/Tzxi
\sphinxbeforeendvarwidth
\end{varwidth}%
&\begin{varwidth}[t]{\sphinxcolwidth{1}{2}}
\sphinxAtStartPar
Tzx实部/虚部
\sphinxbeforeendvarwidth
\end{varwidth}%
\\
\sphinxhline\begin{varwidth}[t]{\sphinxcolwidth{1}{2}}
\sphinxAtStartPar
MT\_Tzyr/Tzyi
\sphinxbeforeendvarwidth
\end{varwidth}%
&\begin{varwidth}[t]{\sphinxcolwidth{1}{2}}
\sphinxAtStartPar
Tzy实部/虚部
\sphinxbeforeendvarwidth
\end{varwidth}%
\\
\sphinxhline\begin{varwidth}[t]{\sphinxcolwidth{1}{2}}
\sphinxAtStartPar
MT\_Rxx/Rxy/Ryx/Ryy
\sphinxbeforeendvarwidth
\end{varwidth}%
&\begin{varwidth}[t]{\sphinxcolwidth{1}{2}}
\sphinxAtStartPar
视电阻率
\sphinxbeforeendvarwidth
\end{varwidth}%
\\
\sphinxhline\begin{varwidth}[t]{\sphinxcolwidth{1}{2}}
\sphinxAtStartPar
MT\_Pxx/Pxy/Pyx/Pyy
\sphinxbeforeendvarwidth
\end{varwidth}%
&\begin{varwidth}[t]{\sphinxcolwidth{1}{2}}
\sphinxAtStartPar
相位
\sphinxbeforeendvarwidth
\end{varwidth}%
\\
\sphinxhline\begin{varwidth}[t]{\sphinxcolwidth{1}{2}}
\sphinxAtStartPar
MT\_Coh*
\sphinxbeforeendvarwidth
\end{varwidth}%
&\begin{varwidth}[t]{\sphinxcolwidth{1}{2}}
\sphinxAtStartPar
相干度
\sphinxbeforeendvarwidth
\end{varwidth}%
\\
\sphinxhline\begin{varwidth}[t]{\sphinxcolwidth{1}{2}}
\sphinxAtStartPar
MT\_SNR*
\sphinxbeforeendvarwidth
\end{varwidth}%
&\begin{varwidth}[t]{\sphinxcolwidth{1}{2}}
\sphinxAtStartPar
信噪比
\sphinxbeforeendvarwidth
\end{varwidth}%
\\
\sphinxhline\begin{varwidth}[t]{\sphinxcolwidth{1}{2}}
\sphinxAtStartPar
MT\_alpha/beta
\sphinxbeforeendvarwidth
\end{varwidth}%
&\begin{varwidth}[t]{\sphinxcolwidth{1}{2}}
\sphinxAtStartPar
相位张量参数
\sphinxbeforeendvarwidth
\end{varwidth}%
\\
\sphinxhline\begin{varwidth}[t]{\sphinxcolwidth{1}{2}}
\sphinxAtStartPar
MT\_PtMax/PtMin
\sphinxbeforeendvarwidth
\end{varwidth}%
&\begin{varwidth}[t]{\sphinxcolwidth{1}{2}}
\sphinxAtStartPar
相位张量最大/最小值
\sphinxbeforeendvarwidth
\end{varwidth}%
\\
\sphinxhline\begin{varwidth}[t]{\sphinxcolwidth{1}{2}}
\sphinxAtStartPar
MT\_Skew1D/Skew2D
\sphinxbeforeendvarwidth
\end{varwidth}%
&\begin{varwidth}[t]{\sphinxcolwidth{1}{2}}
\sphinxAtStartPar
相位张量偏斜度
\sphinxbeforeendvarwidth
\end{varwidth}%
\\
\sphinxbottomrule
\end{tabulary}
\sphinxtableafterendhook\par
\sphinxattableend\end{savenotes}

\sphinxstepscope


\section{附录C: 文件格式规范}
\label{\detokenize{appendices/appendixC:c}}\label{\detokenize{appendices/appendixC::doc}}

\subsection{工程文件结构}
\label{\detokenize{appendices/appendixC:id1}}
\sphinxAtStartPar
MTDP工程使用加密格式存储，扩展名为 \sphinxcode{\sphinxupquote{.MTDPE}}。


\subsubsection{工程目录结构}
\label{\detokenize{appendices/appendixC:id2}}
\begin{sphinxVerbatim}[commandchars=\\\{\}]
ProjectName/
├── ProjectName.MTDPE        \PYGZsh{} 工程主文件（加密）
├── Configurations/          \PYGZsh{} 配置文件
│   ├── MTDP.SETTING        \PYGZsh{} 软件设置
│   ├── ProjectInfos.XML    \PYGZsh{} 工程信息列表
│   └── Lang/               \PYGZsh{} 语言文件
├── CLB/                     \PYGZsh{} 电场盒标定
├── CLC/                     \PYGZsh{} 磁传感器标定
└── temp/                    \PYGZsh{} 临时文件
\end{sphinxVerbatim}


\bigskip\hrule\bigskip



\subsection{EDI文件格式}
\label{\detokenize{appendices/appendixC:edi}}
\sphinxAtStartPar
EDI（Electrical Data Interchange）是MT数据的标准交换格式。


\subsubsection{EDI文件结构示例}
\label{\detokenize{appendices/appendixC:id3}}
\begin{sphinxVerbatim}[commandchars=\\\{\}]
\PYG{o}{\PYGZgt{}}\PYG{n}{HEAD}
  \PYG{n}{SITE} \PYG{n}{NAME}\PYG{o}{=}\PYG{l+s+s2}{\PYGZdq{}}\PYG{l+s+s2}{MT001}\PYG{l+s+s2}{\PYGZdq{}}
  \PYG{n}{LATITUDE}\PYG{o}{=}\PYG{l+m+mi}{40}\PYG{p}{:}\PYG{l+m+mi}{07}\PYG{p}{:}\PYG{l+m+mf}{24.2} \PYG{n}{N}
  \PYG{n}{LONGITUDE}\PYG{o}{=}\PYG{l+m+mi}{116}\PYG{p}{:}\PYG{l+m+mi}{34}\PYG{p}{:}\PYG{l+m+mf}{04.1} \PYG{n}{E}
  \PYG{n}{ELEVATION}\PYG{o}{=}\PYG{l+m+mf}{150.5}
\PYG{o}{\PYGZgt{}}\PYG{n}{END}

\PYG{o}{\PYGZgt{}}\PYG{n}{FREQ} \PYG{n}{UNITS}\PYG{o}{=}\PYG{n}{HZ}
  \PYG{l+m+mf}{0.001} \PYG{l+m+mf}{0.002} \PYG{l+m+mf}{0.005} \PYG{l+m+mf}{0.01} \PYG{l+m+mf}{0.02} \PYG{l+m+mf}{0.05} \PYG{l+m+mf}{0.1} \PYG{l+m+mf}{0.2} \PYG{l+m+mf}{0.5} \PYG{l+m+mf}{1.0} \PYG{l+m+mf}{2.0} \PYG{l+m+mf}{5.0} \PYG{l+m+mf}{10.0}
\PYG{o}{\PYGZgt{}}\PYG{n}{END}

\PYG{o}{\PYGZgt{}}\PYG{n}{ZXX} \PYG{n}{UNITS}\PYG{o}{=}\PYG{n}{MV}\PYG{o}{/}\PYG{n}{KM}\PYG{o}{/}\PYG{n}{NT}
  \PYG{o}{/}\PYG{o}{/} \PYG{o}{|}\PYG{n}{Zxx}\PYG{o}{|}\PYG{p}{,} \PYG{n}{arg}\PYG{p}{(}\PYG{n}{Zxx}\PYG{p}{)}\PYG{p}{,} \PYG{o}{|}\PYG{n}{Zxx}\PYG{o}{|}\PYG{n}{\PYGZus{}err}\PYG{p}{,} \PYG{n}{arg}\PYG{p}{(}\PYG{n}{Zxx}\PYG{p}{)}\PYG{n}{\PYGZus{}err}
  \PYG{l+m+mf}{0.00123}\PYG{p}{,} \PYG{l+m+mf}{45.2}\PYG{p}{,} \PYG{l+m+mf}{0.00012}\PYG{p}{,} \PYG{l+m+mf}{2.5}
  \PYG{l+m+mf}{0.00234}\PYG{p}{,} \PYG{l+m+mf}{43.8}\PYG{p}{,} \PYG{l+m+mf}{0.00023}\PYG{p}{,} \PYG{l+m+mf}{3.1}
  \PYG{o}{.}\PYG{o}{.}\PYG{o}{.}
\PYG{o}{\PYGZgt{}}\PYG{n}{END}

\PYG{o}{\PYGZgt{}}\PYG{n}{ZXY} \PYG{n}{UNITS}\PYG{o}{=}\PYG{n}{MV}\PYG{o}{/}\PYG{n}{KM}\PYG{o}{/}\PYG{n}{NT}
  \PYG{o}{/}\PYG{o}{/} \PYG{o}{|}\PYG{n}{Zxy}\PYG{o}{|}\PYG{p}{,} \PYG{n}{arg}\PYG{p}{(}\PYG{n}{Zxy}\PYG{p}{)}\PYG{p}{,} \PYG{o}{|}\PYG{n}{Zxy}\PYG{o}{|}\PYG{n}{\PYGZus{}err}\PYG{p}{,} \PYG{n}{arg}\PYG{p}{(}\PYG{n}{Zxy}\PYG{p}{)}\PYG{n}{\PYGZus{}err}
  \PYG{o}{.}\PYG{o}{.}\PYG{o}{.}
\PYG{o}{\PYGZgt{}}\PYG{n}{END}
\end{sphinxVerbatim}


\bigskip\hrule\bigskip



\subsection{各仪器原始数据格式摘要}
\label{\detokenize{appendices/appendixC:id4}}

\subsubsection{Phoenix TBL格式}
\label{\detokenize{appendices/appendixC:phoenix-tbl}}
\begin{sphinxVerbatim}[commandchars=\\\{\}]
\PYG{n}{SITE\PYGZus{}NAME}\PYG{p}{:} \PYG{n}{MT001}
\PYG{n}{LATITUDE}\PYG{p}{:} \PYG{l+m+mf}{40.1234}
\PYG{n}{LONGITUDE}\PYG{p}{:} \PYG{l+m+mf}{116.5678}
\PYG{n}{ELEVATION}\PYG{p}{:} \PYG{l+m+mf}{150.5}
\PYG{n}{SAMPLING\PYGZus{}RATE}\PYG{p}{:} \PYG{l+m+mi}{256}
\PYG{n}{NUM\PYGZus{}CHANNELS}\PYG{p}{:} \PYG{l+m+mi}{5}
\PYG{n}{START\PYGZus{}TIME}\PYG{p}{:} \PYG{l+m+mi}{2024}\PYG{o}{\PYGZhy{}}\PYG{l+m+mi}{01}\PYG{o}{\PYGZhy{}}\PYG{l+m+mi}{15} \PYG{l+m+mi}{08}\PYG{p}{:}\PYG{l+m+mi}{00}\PYG{p}{:}\PYG{l+m+mi}{00}
\PYG{n}{END\PYGZus{}TIME}\PYG{p}{:} \PYG{l+m+mi}{2024}\PYG{o}{\PYGZhy{}}\PYG{l+m+mi}{01}\PYG{o}{\PYGZhy{}}\PYG{l+m+mi}{16} \PYG{l+m+mi}{08}\PYG{p}{:}\PYG{l+m+mi}{00}\PYG{p}{:}\PYG{l+m+mi}{00}
\end{sphinxVerbatim}


\subsubsection{Phoenix JSON格式}
\label{\detokenize{appendices/appendixC:phoenix-json}}
\begin{sphinxVerbatim}[commandchars=\\\{\}]
\PYG{p}{\PYGZob{}}
\PYG{+w}{  }\PYG{n+nt}{\PYGZdq{}site\PYGZdq{}}\PYG{p}{:}\PYG{+w}{ }\PYG{p}{\PYGZob{}}
\PYG{+w}{    }\PYG{n+nt}{\PYGZdq{}name\PYGZdq{}}\PYG{p}{:}\PYG{+w}{ }\PYG{l+s+s2}{\PYGZdq{}MT001\PYGZdq{}}\PYG{p}{,}
\PYG{+w}{    }\PYG{n+nt}{\PYGZdq{}coordinates\PYGZdq{}}\PYG{p}{:}\PYG{+w}{ }\PYG{p}{\PYGZob{}}
\PYG{+w}{      }\PYG{n+nt}{\PYGZdq{}latitude\PYGZdq{}}\PYG{p}{:}\PYG{+w}{ }\PYG{l+m+mf}{40.1234}\PYG{p}{,}
\PYG{+w}{      }\PYG{n+nt}{\PYGZdq{}longitude\PYGZdq{}}\PYG{p}{:}\PYG{+w}{ }\PYG{l+m+mf}{116.5678}\PYG{p}{,}
\PYG{+w}{      }\PYG{n+nt}{\PYGZdq{}elevation\PYGZdq{}}\PYG{p}{:}\PYG{+w}{ }\PYG{l+m+mf}{150.5}
\PYG{+w}{    }\PYG{p}{\PYGZcb{}}
\PYG{+w}{  }\PYG{p}{\PYGZcb{},}
\PYG{+w}{  }\PYG{n+nt}{\PYGZdq{}acquisition\PYGZdq{}}\PYG{p}{:}\PYG{+w}{ }\PYG{p}{\PYGZob{}}
\PYG{+w}{    }\PYG{n+nt}{\PYGZdq{}sampling\PYGZus{}rate\PYGZdq{}}\PYG{p}{:}\PYG{+w}{ }\PYG{l+m+mi}{256}\PYG{p}{,}
\PYG{+w}{    }\PYG{n+nt}{\PYGZdq{}start\PYGZus{}time\PYGZdq{}}\PYG{p}{:}\PYG{+w}{ }\PYG{l+s+s2}{\PYGZdq{}2024\PYGZhy{}01\PYGZhy{}15T08:00:00Z\PYGZdq{}}\PYG{p}{,}
\PYG{+w}{    }\PYG{n+nt}{\PYGZdq{}duration\PYGZus{}hours\PYGZdq{}}\PYG{p}{:}\PYG{+w}{ }\PYG{l+m+mi}{24}
\PYG{+w}{  }\PYG{p}{\PYGZcb{}}
\PYG{p}{\PYGZcb{}}
\end{sphinxVerbatim}


\subsubsection{Metronix ATM格式}
\label{\detokenize{appendices/appendixC:metronix-atm}}
\sphinxAtStartPar
ATM文件为二进制格式，包含：
\begin{itemize}
\item {} 
\sphinxAtStartPar
文件头：标识信息、版本号

\item {} 
\sphinxAtStartPar
通道配置：通道数、采样率

\item {} 
\sphinxAtStartPar
数据块：时间戳、采样数据

\end{itemize}


\subsubsection{LEMI格式}
\label{\detokenize{appendices/appendixC:lemi}}\begin{itemize}
\item {} 
\sphinxAtStartPar
二进制时间序列

\item {} 
\sphinxAtStartPar
文本头文件

\item {} 
\sphinxAtStartPar
支持多种采样率

\end{itemize}


\subsubsection{ATTS格式}
\label{\detokenize{appendices/appendixC:atts}}\begin{itemize}
\item {} 
\sphinxAtStartPar
自定义二进制格式

\item {} 
\sphinxAtStartPar
高精度时间戳

\item {} 
\sphinxAtStartPar
多通道同步采集

\end{itemize}


\bigskip\hrule\bigskip



\subsection{标定文件格式}
\label{\detokenize{appendices/appendixC:id5}}

\subsubsection{CLB文件（电场盒标定）}
\label{\detokenize{appendices/appendixC:clb}}
\sphinxAtStartPar
记录电场采集盒的频率响应：
\begin{itemize}
\item {} 
\sphinxAtStartPar
频率点

\item {} 
\sphinxAtStartPar
幅度响应

\item {} 
\sphinxAtStartPar
相位响应

\end{itemize}


\subsubsection{CLC文件（磁传感器标定）}
\label{\detokenize{appendices/appendixC:clc}}
\sphinxAtStartPar
记录磁传感器的频率响应：
\begin{itemize}
\item {} 
\sphinxAtStartPar
传感器序列号

\item {} 
\sphinxAtStartPar
频率点

\item {} 
\sphinxAtStartPar
复数响应值

\end{itemize}


\bigskip\hrule\bigskip



\subsection{处理方案文件}
\label{\detokenize{appendices/appendixC:id6}}
\sphinxAtStartPar
处理方案可导出为XML或JSON格式，包含：
\begin{itemize}
\item {} 
\sphinxAtStartPar
窗函数设置

\item {} 
\sphinxAtStartPar
FFT参数

\item {} 
\sphinxAtStartPar
频率表

\item {} 
\sphinxAtStartPar
估计方法配置

\item {} 
\sphinxAtStartPar
自动筛选参数

\end{itemize}

\sphinxstepscope


\section{附录D: 常见问题FAQ}
\label{\detokenize{appendices/appendixD:d-faq}}\label{\detokenize{appendices/appendixD::doc}}

\subsection{安装与授权问题}
\label{\detokenize{appendices/appendixD:id1}}
\sphinxAtStartPar
\sphinxstylestrong{Q1: 软件无法启动，提示缺少DLL文件}

\sphinxAtStartPar
A: 请检查是否正确安装了所有依赖库。尝试重新安装软件，确保安装目录完整。

\sphinxAtStartPar
\sphinxstylestrong{Q2: 授权验证失败}

\sphinxAtStartPar
A: 请确认授权码输入正确，联系技术支持获取帮助。

\sphinxAtStartPar
\sphinxstylestrong{Q3: 软件运行缓慢}

\sphinxAtStartPar
A: 检查系统资源占用，关闭其他占用内存的程序。对于大规模数据处理，建议增加内存或使用更高配置的计算机。


\bigskip\hrule\bigskip



\subsection{数据导入问题}
\label{\detokenize{appendices/appendixD:id2}}
\sphinxAtStartPar
\sphinxstylestrong{Q4: 导入Phoenix数据时提示"文件格式不识别"}

\sphinxAtStartPar
A: 请确认数据格式与选择的导入类型匹配（TBL或JSON）。检查数据文件是否完整，必要时重新从仪器导出数据。

\sphinxAtStartPar
\sphinxstylestrong{Q5: 导入后时间序列显示异常}

\sphinxAtStartPar
A: 检查采样率设置是否正确，确认通道映射是否匹配。可能需要手动调整导入参数。

\sphinxAtStartPar
\sphinxstylestrong{Q6: 批量导入时部分测点失败}

\sphinxAtStartPar
A: 查看错误日志了解具体原因。常见原因包括：文件损坏、格式不一致、磁盘空间不足。


\bigskip\hrule\bigskip



\subsection{处理结果异常}
\label{\detokenize{appendices/appendixD:id3}}
\sphinxAtStartPar
\sphinxstylestrong{Q7: FFT计算后数据点很少}

\sphinxAtStartPar
A: 检查FFT长度设置是否过大（相对于数据长度）。尝试减小FFT长度或增加重叠率。

\sphinxAtStartPar
\sphinxstylestrong{Q8: 视电阻率曲线严重离散}

\sphinxAtStartPar
A: 可能是数据质量问题。尝试：
\begin{enumerate}
\sphinxsetlistlabels{\arabic}{enumi}{enumii}{}{.}%
\item {} 
\sphinxAtStartPar
应用Robust估计

\item {} 
\sphinxAtStartPar
配置参考道

\item {} 
\sphinxAtStartPar
检查是否有强噪声干扰

\item {} 
\sphinxAtStartPar
使用数据筛选功能

\end{enumerate}

\sphinxAtStartPar
\sphinxstylestrong{Q9: 相位曲线超出合理范围}

\sphinxAtStartPar
A: 检查通道极性是否正确，确认标定是否正确应用。可能存在系统响应校正问题。

\sphinxAtStartPar
\sphinxstylestrong{Q10: 智能筛选运行时间过长}

\sphinxAtStartPar
A: 遗传算法计算量较大。可尝试减少迭代次数、减小种群规模，或先用手动筛选预处理。


\bigskip\hrule\bigskip



\subsection{导出问题}
\label{\detokenize{appendices/appendixD:id4}}
\sphinxAtStartPar
\sphinxstylestrong{Q11: 导出的EDI文件无法被其他软件读取}

\sphinxAtStartPar
A: 检查EDI格式是否符合标准。某些软件可能不支持所有EDI变体，尝试使用更通用的格式设置。

\sphinxAtStartPar
\sphinxstylestrong{Q12: 批量导出时文件名重复}

\sphinxAtStartPar
A: 检查导出设置中的文件命名规则，确保包含测点名或其他唯一标识。

\sphinxAtStartPar
\sphinxstylestrong{Q13: KML文件在Google Earth中位置不对}

\sphinxAtStartPar
A: 检查测点的经纬度坐标是否正确设置，确认坐标单位（度/度分秒）是否匹配。


\bigskip\hrule\bigskip



\subsection{其他问题}
\label{\detokenize{appendices/appendixD:id5}}
\sphinxAtStartPar
\sphinxstylestrong{Q14: 如何恢复误删除的数据？}

\sphinxAtStartPar
A: 使用撤销功能（Ctrl+Z）或从历史记录恢复。如果已保存工程，尝试从备份恢复。

\sphinxAtStartPar
\sphinxstylestrong{Q15: 处理参数可以保存吗？}

\sphinxAtStartPar
A: 可以使用"保存处理模板"功能保存常用参数设置，下次使用时快速加载。

\sphinxAtStartPar
\sphinxstylestrong{Q16: 软件支持哪些操作系统？}

\sphinxAtStartPar
A: 目前仅支持Windows 10/11 64位系统。

\sphinxAtStartPar
\sphinxstylestrong{Q17: 如何获取技术支持？}

\sphinxAtStartPar
A: 请联系MTDP开发团队，提供详细的问题描述和必要的数据文件。

\sphinxstepscope


\section{附录E: 更新日志}
\label{\detokenize{appendices/appendixE:e}}\label{\detokenize{appendices/appendixE::doc}}

\subsection{版本 1.9.4 (当前版本)}
\label{\detokenize{appendices/appendixE:id1}}
\sphinxAtStartPar
\sphinxstylestrong{发布日期：} 2026年2月2日

\sphinxAtStartPar
\sphinxstylestrong{主要功能：}


\subsubsection{多仪器支持}
\label{\detokenize{appendices/appendixE:id2}}\begin{itemize}
\item {} 
\sphinxAtStartPar
Phoenix MTU\sphinxhyphen{}5/5A TBL格式

\item {} 
\sphinxAtStartPar
Phoenix MTU\sphinxhyphen{}5C/8A JSON格式

\item {} 
\sphinxAtStartPar
Metronix ADU\sphinxhyphen{}06/ADU\sphinxhyphen{}07/MMS系列

\item {} 
\sphinxAtStartPar
LEMI长周期仪器

\item {} 
\sphinxAtStartPar
RMT/CASRMT格式

\item {} 
\sphinxAtStartPar
Aether (ATTS) 格式

\item {} 
\sphinxAtStartPar
EDI格式导入

\end{itemize}


\subsubsection{数据处理}
\label{\detokenize{appendices/appendixE:id3}}\begin{itemize}
\item {} 
\sphinxAtStartPar
FFT/DFT频谱计算

\item {} 
\sphinxAtStartPar
多锥谱分析（Multi\sphinxhyphen{}Taper）

\item {} 
\sphinxAtStartPar
系统响应标定

\item {} 
\sphinxAtStartPar
工频陷波滤波

\item {} 
\sphinxAtStartPar
高通/低通滤波

\item {} 
\sphinxAtStartPar
降采样处理（2x\sphinxhyphen{}4800x）

\end{itemize}


\subsubsection{稳健估计}
\label{\detokenize{appendices/appendixE:id4}}\begin{itemize}
\item {} 
\sphinxAtStartPar
Regression\sphinxhyphen{}M估计

\item {} 
\sphinxAtStartPar
Repeated Median估计

\item {} 
\sphinxAtStartPar
Bounded Influence估计

\item {} 
\sphinxAtStartPar
Huber/Thomson预加权

\item {} 
\sphinxAtStartPar
Robust+AI估计

\end{itemize}


\subsubsection{参考道技术}
\label{\detokenize{appendices/appendixE:id5}}\begin{itemize}
\item {} 
\sphinxAtStartPar
远参考道处理

\item {} 
\sphinxAtStartPar
本地参考道

\item {} 
\sphinxAtStartPar
反向参考估计

\end{itemize}


\subsubsection{智能筛选}
\label{\detokenize{appendices/appendixE:id6}}\begin{itemize}
\item {} 
\sphinxAtStartPar
多参数自动筛选

\item {} 
\sphinxAtStartPar
马氏距离筛选

\item {} 
\sphinxAtStartPar
Rhoplus参考曲线

\item {} 
\sphinxAtStartPar
遗传算法筛选

\end{itemize}


\subsubsection{导出功能}
\label{\detokenize{appendices/appendixE:id7}}\begin{itemize}
\item {} 
\sphinxAtStartPar
EDI格式（标准/Spe/ZT/MTpy/PLT）

\item {} 
\sphinxAtStartPar
KML/KMZ地理信息

\item {} 
\sphinxAtStartPar
时间序列导出（GMT/ATGMT格式）

\end{itemize}


\bigskip\hrule\bigskip



\subsection{版本 1.5}
\label{\detokenize{appendices/appendixE:id8}}
\sphinxAtStartPar
\sphinxstylestrong{主要功能：}
\begin{itemize}
\item {} 
\sphinxAtStartPar
Phoenix MTU\sphinxhyphen{}5/5A TBL格式支持

\item {} 
\sphinxAtStartPar
Metronix、LEMI基本支持

\item {} 
\sphinxAtStartPar
FFT计算

\item {} 
\sphinxAtStartPar
手动数据筛选

\item {} 
\sphinxAtStartPar
Robust估计

\item {} 
\sphinxAtStartPar
EDI导出

\end{itemize}


\bigskip\hrule\bigskip



\subsection{版本 1.0}
\label{\detokenize{appendices/appendixE:id9}}
\sphinxAtStartPar
\sphinxstylestrong{初始版本功能：}
\begin{itemize}
\item {} 
\sphinxAtStartPar
基本工程管理

\item {} 
\sphinxAtStartPar
Phoenix数据导入

\item {} 
\sphinxAtStartPar
FFT处理

\item {} 
\sphinxAtStartPar
基本数据筛选

\item {} 
\sphinxAtStartPar
EDI导出

\end{itemize}


\bigskip\hrule\bigskip



\subsection{已知问题}
\label{\detokenize{appendices/appendixE:id10}}\begin{itemize}
\item {} 
\sphinxAtStartPar
超大数据集（>10GB）时内存占用较高

\item {} 
\sphinxAtStartPar
某些显卡驱动可能与图表渲染冲突

\item {} 
\sphinxAtStartPar
极低频率（<0.001Hz）处理可能需要较长时间

\end{itemize}


\bigskip\hrule\bigskip



\subsection{系统要求}
\label{\detokenize{appendices/appendixE:id11}}

\begin{savenotes}\sphinxattablestart
\sphinxthistablewithglobalstyle
\centering
\begin{tabulary}{\linewidth}[t]{TT}
\sphinxtoprule
\begin{varwidth}[t]{\sphinxcolwidth{1}{2}}
\sphinxstyletheadfamily \sphinxAtStartPar
组件
\sphinxbeforeendvarwidth
\end{varwidth}%
&\begin{varwidth}[t]{\sphinxcolwidth{1}{2}}
\sphinxstyletheadfamily \sphinxAtStartPar
要求
\sphinxbeforeendvarwidth
\end{varwidth}%
\\
\sphinxmidrule
\sphinxtableatstartofbodyhook\begin{varwidth}[t]{\sphinxcolwidth{1}{2}}
\sphinxAtStartPar
操作系统
\sphinxbeforeendvarwidth
\end{varwidth}%
&\begin{varwidth}[t]{\sphinxcolwidth{1}{2}}
\sphinxAtStartPar
Windows 10/11 64位
\sphinxbeforeendvarwidth
\end{varwidth}%
\\
\sphinxhline\begin{varwidth}[t]{\sphinxcolwidth{1}{2}}
\sphinxAtStartPar
处理器
\sphinxbeforeendvarwidth
\end{varwidth}%
&\begin{varwidth}[t]{\sphinxcolwidth{1}{2}}
\sphinxAtStartPar
Intel/AMD 64位处理器，4核心以上推荐
\sphinxbeforeendvarwidth
\end{varwidth}%
\\
\sphinxhline\begin{varwidth}[t]{\sphinxcolwidth{1}{2}}
\sphinxAtStartPar
内存
\sphinxbeforeendvarwidth
\end{varwidth}%
&\begin{varwidth}[t]{\sphinxcolwidth{1}{2}}
\sphinxAtStartPar
8GB以上，16GB推荐
\sphinxbeforeendvarwidth
\end{varwidth}%
\\
\sphinxhline\begin{varwidth}[t]{\sphinxcolwidth{1}{2}}
\sphinxAtStartPar
硬盘
\sphinxbeforeendvarwidth
\end{varwidth}%
&\begin{varwidth}[t]{\sphinxcolwidth{1}{2}}
\sphinxAtStartPar
SSD推荐，至少10GB可用空间
\sphinxbeforeendvarwidth
\end{varwidth}%
\\
\sphinxhline\begin{varwidth}[t]{\sphinxcolwidth{1}{2}}
\sphinxAtStartPar
显卡
\sphinxbeforeendvarwidth
\end{varwidth}%
&\begin{varwidth}[t]{\sphinxcolwidth{1}{2}}
\sphinxAtStartPar
支持OpenGL 3.0以上
\sphinxbeforeendvarwidth
\end{varwidth}%
\\
\sphinxbottomrule
\end{tabulary}
\sphinxtableafterendhook\par
\sphinxattableend\end{savenotes}


\bigskip\hrule\bigskip



\subsection{技术支持}
\label{\detokenize{appendices/appendixE:id12}}
\sphinxAtStartPar
如需技术支持，请联系MTDP开发团队，提供：
\begin{itemize}
\item {} 
\sphinxAtStartPar
软件版本号

\item {} 
\sphinxAtStartPar
问题描述

\item {} 
\sphinxAtStartPar
相关数据文件（如适用）

\item {} 
\sphinxAtStartPar
系统配置信息

\end{itemize}

\sphinxstepscope


\section{附录F: 图表操作详解}
\label{\detokenize{appendices/appendixF:f}}\label{\detokenize{appendices/appendixF::doc}}
\sphinxAtStartPar
本附录详细介绍MTDP中图表的各种操作方法。


\bigskip\hrule\bigskip



\subsection{F.1 图表右键菜单}
\label{\detokenize{appendices/appendixF:f-1}}
\sphinxAtStartPar
在任意图表上点击右键，弹出以下菜单：


\subsubsection{F.1.1 显示设置}
\label{\detokenize{appendices/appendixF:f-1-1}}

\begin{savenotes}\sphinxattablestart
\sphinxthistablewithglobalstyle
\centering
\begin{tabulary}{\linewidth}[t]{TT}
\sphinxtoprule
\begin{varwidth}[t]{\sphinxcolwidth{1}{2}}
\sphinxstyletheadfamily \sphinxAtStartPar
菜单项
\sphinxbeforeendvarwidth
\end{varwidth}%
&\begin{varwidth}[t]{\sphinxcolwidth{1}{2}}
\sphinxstyletheadfamily \sphinxAtStartPar
功能
\sphinxbeforeendvarwidth
\end{varwidth}%
\\
\sphinxmidrule
\sphinxtableatstartofbodyhook\begin{varwidth}[t]{\sphinxcolwidth{1}{2}}
\sphinxAtStartPar
图表编辑器
\sphinxbeforeendvarwidth
\end{varwidth}%
&\begin{varwidth}[t]{\sphinxcolwidth{1}{2}}
\sphinxAtStartPar
打开完整的图表设置对话框
\sphinxbeforeendvarwidth
\end{varwidth}%
\\
\sphinxhline\begin{varwidth}[t]{\sphinxcolwidth{1}{2}}
\sphinxAtStartPar
显示图例
\sphinxbeforeendvarwidth
\end{varwidth}%
&\begin{varwidth}[t]{\sphinxcolwidth{1}{2}}
\sphinxAtStartPar
显示/隐藏图例
\sphinxbeforeendvarwidth
\end{varwidth}%
\\
\sphinxhline\begin{varwidth}[t]{\sphinxcolwidth{1}{2}}
\sphinxAtStartPar
显示命令栏
\sphinxbeforeendvarwidth
\end{varwidth}%
&\begin{varwidth}[t]{\sphinxcolwidth{1}{2}}
\sphinxAtStartPar
显示/隐藏图表工具栏
\sphinxbeforeendvarwidth
\end{varwidth}%
\\
\sphinxhline\begin{varwidth}[t]{\sphinxcolwidth{1}{2}}
\sphinxAtStartPar
全屏显示
\sphinxbeforeendvarwidth
\end{varwidth}%
&\begin{varwidth}[t]{\sphinxcolwidth{1}{2}}
\sphinxAtStartPar
图表全屏模式
\sphinxbeforeendvarwidth
\end{varwidth}%
\\
\sphinxbottomrule
\end{tabulary}
\sphinxtableafterendhook\par
\sphinxattableend\end{savenotes}


\subsubsection{F.1.2 坐标轴设置}
\label{\detokenize{appendices/appendixF:f-1-2}}

\begin{savenotes}\sphinxattablestart
\sphinxthistablewithglobalstyle
\centering
\begin{tabulary}{\linewidth}[t]{TT}
\sphinxtoprule
\begin{varwidth}[t]{\sphinxcolwidth{1}{2}}
\sphinxstyletheadfamily \sphinxAtStartPar
菜单项
\sphinxbeforeendvarwidth
\end{varwidth}%
&\begin{varwidth}[t]{\sphinxcolwidth{1}{2}}
\sphinxstyletheadfamily \sphinxAtStartPar
功能
\sphinxbeforeendvarwidth
\end{varwidth}%
\\
\sphinxmidrule
\sphinxtableatstartofbodyhook\begin{varwidth}[t]{\sphinxcolwidth{1}{2}}
\sphinxAtStartPar
自动调整全部
\sphinxbeforeendvarwidth
\end{varwidth}%
&\begin{varwidth}[t]{\sphinxcolwidth{1}{2}}
\sphinxAtStartPar
X、Y轴自动调整到最佳范围
\sphinxbeforeendvarwidth
\end{varwidth}%
\\
\sphinxhline\begin{varwidth}[t]{\sphinxcolwidth{1}{2}}
\sphinxAtStartPar
自动调整X轴
\sphinxbeforeendvarwidth
\end{varwidth}%
&\begin{varwidth}[t]{\sphinxcolwidth{1}{2}}
\sphinxAtStartPar
仅X轴自动调整
\sphinxbeforeendvarwidth
\end{varwidth}%
\\
\sphinxhline\begin{varwidth}[t]{\sphinxcolwidth{1}{2}}
\sphinxAtStartPar
自动调整Y轴
\sphinxbeforeendvarwidth
\end{varwidth}%
&\begin{varwidth}[t]{\sphinxcolwidth{1}{2}}
\sphinxAtStartPar
仅Y轴自动调整
\sphinxbeforeendvarwidth
\end{varwidth}%
\\
\sphinxhline\begin{varwidth}[t]{\sphinxcolwidth{1}{2}}
\sphinxAtStartPar
恢复坐标轴
\sphinxbeforeendvarwidth
\end{varwidth}%
&\begin{varwidth}[t]{\sphinxcolwidth{1}{2}}
\sphinxAtStartPar
恢复到保存的坐标范围
\sphinxbeforeendvarwidth
\end{varwidth}%
\\
\sphinxhline\begin{varwidth}[t]{\sphinxcolwidth{1}{2}}
\sphinxAtStartPar
记录坐标轴
\sphinxbeforeendvarwidth
\end{varwidth}%
&\begin{varwidth}[t]{\sphinxcolwidth{1}{2}}
\sphinxAtStartPar
保存当前坐标范围
\sphinxbeforeendvarwidth
\end{varwidth}%
\\
\sphinxhline\begin{varwidth}[t]{\sphinxcolwidth{1}{2}}
\sphinxAtStartPar
X轴对数
\sphinxbeforeendvarwidth
\end{varwidth}%
&\begin{varwidth}[t]{\sphinxcolwidth{1}{2}}
\sphinxAtStartPar
X轴切换为对数坐标
\sphinxbeforeendvarwidth
\end{varwidth}%
\\
\sphinxhline\begin{varwidth}[t]{\sphinxcolwidth{1}{2}}
\sphinxAtStartPar
Y轴对数
\sphinxbeforeendvarwidth
\end{varwidth}%
&\begin{varwidth}[t]{\sphinxcolwidth{1}{2}}
\sphinxAtStartPar
Y轴切换为对数坐标
\sphinxbeforeendvarwidth
\end{varwidth}%
\\
\sphinxbottomrule
\end{tabulary}
\sphinxtableafterendhook\par
\sphinxattableend\end{savenotes}


\subsubsection{F.1.3 缩放功能}
\label{\detokenize{appendices/appendixF:f-1-3}}

\begin{savenotes}\sphinxattablestart
\sphinxthistablewithglobalstyle
\centering
\begin{tabulary}{\linewidth}[t]{TT}
\sphinxtoprule
\begin{varwidth}[t]{\sphinxcolwidth{1}{2}}
\sphinxstyletheadfamily \sphinxAtStartPar
菜单项
\sphinxbeforeendvarwidth
\end{varwidth}%
&\begin{varwidth}[t]{\sphinxcolwidth{1}{2}}
\sphinxstyletheadfamily \sphinxAtStartPar
功能
\sphinxbeforeendvarwidth
\end{varwidth}%
\\
\sphinxmidrule
\sphinxtableatstartofbodyhook\begin{varwidth}[t]{\sphinxcolwidth{1}{2}}
\sphinxAtStartPar
放大X轴
\sphinxbeforeendvarwidth
\end{varwidth}%
&\begin{varwidth}[t]{\sphinxcolwidth{1}{2}}
\sphinxAtStartPar
放大X轴（缩小显示范围）
\sphinxbeforeendvarwidth
\end{varwidth}%
\\
\sphinxhline\begin{varwidth}[t]{\sphinxcolwidth{1}{2}}
\sphinxAtStartPar
放大Y轴
\sphinxbeforeendvarwidth
\end{varwidth}%
&\begin{varwidth}[t]{\sphinxcolwidth{1}{2}}
\sphinxAtStartPar
放大Y轴
\sphinxbeforeendvarwidth
\end{varwidth}%
\\
\sphinxhline\begin{varwidth}[t]{\sphinxcolwidth{1}{2}}
\sphinxAtStartPar
缩小X轴
\sphinxbeforeendvarwidth
\end{varwidth}%
&\begin{varwidth}[t]{\sphinxcolwidth{1}{2}}
\sphinxAtStartPar
缩小X轴（增大显示范围）
\sphinxbeforeendvarwidth
\end{varwidth}%
\\
\sphinxhline\begin{varwidth}[t]{\sphinxcolwidth{1}{2}}
\sphinxAtStartPar
缩小Y轴
\sphinxbeforeendvarwidth
\end{varwidth}%
&\begin{varwidth}[t]{\sphinxcolwidth{1}{2}}
\sphinxAtStartPar
缩小Y轴
\sphinxbeforeendvarwidth
\end{varwidth}%
\\
\sphinxbottomrule
\end{tabulary}
\sphinxtableafterendhook\par
\sphinxattableend\end{savenotes}


\subsubsection{F.1.4 导出功能}
\label{\detokenize{appendices/appendixF:f-1-4}}

\begin{savenotes}\sphinxattablestart
\sphinxthistablewithglobalstyle
\centering
\begin{tabulary}{\linewidth}[t]{TTT}
\sphinxtoprule
\begin{varwidth}[t]{\sphinxcolwidth{1}{3}}
\sphinxstyletheadfamily \sphinxAtStartPar
菜单项
\sphinxbeforeendvarwidth
\end{varwidth}%
&\begin{varwidth}[t]{\sphinxcolwidth{1}{3}}
\sphinxstyletheadfamily \sphinxAtStartPar
格式
\sphinxbeforeendvarwidth
\end{varwidth}%
&\begin{varwidth}[t]{\sphinxcolwidth{1}{3}}
\sphinxstyletheadfamily \sphinxAtStartPar
说明
\sphinxbeforeendvarwidth
\end{varwidth}%
\\
\sphinxmidrule
\sphinxtableatstartofbodyhook\begin{varwidth}[t]{\sphinxcolwidth{1}{3}}
\sphinxAtStartPar
导出到剪贴板
\sphinxbeforeendvarwidth
\end{varwidth}%
&\begin{varwidth}[t]{\sphinxcolwidth{1}{3}}
\sphinxAtStartPar
BMP
\sphinxbeforeendvarwidth
\end{varwidth}%
&\begin{varwidth}[t]{\sphinxcolwidth{1}{3}}
\sphinxAtStartPar
复制到剪贴板
\sphinxbeforeendvarwidth
\end{varwidth}%
\\
\sphinxhline\begin{varwidth}[t]{\sphinxcolwidth{1}{3}}
\sphinxAtStartPar
导出为PDF
\sphinxbeforeendvarwidth
\end{varwidth}%
&\begin{varwidth}[t]{\sphinxcolwidth{1}{3}}
\sphinxAtStartPar
PDF
\sphinxbeforeendvarwidth
\end{varwidth}%
&\begin{varwidth}[t]{\sphinxcolwidth{1}{3}}
\sphinxAtStartPar
矢量格式，适合打印
\sphinxbeforeendvarwidth
\end{varwidth}%
\\
\sphinxhline\begin{varwidth}[t]{\sphinxcolwidth{1}{3}}
\sphinxAtStartPar
导出为图片
\sphinxbeforeendvarwidth
\end{varwidth}%
&\begin{varwidth}[t]{\sphinxcolwidth{1}{3}}
\sphinxAtStartPar
PNG/JPG
\sphinxbeforeendvarwidth
\end{varwidth}%
&\begin{varwidth}[t]{\sphinxcolwidth{1}{3}}
\sphinxAtStartPar
位图格式
\sphinxbeforeendvarwidth
\end{varwidth}%
\\
\sphinxhline\begin{varwidth}[t]{\sphinxcolwidth{1}{3}}
\sphinxAtStartPar
导出为SVG
\sphinxbeforeendvarwidth
\end{varwidth}%
&\begin{varwidth}[t]{\sphinxcolwidth{1}{3}}
\sphinxAtStartPar
SVG
\sphinxbeforeendvarwidth
\end{varwidth}%
&\begin{varwidth}[t]{\sphinxcolwidth{1}{3}}
\sphinxAtStartPar
矢量格式
\sphinxbeforeendvarwidth
\end{varwidth}%
\\
\sphinxhline\begin{varwidth}[t]{\sphinxcolwidth{1}{3}}
\sphinxAtStartPar
导出为HTML
\sphinxbeforeendvarwidth
\end{varwidth}%
&\begin{varwidth}[t]{\sphinxcolwidth{1}{3}}
\sphinxAtStartPar
HTML
\sphinxbeforeendvarwidth
\end{varwidth}%
&\begin{varwidth}[t]{\sphinxcolwidth{1}{3}}
\sphinxAtStartPar
网页格式
\sphinxbeforeendvarwidth
\end{varwidth}%
\\
\sphinxhline\begin{varwidth}[t]{\sphinxcolwidth{1}{3}}
\sphinxAtStartPar
导出数据
\sphinxbeforeendvarwidth
\end{varwidth}%
&\begin{varwidth}[t]{\sphinxcolwidth{1}{3}}
\sphinxAtStartPar
TXT/CSV
\sphinxbeforeendvarwidth
\end{varwidth}%
&\begin{varwidth}[t]{\sphinxcolwidth{1}{3}}
\sphinxAtStartPar
导出数据点
\sphinxbeforeendvarwidth
\end{varwidth}%
\\
\sphinxbottomrule
\end{tabulary}
\sphinxtableafterendhook\par
\sphinxattableend\end{savenotes}


\bigskip\hrule\bigskip



\subsection{F.2 坐标轴操作}
\label{\detokenize{appendices/appendixF:f-2}}

\subsubsection{F.2.1 点击坐标轴设置范围}
\label{\detokenize{appendices/appendixF:f-2-1}}\begin{enumerate}
\sphinxsetlistlabels{\arabic}{enumi}{enumii}{}{.}%
\item {} 
\sphinxAtStartPar
点击图表的坐标轴（X轴或Y轴）

\item {} 
\sphinxAtStartPar
在弹出对话框中设置：
\begin{itemize}
\item {} 
\sphinxAtStartPar
最小值

\item {} 
\sphinxAtStartPar
最大值

\item {} 
\sphinxAtStartPar
是否自动

\end{itemize}

\item {} 
\sphinxAtStartPar
点击确定

\end{enumerate}


\subsubsection{F.2.2 对数/线性坐标切换}
\label{\detokenize{appendices/appendixF:f-2-2}}\begin{itemize}
\item {} 
\sphinxAtStartPar
\sphinxstylestrong{X轴对数}：右键菜单 → X轴对数

\item {} 
\sphinxAtStartPar
\sphinxstylestrong{Y轴对数}：右键菜单 → Y轴对数

\end{itemize}
\begin{quote}

\sphinxAtStartPar
💡 \sphinxstylestrong{提示}：视电阻率曲线通常使用双对数坐标。
\end{quote}


\bigskip\hrule\bigskip



\subsection{F.3 图表缩放与平移}
\label{\detokenize{appendices/appendixF:f-3}}

\subsubsection{F.3.1 鼠标操作}
\label{\detokenize{appendices/appendixF:f-3-1}}

\begin{savenotes}\sphinxattablestart
\sphinxthistablewithglobalstyle
\centering
\begin{tabulary}{\linewidth}[t]{TT}
\sphinxtoprule
\begin{varwidth}[t]{\sphinxcolwidth{1}{2}}
\sphinxstyletheadfamily \sphinxAtStartPar
操作
\sphinxbeforeendvarwidth
\end{varwidth}%
&\begin{varwidth}[t]{\sphinxcolwidth{1}{2}}
\sphinxstyletheadfamily \sphinxAtStartPar
功能
\sphinxbeforeendvarwidth
\end{varwidth}%
\\
\sphinxmidrule
\sphinxtableatstartofbodyhook\begin{varwidth}[t]{\sphinxcolwidth{1}{2}}
\sphinxAtStartPar
滚轮向上
\sphinxbeforeendvarwidth
\end{varwidth}%
&\begin{varwidth}[t]{\sphinxcolwidth{1}{2}}
\sphinxAtStartPar
放大
\sphinxbeforeendvarwidth
\end{varwidth}%
\\
\sphinxhline\begin{varwidth}[t]{\sphinxcolwidth{1}{2}}
\sphinxAtStartPar
滚轮向下
\sphinxbeforeendvarwidth
\end{varwidth}%
&\begin{varwidth}[t]{\sphinxcolwidth{1}{2}}
\sphinxAtStartPar
缩小
\sphinxbeforeendvarwidth
\end{varwidth}%
\\
\sphinxhline\begin{varwidth}[t]{\sphinxcolwidth{1}{2}}
\sphinxAtStartPar
左键拖动
\sphinxbeforeendvarwidth
\end{varwidth}%
&\begin{varwidth}[t]{\sphinxcolwidth{1}{2}}
\sphinxAtStartPar
平移图表
\sphinxbeforeendvarwidth
\end{varwidth}%
\\
\sphinxhline\begin{varwidth}[t]{\sphinxcolwidth{1}{2}}
\sphinxAtStartPar
双击
\sphinxbeforeendvarwidth
\end{varwidth}%
&\begin{varwidth}[t]{\sphinxcolwidth{1}{2}}
\sphinxAtStartPar
恢复原始视图
\sphinxbeforeendvarwidth
\end{varwidth}%
\\
\sphinxbottomrule
\end{tabulary}
\sphinxtableafterendhook\par
\sphinxattableend\end{savenotes}


\subsubsection{F.3.2 键盘快捷键}
\label{\detokenize{appendices/appendixF:f-3-2}}

\begin{savenotes}\sphinxattablestart
\sphinxthistablewithglobalstyle
\centering
\begin{tabulary}{\linewidth}[t]{TT}
\sphinxtoprule
\begin{varwidth}[t]{\sphinxcolwidth{1}{2}}
\sphinxstyletheadfamily \sphinxAtStartPar
快捷键
\sphinxbeforeendvarwidth
\end{varwidth}%
&\begin{varwidth}[t]{\sphinxcolwidth{1}{2}}
\sphinxstyletheadfamily \sphinxAtStartPar
功能
\sphinxbeforeendvarwidth
\end{varwidth}%
\\
\sphinxmidrule
\sphinxtableatstartofbodyhook\begin{varwidth}[t]{\sphinxcolwidth{1}{2}}
\sphinxAtStartPar
Ctrl
\sphinxbeforeendvarwidth
\end{varwidth}%
&\begin{varwidth}[t]{\sphinxcolwidth{1}{2}}
\sphinxAtStartPar
按住时拖动平移
\sphinxbeforeendvarwidth
\end{varwidth}%
\\
\sphinxhline\begin{varwidth}[t]{\sphinxcolwidth{1}{2}}
\sphinxAtStartPar
Ctrl + +
\sphinxbeforeendvarwidth
\end{varwidth}%
&\begin{varwidth}[t]{\sphinxcolwidth{1}{2}}
\sphinxAtStartPar
放大
\sphinxbeforeendvarwidth
\end{varwidth}%
\\
\sphinxhline\begin{varwidth}[t]{\sphinxcolwidth{1}{2}}
\sphinxAtStartPar
Ctrl + \sphinxhyphen{}
\sphinxbeforeendvarwidth
\end{varwidth}%
&\begin{varwidth}[t]{\sphinxcolwidth{1}{2}}
\sphinxAtStartPar
缩小
\sphinxbeforeendvarwidth
\end{varwidth}%
\\
\sphinxhline\begin{varwidth}[t]{\sphinxcolwidth{1}{2}}
\sphinxAtStartPar
Ctrl + 0
\sphinxbeforeendvarwidth
\end{varwidth}%
&\begin{varwidth}[t]{\sphinxcolwidth{1}{2}}
\sphinxAtStartPar
恢复原始大小
\sphinxbeforeendvarwidth
\end{varwidth}%
\\
\sphinxbottomrule
\end{tabulary}
\sphinxtableafterendhook\par
\sphinxattableend\end{savenotes}


\bigskip\hrule\bigskip



\subsection{F.4 数据筛选图表操作}
\label{\detokenize{appendices/appendixF:f-4}}
\sphinxAtStartPar
在数据筛选界面的图表支持多种选择方式：


\subsubsection{F.4.1 选择模式}
\label{\detokenize{appendices/appendixF:f-4-1}}
\sphinxAtStartPar
通过下拉框或菜单切换选择模式：


\begin{savenotes}\sphinxattablestart
\sphinxthistablewithglobalstyle
\centering
\begin{tabulary}{\linewidth}[t]{TTT}
\sphinxtoprule
\begin{varwidth}[t]{\sphinxcolwidth{1}{3}}
\sphinxstyletheadfamily \sphinxAtStartPar
模式
\sphinxbeforeendvarwidth
\end{varwidth}%
&\begin{varwidth}[t]{\sphinxcolwidth{1}{3}}
\sphinxstyletheadfamily \sphinxAtStartPar
说明
\sphinxbeforeendvarwidth
\end{varwidth}%
&\begin{varwidth}[t]{\sphinxcolwidth{1}{3}}
\sphinxstyletheadfamily \sphinxAtStartPar
操作方法
\sphinxbeforeendvarwidth
\end{varwidth}%
\\
\sphinxmidrule
\sphinxtableatstartofbodyhook\begin{varwidth}[t]{\sphinxcolwidth{1}{3}}
\sphinxAtStartPar
点选择
\sphinxbeforeendvarwidth
\end{varwidth}%
&\begin{varwidth}[t]{\sphinxcolwidth{1}{3}}
\sphinxAtStartPar
选择单个数据点
\sphinxbeforeendvarwidth
\end{varwidth}%
&\begin{varwidth}[t]{\sphinxcolwidth{1}{3}}
\sphinxAtStartPar
点击数据点
\sphinxbeforeendvarwidth
\end{varwidth}%
\\
\sphinxhline\begin{varwidth}[t]{\sphinxcolwidth{1}{3}}
\sphinxAtStartPar
矩形选择
\sphinxbeforeendvarwidth
\end{varwidth}%
&\begin{varwidth}[t]{\sphinxcolwidth{1}{3}}
\sphinxAtStartPar
框选矩形区域
\sphinxbeforeendvarwidth
\end{varwidth}%
&\begin{varwidth}[t]{\sphinxcolwidth{1}{3}}
\sphinxAtStartPar
拖动绘制矩形
\sphinxbeforeendvarwidth
\end{varwidth}%
\\
\sphinxhline\begin{varwidth}[t]{\sphinxcolwidth{1}{3}}
\sphinxAtStartPar
矩形反选
\sphinxbeforeendvarwidth
\end{varwidth}%
&\begin{varwidth}[t]{\sphinxcolwidth{1}{3}}
\sphinxAtStartPar
选择矩形外的点
\sphinxbeforeendvarwidth
\end{varwidth}%
&\begin{varwidth}[t]{\sphinxcolwidth{1}{3}}
\sphinxAtStartPar
拖动绘制矩形
\sphinxbeforeendvarwidth
\end{varwidth}%
\\
\sphinxhline\begin{varwidth}[t]{\sphinxcolwidth{1}{3}}
\sphinxAtStartPar
多边形选择
\sphinxbeforeendvarwidth
\end{varwidth}%
&\begin{varwidth}[t]{\sphinxcolwidth{1}{3}}
\sphinxAtStartPar
选择多边形区域
\sphinxbeforeendvarwidth
\end{varwidth}%
&\begin{varwidth}[t]{\sphinxcolwidth{1}{3}}
\sphinxAtStartPar
依次点击顶点
\sphinxbeforeendvarwidth
\end{varwidth}%
\\
\sphinxhline\begin{varwidth}[t]{\sphinxcolwidth{1}{3}}
\sphinxAtStartPar
多边形反选
\sphinxbeforeendvarwidth
\end{varwidth}%
&\begin{varwidth}[t]{\sphinxcolwidth{1}{3}}
\sphinxAtStartPar
选择多边形外的点
\sphinxbeforeendvarwidth
\end{varwidth}%
&\begin{varwidth}[t]{\sphinxcolwidth{1}{3}}
\sphinxAtStartPar
依次点击顶点
\sphinxbeforeendvarwidth
\end{varwidth}%
\\
\sphinxhline\begin{varwidth}[t]{\sphinxcolwidth{1}{3}}
\sphinxAtStartPar
X轴范围
\sphinxbeforeendvarwidth
\end{varwidth}%
&\begin{varwidth}[t]{\sphinxcolwidth{1}{3}}
\sphinxAtStartPar
选择X轴范围内的点
\sphinxbeforeendvarwidth
\end{varwidth}%
&\begin{varwidth}[t]{\sphinxcolwidth{1}{3}}
\sphinxAtStartPar
拖动设置范围
\sphinxbeforeendvarwidth
\end{varwidth}%
\\
\sphinxhline\begin{varwidth}[t]{\sphinxcolwidth{1}{3}}
\sphinxAtStartPar
Y轴范围
\sphinxbeforeendvarwidth
\end{varwidth}%
&\begin{varwidth}[t]{\sphinxcolwidth{1}{3}}
\sphinxAtStartPar
选择Y轴范围内的点
\sphinxbeforeendvarwidth
\end{varwidth}%
&\begin{varwidth}[t]{\sphinxcolwidth{1}{3}}
\sphinxAtStartPar
拖动设置范围
\sphinxbeforeendvarwidth
\end{varwidth}%
\\
\sphinxhline\begin{varwidth}[t]{\sphinxcolwidth{1}{3}}
\sphinxAtStartPar
X轴范围反选
\sphinxbeforeendvarwidth
\end{varwidth}%
&\begin{varwidth}[t]{\sphinxcolwidth{1}{3}}
\sphinxAtStartPar
选择X轴范围外的点
\sphinxbeforeendvarwidth
\end{varwidth}%
&\begin{varwidth}[t]{\sphinxcolwidth{1}{3}}
\sphinxAtStartPar
拖动设置范围
\sphinxbeforeendvarwidth
\end{varwidth}%
\\
\sphinxhline\begin{varwidth}[t]{\sphinxcolwidth{1}{3}}
\sphinxAtStartPar
Y轴范围反选
\sphinxbeforeendvarwidth
\end{varwidth}%
&\begin{varwidth}[t]{\sphinxcolwidth{1}{3}}
\sphinxAtStartPar
选择Y轴范围外的点
\sphinxbeforeendvarwidth
\end{varwidth}%
&\begin{varwidth}[t]{\sphinxcolwidth{1}{3}}
\sphinxAtStartPar
拖动设置范围
\sphinxbeforeendvarwidth
\end{varwidth}%
\\
\sphinxbottomrule
\end{tabulary}
\sphinxtableafterendhook\par
\sphinxattableend\end{savenotes}


\subsubsection{F.4.2 完成选择}
\label{\detokenize{appendices/appendixF:f-4-2}}

\begin{savenotes}\sphinxattablestart
\sphinxthistablewithglobalstyle
\centering
\begin{tabulary}{\linewidth}[t]{TT}
\sphinxtoprule
\begin{varwidth}[t]{\sphinxcolwidth{1}{2}}
\sphinxstyletheadfamily \sphinxAtStartPar
操作
\sphinxbeforeendvarwidth
\end{varwidth}%
&\begin{varwidth}[t]{\sphinxcolwidth{1}{2}}
\sphinxstyletheadfamily \sphinxAtStartPar
功能
\sphinxbeforeendvarwidth
\end{varwidth}%
\\
\sphinxmidrule
\sphinxtableatstartofbodyhook\begin{varwidth}[t]{\sphinxcolwidth{1}{2}}
\sphinxAtStartPar
Enter
\sphinxbeforeendvarwidth
\end{varwidth}%
&\begin{varwidth}[t]{\sphinxcolwidth{1}{2}}
\sphinxAtStartPar
完成当前选择
\sphinxbeforeendvarwidth
\end{varwidth}%
\\
\sphinxhline\begin{varwidth}[t]{\sphinxcolwidth{1}{2}}
\sphinxAtStartPar
空格
\sphinxbeforeendvarwidth
\end{varwidth}%
&\begin{varwidth}[t]{\sphinxcolwidth{1}{2}}
\sphinxAtStartPar
完成当前选择
\sphinxbeforeendvarwidth
\end{varwidth}%
\\
\sphinxhline\begin{varwidth}[t]{\sphinxcolwidth{1}{2}}
\sphinxAtStartPar
Backspace
\sphinxbeforeendvarwidth
\end{varwidth}%
&\begin{varwidth}[t]{\sphinxcolwidth{1}{2}}
\sphinxAtStartPar
清除当前选择区域
\sphinxbeforeendvarwidth
\end{varwidth}%
\\
\sphinxbottomrule
\end{tabulary}
\sphinxtableafterendhook\par
\sphinxattableend\end{savenotes}


\subsubsection{F.4.3 撤销与重做}
\label{\detokenize{appendices/appendixF:f-4-3}}

\begin{savenotes}\sphinxattablestart
\sphinxthistablewithglobalstyle
\centering
\begin{tabulary}{\linewidth}[t]{TT}
\sphinxtoprule
\begin{varwidth}[t]{\sphinxcolwidth{1}{2}}
\sphinxstyletheadfamily \sphinxAtStartPar
快捷键
\sphinxbeforeendvarwidth
\end{varwidth}%
&\begin{varwidth}[t]{\sphinxcolwidth{1}{2}}
\sphinxstyletheadfamily \sphinxAtStartPar
功能
\sphinxbeforeendvarwidth
\end{varwidth}%
\\
\sphinxmidrule
\sphinxtableatstartofbodyhook\begin{varwidth}[t]{\sphinxcolwidth{1}{2}}
\sphinxAtStartPar
Ctrl + Z
\sphinxbeforeendvarwidth
\end{varwidth}%
&\begin{varwidth}[t]{\sphinxcolwidth{1}{2}}
\sphinxAtStartPar
撤销上一步选择
\sphinxbeforeendvarwidth
\end{varwidth}%
\\
\sphinxhline\begin{varwidth}[t]{\sphinxcolwidth{1}{2}}
\sphinxAtStartPar
Ctrl + Y
\sphinxbeforeendvarwidth
\end{varwidth}%
&\begin{varwidth}[t]{\sphinxcolwidth{1}{2}}
\sphinxAtStartPar
重做
\sphinxbeforeendvarwidth
\end{varwidth}%
\\
\sphinxbottomrule
\end{tabulary}
\sphinxtableafterendhook\par
\sphinxattableend\end{savenotes}


\subsubsection{F.4.4 操作示例}
\label{\detokenize{appendices/appendixF:f-4-4}}
\sphinxAtStartPar
\sphinxstylestrong{矩形选择：}
\begin{enumerate}
\sphinxsetlistlabels{\arabic}{enumi}{enumii}{}{.}%
\item {} 
\sphinxAtStartPar
选择"矩形选择"模式

\item {} 
\sphinxAtStartPar
在图表上按住左键拖动

\item {} 
\sphinxAtStartPar
松开鼠标完成矩形

\item {} 
\sphinxAtStartPar
按 Enter 确认选择

\end{enumerate}

\sphinxAtStartPar
\sphinxstylestrong{多边形选择：}
\begin{enumerate}
\sphinxsetlistlabels{\arabic}{enumi}{enumii}{}{.}%
\item {} 
\sphinxAtStartPar
选择"多边形选择"模式

\item {} 
\sphinxAtStartPar
依次点击多边形顶点

\item {} 
\sphinxAtStartPar
按 Enter 闭合多边形并确认

\end{enumerate}


\bigskip\hrule\bigskip



\subsection{F.5 添加临时曲线}
\label{\detokenize{appendices/appendixF:f-5}}
\sphinxAtStartPar
可以在现有图表上叠加显示外部数据：
\begin{enumerate}
\sphinxsetlistlabels{\arabic}{enumi}{enumii}{}{.}%
\item {} 
\sphinxAtStartPar
右键菜单 → 添加临时曲线

\item {} 
\sphinxAtStartPar
选择数据文件（TXT格式）

\item {} 
\sphinxAtStartPar
数据格式：每行两个数值，空格或制表符分隔

\item {} 
\sphinxAtStartPar
曲线自动添加到图表

\end{enumerate}

\sphinxAtStartPar
\sphinxstylestrong{数据文件示例：}

\begin{sphinxVerbatim}[commandchars=\\\{\}]
\PYG{l+m+mf}{0.01}  \PYG{l+m+mi}{100}
\PYG{l+m+mf}{0.02}  \PYG{l+m+mi}{95}
\PYG{l+m+mf}{0.05}  \PYG{l+m+mi}{90}
\PYG{l+m+mf}{0.1}   \PYG{l+m+mi}{85}
\end{sphinxVerbatim}


\bigskip\hrule\bigskip



\subsection{F.6 图表编辑器}
\label{\detokenize{appendices/appendixF:f-6}}
\sphinxAtStartPar
右键菜单 → 图表编辑器，打开完整的图表设置对话框：


\subsubsection{F.6.1 常用设置}
\label{\detokenize{appendices/appendixF:f-6-1}}

\begin{savenotes}\sphinxattablestart
\sphinxthistablewithglobalstyle
\centering
\begin{tabulary}{\linewidth}[t]{TT}
\sphinxtoprule
\begin{varwidth}[t]{\sphinxcolwidth{1}{2}}
\sphinxstyletheadfamily \sphinxAtStartPar
分类
\sphinxbeforeendvarwidth
\end{varwidth}%
&\begin{varwidth}[t]{\sphinxcolwidth{1}{2}}
\sphinxstyletheadfamily \sphinxAtStartPar
可设置项
\sphinxbeforeendvarwidth
\end{varwidth}%
\\
\sphinxmidrule
\sphinxtableatstartofbodyhook\begin{varwidth}[t]{\sphinxcolwidth{1}{2}}
\sphinxAtStartPar
标题
\sphinxbeforeendvarwidth
\end{varwidth}%
&\begin{varwidth}[t]{\sphinxcolwidth{1}{2}}
\sphinxAtStartPar
标题文本、字体、颜色
\sphinxbeforeendvarwidth
\end{varwidth}%
\\
\sphinxhline\begin{varwidth}[t]{\sphinxcolwidth{1}{2}}
\sphinxAtStartPar
坐标轴
\sphinxbeforeendvarwidth
\end{varwidth}%
&\begin{varwidth}[t]{\sphinxcolwidth{1}{2}}
\sphinxAtStartPar
范围、刻度、标签、网格
\sphinxbeforeendvarwidth
\end{varwidth}%
\\
\sphinxhline\begin{varwidth}[t]{\sphinxcolwidth{1}{2}}
\sphinxAtStartPar
图例
\sphinxbeforeendvarwidth
\end{varwidth}%
&\begin{varwidth}[t]{\sphinxcolwidth{1}{2}}
\sphinxAtStartPar
位置、字体、可见性
\sphinxbeforeendvarwidth
\end{varwidth}%
\\
\sphinxhline\begin{varwidth}[t]{\sphinxcolwidth{1}{2}}
\sphinxAtStartPar
系列
\sphinxbeforeendvarwidth
\end{varwidth}%
&\begin{varwidth}[t]{\sphinxcolwidth{1}{2}}
\sphinxAtStartPar
线型、颜色、标记
\sphinxbeforeendvarwidth
\end{varwidth}%
\\
\sphinxhline\begin{varwidth}[t]{\sphinxcolwidth{1}{2}}
\sphinxAtStartPar
背景
\sphinxbeforeendvarwidth
\end{varwidth}%
&\begin{varwidth}[t]{\sphinxcolwidth{1}{2}}
\sphinxAtStartPar
颜色、渐变、边框
\sphinxbeforeendvarwidth
\end{varwidth}%
\\
\sphinxbottomrule
\end{tabulary}
\sphinxtableafterendhook\par
\sphinxattableend\end{savenotes}


\subsubsection{F.6.2 保存设置}
\label{\detokenize{appendices/appendixF:f-6-2}}
\sphinxAtStartPar
图表设置会自动保存，下次打开时恢复。


\bigskip\hrule\bigskip



\subsection{F.7 多图表对比}
\label{\detokenize{appendices/appendixF:f-7}}

\subsubsection{F.7.1 多测点叠加}
\label{\detokenize{appendices/appendixF:f-7-1}}\begin{enumerate}
\sphinxsetlistlabels{\arabic}{enumi}{enumii}{}{.}%
\item {} 
\sphinxAtStartPar
选择多个测点（Ctrl+点击）

\item {} 
\sphinxAtStartPar
右键选择"对比曲线"

\item {} 
\sphinxAtStartPar
所有曲线显示在同一图表

\end{enumerate}


\subsubsection{F.7.2 图表工具栏}
\label{\detokenize{appendices/appendixF:f-7-2}}
\sphinxAtStartPar
显示图表工具栏后可快速访问常用功能：
\begin{itemize}
\item {} 
\sphinxAtStartPar
缩放工具

\item {} 
\sphinxAtStartPar
平移工具

\item {} 
\sphinxAtStartPar
复制

\item {} 
\sphinxAtStartPar
打印

\item {} 
\sphinxAtStartPar
保存

\end{itemize}


\bigskip\hrule\bigskip



\subsection{F.8 常见问题}
\label{\detokenize{appendices/appendixF:f-8}}
\sphinxAtStartPar
\sphinxstylestrong{Q: 如何恢复误操作的坐标范围？}

\sphinxAtStartPar
A: 右键菜单 → 恢复坐标轴，或双击图表恢复。

\sphinxAtStartPar
\sphinxstylestrong{Q: 如何同时调整多个图表的坐标？}

\sphinxAtStartPar
A: 目前需要分别调整每个图表。

\sphinxAtStartPar
\sphinxstylestrong{Q: 导出的图片模糊怎么办？}

\sphinxAtStartPar
A: 导出为PDF或SVG矢量格式，可无限放大。

\sphinxstepscope


\section{贡献指南}
\label{\detokenize{appendix/contributing:id1}}\label{\detokenize{appendix/contributing::doc}}
\sphinxAtStartPar
感谢您对 MTDataPro 文档的关注与支持！


\subsection{如何贡献}
\label{\detokenize{appendix/contributing:id2}}\begin{enumerate}
\sphinxsetlistlabels{\arabic}{enumi}{enumii}{}{.}%
\item {} 
\sphinxAtStartPar
\sphinxstylestrong{发现问题}：浏览 \sphinxhref{https://github.com/EMWPJ/MTDataPro-docs/issues}{Issue 列表} (https://github.com/EMWPJ/MTDataPro\sphinxhyphen{}docs/issues) 或报告新问题

\item {} 
\sphinxAtStartPar
\sphinxstylestrong{Fork 仓库}：创建您自己的副本

\item {} 
\sphinxAtStartPar
\sphinxstylestrong{进行修改}：编辑 Markdown 源文件

\item {} 
\sphinxAtStartPar
\sphinxstylestrong{本地预览}：构建文档验证修改

\item {} 
\sphinxAtStartPar
\sphinxstylestrong{提交 Pull Request}：您的修改将被审核并合并

\end{enumerate}


\subsection{文档结构}
\label{\detokenize{appendix/contributing:id3}}
\begin{sphinxVerbatim}[commandchars=\\\{\}]
docs/
├── source/
│   ├── conf.py           \PYGZsh{} Sphinx 配置
│   ├── index.md          \PYGZsh{} 主页面
│   ├── chapters/         \PYGZsh{} 章节文件 (chapter1\PYGZhy{}10)
│   ├── appendices/       \PYGZsh{} 附录文件 (appendixA\PYGZhy{}F)
│   ├── intro/            \PYGZsh{} 入门指南
│   ├── tutorial/         \PYGZsh{} 进阶教程
│   └── gallery/          \PYGZsh{} 实例展示
└── requirements.txt      \PYGZsh{} Python 依赖
\end{sphinxVerbatim}


\subsection{本地构建}
\label{\detokenize{appendix/contributing:id4}}
\begin{sphinxVerbatim}[commandchars=\\\{\}]
\PYG{c+c1}{\PYGZsh{} 安装依赖}
pip\PYG{+w}{ }install\PYG{+w}{ }\PYGZhy{}r\PYG{+w}{ }requirements.txt

\PYG{c+c1}{\PYGZsh{} 构建 HTML 文档}
sphinx\PYGZhy{}build\PYG{+w}{ }\PYGZhy{}b\PYG{+w}{ }html\PYG{+w}{ }\PYG{n+nb}{source}\PYG{+w}{ }\PYGZus{}build/html

\PYG{c+c1}{\PYGZsh{} 预览文档}
open\PYG{+w}{ }\PYGZus{}build/html/index.html
\end{sphinxVerbatim}


\subsection{写作规范}
\label{\detokenize{appendix/contributing:id5}}\begin{itemize}
\item {} 
\sphinxAtStartPar
使用清晰、简洁的语言

\item {} 
\sphinxAtStartPar
适当包含代码示例

\item {} 
\sphinxAtStartPar
用户面向的内容使用中文编写

\item {} 
\sphinxAtStartPar
保持标题结构的一致性

\item {} 
\sphinxAtStartPar
公式使用 LaTeX 格式

\item {} 
\sphinxAtStartPar
Mermaid 图表用于流程说明

\end{itemize}


\subsection{注意事项}
\label{\detokenize{appendix/contributing:id6}}\begin{itemize}
\item {} 
\sphinxAtStartPar
提交前请确保文档可以正常构建

\item {} 
\sphinxAtStartPar
确保所有链接都有效

\item {} 
\sphinxAtStartPar
图片请使用相对路径引用

\item {} 
\sphinxAtStartPar
代码块请指定正确的语言以获得语法高亮

\end{itemize}


\bigskip\hrule\bigskip


\sphinxAtStartPar
© 版权所有 2026, 王培杰。最后更新于 2026 年 03 月 29 日。

\sphinxAtStartPar
利用 \sphinxhref{https://www.sphinx-doc.org/}{Sphinx} (https://www.sphinx\sphinxhyphen{}doc.org/) 构建，使用的 \sphinxhref{https://github.com/readthedocs/sphinx\_rtd\_theme}{sphinx\_rtd\_theme 主题} (https://github.com/readthedocs/sphinx\_rtd\_theme) 由 \sphinxhref{https://readthedocs.org}{Read the Docs} (https://readthedocs.org) 开发。



\renewcommand{\indexname}{索引}
\printindex
\end{document}